\documentclass[11pt,oneside]{book}
\usepackage[margin=1in]{geometry}
\usepackage{amsmath,amssymb,amsthm,mathtools}
\usepackage{booktabs,longtable,array}
\usepackage{enumitem}
\usepackage[colorlinks=true,linkcolor=blue,urlcolor=blue,citecolor=blue]{hyperref}
\usepackage[nameinlink,noabbrev]{cleveref}

\newtheorem{theorem}{Theorem}[chapter]
\newtheorem{lemma}[theorem]{Lemma}
\newtheorem{proposition}[theorem]{Proposition}
\newtheorem{corollary}[theorem]{Corollary}
\newtheorem{assessment}[theorem]{Assessment}
\newtheorem{firewall}[theorem]{Claim firewall}
\newtheorem{openproblem}[theorem]{Open problem}
\newtheorem{record}[theorem]{Record}
\newtheorem{programme}[theorem]{Programme}
\newtheorem{audit}[theorem]{Audit}
\newtheorem{correction}[theorem]{Correction}
\newtheorem{warning}[theorem]{Warning}
\theoremstyle{definition}
\newtheorem{definition}[theorem]{Definition}
\theoremstyle{remark}
\newtheorem{remark}[theorem]{Remark}

\newcommand{\MGclosed}{\textsc{closed}}
\newcommand{\MGopen}{\textsc{open}}
\newcommand{\MGpartial}{\textsc{partial}}
\newcommand{\opnorm}[1]{\lVert#1\rVert_{\mathrm{op}}}
\newcommand{\Tr}{\operatorname{Tr}}

\title{Renewal Yang--Mills Mass-Gap Research Notes\\
Fifth Series, Version V1}
\author{Compact proof-indexed continuation after archives f1--f4}
\date{September 2026}

\begin{document}
\maketitle
\tableofcontents

\chapter{Context, cumulative proof ledger, and continuation}
\label{chap:newV1-context}

\section{Purpose and claim boundary}
\label{sec:newV1-purpose}

This is the compact fifth-series continuation of a programme whose
ultimate target is a rigorous construction of four-dimensional
quantum Yang--Mills theory together with a strictly positive
Hamiltonian spectral gap above the vacuum.  The preceding fourth
series reached V100 and was frozen because the active source had
again become too large to navigate reliably.  The mathematics is
not being restarted: this note records the main proved interfaces,
states the exact live obstruction, and continues from it.

The current results concern concentrated one-link Wilson measures,
their representation-valued cumulants, and the marked-tree bounds
needed by a later multiscale construction.  They do not yet construct
the continuum measure or its Hilbert space.  No continuum mass gap is
assumed or claimed.

\begin{record}[Exact frozen source archives]
\label{rec:newV1-archives}
\begin{center}
\begin{tabular}{p{0.58\textwidth}rr}
\toprule
archive & source lines & bytes\\
\midrule
\texttt{Renewal\_Yang\_Mills\_Mass\_Gap\_Research\_Notes\_V135\_f1.tex}
& \(88\,517\) & \(3\,083\,467\)\\
\texttt{Renewal\_Yang\_Mills\_Mass\_Gap\_Research\_Notes\_V100\_f2.tex}
& \(49\,581\) & \(1\,706\,928\)\\
\texttt{Renewal\_Yang\_Mills\_Mass\_Gap\_Research\_Notes\_V100\_f3.tex}
& \(48\,381\) & \(1\,720\,945\)\\
\texttt{Renewal\_Yang\_Mills\_Mass\_Gap\_Research\_Notes\_V100\_f4.tex}
& \(44\,712\) & \(1\,483\,037\)\\
\bottomrule
\end{tabular}
\end{center}
Their SHA--256 digests, in the same order, are
\[
\begin{split}
&\texttt{82f6bbffeecdad3c5c63551660b19a0bdc7f3d6dffe4c68b4e9bba3a783a2b71},\\
&\texttt{6444a1044e74dfcb1a72ed9c7248bc0f234efc4a0e4248c736cd7db212dfca9a},\\
&\texttt{a5e05a144c753426d8a461c8290e9538a4cd19d8672bab8b08b509d82faf892f},\\
&\texttt{069cdb5806423cb3a871776eb571b63919feb56d4dd638be694c1991ccb0ae16}.
\end{split}
\]
Typewriter theorem names below are proof indices into these immutable
archives, not local cross-references.
\end{record}

\begin{record}[Version and audit protocol]
\label{rec:newV1-protocol}
Every successor copies the complete active note, appends one proof
chapter, passes the append-prefix, terminator, environment, label,
and reference validator, and then replaces its predecessor.  Frozen
archives remain.  At each version divisible by five, the newest
active Standard-Model--Einstein and Navier--Stokes notes are
inspected read-only.  At V100 the active volume is frozen under the
next suffix \(\texttt{\_fX}\), and another contextual V1 begins.
\end{record}

\section{The long implication chain}
\label{sec:newV1-route}

The intended route is
\begin{equation}
\begin{split}
\text{one-link connected Wilson cards}
&\Longrightarrow \text{marked-tree atlas (MTA)}
\Longrightarrow \text{weighted uniform fusion control (WUFC)}
\\
&\Longrightarrow \text{connected polymer/source majorant (CPM)}
\Longrightarrow \text{cutoff-uniform finite-volume measures}
\\
&\Longrightarrow \text{cutoff removal and reflection positivity}
\Longrightarrow \text{Osterwalder--Schrader reconstruction}
\\
&\Longrightarrow \text{positive physical Hamiltonian gap}.
\end{split}
\label{eq:newV1-route}
\end{equation}
The live work remains in the first implication.  The recent advance
is that every fixed arity is now controlled, but the constants have
not yet been proved to have the uniform factorial growth required to
sum all arities.

\section{Major results obtained before this volume}
\label{sec:newV1-summary}

\subsection{Representation normalization and connected combinatorics}

Archive f1 proves complete Peter--Weyl output pinching and removes
spurious output-dimension losses.  It introduces normalized row-pair
stocks and proves the binary and ternary marked-tree bounds.  Archive
f2 proves exact first-connectivity factorization: every connected
coordinate word has a unique first coordinate at which its cumulative
row partition becomes connected.  The proper prefix factors over its
blocks, the joining letter is a complete cumulant of the block
products, and the suffix is an unrestricted complete moment.  This
allocates the final resolvent payer once rather than once per block.

The same archives contain sharp obstruction tests.  A bare
one-Casimir-per-input quintic estimate fails on a coherent
\(SU(3)\) family; individual merge differences leak endpoint
covariance; and an equal-spin \(SU(2)\) heat square retains a singlet
with product-state capacity
\[
 (2j+1)^{-1}\asymp(1+C_2(j))^{-1/2}.
\]
These examples prevent later arguments from purchasing an
unavailable inverse Casimir power.

\subsection{Restored quartic atlas and the former quintic obstruction}

Archive f3 repairs the normalized \(3+1\) and \(2+1+1\) source
atlases by globally rescued, coordinate-summable square
decompositions.  Together with the established \(2+2\) sources it
proves the complete quartic MTA.  It then isolates the former
all-thermal quintic bottleneck rather than hiding it behind a
formal tree expansion.

\subsection{Symmetry-first projected tails}

Archive f4 replaces direct cumulant partitions by a centered subset
recursion.  For the one-link Wilson generator \(A_\alpha\), write
\[
 Q_\alpha F=F-\mathbb E_{\nu_\alpha}F,\qquad
 R_\alpha=A_\alpha^{-1}Q_\alpha,
\]
and let \(\Gamma\) be its carré du champ.  The centered contraction is
\begin{equation}
 \mathcal D_\alpha(F,G)
 :=
 Q_\alpha\Gamma(F,R_\alpha G).
 \label{eq:newV1-D}
\end{equation}
For a nonempty label set \(I\), the fully symmetrized projected tail
obeys the exact recursion
\begin{equation}
 \mathfrak P_\alpha(I)
 =
 \sum_{i\in I}
 \mathcal D_\alpha
 \bigl(F_i,\mathfrak P_\alpha(I\setminus\{i\})\bigr).
 \label{eq:newV1-subset}
\end{equation}
The exact covariance-centering identity and the negative norm
\[
 \|H\|_{-1,c}:=\|dR_\alpha H\|_{2,c}
\]
convert the final outer pairing into one terminal tail estimate.
The resulting reduced symmetrized projected-chain certificate
(RSPCC) at \(r=|I|\) is
\begin{equation}
 \|\mathfrak P_\alpha(I)\|_{-1,c}
 \le
 {r!K^r\over\sqrt\alpha}
 \left(\prod_{i\in I}b_i\right)
 \left(\sum_{i\in I}b_i\right)^{r-1}.
 \label{eq:newV1-RSPCC}
\end{equation}
Here
\[
 b_i=\sqrt{s_\alpha A_i}<1,\qquad
 s_\alpha\asymp\alpha^{-1},\qquad
 A_i\ge1+C_2(R_i).
\]
Whenever \eqref{eq:newV1-RSPCC} holds uniformly in \(r\), the Prüfer
identity produces the desired local labelled-tree card.

\subsection{One-link analytic estimates}

Archive f4 proves a Wilson spectral gap of order \(\alpha\) and, for
every fixed \(2\le p<\infty\), the complete dimension-free gradient
resolvent estimate
\begin{equation}
 \|dR_\alpha G\|_{p,c}
 \le {C_p\over\sqrt\alpha}\|Q_\alpha G\|_{p,c}.
 \label{eq:newV1-FWRC}
\end{equation}
The proof combines the exact global Bakry curvature lower bound
\(-\kappa\alpha\), a reverse semigroup Poincaré inequality, and
mean-zero \(L^p\) decay.  It does not infer positive curvature near
the antipode.

Global quaternion coordinates
\[
 U=\phi I+2\sum_{a=1}^3x_a e_a,\qquad
 \phi^2+|x|^2=1,
\]
avoid the logarithmic cut locus.  Exact linear, quadratic, and cubic
generator identities yield triangular resolvent normal forms.  In
particular, coordinate pairs become
\(\alpha^{-1}\) times centered quadratics plus a smaller remainder,
and all-linear three-tails become
\(\alpha^{-2}\) times a linear coordinate plus a smaller remainder.

\subsection{Tangent filtration and the current fixed-arity frontier}

After scaling \(y=\sqrt\alpha Z\), the generator converges to
\[
 A_0={1\over4}(-\Delta+y\cdot\nabla)
\]
and the contraction converges to the Gaussian
Ornstein--Uhlenbeck contraction.  Archive f4 proves the filtered
expansion
\begin{equation}
 \deg T_j(p_1,\ldots,p_r)
 \le
 D-2(r-1)+2j,
 \qquad D=\sum_i\deg p_i,
 \label{eq:newV1-filtration}
\end{equation}
for the coefficient of \(\alpha^{-j}\) in an ordered projected
tail.  Negative degree vanishes, and degree zero is removed by the
outer centering.

This closes the RSPCC with the sharp tree power for label sets of
size at most five and gives local all-thermal tree cards through
cumulant arity six.  The last two exact archival endpoints are
\texttt{thm:V99-complete-four-RSPCC} and
\texttt{thm:V100-complete-five-RSPCC}.

\begin{firewall}[Summary does not upgrade the result]
\label{fire:newV1-summary-boundary}
No result above controls the RSPCC constant uniformly at arbitrary
arity.  Nor does it prove anisotropic scalar optimization, the
complete selected gate, all-order MTA, WUFC, CPM, cutoff removal,
reflection positivity, Osterwalder--Schrader reconstruction, a
continuum Yang--Mills measure, or a positive continuum Hamiltonian
gap.  The full Yang--Mills mass gap has not been proved.
\end{firewall}

\chapter{V1 continuation: every fixed arity closes}
\label{chap:newV1-fixed-arity}

\section{Inherited proved interfaces}
\label{sec:newV1-inherited}

\begin{record}[Interfaces used in this chapter]
\label{rec:newV1-inherited}
The following statements are proved in archive f4 and used with
their original hypotheses:
\begin{enumerate}[leftmargin=3em]
\item the exact subset recursion
      \texttt{thm:V88-subset-recursion};
\item the complete fixed-\(p\) gradient resolvent
      \texttt{thm:V96-FWRC};
\item fixed logarithmic moments and complete Casimir row bounds;
\item the fixed-order tangent filtration
      \texttt{lem:V100-filtered-chain}; and
\item the terminal \(L^2\)-to-negative gap estimate.
\end{enumerate}
They are cited as inherited theorems, not reproved from assumptions.
\end{record}

\section{Taylor and Hölder cards at arbitrary fixed arity}
\label{sec:newV1-fixed-cards}

Fix \(r\ge2\).  Let \(L_R=dR(Z)\) and define
\begin{equation}
 \mathsf P_{R,d}:=
 Q_\alpha {L_R^d\over d!},
 \qquad 1\le d\le r-1,
 \label{eq:newV1-Pd}
\end{equation}
together with the exact centered remainder
\begin{equation}
 \mathsf N_{R,r}
 :=
 Q_\alpha D^R-\sum_{d=1}^{r-1}\mathsf P_{R,d}.
 \label{eq:newV1-Nr}
\end{equation}

\begin{lemma}[Fixed-\(r\) Taylor cards]
\label{lem:newV1-Taylor-cards}
For every fixed \(r\) and every fixed
\(2\le p\le2^r\),
\begin{align}
 \|\mathsf P_{R,d}\|_{p,c}
 &\le C_{r,p}b_R^d,
 &
 \|d\mathsf P_{R,d}\|_{p,c}
 &\le C_{r,p}\sqrt\alpha\,b_R^d,
 \qquad 1\le d\le r-1,
 \label{eq:newV1-Pd-cards}\\
 \|\mathsf N_{R,r}\|_{p,c}
 &\le C_{r,p}b_R^r,
 &
 \|d\mathsf N_{R,r}\|_{p,c}
 &\le C_{r,p}\sqrt\alpha\,b_R^r.
 \label{eq:newV1-Nr-cards}
\end{align}
The homogeneous coefficient rows obey
\begin{equation}
 |\mathsf P_{R,d}|_{\rm row}
 \le C_d\alpha^{d/2}b_R^d.
 \label{eq:newV1-Pd-row}
\end{equation}
All constants are independent of the representation and target
dimensions.
\end{lemma}

\begin{proof}
Taylor's integral remainder for \(e^{L_R}\) gives
\[
 \left\|
 e^{L_R}-\sum_{d=0}^{r-1}{L_R^d\over d!}
 \right\|
 \le {1\over(r-1)!}\|L_R\|^r
\]
in the unitary integral form.  After differentiation, every term
contains one \(dL_R\) and at least \(r-1\) undifferentiated
\(L_R\)'s.  Complete row--column Hölder, the identity
\(\sum_aT_a^*T_a=C_2(R)I\), and the inherited fixed logarithmic
moments give
\[
 \|d\mathsf N_{R,r}\|_{p,c}
 \le
 C_{r,p}\sqrt{C_2(R)}
 \left({C_2(R)\over\alpha}\right)^{(r-1)/2}
 \le C_{r,p}\sqrt\alpha\,b_R^r.
\]
The value bound is the same calculation without the differentiated
factor.  For a homogeneous degree-\(d\) piece, use \(d\) Casimir
rows and the \(d\)-th logarithmic moment to obtain
\eqref{eq:newV1-Pd-cards}; omit the moment scaling to obtain
\eqref{eq:newV1-Pd-row}.  Centering costs at most two in value norms
and commutes with differentiation.
\end{proof}

\begin{lemma}[Fixed-\(r\) degree-preserving chain]
\label{lem:newV1-degree-chain}
Assign degree \(d\) to \(\mathsf P_{R,d}\) and degree \(r\) to
\(\mathsf N_{R,r}\).  For any such \(r\) inputs and every ordering
\(\pi\),
\begin{equation}
 \|\mathcal C^\pi_{\alpha,r}\|_{-1,c}
 \le
 {C_r\over\sqrt\alpha}\prod_{i=1}^r b_i^{d_i}.
 \label{eq:newV1-degree-chain}
\end{equation}
\end{lemma}

\begin{proof}
Begin with the innermost pair and apply Hölder and
\eqref{eq:newV1-FWRC} at exponent \(2^r\); the centered output lies
in \(L^{2^{r-1}}\).  At each subsequent contraction halve the
exponent.  Lemma~\ref{lem:newV1-Taylor-cards} supplies the input
derivative factor \(\sqrt\alpha b_i^{d_i}\), and the resolvent
gradient supplies \(\alpha^{-1/2}\), so the scale powers cancel.
After \(r-1\) contractions the output is centered in \(L^2\).
The inherited gap estimate supplies the final
\(\alpha^{-1/2}\).  Complete row--column Hölder prevents any target
dimension factor.
\end{proof}

\begin{corollary}[Every high-excess sector]
\label{cor:newV1-high-excess}
If
\[
 E:=\sum_{i=1}^r(d_i-1)\ge r-1,
\]
then the corresponding degree sector satisfies
\begin{equation}
 \|\mathfrak P_\alpha(P_1,\ldots,P_r)\|_{-1,c}
 \le
 {C_r\over\sqrt\alpha}
 \left(\prod_i b_i\right)
 \left(\sum_i b_i\right)^{r-1}.
 \label{eq:newV1-high-excess}
\end{equation}
\end{corollary}

\begin{proof}
Sum \eqref{eq:newV1-degree-chain} over the \(r!\) orderings.
After extracting one \(b_i\) at every label, retain any \(r-1\)
excess factors, bound their product by
\((\sum_i b_i)^{r-1}\), and discard the rest using \(b_i<1\).
\end{proof}

\section{All low-excess sectors}
\label{sec:newV1-low-excess}

For \(0\le E\le r-2\), define
\begin{equation}
 k_r(E):=
 \left\lceil {r-1-E\over2}\right\rceil.
 \label{eq:newV1-kr}
\end{equation}

\begin{lemma}[First possible tangent-defect order]
\label{lem:newV1-first-defect}
Let \(d_i\ge1\), \(D=\sum_i d_i=r+E\), and
\(0\le E\le r-2\).  Every scaled projected-chain coefficient of
order \(\alpha^{-j}\) vanishes for \(j<k_r(E)\).
\end{lemma}

\begin{proof}
The inherited filtration \eqref{eq:newV1-filtration} gives
\[
 \deg T_j
 \le
 r+E-2(r-1)+2j
 =
 E-r+2+2j.
\]
For \(j<\lceil(r-1-E)/2\rceil\), the right side is at most zero.
Negative-degree coefficients vanish.  A degree-zero coefficient is
constant and is removed by the final centered contraction.  Hence
all such \(T_j\) vanish.
\end{proof}

\begin{lemma}[Fixed-\(r\) low-sector tangent estimate]
\label{lem:newV1-low-tangent}
For homogeneous logarithmic polynomials \(p_i\) of degrees \(d_i\)
as in Lemma~\ref{lem:newV1-first-defect},
\begin{equation}
 \|\mathfrak P_\alpha(Q_\alpha p_1,\ldots,Q_\alpha p_r)\|_{-1,c}
 \le
 C_r
 \alpha^{-(D+1)/2-k_r(E)}
 \prod_i|p_i|_{\rm row}.
 \label{eq:newV1-low-tangent}
\end{equation}
\end{lemma}

\begin{proof}
Use the inherited filtered expansion through order
\(k_r(E)-1\).  Lemma~\ref{lem:newV1-first-defect} removes every
displayed coefficient, leaving a scaled \(L^2\) remainder of order
\(\alpha^{-k_r(E)}\).  Returning from \(y=\sqrt\alpha Z\) supplies
\(\alpha^{-D/2}\), and the terminal gap supplies
\(\alpha^{-1/2}\).  The fixed \(r!\) orderings are absorbed into
\(C_r\).
\end{proof}

\begin{theorem}[Sharp tree power at every fixed arity]
\label{thm:newV1-every-fixed-arity}
For every fixed \(r\ge2\), there is a finite constant \(C_r\),
independent of \(\alpha\), the representations, and all target
dimensions, such that
\begin{equation}
 \boxed{\qquad
 \|\mathfrak P_\alpha(D^{R_1},\ldots,D^{R_r})\|_{-1,c}
 \le
 {C_r\over\sqrt\alpha}
 \left(\prod_{i=1}^r b_i\right)
 \left(\sum_{i=1}^r b_i\right)^{r-1}.
 \qquad}
 \label{eq:newV1-every-fixed-arity}
\end{equation}
\end{theorem}

\begin{proof}
Expand every input by
\eqref{eq:newV1-Pd}--\eqref{eq:newV1-Nr}.  A sector containing a
remainder has a degree-\(r\) label and therefore at least \(r-1\)
excess powers; it is covered by
\Cref{cor:newV1-high-excess}.  The same is true of every polynomial
sector with \(E\ge r-1\).

It remains to consider \(E\le r-2\).  Combine
\Cref{lem:newV1-low-tangent} with
\eqref{eq:newV1-Pd-row}.  The coefficient powers
\(\alpha^{d_i/2}\) cancel \(\alpha^{-D/2}\), giving
\[
 {C_r\over\sqrt\alpha}\,
 \alpha^{-k_r(E)}
 \prod_i b_i^{d_i}.
\]
After extracting one base factor at each label, the excess product
is at most \(B^E\), where \(B=\sum_i b_i\).  Since
\(\alpha^{-1}\le C_rB^2\) and \(B\le r\),
\[
 \alpha^{-k_r(E)}B^E
 \le
 C_r B^{2k_r(E)+E}
 \le
 C'_rB^{r-1}.
\]
The last inequality uses
\(2k_r(E)+E\ge r-1\); its excess over \(r-1\) is at most one and is
absorbed using \(B\le r\).  Sum the finite Taylor sector ledger.
\end{proof}

\begin{corollary}[Every fixed cumulant arity has a local tree card]
\label{cor:newV1-fixed-cumulants}
For each fixed \(m\ge2\), the \(m\)-th all-thermal Wilson cumulant
admits a labelled-tree bound with a finite arity-dependent constant
\(C_m\), uniform in \(\alpha\), representations, and target
dimensions.
\end{corollary}

\begin{proof}
Apply
\Cref{thm:newV1-every-fixed-arity} with \(r=m-1\) to every terminal
tail in the inherited symmetry-first cumulant compression.  Use the
outer negative pairing and the Prüfer identity.  All constants are
finite because \(m\) is fixed.
\end{proof}

\section{The exact remaining obstruction}
\label{sec:newV1-uniform-obstruction}

\begin{assessment}[Pointwise-in-arity is not all-order]
\label{ass:newV1-not-uniform}
\Cref{thm:newV1-every-fixed-arity} does not prove
\eqref{eq:newV1-RSPCC}.  Its proof permits:
\begin{enumerate}[leftmargin=3em]
\item fixed-\(p\) resolvent constants at \(p=2^r\);
\item filtered expansion constants through order about \(r/2\);
\item \(r^r\)-scale Taylor-sector counts; and
\item an uncontrolled \(r!\) permutation sum before cancellation.
\end{enumerate}
No bound of the form \(C_r\le r!K^r\) has been established.  This
constant-growth problem is now the first local obstruction; there is
no remaining exceptional finite arity.
\end{assessment}

\begin{openproblem}[Uniform analytic projected-chain card]
\label{open:newV1-UAPC}
Construct complete graded norms
\(\|\cdot\|_{\rho,c}\), \(0<\rho\le\rho_0\), and constants
\(K,c>0\) such that the following three statements hold
simultaneously:
\begin{enumerate}[leftmargin=3em]
\item the Wilson gradient resolvent has an analytic-radius estimate
\[
 \|dR_\alpha Q_\alpha G\|_{\rho',c}
 \le {K\over\sqrt\alpha(\rho-\rho')^{1/2}}
 \|G\|_{\rho,c};
\]
\item the centered contraction is a bounded bilinear map with only
one Cauchy radius loss, uniformly across all chaos degrees; and
\item the complete representation exponential and every filtered
tangent-defect coefficient have exponential generating-function
majorants, with the final centering performed before the norm.
\end{enumerate}
The radius losses must be schedulable across \(r-1\) contractions
with total cost at most \(r!K^r\), not \(K^{2^r}\).
\end{openproblem}

\begin{proposition}[Why the current proof cannot certify the analytic card]
\label{prop:newV1-current-method-fails}
The fixed-\(p\) estimate \eqref{eq:newV1-FWRC} and ordinary Hölder
alone do not imply the uniform analytic projected-chain card.
\end{proposition}

\begin{proof}
An \(r\)-input nested product ending in \(L^2\) forces the exponent
sequence
\[
 2^r,\ 2^{r-1},\ldots,8,4,2
\]
under repeated symmetric Hölder.  The inherited proof of
\eqref{eq:newV1-FWRC} yields a finite constant for each fixed \(p\)
but no estimate of \(C_p\) along \(p=2^r\).  Even an exponential
bound \(C_p\le e^{cp}\) would become doubly exponential in \(r\),
far larger than \(r!K^r\).  Therefore a proof using only those two
ingredients cannot establish the required uniform card; a
chaos-aware multiplication estimate or an assembled-chain
square-function estimate is genuinely new input.
\end{proof}

\begin{programme}[V2: upstream analytic-norm construction]
\label{prog:newV1-V2}
\begin{enumerate}[leftmargin=3em]
\item Put the Wilson diffusion in scaled radial chaos coordinates
      and define a complete Fock majorant whose degree weights
      include the representation factorial \(1/d!\).
\item Prove multiplication with additive, rather than doubling,
      chaos degree and quantify the smallest unavoidable analytic
      radius loss.
\item Revisit the semigroup proof of
      \eqref{eq:newV1-FWRC} in that majorant; do not infer an
      all-\(p\) bound from the fixed-\(p\) statement.
\item Expand the entire symmetrized subset recursion before applying
      the norm, preserving final centering and tangent-degree
      cancellation.
\item If an analytic resolvent card fails, exhibit the first chaos
      degree or radial coefficient causing failure and move upstream
      to a weighted carré-du-champ estimate.
\end{enumerate}
\end{programme}

\begin{firewall}[Claim boundary after fifth-series V1]
\label{fire:newV1-claim-boundary}
This note proves the sharp incidence power at every separately fixed
arity, with an arity-dependent constant.  It does not prove the
uniform factorial constant required for the all-order RSPCC.
Consequently it does not yet prove the all-order local tree card,
anisotropic scalar optimization, the selected gate, MTA, WUFC, CPM,
cutoff removal, reflection positivity,
Osterwalder--Schrader reconstruction, construction of the continuum
Yang--Mills measure, or a positive continuum Hamiltonian gap.  The
full Yang--Mills mass gap has not been proved.
\end{firewall}

\section{Fifth-series V1 result ledger}
\label{sec:newV1-results}

\begin{theorem}[V1 result ledger]
\label{thm:newV1-results}
This restarted V1:
\begin{enumerate}[label=\textup{(V1.\arabic*)},leftmargin=5.8em]
\item records exact fingerprints for all four frozen proof archives;
\item summarizes the cumulative mathematical route and claim
      boundary;
\item proves Taylor and derivative cards at arbitrary fixed arity;
\item proves the filtered low-sector estimate at arbitrary fixed
      arity;
\item closes the sharp tree incidence power for every separately
      fixed projected-tail and cumulant arity; and
\item isolates uniform factorial constant growth as the first
      remaining local obstruction.
\end{enumerate}
\end{theorem}

\begin{proof}
Use
\Cref{lem:newV1-Taylor-cards,
lem:newV1-degree-chain,
cor:newV1-high-excess,
lem:newV1-first-defect,
lem:newV1-low-tangent,
thm:newV1-every-fixed-arity,
cor:newV1-fixed-cumulants,
prop:newV1-current-method-fails}.
\end{proof}

\paragraph{Rollover record.}
This file begins after the validated fourth-series V100 was frozen
as
\texttt{Renewal\_Yang\_Mills\_Mass\_Gap\_Research\_Notes\_V100\_f4.tex}.
That archive has \(44\,712\) logical source lines,
\(1\,483\,037\) bytes, and SHA--256
\[
 \texttt{069cdb5806423cb3a871776eb571b63919feb56d4dd638be694c1991ccb0ae16}.
\]
The next compulsory companion audit is V5.

% =============================================================================
% BEGIN APPENDED FIFTH-SERIES V2 MATERIAL
% =============================================================================

\chapter{V2: Gaussian Fock majorants and the radius-loss audit}
\label{chap:newV2-Fock}

\section{A complete coefficient norm for Gaussian chaos}
\label{sec:newV2-chaos-norm}

Let \(H=\mathbb R^3\).  For a finite-dimensional Hilbert target
\(\mathcal V\), write a polynomial in its Wiener--Itô expansion
\[
 F(Y)=\sum_{n=0}^N I_n(f_n),
 \qquad
 f_n\in H^{\odot n}\otimes\mathcal V.
\]
The normalization is
\[
 \|I_n(f_n)\|_{L^2(\gamma_3;\mathcal V)}^2
 =n!\|f_n\|^2.
\]
For \(\rho>0\), define the coefficient majorant
\begin{equation}
 \|F\|_{\mathcal A_\rho,c}
 :=
 \sup_{\|v\|=1}
 \sum_{n\ge0}n!\rho^n\|f_n v\|.
 \label{eq:newV2-A-rho}
\end{equation}
For scalar or Hilbert-valued coefficients the supremum is omitted.
The completion of polynomials in this norm is denoted
\(\mathcal A_\rho\).  Tensor spectator legs are included in the
target \(\mathcal V\).

For a one-form chaos
\[
 \omega=\sum_{n\ge0}I_n(w_n),
 \qquad
 w_n\in H^{\odot n}\otimes H\otimes\mathcal V,
\]
use the same norm with the extra \(H\) factor in the Hilbert target.

\begin{lemma}[Exact chaos derivative weight]
\label{lem:newV2-derivative-weight}
For every polynomial \(F\),
\begin{equation}
 \|dF\|_{\mathcal A_\rho,c}
 \le {1\over\rho}\|F-\mathbb E_{\gamma_3}F\|_{\mathcal A_\rho,c}.
 \label{eq:newV2-derivative-weight}
\end{equation}
\end{lemma}

\begin{proof}
The derivative of \(I_n(f_n)\) has kernel \(nf_n\) and chaos
degree \(n-1\).  Its one-form contribution to the left side is
\[
 (n-1)!\rho^{n-1}n\|f_n v\|
 =
 {1\over\rho}n!\rho^n\|f_nv\|.
\]
Sum over \(n\ge1\) and take the complete column supremum.
\end{proof}

\section{Pointwise multiplication with an explicit radius loss}
\label{sec:newV2-product}

\begin{theorem}[Fock product card]
\label{thm:newV2-Fock-product}
If \(0<2\rho'\le\rho\), then
\begin{equation}
 \boxed{\qquad
 \|FG\|_{\mathcal A_{\rho'},c}
 \le
 e^{\rho^{-2}}
 \|F\|_{\mathcal A_\rho,c}
 \|G\|_{\mathcal A_\rho,c}.
 \qquad}
 \label{eq:newV2-Fock-product}
\end{equation}
The estimate is complete and independent of coefficient and
spectator dimensions.
\end{theorem}

\begin{proof}
It suffices to prove the Hilbert coefficient estimate after fixing
a unit input column.  The Gaussian product formula is
\begin{equation}
 I_m(f_m)I_n(g_n)
 =
 \sum_{k=0}^{\min(m,n)}
 k!\binom mk\binom nk
 I_{m+n-2k}
 \bigl(\operatorname{Sym}(f_m\otimes_k g_n)\bigr).
 \label{eq:newV2-product-formula}
\end{equation}
Contraction and symmetrization have norm at most one.
Put
\[
 a=m-k,\qquad b=n-k,\qquad \ell=a+b.
\]
After division by the two input weights, the contribution of this
term to the output weight is
\begin{align}
 &{\ell!\rho'^\ell\,
 k!\binom{a+k}{k}\binom{b+k}{k}
 \over
 (a+k)!\rho^{a+k}(b+k)!\rho^{b+k}}
 \nonumber\\
 &\hspace{4em}=
 {\ell!\over a!b!k!}
 {\rho'^{a+b}\over\rho^{a+b+2k}}
 \le
 \left({2\rho'\over\rho}\right)^{a+b}
 {\rho^{-2k}\over k!}.
 \label{eq:newV2-weight-ratio}
\end{align}
The first factor is at most one.  Sum the second over \(k\) to
obtain \(e^{\rho^{-2}}\); enlarging the sums in \(a,b\) factors the
two input majorants.  The argument uses only Hilbert contraction
and therefore remains valid in complete column norm with spectator
legs.
\end{proof}

\begin{corollary}[Centered product and carré-du-champ product]
\label{cor:newV2-centered-product}
Under the hypotheses of
\Cref{thm:newV2-Fock-product},
\begin{equation}
 \|Q_0(FG)\|_{\mathcal A_{\rho'},c}
 \le
 2e^{\rho^{-2}}
 \|F\|_{\mathcal A_\rho,c}
 \|G\|_{\mathcal A_\rho,c}.
 \label{eq:newV2-centered-product}
\end{equation}
The same estimate holds for the Hilbert contraction of two
one-form chaoses.
\end{corollary}

\begin{proof}
Deleting the zeroth chaos is contractive in
\(\mathcal A_{\rho'}\); the displayed factor two also covers the
Wilson-centering convention used after comparison.  The proof of
\Cref{thm:newV2-Fock-product} treats an extra finite-dimensional
Hilbert factor as a spectator, so it applies to one-form
contraction.
\end{proof}

\section{The exact tangent resolvent card}
\label{sec:newV2-OU-resolvent}

Recall
\[
 A_0={1\over4}(-\Delta+y\cdot\nabla),
 \qquad R_0=A_0^{-1}Q_0.
\]

\begin{theorem}[Fock gradient-resolvent card]
\label{thm:newV2-Fock-resolvent}
For every \(\rho>0\),
\begin{equation}
 \boxed{\qquad
 \|dR_0F\|_{\mathcal A_\rho,c}
 \le {4\over\rho}
 \|Q_0F\|_{\mathcal A_\rho,c}.
 \qquad}
 \label{eq:newV2-Fock-resolvent}
\end{equation}
\end{theorem}

\begin{proof}
On the \(n\)-th nonconstant chaos, \(A_0=n/4\), so
\(R_0I_n(f_n)=(4/n)I_n(f_n)\).  Its differentiated contribution is
\[
 (n-1)!\rho^{n-1}n{4\over n}\|f_nv\|
 =
 {4\over n\rho}n!\rho^n\|f_nv\|
 \le
 {4\over\rho}n!\rho^n\|f_nv\|.
\]
Sum and take the column supremum.
\end{proof}

\begin{corollary}[One Gaussian centered contraction]
\label{cor:newV2-one-D0}
If \(0<2\rho'\le\rho\), then
\begin{equation}
 \|\mathcal D_0(F,G)\|_{\mathcal A_{\rho'},c}
 \le
 {8e^{\rho^{-2}}\over\rho^2}
 \|Q_0F\|_{\mathcal A_\rho,c}
 \|Q_0G\|_{\mathcal A_\rho,c}.
 \label{eq:newV2-one-D0}
\end{equation}
\end{corollary}

\begin{proof}
Use
\[
 \mathcal D_0(F,G)
 =
 Q_0\Gamma_0(F,R_0G)
 =
 {1\over4}Q_0\langle dF,dR_0G\rangle.
\]
Apply
\Cref{lem:newV2-derivative-weight,
thm:newV2-Fock-resolvent,
cor:newV2-centered-product}.  The factors are
\(\rho^{-1}\), \(4\rho^{-1}\), \(1/4\), and at most
\(2e^{\rho^{-2}}\).
\end{proof}

\section{Wilson perturbations in analytic radius}
\label{sec:newV2-Wilson-radius}

\begin{lemma}[Polynomial differential operators cost Cauchy radius]
\label{lem:newV2-Cauchy-loss}
Let \(E\) be a differential operator on Gaussian polynomials whose
coefficient of every derivative has degree at most the derivative
order plus \(q\), with fixed coefficient norm.  For
\(0<\rho'<\rho\),
\begin{equation}
 \|EF\|_{\mathcal A_{\rho'},c}
 \le
 {C_E\over(\rho-\rho')^{q+s}}
 \|F\|_{\mathcal A_\rho,c},
 \label{eq:newV2-Cauchy-loss}
\end{equation}
where \(s\) is the differential order.  Only the existence of a
fixed polynomial power is used below.
\end{lemma}

\begin{proof}
Acting on chaos degree \(n\), the derivative and multiplication
coefficients are bounded by a fixed polynomial \(C_E(1+n)^{q+s}\).
Let \(t=\rho'/\rho<1\).  The elementary bound
\[
 \sup_{n\ge0}(1+n)^a t^n
 \le {C_a\over(1-t)^a}
\]
with \(a=q+s\) absorbs that polynomial.  Since
\(1-t=(\rho-\rho')/\rho\), the harmless fixed powers of \(\rho\)
are absorbed in \(C_E\) on a prescribed compact radius interval.
Sum the chaos coefficient majorant.
\end{proof}

\begin{proposition}[First Wilson analytic resolvent defect]
\label{prop:newV2-first-Wilson-defect}
Fix \(0<\rho'<\rho\) in a compact positive radius interval.  In
the scaled logarithmic chart, for polynomial \(F\) and sufficiently
large \(\alpha\),
\begin{equation}
 \|d(\widehat R_\alpha-R_0)F\|_{\mathcal A_{\rho'},c}
 \le
 {C\over\alpha(\rho-\rho')^q}
 \|Q_0F\|_{\mathcal A_\rho,c}
 +O(e^{-c\alpha})\|F\|_{\mathcal A_\rho,c}
 \label{eq:newV2-first-Wilson-defect}
\end{equation}
for a fixed integer \(q\).  The constant is complete and
dimension independent.
\end{proposition}

\begin{proof}
The scaled generator expansion proved in archive f4 V99--V100 is
\[
 \widehat A_\alpha=A_0+\alpha^{-1}E_1
 +O(\alpha^{-2}),
\]
where \(E_1\) is a fixed polynomial differential operator and the
remainder obeys the same statement with a larger fixed Cauchy loss.
Use the resolvent identity
\[
 \widehat R_\alpha-R_0
 =
 -\widehat R_\alpha
 (\widehat A_\alpha-A_0)R_0
\]
on centered cutoff polynomials.  Apply
\Cref{thm:newV2-Fock-resolvent,lem:newV2-Cauchy-loss} and the
scaled Wilson gradient-resolvent theorem
\texttt{thm:V96-FWRC}, first on finite chaos truncations and then
by closure.  The identity-neighbourhood cutoff error is
\(O(e^{-c\alpha})\) in every fixed analytic radius strictly inside
the input radius, by the Wilson tail estimate.  This proves
\eqref{eq:newV2-first-Wilson-defect}.
\end{proof}

\section{The decisive audit: a binary Fock norm is still too local}
\label{sec:newV2-radius-audit}

\begin{proposition}[Repeated binary use collapses the radius]
\label{prop:newV2-radius-collapse}
An \(r\)-input projected chain estimated only by repeated use of
\Cref{cor:newV2-one-D0} requires a radius sequence satisfying
\[
 2\rho_{j+1}\le\rho_j
\]
at every contraction.  Hence
\[
 \rho_{r-1}\le2^{-(r-1)}\rho_0,
\]
and the derivative factors alone can grow at least like
\[
 \prod_{j=0}^{r-2}\rho_j^{-2}
 \ge
 \rho_0^{-2(r-1)}2^{(r-1)(r-2)}.
 \label{eq:newV2-radius-collapse}
\]
This is superfactorial and cannot prove an \(r!K^r\) RSPCC.
\end{proposition}

\begin{proof}
The product theorem requires the displayed radius halving.  The
contraction card contributes \(\rho_j^{-2}\) at the \(j\)-th use.
The largest admissible output radii are
\(\rho_j=2^{-j}\rho_0\).  Their product is
\[
 \prod_{j=0}^{r-2}(2^j/\rho_0)^2
 =
 \rho_0^{-2(r-1)}
 2^{2\sum_{j=0}^{r-2}j},
\]
which is \eqref{eq:newV2-radius-collapse}.  Since
\(\log(r!K^r)=O(r\log r)\) but the logarithm of the last factor is
of order \(r^2\), no fixed \(K\) absorbs it.
\end{proof}

\begin{assessment}[What V2 has ruled out]
\label{ass:newV2-ruled-out}
The coefficient Fock norm removes the uncontrolled \(C_{2^r}\)
constants and gives an exact tangent resolvent estimate.  It does
not solve the all-order problem because applying its binary product
card after each contraction destroys analytic radius geometrically.
Thus the failure in fifth-series V1 is not merely an unfortunate
choice of \(L^p\) norms.  Any successful estimate must assemble the
connected \(r\)-label diagram before applying the analytic norm, or
use a norm whose connected product is bounded without binary radius
halving.
\end{assessment}

\begin{openproblem}[Assembled connected Fock card]
\label{open:newV2-assembled-Fock}
Prove directly, without iterating
\eqref{eq:newV2-one-D0}, that the fully symmetrized Gaussian tail
has a connected-diagram estimate
\begin{equation}
 \|\mathfrak P_0(F_1,\ldots,F_r)\|_{\mathcal A_{\rho'},c}
 \le
 r!K^r
 \sum_{\tau\in\mathfrak T_r}
 \prod_{ij\in E(\tau)}w_{ij}
 \label{eq:newV2-assembled-Fock}
\end{equation}
with one total radius loss \(\rho-\rho'\), not one halving per
edge.  The weights \(w_{ij}\) must be compatible with the Casimir
row-pair stocks after the representation exponential is inserted.
\end{openproblem}

\begin{programme}[V3: rooted connected Gaussian diagrams]
\label{prog:newV2-V3}
\begin{enumerate}[leftmargin=3em]
\item Expand the fully symmetrized projected tail in Hermite
      contractions before taking any norm.
\item Prove that final centering leaves only diagrams whose label
      graph is connected and rooted at the uncontracted chaos legs.
\item Sum contraction orders with the \(1/d!\) representation
      coefficients before bounding coefficient tensors.
\item Seek a Penrose or Prüfer allocation that charges every
      connected diagram to one labelled tree with only exponential
      multiplicity.
\item Only after the assembled Gaussian card is proved, insert one
      Wilson generator or measure defect and test whether its
      Cauchy radius cost remains summable.
\end{enumerate}
\end{programme}

\begin{firewall}[A tangent Fock norm is not the all-order Wilson card]
\label{fire:newV2-boundary}
\Cref{thm:newV2-Fock-product,thm:newV2-Fock-resolvent} are exact
Gaussian analytic estimates, and
\Cref{prop:newV2-first-Wilson-defect} controls one perturbative
defect with radius loss.  V2 does not prove the assembled connected
card \eqref{eq:newV2-assembled-Fock}, uniform Wilson defect
summation, the all-order RSPCC, all-order local MTA, WUFC, CPM,
cutoff removal, reflection positivity,
Osterwalder--Schrader reconstruction, a continuum Yang--Mills
measure, or a positive Hamiltonian gap.  The full Yang--Mills mass
gap has not been proved.
\end{firewall}

\section{Fifth-series V2 result ledger}
\label{sec:newV2-results}

\begin{theorem}[V2 result ledger]
\label{thm:newV2-results}
V2:
\begin{enumerate}[label=\textup{(V2.\arabic*)},leftmargin=5.8em]
\item defines a complete Gaussian chaos coefficient majorant;
\item proves its exact derivative weight and pointwise product card;
\item proves the exact Ornstein--Uhlenbeck gradient-resolvent card;
\item derives a first Wilson analytic resolvent-defect estimate;
\item proves that repeated binary use collapses radius
      superfactorially; and
\item isolates the assembled connected Fock card as the next
      upstream acquisition.
\end{enumerate}
\end{theorem}

\begin{proof}
Use
\Cref{lem:newV2-derivative-weight,
thm:newV2-Fock-product,
thm:newV2-Fock-resolvent,
cor:newV2-one-D0,
lem:newV2-Cauchy-loss,
prop:newV2-first-Wilson-defect,
prop:newV2-radius-collapse}.
\end{proof}

\paragraph{Parent-version record.}
V2 was formed from
\texttt{Renewal\_Yang\_Mills\_Mass\_Gap\_Research\_Notes\_V1.tex}.
The V1 parent had \(682\) logical source lines,
\(23\,665\) bytes, and SHA--256
\[
 \texttt{eeaf1dfa6b36dab5fa43decab4f46987bec9fc0ea12a86f99e906e7fbabb6336}.
\]
The next compulsory companion audit is V5.

% =============================================================================
% END APPENDED FIFTH-SERIES V2 MATERIAL
% =============================================================================

% =============================================================================
% BEGIN APPENDED FIFTH-SERIES V3 MATERIAL
% =============================================================================

\chapter{V3: Forest interpolation closes the assembled Gaussian card}
\label{chap:newV3-forest}

\section{Replicated Gaussian interpolation}
\label{sec:newV3-replicas}

Let \(I=\{1,\ldots,r\}\).  For a positive semidefinite correlation
matrix \(C=(C_{ij})_{i,j\in I}\) with \(C_{ii}=1\), let
\(\mathbb E_C\) denote expectation of \(H\)-valued Gaussian replicas
\((Y_i)_{i\in I}\) satisfying
\[
 \mathbb E_C(Y_i^aY_j^b)=C_{ij}\delta_{ab}.
\]
Functions attached to distinct replicas take values on distinct
tensor legs and therefore commute.

For a forest \(F\) on \(I\) and parameters
\(\mathbf s=(s_e)_{e\in F}\in[0,1]^F\), define
\begin{equation}
 C^F_{ij}(\mathbf s)
 :=
 \begin{cases}
 1,&i=j,\\
 \min\{s_e:e\text{ lies on the unique }F\text{-path from }i
 \text{ to }j\},&i,j\text{ are }F\text{-connected},\\
 0,&\text{otherwise}.
 \end{cases}
 \label{eq:newV3-forest-covariance}
\end{equation}

\begin{lemma}[Positivity of the forest covariance]
\label{lem:newV3-forest-positive}
For every forest \(F\) and every \(\mathbf s\in[0,1]^F\),
\(C^F(\mathbf s)\) is positive semidefinite.
\end{lemma}

\begin{proof}
For \(t\in[0,1]\), delete from \(F\) every edge \(e\) with
\(s_e<t\), and let \(P_t\) be the matrix whose \(ij\) entry is one
when \(i,j\) lie in the same remaining component and zero otherwise.
Each \(P_t\) is a block-diagonal sum of all-ones matrices and is
positive semidefinite.  Moreover
\[
 C^F(\mathbf s)=\int_0^1P_t\,dt
\]
entrywise, including the diagonal.  The integral is positive
semidefinite.
\end{proof}

For an edge \(e=\{i,j\}\), set
\begin{equation}
 \Delta_e:=\sum_{a=1}^{\dim H}\partial_{y_i^a}\partial_{y_j^a}.
 \label{eq:newV3-edge-derivative}
\end{equation}

\begin{theorem}[BKAR connected forest formula]
\label{thm:newV3-BKAR}
Let \(f_i:H\to\operatorname{End}(V_i)\) be smooth with bounded
derivatives through order \(r-1\).  Then
\begin{equation}
 \boxed{\qquad
 \kappa_r^{\gamma_H}(f_1(Y),\ldots,f_r(Y))
 =
 \sum_{T\in\mathfrak T_r}
 \int_{[0,1]^{E(T)}}\!
 \mathbb E_{C^T(\mathbf s)}
 \left[
 \left(\prod_{e\in E(T)}\Delta_e\right)
 \prod_{i=1}^r f_i(Y_i)
 \right]\,d\mathbf s .
 \qquad}
 \label{eq:newV3-BKAR}
\end{equation}
The identity is valid entrywise and hence for the tensor-valued
functions used here.
\end{theorem}

\begin{proof}
For completeness, start with the replicated moment
\[
 M(C):=\mathbb E_C\prod_i f_i(Y_i).
\]
Gaussian differentiation gives, for \(i\ne j\),
\begin{equation}
 {\partial M\over\partial C_{ij}}
 =
 \mathbb E_C\left[\Delta_{\{i,j\}}\prod_i f_i(Y_i)\right].
 \label{eq:newV3-Gaussian-differentiation}
\end{equation}
Apply the fundamental theorem of calculus successively to the
off-diagonal entries, using the forest interpolation
\eqref{eq:newV3-forest-covariance}.  The resulting forest formula is
\begin{equation}
 M(\mathbf 1)
 =
 \sum_{F\text{ forest on }I}
 \int_{[0,1]^{E(F)}}\!
 \mathbb E_{C^F(\mathbf s)}
 \left[
 \left(\prod_{e\in E(F)}\Delta_e\right)
 \prod_i f_i(Y_i)
 \right]d\mathbf s .
 \label{eq:newV3-BKAR-moment}
\end{equation}
One proves this finite identity first for Gaussian-polynomial
functions by induction on the number of off-diagonal variables;
bounded-derivative approximation then gives the stated class.
Lemma~\ref{lem:newV3-forest-positive} ensures that every interpolated
expectation is a genuine Gaussian expectation.

If \(F\) has components \(B_1,\ldots,B_k\), both its covariance and
its differential operator factor over those components.  Therefore
the \(F\)-term in \eqref{eq:newV3-BKAR-moment} is a product of
connected component functionals.  Möbius inversion on set
partitions says that the joint cumulant is precisely the connected
component functional.  The connected forests on \(I\) are exactly
the labelled spanning trees, proving \eqref{eq:newV3-BKAR}.
\end{proof}

\begin{remark}[Why this avoids the V2 radius collapse]
\label{rem:newV3-no-binary-radius}
All \(r-1\) edge derivatives in \eqref{eq:newV3-BKAR} act before a
single expectation and a single norm.  No intermediate product is
placed in a smaller analytic radius.  Connectedness is imposed by
Möbius inversion, not recovered after \(r-1\) applications of a
binary estimate.
\end{remark}

\section{Complete derivative rows of representation exponentials}
\label{sec:newV3-representation-derivatives}

For an irreducible unitary representation \(R\), let
\(T_{R,a}=dR(e_a)\), normalized by
\begin{equation}
 \sum_aT_{R,a}^*T_{R,a}=C_2(R)I.
 \label{eq:newV3-Casimir}
\end{equation}
For \(z\in\mathfrak{su}(2)\), put
\[
 F_R(z):=D^R(e^z).
\]

\begin{lemma}[Complete symmetrized derivative row]
\label{lem:newV3-derivative-row}
For every \(d\ge1\),
\begin{equation}
 \sup_z
 \left\|
 \left(
 \partial_{a_1}\cdots\partial_{a_d}F_R(z)
 \right)_{a_1,\ldots,a_d}
 \right\|_{\rm row}
 \le
 C_0^d C_2(R)^{d/2}.
 \label{eq:newV3-derivative-row}
\end{equation}
The bound is independent of \(\dim R\).
\end{lemma}

\begin{proof}
The \(d\)-th derivative of the exponential is the symmetrized
Duhamel integral over the \(d\)-simplex:
\[
 \sum_{\sigma\in S_d}
 \int_{0<t_1<\cdots<t_d<1}
 U_0T_{R,a_{\sigma(1)}}U_1\cdots
 T_{R,a_{\sigma(d)}}U_d\,d\mathbf t,
\]
where all \(U_j\) are unitary representation matrices.  The simplex
has volume \(1/d!\), which cancels the \(d!\) permutations.
Unitary factors do not change column norms.  Apply complete
row--column Cauchy--Schwarz successively to the \(d\) generator
slots and use \eqref{eq:newV3-Casimir} at each slot.  A fixed change
between the chosen orthonormal Lie basis and exponential coordinates
is absorbed in \(C_0^d\).  No trace or output-dimension estimate is
used.
\end{proof}

\begin{corollary}[Scaled derivative weight]
\label{cor:newV3-scaled-row}
For
\[
 f_{R,\alpha}(y):=D^R(e^{y/\sqrt\alpha}),
\]
one has
\begin{equation}
 \sup_y\|\nabla^d f_{R,\alpha}(y)\|_{\rm row}
 \le
 C_1^d b_R^d,
 \label{eq:newV3-scaled-row}
\end{equation}
where \(b_R=\sqrt{s_\alpha A_R}\),
\(A_R\ge1+C_2(R)\), and \(s_\alpha\asymp\alpha^{-1}\).
\end{corollary}

\begin{proof}
Each derivative contributes \(\alpha^{-1/2}\).  Apply
\Cref{lem:newV3-derivative-row} and
\(\alpha^{-1/2}\sqrt{C_2(R)}\le Cb_R\).
\end{proof}

\section{The all-order Gaussian tangent tree card}
\label{sec:newV3-Gaussian-card}

\begin{theorem}[Complete Gaussian representation tree card]
\label{thm:newV3-Gaussian-tree-card}
For every \(r\ge2\) and all all-thermal irreducible unitary
representations,
\begin{equation}
 \boxed{\qquad
 \left\|
 \kappa_r^{\gamma_3}
 \bigl(
 f_{R_1,\alpha}(Y),\ldots,f_{R_r,\alpha}(Y)
 \bigr)
 \right\|_{\rm op}
 \le
 C_1^{2r-2}
 \sum_{T\in\mathfrak T_r}
 \prod_{i=1}^r b_i^{d_i(T)}.
 \qquad}
 \label{eq:newV3-Gaussian-tree-card}
\end{equation}
The estimate is uniform in \(r,\alpha\), representation dimensions,
and tensor output dimension.
\end{theorem}

\begin{proof}
Insert the functions \(f_{R_i,\alpha}\) in
\eqref{eq:newV3-BKAR}.  At a vertex \(i\), the edge derivatives
assemble into the \(d_i(T)\)-th derivative row of that one function.
Complete row--column Cauchy--Schwarz along the tree edges and
\eqref{eq:newV3-scaled-row} give
\[
 C_1^{\sum_i d_i(T)}
 \prod_i b_i^{d_i(T)}
 =
 C_1^{2r-2}\prod_i b_i^{d_i(T)}.
\]
Every interpolated Gaussian measure is a probability measure and
the remaining representation factors are unitary, so its
expectation has operator norm at most the same row contraction
bound.  The parameter cube has volume one.  Sum over labelled
trees.  At no stage is a tensor trace or matrix-dimension factor
introduced.
\end{proof}

\begin{corollary}[Uniform tangent all-order MTA]
\label{cor:newV3-Gaussian-MTA}
The complete all-order local marked-tree atlas holds for the
Gaussian tangent model, with exponential constant
\(K_G=C_1^2\).
\end{corollary}

\begin{proof}
Equation \eqref{eq:newV3-Gaussian-tree-card} is exactly the local
tree card: every labelled tree appears once and receives its
Casimir incidence weight.  Absorb \(C_1^{2r-2}\) into \(K_G^r\).
\end{proof}

\section{What remains when the tangent law is replaced by Wilson}
\label{sec:newV3-Wilson-gap}

\begin{assessment}[The Gaussian constant-growth problem is closed]
\label{ass:newV3-Gaussian-closed}
V2's superfactorial radius loss was an artifact of norming binary
contractions before assembly.  The forest formula proves that the
fully assembled Gaussian cumulant has exactly one sum over labelled
trees and only an exponential derivative constant.  Thus neither
Gaussian Wick contractions nor representation derivative rows are
the source of the remaining arity-growth obstruction.
\end{assessment}

\begin{definition}[Uniform Wilson transport card]
\label{def:newV3-UWTC}
Let \(\mathcal T_\alpha:\mathbb R^3\to\mathfrak{su}(2)\) be a
measurable map pushing \(\gamma_3\) to the global logarithmic Wilson
law.  The uniform Wilson transport card (UWTC) is the assertion that
one may choose \(\mathcal T_\alpha\) so that, for every \(d\ge1\),
\begin{equation}
 \left\|
 \nabla^d
 \left[
 D^R\bigl(e^{\mathcal T_\alpha(y)}\bigr)
 \right]
 \right\|_{\rm row}
 \le
 d!K_T^d b_R^d
 \label{eq:newV3-UWTC}
\end{equation}
in every replicated Gaussian expectation occurring in
\eqref{eq:newV3-BKAR}, with a common \(K_T\) independent of
\(\alpha,R,d\).  An integrated version with the same treewise
product bound is sufficient; a global pointwise bound is not
required.
\end{definition}

\begin{proposition}[UWTC implies the all-order Wilson local card]
\label{prop:newV3-UWTC-implies-card}
If the UWTC holds and its vertex factorials satisfy
\begin{equation}
 \prod_{i=1}^r d_i(T)!\le(r-1)!
 \quad\text{for every }T\in\mathfrak T_r,
 \label{eq:newV3-tree-factorials}
\end{equation}
then the all-order all-thermal Wilson local tree card holds with
constant \(r!K^r\).
\end{proposition}

\begin{proof}
Push every Wilson expectation to \(\gamma_3\) by
\(\mathcal T_\alpha\) and apply
\Cref{thm:newV3-BKAR}.  The UWTC bounds the \(T\)-term by
\[
 K_T^{2r-2}
 \left(\prod_i d_i(T)!\right)
 \prod_i b_i^{d_i(T)}.
\]
Use \eqref{eq:newV3-tree-factorials} and absorb
\(K_T^{2r-2}\) into \(K^r\).

It remains to justify the elementary factorial inequality.
Remove a leaf and its incident edge.  If its neighbour has degree
\(d\), the product of degree factorials changes by the factor \(d\).
Inductively expose leaves in an order obtained by a Prüfer word.
At the \(k\)-th removal this factor is at most \(k\); their product
is at most \((r-1)!\).  This proves
\eqref{eq:newV3-tree-factorials}.
\end{proof}

\begin{remark}[Why the transport card is not assumed]
\label{rem:newV3-transport-not-proved}
The radial Wilson law admits a monotone radial transport from
\(\gamma_3\), but its derivatives can deteriorate near the
antipodal boundary where the target density vanishes.  Exponential
smallness of that region is enough for each fixed derivative order;
it has not yet been shown to compensate uniformly for all orders
in the replicated forest expectations.  The UWTC is therefore a
precise acquisition target, not a theorem of V3.
\end{remark}

\begin{programme}[V4: radial transport derivative ledger]
\label{prog:newV3-V4}
\begin{enumerate}[leftmargin=3em]
\item Write the exact radial quantile equation taking the
      three-dimensional Gaussian radius to the scaled Wilson
      logarithmic radius.
\item Prove uniform first and second radial derivative bounds on the
      central region and quantify the antipodal tail contribution.
\item Derive the general inverse-function recursion for the
      \(d\)-th radial derivative with factorially weighted
      coefficients.
\item Test whether the Gaussian tail compensates the boundary
      denominators with only \(d!K^d\) growth in every forest
      interpolation.
\item If pointwise transport derivatives fail, retain the integrated
      treewise UWTC and avoid strengthening it.
\end{enumerate}
\end{programme}

\begin{firewall}[Gaussian MTA is not Wilson MTA]
\label{fire:newV3-boundary}
\Cref{thm:newV3-Gaussian-tree-card} proves the all-order local card
for the Gaussian tangent model.  It does not compare all derivative
orders of the Wilson law uniformly, and the UWTC remains unproved.
Therefore V3 does not establish the all-order Wilson local MTA,
WUFC, CPM, cutoff removal, reflection positivity,
Osterwalder--Schrader reconstruction, a continuum Yang--Mills
measure, or a positive Hamiltonian gap.  The full Yang--Mills mass
gap has not been proved.
\end{firewall}

\section{Fifth-series V3 result ledger}
\label{sec:newV3-results}

\begin{theorem}[V3 result ledger]
\label{thm:newV3-results}
V3:
\begin{enumerate}[label=\textup{(V3.\arabic*)},leftmargin=5.8em]
\item proves positivity of the forest interpolation covariances;
\item proves the connected Gaussian BKAR tree formula;
\item proves complete dimension-free derivative rows for unitary
      representation exponentials;
\item proves the all-order Gaussian tangent local tree card with
      exponential constants;
\item shows that the V2 radius collapse disappears after connected
      assembly; and
\item reduces the Wilson deformation to the explicit UWTC radial
      transport derivative problem.
\end{enumerate}
\end{theorem}

\begin{proof}
Use
\Cref{lem:newV3-forest-positive,
thm:newV3-BKAR,
lem:newV3-derivative-row,
cor:newV3-scaled-row,
thm:newV3-Gaussian-tree-card,
prop:newV3-UWTC-implies-card}.
\end{proof}

\paragraph{Parent-version record.}
V3 was formed from
\texttt{Renewal\_Yang\_Mills\_Mass\_Gap\_Research\_Notes\_V2.tex}.
The V2 parent had \(1\,118\) logical source lines,
\(37\,011\) bytes, and SHA--256
\[
 \texttt{7d4db79d1e06171a1810d96471f29bb632ed644e8c4053fd77f7870ef05a6024}.
\]
The next compulsory companion audit is V5.

% =============================================================================
% END APPENDED FIFTH-SERIES V3 MATERIAL
% =============================================================================

% =============================================================================
% BEGIN APPENDED FIFTH-SERIES V4 MATERIAL
% =============================================================================

\chapter{V4: Archive correction, high-arity closure, and rare-error
absorption}
\label{chap:newV4-arity-split}

\section{Nonduplication correction after the transport audit}
\label{sec:newV4-correction}

Before following the V3 transport programme, archives f3--f4 were
searched for the complete earlier transport history.  The following
facts were already proved there:
\begin{enumerate}[leftmargin=3em]
\item \texttt{thm:V71-pushforward} constructs the exact monotone
      radial Gaussian-to-Wilson transport;
\item \texttt{thm:V71-RTDC-implies-card} gives the same conditional
      forest compiler as V3;
\item \texttt{thm:V72-RTDC-false} proves that its pointwise
      derivative certificate fails already at first order because
      of the square-root antipodal caustic;
\item \texttt{thm:V77-no-SEDE} proves that a common absolute
      sub-Gaussian derivative envelope fails at second order; and
\item \texttt{fire:V77-WRTTC-open} leaves only an integrated,
      cancellation-sensitive transport card open.
\end{enumerate}

\begin{record}[Status of the V3 claims]
\label{rec:newV4-V3-status}
The Gaussian BKAR formula and Gaussian tree card in V3 are correct,
but they compactly reprove the archived Gaussian forest interface
rather than add a new Wilson theorem.  The integrated option in the
V3 UWTC is the archived WRTTC under new notation.  Its pointwise
interpretation is unavailable.  No later argument may cite V3 as
having proved a Wilson transport derivative card.
\end{record}

\begin{assessment}[Direction correction]
\label{ass:newV4-direction}
V4 does not differentiate the radial quantile again.  It instead
asks how much of the all-order local card can be removed before any
transport or intrinsic derivative estimate is attempted.  Two
elementary but uniform mechanisms suffice:
\begin{enumerate}[leftmargin=3em]
\item at arity comparable with or larger than \(\alpha\), the
      number of labelled trees and the universal lower bound on
      every \(b_i\) dominate the trivial cumulant estimate; and
\item at the complementary arities, an \(e^{-c\alpha}\) error in
      the carriers is smaller than the complete tree scale after a
      fixed enlargement of the exponential base.
\end{enumerate}
\end{assessment}

\section{A universal bounded-variable cumulant estimate}
\label{sec:newV4-trivial-cumulant}

\begin{lemma}[One-factorial trivial cumulant bound]
\label{lem:newV4-trivial-cumulant}
Let \(X_1,\ldots,X_r\) be compatible operator-valued random
variables on distinct tensor legs with
\(\|X_i\|_{\rm op}\le1\) almost surely.  Then
\begin{equation}
 \|\kappa_r(X_1,\ldots,X_r)\|_{\rm op}
 \le r^r\le e^r r!.
 \label{eq:newV4-trivial-cumulant}
\end{equation}
\end{lemma}

\begin{proof}
The moment--cumulant formula gives
\[
 \kappa_r(X_1,\ldots,X_r)
 =
 \sum_{\pi\in\Pi_r}
 (|\pi|-1)!(-1)^{|\pi|-1}
 \prod_{B\in\pi}\mathbb E\prod_{i\in B}X_i.
\]
Every block moment has norm at most one.  Group partitions by
\(k=|\pi|\).  Since
\[
 S(r,k)\le{k^r\over k!},
\]
the norm is at most
\[
 \sum_{k=1}^r S(r,k)(k-1)!
 \le\sum_{k=1}^r k^{r-1}
 \le r^r.
\]
The last inequality in \eqref{eq:newV4-trivial-cumulant} follows
from the elementary Stirling lower bound
\(r!\ge(r/e)^r\).
\end{proof}

For the all-thermal scale, put
\begin{equation}
 \mathcal W_r(\mathbf b)
 :=
 \sum_{T\in\mathfrak T_r}
 \prod_{i=1}^r b_i^{d_i(T)}.
 \label{eq:newV4-tree-weight}
\end{equation}
The inherited definitions give constants \(c_*>0\) and
\(\alpha_*\) such that
\begin{equation}
 b_i\ge {c_*\over\sqrt\alpha}
 \label{eq:newV4-b-lower}
\end{equation}
for \(\alpha\ge\alpha_*\), because
\(A_i\ge1\) and \(s_\alpha\asymp\alpha^{-1}\).

\begin{lemma}[Uniform lower tree mass]
\label{lem:newV4-tree-lower}
For \(r\ge2\),
\begin{equation}
 \mathcal W_r(\mathbf b)
 \ge
 r^{r-2}
 \left({c_*^2\over\alpha}\right)^{r-1}
 =
 {1\over r}
 \left({c_*^2r\over\alpha}\right)^{r-1}.
 \label{eq:newV4-tree-lower}
\end{equation}
\end{lemma}

\begin{proof}
Cayley's formula gives \(r^{r-2}\) labelled trees.  Every tree has
total degree \(2r-2\), so
\eqref{eq:newV4-b-lower} bounds each of its weights below by
\((c_*/\sqrt\alpha)^{2r-2}\).
\end{proof}

\begin{theorem}[All sufficiently high arities are closed trivially]
\label{thm:newV4-high-arity}
Fix \(\delta>0\).  There is \(K_\delta<\infty\) such that for every
\(\alpha\ge\alpha_*\), every integer \(r\ge\delta\alpha\), and all
all-thermal unitary representation carriers,
\begin{equation}
 \boxed{\qquad
 \|\kappa_r^{\nu_\alpha}
 (D^{R_1},\ldots,D^{R_r})\|_{\rm op}
 \le
 r!K_\delta^r\mathcal W_r(\mathbf b).
 \qquad}
 \label{eq:newV4-high-arity}
\end{equation}
Moreover the cumulant admits an exact decomposition over labelled
trees with the same per-tree constant.
\end{theorem}

\begin{proof}
Representation matrices are unitary, so
\Cref{lem:newV4-trivial-cumulant} applies.  If
\(r\ge\delta\alpha\), then
\[
 \mathcal W_r(\mathbf b)
 \ge
 {1\over r}(c_*^2\delta)^{r-1}.
\]
Choose \(K_\delta\) so large that
\[
 K_\delta^r{(c_*^2\delta)^{r-1}\over r}
 \ge e^r
\]
for every \(r\ge2\); finitely many small \(r\) are absorbed by one
further enlargement.  Then
\eqref{eq:newV4-trivial-cumulant} gives
\eqref{eq:newV4-high-arity}.

For an exact tree decomposition, set
\[
 K_{r,T}:=
 {\prod_i b_i^{d_i(T)}\over\mathcal W_r(\mathbf b)}
 \kappa_r.
\]
The coefficients sum to one.  The total estimate implies
\[
 \|K_{r,T}\|_{\rm op}
 \le r!K_\delta^r\prod_i b_i^{d_i(T)}
\]
for every \(T\).
\end{proof}

\section{Stability of cumulants under rare or small modifications}
\label{sec:newV4-rare-stability}

\begin{lemma}[Cumulant perturbation in \(L^1\)]
\label{lem:newV4-cumulant-perturbation}
Let \(X_i,Y_i\) be compatible operator-valued variables on distinct
tensor legs with
\[
 \|X_i\|_{\rm op},\|Y_i\|_{\rm op}\le1
\]
almost surely, and put
\[
 \eta:=\max_i\mathbb E\|X_i-Y_i\|_{\rm op}.
\]
Then
\begin{equation}
 \|\kappa_r(X_1,\ldots,X_r)
 -\kappa_r(Y_1,\ldots,Y_r)\|_{\rm op}
 \le
 \eta\,r!K_0^r
 \label{eq:newV4-cumulant-perturbation}
\end{equation}
for a universal \(K_0\).
\end{lemma}

\begin{proof}
For a block \(B\), telescope its two products:
\[
 \left\|
 \mathbb E\prod_{i\in B}X_i
 -
 \mathbb E\prod_{i\in B}Y_i
 \right\|
 \le |B|\eta.
\]
For a partition \(\pi\), telescope the product of its block
moments.  Its difference is at most
\(\sum_{B\in\pi}|B|\eta=r\eta\).  The moment--cumulant formula and
the estimate in
\Cref{lem:newV4-trivial-cumulant} therefore give
\[
 r\eta\sum_{k=1}^rS(r,k)(k-1)!
 \le\eta r^{r+1}.
\]
Stirling's bound and absorption of the polynomial \(r\) into a
fixed exponential yield
\eqref{eq:newV4-cumulant-perturbation}.
\end{proof}

\begin{corollary}[A common rare event]
\label{cor:newV4-rare-event}
If \(X_i=Y_i\) outside one event \(E\) and both families are
unitarily bounded, then
\[
 \|\kappa_r(\mathbf X)-\kappa_r(\mathbf Y)\|
 \le
 2\mathbb P(E)r!K_0^r.
\]
\end{corollary}

\begin{proof}
Here
\(\mathbb E\|X_i-Y_i\|\le2\mathbb P(E)\); apply
\Cref{lem:newV4-cumulant-perturbation}.
\end{proof}

\begin{lemma}[Exponential rarity is below the tree scale]
\label{lem:newV4-rarity-below-tree}
Fix \(c,\delta>0\).  There is \(K_{c,\delta}<\infty\) such that,
for all \(\alpha\ge\alpha_*\) and
\(2\le r\le\delta\alpha\),
\begin{equation}
 e^{-c\alpha}K_0^r
 \le
 K_{c,\delta}^r\mathcal W_r(\mathbf b).
 \label{eq:newV4-rarity-below-tree}
\end{equation}
\end{lemma}

\begin{proof}
By \Cref{lem:newV4-tree-lower}, it is enough to prove
\[
 e^{-c\alpha}K_0^r
 \le
 {K^r\over r}
 \left({c_*^2r\over\alpha}\right)^{r-1}.
\]
Put \(x=r/\alpha\).  Up to the harmless factor \(r^{-1}\), the
logarithm of the right side is
\[
 \alpha x\log(Kc_*^2x)+O(r).
\]
For \(a>0\),
\[
 \inf_{x>0}x\log(ax)=-{1\over ea}.
\]
Choose \(K\) so large that
\((eKc_*^2)^{-1}<c/4\) and also absorb \(K_0^r\) and the
polynomial factor \(r\).  The right side is then bounded below by
\(e^{-c\alpha/2}\) times the absorbed factors, uniformly for
\(0<x\le\delta\).  This proves
\eqref{eq:newV4-rarity-below-tree}.
\end{proof}

\begin{theorem}[All-order transfer from an exponentially accurate model]
\label{thm:newV4-model-transfer}
Suppose that for every \(\alpha\) and representation \(R\) there is
a unitarily bounded model carrier \(\widetilde D_\alpha^R\), on the
same probability space as \(D^R\), such that
\begin{align}
 \mathbb E\|D^R-\widetilde D_\alpha^R\|_{\rm op}
 &\le e^{-c\alpha},
 \label{eq:newV4-model-error}\\
 \|\kappa_r(\widetilde D_\alpha^{R_1},\ldots,
 \widetilde D_\alpha^{R_r})\|_{\rm op}
 &\le r!K_1^r\mathcal W_r(\mathbf b)
 \quad\text{for }2\le r\le\delta\alpha.
 \label{eq:newV4-model-card}
\end{align}
Then the exact Wilson carriers satisfy the all-order local tree card
for every \(r\ge2\).
\end{theorem}

\begin{proof}
For \(r\le\delta\alpha\), combine
\Cref{lem:newV4-cumulant-perturbation,
lem:newV4-rarity-below-tree} with
\eqref{eq:newV4-model-card}.  For \(r\ge\delta\alpha\), apply
\Cref{thm:newV4-high-arity}.  Increase one fixed exponential base
to cover the boundary and the bounded range
\(\alpha<\alpha_*\).  The proportional allocation used in
\Cref{thm:newV4-high-arity} converts each total estimate into an
exact labelled-tree decomposition.
\end{proof}

\section{The surviving approximation problem}
\label{sec:newV4-approximation}

\begin{definition}[Central analytic approximation card]
\label{def:newV4-CAAC}
A central analytic approximation card (CAAC) is a family of
unitarily bounded functions
\(\widetilde f_{R,\alpha}:\mathbb R^3\to\operatorname{End}(V_R)\)
such that:
\begin{align}
 \mathbb E_{\gamma_3}
 \|D^R(\Phi_\alpha(Y))-\widetilde f_{R,\alpha}(Y)\|_{\rm op}
 &\le e^{-c\alpha},
 \label{eq:newV4-CAAC-error}\\
 \sup_y\|\nabla^d\widetilde f_{R,\alpha}(y)\|_{\rm row}
 &\le d!K_A^d b_R^d
 \qquad(d\ge1).
 \label{eq:newV4-CAAC-derivative}
\end{align}
An integrated BKAR version of
\eqref{eq:newV4-CAAC-derivative} is sufficient.
\end{definition}

\begin{corollary}[CAAC implies the all-order Wilson local card]
\label{cor:newV4-CAAC-implies-card}
The CAAC implies the all-order all-thermal Wilson local marked-tree
card.
\end{corollary}

\begin{proof}
Apply the Gaussian forest theorem
\eqref{eq:newV3-BKAR} to the model carriers.  The proof of
\Cref{prop:newV3-UWTC-implies-card} and the tree-degree factorial
inequality give
\eqref{eq:newV4-model-card}.  The exact radial pushforward identifies
the \(L^1\) error in
\eqref{eq:newV4-CAAC-error} with the Wilson carrier error.  Apply
\Cref{thm:newV4-model-transfer}.
\end{proof}

\begin{proposition}[Exact analytic splicing is impossible]
\label{prop:newV4-no-exact-splice}
One cannot prove the CAAC by choosing a globally real-analytic
carrier which agrees exactly with
\(D^R\circ\Phi_\alpha\) on a nonempty central ball and differs only
past a pre-caustic radius, unless the original radial carrier itself
admits that global analytic continuation.
\end{proposition}

\begin{proof}
The derivative bound
\eqref{eq:newV4-CAAC-derivative} implies real analyticity with a
uniform positive radius.  Two real-analytic functions agreeing on a
nonempty open subset of the connected domain \(\mathbb R^3\) agree
everywhere by the identity theorem.  The original radial transport
has the caustic behavior proved in archived V72 and V77, so an exact
analytic cutoff does not remove it.  The approximation must be
genuinely approximate in \(L^1\), or use the integrated forest card
directly.
\end{proof}

\begin{assessment}[The narrowed all-order gate]
\label{ass:newV4-narrowed-gate}
The all-order problem no longer includes:
\begin{enumerate}[leftmargin=3em]
\item any arity \(r\ge\delta\alpha\), which is closed by
      \Cref{thm:newV4-high-arity};
\item any \(e^{-c\alpha}\) carrier error, which is absorbed by
      \Cref{thm:newV4-model-transfer}; or
\item the Gaussian forest combinatorics, already controlled with
      one global factorial.
\end{enumerate}
What remains is a central analytic approximation or an equivalent
integrated forest estimate for \(2\le r\le\delta\alpha\).  This is
strictly weaker than the disproved pointwise RTDC and SEDE, but it
has not been proved.
\end{assessment}

\begin{programme}[V5: companion audit and central approximation]
\label{prog:newV4-V5}
\begin{enumerate}[leftmargin=3em]
\item Perform the compulsory read-only SM/NS audit and check whether
      either programme has an exponentially accurate
      symmetry-preserving approximation mechanism.
\item Work in the pre-caustic scaled ball and quantify the analytic
      continuation radius of the exact radial carrier.
\item Construct a bounded analytic approximant without exact
      splicing, aiming for \(e^{-c\alpha}\) Gaussian \(L^1\) error.
\item Track representation derivative rows before taking operator
      norms, retaining unitarity on the approximating family.
\item If bounded analytic approximation is impossible, formulate
      and test the integrated BKAR card directly on the caustic
      layer, without a common derivative envelope.
\end{enumerate}
\end{programme}

\begin{firewall}[The arity split does not prove the central card]
\label{fire:newV4-boundary}
\Cref{thm:newV4-high-arity,thm:newV4-model-transfer} are
unconditional, but the CAAC is not proved.  Therefore V4 does not
establish the all-order Wilson local MTA, WUFC, CPM, cutoff removal,
reflection positivity, Osterwalder--Schrader reconstruction, a
continuum Yang--Mills measure, or a positive Hamiltonian gap.  The
full Yang--Mills mass gap has not been proved.
\end{firewall}

\section{Fifth-series V4 result ledger}
\label{sec:newV4-results}

\begin{theorem}[V4 result ledger]
\label{thm:newV4-results}
V4:
\begin{enumerate}[label=\textup{(V4.\arabic*)},leftmargin=5.8em]
\item records the archived radial-transport no-go results and
      prevents reuse of a disproved pointwise route;
\item proves a one-factorial trivial operator-cumulant bound;
\item proves a uniform lower bound on total labelled-tree weight;
\item closes every arity \(r\ge\delta\alpha\) without transport;
\item proves stability of arbitrary-order cumulants under
      exponentially small \(L^1\) carrier errors;
\item proves the all-order transfer theorem from an exponentially
      accurate central model; and
\item isolates the CAAC as the remaining low-relative-arity gate.
\end{enumerate}
\end{theorem}

\begin{proof}
Use
\Cref{lem:newV4-trivial-cumulant,
lem:newV4-tree-lower,
thm:newV4-high-arity,
lem:newV4-cumulant-perturbation,
lem:newV4-rarity-below-tree,
thm:newV4-model-transfer,
cor:newV4-CAAC-implies-card,
prop:newV4-no-exact-splice}.
\end{proof}

\paragraph{Parent-version record.}
V4 was formed from
\texttt{Renewal\_Yang\_Mills\_Mass\_Gap\_Research\_Notes\_V3.tex}.
The V3 parent had \(1\,538\) logical source lines,
\(51\,011\) bytes, and SHA--256
\[
 \texttt{0194be468080c1d7960762f1aa0e125985cb9b7b7bac4c6544b15c22d454b566}.
\]
The next compulsory companion audit is V5.

% =============================================================================
% END APPENDED FIFTH-SERIES V4 MATERIAL
% =============================================================================

% =============================================================================
% BEGIN APPENDED FIFTH-SERIES V5 MATERIAL
% =============================================================================

\chapter{V5: Companion audit and an arity-adapted Wilson density
expansion}
\label{chap:newV5-density}

\section{Compulsory companion audit}
\label{sec:newV5-audit}

The newest active companion files were inspected read-only.  At the
time of inspection they were
\begin{equation}
\begin{split}
&\texttt{Renewal\_SM\_Einstein\_Continuum\_Research\_Notes\_V100.tex},\\
&37\,768\text{ logical source lines},\qquad
1\,322\,035\text{ bytes},\\
&\operatorname{SHA256}=
\texttt{7072c98999eac60f2378c35b8f137f2c40347eb1744894a14e48e4ae4b81f386},
\end{split}
\label{eq:newV5-SM-fingerprint}
\end{equation}
and
\begin{equation}
\begin{split}
&\texttt{Renewal\_NS\_Stability\_Research\_Notes\_V31.tex},\\
&23\,052\text{ logical source lines},\qquad
887\,537\text{ bytes},\\
&\operatorname{SHA256}=
\texttt{f1121b62238ad5491bb7cf03635ea9934942edf573ed8faaa9ee79e1d38a6696}.
\end{split}
\label{eq:newV5-NS-fingerprint}
\end{equation}
Neither file was modified and neither is treated as an instruction.

SM V95 proves a two-regime summability lemma: a low-band Taylor
gain and an owned high-shell decay may be balanced only when their
exponents leave a nonempty crossover interval.  It also insists that
each remote annulus have one owner.  SM V100 proves a Gaussian BKAR
card on a massive reference but explicitly leaves transport to the
coupled positive law open.  NS V31 supplies the same accounting
warning in a different problem: repeatedly charging one remote
mode recreates the critical weight.

\begin{assessment}[Direction selected by the V5 audit]
\label{ass:newV5-audit-direction}
The YM analogue is an order/arity split.  V4 already owns the
high-arity regime \(r\ge\delta\alpha\) and every
\(e^{-c\alpha}\) error.  In the complementary regime, the Wilson
density should be expanded only to the order demanded by the
connected half-edge deficit.  Its coefficients are allowed
Gevrey-one growth \(K^NN!\): when \(N\lesssim r\), the factor
\(\alpha^{-N}N!\) is paid by the existing scalar tree power
\((\sum_i b_i)^{2N}\).  No all-order derivative of the antipodal
transport is needed.
\end{assessment}

\section{Exact scaled density and its two correction series}
\label{sec:newV5-exact-density}

Put \(\epsilon=\alpha^{-1}\) and
\(Y_\alpha=\sqrt\alpha Z\).  Relative to \(\gamma_3\), define the
unnormalized density
\begin{equation}
 \widetilde h_\epsilon(y)
 :=
 {\bf1}_{\{\sqrt\epsilon|y|<\pi\}}
 \exp\{\mathcal W_\epsilon(y)+\mathcal H_\epsilon(y)\},
 \label{eq:newV5-unnormalized}
\end{equation}
where
\begin{align}
 \mathcal W_\epsilon(y)
 &:=
 \epsilon^{-1}
 \{\cos(\sqrt\epsilon|y|)-1\}
 +{|y|^2\over2},
 \label{eq:newV5-W}\\
 \mathcal H_\epsilon(y)
 &:=
 2\log
 \left(
 {\sin(\sqrt\epsilon|y|)
  \over\sqrt\epsilon|y|}
 \right).
 \label{eq:newV5-H}
\end{align}
Then
\begin{equation}
 h_\epsilon(y)
 :=
 {d\mathcal L(Y_\alpha)\over d\gamma_3}(y)
 =
 {\widetilde h_\epsilon(y)\over
  \mathbb E_{\gamma_3}\widetilde h_\epsilon}.
 \label{eq:newV5-normalized-density}
\end{equation}
On \(\sqrt\epsilon|y|<\pi\), the exact series are
\begin{align}
 \mathcal W_\epsilon(y)
 &=
 \sum_{j\ge1}
 {(-1)^{j+1}\epsilon^j|y|^{2j+2}\over(2j+2)!},
 \label{eq:newV5-W-series}\\
 \mathcal H_\epsilon(y)
 &=
 -2\sum_{j\ge1}
 {\zeta(2j)\over j\pi^{2j}}
 \epsilon^j|y|^{2j}.
 \label{eq:newV5-H-series}
\end{align}
The second formula retains the factorial Haar derivatives identified
in archived V86, but its Taylor \emph{coefficients} have only the
displayed \(j^{-1}\pi^{-2j}\) size.  This distinction is essential.

\section{An integrated Gevrey estimate}
\label{sec:newV5-Gevrey}

For a radial polynomial \(P(y)=\sum_\ell c_\ell|y|^{2\ell}\), set
\begin{equation}
 \|P\|_{\mathfrak G}
 :=
 \sum_\ell |c_\ell|\,
 \mathbb E_{\gamma_3}|Y|^{2\ell}.
 \label{eq:newV5-G-norm}
\end{equation}
It dominates \(\|P\|_{L^1(\gamma_3)}\) and is submultiplicative up
to the Gaussian moment convolution recorded below.

\begin{lemma}[Gaussian radial moment convolution]
\label{lem:newV5-moment-convolution}
There is \(C<\infty\) such that for nonnegative integers
\(\ell_1,\ldots,\ell_k\), with \(L=\sum_q\ell_q\),
\begin{equation}
 \mathbb E_{\gamma_3}|Y|^{2L}
 \le
 C^{L+k}L!,
 \qquad
 L!\le k^L\prod_{q=1}^k(\ell_q+1)!.
 \label{eq:newV5-moment-convolution}
\end{equation}
\end{lemma}

\begin{proof}
The radial square is chi-square with three degrees of freedom, so
\[
 \mathbb E|Y|^{2L}
 =
 2^L{\Gamma(L+3/2)\over\Gamma(3/2)}
 \le C^{L+1}L!.
\]
For the second inequality, the multinomial coefficient
\(L!/\prod\ell_q!\) is at most \(k^L\), and
\(\ell_q!\le(\ell_q+1)!\).
\end{proof}

\begin{lemma}[Unnormalized coefficient majorant]
\label{lem:newV5-unnormalized-coefficients}
There are radial even polynomials \(U_j\), \(j\ge0\), with
\[
 U_0=1,\qquad\deg U_j\le4j,
\]
such that the formal central expansion is
\begin{equation}
 \exp(\mathcal W_\epsilon+\mathcal H_\epsilon)
 \sim\sum_{j\ge0}\epsilon^jU_j,
 \label{eq:newV5-U-expansion}
\end{equation}
and
\begin{equation}
 \|U_j\|_{\mathfrak G}\le K_0^j j!.
 \label{eq:newV5-U-majorant}
\end{equation}
\end{lemma}

\begin{proof}
Write
\[
 \mathcal W_\epsilon+\mathcal H_\epsilon
 =
 \sum_{j\ge1}\epsilon^jV_j.
\]
Equations
\eqref{eq:newV5-W-series}--\eqref{eq:newV5-H-series}
show that \(V_j\) is the sum of a degree-\((2j+2)\) monomial with
coefficient \((2j+2)!^{-1}\) and a degree-\(2j\) monomial with
coefficient at most \(C/(j\pi^{2j})\).  The exponential formula gives
\begin{equation}
 U_j=
 \sum_{k=1}^j{1\over k!}
 \sum_{\substack{j_1+\cdots+j_k=j\\j_q\ge1}}
 V_{j_1}\cdots V_{j_k}.
 \label{eq:newV5-U-partitions}
\end{equation}
Every product has degree at most
\[
 2j+2k\le4j.
\]
Apply
\Cref{lem:newV5-moment-convolution}.  For the Haar choices, the
Gaussian \(L!\) is bounded by \(K^jj!\) times the product of the
displayed \(j_q^{-1}\pi^{-2j_q}\) coefficients.  For an action
choice, its \((2j_q+2)!\) denominator is stronger than the
corresponding radial moment factorial.  Mixed choices obey the same
bound after increasing \(K\).  The number of compositions of \(j\)
and the outer \(1/k!\) cost only another exponential base.  Summing
\eqref{eq:newV5-U-partitions} proves
\eqref{eq:newV5-U-majorant}.
\end{proof}

\begin{theorem}[Arity-adapted integrated density expansion]
\label{thm:newV5-density-expansion}
There exist constants \(K,c,\delta>0\) and centered radial even
polynomials \(V_j\), \(j\ge1\), with
\[
 \deg V_j\le4j,
 \qquad
 \mathbb E_{\gamma_3}V_j=0,
 \qquad
 \|V_j\|_{\mathfrak G}\le K^jj!,
\]
such that for every integer \(N\ge1\) satisfying
\begin{equation}
 N\epsilon\le\delta,
 \label{eq:newV5-order-window}
\end{equation}
one has
\begin{equation}
 \boxed{\qquad
 \left\|
 h_\epsilon-
 \left(1+\sum_{j=1}^{N-1}\epsilon^jV_j\right)
 \right\|_{L^1(\gamma_3)}
 \le
 \left(KN\epsilon\right)^N+Ke^{-c/\epsilon}.
 \qquad}
 \label{eq:newV5-density-remainder}
\end{equation}
The constants are independent of \(N\) and \(\epsilon\).
\end{theorem}

\begin{proof}
We separate coefficient growth, integrated Taylor remainder, and
normalization.

\emph{Coefficient growth.}
Lemma~\ref{lem:newV5-unnormalized-coefficients} gives
\(\|U_j\|_{\mathfrak G}\le K_0^jj!\).

\emph{Integrated remainder.}
Expand the two exact series
\eqref{eq:newV5-W-series}--\eqref{eq:newV5-H-series} and then their
exponential, but stop after total \(\epsilon\)-order \(N-1\).
Taylor's integral remainder has the same composition formula as
\eqref{eq:newV5-U-partitions}, with total order at least \(N\).
Integrate termwise and use
\Cref{lem:newV5-moment-convolution}.  Group a term of total order
\(q\ge N\) by the first \(N\) units of its composition.  The
factorial moment bound and \(q!\le q^{q-N}N!\binom qN\), together
with the exact action factorials and the Haar factor
\(\pi^{-2q}\), give the convergent majorant
\[
 \sum_{q\ge N}(K_1q\epsilon)^q
 \le (K_2N\epsilon)^N
\]
after choosing \(\delta\) small.  This is the standard
Gevrey-one remainder estimate obtained directly from the explicit
radial series; it does not assert convergence of the full series at
fixed \(\epsilon\).

The indicator boundary in
\eqref{eq:newV5-unnormalized} contributes only the Wilson and
Gaussian radial tails at \(|y|\asymp\epsilon^{-1/2}\).
The inequality
\[
 1-\cos\theta\ge{2\theta^2\over\pi^2}
\]
and the Gaussian tail bound make their integrated contribution at
most \(Ke^{-c/\epsilon}\).  Polynomial terms of degree at most
\(4N\) have the same tail bound with a smaller \(c\), uniformly
when \(N\epsilon\le\delta\).  Thus
\begin{equation}
 \left\|
 \widetilde h_\epsilon-
 \sum_{j=0}^{N-1}\epsilon^jU_j
 \right\|_1
 \le
 (K_3N\epsilon)^N+Ke^{-c/\epsilon}.
 \label{eq:newV5-unnormalized-remainder}
\end{equation}

\emph{Normalization.}
Let \(a_j=\mathbb E_{\gamma_3}U_j\).  Then
\(|a_j|\le K_0^jj!\) and \(a_0=1\).
Formal inversion of
\(1+\sum_{j\ge1}a_j\epsilon^j\) has coefficients bounded by
\(K_4^jj!\), by the same composition argument.  Multiply its
truncation by
\eqref{eq:newV5-unnormalized-remainder}.  The product coefficients
are the declared \(1,V_1,\ldots,V_{N-1}\); exact normalization makes
every \(V_j\) centered.  The convolution of two Gevrey-one
sequences remains bounded by \(K^jj!\), and the two remainders retain
the right side of
\eqref{eq:newV5-density-remainder}.  This completes the proof.
\end{proof}

\begin{remark}[Consistency with the fixed-order archives]
\label{rem:newV5-fixed-orders}
The first coefficient is
\[
 V_1(y)
 =
 {|y|^4\over24}-{|y|^2\over3}-c_0,
\]
and the second has degree at most eight.  Thus
\Cref{thm:newV5-density-expansion} extends the fixed-order V84--V85
density comparison in the only direction needed here: order may grow
with arity, but only while \(N/\alpha\) stays below a fixed constant.
It does not contradict the V86 no-go for an exponential absolute
majorant of all derivatives of the logarithmic Haar Jacobian.
\end{remark}

\section{Connected source expansion with owned defect order}
\label{sec:newV5-connected-source}

For bounded compatible carriers \(F_i\), introduce formal sources
\(\mathbf t=(t_1,\ldots,t_r)\) and
\begin{equation}
 \mathcal K_\epsilon(\mathbf t)
 :=
 \log
 \mathbb E_{\gamma_3}
 \left[
 h_\epsilon(Y)
 \exp\left(\sum_{i=1}^rt_iF_i(Y)\right)
 \right].
 \label{eq:newV5-source-log}
\end{equation}
The coefficient of \(t_1\cdots t_r\) is the exact Wilson cumulant.

\begin{proposition}[Defect order is assigned once]
\label{prop:newV5-owned-defect}
Insert the truncated density from
\Cref{thm:newV5-density-expansion} into
\eqref{eq:newV5-source-log} and expand formally through
\(\epsilon^{N-1}\).  The coefficient of
\(\epsilon^jt_1\cdots t_r\) is a finite sum of connected Gaussian
cumulants containing the \(r\) carrier vertices and density-defect
vertices whose positive orders sum exactly to \(j\).  No defect order
is charged more than once.
\end{proposition}

\begin{proof}
Regard each term \(\epsilon^jV_j\) as a labelled auxiliary source
and use
\[
 \log\mathbb E_{\gamma_3}
 \exp\left(
 \sum_i t_iF_i+\sum_{j\ge1}u_jV_j
 \right).
\]
Its multivariate Taylor coefficients are joint Gaussian cumulants,
hence connected by the moment--cumulant formula.  Recovering the
density polynomial amounts to substituting the finite polynomial in
\(\epsilon\) for the auxiliary sources and collecting total order.
Ordinary coefficient extraction gives each composition of \(j\)
once.  This is the defect-order analogue of the one-use shell
identity in the SM audit.
\end{proof}

\begin{assessment}[What remains after the density acquisition]
\label{ass:newV5-next-gate}
The density itself now has the required arity-adapted Gevrey
control.  To obtain the Wilson tree card, one must apply the Gaussian
forest formula to the connected carrier-plus-defect cumulants and
show that:
\begin{enumerate}[leftmargin=3em]
\item an order-\(j\) density defect supplies the same two incidence
      powers as \(\alpha^{-j}\);
\item its polynomial derivative rows cost at most \(K^jj!\);
\item all auxiliary defect vertices can be contracted onto one tree
      on the original carrier labels with exponential multiplicity;
      and
\item the remainder in
      \eqref{eq:newV5-density-remainder} lies below the V4 tree scale
      at the chosen \(N\).
\end{enumerate}
The last item is scalar optimization; the first three form the live
connected compiler.
\end{assessment}

\begin{programme}[V6: defect-vertex forest compiler]
\label{prog:newV5-V6}
\begin{enumerate}[leftmargin=3em]
\item Differentiate each radial polynomial \(V_j\) inside the BKAR
      formula and prove a factorial row bound indexed by its incident
      degree.
\item Use the connected half-edge criterion to determine the minimum
      total defect order for every carrier Taylor-degree pattern.
\item Contract auxiliary defect vertices to a canonical spanning
      tree on the carrier labels, recording the multiplicity.
\item Sum compositions of total defect order before taking norms.
\item Optimize \(N\) against \(r/\alpha\) and invoke V4 only after
      the connected density polynomial has been assembled.
\end{enumerate}
\end{programme}

\begin{firewall}[A density expansion is not the Wilson tree card]
\label{fire:newV5-boundary}
\Cref{thm:newV5-density-expansion} proves a uniform integrated
Gevrey expansion in the arity-adapted window, not the defect-vertex
forest compiler.  Therefore V5 does not establish the all-order
Wilson local MTA, WUFC, CPM, cutoff removal, reflection positivity,
Osterwalder--Schrader reconstruction, a continuum Yang--Mills
measure, or a positive Hamiltonian gap.  The full Yang--Mills mass
gap has not been proved.
\end{firewall}

\section{Fifth-series V5 result ledger}
\label{sec:newV5-results}

\begin{theorem}[V5 result ledger]
\label{thm:newV5-results}
V5:
\begin{enumerate}[label=\textup{(V5.\arabic*)},leftmargin=5.8em]
\item performs the compulsory SM V100 and NS V31 audit;
\item writes the exact normalized scaled Wilson density;
\item proves a Gevrey-one majorant for every integrated density
      coefficient;
\item proves a uniform truncation remainder for
      \(N\le\delta\alpha\);
\item reconciles this estimate with the archived Haar-derivative
      no-go; and
\item proves one-use ownership of total density-defect order in the
      connected source expansion.
\end{enumerate}
\end{theorem}

\begin{proof}
Use
\Cref{lem:newV5-moment-convolution,
lem:newV5-unnormalized-coefficients,
thm:newV5-density-expansion,
prop:newV5-owned-defect}.
\end{proof}

\paragraph{Parent-version record.}
V5 was formed from
\texttt{Renewal\_Yang\_Mills\_Mass\_Gap\_Research\_Notes\_V4.tex}.
The V4 parent had \(2\,009\) logical source lines,
\(66\,244\) bytes, and SHA--256
\[
 \texttt{8f64e955a6fd6edda9d26eb063aca52e0f543997a099a2486b935363ade70ca3}.
\]
The next compulsory companion audit is V10.

% =============================================================================
% END APPENDED FIFTH-SERIES V5 MATERIAL
% =============================================================================

% =============================================================================
% BEGIN APPENDED FIFTH-SERIES V6 MATERIAL
% =============================================================================

\chapter{V6: Extended Steiner trees and exact defect-incidence payment}
\label{chap:newV6-Steiner}

\section{Elementary correction vertices}
\label{sec:newV6-elementary}

The normalized polynomials \(V_j\) of V5 are convenient for
integrated remainders but hide the incidence content.  Before
normalization, the logarithmic correction has the elementary form
\begin{equation}
 \mathcal W_\epsilon+\mathcal H_\epsilon
 =
 \sum_{n\ge2}\epsilon^{n-1}\mathsf A_n
 +
 \sum_{n\ge1}\epsilon^n\mathsf H_n,
 \label{eq:newV6-elementary-log}
\end{equation}
where
\begin{align}
 \mathsf A_n(y)
 &:=
 {(-1)^n|y|^{2n}\over(2n)!},
 \label{eq:newV6-An}\\
 \mathsf H_n(y)
 &:=
 -{2\zeta(2n)\over n\pi^{2n}}|y|^{2n}.
 \label{eq:newV6-Hn}
\end{align}
Thus both vertices have polynomial capacity \(k=2n\), while their
activity orders are
\begin{equation}
 q(\mathsf A_n)=n-1={k-2\over2},
 \qquad
 q(\mathsf H_n)=n={k\over2}.
 \label{eq:newV6-activity-orders}
\end{equation}

Expanding a normalized tilted Gaussian source in the elementary
activities gives connected Gaussian cumulants containing:
\begin{enumerate}[leftmargin=3em]
\item \(r\) labelled carrier vertices;
\item \(s\) initially labelled correction vertices, divided by
      \(s!\); and
\item one BKAR spanning tree on all \(r+s\) vertices.
\end{enumerate}
Correction labels are temporary devices for the exponential formula.

\section{Suppressing auxiliary vertices}
\label{sec:newV6-suppression}

\begin{lemma}[Auxiliary-vertex suppression]
\label{lem:newV6-suppression}
Let \(\mathcal T\) be a tree with \(r\ge2\) distinguished carrier
vertices and \(s\) auxiliary vertices.  There is a carrier tree
\(T\in\mathfrak T_r\) obtained by successively deleting auxiliary
leaves and suppressing the remaining auxiliary vertices.  If an
auxiliary vertex \(a\) has degree \(v_a\) at the time it is
suppressed, the total number of new carrier-side incidences required
by the operation is at most
\begin{equation}
 L(\mathcal T):=
 \sum_{a=1}^s(v_a-2)_+.
 \label{eq:newV6-L}
\end{equation}
More precisely, the degrees of \(T\) can be obtained from the
degrees originally carried by distinguished vertices by adding at
most \(L(\mathcal T)\) unit incidences and deleting any surplus
incidences produced by auxiliary leaves.
\end{lemma}

\begin{proof}
Delete every auxiliary leaf and its incident edge.  This can only
remove one incidence from its neighbour and therefore creates no
missing carrier incidence.

Now choose an auxiliary vertex \(a\) of degree \(v\ge2\).  Remove
\(a\) and its \(v\) incident edges.  Connect its \(v\) neighbours by
any tree, using \(v-1\) new edges.  Connectivity is preserved and no
cycle is created: contracting the new neighbour tree to one point
recovers the original tree with \(a\) contracted.  Before the
operation the neighbours had \(v\) incidences toward \(a\); after it
they have \(2(v-1)\) incidences among themselves.  The aggregate
increase is \(v-2\).  It is zero at \(v=2\) and at most
\((v-2)_+\) in all cases.

Iterate until no auxiliary vertex remains.  Every intermediate graph
is a tree, so the final graph is a tree on the distinguished
vertices.  Summing the incidence increases proves the claim.
\end{proof}

\begin{lemma}[Exact degree identity before pruning]
\label{lem:newV6-degree-identity}
If no auxiliary leaves are deleted and the auxiliary degrees in the
original extended tree are \(v_a\), then the aggregate distinguished
degree is
\begin{equation}
 D_{\rm car}
 =
 2r-2-\sum_{a=1}^s(v_a-2).
 \label{eq:newV6-degree-identity}
\end{equation}
\end{lemma}

\begin{proof}
The extended tree has \(r+s-1\) edges and total degree
\(2(r+s-1)\).  Subtracting the auxiliary degree sum gives
\[
 D_{\rm car}
 =
 2(r+s-1)-\sum_av_a
 =
 2r-2-\sum_a(v_a-2).
\]
\end{proof}

\section{Activities pay every missing incidence}
\label{sec:newV6-payment}

\begin{theorem}[Defect-incidence payment]
\label{thm:newV6-incidence-payment}
Consider a nonzero BKAR tree term built from the elementary
corrections \eqref{eq:newV6-An}--\eqref{eq:newV6-Hn}.  Let \(v_a\)
be the degree of correction vertex \(a\).  Its total activity obeys
\begin{equation}
 \prod_{a=1}^s\epsilon^{q_a}
 \le
 \epsilon^{L(\mathcal T)/2},
 \qquad
 L(\mathcal T)=\sum_a(v_a-2)_+.
 \label{eq:newV6-activity-pays}
\end{equation}
Consequently, for any assignment of the \(L(\mathcal T)\) new
incidences in
\Cref{lem:newV6-suppression} to carrier labels
\(i_1,\ldots,i_L\),
\begin{equation}
 \prod_a\epsilon^{q_a}
 \le
 C^{L(\mathcal T)}
 \prod_{\ell=1}^{L(\mathcal T)}b_{i_\ell}.
 \label{eq:newV6-b-payment}
\end{equation}
Thus auxiliary correction vertices never create a deficit in the
power of the carrier tree weight.
\end{theorem}

\begin{proof}
A degree-\(v_a\) derivative of a polynomial of degree \(k_a\)
vanishes unless \(v_a\le k_a\).  For an action vertex,
\[
 q_a={k_a-2\over2}
 \ge{(v_a-2)_+\over2}.
\]
For a Haar vertex,
\[
 q_a={k_a\over2}
 \ge{(v_a-2)_+\over2}.
\]
Multiply these inequalities.  Since \(0<\epsilon\le1\), the larger
activity exponent gives the smaller factor and proves
\eqref{eq:newV6-activity-pays}.

The inherited all-thermal lower bound is
\(\sqrt\epsilon\le Cb_i\) for every carrier label \(i\).
Apply it once for each assigned new incidence to obtain
\eqref{eq:newV6-b-payment}.
\end{proof}

\begin{corollary}[Haar vertices have one reserve pair]
\label{cor:newV6-Haar-reserve}
At a Haar vertex of capacity \(k=2n\), even maximal valence
\(v=k\) uses only \(\epsilon^{n-1}\) for the
\((v-2)\) missing incidences.  One complete factor
\(\epsilon\) remains.  At an action vertex of maximal valence, the
activity is exactly exhausted.
\end{corollary}

\begin{proof}
Substitute \(v=k=2n\) in
\eqref{eq:newV6-activity-orders}.
\end{proof}

\section{The unlabeled auxiliary-tree census}
\label{sec:newV6-census}

\begin{lemma}[One factorial after auxiliary labels are removed]
\label{lem:newV6-auxiliary-census}
Fix \(r\ge2\) and total correction order \(j\le r\).  Sum over:
\begin{enumerate}[leftmargin=3em]
\item \(s\ge0\) elementary correction vertices;
\item action/Haar types and positive activity orders summing to
      \(j\); and
\item spanning trees on the \(r+s\) carrier and temporarily labelled
      correction vertices,
\end{enumerate}
including the exponential factor \(1/s!\).  The number of resulting
decorated tree histories is at most
\begin{equation}
 r!K_{\rm cen}^{\,r+j}
 \label{eq:newV6-census}
\end{equation}
for a universal \(K_{\rm cen}\).
\end{lemma}

\begin{proof}
Every elementary correction has positive activity order, so
\(s\le j\le r\).  For fixed \(s,j\), the number of compositions of
\(j\) into \(s\) positive parts is
\(\binom{j-1}{s-1}\), and each part has at most two types.
Cayley's formula gives \((r+s)^{r+s-2}\) extended labelled trees.
Thus the census is bounded by
\[
 \sum_{s=1}^j
 {2^s\over s!}\binom{j-1}{s-1}
 (r+s)^{r+s-2}.
\]
Because \(s\le r\), Stirling's lower bound gives
\begin{align*}
 {(r+s)^{r+s}\over s!\,r!}
 &\le
 { (2r)^{r+s}e^{r+s}\over s^sr^r}\\
 &=
 (2e)^{r+s}\left({r\over s}\right)^s
 \le K_1^{r+s}.
\end{align*}
The last inequality uses
\((r/s)^s\le e^{r/e}\).  Also
\(\binom{j-1}{s-1}\le2^j\), and the sum has at most \(r\)
terms.  Absorb all fixed exponentials and the final polynomial
factor into \(K_{\rm cen}^{r+j}\), proving
\eqref{eq:newV6-census}.
\end{proof}

\begin{assessment}[Topology and census are no longer bottlenecks]
\label{ass:newV6-topology-closed}
\Cref{thm:newV6-incidence-payment} proves the sharp scalar incidence
balance before norms, and
\Cref{lem:newV6-auxiliary-census} proves that auxiliary labels and
extended-tree choices cost only one carrier factorial times an
exponential.  A second factorial cannot arise from merely labelling
density defects.
\end{assessment}

\section{Action vertices close; the Haar amplitude is isolated}
\label{sec:newV6-amplitudes}

\begin{lemma}[Integrated action-vertex derivative]
\label{lem:newV6-action-amplitude}
For \(0\le v\le2n\),
\begin{equation}
 \mathbb E_{\gamma_3}
 \|\nabla^v\mathsf A_n(Y)\|
 \le K_A^{n+v}.
 \label{eq:newV6-action-amplitude}
\end{equation}
The same estimate, with another exponential base, holds inside a
joint Gaussian moment with any fixed total number of other
polynomial legs after Wick expansion.
\end{lemma}

\begin{proof}
Differentiating \(|y|^{2n}\) \(v\) times produces at most
\(K^v(2n)!/(2n-v)!\) tensor monomials of degree \(2n-v\).
The coefficient \((2n)!^{-1}\) in
\eqref{eq:newV6-An} leaves \((2n-v)!^{-1}\).  Gaussian moments obey
\[
 \mathbb E|Y|^{2n-v}
 \le K^{2n-v}\sqrt{(2n-v)!}.
\]
Their ratio is at most an exponential after discarding the remaining
inverse square-root factorial.  In a joint Gaussian moment, Wick's
formula replaces scalar moments by pairings.  The action factorial
denominator cancels the choices of legs at this vertex, and every
additional pairing costs only a fixed covariance and
finite-dimensional index factor.
\end{proof}

\begin{corollary}[The action-only defect forest compiler]
\label{cor:newV6-action-only}
If the Haar correction \(\mathcal H_\epsilon\) is omitted, every
truncated connected defect expansion of total order \(j\le r\)
obeys the required carrier-tree incidence powers with at most
\(r!K^{r+j}\) multiplicity and amplitude.
\end{corollary}

\begin{proof}
Combine
\Cref{lem:newV6-suppression,
thm:newV6-incidence-payment,
lem:newV6-auxiliary-census,
lem:newV6-action-amplitude}.
Suppress every correction vertex, allocate its activity to the new
carrier incidences, and sum the exponentially many histories.
\end{proof}

For a Haar vertex, differentiation gives the exact radial leading
coefficient
\begin{equation}
 \|\nabla^{2n}\mathsf H_n(0)\|
 =
 {2(2n)!\zeta(2n)\over n\pi^{2n}}.
 \label{eq:newV6-Haar-full-derivative}
\end{equation}
This is the archived V86 factorial.  It cannot be bounded by
\(K^{2n}\) vertexwise.

\begin{definition}[Assembled Haar vertex card]
\label{def:newV6-AHVC}
The assembled Haar vertex card (AHVC) is the assertion that, after:
\begin{enumerate}[leftmargin=3em]
\item multiplying every Haar vertex by its full activity
      \(\epsilon^n\);
\item allocating only
      \(\epsilon^{(v-2)_+/2}\) to topology as in
      \Cref{thm:newV6-incidence-payment};
\item retaining the reserve powers from
      \Cref{cor:newV6-Haar-reserve};
\item multiplying by \(1/s!\) from the correction exponential; and
\item summing Wick pairings and extended trees at fixed total order
      \(j\le r\le\delta\alpha\),
\end{enumerate}
the residual amplitude is at most \(r!K_H^{r+j}\).
\end{definition}

\begin{proposition}[AHVC completes the V5 defect compiler]
\label{prop:newV6-AHVC-sufficient}
If the AHVC holds for some sufficiently small fixed \(\delta>0\),
then the arity-adapted density expansion of V5 gives the local Wilson
tree card for every \(2\le r\le\delta\alpha\).  Together with V4,
this implies the all-order all-thermal one-link local MTA.
\end{proposition}

\begin{proof}
The action-only vertices are controlled by
\Cref{cor:newV6-action-only}; mixed histories use the AHVC for their
Haar subfamily.  The incidence payment produces the carrier-tree
weight, and the auxiliary census produces only one global
factorial.  Truncate at \(N=r\).  The V5 density remainder is
\[
 (Kr/\alpha)^r+Ke^{-c\alpha}.
\]
By
\Cref{lem:newV4-tree-lower,lem:newV4-rarity-below-tree}, this is
below \(K_1^r\mathcal W_r(\mathbf b)\) after increasing \(K_1\);
the harmless factor \(r^2/\alpha\) in the first comparison is
absorbed exponentially.  The cumulant perturbation lemma transfers
the estimate to the exact density.  V4 handles
\(r\ge\delta\alpha\).
\end{proof}

\begin{assessment}[One precise local amplitude gate remains]
\label{ass:newV6-one-gate}
The defect compiler has been reduced to the AHVC.  This is narrower
than the V5 four-part programme:
\begin{enumerate}[leftmargin=3em]
\item correction topology and incidence payment are proved;
\item auxiliary-tree multiplicity is proved exponential times one
      factorial;
\item action vertices are completely controlled; and
\item truncation and high-arity transfer are already owned.
\end{enumerate}
Only the factorial Haar derivative must be paired quantitatively
with its reserve activity, the exponential \(1/s!\), and the global
tree factorial.  Bounding it vertex by vertex would repeat the
disproved V86 route.
\end{assessment}

\begin{programme}[V7: assembled Haar factorial allocation]
\label{prog:newV6-V7}
\begin{enumerate}[leftmargin=3em]
\item For a fixed extended tree, order Haar vertices by leaf
      removal and assign each derivative factorial to the single
      global degree-factorial ledger.
\item Prove
      \(\prod_v d_v!\le(r+s-1)!\) and combine it with \(1/s!\).
\item Use one reserve \(\epsilon\) per Haar vertex to pay the
      remaining binomial factor between \((r+s-1)!/s!\) and \(r!\)
      in the window \(r\le\delta\alpha\).
\item Sum Haar orders only after this allocation, retaining
      \(\pi^{-2n}/n\).
\item Test the extremal one-Haar-star and many-quadratic-Haar paths
      before declaring the AHVC.
\end{enumerate}
\end{programme}

\begin{firewall}[Incidence balance is not amplitude summation]
\label{fire:newV6-boundary}
V6 proves the topology, census, and action-only compiler, but not
the AHVC.  Hence it does not yet prove the all-order Wilson local
MTA, WUFC, CPM, cutoff removal, reflection positivity,
Osterwalder--Schrader reconstruction, a continuum Yang--Mills
measure, or a positive Hamiltonian gap.  The full Yang--Mills mass
gap has not been proved.
\end{firewall}

\section{Fifth-series V6 result ledger}
\label{sec:newV6-results}

\begin{theorem}[V6 result ledger]
\label{thm:newV6-results}
V6:
\begin{enumerate}[label=\textup{(V6.\arabic*)},leftmargin=5.8em]
\item restores the elementary action and Haar correction vertices;
\item proves the auxiliary-tree suppression and exact carrier-degree
      identity;
\item proves that every correction activity pays its missing
      carrier incidences;
\item proves a one-factorial auxiliary-tree census;
\item closes the complete action-only defect compiler; and
\item isolates the assembled Haar vertex card as the sole remaining
      local all-order amplitude gate.
\end{enumerate}
\end{theorem}

\begin{proof}
Use
\Cref{lem:newV6-suppression,
lem:newV6-degree-identity,
thm:newV6-incidence-payment,
lem:newV6-auxiliary-census,
lem:newV6-action-amplitude,
cor:newV6-action-only,
prop:newV6-AHVC-sufficient}.
\end{proof}

\paragraph{Parent-version record.}
V6 was formed from
\texttt{Renewal\_Yang\_Mills\_Mass\_Gap\_Research\_Notes\_V5.tex}.
The V5 parent had \(2\,470\) logical source lines,
\(81\,365\) bytes, and SHA--256
\[
 \texttt{7799d4a189b42e94cbdcd4742274b5a654578dab0966d12e41f408084ad0f973}.
\]
The next compulsory companion audit is V10.

% =============================================================================
% END APPENDED FIFTH-SERIES V6 MATERIAL
% =============================================================================

% =============================================================================
% BEGIN APPENDED FIFTH-SERIES V7 MATERIAL
% =============================================================================

\chapter{V7: Assembled Haar allocation and the all-order local card}
\label{chap:newV7-Haar}

\section{Adaptive Hölder for correlated radial polynomials}
\label{sec:newV7-adaptive-Holder}

\begin{lemma}[Total-degree Hölder allocation]
\label{lem:newV7-total-degree-Holder}
Let \((Y_a)_{a=1}^s\) be jointly Gaussian \(\mathbb R^3\)-valued
variables, each with standard marginal law; their cross-covariances
may be arbitrary.  Let \(\mu_a\ge0\) be integers or half-integers
and put \(M=\sum_a\mu_a\).  Then
\begin{equation}
 \mathbb E\prod_{a=1}^s(1+|Y_a|)^{2\mu_a}
 \le
 K^{M+s}(1+M)^M.
 \label{eq:newV7-total-degree-Holder}
\end{equation}
The estimate is independent of the joint covariance and remains
valid for degenerate Gaussian replicas.
\end{lemma}

\begin{proof}
Discard indices with \(\mu_a=0\).  If \(M>0\), choose
\[
 p_a={M\over\mu_a},
 \qquad
 \sum_a{1\over p_a}=1.
\]
Hölder's inequality uses only the marginal laws.  Standard Gaussian
moments give
\[
 \|(1+|Y|)^{2\mu_a}\|_{L^{p_a}}
 \le
 K^{\mu_a+1}
 (1+\mu_ap_a)^{\mu_a}
 =
 K^{\mu_a+1}(1+M)^{\mu_a}.
\]
Multiply over \(a\).  The case \(M=0\) is immediate.
\end{proof}

\begin{remark}[Why equal-exponent Hölder is forbidden]
\label{rem:newV7-no-equal-Holder}
Using exponent \(s\) at every vertex would create
\(\prod_a(sn_a)^{n_a}\), which can hide an unnecessary
superfactorial.  The allocation in
\Cref{lem:newV7-total-degree-Holder} makes the common moment scale
the total residual degree \(M\), which is already owned by the total
defect order.
\end{remark}

\section{One Haar vertex after topological payment}
\label{sec:newV7-one-Haar}

\begin{lemma}[Haar derivative before activity]
\label{lem:newV7-Haar-derivative}
Let \(0\le v\le2n\) and put
\[
 \mu={2n-v\over2}.
\]
Then
\begin{equation}
 \|\nabla^v\mathsf H_n(y)\|
 \le
 v!K_0^{n+v}(1+|y|)^{2\mu}.
 \label{eq:newV7-Haar-derivative}
\end{equation}
The tensor norm includes all derivative-coordinate rows and has no
ambient-dimension dependence.
\end{lemma}

\begin{proof}
From \eqref{eq:newV6-Hn},
\[
 \left|{2\zeta(2n)\over n\pi^{2n}}\right|
 \le {K_1^n\over n}.
\]
The \(v\)-th derivative of \(|y|^{2n}\) is a sum of tensor monomials
of degree \(2n-v\).  Factoring \(v!\), their total coefficient is
bounded by
\[
 v!\binom{2n}{v}K_2^v
 \le v!K_3^{n+v},
\]
because \(\binom{2n}{v}\le4^n\).  Replacing
\(|y|^{2n-v}\) by \((1+|y|)^{2\mu}\) proves the claim.
\end{proof}

\begin{lemma}[Residual activity exponent]
\label{lem:newV7-residual-activity}
After allocating
\(\epsilon^{(v-2)_+/2}\) to the missing incidences at a Haar vertex,
its unused activity is at least \(\epsilon^\mu\):
\begin{equation}
 \epsilon^{n-(v-2)_+/2}\le\epsilon^\mu,
 \qquad 0<\epsilon\le1.
 \label{eq:newV7-residual-activity}
\end{equation}
For \(v\ge2\), one additional full factor \(\epsilon\) remains.
\end{lemma}

\begin{proof}
If \(v=0\), then \(\mu=n\) and equality holds.  If \(v=1\), then
\(\mu=n-\frac12\) and \(\epsilon^n\le\epsilon^\mu\).
If \(v\ge2\), then
\[
 n-{v-2\over2}
 =
 {2n-v\over2}+1
 =
 \mu+1.
\]
\end{proof}

\section{Proof of the assembled Haar vertex card}
\label{sec:newV7-AHVC}

\begin{theorem}[Assembled Haar vertex card]
\label{thm:newV7-AHVC}
There are \(\delta>0\) and \(K_H<\infty\) such that the AHVC of
\Cref{def:newV6-AHVC} holds whenever
\[
 j\le r\le\delta\alpha.
\]
The constant is uniform in \(\alpha,r,j\), all representations,
and all tensor target dimensions.
\end{theorem}

\begin{proof}
Fix a connected extended BKAR tree with \(s_H\) Haar vertices.
At Haar vertex \(a\), let its order be \(n_a\), tree valence
\(v_a\le2n_a\), and residual polynomial half-degree
\[
 \mu_a={2n_a-v_a\over2}.
\]
Apply
\Cref{lem:newV7-Haar-derivative,
lem:newV7-residual-activity}.  After the topological activity
payment of V6, the product of Haar factors is bounded by
\begin{equation}
 \left(\prod_{a=1}^{s_H}v_a!\right)
 K_0^{\sum_a(n_a+v_a)}
 \epsilon^M
 \prod_{a=1}^{s_H}(1+|Y_a|)^{2\mu_a},
 \qquad
 M:=\sum_a\mu_a.
 \label{eq:newV7-Haar-product}
\end{equation}
We have discarded the additional reserve \(\epsilon\) at every
vertex of valence at least two.

All BKAR replica marginals are standard Gaussian by construction.
Lemma~\ref{lem:newV7-total-degree-Holder} therefore bounds the
expectation of the last product by
\[
 K^{M+s_H}(1+M)^M.
\]
Furthermore
\[
 M\le\sum_an_a\le j\le r,
 \qquad
 \epsilon r\le\delta.
\]
Choose \(\delta\) so small that \(K\delta<1/2\).  Then
\begin{equation}
 \epsilon^M K^M(1+M)^M
 \le
 K_1^M\{\epsilon(1+M)\}^M
 \le K_2^r.
 \label{eq:newV7-residual-moment-paid}
\end{equation}
The cases \(M=0\) and \(M<1\) are covered by enlarging \(K_2\).
Since \(\sum n_a\le j\) and
\(\sum v_a\le2(r+s-1)\), all nonfactorial terms in
\eqref{eq:newV7-Haar-product} are at most \(K_3^{r+j}\).

It remains to place the vertex factorials.  For every tree on
\(r+s\) vertices, the degree-factorial inequality gives
\[
 \prod_{x=1}^{r+s}d_x!\le(r+s-1)!.
\]
Carrier derivative rows have no factorial loss, so inserting their
degree factorials only enlarges the bound.  Hence
\[
 \prod_{a=1}^{s_H}v_a!\le(r+s-1)!.
\]
The elementary correction exponential supplies \(1/s!\), where
\(s\ge s_H\) is the number of all correction vertices.  Because
\(s\le j\le r\),
\begin{align}
 {(r+s-1)!\over s!}
 &=
 (r-1)!\binom{r+s-1}{s}
 \nonumber\\
 &\le
 r!\,2^{r+s}
 \le r!\,4^r.
 \label{eq:newV7-global-factorial}
\end{align}
Action vertices have already been controlled by
\Cref{lem:newV6-action-amplitude}; mixed Gaussian moments use the
same total-degree allocation after including their residual
polynomial degrees.

Finally sum types, order compositions, and extended trees.
\Cref{lem:newV6-auxiliary-census} shows that, after the same
\(1/s!\) is retained, this costs only another exponential in
\(r+j\).  Combining
\eqref{eq:newV7-residual-moment-paid} and
\eqref{eq:newV7-global-factorial} proves the AHVC.
\end{proof}

\begin{assessment}[Extremal configurations]
\label{ass:newV7-extremes}
The proof includes both extremal tests from the V6 programme.
\begin{enumerate}[leftmargin=3em]
\item For one Haar star of valence \(v=2n\), \(M=0\); the full
      derivative factorial is placed in the single global
      \((r+s-1)!\) ledger, and the Haar reserve leaves one unused
      \(\epsilon\).
\item For a path of quadratic Haar vertices, every internal vertex
      has \(n=1,v=2\), again \(M=0\), and each vertex leaves one
      reserve \(\epsilon\).  The exponential \(1/s!\) prevents the
      path orderings from generating a second factorial.
\end{enumerate}
Neither case is visible to a vertexwise exponential derivative
bound, which explains the V86 no-go without obstructing the assembled
card.
\end{assessment}

\section{The all-order all-thermal local Wilson card}
\label{sec:newV7-local-card}

\begin{theorem}[All-order all-thermal one-link MTA]
\label{thm:newV7-all-order-local-MTA}
There is \(K<\infty\) such that, for every
\(\alpha\) in the inherited Wilson concentration regime, every
\(r\ge2\), and all all-thermal irreducible unitary representations,
the exact one-link Wilson cumulant has a decomposition
\begin{equation}
 \kappa_r^{\nu_\alpha}
 (D^{R_1},\ldots,D^{R_r})
 =
 \sum_{T\in\mathfrak T_r}K_{\alpha,T}
 \label{eq:newV7-local-decomposition}
\end{equation}
with
\begin{equation}
 \boxed{\qquad
 \|K_{\alpha,T}\|_{\rm op}
 \le
 r!K^r
 \prod_{i=1}^rb_i^{d_i(T)}.
 \qquad}
 \label{eq:newV7-local-tree-card}
\end{equation}
The constant is independent of \(r,\alpha\), representation
dimensions, and tensor output dimension.
\end{theorem}

\begin{proof}
Choose the \(\delta\) from
\Cref{thm:newV7-AHVC}.  For
\(2\le r\le\delta\alpha\), the AHVC supplies the sole hypothesis
left open in
\Cref{prop:newV6-AHVC-sufficient}.  That proposition combines the
action compiler, exact incidence payment, one-factorial census, V5
density remainder, and V4 cumulant perturbation to give
\eqref{eq:newV7-local-tree-card}.

For \(r\ge\delta\alpha\), use
\Cref{thm:newV4-high-arity}, which gives both the total estimate and
the proportional exact tree decomposition.  Enlarge \(K\) to cover
the finite overlap and the inherited bounded-\(\alpha\) regime.
All component estimates are complete row--column estimates, so the
constant has the stated dimension independence.
\end{proof}

\begin{corollary}[The local all-order obstruction is removed]
\label{cor:newV7-local-obstruction-removed}
The first arrow in
\eqref{eq:newV1-route}, restricted to the all-thermal one-link
sector, now holds at every cumulant arity with the required
factorial-exponential constant.
\end{corollary}

\begin{proof}
Equation \eqref{eq:newV7-local-tree-card} is precisely the
all-thermal local marked-tree atlas.
\end{proof}

\section{Return to the complete selected gate}
\label{sec:newV7-next}

\begin{assessment}[What local sectors remain]
\label{ass:newV7-remaining-local}
The all-thermal sector is now uniform in arity.  The frozen archives
also contain:
\begin{enumerate}[leftmargin=3em]
\item exact high-\(b\) matrix projected-chain estimates when the
      total large parameter exceeds the declared threshold;
\item binary and ternary arbitrary-representation cards;
\item selected quartic and fixed quintic transition-channel cards;
      and
\item scalar anisotropic reductions with explicit unproved
      optimization gates.
\end{enumerate}
The next task is not another all-thermal expansion.  It is to merge
\eqref{eq:newV7-local-tree-card} with those high-channel estimates
without charging a mixed tree twice, and to prove the remaining
anisotropic scalar inequality uniformly in arity.
\end{assessment}

\begin{programme}[V8: thermal/high-channel selected merger]
\label{prog:newV7-V8}
\begin{enumerate}[leftmargin=3em]
\item Restate the exact frozen selected-gate partition in current
      notation and fingerprint its proved versus open cells.
\item Assign every labelled tree one first high-\(b\) vertex or
      declare it all-thermal, so the two local cards are disjoint.
\item Apply the V7 card only to the all-thermal component left after
      conditioning on the high vertex.
\item Reduce the mixed scalar weights to the smallest explicit
      anisotropic inequality and test star and path extrema.
\item Do not enter WUFC until the complete selected local gate is
      genuinely closed.
\end{enumerate}
\end{programme}

\begin{firewall}[A local MTA is not the continuum mass gap]
\label{fire:newV7-boundary}
\Cref{thm:newV7-all-order-local-MTA} proves the all-order
all-thermal one-link local card.  It does not yet prove the complete
mixed selected gate, the global MTA compiler, WUFC, CPM, cutoff
removal, reflection positivity, Osterwalder--Schrader reconstruction,
a continuum Yang--Mills measure, or a positive Hamiltonian gap.
The full Yang--Mills mass gap has not been proved.
\end{firewall}

\section{Fifth-series V7 result ledger}
\label{sec:newV7-results}

\begin{theorem}[V7 result ledger]
\label{thm:newV7-results}
V7:
\begin{enumerate}[label=\textup{(V7.\arabic*)},leftmargin=5.8em]
\item proves the total-degree adaptive Hölder allocation for
      correlated Gaussian replicas;
\item proves the Haar derivative and residual-activity bounds;
\item proves the assembled Haar vertex card;
\item verifies the Haar-star and quadratic-Haar-path extrema;
\item proves the uniform all-order all-thermal one-link local MTA;
      and
\item moves the live programme to the mixed selected-gate merger.
\end{enumerate}
\end{theorem}

\begin{proof}
Use
\Cref{lem:newV7-total-degree-Holder,
lem:newV7-Haar-derivative,
lem:newV7-residual-activity,
thm:newV7-AHVC,
thm:newV7-all-order-local-MTA,
cor:newV7-local-obstruction-removed}.
\end{proof}

\paragraph{Parent-version record.}
V7 was formed from
\texttt{Renewal\_Yang\_Mills\_Mass\_Gap\_Research\_Notes\_V6.tex}.
The V6 parent had \(2\,913\) logical source lines,
\(96\,516\) bytes, and SHA--256
\[
 \texttt{c28e4c7d221f2920dc1f9e46ff7762509a82c9a4e7433f9707054f65328f86aa}.
\]
The next compulsory companion audit is V10.

% =============================================================================
% END APPENDED FIFTH-SERIES V7 MATERIAL
% =============================================================================

% =============================================================================
% BEGIN APPENDED FIFTH-SERIES V8 MATERIAL
% =============================================================================

\chapter{V8: The selected-junction merger and the one-spike test}
\label{chap:newV8-selected}

\section{The exact inherited partition}
\label{sec:newV8-partition}

Fix a selected \(m\)-junction, \(m\ge3\), and retain the notation
\[
 a_i=1+C_2(R_i),\qquad
 P_i=\sum_{f\in\mathcal P_i}(1+C_2(R_f)),\qquad
 \Xi_i=a_i+P_i .
\]
Write
\[
 s=s_\alpha,\qquad
 t_i=sa_i,\qquad p_i=sP_i,\qquad
 x_i={a_i\over\Xi_i}={t_i\over t_i+p_i},
\]
and
\[
 q_i=\prod_{f\in\mathcal P_i}\eta_{R_f}(\alpha),
 \qquad q_P=\prod_{i=1}^m q_i,
 \qquad
 H=\sum_i\sqrt{t_i},\qquad S=\sum_i x_i .
\]
For a labelled tree \(\tau\), let
\[
 B_\tau=\prod_i\Xi_i^{d_i(\tau)-1},
 \qquad
 {\cal T}_\tau
 =\left(\prod_i a_i\right)
  \prod_i x_i^{d_i(\tau)-1}.
\]
The exact Prüfer identities in archived fourth-series V64 are
\begin{align}
 \sum_{\tau\in\mathfrak T_m}{\cal T}_\tau
 &=\left(\prod_i a_i\right)S^{m-2},
 \label{eq:newV8-target-Pruefer}\\
 \sum_{\tau\in\mathfrak T_m}
 \prod_i t_i^{d_i(\tau)/2}
 &=\left(\prod_i\sqrt{t_i}\right)H^{m-2}.
 \label{eq:newV8-local-Pruefer}
\end{align}

Choose the exponential threshold
\[
 D_m=\max\{(4K_{\rm F})^m,L^{m-2}\}\le K_D^m
\]
from archived V63.  Its exact re-fibration proves the desired
selected marked-tree bound outside
\begin{equation}
 \Xi_T\le D_mY,\qquad
 B_\tau\le D_mY^{m-2},
 \qquad Y=\sum_iP_i.
 \label{eq:newV8-private-residual}
\end{equation}
Inside this residual, archived V64 closes every fibre for which
\(P_i\le a_i\) for at least one row.  Thus the residual still needing
the scalar compiler satisfies
\begin{equation}
 p_i>t_i,\qquad 0<x_i<\frac12
 \quad(1\le i\le m).
 \label{eq:newV8-hard-rows}
\end{equation}
These are exact partitions of the same selected operator; no
absolute private-moment comparison is used to move between them.

\begin{record}[Selected-gate status after V7]
\label{rec:newV8-status}
The selected-junction cells are now as follows.
\begin{center}
\begin{tabular}{p{0.38\textwidth}p{0.18\textwidth}p{0.32\textwidth}}
\toprule
cell & status & proof interface\\
\midrule
\(\Xi_T>D_mY\) & \(\MGclosed\) & archived V63 output star\\
\(B_\tau>D_mY^{m-2}\) & \(\MGclosed\) & archived V63 branching star\\
some \(P_i\le a_i\) & \(\MGclosed\) & archived V64 weak-row star\\
hard rows, all thermal, bounded \(t_{\max}/t_{\min}\)
& \(\MGclosed\) & V7 local card plus the proof below\\
hard rows, unbounded anisotropy & \(\MGopen\) &
rowwise Wilson-tail optimization\\
hard rows with a thermally hot selected row & \(\MGopen\) &
mixed hot/private compiler\\
\bottomrule
\end{tabular}
\end{center}
\end{record}

\section{Unconditional closure of the balanced all-thermal cell}
\label{sec:newV8-balanced}

The frozen first private debit gives a fixed \(C_1\) such that
\begin{equation}
 p_iq_i\le C_1,\qquad 0\le q_i\le1.
 \label{eq:newV8-first-debit}
\end{equation}
If \eqref{eq:newV8-hard-rows} holds, then, for
\(0\le\gamma\le1\),
\begin{equation}
 q_i\le C_2\left({x_i\over t_i}\right)^\gamma .
 \label{eq:newV8-fractional-debit}
\end{equation}
Indeed \(p_i=t_i(x_i^{-1}-1)\ge t_i/(2x_i)\), so
\(q_i\le\min\{1,2C_1x_i/t_i\}\), and
\(\min\{1,u\}\le u^\gamma\).

\begin{theorem}[Balanced selected-junction closure]
\label{thm:newV8-balanced}
Fix \(B<\infty\).  Suppose that
\eqref{eq:newV8-private-residual} and
\eqref{eq:newV8-hard-rows} hold, that all selected rows are thermal,
\[
 0<t_i<1,
\]
and that
\begin{equation}
 t_{\max}\le B\,t_{\min}.
 \label{eq:newV8-balance}
\end{equation}
Then both final-resolvent branches of the complete selected
\(m\)-junction obey the summed marked-tree target with constant
\[
 m!K_B^m.
\]
Thus the balanced all-thermal cell in
\Cref{rec:newV8-status} is unconditional.
\end{theorem}

\begin{proof}
By \Cref{thm:newV7-all-order-local-MTA}, the selected one-link
cumulant has an exact tree decomposition whose summed operator norm
is at most
\begin{equation}
 m!K_0^m
 \left(\prod_i\sqrt{t_i}\right)H^{m-2}.
 \label{eq:newV8-summed-local}
\end{equation}
The summed target is
\begin{equation}
 s^{-m}\left(\prod_it_i\right)S^{m-2}
 \label{eq:newV8-summed-target}
\end{equation}
by \eqref{eq:newV8-target-Pruefer}.

For low final output, the payer is at most \(C/s\).  Put
\(\gamma=(m-2)/m\).  From
\eqref{eq:newV8-fractional-debit} and the
arithmetic--geometric mean inequality,
\[
 q_P\le C_2^m\prod_i(x_i/t_i)^\gamma,
 \qquad
 \prod_i x_i^\gamma
 \le (S/m)^{m-2}.
\]
After dividing the paid version of
\eqref{eq:newV8-summed-local} by
\eqref{eq:newV8-summed-target}, the remaining scalar ratio is at
most
\begin{equation}
 C^m
 {s^{m-1}(H/m)^{m-2}
  \over\prod_i t_i^{1/2+\gamma}}.
 \label{eq:newV8-low-ratio-balanced}
\end{equation}
Since \(H/m\le\sqrt{t_{\max}}\), \(s\le t_{\min}\), and
\(\sum_i(1/2+\gamma)=3m/2-2\), the last expression is at most
\((C\sqrt B)^m\).  The powers of \(t_{\min}\) cancel exactly:
\[
 (m-1)+{m-2\over2}
 = {3m\over2}-2.
\]

For high final output, the residual inequality gives
\[
 \Xi_T\le D_mY={D_m\over s}\sum_kp_k.
\]
For a fixed summand use \(p_kq_k\le C_1\), and on the other
\(m-1\) rows use
\eqref{eq:newV8-fractional-debit} with
\(\delta=(m-2)/(m-1)\).  Then
\[
 p_kq_P
 \le C^m\prod_{i\ne k}(x_i/t_i)^\delta,
 \qquad
 \prod_{i\ne k}x_i^\delta
 \le \left({S\over m-1}\right)^{m-2}.
\]
After summing \(k\), the scalar ratio is at most
\[
 C^mD_m
 {s^{m-1}t_{\max}^{(m-2)/2}
  \over t_{\min}^{m/2+m-2}}.
\]
The same identity of exponents cancels the \(t_{\min}\) powers.
The factors \(D_m\le K_D^m\), \(B^{(m-2)/2}\), the \(m\) choices
of \(k\), and the harmless replacement of \(m-1\) by \(m\) are all
exponential per input.  This proves the high branch and the theorem.
\end{proof}

\begin{corollary}[A rigorous reduction of the all-thermal residual]
\label{cor:newV8-reduction}
In the all-thermal selected sector, the only unclosed part of
\eqref{eq:newV8-private-residual} consists of hard rows for which
\[
 \sup_m {t_{\max}\over t_{\min}}=\infty .
\]
The local connected operator estimate is no longer an obstruction
there; only its scalar coupling to the private bank remains.
\end{corollary}

\begin{proof}
The local card is \Cref{thm:newV7-all-order-local-MTA}.  The
output-dominant, branching-dominant, and weak-row complements were
removed by the exact inherited partition, and
\Cref{thm:newV8-balanced} removes every fixed balance class.
\end{proof}

\section{The exact scalar ratios}
\label{sec:newV8-ratios}

For later optimization, it is useful not to hide the two remaining
quantities behind the informal term \emph{compiler}.  Division of
the paid local sum \eqref{eq:newV8-summed-local} by the target
\eqref{eq:newV8-summed-target} gives, up to a fixed exponential base,
\begin{align}
 {\cal R}^{\rm lo}_m
 &:=
 {q_Ps^{m-1}H^{m-2}
  \over(\prod_i\sqrt{t_i})S^{m-2}},
 \label{eq:newV8-low-ratio}\\
 {\cal R}^{\rm hi}_m
 &:=
 D_m\left(\sum_i p_i\right){\cal R}^{\rm lo}_m .
 \label{eq:newV8-high-ratio}
\end{align}
The desired anisotropic lemma is
\begin{equation}
 {\cal R}^{\rm lo}_m+{\cal R}^{\rm hi}_m\le K^m
 \label{eq:newV8-anisotropic-target}
\end{equation}
for exact private Wilson multipliers.  The false route exposed in
archived V65 replaces those multipliers only by
\(p_iq_i\le C_1\); that weaker axiom does not imply
\eqref{eq:newV8-anisotropic-target}.

\section{Uniform repair of the sharp one-spike counterpacket}
\label{sec:newV8-one-spike}

The archived V66 \(SU(2)\) Bessel estimate says that
\[
 q_i\le
 \exp\{-c\Psi_s(p_i)\},
 \qquad
 \Psi_s(p)={1\over s}h(c_*\sqrt{sp}),
\]
where
\[
 h(r)=r\operatorname{arsinh}r-\sqrt{1+r^2}+1.
\]
The following strengthens the old fixed-\(m\) regression test to the
uniform all-arity constant required by the selected gate.

\begin{lemma}[Exponential absorption at the anisotropic scale]
\label{lem:newV8-absorption}
For every \(c>0\), \(m\ge3\), \(0<s\le1\), and \(0<\theta\le1\),
\begin{equation}
 e^{-c/s}
 \left(1+{\sqrt\theta\over m\sqrt s}\right)^{m-2}
 \le \exp\{1/(4c)\}.
 \label{eq:newV8-absorption}
\end{equation}
Moreover
\begin{equation}
 (s^{-1}+m)e^{-c/(2s)}\le C_c\,2^m.
 \label{eq:newV8-polynomial-absorption}
\end{equation}
\end{lemma}

\begin{proof}
Put \(y=\sqrt\theta/(m\sqrt s)\).  Since \(\theta\le1\),
\[
 -{c\over s}+(m-2)\log(1+y)
 \le -cm^2y^2+my
 \le {1\over4c}.
\]
This proves the first assertion.  For the second,
\(s^{-1}e^{-c/(2s)}\) is bounded on \((0,1]\), while
\(m\le2^m\).
\end{proof}

\begin{theorem}[Uniform \(SU(2)\) one-spike closure]
\label{thm:newV8-one-spike}
Fix \(0<\theta_0<\theta_1<1\).  Consider the exact anisotropic
packet
\begin{equation}
 t_1=\theta\in[\theta_0,\theta_1],\qquad
 t_i=s\ (i\ge2),\qquad
 x_i=s\ (1\le i\le m),
 \label{eq:newV8-one-spike-data}
\end{equation}
with \(0<s\le s_0<1/2\), and suppose the private rows are actual
\(SU(2)\) Wilson banks.  Then there is a fixed \(K\), independent of
\(m,s,\theta\), such that
\begin{equation}
 {\cal R}^{\rm lo}_m+{\cal R}^{\rm hi}_m\le K^m.
 \label{eq:newV8-one-spike-bound}
\end{equation}
\end{theorem}

\begin{proof}
The first row has
\[
 p_1=\theta{s^{-1}-1},\qquad
 sp_1=\theta(1-s).
\]
After decreasing \(s_0\), the compact bounds on \(\theta\) and the
positivity of \(h\) away from zero give
\[
 q_1\le e^{-c_0/s}.
\]
We only use \(q_i\le1\) on the other rows.  Directly from
\eqref{eq:newV8-one-spike-data},
\[
 H=\sqrt\theta+(m-1)\sqrt s,\qquad
 S=ms,\qquad
 \prod_i\sqrt{t_i}=\sqrt\theta\,s^{(m-1)/2}.
\]
Therefore
\begin{align*}
 {\cal R}^{\rm lo}_m
 &\le
 {e^{-c_0/s}\over\sqrt\theta}\,
 s^{-(m-3)/2}
 \left({H\over m}\right)^{m-2}\\
 &\le
 {\sqrt s\over\sqrt{\theta_0}}\,
 e^{-c_0/s}
 \left(1+{\sqrt\theta\over m\sqrt s}\right)^{m-2}.
\end{align*}
The first part of \Cref{lem:newV8-absorption} bounds this uniformly.

Also
\[
 \sum_i p_i
 =\theta(s^{-1}-1)+(m-1)(1-s)
 \le s^{-1}+m .
\]
Split \(e^{-c_0/s}\) into two equal factors.  Use the first part of
\Cref{lem:newV8-absorption} on one and
\eqref{eq:newV8-polynomial-absorption} on the other.  Since
\(D_m\le K_D^m\), this proves
\({\cal R}^{\rm hi}_m\le K^m\) as well.
\end{proof}

\begin{assessment}[What the one-spike theorem establishes]
\label{ass:newV8-one-spike-scope}
The packet \eqref{eq:newV8-one-spike-data} is the sharp geometry of
the archived first-debit countermodel: one selected thermal scale is
macroscopic, the others are at the regulator floor, and every
private ratio is hard.  The theorem shows that the exact Wilson tail
repairs this obstruction with an exponential-per-input constant,
including the high-output payer.  It does not prove the arbitrary
multi-spike optimization, because several scale clusters can share
both \(H\) and \(S\) without satisfying
\eqref{eq:newV8-one-spike-data}.
\end{assessment}

\section{V9 programme and claim boundary}
\label{sec:newV8-next}

\begin{programme}[V9: multiscale anisotropic rows]
\label{prog:newV8-V9}
\begin{enumerate}[leftmargin=3em]
\item Order the thermal scales and divide them into dyadic clusters
      without paying the number of nonempty scale bins.
\item Assign each cluster its exact share of \(S=\sum x_i\), so a
      row with a small private ratio cannot be charged as though all
      of \(S\) were small.
\item Apply the \(SU(2)\) Wilson rate only to scale clusters whose
      \(x\)-share is deficient; use
      \Cref{lem:newV8-absorption} at the cluster cardinality.
\item Prove both ratios
      \eqref{eq:newV8-low-ratio}--\eqref{eq:newV8-high-ratio} with
      one common exponential base, or exhibit a genuine exact
      Wilson packet which prevents it.
\item Keep thermally hot selected rows separate until the
      all-thermal scalar lemma is complete.
\end{enumerate}
\end{programme}

\begin{firewall}[The selected gate and mass gap remain open]
\label{fire:newV8-boundary}
\Cref{thm:newV8-balanced} closes the bounded-anisotropy all-thermal
selected residual, and \Cref{thm:newV8-one-spike} closes the sharp
one-spike \(SU(2)\) counterpacket uniformly in arity.  V8 does not
prove the arbitrary anisotropic scalar estimate
\eqref{eq:newV8-anisotropic-target}, the hard thermally hot cell, the
complete selected gate, global MTA, WUFC, CPM, cutoff removal,
reflection positivity, Osterwalder--Schrader reconstruction, a
continuum Yang--Mills measure, or a positive Hamiltonian gap.  The
full Yang--Mills mass gap has not been proved.
\end{firewall}

\begin{theorem}[V8 result ledger]
\label{thm:newV8-results}
V8:
\begin{enumerate}[label=\textup{(V8.\arabic*)},leftmargin=5.8em]
\item reconstructs the exact selected-junction partition;
\item combines the V7 local card with the archived scalar compiler
      to close the balanced all-thermal residual unconditionally;
\item writes the low- and high-output anisotropic ratios explicitly;
\item proves a regulator-scale exponential absorption lemma; and
\item closes the exact \(SU(2)\) one-spike counterpacket uniformly
      in arity in both output branches.
\end{enumerate}
\end{theorem}

\begin{proof}
Use
\Cref{thm:newV8-balanced,
cor:newV8-reduction,
lem:newV8-absorption,
thm:newV8-one-spike}.
\end{proof}

\paragraph{Parent-version record.}
V8 was formed from
\texttt{Renewal\_Yang\_Mills\_Mass\_Gap\_Research\_Notes\_V7.tex}.
The V7 parent had \(3\,299\) logical source lines,
\(108\,866\) bytes, and SHA--256
\[
 \texttt{ae00ee159af2f85646ede3d78fc859dc6d941bfd530995014eb447f9fd678797}.
\]
The next compulsory companion audit is V10.

% =============================================================================
% END APPENDED FIFTH-SERIES V8 MATERIAL
% =============================================================================

% =============================================================================
% BEGIN APPENDED FIFTH-SERIES V9 MATERIAL
% =============================================================================

\chapter{V9: Weighted-load closure of all thermal anisotropy}
\label{chap:newV9-anisotropy}

\section{A usable consequence of the exact Wilson rate}
\label{sec:newV9-rate}

Retain the \(SU(2)\) private-row estimate from archived
fourth-series V66:
\[
 q_i\le\exp\{-c\Psi_s(p_i)\},
 \qquad
 \Psi_s(p)={1\over s}h(c_*\sqrt{sp}),
\]
with
\[
 h(r)=r\operatorname{arsinh}r-\sqrt{1+r^2}+1.
\]

\begin{lemma}[Two-regime lower bound for the Bessel rate]
\label{lem:newV9-rate-lower}
There is \(c_0>0\) such that, for every \(0<s\le1\) and \(p\ge0\),
\begin{equation}
 \Psi_s(p)\ge
 c_0\min\left\{p,\sqrt{p/s}\right\}.
 \label{eq:newV9-rate-lower}
\end{equation}
Consequently every private row obeys
\begin{equation}
 q_i\le
 \exp\left\{-c_1
 \min\left(p_i,\sqrt{p_i/s}\right)\right\}.
 \label{eq:newV9-row-tail}
\end{equation}
\end{lemma}

\begin{proof}
Since \(h'(r)=\operatorname{arsinh}r\), for \(0\le r\le1\),
\[
 h(r)=\int_0^r\operatorname{arsinh}v\,dv
 \ge {1\over\sqrt2}\int_0^r v\,dv
 ={r^2\over2\sqrt2}.
\]
For \(r\ge1\), positivity of \(h(1)\) and monotonicity of
\(\operatorname{arsinh}\) give
\[
 h(r)
 =h(1)+\int_1^r\operatorname{arsinh}v\,dv
 \ge h(1)+(r-1)\operatorname{arsinh}1
 \ge c_2r.
\]
Thus \(h(r)\ge c_3\min(r^2,r)\).  Substitute
\(r=c_*\sqrt{sp}\) and absorb fixed powers of \(c_*\):
\[
 {1\over s}\min\{sp,\sqrt{sp}\}
 =\min\{p,\sqrt{p/s}\}.
\]
The last assertion follows from the archived bank-tail theorem.
\end{proof}

\section{The rowwise fractional-load transform}
\label{sec:newV9-row}

The next lemma is the analytic core.  It retains the crossover at
\(p=s^{-1}\); replacing it by a polynomial Casimir moment would
reintroduce the false second factorial from archived V62.

\begin{lemma}[Uniform fractional-load transform]
\label{lem:newV9-load}
Suppose
\[
 0<s\le1,\qquad t\ge s,\qquad p\ge0,\qquad
 q\le e^{-c\min\{p,\sqrt{p/s}\}}.
\]
For every \(\gamma\ge0\), put
\[
 A_\gamma=\gamma+s\gamma^2.
\]
Then
\begin{equation}
 q\left(1+{p\over t}\right)^\gamma
 \le
 e^{C\gamma}
 \left(1+{A_\gamma\over t}\right)^\gamma ,
 \label{eq:newV9-load}
\end{equation}
where \(C\) depends only on \(c\).
\end{lemma}

\begin{proof}
The assertion is immediate for \(\gamma=0\), so assume
\(\gamma>0\).  Since
\[
 1+{p\over t}
 \le
 \left(1+{A_\gamma\over t}\right)
 \left(1+{p\over A_\gamma}\right),
\]
it is enough to prove
\begin{equation}
 \sup_{p\ge0}
 \left\{
 \gamma\log\left(1+{p\over A_\gamma}\right)
 -c\min\left(p,\sqrt{p/s}\right)
 \right\}
 \le C\gamma.
 \label{eq:newV9-Legendre}
\end{equation}

If \(p\le s^{-1}\), the minimum is \(p\).  Put
\(z=p/A_\gamma\).  Since \(A_\gamma\ge\gamma\),
\[
 \gamma\log(1+z)-cp
 \le
 \gamma\{\log(1+z)-cz\}
 \le C_1\gamma.
\]
If \(p\ge s^{-1}\), put
\[
 y={\sqrt{p/s}\over\gamma},
 \qquad p=s\gamma^2y^2.
\]
Because \(A_\gamma\ge s\gamma^2\),
\[
 \gamma\log\left(1+{p\over A_\gamma}\right)
 -c\sqrt{p/s}
 \le
 \gamma\{\log(1+y^2)-cy\}
 \le C_2\gamma.
\]
Both scalar suprema are finite.  This proves
\eqref{eq:newV9-Legendre} and hence
\eqref{eq:newV9-load}.
\end{proof}

\begin{corollary}[A split tail pays one high-output mass]
\label{cor:newV9-first-mass}
Under the hypotheses of \Cref{lem:newV9-load},
\begin{equation}
 pq^{1/2}\le C.
 \label{eq:newV9-first-mass}
\end{equation}
Moreover \(q^{1/2}\) satisfies the hypothesis of
\Cref{lem:newV9-load} with \(c\) replaced by \(c/2\).
\end{corollary}

\begin{proof}
In the range \(p\le s^{-1}\),
\(pq^{1/2}\le pe^{-cp/2}\le C\).  In the range
\(p\ge s^{-1}\), write \(y=\sqrt{p/s}\).  Then
\[
 pq^{1/2}\le sy^2e^{-cy/2}.
\]
Here \(p\ge s^{-1}\) implies \(y\ge s^{-1}\), hence
\(s\le y^{-1}\); the last display is bounded by
\(ye^{-cy/2}\).  The second assertion is immediate.
\end{proof}

\section{Weighted AM--GM removes every scale bin}
\label{sec:newV9-weighted}

\begin{theorem}[Anisotropy-uniform low-output scalar bound]
\label{thm:newV9-low}
Let \(m\ge3\), \(0<s\le t_i<1\), \(p_i\ge0\), and
\[
 x_i={t_i\over t_i+p_i},\qquad
 q_i\le
 \exp\left\{-c
 \min\left(p_i,\sqrt{p_i/s}\right)\right\}.
\]
With
\[
 H=\sum_i\sqrt{t_i},\qquad S=\sum_ix_i,\qquad
 q_P=\prod_iq_i,
\]
the exact low-output ratio from V8 satisfies
\begin{equation}
 {\cal R}^{\rm lo}_m
 =
 {q_Ps^{m-1}H^{m-2}
  \over(\prod_i\sqrt{t_i})S^{m-2}}
 \le K^m.
 \label{eq:newV9-low}
\end{equation}
The constant is independent of \(m,s\), the anisotropy of the
\(t_i\), the private masses, and the support sizes.
\end{theorem}

\begin{proof}
Put \(n=m-2\) and choose the probability weights
\[
 w_i={\sqrt{t_i}\over H},
 \qquad
 \gamma_i=nw_i.
\]
Thus \(\sum_i\gamma_i=n\).  Weighted AM--GM applied to
\(x_i/w_i\) gives
\[
 S=\sum_iw_i{x_i\over w_i}
 \ge\prod_i\left({x_i\over w_i}\right)^{w_i},
\]
and hence
\begin{equation}
 S^{-n}\le
 \prod_i w_i^{\gamma_i}x_i^{-\gamma_i}.
 \label{eq:newV9-weighted-AMGM}
\end{equation}
Since \(H^n\prod_iw_i^{\gamma_i}
=\prod_it_i^{\gamma_i/2}\), substitution in the scalar ratio yields
\begin{equation}
 {\cal R}^{\rm lo}_m
 \le
 s^{m-1}\prod_i
 \left[
 t_i^{(\gamma_i-1)/2}
 q_ix_i^{-\gamma_i}
 \right].
 \label{eq:newV9-factorized}
\end{equation}

Apply \Cref{lem:newV9-load} row by row and use
\(x_i^{-1}=1+p_i/t_i\).  Up to \(e^{Cn}\), the \(i\)-th factor
carrying exponent \(\gamma_i\) is
\[
 \left[
 \sqrt{t_i}
 +{\gamma_i\over\sqrt{t_i}}
 +{s\gamma_i^2\over\sqrt{t_i}}
 \right]^{\gamma_i}.
\]
The chosen loads make the two potentially dangerous quotients
uniform across all rows:
\[
 {\gamma_i\over\sqrt{t_i}}={n\over H},
 \qquad
 {s\gamma_i^2\over\sqrt{t_i}}
 ={sn^2\sqrt{t_i}\over H^2}.
\]
Because \(t_i<1\), \(H\ge m\sqrt s\), and \(n\le m\),
\[
 \sqrt{t_i}+{n\over H}
 +{sn^2\sqrt{t_i}\over H^2}
 \le 2+{1\over\sqrt s}
 \le {3\over\sqrt s}.
\]
Using \(\sum_i\gamma_i=n\) in
\eqref{eq:newV9-factorized} therefore gives
\[
 {\cal R}^{\rm lo}_m
 \le
 C^m s^{m-1}
 \left(\prod_it_i^{-1/2}\right)s^{-n/2}.
\]
Finally \(t_i\ge s\), so
\[
 s^{m-1}
 \left(\prod_it_i^{-1/2}\right)s^{-n/2}
 \le
 s^{m-1-m/2-(m-2)/2}=1.
\]
This proves \eqref{eq:newV9-low}.
\end{proof}

\begin{theorem}[Anisotropy-uniform high-output scalar bound]
\label{thm:newV9-high}
Under the hypotheses of \Cref{thm:newV9-low}, if
\(D_m\le K_D^m\), then
\begin{equation}
 {\cal R}^{\rm hi}_m
 =
 D_m\left(\sum_i p_i\right){\cal R}^{\rm lo}_m
 \le K_1^m.
 \label{eq:newV9-high}
\end{equation}
\end{theorem}

\begin{proof}
Expand the mass sum.  In its \(k\)-th term write
\[
 p_kq_P
 =(p_kq_k^{1/2})
  \left(q_k^{1/2}\prod_{i\ne k}q_i\right).
\]
The first factor is uniformly bounded by
\Cref{cor:newV9-first-mass}.  Every multiplier in the second factor
satisfies the same two-regime tail with \(c\) replaced, if necessary,
by \(c/2\).  Therefore the proof of
\Cref{thm:newV9-low} applies verbatim and bounds each \(k\)-term by
\(C^m\).  Sum the \(m\) choices of \(k\), use
\(m\le2^m\), and absorb \(D_m\le K_D^m\).
\end{proof}

\begin{remark}[Why no dyadic-bin count appears]
\label{rem:newV9-no-bins}
The load
\(\gamma_i=(m-2)\sqrt{t_i}/H\) is a continuous Prüfer allocation.
It assigns to a scale precisely its share of the local tree
polynomial.  The identities
\[
 {\gamma_i\over\sqrt{t_i}}={m-2\over H},
 \qquad
 {s\gamma_i^2\over\sqrt{t_i}}
 ={s(m-2)^2\sqrt{t_i}\over H^2}
\]
replace the planned dyadic decomposition, so neither
\(\log(t_{\max}/t_{\min})\) nor the number of occupied scale bins is
ever paid.
\end{remark}

\section{Closure of the all-thermal selected gate}
\label{sec:newV9-gate}

\begin{theorem}[Complete all-order all-thermal \(SU(2)\) selected gate]
\label{thm:newV9-all-thermal-selected}
For the exact \(SU(2)\) Wilson multiplier and every selected
\(m\)-junction whose selected rows are all thermal, the complete
paid junction obeys the required summed marked-tree bound with
constant
\begin{equation}
 m!K_{\rm sel}^m,
 \label{eq:newV9-selected-bound}
\end{equation}
uniformly in \(m,\alpha\), all representation labels, output
dimensions, private supports, and thermal anisotropy.
\end{theorem}

\begin{proof}
Use the exact partition recorded in V8.  The output-dominant and
branching-dominant cells are the proved star re-fibrations of
archived V63.  A row with \(P_i\le a_i\) is the proved weak-row star
of archived V64.  In the remaining hard-row residual,
\Cref{thm:newV7-all-order-local-MTA} supplies the exact local
tree decomposition.  Its low- and high-output scalar quotients are
exactly \eqref{eq:newV8-low-ratio} and
\eqref{eq:newV8-high-ratio}; they are bounded by
\Cref{thm:newV9-low,thm:newV9-high}.  All partitions are exact
decompositions or reassignments of the same selected junction, so
their disjoint sum proves \eqref{eq:newV9-selected-bound}.
\end{proof}

\begin{assessment}[New live boundary]
\label{ass:newV9-boundary}
The unbounded all-thermal anisotropy listed as open in V8 is now
closed for \(SU(2)\).  The remaining selected local sector is the
hard residual containing at least one thermally hot selected row.
Archived V61 bounds the hot local cumulant by a star, but archived
V64 correctly observed that this alone does not couple the same star
to an arbitrary hard private bank.  The next proof must merge the
hot displacement center with the rowwise Wilson tail without using
the all-thermal tree card on that hot row.
\end{assessment}

\section{V10 programme and claim boundary}
\label{sec:newV9-next}

\begin{programme}[V10: companion audit and the hot/private merger]
\label{prog:newV9-V10}
\begin{enumerate}[leftmargin=3em]
\item Perform the mandatory read-only audit of the newest active
      Standard-Model--Einstein and Navier--Stokes notes.
\item Restate the exact archived V61 hot-star estimate, including
      the center choice and both final-output payers.
\item Split the hard private multiplier at the hot center and apply
      the row transform \eqref{eq:newV9-load} only to the remaining
      selected rows.
\item Determine whether a maximal-\(t_i\) star has enough target
      mass to pay the center's private ratio; if not, introduce an
      exact hot-center re-fibration rather than a higher Casimir
      moment.
\item Do not enter the global WUFC compiler until the hard
      hot/private cell is closed.
\end{enumerate}
\end{programme}

\begin{firewall}[The continuum mass gap is not yet proved]
\label{fire:newV9-boundary}
\Cref{thm:newV9-all-thermal-selected} proves the complete all-order
all-thermal \(SU(2)\) selected gate.  It does not prove the hard
hot/private cell, the complete selected gate, global MTA, WUFC,
CPM, cutoff removal, reflection positivity,
Osterwalder--Schrader reconstruction, a continuum Yang--Mills
measure, or a positive Hamiltonian gap.  The full Yang--Mills mass
gap has not been proved.
\end{firewall}

\begin{theorem}[V9 result ledger]
\label{thm:newV9-results}
V9:
\begin{enumerate}[label=\textup{(V9.\arabic*)},leftmargin=5.8em]
\item derives a two-regime lower bound from the exact \(SU(2)\)
      Bessel rate;
\item proves the uniform fractional-load transform;
\item proves anisotropy-uniform low- and high-output scalar bounds;
\item removes every scale-bin counting loss by a weighted AM--GM
      allocation; and
\item closes the complete all-order all-thermal \(SU(2)\) selected
      gate.
\end{enumerate}
\end{theorem}

\begin{proof}
Use
\Cref{lem:newV9-rate-lower,
lem:newV9-load,
cor:newV9-first-mass,
thm:newV9-low,
thm:newV9-high,
thm:newV9-all-thermal-selected}.
\end{proof}

\paragraph{Parent-version record.}
V9 was formed from
\texttt{Renewal\_Yang\_Mills\_Mass\_Gap\_Research\_Notes\_V8.tex}.
The V8 parent had \(3\,743\) logical source lines,
\(122\,542\) bytes, and SHA--256
\[
 \texttt{8f27948afc0c567038ac2c8ab6aed754dca33f3211f515de207a5a61bd1b052c}.
\]
The next compulsory companion audit is V10.

% =============================================================================
% END APPENDED FIFTH-SERIES V9 MATERIAL
% =============================================================================

% =============================================================================
% BEGIN APPENDED FIFTH-SERIES V10 MATERIAL
% =============================================================================

\chapter{V10: Companion audit and closure of the hot/private cell}
\label{chap:newV10-hot}

\section{Mandatory read-only companion audit}
\label{sec:newV10-audit}

The newest active companion files inspected at this checkpoint were
\begin{center}
\begin{tabular}{p{0.57\textwidth}rr}
\toprule
file & logical lines & bytes\\
\midrule
\texttt{Renewal\_SM\_Einstein\_Continuum\_Research\_Notes\_V8.tex}
& \(2\,508\) & \(102\,599\)\\
\texttt{Renewal\_NS\_Stability\_Research\_Notes\_V31.tex}
& \(23\,052\) & \(887\,537\)\\
\bottomrule
\end{tabular}
\end{center}
Their SHA--256 digests, in the same order, were
\[
\begin{split}
&\texttt{40ee23e210ffc756395cf1280cd00b8758374fa1d3af4d2a54ac78908eeae188},\\
&\texttt{f1121b62238ad5491bb7cf03635ea9934942edf573ed8faaa9ee79e1d38a6696}.
\end{split}
\]
Neither file was modified.

The SM V8 note proves, on its positive massive-safe compact sector,
that a coercive compact density has a uniform \(L^p\) interval above
one.  Sparse coloring makes the replicated covariance exponent lie
inside that interval, so exact square-free compact remainders
multiply at arbitrary arity.  It explicitly excludes unbounded
Higgs owners, determinant-zero crossings, massless covariance, and
a cutoff-dependent correlation length.

The NS V31 note proves an exact dyadic first-moment reconstruction
and a scale-matched Duhamel formation ladder.  Geometric parent
summation removes dyadic multiplicity, but its universal estimate
returns the enstrophy-square continuation stock.  It therefore keeps
the recent-entry correlation and the already-paid persistent
intervals separate.

\begin{assessment}[Type-compatible audit transfer]
\label{ass:newV10-audit}
No SM integrability exponent, covariance constant, or NS
formation estimate is imported into Yang--Mills.  Two structural
rules are retained:
\begin{enumerate}[leftmargin=3em]
\item match the exact tail exponent to the assembled product before
      replacing it by an absolute moment; and
\item allocate a continuously shared stock once, without separately
      charging each scale bin or output branch.
\end{enumerate}
For the present one-link problem these rules point to the exact
private Wilson tail and to a single capped continuous row weight.
\end{assessment}

\begin{firewall}[No cross-program theorem transfer]
\label{fire:newV10-no-transfer}
The companion results concern different measures, operators, and
target topologies.  They imply neither a Wilson selected-junction
bound nor WUFC nor a Yang--Mills spectral gap.
\end{firewall}

\section{The hot centered numerator and the summed target}
\label{sec:newV10-numerator}

Retain the selected-junction variables from V8 and put
\begin{equation}
 c_i:=\min\{1,\sqrt{t_i}\},\qquad
 C_0:=\sum_{i=1}^m c_i,\qquad n:=m-2.
 \label{eq:newV10-capped-scales}
\end{equation}
Assume that at least one selected row is thermally hot.  Choose one
such row \(h\), so \(t_h\ge1\) and \(c_h=1\).

The all-order centered displacement theorem in archived
fourth-series V61 gives
\begin{equation}
 \|K_m\|_{\rm op}
 \le m!K_{\rm h}^m\prod_{i\ne h}c_i.
 \label{eq:newV10-centered}
\end{equation}
Unlike the archived hot-star use of this estimate, we shall not
assign the selected junction to that local star.  Its target would
contain \(x_h^{m-2}\) and would ask one private row for a growing
Casimir moment.  Instead we retain the exact full tree sum
\begin{equation}
 {\cal T}_{\rm sum}
 :=\sum_{\tau\in\mathfrak T_m}{\cal T}_\tau
 =\left(\prod_ia_i\right)S^n,
 \qquad S=\sum_ix_i,
 \label{eq:newV10-target-sum}
\end{equation}
and assign the complete numerator proportionally only after the
summed inequality is proved.

\begin{lemma}[Exact hot/private scalar quotients]
\label{lem:newV10-quotients}
On the private-dominant residual
\[
 \Xi_T\le D_mY,\qquad D_m\le K_D^m,
\]
division of the paid estimate
\eqref{eq:newV10-centered} by
\eqref{eq:newV10-target-sum} leaves, apart from
\(m!K^m\), the low-output quotient
\begin{equation}
 {\cal H}^{\rm lo}_{m,h}
 :=
 {q_Ps^{m-1}\prod_{i\ne h}c_i
  \over(\prod_it_i)S^n},
 \label{eq:newV10-hot-low}
\end{equation}
and the high-output quotient
\begin{equation}
 {\cal H}^{\rm hi}_{m,h}
 :=
 D_m\left(\sum_i p_i\right){\cal H}^{\rm lo}_{m,h}.
 \label{eq:newV10-hot-high}
\end{equation}
\end{lemma}

\begin{proof}
Since \(a_i=t_i/s\),
\[
 {\cal T}_{\rm sum}
 =s^{-m}\left(\prod_it_i\right)S^n.
\]
The low-output payer is at most \(C/s\), which gives
\eqref{eq:newV10-hot-low}.  On high output,
\[
 \Xi_T\le D_mY={D_m\over s}\sum_i p_i.
\]
This adds exactly the factor in
\eqref{eq:newV10-hot-high}.
\end{proof}

\section{Capped weighted-load allocation}
\label{sec:newV10-load}

\begin{theorem}[Anisotropy-uniform hot/private low branch]
\label{thm:newV10-hot-low}
Assume \(0<s\le1\), \(t_i\ge s\), at least one \(t_h\ge1\), and
the exact \(SU(2)\) private-row tails
\[
 q_i\le
 \exp\left\{-c
 \min\left(p_i,\sqrt{p_i/s}\right)\right\}.
\]
Then
\begin{equation}
 {\cal H}^{\rm lo}_{m,h}\le K^m,
 \label{eq:newV10-hot-low-bound}
\end{equation}
uniformly in the number and sizes of both hot and thermal rows and
in all private masses.
\end{theorem}

\begin{proof}
Use the probability weights
\begin{equation}
 w_i={c_i\over C_0},\qquad
 \gamma_i=nw_i.
 \label{eq:newV10-hot-weights}
\end{equation}
Weighted AM--GM gives, exactly as in V9,
\[
 S^{-n}\le
 \prod_iw_i^{\gamma_i}x_i^{-\gamma_i}.
\]
Apply the V9 fractional-load transform:
\[
 q_ix_i^{-\gamma_i}
 \le
 e^{C\gamma_i}
 \left(1+{\gamma_i+s\gamma_i^2\over t_i}\right)^{\gamma_i}.
\]
The product of the exponential factors is \(e^{Cn}\).

For a thermal row, \(t_i<1\) and \(c_i=\sqrt{t_i}\), whence
\begin{align}
 w_i\left(1+{\gamma_i+s\gamma_i^2\over t_i}\right)
 &=
 {\sqrt{t_i}\over C_0}
 +{n\over C_0^2}
 +{sn^2\sqrt{t_i}\over C_0^3}.
 \label{eq:newV10-thermal-bracket}
\end{align}
For a hot row, \(t_i\ge1\) and \(c_i=1\), so
\begin{align}
 w_i\left(1+{\gamma_i+s\gamma_i^2\over t_i}\right)
 &\le
 {1\over C_0}
 +{n\over C_0^2}
 +{sn^2\over C_0^3}.
 \label{eq:newV10-hot-bracket}
\end{align}
Every \(t_i\ge s\) implies \(c_i\ge\sqrt s\).  Therefore
\[
 C_0\ge m\sqrt s,\qquad C_0\ge1,
\]
and both brackets are at most
\begin{equation}
 {1\over C_0}
 \left(1+{n\over C_0}+{sn^2\over C_0^2}\right)
 \le {3\over C_0\sqrt s}.
 \label{eq:newV10-common-bracket}
\end{equation}

Let \(k\) be the number of thermal rows different from \(h\).
In the prefactor of \eqref{eq:newV10-hot-low},
\[
 {1\over t_h}\le1,\qquad
 {c_i\over t_i}
 =
 \begin{cases}
  t_i^{-1/2}\le s^{-1/2},&t_i<1,\\
  t_i^{-1}\le1,&t_i\ge1.
 \end{cases}
\]
Combining this with \eqref{eq:newV10-common-bracket} and
\(\sum_i\gamma_i=n\) gives
\begin{align}
 {\cal H}^{\rm lo}_{m,h}
 &\le
 C^m
 s^{m-1}s^{-k/2}
 (C_0\sqrt s)^{-n}
 \nonumber\\
 &\le
 C^m s^{m-1-k/2-n/2}.
 \label{eq:newV10-final-power}
\end{align}
Here \(C_0^{-n}\le1\).  Since \(k\le m-1\) and \(n=m-2\),
\[
 m-1-{k\over2}-{n\over2}
 ={m-k\over2}\ge{1\over2}.
\]
As \(s\le1\), the final power is at most one.  This proves
\eqref{eq:newV10-hot-low-bound}.
\end{proof}

\begin{theorem}[Anisotropy-uniform hot/private high branch]
\label{thm:newV10-hot-high}
Under the hypotheses of \Cref{thm:newV10-hot-low},
\begin{equation}
 {\cal H}^{\rm hi}_{m,h}\le K_1^m.
 \label{eq:newV10-hot-high-bound}
\end{equation}
\end{theorem}

\begin{proof}
Expand \(\sum_kp_k\).  In its \(k\)-th summand split only the
corresponding private multiplier:
\[
 p_kq_P
 =(p_kq_k^{1/2})
  \left(q_k^{1/2}\prod_{i\ne k}q_i\right).
\]
By the V9 split-tail lemma, the first factor is uniformly bounded,
and the multipliers in the second factor satisfy the same two-regime
tail with a fixed smaller constant.  Apply
\Cref{thm:newV10-hot-low} to each summand.  The \(m\) choices and
\(D_m\le K_D^m\) are absorbed into the exponential base.
\end{proof}

\section{Exact proportional tree assignment}
\label{sec:newV10-proportional}

\begin{lemma}[A summed positive target gives an exact atlas]
\label{lem:newV10-proportional}
Suppose a complete operator \(K_m\) and positive scalar targets
\({\cal T}_\tau\) satisfy
\[
 q\,\|\mathcal P K_m\|_{\rm op}
 \le m!K^m\sum_{\sigma\in\mathfrak T_m}{\cal T}_\sigma
\]
for the relevant payer \(\mathcal P\).  Define
\begin{equation}
 K_{m,\tau}
 :={ {\cal T}_\tau\over
       \sum_\sigma{\cal T}_\sigma}\,K_m .
 \label{eq:newV10-proportional-assignment}
\end{equation}
Then
\[
 K_m=\sum_\tau K_{m,\tau},
 \qquad
 q\,\|\mathcal P K_{m,\tau}\|_{\rm op}
 \le m!K^m{\cal T}_\tau .
\]
\end{lemma}

\begin{proof}
All \(a_i,\Xi_i\) are positive, hence every
\({\cal T}_\tau>0\).  The coefficients in
\eqref{eq:newV10-proportional-assignment} sum to one.  Multiply the
summed norm inequality by the corresponding positive coefficient.
\end{proof}

\begin{theorem}[Complete all-order hot/private \(SU(2)\) selected cell]
\label{thm:newV10-hot-cell}
Every selected \(SU(2)\) junction containing at least one thermally
hot selected row obeys the exact marked-tree bound with constant
\begin{equation}
 m!K_{\rm hot}^m
 \label{eq:newV10-hot-cell}
\end{equation}
in both final-output branches, uniformly in arity,
representations, private supports, and output dimensions.
\end{theorem}

\begin{proof}
Outside the private-dominant residual, use the exact output- and
branching-dominant re-fibrations of archived V63.  A residual with a
weak private row is the archived V64 star theorem.  On the remaining
hard residual, use \eqref{eq:newV10-centered};
\Cref{lem:newV10-quotients,thm:newV10-hot-low,
thm:newV10-hot-high} prove its paid norm is bounded by
\(m!K^m{\cal T}_{\rm sum}\).  Apply
\Cref{lem:newV10-proportional}.  The inherited Fourier
normalization and raw pinching make the estimate dimension-free.
\end{proof}

\begin{theorem}[Complete all-order \(SU(2)\) selected gate]
\label{thm:newV10-complete-selected}
Every selected one-coordinate \(SU(2)\) junction, at every arity,
obeys the exact paid marked-tree atlas with factorial-exponential
constant
\begin{equation}
 \boxed{\qquad m!K_{\rm sel}^m.\qquad}
 \label{eq:newV10-complete-selected}
\end{equation}
\end{theorem}

\begin{proof}
If all selected rows are thermal, use
\Cref{thm:newV9-all-thermal-selected}.  If at least one is hot, use
\Cref{thm:newV10-hot-cell}.  These two cases are disjoint and
exhaustive.
\end{proof}

\begin{assessment}[Why this does not repeat the false private moment]
\label{ass:newV10-no-false-moment}
The proof never asks one row to pay
\((sP_i)^{m-2}q_i\).  The continuous load
\(\gamma_i=(m-2)c_i/C_0\) is assigned before taking row norms.
Its total is exactly \(m-2\), and the two-regime Wilson transform
pays each row at that fractional load.  The high-output mass uses a
single split half-tail and is not charged again in the low-output
allocation.
\end{assessment}

\section{V11 programme and claim boundary}
\label{sec:newV10-next}

\begin{programme}[V11: insertion into the global marked-tree compiler]
\label{prog:newV10-V11}
\begin{enumerate}[leftmargin=3em]
\item Reopen the exact first-connectivity source factorization and
      identify each occurrence of the now-proved selected gate.
\item Verify that proper-prefix block moments, the unique joining
      cumulant, and the unrestricted suffix share one
      factorial-exponential generating function.
\item Insert \eqref{eq:newV10-complete-selected} without paying a
      second labelled-tree or set-partition factorial.
\item Prove the all-order global MTA if the remaining nonselected
      source cells are already closed; otherwise isolate the smallest
      missing cell and go upstream.
\item Enter WUFC only after the complete global MTA constant has been
      checked.
\end{enumerate}
\end{programme}

\begin{firewall}[A complete selected gate is not the mass gap]
\label{fire:newV10-boundary}
\Cref{thm:newV10-complete-selected} closes the complete all-order
selected one-coordinate gate for \(SU(2)\).  It does not by itself
prove the global MTA compiler, WUFC, CPM, cutoff removal, reflection
positivity, Osterwalder--Schrader reconstruction, a continuum
Yang--Mills measure, or a positive Hamiltonian gap.  The full
Yang--Mills mass gap has not been proved.
\end{firewall}

\begin{theorem}[V10 result ledger]
\label{thm:newV10-results}
V10:
\begin{enumerate}[label=\textup{(V10.\arabic*)},leftmargin=6.3em]
\item completes the mandatory read-only SM/NS audit;
\item derives the exact hot/private low- and high-output quotients;
\item proves their anisotropy-uniform bounds by capped continuous
      load allocation;
\item avoids every growing private Casimir moment;
\item proves an exact proportional tree assignment; and
\item closes the complete all-order \(SU(2)\) selected gate.
\end{enumerate}
\end{theorem}

\begin{proof}
Use
\Cref{ass:newV10-audit,
lem:newV10-quotients,
thm:newV10-hot-low,
thm:newV10-hot-high,
lem:newV10-proportional,
thm:newV10-hot-cell,
thm:newV10-complete-selected}.
\end{proof}

\paragraph{Parent-version record.}
V10 was formed from
\texttt{Renewal\_Yang\_Mills\_Mass\_Gap\_Research\_Notes\_V9.tex}.
The V9 parent had \(4\,170\) logical source lines,
\(134\,582\) bytes, and SHA--256
\[
 \texttt{386cf89d2b4f9f60ea2cf8d480d6d1d105547b27880663eefc337953b50bf8b7}.
\]
The next compulsory companion audit is V15.

% =============================================================================
% END APPENDED FIFTH-SERIES V10 MATERIAL
% =============================================================================

% =============================================================================
% BEGIN APPENDED FIFTH-SERIES V11 MATERIAL
% =============================================================================

\chapter{V11: The all-order first-connectivity compiler}
\label{chap:newV11-global}

\section{Exact joining at arbitrary quotient arity}
\label{sec:newV11-joining}

Let \(I=[m]\), let \(\rho=\{C_1,\ldots,C_k\}\) be the cumulative
row partition immediately before the first-connectivity coordinate
\(r\), and put
\[
 Y_\nu=\bigotimes_{i\in C_\nu}A_{i,r}
 \qquad(1\le\nu\le k),
\]
with identity factors on inactive row legs.  As in the archived
first-connectivity formula, \(L_{r,\pi}\) denotes the product of
complete local cumulants over the blocks of \(\pi\).

\begin{theorem}[All-arity relative product-cumulant identity]
\label{thm:newV11-relative-cumulant}
For every proper predecessor partition \(\rho<\widehat1\),
\begin{equation}
 \boxed{\qquad
 \sum_{\substack{\pi\in\Pi(I)\\
                  \pi\vee\rho=\widehat1}}
 L_{r,\pi}
 =
 \kappa_k^{\nu_\alpha}(Y_1,\ldots,Y_k).
 \qquad}
 \label{eq:newV11-relative-cumulant}
\end{equation}
Thus all local partitions which first join the \(k\) predecessor
blocks are assembled into one complete quotient cumulant before any
norm is taken.
\end{theorem}

\begin{proof}
Expand the moment of every product \(Y_\nu\) into local cumulants of
the individual row variables.  For a partition
\(\sigma\in\Pi([k])\), the product of the moments over its blocks
contains exactly the row partitions \(\pi\in\Pi(I)\) satisfying
\[
 \pi\le \widetilde\sigma,
\]
where \(\widetilde\sigma\) is obtained by merging the
\(\rho\)-blocks whose quotient labels lie in a common
\(\sigma\)-block.  Applying the cumulant Möbius coefficient
\(\mu(\sigma,\widehat1)\), the coefficient of a fixed \(\pi\) is
\[
 \sum_{\sigma:\ \pi\vee\rho\le\widetilde\sigma}
 \mu(\sigma,\widehat1).
\]
Möbius cancellation makes this coefficient one when
\(\pi\vee\rho=\widehat1\), and zero otherwise.  This proves
\eqref{eq:newV11-relative-cumulant}.
\end{proof}

\begin{firewall}[No iterated quotient inverses]
\label{fire:newV11-no-iterated}
Equation \eqref{eq:newV11-relative-cumulant} is a signed algebraic
identity.  Replacing its right side by a sequence of binary
covariances chooses a bracketing, omits Möbius terms, and can insert
several resolvent payers.  V11 uses the single complete quotient
cumulant and the one final payer supplied by the selected gate.
\end{firewall}

\section{Reducible block-products and row endpoints}
\label{sec:newV11-reducible}

\begin{lemma}[All-order selected gate under block-product amplification]
\label{lem:newV11-reducible}
The complete selected gate
\eqref{eq:newV10-complete-selected} remains valid when its \(k\)
inputs are finite reducible tensor-product \(SU(2)\)
representations \(Y_\nu\), carrying arbitrary trace-class prefix
coefficients.  Its constant is
\[
 k!K^kK_0^{\sum_\nu|C_\nu|}
 \le k!K_1^{m},
\]
and its raw input and output sums have no representation-multiplicity
factor.
\end{lemma}

\begin{proof}
Decompose each block-product representation orthogonally into
irreducible isotypic summands.  Apply
\Cref{thm:newV10-complete-selected} to every tuple of irreducible
summands.  Its constant is uniform in representations, physical
multiplicities, output dimensions, and private support.
Orthogonal compression, isotypic pinching, and multiplicity partial
trace are trace-norm contractions.  Complete raw pinching therefore
sums the block coefficients to at most the trace norm of the
unresolved prefix coefficient.

For a block \(C_\nu\), fixed-group Casimir fusion and
Cauchy--Schwarz give
\[
 A(Y_\nu)
 \le K_0^{|C_\nu|}\sum_{i\in C_\nu}a_{i,r}.
\]
Multiplying over the quotient-tree incidences introduces at most an
exponential \(K_0^m\), not a multiplicity count.  This proves the
stated constant.
\end{proof}

\begin{lemma}[Exact row lift of internal and quotient trees]
\label{lem:newV11-tree-lift}
Choose a spanning tree \(T_\nu\) inside every predecessor block
\(C_\nu\), interpreting a singleton as the empty tree.  Choose a
tree \(T_\rho\) on the \(k\) quotient vertices.  For every quotient
edge \(C_\nu C_\lambda\), expand its nonnegative block mass and
choose one endpoint pair \(ij\) with
\(i\in C_\nu,j\in C_\lambda\).  Then
\[
 T=\left(\bigcup_\nu T_\nu\right)
 \cup\{\hbox{the chosen quotient-edge lifts}\}
\]
is a labelled spanning tree on \(I\).
\end{lemma}

\begin{proof}
The internal trees have
\[
 \sum_\nu(|C_\nu|-1)=m-k
\]
edges.  The quotient tree has \(k-1\) edges.  After lifting, its
edges join distinct predecessor blocks according to a connected
acyclic quotient graph.  The union is therefore connected and has
\((m-k)+(k-1)=m-1\) edges, hence is a tree.
\end{proof}

\begin{lemma}[Endpoint expansion has no hidden block count]
\label{lem:newV11-endpoints}
Every quotient-tree pair stock is bounded by a nonnegative endpoint
sum
\begin{equation}
 \Theta_{C_\nu C_\lambda,r}
 \le
 K_0^{|C_\nu|+|C_\lambda|}
 \sum_{\substack{i\in C_\nu\\j\in C_\lambda}}
 \theta_{ij,r}.
 \label{eq:newV11-endpoints}
\end{equation}
If all endpoint sums are left uncollapsed, their product is exactly a
sum over lifted row trees from
\Cref{lem:newV11-tree-lift}; no power of \(m\) is paid outside that
tree sum.
\end{lemma}

\begin{proof}
Apply the Casimir fusion estimate in the proof of
\Cref{lem:newV11-reducible} to the two block-products and expand the
product of their row-mass sums.  The original row normalizations turn
each summand into the corresponding normalized pair mark
\(\theta_{ij,r}\).  Distributivity expands a quotient-tree monomial
into endpoint choices.  By
\Cref{lem:newV11-tree-lift}, every choice is already an ordinary
row-tree monomial, so the endpoint count is part of the displayed
tree sum rather than a separate constant.
\end{proof}

\section{Ordered coordinates are support-uniform}
\label{sec:newV11-coordinate}

For every row pair \(ij\), write
\[
 \Theta_{ij}=\sum_e\theta_{ij,e},
\qquad \theta_{ij,e}\ge0.
\]
These are the normalized global pair stocks inherited from the
marked-link probability budget.

\begin{lemma}[Ordered diagonal domination]
\label{lem:newV11-diagonal}
Let a lifted row tree \(T\) be split into internal prefix edges
\(E_{\rm p}\) and quotient joining edges \(E_{\rm j}\), with
\(E_{\rm j}\ne\varnothing\).  Then
\begin{equation}
 \sum_r
 \left[
 \prod_{ij\in E_{\rm p}}
 \left(\sum_{e<r}\theta_{ij,e}\right)
 \right]
 \left[
 \prod_{ij\in E_{\rm j}}\theta_{ij,r}
 \right]
 \le
 \prod_{ij\in E(T)}\Theta_{ij}.
 \label{eq:newV11-diagonal}
\end{equation}
The constant is one and is independent of support size.
\end{lemma}

\begin{proof}
Expand the left side into coordinate assignments to all edges.  It
contains only assignments for which every joining edge has the common
coordinate \(r\) and every prefix edge lies strictly before \(r\).
The right side is the sum of the same nonnegative monomials over all
independent coordinate assignments, without either restriction.
Thus the left indexing set is a subset of the right indexing set.
\end{proof}

\begin{remark}[Suffix contractions and a suffix payer]
\label{rem:newV11-suffix}
The first-connectivity suffix is formed as a complete moment before
norms and is contractive.  If the unique payer is represented at a
suffix coordinate, the inherited payer allocation keeps that one
factor and contracts all other suffix moments.  This adds one
nonnegative marked coordinate to the same assignment sum; dropping
its order constraint is again bounded by the corresponding full
global stock.  No number of suffix coordinates is paid.
\end{remark}

\section{The one-factorial partition census}
\label{sec:newV11-factorial}

\begin{lemma}[Exact partition-factorial sum]
\label{lem:newV11-partition-factorial}
For every \(m\ge1\),
\begin{equation}
 \boxed{\qquad
 \sum_{\rho\in\Pi([m])}
 |\rho|!\prod_{B\in\rho}|B|!
 =
 m!\,2^{m-1}.
 \qquad}
 \label{eq:newV11-partition-factorial}
\end{equation}
\end{lemma}

\begin{proof}
For a partition \(\rho\) with \(k\) blocks, the factor \(k!\) orders
the blocks and the factors \(|B|!\) order the elements within every
block.  Such data are equivalent to a permutation of \([m]\) together
with \(k-1\) cut positions among its \(m-1\) gaps.  Hence the
contribution at fixed \(k\) is
\[
 m!\binom{m-1}{k-1}.
\]
Sum \(k=1,\ldots,m\).
\end{proof}

\begin{corollary}[No second factorial in recursive connectivity]
\label{cor:newV11-one-factorial}
Suppose a predecessor partition \(\rho\) has \(k\) blocks and that a
connected prefix block \(B\) costs at most
\(|B|!K^{|B|}\), while the quotient joining gate costs at most
\(k!K^k\).  Summing over all predecessor partitions costs at most
\[
 m!(2K_1)^m.
\]
\end{corollary}

\begin{proof}
Since \(k\le m\) and \(\sum_B|B|=m\), all exponential bases are
bounded by \(K_1^m\).  Apply
\Cref{lem:newV11-partition-factorial}.
\end{proof}

\section{Strong induction for the global atlas}
\label{sec:newV11-induction}

\begin{theorem}[All-order global \(SU(2)\) marked-tree atlas]
\label{thm:newV11-global-MTA}
There is \(K_{\rm MTA}<\infty\) such that, for every \(m\ge2\),
arbitrary finite row supports, and trace-class row-factorized Fourier
coefficients, the complete once-paid connected \(m\)-row physical
carrier satisfies
\begin{equation}
 \boxed{\qquad
 \kappa_\otimes^{\rm abs}(\mathcal C_m^{\rm pay})
 \le
 m!K_{\rm MTA}^m
 \sum_{T\in\mathfrak T_m}
 \prod_{ij\in E(T)}\Theta_{ij}.
 \qquad}
 \label{eq:newV11-global-MTA}
\end{equation}
The estimate is uniform in volume, coordinate supports, input
representations, output multiplicities, output-duality orientation,
and payer placement.
\end{theorem}

\begin{proof}
We use strong induction on \(m\).  The binary case is the inherited
support-uniform covariance MTA.  Assume
\eqref{eq:newV11-global-MTA} through order \(m-1\).

Apply the exact linked-word identity and disintegrate at the unique
first-connectivity coordinate \(r\).  If the predecessor join is
\(\rho=\{C_1,\ldots,C_k\}<\widehat1\), its prefix carrier factors
exactly into complete connected carriers on the nonsingleton blocks
and complete one-row moments on the singleton blocks.  For every
nonsingleton \(C_\nu\), apply the inductive once-paid theorem and
then the inherited scalar-inclusive payer-stripping principle.  This
removes the lower-order payer, includes scalar prefix outputs, and
leaves the same internal tree sum with one harmless numerator factor
\(s\).  Singleton prefix moments and the complete suffix are
contractions.

Assemble the entire local family which joins the \(k\) predecessor
blocks by \Cref{thm:newV11-relative-cumulant}.  Apply the complete
selected gate
\Cref{thm:newV10-complete-selected} through the reducible
amplification of \Cref{lem:newV11-reducible}.  This joining theorem
contains the sole final payer and supplies a quotient tree on the
\(k\) blocks.

Lift the internal and quotient trees by
\Cref{lem:newV11-tree-lift,lem:newV11-endpoints}.  Sum the unique
joining coordinate, all prefix coordinates, and the possible suffix
payer by \Cref{lem:newV11-diagonal,rem:newV11-suffix}.  The result is
the global tree monomial on the right of
\eqref{eq:newV11-global-MTA}, with no support factor.

For this fixed predecessor partition, the factorial constant is at
most
\[
 k!K^k\prod_{B\in\rho}|B|!K^{|B|}.
\]
The factors \(s\le1\) left by the nonsingleton prefix numerators may
be discarded.  Finally sum all predecessor partitions and use
\Cref{cor:newV11-one-factorial}.  The exact first-connectivity
classes are disjoint before this last triangle inequality, so no word
is counted twice.  Enlarging a common exponential base closes the
induction.
\end{proof}

\begin{corollary}[The global MTA obstruction is removed for \(SU(2)\)]
\label{cor:newV11-MTA-removed}
The first implication in the route
\eqref{eq:newV1-route}, from one-link connected Wilson cards to the
all-order global marked-tree atlas, now holds with the required
factorial-exponential constant.
\end{corollary}

\begin{proof}
The local input is
\Cref{thm:newV10-complete-selected}; the global conclusion is
\Cref{thm:newV11-global-MTA}.
\end{proof}

\begin{assessment}[Scope and next interface]
\label{ass:newV11-scope}
The proof is specific to the exact \(SU(2)\) Wilson tail used in
V9--V10.  A general compact semisimple gauge group requires an
analogous support-uniform highest-weight Fourier tail before the same
compiler applies.  For \(SU(2)\), the next live interface is WUFC:
one must insert \eqref{eq:newV11-global-MTA} into the rooted spatial
fusion operator, verify its norm and activity conventions, and check
that the factor \(m!\) is paired with the exponential generating
coefficient rather than summed a second time.
\end{assessment}

\section{V12 programme and claim boundary}
\label{sec:newV11-next}

\begin{programme}[V12: weighted uniform fusion control]
\label{prog:newV11-V12}
\begin{enumerate}[leftmargin=3em]
\item Restate the exact archived MTA-to-WUFC compiler and its rooted
      Banach norm, including the activity factorial.
\item Insert \eqref{eq:newV11-global-MTA} and cancel its \(m!\)
      against the exponential generating coefficient.
\item Use the mark probability budget to sum spatial attachment
      sites without a support or volume count.
\item Compute the explicit small-activity condition under which the
      rooted fusion operator is a contraction.
\item If WUFC needs an estimate not contained in the global MTA,
      isolate that estimate rather than declaring CPM.
\end{enumerate}
\end{programme}

\begin{firewall}[A global MTA is not the continuum mass gap]
\label{fire:newV11-boundary}
\Cref{thm:newV11-global-MTA} proves the all-order global \(SU(2)\)
marked-tree atlas.  It does not yet prove WUFC, CPM, a convergent
cutoff-uniform polymer measure, cutoff removal, reflection
positivity, Osterwalder--Schrader reconstruction, a continuum
Yang--Mills measure, or a positive Hamiltonian gap.  The full
Yang--Mills mass gap has not been proved.
\end{firewall}

\begin{theorem}[V11 result ledger]
\label{thm:newV11-results}
V11:
\begin{enumerate}[label=\textup{(V11.\arabic*)},leftmargin=6.5em]
\item proves the all-arity relative product-cumulant joining
      identity;
\item proves reducible block-product stability and exact row-tree
      lifting;
\item proves support-uniform ordered diagonal summation;
\item proves the exact one-factorial partition census; and
\item proves the all-order global \(SU(2)\) marked-tree atlas by
      strong induction.
\end{enumerate}
\end{theorem}

\begin{proof}
Use
\Cref{thm:newV11-relative-cumulant,
lem:newV11-reducible,
lem:newV11-tree-lift,
lem:newV11-endpoints,
lem:newV11-diagonal,
lem:newV11-partition-factorial,
thm:newV11-global-MTA}.
\end{proof}

\paragraph{Parent-version record.}
V11 was formed from
\texttt{Renewal\_Yang\_Mills\_Mass\_Gap\_Research\_Notes\_V10.tex}.
The V10 parent had \(4\,599\) logical source lines,
\(148\,157\) bytes, and SHA--256
\[
 \texttt{c363a929168a36284ef65ba810efaa12459b8d7f41fe206c3244c90091ea8b67}.
\]
The next compulsory companion audit is V15.

% =============================================================================
% END APPENDED FIFTH-SERIES V11 MATERIAL
% =============================================================================

% =============================================================================
% BEGIN APPENDED FIFTH-SERIES V12 MATERIAL
% =============================================================================

\chapter{V12: Factorial-grade correction and exact Haar resummation}
\label{chap:newV12-correction}

\section{The archived MTA normalization}
\label{sec:newV12-normalization}

The exact archived spatial compiler uses the following two grades.
Its marked-tree atlas requires, for every fixed labelled tree,
\begin{equation}
 \|\mathcal K_{m,T}\|
 \le K_0^{m-1}\mathcal W_T,
 \label{eq:newV12-strong-MTA}
\end{equation}
with no factorial.  Rooting the tree and summing its markings costs
one per input, while the number of rooted labelled trees obeys
\begin{equation}
 m^{m-1}\le e^{m-1}m!.
 \label{eq:newV12-Cayley-factorial}
\end{equation}
This produces the WUFC grade
\begin{equation}
 \|\mathcal RQ_0\kappa_m(g_1,\ldots,g_m)\|_{\rm FI}
 \le m!K^{m-1}\prod_i\|g_i\|_{\rm FI}.
 \label{eq:newV12-WUFC-grade}
\end{equation}
Thus the single allowed factorial is supplied by rooted Cayley
enumeration, not by the fixed-tree analytic estimate.

\begin{definition}[Factorial-graded marked-tree atlas]
\label{def:newV12-FGMTA}
Call an estimate an FGMTA if its fixed-tree grade is
\begin{equation}
 \|\mathcal K_{m,T}\|
 \le m!K^m\mathcal W_T.
 \label{eq:newV12-FGMTA}
\end{equation}
\end{definition}

\begin{theorem}[Correction of the V7--V11 claim grade]
\label{thm:newV12-correction}
The conclusions proved in V7, V9, V10, and V11 have the FGMTA grade
\eqref{eq:newV12-FGMTA}, not the strong MTA grade
\eqref{eq:newV12-strong-MTA}.  In particular:
\begin{enumerate}[leftmargin=3em]
\item the estimates and exact decompositions in those versions
      remain valid with their displayed \(m!K^m\) constants;
\item the V11 assertion that the global MTA obstruction is removed
      is withdrawn; and
\item the archived MTA-to-WUFC compiler cannot yet be invoked.
\end{enumerate}
\end{theorem}

\begin{proof}
The V7 local tree estimate explicitly contains \(r!K^r\).
The V9 and V10 scalar compilers multiply that estimate only by
exponential-per-input constants, so they retain the same factorial.
In V11 a predecessor partition
\(\rho=\{B_1,\ldots,B_k\}\) contributes the joining factorial \(k!\)
and the prefix factorials \(\prod_j|B_j|!\).  The exact V11 census is
\[
 \sum_{\rho\in\Pi([m])}|\rho|!\prod_{B\in\rho}|B|!
 =m!2^{m-1}.
\]
Thus the displayed conclusion of V11 is exactly FGMTA.

If \eqref{eq:newV12-FGMTA} is substituted into the archived spatial
compiler, the rooted-tree count
\eqref{eq:newV12-Cayley-factorial} yields a bound of order
\[
 (m!)^2K_1^m,
\]
not \eqref{eq:newV12-WUFC-grade}.  Since no fixed exponential base
absorbs \(m!\), the strong inference is invalid.
\end{proof}

\begin{corollary}[An exponential generating coefficient does not
repair the operator gate]
\label{cor:newV12-EGF}
Multiplying the FGMTA estimate by the usual exponential generating
coefficient \(1/m!\) cancels only its analytic factorial.  Rooted
Cayley enumeration leaves a coefficient of order \(m!K^m\), whose
ordinary power series has zero radius of convergence.  Therefore the
missing factorial cannot be deferred to CPM.
\end{corollary}

\begin{proof}
After the \(1/m!\) cancellation, the remaining rooted-tree count is
bounded above at the factorial scale and is of that scale up to an
exponential factor.  The ratio test for
\(\sum_m m!K^mz^m\) gives zero radius.
\end{proof}

\begin{firewall}[Correct live boundary]
\label{fire:newV12-correct-boundary}
V11 proves a support-uniform, all-order FGMTA for \(SU(2)\).  It does
not prove the strong MTA, WUFC, or CPM.  Every subsequent claim uses
this corrected grade.
\end{firewall}

\section{Why the logarithmic Haar vertex cannot remove the factorial}
\label{sec:newV12-log-no-go}

In the scaled \(SU(2)\) exponential chart, V6--V7 used logarithmic
Haar vertices
\begin{equation}
 H_n(y)
 =-{2\zeta(2n)\over n\pi^{2n}}|y|^{2n},
 \qquad n\ge1,
 \label{eq:newV12-log-Haar}
\end{equation}
with activity \(\epsilon^n\), where
\(\epsilon\asymp\alpha^{-1}\).

\begin{proposition}[Isolated logarithmic Haar derivatives are not
exponential-grade]
\label{prop:newV12-log-no-go}
There is no fixed \(K\) such that
\begin{equation}
 \epsilon^n
 \|\nabla^{2n}H_n(0)\|_{\rm row}
 \le K^{2n}
 \label{eq:newV12-false-log-bound}
\end{equation}
for all sufficiently small \(\epsilon\) and all
\(n\le\delta/\epsilon\), with fixed \(\delta>0\).
\end{proposition}

\begin{proof}
Restrict to one coordinate line \(y=te_1\).  Then
\[
 \left|{d^{2n}\over dt^{2n}}H_n(te_1)\bigg|_{t=0}\right|
 =
 {2\zeta(2n)\over n\pi^{2n}}(2n)!
 \ge {2\over n\pi^{2n}}(2n)!.
\]
Choose \(n=\lfloor\delta/(2\epsilon)\rfloor\).  Stirling's lower
bound gives
\[
 \left\{\epsilon^n(2n)!\right\}^{1/(2n)}
 \ge c\sqrt\epsilon\,n
 \ge {c_\delta\over\sqrt\epsilon},
\]
which diverges as \(\epsilon\downarrow0\).  The polynomial factor
\(n^{-1}\pi^{-2n}\) changes only a fixed exponential base and cannot
restore \eqref{eq:newV12-false-log-bound}.
\end{proof}

\begin{assessment}[Meaning of the no-go]
\label{ass:newV12-log-no-go}
This proposition does not disprove the strong Wilson tree card.
It proves that the V7 operation of bounding each logarithmic Haar
vertex separately cannot establish it.  The logarithm has converted
the factorial denominators of the exact Jacobian into vertices with
large isolated derivatives; cancellations between different
logarithmic orders must be restored before taking norms.
\end{assessment}

\section{The exact Haar factor restores the denominators}
\label{sec:newV12-exact-Haar}

Up to a fixed convention-dependent radial rescaling, the normalized
\(SU(2)\) Haar Jacobian in exponential coordinates is
\begin{equation}
 J_\epsilon(y)
 =
 \left\{
 {\sin(\sqrt\epsilon\,|y|)
  \over\sqrt\epsilon\,|y|}
 \right\}^{\!2}.
 \label{eq:newV12-Haar-J}
\end{equation}
The value at the origin is defined by continuity.

\begin{lemma}[Exact factorial-denominator series]
\label{lem:newV12-Haar-series}
The entire radial series is
\begin{equation}
 J_\epsilon(y)
 =
 \sum_{n=0}^{\infty}
 {(-1)^n2^{2n+1}\over(2n+2)!}
 \epsilon^n|y|^{2n}.
 \label{eq:newV12-Haar-series}
\end{equation}
\end{lemma}

\begin{proof}
Use
\[
 \left({\sin z\over z}\right)^2
 ={1-\cos(2z)\over2z^2}
\]
and expand the cosine series.  Set \(z=\sqrt\epsilon\,|y|\).
\end{proof}

\begin{lemma}[Radial derivative row]
\label{lem:newV12-radial-row}
In fixed dimension three, for \(0\le v\le2n\),
\begin{equation}
 \|\nabla^v|y|^{2n}\|_{\rm row}
 \le
 C^v{(2n)!\over(2n-v)!}(1+|y|)^{2n-v}.
 \label{eq:newV12-radial-row}
\end{equation}
\end{lemma}

\begin{proof}
Differentiate the product of \(2n\) linear coordinate factors after
polarizing \((y\cdot y)^n\).  Choosing the differentiated slots gives
the falling factorial \((2n)!/(2n-v)!\).  Contraction of the fixed
Euclidean metric tensors changes the complete row norm by at most
\(C^v\).  Every undifferentiated slot is bounded by \(1+|y|\).
\end{proof}

\begin{theorem}[Factorial-free integrated exact-Haar derivative card]
\label{thm:newV12-exact-Haar-card}
There are \(C<\infty\) and \(\epsilon_0>0\) such that, for every
\(0<\epsilon\le\epsilon_0\), every \(v\ge0\), and every centered
Gaussian \(Y\in\mathbb R^3\) with covariance bounded by the identity,
\begin{equation}
 \boxed{\qquad
 \mathbb E
 \|\nabla^vJ_\epsilon(Y)\|_{\rm row}
 \le
 C^v\epsilon^{v/2}.
 \qquad}
 \label{eq:newV12-exact-Haar-card}
\end{equation}
For \(v=0\), the right side is interpreted after enlarging the fixed
constant.
\end{theorem}

\begin{proof}
For \(q\ge0\), fixed-dimensional Gaussian moments satisfy
\begin{equation}
 \mathbb E(1+|Y|)^q\le C_0^qq^{q/2}.
 \label{eq:newV12-Gaussian-moment}
\end{equation}
Apply \Cref{lem:newV12-radial-row} termwise in
\eqref{eq:newV12-Haar-series}.  Only \(n\ge\lceil v/2\rceil\)
contribute.  With \(q=2n-v\), the \(n\)-th expected derivative is at
most
\begin{align}
 &{2^{2n+1}\epsilon^n\over(2n+2)!}
 C^v{(2n)!\over q!}
 C_0^qq^{q/2}
 \nonumber\\
 &\qquad\le
 C_1^{n+v}\epsilon^n.
 \label{eq:newV12-Haar-term}
\end{align}
Indeed
\[
 {(2n)!\over(2n+2)!}
 ={1\over(2n+1)(2n+2)}
\]
and \(q^{q/2}/q!\le e^q\) by the elementary Stirling lower bound;
the cases \(q=0,1\) are absorbed into the constant.

Choose \(\epsilon_0\) so that \(C_1\epsilon_0<1/2\).  Summing
\eqref{eq:newV12-Haar-term} from
\(\lceil v/2\rceil\) to infinity gives
\[
 C_2^v
 \sum_{n\ge\lceil v/2\rceil}(C_1\epsilon)^n
 \le C_3^v\epsilon^{v/2}.
\]
Termwise differentiation and expectation are justified by this
absolute majorant.  The covariance hypothesis gives
\eqref{eq:newV12-Gaussian-moment} uniformly.
\end{proof}

\begin{corollary}[The maximal-valence Haar star is exponential-grade]
\label{cor:newV12-Haar-star}
For \(v=2n\), the exact Haar factor contributes at most
\[
 (C\sqrt\epsilon)^{2n}
\]
to the integrated derivative row.  Thus the maximal-valence star
which forced the \(r!\) ledger in V7 has no factorial when all
logarithmic Haar orders are reassembled first.
\end{corollary}

\begin{proof}
Apply \Cref{thm:newV12-exact-Haar-card} with \(v=2n\).
\end{proof}

\begin{assessment}[What is still needed for strong MTA]
\label{ass:newV12-remaining}
The exact-Haar card removes one identified source of the analytic
factorial.  A strong Wilson tree card still requires:
\begin{enumerate}[leftmargin=3em]
\item an analogous exponential-grade assembled card for the exact
      action correction and normalization;
\item a BKAR formula in which exact Haar and action factors are
      differentiated only after their internal order sums are
      assembled;
\item proportional tree assignment at the strong grade, without the
      V11 triangle inequality over predecessor partitions; and
\item the private-tail scalar compiler with those new
      factorial-free local constants.
\end{enumerate}
The Gaussian tangent card already has the correct fixed-tree
exponential grade.  The live obstruction is therefore the
factorial-free Wilson-density transfer, not spatial support
summation.
\end{assessment}

\section{V13 programme and claim boundary}
\label{sec:newV12-next}

\begin{programme}[V13: exact-density strong local card]
\label{prog:newV12-V13}
\begin{enumerate}[leftmargin=3em]
\item Write the exact scaled Wilson density as one assembled
      action factor times \(J_\epsilon\), without logarithmic Haar
      vertices.
\item Prove a complete derivative-row convolution whose constants
      are exponential in total derivative order.
\item Insert that convolution into the Gaussian BKAR formula and
      retain one labelled-tree sum with no \(m!\) per tree.
\item Treat the density normalization as a scalar quotient after
      the connected numerator has been assembled.
\item If the action factor fails the exponential derivative card,
      identify its first exact coefficient obstruction rather than
      returning to the logarithmic vertex expansion.
\end{enumerate}
\end{programme}

\begin{firewall}[Corrected claim boundary after V12]
\label{fire:newV12-boundary}
V12 corrects the V7--V11 factorial grade and proves the
factorial-free integrated derivative card for the exact \(SU(2)\)
Haar Jacobian.  It does not yet prove the exact-action card, strong
local or global MTA, WUFC, CPM, cutoff removal, reflection positivity,
Osterwalder--Schrader reconstruction, a continuum Yang--Mills
measure, or a positive Hamiltonian gap.  The full Yang--Mills mass
gap has not been proved.
\end{firewall}

\begin{theorem}[V12 result ledger]
\label{thm:newV12-results}
V12:
\begin{enumerate}[label=\textup{(V12.\arabic*)},leftmargin=6.5em]
\item identifies the exact strong-MTA and WUFC factorial grades;
\item reclassifies V7--V11 as a valid FGMTA and withdraws the
      overstrong WUFC implication;
\item proves that isolated logarithmic Haar derivatives cannot
      remove the factorial;
\item restores the exact Haar factorial denominators; and
\item proves the all-order factorial-free integrated exact-Haar
      derivative card.
\end{enumerate}
\end{theorem}

\begin{proof}
Use
\Cref{thm:newV12-correction,
prop:newV12-log-no-go,
lem:newV12-Haar-series,
lem:newV12-radial-row,
thm:newV12-exact-Haar-card,
cor:newV12-Haar-star}.
\end{proof}

\paragraph{Parent-version record.}
V12 was formed from
\texttt{Renewal\_Yang\_Mills\_Mass\_Gap\_Research\_Notes\_V11.tex}.
The V11 parent had \(5\,026\) logical source lines,
\(163\,451\) bytes, and SHA--256
\[
 \texttt{0501006a043aa1b10426ab797b9002bf0ee87adb38bc57ec489c681d81cbfc2f}.
\]
The next compulsory companion audit is V15.

% =============================================================================
% END APPENDED FIFTH-SERIES V12 MATERIAL
% =============================================================================

% =============================================================================
% BEGIN APPENDED FIFTH-SERIES V13 MATERIAL
% =============================================================================

\chapter{V13: The exact-density \(L^p\) boundary and intrinsic redirection}
\label{chap:newV13-density}

\section{The exact scaled \(SU(2)\) density}
\label{sec:newV13-exact-density}

Take the radial convention in which the Wilson action is
\(1-\cos\theta\), the exponential-chart boundary is
\(\theta=\pi\), and \(\epsilon=\alpha^{-1}\).  Relative to standard
Gaussian measure
\[
 d\gamma_3(y)=(2\pi)^{-3/2}e^{-|y|^2/2}\,dy,
\]
the unnormalized scaled density is
\begin{equation}
 \widetilde h_\epsilon(y)
 =
 J_\epsilon(y)A_\epsilon(y)
 \mathbf1_{\{\sqrt\epsilon|y|<\pi\}},
 \label{eq:newV13-density}
\end{equation}
where
\begin{align}
 J_\epsilon(y)
 &=
 \left\{
 {\sin(\sqrt\epsilon|y|)
  \over\sqrt\epsilon|y|}
 \right\}^2,
 \label{eq:newV13-J}\\
 A_\epsilon(y)
 &=
 \exp\left\{
 -{1-\cos(\sqrt\epsilon|y|)\over\epsilon}
 +{|y|^2\over2}
 \right\}.
 \label{eq:newV13-A}
\end{align}
The normalized density is
\[
 h_\epsilon
 =\widetilde h_\epsilon/
   \mathbb E_{\gamma_3}\widetilde h_\epsilon.
\]
The denominator is bounded above and below by positive constants for
all sufficiently small \(\epsilon\), by the standard Laplace
comparison at the unique Wilson well.

\begin{lemma}[Exact cancellation of the Gaussian reference]
\label{lem:newV13-reference-cancellation}
For every nonnegative measurable \(F\),
\begin{equation}
 \int F(y)\widetilde h_\epsilon(y)\,d\gamma_3(y)
 =
 (2\pi)^{-3/2}
 \int_{\sqrt\epsilon|y|<\pi}
 F(y)J_\epsilon(y)
 e^{-\{1-\cos(\sqrt\epsilon|y|)\}/\epsilon}\,dy.
 \label{eq:newV13-reference-cancellation}
\end{equation}
\end{lemma}

\begin{proof}
Multiply \eqref{eq:newV13-A} by the Gaussian density.  The two
quadratic exponents cancel exactly.
\end{proof}

\begin{assessment}[Why the action factor must not be estimated alone]
\label{ass:newV13-action-alone}
The factor \(A_\epsilon\) grows like \(e^{|y|^2/2}\) near the chart
scale.  Its integrability comes from the Gaussian cancellation in
\eqref{eq:newV13-reference-cancellation}, not from a uniform
pointwise bound.  Separating the two factors before choosing the
replica exponent can therefore destroy the exact coercivity.
\end{assessment}

\section{Sharp uniform \(L^p\) window}
\label{sec:newV13-Lp}

Put
\begin{equation}
 p_*={\pi^2\over\pi^2-4}>1.
 \label{eq:newV13-pstar}
\end{equation}

\begin{theorem}[Uniform \(L^p\) window and its sharp upper boundary]
\label{thm:newV13-Lp-window}
For every fixed \(1\le p<p_*\),
\begin{equation}
 \sup_{0<\epsilon\le\epsilon_0}
 \|h_\epsilon\|_{L^p(\gamma_3)}<\infty.
 \label{eq:newV13-Lp-upper}
\end{equation}
For every fixed \(p>p_*\),
\begin{equation}
 \|h_\epsilon\|_{L^p(\gamma_3)}
 \longrightarrow\infty
 \qquad(\epsilon\downarrow0),
 \label{eq:newV13-Lp-divergence}
\end{equation}
indeed at an exponential rate in \(\epsilon^{-1}\).
\end{theorem}

\begin{proof}
Normalization changes the estimates only by fixed constants.  Since
\[
 1-\cos\theta\ge {2\theta^2\over\pi^2}
 \qquad(|\theta|\le\pi)
\]
and \(0\le J_\epsilon\le1\), the \(p\)-th power of
\(\widetilde h_\epsilon\) times the Gaussian density is bounded by a
fixed constant times
\begin{align}
 \exp\left\{
 -{p\over\epsilon}(1-\cos(\sqrt\epsilon|y|))
 +{p-1\over2}|y|^2
 \right\}
 &\le
 e^{-c_p|y|^2},
 \label{eq:newV13-Lp-coercivity}
\end{align}
where
\[
 c_p={2p\over\pi^2}-{p-1\over2}>0
\]
exactly when \(p<p_*\).  Integration proves
\eqref{eq:newV13-Lp-upper}.

Now let \(p>p_*\).  The continuous function
\[
 \Phi_p(\theta)
 ={p-1\over2}\theta^2-p(1-\cos\theta)
\]
satisfies
\[
 \Phi_p(\pi)
 ={p-1\over2}\pi^2-2p>0.
\]
Choose a compact interval
\([a,b]\subset(0,\pi)\), sufficiently close to \(\pi\), on which
\(\Phi_p\ge c_0>0\).  On the shell
\[
 a\le\sqrt\epsilon|y|\le b,
\]
the Jacobian \(J_\epsilon\) is bounded below by a positive constant
depending only on \(a,b\).  The radial integral of
\(\widetilde h_\epsilon^p\,d\gamma_3\) over this shell is therefore
bounded below by
\[
 c\,\epsilon^{-3/2}e^{c_0/\epsilon}.
\]
The normalization remains bounded, proving
\eqref{eq:newV13-Lp-divergence}.
\end{proof}

\begin{remark}[The endpoint is not needed]
\label{rem:newV13-endpoint}
The theorem deliberately leaves \(p=p_*\) undecided.  The all-arity
obstruction uses arbitrarily large integer exponents, so the strict
upper boundary is sufficient.
\end{remark}

\section{Fully correlated replicas leave the window}
\label{sec:newV13-replicas}

\begin{lemma}[The coincident-replica exponent is the arity]
\label{lem:newV13-coincident}
Let \(Y_1=\cdots=Y_r=Y\) be coincident standard Gaussian replicas.
For every nonnegative \(f\),
\begin{equation}
 \mathbb E\prod_{i=1}^r f(Y_i)
 =\|f\|_{L^r(\gamma_3)}^r.
 \label{eq:newV13-coincident}
\end{equation}
Consequently no product inequality based only on a uniform
\(L^{p_0}\) norm with fixed \(p_0<\infty\) can control coincident
replicas at all arities.
\end{lemma}

\begin{proof}
The first equality is the definition of the \(L^r\) norm.  If
\(r>p_0\), an \(L^{p_0}\) bound alone does not imply an \(L^r\)
bound on a probability space; the exact density family in
\Cref{thm:newV13-Lp-window} realizes this failure uniformly as
\(\epsilon\downarrow0\).
\end{proof}

\begin{theorem}[No all-arity Gaussian density-product transfer]
\label{thm:newV13-no-density-transfer}
There is no proof of the strong local Wilson MTA which uses only:
\begin{enumerate}[leftmargin=3em]
\item the Gaussian BKAR representation;
\item absolute multiplication of one copy of \(h_\epsilon\) for
      every selected replica; and
\item a uniform \(L^{p_0}(\gamma_3)\) bound on \(h_\epsilon\) at one
      fixed exponent \(p_0>1\).
\end{enumerate}
Such an argument already fails on the coincident-replica endpoint for
every integer \(r>\max\{p_0,p_*\}\).
\end{theorem}

\begin{proof}
At the coincident endpoint, the absolute density product is
\[
 \mathbb E_{\gamma_3}h_\epsilon(Y)^r
 =\|h_\epsilon\|_r^r
\]
by \Cref{lem:newV13-coincident}.  For \(r>p_*\), this quantity
diverges exponentially in \(\epsilon^{-1}\) by
\Cref{thm:newV13-Lp-window}.  No constant \(K^r\), uniform in
\(\epsilon\), can bound it.  A fixed \(L^{p_0}\) estimate supplies no
additional information at exponent \(r>p_0\).
\end{proof}

\begin{firewall}[Why sparse coloring does not alter the one-link test]
\label{fire:newV13-no-color}
Sparse coloring separates distinct spatial anchors.  The replicas in
\Cref{lem:newV13-coincident} are several selected observables at the
same Wilson link.  They cannot be spatially recolored without changing
the local cumulant.  Thus the SM V8 sparse-color mechanism from the
V10 audit is type-incompatible with this one-link obstruction.
\end{firewall}

\section{A complementary pointwise action no-go}
\label{sec:newV13-entire}

\begin{proposition}[The exact action correction is not of exponential
derivative type]
\label{prop:newV13-action-entire}
Fix \(\epsilon>0\) and consider the entire one-dimensional action
correction
\[
 A_\epsilon(z)
 =
 \exp\left\{
 {\cos(\sqrt\epsilon z)-1\over\epsilon}
 +{z^2\over2}
 \right\}.
\]
There is no finite \(K_\epsilon\) such that
\begin{equation}
 |A_\epsilon^{(v)}(0)|\le K_\epsilon^v
 \qquad(v\ge0).
 \label{eq:newV13-false-entire}
\end{equation}
\end{proposition}

\begin{proof}
If \eqref{eq:newV13-false-entire} held, the Taylor series would give
\[
 |A_\epsilon(z)|\le e^{K_\epsilon|z|}
\]
for every complex \(z\).  On the imaginary axis \(z=iy\),
\[
 |A_\epsilon(iy)|
 =
 \exp\left\{
 {\cosh(\sqrt\epsilon y)-1\over\epsilon}
 -{y^2\over2}
 \right\},
\]
which grows faster than \(e^{K_\epsilon|y|}\) for every fixed
\(K_\epsilon\).  This contradiction proves the claim.
\end{proof}

\begin{assessment}[Scope of the action no-go]
\label{ass:newV13-action-no-go}
The proposition rules out a global pointwise exponential derivative
row for the separated action factor.  It does not rule out a
cancellation-preserving integrated identity for the complete Wilson
law.  The compact chart cutoff and exact Gaussian cancellation must be
retained in any such identity.
\end{assessment}

\section{The corrected strong-card target}
\label{sec:newV13-intrinsic}

The Gaussian tangent forest formula has the correct strong fixed-tree
grade because it differentiates only the representation observables.
\Cref{thm:newV13-no-density-transfer} shows that the Wilson law cannot
be inserted as an absolute product of replica densities.  The
remaining route must keep the Wilson measure intrinsic.

\begin{definition}[Intrinsic strong Wilson forest card]
\label{def:newV13-ISWFC}
An intrinsic strong Wilson forest card is an exact decomposition
\begin{equation}
 \kappa_m^{\nu_\alpha}
 (D^{R_1},\ldots,D^{R_m})
 =
 \sum_{T\in\mathfrak T_m}{\cal I}_{\alpha,T}
 \label{eq:newV13-ISWFC-decomposition}
\end{equation}
such that
\begin{equation}
 \|{\cal I}_{\alpha,T}\|_{\rm op}
 \le K^{m-1}\prod_i b_i^{d_i(T)}
 \label{eq:newV13-ISWFC}
\end{equation}
for every all-thermal input family, with the Wilson law, its
normalization, and all internal score or resolvent sums assembled
before the norm.
\end{definition}

\begin{theorem}[ISWFC is the missing local factorial-removal
certificate]
\label{thm:newV13-ISWFC-sufficient}
Assume \eqref{eq:newV13-ISWFC} and its hot/private selected analogue,
with the rowwise Wilson-tail scalar compiler at exponential grade.
Then the first-connectivity construction of V11 can be replaced by a
strong global MTA construction, after the predecessor partition
family is kept in its exact relative cumulant rather than bounded
partition by partition.
\end{theorem}

\begin{proof}
The local fixed-tree factor is then \(K^{m-1}\), the grade required by
\eqref{eq:newV12-strong-MTA}.  Exact relative product-cumulant
assembly is V11, reducible block-product amplification and row-tree
lifting are dimension-free, and ordered coordinate summation has
constant one.  The only factorial in the V11 proof came from applying
the triangle inequality separately to predecessor partitions and
from the FGMTA local input.  Retaining the exact relative cumulant
forest removes the first operation, while
\eqref{eq:newV13-ISWFC} removes the second.  The resulting fixed-tree
constant is exponential, so the archived spatial compiler supplies
the single WUFC factorial through Cayley enumeration.
\end{proof}

\begin{assessment}[Exact remaining analytic problem]
\label{ass:newV13-remaining}
V13 has ruled out both separated logarithmic-Haar vertices and
absolute Gaussian density products.  The viable upstream objects are
the intrinsic Wilson projected-tail recursion and an intrinsic
Wilson forest interpolation.  The frozen archive already supplies
the exact symmetry-first subset recursion and negative-norm terminal
pairing; the new target is to improve their \(m!K^m\) estimate to the
fixed-tree exponential grade, or to derive
\eqref{eq:newV13-ISWFC} directly from the Wilson diffusion.
\end{assessment}

\section{V14 programme and claim boundary}
\label{sec:newV13-next}

\begin{programme}[V14: intrinsic forest identity]
\label{prog:newV13-V14}
\begin{enumerate}[leftmargin=3em]
\item Start from the exact Wilson covariance identity
      \(\operatorname{Cov}(F,G)=
      \mathbb E\Gamma(F,R_\alpha G)\).
\item Introduce a forest interpolation through conditional
      expectations or resolvent contractions of the Wilson law,
      without multiplying replica densities.
\item Verify at arities two and three that each forest edge supplies
      one representation derivative row and no repeated inverse.
\item Prove positivity or an absolute contraction bound for every
      interpolated functional before asserting
      \eqref{eq:newV13-ISWFC}.
\item If a Markov forest identity fails, return to the exact
      symmetrized projected-tail sum and identify its first
      non-exponential coefficient.
\end{enumerate}
\end{programme}

\begin{firewall}[Claim boundary after V13]
\label{fire:newV13-boundary}
V13 proves the sharp strict \(L^p\) boundary and rules out an
all-arity absolute Gaussian density-product transfer.  It defines a
sufficient intrinsic strong Wilson forest card but does not prove
that card, strong local or global MTA, WUFC, CPM, cutoff removal,
reflection positivity, Osterwalder--Schrader reconstruction, a
continuum Yang--Mills measure, or a positive Hamiltonian gap.  The
full Yang--Mills mass gap has not been proved.
\end{firewall}

\begin{theorem}[V13 result ledger]
\label{thm:newV13-results}
V13:
\begin{enumerate}[label=\textup{(V13.\arabic*)},leftmargin=6.5em]
\item writes the exact scaled \(SU(2)\) Wilson density and Gaussian
      cancellation;
\item proves its uniform \(L^p\) window and strict sharp upper
      boundary \(p_*\);
\item proves the coincident-replica obstruction to fixed-exponent
      density multiplication;
\item proves a complementary pointwise exact-action derivative
      no-go; and
\item identifies the intrinsic strong Wilson forest card as the
      factorial-removal certificate.
\end{enumerate}
\end{theorem}

\begin{proof}
Use
\Cref{lem:newV13-reference-cancellation,
thm:newV13-Lp-window,
lem:newV13-coincident,
thm:newV13-no-density-transfer,
prop:newV13-action-entire,
thm:newV13-ISWFC-sufficient}.
\end{proof}

\paragraph{Parent-version record.}
V13 was formed from
\texttt{Renewal\_Yang\_Mills\_Mass\_Gap\_Research\_Notes\_V12.tex}.
The V12 parent had \(5\,409\) logical source lines,
\(176\,208\) bytes, and SHA--256
\[
 \texttt{e49d5fd3495e82f6cc6813e1d1e494bfab5eab0a385ed1d76c5683ce6deca422}.
\]
The next compulsory companion audit is V15.

% =============================================================================
% END APPENDED FIFTH-SERIES V13 MATERIAL
% =============================================================================

% =============================================================================
% BEGIN APPENDED FIFTH-SERIES V14 MATERIAL
% =============================================================================

\chapter{V14: A density-free random-cluster forest identity}
\label{chap:newV14-forest}

\section{Replica coupling by a Bernoulli graph}
\label{sec:newV14-coupling}

Let \((\Omega,\nu)\) be any probability space and let
\(f_i:\Omega\to\operatorname{End}(V_i)\), \(i\in[m]\), be bounded.
For a graph \(G\) on \([m]\), with connected-component partition
\(\pi(G)\), define
\begin{equation}
 M_G(f_1,\ldots,f_m)
 :=
 \bigotimes_{B\in\pi(G)}
 \mathbb E_\nu\bigotimes_{i\in B}f_i.
 \label{eq:newV14-graph-moment}
\end{equation}
All tensor legs retain their canonical label order.

For
\(\boldsymbol x=(x_{ij})_{1\le i<j\le m}\in[0,1]^{\binom m2}\),
let \(G_{\boldsymbol x}\) be the random graph whose edges are
independent and satisfy
\(\mathbb P(ij\in G_{\boldsymbol x})=x_{ij}\).  Put
\begin{equation}
 F_m(\boldsymbol x)
 :=
 \mathbb E_{G_{\boldsymbol x}}M_{G_{\boldsymbol x}}.
 \label{eq:newV14-interpolated-moment}
\end{equation}
This is a multilinear polynomial in the edge variables.

\begin{lemma}[Endpoint and factorization properties]
\label{lem:newV14-endpoints}
The interpolation satisfies
\begin{align}
 F_m(\boldsymbol1)
 &=\mathbb E_\nu\bigotimes_{i=1}^mf_i,
 \label{eq:newV14-common-endpoint}\\
 F_m(\boldsymbol0)
 &=\bigotimes_{i=1}^m\mathbb E_\nu f_i.
 \label{eq:newV14-independent-endpoint}
\end{align}
More generally, if \(\rho\in\Pi([m])\) and
\(x_{ij}=0\) whenever \(i,j\) belong to different \(\rho\)-blocks,
then \(F_m(\boldsymbol x)\) factors as a tensor product of the
corresponding interpolations on the blocks of \(\rho\).
\end{lemma}

\begin{proof}
At \(\boldsymbol1\), the random graph is complete and has one
component.  At \(\boldsymbol0\), it has \(m\) singleton components.
Under the block-zero hypothesis no random edge crosses two
\(\rho\)-blocks.  Independence of all remaining Bernoulli edges and
the definition \eqref{eq:newV14-graph-moment} then factor the graph
sum block by block.
\end{proof}

\section{The universal forest formula}
\label{sec:newV14-BKAR}

For a forest \(F\) and \(\boldsymbol s\in[0,1]^{E(F)}\), define
\[
 x^F_{ij}(\boldsymbol s)
 =
 \begin{cases}
 \min\{s_e:e\text{ lies on the }F\text{-path from }i\text{ to }j\},
 &i,j\text{ are \(F\)-connected},\\
 0,&\text{otherwise}.
 \end{cases}
\]
For an edge set \(A\), write
\(\partial_A=\prod_{e\in A}\partial_{x_e}\).

\begin{theorem}[Random-cluster BKAR moment formula]
\label{thm:newV14-BKAR-moment}
\begin{equation}
 F_m(\boldsymbol1)
 =
 \sum_{F\ {\rm forest\ on}\ [m]}
 \int_{[0,1]^{E(F)}}
 \partial_{E(F)}F_m
 \bigl(\boldsymbol x^F(\boldsymbol s)\bigr)
 \,d\boldsymbol s .
 \label{eq:newV14-BKAR-moment}
\end{equation}
\end{theorem}

\begin{proof}
This is the finite multivariable fundamental-theorem-of-calculus
forest formula.  For completeness, induct on the number of vertices.
Apply the fundamental theorem of calculus to the edge variables in
any fixed order, stopping a branch when its selected edges contain a
cycle.  The minimum-along-path prescription records the last
parameter at which two components first become joined.  Branches with
the same acyclic selected edge set combine into the cube integral
displayed above.  Since \(F_m\) is a polynomial, no limiting argument
is required.  Evaluating the identity on edge monomials verifies that
both sides equal one when the monomial edge set is connected inside
each component of the selected forest and equal zero otherwise.
\end{proof}

\begin{theorem}[Density-free cumulant forest identity]
\label{thm:newV14-cumulant-forest}
For every \(m\ge2\),
\begin{equation}
 \boxed{\qquad
 \kappa_m^\nu(f_1,\ldots,f_m)
 =
 \sum_{T\in\mathfrak T_m}
 \int_{[0,1]^{E(T)}}
 \partial_{E(T)}F_m
 \bigl(\boldsymbol x^T(\boldsymbol s)\bigr)
 \,d\boldsymbol s.
 \qquad}
 \label{eq:newV14-cumulant-forest}
\end{equation}
The formula uses only moments of the original law \(\nu\); no
Radon--Nikodym density is raised to a replica-dependent power.
\end{theorem}

\begin{proof}
For a forest \(F\) with components
\(B_1,\ldots,B_k\), the block-zero property in
\Cref{lem:newV14-endpoints} and the fact that all derivatives lie
inside forest components give
\[
 \partial_{E(F)}F_m(\boldsymbol x^F)
 =
 \bigotimes_{\ell=1}^k
 \partial_{E(F|_{B_\ell})}
 F_{B_\ell}(\boldsymbol x^{F|_{B_\ell}}).
\]
Thus \eqref{eq:newV14-BKAR-moment} is a moment expansion whose term
indexed by \(F\) factors over the component partition \(\pi(F)\).
Möbius inversion on \(\Pi([m])\) extracts its connected component.
The connected forests are exactly the labelled spanning trees,
which proves \eqref{eq:newV14-cumulant-forest}.
\end{proof}

\begin{corollary}[The Wilson interpolation stays intrinsic]
\label{cor:newV14-Wilson-intrinsic}
Taking \(\nu=\nu_\alpha\), every forest integrand in
\eqref{eq:newV14-cumulant-forest} is a finite signed combination of
ordinary \(SU(2)\) Wilson moments.  The \(L^p\) obstruction in V13
does not occur.
\end{corollary}

\begin{proof}
Each graph moment in
\eqref{eq:newV14-interpolated-moment} assigns one independent
\(\nu_\alpha\)-distributed group variable to every graph component.
There is never a product of density ratios on one Gaussian variable.
\end{proof}

\section{Every edge derivative is an exact component merge}
\label{sec:newV14-merge}

\begin{lemma}[Single-edge derivative]
\label{lem:newV14-edge-derivative}
Fix \(e=ij\).  Let \(G^{-e}_{\boldsymbol x}\) be the Bernoulli graph
with the edge \(e\) omitted and all other edge probabilities
unchanged.  Then
\begin{equation}
 \partial_{x_e}F_m(\boldsymbol x)
 =
 \mathbb E_{G^{-e}_{\boldsymbol x}}
 \left[
 M_{G^{-e}_{\boldsymbol x}+e}
 -M_{G^{-e}_{\boldsymbol x}}
 \right].
 \label{eq:newV14-edge-derivative}
\end{equation}
If \(i,j\) are already in the same component, the bracket is zero.
If they belong to distinct components \(B,C\), the bracket is
\begin{equation}
 \left\{
 \mathbb E_\nu(F_BF_C)
 -\mathbb E_\nu F_B\,\mathbb E_\nu F_C
 \right\}
 \bigotimes_{D\ne B,C}\mathbb E_\nu F_D,
 \label{eq:newV14-component-covariance}
\end{equation}
where \(F_B=\bigotimes_{i\in B}f_i\).
\end{lemma}

\begin{proof}
Condition on every Bernoulli edge except \(e\).  The conditional
expectation is
\[
 x_eM_{G^{-e}+e}+(1-x_e)M_{G^{-e}}.
\]
Differentiate.  Adding \(e\) changes nothing when its endpoints are
already connected.  Otherwise it replaces the two moment factors for
\(B,C\) by the single moment of their product, giving
\eqref{eq:newV14-component-covariance}.
\end{proof}

\begin{corollary}[Iterated derivatives are merge histories]
\label{cor:newV14-merge-history}
For a tree \(T\), the derivative
\(\partial_{E(T)}F_m\) is an average of signed histories in which
each effective tree edge either vanishes because its endpoints were
already connected by auxiliary Bernoulli edges, or merges two current
component-products by the complete covariance
\eqref{eq:newV14-component-covariance}.
\end{corollary}

\begin{proof}
Apply \Cref{lem:newV14-edge-derivative} successively, conditioning at
each step on all still-undifferentiated Bernoulli edges.  Mixed
partials commute because \(F_m\) is a polynomial.
\end{proof}

\begin{firewall}[A merge history is not a product of binary paid
estimates]
\label{fire:newV14-no-paid-iteration}
Although every derivative is a covariance, applying the once-paid
binary WUFC estimate at every edge would insert \(m-1\) resolvents.
The forest integrand has only the one final payer.  Internal merge
edges require projection-complete numerator estimates; the final
edge alone may carry the output resolvent.
\end{firewall}

\section{The ternary integrand exposes the first analytic test}
\label{sec:newV14-ternary}

For \(m=3\), abbreviate
\[
 a=M_1M_2M_3,\quad
 b_{12}=M_{12}M_3,\quad
 b_{13}=M_{13}M_2,\quad
 b_{23}=M_{23}M_1,\quad
 c=M_{123},
\]
where \(M_B=\mathbb E_\nu\prod_{i\in B}f_i\).

\begin{lemma}[Exact two-edge derivative]
\label{lem:newV14-ternary-integrand}
For the tree with edges \(12,13\), treating \(x_{23}=z\) as fixed,
\begin{equation}
 \partial_{x_{12}}\partial_{x_{13}}F_3
 =
 (1-z)(c-b_{12}-b_{13}+a)
 +z(b_{23}-c).
 \label{eq:newV14-ternary-integrand}
\end{equation}
In the BKAR integrand,
\[
 z=\min\{s_{12},s_{13}\}.
\]
\end{lemma}

\begin{proof}
Enumerating the eight graphs on three vertices gives
\begin{align*}
 F_3
 &=
 (1-x)(1-y)(1-z)a
 +x(1-y)(1-z)b_{12}\\
 &\quad
 +y(1-x)(1-z)b_{13}
 +z(1-x)(1-y)b_{23}\\
 &\quad
 +(xy+xz+yz-2xyz)c,
\end{align*}
where \(x=x_{12}\), \(y=x_{13}\).  Differentiate in \(x,y\) and
collect the coefficients of \(1-z\) and \(z\).
\end{proof}

\begin{assessment}[The first nontrivial strong-card row]
\label{ass:newV14-ternary}
Formula \eqref{eq:newV14-ternary-integrand} is not the product of two
separately normed covariances.  It is an interpolation between two
complete signed merge remainders.  The established ternary Wilson MTA
proves the integral after all three trees are assembled, but a
factorial-free all-order proof needs a uniform estimate of the
individual forest integrands in this cancellation-preserving form.
\end{assessment}

\section{The conditional merge certificate}
\label{sec:newV14-certificate}

\begin{definition}[Projection-complete conditional merge card]
\label{def:newV14-PCMC}
The PCMC holds if, for every labelled tree \(T\), every
\(\boldsymbol s\), every Bernoulli auxiliary graph appearing in
\Cref{cor:newV14-merge-history}, and every assignment of irreducible
Wilson representation inputs, the complete forest derivative after
one final payer satisfies
\begin{equation}
 \left\|
 \mathcal P_{\rm final}
 \partial_{E(T)}F_m
 \bigl(\boldsymbol x^T(\boldsymbol s)\bigr)
 \right\|_{\rm op}
 \le
 K^{m-1}\prod_i b_i^{d_i(T)}.
 \label{eq:newV14-PCMC}
\end{equation}
All internal component merges are kept as projection-complete
numerators and the Bernoulli expectation is taken before the norm.
\end{definition}

\begin{theorem}[PCMC implies the intrinsic strong Wilson forest card]
\label{thm:newV14-PCMC-sufficient}
If \eqref{eq:newV14-PCMC} holds, then the ISWFC
\eqref{eq:newV13-ISWFC} holds with the same exponential base up to a
fixed factor.
\end{theorem}

\begin{proof}
Apply the final payer to the exact identity
\eqref{eq:newV14-cumulant-forest}.  For a fixed tree, use
\eqref{eq:newV14-PCMC} inside its parameter integral.  The cube has
volume one.  This gives one term with the required weight for each
labelled tree and no additional tree, partition, or permutation
count.
\end{proof}

\begin{proposition}[Binary PCMC is already proved]
\label{prop:newV14-binary-PCMC}
At \(m=2\), the PCMC is exactly the inherited marked binary Wilson
quotient theorem and has the strong exponential grade.
\end{proposition}

\begin{proof}
There is one tree and
\[
 \partial_{x_{12}}F_2
 =M_{12}-M_1M_2.
\]
After the unique output payer and exact Fourier normalization, this
is the marked binary quotient.  Its two endpoint marks give
\(b_1b_2\) with a representation- and support-uniform constant.
\end{proof}

\begin{assessment}[What V14 changes]
\label{ass:newV14-progress}
V13 left open whether an intrinsic forest interpolation exists.  V14
constructs one for every probability law and proves its exact
component-merge derivative.  The strong local problem is now the
explicit PCMC integrand estimate.  Its advantages are that the
Wilson law remains a probability measure, the tree census appears
exactly once, and there is no logarithmic-Haar or density-product
factorial.  Its unresolved point is analytic: projection-complete
merge remainders such as
\eqref{eq:newV14-ternary-integrand} must retain all incidence powers
under one final payer.
\end{assessment}

\section{V15 programme and claim boundary}
\label{sec:newV14-next}

\begin{programme}[V15: companion audit and ternary PCMC]
\label{prog:newV14-V15}
\begin{enumerate}[leftmargin=3em]
\item Perform the mandatory read-only SM/NS companion audit.
\item Insert the exact Wilson moments into
      \eqref{eq:newV14-ternary-integrand} and rewrite both brackets
      as centered representation products.
\item Prove the ternary PCMC uniformly in the BKAR parameter
      \(z\in[0,1]\), retaining the final payer once.
\item Compare the resulting integrand bound with the archived
      complete ternary marked quotient, without inferring an
      integrand estimate from its integral.
\item Identify the induction norm on component-products needed for
      arbitrary tree pruning.
\end{enumerate}
\end{programme}

\begin{firewall}[Claim boundary after V14]
\label{fire:newV14-boundary}
V14 proves an exact density-free random-cluster forest identity and
reduces the strong local Wilson card to the PCMC.  It proves the PCMC
only at arity two.  It does not prove the ternary or all-order PCMC,
strong local or global MTA, WUFC, CPM, cutoff removal, reflection
positivity, Osterwalder--Schrader reconstruction, a continuum
Yang--Mills measure, or a positive Hamiltonian gap.  The full
Yang--Mills mass gap has not been proved.
\end{firewall}

\begin{theorem}[V14 result ledger]
\label{thm:newV14-results}
V14:
\begin{enumerate}[label=\textup{(V14.\arabic*)},leftmargin=6.5em]
\item constructs the Bernoulli random-cluster replica interpolation;
\item proves the density-free cumulant BKAR forest identity;
\item proves that every edge derivative is an exact complete
      component covariance;
\item computes the first ternary forest integrand explicitly;
\item defines the projection-complete conditional merge card; and
\item proves that PCMC implies ISWFC and that binary PCMC is already
      closed.
\end{enumerate}
\end{theorem}

\begin{proof}
Use
\Cref{lem:newV14-endpoints,
thm:newV14-BKAR-moment,
thm:newV14-cumulant-forest,
lem:newV14-edge-derivative,
lem:newV14-ternary-integrand,
thm:newV14-PCMC-sufficient,
prop:newV14-binary-PCMC}.
\end{proof}

\paragraph{Parent-version record.}
V14 was formed from
\texttt{Renewal\_Yang\_Mills\_Mass\_Gap\_Research\_Notes\_V13.tex}.
The V13 parent had \(5\,829\) logical source lines,
\(190\,034\) bytes, and SHA--256
\[
 \texttt{927aaae4e2f95473f2922411a5614002b9471f36946ae6e1f7475c0b8f5eeee7}.
\]
The next compulsory companion audit is V15.

% =============================================================================
% END APPENDED FIFTH-SERIES V14 MATERIAL
% =============================================================================

% =============================================================================
% BEGIN APPENDED FIFTH-SERIES V15 MATERIAL
% =============================================================================

\chapter{V15: Companion audit and the failure of raw forest trees}
\label{chap:newV15-raw-tree}

\section{Mandatory read-only companion audit}
\label{sec:newV15-audit}

The newest active companion notes inspected at this checkpoint were
\begin{center}
\begin{tabular}{p{0.63\textwidth}r}
\hline
file & bytes\\
\hline
\texttt{Renewal\_SM\_Einstein\_Continuum\_Research\_Notes\_V15.tex}
& \(192\,863\)\\
\texttt{Renewal\_NS\_Stability\_Research\_Notes\_V31.tex}
& \(887\,537\)\\
\hline
\end{tabular}
\end{center}
Their logical source-line counts and SHA--256 digests were,
respectively,
\begin{align}
 4758,\qquad&
 \texttt{c4d163346216c0f7048db6dbbcaacf040a6217453f9b4447658d37ef63e92a60},
 \label{eq:newV15-SM-audit}\\
 23052,\qquad&
 \texttt{f1121b62238ad5491bb7cf03635ea9934942edf573ed8faaa9ee79e1d38a6696}.
 \label{eq:newV15-NS-audit}
\end{align}
Neither companion file was modified.

The SM V15 note proves the exact degree-weighted Cayley census
\[
 \sum_{T\in\mathfrak T_N}\prod_i d_i(T)!
 =(N-2)!\binom{3N-3}{N-2}\le 8^N N!.
\]
This is a useful warning and a useful compiler: derivative-degree
factorials and labelled-tree multiplicity must be counted jointly.
It repairs an adjacent-cutoff Gevrey expansion because its unordered
owner coefficient cancels the resulting owner factorial.  It does
\emph{not} repair the present strong fixed-tree norm, where no
analytic factorial is allowed before the spatial Cayley sum.
The same note explicitly preserves that grade distinction.

The NS V31 note proves an exact dyadic first-moment identity and then
shows that its available global payment returns the continuation
quantity one is trying to prove.  Its transferable discipline is
therefore structural: keep the restricted, stopping-time-correlated
object intact until its cancellation is used.  In the present
problem, the analogue is to keep the sum of replica trees intact
until its lower-cumulant leakage has canceled.  Taking the norm of
each raw tree first can be circularly stronger than the cumulant
estimate itself.

\begin{firewall}[No analytic theorem crosses the companion audit]
\label{fire:newV15-audit}
The SM and NS notes concern different measures and target
statements.  Only the two proof disciplines just stated are used
below.  No SM determinant estimate and no Navier--Stokes formation
estimate is imported into the Wilson law.
\end{firewall}

\section{The exact integrated ternary tree}
\label{sec:newV15-integrated-tree}

Retain the notation of V14.  For the tree
\(\tau_1=\{12,13\}\), put
\[
 {\cal J}_1
 :=
 \int_{[0,1]^2}
 \partial_{x_{12}}\partial_{x_{13}}F_3
 \bigl(\boldsymbol x^{\tau_1}(\boldsymbol s)\bigr)
 \,d\boldsymbol s .
\]
Define the complete binary covariances
\[
 C_{ij}:=M_{ij}-M_iM_j
\]
and the complete third cumulant
\[
 K_3:=M_{123}-M_{12}M_3-M_{13}M_2-M_{23}M_1
       +2M_1M_2M_3.
\]

\begin{lemma}[Parameter integrals of the BKAR minimum]
\label{lem:newV15-minimum}
If \(z=\min(s,t)\) on the unit square, then
\begin{equation}
 \int_{[0,1]^2}z\,ds\,dt={1\over3},
 \qquad
 \int_{[0,1]^2}(1-z)\,ds\,dt={2\over3}.
 \label{eq:newV15-minimum}
\end{equation}
\end{lemma}

\begin{proof}
For \(0\le u\le1\),
\(\mathbb P(\min(s,t)>u)=(1-u)^2\).  The layer-cake identity gives
\[
 \int_{[0,1]^2}\min(s,t)\,ds\,dt
 =\int_0^1(1-u)^2\,du={1\over3}.
\]
The second equality follows because the square has volume one.
\end{proof}

\begin{lemma}[Exact lower-cumulant leakage of one tree]
\label{lem:newV15-tree-leakage}
For \(\{i,j,k\}=\{1,2,3\}\), let \({\cal J}_i\) be the integrated
raw forest term for the tree centered at \(i\).  Then
\begin{equation}
 \boxed{\quad
 {\cal J}_i
 ={1\over3}K_3
  +{2\over3}M_iC_{jk}
  -{1\over3}M_jC_{ik}
  -{1\over3}M_kC_{ij}.
 \quad}
 \label{eq:newV15-tree-leakage}
\end{equation}
In particular,
\begin{equation}
 {\cal J}_1+{\cal J}_2+{\cal J}_3=K_3.
 \label{eq:newV15-leakage-cancels}
\end{equation}
\end{lemma}

\begin{proof}
For \(i=1\), insert
\Cref{lem:newV14-ternary-integrand} and
\Cref{lem:newV15-minimum}:
\begin{align*}
 {\cal J}_1
 &={2\over3}(M_{123}-M_{12}M_3-M_{13}M_2+M_1M_2M_3)\\
 &\quad+{1\over3}(M_{23}M_1-M_{123})\\
 &={1\over3}M_{123}
   -{2\over3}M_{12}M_3
   -{2\over3}M_{13}M_2
   +{1\over3}M_{23}M_1
   +{2\over3}M_1M_2M_3.
\end{align*}
Expanding the right side of
\eqref{eq:newV15-tree-leakage} gives exactly this expression.
The other two identities follow by permutation.  On summing them,
the coefficient of each \(M_iC_{jk}\) is
\(2/3-1/3-1/3=0\), proving
\eqref{eq:newV15-leakage-cancels}.
\end{proof}

\begin{assessment}[Why a raw tree is not a connected object]
\label{ass:newV15-raw-not-connected}
The connectedness projection in
\Cref{thm:newV14-cumulant-forest} acts on the \emph{sum} of the
three ternary trees.  Equation
\eqref{eq:newV15-tree-leakage} shows that an individual raw tree
still contains binary covariances.  Those terms have only two
incidences and cannot satisfy a four-incidence ternary tree weight
uniformly at weak coupling.  Their cancellation occurs only in
\eqref{eq:newV15-leakage-cancels}.
\end{assessment}

\section{A representation-exact counterexample to PCMC}
\label{sec:newV15-counterexample}

We now verify the preceding warning inside the actual \(SU(2)\)
Wilson law.  Let
\[
 V=\mathbb C^2,\qquad
 X(U)=D^{1/2}(U),
\]
and put one copy of \(X\) on each of the three labelled tensor
legs.  Write
\[
 q=q_{\alpha,1/2},\qquad
 r=q_{\alpha,1},\qquad
 t=q_{\alpha,3/2}.
\]
The fixed-representation Wilson multiplier expansion inherited from
the exact one-link analysis is
\begin{equation}
 q_{\alpha,j}
 =1-s_\alpha j(j+1)+O_j(s_\alpha^2),
 \qquad s_\alpha\downarrow0.
 \label{eq:newV15-multiplier-expansion}
\end{equation}
Thus
\begin{align}
 q&=1-\frac34s_\alpha+O(s_\alpha^2),&
 r&=1-2s_\alpha+O(s_\alpha^2),&
 t&=1-\frac{15}{4}s_\alpha+O(s_\alpha^2).
 \label{eq:newV15-three-multipliers}
\end{align}

\begin{theorem}[Pointwise ternary PCMC is false]
\label{thm:newV15-pointwise-PCMC-false}
Let \(\tau_1=\{12,13\}\), and evaluate its V14 forest
integrand at any parameters for which
\(z=\min(s_{12},s_{13})=1\).  On the total-spin \(3/2\) subspace,
\begin{equation}
 P_{3/2}
 \partial_{x_{12}}\partial_{x_{13}}F_3
 P_{3/2}
 =(qr-t)P_{3/2}
 =\bigl(s_\alpha+O(s_\alpha^2)\bigr)P_{3/2}.
 \label{eq:newV15-pointwise-lower}
\end{equation}
Consequently no constant \(K\), uniform in \(\alpha\), can bound
this raw integrand by
\[
 K^2b_1^2b_2b_3
\]
for the tree \(\tau_1\), where
\(b_i=\sqrt{s_\alpha(1+C_2(1/2))}\).
The pointwise PCMC of V14 is therefore false already at arity three.
\end{theorem}

\begin{proof}
Clebsch--Gordan decomposition gives
\[
 \mathbb E_{\nu_\alpha}(X_2\otimes X_3)
 =P_0^{23}+rP_1^{23},
\qquad
 \mathbb E_{\nu_\alpha}X_1=qI.
\]
On the completely symmetric total-spin \(3/2\) subspace, the
\((2,3)\) pair is triplet and
\[
 M_{23}M_1=qrP_{3/2},
\qquad
 M_{123}=tP_{3/2}.
\]
At \(z=1\), \Cref{lem:newV14-ternary-integrand} is
\(M_{23}M_1-M_{123}\).  This proves the first equality.
Using \eqref{eq:newV15-three-multipliers},
\[
 qr-t
 =\left(1-\frac{11}{4}s_\alpha+O(s_\alpha^2)\right)
  -\left(1-\frac{15}{4}s_\alpha+O(s_\alpha^2)\right)
 =s_\alpha+O(s_\alpha^2).
\]
All three input marks are equal and
\[
 b_1^2b_2b_3
 =b^4
 =s_\alpha^2\left(1+C_2(1/2)\right)^2.
\]
The ratio of \eqref{eq:newV15-pointwise-lower} to this weight
diverges like \(s_\alpha^{-1}\).  Projection onto \(P_{3/2}\) is a
legal nontrivial output compression.  The usual final
mass--resolvent payer is a nonzero scalar on this fixed output and
cannot annihilate the displayed component.
\end{proof}

The failure is not caused merely by requiring uniformity in the
interpolation parameter.  It remains after integrating one labelled
tree.

\begin{theorem}[An integrated raw tree also loses one Wilson clock]
\label{thm:newV15-integrated-PCMC-false}
For the same three fundamental inputs,
\begin{equation}
 \|P_{1/2}{\cal J}_1P_{1/2}\|_{\rm op}
 \ge {1\over2}s_\alpha
 \label{eq:newV15-integrated-lower}
\end{equation}
for all sufficiently small \(s_\alpha\).  Hence an
integrated fixed-tree estimate
\[
 \|{\cal J}_1\|_{\rm op}\le K^2b_1^2b_2b_3
\]
is also false uniformly in weak coupling.
\end{theorem}

\begin{proof}
Let \(S_{ij}\) be the flip of tensor legs \(i,j\).  On
\(V\otimes V\),
\[
 P_0^{ij}={I-S_{ij}\over2},
\qquad
 P_1^{ij}={I+S_{ij}\over2}.
\]
Therefore the pair moment is
\begin{equation}
 M_{ij}
 =A I+B S_{ij},
\qquad
 A={1+r\over2},
\qquad
 B={r-1\over2}.
 \label{eq:newV15-pair-flip}
\end{equation}
The total-spin \(1/2\) isotypic subspace has the standard
two-dimensional \(S_3\) multiplicity representation.  On that
representation,
\begin{equation}
 S_{12}+S_{13}+S_{23}=0.
 \label{eq:newV15-standard-relation}
\end{equation}
Moreover \(M_{123}=qI\) there and \(M_i=qI\).  Substitution of
\eqref{eq:newV15-pair-flip} into the explicit expression in the
proof of \Cref{lem:newV15-tree-leakage} gives
\begin{equation}
 P_{1/2}{\cal J}_1P_{1/2}
 =
 \lambda_\alpha I+qB S_{23},
 \qquad
 \lambda_\alpha={q\over3}+{2q^3\over3}-qA.
 \label{eq:newV15-integrated-exact}
\end{equation}
Indeed, the permutation part is
\[
 {qB\over3}(-2S_{12}-2S_{13}+S_{23})
 =qB S_{23}
\]
by \eqref{eq:newV15-standard-relation}.
The flip \(S_{23}\) is a self-adjoint unitary on the multiplicity
space, so
\[
 \|\lambda_\alpha I+qB S_{23}\|_{\rm op}\ge |qB|.
\]
Equations \eqref{eq:newV15-three-multipliers} and
\eqref{eq:newV15-pair-flip} give
\[
 qB=-s_\alpha+O(s_\alpha^2).
\]
This proves \eqref{eq:newV15-integrated-lower}.  The proposed right
side is \(K^2b^4=O(K^2s_\alpha^2)\), which cannot dominate it for
fixed \(K\).
\end{proof}

\begin{firewall}[Append-only correction to the V14 direction]
\label{fire:newV15-PCMC-correction}
The random-cluster identities of V14, including
\Cref{thm:newV14-cumulant-forest}, remain exact.  The implication
PCMC \(\Rightarrow\) ISWFC also remains a valid conditional
statement.  However,
\Cref{thm:newV15-pointwise-PCMC-false,
thm:newV15-integrated-PCMC-false}
prove that the PCMC hypothesis as stated is false.  Thus the V14
programme of estimating raw forest trees separately is retired.
\end{firewall}

\section{The ternary compensation that does work}
\label{sec:newV15-compensation}

The obstruction is a decomposition defect, not a failure of the
complete ternary Wilson estimate.  The complete cumulant has the
four-incidence cancellation already proved in the frozen archive.
We record the exact balancing step because it is the model required
at arbitrary arity.

\begin{lemma}[Positive redistribution]
\label{lem:newV15-redistribution}
Let \(K\) be an element of a normed vector space and let
\((w_a)_{a\in A}\) be finitely many nonnegative numbers with
\(W=\sum_aw_a>0\).  The following are equivalent:
\begin{enumerate}[label=\textup{(\roman*)},leftmargin=3.5em]
\item \(\|K\|\le CW\);
\item there is an exact decomposition \(K=\sum_aK_a\) with
      \(\|K_a\|\le Cw_a\) for every \(a\).
\end{enumerate}
One valid choice is
\begin{equation}
 K_a={w_a\over W}K.
 \label{eq:newV15-positive-redistribution}
\end{equation}
\end{lemma}

\begin{proof}
The displayed choice proves \textup{(i)} implies \textup{(ii)}.
The triangle inequality proves the converse.
\end{proof}

For the three labelled trees set
\[
 w_i:=b_i^2b_jb_k,
\qquad
 W=w_1+w_2+w_3.
\]
The archived all-thermal ternary theorem gives
\begin{equation}
 \|K_3\|_{\rm op}\le C W.
 \label{eq:newV15-complete-ternary}
\end{equation}

\begin{theorem}[Compensated ternary forest card]
\label{thm:newV15-compensated-ternary}
Define
\begin{equation}
 \widehat{\cal J}_i:={w_i\over W}K_3,
\qquad
 \Delta_i:=\widehat{\cal J}_i-{\cal J}_i.
 \label{eq:newV15-balanced-trees}
\end{equation}
Then
\begin{align}
 K_3&=\sum_{i=1}^3\widehat{\cal J}_i,
 &
 \sum_{i=1}^3\Delta_i&=0,
 \label{eq:newV15-balanced-sum}\\
 \|\widehat{\cal J}_i\|_{\rm op}
 &\le Cb_i^2b_jb_k.
 \label{eq:newV15-balanced-bound}
\end{align}
Thus the ternary intrinsic strong tree card is valid, but its tree
terms are compensated after the complete cumulant cancellation;
they are not the individual random-cluster forest integrals.
\end{theorem}

\begin{proof}
The first identity and the bound are
\Cref{lem:newV15-redistribution} applied to
\eqref{eq:newV15-complete-ternary}.  The second identity follows
from \eqref{eq:newV15-leakage-cancels}.  All operations take place
before an output norm is summed.  The arbitrary-support version is
the already established global ternary marked-tree atlas; this
theorem isolates only the new relation between that atlas and the
raw V14 forest formula.
\end{proof}

\begin{assessment}[The correct all-order target]
\label{ass:newV15-correct-target}
At arbitrary arity, an exact forest formula is useful only if its
lower-cumulant leakage can be balanced across trees with an
exponential, support-uniform constant.  By
\Cref{lem:newV15-redistribution}, existence of some tree
decomposition is equivalent to the aggregate weighted estimate
\[
 \left\|\mathcal P_{\rm final}K_m\right\|
 \le
 K^{m-1}
 \sum_{T\in\mathfrak T_m}\prod_i b_i^{d_i(T)}.
\]
The inherited Prüfer identity rewrites the weight sum as
\[
 \left(\prod_i b_i\right)
 \left(\sum_i b_i\right)^{m-2}.
\]
What is still missing is not an arbitrary redistribution, which is
automatic once this bound is known, but a cancellation-preserving
proof of the aggregate estimate with an exponential base and a
balancing operation stable under predecessor products and the
single final payer.  The raw Bernoulli-tree triangle inequality
cannot provide it.
\end{assessment}

\section{V16 acquisition order and claim boundary}
\label{sec:newV15-next}

\begin{programme}[V16: cumulant-level tree balancing]
\label{prog:newV15-V16}
\begin{enumerate}[leftmargin=3em]
\item Compute the general lower-cumulant leakage of a raw
      random-cluster tree by Möbius inversion over edge cuts.
\item Define a tree-balancing operator on the complete family of raw
      forest integrals which annihilates every proper-partition
      component before norms are taken.
\item Require the balancing operator to commute with spectator
      amplification, isotypic compression, and the final Wilson
      payer.
\item Test its operator norm on stars and paths.  An exponential
      bound is admissible; a Bell, ordered-partition, or extra
      factorial bound is not.
\item Compare the balanced operator with the symmetry-first
      projected-tail recursion.  If their coefficients coincide,
      use the exact subset recursion rather than maintaining two
      expansions.
\item If the balancing norm is already factorial on a star, record
      the obstruction and return upstream to the generator-level
      projected recursion.
\end{enumerate}
\end{programme}

\begin{firewall}[Claim boundary after V15]
\label{fire:newV15-boundary}
V15 proves that the pointwise and integrated raw-tree PCMC proposed
in V14 are false at arity three, and it restores the valid ternary
tree card by exact all-tree compensation.  It does not prove an
all-order balancing bound, strong all-order local or global MTA,
WUFC, CPM, cutoff removal, reflection positivity,
Osterwalder--Schrader reconstruction, a continuum Yang--Mills
measure, or a positive Hamiltonian gap.  The full Yang--Mills mass
gap has not been proved.
\end{firewall}

\begin{theorem}[V15 result ledger]
\label{thm:newV15-results}
V15:
\begin{enumerate}[label=\textup{(V15.\arabic*)},leftmargin=6.5em]
\item performs the compulsory read-only SM/NS audit;
\item computes the exact integrated ternary random-cluster tree and
      its binary-covariance leakage;
\item gives an actual fundamental-representation counterexample to
      pointwise PCMC;
\item proves that even one parameter-integrated raw tree loses one
      Wilson clock;
\item gives the exact zero-sum compensation and recovers the proved
      ternary strong tree card; and
\item replaces raw-tree estimation by a cumulant-level balancing
      problem at arbitrary arity.
\end{enumerate}
\end{theorem}

\begin{proof}
Use
\Cref{lem:newV15-tree-leakage,
thm:newV15-pointwise-PCMC-false,
thm:newV15-integrated-PCMC-false,
lem:newV15-redistribution,
thm:newV15-compensated-ternary}.
\end{proof}

\paragraph{Parent-version record.}
V15 was formed from
\texttt{Renewal\_Yang\_Mills\_Mass\_Gap\_Research\_Notes\_V14.tex}.
The V14 parent had \(6\,255\) logical source lines,
\(204\,459\) bytes, and SHA--256
\[
 \texttt{6bf739306425190963302405402e3282d37e4be82789e390fb1f159004fcfd3a}.
\]
The next compulsory companion audit is V20.

% =============================================================================
% END APPENDED FIFTH-SERIES V15 MATERIAL
% =============================================================================

% =============================================================================
% BEGIN APPENDED FIFTH-SERIES V16 MATERIAL
% =============================================================================

\chapter{V16: Partition leakage and the exact balancing projection}
\label{chap:newV16-balancing}

\section{Cumulant coordinates of the random-cluster interpolation}
\label{sec:newV16-cumulant-coordinates}

V15 showed at arity three that a raw replica-tree integral is not
itself connected: proper-block cumulants occur in it and cancel only
after all labelled trees are summed.  We now compute this phenomenon
at every arity.

For \(\sigma\in\Pi([m])\), define the cumulant tensor
\begin{equation}
 {\mathfrak K}_\sigma
 :=
 \bigotimes_{B\in\sigma}
 \kappa^\nu_{|B|}\bigl((f_i)_{i\in B}\bigr).
 \label{eq:newV16-partition-cumulant}
\end{equation}
The tensor legs retain their original label order.  Write
\(\sigma\preceq\rho\) when every block of \(\sigma\) is contained
in a block of \(\rho\).  For the Bernoulli graph
\(G_{\boldsymbol x}\) of V14, put
\begin{equation}
 \Phi_\sigma(\boldsymbol x)
 :=
 \mathbb P_{\boldsymbol x}
 \{\sigma\preceq\pi(G_{\boldsymbol x})\}.
 \label{eq:newV16-connectivity-polynomial}
\end{equation}
Thus every block of \(\sigma\) must lie inside one random-graph
component, while different \(\sigma\)-blocks are allowed to merge.

\begin{theorem}[Exact cumulant-coordinate expansion]
\label{thm:newV16-cumulant-coordinate}
For every probability law \(\nu\), every bounded input family, and
every \(\boldsymbol x\),
\begin{equation}
 \boxed{\qquad
 F_m(\boldsymbol x)
 =
 \sum_{\sigma\in\Pi([m])}
 \Phi_\sigma(\boldsymbol x)\,
 {\mathfrak K}_\sigma.
 \qquad}
 \label{eq:newV16-cumulant-coordinate}
\end{equation}
\end{theorem}

\begin{proof}
Fix a deterministic graph \(G\), with component partition
\(\rho=\pi(G)\).  On every component \(C\in\rho\), the ordinary
moment--cumulant formula gives
\[
 \mathbb E_\nu\bigotimes_{i\in C}f_i
 =
 \sum_{\sigma_C\in\Pi(C)}
 \bigotimes_{B\in\sigma_C}
 \kappa^\nu_{|B|}\bigl((f_i)_{i\in B}\bigr).
\]
Multiplying these identities over the components \(C\) shows
\[
 M_G
 =
 \sum_{\sigma\preceq\pi(G)}{\mathfrak K}_\sigma.
\]
Average over \(G_{\boldsymbol x}\), interchange the two finite
sums, and recognize
\eqref{eq:newV16-connectivity-polynomial}.
\end{proof}

For \(T\in\mathfrak T_m\), let
\begin{align}
 {\cal J}_T
 &:=
 \int_{[0,1]^{E(T)}}
 \partial_{E(T)}F_m
 \bigl(\boldsymbol x^T(\boldsymbol s)\bigr)
 \,d\boldsymbol s,
 \label{eq:newV16-raw-tree}\\
 a_T(\sigma)
 &:=
 \int_{[0,1]^{E(T)}}
 \partial_{E(T)}\Phi_\sigma
 \bigl(\boldsymbol x^T(\boldsymbol s)\bigr)
 \,d\boldsymbol s.
 \label{eq:newV16-leakage-coefficient}
\end{align}

\begin{theorem}[General proper-partition leakage]
\label{thm:newV16-general-leakage}
Every raw tree has the universal expansion
\begin{equation}
 {\cal J}_T
 =
 \sum_{\sigma\in\Pi([m])}
 a_T(\sigma){\mathfrak K}_\sigma.
 \label{eq:newV16-raw-cumulant-expansion}
\end{equation}
Its coefficients obey
\begin{equation}
 \boxed{\qquad
 \sum_{T\in\mathfrak T_m}a_T(\sigma)
 =
 \begin{cases}
 1,&\sigma=\{[m]\},\\
 0,&\sigma\ne\{[m]\}.
 \end{cases}
 \qquad}
 \label{eq:newV16-leakage-zero-sum}
\end{equation}
For the discrete partition \(\widehat0\),
\begin{equation}
 a_T(\widehat0)=0
 \qquad(m\ge2).
 \label{eq:newV16-discrete-zero}
\end{equation}
Thus all nontrivial leakage comes from proper partitions having at
least one nonsingleton block, and each such leakage family has
exactly zero total over the labelled trees.
\end{theorem}

\begin{proof}
Differentiate
\eqref{eq:newV16-cumulant-coordinate} and integrate to obtain
\eqref{eq:newV16-raw-cumulant-expansion}.  Summing it over \(T\)
and applying the exact V14 cumulant forest identity gives
\[
 \sum_T{\cal J}_T
 =\kappa_m^\nu(f_1,\ldots,f_m)
 ={\mathfrak K}_{\{[m]\}}.
\]
The calculation is an identity in the universal
moment--cumulant polynomial algebra.  Equating the formal
\({\mathfrak K}_\sigma\) coefficients proves
\eqref{eq:newV16-leakage-zero-sum}.  Finally
\(\Phi_{\widehat0}\equiv1\), since every singleton is contained in
a graph component.  Its derivative of positive order is zero,
proving \eqref{eq:newV16-discrete-zero}.
\end{proof}

\begin{corollary}[V15 is the first proper-block instance]
\label{cor:newV16-V15-instance}
At \(m=3\), the only possible leakage partitions have type
\(2+1\).  Equation \eqref{eq:newV16-leakage-zero-sum} is exactly
the cancellation of the three binary-covariance rows in
\eqref{eq:newV15-leakage-cancels}.
\end{corollary}

\begin{proof}
The discrete \(1+1+1\) coefficient vanishes by
\eqref{eq:newV16-discrete-zero}; the top \(3\)-block is the desired
third cumulant.  The three remaining partitions are
\(12|3,13|2,23|1\), whose tensors are
\(C_{12}M_3,C_{13}M_2,C_{23}M_1\).
\end{proof}

\section{The rank-one balancing projection}
\label{sec:newV16-projection}

Let \(w_T>0\) be any desired family of tree weights and put
\[
 W_m:=\sum_{T\in\mathfrak T_m}w_T.
\]
For a family \(\boldsymbol A=(A_T)_{T\in\mathfrak T_m}\) in a
vector space, define
\begin{equation}
 (\mathscr B_w\boldsymbol A)_T
 :={w_T\over W_m}\sum_{S\in\mathfrak T_m}A_S.
 \label{eq:newV16-balancing-projection}
\end{equation}

\begin{theorem}[Exact tree-balancing projection]
\label{thm:newV16-balancing-projection}
The map \(\mathscr B_w\) has the following properties.
\begin{enumerate}[label=\textup{(\roman*)},leftmargin=3.7em]
\item It preserves the total:
      \[
      \sum_T(\mathscr B_w\boldsymbol A)_T=\sum_TA_T.
      \]
\item It is idempotent:
      \(\mathscr B_w^2=\mathscr B_w\).
\item Its kernel is the space of zero-total families.
\item Its range consists of the families
      \(A_T=w_TA/W_m\) for a single vector \(A\).
\item For every common linear map \(L\),
      \[
      L(\mathscr B_w\boldsymbol A)
      =\mathscr B_w(L\boldsymbol A).
      \]
\item No star or path coefficient is amplified:
      \[
      0<{w_T\over W_m}\le1,
      \qquad
      \sum_T{w_T\over W_m}=1.
      \]
\end{enumerate}
\end{theorem}

\begin{proof}
Each assertion follows directly from
\eqref{eq:newV16-balancing-projection}.  For idempotence, use
preservation of the total in the second application.  A family is
killed precisely when its total is zero.  Common linear maps pass
through the finite sum and scalar coefficients.  The last assertion
uses positivity of the weights.
\end{proof}

\begin{corollary}[All proper cumulant blocks are removed exactly]
\label{cor:newV16-balancing-removes-leakage}
Apply \(\mathscr B_w\) to the raw family
\(\boldsymbol{\cal J}=({\cal J}_T)_T\).  Then
\begin{equation}
 (\mathscr B_w\boldsymbol{\cal J})_T
 ={w_T\over W_m}
 \kappa_m^\nu(f_1,\ldots,f_m).
 \label{eq:newV16-balanced-cumulant}
\end{equation}
In the cumulant-coordinate expansion, every proper-partition row is
annihilated and the top-partition coefficient becomes \(w_T/W_m\).
\end{corollary}

\begin{proof}
The V14 identity gives
\(\sum_S{\cal J}_S=\kappa_m^\nu\).  Substitute this into
\eqref{eq:newV16-balancing-projection}.  Alternatively,
\Cref{thm:newV16-general-leakage} puts every proper-partition
coefficient family in the kernel and gives total one to the top
partition.
\end{proof}

\begin{corollary}[Compatibility with common Wilson output operations]
\label{cor:newV16-output-compatibility}
The balancing projection commutes with every output operation that
is linear and independent of the tree label, including:
\begin{enumerate}[label=\textup{(\alph*)},leftmargin=3.5em]
\item tensoring by a fixed spectator operator;
\item a fixed isotypic compression;
\item raw multiplicity pinching; and
\item the single final reduced Wilson resolvent on the assembled
      output.
\end{enumerate}
\end{corollary}

\begin{proof}
Each listed operation is a common linear map on the output family.
Apply part \textup{(v)} of
\Cref{thm:newV16-balancing-projection}.  No statement is made here
about a tree-dependent intermediate resolvent, which remains
forbidden.
\end{proof}

\begin{firewall}[Nonlinear predecessor assembly is not automatic]
\label{fire:newV16-product-compatibility}
The preceding corollary concerns common linear output maps.
Multiplication of separately balanced predecessor components is a
bilinear operation on different tree families.  Its compatibility
with a later quotient-tree weight does not follow from
\Cref{thm:newV16-balancing-projection}.  Any use inside the V11
relative product-cumulant assembly must therefore balance only after
the complete relative cumulant is formed, or prove a separate
block-factorization theorem.
\end{firewall}

\section{Equivalence with the aggregate strong card}
\label{sec:newV16-aggregate}

For all-thermal Wilson inputs take
\begin{equation}
 w_T:=\prod_{i=1}^m b_i^{d_i(T)}.
 \label{eq:newV16-Wilson-tree-weight}
\end{equation}
The exact Prüfer identity inherited in the active note gives
\begin{equation}
 W_m
 =
 \sum_{T\in\mathfrak T_m}\prod_i b_i^{d_i(T)}
 =
 \left(\prod_i b_i\right)
 \left(\sum_i b_i\right)^{m-2}.
 \label{eq:newV16-Prufer-weight}
\end{equation}

\begin{theorem}[Strong tree decomposition is equivalent to one
aggregate estimate]
\label{thm:newV16-aggregate-equivalence}
For a fixed \(m\), the following statements are equivalent up to the
same constant \(C_m\).
\begin{enumerate}[label=\textup{(\roman*)},leftmargin=3.7em]
\item There is an exact decomposition
      \(K_m=\sum_TK_{m,T}\) satisfying
      \[
      \|K_{m,T}\|\le C_mw_T
      \quad\text{for every }T.
      \]
\item The complete cumulant satisfies
      \begin{equation}
      \|K_m\|
      \le C_mW_m
      =
      C_m
      \left(\prod_i b_i\right)
      \left(\sum_i b_i\right)^{m-2}.
      \label{eq:newV16-aggregate-card}
      \end{equation}
\end{enumerate}
If \textup{(ii)} holds, the balanced family
\[
 K_{m,T}:={w_T\over W_m}K_m
\]
proves \textup{(i)}.
\end{theorem}

\begin{proof}
Sum the bounds in \textup{(i)} and use the triangle inequality to
obtain \textup{(ii)}.  Conversely use
\Cref{lem:newV15-redistribution} and
\eqref{eq:newV16-Prufer-weight}.
\end{proof}

\begin{assessment}[What the random-cluster formula does and does not
buy]
\label{ass:newV16-random-cluster-scope}
The random-cluster formula gives an exact, density-free
representation and identifies the entire proper-partition leakage
kernel.  The balancing projection removes that kernel with no
combinatorial loss and commutes with the final linear Wilson
operations.  However, after balancing,
\eqref{eq:newV16-balanced-cumulant} is simply the complete cumulant
redistributed by its desired weights.  It does not prove the
aggregate analytic estimate
\eqref{eq:newV16-aggregate-card}.

Consequently the V13 ISWFC definition, if it asks only for the
existence of some exact tree decomposition, is equivalent to the
aggregate bound and not a stronger constructive certificate.  The
hard problem remains to replace the current factorial-grade
aggregate estimate by
\[
 C_m\le K^{m-1}.
\]
Raw forest-edge iteration cannot do this because V15 proves that its
individual terms have the wrong weak-coupling order.
\end{assessment}

\section{Comparison with the symmetry-first recursion}
\label{sec:newV16-subset-comparison}

The archived projected-tail recursion
\[
 \mathfrak P_\alpha(I)
 =
 \sum_{i\in I}
 \mathcal D_\alpha
 \bigl(F_i,\mathfrak P_\alpha(I\setminus\{i\})\bigr)
\]
centers every contraction before proceeding and therefore removes
proper-partition rows dynamically.  The rank-one projection
\(\mathscr B_w\) removes the same rows only after the complete raw
tree family is assembled.  These are not termwise identical
expansions, but they have the same essential cancellation boundary:
no proper cumulant block may be normed before its zero-sum assembly.

\begin{proposition}[The balancing map itself has exponential grade]
\label{prop:newV16-balancing-grade}
The coefficients of \(\mathscr B_w\) have total variation one,
independently of \(m\).  Hence this algebraic compensation introduces
neither a factorial nor an exponential loss.  Any factorial in a
proof of \eqref{eq:newV16-aggregate-card} must arise from the
analytic estimate of the complete cumulant or from expanding it
before balancing, not from the balancing map.
\end{proposition}

\begin{proof}
For each output tree, the map uses the nonnegative scalar
\(w_T/W_m\), and these scalars sum to one over \(T\).  This is exactly
part \textup{(vi)} of
\Cref{thm:newV16-balancing-projection}.
\end{proof}

\begin{assessment}[Direction selected]
\label{ass:newV16-direction}
The density-free forest route has now reached its exact limit: it
supplies algebraic connectivity but not the missing analytic clock.
The next attack returns upstream to the symmetry-first projected
recursion.  Its current \(m!\) was produced by summing ordered
contraction histories separately.  The target is to group histories
by their unordered subset state before applying Hölder and the
resolvent estimate, and to determine whether the recursion admits a
majorant with exponential rather than factorial constants.
\end{assessment}

\section{V17 programme and claim boundary}
\label{sec:newV16-next}

\begin{programme}[V17: unordered subset majorant]
\label{prog:newV16-V17}
\begin{enumerate}[leftmargin=3em]
\item Normalize the projected tail by the exact number of orderings
      already present in its subset recursion, without changing the
      cumulant compression identity.
\item Derive the sharp scalar recurrence for the normalized
      negative norm before inserting any common upper bound.
\item Use the exact weighted subset convolution
      \[
      \sum_{i\in I}b_i
      \prod_{j\ne i}b_j
      \left(\sum_{j\ne i}b_j\right)^{|I|-2}
      \]
      and compare it with the Prüfer weight
      \((\prod_i b_i)(\sum_i b_i)^{|I|-1}\).
\item Keep the cyclic or subset sum inside the column norm whenever
      the tangent leading symbol cancels; do not apply the
      triangle inequality to its ordered histories.
\item Prove the recurrence first on the exact Gaussian
      Ornstein--Uhlenbeck model, then isolate the Wilson generator
      defect needed to transfer it.
\item If the normalized recursion still has factorial growth,
      exhibit the first star or path family forcing it and return to
      a semigroup interpolation of the complete cumulant.
\end{enumerate}
\end{programme}

\begin{firewall}[Claim boundary after V16]
\label{fire:newV16-boundary}
V16 computes the all-arity proper-partition leakage of the
random-cluster forest family and removes it by an exact norm-one
balancing projection.  It proves the equivalence between a strong
tree decomposition and the aggregate Prüfer-weighted cumulant
estimate.  It does not prove that aggregate estimate with an
exponential base, all-order strong MTA, WUFC, CPM, cutoff removal,
reflection positivity, Osterwalder--Schrader reconstruction, a
continuum Yang--Mills measure, or a positive Hamiltonian gap.  The
full Yang--Mills mass gap has not been proved.
\end{firewall}

\begin{theorem}[V16 result ledger]
\label{thm:newV16-results}
V16:
\begin{enumerate}[label=\textup{(V16.\arabic*)},leftmargin=6.5em]
\item expands the random-cluster interpolation exactly in
      partition-cumulant coordinates;
\item proves that every proper-partition coefficient has zero total
      across labelled trees;
\item constructs the exact rank-one tree-balancing projection;
\item proves its compatibility with common linear Wilson output
      operations and states the nonlinear product firewall;
\item proves the equivalence of the strong tree card and the
      aggregate Prüfer-weighted bound; and
\item returns the factorial-removal attack to the unordered
      symmetry-first subset recursion.
\end{enumerate}
\end{theorem}

\begin{proof}
Use
\Cref{thm:newV16-cumulant-coordinate,
thm:newV16-general-leakage,
thm:newV16-balancing-projection,
cor:newV16-output-compatibility,
thm:newV16-aggregate-equivalence,
prop:newV16-balancing-grade}.
\end{proof}

\paragraph{Parent-version record.}
V16 was formed from
\texttt{Renewal\_Yang\_Mills\_Mass\_Gap\_Research\_Notes\_V15.tex}.
The V15 parent had \(6\,785\) logical source lines,
\(221\,304\) bytes, and SHA--256
\[
 \texttt{800a5a73a79c6b36d8a1ad164b43acfa240f0032869600cd840f41d2459e9c63}.
\]
The next compulsory companion audit is V20.

% =============================================================================
% END APPENDED FIFTH-SERIES V16 MATERIAL
% =============================================================================

% =============================================================================
% BEGIN APPENDED FIFTH-SERIES V17 MATERIAL
% =============================================================================

\chapter{V17: Normalized subset recursion and the missing clock}
\label{chap:newV17-subset}

\section{The strong certificate required by V16}
\label{sec:newV17-strong-certificate}

For a nonempty label set \(I\), retain the exact symmetrized
projected tail
\[
 \mathfrak P_\alpha(I)
 =
 \sum_{\pi\in S_I}
 \mathcal P_\alpha
 (D^{R_{\pi(1)}},\ldots,D^{R_{\pi(|I|)}})
\]
and put
\[
 b_i=\sqrt{s_\alpha(1+C_2(R_i))}<1,
 \qquad
 B_I=\sum_{i\in I}b_i.
\]
The factorial-graded RSPCC recalled in V1 and corrected in V12 is
not sufficient for the strong fixed-tree grade.  The exact stronger
target is the following.

\begin{definition}[Strong RSPCC]
\label{def:newV17-SRSPCC}
The strong resolvent-norm symmetrized projected-chain certificate
holds if
\begin{equation}
 \boxed{\qquad
 \|\mathfrak P_\alpha(I)\|_{-1,c}
 \le
 {K^{|I|}\over\sqrt\alpha}
 \left(\prod_{i\in I}b_i\right)
 B_I^{|I|-1}
 \qquad}
 \label{eq:newV17-SRSPCC}
\end{equation}
with \(K\) independent of \(I,\alpha\), all representations, and
spectator legs.
\end{definition}

\begin{theorem}[Strong RSPCC gives the aggregate strong card]
\label{thm:newV17-SRSPCC-sufficient}
If \eqref{eq:newV17-SRSPCC} holds, then
\begin{equation}
 \left\|
 \kappa_m^{\nu_\alpha}
 (D^{R_1},\ldots,D^{R_m})
 \right\|_{\rm op}
 \le
 K_1^{m-1}
 \left(\prod_{i=1}^mb_i\right)
 \left(\sum_{i=1}^mb_i\right)^{m-2}.
 \label{eq:newV17-aggregate-strong}
\end{equation}
Consequently the V16 aggregate card, and hence the strong
all-thermal local tree decomposition, follow without an analytic
factorial.
\end{theorem}

\begin{proof}
Use the exact symmetry-first compression
\[
 \kappa_m
 =
 {1\over m}\sum_{i=1}^m
 \mathbb E_{\nu_\alpha}
 \mathcal B_\alpha
 \left(
 D^{R_i},
 \mathfrak P_\alpha([m]\setminus\{i\})
 \right).
\]
The inherited terminal negative-norm pairing gives
\[
 \left\|
 \mathbb E\mathcal B_\alpha(D^{R_i},G)
 \right\|_{\rm op}
 \le C\sqrt\alpha\,b_i\|G\|_{-1,c}.
\]
Apply \eqref{eq:newV17-SRSPCC} to the set
\(I_i=[m]\setminus\{i\}\):
\begin{align*}
 \|\kappa_m\|_{\rm op}
 &\le
 {CK^{m-1}\over m}
 \left(\prod_{j=1}^mb_j\right)
 \sum_{i=1}^m
 (B_{[m]}-b_i)^{m-2}\\
 &\le
 CK^{m-1}
 \left(\prod_jb_j\right)
 B_{[m]}^{m-2}.
\end{align*}
Absorb the fixed \(C\) in the exponential base.  Now apply
\Cref{thm:newV16-aggregate-equivalence}.
\end{proof}

\begin{firewall}[Exact grade correction]
\label{fire:newV17-grade}
The earlier RSPCC has \(|I|!K^{|I|}\) on its right side and yields
the FGMTA grade identified in V12.  The sufficient certificate
\eqref{eq:newV17-SRSPCC} has no \(|I|!\).  These statements must not
be conflated.
\end{firewall}

\section{Removing the visible ordering factorial}
\label{sec:newV17-normalized}

Define the normalized tail
\begin{equation}
 \widehat{\mathfrak P}_\alpha(I)
 :={1\over |I|!}\mathfrak P_\alpha(I).
 \label{eq:newV17-normalized-tail}
\end{equation}

\begin{lemma}[Exact averaged subset recursion]
\label{lem:newV17-averaged-recursion}
For a singleton,
\(\widehat{\mathfrak P}_\alpha(\{i\})=Q_\alpha D^{R_i}\).
For \(r=|I|\ge2\),
\begin{equation}
 \boxed{\qquad
 \widehat{\mathfrak P}_\alpha(I)
 =
 {1\over r}\sum_{i\in I}
 \mathcal D_\alpha
 \left(
 D^{R_i},
 \widehat{\mathfrak P}_\alpha(I\setminus\{i\})
 \right).
 \qquad}
 \label{eq:newV17-averaged-recursion}
\end{equation}
\end{lemma}

\begin{proof}
Divide the exact subset recursion by \(r!\) and use
\[
 \mathfrak P_\alpha(I\setminus\{i\})
 =(r-1)!\widehat{\mathfrak P}_\alpha(I\setminus\{i\}).
\]
\end{proof}

\begin{lemma}[The factorial moves into cumulant compression]
\label{lem:newV17-normalized-compression}
For \(m\ge2\),
\begin{equation}
 \kappa_m
 =
 {(m-1)!\over m}
 \sum_{i=1}^m
 \mathbb E\mathcal B_\alpha
 \left(
 D^{R_i},
 \widehat{\mathfrak P}_\alpha([m]\setminus\{i\})
 \right).
 \label{eq:newV17-normalized-compression}
\end{equation}
\end{lemma}

\begin{proof}
Substitute
\(\mathfrak P_\alpha(I_i)=(m-1)!
\widehat{\mathfrak P}_\alpha(I_i)\)
in the exact symmetry-first cumulant formula.
\end{proof}

\begin{corollary}[The normalized strong target]
\label{cor:newV17-normalized-target}
The SRSPCC is equivalent to
\begin{equation}
 \|\widehat{\mathfrak P}_\alpha(I)\|_{-1,c}
 \le
 {K^{|I|}\over |I|!\sqrt\alpha}
 \left(\prod_{i\in I}b_i\right)
 B_I^{|I|-1}.
 \label{eq:newV17-normalized-target}
\end{equation}
Thus division by the number of orderings does not by itself remove
the required factorial gain: the normalized tail must be smaller by
the same factor because
\eqref{eq:newV17-normalized-compression} puts it back at the outer
compression.
\end{corollary}

\begin{proof}
Divide \eqref{eq:newV17-SRSPCC} by \(|I|!\).
\end{proof}

\section{The sharp scalar recurrence}
\label{sec:newV17-scalar}

The exact averaged recursion removes a purely combinatorial
factorial under a termwise estimate, but it does not create the
Wilson clocks required by the Prüfer weight.

\begin{proposition}[Termwise normalized recurrence]
\label{prop:newV17-termwise}
Suppose only the inherited atomic-left inequality is used separately
on the \(r\) summands in
\eqref{eq:newV17-averaged-recursion}.  Then
\begin{equation}
 \|\widehat{\mathfrak P}_\alpha(I)\|_{-1,c}
 \le
 C^{|I|}\prod_{i\in I}b_i
 \label{eq:newV17-normalized-triangle}
\end{equation}
up to the fixed singleton normalization.  This bound is exponential
in arity but contains none of the \(B_I^{|I|-1}\) gain.
\end{proposition}

\begin{proof}
The singleton bound has the displayed form.  If it holds on every
set of size \(r-1\), then
\begin{align*}
 \|\widehat{\mathfrak P}_\alpha(I)\|_{-1,c}
 &\le
 {C\over r}\sum_{i\in I}
 b_i
 \|\widehat{\mathfrak P}_\alpha(I\setminus\{i\})\|_{-1,c}\\
 &\le
 {C^r\over r}\sum_{i\in I}
 b_i\prod_{j\ne i}b_j
 =C^r\prod_{j\in I}b_j.
\end{align*}
\end{proof}

\begin{lemma}[Exact weighted subset convolution]
\label{lem:newV17-subset-convolution}
For \(r=|I|\ge2\),
\begin{equation}
 \sum_{i\in I}
 b_i
 \left(\prod_{j\ne i}b_j\right)
 (B_I-b_i)^{r-2}
 =
 \left(\prod_{j\in I}b_j\right)
 \sum_{i\in I}(B_I-b_i)^{r-2}.
 \label{eq:newV17-subset-convolution}
\end{equation}
Moreover,
\begin{equation}
 \sum_{i\in I}(B_I-b_i)^{r-2}
 \le rB_I^{r-2}.
 \label{eq:newV17-subset-upper}
\end{equation}
When all \(b_i=b>0\), the ratio of the left scalar sum in
\eqref{eq:newV17-subset-upper} to \(B_I^{r-1}\) is
\begin{equation}
 {r((r-1)b)^{r-2}\over (rb)^{r-1}}
 =
 {1\over b}
 \left(1-{1\over r}\right)^{r-2}.
 \label{eq:newV17-missing-clock-ratio}
\end{equation}
\end{lemma}

\begin{proof}
In each summand,
\(b_i\prod_{j\ne i}b_j=\prod_jb_j\), proving the identity.
Since \(B_I-b_i\le B_I\), summing gives
\eqref{eq:newV17-subset-upper}.  Direct substitution of
\(B_I=rb\) proves \eqref{eq:newV17-missing-clock-ratio}.
\end{proof}

\begin{theorem}[Normalization alone cannot prove SRSPCC]
\label{thm:newV17-normalization-no-go}
Assume inductively the normalized strong target
\eqref{eq:newV17-normalized-target} on all proper subsets and apply
the atomic-left bound separately to the summands of
\eqref{eq:newV17-averaged-recursion}.  The resulting scalar
condition would require
\begin{equation}
 \sum_{i\in I}(B_I-b_i)^{r-2}
 \le K B_I^{r-1}
 \label{eq:newV17-false-scalar}
\end{equation}
with a constant independent of the marks.  This is false already
when all \(b_i=b\downarrow0\).  Hence averaging the orderings removes
the visible \(r!\) but still loses exactly one small Wilson clock at
each subset step.
\end{theorem}

\begin{proof}
Insert \eqref{eq:newV17-normalized-target} for
\(I\setminus\{i\}\) into the termwise estimate.  Since
\(r(r-1)!=r!\), the desired conclusion reduces precisely to
\eqref{eq:newV17-false-scalar}.  Equation
\eqref{eq:newV17-missing-clock-ratio} shows that the best constant in
that inequality diverges like \(b^{-1}\).
\end{proof}

\begin{assessment}[Factorial and clock losses are independent]
\label{ass:newV17-two-losses}
The averaged recursion proves that the elementary ordering count can
be normalized to exponential grade.  The normalized cumulant
compression puts that factorial back unless the analytic tail is
smaller by \(1/r!\).  Independently,
\Cref{thm:newV17-normalization-no-go} shows that termwise estimates
lose one power of \(B_I\) at every recursion.  A valid strong proof
must therefore use the assembled subset cancellation both to supply
the missing clock and to control the normalized tail at the
\(1/r!\) scale required by the outer cumulant formula.
\end{assessment}

\section{The tangent cancellation and the Wilson defect}
\label{sec:newV17-tangent}

The missing clock has a known local source.  In the Gaussian
Ornstein--Uhlenbeck tangent model, the contraction of two linear
representation forms is constant.  The outer centering in
\(\mathcal D_0\) kills it exactly.  This is the inherited result
\texttt{prop:V89-Gaussian-linear}.

\begin{proposition}[All-linear tangent tails vanish]
\label{prop:newV17-linear-tangent}
If every input to the tangent subset recursion is linear, then
\begin{equation}
 \mathfrak P_0(I)=0
 \qquad(|I|\ge2).
 \label{eq:newV17-linear-tangent-zero}
\end{equation}
\end{proposition}

\begin{proof}
For \(|I|=2\), each projected contraction is the centered part of a
constant and is zero by the inherited linear-pair identity.  If the
claim holds for sets of size \(r-1\), every term in the exact subset
recursion at size \(r\) has zero second argument.  Induction proves
the result.
\end{proof}

\begin{assessment}[Meaning of the extra clock]
\label{ass:newV17-clock-source}
Equation \eqref{eq:newV17-linear-tangent-zero} shows that the
termwise bound is estimating a leading symbol whose assembled value
is exactly zero.  Every nonzero Wilson contribution must contain at
least one nonlinear representation remainder or one
Wilson-versus-tangent generator defect.  Either can in principle
supply the missing factor \(B_I\).  The fixed-arity archive verifies
this mechanism through the proved subset frontier, but its constants
have not been controlled uniformly in \(r\).
\end{assessment}

\begin{definition}[Uniform assembled clock-gain gate]
\label{def:newV17-UACG}
The UACG is the assertion that the complete one-form subset sum
satisfies an arity-uniform estimate which, after insertion of the
SRSPCC bounds on proper subsets, gains one additional factor \(B_I\)
relative to the termwise scalar convolution
\eqref{eq:newV17-subset-convolution}, with an exponential constant
and no factorial.
\end{definition}

\begin{theorem}[UACG plus factorial-normalized propagation is the
remaining local gate]
\label{thm:newV17-UACG-reduction}
Suppose the UACG holds at every subset size and its recurrence for
the normalized tails has coefficients whose total majorant is
bounded by \(K^r/r!\).  Then the SRSPCC holds, hence so do the
aggregate strong card and the all-thermal strong local tree card.
\end{theorem}

\begin{proof}
The UACG supplies the missing \(B_I\) at the step isolated in
\Cref{thm:newV17-normalization-no-go}.  The stated normalized
coefficient majorant supplies the \(1/r!\) in
\eqref{eq:newV17-normalized-target}.  Induction on \(|I|\) proves
that target.  Apply
\Cref{thm:newV17-SRSPCC-sufficient}.
\end{proof}

\begin{firewall}[The UACG is not proved]
\label{fire:newV17-UACG-open}
\Cref{thm:newV17-UACG-reduction} is conditional.  The tangent
identity proves exact cancellation only for the leading all-linear
symbol.  It does not by itself control the complete Wilson
generator defect uniformly at arbitrary arity, nor does it prove the
\(1/r!\) normalized coefficient majorant.
\end{firewall}

\section{V18 programme and claim boundary}
\label{sec:newV17-next}

\begin{programme}[V18: Duhamel placement of the first defect]
\label{prog:newV17-V18}
\begin{enumerate}[leftmargin=3em]
\item Expand the complete normalized subset tail around the
      all-linear tangent recursion by a telescoping Duhamel identity,
      marking the first nonlinear or generator-defect occurrence.
\item Keep the sum over its possible locations assembled.  Do not
      bound the location count by \(r\) before using the
      \(1/r\) in \eqref{eq:newV17-averaged-recursion}.
\item Prove that the first defect supplies one of the marks in
      \(B_I\) and that the remaining normalized prefix and suffix
      weights form a binomial convolution, not an ordered
      permutation count.
\item Determine whether the binomial denominators yield the needed
      \(1/r!\) scale or only an exponential normalized scale.
\item Work first in the exact Gaussian model with exponential
      representation observables, where the generator is fixed and
      only representation nonlinearities remain.
\item Transfer to the Wilson law only after isolating a
      dimension-free one-defect resolvent estimate; avoid a global
      density comparison, ruled out in V13.
\end{enumerate}
\end{programme}

\begin{firewall}[Claim boundary after V17]
\label{fire:newV17-boundary}
V17 identifies the exact factorial-free SRSPCC sufficient for the
strong local card, derives the normalized subset recursion and its
outer compression, and proves that normalization plus termwise
estimation still loses one Wilson clock.  It proves the vanishing of
the all-linear tangent tails and isolates the UACG as the analytic
defect estimate.  It does not prove UACG, SRSPCC, all-order strong
MTA, WUFC, CPM, cutoff removal, reflection positivity,
Osterwalder--Schrader reconstruction, a continuum Yang--Mills
measure, or a positive Hamiltonian gap.  The full Yang--Mills mass
gap has not been proved.
\end{firewall}

\begin{theorem}[V17 result ledger]
\label{thm:newV17-results}
V17:
\begin{enumerate}[label=\textup{(V17.\arabic*)},leftmargin=6.5em]
\item defines the factorial-free SRSPCC and proves that it implies
      the V16 aggregate strong card;
\item derives the exact averaged subset recursion and normalized
      cumulant compression;
\item proves the sharp weighted subset convolution;
\item proves that ordering normalization alone loses exactly one
      Wilson clock per subset step;
\item proves all higher all-linear tangent tails vanish; and
\item reduces the next analytic step to a first-defect assembled
      clock gain with factorial-normalized propagation.
\end{enumerate}
\end{theorem}

\begin{proof}
Use
\Cref{thm:newV17-SRSPCC-sufficient,
lem:newV17-averaged-recursion,
lem:newV17-normalized-compression,
lem:newV17-subset-convolution,
thm:newV17-normalization-no-go,
prop:newV17-linear-tangent,
thm:newV17-UACG-reduction}.
\end{proof}

\paragraph{Parent-version record.}
V17 was formed from
\texttt{Renewal\_Yang\_Mills\_Mass\_Gap\_Research\_Notes\_V16.tex}.
The V16 parent had \(7\,267\) logical source lines,
\(237\,766\) bytes, and SHA--256
\[
 \texttt{5bf67f775d8d95839630a78d62ce662de2460cd6732de449316a28af812ed1a5}.
\]
The next compulsory companion audit is V20.

% =============================================================================
% END APPENDED FIFTH-SERIES V17 MATERIAL
% =============================================================================

% =============================================================================
% BEGIN APPENDED FIFTH-SERIES V18 MATERIAL
% =============================================================================

\chapter{V18: First-defect Duhamel expansion and its exact limit}
\label{chap:newV18-Duhamel}

\section{Ordered contractions as products of linear maps}
\label{sec:newV18-ordered}

V17 proposed marking the first nonlinear or generator-defect
occurrence in the normalized subset recursion.  Before estimating
that occurrence, one must determine whether the resulting
combinatorics can possibly supply the factorial-normalized scale in
\eqref{eq:newV17-normalized-target}.  This chapter carries out that
test exactly.

Fix an ordered label list
\(\pi=(\pi(1),\ldots,\pi(r))\).  After the innermost input has been
chosen, an ordered projected chain is a composition of \(r-1\)
linear contraction maps acting on that terminal input.  Abstractly
write
\begin{equation}
 {\cal C}^{W}_\pi
 =
 A_{\pi,1}\cdots A_{\pi,r-1}v_{\pi,r},
 \qquad
 {\cal C}^{0}_\pi
 =
 A^0_{\pi,1}\cdots A^0_{\pi,r-1}v^0_{\pi,r}.
 \label{eq:newV18-two-chains}
\end{equation}
The superscripts \(W,0\) denote the Wilson and chosen tangent
reference chains.  All maps act in their displayed order; no
commutativity is assumed.

\begin{lemma}[Noncommutative telescoping identity]
\label{lem:newV18-telescoping}
For linear maps \(A_1,\ldots,A_n,B_1,\ldots,B_n\) with compatible
domains,
\begin{equation}
 A_1\cdots A_n-B_1\cdots B_n
 =
 \sum_{k=1}^n
 A_1\cdots A_{k-1}
 (A_k-B_k)
 B_{k+1}\cdots B_n.
 \label{eq:newV18-telescoping}
\end{equation}
\end{lemma}

\begin{proof}
Insert and subtract successively
\[
 A_1\cdots A_kB_{k+1}\cdots B_n,
 \qquad k=0,\ldots,n.
\]
The resulting sum telescopes and preserves the order of every
factor.
\end{proof}

Terminal-vector replacement adds one further row:
\begin{align}
 {\cal C}^{W}_\pi-{\cal C}^{0}_\pi
 &=
 \sum_{k=1}^{r-1}
 A_{\pi,1}\cdots A_{\pi,k-1}
 (A_{\pi,k}-A^0_{\pi,k})
 A^0_{\pi,k+1}\cdots A^0_{\pi,r-1}v^0_{\pi,r}
 \notag\\
 &\quad+
 A_{\pi,1}\cdots A_{\pi,r-1}
 (v_{\pi,r}-v^0_{\pi,r}).
 \label{eq:newV18-first-defect-expansion}
\end{align}
This is an exact one-defect Duhamel expansion: Wilson maps lie before
the marked defect and reference maps lie after it.

\section{The total-variation census}
\label{sec:newV18-census}

Let
\[
 \widehat{\mathfrak P}^{W}_\alpha(I)
 ={1\over r!}\sum_{\pi\in S_I}{\cal C}^{W}_\pi,
 \qquad
 \widehat{\mathfrak P}^{0}_\alpha(I)
 ={1\over r!}\sum_{\pi\in S_I}{\cal C}^{0}_\pi.
\]

\begin{theorem}[Exact first-defect coefficient mass]
\label{thm:newV18-defect-mass}
After substituting
\eqref{eq:newV18-first-defect-expansion} into the normalized
permutation average, there are exactly \(r!\,r\) marked rows,
counted with multiplicity, and each has coefficient \(1/r!\).
Consequently their total coefficient variation is exactly
\begin{equation}
 {r!\,r\over r!}=r.
 \label{eq:newV18-total-variation}
\end{equation}
For the unnormalized symmetrized tail the corresponding variation
is \(r!\,r\).
\end{theorem}

\begin{proof}
There are \(r!\) permutations.  For each permutation,
\eqref{eq:newV18-first-defect-expansion} has \(r-1\) map-defect
locations and one terminal-defect location.  Every coefficient in
the normalized average is \(1/r!\).  Summing their absolute values
proves \eqref{eq:newV18-total-variation}.  Multiplication by \(r!\)
gives the unnormalized statement.
\end{proof}

\begin{corollary}[A first-defect triangle bound cannot supply
\(1/r!\)]
\label{cor:newV18-no-factorial}
Suppose the tangent average vanishes and every marked row in
\eqref{eq:newV18-first-defect-expansion} is bounded separately by
the same exponential-per-label majorant \(K^r{\cal W}_I\).
Then the Duhamel triangle inequality gives only
\begin{equation}
 \|\widehat{\mathfrak P}^{W}_\alpha(I)\|
 \le rK^r{\cal W}_I.
 \label{eq:newV18-Duhamel-triangle}
\end{equation}
It cannot give the \(K_1^r{\cal W}_I/r!\) scale required by
\eqref{eq:newV17-normalized-target}.
\end{corollary}

\begin{proof}
The reference average is zero by hypothesis.  Multiply the common
row bound by the exact total variation
\eqref{eq:newV18-total-variation}.  Since
\(r\,r!\) cannot be absorbed into a fixed exponential base after
comparison with \(1/r!\), the required normalized scale does not
follow.
\end{proof}

\begin{firewall}[Scope of the census no-go]
\label{fire:newV18-census-scope}
\Cref{cor:newV18-no-factorial} does not disprove SRSPCC.  It proves
that marking one defect and then taking norms row by row is
insufficient.  Further cancellation among defect locations and
permutations, or a cumulant-level representation which never
creates the \(r!\) histories, is necessary.
\end{firewall}

\section{The missing clock is present but not enough}
\label{sec:newV18-clock}

Assume a marked nonlinear observable remainder at label \(j\) has
one extra factor \(b_j\), and a generator defect has an analogous
mark bounded by \(CB_I\).  Then every Duhamel row can indeed improve
the termwise V17 homogeneity by one Wilson clock.

\begin{proposition}[Clock gain and factorial gain are orthogonal]
\label{prop:newV18-orthogonal-gains}
Under the preceding one-defect bound, the normalized first-defect
triangle estimate can have the homogeneity
\begin{equation}
 rK^r
 \left(\prod_{i\in I}b_i\right)B_I^{r-1},
 \label{eq:newV18-correct-homogeneity}
\end{equation}
but still lacks the factor \(1/r!\) in
\eqref{eq:newV17-normalized-target}.  Thus a correct clock power
does not repair the ordering grade.
\end{proposition}

\begin{proof}
The marked defect supplies the one factor of \(B_I\) missing in
\Cref{thm:newV17-normalization-no-go}.  The other incidence powers
are those of the proper-subset majorant.  The coefficient census is
unchanged and contributes \(r\), not \(1/r!\), by
\Cref{thm:newV18-defect-mass}.
\end{proof}

\begin{assessment}[The proposed V18 route is too local]
\label{ass:newV18-route-limit}
The V17 UACG separated the correct two obligations, but a
first-defect expansion of individual ordered chains addresses only
the clock obligation.  It cannot address the factorial-normalized
obligation after rowwise norms.  Continuing to sharpen a
one-defect row estimate would therefore solve the wrong intermediate
problem.
\end{assessment}

\section{A binomial regrouping still does not yield factorial decay}
\label{sec:newV18-binomial}

One may group marked rows by the unordered set \(S\) of labels
preceding the defect.  This reduces storage from permutations to
subsets, but its coefficient variation remains polynomial rather
than factorially small.

\begin{lemma}[Predecessor-set census]
\label{lem:newV18-predecessor-census}
Fix a defect label \(j\) and an integer \(0\le k\le r-1\).
For a fixed set
\(S\subset I\setminus\{j\}\) with \(|S|=k\), the number of
permutations in which precisely the labels of \(S\) precede \(j\)
is
\begin{equation}
 k!(r-k-1)!.
 \label{eq:newV18-fixed-set-count}
\end{equation}
After normalization by \(r!\), the total coefficient of that fixed
predecessor set is
\begin{equation}
 {1\over r\binom{r-1}{k}}.
 \label{eq:newV18-fixed-set-weight}
\end{equation}
Summing over all \(\binom{r-1}{k}\) predecessor sets of size \(k\)
gives \(1/r\), and summing over the \(r\) possible values of \(k\)
gives one for the fixed defect label.
\end{lemma}

\begin{proof}
Order the labels in \(S\) in \(k!\) ways, place \(j\) next, and
order the remaining \(r-k-1\) labels in
\((r-k-1)!\) ways.  Division by
\[
 r!=r\binom{r-1}{k}k!(r-k-1)!
\]
proves \eqref{eq:newV18-fixed-set-weight}.  The two sums are
immediate.
\end{proof}

\begin{corollary}[Subset regrouping has unit mass per defect label]
\label{cor:newV18-subset-mass}
Even after exact binomial regrouping, the normalized coefficient
mass is one for each possible defect label and \(r\) after the
defect label is summed.  Binomial denominators prevent an additional
combinatorial explosion, but they do not supply \(1/r!\).
\end{corollary}

\begin{proof}
This is the final assertion of
\Cref{lem:newV18-predecessor-census}, summed over \(j\in I\).
\end{proof}

\section{Return to the complete cumulant}
\label{sec:newV18-cumulant}

The V16 aggregate card is a bound on the complete cumulant, not on a
particular projected-tail norm.  SRSPCC is sufficient but may be
strictly stronger because its negative norm tests the tail against a
larger class than the finite family of outer representation rows in
the cumulant compression.

\begin{proposition}[The reverse implication is unavailable]
\label{prop:newV18-no-converse}
The aggregate cumulant estimate
\eqref{eq:newV17-aggregate-strong} does not, by the identities proved
so far, imply SRSPCC.  Recovering the negative norm of
\(\mathfrak P_\alpha(I)\) would require a uniform dual density of
outer representation derivatives in the complete negative-norm
test space, and no such theorem has been proved.
\end{proposition}

\begin{proof}
The terminal pairing maps a projected tail \(G\) to the restricted
family
\[
 G\longmapsto
 \mathbb E\mathcal B_\alpha(D^R,G).
\]
The inherited estimate bounds these functionals by
\(\sqrt\alpha b_R\|G\|_{-1,c}\), which is one-sided.  None of the
archived results supplies the reverse norming inequality uniformly
over \(R\), spectator legs, and output compression.  Therefore a
bound on all displayed cumulant pairings cannot be inverted to
SRSPCC.
\end{proof}

\begin{assessment}[Upstream correction]
\label{ass:newV18-upstream}
V17's SRSPCC remains a valid sufficient condition, but
\Cref{thm:newV18-defect-mass} shows that its proposed
first-defect proof does not remove the factorial.  The more economical
target is the aggregate cumulant itself.  Its connected-graph
representation already has the correct census in the Gaussian
tangent model; the Wilson problem is to compare complete connected
objects, rather than their ordered projected tails.
\end{assessment}

\section{The next exact comparison object}
\label{sec:newV18-next-object}

Let \(K_m^W\) denote the complete intrinsic Wilson cumulant and let
\(K_m^0\) denote the tangent Gaussian cumulant under the canonical
scaled representation linearization.  The frozen Gaussian forest
theorem gives \(K_m^0\) the strong aggregate Prüfer grade.  Define
the complete connected defect
\begin{equation}
 {\cal E}_m
 :=K_m^W-\mathcal U_\alpha K_m^0,
 \label{eq:newV18-connected-defect}
\end{equation}
where \(\mathcal U_\alpha\) is only a notation for the fixed scaled
coordinate identification on the principal chart.  It is not a
global density transport.

\begin{definition}[Connected Wilson--tangent defect card]
\label{def:newV18-CWTDC}
The CWTDC is
\begin{equation}
 \|{\cal E}_m\|_{\rm op}
 \le
 K^{m-1}
 \left(\prod_i b_i\right)
 \left(\sum_i b_i\right)^{m-2},
 \label{eq:newV18-CWTDC}
\end{equation}
with the chart complement, exact Haar factor, and final Wilson payer
included before the norm.
\end{definition}

\begin{theorem}[CWTDC is sufficient for the all-thermal aggregate
card]
\label{thm:newV18-CWTDC-sufficient}
If \eqref{eq:newV18-CWTDC} holds and the canonical tangent
identification preserves the inherited Gaussian strong card with an
exponential base, then
\eqref{eq:newV17-aggregate-strong} holds.
\end{theorem}

\begin{proof}
Use \eqref{eq:newV18-connected-defect} and the triangle inequality
only between the two \emph{complete connected} objects.  Bound the
Gaussian term by its inherited strong aggregate card and the defect
by \eqref{eq:newV18-CWTDC}.  Enlarge the common exponential base.
\end{proof}

\begin{firewall}[No density transfer is reintroduced]
\label{fire:newV18-no-density}
The CWTDC must not be proved by multiplying one Wilson-to-Gaussian
density ratio per replica.  V13 proves that such a product leaves
the uniform \(L^p\) window.  A valid comparison must use one
intrinsic Wilson integration, a generator interpolation, or an
exact connected-coordinate identity in which the density is
assembled once.
\end{firewall}

\section{V19 programme and claim boundary}
\label{sec:newV18-next}

\begin{programme}[V19: connected generator interpolation]
\label{prog:newV18-V19}
\begin{enumerate}[leftmargin=3em]
\item Define a one-parameter intrinsic generator interpolation
      between the scaled tangent diffusion and the exact Wilson
      one-link diffusion, with its invariant law carried along once.
\item Differentiate the \emph{complete cumulant} with respect to the
      interpolation parameter before expanding any replica product.
\item Express the derivative by one marked generator-defect
      insertion in a connected cumulant, not by \(m!\) ordered
      chains.
\item Prove the exact cumulant derivative identity first for bounded
      finite-dimensional test matrices.
\item Determine whether the marked defect supplies one extra Wilson
      clock and whether the remaining connected Gaussian forest
      retains its exponential base.
\item Treat the antipodal chart complement intrinsically; do not
      use a replica-dependent density power.
\item If a global generator interpolation cannot retain a uniform
      gap, localize it only after the complete connected derivative
      has been formed.
\end{enumerate}
\end{programme}

\begin{firewall}[Claim boundary after V18]
\label{fire:newV18-boundary}
V18 proves the exact noncommutative one-defect telescoping identity,
the permutation and predecessor-set coefficient censuses, and the
failure of a rowwise first-defect proof to reach the
factorial-normalized SRSPCC scale.  It replaces that overstrong
route by the complete connected Wilson--tangent defect card, proved
only as a sufficient condition.  It does not prove CWTDC, the
aggregate strong card, all-order strong MTA, WUFC, CPM, cutoff
removal, reflection positivity, Osterwalder--Schrader reconstruction,
a continuum Yang--Mills measure, or a positive Hamiltonian gap.  The
full Yang--Mills mass gap has not been proved.
\end{firewall}

\begin{theorem}[V18 result ledger]
\label{thm:newV18-results}
V18:
\begin{enumerate}[label=\textup{(V18.\arabic*)},leftmargin=6.5em]
\item gives the exact noncommutative first-defect Duhamel identity;
\item proves its normalized total coefficient variation is \(r\),
      not \(1/r!\);
\item proves that clock gain and factorial gain are independent;
\item computes the exact binomial predecessor-set weights;
\item explains why aggregate cumulant control cannot presently be
      inverted to SRSPCC; and
\item identifies a complete connected Wilson--tangent defect card
      as the next sufficient object.
\end{enumerate}
\end{theorem}

\begin{proof}
Use
\Cref{lem:newV18-telescoping,
thm:newV18-defect-mass,
prop:newV18-orthogonal-gains,
lem:newV18-predecessor-census,
prop:newV18-no-converse,
thm:newV18-CWTDC-sufficient}.
\end{proof}

\paragraph{Parent-version record.}
V18 was formed from
\texttt{Renewal\_Yang\_Mills\_Mass\_Gap\_Research\_Notes\_V17.tex}.
The V17 parent had \(7\,738\) logical source lines,
\(252\,899\) bytes, and SHA--256
\[
 \texttt{1ef1f7ad1e974762950e424d324ffbd94f35dfb90dcc1bd34b6fd8e5226bdd88}.
\]
The next compulsory companion audit is V20.

% =============================================================================
% END APPENDED FIFTH-SERIES V18 MATERIAL
% =============================================================================

% =============================================================================
% BEGIN APPENDED FIFTH-SERIES V19 MATERIAL
% =============================================================================

\chapter{V19: A score insertion for the complete connected cumulant}
\label{chap:newV19-score}

\section{Differentiating a probability law once}
\label{sec:newV19-law}

V18 showed that a Duhamel expansion of ordered contraction histories
creates \(r!\) rows before the relevant cancellation.  There is a
different exact derivative identity which acts directly on the
complete cumulant and introduces only one marked observable.

Let \((\Omega,\lambda)\) be a finite measure space, and let
\[
 d\nu_t=w_t\,d\lambda,
 \qquad t\in[0,1],
\]
be a \(C^1\) family of strictly positive probability densities.
Assume that \(w_t,\partial_tw_t\) are bounded and that all test
functions below are bounded.  Define the centered score
\begin{equation}
 S_t
 :={\partial\over\partial t}\log w_t.
 \label{eq:newV19-score}
\end{equation}
Normalization gives
\begin{equation}
 \mathbb E_tS_t=0.
 \label{eq:newV19-score-centered}
\end{equation}

\begin{lemma}[Derivative of one moment]
\label{lem:newV19-moment-derivative}
For every bounded, \(t\)-independent scalar or operator-valued
function \(F\),
\begin{equation}
 {d\over dt}\mathbb E_tF
 =
 \mathbb E_t(FS_t).
 \label{eq:newV19-moment-derivative}
\end{equation}
\end{lemma}

\begin{proof}
Differentiate under the integral:
\[
 {d\over dt}\int Fw_t\,d\lambda
 =
 \int Fw_t\,\partial_t\log w_t\,d\lambda.
\]
The hypotheses justify the operation.  Taking \(F=1\) proves
\eqref{eq:newV19-score-centered}.
\end{proof}

\section{The score--cumulant identity}
\label{sec:newV19-cumulant}

For tensor-leg ordered operator inputs, use the classical
partition cumulant
\begin{equation}
 \kappa_m^t(f_1,\ldots,f_m)
 =
 \sum_{\pi\in\Pi([m])}
 (|\pi|-1)!(-1)^{|\pi|-1}
 \bigotimes_{B\in\pi}
 \mathbb E_t\bigotimes_{i\in B}f_i.
 \label{eq:newV19-partition-cumulant}
\end{equation}
The scalar score commutes with all operator legs.

\begin{theorem}[Complete connected score insertion]
\label{thm:newV19-score-cumulant}
For every \(m\ge1\),
\begin{equation}
 \boxed{\qquad
 {d\over dt}\kappa_m^t(f_1,\ldots,f_m)
 =
 \kappa_{m+1}^t(f_1,\ldots,f_m,S_t).
 \qquad}
 \label{eq:newV19-score-cumulant}
\end{equation}
No ordered contraction history and no replica density product occurs.
\end{theorem}

\begin{proof}
Differentiate \eqref{eq:newV19-partition-cumulant}.
For a partition \(\pi\), the product rule chooses one block
\(B\in\pi\); by \Cref{lem:newV19-moment-derivative}, its moment is
replaced by the moment with \(S_t\) inserted.

Now expand the right side of
\eqref{eq:newV19-score-cumulant} by partitions of
\([m]\cup\{*\}\), where \(*\) labels \(S_t\).
Every partition in which \(\{*\}\) is a singleton vanishes by
\eqref{eq:newV19-score-centered}.  Every remaining partition is
obtained uniquely by taking a partition \(\pi\) of \([m]\),
choosing one block \(B\in\pi\), and adjoining \(*\) to \(B\).
The number of blocks remains \(|\pi|\), so its Möbius coefficient is
exactly
\((|\pi|-1)!(-1)^{|\pi|-1}\), the coefficient produced by
differentiating the \(\pi\) term.  This proves the identity in the
operator-valued tensor algebra.
\end{proof}

\begin{corollary}[Moving observables]
\label{cor:newV19-moving-observables}
If the inputs \(f_i(t)\) are also \(C^1\), then
\begin{align}
 {d\over dt}\kappa_m^t(f_1(t),\ldots,f_m(t))
 &=
 \kappa_{m+1}^t(f_1(t),\ldots,f_m(t),S_t)
 \notag\\
 &\quad+
 \sum_{i=1}^m
 \kappa_m^t
 (f_1(t),\ldots,\dot f_i(t),\ldots,f_m(t)).
 \label{eq:newV19-moving-cumulant}
\end{align}
\end{corollary}

\begin{proof}
Apply \Cref{thm:newV19-score-cumulant} to the law dependence and
multilinearity to the input dependence.
\end{proof}

\section{An intrinsic Wilson--heat path}
\label{sec:newV19-path}

To avoid a chart boundary and to keep the density assembled once,
choose a strictly positive central heat-kernel density
\(h_{\tau_\alpha}\) on \(SU(2)\), with time calibrated so that its
fixed-representation multiplier has the same first Casimir
coefficient as the Wilson multiplier.  Let
\[
 w_\alpha^W(U)
\]
be the normalized one-link Wilson density with respect to Haar
measure.  Define the geometric interpolation
\begin{align}
 d\nu_{\alpha,t}(U)
 &:=
 Z_{\alpha,t}^{-1}
 h_{\tau_\alpha}(U)^{1-t}
 w_\alpha^W(U)^t\,dU,
 \label{eq:newV19-geometric-path}\\
 L_\alpha(U)
 &:=
 \log{w_\alpha^W(U)\over h_{\tau_\alpha}(U)}.
 \label{eq:newV19-log-ratio}
\end{align}

\begin{lemma}[Exact score of the geometric path]
\label{lem:newV19-geometric-score}
For every fixed \(\alpha\), the path
\eqref{eq:newV19-geometric-path} is a smooth family of central
probability measures and
\begin{equation}
 S_{\alpha,t}
 =
 L_\alpha-\mathbb E_{\alpha,t}L_\alpha.
 \label{eq:newV19-geometric-score}
\end{equation}
Its endpoints are the calibrated heat-kernel law and the Wilson
law.
\end{lemma}

\begin{proof}
Both densities are smooth and strictly positive on the compact
group, so every geometric interpolant is positive and integrable.
Differentiate
\[
 \log w_{\alpha,t}
 =(1-t)\log h_{\tau_\alpha}
  +t\log w_\alpha^W-\log Z_{\alpha,t}.
\]
The derivative of \(\log Z_{\alpha,t}\) is
\(\mathbb E_{\alpha,t}L_\alpha\).  This gives
\eqref{eq:newV19-geometric-score}.  The endpoint statement is
immediate.
\end{proof}

\begin{theorem}[Exact Wilson--heat connected comparison]
\label{thm:newV19-WH-comparison}
For fixed bounded representation inputs,
\begin{equation}
 \boxed{\qquad
 \kappa_m^{W}(f_1,\ldots,f_m)
 -
 \kappa_m^{H}(f_1,\ldots,f_m)
 =
 \int_0^1
 \kappa_{m+1}^{\nu_{\alpha,t}}
 (f_1,\ldots,f_m,S_{\alpha,t})\,dt.
 \qquad}
 \label{eq:newV19-WH-comparison}
\end{equation}
\end{theorem}

\begin{proof}
Apply \Cref{thm:newV19-score-cumulant} to the path
\eqref{eq:newV19-geometric-path} and integrate in \(t\).
\end{proof}

\begin{assessment}[Why the V13 obstruction is absent]
\label{ass:newV19-one-density}
At each \(t\), the right side of
\eqref{eq:newV19-WH-comparison} is one cumulant under one probability
law.  There is no multiplication of \(m\) Radon--Nikodym densities
on coincident Gaussian replicas.  Thus the sharp fixed-\(L^p\)
obstruction of V13 is not triggered.
\end{assessment}

\section{Census and the new analytic card}
\label{sec:newV19-card}

Equation \eqref{eq:newV19-WH-comparison} has one integral and one
extra cumulant input, independently of \(m\).  It therefore avoids
the \(r!\,r\) first-defect census of V18.  This algebraic improvement
does not by itself estimate the score insertion.

\begin{definition}[Connected score-insertion card]
\label{def:newV19-CSIC}
The CSIC holds if, uniformly in \(t\in[0,1]\),
\begin{equation}
 \left\|
 \kappa_{m+1}^{\nu_{\alpha,t}}
 (D^{R_1},\ldots,D^{R_m},S_{\alpha,t})
 \right\|_{\rm op}
 \le
 K^m
 \left(\prod_{i=1}^mb_i\right)
 \left(\sum_{i=1}^mb_i\right)^{m-2}
 \label{eq:newV19-CSIC}
\end{equation}
for \(m\ge2\), with the exact output normalization and final payer
inserted in the same way as in the target Wilson card.
\end{definition}

\begin{theorem}[Heat strong card plus CSIC gives the Wilson strong
aggregate card]
\label{thm:newV19-CSIC-sufficient}
Assume:
\begin{enumerate}[label=\textup{(\roman*)},leftmargin=3.7em]
\item the calibrated heat-kernel cumulant has the aggregate strong
      Prüfer bound with exponential base; and
\item the CSIC \eqref{eq:newV19-CSIC} holds.
\end{enumerate}
Then the Wilson cumulant satisfies
\eqref{eq:newV17-aggregate-strong}.
\end{theorem}

\begin{proof}
Take norms in the exact comparison
\eqref{eq:newV19-WH-comparison} only after each complete connected
cumulant has been formed.  The heat term is bounded by assumption
\textup{(i)}.  The integral of the score term is bounded by
assumption \textup{(ii)}, since the parameter interval has length
one.  Absorb their two exponential bases into one.
\end{proof}

\begin{firewall}[No circular use of the all-order card]
\label{fire:newV19-CSIC-circular}
The CSIC may not be proved by applying the desired all-order
Wilson tree card at arity \(m+1\) with \(S_{\alpha,t}\) treated as
an ordinary extra input.  That would assume the conclusion at a
higher arity.  The score must instead be used through its explicit
central scalar structure and the quadratic matching of the two
endpoint laws.
\end{firewall}

\section{Fourier matching and the first score identity}
\label{sec:newV19-matching}

The calibration of \(\tau_\alpha\) matches
\[
 q_{\alpha,R}^W
 =1-s_\alpha C_2(R)+O_R(s_\alpha^2)
\]
with the heat multiplier
\[
 q_{\tau_\alpha,R}^H
 =e^{-\tau_\alpha C_2(R)}
 =1-\tau_\alpha C_2(R)+O_R(\tau_\alpha^2)
\]
by taking \(\tau_\alpha=s_\alpha+O(s_\alpha^2)\) at first order.

\begin{proposition}[Integrated fixed-mode score cancellation]
\label{prop:newV19-fixed-score}
For every fixed irreducible representation \(R\), the calibrated
path satisfies
\begin{equation}
 \int_0^1
 \kappa_2^{\nu_{\alpha,t}}
 (D^R,S_{\alpha,t})\,dt
 =
 O_R(s_\alpha^2)I_R.
 \label{eq:newV19-fixed-score-cancellation}
\end{equation}
Thus the complete score insertion has no first-order Fourier
component against any fixed representation matrix.
\end{proposition}

\begin{proof}
Apply \Cref{thm:newV19-WH-comparison} with \(m=1\).  Its left side
is
\[
 (q_{\alpha,R}^W-q_{\tau_\alpha,R}^H)I_R.
\]
The two displayed multiplier expansions and
\(\tau_\alpha=s_\alpha+O(s_\alpha^2)\) bound this difference by
\(O_R(s_\alpha^2)\).  This is exactly
\eqref{eq:newV19-fixed-score-cancellation}.
\end{proof}

\begin{firewall}[Uniform score control remains open]
\label{fire:newV19-score-uniform}
\Cref{prop:newV19-fixed-score} does not prove a pointwise or
normwise score estimate.  It is an integrated test against one fixed
Fourier mode; its constant is not uniform on growing representation
scale and it permits cancellation in \(t\).  It also gives no
uniform spectral gap for the interpolating laws.  Those are the next
analytic requirements.
\end{firewall}

\section{The interpolation-gap bottleneck}
\label{sec:newV19-gap}

The score identity is exact for any positive path, but estimates
which use a resolvent or covariance representation along the path
require a uniform Poincaré inequality.

\begin{definition}[Wilson--heat path gap]
\label{def:newV19-WHPG}
The WHPG is the existence of constants \(c,\alpha_0>0\) such that
the Markov generator symmetric in
\(L^2(\nu_{\alpha,t})\) satisfies
\begin{equation}
 \operatorname{gap}(\nu_{\alpha,t})
 \ge c\alpha
 \qquad
 (\alpha\ge\alpha_0,\ 0\le t\le1).
 \label{eq:newV19-WHPG}
\end{equation}
\end{definition}

\begin{assessment}[Exact live bottleneck]
\label{ass:newV19-bottleneck}
V19 removes the factorial census from the comparison identity.  The
remaining upstream analytic problem is now explicit:
\begin{enumerate}[leftmargin=3em]
\item prove WHPG for the geometric Wilson--heat path, or replace the
      path by one for which a uniform gap is known;
\item prove a representation-uniform centered score row with the
      extra matched clock; and
\item insert that row into a complete connected cumulant without
      appealing to the all-order card recursively.
\end{enumerate}
The first item is logically prior because every proposed
score-resolvent estimate depends on it.
\end{assessment}

\section{V20 programme and claim boundary}
\label{sec:newV19-next}

\begin{programme}[V20: audit and interpolation-gap test]
\label{prog:newV19-V20}
\begin{enumerate}[leftmargin=3em]
\item Perform the compulsory read-only SM/NS companion audit.
\item Write the radial potential of
      \(\nu_{\alpha,t}\) exactly in the class-angle coordinate.
\item Test a one-dimensional Poincaré or Hardy criterion uniformly
      in \(t\), retaining the Haar radial factor at both endpoints.
\item Separate the identity well, transition annulus, and antipodal
      cap before taking suprema.
\item If the geometric path develops a secondary bottleneck, replace
      it by an invariant-law path with an explicit monotone radial
      drift and repeat the score identity.
\item Only after WHPG is proved, derive the centered score
      resolvent row needed by CSIC.
\end{enumerate}
\end{programme}

\begin{firewall}[Claim boundary after V19]
\label{fire:newV19-boundary}
V19 proves the complete score--cumulant derivative identity and an
exact one-density Wilson--heat comparison.  It removes the ordered
history census but does not prove the Wilson--heat path gap, CSIC,
the heat strong card in the required output norm, the Wilson
aggregate strong card, all-order strong MTA, WUFC, CPM, cutoff
removal, reflection positivity, Osterwalder--Schrader reconstruction,
a continuum Yang--Mills measure, or a positive Hamiltonian gap.  The
full Yang--Mills mass gap has not been proved.
\end{firewall}

\begin{theorem}[V19 result ledger]
\label{thm:newV19-results}
V19:
\begin{enumerate}[label=\textup{(V19.\arabic*)},leftmargin=6.5em]
\item proves the complete connected score-insertion identity for
      tensor-leg ordered operator cumulants;
\item constructs an intrinsic geometric Wilson--heat path on
      \(SU(2)\);
\item proves the exact Wilson--heat cumulant comparison with one
      centered score insertion;
\item defines the noncircular CSIC and proves its sufficiency
      together with a heat strong card;
\item proves integrated fixed-Fourier-mode score cancellation; and
\item identifies the uniform interpolation gap as the first
      upstream analytic gate.
\end{enumerate}
\end{theorem}

\begin{proof}
Use
\Cref{thm:newV19-score-cumulant,
lem:newV19-geometric-score,
thm:newV19-WH-comparison,
thm:newV19-CSIC-sufficient,
prop:newV19-fixed-score}.
\end{proof}

\paragraph{Parent-version record.}
V19 was formed from
\texttt{Renewal\_Yang\_Mills\_Mass\_Gap\_Research\_Notes\_V18.tex}.
The V18 parent had \(8\,164\) logical source lines,
\(268\,176\) bytes, and SHA--256
\[
 \texttt{c1c0551d77f1556a9a36bb9f53df5a4a42d5393af98b1151ae0571015e4accc7}.
\]
The next compulsory companion audit is V20.

% =============================================================================
% END APPENDED FIFTH-SERIES V19 MATERIAL
% =============================================================================

% =============================================================================
% BEGIN APPENDED FIFTH-SERIES V20 MATERIAL
% =============================================================================

\chapter{V20: Companion audit and the Wilson--heat path gap}
\label{chap:newV20-gap}

\section{Mandatory read-only companion audit}
\label{sec:newV20-audit}

The newest active companion notes inspected at this checkpoint were
\begin{center}
\begin{tabular}{p{0.63\textwidth}r}
\hline
file & bytes\\
\hline
\texttt{Renewal\_SM\_Einstein\_Continuum\_Research\_Notes\_V23.tex}
& \(296\,732\)\\
\texttt{Renewal\_NS\_Stability\_Research\_Notes\_V31.tex}
& \(887\,537\)\\
\hline
\end{tabular}
\end{center}
Their logical source-line counts and SHA--256 digests were
\begin{align}
 7430,\qquad&
 \texttt{11030a7ecdecfbe587e20163f3781f7f70a27df1fbd9d2d650c7d60087990749},
 \label{eq:newV20-SM-audit}\\
 23052,\qquad&
 \texttt{f1121b62238ad5491bb7cf03635ea9934942edf573ed8faaa9ee79e1d38a6696}.
 \label{eq:newV20-NS-audit}
\end{align}
Neither companion file was modified.

The new SM V16 calculation keeps the exact scaled \(SU(2)\)
action--Haar factor assembled.  Its owned coefficients have a valid
Gevrey-one integrated card, while the quartic action term proves that
a factorial-free pointwise all-derivative card is false.  This
supports the present decision to differentiate the complete
probability law rather than its logarithmic action vertices.

SM V20 records the V15 Yang--Mills raw-tree counterexample and
restricts its own Gaussian compiler to square-free adjacent owners.
SM V23 proves that an exact endpoint difference and a formal Ward
identity are not interchangeable.  In the present problem, the
fixed-mode identity
\eqref{eq:newV19-fixed-score-cancellation} is likewise not a
normwise score estimate and cannot be promoted to CSIC by analogy.

The NS V31 file is unchanged.  Its exact first-moment calculation
again shows that a multiplicity-free representation can still be
analytically circular if its payment is the target continuation
stock.  Here, using the desired all-order card at arity \(m+1\) to
bound the score insertion would be the corresponding circular step.

\begin{firewall}[Audit type boundary]
\label{fire:newV20-audit}
No determinant, metric, or fluid estimate is imported into the
one-link Wilson law.  The audit changes only the order in which the
Yang--Mills objects are assembled and normed.
\end{firewall}

\section{Exact radial form of the geometric path}
\label{sec:newV20-radial}

Identify \(SU(2)\) with the unit three-sphere and write a conjugacy
class as
\[
 U=\cos\theta\,I+i\sin\theta\,\boldsymbol n\cdot\boldsymbol\sigma,
 \qquad
 0\le\theta\le\pi,\quad \boldsymbol n\in S^2.
\]
Normalize the bi-invariant metric so that \(\theta=d(I,U)\) and
\(-\Delta\) has eigenvalue \(4C_2(R)\) on representation \(R\).
The radial Haar factor is a fixed constant times
\(\sin^2\theta\,d\theta\).

Let \(h_\tau(\theta)\) be the heat-kernel density with respect to
Haar measure.  Its representation multiplier is
\[
 e^{-4\tau C_2(R)}.
\]
Since the Wilson multiplier is
\[
 1-{2\over\alpha}C_2(R)+O_R(\alpha^{-2}),
\]
the first-order calibrated heat time is
\begin{equation}
 \tau_\alpha={1\over2\alpha}+O(\alpha^{-2}).
 \label{eq:newV20-calibrated-time}
\end{equation}

\begin{lemma}[Radial density of every interpolant]
\label{lem:newV20-radial-density}
After angular integration, the geometric path
\eqref{eq:newV19-geometric-path} has probability density
\begin{equation}
 p_{\alpha,t}(\theta)
 =
 {1\over Z_{\alpha,t}^{\rm rad}}
 \sin^2\theta\,
 h_{\tau_\alpha}(\theta)^{1-t}
 e^{t\alpha\cos\theta},
 \qquad 0<\theta<\pi.
 \label{eq:newV20-radial-density}
\end{equation}
The full measure is
\[
 p_{\alpha,t}(\theta)\,d\theta\,d\omega,
\]
where \(d\omega\) is normalized uniform measure on \(S^2\).
\end{lemma}

\begin{proof}
The Wilson density relative to Haar is proportional to
\(e^{\alpha\cos\theta}\).  Insert it and the central heat density
into \eqref{eq:newV19-geometric-path}.  Haar disintegration supplies
the factor \(\sin^2\theta\); all constants independent of \(\theta\)
are absorbed into \(Z_{\alpha,t}^{\rm rad}\).
\end{proof}

\begin{corollary}[Exact radial score]
\label{cor:newV20-radial-score}
The centered score is
\begin{equation}
 S_{\alpha,t}(\theta)
 =
 \alpha\cos\theta-\log h_{\tau_\alpha}(\theta)
 -
 \mathbb E_{\alpha,t}
 \left[
 \alpha\cos\theta-\log h_{\tau_\alpha}(\theta)
 \right].
 \label{eq:newV20-radial-score}
\end{equation}
\end{corollary}

\begin{proof}
This is \eqref{eq:newV19-geometric-score} after canceling the
endpoint normalizing constants, which disappear under centering.
\end{proof}

\section{Why bounded-density comparison is exponentially weak}
\label{sec:newV20-comparison-no-go}

For a probability density \(w=e^L/Z\) relative to a reference
measure with spectral gap \(\lambda_0\), the
Holley--Stroock oscillation comparison gives only a gap of order
\[
 \lambda_0e^{-\operatorname{osc}L}.
\]
The endpoint log ratio in V19 has oscillation linear in \(\alpha\).

\begin{theorem}[Linear oscillation of the matched endpoint ratio]
\label{thm:newV20-ratio-oscillation}
With the metric and heat-time normalization above,
\begin{equation}
 \operatorname{osc}_{SU(2)}
 \log{w_\alpha^W\over h_{\tau_\alpha}}
 \ge
 \left({\pi^2\over2}-2+o(1)\right)\alpha.
 \label{eq:newV20-ratio-oscillation}
\end{equation}
In particular, a global bounded-density comparison between the heat
and Wilson endpoint measures cannot prove a gap of order \(\alpha\).
\end{theorem}

\begin{proof}
The endpoint normalizing constants cancel when two class angles are
subtracted.  For the Wilson density,
\begin{equation}
 \log{w_\alpha^W(\pi)\over w_\alpha^W(0)}
 =\alpha(\cos\pi-\cos0)=-2\alpha.
 \label{eq:newV20-Wilson-endpoint-ratio}
\end{equation}
Varadhan's compact-manifold heat-kernel asymptotic gives
\[
 \log{h_{\tau_\alpha}(\pi)\over h_{\tau_\alpha}(0)}
 =
 -{\pi^2\over4\tau_\alpha}+o(\alpha)
 =
 -{\pi^2\over2}\alpha+o(\alpha),
\]
where \eqref{eq:newV20-calibrated-time} was used.  Subtracting the
heat log ratio from
\eqref{eq:newV20-Wilson-endpoint-ratio} gives
\[
 \log{(w_\alpha^W/h_{\tau_\alpha})(\pi)
 \over
 (w_\alpha^W/h_{\tau_\alpha})(0)}
 =
 \left({\pi^2\over2}-2+o(1)\right)\alpha.
\]
The oscillation dominates the absolute value of this difference.
Since the coefficient is positive, the comparison lower bound is at
best exponentially small in \(\alpha\), even if the reference gap
is \(O(\alpha)\).
\end{proof}

\begin{firewall}[Oscillation failure is not a path-gap
counterexample]
\label{fire:newV20-oscillation-scope}
\Cref{thm:newV20-ratio-oscillation} proves that a crude global
density-ratio theorem is too weak.  It does not show that the actual
interpolating measures have a small gap.  Both endpoint densities
have one dominant well at the identity; a shape-sensitive
one-dimensional argument may still prove WHPG.
\end{firewall}

\section{Exact spherical-harmonic reduction}
\label{sec:newV20-harmonic}

Let \({\cal E}_{\alpha,t}\) be the gradient Dirichlet form of the
reversible diffusion with invariant law
\(\nu_{\alpha,t}\):
\[
 {\cal E}_{\alpha,t}(f,f)
 =
 \int_{SU(2)}|\nabla f|^2\,d\nu_{\alpha,t}.
\]
Expand
\[
 f(\theta,\omega)
 =
 \sum_{\ell=0}^\infty\sum_{k=1}^{2\ell+1}
 g_{\ell,k}(\theta)Y_{\ell,k}(\omega),
\]
where
\(-\Delta_{S^2}Y_{\ell,k}=\ell(\ell+1)Y_{\ell,k}\).

\begin{lemma}[Modewise Rayleigh quotient]
\label{lem:newV20-mode-Rayleigh}
For every smooth \(f\),
\begin{align}
 \|f\|_{L^2(\nu_{\alpha,t})}^2
 &=
 \sum_{\ell,k}
 \int_0^\pi|g_{\ell,k}(\theta)|^2
 p_{\alpha,t}(\theta)\,d\theta,
 \label{eq:newV20-mode-L2}\\
 {\cal E}_{\alpha,t}(f,f)
 &=
 \sum_{\ell,k}
 \int_0^\pi
 \left(
 |g'_{\ell,k}(\theta)|^2
 {\ell(\ell+1)\over\sin^2\theta}
 |g_{\ell,k}(\theta)|^2
 \right)
 p_{\alpha,t}(\theta)\,d\theta.
 \label{eq:newV20-mode-energy}
\end{align}
The mean-zero restriction affects only the \(\ell=0\) mode.
\end{lemma}

\begin{proof}
In geodesic polar coordinates on the three-sphere,
\[
 |\nabla f|^2
 =
 |\partial_\theta f|^2
 +{\left|\nabla_{S^2}f\right|^2\over\sin^2\theta}.
\]
Insert the orthonormal spherical-harmonic expansion and use
Parseval on \(S^2\).  All \(\ell\ge1\) harmonics have zero angular
mean.
\end{proof}

Define the radial Poincaré eigenvalue
\begin{equation}
 \lambda_{\alpha,t}^{\rm rad}
 :=
 \inf_{\substack{g\not\equiv0\\\int gp_{\alpha,t}=0}}
 {\int_0^\pi|g'|^2p_{\alpha,t}\,d\theta
  \over
  \int_0^\pi|g|^2p_{\alpha,t}\,d\theta}
 \label{eq:newV20-radial-gap}
\end{equation}
and, for \(\ell\ge1\),
\begin{equation}
 \lambda_{\alpha,t}^{(\ell)}
 :=
 \inf_{g\not\equiv0}
 {\int_0^\pi
   \left(|g'|^2+
   {\ell(\ell+1)\over\sin^2\theta}|g|^2\right)
   p_{\alpha,t}\,d\theta
  \over
  \int_0^\pi|g|^2p_{\alpha,t}\,d\theta}.
 \label{eq:newV20-angular-gap}
\end{equation}

\begin{theorem}[Exact two-scalar reduction of WHPG]
\label{thm:newV20-gap-reduction}
The spectral gap of \(\nu_{\alpha,t}\) is
\begin{equation}
 \operatorname{gap}(\nu_{\alpha,t})
 =
 \min\left\{
 \lambda_{\alpha,t}^{\rm rad},
 \inf_{\ell\ge1}\lambda_{\alpha,t}^{(\ell)}
 \right\}
 =
 \min\left\{
 \lambda_{\alpha,t}^{\rm rad},
 \lambda_{\alpha,t}^{(1)}
 \right\}.
 \label{eq:newV20-gap-reduction}
\end{equation}
Consequently WHPG is equivalent to the two uniform inequalities
\begin{equation}
 \lambda_{\alpha,t}^{\rm rad}\ge c\alpha,
 \qquad
 \lambda_{\alpha,t}^{(1)}\ge c\alpha.
 \label{eq:newV20-two-gap-tests}
\end{equation}
\end{theorem}

\begin{proof}
\Cref{lem:newV20-mode-Rayleigh} makes the variance and energy
orthogonal sums of the displayed one-dimensional mode forms.  The
bottom Rayleigh quotient is therefore the infimum of their bottoms,
with mean zero imposed on the radial constant mode.  Since
\(\ell(\ell+1)\) increases with \(\ell\), the quadratic form for
\(\ell\ge1\) is bounded below by the one for \(\ell=1\), and equality
of the infimum is attained at \(\ell=1\).  This proves
\eqref{eq:newV20-gap-reduction} and the equivalence.
\end{proof}

\section{The radial Hardy test}
\label{sec:newV20-Hardy}

Let \(m_{\alpha,t}\) be a median of \(p_{\alpha,t}\).  Define
\begin{align}
 {\cal H}_{\alpha,t}^-
 &:=
 \sup_{0<x<m_{\alpha,t}}
 \left(\int_0^xp_{\alpha,t}(\theta)\,d\theta\right)
 \left(\int_x^{m_{\alpha,t}}
 {d\theta\over p_{\alpha,t}(\theta)}\right),
 \label{eq:newV20-Hardy-minus}\\
 {\cal H}_{\alpha,t}^+
 &:=
 \sup_{m_{\alpha,t}<x<\pi}
 \left(\int_x^\pi p_{\alpha,t}(\theta)\,d\theta\right)
 \left(\int_{m_{\alpha,t}}^x
 {d\theta\over p_{\alpha,t}(\theta)}\right).
 \label{eq:newV20-Hardy-plus}
\end{align}

\begin{theorem}[Exact radial acquisition criterion]
\label{thm:newV20-Hardy-criterion}
There are universal numerical constants \(c,C>0\) such that
\begin{equation}
 c\max\{{\cal H}_{\alpha,t}^-,
        {\cal H}_{\alpha,t}^+\}
 \le
 {1\over\lambda_{\alpha,t}^{\rm rad}}
 \le
 C\max\{{\cal H}_{\alpha,t}^-,
        {\cal H}_{\alpha,t}^+\}.
 \label{eq:newV20-Hardy-criterion}
\end{equation}
Hence the radial half of WHPG is equivalent, up to universal
constants, to
\begin{equation}
 \sup_{0\le t\le1}
 \max\{{\cal H}_{\alpha,t}^-,
        {\cal H}_{\alpha,t}^+\}
 \le {C\over\alpha}.
 \label{eq:newV20-Hardy-target}
\end{equation}
\end{theorem}

\begin{proof}
This is the one-dimensional Muckenhoupt--Hardy characterization of
the Poincaré constant on an interval, applied separately on the two
sides of a median and then glued by the median decomposition of the
variance.  Its proof uses Cauchy--Schwarz on
\[
 g(x)-g(m)=\int_m^xg'(\theta)\,d\theta
\]
with weights \(p_{\alpha,t}^{1/2}\) and
\(p_{\alpha,t}^{-1/2}\), followed by Fubini; the converse tests
piecewise antiderivatives of \(p_{\alpha,t}^{-1}\) truncated at the
maximizing point.  These two steps give fixed numerical constants
independent of the density.
\end{proof}

\begin{assessment}[The gap problem is now scalar and
nonperturbative]
\label{ass:newV20-gap-position}
The geometric path cannot be controlled by global density
oscillation, but its gap question is now exactly the pair of
one-dimensional estimates
\eqref{eq:newV20-two-gap-tests}.  The radial estimate has the
explicit Hardy form
\eqref{eq:newV20-Hardy-target}; the angular estimate is the bottom
of the killed radial form
\eqref{eq:newV20-angular-gap} at \(\ell=1\).
Both retain the antipodal Haar zero and the complete heat density,
so no chart complement is discarded.
\end{assessment}

\section{V21 programme and claim boundary}
\label{sec:newV20-next}

\begin{programme}[V21: uniform radial and angular gap bounds]
\label{prog:newV20-V21}
\begin{enumerate}[leftmargin=3em]
\item Establish two-sided, \(t\)-uniform estimates for the radial
      density \eqref{eq:newV20-radial-density} on the identity well
      \(\theta\lesssim\alpha^{-1/2}\), the transition interval, and
      the antipodal cap.
\item Locate a median at
      \(m_{\alpha,t}\asymp\alpha^{-1/2}\) uniformly in \(t\).
\item Insert those estimates into
      \eqref{eq:newV20-Hardy-minus}--%
      \eqref{eq:newV20-Hardy-plus} and prove or refute the
      \(C/\alpha\) target.
\item For the \(\ell=1\) form, split at
      \(\theta=c\alpha^{-1/2}\): use
      \(2/\sin^2\theta\) inside and radial transport back to the well
      outside.
\item If the heat kernel's cut-locus prefactor prevents a uniform
      estimate, replace the heat endpoint by the explicit geodesic
      Gaussian reference and repeat the exact score identity.
\item Do not infer score control from the gap alone.  After WHPG,
      return separately to the centered score-insertion card.
\end{enumerate}
\end{programme}

\begin{firewall}[Claim boundary after V20]
\label{fire:newV20-boundary}
V20 performs the compulsory audit, writes the exact radial
Wilson--heat interpolation, proves that global density comparison is
exponentially too weak, and reduces WHPG to two explicit
one-dimensional spectral inequalities.  It does not prove those
uniform inequalities, WHPG, CSIC, the heat strong card, the Wilson
aggregate strong card, all-order strong MTA, WUFC, CPM, cutoff
removal, reflection positivity, Osterwalder--Schrader reconstruction,
a continuum Yang--Mills measure, or a positive Hamiltonian gap.  The
full Yang--Mills mass gap has not been proved.
\end{firewall}

\begin{theorem}[V20 result ledger]
\label{thm:newV20-results}
V20:
\begin{enumerate}[label=\textup{(V20.\arabic*)},leftmargin=6.5em]
\item performs the mandatory read-only SM/NS audit;
\item gives the exact radial density and score of the geometric
      Wilson--heat path;
\item proves linear-in-\(\alpha\) endpoint log-ratio oscillation and
      the failure of crude bounded-density gap comparison;
\item derives the exact spherical-harmonic Dirichlet decomposition;
\item reduces WHPG to the radial and \(\ell=1\) scalar gaps; and
\item states the explicit Muckenhoupt--Hardy quantity equivalent to
      the radial gap.
\end{enumerate}
\end{theorem}

\begin{proof}
Use
\Cref{lem:newV20-radial-density,
cor:newV20-radial-score,
thm:newV20-ratio-oscillation,
lem:newV20-mode-Rayleigh,
thm:newV20-gap-reduction,
thm:newV20-Hardy-criterion}.
\end{proof}

\paragraph{Parent-version record.}
V20 was formed from
\texttt{Renewal\_Yang\_Mills\_Mass\_Gap\_Research\_Notes\_V19.tex}.
The V19 parent had \(8\,602\) logical source lines,
\(282\,336\) bytes, and SHA--256
\[
 \texttt{936d939fdc70f9c609cc19cce4b906b1a484b2b46e3821c020d41b5b8ad3cfb2}.
\]
The next compulsory companion audit is V25.

% =============================================================================
% END APPENDED FIFTH-SERIES V20 MATERIAL
% =============================================================================

% =============================================================================
% BEGIN APPENDED FIFTH-SERIES V21 MATERIAL
% =============================================================================

\chapter{V21: A gapped geodesic reference path and its exact score}
\label{chap:newV21-geodesic}

\section{Replacing the heat endpoint}
\label{sec:newV21-reference}

V20 reduces the Wilson--heat path gap to explicit scalar estimates,
but the heat-kernel cut-locus prefactor is unnecessary for the score
identity.  We now choose a reference with the same tangent quadratic
form and an elementary radial potential.

Define the geodesic Gaussian probability law on \(SU(2)\) by
\begin{equation}
 d\gamma_\alpha^G(U)
 =
 {1\over Z_\alpha^G}
 \exp\left(-{\alpha\over2}\theta(U)^2\right)dU,
 \qquad
 \theta(U)=d(I,U)\in[0,\pi].
 \label{eq:newV21-geodesic-Gaussian}
\end{equation}
The cut locus is the single Haar-null point \(-I\), so the weak
gradient Dirichlet form is well defined.  Interpolate geometrically
between \(\gamma_\alpha^G\) and the Wilson law:
\begin{equation}
 d\widetilde\nu_{\alpha,u}(U)
 =
 {1\over\widetilde Z_{\alpha,u}}
 e^{-\alpha V_u(\theta(U))}\,dU,
 \qquad
 V_u(\theta)
 :=
 (1-u){\theta^2\over2}
 +u(1-\cos\theta),
 \label{eq:newV21-geodesic-path}
\end{equation}
where \(0\le u\le1\).

\begin{lemma}[Exact path score]
\label{lem:newV21-score}
Put
\begin{equation}
 R(\theta)
 :={\theta^2\over2}-(1-\cos\theta).
 \label{eq:newV21-score-profile}
\end{equation}
The centered score of
\eqref{eq:newV21-geodesic-path} is
\begin{equation}
 \widetilde S_{\alpha,u}
 =
 \alpha\left(
 R-\mathbb E_{\alpha,u}R
 \right).
 \label{eq:newV21-exact-score}
\end{equation}
The complete cumulants obey
\begin{equation}
 \kappa_m^W(f_1,\ldots,f_m)
 -
 \kappa_m^G(f_1,\ldots,f_m)
 =
 \int_0^1
 \kappa_{m+1}^{\widetilde\nu_{\alpha,u}}
 (f_1,\ldots,f_m,\widetilde S_{\alpha,u})\,du.
 \label{eq:newV21-WG-comparison}
\end{equation}
\end{lemma}

\begin{proof}
The log density in \eqref{eq:newV21-geodesic-path} has
\(u\)-derivative
\[
 \alpha\left({\theta^2\over2}-(1-\cos\theta)\right)
 -\partial_u\log\widetilde Z_{\alpha,u}.
\]
The normalizing derivative is \(\alpha\mathbb E_{\alpha,u}R\),
proving \eqref{eq:newV21-exact-score}.  Apply
\Cref{thm:newV19-score-cumulant} and integrate in \(u\) to obtain
\eqref{eq:newV21-WG-comparison}.
\end{proof}

\section{Uniform one-well geometry}
\label{sec:newV21-one-well}

\begin{lemma}[Uniform potential inequalities]
\label{lem:newV21-potential}
For \(0\le\theta\le\pi\) and \(0\le u\le1\),
\begin{align}
 {2\over\pi^2}\theta^2
 &\le V_u(\theta)\le{\theta^2\over2},
 \label{eq:newV21-potential-quadratic}\\
 V_u'(\theta)
 &=(1-u)\theta+u\sin\theta>0
 \qquad(0<\theta<\pi).
 \label{eq:newV21-potential-monotone}
\end{align}
On \(0\le\theta\le\pi/2\),
\begin{equation}
 V_u'(\theta)\ge {2\over\pi}\theta.
 \label{eq:newV21-inner-slope}
\end{equation}
If \(\theta=\pi-\delta\), \(0\le\delta\le\pi/2\), then
\begin{equation}
 V_u'(\pi-\delta)
 \ge c\bigl((1-u)+\delta\bigr)
 \label{eq:newV21-cap-slope}
\end{equation}
with a numerical \(c>0\).
\end{lemma}

\begin{proof}
The elementary chord inequalities
\[
 {2\theta^2\over\pi^2}
 \le1-\cos\theta\le{\theta^2\over2}
\]
prove \eqref{eq:newV21-potential-quadratic}.  Differentiation gives
\eqref{eq:newV21-potential-monotone}.  On the half interval,
\(\sin\theta\ge2\theta/\pi\), proving
\eqref{eq:newV21-inner-slope}.

For the cap, if \(u\le1/2\), then
\((1-u)(\pi-\delta)\ge\pi/4\), which dominates a fixed multiple of
\((1-u)+\delta\).  If \(u>1/2\), use
\(\sin\delta\ge2\delta/\pi\) and
\((1-u)(\pi-\delta)\ge(\pi/2)(1-u)\).  This proves
\eqref{eq:newV21-cap-slope}.
\end{proof}

After angular integration, use the unnormalized radial density
\begin{equation}
 q_{\alpha,u}(\theta)
 :=
 \sin^2\theta\,e^{-\alpha V_u(\theta)}.
 \label{eq:newV21-q}
\end{equation}

\begin{lemma}[Uniform normalization and moments]
\label{lem:newV21-moments}
For every fixed \(k\ge0\), uniformly in \(u\in[0,1]\) and
\(\alpha\ge1\),
\begin{align}
 c\alpha^{-3/2}
 &\le
 \int_0^\pi q_{\alpha,u}(\theta)\,d\theta
 \le C\alpha^{-3/2},
 \label{eq:newV21-normalization}\\
 \mathbb E_{\alpha,u}\theta^k
 &\le C_k\alpha^{-k/2}.
 \label{eq:newV21-moment}
\end{align}
\end{lemma}

\begin{proof}
For the upper bound, use
\(\sin\theta\le\theta\) and the lower potential bound:
\[
 \int_0^\pi q_{\alpha,u}
 \le
 \int_0^\infty
 \theta^2e^{-(2/\pi^2)\alpha\theta^2}\,d\theta
 \le C\alpha^{-3/2}.
\]
For \(0\le\theta\le\alpha^{-1/2}\),
\(\sin\theta\ge c\theta\) and
\(V_u(\theta)\le\theta^2/2\).  Integration on this interval gives
the lower bound.  Inserting \(\theta^k\) in the upper calculation
and dividing by the lower normalization bound proves
\eqref{eq:newV21-moment}.
\end{proof}

\section{A uniform radial Poincaré bound}
\label{sec:newV21-radial-gap}

Choose a fixed large \(R>4\) and put
\begin{equation}
 a_\alpha={R\over\sqrt\alpha}.
 \label{eq:newV21-anchor}
\end{equation}
Increase the fixed lower threshold \(\alpha_0\) so that
\(a_\alpha\le\pi/4\).

\begin{lemma}[The anchor has two-sided mass]
\label{lem:newV21-anchor-mass}
There is \(\eta_R>0\), independent of \(u\) and
\(\alpha\ge\alpha_0\), such that
\begin{equation}
 \widetilde\nu_{\alpha,u}
 \{\theta\le a_\alpha\}\ge\eta_R,
 \qquad
 \widetilde\nu_{\alpha,u}
 \{\theta\ge a_\alpha\}\ge\eta_R.
 \label{eq:newV21-anchor-mass}
\end{equation}
\end{lemma}

\begin{proof}
On every fixed scaled interval
\(\theta=v/\sqrt\alpha\), Taylor's theorem gives, uniformly in
\(u\),
\[
 \alpha V_u(v/\sqrt\alpha)
 ={v^2\over2}+O_R(\alpha^{-1}),
\qquad
 \sin(v/\sqrt\alpha)\asymp {v\over\sqrt\alpha}.
\]
The interval \(1\le v\le2\) lies below \(R\) and supplies a fixed
positive fraction of the normalization to the left.  The interval
\(R\le v\le R+1\) supplies a fixed positive fraction to the right.
Use \eqref{eq:newV21-normalization}.
\end{proof}

\begin{lemma}[Uniform descending log slope]
\label{lem:newV21-log-slope}
If \(R\) is chosen sufficiently large, there is \(c_R>0\) such that
\begin{equation}
 -{d\over d\theta}\log q_{\alpha,u}(\theta)
 =
 \alpha V_u'(\theta)-2\cot\theta
 \ge c_R\sqrt\alpha
 \label{eq:newV21-log-slope}
\end{equation}
for every
\(a_\alpha\le\theta<\pi\), every \(u\in[0,1]\), and
\(\alpha\ge\alpha_0\).
\end{lemma}

\begin{proof}
For \(a_\alpha\le\theta\le\pi/2\), use
\eqref{eq:newV21-inner-slope} and
\(\cot\theta\le1/\theta\):
\[
 \alpha V_u'(\theta)-2\cot\theta
 \ge {2\alpha\theta\over\pi}-{2\over\theta}.
\]
The right side is increasing in \(\theta\), and at
\(\theta=R/\sqrt\alpha\) it equals
\(\sqrt\alpha(2R/\pi-2/R)\), which has the required lower bound for
large fixed \(R\).

Now let \(\pi/2\le\theta<\pi\) and
\(\delta=\pi-\theta\).  If \(u\le1/2\),
\Cref{lem:newV21-potential} gives \(V_u'\ge c\), hence a lower
bound \(c\alpha\).  If \(u>1/2\), then
\(V_u'\ge c\delta\).  For \(0<\delta\le\pi/4\),
\(-\cot\theta=\cot\delta\ge c/\delta\), so
\[
 \alpha V_u'-2\cot\theta
 \ge c(\alpha\delta+\delta^{-1})
 \ge c\sqrt\alpha.
\]
For \(\pi/4\le\delta\le\pi/2\), the term
\(\alpha V_u'\) alone is at least \(c\alpha\).
\end{proof}

\begin{lemma}[Two-sided Hardy products at the anchor]
\label{lem:newV21-Hardy-products}
With unnormalized density \(q=q_{\alpha,u}\),
\begin{align}
 \sup_{0<x<a_\alpha}
 \left(\int_0^xq\right)
 \left(\int_x^{a_\alpha}{d\theta\over q(\theta)}\right)
 &\le {C_R\over\alpha},
 \label{eq:newV21-Hardy-left}\\
 \sup_{a_\alpha<x<\pi}
 \left(\int_x^\pi q\right)
 \left(\int_{a_\alpha}^x{d\theta\over q(\theta)}\right)
 &\le {C_R\over\alpha}.
 \label{eq:newV21-Hardy-right}
\end{align}
\end{lemma}

\begin{proof}
On \([0,a_\alpha]\), the potential exponent is bounded by
\(R^2/2\) and \(\sin\theta\asymp\theta\).  Hence, with constants
depending only on \(R\),
\[
 c_R\theta^2\le q(\theta)\le C_R\theta^2.
\]
Therefore the left product is at most
\[
 C_Rx^3\int_x^{a_\alpha}\theta^{-2}\,d\theta
 \le C_Rx^2\le {C_R\over\alpha}.
\]

Put \(k=c_R\sqrt\alpha\).  By
\Cref{lem:newV21-log-slope}, for \(y\ge x\ge a_\alpha\),
\[
 q(y)\le q(x)e^{-k(y-x)},
\]
and for \(a_\alpha\le y\le x\),
\[
 {1\over q(y)}
 \le {1\over q(x)}e^{-k(x-y)}.
\]
Consequently
\[
 \int_x^\pi q\le {q(x)\over k},
 \qquad
 \int_{a_\alpha}^x{d\theta\over q(\theta)}
 \le {1\over kq(x)}.
\]
Their product is at most \(k^{-2}\le C_R/\alpha\).
\end{proof}

\begin{theorem}[Uniform radial gap on the geodesic path]
\label{thm:newV21-radial-gap}
There is \(c>0\) such that
\begin{equation}
 \lambda_{\alpha,u}^{\rm rad}\ge c\alpha
 \qquad
 (0\le u\le1,\ \alpha\ge\alpha_0).
 \label{eq:newV21-radial-gap}
\end{equation}
\end{theorem}

\begin{proof}
The two-sided anchor mass
\eqref{eq:newV21-anchor-mass} permits the
one-dimensional Poincaré inequality to be glued at \(a_\alpha\)
instead of an exact median, with a constant depending only on
\(\eta_R\).  The Hardy products are unchanged by normalization of
\(q\), because its normalizing constant cancels between the two
factors.  Apply
\Cref{lem:newV21-Hardy-products} and the same
Muckenhoupt--Hardy proof as in
\Cref{thm:newV20-Hardy-criterion}.
\end{proof}

\section{The angular gap}
\label{sec:newV21-angular-gap}

\begin{lemma}[Dirichlet Hardy inequality on the descending side]
\label{lem:newV21-Dirichlet-Hardy}
If \(h(a_\alpha)=0\), then
\begin{equation}
 \int_{a_\alpha}^{\pi}|h|^2q_{\alpha,u}\,d\theta
 \le {C\over\alpha}
 \int_{a_\alpha}^{\pi}|h'|^2q_{\alpha,u}\,d\theta.
 \label{eq:newV21-Dirichlet-Hardy}
\end{equation}
\end{lemma}

\begin{proof}
The weighted Hardy constant for a Dirichlet boundary at
\(a_\alpha\) is bounded by
\[
 4\sup_{a_\alpha<x<\pi}
 \left(\int_x^\pi q\right)
 \left(\int_{a_\alpha}^xq^{-1}\right).
\]
Use \eqref{eq:newV21-Hardy-right}.
\end{proof}

\begin{lemma}[The anchor trace is paid by local angular energy]
\label{lem:newV21-trace}
For every absolutely continuous \(g\),
\begin{align}
 |g(a_\alpha)|^2
 \int_{a_\alpha}^{\pi}q_{\alpha,u}\,d\theta
 \le {C_R\over\alpha}
 \int_{a_\alpha/2}^{a_\alpha}
 \left(
 |g'|^2+{|g|^2\over\sin^2\theta}
 \right)q_{\alpha,u}\,d\theta.
 \label{eq:newV21-trace}
\end{align}
\end{lemma}

\begin{proof}
Average the elementary trace bound
\[
 |g(a_\alpha)|^2
 \le2|g(y)|^2
 +2(a_\alpha-y)
 \int_y^{a_\alpha}|g'(\theta)|^2\,d\theta
\]
over \(y\in[a_\alpha/2,a_\alpha]\).  On this interval,
\[
 q_{\alpha,u}(\theta)\asymp_R a_\alpha^2,
 \qquad
 \sin^2\theta\asymp_R a_\alpha^2.
\]
Also, \Cref{lem:newV21-log-slope} gives
\[
 \int_{a_\alpha}^{\pi}q_{\alpha,u}
 \le C_R{q_{\alpha,u}(a_\alpha)\over\sqrt\alpha}
 \le C_R{a_\alpha^2\over\sqrt\alpha}.
\]
After multiplication, the averaged value term and derivative term
are respectively bounded by
\[
 {C_R\over\alpha}
 \int_{a_\alpha/2}^{a_\alpha}
 {|g|^2\over\sin^2\theta}q_{\alpha,u},
 \qquad
 {C_R\over\alpha}
 \int_{a_\alpha/2}^{a_\alpha}
 |g'|^2q_{\alpha,u}.
\]
\end{proof}

\begin{theorem}[Uniform angular gap and WHPG for the geodesic path]
\label{thm:newV21-full-gap}
For the \(\ell=1\) radial form,
\begin{equation}
 \lambda_{\alpha,u}^{(1)}\ge c\alpha.
 \label{eq:newV21-angular-gap}
\end{equation}
Consequently the full reversible Dirichlet form of
\(\widetilde\nu_{\alpha,u}\) has spectral gap
\begin{equation}
 \boxed{\qquad
 \operatorname{gap}(\widetilde\nu_{\alpha,u})
 \ge c\alpha
 \qquad
 (0\le u\le1,\ \alpha\ge\alpha_0).
 \qquad}
 \label{eq:newV21-full-gap}
\end{equation}
\end{theorem}

\begin{proof}
On \(0\le\theta\le a_\alpha\),
\[
 {2\over\sin^2\theta}
 \ge {2\over a_\alpha^2}
 ={2\alpha\over R^2},
\]
so the angular potential pays the \(L^2\) mass there at the required
scale.

On \([a_\alpha,\pi]\), write
\[
 g=(g-g(a_\alpha))+g(a_\alpha).
\]
\Cref{lem:newV21-Dirichlet-Hardy} pays the first term by
\(\alpha^{-1}\int|g'|^2q\).
\Cref{lem:newV21-trace} pays the constant term by
\(\alpha^{-1}\) times the derivative and angular energy on
\([a_\alpha/2,a_\alpha]\).  Combining the inside and outside
estimates gives
\[
 \int_0^\pi|g|^2q
 \le {C\over\alpha}
 \int_0^\pi
 \left(
 |g'|^2+{2\over\sin^2\theta}|g|^2
 \right)q.
\]
This is \eqref{eq:newV21-angular-gap}.  Higher angular modes have a
larger potential.  Combine with
\Cref{thm:newV21-radial-gap} and the exact mode reduction
\eqref{eq:newV20-gap-reduction}.
\end{proof}

\section{Uniform score rows}
\label{sec:newV21-score-rows}

\begin{lemma}[Quartic score profile]
\label{lem:newV21-quartic-profile}
For \(0\le\theta\le\pi\),
\begin{equation}
 0\le R(\theta)\le{\theta^4\over24},
 \qquad
 0\le R'(\theta)=\theta-\sin\theta\le{\theta^3\over6}.
 \label{eq:newV21-quartic-profile}
\end{equation}
\end{lemma}

\begin{proof}
The alternating Taylor bounds
\[
 1-\cos\theta\ge{\theta^2\over2}-{\theta^4\over24},
 \qquad
 \sin\theta\ge\theta-{\theta^3\over6}
\]
hold on \([0,\pi]\).  The opposite inequalities
\(1-\cos\theta\le\theta^2/2\) and
\(\sin\theta\le\theta\) give nonnegativity.
\end{proof}

\begin{theorem}[Value and gradient score bounds]
\label{thm:newV21-score-rows}
For every fixed \(1\le p<\infty\),
\begin{align}
 \|\widetilde S_{\alpha,u}\|_{L^p}
 &\le {C_p\over\alpha},
 \label{eq:newV21-score-value}\\
 \|\nabla\widetilde S_{\alpha,u}\|_{L^p}
 &\le {C_p\over\sqrt\alpha},
 \label{eq:newV21-score-gradient}
\end{align}
uniformly in \(u\in[0,1]\) and \(\alpha\ge\alpha_0\).
\end{theorem}

\begin{proof}
Centering costs at most two in \(L^p\).  By
\Cref{lem:newV21-quartic-profile,lem:newV21-moments},
\[
 \|\widetilde S_{\alpha,u}\|_p
 \le2\alpha\|R\|_p
 \le C_p\alpha
 \left(\mathbb E\theta^{4p}\right)^{1/p}
 \le {C_p\over\alpha}.
\]
The weak gradient has magnitude
\(\alpha|R'(\theta)|\) away from the Haar-null cut point.  Hence
\[
 \|\nabla\widetilde S_{\alpha,u}\|_p
 \le C_p\alpha
 \left(\mathbb E\theta^{3p}\right)^{1/p}
 \le {C_p\over\sqrt\alpha}.
\]
\end{proof}

\begin{assessment}[What has been acquired]
\label{ass:newV21-acquired}
Replacing the heat endpoint by the geodesic Gaussian closes the
first upstream gate from V19: the whole interpolation path has a
uniform gap of order \(\alpha\).  Its exact centered score starts at
quartic order in the identity well and has the uniform fixed-\(p\)
rows \eqref{eq:newV21-score-value}--%
\eqref{eq:newV21-score-gradient}.  The remaining issue is
all-arity connected insertion.  Fixed-\(p\) bounds alone cannot be
multiplied over arbitrary coincident replicas, and CSIC may not be
obtained by invoking the desired arity-\(m+1\) card.
\end{assessment}

\section{V22 programme and claim boundary}
\label{sec:newV21-next}

\begin{programme}[V22: geodesic reference card and score insertion]
\label{prog:newV21-V22}
\begin{enumerate}[leftmargin=3em]
\item Express \(\gamma_\alpha^G\) in the scaled principal
      logarithmic chart.  Its density relative to the tangent
      Gaussian is only the exact Haar Jacobian times the remote
      radius cutoff.
\item Use \(0\le J_\alpha\le1\) and the V12 integrated exact-Haar
      derivative card to test a strong geodesic-reference cumulant
      bound without the Wilson action growth.
\item Keep the radius-cutoff boundary as one exponentially remote
      remainder; do not differentiate its indicator.
\item Differentiate the complete connected cumulant along
      \eqref{eq:newV21-geodesic-path} and insert the single explicit
      score \eqref{eq:newV21-exact-score}.
\item Use the quartic score profile and the path gap to seek a
      leaf-marked connected estimate, without applying the
      all-order card at arity \(m+1\).
\item If fixed-\(p\) score Hölder constants grow factorially, return
      to the exact radial series for \(R\) before taking norms.
\end{enumerate}
\end{programme}

\begin{firewall}[Claim boundary after V21]
\label{fire:newV21-boundary}
V21 constructs an explicit one-density geodesic
Gaussian--Wilson path, proves its uniform radial, angular, and full
spectral gaps of order \(\alpha\), and proves uniform fixed-\(p\)
value and gradient rows for its exact quartic centered score.  It
does not prove the geodesic-reference strong card, CSIC, the Wilson
aggregate strong card, all-order strong MTA, WUFC, CPM, cutoff
removal, reflection positivity, Osterwalder--Schrader reconstruction,
a continuum Yang--Mills measure, or a positive Hamiltonian gap.  The
full Yang--Mills mass gap has not been proved.
\end{firewall}

\begin{theorem}[V21 result ledger]
\label{thm:newV21-results}
V21:
\begin{enumerate}[label=\textup{(V21.\arabic*)},leftmargin=6.5em]
\item replaces the heat endpoint by an explicit geodesic Gaussian
      reference and writes the exact quartic score;
\item proves uniform quadratic confinement and all radial moments;
\item proves the radial path gap by two-sided Hardy products;
\item proves the \(\ell=1\) angular gap by a descending-side
      Dirichlet Hardy estimate and a local trace payment;
\item obtains the full path gap \(c\alpha\); and
\item proves the fixed-\(p\) score value and gradient rows
      \(O(\alpha^{-1})\) and \(O(\alpha^{-1/2})\).
\end{enumerate}
\end{theorem}

\begin{proof}
Use
\Cref{lem:newV21-score,
lem:newV21-potential,
lem:newV21-moments,
thm:newV21-radial-gap,
thm:newV21-full-gap,
thm:newV21-score-rows}.
\end{proof}

\paragraph{Parent-version record.}
V21 was formed from
\texttt{Renewal\_Yang\_Mills\_Mass\_Gap\_Research\_Notes\_V20.tex}.
The V20 parent had \(9\,087\) logical source lines,
\(297\,941\) bytes, and SHA--256
\[
 \texttt{6bb15b15f550c0788a91af137531fd77ffb335e1f9c8741998115e17bd04dd7d}.
\]
The next compulsory companion audit is V25.

% =============================================================================
% END APPENDED FIFTH-SERIES V21 MATERIAL
% =============================================================================

% =============================================================================
% BEGIN APPENDED FIFTH-SERIES V22 MATERIAL
% =============================================================================

\chapter{V22: Adaptive score rows and the geodesic chart boundary}
\label{chap:newV22-reference}

\section{All-\(p\) concentration on the gapped path}
\label{sec:newV22-all-p}

V21 proved fixed-\(p\) score rows.  A connected cumulant of arity
\(m+1\) may require a Hölder exponent growing with \(m\), so the
dependence on \(p\) must be explicit.

\begin{lemma}[Sharp-order radial moments]
\label{lem:newV22-radial-moments}
There is a numerical \(C\) such that for every \(q\ge1\),
every \(u\in[0,1]\), and every \(\alpha\ge\alpha_0\),
\begin{equation}
 \left(
 \mathbb E_{\alpha,u}\theta^q
 \right)^{1/q}
 \le C\sqrt{q\over\alpha}.
 \label{eq:newV22-radial-moments}
\end{equation}
\end{lemma}

\begin{proof}
Use
\(\sin\theta\le\theta\),
\(V_u(\theta)\ge(2/\pi^2)\theta^2\), and the lower normalization
bound \eqref{eq:newV21-normalization}:
\begin{align*}
 \mathbb E_{\alpha,u}\theta^q
 &\le
 C\alpha^{3/2}
 \int_0^\infty
 \theta^{q+2}e^{-c\alpha\theta^2}\,d\theta\\
 &=
 Cc^{-(q+3)/2}\alpha^{-q/2}
 \Gamma\left({q+3\over2}\right).
\end{align*}
The elementary Gamma bound
\(\Gamma((q+3)/2)^{1/q}\le C\sqrt q\)
proves the claim after enlarging \(C\) for bounded \(q\).
\end{proof}

\begin{theorem}[Adaptive score rows]
\label{thm:newV22-adaptive-score}
For every \(p\ge1\),
\begin{align}
 \|\widetilde S_{\alpha,u}\|_{L^p}
 &\le {Cp^2\over\alpha},
 \label{eq:newV22-score-value}\\
 \|\nabla\widetilde S_{\alpha,u}\|_{L^p}
 &\le {Cp^{3/2}\over\sqrt\alpha},
 \label{eq:newV22-score-gradient}
\end{align}
uniformly in \(u\in[0,1]\) and
\(\alpha\ge\alpha_0\).
\end{theorem}

\begin{proof}
Use the exact centered score
\eqref{eq:newV21-exact-score} and the quartic profile
\eqref{eq:newV21-quartic-profile}.  Centering costs at most two:
\[
 \|\widetilde S_{\alpha,u}\|_p
 \le {C\alpha}\|\theta^4\|_p
 \le
 C\alpha
 \left(C\sqrt{4p/\alpha}\right)^4
 \le {Cp^2\over\alpha}.
\]
Similarly,
\[
 \|\nabla\widetilde S_{\alpha,u}\|_p
 \le C\alpha\|\theta^3\|_p
 \le {Cp^{3/2}\over\sqrt\alpha}.
\]
\end{proof}

\begin{corollary}[Adaptive Hölder costs only an exponential base]
\label{cor:newV22-adaptive-exponential}
If the score is placed in a Hölder estimate with any exponent
\(p\le c(m+1)\), then its polynomial factor \(p^2\) or
\(p^{3/2}\) can be absorbed into \(K^m\), uniformly in \(m\).
Thus the score row itself introduces no factorial.
\end{corollary}

\begin{proof}
For \(m\ge1\), every fixed power of \(m+1\) is bounded by
\(C^m\) after increasing a numerical \(C\).  Apply
\Cref{thm:newV22-adaptive-score}.
\end{proof}

\begin{firewall}[An adaptive row is not CSIC]
\label{fire:newV22-row-not-CSIC}
The score occurs once, so
\Cref{cor:newV22-adaptive-exponential} removes one possible
high-arity loss.  It does not organize the remaining \(m\)
representation inputs into a connected Prüfer weight and therefore
does not prove CSIC.
\end{firewall}

\section{Exact scaled form of the geodesic reference}
\label{sec:newV22-scaled-reference}

On the principal exponential chart write
\[
 U=\exp(X),
 \qquad
 X=\theta\boldsymbol n,
 \qquad
 |X|=\theta<\pi,
\]
and scale \(Y=\sqrt\alpha X\).  The \(SU(2)\) Haar Jacobian is
\begin{equation}
 J_\alpha(Y)
 =
 \left(
 {\sin(|Y|/\sqrt\alpha)\over |Y|/\sqrt\alpha}
 \right)^2.
 \label{eq:newV22-Haar-Jacobian}
\end{equation}
Let \(\gamma_3\) denote standard Gaussian measure on
\(\mathbb R^3\), and put
\begin{equation}
 z_\alpha
 :=
 \mathbb E_{\gamma_3}
 \left[
 J_\alpha(Y)
 \mathbf1_{\{|Y|<\pi\sqrt\alpha\}}
 \right].
 \label{eq:newV22-z}
\end{equation}

\begin{lemma}[Exact Gaussian density and uniform \(L^\infty\) bound]
\label{lem:newV22-reference-density}
The scaled geodesic Gaussian law has density
\begin{equation}
 h_\alpha^G(Y)
 =
 z_\alpha^{-1}
 J_\alpha(Y)
 \mathbf1_{\{|Y|<\pi\sqrt\alpha\}}
 \label{eq:newV22-reference-density}
\end{equation}
relative to \(\gamma_3\).  Moreover, for
\(\alpha\ge\alpha_0\),
\begin{equation}
 0\le h_\alpha^G\le C,
 \qquad
 c\le z_\alpha\le1.
 \label{eq:newV22-reference-Linfty}
\end{equation}
\end{lemma}

\begin{proof}
The exponential-coordinate Haar formula is
\[
 dU=CJ_\alpha(Y)\alpha^{-3/2}\,dY.
\]
Multiplication by the geodesic factor
\(e^{-|Y|^2/2}\) and normalization gives
\eqref{eq:newV22-reference-density}.  Since
\(|\sin r|\le|r|\), \(0\le J_\alpha\le1\), so
\(z_\alpha\le1\).  On \(|Y|\le1\), for large \(\alpha\),
\(J_\alpha(Y)\ge1/2\) and the chart indicator is one.  Hence
\[
 z_\alpha\ge {1\over2}\gamma_3\{|Y|\le1\}=:c>0.
\]
The density upper bound follows.
\end{proof}

\begin{corollary}[Coincident replicas remain bounded]
\label{cor:newV22-coincident-reference}
For every integer \(r\ge1\),
\begin{equation}
 \mathbb E_{\gamma_3}
 \left(h_\alpha^G(Y)\right)^r
 \le C^{r-1}.
 \label{eq:newV22-coincident-reference}
\end{equation}
Thus the finite uniform \(L^p\) obstruction of the exact Wilson
density in V13 is absent for the geodesic reference.
\end{corollary}

\begin{proof}
Since \(h_\alpha^G\le C\) and
\(\mathbb E_{\gamma_3}h_\alpha^G=1\),
\[
 \mathbb E(h_\alpha^G)^r
 \le C^{r-1}\mathbb Eh_\alpha^G.
\]
\end{proof}

\section{The boundary regularity obstruction}
\label{sec:newV22-boundary}

The bounded density solves the coincident-replica problem, but a
Gaussian covariance interpolation differentiates each vertex
according to its tree degree.  The chart indicator cannot be
ignored at unbounded degree.

Restrict to a radial line and put
\[
 R_\alpha=\pi\sqrt\alpha,
 \qquad
 j_\alpha(r)
 =
 \left(
 {\sin(r/\sqrt\alpha)\over r/\sqrt\alpha}
 \right)^2,
 \qquad
 k_\alpha(r)=j_\alpha(r)\mathbf1_{\{r<R_\alpha\}}.
\]

\begin{lemma}[Exact boundary jet]
\label{lem:newV22-boundary-jet}
As \(r\uparrow R_\alpha\),
\begin{equation}
 j_\alpha(r)
 =
 {(R_\alpha-r)^2\over\pi^2\alpha}
 +O\left(
 {(R_\alpha-r)^3\over\alpha^{3/2}}
 \right).
 \label{eq:newV22-boundary-jet}
\end{equation}
Consequently the extension \(k_\alpha\) is \(C^1\) but not \(C^2\)
at \(R_\alpha\).  Its third distributional derivative contains a
nonzero multiple of the surface delta at the chart boundary.
\end{lemma}

\begin{proof}
Write \(r=R_\alpha-\delta\).  Then
\[
 \sin(r/\sqrt\alpha)
 =\sin(\pi-\delta/\sqrt\alpha)
 ={\delta\over\sqrt\alpha}
 +O(\delta^3\alpha^{-3/2}),
\]
while
\[
 {r\over\sqrt\alpha}
 =\pi-{\delta\over\sqrt\alpha}.
\]
Squaring the quotient proves
\eqref{eq:newV22-boundary-jet}.  The value and first derivative
match the zero exterior extension; the second interior derivative
has nonzero boundary limit while the exterior second derivative is
zero.  Differentiating this jump once in distributions produces a
nonzero delta term.
\end{proof}

\begin{theorem}[Direct smooth Gaussian BKAR is unavailable at all
degrees]
\label{thm:newV22-direct-BKAR-unavailable}
The usual smooth Gaussian covariance-derivative formula cannot be
applied directly to
\(D^R(\exp(Y/\sqrt\alpha))h_\alpha^G(Y)\)
at vertices of arbitrary tree degree.  Starting at degree three,
weak derivatives include chart-boundary distributions which are not
covered by the V12 interior exact-Haar derivative card.
\end{theorem}

\begin{proof}
A vertex of degree \(d\) in the Gaussian BKAR formula receives
\(d\) derivatives of its complete vertex function.  The
representation factor is smooth, but
\Cref{lem:newV22-boundary-jet} shows that the density factor has a
surface distribution in its third weak derivative.  The V12 card
estimates ordinary derivatives of the entire function
\(J_\alpha\) under a Gaussian expectation; it contains no surface
distribution.  Thus its hypotheses do not justify the degree-three
or higher covariance derivative of the truncated density.
\end{proof}

\begin{firewall}[A remote boundary is not automatically harmless]
\label{fire:newV22-remote-boundary}
The Gaussian surface weight at
\(|Y|=\pi\sqrt\alpha\) is exponentially small in \(\alpha\), but
arbitrary arity permits arbitrary normal derivative order.  No
uniform exponential-in-degree surface-jet estimate has been proved.
Discarding the boundary solely because it is remote would repeat the
fixed-arity-to-all-arity error corrected in V12.
\end{firewall}

\section{An intrinsic reference alternative}
\label{sec:newV22-intrinsic}

The geodesic reference remains a compact-group probability measure
with the path gap proved in V21.  Its cumulants can therefore be
kept intrinsic, avoiding differentiation of the chart indicator.
The same score identity then compares two intrinsic complete
cumulants.

\begin{definition}[Geodesic intrinsic aggregate card]
\label{def:newV22-GIAC}
The GIAC is
\begin{equation}
 \left\|
 \kappa_m^{\gamma_\alpha^G}
 (D^{R_1},\ldots,D^{R_m})
 \right\|_{\rm op}
 \le
 K^{m-1}
 \left(\prod_i b_i\right)
 \left(\sum_i b_i\right)^{m-2}.
 \label{eq:newV22-GIAC}
\end{equation}
\end{definition}

\begin{definition}[Geodesic-path connected score card]
\label{def:newV22-GCSIC}
The GCSIC is the analogous bound
\begin{equation}
 \left\|
 \kappa_{m+1}^{\widetilde\nu_{\alpha,u}}
 (D^{R_1},\ldots,D^{R_m},\widetilde S_{\alpha,u})
 \right\|_{\rm op}
 \le
 K^m
 \left(\prod_i b_i\right)
 \left(\sum_i b_i\right)^{m-2},
 \label{eq:newV22-GCSIC}
\end{equation}
uniformly in \(u\).
\end{definition}

\begin{theorem}[GIAC and GCSIC close the Wilson all-thermal
aggregate card]
\label{thm:newV22-two-cards}
If \eqref{eq:newV22-GIAC} and
\eqref{eq:newV22-GCSIC} hold with the exact output payer, then the
Wilson aggregate strong card
\eqref{eq:newV17-aggregate-strong} holds.
\end{theorem}

\begin{proof}
Use the exact one-density comparison
\eqref{eq:newV21-WG-comparison}.  Apply GIAC to its geodesic
endpoint and GCSIC inside the parameter integral.  Take the triangle
inequality only between these two complete connected objects and
enlarge the exponential base.
\end{proof}

\begin{assessment}[Current position]
\label{ass:newV22-position}
The geodesic path now has three rigorous advantages:
\begin{enumerate}[leftmargin=3em]
\item a uniform gap \(c\alpha\);
\item an explicit quartic score with adaptive polynomial-in-\(p\)
      rows; and
\item a uniformly bounded Gaussian density, even at coincident
      replica arity.
\end{enumerate}
The direct Gaussian forest proof nevertheless encounters the hard
chart boundary at degree three.  The live choice is to prove an
exponential surface-jet resummation, or to remain intrinsic and
prove GIAC and GCSIC by a compact-group generator method.
\end{assessment}

\section{V23 programme and claim boundary}
\label{sec:newV22-next}

\begin{programme}[V23: boundary resummation test]
\label{prog:newV22-V23}
\begin{enumerate}[leftmargin=3em]
\item Compute the complete distributional normal-jet series of
      \(J_\alpha\mathbf1_{\{|Y|<\pi\sqrt\alpha\}}\) under Gaussian
      covariance differentiation.
\item Sum the surface terms before taking norms and test their
      dependence on the vertex degree.
\item Use the fact that every representation matrix is constant on
      the antipodal boundary, up to its central parity, before
      differentiating tangentially.
\item Determine whether centering kills the first surviving surface
      row in the complete cumulant.
\item If the surface-jet generating function has superexponential
      degree growth, abandon the chart forest and use the intrinsic
      gapped generator path.
\item Retain the adaptive score rows
      \eqref{eq:newV22-score-value}--%
      \eqref{eq:newV22-score-gradient}; they are valid in either
      route.
\end{enumerate}
\end{programme}

\begin{firewall}[Claim boundary after V22]
\label{fire:newV22-boundary}
V22 proves adaptive all-\(p\) score bounds, the exact uniformly
bounded Gaussian density of the geodesic reference, and the first
chart-boundary distributional obstruction to a direct all-degree
Gaussian BKAR proof.  It does not prove a surface-jet resummation,
GIAC, GCSIC, the Wilson aggregate strong card, all-order strong MTA,
WUFC, CPM, cutoff removal, reflection positivity,
Osterwalder--Schrader reconstruction, a continuum Yang--Mills
measure, or a positive Hamiltonian gap.  The full Yang--Mills mass
gap has not been proved.
\end{firewall}

\begin{theorem}[V22 result ledger]
\label{thm:newV22-results}
V22:
\begin{enumerate}[label=\textup{(V22.\arabic*)},leftmargin=6.5em]
\item proves the sharp-order all-\(p\) radial moment bound;
\item proves adaptive score value and gradient rows with only
      polynomial \(p\)-growth;
\item writes the exact scaled geodesic-reference density and proves
      its uniform \(L^\infty\) bound;
\item proves coincident density replicas cost only an exponential;
\item computes the boundary jet and identifies the first surface
      distribution at derivative degree three; and
\item reduces Wilson closure to GIAC plus GCSIC, or a valid boundary
      resummation proving those cards.
\end{enumerate}
\end{theorem}

\begin{proof}
Use
\Cref{lem:newV22-radial-moments,
thm:newV22-adaptive-score,
lem:newV22-reference-density,
cor:newV22-coincident-reference,
lem:newV22-boundary-jet,
thm:newV22-direct-BKAR-unavailable,
thm:newV22-two-cards}.
\end{proof}

\paragraph{Parent-version record.}
V22 was formed from
\texttt{Renewal\_Yang\_Mills\_Mass\_Gap\_Research\_Notes\_V21.tex}.
The V21 parent had \(9\,700\) logical source lines,
\(315\,512\) bytes, and SHA--256
\[
 \texttt{ac983820feabee4e2dafa1dbd4a45833d57fbb643f097045df657d7b958c9fbe}.
\]
The next compulsory companion audit is V25.

% =============================================================================
% END APPENDED FIFTH-SERIES V22 MATERIAL
% =============================================================================

% =============================================================================
% BEGIN APPENDED FIFTH-SERIES V23 MATERIAL
% =============================================================================

\chapter{V23: Complete boundary jets and abandonment of the chart forest}
\label{chap:newV23-boundary}

\section{Distributional derivative formula}
\label{sec:newV23-distribution}

V22 found the first surface delta at derivative degree three.  We
now compute every boundary row and test the complete surface sum
against the translated Gaussian kernels which arise after
conditioning correlated replicas.

Let \(R>0\), let \(j\) be smooth on a neighborhood of
\((-\infty,R]\), and put
\[
 k(r)=j(r)\mathbf1_{\{r<R\}}.
\]
Write \(\delta_R^{(a)}\) for the \(a\)-th distributional derivative
of the point mass at \(R\).

\begin{lemma}[Complete one-sided boundary derivative]
\label{lem:newV23-distribution}
For every integer \(n\ge1\),
\begin{equation}
 D^nk
 =
 j^{(n)}\mathbf1_{\{r<R\}}
 -
 \sum_{a=0}^{n-1}
 j^{(n-1-a)}(R^-)\delta_R^{(a)}.
 \label{eq:newV23-distribution}
\end{equation}
If \(j(R^-)=j'(R^-)=0\), this becomes
\begin{equation}
 D^nk
 =
 j^{(n)}\mathbf1_{\{r<R\}}
 -
 \sum_{q=2}^{n-1}
 j^{(q)}(R^-)\delta_R^{(n-1-q)}.
 \label{eq:newV23-distribution-reduced}
\end{equation}
\end{lemma}

\begin{proof}
Since
\(D\mathbf1_{\{r<R\}}=-\delta_R\), the \(n=1\) formula is the
distributional product rule.  Differentiating the displayed formula
once changes the interior term to
\[
 j^{(n+1)}\mathbf1_{\{r<R\}}-j^{(n)}(R^-)\delta_R
\]
and raises every existing delta derivative by one.  This is the
formula at \(n+1\).  Induction proves
\eqref{eq:newV23-distribution}; the two vanishing boundary values
give \eqref{eq:newV23-distribution-reduced}.
\end{proof}

\section{Every ordinary Haar boundary jet is exponential}
\label{sec:newV23-Haar-jets}

Return to
\[
 R_\alpha=\pi\sqrt\alpha,
 \qquad
 j_\alpha(r)
 =
 \left(
 {\sin(r/\sqrt\alpha)\over r/\sqrt\alpha}
 \right)^2.
\]
Let
\[
 H(z)=\left({\sin z\over z}\right)^2.
\]
Then \(j_\alpha(r)=H(r/\sqrt\alpha)\).

\begin{lemma}[Fourier representation of the squared sinc]
\label{lem:newV23-sinc-Fourier}
For every complex \(z\),
\begin{equation}
 H(z)
 =
 \int_{-1}^{1}(1-|t|)e^{2itz}\,dt.
 \label{eq:newV23-sinc-Fourier}
\end{equation}
Consequently,
\begin{equation}
 |H^{(q)}(x)|\le2^q
 \qquad(x\in\mathbb R,\ q\ge0).
 \label{eq:newV23-H-derivative}
\end{equation}
\end{lemma}

\begin{proof}
Direct integration of the even function on the right gives
\[
 2\int_0^1(1-t)\cos(2zt)\,dt
 ={1-\cos(2z)\over2z^2}
 =\left({\sin z\over z}\right)^2.
\]
Both sides are entire, so the identity holds for complex \(z\).
Differentiate under the finite integral:
\[
 |H^{(q)}(x)|
 \le
 \int_{-1}^1(1-|t|)(2|t|)^q\,dt
 \le2^q.
\]
\end{proof}

\begin{corollary}[Exponential ordinary boundary jets]
\label{cor:newV23-boundary-jets}
For every \(q\ge0\),
\begin{equation}
 |j_\alpha^{(q)}(R_\alpha^-)|
 \le
 \left({2\over\sqrt\alpha}\right)^q.
 \label{eq:newV23-boundary-jets}
\end{equation}
Moreover,
\begin{equation}
 j_\alpha''(R_\alpha^-)
 ={2\over\pi^2\alpha}.
 \label{eq:newV23-second-jet}
\end{equation}
\end{corollary}

\begin{proof}
The chain rule and
\eqref{eq:newV23-H-derivative} give the first assertion.  The second
is the exact quadratic coefficient computed in
\eqref{eq:newV22-boundary-jet}.
\end{proof}

\begin{assessment}[The ordinary jet is not the problem]
\label{ass:newV23-ordinary-jet}
The one-sided derivatives of the Haar Jacobian have precisely the
desired exponential dependence on derivative degree.  The possible
failure lies in the order of the delta derivative left after
\eqref{eq:newV23-distribution-reduced} is paired with a conditional
Gaussian density.
\end{assessment}

\section{The complete surface sum against a translated Gaussian}
\label{sec:newV23-Gaussian-test}

Let
\[
 \varphi_\mu(r)
 ={1\over\sqrt{2\pi}}e^{-(r-\mu)^2/2}.
\]
At the boundary-centered translate \(\mu=R_\alpha\),
\begin{align}
 \varphi_{R_\alpha}^{(2h)}(R_\alpha)
 &=
 {(-1)^h\over\sqrt{2\pi}}
 { (2h)!\over2^hh!},
 \label{eq:newV23-Gaussian-even}\\
 \varphi_{R_\alpha}^{(2h+1)}(R_\alpha)&=0.
 \label{eq:newV23-Gaussian-odd}
\end{align}
Put
\[
 A_h:={ (2h)!\over2^hh!}=(2h-1)!!.
\]

\begin{lemma}[Double-factorial ratios]
\label{lem:newV23-ratios}
For \(1\le s\le h\),
\begin{equation}
 {A_{h-s}\over A_h}
 =
 {1\over
 (2h-1)(2h-3)\cdots(2h-2s+1)}.
 \label{eq:newV23-ratios}
\end{equation}
In particular, for each fixed \(s\),
\[
 {A_{h-s}\over A_h}=O_s(h^{-s}).
\]
\end{lemma}

\begin{proof}
Cancel the last \(s\) factors in
\[
 A_h=(2h-1)(2h-3)\cdots1.
\]
\end{proof}

For \(n=2h+3\), define the complete surface part of
\(\langle D^nk_\alpha,\varphi_{R_\alpha}\rangle\) by
\begin{equation}
 {\cal S}_{n,\alpha}
 :=
 -
 \sum_{q=2}^{n-1}
 j_\alpha^{(q)}(R_\alpha^-)
 \left\langle
 \delta_{R_\alpha}^{(n-1-q)},
 \varphi_{R_\alpha}
 \right\rangle.
 \label{eq:newV23-surface-sum}
\end{equation}

\begin{theorem}[The summed surface row is superexponential]
\label{thm:newV23-surface-growth}
For every fixed \(\alpha\ge\alpha_0\), there is
\(h_0(\alpha)\) such that, for all \(h\ge h_0(\alpha)\),
\begin{equation}
 |{\cal S}_{2h+3,\alpha}|
 \ge
 {c\over\alpha}A_h.
 \label{eq:newV23-surface-growth}
\end{equation}
Consequently there is no finite \(K_\alpha\) satisfying
\[
 |{\cal S}_{n,\alpha}|\le K_\alpha^n
 \qquad(n\ge3).
\]
\end{theorem}

\begin{proof}
The distribution rule
\(\langle\delta_R^{(a)},\phi\rangle=(-1)^a\phi^{(a)}(R)\)
and
\eqref{eq:newV23-Gaussian-odd} show that only even \(q=2s\)
contribute.  The \(q=2\) term has magnitude
\begin{equation}
 {2\over\pi^2\alpha\sqrt{2\pi}}A_h
 \label{eq:newV23-leading-surface}
\end{equation}
by \eqref{eq:newV23-second-jet}.

For \(s\ge2\), use
\eqref{eq:newV23-boundary-jets} and
\eqref{eq:newV23-ratios}.  Relative to
\eqref{eq:newV23-leading-surface}, the \(q=2s\) term is at most
\[
 C\left({4\over\alpha}\right)^{s-1}
 {A_{h+1-s}\over A_h}.
\]
For \(2\le s\le h/2\), this is bounded by
\[
 C\left({C\over\alpha h}\right)^{s-1}.
\]
Its sum tends to zero as \(h\to\infty\).
For \(h/2<s\le h+1\), use the crude bound
\[
 A_{h+1-s}\le A_{\lceil h/2\rceil},
\]
while
\[
 {A_{\lceil h/2\rceil}\over A_h}
 \le(c h)^{-h/2}.
\]
This dominates the at most \(h\) exponential jet factors for fixed
\(\alpha\), so the tail sum also tends to zero relative to
\eqref{eq:newV23-leading-surface}.  For large \(h\), all
nonleading surface terms have total magnitude at most one half of
the leading term.  This proves
\eqref{eq:newV23-surface-growth}.

Finally
\[
 A_h^{1/(2h)}
 \asymp\sqrt h
\]
by Stirling's formula.  Thus no exponential \(K_\alpha^n\) can
dominate the lower bound.
\end{proof}

\begin{theorem}[The complete distributional derivative has the same
obstruction]
\label{thm:newV23-complete-growth}
Let
\[
 {\cal I}_{n,\alpha}
 :=
 \left\langle
 D^nk_\alpha,\varphi_{R_\alpha}
 \right\rangle.
\]
Then, for fixed \(\alpha\) and all sufficiently large \(h\),
\begin{equation}
 |{\cal I}_{2h+3,\alpha}|
 \ge {c\over\alpha}A_h.
 \label{eq:newV23-complete-growth}
\end{equation}
\end{theorem}

\begin{proof}
By \eqref{eq:newV23-distribution-reduced},
\[
 {\cal I}_{n,\alpha}
 =
 \int_{-\infty}^{R_\alpha}
 j_\alpha^{(n)}(r)\varphi_{R_\alpha}(r)\,dr
 +{\cal S}_{n,\alpha}.
\]
The Fourier representation gives the pointwise bound
\[
 |j_\alpha^{(n)}(r)|
 \le(2/\sqrt\alpha)^n.
\]
Hence the interior integral is at most
\((2/\sqrt\alpha)^n\), an exponential in \(n\).  The lower bound
\eqref{eq:newV23-surface-growth} is superexponential and therefore
dominates this interior term for all sufficiently large \(h\).
\end{proof}

\begin{firewall}[Precise scope of the translated test]
\label{fire:newV23-test-scope}
The translated Gaussian test is the conditional kernel encountered
when one Gaussian replica is conditioned near the chart boundary.
\Cref{thm:newV23-complete-growth} rules out any direct chart-BKAR
argument requiring a derivative-row estimate uniform in the
conditional mean.  It does not by itself rule out an as-yet unknown
global cancellation which integrates all conditional means and all
trees before taking norms.
\end{firewall}

\section{Antipodal representation values do not center away}
\label{sec:newV23-antipode}

For spin \(j\), let
\(\varepsilon_j=(-1)^{2j}\).  The central element \(-I\) acts as
\[
 D^j(-I)=\varepsilon_jI.
\]
The Wilson or geodesic mean is \(q_{\alpha,j}I\).

\begin{lemma}[Centered antipodal value]
\label{lem:newV23-centered-antipode}
At the chart boundary,
\begin{equation}
 QD^j(-I)
 =
 (\varepsilon_j-q_{\alpha,j})I.
 \label{eq:newV23-centered-antipode}
\end{equation}
For every fixed half-integer \(j\),
\begin{equation}
 |\varepsilon_j-q_{\alpha,j}|
 \longrightarrow2
 \qquad(\alpha\to\infty).
 \label{eq:newV23-half-antipode}
\end{equation}
Thus centering does not annihilate the leading boundary value
uniformly over irreducible inputs.
\end{lemma}

\begin{proof}
The first identity is the definition of centering and Schur
centrality of the mean.  For half-integer spin,
\(\varepsilon_j=-1\), while concentration at the identity gives
\(q_{\alpha,j}\to1\).  This proves
\eqref{eq:newV23-half-antipode}.
\end{proof}

\begin{assessment}[Decision on the chart route]
\label{ass:newV23-chart-decision}
The ordinary Haar jets are exponential, but the complete surface
distribution paired with a boundary-centered conditional Gaussian
has double-factorial growth.  The leading row is generated by the
nonzero second Haar jet and is not removed by centering
half-integer representation inputs.  Therefore the direct
principal-chart Gaussian BKAR route cannot supply the uniform
all-degree vertex card required by GIAC.  In accordance with the V22
programme, the attack now abandons that chart-differentiation route
and retains the compact-group law intrinsically.
\end{assessment}

\section{The intrinsic score leaf}
\label{sec:newV23-score-leaf}

Let \(A_{\alpha,u}\) be the nonnegative generator associated with
the V21 Dirichlet form, and let
\[
 R_{\alpha,u}=A_{\alpha,u}^{-1}Q_{\alpha,u}.
\]
The gap \eqref{eq:newV21-full-gap} and
\eqref{eq:newV22-score-value} immediately give a strong negative
norm for the one score insertion.

\begin{proposition}[Negative norm of the exact score]
\label{prop:newV23-score-negative}
Uniformly in \(u\in[0,1]\),
\begin{equation}
 \|\widetilde S_{\alpha,u}\|_{-1}
 :=
 \|dR_{\alpha,u}\widetilde S_{\alpha,u}\|_2
 \le {C\over\alpha^{3/2}}.
 \label{eq:newV23-score-negative}
\end{equation}
\end{proposition}

\begin{proof}
For every centered \(F\), the spectral theorem and the gap
\(\lambda_{\alpha,u}\ge c\alpha\) give
\[
 \|dR_{\alpha,u}F\|_2^2
 =
 \langle F,R_{\alpha,u}F\rangle
 \le {1\over c\alpha}\|F\|_2^2.
\]
Apply this with \(F=\widetilde S_{\alpha,u}\) and use
\(\|F\|_2\le C/\alpha\) from
\eqref{eq:newV22-score-value}.
\end{proof}

\begin{assessment}[The next viable payer]
\label{ass:newV23-next-payer}
The score can be placed as a distinguished leaf and paid by
\eqref{eq:newV23-score-negative} before any representation
resolvent is used.  This is not circular: it follows from the proved
path gap and explicit score row.  What remains is to express
\(\kappa_{m+1}(D^{R_1},\ldots,D^{R_m},\widetilde S)\) as one
covariance between the score and a complete cumulant Wick
polynomial, then differentiate that polynomial without expanding
its proper partitions separately.
\end{assessment}

\section{V24 programme and claim boundary}
\label{sec:newV23-next}

\begin{programme}[V24: cumulant Wick polynomial and score-leaf
pairing]
\label{prog:newV23-V24}
\begin{enumerate}[leftmargin=3em]
\item Define the complete cumulant Wick polynomial
      \(\mathcal W_m(f_1,\ldots,f_m)\) by Möbius inversion with one
      unaveraged product block.
\item Prove exactly that
      \[
      \kappa_{m+1}(f_1,\ldots,f_m,\widetilde S)
      =
      \mathbb E[
      \widetilde S\,\mathcal W_m(f_1,\ldots,f_m)].
      \]
\item Apply the covariance--resolvent identity once, on the score
      leaf, using \eqref{eq:newV23-score-negative}.
\item Differentiate the \emph{complete} Wick polynomial before
      taking norms and derive its exact subset recursion.
\item Test whether its derivative has the aggregate Prüfer weight
      with an exponential base.  Do not estimate its partition
      terms separately.
\item If the derivative recursion regenerates a Bell factorial,
      identify the first uncanceled proper-partition row and return
      to the connected score identity before inserting a resolvent.
\end{enumerate}
\end{programme}

\begin{firewall}[Claim boundary after V23]
\label{fire:newV23-boundary}
V23 computes the complete distributional boundary series, proves
superexponential double-factorial growth against a
boundary-centered Gaussian kernel, and shows that centering does not
remove it uniformly.  It therefore abandons the direct chart forest
and proves instead the intrinsic score negative-norm bound
\(C\alpha^{-3/2}\).  It does not prove the Wick-polynomial derivative
card, GIAC, GCSIC, the Wilson aggregate strong card, all-order strong
MTA, WUFC, CPM, cutoff removal, reflection positivity,
Osterwalder--Schrader reconstruction, a continuum Yang--Mills
measure, or a positive Hamiltonian gap.  The full Yang--Mills mass
gap has not been proved.
\end{firewall}

\begin{theorem}[V23 result ledger]
\label{thm:newV23-results}
V23:
\begin{enumerate}[label=\textup{(V23.\arabic*)},leftmargin=6.5em]
\item gives the complete distributional derivative formula for the
      truncated Haar density;
\item proves every ordinary Haar boundary jet has exponential
      degree dependence;
\item proves the summed surface distribution and complete weak
      derivative have double-factorial translated-Gaussian growth;
\item proves half-integer antipodal values survive centering;
\item rigorously retires the direct chart-BKAR route; and
\item proves the intrinsic exact-score negative norm
      \(O(\alpha^{-3/2})\).
\end{enumerate}
\end{theorem}

\begin{proof}
Use
\Cref{lem:newV23-distribution,
cor:newV23-boundary-jets,
thm:newV23-surface-growth,
thm:newV23-complete-growth,
lem:newV23-centered-antipode,
prop:newV23-score-negative}.
\end{proof}

\paragraph{Parent-version record.}
V23 was formed from
\texttt{Renewal\_Yang\_Mills\_Mass\_Gap\_Research\_Notes\_V22.tex}.
The V22 parent had \(10\,146\) logical source lines,
\(329\,198\) bytes, and SHA--256
\[
 \texttt{08d506df506aec9bc138e8b7f4bae66b236dbe2f5eb3938b52d20464fd4a5b5a}.
\]
The next compulsory companion audit is V25.

% =============================================================================
% END APPENDED FIFTH-SERIES V23 MATERIAL
% =============================================================================

% =============================================================================
% BEGIN APPENDED FIFTH-SERIES V24 MATERIAL
% =============================================================================

\chapter{V24: Cumulant Wick polynomials and the score-leaf no-go}
\label{chap:newV24-Wick}

\section{Subset-convolution inverse of the moments}
\label{sec:newV24-inverse}

Let \(I\) be a finite label set.  Put
\[
 F_A(U):=\bigotimes_{i\in A}f_i(U),
 \qquad
 M_A:=\mathbb EF_A,
 \qquad
 F_\varnothing=M_\varnothing=I.
\]
Operators supported on disjoint tensor legs commute.  Define
\begin{equation}
 c_\varnothing:=I,
 \qquad
 c_J
 :=
 \sum_{\rho\in\Pi(J)}
 (-1)^{|\rho|}|\rho|!
 \bigotimes_{B\in\rho}M_B
 \quad(J\ne\varnothing).
 \label{eq:newV24-c}
\end{equation}

\begin{lemma}[Exact subset-convolution inverse]
\label{lem:newV24-convolution-inverse}
For every finite \(J\),
\begin{equation}
 \sum_{A\subseteq J}M_Ac_{J\setminus A}
 =
 \begin{cases}
 I,&J=\varnothing,\\
 0,&J\ne\varnothing.
 \end{cases}
 \label{eq:newV24-convolution-inverse}
\end{equation}
\end{lemma}

\begin{proof}
Introduce commuting nilpotent variables \(z_i\), \(z_i^2=0\), and
the finite moment polynomial
\[
 {\cal M}(\boldsymbol z)
 :=\sum_{A\subseteq J}M_Az_A,
 \qquad z_A:=\prod_{i\in A}z_i.
\]
Writing \({\cal M}=I+X\), its inverse is the finite series
\[
 {\cal M}^{-1}
 =\sum_{k\ge0}(-1)^kX^k.
\]
The coefficient of \(z_J\) in \(X^k\) is the sum over ordered
partitions of \(J\) into \(k\) nonempty blocks.  Every unordered
partition has \(k!\) orderings, so this coefficient is precisely
\[
 \sum_{\substack{\rho\in\Pi(J)\\|\rho|=k}}
 k!\bigotimes_{B\in\rho}M_B.
\]
Thus the coefficient of \(z_J\) in \({\cal M}^{-1}\) is \(c_J\).
The coefficient identity in
\({\cal M}{\cal M}^{-1}=I\) is
\eqref{eq:newV24-convolution-inverse}.
\end{proof}

\section{The complete cumulant Wick polynomial}
\label{sec:newV24-Wick}

\begin{definition}[Multilinear cumulant Wick polynomial]
\label{def:newV24-Wick}
Define
\begin{equation}
 {\cal W}_I(U)
 :=
 \sum_{A\subseteq I}F_A(U)c_{I\setminus A}.
 \label{eq:newV24-Wick}
\end{equation}
\end{definition}

\begin{lemma}[Centering and constant annihilation]
\label{lem:newV24-Wick-properties}
If \(I\ne\varnothing\), then
\begin{equation}
 \mathbb E{\cal W}_I=0.
 \label{eq:newV24-Wick-centered}
\end{equation}
If any input \(f_i\) is constant, then
\begin{equation}
 {\cal W}_I=0.
 \label{eq:newV24-Wick-constant}
\end{equation}
\end{lemma}

\begin{proof}
Taking expectation in \eqref{eq:newV24-Wick} gives
\eqref{eq:newV24-convolution-inverse}.

For the second assertion, write \(f_i=a\), so every moment containing
\(i\) is \(a\) times the moment with \(i\) removed.  Equivalently,
the moment polynomial factors as
\[
 {\cal M}_I(\boldsymbol z)
 =(I+az_i){\cal M}_{I\setminus\{i\}}(\boldsymbol z).
\]
The random polynomial
\(\sum_AF_Az_A\) has the same factor.  Their quotient, whose
\(z_I\)-coefficient is \({\cal W}_I\), is therefore independent of
\(z_i\).  That coefficient is zero.
\end{proof}

\begin{theorem}[Exact score--Wick pairing]
\label{thm:newV24-score-Wick}
Let \(S\) be a centered scalar observable.  Then
\begin{equation}
 \boxed{\qquad
 \kappa_{|I|+1}\bigl((f_i)_{i\in I},S\bigr)
 =
 \mathbb E\left[S{\cal W}_I\right].
 \qquad}
 \label{eq:newV24-score-Wick}
\end{equation}
\end{theorem}

\begin{proof}
Expand \({\cal W}_I\) using
\eqref{eq:newV24-c}:
\[
 \mathbb E[S{\cal W}_I]
 =
 \sum_{A\subseteq I}
 \sum_{\rho\in\Pi(I\setminus A)}
 (-1)^{|\rho|}|\rho|!\,
 \mathbb E(SF_A)
 \bigotimes_{B\in\rho}M_B.
\]
The term \(A=\varnothing\) vanishes because \(\mathbb ES=0\).
Every remaining pair \((A,\rho)\) specifies uniquely a partition of
\(I\cup\{*\}\): the block containing the score label \(*\) is
\(A\cup\{*\}\), and the other blocks are \(\rho\).  This partition
has \(1+|\rho|\) blocks, so its cumulant Möbius coefficient is
\[
 (-1)^{|\rho|}|\rho|!.
\]
The displayed sum is therefore the complete partition formula for
the cumulant on the left.
\end{proof}

\section{The derivative recursion}
\label{sec:newV24-derivative}

\begin{theorem}[Exact Appell derivative identity]
\label{thm:newV24-Appell}
For every nonempty \(I\),
\begin{equation}
 \boxed{\qquad
 d{\cal W}_I
 =
 \sum_{i\in I}
 (df_i)\otimes{\cal W}_{I\setminus\{i\}}.
 \qquad}
 \label{eq:newV24-Appell}
\end{equation}
The tensor factors are restored to canonical label order.
\end{theorem}

\begin{proof}
The coefficients \(c_J\) are deterministic, so differentiation of
\eqref{eq:newV24-Wick} gives
\[
 d{\cal W}_I
 =
 \sum_{A\subseteq I}
 \sum_{i\in A}
 (df_i)\otimes F_{A\setminus\{i\}}c_{I\setminus A}.
\]
For fixed \(i\), put \(B=A\setminus\{i\}\).  The inner sum becomes
\[
 \sum_{B\subseteq I\setminus\{i\}}
 F_Bc_{(I\setminus\{i\})\setminus B}
 =
 {\cal W}_{I\setminus\{i\}},
\]
which proves the identity.
\end{proof}

\begin{corollary}[Exact intrinsic score-leaf formula]
\label{cor:newV24-score-leaf}
For the V21 path score,
\begin{equation}
 \kappa_{m+1}
 (f_1,\ldots,f_m,\widetilde S)
 =
 \mathbb E\,
 \Gamma\left(
 R\widetilde S,
 {\cal W}_{[m]}
 \right)
 =
 \sum_{i=1}^m
 \mathbb E\,
 \Gamma\left(
 R\widetilde S,
 f_i
 \right)
 {\cal W}_{[m]\setminus\{i\}},
 \label{eq:newV24-score-leaf}
\end{equation}
with the operator order inherited from
\eqref{eq:newV24-Appell}.
\end{corollary}

\begin{proof}
The score is centered, so
\(\mathbb E[\widetilde S{\cal W}_{[m]}]\) is a covariance.
Apply the reversible covariance--resolvent identity to the score
slot, then use the diffusion product rule and
\Cref{thm:newV24-Appell}.  The scalar score gradient commutes with
the tensor legs.
\end{proof}

\section{What the score negative norm can prove}
\label{sec:newV24-negative}

Cauchy--Schwarz in
\eqref{eq:newV24-score-leaf} gives
\begin{equation}
 \left\|
 \kappa_{m+1}(f_1,\ldots,f_m,\widetilde S)
 \right\|
 \le
 \|\widetilde S\|_{-1}
 \|d{\cal W}_{[m]}\|_{2,c}.
 \label{eq:newV24-CS}
\end{equation}
By V23, the first factor is \(O(\alpha^{-3/2})\).
The tempting derivative card would be
\begin{equation}
 \|d{\cal W}_{[m]}\|_{2,c}
 \le
 K^m\sqrt\alpha
 \left(\prod_i b_i\right)
 \left(\sum_i b_i\right)^{m-2}.
 \label{eq:newV24-false-derivative-card}
\end{equation}
Together with the all-thermal final payer of order \(O(\alpha)\),
this would prove GCSIC.  It is, however, false.

\begin{theorem}[The normwise Wick derivative card fails at three
labels]
\label{thm:newV24-derivative-no-go}
Consider one standard Gaussian scalar \(X\), and take three
centered linear inputs
\[
 f_i=b_iX.
\]
Then
\begin{align}
 {\cal W}_{123}
 &=b_1b_2b_3(X^3-3X),
 \label{eq:newV24-Wick-cubic}\\
 {d\over dX}{\cal W}_{123}
 &=3b_1b_2b_3(X^2-1),
 \label{eq:newV24-Wick-cubic-derivative}\\
 \left\|
 {d\over dX}{\cal W}_{123}
 \right\|_2
 &=3\sqrt2\,b_1b_2b_3.
 \label{eq:newV24-Wick-cubic-norm}
\end{align}
For \(b_1=b_2=b_3=b\downarrow0\), this cannot be bounded by
\[
 K^3b_1b_2b_3(b_1+b_2+b_3).
\]
Thus \eqref{eq:newV24-false-derivative-card} cannot hold uniformly,
even in the tangent Gaussian model.
\end{theorem}

\begin{proof}
The subset inverse coefficients for centered Gaussian linear inputs
give the ordinary third Hermite polynomial:
\[
 X^3-3\mathbb EX^2\,X=X^3-3X.
\]
Differentiate and use
\(\|X^2-1\|_2=\sqrt2\).
For equal marks, the ratio of
\eqref{eq:newV24-Wick-cubic-norm} to the proposed right side is
\(\sqrt2/(Kb)\), which diverges as \(b\downarrow0\).
\end{proof}

\begin{firewall}[The score cumulant is not disproved]
\label{fire:newV24-no-go-scope}
\Cref{thm:newV24-derivative-no-go} disproves only the use of
Cauchy--Schwarz after replacing the score by its negative norm.
It does not disprove GCSIC.  The exact score has an odd gradient
whose tangent pairing with the even quadratic
\eqref{eq:newV24-Wick-cubic-derivative} vanishes.  The norm estimate
\eqref{eq:newV24-CS} discards precisely that cancellation.
\end{firewall}

\section{Parity of the first tangent score pairing}
\label{sec:newV24-parity}

In scaled tangent coordinates the leading score is a centered even
quartic radial polynomial.  Its Ornstein--Uhlenbeck resolvent is
even and its gradient is odd.  For linear Gaussian inputs,
\({\cal W}_{[m]}\) is a homogeneous-plus-lower Wick polynomial of
parity \(m\), so its gradient has parity \(m-1\).

\begin{proposition}[Odd input arity cancels at the score leaf]
\label{prop:newV24-parity}
In the all-linear Gaussian tangent sector,
\begin{equation}
 \mathbb E\,
 \Gamma_0\left(
 R_0S_4,
 {\cal W}_{[m]}
 \right)=0
 \qquad(m\ \hbox{odd}),
 \label{eq:newV24-parity}
\end{equation}
where \(S_4\) is any centered even quartic polynomial.
\end{proposition}

\begin{proof}
The gradient of \(R_0S_4\) is odd.  The gradient of the
\(m\)-th Gaussian Wick polynomial has parity \(m-1\).  If \(m\) is
odd, the latter is even.  Their scalar contraction is odd and has
zero expectation under the centered Gaussian law.
\end{proof}

\begin{assessment}[The correct score card is a pairing card]
\label{ass:newV24-pairing}
The exact Wick polynomial eliminates constants and exposes every
label before differentiation, but its derivative norm does not
contain the full connected-tree power.  The missing power is
restored only in the pairing with the explicit quartic score
resolvent.  The next target must therefore estimate the complete
bilinear form in \eqref{eq:newV24-score-leaf}, not its two factors
separately.
\end{assessment}

\section{A cancellation-preserving target}
\label{sec:newV24-target}

\begin{definition}[Score--Wick pairing card]
\label{def:newV24-SWPC}
The SWPC is the estimate
\begin{equation}
 \left\|
 \mathbb E_{\alpha,u}
 \Gamma\left(
 R_{\alpha,u}\widetilde S_{\alpha,u},
 {\cal W}_{[m]}(D^{R_1},\ldots,D^{R_m})
 \right)
 \right\|_{\rm op}
 \le
 {K^m\over\alpha}
 \left(\prod_i b_i\right)
 \left(\sum_i b_i\right)^{m-2},
 \label{eq:newV24-SWPC}
\end{equation}
uniformly in \(u\), before the final \(O(\alpha)\) output payer.
\end{definition}

\begin{theorem}[SWPC implies GCSIC]
\label{thm:newV24-SWPC-sufficient}
The SWPC, followed by the inherited all-thermal final
mass--resolvent bound \(O(\alpha)\), implies
\eqref{eq:newV22-GCSIC}.
\end{theorem}

\begin{proof}
The left side of \eqref{eq:newV24-SWPC} is exactly the unresolvented
score cumulant by
\Cref{cor:newV24-score-leaf}.  Multiply by the final payer and
cancel its factor \(\alpha\) against the displayed
\(\alpha^{-1}\).
\end{proof}

\section{V25 programme and claim boundary}
\label{sec:newV24-next}

\begin{programme}[V25: audit and the tangent Score--Wick pairing]
\label{prog:newV24-V25}
\begin{enumerate}[leftmargin=3em]
\item Perform the compulsory read-only SM/NS companion audit.
\item Compute the SWPC exactly in the scalar Gaussian Weyl model,
      retaining the quartic score before the expectation.
\item Express the result as a connected graph polynomial and test
      the Prüfer weight with an exponential base.
\item Extend the calculation to matrix linearizations using
      row--column Gaussian contraction before taking operator norms.
\item Separate the geodesic-Haar and Wilson-generator defects from
      the exact tangent pairing.
\item Do not use \eqref{eq:newV24-CS}; it is ruled out by
      \Cref{thm:newV24-derivative-no-go}.
\end{enumerate}
\end{programme}

\begin{firewall}[Claim boundary after V24]
\label{fire:newV24-boundary}
V24 constructs the exact cumulant Wick polynomial, proves its
Appell derivative identity and the one-resolvent score-leaf formula,
and proves that a normwise derivative card is false at three
Gaussian labels.  It identifies the cancellation-preserving SWPC
but does not prove it, GIAC, GCSIC, the Wilson aggregate strong
card, all-order strong MTA, WUFC, CPM, cutoff removal, reflection
positivity, Osterwalder--Schrader reconstruction, a continuum
Yang--Mills measure, or a positive Hamiltonian gap.  The full
Yang--Mills mass gap has not been proved.
\end{firewall}

\begin{theorem}[V24 result ledger]
\label{thm:newV24-results}
V24:
\begin{enumerate}[label=\textup{(V24.\arabic*)},leftmargin=6.5em]
\item constructs the exact subset-convolution inverse of all
      moments;
\item defines centered multilinear cumulant Wick polynomials and
      proves their constant-annihilation property;
\item proves the exact score--Wick pairing and Appell derivative
      recursion;
\item derives the intrinsic one-score-resolvent formula;
\item disproves the separated Wick-derivative norm card at arity
      three; and
\item defines the cancellation-preserving SWPC sufficient for
      GCSIC.
\end{enumerate}
\end{theorem}

\begin{proof}
Use
\Cref{lem:newV24-convolution-inverse,
lem:newV24-Wick-properties,
thm:newV24-score-Wick,
thm:newV24-Appell,
cor:newV24-score-leaf,
thm:newV24-derivative-no-go,
thm:newV24-SWPC-sufficient}.
\end{proof}

\paragraph{Parent-version record.}
V24 was formed from
\texttt{Renewal\_Yang\_Mills\_Mass\_Gap\_Research\_Notes\_V23.tex}.
The V23 parent had \(10\,658\) logical source lines,
\(343\,899\) bytes, and SHA--256
\[
 \texttt{138b0be3c3ed5f1eb3bfeb9f3acfc95a7d240c3599eedbfe8d29c4bd8f2a04e5}.
\]
The next compulsory companion audit is V25.

% =============================================================================
% END APPENDED FIFTH-SERIES V24 MATERIAL
% =============================================================================

% =============================================================================
% BEGIN APPENDED FIFTH-SERIES V25 MATERIAL
% =============================================================================

\chapter{V25: Quartic-score clock census and the payer-matching
obstruction}
\label{chap:newV25-quartic}

\section{Compulsory read-only companion audit}
\label{sec:newV25-audit}

The newest active companion files inspected before this append were
\begin{center}
\begin{tabular}{p{0.62\textwidth}r r}
\hline
file & bytes & logical lines\\
\hline
\texttt{Renewal\_SM\_Einstein\_Continuum\_Research\_Notes\_V30.tex}
& \(385\,907\) & \(9\,847\)\\
\texttt{Renewal\_NS\_Stability\_Research\_Notes\_V31.tex}
& \(887\,537\) & \(23\,052\)\\
\hline
\end{tabular}
\end{center}
Their SHA--256 digests were respectively
\begin{align}
 &\texttt{8c5c804faf73917d9a11bff13b9861d2c8654ac2c39e10a8e0ccee1f03b2ca42},
 \label{eq:newV25-SM-hash}\\
 &\texttt{f1121b62238ad5491bb7cf03635ea9934942edf573ed8faaa9ee79e1d38a6696}.
 \label{eq:newV25-NS-hash}
\end{align}
Neither companion file was modified.

SM V30 independently derives a complete cumulant Wick polynomial,
the exact score--Wick identity, and a one-negative-norm score leaf.
It correctly leaves its Ward-first Wick derivative card open.  The
YM V24 Gaussian cubic test shows why that open hypothesis must not
be replaced by a normwise derivative estimate without additional
cancellation.

NS V31 proves an exact dyadic unmarking identity.  Its normalized
probability mark is harmless, but returning to the raw critical
coarea costs exactly one first radial moment; heat formation and
high-parent lifting then return the original continuation-class
stock.  The transferable lesson is only an order of operations:
retain the complete correlated object through the scale-restoring
operation.  No Navier--Stokes estimate, sign, or continuation claim
is imported here.

\begin{firewall}[No companion theorem is transferred]
\label{fire:newV25-no-transfer}
The SM score acts on a massless gauge--matter--gravity shell and the
NS first moment acts on a parabolic Fourier tail.  Neither gives a
Wilson one-link cumulant estimate, a compact-group output-payer
identity, or a Yang--Mills mass gap.  The common information is the
diagnostic warning that a separately bounded restoring operator can
consume the scale gained by a centered insertion.
\end{firewall}

\section{The exact scalar Gaussian score graph}
\label{sec:newV25-Gaussian-graph}

Let \(X\) be a standard real Gaussian and let
\[
 H_4(X):=X^4-6X^2+3.
\]
For \(m\ge1\), let \(\mathfrak C_m^{(4)}\) be the set of connected
loop-free multigraphs on vertices
\(\{0,1,\ldots,m\}\), described by edge multiplicities
\[
 n_{ij}\in\mathbb N_0,\qquad 0\le i<j\le m,
\]
whose distinguished root has degree
\[
 \sum_{j=1}^m n_{0j}=4.
\]
Parallel edges are allowed.

\begin{theorem}[Exact quartic-root connected graph formula]
\label{thm:newV25-exact-graph}
For all real \(t_1,\ldots,t_m\),
\begin{align}
 &\kappa_{m+1}
 \left(
 e^{it_1X},\ldots,e^{it_mX},H_4(X)
 \right)
 \notag\\
 &\quad=
 e^{-\frac12\sum_{j=1}^m t_j^2}
 \sum_{\boldsymbol n\in\mathfrak C_m^{(4)}}
 {4!\over
   \displaystyle
   \prod_{j=1}^m n_{0j}!
   \prod_{1\le i<j\le m}n_{ij}!}
 \prod_{j=1}^m(it_j)^{n_{0j}}
 \prod_{1\le i<j\le m}(-t_it_j)^{n_{ij}}.
 \label{eq:newV25-exact-graph}
\end{align}
The series is absolutely convergent on every compact subset of
\(\mathbb R^m\).
\end{theorem}

\begin{proof}
Use Gaussian normal ordering
\[
 e^{itX}=e^{-t^2/2}:e^{itX}:,
 \qquad
 H_4(X)=:X^4:.
\]
Multilinearity extracts the displayed product of one-vertex
factors.  Wick contraction between two normal-ordered exponential
vertices \(i,j\) has weight
\((it_i)(it_j)=-t_it_j\).  Summing \(n_{ij}\) parallel contractions
gives the exponential-series divisor \(n_{ij}!\).

The four labelled legs at the root must contract with exponential
vertices.  If \(n_{0j}\) of them meet vertex \(j\), their allocation
has multiplicity
\[
 {4!\over\prod_jn_{0j}!}
\]
and weight \(\prod_j(it_j)^{n_{0j}}\).
The moment expansion is therefore the sum over all such
multigraphs, connected or not.  The moment--cumulant Möbius
inversion retains exactly the graphs connected on the complete
label set \(\{0,\ldots,m\}\), which proves
\eqref{eq:newV25-exact-graph}.

Absolute convergence follows by discarding connectedness and
summing every family of parallel edges.  The root sum is finite,
and the remaining sum is bounded on compact sets by
\[
 4!\left(\sum_j|t_j|\right)^4
 \exp\left(\sum_{i<j}|t_it_j|\right).
\]
\end{proof}

\begin{lemma}[Exact clock degree of a quartic-root graph]
\label{lem:newV25-clock-degree}
Every monomial in \eqref{eq:newV25-exact-graph} is divisible by
\(\prod_{j=1}^m t_j\).  If \(E\) is the number of input--input
edges counted with multiplicity, its total \(t\)-degree is
\[
 4+2E.
 \label{eq:newV25-graph-degree}
\]
Connectivity forces
\[
 E\ge\max\{0,m-4\}.
 \label{eq:newV25-edge-lower}
\]
Consequently the least possible input degree is
\[
 d_m^{\min}
 =
 4+2\max\{0,m-4\}.
 \label{eq:newV25-min-degree}
\]
\end{lemma}

\begin{proof}
Every input vertex belongs to a connected graph and therefore has
positive degree.  Its edge factors supply at least one copy of its
own \(t_j\).  The four root edges supply four input half-edges and
every input--input edge supplies two, proving
\eqref{eq:newV25-graph-degree}.  A connected multigraph on
\(m+1\) vertices has at least \(m\) edges.  There are four root
edges, so \(4+E\ge m\), which gives
\eqref{eq:newV25-edge-lower}.  Both lower bounds can be attained:
for \(m\le4\), distribute the four root legs so that every input is
hit; for \(m>4\), hit four inputs and attach each remaining input by
one input--input edge.
\end{proof}

\begin{assessment}[The quartic score is exactly two clocks]
\label{ass:newV25-two-clocks}
At fixed representation,
\(b\asymp\alpha^{-1/2}\).  The score coefficient
\(\alpha^{-1}\) therefore has the homogeneity of two Wilson
clocks.  Adding those two virtual clocks to
\eqref{eq:newV25-min-degree} gives
\[
 6+2\max\{0,m-4\}\ge2m-2,
\]
with equality for every \(m\ge4\).  Thus the quartic score can
supply the ordinary input-tree homogeneity, but it has no further
generic \(\alpha^{-1}\) left for a separately estimated
\(O(\alpha)\) output payer.
\end{assessment}

\section{The arity-four coefficient}
\label{sec:newV25-arity-four}

\begin{corollary}[The first arity-four monomial is nonzero]
\label{cor:newV25-four-leading}
As \(\boldsymbol t\to0\),
\begin{equation}
 \kappa_5
 \left(
 e^{it_1X},e^{it_2X},e^{it_3X},e^{it_4X},H_4(X)
 \right)
 =
 24t_1t_2t_3t_4+O_{\rm an}(|\boldsymbol t|^6).
 \label{eq:newV25-four-leading}
\end{equation}
Equivalently,
\[
 \left.
 \partial_{t_1}\partial_{t_2}\partial_{t_3}\partial_{t_4}
 \kappa_5
 \left(
 e^{it_1X},\ldots,e^{it_4X},H_4(X)
 \right)
 \right|_{\boldsymbol t=0}
 =24.
\]
\end{corollary}

\begin{proof}
At total degree four, the only connected graph has one root edge
to each input and no input--input edge.  Its coefficient in
\eqref{eq:newV25-exact-graph} is
\[
 4!i^4=24.
\]
Every other graph has an input--input edge and hence adds two to
the degree.  The exterior Gaussian factor has constant term one
and only even corrections.
\end{proof}

The same coefficient can be checked without diagrams.  If
\(Y=(Y_1,Y_2,Y_3)\) is standard Gaussian and
\[
 Q_4(Y):=|Y|^4-15,
\]
then exponential tilting in the first coordinate gives
\begin{equation}
 \mathbb E\left[
 e^{hY_1-h^2/2}Q_4(Y)
 \right]
 =h^4+10h^2.
 \label{eq:newV25-radial-tilt}
\end{equation}

\begin{lemma}[Radial quartic cumulant]
\label{lem:newV25-radial-cumulant}
\begin{equation}
 \kappa_5^{\gamma_3}
 (Y_1,Y_1,Y_1,Y_1,Q_4)=24.
 \label{eq:newV25-radial-cumulant}
\end{equation}
\end{lemma}

\begin{proof}
For \(Y\sim N(he_1,I_3)\),
\[
 \mathbb E|Y|^4=h^4+10h^2+15,
\]
which proves \eqref{eq:newV25-radial-tilt}.  Differentiate the
joint cumulant generating function once in the \(Q_4\) source and
four times in \(h\) at zero.  The fourth derivative of
\(h^4+10h^2\) is \(24\).
\end{proof}

\section{Transfer to the actual \(SU(2)\) geodesic--Wilson path}
\label{sec:newV25-SU2-transfer}

Write
\[
 U=\cos\theta\,I+i\sin\theta\,\boldsymbol n\cdot\boldsymbol\sigma,
 \qquad
 Y=\sqrt\alpha\,\theta\boldsymbol n,
\]
and define the bounded fundamental matrix functional
\begin{equation}
 g_\alpha(U)
 :={1\over2i}\operatorname{tr}(\sigma_1U)
 =n_1\sin\theta.
 \label{eq:newV25-fundamental-coordinate}
\end{equation}
It is a fixed norm-bounded linear functional of the fundamental
representation matrix.

\begin{lemma}[Uniform fixed-moment tangent expansion]
\label{lem:newV25-path-expansion}
For every fixed \(p<\infty\), uniformly in \(u\in[0,1]\),
\begin{align}
 g_\alpha
 &={Y_1\over\sqrt\alpha}
   +O_{L^p}(\alpha^{-3/2}),
 \label{eq:newV25-g-expansion}\\
 \widetilde S_{\alpha,u}
 &={|Y|^4-15\over24\alpha}
   +O_{L^p}(\alpha^{-2}).
 \label{eq:newV25-score-expansion}
\end{align}
Moreover, every fixed polynomial moment of \(Y\) under
\(\widetilde\nu_{\alpha,u}\) differs from its standard
three-dimensional Gaussian value by \(O(\alpha^{-1})\).
\end{lemma}

\begin{proof}
The elementary bound
\[
 \left|{\sin z\over z}-1\right|\le Cz^2
 \qquad(0\le z\le\pi)
\]
and the uniform scaled moments from V21 give
\eqref{eq:newV25-g-expansion}.  Taylor's theorem on
\([0,\pi]\) gives
\[
 R(\theta)={\theta^4\over24}+O(\theta^6).
\]
Hence
\[
 \alpha R(|Y|/\sqrt\alpha)
 ={|Y|^4\over24\alpha}
  +O_{L^p}(\alpha^{-2}).
\]

It remains to center this expansion.  On every fixed scaled ball,
the Haar factor and path density have expansions
\[
 \left({\sin(|Y|/\sqrt\alpha)\over|Y|/\sqrt\alpha}\right)^2
 =1+O\left({|Y|^2\over\alpha}\right),
\]
\[
 e^{-\alpha V_u(|Y|/\sqrt\alpha)}
 =
 e^{-|Y|^2/2}
 \left[
 1+O\left({|Y|^4+|Y|^6\over\alpha}\right)
 \right],
\]
uniformly in \(u\).  To justify integration, split at
\(|Y|=\alpha^{1/8}\).  On the inner part the displayed expansions
are dominated by a fixed Gaussian times a polynomial.  On the
outer part the uniform confinement
\(V_u(\theta)\ge2\theta^2/\pi^2\) makes every fixed polynomial
moment \(O(e^{-c\alpha^{1/4}})\).  Normalization has the same
expansion.  Thus fixed polynomial moments differ from Gaussian
moments by \(O(\alpha^{-1})\), in particular
\[
 \mathbb E_{\alpha,u}|Y|^4=15+O(\alpha^{-1}).
\]
Centering proves \eqref{eq:newV25-score-expansion}.
\end{proof}

\begin{theorem}[Actual path score saturates the unpaid tree scale]
\label{thm:newV25-path-saturation}
Uniformly for \(u\in[0,1]\),
\begin{equation}
 \kappa_5^{\widetilde\nu_{\alpha,u}}
 (g_\alpha,g_\alpha,g_\alpha,g_\alpha,
  \widetilde S_{\alpha,u})
 =
 {1\over\alpha^3}+O(\alpha^{-4}).
 \label{eq:newV25-path-saturation}
\end{equation}
In particular its absolute value is at least
\(\frac12\alpha^{-3}\) for all sufficiently large \(\alpha\).
\end{theorem}

\begin{proof}
Cumulants of fixed arity are finite linear combinations of
moments.  Insert
\eqref{eq:newV25-g-expansion}--%
\eqref{eq:newV25-score-expansion}, use H\"older's inequality and
the uniform moments, and then use the \(O(\alpha^{-1})\)
moment convergence in
\Cref{lem:newV25-path-expansion}.  Multilinearity gives
\begin{align*}
 &\kappa_5^{\widetilde\nu_{\alpha,u}}
 (g_\alpha,g_\alpha,g_\alpha,g_\alpha,
  \widetilde S_{\alpha,u})\\
 &\qquad=
 {1\over24\alpha^3}
 \kappa_5^{\gamma_3}
 (Y_1,Y_1,Y_1,Y_1,Q_4)
 +O(\alpha^{-4}).
\end{align*}
Apply \Cref{lem:newV25-radial-cumulant}.
\end{proof}

\section{The V24 pairing card is false}
\label{sec:newV25-SWPC-no-go}

For four fundamental inputs the inherited thermal marks obey
\[
 b_i=b_{\rm f}
 =\sqrt{s_\alpha(1+C_2(1/2))}
 \asymp\alpha^{-1/2}.
\]

\begin{theorem}[Arity-four no-go for SWPC]
\label{thm:newV25-SWPC-false}
The Score--Wick pairing card
\eqref{eq:newV24-SWPC} is false, already for four fundamental
\(SU(2)\) inputs along the geodesic--Wilson path.
\end{theorem}

\begin{proof}
Let
\[
 {\cal K}_{\alpha,u}
 :=
 \kappa_5^{\widetilde\nu_{\alpha,u}}
 (U,U,U,U,\widetilde S_{\alpha,u})
\]
be the four-fold fundamental tensor score cumulant before the
final output payer.  Applying the fixed bounded functional
\[
 A\longmapsto {1\over2i}\operatorname{tr}(\sigma_1A)
\]
to all four tensor legs gives exactly the scalar cumulant in
\eqref{eq:newV25-path-saturation}.  Hence the operator cross norm
of \({\cal K}_{\alpha,u}\) is at least \(c\alpha^{-3}\), with a
fixed \(c>0\).

By the exact score--Wick and score-leaf identities of V24, this is
also the corresponding contraction of the left side of
\eqref{eq:newV24-SWPC}.  At \(m=4\), the proposed right side is
\[
 {K^4\over\alpha}
 b_{\rm f}^4(4b_{\rm f})^2
 =O(K^4\alpha^{-4}).
\]
No \(K\) independent of \(\alpha\) can dominate the lower bound.
\end{proof}

\begin{assessment}[Append-only correction to V24]
\label{ass:newV25-V24-correction}
The algebraic results of V24 remain valid: the subset inverse, Wick
polynomial, Appell derivative, score--Wick identity, and intrinsic
score-leaf formula are exact.  Its conditional implication
SWPC \(\Rightarrow\) GCSIC is also logically valid.  However,
\Cref{thm:newV25-SWPC-false} proves that the SWPC hypothesis itself
is false.  The V24 statement that odd arity cancels remains correct
at \(m=3\); the first saturating even test occurs at \(m=4\).
\end{assessment}

\section{Corrected card hierarchy}
\label{sec:newV25-cards}

\begin{definition}[Unpaid quartic-score aggregate card]
\label{def:newV25-UQSAC}
The UQSAC is the scale-compatible estimate
\begin{equation}
 \left\|
 \kappa_{m+1}^{\widetilde\nu_{\alpha,u}}
 (D^{R_1},\ldots,D^{R_m},\widetilde S_{\alpha,u})
 \right\|_{\rm op}
 \le
 K^m
 \left(\prod_i b_i\right)
 \left(\sum_i b_i\right)^{m-2},
 \label{eq:newV25-UQSAC}
\end{equation}
before applying the final output payer.
\end{definition}

\begin{definition}[Payer-matched score insertion card]
\label{def:newV25-PMSIC}
If \(\mathcal P_{\rm out}\) denotes the exact selected output
normalization and mass--resolvent operation, the PMSIC is
\begin{equation}
 \left\|
 \mathcal P_{\rm out}
 \kappa_{m+1}^{\widetilde\nu_{\alpha,u}}
 (D^{R_1},\ldots,D^{R_m},\widetilde S_{\alpha,u})
 \right\|_{\rm op}
 \le
 K^m
 \left(\prod_i b_i\right)
 \left(\sum_i b_i\right)^{m-2}.
 \label{eq:newV25-PMSIC}
\end{equation}
The payer is applied to the complete score cumulant before any
operator norm.
\end{definition}

\begin{proposition}[What the tangent graph does and does not prove]
\label{prop:newV25-card-status}
The graph census
\Cref{thm:newV25-exact-graph,lem:newV25-clock-degree} is compatible
with the homogeneity of UQSAC and sharp at four equal fixed
representations.  It is incompatible with any uniform implication
of the form
\[
 \hbox{UQSAC}
 \quad+\quad
 \|\mathcal P_{\rm out}\|=O(\alpha)
 \quad\Longrightarrow\quad
 \hbox{PMSIC}.
 \label{eq:newV25-false-payer-separation}
\]
It neither proves nor disproves PMSIC, because the exact output
projection may cancel the saturating quartic tensor in some
sectors.
\end{proposition}

\begin{proof}
The first statement is
\Cref{ass:newV25-two-clocks,thm:newV25-path-saturation}.  A
separate \(O(\alpha)\) norm bound applied to the saturating
\(\alpha^{-3}\) arity-four cumulant gives only
\(O(\alpha^{-2})\), whereas the four-input tree weight is
\(b_{\rm f}^6\asymp\alpha^{-3}\).  Therefore such a separated
argument loses one Wilson clock pair.  This calculation concerns an
upper bound on the payer norm; deciding PMSIC requires its actual
action on each Peter--Weyl output sector.
\end{proof}

\begin{firewall}[Do not call the universal payer after the score norm]
\label{fire:newV25-no-separated-payer}
From V25 onward, the score comparison may be estimated either as an
unpaid cumulant at the ordinary tree scale or with the exact output
payer already inserted.  One must not demand an extra
\(\alpha^{-1}\) from the score pairing and then multiply by a
universal \(O(\alpha)\) payer.  The former demand is false by
\Cref{thm:newV25-SWPC-false}; the latter multiplication discards
the only possible output-sector cancellation.
\end{firewall}

\section{V26 programme and claim boundary}
\label{sec:newV25-next}

\begin{programme}[V26: exact fundamental output-sector test]
\label{prog:newV25-V26}
\begin{enumerate}[leftmargin=3em]
\item Decompose the four-fold fundamental tensor cumulant in
      \eqref{eq:newV25-path-saturation} into total-spin
      Peter--Weyl output sectors before taking norms.
\item Write the exact final mass--resolvent eigenvalue on every
      sector which receives the radial quartic tensor.
\item Determine whether the saturating
      \(\alpha^{-3}\) coefficient lies in the kernel of the payer,
      receives an \(O(1)\) multiplier, or receives the dangerous
      \(O(\alpha)\) multiplier.
\item If the dangerous sector survives, correct the all-thermal
      strong-card target itself rather than imposing another false
      sufficient card.
\item If it cancels, express that cancellation as a
      projection-complete payer-matched identity stable under
      spectator amplification.
\item Retain the complete quartic-root graph formula when passing
      to arbitrary representations; do not estimate its score and
      payer factors separately.
\end{enumerate}
\end{programme}

\begin{firewall}[Claim boundary after V25]
\label{fire:newV25-boundary}
V25 performs the compulsory companion audit, proves an exact
connected multigraph formula for a quartic Gaussian score, computes
its sharp clock degree, transfers the nonzero arity-four coefficient
to the actual \(SU(2)\) geodesic--Wilson path, and disproves the V24
SWPC.  It defines the corrected unpaid and payer-matched cards but
does not prove either uniformly at all arities.  In particular it
does not prove PMSIC, GIAC, GCSIC, the Wilson aggregate strong card,
all-order strong MTA, WUFC, CPM, cutoff removal, reflection
positivity, Osterwalder--Schrader reconstruction, a continuum
Yang--Mills measure, or a positive Hamiltonian gap.  The full
Yang--Mills mass gap has not been proved.
\end{firewall}

\begin{theorem}[V25 result ledger]
\label{thm:newV25-results}
V25:
\begin{enumerate}[label=\textup{(V25.\arabic*)},leftmargin=6.5em]
\item completes the required read-only SM/NS audit;
\item derives the exact connected multigraph series with a
      degree-four score root;
\item proves that the score coefficient supplies exactly two
      virtual Wilson clocks;
\item computes the nonzero arity-four tangent coefficient;
\item proves the uniform actual-path asymptotic
      \(\alpha^{-3}+O(\alpha^{-4})\);
\item disproves the extra-\(\alpha^{-1}\) SWPC of V24; and
\item reduces the next gate to the exact output-sector action of the
      unique final payer.
\end{enumerate}
\end{theorem}

\begin{proof}
Use
\Cref{thm:newV25-exact-graph,
lem:newV25-clock-degree,
cor:newV25-four-leading,
lem:newV25-radial-cumulant,
lem:newV25-path-expansion,
thm:newV25-path-saturation,
thm:newV25-SWPC-false,
prop:newV25-card-status}.
\end{proof}

\paragraph{Parent-version record.}
V25 was formed from
\texttt{Renewal\_Yang\_Mills\_Mass\_Gap\_Research\_Notes\_V24.tex}.
The V24 parent had \(11\,120\) logical source lines,
\(357\,163\) bytes, and SHA--256
\[
 \texttt{f039186e44fce88f074b78363b4f69566ade9b66eee7cae1ab1b922d119f9dc4}.
\]
The next compulsory companion audit is V30.

% =============================================================================
% END APPENDED FIFTH-SERIES V25 MATERIAL
% =============================================================================

% =============================================================================
% BEGIN APPENDED FIFTH-SERIES V26 MATERIAL
% =============================================================================

\chapter{V26: Total-spin output sectors and the corrected raw
comparison}
\label{chap:newV26-sectors}

\section{Why the output sector must be computed}
\label{sec:newV26-purpose}

V25 proves that the unresolvented four-input score cumulant has
exact size \(\alpha^{-3}\), the same scale as the four-input
Pr\"ufer weight.  It leaves one precise alternative: perhaps the
dangerous tensor lies in the trivial output deleted by \(Q_0\), or
perhaps a nontrivial output survives and the inherited
\(O(\alpha)\) payer really loses one clock pair.  This chapter
settles that alternative in the four-fold fundamental tensor
product.

No companion audit is due at V26.  The newest audit remains the V25
read-only inspection recorded in
\Cref{sec:newV25-audit}.

\section{The invariant quartic tensor}
\label{sec:newV26-invariant}

Let
\[
 {\cal H}:=(\mathbb C^2)^{\otimes4},
\]
let \(\sigma_a^{(i)}\) denote the Pauli matrix \(\sigma_a\) on
factor \(i\), and put
\begin{equation}
 S_{ij}:=\sum_{a=1}^3\sigma_a^{(i)}\sigma_a^{(j)}.
 \label{eq:newV26-Sij}
\end{equation}
Define the completely symmetric invariant
\begin{equation}
 {\cal T}_4
 :=
 S_{12}S_{34}+S_{13}S_{24}+S_{14}S_{23}.
 \label{eq:newV26-T4}
\end{equation}

\begin{theorem}[Operator tangent expansion of the path score]
\label{thm:newV26-operator-expansion}
Let
\begin{equation}
 {\cal K}_{\alpha,u}
 :=
 \kappa_5^{\widetilde\nu_{\alpha,u}}
 (U,U,U,U,\widetilde S_{\alpha,u})
 \in\operatorname{End}({\cal H}).
 \label{eq:newV26-K}
\end{equation}
Then, uniformly in \(u\in[0,1]\),
\begin{equation}
 \boxed{\qquad
 {\cal K}_{\alpha,u}
 ={1\over3\alpha^3}{\cal T}_4
 +O_{\rm op}(\alpha^{-4}).
 \qquad}
 \label{eq:newV26-K-expansion}
\end{equation}
\end{theorem}

\begin{proof}
In the scaled chart,
\[
 U
 =I+{i\over\sqrt\alpha}\sum_{a=1}^3Y_a\sigma_a
 +O_{L^p}(\alpha^{-1})
\]
for every fixed \(p\).  A cumulant vanishes if one of its arguments
is constant.  Thus the first nonzero contribution with all four
input labels visible takes the linear term from every copy of
\(U\).  Terms of total odd Gaussian degree vanish, and the first
remaining correction has two more half-powers of \(\alpha^{-1/2}\).
By \eqref{eq:newV25-score-expansion},
\[
 \widetilde S_{\alpha,u}
 ={|Y|^4-15\over24\alpha}
 +O_{L^p}(\alpha^{-2}).
\]
The uniform moment comparison of
\Cref{lem:newV25-path-expansion} therefore gives
\begin{align*}
 {\cal K}_{\alpha,u}
 &=
 {1\over24\alpha^3}
 \sum_{a,b,c,d=1}^3
 \kappa_5^{\gamma_3}
 (Y_a,Y_b,Y_c,Y_d,|Y|^4-15)
 \sigma_a^{(1)}\sigma_b^{(2)}
 \sigma_c^{(3)}\sigma_d^{(4)}
 +O_{\rm op}(\alpha^{-4}).
\end{align*}
The fourth derivative at zero of \(|h|^4\) is
\[
 8\left(
 \delta_{ab}\delta_{cd}
 +\delta_{ac}\delta_{bd}
 +\delta_{ad}\delta_{bc}
 \right).
\]
This is the displayed Gaussian cumulant, by the same exponential
tilt used in \Cref{lem:newV25-radial-cumulant}.  Since
\(8/24=1/3\), the three contractions are precisely
\eqref{eq:newV26-T4}.  Finite-dimensional entrywise remainder
bounds imply the stated operator-norm remainder.
\end{proof}

\section{Exact total-spin diagonalization}
\label{sec:newV26-diagonalization}

Let \(\Pi_{ij}\) swap tensor factors \(i,j\).  The two-spin identity
\begin{equation}
 S_{ij}=2\Pi_{ij}-I
 \label{eq:newV26-swap}
\end{equation}
turns \({\cal T}_4\) into a central element of the permutation
algebra.  Let \(P_j\) be the total-spin-\(j\) projector in
\[
 \left({1\over2}\right)^{\otimes4},
 \qquad j=0,1,2.
\]

\begin{theorem}[The quartic score survives nontrivial output]
\label{thm:newV26-sector-eigenvalues}
The exact spectral decomposition is
\begin{equation}
 \boxed{\qquad
 {\cal T}_4=15P_0-5P_1+3P_2.
 \qquad}
 \label{eq:newV26-T4-spectrum}
\end{equation}
Consequently
\begin{equation}
 {\cal K}_{\alpha,u}
 =
 {1\over\alpha^3}
 \left(
 5P_0-{5\over3}P_1+P_2
 \right)
 +O_{\rm op}(\alpha^{-4}).
 \label{eq:newV26-K-spectrum}
\end{equation}
In particular \(Q_0=I-P_0\) does not delete the leading tensor.
\end{theorem}

\begin{proof}
Let \(C_{(2)}\) be the sum of the six transpositions in
\(\mathbb C[S_4]\), and let \(C_{(2,2)}\) be the sum of the three
double transpositions.  Expanding \eqref{eq:newV26-swap} in the
three disjoint pairings gives
\begin{equation}
 {\cal T}_4
 =4C_{(2,2)}-2C_{(2)}+3I.
 \label{eq:newV26-class-sum}
\end{equation}

Schur--Weyl duality pairs total spins \(j=2,1,0\) with the
\(S_4\) partitions \([4],[3,1],[2,2]\), respectively.  On an
irreducible \(S_4\) module \(\lambda\), a class sum \(C_{\cal C}\)
acts by
\[
 {|{\cal C}|\chi_\lambda({\cal C})\over d_\lambda}.
\]
The required character rows are
\[
\begin{array}{c|ccc}
\lambda&d_\lambda&\chi_\lambda((2))&
                    \chi_\lambda((2,2))\\
\hline
[4]&1&1&1\\
[3,1]&3&1&-1\\
[2,2]&2&0&2.
\end{array}
\]
Thus \(C_{(2)}\) has eigenvalues \(6,2,0\), while
\(C_{(2,2)}\) has eigenvalues \(3,-1,3\).  Substitution into
\eqref{eq:newV26-class-sum} gives
\[
 3,\quad-5,\quad15
\]
on \(j=2,1,0\).  This proves
\eqref{eq:newV26-T4-spectrum}; combine with
\eqref{eq:newV26-K-expansion} for
\eqref{eq:newV26-K-spectrum}.
\end{proof}

\begin{corollary}[The integrated Wilson--geodesic defect survives]
\label{cor:newV26-integrated-defect}
Let
\[
 {\cal E}_{\alpha,4}
 :=
 \kappa_4^W(U,U,U,U)
 -\kappa_4^G(U,U,U,U).
\]
Then
\begin{equation}
 {\cal E}_{\alpha,4}
 =
 {1\over\alpha^3}
 \left(
 5P_0-{5\over3}P_1+P_2
 \right)
 +O_{\rm op}(\alpha^{-4}).
 \label{eq:newV26-integrated-defect}
\end{equation}
\end{corollary}

\begin{proof}
Integrate the exact score identity
\eqref{eq:newV21-WG-comparison} in \(u\) and use the uniform
expansion \eqref{eq:newV26-K-spectrum}.
\end{proof}

\section{Action of the inherited one-coordinate payer}
\label{sec:newV26-payer}

The frozen archives define the nontrivial output payer on a
Peter--Weyl sector \(j\) by
\begin{equation}
 w_{\alpha,j}
 :={\Xi_j\over1-q_{\alpha,j}},
 \qquad j>0,
 \label{eq:newV26-payer}
\end{equation}
where \(\Xi_j>0\) is the fixed output Casimir mass and
\(q_{\alpha,j}\) is the central Wilson Fourier multiplier.  At
fixed \(j\),
\begin{equation}
 1-q_{\alpha,j}
 =s_\alpha C_2(j)+O(s_\alpha^2),
 \qquad
 s_\alpha\asymp\alpha^{-1}.
 \label{eq:newV26-fixed-gap}
\end{equation}

\begin{lemma}[The nontrivial fixed-spin payer is genuinely large]
\label{lem:newV26-payer-sharp}
For \(j=1,2\),
\begin{equation}
 c_j\alpha\le w_{\alpha,j}\le C_j\alpha
 \label{eq:newV26-payer-linear}
\end{equation}
for all sufficiently large \(\alpha\).
\end{lemma}

\begin{proof}
Insert \eqref{eq:newV26-fixed-gap} into
\eqref{eq:newV26-payer}.  Both \(C_2(j)\) and \(\Xi_j\) are fixed
strictly positive numbers, while
\(s_\alpha\asymp\alpha^{-1}\).
\end{proof}

\begin{theorem}[The standalone paid score card is false]
\label{thm:newV26-paid-score-false}
Under the inherited payer convention
\eqref{eq:newV26-payer},
\begin{equation}
 \left\|
 Q_0{\Xi\over1-q_\alpha}
 {\cal K}_{\alpha,u}
 \right\|_{\rm op}
 \ge c\alpha^{-2}
 \label{eq:newV26-paid-lower}
\end{equation}
uniformly in \(u\) for all sufficiently large \(\alpha\).
The four-input Pr\"ufer target is
\[
 b_{\rm f}^4(4b_{\rm f})^2
 \asymp\alpha^{-3}.
\]
Therefore the PMSIC \eqref{eq:newV25-PMSIC} is false if
\(\mathcal P_{\rm out}\) is the inherited final payer, and the
V22 GCSIC is false under its stated convention that this payer is
included.
\end{theorem}

\begin{proof}
On \(P_2\), the leading eigenvalue of
\({\cal K}_{\alpha,u}\) is
\(\alpha^{-3}+O(\alpha^{-4})\).  On \(P_1\), it is
\(-5/(3\alpha^3)+O(\alpha^{-4})\).  Neither sector is deleted by
\(Q_0\).  Multiply either one by
\eqref{eq:newV26-payer-linear} to obtain
\eqref{eq:newV26-paid-lower}.  The target calculation uses
\(b_{\rm f}\asymp\alpha^{-1/2}\).
\end{proof}

\begin{assessment}[The payer ambiguity is resolved]
\label{ass:newV26-payer-resolution}
V25 conservatively left open whether output projection cancels the
quartic tensor.  \Cref{thm:newV26-sector-eigenvalues} shows that it
does not: both spin one and spin two survive.  Hence the
\(O(\alpha)\) payer loss is not an artefact of taking its norm.
It is an exact eigenvalue on the first saturating score sector.
\end{assessment}

\section{Append-only normalization correction}
\label{sec:newV26-correction}

The intrinsic strong Wilson forest card
\eqref{eq:newV13-ISWFC} is a statement about the raw one-link
cumulant.  It contains the Wilson law and any internal score or
resolvent sums, but it does not apply the later
first-connectivity output mass and reduced resolvent.  The
score-comparison reduction must use the same convention.

\begin{definition}[Raw geodesic score card]
\label{def:newV26-RGSC}
The RGSC is exactly the unpaid card
\eqref{eq:newV25-UQSAC}:
\[
 \left\|
 \kappa_{m+1}^{\widetilde\nu_{\alpha,u}}
 (D^{R_1},\ldots,D^{R_m},\widetilde S_{\alpha,u})
 \right\|_{\rm op}
 \le
 K^m
 \left(\prod_i b_i\right)
 \left(\sum_i b_i\right)^{m-2}.
\]
\end{definition}

\begin{theorem}[Correct raw comparison theorem]
\label{thm:newV26-raw-comparison}
If GIAC \eqref{eq:newV22-GIAC} and RGSC hold uniformly, then the
raw Wilson aggregate obeys
\begin{equation}
 \left\|
 \kappa_m^W(D^{R_1},\ldots,D^{R_m})
 \right\|_{\rm op}
 \le
 K_1^m
 \left(\prod_i b_i\right)
 \left(\sum_i b_i\right)^{m-2}.
 \label{eq:newV26-raw-Wilson}
\end{equation}
Consequently the balancing equivalence of V16 yields the intrinsic
strong Wilson forest card.  No final first-connectivity payer is
used in this local implication.
\end{theorem}

\begin{proof}
Use the exact complete-cumulant identity
\[
 \kappa_m^W
 =
 \kappa_m^G
 +\int_0^1
 \kappa_{m+1}^{\widetilde\nu_{\alpha,u}}
 (D^{R_1},\ldots,D^{R_m},\widetilde S_{\alpha,u})\,du.
\]
Apply GIAC to the first complete object and RGSC to the second,
then take the triangle inequality between those two objects only.
The \(u\)-interval has length one.  Enlarge the exponential base
and apply the exact Pr\"ufer-weight balancing theorem
\Cref{thm:newV16-balancing-projection}.
\end{proof}

\begin{firewall}[A local card is a numerator, not a paid junction]
\label{fire:newV26-local-not-paid}
The later global first-connectivity construction may apply one
output mass and one reduced resolvent only after a joining
covariance or another numerator supplies its cancelling
\(s_\alpha\) clock.  Applying that payer directly to the raw local
card changes its scale and is forbidden by
\Cref{thm:newV26-paid-score-false}.  Thus the payer references in
V22 GCSIC, V24 SWPC, and V25 PMSIC are retired as local sufficient
conditions.  Their raw algebraic score identities remain valid.
\end{firewall}

\section{The complete lowest graph layer}
\label{sec:newV26-lowest-layer}

The scalar tangent graph has a useful all-arity result beyond power
counting.  Let \(B=\sum_{i=1}^m b_i\), and assume the inherited
thermal floor
\begin{equation}
 b_i\ge {c_*\over\sqrt\alpha}.
 \label{eq:newV26-mark-floor}
\end{equation}

\begin{lemma}[Weighted root-degree-four tree census]
\label{lem:newV26-root-tree-census}
For \(m\ge4\),
\begin{equation}
 \sum_{\substack{T\in\mathfrak T_{\{0,\ldots,m\}}\\
                  d_0(T)=4}}
 \prod_{i=1}^m b_i^{d_i(T)}
 =
 {m-1\choose3}
 \left(\prod_{i=1}^m b_i\right)B^{m-4}.
 \label{eq:newV26-root-tree-census}
\end{equation}
\end{lemma}

\begin{proof}
The weighted Cayley--Pr\"ufer identity on \(m+1\) vertices is
\[
 \sum_T\prod_{i=0}^m x_i^{d_i(T)}
 =
 \left(\prod_{i=0}^m x_i\right)
 \left(\sum_{i=0}^m x_i\right)^{m-1}.
\]
Set \(x_i=b_i\) for \(i\ge1\) and extract the coefficient of
\(x_0^4\).  The exterior \(x_0\) leaves the coefficient of
\(x_0^3\) in \((x_0+B)^{m-1}\), which is the right side of
\eqref{eq:newV26-root-tree-census}.
\end{proof}

\begin{theorem}[The first saturating score layer has exponential
Pr\"ufer grade]
\label{thm:newV26-lowest-layer}
For \(m\ge4\), take \(|t_i|\le C_0b_i\) in the scalar Gaussian
formula \eqref{eq:newV25-exact-graph}, and multiply the quartic
score by \(1/(24\alpha)\).  The complete homogeneous layer of
least degree \(2m-4\) obeys
\begin{equation}
 \left|\mathcal L_m(\boldsymbol t)\right|
 \le
 K^m
 \left(\prod_{i=1}^m b_i\right)B^{m-2}
 \label{eq:newV26-lowest-layer}
\end{equation}
with \(K\) independent of \(m,\alpha\), and the marks.
\end{theorem}

\begin{proof}
By \Cref{lem:newV25-clock-degree}, the layer of degree
\(2m-4\) has \(m-4\) input--input edges.  Together with the four
root edges it has \(m\) edges on \(m+1\) vertices.  Connectivity
therefore makes it a tree, and the root has degree four.  In
\eqref{eq:newV25-exact-graph} every edge multiplicity is one, so
the root allocation contributes \(4!=24\), which cancels the
score normalization.  Hence
\[
 |\mathcal L_m(\boldsymbol t)|
 \le
 {C_0^{2m-4}\over\alpha}
 \sum_{\substack{T\in\mathfrak T_{\{0,\ldots,m\}}\\d_0(T)=4}}
 \prod_{i=1}^m b_i^{d_i(T)}.
\]
Apply \Cref{lem:newV26-root-tree-census}.  From
\eqref{eq:newV26-mark-floor},
\[
 B\ge {c_*m\over\sqrt\alpha},
 \qquad
 {1\over\alpha}\le {B^2\over c_*^2m^2}.
\]
Therefore the last display is at most
\[
 {C_0^{2m-4}\over c_*^2}
 {{m-1\choose3}\over m^2}
 \left(\prod_i b_i\right)B^{m-2}.
\]
The remaining factor is \(O(m)\), hence is bounded by an
exponential \(K^m\).
\end{proof}

\begin{assessment}[What remains after the exact lowest layer]
\label{ass:newV26-remaining}
The first graph layer is not the all-order obstruction: its
coefficient census is exactly one root-degree-four Pr\"ufer
coefficient and has exponential grade.  The unresolved scalar
problem is the cancellation-preserving resummation of the higher
input--input multigraph layers.  Taking their absolute exponential
series gives a factor of order
\(\exp(\sum_{i<j}b_ib_j)\), which is not bounded by \(K^m\) when
many marks are order one.  The signs and the unit-modulus Weyl
vertices must therefore be retained before the norm.
\end{assessment}

\section{V27 programme and claim boundary}
\label{sec:newV26-next}

\begin{programme}[V27: signed cycle resummation]
\label{prog:newV26-V27}
\begin{enumerate}[leftmargin=3em]
\item Start from the exact graph formula
      \eqref{eq:newV25-exact-graph} and sum all parallel
      input--input edges before absolute values.
\item Express the result through a connected tree formula with
      edge functions \(e^{-t_it_j}-1\), retaining their signs.
\item Prove or sharply test an exponential-grade tree-graph bound
      under the thermal floor
      \eqref{eq:newV26-mark-floor}.
\item Test equal order-one marks and alternating signs; these are
      precisely where the absolute multigraph exponential is too
      large.
\item Lift only a proved scalar cancellation to matrix
      representation rows, using complete row--column contraction
      before operator norm.
\item Keep the local comparison raw.  Do not attach the global
      final payer to GIAC or RGSC.
\end{enumerate}
\end{programme}

\begin{firewall}[Claim boundary after V26]
\label{fire:newV26-boundary}
V26 computes the complete leading four-fundamental score tensor,
diagonalizes it on total-spin sectors, proves that nontrivial
spin-one and spin-two outputs survive, and thereby disproves the
standalone paid GCSIC/PMSIC normalization.  It restores the correct
raw comparison theorem and proves an all-arity exponential
Pr\"ufer bound for the first saturating quartic-root graph layer.
It does not resum the higher graph layers or prove RGSC, GIAC, the
Wilson aggregate strong card, all-order strong MTA, WUFC, CPM,
cutoff removal, reflection positivity,
Osterwalder--Schrader reconstruction, a continuum Yang--Mills
measure, or a positive Hamiltonian gap.  The full Yang--Mills mass
gap has not been proved.
\end{firewall}

\begin{theorem}[V26 result ledger]
\label{thm:newV26-results}
V26:
\begin{enumerate}[label=\textup{(V26.\arabic*)},leftmargin=6.5em]
\item computes the operator-valued tangent score coefficient
      \({\cal T}_4/(3\alpha^3)\);
\item proves
      \({\cal T}_4=15P_0-5P_1+3P_2\);
\item proves the integrated defect survives \(Q_0\);
\item proves the inherited fixed-spin payer grows like \(\alpha\)
      and makes the standalone paid score card false;
\item corrects the comparison theorem to the raw ISWFC
      normalization;
\item evaluates the exact root-degree-four weighted tree census;
      and
\item proves the complete least-degree score layer has the desired
      exponential Pr\"ufer grade.
\end{enumerate}
\end{theorem}

\begin{proof}
Use
\Cref{thm:newV26-operator-expansion,
thm:newV26-sector-eigenvalues,
cor:newV26-integrated-defect,
lem:newV26-payer-sharp,
thm:newV26-paid-score-false,
thm:newV26-raw-comparison,
lem:newV26-root-tree-census,
thm:newV26-lowest-layer}.
\end{proof}

\paragraph{Parent-version record.}
V26 was formed from
\texttt{Renewal\_Yang\_Mills\_Mass\_Gap\_Research\_Notes\_V25.tex}.
The V25 parent had \(11\,730\) logical source lines,
\(376\,620\) bytes, and SHA--256
\[
 \texttt{048b37af3b11bb3f0b8249423f34f5b1724309e0a807c213bf44f65528521ce8}.
\]
The next compulsory companion audit is V30.

% =============================================================================
% END APPENDED FIFTH-SERIES V26 MATERIAL
% =============================================================================

% =============================================================================
% BEGIN APPENDED FIFTH-SERIES V27 MATERIAL
% =============================================================================

\chapter{V27: Signed cycle resummation in the scalar tangent model}
\label{chap:newV27-resummation}

\section{Parallel edges must be summed first}
\label{sec:newV27-simple-graphs}

The absolute multigraph estimate rejected at the end of V26 treats
every parallel contraction separately.  Gaussian normal ordering
allows those contractions to be summed exactly.  Put
\begin{equation}
 f_{ij}:=e^{-t_it_j}-1,
 \qquad 1\le i<j\le m.
 \label{eq:newV27-Mayer-edge}
\end{equation}
For a root allocation
\[
 \boldsymbol k=(k_1,\ldots,k_m)\in\mathbb N_0^m,
 \qquad
 \sum_i k_i=4,
\]
write
\[
 A(\boldsymbol k):=\{i:k_i>0\}.
\]
Let \(\mathfrak G_A\) be the family of simple graphs on
\([m]\) for which every connected component contains at least one
vertex of \(A\).

\begin{theorem}[Exact signed simple-graph resummation]
\label{thm:newV27-simple-resummation}
For all real \(t_1,\ldots,t_m\),
\begin{align}
 &\kappa_{m+1}
 \left(
 e^{it_1X},\ldots,e^{it_mX},H_4(X)
 \right)
 \notag\\
 &\quad=
 e^{-\frac12\sum_i t_i^2}
 \sum_{\substack{\boldsymbol k\in\mathbb N_0^m\\
                  \sum_i k_i=4}}
 {4!\over\prod_i k_i!}
 \prod_i(it_i)^{k_i}
 \sum_{G\in\mathfrak G_{A(\boldsymbol k)}}
 \prod_{ij\in E(G)}f_{ij}.
 \label{eq:newV27-simple-resummation}
\end{align}
\end{theorem}

\begin{proof}
Start from \eqref{eq:newV25-exact-graph} and fix the four root-edge
multiplicities \(n_{0i}=k_i\).  For a fixed input pair \(ij\),
summing a positive number of parallel edges gives
\[
 \sum_{n=1}^{\infty}{(-t_it_j)^n\over n!}
 =e^{-t_it_j}-1=f_{ij}.
\]
Pairs carrying no edge contribute one.  Connectivity of the full
multigraph on \(\{0,\ldots,m\}\) depends only on the support of
the input edges: after the root and its incident edges are removed,
every input component must meet
\(A(\boldsymbol k)\).  This is exactly
\(G\in\mathfrak G_{A(\boldsymbol k)}\), proving the formula.
\end{proof}

\begin{assessment}[Where the cancellation is stored]
\label{ass:newV27-cancellation-location}
The dangerous factor
\(\exp(\sum_{i<j}|t_it_j|)\) arises only if the Taylor series of
every \(f_{ij}\) is made positive before graph assembly.
\Cref{thm:newV27-simple-resummation} retains instead the signed
Mayer edge \(e^{-t_it_j}-1\).  Its pair potential
\(t_it_j\) is not nonnegative, but it is uniformly stable; this is
the exact input needed by the tree-graph inequality.
\end{assessment}

\section{A rooted stable tree-graph inequality}
\label{sec:newV27-rooted-Penrose}

\begin{definition}[Stability]
\label{def:newV27-stability}
A real pair array \(v_{ij}\) on \([m]\) is \(B\)-stable if, for
every \(I\subseteq[m]\),
\begin{equation}
 \sum_{\substack{i<j\\i,j\in I}}v_{ij}
 \ge-B|I|.
 \label{eq:newV27-stability}
\end{equation}
\end{definition}

For \(\varnothing\ne R\subseteq[m]\), let
\(\mathfrak F_R\) denote the forests on \([m]\) in which every
component contains exactly one vertex of \(R\).  Equivalently, each
component is oriented toward its unique root in \(R\).

\begin{lemma}[Rooted Penrose bound]
\label{lem:newV27-rooted-Penrose}
Let \(v\) be \(B\)-stable and put \(f_{ij}=e^{-v_{ij}}-1\).
For every nonempty \(A\subseteq[m]\),
\begin{align}
 \left|
 \sum_{G\in\mathfrak G_A}\prod_{ij\in E(G)}f_{ij}
 \right|
 &\le
 e^{2Bm}
 \sum_{\varnothing\ne R\subseteq A}
 \ \sum_{F\in\mathfrak F_R}
 \prod_{ij\in E(F)}
 \left(1-e^{-|v_{ij}|}\right).
 \label{eq:newV27-rooted-Penrose}
\end{align}
\end{lemma}

\begin{proof}
Adjoin a ghost vertex \(0\) and the fixed zero-cost edges
\(\{0a:a\in A\}\).  An input graph \(G\) belongs to
\(\mathfrak G_A\) exactly when this augmented graph is connected.
Choose a total edge order with the ghost edges first and apply
Kruskal's deterministic partition scheme to the connected augmented
graphs.  Removing the ghost from the selected spanning tree leaves
a forest.  The ghost edges retained by the tree select a nonempty
set \(R\subseteq A\), and every remaining component contains
exactly one member of \(R\).  Thus the tree part lies in
\(\mathfrak F_R\).

For a fixed selected forest, summation of all graph edges admitted
by its Penrose interval produces exponential factors
\(e^{-\sum v_{ij}}\).  The usual deletion sequence exposes at most
two induced input subgraphs at each vertex.  Applying
\eqref{eq:newV27-stability} to those disjoint induced sets bounds
the product of all adverse exponential factors by \(e^{2Bm}\).
On a selected forest edge,
\[
 |e^{-v}-1|e^{-\max\{-v,0\}}
 =1-e^{-|v|}.
\]
Multiplying these bounds and summing the disjoint Kruskal intervals
gives \eqref{eq:newV27-rooted-Penrose}.  The fixed ghost edges have
unit weight and introduce no stability term.
\end{proof}

\begin{remark}[Scope of the rooted form]
\label{rem:newV27-rooted-scope}
For \(|A|=1\), \Cref{lem:newV27-rooted-Penrose} is the ordinary
connected Penrose tree bound.  For \(A=[m]\), the left side is the
complete graph partition product \(e^{-\sum_{i<j}v_{ij}}\); the
stability factor controls it even though the empty rooted forest is
present.  Thus the two extreme cases are included without a hidden
connectivity assumption.
\end{remark}

\section{Stability of the Gaussian rank-one potential}
\label{sec:newV27-stability}

\begin{lemma}[Uniform stability]
\label{lem:newV27-rank-one-stability}
If \(|t_i|\le C_0\), then the pair potential
\[
 v_{ij}=t_it_j
\]
is \(B_0\)-stable with \(B_0=C_0^2/2\).
\end{lemma}

\begin{proof}
For every \(I\subseteq[m]\),
\begin{align*}
 \sum_{\substack{i<j\\i,j\in I}}t_it_j
 &={1\over2}
 \left[
 \left(\sum_{i\in I}t_i\right)^2
 -\sum_{i\in I}t_i^2
 \right]\\
 &\ge-{1\over2}\sum_{i\in I}t_i^2
 \ge-{C_0^2\over2}|I|.
\end{align*}
\end{proof}

\begin{lemma}[Rooted forest majorant]
\label{lem:newV27-forest-majorant}
Assume \(|t_i|\le C_0b_i\) and \(0<b_i<1\).
For nonempty \(R\subseteq[m]\),
\begin{equation}
 \sum_{F\in\mathfrak F_R}
 \prod_{ij\in E(F)}
 \left(1-e^{-|t_it_j|}\right)
 \le
 C_1^{m-|R|}
 \left(\prod_{i\notin R}b_i\right)
 B^{m-|R|},
 \qquad B=\sum_i b_i.
 \label{eq:newV27-forest-majorant}
\end{equation}
\end{lemma}

\begin{proof}
Use
\[
 1-e^{-|t_it_j|}
 \le |t_it_j|
 \le C_0^2b_ib_j.
\]
Orient each forest edge toward the unique root of its component.
Every \(i\notin R\) then chooses one parent \(p(i)\ne i\), and its
edge contributes \(C_0^2b_ib_{p(i)}\).  Drop the acyclicity
condition and all restrictions on the parent choices.  The enlarged
sum factorizes:
\[
 \prod_{i\notin R}
 \left(
 C_0^2b_i\sum_{j\ne i}b_j
 \right)
 \le
 C_0^{2(m-|R|)}
 \left(\prod_{i\notin R}b_i\right)
 B^{m-|R|}.
\]
\end{proof}

\section{The scalar all-order raw score card}
\label{sec:newV27-scalar-card}

\begin{theorem}[Scalar Gaussian RGSC]
\label{thm:newV27-scalar-RGSC}
Let \(m\ge2\), \(X\sim N(0,1)\), and
\[
 S_\alpha^{(4)}:={H_4(X)\over24\alpha}.
\]
Assume
\begin{equation}
 {c_*\over\sqrt\alpha}\le b_i<1,
 \qquad
 |t_i|\le C_0b_i.
 \label{eq:newV27-scalar-hypotheses}
\end{equation}
Then
\begin{equation}
 \boxed{\qquad
 \left|
 \kappa_{m+1}
 \left(
 e^{it_1X},\ldots,e^{it_mX},S_\alpha^{(4)}
 \right)
 \right|
 \le
 K^m
 \left(\prod_{i=1}^m b_i\right)
 \left(\sum_{i=1}^m b_i\right)^{m-2},
 \qquad}
 \label{eq:newV27-scalar-RGSC}
\end{equation}
where \(K\) depends only on \(c_*,C_0\).
\end{theorem}

\begin{proof}
Use the exact resummation
\eqref{eq:newV27-simple-resummation}; its exterior Gaussian factor
has modulus at most one.  Fix a root allocation
\(\boldsymbol k\), put \(A=A(\boldsymbol k)\), and use
\Cref{lem:newV27-rooted-Penrose,
lem:newV27-rank-one-stability,
lem:newV27-forest-majorant}.
For every nonempty \(R\subseteq A\), let \(r=|R|\).  Apart from an
exponential \(K_0^m\), its contribution after the score
normalization is at most
\begin{align}
 {1\over\alpha}
 \left(\prod_{a\in A}b_a^{k_a}\right)
 \left(\prod_{i\notin R}b_i\right)
 B^{m-r}.
 \label{eq:newV27-one-root-set}
\end{align}
Since \(R\subseteq A\) and \(\sum_{a\in A}k_a=4\), this is bounded
by
\begin{align}
 {1\over\alpha}
 \left(\prod_{i=1}^m b_i\right)
 B^{4-r}B^{m-r}
 &=
 \left(\prod_i b_i\right)B^{m-2}
 \left[
 {1\over\alpha}B^{6-2r}
 \right].
 \label{eq:newV27-ratio}
\end{align}
There are only four cases:
\begin{align*}
 r=1:&\quad \alpha^{-1}B^4\le m^4,\\
 r=2:&\quad \alpha^{-1}B^2\le m^2,\\
 r=3:&\quad \alpha^{-1}\le1,\\
 r=4:&\quad
 \alpha^{-1}B^{-2}
 \le(c_*^2m^2)^{-1}.
\end{align*}
The first three use \(B<m\) and \(\alpha\ge1\); the last uses
\(B\ge c_*m/\sqrt\alpha\).

The number of root allocations is
\({m+3\choose3}\), and each has at most \(2^4-1\) nonempty
root subsets.  All displayed polynomial factors in \(m\), as well
as the fixed multinomial coefficient
\(4!/\prod k_i!\), are absorbed into a larger exponential \(K^m\).
This proves \eqref{eq:newV27-scalar-RGSC}.
\end{proof}

\begin{corollary}[All higher scalar cycles are resummed]
\label{cor:newV27-cycles-closed}
The obstruction identified in
\Cref{ass:newV26-remaining} is closed in the scalar tangent Weyl
model.  No absolute factor
\(\exp(\sum_{i<j}|t_it_j|)\) occurs; stability costs only an
exponential per label.
\end{corollary}

\begin{proof}
\Cref{thm:newV27-scalar-RGSC} includes every input--input graph
layer in \eqref{eq:newV25-exact-graph}, not only the least-degree
tree layer treated in V26.
\end{proof}

\section{Why this is not yet the compact-group card}
\label{sec:newV27-matrix-boundary}

\begin{firewall}[Scalar stability does not imply matrix stability]
\label{fire:newV27-no-matrix-transfer}
For scalar Weyl inputs every pair contraction is the rank-one real
number \(t_it_j\), and
\Cref{lem:newV27-rank-one-stability} controls the negative part of
their total sum.  Matrix representation derivatives produce
noncommuting row--column contractions and several Lie algebra
coordinates.  Replacing them by their operator norms before the
Penrose sum destroys the signed stability identity.  Therefore
\Cref{thm:newV27-scalar-RGSC} is a proved tangent model, not RGSC
for arbitrary \(SU(2)\) representations.
\end{firewall}

\begin{assessment}[The precise matrix lift]
\label{ass:newV27-matrix-lift}
The next analytic object is a noncommutative stable Mayer row.  One
must normal-order the complete matrix exponential vertices, sum
parallel Gaussian contractions into a time-ordered pair operator,
and prove that its negative quadratic form is bounded by a constant
times the number of labels before taking the final operator norm.
Only after that stability statement is proved may the rooted
Penrose forest census of V27 be reused.
\end{assessment}

\section{V28 programme and claim boundary}
\label{sec:newV27-next}

\begin{programme}[V28: commuting Cartan and matrix stability tests]
\label{prog:newV27-V28}
\begin{enumerate}[leftmargin=3em]
\item Lift \Cref{thm:newV27-scalar-RGSC} to simultaneously
      diagonalizable Cartan representation inputs, retaining all
      weight signs.
\item Write the exact pair-contraction operator for two
      noncommuting \(SU(2)\) matrix exponentials and compare it with
      the scalar Mayer edge.
\item Test stability first on four fundamental inputs, where the
      total-spin diagonalization of V26 is explicit.
\item Formulate the required completely bounded stability
      inequality on a tensor product of arbitrary representation
      rows.
\item If noncommutativity violates scalar stability, retain the
      ordered Gaussian covariance interpolation rather than taking
      a pairwise exponential.
\item Keep the comparison raw and preserve the exact root score
      through the complete matrix contraction.
\end{enumerate}
\end{programme}

\begin{firewall}[Claim boundary after V27]
\label{fire:newV27-boundary}
V27 exactly resums all parallel Gaussian edges into signed Mayer
edges, proves a rooted stable Penrose inequality, verifies stability
of the scalar rank-one Gaussian potential, and proves the
factorial-free all-order raw quartic-score card for scalar Weyl
inputs.  It does not prove noncommutative matrix stability, RGSC
for arbitrary \(SU(2)\) representations, GIAC, the Wilson
aggregate strong card, all-order strong MTA, WUFC, CPM, cutoff
removal, reflection positivity, Osterwalder--Schrader
reconstruction, a continuum Yang--Mills measure, or a positive
Hamiltonian gap.  The full Yang--Mills mass gap has not been
proved.
\end{firewall}

\begin{theorem}[V27 result ledger]
\label{thm:newV27-results}
V27:
\begin{enumerate}[label=\textup{(V27.\arabic*)},leftmargin=6.5em]
\item derives the exact signed simple-graph resummation;
\item proves the rooted stable Penrose bound;
\item proves uniform stability of \(v_{ij}=t_it_j\);
\item bounds the full rooted forest family without a Cayley
      factorial;
\item proves the scalar Gaussian RGSC at every arity with an
      exponential base; and
\item isolates noncommutative matrix stability as the next gate.
\end{enumerate}
\end{theorem}

\begin{proof}
Use
\Cref{thm:newV27-simple-resummation,
lem:newV27-rooted-Penrose,
lem:newV27-rank-one-stability,
lem:newV27-forest-majorant,
thm:newV27-scalar-RGSC,
cor:newV27-cycles-closed}.
\end{proof}

\paragraph{Parent-version record.}
V27 was formed from
\texttt{Renewal\_Yang\_Mills\_Mass\_Gap\_Research\_Notes\_V26.tex}.
The V26 parent had \(12\,295\) logical source lines,
\(393\,873\) bytes, and SHA--256
\[
 \texttt{773402da87afd686b464c64069b23b9a321bffd7aacc089570b69d5f7cc1d79e}.
\]
The next compulsory companion audit is V30.

% =============================================================================
% END APPENDED FIFTH-SERIES V27 MATERIAL
% =============================================================================

% =============================================================================
% BEGIN APPENDED FIFTH-SERIES V28 MATERIAL
% =============================================================================

\chapter{V28: Cartan closure and noncommutative Casimir stability}
\label{chap:newV28-stability}

\section{The radial quartic score seen from one Cartan coordinate}
\label{sec:newV28-conditional-score}

Let \(Y=(Y_1,Y_2,Y_3)\sim\gamma_3\) and put \(X=Y_3\).
The leading radial score polynomial from V25 is
\[
 Q_4(Y)=|Y|^4-15.
\]

\begin{lemma}[Exact Cartan conditional score]
\label{lem:newV28-conditional-score}
\begin{equation}
 \mathbb E[Q_4(Y)\mid X]
 =H_4(X)+10H_2(X),
 \label{eq:newV28-conditional-score}
\end{equation}
where
\[
 H_2(X)=X^2-1,
 \qquad
 H_4(X)=X^4-6X^2+3.
\]
\end{lemma}

\begin{proof}
Write \(R_\perp^2=Y_1^2+Y_2^2\).  Then
\[
 |Y|^4=X^4+2X^2R_\perp^2+R_\perp^4.
\]
For a two-dimensional standard Gaussian,
\(\mathbb ER_\perp^2=2\) and
\(\mathbb ER_\perp^4=8\).  Therefore
\[
 \mathbb E[Q_4\mid X]
 =X^4+4X^2-7
 =H_4(X)+10H_2(X).
\]
\end{proof}

\section{The degree-two companion graph}
\label{sec:newV28-H2}

For \(d\in\{2,4\}\), let \(\mathfrak C_m^{(d)}\) be the connected
multigraph family of V25 with distinguished root degree \(d\).
The normal-ordering proof of
\Cref{thm:newV25-exact-graph} applies verbatim to \(H_d\).

\begin{proposition}[Scalar \(H_2\) raw score card]
\label{prop:newV28-H2-card}
Under \eqref{eq:newV27-scalar-hypotheses},
\begin{equation}
 \left|
 \kappa_{m+1}
 \left(
 e^{it_1X},\ldots,e^{it_mX},{H_2(X)\over\alpha}
 \right)
 \right|
 \le
 K^m
 \left(\prod_i b_i\right)B^{m-2},
 \qquad m\ge2.
 \label{eq:newV28-H2-card}
\end{equation}
\end{proposition}

\begin{proof}
Repeat the exact parallel-edge resummation of
\Cref{thm:newV27-simple-resummation} with root allocations
\(\sum_i k_i=2\).  Apply the rooted Penrose and forest bounds
\Cref{lem:newV27-rooted-Penrose,
lem:newV27-rank-one-stability,
lem:newV27-forest-majorant}.

For a selected nonempty \(R\subseteq A(\boldsymbol k)\), put
\(r=|R|\).  Here \(r\in\{1,2\}\).  The analogue of
\eqref{eq:newV27-ratio} is the target Pr\"ufer weight times
\[
 {1\over\alpha}B^{4-2r}.
\]
For \(r=1\), this is at most \(m^2\), because \(B<m\) and
\(\alpha\ge1\).  For \(r=2\), it is at most one.  The number of
root allocations is \(O(m^2)\).  All polynomial factors are
absorbed into \(K^m\).
\end{proof}

\begin{corollary}[Conditional radial scalar card]
\label{cor:newV28-radial-Cartan-card}
Under the same hypotheses,
\begin{equation}
 \left|
 \kappa_{m+1}
 \left(
 e^{it_1X},\ldots,e^{it_mX},
 {H_4(X)+10H_2(X)\over24\alpha}
 \right)
 \right|
 \le
 K^m
 \left(\prod_i b_i\right)B^{m-2}.
 \label{eq:newV28-radial-Cartan-card}
\end{equation}
\end{corollary}

\begin{proof}
Combine
\Cref{thm:newV27-scalar-RGSC,prop:newV28-H2-card} and enlarge the
exponential base to absorb the fixed coefficient \(10\).
\end{proof}

\section{Simultaneously diagonalizable representation rows}
\label{sec:newV28-Cartan-lift}

For each label \(i\), let \(J_i^3\) be the Hermitian Cartan
generator of a finite-dimensional unitary \(SU(2)\)
representation.  Suppose
\[
 \|J_i^3\|\le c_0\sqrt{C_2(R_i)}
\]
and put
\begin{equation}
 F_i(X):=\exp\left(i{\sqrt{s_\alpha}}J_i^3X\right).
 \label{eq:newV28-Cartan-input}
\end{equation}
The matrices \(F_i\) act on distinct tensor factors.

\begin{theorem}[Complete Cartan matrix RGSC]
\label{thm:newV28-Cartan-RGSC}
Let
\[
 S_{\alpha}^{\rm Car}
 :={H_4(X)+10H_2(X)\over24\alpha}.
\]
Then
\begin{equation}
 \left\|
 \kappa_{m+1}
 (F_1,\ldots,F_m,S_{\alpha}^{\rm Car})
 \right\|_{\rm op}
 \le
 K^m
 \left(\prod_i b_i\right)B^{m-2},
 \qquad m\ge2,
 \label{eq:newV28-Cartan-RGSC}
\end{equation}
uniformly in the representations and their dimensions.
\end{theorem}

\begin{proof}
Choose an orthonormal weight basis for each \(J_i^3\).  On a joint
weight vector \(\boldsymbol\mu=(\mu_1,\ldots,\mu_m)\), the tensor
cumulant is the scalar cumulant in
\eqref{eq:newV28-radial-Cartan-card} with
\[
 t_i=\sqrt{s_\alpha}\,\mu_i.
\]
The Casimir row bound gives
\[
 |t_i|
 \le c_0\sqrt{s_\alpha C_2(R_i)}
 \le c_0 b_i.
\]
The inherited thermal floor supplies
\(b_i\ge c_*/\sqrt\alpha\).  Thus
\Cref{cor:newV28-radial-Cartan-card} bounds every diagonal entry by
the same right side, independently of \(\boldsymbol\mu\) and the
number of weights.  The operator norm of a diagonal matrix is the
supremum of its diagonal entries, proving the claim.
\end{proof}

\begin{assessment}[A genuine matrix sector is closed]
\label{ass:newV28-Cartan-closed}
\Cref{thm:newV28-Cartan-RGSC} is not merely a one-dimensional
scalar example.  It is dimension-free for arbitrary reducible
tensor products of representation weights, provided all Gaussian
Lie algebra directions are compressed to a common Cartan
coordinate.  The remaining obstruction is precisely the
noncommutation of different Lie algebra directions at a shared
representation vertex.
\end{assessment}

\section{The exact noncommutative stability identity}
\label{sec:newV28-operator-stability}

Let \(J_i^a\), \(a=1,2,3\), be Hermitian generators on the
\(i\)-th representation space, normalized by
\[
 C_i:=\sum_{a=1}^3(J_i^a)^2=C_2(R_i)I.
\]
On the full tensor product define
\begin{align}
 \Omega_{ij}
 &:=\sum_{a=1}^3J_i^aJ_j^a,
 \label{eq:newV28-Omega}\\
 C_I^{\rm tot}
 &:=\sum_{a=1}^3
 \left(\sum_{i\in I}J_i^a\right)^2
 \qquad(I\subseteq[m]).
 \label{eq:newV28-total-Casimir}
\end{align}

\begin{theorem}[Complete operator stability on every subset]
\label{thm:newV28-operator-stability}
For every \(I\subseteq[m]\),
\begin{equation}
 \boxed{\qquad
 2s_\alpha
 \sum_{\substack{i<j\\i,j\in I}}\Omega_{ij}
 =
 s_\alpha C_I^{\rm tot}
 -s_\alpha\sum_{i\in I}C_i
 \ge
 -\sum_{i\in I}b_i^2\,I
 \ge-|I|\,I.
 \qquad}
 \label{eq:newV28-operator-stability}
\end{equation}
Thus the full noncommutative pair family is operator-stable with a
constant independent of the number and dimensions of the
representations.
\end{theorem}

\begin{proof}
Expand the square in
\eqref{eq:newV28-total-Casimir}:
\[
 C_I^{\rm tot}
 =
 \sum_{i\in I}C_i
 +2\sum_{\substack{i<j\\i,j\in I}}\Omega_{ij}.
\]
The total Casimir is a sum of Hermitian squares and is therefore
positive semidefinite.  Since
\[
 s_\alpha C_i
 \le b_i^2
\]
by the definition
\(b_i^2=s_\alpha(1+C_2(R_i))\), the first lower bound follows.
Finally \(b_i<1\) gives the second.
\end{proof}

\begin{corollary}[Stable total heat interaction]
\label{cor:newV28-heat-stability}
For every \(I\subseteq[m]\),
\begin{equation}
 \left\|
 \exp\left(
 -2s_\alpha\sum_{\substack{i<j\\i,j\in I}}\Omega_{ij}
 \right)
 \right\|_{\rm op}
 \le e^{|I|}.
 \label{eq:newV28-heat-stability}
\end{equation}
\end{corollary}

\begin{proof}
The exponent is Hermitian.  Apply functional calculus to
\eqref{eq:newV28-operator-stability}.
\end{proof}

\begin{proposition}[Pair row size]
\label{prop:newV28-pair-row}
\begin{equation}
 \|2s_\alpha\Omega_{ij}\|_{\rm op}
 \le
 2b_ib_j.
 \label{eq:newV28-pair-row}
\end{equation}
\end{proposition}

\begin{proof}
The row--column Cauchy--Schwarz inequality gives
\[
 \|\Omega_{ij}\|
 \le
 \left\|\sum_a(J_i^a)^2\right\|^{1/2}
 \left\|\sum_a(J_j^a)^2\right\|^{1/2}
 =
 \sqrt{C_2(R_i)C_2(R_j)}.
\]
Multiply by \(2s_\alpha\) and use
\(s_\alpha C_2(R_i)\le b_i^2\).
\end{proof}

\section{The noncommutative ordering gate}
\label{sec:newV28-ordering}

\begin{firewall}[Operator stability is not yet a quantum Penrose
theorem]
\label{fire:newV28-no-quantum-Penrose}
The scalar proof of
\Cref{lem:newV27-rooted-Penrose} partitions a product of commuting
Mayer factors and then rearranges exponential weights inside each
graph interval.  The operators \(\Omega_{ij}\) fail to commute when
their edges share a vertex.  Although
\Cref{thm:newV28-operator-stability} controls the complete
exponential of their sum, it does not justify replacing a
time-ordered Duhamel word by an unordered graph product.  A quantum
or time-ordered tree-graph identity is still required.
\end{firewall}

\begin{lemma}[Disjoint pair interactions commute]
\label{lem:newV28-disjoint-commute}
If \(\{i,j\}\cap\{k,\ell\}=\varnothing\), then
\[
 [\Omega_{ij},\Omega_{k\ell}]=0.
\]
If the edges share a vertex, their commutator is generally nonzero
and contains one Lie bracket at that vertex.
\end{lemma}

\begin{proof}
Operators on disjoint tensor factors commute.  For example,
\[
 [\Omega_{ij},\Omega_{ik}]
 =
 \sum_{a,b}[J_i^a,J_i^b]J_j^aJ_k^b,
\]
which is nonzero in a nontrivial \(SU(2)\) representation.
\end{proof}

\begin{assessment}[The exact acquisition]
\label{ass:newV28-acquisition}
V28 acquires both scalar ingredients of a tree-graph proof in their
matrix analogues: the individual pair row
\eqref{eq:newV28-pair-row} and the subset-uniform stability
\eqref{eq:newV28-operator-stability}.  The missing step is no
longer an inequality about Casimir size.  It is an ordering theorem
which converts a complete time-ordered exponential into rooted
forests while applying the stability estimate to complete partial
sums, not to separately normed edges.
\end{assessment}

\section{V29 programme and claim boundary}
\label{sec:newV28-next}

\begin{programme}[V29: time-ordered quantum forest formula]
\label{prog:newV28-V29}
\begin{enumerate}[leftmargin=3em]
\item Use Duhamel's formula to differentiate
      \(\exp[-s_\alpha C_I^{\rm tot}]\) with respect to pair
      couplings \(x_{ij}\in[0,1]\).
\item Apply the BKAR minimum interpolation to the complete
      self-adjoint Hamiltonian
      \(\sum x_{ij}\Omega_{ij}\), retaining simplex time ordering.
\item Prove that each forest derivative has one pair row per edge
      and a complete exponential remainder controlled by
      \eqref{eq:newV28-operator-stability}.
\item Determine whether commutator terms appear outside the
      time-ordered simplex.  If they do, keep the simplex integral;
      do not commute the insertions.
\item Add the degree-two and degree-four score roots only after the
      unrooted quantum forest bound is established.
\item Test the formula on the spin decomposition of four
      fundamentals from V26.
\end{enumerate}
\end{programme}

\begin{firewall}[Claim boundary after V28]
\label{fire:newV28-boundary}
V28 proves the exact conditional radial score, closes the
all-order degree-two scalar graph, proves the dimension-free Cartan
matrix RGSC, and establishes subset-uniform noncommutative Casimir
stability plus the sharp pair-row bound.  It does not prove a
time-ordered quantum Penrose formula, full matrix RGSC, GIAC, the
Wilson aggregate strong card, all-order strong MTA, WUFC, CPM,
cutoff removal, reflection positivity,
Osterwalder--Schrader reconstruction, a continuum Yang--Mills
measure, or a positive Hamiltonian gap.  The full Yang--Mills mass
gap has not been proved.
\end{firewall}

\begin{theorem}[V28 result ledger]
\label{thm:newV28-results}
V28:
\begin{enumerate}[label=\textup{(V28.\arabic*)},leftmargin=6.5em]
\item computes
      \(\mathbb E[|Y|^4-15\mid Y_3]=H_4+10H_2\);
\item proves the all-order scalar \(H_2\) rooted-graph card;
\item lifts the complete scalar result to arbitrary
      simultaneously diagonalizable representation weights;
\item proves the exact noncommutative subset stability identity;
\item proves the representation-uniform pair-row bound; and
\item isolates time ordering, rather than Casimir growth, as the
      next matrix gate.
\end{enumerate}
\end{theorem}

\begin{proof}
Use
\Cref{lem:newV28-conditional-score,
prop:newV28-H2-card,
cor:newV28-radial-Cartan-card,
thm:newV28-Cartan-RGSC,
thm:newV28-operator-stability,
cor:newV28-heat-stability,
prop:newV28-pair-row,
lem:newV28-disjoint-commute}.
\end{proof}

\paragraph{Parent-version record.}
V28 was formed from
\texttt{Renewal\_Yang\_Mills\_Mass\_Gap\_Research\_Notes\_V27.tex}.
The V27 parent had \(12\,724\) logical source lines,
\(407\,563\) bytes, and SHA--256
\[
 \texttt{bacc59b2cd70ee50a3b1e5890e90939a2dc74eaeafeb714a7733fad3fa56ddb3}.
\]
The next compulsory companion audit is V30.

% =============================================================================
% END APPENDED FIFTH-SERIES V28 MATERIAL
% =============================================================================

% =============================================================================
% BEGIN APPENDED FIFTH-SERIES V29 MATERIAL
% =============================================================================

\chapter{V29: Time-ordered quantum forests and the strong heat
endpoint}
\label{chap:newV29-quantum-forest}

\section{The calibrated heat moment}
\label{sec:newV29-heat-moment}

Let \(\mu_s^H\) be the central \(SU(2)\) heat-kernel law normalized
so that its Fourier multiplier in representation \(R\) is
\[
 q_{s,R}^H=e^{-sC_2(R)}.
\]
For inputs \(R_1,\ldots,R_m\), use the Hermitian generators and
pair interactions of V28:
\[
 C_i=C_2(R_i)I,
 \qquad
 \Omega_{ij}=\sum_{a=1}^3J_i^aJ_j^a.
\]

\begin{lemma}[Exact tensor heat moment]
\label{lem:newV29-heat-moment}
For every \(I\subseteq[m]\),
\begin{align}
 M_I^H
 &:=
 \mathbb E_{\mu_s^H}
 \bigotimes_{i\in I}D^{R_i}(U)
 \notag\\
 &=
 e^{-sC_I^{\rm tot}}
 =
 e^{-s\sum_{i\in I}C_i}
 \exp\left(
 -2s\sum_{\substack{i<j\\i,j\in I}}\Omega_{ij}
 \right).
 \label{eq:newV29-heat-moment}
\end{align}
\end{lemma}

\begin{proof}
The tensor product is the representation
\(\bigotimes_{i\in I}R_i\), whose infinitesimal generators are
\(\sum_{i\in I}J_i^a\).  The heat-kernel Fourier transform is
functional calculus of its total Casimir, giving the first
exponential.  Expand the total Casimir as in
\eqref{eq:newV28-operator-stability}.  Each \(C_i\) is scalar on
its irreducible input factor and therefore commutes with every
\(\Omega_{jk}\), which proves the factorization.
\end{proof}

\section{BKAR minimum matrices preserve operator stability}
\label{sec:newV29-BKAR-stability}

Let \(T\) be a labelled tree on \([m]\), assign
\(\boldsymbol u=(u_e)_{e\in E(T)}\in[0,1]^{m-1}\), and define
\begin{equation}
 x_{ij}^T(\boldsymbol u)
 :=
 \min_{e\in{\rm path}_T(i,j)}u_e
 \quad(i\ne j),
 \qquad
 x_{ii}^T=1.
 \label{eq:newV29-BKAR-min}
\end{equation}

\begin{lemma}[Threshold-block layer cake]
\label{lem:newV29-threshold-layer}
For \(0\le t\le1\), delete from \(T\) all edges with \(u_e<t\)
and let \(\mathcal C_t\) be the resulting component partition.
Then
\begin{equation}
 x_{ij}^T(\boldsymbol u)
 =
 \int_0^1
 \mathbf1_{\{i,j\text{ belong to the same }C\in\mathcal C_t\}}
 \,dt.
 \label{eq:newV29-threshold-layer}
\end{equation}
\end{lemma}

\begin{proof}
Vertices \(i,j\) remain in the same threshold component exactly
when every edge on their unique tree path has \(u_e\ge t\), that
is, when \(t\le\min_{e\in{\rm path}_T(i,j)}u_e\).  The measure of
this interval is \eqref{eq:newV29-BKAR-min}.
\end{proof}

\begin{theorem}[Operator stability at every BKAR point]
\label{thm:newV29-BKAR-stability}
Let \(b_i^2=s(1+C_2(R_i))<1\).  At every tree minimum matrix,
\begin{equation}
 2s\sum_{i<j}x_{ij}^T(\boldsymbol u)\Omega_{ij}
 \ge
 -\sum_{i=1}^m b_i^2\,I
 \ge-mI.
 \label{eq:newV29-BKAR-stability}
\end{equation}
The same conclusion holds for a forest minimum matrix, block by
block.
\end{theorem}

\begin{proof}
Use \Cref{lem:newV29-threshold-layer} to write
\begin{align*}
 2s\sum_{i<j}x_{ij}^T\Omega_{ij}
 &=
 \int_0^1
 \sum_{C\in\mathcal C_t}
 2s\sum_{\substack{i<j\\i,j\in C}}\Omega_{ij}\,dt.
\end{align*}
Apply the subset operator stability
\eqref{eq:newV28-operator-stability} to every block \(C\):
\[
 2s\sum_{i<j\in C}\Omega_{ij}
 \ge-\sum_{i\in C}b_i^2I.
\]
The blocks partition \([m]\), so their right sides sum to
\(-\sum_i b_i^2I\), independently of \(t\).  Integrate.  A forest
has the same threshold representation inside its permanent
components.
\end{proof}

\section{Mixed Duhamel derivatives without a factorial}
\label{sec:newV29-Duhamel}

For a coupling array \(\boldsymbol x\), set
\begin{equation}
 H(\boldsymbol x)
 :=
 2s\sum_{i<j}x_{ij}\Omega_{ij},
 \qquad
 {\cal U}(\boldsymbol x):=e^{-H(\boldsymbol x)}.
 \label{eq:newV29-Hx}
\end{equation}

\begin{lemma}[Exact time-ordered mixed derivative]
\label{lem:newV29-Duhamel}
For distinct edges \(e_1,\ldots,e_k\), write
\(V_e=2s\Omega_e\).  Then
\begin{align}
 \partial_{e_1}\cdots\partial_{e_k}{\cal U}(\boldsymbol x)
 &=
 (-1)^k
 \sum_{\pi\in S_k}
 \int_{\Delta_k}
 e^{-\tau_0H}
 V_{e_{\pi(1)}}e^{-\tau_1H}
 \cdots
 V_{e_{\pi(k)}}e^{-\tau_kH}
 \,d\boldsymbol\tau,
 \label{eq:newV29-Duhamel}
\end{align}
where
\[
 \Delta_k
 :=
 \{\tau_j\ge0:\ \tau_0+\cdots+\tau_k=1\}
\]
has volume \(1/k!\).
\end{lemma}

\begin{proof}
The first derivative is Duhamel's formula
\[
 \partial_e e^{-H}
 =
 -\int_0^1e^{-(1-r)H}V_e e^{-rH}\,dr.
\]
Iterate it.  Each new derivative can be inserted into any one of
the existing exponential intervals.  Ordering the \(k\) insertions
gives the sum over \(S_k\), and the interval lengths give the
simplex coordinates.  This proves the formula by induction.
\end{proof}

\begin{theorem}[Quantum tree derivative bound]
\label{thm:newV29-tree-derivative}
At a BKAR minimum matrix and for every set \(E\) of distinct
edges,
\begin{equation}
 \left\|
 \partial_E{\cal U}(\boldsymbol x^T(\boldsymbol u))
 \right\|_{\rm op}
 \le
 e^m\prod_{ij\in E}2b_ib_j.
 \label{eq:newV29-tree-derivative}
\end{equation}
There is no \(k!\) factor.
\end{theorem}

\begin{proof}
By \Cref{thm:newV29-BKAR-stability},
\[
 H(\boldsymbol x^T)\ge-mI,
\]
so functional calculus gives
\[
 \|e^{-\tau H(\boldsymbol x^T)}\|_{\rm op}
 \le e^{\tau m}.
\]
For every Duhamel word in
\eqref{eq:newV29-Duhamel}, the semigroup factors therefore
contribute at most
\[
 e^{m(\tau_0+\cdots+\tau_k)}=e^m.
\]
The pair-row bound \eqref{eq:newV28-pair-row} gives
\(\|V_{ij}\|\le2b_ib_j\).  There are \(k!\) ordered words, while
\(\operatorname{vol}(\Delta_k)=1/k!\).  Their factors cancel
exactly, proving \eqref{eq:newV29-tree-derivative}.
\end{proof}

\begin{assessment}[Time ordering is harmless when retained]
\label{ass:newV29-ordering-closed}
V28 isolated the commutators of overlapping pair edges.  They do
not need to be expanded.  The exact Duhamel simplex retains their
order, and submultiplicativity uses only the complete exponential
stability between insertions.  Thus noncommutativity introduces
neither a commutator census nor an analytic factorial.
\end{assessment}

\section{Exact heat cumulant forest}
\label{sec:newV29-heat-forest}

For \(0\le x_{ij}\le1\), define
\begin{equation}
 F_m(\boldsymbol x)
 :=
 e^{-s\sum_iC_i}
 \exp\left(-2s\sum_{i<j}x_{ij}\Omega_{ij}\right).
 \label{eq:newV29-Fm}
\end{equation}
If \(\rho\) is a partition of \([m]\), let
\(\boldsymbol x^\rho\) have \(x_{ij}=1\) within blocks and zero
between blocks.

\begin{lemma}[Partition endpoints are exact moment products]
\label{lem:newV29-partition-endpoints}
\begin{equation}
 F_m(\boldsymbol x^\rho)
 =
 \bigotimes_{B\in\rho}M_B^H.
 \label{eq:newV29-partition-endpoints}
\end{equation}
\end{lemma}

\begin{proof}
At \(\boldsymbol x^\rho\), the Hamiltonian in
\eqref{eq:newV29-Fm} is a sum over disjoint tensor blocks.
The block operators commute and their exponential factorizes.
Apply \Cref{lem:newV29-heat-moment} on every block.
\end{proof}

\begin{theorem}[Exact time-ordered heat cumulant forest]
\label{thm:newV29-heat-forest}
For \(m\ge2\),
\begin{equation}
 \kappa_m^H(D^{R_1},\ldots,D^{R_m})
 =
 \sum_{T\in\mathfrak T_m}
 \int_{[0,1]^{m-1}}
 \partial_{E(T)}
 F_m(\boldsymbol x^T(\boldsymbol u))
 \,d\boldsymbol u.
 \label{eq:newV29-heat-forest}
\end{equation}
\end{theorem}

\begin{proof}
Apply the multivariate BKAR forest formula to
\(F_m\).  By
\Cref{lem:newV29-partition-endpoints}, its value at a partition
matrix is the tensor product of the heat moments on that partition.
The Möbius alternating sum over partition endpoints is exactly the
complete tensor cumulant.  In the forest formula, that Möbius sum
annihilates every disconnected forest; the surviving connected
forests are precisely the labelled spanning trees.  Their minimum
matrices and edge derivatives give
\eqref{eq:newV29-heat-forest}.  All derivatives are justified in
finite dimension by
\Cref{lem:newV29-Duhamel}.
\end{proof}

\begin{theorem}[Strong heat-kernel forest card]
\label{thm:newV29-strong-heat-card}
Every tree term in \eqref{eq:newV29-heat-forest} satisfies
\begin{equation}
 \left\|
 \int_{[0,1]^{m-1}}
 \partial_{E(T)}
 F_m(\boldsymbol x^T(\boldsymbol u))
 \,d\boldsymbol u
 \right\|_{\rm op}
 \le
 (2e)^{m}
 \prod_{i=1}^m b_i^{d_i(T)}.
 \label{eq:newV29-strong-heat-card}
\end{equation}
Consequently
\begin{equation}
 \left\|
 \kappa_m^H(D^{R_1},\ldots,D^{R_m})
 \right\|_{\rm op}
 \le
 (2e)^m
 \left(\prod_i b_i\right)
 \left(\sum_i b_i\right)^{m-2}.
 \label{eq:newV29-heat-aggregate}
\end{equation}
Both constants are uniform in \(m,s\), the representations, and
all representation dimensions, subject only to \(b_i<1\).
\end{theorem}

\begin{proof}
The scalar factor \(e^{-s\sum_iC_i}\) in
\eqref{eq:newV29-Fm} has norm at most one and is independent of
\(\boldsymbol x\).  Apply
\Cref{thm:newV29-tree-derivative} with \(k=m-1\):
\[
 \|\partial_{E(T)}F_m\|
 \le
 e^m2^{m-1}
 \prod_{ij\in E(T)}b_ib_j
 =
 e^m2^{m-1}
 \prod_i b_i^{d_i(T)}.
\]
The parameter cube has volume one, proving
\eqref{eq:newV29-strong-heat-card} after a harmless enlargement.
Sum over \(T\) and use the weighted Pr\"ufer identity
\[
 \sum_{T\in\mathfrak T_m}\prod_i b_i^{d_i(T)}
 =
 \left(\prod_i b_i\right)
 \left(\sum_i b_i\right)^{m-2}.
\]
\end{proof}

\begin{corollary}[The former conditional heat endpoint is closed]
\label{cor:newV29-heat-endpoint-closed}
The heat strong-card hypothesis used conditionally in V19 is now a
theorem at the raw local normalization.  Its proof is
nonperturbative in \(sC_2(R_i)\) and remains uniform for growing
representations in the thermal regime \(b_i<1\).
\end{corollary}

\begin{proof}
\Cref{thm:newV29-strong-heat-card} is exactly the required
fixed-tree estimate with an exponential base, and
\eqref{eq:newV29-heat-aggregate} is its aggregate form.
\end{proof}

\section{What this changes in the comparison route}
\label{sec:newV29-route}

\begin{assessment}[One of the two raw endpoint cards is acquired]
\label{ass:newV29-route}
The Wilson programme no longer needs GIAC specifically if a
Wilson--heat comparison can be made with a raw connected score card:
the calibrated heat endpoint already has the exact all-order strong
forest bound.  The live comparison gate is therefore reduced to a
single family of path score cumulants.  The old V19 heat path was
abandoned because its uniform gap was not established; V29 does not
silently restore that gap.  It supplies only the endpoint analytic
card.
\end{assessment}

\begin{firewall}[Endpoint closure is not path closure]
\label{fire:newV29-no-path}
\Cref{thm:newV29-strong-heat-card} does not prove a uniform
spectral gap for a Wilson--heat interpolation, an intrinsic
negative-norm score row on that path, or a raw score insertion
card.  Nor may the later first-connectivity payer be attached to
\eqref{eq:newV29-heat-aggregate}; V26 proves that such a local
normalization is false.
\end{firewall}

\section{V30 programme and claim boundary}
\label{sec:newV29-next}

\begin{programme}[V30: audit and a gapped heat--Wilson bridge]
\label{prog:newV29-V30}
\begin{enumerate}[leftmargin=3em]
\item Perform the compulsory read-only audit of the newest active
      SM--Einstein and Navier--Stokes notes.
\item Compare the heat-kernel one-link density with the monotone
      geodesic path of V21 without differentiating a principal-chart
      boundary.
\item Seek a two-stage bridge: heat to a monotone intrinsic
      reference with an exact Fourier endpoint, then reference to
      Wilson through the proved V21 gap.
\item Alternatively construct a heat--Wilson potential path whose
      radial and angular gaps admit the V21 Hardy proof.
\item Keep every score comparison at the raw local normalization.
\item Reuse \Cref{thm:newV29-strong-heat-card} as the endpoint;
      do not re-prove it by a Gaussian chart expansion.
\end{enumerate}
\end{programme}

\begin{firewall}[Claim boundary after V29]
\label{fire:newV29-boundary}
V29 proves that BKAR minimum matrices preserve the noncommutative
Casimir stability bound, derives the exact mixed Duhamel formula,
proves a factorial-free quantum tree derivative bound, and obtains
the full all-order strong forest card for the calibrated
\(SU(2)\) heat-kernel endpoint.  It does not prove a gapped
heat--Wilson bridge, a path score card, the Wilson aggregate strong
card, all-order strong MTA, WUFC, CPM, cutoff removal, reflection
positivity, Osterwalder--Schrader reconstruction, a continuum
Yang--Mills measure, or a positive Hamiltonian gap.  The full
Yang--Mills mass gap has not been proved.
\end{firewall}

\begin{theorem}[V29 result ledger]
\label{thm:newV29-results}
V29:
\begin{enumerate}[label=\textup{(V29.\arabic*)},leftmargin=6.5em]
\item writes the exact tensor heat moment as a total-Casimir
      exponential;
\item proves the threshold-block representation of every BKAR
      minimum matrix;
\item proves subset stability at every interpolation point;
\item derives time-ordered mixed derivatives with exact
      permutation--simplex cancellation;
\item proves the exact heat cumulant forest identity; and
\item proves the representation- and arity-uniform strong
      heat-kernel forest card.
\end{enumerate}
\end{theorem}

\begin{proof}
Use
\Cref{lem:newV29-heat-moment,
lem:newV29-threshold-layer,
thm:newV29-BKAR-stability,
lem:newV29-Duhamel,
thm:newV29-tree-derivative,
lem:newV29-partition-endpoints,
thm:newV29-heat-forest,
thm:newV29-strong-heat-card}.
\end{proof}

\paragraph{Parent-version record.}
V29 was formed from
\texttt{Renewal\_Yang\_Mills\_Mass\_Gap\_Research\_Notes\_V28.tex}.
The V28 parent had \(13\,145\) logical source lines,
\(419\,765\) bytes, and SHA--256
\[
 \texttt{d8fe61e8c2edc16be47ddd6318140d6d38aba8d3d81425cf2b010337d4dee66b}.
\]
The next compulsory companion audit is V30.

% =============================================================================
% END APPENDED FIFTH-SERIES V29 MATERIAL
% =============================================================================

% =============================================================================
% BEGIN APPENDED FIFTH-SERIES V30 MATERIAL
% =============================================================================

\chapter{V30: Compulsory audit and concentration of the
heat--Wilson bridge}
\label{chap:newV30-bridge}

\section{Read-only companion audit}
\label{sec:newV30-audit}

The newest active companion files inspected before this append were
\begin{center}
\begin{tabular}{p{0.62\textwidth}r r}
\hline
file & bytes & logical lines\\
\hline
\texttt{Renewal\_SM\_Einstein\_Continuum\_Research\_Notes\_V38.tex}
& \(474\,307\) & \(12\,182\)\\
\texttt{Renewal\_NS\_Stability\_Research\_Notes\_V31.tex}
& \(887\,537\) & \(23\,052\)\\
\hline
\end{tabular}
\end{center}
Their SHA--256 digests were
\begin{align}
 &\texttt{54e698287185e277beb0ca3ad6649ebd0397f3874577519044b1de86f006f9ce},
 \label{eq:newV30-SM-hash}\\
 &\texttt{f1121b62238ad5491bb7cf03635ea9934942edf573ed8faaa9ee79e1d38a6696}.
 \label{eq:newV30-NS-hash}
\end{align}
Neither companion file was modified.

SM V35 independently imports the V25 quartic-score saturation test
and concludes that a large restoring operation cannot be bounded
after a separate score norm.  Its payer-complete regional card
remains open.  YM V26 goes further in the one-link problem:
the saturating tensor survives in spin one and spin two, so the
standalone paid local card is actually false under the inherited
Peter--Weyl payer.  SM V36--V38 then address boundary Hodge
commutators, Schur gluing, and anomaly arithmetic; those are
specific to the gauge--matter--gravity programme and do not supply
a Wilson one-link bridge.

NS V31 is unchanged.  Its exact heat-formation ladder remains a
useful warning: a multiplicity-free heat representation can return
the critical continuation stock when the restoring first moment is
applied after norms.  No fluid formation, flux, or ancestry estimate
is transferred to Yang--Mills.

\begin{firewall}[Audit type boundary]
\label{fire:newV30-audit-boundary}
The companion files support the order of operations used in
V26--V29 but prove no compact-group heat-kernel logarithmic slope,
Wilson score row, all-order matrix cumulant, or Yang--Mills
continuum statement.
\end{firewall}

\section{Exact image form of the \(SU(2)\) heat endpoint}
\label{sec:newV30-image}

Retain the V20 metric convention: the class angle is
\(\theta\in[0,\pi]\), and the heat multiplier is
\(e^{-4\tau C_2(R)}\).  Poisson summation of the character series
gives the following radial form.

\begin{lemma}[Wrapped image formula]
\label{lem:newV30-image}
For \(0<\theta<\pi\),
\begin{equation}
 h_\tau(\theta)
 =
 {c_\tau\over\sin\theta}
 \sum_{k\in\mathbb Z}
 (\theta+2\pi k)
 \exp\left(
 -{(\theta+2\pi k)^2\over4\tau}
 \right),
 \label{eq:newV30-image}
\end{equation}
where \(c_\tau>0\) depends only on \(\tau\) and the normalization of
Haar measure.  The numerator has simple zeros at \(0,\pi\), and the
quotient extends to a smooth strictly positive class function at
both endpoints.
\end{lemma}

\begin{proof}
The character expansion is
\[
 h_\tau(\theta)
 =
 \sum_{n=0}^\infty
 (n+1)e^{-\tau n(n+2)}
 {\sin((n+1)\theta)\over\sin\theta}.
\]
Write \(n(n+2)=(n+1)^2-1\), extend the sine series oddly to
\(\mathbb Z\), and apply the one-dimensional Poisson summation
formula to the derivative of the periodized Gaussian.  All
\(\theta\)-independent factors, including \(e^\tau\), form
\(c_\tau\), giving \eqref{eq:newV30-image}.  Oddness at \(0\) and
pairing \(k\) with \(-k-1\) at \(\pi\) give the endpoint zeros.
The character series is smooth and positive for \(\tau>0\), so the
removable quotients have the stated extensions.
\end{proof}

\begin{lemma}[Uniform central heat scaling]
\label{lem:newV30-central-scaling}
For every fixed \(R<\infty\), there are \(0<c_R<C_R<\infty\)
such that, for all sufficiently small \(\tau\) and
\(0\le\theta\le R\sqrt\tau\),
\begin{equation}
 c_R\tau^{-3/2}e^{-\theta^2/(4\tau)}
 \le h_\tau(\theta)
 \le
 C_R\tau^{-3/2}e^{-\theta^2/(4\tau)}.
 \label{eq:newV30-central-scaling}
\end{equation}
\end{lemma}

\begin{proof}
In \eqref{eq:newV30-image}, the \(k=0\) image is
\[
 {\theta\over\sin\theta}e^{-\theta^2/(4\tau)}
\]
and \(\theta/\sin\theta\) is bounded above and below on the stated
shrinking interval.  Every \(k\ne0\) image is
\(O_R(e^{-c/\tau})\) relative to it.  Normalization follows by
integrating against
\(\sin^2\theta\,d\theta\): after
\(\theta=\sqrt\tau\,r\), the leading radial mass is a fixed
positive multiple of
\(\tau^{3/2}\int_0^\infty r^2e^{-r^2/4}\,dr\).
Hence \(c_\tau\asymp\tau^{-3/2}\), proving the two-sided bound.
\end{proof}

\begin{lemma}[Heat distance moments]
\label{lem:newV30-heat-moments}
There is a numerical \(C\) such that for every \(q\ge1\) and
\(0<\tau\le\tau_0\),
\begin{equation}
 \left(
 \int_{SU(2)}d(I,U)^q h_\tau(U)\,dU
 \right)^{1/q}
 \le C\sqrt{q\tau}.
 \label{eq:newV30-heat-moments}
\end{equation}
\end{lemma}

\begin{proof}
Multiply \eqref{eq:newV30-image} by the radial Haar factor
\(\sin^2\theta\).  Pair the images adjacent to each endpoint and
use the triangle inequality only after that pairing.  The result is
bounded by
\[
 C\tau^{-3/2}\theta^2e^{-c\theta^2/\tau}
 +Ce^{-c/\tau}
\]
on \(0\le\theta\le\pi\).  The second term is harmless after
integration.  The first gives
\[
 C\tau^{-3/2}
 \int_0^\infty
 \theta^{q+2}e^{-c\theta^2/\tau}\,d\theta
 \le
 C^q\tau^{q/2}
 \Gamma\left({q+3\over2}\right).
\]
Taking the \(q\)-th root gives
\eqref{eq:newV30-heat-moments} while
\(q\tau\le1\).  If \(q\tau>1\), enlarge \(C\) and use
\(\theta\le\pi\).
\end{proof}

\section{Uniform overlap and moments of the geometric bridge}
\label{sec:newV30-overlap}

Let \(w_\alpha\) be the normalized Wilson density relative to Haar,
let
\[
 \tau_\alpha={1\over2\alpha}+O(\alpha^{-2}),
\]
and define the geometric bridge
\begin{align}
 d\nu_{\alpha,t}
 &:=
 {h_{\tau_\alpha}^{\,1-t}w_\alpha^t\over
   {\cal Z}_{\alpha,t}}\,dU,
 \label{eq:newV30-geometric-bridge}\\
 {\cal Z}_{\alpha,t}
 &:=
 \int h_{\tau_\alpha}^{\,1-t}w_\alpha^t\,dU.
 \label{eq:newV30-Hellinger}
\end{align}
This is the V19--V20 path, now with the V29 endpoint card available.

\begin{theorem}[Uniform Hellinger overlap]
\label{thm:newV30-overlap}
There is \(c>0\) such that
\begin{equation}
 c\le{\cal Z}_{\alpha,t}\le1
 \label{eq:newV30-overlap}
\end{equation}
for every \(t\in[0,1]\) and all sufficiently large \(\alpha\).
\end{theorem}

\begin{proof}
The upper bound is H\"older's inequality and normalization of both
endpoint densities.  On the ball
\(\theta\le R/\sqrt\alpha\),
\Cref{lem:newV30-central-scaling} and the Wilson local scaling from
V21 give, for fixed sufficiently large \(R\),
\[
 h_{\tau_\alpha}(\theta)
 \asymp_R
 w_\alpha(\theta)
 \asymp_R
 \alpha^{3/2}e^{-\alpha\theta^2/2}.
\]
The exponent calibration uses
\(\tau_\alpha=(2\alpha)^{-1}+O(\alpha^{-2})\).
Their weighted geometric mean has the same lower bound uniformly
in \(t\).  Integrating it against
\(\sin^2\theta\,d\theta\,d\omega\) on this ball and scaling
\(r=\sqrt\alpha\,\theta\) gives a fixed positive constant.
\end{proof}

\begin{theorem}[Adaptive bridge concentration]
\label{thm:newV30-bridge-moments}
There is a numerical \(C\) such that, uniformly in
\(t\in[0,1]\), \(q\ge1\), and sufficiently large \(\alpha\),
\begin{equation}
 \left(
 \mathbb E_{\alpha,t}\theta^q
 \right)^{1/q}
 \le C\sqrt{q\over\alpha}.
 \label{eq:newV30-bridge-moments}
\end{equation}
In particular, the bridge has a median at
\(O(\alpha^{-1/2})\), uniformly in \(t\).
\end{theorem}

\begin{proof}
By \eqref{eq:newV30-Hellinger} and H\"older,
\begin{align*}
 {\cal Z}_{\alpha,t}
 \mathbb E_{\alpha,t}\theta^q
 &=
 \int
 (\theta^qh_{\tau_\alpha})^{1-t}
 (\theta^qw_\alpha)^t\,dU\\
 &\le
 \left(\mathbb E_{h_{\tau_\alpha}}\theta^q\right)^{1-t}
 \left(\mathbb E_{w_\alpha}\theta^q\right)^t.
\end{align*}
The heat factor is bounded by
\Cref{lem:newV30-heat-moments}; the Wilson factor is bounded by
\Cref{lem:newV22-radial-moments}.  Divide by the uniform lower
bound in \eqref{eq:newV30-overlap} and take \(q\)-th roots,
absorbing \(c^{-1/q}\) into \(C\).  Markov's inequality at a fixed
large multiple of \(\alpha^{-1/2}\) places at least half the mass
to its left, proving the median statement.
\end{proof}

\begin{assessment}[The density-ratio obstruction is bypassed]
\label{ass:newV30-overlap-progress}
V20 proves that the endpoint log-density ratio has oscillation
linear in \(\alpha\), so a global Holley--Stroock comparison is
exponentially weak.  \Cref{thm:newV30-overlap,
thm:newV30-bridge-moments} show that this remote antipodal
oscillation does not destroy the normalization or concentration of
the actual geometric bridge.  The remaining gap question is a
tail-transport estimate, not a mass-overlap estimate.
\end{assessment}

\section{The single remaining gap profile}
\label{sec:newV30-log-slope}

\begin{definition}[Heat logarithmic-slope card]
\label{def:newV30-HLSC}
The HLSC holds if there are fixed \(R,c>0\) such that, for all
sufficiently small \(\tau\),
\begin{align}
 -\partial_\theta\log h_\tau(\theta)
 &\ge c{\theta\over\tau},
 &&R\sqrt\tau\le\theta\le{\pi\over2},
 \label{eq:newV30-HLSC-inner}\\
 -\partial_\theta\log h_\tau(\theta)
 &\ge c{\pi-\theta\over\tau},
 &&{\pi\over2}\le\theta<\pi.
 \label{eq:newV30-HLSC-cap}
\end{align}
\end{definition}

\begin{theorem}[HLSC closes the old Wilson--heat path gap]
\label{thm:newV30-HLSC-sufficient}
If HLSC holds, then
\begin{equation}
 \operatorname{gap}(\nu_{\alpha,t})\ge c_1\alpha
 \label{eq:newV30-bridge-gap}
\end{equation}
uniformly in \(t\in[0,1]\) for all sufficiently large \(\alpha\).
\end{theorem}

\begin{proof}
The unnormalized radial density is
\[
 p_{\alpha,t}(\theta)
 =
 \sin^2\theta\,
 h_{\tau_\alpha}(\theta)^{1-t}
 e^{t\alpha\cos\theta}.
\]
For \(\theta\ge R/\sqrt\alpha\),
\begin{align*}
 -\partial_\theta\log p_{\alpha,t}
 &=
 (1-t)(-\partial_\theta\log h_{\tau_\alpha})
 +t\alpha\sin\theta-2\cot\theta.
\end{align*}
On \([R/\sqrt\alpha,\pi/2]\), HLSC,
\(\tau_\alpha\asymp\alpha^{-1}\), and
\(\sin\theta\ge2\theta/\pi\) give
\[
 -\partial_\theta\log p_{\alpha,t}
 \ge c\alpha\theta-{2\over\theta}
 \ge c_R\sqrt\alpha
\]
after increasing \(R\).

On \([\pi/2,\pi)\), put \(\delta=\pi-\theta\).
HLSC and \(\sin\theta=\sin\delta\) give a lower bound
\(c\alpha\delta-2\cot\theta\).
For small \(\delta\), the last term is
\(2\cot\delta\ge c/\delta\); for
\(\delta\) bounded below it is dominated by
\(c\alpha\delta\).  Hence the whole descending side again has
logarithmic slope at least \(c\sqrt\alpha\).

The adaptive concentration
\eqref{eq:newV30-bridge-moments} supplies the two-sided anchor mass
at \(R/\sqrt\alpha\).  The preceding descending slope gives the
right Hardy product \(O(\alpha^{-1})\), exactly as in
\Cref{lem:newV21-Hardy-products}.  On the left of the anchor,
the central heat and Wilson estimates give
\(p_{\alpha,t}(\theta)\asymp_R\theta^2\), up to normalization, so
the left Hardy product is also \(O(\alpha^{-1})\).
Thus
\(\lambda_{\alpha,t}^{\rm rad}\ge c\alpha\) by
\Cref{thm:newV20-Hardy-criterion}.

For the angular \(\ell=1\) form, use
\(2/\sin^2\theta\ge c\alpha\) inside the anchor and the same
descending-side Dirichlet Hardy plus anchor trace argument as
\Cref{lem:newV21-Dirichlet-Hardy,
lem:newV21-trace}.  This gives
\(\lambda_{\alpha,t}^{(1)}\ge c\alpha\).  Apply the exact harmonic
reduction \eqref{eq:newV20-gap-reduction}.
\end{proof}

\begin{firewall}[HLSC remains unproved]
\label{fire:newV30-HLSC-open}
\Cref{lem:newV30-image} makes HLSC an explicit one-dimensional
inequality for a wrapped Gaussian derivative.  V30 proves neither
\eqref{eq:newV30-HLSC-inner} nor
\eqref{eq:newV30-HLSC-cap}; cancellations between the two shortest
images become singular on the antipodal scale
\(\pi-\theta\asymp\tau\).  Monotonicity of the heat kernel alone is
insufficient because the proof requires the displayed quantitative
slope.  No Wilson--heat path gap is claimed.
\end{firewall}

\section{V31 programme and claim boundary}
\label{sec:newV30-next}

\begin{programme}[V31: two-image logarithmic slope]
\label{prog:newV30-V31}
\begin{enumerate}[leftmargin=3em]
\item Differentiate the exact image numerator in
      \eqref{eq:newV30-image}.
\item On \(0\le\theta\le\pi/2\), dominate every nonzero image by
      the \(k=0\) image and prove
      \eqref{eq:newV30-HLSC-inner}.
\item On the antipodal cap, pair \(k=0\) with \(k=-1\) before any
      absolute estimate and use
      \(\delta=(\pi-\theta)/\tau\).
\item Prove the cap slope first for
      \(\pi-\theta\le c\tau\), then for
      \(c\tau\le\pi-\theta\le\pi/2\), where one shortest image
      dominates.
\item Insert the resulting HLSC into
      \Cref{thm:newV30-HLSC-sufficient}; do not repeat the Hardy
      and angular trace arguments.
\item Only after the path gap is proved, return to its raw complete
      score insertion, using the V29 heat endpoint card.
\end{enumerate}
\end{programme}

\begin{firewall}[Claim boundary after V30]
\label{fire:newV30-boundary}
V30 performs the compulsory companion audit, derives the exact
wrapped image form of the \(SU(2)\) heat kernel, proves central
heat scaling and adaptive heat moments, and proves uniform
Hellinger overlap plus adaptive concentration along the full
geometric heat--Wilson bridge.  It reduces the path gap to the
explicit HLSC and proves that HLSC is sufficient, but does not prove
HLSC itself.  It therefore does not prove the Wilson--heat path
gap, its raw score card, the Wilson aggregate strong card,
all-order strong MTA, WUFC, CPM, cutoff removal, reflection
positivity, Osterwalder--Schrader reconstruction, a continuum
Yang--Mills measure, or a positive Hamiltonian gap.  The full
Yang--Mills mass gap has not been proved.
\end{firewall}

\begin{theorem}[V30 result ledger]
\label{thm:newV30-results}
V30:
\begin{enumerate}[label=\textup{(V30.\arabic*)},leftmargin=6.5em]
\item completes the required read-only SM/NS audit;
\item derives the wrapped Gaussian heat-kernel formula;
\item proves central heat scaling and all adaptive distance moments;
\item proves uniform endpoint Hellinger overlap;
\item proves adaptive concentration and a uniform identity-scale
      median for every bridge measure; and
\item reduces the full path gap to one quantitative heat
      logarithmic-slope card.
\end{enumerate}
\end{theorem}

\begin{proof}
Use
\Cref{lem:newV30-image,
lem:newV30-central-scaling,
lem:newV30-heat-moments,
thm:newV30-overlap,
thm:newV30-bridge-moments,
thm:newV30-HLSC-sufficient}.
\end{proof}

\paragraph{Parent-version record.}
V30 was formed from
\texttt{Renewal\_Yang\_Mills\_Mass\_Gap\_Research\_Notes\_V29.tex}.
The V29 parent had \(13\,605\) logical source lines,
\(433\,724\) bytes, and SHA--256
\[
 \texttt{533daa54c78a98e56657c62a0b26b15150383cad003a1e76319dd333f2dc1782}.
\]
The next compulsory companion audit is V35.

% =============================================================================
% END APPENDED FIFTH-SERIES V30 MATERIAL
% =============================================================================

% =============================================================================
% BEGIN APPENDED FIFTH-SERIES V31 MATERIAL
% =============================================================================

\chapter{V31: Two-image heat slope and the gapped
heat--Wilson bridge}
\label{chap:newV31-slope}

\section{The image numerator}
\label{sec:newV31-numerator}

Put
\begin{equation}
 N_\tau(\theta)
 :=
 \sum_{k\in\mathbb Z}
 (\theta+2\pi k)
 e^{-(\theta+2\pi k)^2/(4\tau)}.
 \label{eq:newV31-N}
\end{equation}
By \eqref{eq:newV30-image},
\[
 h_\tau(\theta)=c_\tau N_\tau(\theta)/\sin\theta.
\]
Only logarithmic derivatives are needed, so \(c_\tau\) disappears.

\begin{lemma}[Shortest-image control on the inner half]
\label{lem:newV31-inner-image}
There are \(c,C>0\) such that, for
\(0<\tau\le\tau_0\) and \(0<\theta\le\pi/2\),
\begin{align}
 N_\tau(\theta)
 &=
 \theta e^{-\theta^2/(4\tau)}
 (1+\varepsilon_\tau(\theta)),
 \label{eq:newV31-inner-N}\\
 |\varepsilon_\tau(\theta)|
 +\sqrt\tau\,|\partial_\theta\varepsilon_\tau(\theta)|
 &\le Ce^{-c/\tau}.
 \label{eq:newV31-inner-error}
\end{align}
\end{lemma}

\begin{proof}
For \(k\ne0\) and \(0<\theta\le\pi/2\),
\[
 |\theta+2\pi k|\ge {3\pi\over2}
\]
except for larger absolute values.  The \(k=0\) distance is at most
\(\pi/2\).  Hence every nonzero image divided by the \(k=0\)
image is bounded by a polynomial in \(\tau^{-1}\) times
\[
 \exp\left[
 -{(3\pi/2)^2-(\pi/2)^2\over4\tau}
 \right]
 =e^{-\pi^2/(2\tau)}.
\]
The same estimate after one derivative adds only a factor
\(C\tau^{-1}\).  Summing the geometric Gaussian tails in
\(|k|\) and weakening the exponent proves
\eqref{eq:newV31-inner-error}.  Near \(\theta=0\), pair \(k\) with
\(-k\) before dividing by the leading odd factor \(\theta\);
the quotient and its derivative obey the same exponentially small
bound.
\end{proof}

\begin{theorem}[Inner heat logarithmic slope]
\label{thm:newV31-inner-slope}
There are fixed \(R,c>0\) such that
\begin{equation}
 -\partial_\theta\log h_\tau(\theta)
 \ge c{\theta\over\tau}
 \label{eq:newV31-inner-slope}
\end{equation}
whenever
\[
 R\sqrt\tau\le\theta\le{\pi\over2},
 \qquad 0<\tau\le\tau_0.
\]
\end{theorem}

\begin{proof}
Use \Cref{lem:newV31-inner-image}:
\begin{align*}
 -\partial_\theta\log h_\tau
 &=
 {\theta\over2\tau}-{1\over\theta}
 +\cot\theta
 -\partial_\theta\log(1+\varepsilon_\tau).
\end{align*}
On \(0<\theta\le\pi/2\),
\[
 \cot\theta\ge {1\over\theta}-C\theta.
\]
The exponentially small error is at most
\(\theta/(12\tau)\) after increasing \(R\) and decreasing
\(\tau_0\).  Thus
\[
 -\partial_\theta\log h_\tau
 \ge{\theta\over2\tau}-C\theta-{\theta\over12\tau}
 \ge{\theta\over3\tau}
\]
for sufficiently small \(\tau_0\).
\end{proof}

\section{The paired antipodal images}
\label{sec:newV31-cap}

Write
\[
 \delta=\pi-\theta,
 \qquad
 a={\pi\delta\over2\tau}.
\]
The two nearest images, \(k=0,-1\), have exact sum
\begin{equation}
 N_\tau^{\rm pair}(\pi-\delta)
 =
 2e^{-(\pi^2+\delta^2)/(4\tau)}
 \left[
 \pi\sinh a-\delta\cosh a
 \right].
 \label{eq:newV31-pair}
\end{equation}

\begin{lemma}[Remote cap images are exponentially negligible]
\label{lem:newV31-cap-remote}
Uniformly for \(0\le\delta\le\pi/2\),
\begin{equation}
 N_\tau(\pi-\delta)
 =
 N_\tau^{\rm pair}(\pi-\delta)
 (1+\rho_\tau(\delta)),
 \label{eq:newV31-cap-pairing}
\end{equation}
where, after the removable value at \(\delta=0\) is taken,
\begin{equation}
 |\rho_\tau(\delta)|
 +\tau|\partial_\delta\rho_\tau(\delta)|
 \le Ce^{-c/\tau}.
 \label{eq:newV31-cap-remote}
\end{equation}
\end{lemma}

\begin{proof}
Pair every image \(k\ge1\) with \(-k-1\).  Its two distances are
at least \(3\pi-\delta\) and \(3\pi+\delta\), while the nearest pair
has distances \(\pi-\delta,\pi+\delta\).  Factoring the common odd
zero at \(\delta=0\), the ratio is bounded by a polynomial in
\(\tau^{-1}\) times \(e^{-2\pi^2/\tau}\).  Differentiation adds one
factor \(C\tau^{-1}\).  Sum the remaining Gaussian image tails and
weaken the exponent.
\end{proof}

\begin{lemma}[Positive scaled cap profile]
\label{lem:newV31-cap-profile}
Let
\[
 \phi(z):={\pi\over2}\coth\left({\pi z\over2}\right)-{1\over z},
 \qquad z>0,
\]
with \(\phi(0)=0\).  For every fixed \(M<\infty\), there is
\(c_M>0\) such that
\begin{equation}
 \phi(z)\ge c_Mz
 \qquad(0\le z\le M).
 \label{eq:newV31-phi-lower}
\end{equation}
\end{lemma}

\begin{proof}
For \(x>0\), \(x\coth x>1\), so \(\phi(z)>0\).
The Taylor expansion at zero is
\[
 \phi(z)={\pi^2\over12}z+O(z^3).
\]
Hence \(\phi(z)/z\) extends continuously and positively to
\([0,M]\).  Its minimum is a positive \(c_M\).
\end{proof}

\begin{lemma}[Very small cap slope]
\label{lem:newV31-small-cap}
For every fixed sufficiently large \(M\), there are
\(\tau_M,c_M>0\) such that
\begin{equation}
 -\partial_\theta\log h_\tau(\theta)
 \ge c_M{\pi-\theta\over\tau}
 \label{eq:newV31-small-cap}
\end{equation}
when
\[
 0<\pi-\theta\le M\tau,
 \qquad 0<\tau\le\tau_M.
\]
\end{lemma}

\begin{proof}
Set \(\delta=\tau z\), \(0\le z\le M\).  Dividing
\eqref{eq:newV31-pair} by
\(\sin\delta\), and removing factors independent of \(\delta\),
gives
\begin{equation}
 e^{-\tau z^2/4}
 {\pi\sinh(\pi z/2)-\tau z\cosh(\pi z/2)
  \over\sin(\tau z)}.
 \label{eq:newV31-scaled-cap-density}
\end{equation}
After division by its value at \(z=0\), this converges in
\(C^2([0,M])\) to
\[
 {\sinh(\pi z/2)\over(\pi/2)z}.
\]
All functions are even in \(z\).  Therefore their logarithmic
derivatives satisfy, uniformly on \([0,M]\),
\[
 \partial_\delta\log h_\tau(\pi-\delta)
 =
 {1\over\tau}\left[\phi(z)+O_M(\tau z)\right]
 +O_M(\tau z).
\]
Use \eqref{eq:newV31-phi-lower} and decrease \(\tau_M\):
\[
 \partial_\delta\log h_\tau(\pi-\delta)
 \ge {c_Mz\over2\tau}.
\]
Since
\(-\partial_\theta=\partial_\delta\) and
\(z/\tau=\delta/\tau^2\ge\delta/\tau\) for \(\tau\le1\), this
proves \eqref{eq:newV31-small-cap}.  The remote-image correction
\eqref{eq:newV31-cap-remote} is exponentially smaller and is
absorbed in the same bound; its odd derivative also vanishes at
\(\delta=0\).
\end{proof}

\begin{lemma}[Dominant-image cap slope]
\label{lem:newV31-large-cap}
For sufficiently large fixed \(M\), there are \(c,\tau_0>0\) such
that
\begin{equation}
 -\partial_\theta\log h_\tau(\theta)
 \ge c{\pi-\theta\over\tau}
 \label{eq:newV31-large-cap}
\end{equation}
when
\[
 M\tau\le\pi-\theta\le{\pi\over2},
 \qquad 0<\tau\le\tau_0.
\]
\end{lemma}

\begin{proof}
On this region the ratio of the \(k=-1\) image to the \(k=0\)
image is bounded, together with its logarithmic derivative error,
by
\[
 C\tau^{-1}e^{-\pi\delta/\tau}
 \le C\tau^{-1}e^{-\pi M}.
\]
Choose \(M\) large enough that this is absorbed relative to the
main \(1/\tau\) slope; all further images are smaller by
\Cref{lem:newV31-cap-remote}.  The leading quotient is
\[
 {\theta\over\sin\theta}e^{-\theta^2/(4\tau)}.
\]
Hence
\begin{align*}
 -\partial_\theta\log h_\tau(\theta)
 &=
 {\theta\over2\tau}-{1\over\theta}+\cot\theta
 +\hbox{absorbed error}\\
 &=
 {\pi-\delta\over2\tau}
 -{1\over\pi-\delta}-\cot\delta
 +\hbox{absorbed error}.
\end{align*}
Here \(0<\delta\le\pi/2\), so
\(\pi-\delta\ge\pi/2\) and
\(\delta\le\pi-\delta\).  If
\(\delta\ge M\tau\), then
\(\cot\delta\le1/\delta\le1/(M\tau)\).
After increasing \(M\), the negative rational terms and image error
use at most half of the leading
\((\pi-\delta)/(2\tau)\).  The remainder is at least
\[
 {c\over\tau}\ge c{\delta\over\tau}.
\]
\end{proof}

\begin{theorem}[Heat logarithmic-slope card]
\label{thm:newV31-HLSC}
The HLSC
\eqref{eq:newV30-HLSC-inner}--%
\eqref{eq:newV30-HLSC-cap} holds.
\end{theorem}

\begin{proof}
The inner inequality is
\Cref{thm:newV31-inner-slope}.  Split the antipodal half into
\(\delta\le M\tau\) and \(\delta\ge M\tau\), then apply
\Cref{lem:newV31-small-cap,lem:newV31-large-cap}.
\end{proof}

\section{Closure of the heat--Wilson path gap}
\label{sec:newV31-gap}

\begin{theorem}[Uniform gapped heat--Wilson bridge]
\label{thm:newV31-path-gap}
For the calibrated geometric bridge
\eqref{eq:newV30-geometric-bridge}, there are \(c,\alpha_0>0\)
such that
\begin{equation}
 \boxed{\qquad
 \operatorname{gap}(\nu_{\alpha,t})\ge c\alpha
 \qquad
 (0\le t\le1,\ \alpha\ge\alpha_0).
 \qquad}
 \label{eq:newV31-path-gap}
\end{equation}
\end{theorem}

\begin{proof}
\Cref{thm:newV31-HLSC} proves the only open hypothesis of
\Cref{thm:newV30-HLSC-sufficient}.  Apply that theorem.
\end{proof}

\begin{assessment}[The endpoint and path gates now coexist]
\label{ass:newV31-acquired}
V29 proves the all-order strong heat endpoint card, and V31 proves
the uniform \(c\alpha\) gap along the exact geometric
heat--Wilson path.  These two facts had not previously been
available simultaneously.  The remaining local comparison gate is
now the raw connected insertion of its exact score
\[
 S_{\alpha,t}
 =
 \alpha\cos\theta-\log h_{\tau_\alpha}(\theta)
 -\mathbb E_{\alpha,t}
 [\alpha\cos\theta-\log h_{\tau_\alpha}(\theta)].
\]
The path gap supplies a legitimate intrinsic inverse for that score
but does not by itself prove the all-arity score cumulant.
\end{assessment}

\section{V32 programme and claim boundary}
\label{sec:newV31-next}

\begin{programme}[V32: exact bridge-score profile]
\label{prog:newV31-V32}
\begin{enumerate}[leftmargin=3em]
\item Insert \eqref{eq:newV30-image} into the exact score
      \eqref{eq:newV20-radial-score}.
\item Calibrate \(\tau_\alpha\) through second order so that the
      quadratic tangent score vanishes uniformly on fixed
      representations.
\item Compute the surviving radial Hermite components.  Determine
      whether the first score root has degree four, as in the
      geodesic path, or degree six after calibration.
\item Prove adaptive \(L^p\) rows for the score and its radial
      gradient using the inner and paired-cap estimates of V31.
\item Form the complete raw score cumulant before applying the path
      inverse or an operator norm.
\item Use the V27 scalar rooted resummation and the V29
      time-ordered matrix forest as the two model endpoints for the
      all-order insertion.
\end{enumerate}
\end{programme}

\begin{firewall}[Claim boundary after V31]
\label{fire:newV31-boundary}
V31 proves the quantitative heat logarithmic-slope card by separate
shortest-image and paired-antipodal estimates and, through the V30
Hardy reduction, proves a uniform spectral gap of order
\(\alpha\) on the full geometric heat--Wilson bridge.  It does not
prove the raw bridge-score insertion card, the Wilson aggregate
strong card, all-order strong MTA, WUFC, CPM, cutoff removal,
reflection positivity, Osterwalder--Schrader reconstruction, a
continuum Yang--Mills measure, or a positive Hamiltonian gap.  The
full Yang--Mills mass gap has not been proved.
\end{firewall}

\begin{theorem}[V31 result ledger]
\label{thm:newV31-results}
V31:
\begin{enumerate}[label=\textup{(V31.\arabic*)},leftmargin=6.5em]
\item proves shortest-image logarithmic slope on the inner half;
\item derives the exact two-image antipodal profile;
\item proves exponential suppression of every remote image;
\item proves the positive scaled cap derivative
      \((\pi/2)\coth(\pi z/2)-1/z\);
\item establishes HLSC on both cap scales; and
\item proves the uniform \(c\alpha\) gap on the full calibrated
      heat--Wilson bridge.
\end{enumerate}
\end{theorem}

\begin{proof}
Use
\Cref{lem:newV31-inner-image,
thm:newV31-inner-slope,
lem:newV31-cap-remote,
lem:newV31-cap-profile,
lem:newV31-small-cap,
lem:newV31-large-cap,
thm:newV31-HLSC,
thm:newV31-path-gap}.
\end{proof}

\paragraph{Parent-version record.}
V31 was formed from
\texttt{Renewal\_Yang\_Mills\_Mass\_Gap\_Research\_Notes\_V30.tex}.
The V30 parent had \(14\,057\) logical source lines,
\(448\,795\) bytes, and SHA--256
\[
 \texttt{31117954717a44b5d82031fdcd32c861dfa2b508a7cfeef041e799626b5b95d2}.
\]
The next compulsory companion audit is V35.

% =============================================================================
% END APPENDED FIFTH-SERIES V31 MATERIAL
% =============================================================================

% =============================================================================
% BEGIN APPENDED FIFTH-SERIES V32 MATERIAL
% =============================================================================

\chapter{V32: Second-order calibration and the exact bridge score}
\label{chap:newV32-score}

\section{Calibration which removes the quadratic score}
\label{sec:newV32-calibration}

Define
\begin{equation}
 \tau_\alpha^\sharp
 :={1\over2\alpha+2/3}.
 \label{eq:newV32-sharp-time}
\end{equation}
Then
\[
 \tau_\alpha^\sharp
 ={1\over2\alpha}-{1\over6\alpha^2}
 +O(\alpha^{-3}),
\]
so it remains in the calibrated class used in V20, V29, and V31.
Let
\begin{equation}
 L_\alpha(\theta)
 :=
 \alpha\cos\theta-\log h_{\tau_\alpha^\sharp}(\theta)
 \label{eq:newV32-L}
\end{equation}
and let
\[
 S_{\alpha,t}^\sharp
 :=
 L_\alpha-\mathbb E_{\alpha,t}L_\alpha
\]
be the exact centered score of the associated geometric
heat--Wilson bridge.

\begin{lemma}[Exact inner score decomposition]
\label{lem:newV32-inner-decomposition}
On \(0\le\theta\le\pi/2\),
\begin{align}
 L_\alpha(\theta)-L_\alpha(0)
 &=
 \alpha R(\theta)
 +B(\theta)
 +E_\alpha(\theta),
 \label{eq:newV32-inner-decomposition}\\
 R(\theta)
 &:=
 {\theta^2\over2}-(1-\cos\theta),
 \label{eq:newV32-R}\\
 B(\theta)
 &:=
 {\theta^2\over6}
 -\log{\theta\over\sin\theta},
 \label{eq:newV32-B}
\end{align}
where, for every derivative order \(j=0,1\),
\begin{equation}
 |\partial_\theta^jE_\alpha(\theta)|
 \le C_j\alpha^{j/2}e^{-c\alpha}.
 \label{eq:newV32-E}
\end{equation}
In particular the coefficient of \(\theta^2\) vanishes exactly.
\end{lemma}

\begin{proof}
By \Cref{lem:newV31-inner-image},
\[
 \log h_\tau(\theta)-\log h_\tau(0)
 =
 -{\theta^2\over4\tau}
 +\log{\theta\over\sin\theta}
 +\log{1+\varepsilon_\tau(\theta)\over
             1+\varepsilon_\tau(0)}.
\]
For \(\tau=\tau_\alpha^\sharp\),
\[
 {1\over4\tau}={\alpha\over2}+{1\over6}.
\]
Insert this and
\(\alpha(\cos\theta-1)\) into
\eqref{eq:newV32-L}; the first two terms become
\(\alpha R+B\).  The last logarithm is \(E_\alpha\), and
\eqref{eq:newV32-E} follows from
\eqref{eq:newV31-inner-error} and
\(\tau_\alpha^\sharp\asymp\alpha^{-1}\).
\end{proof}

\begin{lemma}[Quartic value and cubic derivative]
\label{lem:newV32-local-profile}
On \(0\le\theta\le\pi/2\),
\begin{align}
 |R(\theta)|+|B(\theta)|
 &\le C\theta^4,
 \label{eq:newV32-local-value}\\
 |R'(\theta)|+|B'(\theta)|
 &\le C\theta^3.
 \label{eq:newV32-local-gradient}
\end{align}
Moreover,
\begin{align}
 R(\theta)
 &={\theta^4\over24}+O(\theta^6),
 \label{eq:newV32-R-series}\\
 B(\theta)
 &=-{\theta^4\over180}+O(\theta^6).
 \label{eq:newV32-B-series}
\end{align}
\end{lemma}

\begin{proof}
The Taylor series are
\[
 1-\cos\theta
 ={\theta^2\over2}-{\theta^4\over24}
 +O(\theta^6),
\]
and
\[
 \log{\theta\over\sin\theta}
 ={\theta^2\over6}+{\theta^4\over180}
 +O(\theta^6).
\]
They prove the local claims near zero.  Divide the left sides of
\eqref{eq:newV32-local-value} by \(\theta^4\), and their
derivatives by \(\theta^3\); the quotients extend continuously to
zero and are bounded on \([0,\pi/2]\).
\end{proof}

\section{Global paired-image score rows}
\label{sec:newV32-global}

\begin{theorem}[Global quartic score profile]
\label{thm:newV32-global-profile}
There is a numerical \(C\) such that, for all sufficiently large
\(\alpha\) and \(0\le\theta\le\pi\),
\begin{align}
 |L_\alpha(\theta)-L_\alpha(0)|
 &\le C\alpha\theta^4,
 \label{eq:newV32-global-value}\\
 |L_\alpha'(\theta)|
 &\le C\alpha\theta^3.
 \label{eq:newV32-global-gradient}
\end{align}
\end{theorem}

\begin{proof}
On \(0\le\theta\le\pi/2\), combine
\Cref{lem:newV32-inner-decomposition,
lem:newV32-local-profile}; the exponentially small error is
absorbed by the right side, including on
\(\theta\lesssim\alpha^{-1/2}\) by its odd endpoint
factorization.

It remains to treat \(\pi/2\le\theta\le\pi\).  Here
\(\theta^4\ge(\pi/2)^4\) and
\(\theta^3\ge(\pi/2)^3\), so it is enough to prove the crude bounds
\[
 |L_\alpha(\theta)-L_\alpha(0)|\le C\alpha,
 \qquad
 |L_\alpha'(\theta)|\le C\alpha.
\]
The value bound follows directly from the paired image formula
\eqref{eq:newV31-pair}, the remote-image estimate
\eqref{eq:newV31-cap-remote}, and
\(\tau_\alpha^\sharp\asymp\alpha^{-1}\): after the common
\(e^{-\pi^2/(4\tau)}\) factor is extracted, every remaining
logarithm is \(O(\alpha+\log\alpha)\).

For the derivative, split
\(\delta=\pi-\theta\) at \(M\tau_\alpha^\sharp\).
On \(\delta\ge M\tau_\alpha^\sharp\), the dominant-image
calculation in
\Cref{lem:newV31-large-cap} gives
\(|\partial_\theta\log h_\tau|\le C/\tau\).
On \(\delta\le M\tau\), differentiate the scaled density
\eqref{eq:newV31-scaled-cap-density}.  Its logarithmic
\(z\)-derivative is uniformly bounded on \(0\le z\le M\), so its
\(\theta\)-derivative is \(O(\tau^{-1})\); evenness removes any
endpoint singularity.  Since
\(|\alpha\sin\theta|\le\alpha\), this proves the crude gradient
bound.
\end{proof}

\begin{theorem}[Adaptive bridge-score rows]
\label{thm:newV32-adaptive-score}
For every \(p\ge1\), uniformly in \(t\in[0,1]\),
\begin{align}
 \|S_{\alpha,t}^\sharp\|_{L^p(\nu_{\alpha,t})}
 &\le {Cp^2\over\alpha},
 \label{eq:newV32-score-value}\\
 \|\nabla S_{\alpha,t}^\sharp\|_{L^p(\nu_{\alpha,t})}
 &\le {Cp^{3/2}\over\sqrt\alpha}.
 \label{eq:newV32-score-gradient}
\end{align}
\end{theorem}

\begin{proof}
Centering costs at most two in \(L^p\).  If \(p\le\alpha\), use
\Cref{thm:newV32-global-profile} and the adaptive bridge moments
\eqref{eq:newV30-bridge-moments}:
\[
 \|S_{\alpha,t}^\sharp\|_p
 \le2C\alpha\|\theta^4\|_p
 \le {Cp^2\over\alpha},
\]
\[
 \|\nabla S_{\alpha,t}^\sharp\|_p
 \le C\alpha\|\theta^3\|_p
 \le {Cp^{3/2}\over\sqrt\alpha}.
\]
If \(p>\alpha\), use the global pointwise bounds
\(|S|\le C\alpha\) and \(|\nabla S|\le C\alpha\).
The proposed right sides are then at least \(C\alpha\), after a
fixed enlargement of \(C\).
\end{proof}

\begin{corollary}[Intrinsic negative score norm]
\label{cor:newV32-score-negative}
Let \(A_{\alpha,t}\) be the reversible bridge generator and
\(R_{\alpha,t}=A_{\alpha,t}^{-1}Q\).  Then
\begin{equation}
 \|dR_{\alpha,t}S_{\alpha,t}^\sharp\|_2
 \le {C\over\alpha^{3/2}}.
 \label{eq:newV32-score-negative}
\end{equation}
\end{corollary}

\begin{proof}
The path gap
\eqref{eq:newV31-path-gap} gives
\[
 \|dR_{\alpha,t}S\|_2^2
 =
 \langle S,R_{\alpha,t}S\rangle
 \le {C\over\alpha}\|S\|_2^2.
\]
Use \eqref{eq:newV32-score-value} at \(p=2\).
\end{proof}

\section{The tangent score remains quartic}
\label{sec:newV32-tangent}

\begin{theorem}[Uniform tangent expansion]
\label{thm:newV32-tangent-score}
In scaled coordinates \(Y=\sqrt\alpha\,\theta\boldsymbol n\),
uniformly in \(t\in[0,1]\) and every fixed \(p<\infty\),
\begin{equation}
 S_{\alpha,t}^\sharp
 =
 {|Y|^4-15\over24\alpha}
 +O_{L^p}(\alpha^{-2}).
 \label{eq:newV32-tangent-score}
\end{equation}
\end{theorem}

\begin{proof}
On every fixed scaled ball,
\Cref{lem:newV32-inner-decomposition,
lem:newV32-local-profile} gives
\[
 L_\alpha(|Y|/\sqrt\alpha)-L_\alpha(0)
 =
 {|Y|^4\over24\alpha}
 +O\left({|Y|^4+|Y|^6\over\alpha^2}\right)
 +O(e^{-c\alpha}).
\]
The heat and Wilson endpoint densities have the same standard
Gaussian scaled limit under
\eqref{eq:newV32-sharp-time}; their geometric means do as well by
the uniform overlap theorem.  The moment bounds
\eqref{eq:newV30-bridge-moments} permit truncation outside a growing
scaled ball, giving
\[
 \mathbb E_{\alpha,t}|Y|^4=15+O(\alpha^{-1}).
\]
Center the preceding expansion.  Its polynomial remainder is
\(O_{L^p}(\alpha^{-2})\), proving the claim.
\end{proof}

\begin{assessment}[Calibration cannot move the first root to degree
six]
\label{ass:newV32-quartic-survives}
The exact quadratic cancellation in
\eqref{eq:newV32-sharp-time} is useful: it prevents a degree-two
score root.  It does not remove the quartic radial coefficient
\(|Y|^4/(24\alpha)\), which is forced by the Wilson action.  Thus
the bridge score has exactly the same first saturating root as V25.
The correct target remains the raw tree-scale score card, not an
extra-clock card.
\end{assessment}

\section{V33 programme and claim boundary}
\label{sec:newV32-next}

\begin{programme}[V33: time-ordered quartic-root forest]
\label{prog:newV32-V33}
\begin{enumerate}[leftmargin=3em]
\item Introduce one auxiliary source coupled to the exact centered
      score \(S_{\alpha,t}^\sharp\) in the bridge moment family.
\item Apply the BKAR forest formula with the score vertex retained
      as a distinguished root.
\item Use the V29 Duhamel simplex for all heat pair edges and the
      V27 rooted Penrose census for the degree-four tangent root.
\item Separate the exact tangent quartic from the
      \(O_{L^p}(\alpha^{-2})\) score remainder before norms, but do
      not separate the final cumulant partitions.
\item Prove a completely bounded root-stability estimate for the
      matrix quartic insertion.
\item Test the construction on the V26 spin-one and spin-two
      saturation sectors.
\end{enumerate}
\end{programme}

\begin{firewall}[Claim boundary after V32]
\label{fire:newV32-boundary}
V32 chooses the exact second-order heat time which cancels the
quadratic bridge score, proves a global quartic value/cubic
gradient profile using the paired heat images, obtains adaptive
all-\(p\) score rows and the intrinsic negative norm
\(O(\alpha^{-3/2})\), and proves the uniform quartic tangent
expansion.  It does not prove the all-arity raw bridge-score card,
the Wilson aggregate strong card, all-order strong MTA, WUFC, CPM,
cutoff removal, reflection positivity,
Osterwalder--Schrader reconstruction, a continuum Yang--Mills
measure, or a positive Hamiltonian gap.  The full Yang--Mills mass
gap has not been proved.
\end{firewall}

\begin{theorem}[V32 result ledger]
\label{thm:newV32-results}
V32:
\begin{enumerate}[label=\textup{(V32.\arabic*)},leftmargin=6.5em]
\item constructs the exact quadratic-cancelling heat calibration;
\item derives the inner score as \(\alpha R+B\) plus an
      exponentially small image correction;
\item proves global quartic value and cubic gradient profiles;
\item proves adaptive all-\(p\) bridge-score rows;
\item proves the intrinsic score negative norm
      \(O(\alpha^{-3/2})\); and
\item identifies the unavoidable centered quartic tangent root.
\end{enumerate}
\end{theorem}

\begin{proof}
Use
\Cref{lem:newV32-inner-decomposition,
lem:newV32-local-profile,
thm:newV32-global-profile,
thm:newV32-adaptive-score,
cor:newV32-score-negative,
thm:newV32-tangent-score}.
\end{proof}

\paragraph{Parent-version record.}
V32 was formed from
\texttt{Renewal\_Yang\_Mills\_Mass\_Gap\_Research\_Notes\_V31.tex}.
The V31 parent had \(14\,478\) logical source lines,
\(460\,841\) bytes, and SHA--256
\[
 \texttt{2da8e4ac7888160d0ef1cde637cf75a124b23050a10da70f3cc30d7e6292a307}.
\]
The next compulsory companion audit is V35.

% =============================================================================
% END APPENDED FIFTH-SERIES V32 MATERIAL
% =============================================================================

% =============================================================================
% BEGIN APPENDED FIFTH-SERIES V33 MATERIAL
% =============================================================================

\chapter{V33: A block-local quartic root and its exact quantum forest}
\label{chap:newV33-block-root}

\section{Purpose and correction of the naive source}
\label{sec:newV33-purpose}

V32 reduced the exact geometric heat--Wilson score to a centered
radial quartic at tangent order.  The next task is not merely to
bound a polynomial of the total Casimir.  A single score belongs to
one block of every moment partition.  Consequently its Fourier
multiplier must be block-additive at partition endpoints.  The
global square of the interpolating Casimir is not block-additive:
it contains products belonging to two different blocks.  This
chapter constructs the missing block-local gate, proves an exact
time-ordered marked-forest identity for the resulting heat spectral
quartic score, and obtains its all-arity raw tree card.

The construction is an exactly solvable model of the tangent score.
It is not yet the exact bridge score
\(S_{\alpha,t}^\sharp\).  The distinction is maintained throughout.

\section{The full interpolating Casimir is positive}
\label{sec:newV33-positive}

For the representations \(R_1,\ldots,R_m\), retain the V29
notation and put
\begin{equation}
 {\cal K}(\boldsymbol x)
 :=
 s\sum_{i=1}^m C_i
 +2s\sum_{i<j}x_{ij}\Omega_{ij}.
 \label{eq:newV33-full-K}
\end{equation}
Thus the V29 heat interpolation is
\[
 F_m(\boldsymbol x)=e^{-{\cal K}(\boldsymbol x)}.
\]
For \(A\subseteq[m]\), write
\[
 {\cal K}_A:=sC_A^{\rm tot}.
\]

\begin{theorem}[Positive layer cake for the full Casimir]
\label{thm:newV33-positive-layer}
At every tree or forest minimum matrix
\(\boldsymbol x^T(\boldsymbol u)\),
\begin{equation}
 {\cal K}(\boldsymbol x^T(\boldsymbol u))
 =
 \int_0^1
 \sum_{A\in{\cal C}_r}{\cal K}_A\,dr
 \ge0,
 \label{eq:newV33-positive-layer}
\end{equation}
where \({\cal C}_r\) is the threshold-component partition of
\Cref{lem:newV29-threshold-layer}.  In particular, every
semigroup interval in a Duhamel word is a contraction:
\begin{equation}
 \left\|e^{-z\tau{\cal K}(\boldsymbol x^T)}\right\|_{\rm op}
 \le1
 \qquad(\tau\ge0,\ \Re z\ge0).
 \label{eq:newV33-contraction}
\end{equation}
Both assertions remain true after arbitrary matrix amplification.
\end{theorem}

\begin{proof}
For each threshold partition,
\[
 \sum_{A\in{\cal C}_r}{\cal K}_A
 =
 s\sum_iC_i+
 2s\sum_{\substack{i<j\\
 i,j\ {\rm in\ the\ same\ threshold\ block}}}\Omega_{ij}.
\]
Integrating in \(r\), the diagonal term stays
\(s\sum_iC_i\), while
\Cref{lem:newV29-threshold-layer} turns every off-diagonal
indicator into \(x_{ij}^T(\boldsymbol u)\).  This proves the
identity.  Every \({\cal K}_A=sC_A^{\rm tot}\) is positive
semidefinite, hence so is its Bochner integral.  Functional
calculus proves \eqref{eq:newV33-contraction}.  Tensoring all
operators with an identity leaves positivity and the norm
unchanged.
\end{proof}

\begin{remark}[Strict strengthening of the V29 stability input]
\label{rem:newV33-strengthening}
V29 bounded the pair Hamiltonian below by \(-mI\) and consequently
paid \(e^m\) for a complete Duhamel word.  Once the scalar self
Casimirs are kept inside the exponential, the same threshold
argument gives the positivity
\eqref{eq:newV33-positive-layer}; the \(e^m\) loss disappears.
No commutation of overlapping pair insertions is used.
\end{remark}

Let
\[
 V_{ij}:=\partial_{x_{ij}}{\cal K}=2s\Omega_{ij},
 \qquad
 w_{ij}:=b_ib_j.
\]
The V28 pair-row bound gives
\begin{equation}
 \|V_{ij}\|_{\rm cb}\le2w_{ij}.
 \label{eq:newV33-pair-cb}
\end{equation}
Here and below the completely bounded norm means the supremum after
tensoring the displayed operator with an arbitrary identity.

\begin{lemma}[Contractive mixed heat derivative]
\label{lem:newV33-contract-derivative}
At a forest minimum matrix, for every set \(E\) of distinct pair
variables,
\begin{equation}
 \left\|\partial_Ee^{-{\cal K}}\right\|_{\rm cb}
 \le
 \prod_{e\in E}2w_e.
 \label{eq:newV33-contract-derivative}
\end{equation}
\end{lemma}

\begin{proof}
Apply the exact V29 Duhamel simplex to \({\cal K}\).  By
\eqref{eq:newV33-contraction}, all exponential intervals have norm
at most one.  The \(|E|!\) ordered words cancel the simplex volume
\(1/|E|!\), and \eqref{eq:newV33-pair-cb} bounds the insertions.
The same argument holds at every amplification.
\end{proof}

\begin{lemma}[Completely bounded spectral insertion]
\label{lem:newV33-spectral-insertion}
For every fixed integer \(q\ge0\),
\begin{equation}
 \left\|
 \partial_E\left({\cal K}^qe^{-{\cal K}}\right)
 \right\|_{\rm cb}
 \le
 q!\,2^q\left({3\over2}\right)^{|E|}
 \prod_{e\in E}2w_e
 \label{eq:newV33-spectral-insertion}
\end{equation}
at every forest minimum matrix.
\end{lemma}

\begin{proof}
For complex \(z\) with \(\Re z>0\),
\[
 {\cal K}^qe^{-{\cal K}}
 =
 (-\partial_z)^qe^{-z{\cal K}}\big|_{z=1}.
\]
The mixed Duhamel formula and
\eqref{eq:newV33-contraction} give
\[
 \|\partial_Ee^{-z{\cal K}}\|_{\rm cb}
 \le |z|^{|E|}\prod_{e\in E}2w_e.
\]
Apply the Banach-valued Cauchy estimate on
\(|z-1|=1/2\).  On this circle
\(\Re z\ge1/2\) and \(|z|\le3/2\), while the Cauchy radius
contributes \(2^qq!\).  This is
\eqref{eq:newV33-spectral-insertion}.
\end{proof}

\section{A partition-local Casimir square}
\label{sec:newV33-local-square}

For ordered \(i,j\), define
\begin{equation}
 A_{ij}:=s\sum_{a=1}^3J_i^aJ_j^a.
 \label{eq:newV33-Aij}
\end{equation}
Thus \(A_{ii}=sC_i\), \(A_{ij}=s\Omega_{ij}\) for
\(i\ne j\), and
\begin{equation}
 \|A_{ij}\|_{\rm cb}\le b_ib_j.
 \label{eq:newV33-Aij-bound}
\end{equation}

For a nonempty set \(S\subseteq[m]\), let \(r(S)=\min S\) and
let
\[
 {\cal A}(S):=\{\{r(S),v\}:v\in S\setminus\{r(S)\}\}
\]
be its canonical star.  Put
\begin{equation}
 \chi_S(\boldsymbol x)
 :=
 \prod_{e\in{\cal A}(S)}x_e,
 \qquad
 \chi_{\{i\}}:=1.
 \label{eq:newV33-gate}
\end{equation}
Define the block-local square
\begin{equation}
 {\cal Q}_m(\boldsymbol x)
 :=
 \sum_{i,j,k,l=1}^m
 \chi_{\{i,j,k,l\}}(\boldsymbol x)
 A_{ij}A_{kl}.
 \label{eq:newV33-Q}
\end{equation}
The particular canonical star is only a compiler gauge.  Its role
is to recognize whether all labels in one quartic monomial belong
to one endpoint block.

\begin{lemma}[Exact endpoint localization]
\label{lem:newV33-Q-endpoint}
If \(\rho\) is a partition of \([m]\), then
\begin{equation}
 {\cal Q}_m(\boldsymbol x^\rho)
 =
 \sum_{B\in\rho}{\cal K}_B^2.
 \label{eq:newV33-Q-endpoint}
\end{equation}
\end{lemma}

\begin{proof}
At a partition matrix,
\(\chi_S(\boldsymbol x^\rho)=1\) exactly when all vertices of
\(S\) lie in one block of \(\rho\), and it is zero otherwise.
Thus \eqref{eq:newV33-Q} retains exactly the ordered quadruples
\(i,j,k,l\in B\), separately for every \(B\).  Their sum is
\[
 \left(\sum_{i,j\in B}A_{ij}\right)^2
 ={\cal K}_B^2.
\]
Sum over the blocks.
\end{proof}

\begin{proposition}[Why the ungated global square is wrong]
\label{prop:newV33-naive-square}
At a partition endpoint,
\begin{equation}
 {\cal K}(\boldsymbol x^\rho)^2
 =
 \sum_{B\in\rho}{\cal K}_B^2
 +2\sum_{\substack{B,C\in\rho\\B<C}}{\cal K}_B{\cal K}_C.
 \label{eq:newV33-naive-square}
\end{equation}
Consequently a source coupled to the global
\({\cal K}^2\) marks either one block or two distinct blocks.  It
cannot yield the cumulant with one quartic score root.
\end{proposition}

\begin{proof}
The block Casimirs act on disjoint tensor factors and commute.
Square their sum.  The second sum in
\eqref{eq:newV33-naive-square} is precisely the product of two
degree-two marked moments.  Such terms belong to two score roots,
whereas the desired joint cumulant contains one score variable.
\end{proof}

\section{Exact heat spectral tangent score}
\label{sec:newV33-spectral-score}

Let \(h_s\) be the positive central heat density whose Fourier
multiplier is \(e^{-sC_2(R)}\).  Define the smooth central signed
density \(n_s^{\rm sp}\) by the Fourier multipliers
\begin{equation}
 \widehat n_s^{\rm sp}(R)
 =
 {s\over3}
 \left[
 (sC_2(R))^2-5sC_2(R)
 \right]e^{-sC_2(R)}.
 \label{eq:newV33-n-multiplier}
\end{equation}
Exponential decay in \(C_2(R)\) makes this a smooth central
function.  Its trivial multiplier is zero, so its integral is
zero.  Hence
\begin{equation}
 S_s^{\rm sp}(U):={n_s^{\rm sp}(U)\over h_s(U)}
 \label{eq:newV33-spectral-score}
\end{equation}
is a centered smooth observable under the heat law.

For a block \(B\), its marked tensor moment is
\begin{equation}
 N_B^{\rm sp}
 :=
 \mathbb E_{\mu_s^H}
 \left[
 S_s^{\rm sp}(U)\bigotimes_{i\in B}D^{R_i}(U)
 \right]
 =
 {s\over3}({\cal K}_B^2-5{\cal K}_B)e^{-{\cal K}_B}.
 \label{eq:newV33-marked-moment}
\end{equation}

\begin{lemma}[Gaussian origin of the multiplier]
\label{lem:newV33-Gaussian-origin}
If \(Y\sim N(0,I_3)\), \(r=|\xi|^2\), and
\(Q_4(Y)=|Y|^4-15\), then
\begin{equation}
 \mathbb E[Q_4(Y)e^{i\xi\cdot Y}]
 =
 (r^2-10r)e^{-r/2}.
 \label{eq:newV33-Gaussian-Fourier}
\end{equation}
If \(r=2{\cal K}\) and
\(\alpha_s=(2s)^{-1}\), multiplication by
\((24\alpha_s)^{-1}\) gives exactly
\[
 {s\over3}({\cal K}^2-5{\cal K})e^{-{\cal K}}.
\]
Thus \eqref{eq:newV33-n-multiplier} is the exact spectral
counterpart of the leading V32 tangent score.
\end{lemma}

\begin{proof}
Apply \((-\Delta_\xi)^2\) to
\(\mathbb E e^{i\xi\cdot Y}=e^{-r/2}\).  In dimension three,
\[
 \mathbb E[|Y|^4e^{i\xi\cdot Y}]
 =
 (r^2-10r+15)e^{-r/2}.
\]
Subtract \(15e^{-r/2}\).  Since
\((24\alpha_s)^{-1}=s/12\), substituting \(r=2{\cal K}\)
gives the last assertion.
\end{proof}

\section{The exact one-root quantum forest}
\label{sec:newV33-root-forest}

Define
\begin{equation}
 {\cal G}_m(\boldsymbol x)
 :=
 {s\over3}
 \left({\cal Q}_m(\boldsymbol x)-5{\cal K}(\boldsymbol x)\right)
 e^{-{\cal K}(\boldsymbol x)}.
 \label{eq:newV33-G}
\end{equation}

\begin{lemma}[Dual-number endpoint factorization]
\label{lem:newV33-dual-factor}
Let \(\varepsilon^2=0\).  At every partition endpoint,
\begin{equation}
 F_m(\boldsymbol x^\rho)
 +\varepsilon{\cal G}_m(\boldsymbol x^\rho)
 =
 \bigotimes_{B\in\rho}
 \left(M_B^H+\varepsilon N_B^{\rm sp}\right)
 \quad\pmod{\varepsilon^2}.
 \label{eq:newV33-dual-factor}
\end{equation}
\end{lemma}

\begin{proof}
Use the unmarked factorization
\eqref{eq:newV29-partition-endpoints}.  By
\Cref{lem:newV33-Q-endpoint},
\[
 {\cal G}_m(\boldsymbol x^\rho)
 =
 \sum_{B\in\rho}
 N_B^{\rm sp}
 \bigotimes_{C\in\rho\setminus\{B\}}M_C^H.
 \]
This is exactly the coefficient of \(\varepsilon\) in the product
on the right of \eqref{eq:newV33-dual-factor}.
\end{proof}

\begin{theorem}[Exact block-local quartic-root forest]
\label{thm:newV33-root-forest}
For \(m\ge2\),
\begin{align}
 &\kappa_{m+1}^{\mu_s^H}
 \left(
 D^{R_1},\ldots,D^{R_m},S_s^{\rm sp}
 \right)
 \notag\\
 &\qquad=
 \sum_{T\in\mathfrak T_m}
 \int_{[0,1]^{m-1}}
 \partial_{E(T)}
 {\cal G}_m(\boldsymbol x^T(\boldsymbol u))
 \,d\boldsymbol u.
 \label{eq:newV33-root-forest}
\end{align}
This is an identity of complete tensor operators.
\end{theorem}

\begin{proof}
Apply the finite multivariable BKAR forest identity in the
dual-number extension of the tensor-operator algebra to
\(F_m+\varepsilon{\cal G}_m\).  The endpoint factorization
\eqref{eq:newV33-dual-factor} makes its connected component the
ordinary cumulant of the perturbed moments
\(M_B^H+\varepsilon N_B^{\rm sp}\).  Extracting the
\(\varepsilon\) coefficient in the partition formula gives
\[
 \sum_{\rho\in\Pi([m])}
 (-1)^{|\rho|-1}(|\rho|-1)!
 \sum_{B\in\rho}
 N_B^{\rm sp}
 \bigotimes_{C\ne B}M_C^H.
\]
Because \(\mathbb E_{\mu_s^H}S_s^{\rm sp}=0\), this is exactly the
joint cumulant on the left of
\eqref{eq:newV33-root-forest}: in a nonzero partition of
\(\{0,1,\ldots,m\}\), the score label \(0\) joins one nonempty
input block \(B\).  The connected forests on the input labels are
the spanning trees, and the coefficient of \(\varepsilon\) in
their integrands is
\(\partial_{E(T)}{\cal G}_m\).
\end{proof}

\section{Root payment and the strong spectral-score card}
\label{sec:newV33-payment}

\begin{lemma}[Canonical-star root payment]
\label{lem:newV33-star-payment}
Let \(\boldsymbol i=(i_1,i_2,i_3,i_4)\),
\(S=\{i_1,i_2,i_3,i_4\}\), and assume
\(b_i^2=s(1+C_2(R_i))<1\).  Then
\begin{equation}
 s\prod_{\nu=1}^4b_{i_\nu}
 \le
 \prod_{e\in{\cal A}(S)}w_e.
 \label{eq:newV33-star-payment}
\end{equation}
Consequently, for every edge set \(A\),
\begin{equation}
 s\|\partial_A{\cal Q}_m(\boldsymbol x)\|_{\rm cb}
 \le
 m^4\prod_{e\in A}w_e
 \label{eq:newV33-Q-derivative}
\end{equation}
on \(0\le x_e\le1\); the left side is zero unless
\(|A|\le3\).
\end{lemma}

\begin{proof}
Let \(q=|S|\) and \(r=\min S\).  Removing repeated factors, whose
marks are at most one, gives
\[
 s\prod_{\nu=1}^4b_{i_\nu}
 \le s\prod_{v\in S}b_v.
\]
The star weight is
\[
 \prod_{e\in{\cal A}(S)}w_e
 =
 b_r^{q-2}\prod_{v\in S}b_v.
\]
For \(q=1,2,3,4\), respectively, it is enough that
\(s\le1,1,b_r,b_r^2\).  All four follow from
\(s\le b_r^2<1\).

Now differentiate \eqref{eq:newV33-Q}.  A gate contributes only
when \(A\subseteq{\cal A}(S)\), and its remaining monomial has
absolute value at most one.  By
\eqref{eq:newV33-Aij-bound} and
\eqref{eq:newV33-star-payment}, the corresponding operator term,
including the exterior \(s\), is bounded by the full star weight.
That weight is at most \(\prod_{e\in A}w_e\), since all omitted
edge weights are below one.  There are \(m^4\) ordered quadruples.
Every gate has degree at most three.
\end{proof}

\begin{theorem}[Fixed-tree completely bounded root card]
\label{thm:newV33-fixed-tree}
There is a numerical \(K_0\) such that, for every labelled tree
\(T\) on \([m]\),
\begin{equation}
 \sup_{\boldsymbol u\in[0,1]^{m-1}}
 \left\|
 \partial_{E(T)}{\cal G}_m
 (\boldsymbol x^T(\boldsymbol u))
 \right\|_{\rm cb}
 \le
 K_0^m\prod_{i=1}^m b_i^{d_i(T)}.
 \label{eq:newV33-fixed-tree}
\end{equation}
The constant is uniform in \(m,s\), all thermal representations,
and all tensor dimensions.
\end{theorem}

\begin{proof}
Write \(k=m-1\).  For the \({\cal Q}_m\) part, Leibniz gives
\[
 \partial_{E(T)}(s{\cal Q}_me^{-{\cal K}})
 =
 \sum_{A\subseteq E(T)}
 s(\partial_A{\cal Q}_m)
 \partial_{E(T)\setminus A}e^{-{\cal K}}.
\]
Only \(|A|\le3\) contributes.  Apply
\Cref{lem:newV33-star-payment,
lem:newV33-contract-derivative}:
\[
 \left\|
 s(\partial_A{\cal Q}_m)
 \partial_{E(T)\setminus A}e^{-{\cal K}}
 \right\|_{\rm cb}
 \le
 m^4\,2^{k-|A|}
 \prod_{e\in E(T)}w_e.
\]
There are at most
\(\sum_{a=0}^3\binom{k}{a}\le(1+k)^3\) such subsets.
The polynomial \(m^4(1+m)^3\) and \(2^k\) are bounded by
an exponential \(K_1^m\).

For the linear term, use
\Cref{lem:newV33-spectral-insertion} with \(q=1\), multiply by
\(5s/3\), and use \(s<1\).  It is bounded by
\(K_2^m\prod_{e\in E(T)}w_e\).
Finally,
\[
 \prod_{e\in E(T)}w_e
 =
 \prod_i b_i^{d_i(T)}.
\]
Absorb the fixed factor \(1/3\) and combine the two estimates.
\end{proof}

\begin{theorem}[All-arity heat spectral RGSC]
\label{thm:newV33-spectral-RGSC}
Under \(b_i^2=s(1+C_2(R_i))<1\),
\begin{align}
 &\left\|
 \kappa_{m+1}^{\mu_s^H}
 (D^{R_1},\ldots,D^{R_m},S_s^{\rm sp})
 \right\|_{\rm cb}
 \notag\\
 &\qquad\le
 K^m
 \left(\prod_{i=1}^m b_i\right)
 \left(\sum_{i=1}^m b_i\right)^{m-2}.
 \label{eq:newV33-spectral-RGSC}
\end{align}
\end{theorem}

\begin{proof}
Insert \eqref{eq:newV33-fixed-tree} into the exact forest
\eqref{eq:newV33-root-forest}, integrate the unit parameter cubes,
and sum trees.  The weighted Pr\"ufer identity used in V29 gives
\[
 \sum_{T\in\mathfrak T_m}\prod_i b_i^{d_i(T)}
 =
 \left(\prod_i b_i\right)
 \left(\sum_i b_i\right)^{m-2}.
\]
\end{proof}

\begin{assessment}[What the block gate accomplishes]
\label{ass:newV33-block-gate}
The derivatives which hit \(\chi_S\) connect up to four input
labels to one score root without inserting extra pair Hamiltonians.
Their missing pair marks are paid exactly by the exterior \(s\)
and the four marks carried by \(A_{ij}A_{kl}\).  This is the
operator analogue of the V26 degree-four root census.  By contrast,
differentiating the ungated global \({\cal K}^2\) forces ordinary
pair insertions and simultaneously retains the two-root
contamination \eqref{eq:newV33-naive-square}.
\end{assessment}

\section{The nontrivial saturation sectors survive}
\label{sec:newV33-sectors}

Let all four inputs be fundamental, let
\({\cal T}_4\) be the V26 pairing tensor, and set
\(\alpha_s=(2s)^{-1}\).

\begin{theorem}[Four-fundamental saturation of the spectral model]
\label{thm:newV33-sector-test}
As \(s\downarrow0\),
\begin{equation}
 \kappa_5^{\mu_s^H}
 (U,U,U,U,S_s^{\rm sp})
 =
 {1\over3\alpha_s^3}{\cal T}_4
 +O_{\rm op}(\alpha_s^{-4}).
 \label{eq:newV33-sector-test}
\end{equation}
Consequently,
\begin{align}
 \kappa_5^{\mu_s^H}
 (U,U,U,U,S_s^{\rm sp})
 &=
 {\alpha_s^{-3}}
 \left(5P_0-{5\over3}P_1+P_2\right)
 +O_{\rm op}(\alpha_s^{-4}).
 \label{eq:newV33-sector-spectrum}
\end{align}
In particular, both the spin-one and spin-two components survive.
For \(s=\tau_\alpha^\sharp\),
\(\alpha_s=\alpha+1/3\), so the leading term agrees with the exact
bridge-score sector expansion through order \(\alpha^{-3}\).
\end{theorem}

\begin{proof}
For fixed representations, all moment operators
\(M_B^H=e^{-sC_B^{\rm tot}}\) and
\eqref{eq:newV33-marked-moment} are analytic in \(s\).
Expand the finite partition formula for the joint cumulant through
order \(s^3\).  Its connected coefficient is the polarization of
the Gaussian multiplier
\eqref{eq:newV33-Gaussian-Fourier}.  Equivalently, it is the Wick
contraction of four fundamental linear tangent terms with
\[
 {Q_4(Y)\over24\alpha_s}.
\]
The exact contraction was evaluated in V26 and equals
\({\cal T}_4/(3\alpha_s^3)\).  Since the representation space and
the partition sum are finite, Taylor's theorem gives an operator
remainder \(O(s^4)=O(\alpha_s^{-4})\).

The V26 diagonalization
\({\cal T}_4=15P_0-5P_1+3P_2\) proves
\eqref{eq:newV33-sector-spectrum}.  Finally,
\eqref{eq:newV32-sharp-time} gives
\(\alpha_s=(2\tau_\alpha^\sharp)^{-1}=\alpha+1/3\).
The V32 tangent expansion and fixed-moment convergence show that
the exact bridge score has the same order-\(\alpha^{-3}\)
contraction.
\end{proof}

\begin{corollary}[The spectral card is sharp at ordinary tree scale]
\label{cor:newV33-sharp}
For four fixed fundamental inputs, the left side of
\eqref{eq:newV33-spectral-RGSC} is bounded below by
\(c\alpha_s^{-3}\), while its Pr\"ufer weight is comparable to
\(\alpha_s^{-3}\).  Thus the exact marked-forest estimate has no
spare uniform heat clock for a separated \(O(\alpha)\) payer.
\end{corollary}

\begin{proof}
Use either nonzero sector in
\eqref{eq:newV33-sector-spectrum}.  Since
\(b_{\rm f}\asymp\alpha_s^{-1/2}\), the four-input aggregate
weight is \(b_{\rm f}^4(4b_{\rm f})^2\asymp\alpha_s^{-3}\).
\end{proof}

\section{What remains for the exact bridge}
\label{sec:newV33-remaining}

\begin{assessment}[Acquired theorem and live transfer gate]
\label{ass:newV33-live-gate}
V33 proves the complete all-arity raw score card for the heat
spectral tangent score \(S_s^{\rm sp}\).  V32 proves only the
fixed-\(p\) approximation
\[
 S_{\alpha,t}^\sharp
 =
 {Q_4(Y)\over24\alpha}+O_{L^p}(\alpha^{-2}).
\]
Fixed-\(p\) control cannot be inserted at arity \(m\) with
\(p\) fixed: the required H\"older exponent grows with \(m\).
Therefore
\Cref{thm:newV33-spectral-RGSC} does not by itself prove the exact
bridge RGSC.  The remaining local gate is now precise: construct
an endpoint-factorizing block-local interpolation for the exact
marked moments whose difference from
\({\cal G}_m\) has an exponential-grade tree derivative bound.
The comparison must be made after block localization, not by
subtracting score variables inside separate cumulant partitions.
\end{assessment}

\begin{programme}[V34: exact-score block-local remainder]
\label{prog:newV33-V34}
\begin{enumerate}[leftmargin=3em]
\item Express the exact bridge marked multiplier on each tensor
      block as \(N_B^{\rm sp}+{\cal R}_{\alpha,t,B}\), retaining
      centering at the block level.
\item Use the global quartic/cubic score rows of V32 and the bridge
      gap of V31 to seek a completely bounded estimate for the
      block remainder, uniform in the total block representation.
\item Build a dual-number interpolation whose partition endpoints
      are
      \(M_{\alpha,t,B}+\varepsilon
      N_{\alpha,t,B}\), using the V33 canonical gate for every
      quartic leading monomial.
\item Differentiate the assembled remainder only after its
      endpoint blocks have been formed; do not apply an
      arity-dependent density-product H\"older estimate.
\item If the exact bridge lacks a usable Fourier semigroup,
      combine the V14 intrinsic random-cluster interpolation with
      the V33 root-payment lemma and isolate the first conditional
      merge estimate that is not already supplied by the V31 gap.
\item At V35, perform the compulsory read-only SM/NS companion
      audit before choosing the next direction.
\end{enumerate}
\end{programme}

\begin{firewall}[Claim boundary after V33]
\label{fire:newV33-boundary}
V33 proves positivity of the full Casimir at every forest minimum,
a completely bounded spectral-insertion estimate, an exact
block-local one-root forest identity, the all-arity raw score card
for the heat spectral tangent score, and its nonzero spin-one and
spin-two saturation.  It does not prove the corresponding card for
the exact geometric heat--Wilson bridge score.  Hence it does not
yet prove the Wilson aggregate strong card, all-order strong MTA,
WUFC, CPM, cutoff removal, reflection positivity,
Osterwalder--Schrader reconstruction, a continuum Yang--Mills
measure, or a positive Hamiltonian gap.  The full Yang--Mills mass
gap has not been proved.
\end{firewall}

\begin{theorem}[V33 result ledger]
\label{thm:newV33-results}
V33:
\begin{enumerate}[label=\textup{(V33.\arabic*)},leftmargin=6.5em]
\item strengthens V29 stability to positivity of the full
      interpolating Casimir;
\item proves contractive and completely bounded time-ordered
      derivative estimates;
\item constructs a canonical partition-local Casimir square;
\item proves that the naive global square contains an inadmissible
      two-root term;
\item constructs the centered heat spectral tangent score;
\item proves its exact dual-number marked-forest identity;
\item proves its uniform all-arity raw score card; and
\item verifies sharp nonzero spin-one and spin-two saturation.
\end{enumerate}
\end{theorem}

\begin{proof}
Use
\Cref{thm:newV33-positive-layer,
lem:newV33-contract-derivative,
lem:newV33-spectral-insertion,
lem:newV33-Q-endpoint,
prop:newV33-naive-square,
lem:newV33-Gaussian-origin,
thm:newV33-root-forest,
lem:newV33-star-payment,
thm:newV33-fixed-tree,
thm:newV33-spectral-RGSC,
thm:newV33-sector-test}.
\end{proof}

\paragraph{Parent-version record.}
V33 was formed from
\texttt{Renewal\_Yang\_Mills\_Mass\_Gap\_Research\_Notes\_V32.tex}.
The V32 parent had \(14\,849\) logical source lines,
\(471\,910\) bytes, and SHA--256
\[
 \texttt{a98720fac11d309a3dca050b2613a2a61f0411329f4ee61253eff7558719a47b}.
\]
The next compulsory companion audit remains V35.

% =============================================================================
% END APPENDED FIFTH-SERIES V33 MATERIAL
% =============================================================================

% =============================================================================
% BEGIN APPENDED FIFTH-SERIES V34 MATERIAL
% =============================================================================

\chapter{V34: The finite spectral-jet compiler and the resummation boundary}
\label{chap:newV34-jet}

\section{Why the next step stays with block-local gates}
\label{sec:newV34-direction}

The intrinsic Bernoulli interpolation of V14 applies to the exact
bridge law, but V15 proves that its individual raw forest trees
contain lower-cumulant leakage and fail the required clock count
already at arity three.  Replacing the unmarked law by a
dual-number marked law does not remove that leakage.  Thus the V34
attack does not return to a raw-tree PCMC.

Instead, this chapter determines how far the V33 block-local
construction extends.  Every finite polynomial spectral jet admits
an exact one-root forest and the strong aggregate card.  The proof
tracks the root degree explicitly.  It then shows that estimating an
infinite analytic jet degree by degree creates a superexponential
arity majorant.  The exact bridge therefore requires a resummed
root generating function before norms, not merely more Taylor
coefficients.

\section{Block-local powers of arbitrary finite degree}
\label{sec:newV34-powers}

Retain \(A_{ij}\), \({\cal K}\), \(b_i\), \(w_{ij}\), and the
canonical gate \(\chi_S\) from V33.  For an integer \(q\ge1\),
define
\begin{equation}
 {\cal Q}_m^{[q]}(\boldsymbol x)
 :=
 \sum_{i_1,j_1,\ldots,i_q,j_q=1}^m
 \chi_{\{i_1,j_1,\ldots,i_q,j_q\}}(\boldsymbol x)
 A_{i_1j_1}\cdots A_{i_qj_q}.
 \label{eq:newV34-Qq}
\end{equation}
Thus \({\cal Q}_m^{[2]}={\cal Q}_m\).  The product order in
\eqref{eq:newV34-Qq} is the displayed order; no commutation of
overlapping pair operators is assumed.

\begin{lemma}[Endpoint localization of every power]
\label{lem:newV34-Qq-endpoint}
For every partition \(\rho\in\Pi([m])\),
\begin{equation}
 {\cal Q}_m^{[q]}(\boldsymbol x^\rho)
 =
 \sum_{B\in\rho}{\cal K}_B^q.
 \label{eq:newV34-Qq-endpoint}
\end{equation}
\end{lemma}

\begin{proof}
At \(\boldsymbol x^\rho\), the gate in one summand is one exactly
when all \(2q\) displayed indices lie in the same block \(B\), and
is zero otherwise.  The retained ordered sum is
\[
 \sum_{i_1,j_1,\ldots,i_q,j_q\in B}
 A_{i_1j_1}\cdots A_{i_qj_q}
 =
 \left(\sum_{i,j\in B}A_{ij}\right)^q
 =
 {\cal K}_B^q.
\]
Sum over \(B\).
\end{proof}

\begin{lemma}[Degree-\(q\) root payment]
\label{lem:newV34-q-payment}
Let
\(\boldsymbol i=(i_1,j_1,\ldots,i_q,j_q)\), let \(S\) be its
support, and assume
\[
 b_i^2=s(1+C_2(R_i))<1.
\]
Then
\begin{equation}
 s^{q-1}\prod_{\nu=1}^q b_{i_\nu}b_{j_\nu}
 \le
 \prod_{e\in{\cal A}(S)}w_e.
 \label{eq:newV34-q-payment}
\end{equation}
Consequently, for every edge set \(D\),
\begin{equation}
 s^{q-1}
 \|\partial_D{\cal Q}_m^{[q]}(\boldsymbol x)\|_{\rm cb}
 \le
 m^{2q}\prod_{e\in D}w_e
 \label{eq:newV34-Qq-derivative}
\end{equation}
on the edge cube.  The derivative vanishes when
\(|D|>2q-1\).
\end{lemma}

\begin{proof}
Let \(r=|S|\) and \(v=\min S\).  Repeated indices contribute
additional marks at most one, so the left side of
\eqref{eq:newV34-q-payment} is at most
\[
 s^{q-1}\prod_{a\in S}b_a.
\]
For \(r\ge2\), the canonical star has weight
\[
 \prod_{e\in{\cal A}(S)}w_e
 =
 b_v^{r-2}\prod_{a\in S}b_a.
\]
Since \(r\le2q\) and \(0<b_v<1\),
\[
 s^{q-1}
 \le b_v^{2q-2}
 \le b_v^{r-2}.
\]
For \(r=1\), the star is empty and the desired inequality follows
from \(s<1\) and \(b_v<1\).

A gate has degree \(r-1\le2q-1\).  If its derivative by \(D\)
does not vanish, then \(D\subseteq{\cal A}(S)\); the residual
monomial is bounded by one.  Equations
\eqref{eq:newV33-Aij-bound} and
\eqref{eq:newV34-q-payment} bound the operator product by the full
star weight.  Deleting the undifferentiated star edges can only
increase this weight, so it is at most
\(\prod_{e\in D}w_e\).  Finally there are \(m^{2q}\) ordered
index lists.
\end{proof}

\section{Exact forests for finite spectral jets}
\label{sec:newV34-finite-jets}

Fix \(d\ge1\) and coefficients
\(a_1(s),\ldots,a_d(s)\).  Assume
\begin{equation}
 A_d:=\sup_{0<s<1}\max_{1\le q\le d}|a_q(s)|<\infty.
 \label{eq:newV34-coefficient-bound}
\end{equation}
Define a smooth centered signed heat density by its multipliers
\begin{equation}
 \widehat n_s^{[d]}(R)
 :=
 e^{-K_R}
 \sum_{q=1}^d a_q(s)s^{q-1}K_R^q,
 \qquad
 K_R=sC_2(R),
 \label{eq:newV34-jet-multiplier}
\end{equation}
and put \(S_s^{[d]}=n_s^{[d]}/h_s\).
The trivial multiplier is zero, so
\(\mathbb E_{\mu_s^H}S_s^{[d]}=0\).

For tensor blocks,
\begin{equation}
 N_B^{[d]}
 =
 e^{-{\cal K}_B}
 \sum_{q=1}^d a_q(s)s^{q-1}{\cal K}_B^q.
 \label{eq:newV34-jet-block}
\end{equation}
Set
\begin{equation}
 {\cal G}_m^{[d]}(\boldsymbol x)
 :=
 \left[
 \sum_{q=1}^d
 a_q(s)s^{q-1}{\cal Q}_m^{[q]}(\boldsymbol x)
 \right]
 e^{-{\cal K}(\boldsymbol x)}.
 \label{eq:newV34-Gd}
\end{equation}

\begin{theorem}[Exact finite-jet one-root forest]
\label{thm:newV34-jet-forest}
For \(m\ge2\),
\begin{align}
 &\kappa_{m+1}^{\mu_s^H}
 (D^{R_1},\ldots,D^{R_m},S_s^{[d]})
 \notag\\
 &\qquad=
 \sum_{T\in\mathfrak T_m}
 \int_{[0,1]^{m-1}}
 \partial_{E(T)}
 {\cal G}_m^{[d]}(\boldsymbol x^T(\boldsymbol u))
 \,d\boldsymbol u.
 \label{eq:newV34-jet-forest}
\end{align}
\end{theorem}

\begin{proof}
By \Cref{lem:newV34-Qq-endpoint},
\[
 {\cal G}_m^{[d]}(\boldsymbol x^\rho)
 =
 \sum_{B\in\rho}
 N_B^{[d]}
 \bigotimes_{C\in\rho\setminus\{B\}}M_C^H.
\]
Thus
\[
 F_m(\boldsymbol x^\rho)
 +\varepsilon{\cal G}_m^{[d]}(\boldsymbol x^\rho)
 =
 \bigotimes_{B\in\rho}
 (M_B^H+\varepsilon N_B^{[d]})
 \pmod{\varepsilon^2}.
\]
The dual-number BKAR argument of
\Cref{thm:newV33-root-forest} applies verbatim.  Its connected
\(\varepsilon\)-coefficient is the joint cumulant with the one
centered score \(S_s^{[d]}\).
\end{proof}

\begin{theorem}[Finite spectral-jet strong card]
\label{thm:newV34-jet-card}
For each fixed \(d\) and \(A_d\), there is
\(K(d,A_d)<\infty\) such that
\begin{align}
 &\left\|
 \kappa_{m+1}^{\mu_s^H}
 (D^{R_1},\ldots,D^{R_m},S_s^{[d]})
 \right\|_{\rm cb}
 \notag\\
 &\qquad\le
 K(d,A_d)^m
 \left(\prod_{i=1}^m b_i\right)
 \left(\sum_{i=1}^m b_i\right)^{m-2}.
 \label{eq:newV34-jet-card}
\end{align}
The constant is uniform in \(m,s\), the thermal
representations, and all tensor dimensions.
\end{theorem}

\begin{proof}
Fix a tree \(T\), put \(k=m-1\), and apply Leibniz to one
degree-\(q\) summand of \eqref{eq:newV34-Gd}.  By
\Cref{lem:newV34-q-payment,
lem:newV33-contract-derivative}, its norm is at most
\begin{align*}
 &|a_q(s)|m^{2q}
 \sum_{a=0}^{\min(k,2q-1)}
 \binom{k}{a}2^{k-a}
 \prod_{e\in E(T)}w_e\\
 &\qquad\le
 A_d\,m^{2q}(1+m)^{2q-1}2^k
 \prod_{e\in E(T)}w_e.
\end{align*}
For \(m\ge2\), the elementary maximum
\[
 \sup_{x\ge1}{\log x\over x}={1\over e}
\]
gives
\[
 m^{2q}(1+m)^{2q-1}
 \le
 2^{2q-1}
 \exp\left({(4q-1)m\over e}\right).
\]
Since \(q\le d\), summing the \(d\) terms and enlarging the base
gives
\[
 \|\partial_{E(T)}{\cal G}_m^{[d]}\|_{\rm cb}
 \le
 K(d,A_d)^m\prod_{e\in E(T)}w_e.
\]
Integrate, sum \(T\), and use the weighted Pr\"ufer identity.
\end{proof}

\begin{corollary}[V33 is the degree-two sharp instance]
\label{cor:newV34-V33-instance}
Take \(d=2\),
\[
 a_2={1\over3},
 \qquad
 a_1=-{5s\over3}.
\]
Then \(S_s^{[2]}=S_s^{\rm sp}\),
\({\cal G}_m^{[2]}={\cal G}_m\), and
\Cref{thm:newV34-jet-card} recovers
\Cref{thm:newV33-spectral-RGSC}.
\end{corollary}

\begin{proof}
Substitute the coefficients in
\eqref{eq:newV34-jet-multiplier}:
\[
 a_1K_R+a_2sK_R^2
 =
 {s\over3}(K_R^2-5K_R).
\]
The block-local interpolation has the identical substitution.
\end{proof}

\section{Why coefficientwise infinite jets do not close}
\label{sec:newV34-infinite}

The preceding proof also quantifies its own limit.  Before
absorbing powers of \(m\), a degree-\(q\) root carries
\[
 m^{2q}(1+m)^{2q-1}.
\]
It may be absorbed into an exponential in \(m\) for fixed \(q\),
but its exponential base grows exponentially in \(q\).

\begin{proposition}[Geometric coefficient decay is insufficient
after termwise norms]
\label{prop:newV34-geometric-no-go}
Suppose an infinite spectral jet has coefficients satisfying only
\[
 |a_q(s)|\le AR^{-q}
\]
with fixed \(A,R>0\), and estimate every degree by the proof of
\Cref{thm:newV34-jet-card} before summing \(q\).  For all
sufficiently large arities \(m\), the resulting positive majorant
diverges.  Therefore a finite analytic radius in the root degree
cannot justify a termwise-norm infinite-jet argument.
\end{proposition}

\begin{proof}
The proof of \Cref{thm:newV34-jet-card} gives, up to a base
independent of \(q\), a factor
\[
 \exp(cqm)
\]
for some numerical \(c>0\).  The coefficientwise majorant contains
\[
 A\sum_{q\ge1}R^{-q}e^{cqm}.
\]
This geometric series diverges as soon as \(e^{cm}\ge R\).
\end{proof}

\begin{proposition}[Even factorial coefficients give a
superexponential arity majorant]
\label{prop:newV34-factorial-no-go}
If one instead has
\[
 |a_q(s)|\le {A^q\over q!},
\]
the same degreewise triangle argument yields a factor bounded only
by
\begin{equation}
 \sum_{q\ge1}{A^q e^{cqm}\over q!}
 =
 e^{Ae^{cm}}-1,
 \label{eq:newV34-double-exponential}
\end{equation}
which cannot be absorbed into \(K^m\) for fixed \(K\).
\end{proposition}

\begin{proof}
The identity is the exponential series.  Its logarithm grows like
\(Ae^{cm}\), whereas the logarithm of \(K^m\) grows linearly in
\(m\).
\end{proof}

\begin{firewall}[Scope of the infinite-jet no-go]
\label{fire:newV34-infinite-scope}
\Cref{prop:newV34-geometric-no-go,
prop:newV34-factorial-no-go}
do not prove that the exact bridge score fails the RGSC.  They prove
that expanding its root into arbitrarily high degrees and then
taking norms degree by degree cannot establish an exponential
arity constant from ordinary analytic or factorial coefficient
decay.  Cancellation inside the complete root generating function
must be retained.
\end{firewall}

\section{The resummed root certificate}
\label{sec:newV34-certificate}

\begin{definition}[Block-local resummed root card]
\label{def:newV34-BRRC}
Let \(N_{\alpha,t,B}\) denote the exact centered bridge-score moment
on a tensor block \(B\).  The block-local resummed root card
\textup{(BRRC)} is the existence of a smooth interpolation
\({\cal G}_{\alpha,t,m}(\boldsymbol x)\) such that:
\begin{enumerate}[label=\textup{(\roman*)},leftmargin=3.5em]
\item at every partition endpoint,
\[
 {\cal G}_{\alpha,t,m}(\boldsymbol x^\rho)
 =
 \sum_{B\in\rho}N_{\alpha,t,B}
 \bigotimes_{C\in\rho\setminus\{B\}}M_{\alpha,t,C};
\]
\item for every tree \(T\),
\[
 \sup_{\boldsymbol u}
 \|\partial_{E(T)}
 {\cal G}_{\alpha,t,m}(\boldsymbol x^T(\boldsymbol u))\|_{\rm cb}
 \le
 K^m\prod_i b_i^{d_i(T)};
\]
\item the interpolation and bound are uniform in
\(t\in[0,1]\), all thermal representations, and all matrix
amplifications.
\end{enumerate}
\end{definition}

\begin{theorem}[BRRC is the exact remaining local comparison
certificate]
\label{thm:newV34-BRRC-sufficient}
If BRRC holds, then the exact geometric heat--Wilson score obeys
RGSC.  Together with the V29 heat endpoint and the exact path
identity, this proves the raw Wilson aggregate strong card and hence
the intrinsic strong Wilson forest card.
\end{theorem}

\begin{proof}
Apply the connected BKAR identity to the dual-number interpolation
\[
 F_{\alpha,t,m}
 +\varepsilon{\cal G}_{\alpha,t,m}.
\]
BRRC(i) identifies its connected \(\varepsilon\)-coefficient with
the exact score cumulant.  BRRC(ii), parameter integration, and the
weighted Pr\"ufer identity prove RGSC.  Integrate the exact
score identity in \(t\) and use
\Cref{thm:newV29-strong-heat-card}; this is the raw comparison
argument of \Cref{thm:newV26-raw-comparison}, with the heat endpoint
in place of GIAC.  The V16 balancing equivalence yields the
intrinsic strong forest card.
\end{proof}

\begin{assessment}[Precise gain over V33]
\label{ass:newV34-gain}
V33 closed one degree-two spectral root.  V34 proves that no finite
spectral order is an obstruction: every finite jet has an exact
block-local forest and the correct all-arity weight.  The remaining
difficulty is neither the root clock nor noncommutative time
ordering.  It is the construction of one resummed interpolation for
the exact bridge marked multiplier, with its entire root-degree sum
kept inside the forest derivative.  BRRC records exactly that
obligation and excludes the false raw Bernoulli-tree PCMC of V15.
\end{assessment}

\section{V35 programme and claim boundary}
\label{sec:newV34-next}

\begin{programme}[V35: audit and resummed root generating function]
\label{prog:newV34-V35}
\begin{enumerate}[leftmargin=3em]
\item Perform the compulsory read-only audit of the newest active
      SM--Einstein and Navier--Stokes notes before choosing the
      resummation architecture.
\item Write the exact bridge marked block moment as
      \(\partial_t q_{\alpha,t}(C_B^{\rm tot})\) by central
      Peter--Weyl functional calculus.
\item Seek a contour or Laplace representation which keeps the full
      scalar multiplier inside a contractive semigroup before the
      V33 canonical gate is differentiated.
\item Require the representing measure to have an exponential-grade
      weighted variation uniformly in the thermal spectral window;
      otherwise record the first failed moment explicitly.
\item Compare a semigroup contour route with cumulant-level
      all-tree compensation, never with the disproved raw-tree PCMC.
\item Test the proposed resummation on the V33 degree-two multiplier
      and on the V26 nontrivial saturation sectors.
\end{enumerate}
\end{programme}

\begin{firewall}[Claim boundary after V34]
\label{fire:newV34-boundary}
V34 proves endpoint localization, root payment, an exact marked
forest, and the all-arity strong raw card for every finite spectral
jet.  It proves that coefficientwise norm summation of an infinite
jet has an unacceptable arity majorant and defines the resummed
BRRC sufficient for the exact comparison.  It does not prove BRRC
for the exact bridge score.  Hence it does not yet prove the Wilson
aggregate strong card, all-order strong MTA, WUFC, CPM, cutoff
removal, reflection positivity, Osterwalder--Schrader
reconstruction, a continuum Yang--Mills measure, or a positive
Hamiltonian gap.  The full Yang--Mills mass gap has not been proved.
\end{firewall}

\begin{theorem}[V34 result ledger]
\label{thm:newV34-results}
V34:
\begin{enumerate}[label=\textup{(V34.\arabic*)},leftmargin=6.5em]
\item constructs the block-local Casimir power of every finite
      degree;
\item proves the exact degree-\(q\) root-payment inequality;
\item proves a dual-number marked forest for every finite spectral
      jet;
\item proves its uniform all-arity completely bounded score card;
\item shows that geometric or factorial coefficient decay does not
      support degreewise norm summation at arbitrary arity; and
\item isolates BRRC as the exact resummed local comparison
      certificate.
\end{enumerate}
\end{theorem}

\begin{proof}
Use
\Cref{lem:newV34-Qq-endpoint,
lem:newV34-q-payment,
thm:newV34-jet-forest,
thm:newV34-jet-card,
prop:newV34-geometric-no-go,
prop:newV34-factorial-no-go,
thm:newV34-BRRC-sufficient}.
\end{proof}

\paragraph{Parent-version record.}
V34 was formed from
\texttt{Renewal\_Yang\_Mills\_Mass\_Gap\_Research\_Notes\_V33.tex}.
The V33 parent had \(15\,605\) logical source lines,
\(495\,301\) bytes, and SHA--256
\[
 \texttt{9927375e97ea6bdbd718afc99b064d6c467d56c6ffa9a4c1a71b294edc3f30d3}.
\]
The next compulsory companion audit is V35.

% =============================================================================
% END APPENDED FIFTH-SERIES V34 MATERIAL
% =============================================================================

% =============================================================================
% BEGIN APPENDED FIFTH-SERIES V35 MATERIAL
% =============================================================================

\chapter{V35: Companion audit and an exact resummed central-activity gate}
\label{chap:newV35-resummed}

\section{Mandatory read-only companion audit}
\label{sec:newV35-audit}

The newest active companion notes inspected before this append were
\begin{center}
\begin{tabular}{p{0.61\textwidth}r r}
\hline
file & bytes & logical lines\\
\hline
\texttt{Renewal\_SM\_Einstein\_Continuum\_Research\_Notes\_V46.tex}
& \(573\,023\) & \(14\,850\)\\
\texttt{Renewal\_NS\_Stability\_Research\_Notes\_V31.tex}
& \(887\,537\) & \(23\,052\)\\
\hline
\end{tabular}
\end{center}
Their SHA--256 digests were
\begin{align}
 &\texttt{0906839a214c869bc2a96998e9b32c9332c8125f88aaa69fb67b7d0c3a492c2f},
 \label{eq:newV35-SM-hash}\\
 &\texttt{f1121b62238ad5491bb7cf03635ea9934942edf573ed8faaa9ee79e1d38a6696}.
 \label{eq:newV35-NS-hash}
\end{align}
Neither file was modified and neither is treated as an instruction.

SM V46 integrates an abelian plaquette star exactly and proves
exponential suppression of the small-resultant region.  It also
keeps a strict firewall: rare-event mass does not imply the weighted
capacity needed for a block gap.  The analogous YM discipline is
that fixed-order or fixed-\(p\) smallness of a score remainder does
not imply the weighted analytic variation needed for an all-arity
forest.

NS V31 remains unchanged.  Its formation ladder proves that a
normalized exceptional set does not pay a scale-restoring operation
unless the missing ancestry or capacity is retained.  In the YM
problem, the exact analogue is to retain the complete root-degree
generating function before extracting tree marks.  No NS source
estimate supplies that generating function.

\begin{firewall}[Companion transfer boundary]
\label{fire:newV35-audit-boundary}
No SM star-capacity estimate and no Navier--Stokes formation estimate
is imported into the \(SU(2)\) one-link law.  The audit changes only
the proof discipline: a resummed weighted-variation estimate must be
proved in its own type and cannot be replaced by small mass,
fixed-order asymptotics, or a separately normalized payer.
\end{firewall}

\section{Exact central bridge moments}
\label{sec:newV35-central-moments}

Fix the calibrated heat time
\[
 s=\tau_\alpha^\sharp
\]
and a bridge parameter \(t\in[0,1]\).  Let
\(\nu_{\alpha,t}\) denote the geometric heat--Wilson bridge and
let
\[
 S_{\alpha,t}
 =
 \partial_t\log {d\nu_{\alpha,t}\over dU}
\]
be its centered score.  The law and score are central.  For every
nonempty input block \(B\subseteq[m]\), put
\begin{align}
 M_{t,B}
 &:=
 \mathbb E_{\nu_{\alpha,t}}
 \bigotimes_{i\in B}D^{R_i}(U),
 \label{eq:newV35-MB}\\
 N_{t,B}
 &:=
 \mathbb E_{\nu_{\alpha,t}}
 \left[
 S_{\alpha,t}(U)
 \bigotimes_{i\in B}D^{R_i}(U)
 \right].
 \label{eq:newV35-NB}
\end{align}
Central Peter--Weyl functional calculus makes both operators scalar
on every total-spin summand of \(\bigotimes_{i\in B}R_i\).

\begin{lemma}[The marked moment is the path derivative]
\label{lem:newV35-marked-derivative}
\begin{equation}
 N_{t,B}=\partial_tM_{t,B}.
 \label{eq:newV35-marked-derivative}
\end{equation}
Moreover \(N_{t,\varnothing}=0\).
\end{lemma}

\begin{proof}
Differentiate the normalized geometric bridge density.  Its
derivative is \(S_{\alpha,t}\,d\nu_{\alpha,t}\).  The group is
compact and the density is smooth and positive, so differentiation
under the finite integral is valid.  The empty-block assertion is
\(\mathbb E_{\nu_{\alpha,t}}S_{\alpha,t}=0\).
\end{proof}

Retain the positive block heat Casimir
\[
 {\cal K}_B=sC_B^{\rm tot}.
\]
Introduce a central dual number \(\varepsilon\), \(\varepsilon^2=0\),
and define the heat-normalized dual block moment
\begin{equation}
 {\cal Y}_{t,B}(\varepsilon)
 :=
 e^{{\cal K}_B}
 \left(M_{t,B}+\varepsilon N_{t,B}\right).
 \label{eq:newV35-YB}
\end{equation}
This is an exact algebraic normalization.  No bound on
\(e^{{\cal K}_B}M_{t,B}\) is asserted here.

\section{Dual cumulant activities}
\label{sec:newV35-activities}

For every nonempty \(B\), define its dual activity
\({\cal A}_{t,B}(\varepsilon)\) by moment--cumulant inversion:
\begin{align}
 {\cal Y}_{t,B}(\varepsilon)
 &=
 \sum_{\pi\in\Pi(B)}
 \bigotimes_{C\in\pi}{\cal A}_{t,C}(\varepsilon),
 \label{eq:newV35-activity-moment}\\
 {\cal A}_{t,B}(\varepsilon)
 &=
 \sum_{\pi\in\Pi(B)}
 (-1)^{|\pi|-1}(|\pi|-1)!
 \bigotimes_{C\in\pi}{\cal Y}_{t,C}(\varepsilon).
 \label{eq:newV35-activity-inverse}
\end{align}
Products over partition blocks act on disjoint tensor factors and
retain the canonical input order.  Write
\[
 {\cal A}_{t,B}(\varepsilon)
 =
 {\cal A}_{t,B}^{(0)}
 +\varepsilon{\cal A}_{t,B}^{(1)}.
\]

For a partition \(\pi\) of \([m]\), define its gate
\begin{equation}
 \chi_\pi(\boldsymbol x)
 :=
 \prod_{B\in\pi}\chi_B(\boldsymbol x),
 \label{eq:newV35-partition-gate}
\end{equation}
where \(\chi_B\) is the V33 canonical-star monomial.  The resummed
activity field is
\begin{equation}
 {\cal R}_{t,m}(\boldsymbol x;\varepsilon)
 :=
 \sum_{\pi\in\Pi([m])}
 \chi_\pi(\boldsymbol x)
 \bigotimes_{B\in\pi}{\cal A}_{t,B}(\varepsilon)
 =
 {\cal R}_{t,m}^{(0)}(\boldsymbol x)
 +\varepsilon{\cal R}_{t,m}^{(1)}(\boldsymbol x).
 \label{eq:newV35-R}
\end{equation}
Unlike the V34 infinite-jet triangle majorant, this is one finite
object assembled from complete block activities before any norm.

\begin{theorem}[Exact partition endpoint factorization]
\label{thm:newV35-R-endpoint}
For every partition \(\rho\in\Pi([m])\),
\begin{equation}
 {\cal R}_{t,m}(\boldsymbol x^\rho;\varepsilon)
 =
 \bigotimes_{B\in\rho}{\cal Y}_{t,B}(\varepsilon).
 \label{eq:newV35-R-endpoint}
\end{equation}
\end{theorem}

\begin{proof}
At \(\boldsymbol x^\rho\), the factor
\(\chi_C\) equals one exactly when \(C\) is contained in one
\(\rho\)-block.  Hence \(\chi_\pi(\boldsymbol x^\rho)=1\)
exactly when \(\pi\) refines \(\rho\).  Therefore
\begin{align*}
 {\cal R}_{t,m}(\boldsymbol x^\rho;\varepsilon)
 &=
 \sum_{\pi\le\rho}
 \bigotimes_{C\in\pi}{\cal A}_{t,C}(\varepsilon)\\
 &=
 \bigotimes_{B\in\rho}
 \left[
 \sum_{\pi_B\in\Pi(B)}
 \bigotimes_{C\in\pi_B}{\cal A}_{t,C}(\varepsilon)
 \right].
\end{align*}
Apply \eqref{eq:newV35-activity-moment} inside every block.
\end{proof}

\section{A universal exact one-root interpolation}
\label{sec:newV35-universal}

Let \({\cal K}(\boldsymbol x)\) be the positive heat Casimir
interpolation from V33 and set
\begin{align}
 {\cal F}_{\alpha,t,m}(\boldsymbol x)
 &:=
 e^{-{\cal K}(\boldsymbol x)}
 {\cal R}_{t,m}^{(0)}(\boldsymbol x),
 \label{eq:newV35-Fexact}\\
 {\cal G}_{\alpha,t,m}(\boldsymbol x)
 &:=
 e^{-{\cal K}(\boldsymbol x)}
 {\cal R}_{t,m}^{(1)}(\boldsymbol x).
 \label{eq:newV35-Gexact}
\end{align}
The order in these products is part of the definition.  At
partition endpoints all factors split over disjoint blocks, so the
order becomes immaterial.

\begin{theorem}[Exact dual endpoint law for every central bridge]
\label{thm:newV35-dual-endpoint}
For every partition \(\rho\),
\begin{equation}
 {\cal F}_{\alpha,t,m}(\boldsymbol x^\rho)
 +\varepsilon{\cal G}_{\alpha,t,m}(\boldsymbol x^\rho)
 =
 \bigotimes_{B\in\rho}
 (M_{t,B}+\varepsilon N_{t,B})
 \pmod{\varepsilon^2}.
 \label{eq:newV35-dual-endpoint}
\end{equation}
In particular, at the fully coupled endpoint the two coefficients
are \(M_{t,[m]}\) and \(N_{t,[m]}\).
\end{theorem}

\begin{proof}
By \Cref{thm:newV35-R-endpoint}, the left side is the
\(\varepsilon\)-expansion of
\[
 e^{-\sum_{B\in\rho}{\cal K}_B}
 \bigotimes_{B\in\rho}
 e^{{\cal K}_B}(M_{t,B}+\varepsilon N_{t,B}).
\]
Operators from distinct blocks commute, and the heat factors cancel
inside each block.  This is the right side of
\eqref{eq:newV35-dual-endpoint}.
\end{proof}

\begin{theorem}[Exact resummed bridge-score forest]
\label{thm:newV35-exact-forest}
For every \(m\ge2\),
\begin{align}
 &\kappa_{m+1}^{\nu_{\alpha,t}}
 (D^{R_1},\ldots,D^{R_m},S_{\alpha,t})
 \notag\\
 &\qquad=
 \sum_{T\in\mathfrak T_m}
 \int_{[0,1]^{m-1}}
 \partial_{E(T)}
 {\cal G}_{\alpha,t,m}
 (\boldsymbol x^T(\boldsymbol u))
 \,d\boldsymbol u.
 \label{eq:newV35-exact-forest}
\end{align}
This identity is valid for the exact geometric bridge at every
\(t\); it contains no score Taylor series and no raw Bernoulli-tree
PCMC.
\end{theorem}

\begin{proof}
Apply the connected multivariable BKAR formula to the dual family
\({\cal F}_{\alpha,t,m}
+\varepsilon{\cal G}_{\alpha,t,m}\).
By \Cref{thm:newV35-dual-endpoint}, its partition endpoints are
the exact products of the dual moments
\(M_{t,B}+\varepsilon N_{t,B}\).  The connected
\(\varepsilon\)-coefficient is therefore
\[
 \sum_{\rho\in\Pi([m])}
 (-1)^{|\rho|-1}(|\rho|-1)!
 \sum_{B\in\rho}N_{t,B}
 \bigotimes_{C\in\rho\setminus\{B\}}M_{t,C}.
\]
Since the score is centered, this is precisely the joint cumulant
on the left.  Extracting the same coefficient from the tree side
gives \eqref{eq:newV35-exact-forest}.
\end{proof}

\begin{assessment}[What is new compared with the random-cluster
forest]
\label{ass:newV35-not-random-cluster}
The V14 random-cluster interpolation expands incomplete graph
components and its individual trees retain lower-cumulant leakage;
V15 proves their desired norm bound false.  The field
\eqref{eq:newV35-R} first forms the complete heat-normalized
cumulant activity of each block, then gates that whole activity.
Thus every lower partition has already been M\"obius-compressed
inside \({\cal A}_{t,B}\).  The construction is exact; its remaining
issue is the norm of the assembled analytic field, not an unproved
claim that raw Bernoulli trees are small.
\end{assessment}

\section{The heat endpoint becomes a Boolean root difference}
\label{sec:newV35-Boolean}

Suppose temporarily that the unmarked law is the heat law, so
\(M_B=e^{-{\cal K}_B}\), and that a centered marked moment has the
form
\begin{equation}
 N_B=\Phi_s({\cal K}_B)e^{-{\cal K}_B},
 \qquad
 \Phi_s(0)=0.
 \label{eq:newV35-Phi-moment}
\end{equation}
Then
\[
 {\cal Y}_B(\varepsilon)=I+\varepsilon\Phi_s({\cal K}_B).
\]

\begin{lemma}[Boolean form of the marked heat activities]
\label{lem:newV35-Boolean}
The unmarked activities satisfy
\[
 {\cal A}_{\{i\}}^{(0)}=I,
 \qquad
 {\cal A}_B^{(0)}=0\quad(|B|\ge2),
\]
and the marked activity is the Boolean subset difference
\begin{equation}
 {\cal A}_B^{(1)}
 =
 \Delta_B\Phi_s
 :=
 \sum_{C\subseteq B}
 (-1)^{|B|-|C|}\Phi_s({\cal K}_C),
 \label{eq:newV35-Boolean}
\end{equation}
where \(\Phi_s({\cal K}_\varnothing)=\Phi_s(0)I=0\).
Consequently
\begin{equation}
 {\cal R}_m^{(1)}(\boldsymbol x)
 =
 \sum_{\varnothing\ne B\subseteq[m]}
 \chi_B(\boldsymbol x)\Delta_B\Phi_s.
 \label{eq:newV35-resummed-Boolean}
\end{equation}
\end{lemma}

\begin{proof}
At order zero, the moment relation
\eqref{eq:newV35-activity-moment} reads
\[
 I=\sum_{\pi\in\Pi(B)}
 \bigotimes_{C\in\pi}{\cal A}_C^{(0)}.
\]
It is solved recursively by singleton identities and zero higher
activities.  At order \(\varepsilon\), all unmarked blocks other
than singletons vanish.  Hence
\[
 \Phi_s({\cal K}_B)
 =
 \sum_{\varnothing\ne C\subseteq B}{\cal A}_C^{(1)}.
\]
Boolean M\"obius inversion gives
\eqref{eq:newV35-Boolean}.  In
\eqref{eq:newV35-R}, a nonzero marked term then consists of one
marked block \(B\) and unmarked singleton activities on its
complement, proving \eqref{eq:newV35-resummed-Boolean}.
\end{proof}

\begin{corollary}[The V33 gate is the quadratic Boolean instance]
\label{cor:newV35-V33-Boolean}
For
\[
 \Phi_s(K)={s\over3}(K^2-5K),
\]
\eqref{eq:newV35-resummed-Boolean} is an endpoint-equivalent
resummed form of the V33 block-local quartic interpolation.  Its
Boolean activity \(\Delta_B\Phi_s\) vanishes for \(|B|>4\), and
the V33 root-payment lemma proves the required tree card.
\end{corollary}

\begin{proof}
Expanding \(K_C^q\) into ordered \(A_{ij}\)-words and applying the
Boolean difference over \(C\subseteq B\) retains exactly those
words whose vertex support is all of \(B\).  A quadratic word has
at most four support vertices.  Summing
\(\chi_B\Delta_B\Phi_s\) over \(B\) is therefore exactly the
support grouping of the V33 gated monomials.
\end{proof}

\section{Weighted complex polydiscs extract all tree marks at once}
\label{sec:newV35-polydisc}

Fix a tree \(T\), a real forest point
\(\boldsymbol x^T(\boldsymbol u)\), and a numerical radius
\(r_0>0\).  Vary only the tree-edge variables and define
\begin{equation}
 {\mathbb D}_{T,r_0}
 :=
 \left\{
 \boldsymbol z:
 |z_e-x_e^T|\le{r_0\over w_e}\ (e\in E(T)),
 \quad
 z_f=x_f^T\ (f\notin E(T))
 \right\}.
 \label{eq:newV35-polydisc}
\end{equation}
All fields above are entire in the finitely many edge variables, so
this large weighted polydisc is legitimate.

\begin{lemma}[Heat factor on the weighted polydisc]
\label{lem:newV35-complex-heat}
For every \(\boldsymbol z\in{\mathbb D}_{T,r_0}\),
\begin{equation}
 \|e^{-{\cal K}(\boldsymbol z)}\|_{\rm cb}
 \le e^{2r_0(m-1)}.
 \label{eq:newV35-complex-heat}
\end{equation}
\end{lemma}

\begin{proof}
Write
\[
 {\cal K}(\boldsymbol z)
 =
 {\cal K}(\boldsymbol x^T)+W,
 \qquad
 W=\sum_{e\in E(T)}(z_e-x_e^T)V_e.
\]
The real base operator is positive by
\Cref{thm:newV33-positive-layer}, and
\eqref{eq:newV33-pair-cb} gives
\[
 \|W\|_{\rm cb}
 \le
 \sum_{e\in E(T)}{r_0\over w_e}\,2w_e
 =2r_0(m-1).
\]
The bounded-perturbation Dyson series around the contraction
semigroup \(e^{-\tau{\cal K}(\boldsymbol x^T)}\) yields
\(\|e^{-({\cal K}+W)}\|\le e^{\|W\|}\).  The same proof holds
after amplification.
\end{proof}

\begin{definition}[Weighted polydisc resummed activity card]
\label{def:newV35-WPRAC}
The WPRAC holds if there are numerical \(r_0,A>0\) such that,
uniformly in \(\alpha,t,m\), all thermal representations, every
tree \(T\), and every forest point,
\begin{equation}
 \sup_{\boldsymbol z\in{\mathbb D}_{T,r_0}}
 \|{\cal R}_{t,m}^{(1)}(\boldsymbol z)\|_{\rm cb}
 \le A^m.
 \label{eq:newV35-WPRAC}
\end{equation}
\end{definition}

\begin{theorem}[WPRAC implies the exact bridge RGSC]
\label{thm:newV35-WPRAC-sufficient}
If WPRAC holds, then
\begin{align}
 &\left\|
 \kappa_{m+1}^{\nu_{\alpha,t}}
 (D^{R_1},\ldots,D^{R_m},S_{\alpha,t})
 \right\|_{\rm cb}
 \notag\\
 &\qquad\le
 K^m
 \left(\prod_i b_i\right)
 \left(\sum_i b_i\right)^{m-2}
 \label{eq:newV35-exact-RGSC}
\end{align}
uniformly along the entire bridge.  Consequently the raw Wilson
aggregate strong card and the intrinsic strong Wilson forest card
follow.
\end{theorem}

\begin{proof}
By \Cref{lem:newV35-complex-heat} and WPRAC,
\[
 \sup_{{\mathbb D}_{T,r_0}}
 \|{\cal G}_{\alpha,t,m}\|_{\rm cb}
 \le
 (Ae^{2r_0})^m.
\]
Apply the multivariate Banach-valued Cauchy estimate at the center
of the polydisc.  Every distinct tree variable contributes the
inverse radius \(w_e/r_0\), with no factorial:
\[
 \|\partial_{E(T)}{\cal G}_{\alpha,t,m}
 (\boldsymbol x^T)\|_{\rm cb}
 \le
 (Ae^{2r_0})^m r_0^{-(m-1)}
 \prod_{e\in E(T)}w_e.
\]
Insert this in \eqref{eq:newV35-exact-forest}, integrate, and use
the weighted Pr\"ufer identity.  This proves
\eqref{eq:newV35-exact-RGSC}.  The exact path comparison and
\Cref{thm:newV29-strong-heat-card,
thm:newV26-raw-comparison} give the two stated consequences.
\end{proof}

\begin{proposition}[Every finite spectral jet passes WPRAC]
\label{prop:newV35-finite-WPRAC}
For each fixed finite jet of V34, the corresponding heat-endpoint
field \({\cal R}_m^{(1)}\) satisfies WPRAC with constants depending
only on the jet degree and coefficient bound.
\end{proposition}

\begin{proof}
By \Cref{lem:newV35-Boolean}, group every polynomial word by its
exact vertex support \(B\).  V34 root payment assigns to that word
the full canonical-star weight
\(\prod_{e\in{\cal A}(B)}w_e\).  On the weighted polydisc,
\[
 w_e|z_e|
 \le w_e+r_0\le1+r_0
 \qquad(e\in E(T)),
\]
whereas off-tree gate variables remain in \([0,1]\).  Hence each
gate times its paying star weight is at most
\((1+r_0)^{|B|-1}\).  The word and degree census is the finite one
in the proof of \Cref{thm:newV34-jet-card}, and is bounded by
\(A_d^m\) for fixed degree \(d\).  Sum the finitely many degrees
before enlarging the exponential base.
\end{proof}

\begin{assessment}[The exact remaining estimate after V35]
\label{ass:newV35-remaining}
The existence of an exact resummed bridge interpolation is no
longer open: it is \eqref{eq:newV35-Gexact}, constructed directly
from the exact central moments.  The remaining local analytic gate
is WPRAC, a single weighted complex-variation estimate for the
already assembled marked activity field.  It neither expands the
score degree by degree nor estimates raw random-cluster trees.
At the heat endpoint it becomes the concrete Boolean-difference
sum \eqref{eq:newV35-resummed-Boolean}; along the bridge it also
contains the heat-normalized unmarked cumulant activities.
\end{assessment}

\section{V36 programme and claim boundary}
\label{sec:newV35-next}

\begin{programme}[V36: binary activity and thermal-window test]
\label{prog:newV35-V36}
\begin{enumerate}[leftmargin=3em]
\item Compute
      \({\cal A}_{t,\{1,2\}}^{(1)}\) exactly in terms of the
      heat-normalized multipliers \(M_{t,B},N_{t,B}\), including
      both marked placements.
\item Use the V31 path gap and V32 negative score norm to prove the
      one-block and binary WPRAC rows uniformly in the thermal
      representation window.
\item Test the weighted complex radius \(r_0/w_{12}\) on the exact
      two-label gate; determine whether the heat normalization
      \(e^{{\cal K}_{12}}\) is harmless or already grows faster than
      exponentially in the two input marks.
\item Derive the general activity recursion from
      \eqref{eq:newV35-activity-moment} and retain its complete
      subset sum before norms.
\item If high total spin defeats heat normalization, change the
      reference time blockwise before proceeding; do not conceal the
      loss inside the WPRAC hypothesis.
\item Keep the next compulsory companion audit at V40.
\end{enumerate}
\end{programme}

\begin{firewall}[Claim boundary after V35]
\label{fire:newV35-boundary}
V35 performs the compulsory companion audit, constructs exact
heat-normalized dual cumulant activities for every central bridge
law, proves their partition endpoint factorization, and obtains an
exact resummed bridge-score forest.  It proves that the weighted
polydisc card WPRAC would yield the exact RGSC and verifies WPRAC
for every finite spectral jet.  It does not prove WPRAC for the
exact bridge activity field.  Hence the Wilson aggregate strong
card, all-order strong MTA, WUFC, CPM, cutoff removal, reflection
positivity, Osterwalder--Schrader reconstruction, a continuum
Yang--Mills measure, and a positive Hamiltonian gap remain open.
The full Yang--Mills mass gap has not been proved.
\end{firewall}

\begin{theorem}[V35 result ledger]
\label{thm:newV35-results}
V35:
\begin{enumerate}[label=\textup{(V35.\arabic*)},leftmargin=6.5em]
\item records the read-only SM V46 and NS V31 audit;
\item defines exact heat-normalized dual bridge moments and their
      cumulant activities;
\item proves the activity field factors into the exact dual block
      moments at every partition endpoint;
\item proves an exact all-\(t\) resummed bridge-score forest;
\item identifies the heat endpoint with a Boolean root-difference
      field and recovers the V33 quartic gate;
\item proves that WPRAC implies the exact bridge RGSC by a
      factorial-free weighted Cauchy estimate; and
\item proves WPRAC for every finite spectral jet.
\end{enumerate}
\end{theorem}

\begin{proof}
Use
\Cref{lem:newV35-marked-derivative,
thm:newV35-R-endpoint,
thm:newV35-dual-endpoint,
thm:newV35-exact-forest,
lem:newV35-Boolean,
lem:newV35-complex-heat,
thm:newV35-WPRAC-sufficient,
prop:newV35-finite-WPRAC}.
\end{proof}

\paragraph{Parent-version record.}
V35 was formed from
\texttt{Renewal\_Yang\_Mills\_Mass\_Gap\_Research\_Notes\_V34.tex}.
The V34 parent had \(16\,122\) logical source lines,
\(511\,068\) bytes, and SHA--256
\[
 \texttt{9fd9b9964d1157054462036e1fa372d76da742584c3011709468f9bcc4ed1345}.
\]
The next compulsory companion audit is V40.

% =============================================================================
% END APPENDED FIFTH-SERIES V35 MATERIAL
% =============================================================================

% =============================================================================
% BEGIN APPENDED FIFTH-SERIES V36 MATERIAL
% =============================================================================

\chapter{V36: Exact binary resummed activity from the bridge negative norm}
\label{chap:newV36-binary}

\section{A complete marked-moment row for any tensor block}
\label{sec:newV36-block-row}

Retain
\[
 s=\tau_\alpha^\sharp={1\over2\alpha+2/3},
 \qquad
 b_i^2=s(1+C_2(R_i))<1,
 \qquad
 B_A=\sum_{i\in A}b_i.
\]
Let \(C_A^{\rm tot}\) be the total Casimir on the tensor block
\(\bigotimes_{i\in A}R_i\), and let
\({\cal K}_A=sC_A^{\rm tot}\).

\begin{lemma}[Upper total-Casimir row]
\label{lem:newV36-total-Casimir}
\begin{equation}
 \|{\cal K}_A\|_{\rm op}\le B_A^2.
 \label{eq:newV36-total-Casimir}
\end{equation}
\end{lemma}

\begin{proof}
Write the total generator row as
\[
 J_A=(\sum_{i\in A}J_i^1,\sum_{i\in A}J_i^2,
                         \sum_{i\in A}J_i^3).
\]
The row triangle inequality gives
\[
 \|C_A^{\rm tot}\|^{1/2}
 =\|J_A\|_{\rm row}
 \le\sum_{i\in A}\|J_i\|_{\rm row}
 =\sum_{i\in A}\sqrt{C_2(R_i)}.
\]
Multiply by \(\sqrt s\) and use
\(\sqrt{sC_2(R_i)}\le b_i\), then square.
\end{proof}

Let \(A_{\alpha,t}\) be the positive reversible bridge generator,
\(R_{\alpha,t}=A_{\alpha,t}^{-1}Q\), and let
\(S_{\alpha,t}\) be the exact centered bridge score.  V32 proves
\[
 \|dR_{\alpha,t}S_{\alpha,t}\|_2
 \le C\alpha^{-3/2}
 \le Cs^{3/2}.
\]
The last inequality merely uses \(s\asymp\alpha^{-1}\).

\begin{theorem}[Completely bounded marked block row]
\label{thm:newV36-marked-block}
For every nonempty block \(A\),
\begin{equation}
 \|N_{t,A}\|_{\rm cb}
 \le CsB_A,
 \label{eq:newV36-marked-block}
\end{equation}
uniformly in \(t\in[0,1]\).  Consequently the normalized moments
\[
 {\cal R}_{t,A}:=e^{{\cal K}_A}M_{t,A},
 \qquad
 {\cal Z}_{t,A}:=e^{{\cal K}_A}N_{t,A}
\]
obey
\begin{align}
 \|{\cal R}_{t,A}\|_{\rm cb}
 &\le e^{B_A^2},
 \label{eq:newV36-R-block}\\
 \|{\cal Z}_{t,A}\|_{\rm cb}
 &\le CsB_Ae^{B_A^2}.
 \label{eq:newV36-Z-block}
\end{align}
\end{theorem}

\begin{proof}
Test \(N_{t,A}\) against arbitrary unit vectors in the tensor
representation and an arbitrary spectator amplification.  The
resulting scalar matrix coefficient \(f\) has pointwise Dirichlet
row bounded by
\[
 |df|\le\|C_A^{\rm tot}\|^{1/2}.
\]
The negative-norm integration-by-parts identity gives
\[
 |\mathbb E_{\nu_{\alpha,t}}[S_{\alpha,t}f]|
 =
 |\langle dR_{\alpha,t}S_{\alpha,t},df\rangle|
 \le
 Cs^{3/2}\|C_A^{\rm tot}\|^{1/2}.
\]
Use \Cref{lem:newV36-total-Casimir} in the form
\[
 \|C_A^{\rm tot}\|^{1/2}\le {B_A\over\sqrt s}.
\]
This proves \eqref{eq:newV36-marked-block}.  The representation
matrices are unitary, so
\(\|M_{t,A}\|_{\rm cb}\le1\).  Functional calculus and
\eqref{eq:newV36-total-Casimir} give
\(\|e^{{\cal K}_A}\|\le e^{B_A^2}\), proving the last two rows.
\end{proof}

\begin{remark}[Why the estimate is genuinely complete]
\label{rem:newV36-complete}
The spectator identity neither changes the total Casimir row nor
the bridge Dirichlet form.  Thus
\eqref{eq:newV36-marked-block} is not a collection of
entry-dependent scalar bounds; it is stable under all matrix
amplifications required in WPRAC.
\end{remark}

\section{Exact binary activities}
\label{sec:newV36-binary-activities}

For brevity write
\[
 R_A={\cal R}_{t,A},
 \qquad
 Z_A={\cal Z}_{t,A}.
\]
The V35 normalized dual moments are
\[
 {\cal Y}_A(\varepsilon)=R_A+\varepsilon Z_A.
\]

\begin{lemma}[Binary activity formula]
\label{lem:newV36-binary-activity}
For two labels,
\begin{align}
 {\cal A}_{\{i\}}^{(0)}&=R_i,
 &
 {\cal A}_{\{i\}}^{(1)}&=Z_i,
 \label{eq:newV36-single-activity}\\
 {\cal A}_{12}^{(0)}
 &=R_{12}-R_1R_2,
 &
 {\cal A}_{12}^{(1)}
 &=Z_{12}-Z_1R_2-R_1Z_2.
 \label{eq:newV36-pair-activity}
\end{align}
The complete marked activity field is
\begin{equation}
 {\cal R}_{t,2}^{(1)}(x)
 =
 Z_1R_2+R_1Z_2
 +x\left(Z_{12}-Z_1R_2-R_1Z_2\right).
 \label{eq:newV36-binary-field}
\end{equation}
\end{lemma}

\begin{proof}
The singleton identities are definitions.  On \(\{1,2\}\), the
moment--activity relation is
\[
 {\cal Y}_{12}
 =
 {\cal A}_{12}+{\cal A}_1{\cal A}_2.
\]
Extract its two dual-number coefficients.  In the resummed field
\eqref{eq:newV35-R}, the singleton partition has gate one and the
one-block partition has gate \(\chi_{\{1,2\}}=x\), which gives
\eqref{eq:newV36-binary-field}.
\end{proof}

\section{The score clock pays the large binary complex radius}
\label{sec:newV36-payment}

\begin{lemma}[Thermal two-mark payment]
\label{lem:newV36-two-mark}
For \(0<s<1\) and \(b_i^2\ge s\), \(b_i<1\),
\begin{equation}
 s(b_1+b_2)\le2b_1b_2.
 \label{eq:newV36-two-mark}
\end{equation}
\end{lemma}

\begin{proof}
The hypotheses give \(s\le b_i^2\le b_i\).  Therefore
\[
 sb_1\le b_1b_2,
 \qquad
 sb_2\le b_1b_2.
\]
Add the two inequalities.
\end{proof}

\begin{theorem}[Exact binary activity carries its edge mark]
\label{thm:newV36-binary-payment}
There is a numerical \(C\) such that
\begin{align}
 \|Z_1R_2+R_1Z_2\|_{\rm cb}
 &\le Cw_{12},
 \label{eq:newV36-binary-base}\\
 \|{\cal A}_{12}^{(1)}\|_{\rm cb}
 &\le Cw_{12}.
 \label{eq:newV36-binary-edge}
\end{align}
The bounds are uniform in the bridge parameter, the coupling, and
both thermal representations.
\end{theorem}

\begin{proof}
For a singleton, \(B_{\{i\}}=b_i<1\); for the pair,
\(B_{\{1,2\}}<2\).  Hence
\Cref{thm:newV36-marked-block} gives
\begin{align*}
 \|R_i\|&\le e,&
 \|Z_i\|&\le Cesb_i,\\
 \|R_{12}\|&\le e^4,&
 \|Z_{12}\|&\le Ce^4s(b_1+b_2).
\end{align*}
Apply \Cref{lem:newV36-two-mark}.  It bounds the two singleton
placements and \(Z_{12}\) by \(Cw_{12}\).
Equation \eqref{eq:newV36-pair-activity} and the triangle inequality
then prove both displayed estimates.
\end{proof}

\begin{theorem}[Binary WPRAC]
\label{thm:newV36-binary-WPRAC}
The exact geometric heat--Wilson bridge satisfies WPRAC at
\(m=2\).  More precisely, for every fixed \(r_0>0\),
\begin{equation}
 \sup_{\substack{x_0\in[0,1]\\
 |z-x_0|\le r_0/w_{12}}}
 \|{\cal R}_{t,2}^{(1)}(z)\|_{\rm cb}
 \le C(1+r_0).
 \label{eq:newV36-binary-WPRAC}
\end{equation}
\end{theorem}

\begin{proof}
For \(z\) in the displayed disk,
\[
 |z|\le1+{r_0\over w_{12}}.
\]
Insert \eqref{eq:newV36-binary-field} and apply
\Cref{thm:newV36-binary-payment}:
\begin{align*}
 \|{\cal R}_{t,2}^{(1)}(z)\|_{\rm cb}
 &\le Cw_{12}+Cw_{12}|z|\\
 &\le Cw_{12}+C(w_{12}+r_0)
 \le C(1+r_0).
\end{align*}
This is \eqref{eq:newV35-WPRAC} at \(m=2\).
\end{proof}

\begin{corollary}[Exact binary bridge-score card]
\label{cor:newV36-binary-score}
\begin{equation}
 \left\|
 \kappa_3^{\nu_{\alpha,t}}
 (D^{R_1},D^{R_2},S_{\alpha,t})
 \right\|_{\rm cb}
 \le Cb_1b_2.
 \label{eq:newV36-binary-score}
\end{equation}
\end{corollary}

\begin{proof}
Apply \Cref{thm:newV35-WPRAC-sufficient} at \(m=2\), or
differentiate \eqref{eq:newV36-binary-field} once and combine with
the contractive heat factor exactly as in its proof.
\end{proof}

\section{Where the first unresolved growth can occur}
\label{sec:newV36-high-block}

\Cref{thm:newV36-marked-block} is uniform as a raw block moment,
but heat normalization alone gives
\begin{equation}
 \|Z_A\|_{\rm cb}
 \le CsB_Ae^{B_A^2}.
 \label{eq:newV36-crude-high-block}
\end{equation}
Since \(B_A<|A|\), the right side can be as large as
\(e^{|A|^2}\).  This is not an exponential arity bound.

\begin{assessment}[The binary closure does not justify termwise
high-block normalization]
\label{ass:newV36-high-block}
At two labels, \(B_A^2<4\), and the negative score norm supplies
the missing \(w_{12}\) exactly.  At growing arity, bounding
\(e^{{\cal K}_A}\), \(M_{t,A}\), and \(N_{t,A}\) separately gives
\eqref{eq:newV36-crude-high-block} and fails WPRAC.  The only
legitimate next question is whether the complete activity
\({\cal A}_{t,A}^{(1)}\) cancels that high-total-spin growth, or
whether the fixed heat normalization must be replaced by a
block-adaptive reference.  V36 makes no such cancellation claim.
\end{assessment}

\begin{proposition}[The first new activity identity is ternary]
\label{prop:newV36-ternary-identity}
For three labels,
\begin{align}
 {\cal A}_{123}^{(1)}
 &=
 Z_{123}
 -\sum_{\substack{1\le i<j\le3\\
                  \{k\}=[3]\setminus\{i,j\}}}
 \left(
 {\cal A}_{ij}^{(1)}R_k
 +{\cal A}_{ij}^{(0)}Z_k
 \right)
 \notag\\
 &\quad
 -Z_1R_2R_3-R_1Z_2R_3-R_1R_2Z_3.
 \label{eq:newV36-ternary-identity}
\end{align}
This is an exact identity before norms.
\end{proposition}

\begin{proof}
At order zero and one in \(\varepsilon\), expand
\[
 {\cal Y}_{123}
 =
 {\cal A}_{123}
 +{\cal A}_{12}{\cal A}_3
 +{\cal A}_{13}{\cal A}_2
 +{\cal A}_{23}{\cal A}_1
 +{\cal A}_1{\cal A}_2{\cal A}_3,
\]
then solve for the marked three-block activity.
\end{proof}

\begin{firewall}[Do not triangle-bound the ternary identity]
\label{fire:newV36-no-ternary-triangle}
The separate terms in \eqref{eq:newV36-ternary-identity} can each
contain the large heat-normalized total-spin factor.  Applying
\eqref{eq:newV36-crude-high-block} term by term would discard the
M\"obius cancellation which the V35 activity construction was
designed to retain.  The identity is a starting point for a
complete spectral-sector computation, not a license for a raw
triangle estimate.
\end{firewall}

\section{V37 programme and claim boundary}
\label{sec:newV36-next}

\begin{programme}[V37: ternary normalized-activity spectrum]
\label{prog:newV36-V37}
\begin{enumerate}[leftmargin=3em]
\item Diagonalize
      \eqref{eq:newV36-ternary-identity} on total-spin sectors,
      retaining the multiplicity matrices produced by different
      fusion orders.
\item Express \(R_A\) and \(Z_A=\partial_tR_A\) through the exact
      bridge Fourier multipliers and determine whether the
      \(e^{B_A^2}\) factors cancel in the complete activity.
\item Prove the ternary WPRAC bound on all thermal representations,
      or exhibit a high-channel counterexample with its exact spin
      labels and multiplier asymptotics.
\item If fixed heat normalization fails, choose a block reference
      time \(s_A\) from the actual multiplier decay and rederive the
      endpoint cancellation without changing the exact cumulant.
\item Preserve the complex radius \(r_0/w_e\); a bound only on the
      real cube does not imply the tree-mark Cauchy estimate.
\item Do not perform the next companion audit before V40.
\end{enumerate}
\end{programme}

\begin{firewall}[Claim boundary after V36]
\label{fire:newV36-boundary}
V36 proves a completely bounded marked-moment row for every tensor
block, computes the exact binary heat-normalized activity, proves
that it carries the missing edge mark, and closes exact binary
WPRAC and the binary bridge-score card.  It does not prove ternary
or all-arity WPRAC.  Hence the exact RGSC, Wilson aggregate strong
card, all-order strong MTA, WUFC, CPM, cutoff removal, reflection
positivity, Osterwalder--Schrader reconstruction, continuum
Yang--Mills measure, and positive Hamiltonian gap remain open.  The
full Yang--Mills mass gap has not been proved.
\end{firewall}

\begin{theorem}[V36 result ledger]
\label{thm:newV36-results}
V36:
\begin{enumerate}[label=\textup{(V36.\arabic*)},leftmargin=6.5em]
\item proves the upper total-Casimir row
      \(\|sC_A^{\rm tot}\|\le B_A^2\);
\item converts the V32 intrinsic negative score norm into the
      complete marked-block estimate \(\|N_{t,A}\|\le CsB_A\);
\item computes the exact binary dual activities;
\item proves that the marked binary activity carries
      \(w_{12}=b_1b_2\);
\item proves exact binary WPRAC and the binary bridge-score card;
\item identifies the superexponential loss in a separated
      high-block heat normalization; and
\item records the exact ternary activity identity which must be
      estimated without termwise norms.
\end{enumerate}
\end{theorem}

\begin{proof}
Use
\Cref{lem:newV36-total-Casimir,
thm:newV36-marked-block,
lem:newV36-binary-activity,
lem:newV36-two-mark,
thm:newV36-binary-payment,
thm:newV36-binary-WPRAC,
cor:newV36-binary-score,
prop:newV36-ternary-identity}.
\end{proof}

\paragraph{Parent-version record.}
V36 was formed from
\texttt{Renewal\_Yang\_Mills\_Mass\_Gap\_Research\_Notes\_V35.tex}.
The V35 parent had \(16\,756\) logical source lines,
\(531\,447\) bytes, and SHA--256
\[
 \texttt{e9eeefd6fb31691988880be539003ed58ed0a95be4a74075bb0ae840870c7a07}.
\]
The next compulsory companion audit remains V40.

% =============================================================================
% END APPENDED FIFTH-SERIES V36 MATERIAL
% =============================================================================

% =============================================================================
% BEGIN APPENDED FIFTH-SERIES V37 MATERIAL
% =============================================================================

\chapter{V37: Uniform three-block expansion and exact ternary WPRAC}
\label{chap:newV37-ternary}

\section{The spectral score as a differentiated heat kernel}
\label{sec:newV37-spectral-score}

Let \(h_s\) be the central heat density with multiplier
\(e^{-sC_2(R)}\).  The V33 spectral score density has multiplier
\[
 {s\over3}\left[(sC_2)^2-5sC_2\right]e^{-sC_2}.
\]
Since
\[
 \partial_sh_s\ \widehat{}\ (R)
 =-C_2(R)e^{-sC_2(R)},
 \qquad
 \partial_s^2h_s\ \widehat{}\ (R)
 =C_2(R)^2e^{-sC_2(R)},
\]
its exact physical-space form is
\begin{equation}
 S_s^{\rm sp}
 =
 {s\over3}
 {s^2\partial_s^2h_s+5s\partial_sh_s\over h_s}.
 \label{eq:newV37-score-differential}
\end{equation}
It is centered because the numerator integrates to zero.

Put
\[
 \alpha_s={1\over2s},
 \qquad
 Y_s=\sqrt{\alpha_s}\,\theta\boldsymbol n.
\]
For \(s=\tau_\alpha^\sharp\), one has
\(\alpha_s=\alpha+1/3\).

\begin{lemma}[Uniform tangent form of the spectral score]
\label{lem:newV37-spectral-tangent}
For every fixed \(p<\infty\),
\begin{equation}
 \left\|
 S_s^{\rm sp}
 -{|Y_s|^4-15\over24\alpha_s}
 \right\|_{L^p(\mu_s^H)}
 \le C_ps^2.
 \label{eq:newV37-spectral-tangent}
\end{equation}
The same estimate holds in \(L^p(\nu_{\alpha,t})\), uniformly in
\(t\), after increasing \(C_p\).
\end{lemma}

\begin{proof}
On \(0\le\theta\le\pi/2\), insert the dominant wrapped image
from \eqref{eq:newV30-image}.  Up to a scalar normalization and an
exponentially small image error,
\[
 h_s(\theta)
 =
 c_s{\theta\over\sin\theta}
 e^{-\theta^2/(4s)}.
\]
For the Euclidean factor
\[
 g_s(y)=(4\pi s)^{-3/2}e^{-|y|^2/(4s)},
\]
direct differentiation gives
\begin{align*}
 {s^2\partial_s^2g_s+5s\partial_sg_s\over g_s}
 &=
 {1\over4}\left(
 \left|{y\over\sqrt{2s}}\right|^4-15
 \right).
\end{align*}
Multiplication by \(s/3\) is therefore
\[
 {s\over12}(|Y_s|^4-15)
 =
 {|Y_s|^4-15\over24\alpha_s}.
\]

The derivatives of the smooth heat normalization and of
\(\theta/\sin\theta\) contribute
\[
 O\bigl(s^2(1+|Y_s|^6)\bigr)
\]
after the exterior factor \(s\) in
\eqref{eq:newV37-score-differential}; differentiated remote images
are \(O(e^{-c/s})\) on the inner half.  On the antipodal half, the
paired-image bounds of V31 and the heat distance moments of V30
give an exponentially small \(L^p\) contribution after a
polynomial loss in \(s^{-1}\).  This proves the heat-law estimate.

For the bridge law, use the uniform moment bound
\eqref{eq:newV30-bridge-moments}.  The same local expansion is a
pointwise identity for \(S_s^{\rm sp}\), and the antipodal bridge
tail is still exponentially small by the Hellinger/concentration
theorem of V30.  Thus the same polynomial remainder is
\(O_{L^p}(s^2)\).
\end{proof}

\begin{lemma}[Spectral and exact bridge scores agree to second
order]
\label{lem:newV37-score-comparison}
For every fixed \(p<\infty\),
\begin{equation}
 \|S_{\alpha,t}-S_s^{\rm sp}\|_{L^p(\nu_{\alpha,t})}
 \le C_ps^2,
 \qquad s=\tau_\alpha^\sharp,
 \label{eq:newV37-score-comparison}
\end{equation}
uniformly in \(t\in[0,1]\).
\end{lemma}

\begin{proof}
V32 gives
\[
 S_{\alpha,t}
 =
 {|Y_\alpha|^4-15\over24\alpha}
 +O_{L^p}(\alpha^{-2}),
 \qquad
 Y_\alpha=\sqrt\alpha\,\theta\boldsymbol n.
\]
Because \(\alpha_s=\alpha+1/3\), replacing
\((\alpha,Y_\alpha)\) by \((\alpha_s,Y_s)\) changes the displayed
quartic by \(O_{L^p}(\alpha^{-2})=O_{L^p}(s^2)\).
Apply \Cref{lem:newV37-spectral-tangent}.
\end{proof}

\section{The bridge and heat laws are close in the needed weighted norm}
\label{sec:newV37-density}

\begin{lemma}[Total and score-weighted variation]
\label{lem:newV37-density-close}
Uniformly in \(t\in[0,1]\),
\begin{align}
 \|\nu_{\alpha,t}-\mu_s^H\|_{\rm TV}
 &\le Cs,
 \label{eq:newV37-density-close}\\
 \int_{SU(2)}|S_s^{\rm sp}|\,
 |d\nu_{\alpha,t}-d\mu_s^H|
 &\le Cs^2.
 \label{eq:newV37-score-weighted-variation}
\end{align}
\end{lemma}

\begin{proof}
In the scaled variable \(Y_s\), split at
\(|Y_s|=s^{-1/8}\).  On the inner region, the V32 logarithmic
profile and \(\alpha_s=\alpha+1/3\) give, after subtraction of a
scalar \(c_\alpha\),
\[
 \ell_\alpha-c_\alpha
 =
 {s\over12}(|Y_s|^4-15)
 +O\left(s^2(1+|Y_s|^6)\right)
\]
in every fixed polynomially weighted \(L^1\) norm of the heat
Gaussian core.  Taylor's formula for the normalized tilt
\[
 {d\nu_{\alpha,t}\over d\mu_s^H}
 =
 {\exp(t\ell_\alpha)\over
  \mathbb E_{\mu_s^H}\exp(t\ell_\alpha)}
\]
therefore gives an \(O(s)\) unweighted \(L^1\) difference and,
using \Cref{lem:newV37-spectral-tangent}, an \(O(s^2)\) difference
after multiplication by \(|S_s^{\rm sp}|\).

On the complementary region, the heat and bridge confinement
estimates of V30 give
\[
 \mu_s^H(|Y_s|>s^{-1/8})
 +\nu_{\alpha,t}(|Y_s|>s^{-1/8})
 \le Ce^{-cs^{-1/4}}.
\]
The wrapped-image expression differentiated in
\eqref{eq:newV37-score-differential} bounds
\(|S_s^{\rm sp}|\) globally by a fixed power of \(s^{-1}\).
Thus both the unweighted and score-weighted tail errors are smaller
than any power of \(s\).  Combining the two regions proves
\eqref{eq:newV37-density-close} and
\eqref{eq:newV37-score-weighted-variation}.
\end{proof}

\begin{remark}[Why an \(L^2\) density ratio is not used]
\label{rem:newV37-no-L2-ratio}
Near the antipode the Wilson endpoint can be exponentially heavier
than the calibrated heat law, even though both tails are
exponentially small.  Therefore an \(L^2(\mu_s^H)\) bound on the
full density ratio is neither needed nor asserted.  The signed
measure estimates above control exactly the bounded Fourier row and
its one spectral-score weight.
\end{remark}

\section{Uniform thermal Fourier expansion on three blocks}
\label{sec:newV37-Fourier}

Define the spectral polynomial
\begin{equation}
 \Phi_s(K):={s\over3}(K^2-5K).
 \label{eq:newV37-Phi}
\end{equation}
For a tensor block \(A\), the heat spectral marked moment is exactly
\[
 N_A^{\rm sp}
 =
 \Phi_s({\cal K}_A)e^{-{\cal K}_A}.
\]

\begin{theorem}[Uniform marked multiplier expansion]
\label{thm:newV37-marked-expansion}
For every tensor block \(A\) with \(|A|\le3\),
\begin{align}
 N_{t,A}
 &=
 \Phi_s({\cal K}_A)e^{-{\cal K}_A}
 +{\cal E}_{t,A}^{(1)},
 \label{eq:newV37-N-expansion}\\
 \|{\cal E}_{t,A}^{(1)}\|_{\rm cb}
 &\le Cs^2.
 \label{eq:newV37-N-error}
\end{align}
Consequently,
\begin{align}
 M_{t,A}
 &=
 e^{-{\cal K}_A}
 +t\Phi_s({\cal K}_A)e^{-{\cal K}_A}
 +{\cal E}_{t,A}^{(0)},
 \label{eq:newV37-M-expansion}\\
 \|{\cal E}_{t,A}^{(0)}\|_{\rm cb}
 &\le Cs^2.
 \label{eq:newV37-M-error}
\end{align}
All constants are uniform in \(t\) and in the thermal
representations.
\end{theorem}

\begin{proof}
Every tensor representation matrix is unitary, so its operator and
complete norms are one.  Split
\begin{align*}
 N_{t,A}-N_A^{\rm sp}
 &=
 \mathbb E_{\nu_{\alpha,t}}
 [(S_{\alpha,t}-S_s^{\rm sp})D^A]\\
 &\quad+
 \int_{SU(2)}S_s^{\rm sp}D^A
 (d\nu_{\alpha,t}-d\mu_s^H).
\end{align*}
The first term is \(O(s^2)\) by
\Cref{lem:newV37-score-comparison}.  The second is \(O(s^2)\)
directly by \eqref{eq:newV37-score-weighted-variation}.  Testing
against amplified unit vectors gives the completely bounded form
\eqref{eq:newV37-N-error}.

At \(t=0\), \(M_{0,A}=e^{-{\cal K}_A}\).  Integrate
\(\partial_tM_{t,A}=N_{t,A}\) from
\Cref{lem:newV35-marked-derivative}; the leading marked moment is
independent of \(t\), and the error remains \(O(s^2)\).  This proves
\eqref{eq:newV37-M-expansion}.
\end{proof}

\begin{corollary}[Heat-normalized three-block expansion]
\label{cor:newV37-normalized-expansion}
For \(|A|\le3\),
\begin{align}
 R_A
 &=
 I+t\Phi_s({\cal K}_A)+O_{\rm cb}(s^2),
 \label{eq:newV37-R-expansion}\\
 Z_A
 &=
 \Phi_s({\cal K}_A)+O_{\rm cb}(s^2).
 \label{eq:newV37-Z-expansion}
\end{align}
The errors are uniform in all thermal representations.
\end{corollary}

\begin{proof}
By \Cref{lem:newV36-total-Casimir},
\[
 \|{\cal K}_A\|\le B_A^2<9.
\]
Multiply
\eqref{eq:newV37-N-expansion} and
\eqref{eq:newV37-M-expansion} by
\(e^{{\cal K}_A}\).  The factor is at most \(e^9\), which is
absorbed into the numerical error constant.
\end{proof}

\section{First-order activity is a Boolean spectral difference}
\label{sec:newV37-activity}

For \(B\subseteq[3]\), define
\[
 \Delta_B\Phi_s
 :=
 \sum_{C\subseteq B}
 (-1)^{|B|-|C|}\Phi_s({\cal K}_C),
 \qquad
 \Phi_s({\cal K}_\varnothing)=0.
\]

\begin{lemma}[Uniform activity expansion through three labels]
\label{lem:newV37-activity-expansion}
For every nonempty \(B\subseteq[3]\),
\begin{align}
 {\cal A}_{t,B}^{(0)}
 &=
 \mathbf1_{\{|B|=1\}}I
 +t\Delta_B\Phi_s
 +O_{\rm cb}(s^2),
 \label{eq:newV37-A0-expansion}\\
 {\cal A}_{t,B}^{(1)}
 &=
 \Delta_B\Phi_s
 +O_{\rm cb}(s^2).
 \label{eq:newV37-A1-expansion}
\end{align}
\end{lemma}

\begin{proof}
Insert
\Cref{cor:newV37-normalized-expansion} into the finite
moment--cumulant inversion
\eqref{eq:newV35-activity-inverse}.  At order zero in \(s\), all
normalized moments are \(I\); their only nonzero activities are the
singleton identities.  The coefficient linear in the perturbation
\(\Phi_s\) therefore obeys
\[
 \Phi_s({\cal K}_B)
 =
 \sum_{\varnothing\ne C\subseteq B}
 {\cal A}_{C,\rm linear},
\]
and Boolean inversion gives \(\Delta_B\Phi_s\).  Every product
containing two perturbations is \(O(s^2)\).  There are only the
partitions of sets of size at most three, so their total remainder
is \(O_{\rm cb}(s^2)\) uniformly.
\end{proof}

\begin{lemma}[Two- and three-block root payment]
\label{lem:newV37-activity-payment}
Let \(r(B)=\min B\).  For \(2\le|B|\le3\),
\begin{align}
 \|\Delta_B\Phi_s\|_{\rm cb}
 &\le
 C\prod_{e\in{\cal A}(B)}w_e,
 \label{eq:newV37-Phi-payment}\\
 \|{\cal A}_{t,B}^{(0)}\|_{\rm cb}
 +\|{\cal A}_{t,B}^{(1)}\|_{\rm cb}
 &\le
 C\prod_{e\in{\cal A}(B)}w_e.
 \label{eq:newV37-activity-payment}
\end{align}
\end{lemma}

\begin{proof}
Expand \(\Phi_s\) into its linear and quadratic Casimir words.
The Boolean difference retains exactly the words whose vertex
support is \(B\).  The V34 degree-\(q\) root-payment lemma, with
\(q=1,2\), proves \eqref{eq:newV37-Phi-payment}.

For \(|B|=2\), the star weight is \(b_ib_j\ge s\), hence it
dominates the \(O(s^2)\) remainder in
\Cref{lem:newV37-activity-expansion}.  For
\(B=\{r,j,k\}\), its canonical-star weight is
\[
 w_{rj}w_{rk}=b_r^2b_jb_k\ge s^2,
\]
because every \(b_i\ge\sqrt s\).  Thus it also dominates the
remainder.  Apply
\eqref{eq:newV37-A0-expansion}--%
\eqref{eq:newV37-A1-expansion}.
\end{proof}

\section{Exact ternary WPRAC}
\label{sec:newV37-WPRAC}

For three labels, the marked resummed activity field is
\begin{align}
 {\cal R}_{t,3}^{(1)}(\boldsymbol x)
 &=
 \sum_{i=1}^3
 Z_i\prod_{j\ne i}R_j
 \notag\\
 &\quad+
 \sum_{\substack{1\le i<j\le3\\
                  \{k\}=[3]\setminus\{i,j\}}}
 x_{ij}
 \left(
 {\cal A}_{ij}^{(1)}R_k
 +{\cal A}_{ij}^{(0)}Z_k
 \right)
 \notag\\
 &\quad+
 x_{12}x_{13}{\cal A}_{123}^{(1)}.
 \label{eq:newV37-ternary-field}
\end{align}
The last gate uses the canonical root \(1=\min[3]\).

\begin{theorem}[Ternary WPRAC]
\label{thm:newV37-ternary-WPRAC}
For every fixed \(r_0>0\), every labelled tree \(T\) on three
vertices, and every real forest point,
\begin{equation}
 \sup_{\boldsymbol z\in{\mathbb D}_{T,r_0}}
 \|{\cal R}_{t,3}^{(1)}(\boldsymbol z)\|_{\rm cb}
 \le C(1+r_0)^2.
 \label{eq:newV37-ternary-WPRAC}
\end{equation}
Hence the exact bridge satisfies WPRAC at \(m=3\).
\end{theorem}

\begin{proof}
The singleton normalized moments are bounded by a numerical
constant using
\Cref{cor:newV37-normalized-expansion}; hence the first line of
\eqref{eq:newV37-ternary-field} is bounded.

By \Cref{lem:newV37-activity-payment}, the coefficient of
\(x_{ij}\) is at most \(Cw_{ij}\).  If \(ij\in E(T)\), then on the
weighted polydisc
\[
 w_{ij}|z_{ij}|\le w_{ij}+r_0\le1+r_0.
\]
If \(ij\notin E(T)\), the variable stays at its real forest value in
\([0,1]\).

The coefficient of the last line is at most
\(Cw_{12}w_{13}\).  Apply the same payment separately to every one
of its canonical gate edges that belongs to \(T\); every off-tree
gate remains in \([0,1]\).  Thus the last line is bounded by
\(C(1+r_0)^2\).  There are only five partition terms, so their sum
obeys \eqref{eq:newV37-ternary-WPRAC}.
\end{proof}

\begin{corollary}[Exact ternary bridge-score card]
\label{cor:newV37-ternary-score}
\begin{equation}
 \left\|
 \kappa_4^{\nu_{\alpha,t}}
 (D^{R_1},D^{R_2},D^{R_3},S_{\alpha,t})
 \right\|_{\rm cb}
 \le
 K^3
 \left(\prod_{i=1}^3b_i\right)
 \left(\sum_{i=1}^3b_i\right).
 \label{eq:newV37-ternary-score}
\end{equation}
\end{corollary}

\begin{proof}
Apply \Cref{thm:newV35-WPRAC-sufficient} at \(m=3\).
\end{proof}

\begin{assessment}[Why the fixed-three proof closes]
\label{ass:newV37-fixed-three}
The high-total-spin heat normalization is harmless on at most three
thermal inputs because \(\|{\cal K}_A\|<9\).  More importantly, the
exact bridge differs from the spectral quartic model by
\(O_{\rm cb}(s^2)\) in every marked block moment.  The smallest
three-vertex star also has size \(s^2\), so that error has exactly
the missing incidence power.  At growing arity the smallest star is
\(s^{m-1}\), while a fixed \(O(s^2)\) remainder is insufficient.
Thus V37 is a true ternary theorem, not an all-order inference.
\end{assessment}

\section{V38 programme and claim boundary}
\label{sec:newV37-next}

\begin{programme}[V38: order-adaptive uniform block expansion]
\label{prog:newV37-V38}
\begin{enumerate}[leftmargin=3em]
\item Extend the differentiated wrapped-kernel calculation to a
      spectral jet of order \(N\), with a remainder uniform for
      \(\|{\cal K}_A\|\le B_A^2\).
\item Choose \(N\) from the activity support size, so the remainder
      begins at the minimum star weight \(s^{|A|-1}\) when all
      marks are at their thermal floor.
\item Sum the complete jet before applying the V34 root-payment
      compiler; do not take norms degree by degree through an
      infinite series.
\item Track the dependence on \(B_A\) and test whether the heat
      normalization produces \(e^{cB_A^2}\) or a canceling
      \(e^{-cB_A^2}\) Fourier tail.
\item If the optimal truncation still has a base growing with
      \(|A|\), replace the fixed heat reference by an adaptive
      reference before attempting WPRAC.
\item Preserve the established binary and ternary cards as exact
      base cases; do not rederive them by raw random-cluster trees.
\end{enumerate}
\end{programme}

\begin{firewall}[Claim boundary after V37]
\label{fire:newV37-boundary}
V37 derives the exact differentiated-heat representation of the
spectral score, proves its \(O(s^2)\) agreement with the geometric
bridge score, obtains uniform marked and unmarked multiplier
expansions on every three-input thermal block, and closes ternary
WPRAC and the exact ternary bridge-score card.  It does not prove
WPRAC at growing arity.  Hence the all-order RGSC, Wilson aggregate
strong card, all-order strong MTA, WUFC, CPM, cutoff removal,
reflection positivity, Osterwalder--Schrader reconstruction,
continuum Yang--Mills measure, and positive Hamiltonian gap remain
open.  The full Yang--Mills mass gap has not been proved.
\end{firewall}

\begin{theorem}[V37 result ledger]
\label{thm:newV37-results}
V37:
\begin{enumerate}[label=\textup{(V37.\arabic*)},leftmargin=6.5em]
\item writes the spectral score as an exact differentiated heat
      kernel;
\item proves its uniform \(O_{L^p}(s^2)\) tangent remainder;
\item proves first-order total variation and second-order
      score-weighted variation between every bridge law and heat;
\item proves marked and unmarked multiplier expansions uniformly on
      all blocks of size at most three;
\item proves that the first activity correction is the Boolean
      spectral difference;
\item proves the exact two- and three-block activity payments; and
\item closes ternary WPRAC and the ternary bridge-score card.
\end{enumerate}
\end{theorem}

\begin{proof}
Use
\Cref{lem:newV37-spectral-tangent,
lem:newV37-score-comparison,
lem:newV37-density-close,
thm:newV37-marked-expansion,
cor:newV37-normalized-expansion,
lem:newV37-activity-expansion,
lem:newV37-activity-payment,
thm:newV37-ternary-WPRAC,
cor:newV37-ternary-score}.
\end{proof}

\paragraph{Parent-version record.}
V37 was formed from
\texttt{Renewal\_Yang\_Mills\_Mass\_Gap\_Research\_Notes\_V36.tex}.
The V36 parent had \(17\,182\) logical source lines,
\(544\,062\) bytes, and SHA--256
\[
 \texttt{eebfa58f356e18179b6355f2e477662975272c85658560f11f6b3a22ced89549}.
\]
The next compulsory companion audit remains V40.

% =============================================================================
% END APPENDED FIFTH-SERIES V37 MATERIAL
% =============================================================================

% =============================================================================
% BEGIN APPENDED FIFTH-SERIES V38 MATERIAL
% =============================================================================

\chapter{V38: Double-factorial entire roots close the infinite-degree forest}
\label{chap:newV38-entire}

\section{Correction of the coefficient threshold}
\label{sec:newV38-threshold}

V34 proves that geometric decay in the spectral degree is
insufficient after degreewise norms, and that coefficients of size
\(A^q/q!\) still produce an \(e^{cm^2}\)-type arity majorant.
That statement remains correct.  It does not cover the stronger
even-entire decay \(A^q/(2q)!\).  On the weighted polydisc of V35,
the latter exactly compensates the \(m^{2q}\) choices of the
ordered indices in a degree-\(q\) Casimir word.  This chapter proves
the corresponding infinite-degree theorem.

\section{The double-factorial spectral class}
\label{sec:newV38-class}

Let
\begin{equation}
 \Phi_s(K)
 :=
 \sum_{q=1}^\infty a_q(s)s^{q-1}K^q
 \label{eq:newV38-Phi}
\end{equation}
with
\begin{equation}
 |a_q(s)|
 \le {A L^{2q}\over(2q)!}
 \qquad(0<s<1,\ q\ge1)
 \label{eq:newV38-double-factorial}
\end{equation}
for fixed \(A,L<\infty\).  The series is entire in \(K\), and
\(\Phi_s(0)=0\).  Its growth is at most
\[
 {A\over s}\exp(L\sqrt{s|K|}),
\]
so \(e^{-K}\Phi_s(K)\) is rapidly decreasing on the positive
Casimir spectrum.

Define the centered smooth heat-law score \(S_s^\Phi\) by the
Fourier multipliers
\begin{equation}
 \widehat{(h_sS_s^\Phi)}(R)
 =
 e^{-K_R}\Phi_s(K_R),
 \qquad
 K_R=sC_2(R).
 \label{eq:newV38-score-multiplier}
\end{equation}
The trivial multiplier vanishes.

For \(q\ge1\), let \({\cal Q}_m^{[q]}\) be the block-local power
from V34 and set
\begin{equation}
 {\cal Q}_m^\Phi(\boldsymbol x)
 :=
 \sum_{q=1}^\infty
 a_q(s)s^{q-1}{\cal Q}_m^{[q]}(\boldsymbol x).
 \label{eq:newV38-QPhi}
\end{equation}

\begin{lemma}[Normal convergence of the resummed root]
\label{lem:newV38-normal-convergence}
For fixed \(m\), representations, and every compact complex
edge-polydisc, the series
\eqref{eq:newV38-QPhi} converges absolutely in completely bounded
norm, locally uniformly with all edge derivatives.  Hence it
defines an entire operator-valued function.
\end{lemma}

\begin{proof}
On a compact polydisc, all gates have a common finite bound
\(H^{2q-1}\).  There are \(m^{2q}\) ordered words and
\[
 \|A_{i_1j_1}\cdots A_{i_qj_q}\|
 \le\prod_{\nu=1}^q b_{i_\nu}b_{j_\nu}
 \le1.
\]
The \(q\)-th term is therefore at most
\[
 A\,{(L mH)^{2q}\over(2q)!}
\]
up to a fixed factor.  The even exponential series converges
normally.  Each fixed edge derivative removes gate factors and
only changes the harmless power of \(H\), so the same argument
applies to all derivatives.  Amplification does not change any
factor norm.
\end{proof}

\begin{lemma}[Exact endpoint localization of the entire root]
\label{lem:newV38-endpoint}
For every partition \(\rho\),
\begin{equation}
 {\cal Q}_m^\Phi(\boldsymbol x^\rho)
 =
 \sum_{B\in\rho}\Phi_s({\cal K}_B).
 \label{eq:newV38-endpoint}
\end{equation}
\end{lemma}

\begin{proof}
By normal convergence, sum the V34 endpoint identity
\[
 {\cal Q}_m^{[q]}(\boldsymbol x^\rho)
 =\sum_{B\in\rho}{\cal K}_B^q
\]
term by term.
\end{proof}

Define
\begin{equation}
 {\cal G}_m^\Phi(\boldsymbol x)
 :=
 {\cal Q}_m^\Phi(\boldsymbol x)e^{-{\cal K}(\boldsymbol x)}.
 \label{eq:newV38-GPhi}
\end{equation}

\begin{theorem}[Exact entire-root forest]
\label{thm:newV38-entire-forest}
For \(m\ge2\),
\begin{align}
 &\kappa_{m+1}^{\mu_s^H}
 (D^{R_1},\ldots,D^{R_m},S_s^\Phi)
 \notag\\
 &\qquad=
 \sum_{T\in\mathfrak T_m}
 \int_{[0,1]^{m-1}}
 \partial_{E(T)}
 {\cal G}_m^\Phi(\boldsymbol x^T(\boldsymbol u))
 \,d\boldsymbol u.
 \label{eq:newV38-entire-forest}
\end{align}
\end{theorem}

\begin{proof}
At a partition endpoint, the heat exponential factors over blocks.
Use \Cref{lem:newV38-endpoint}:
\[
 {\cal G}_m^\Phi(\boldsymbol x^\rho)
 =
 \sum_{B\in\rho}
 \left[
 \Phi_s({\cal K}_B)e^{-{\cal K}_B}
 \right]
 \bigotimes_{C\in\rho\setminus\{B\}}e^{-{\cal K}_C}.
\]
This is the coefficient of one centered marked heat moment in the
dual endpoint product.  The dual-number BKAR argument of V33 and
V35 gives \eqref{eq:newV38-entire-forest}.  Normal convergence
from \Cref{lem:newV38-normal-convergence} justifies all
differentiation and integration.
\end{proof}

\section{The infinite root satisfies WPRAC}
\label{sec:newV38-WPRAC}

\begin{theorem}[Double-factorial weighted-polydisc bound]
\label{thm:newV38-polydisc}
Fix \(r_0>0\).  For every labelled tree \(T\), every forest point,
and every thermal representation family,
\begin{equation}
 \sup_{\boldsymbol z\in{\mathbb D}_{T,r_0}}
 \|{\cal Q}_m^\Phi(\boldsymbol z)\|_{\rm cb}
 \le
 A\exp\bigl(L(1+r_0)m\bigr)
 \label{eq:newV38-polydisc}
\end{equation}
after a harmless enlargement of \(A,L\).
\end{theorem}

\begin{proof}
Consider one ordered word of degree \(q\) and let \(B\) be its
vertex support.  The degree-\(q\) root payment
\eqref{eq:newV34-q-payment} bounds its operator coefficient,
including \(s^{q-1}\), by the full canonical-star weight
\[
 \prod_{e\in{\cal A}(B)}w_e.
\]
If a star edge lies in \(T\), then on the weighted polydisc
\[
 w_e|z_e|
 \le w_e+r_0\le1+r_0.
\]
If it does not lie in \(T\), its variable stays at a real forest
value in \([0,1]\).  Thus the paying star weight times its complete
gate has modulus at most
\((1+r_0)^{2q-1}\).

There are \(m^{2q}\) ordered index words.  Summing them before
taking the degree sum and applying
\eqref{eq:newV38-double-factorial} gives
\begin{align*}
 \|{\cal Q}_m^\Phi(\boldsymbol z)\|_{\rm cb}
 &\le
 A\sum_{q\ge1}
 {L^{2q}m^{2q}(1+r_0)^{2q-1}\over(2q)!}\\
 &\le
 A\exp\bigl(L(1+r_0)m\bigr).
\end{align*}
This is uniform in \(\boldsymbol z\) and all amplifications.
\end{proof}

\begin{corollary}[Infinite-degree heat score card]
\label{cor:newV38-entire-card}
Every score in the class
\eqref{eq:newV38-double-factorial} satisfies WPRAC and
\begin{align}
 &\left\|
 \kappa_{m+1}^{\mu_s^H}
 (D^{R_1},\ldots,D^{R_m},S_s^\Phi)
 \right\|_{\rm cb}
 \notag\\
 &\qquad\le
 K^m
 \left(\prod_i b_i\right)
 \left(\sum_i b_i\right)^{m-2}.
 \label{eq:newV38-entire-card}
\end{align}
\end{corollary}

\begin{proof}
\Cref{thm:newV38-polydisc} is WPRAC for the heat endpoint, where
the unmarked normalized activities are singleton identities and
the marked activity field is the Boolean support grouping of
\({\cal Q}_m^\Phi\).  Apply
\Cref{lem:newV35-complex-heat} and the weighted Cauchy proof of
\Cref{thm:newV35-WPRAC-sufficient}, or apply it directly to
\eqref{eq:newV38-entire-forest}.
\end{proof}

\section{Sharpness of the coefficient order for this compiler}
\label{sec:newV38-sharpness}

\begin{proposition}[General gamma-order census]
\label{prop:newV38-gamma-order}
Suppose the coefficient majorant has the form
\[
 |a_q(s)|\le {A L^{2q}\over\Gamma(\beta q+1)}
\]
with \(\beta>0\).  The ordered-word polydisc census is bounded by
\[
 \sum_{q\ge1}{(Cm)^{2q}\over\Gamma(\beta q+1)}
 =
 \exp\bigl(O(m^{2/\beta})\bigr).
 \label{eq:newV38-gamma-growth}
\]
Consequently this method has an exponential-in-\(m\) bound when
\(\beta\ge2\).  When \(\beta<2\), its positive coefficient
majorant is superexponential in \(m\).
\end{proposition}

\begin{proof}
Stirling's formula gives
\[
 \log\Gamma(\beta q+1)
 =
 \beta q\log(\beta q)-\beta q+O(\log(q+1)).
\]
Maximizing
\[
 2q\log(Cm)-\log\Gamma(\beta q+1)
\]
places \(q\) at order \(m^{2/\beta}\), and the maximal logarithm
has the same order.  Summing the log-concave tails changes only the
constant in the exponent.  If \(\beta\ge2\), this is \(O(m)\);
if \(\beta<2\), it grows faster than linearly.
\end{proof}

\begin{corollary}[The V34 factorial no-go and V38 closure are
compatible]
\label{cor:newV38-compatible}
The coefficient \(1/q!\) corresponds to \(\beta=1\) and gives
\(e^{O(m^2)}\), as in V34.  The coefficient \(1/(2q)!\) has
\(\beta=2\) and gives \(e^{O(m)}\), as in
\Cref{thm:newV38-polydisc}.
\end{corollary}

\section{A nonpolynomial model closed by the theorem}
\label{sec:newV38-model}

For \(\lambda>0\), set
\begin{equation}
 \Phi_{s,\lambda}^{\rm ev}(K)
 :=
 {1\over s}
 \left[
 \cosh(\lambda\sqrt{sK})
 -1-{\lambda^2sK\over2}
 \right].
 \label{eq:newV38-even-model}
\end{equation}
The removable value at \(sK=0\) is understood.

\begin{proposition}[Calibrated even-entire model]
\label{prop:newV38-even-model}
\begin{equation}
 \Phi_{s,\lambda}^{\rm ev}(K)
 =
 \sum_{q=2}^\infty
 {\lambda^{2q}\over(2q)!}s^{q-1}K^q.
 \label{eq:newV38-even-series}
\end{equation}
Its degree-two root is nonzero, it has infinitely many higher
degrees, and its exact heat score satisfies the all-arity card
\eqref{eq:newV38-entire-card}.
\end{proposition}

\begin{proof}
Expand the hyperbolic cosine and remove its constant and quadratic
terms.  The coefficients satisfy
\eqref{eq:newV38-double-factorial} with \(A=1\), \(L=\lambda\).
Apply \Cref{cor:newV38-entire-card}.
\end{proof}

\begin{assessment}[What remains for the geometric bridge]
\label{ass:newV38-remaining}
V38 removes infinite root degree as an obstruction when the
complete heat-normalized spectral coefficients have even-entire
\((2q)!\) decay.  The physical Taylor series of the cosine action
has this formal decay before the Gaussian/heat transform.  It is
not yet proved that the \emph{complete} calibrated multiplier,
including the logarithmic Haar term, paired heat images,
normalization, and the \(t\)-dependent unmarked activities, obeys
\eqref{eq:newV38-double-factorial}.  That assertion must be proved
after the cancellations are assembled; it cannot be inferred from
the physical-space cosine series alone.
\end{assessment}

\begin{definition}[Exact even-entire multiplier card]
\label{def:newV38-EEMC}
The heat-endpoint EEMC is the assertion that its exact
heat-normalized marked multiplier admits
\eqref{eq:newV38-Phi} with
\eqref{eq:newV38-double-factorial}, uniformly in \(s\).  The
bridge EEMC is the corresponding statement for the complete
normalized dual activities of V35, after their unmarked and marked
parts have been assembled.
\end{definition}

\section{V39 programme and claim boundary}
\label{sec:newV38-next}

\begin{programme}[V39: exact multiplier coefficient test]
\label{prog:newV38-V39}
\begin{enumerate}[leftmargin=3em]
\item Insert the exact decomposition
      \(L_\alpha-L_\alpha(0)=\alpha R+B+E_\alpha\) from V32 into
      the heat Fourier transform, before centering.
\item Use radial heat differentiation to compute the spectral
      coefficient sequence of the entire cosine remainder
      \(\alpha R\), retaining the \((2q)!\) denominator.
\item Treat \(B=\theta^2/6-\log(\theta/\sin\theta)\) as one
      assembled analytic function and determine its exact gamma
      order after the heat transform.
\item Bound the paired-image term by a separate exponentially small
      activity on the full thermal spectral window; do not replace a
      derivative bound by tail mass alone.
\item Decide the heat-endpoint EEMC.  If the Haar term has
      \(\beta<2\), absorb it into a modified reference rather than
      iterating a failed coefficient estimate.
\item Only after the endpoint test, propagate the multiplier class
      along \(t\) through the exact dual activity recursion of V35.
\end{enumerate}
\end{programme}

\begin{firewall}[Claim boundary after V38]
\label{fire:newV38-boundary}
V38 proves normal convergence, exact endpoint localization, an
exact forest, WPRAC, and the all-arity raw score card for every
double-factorial entire heat spectral root.  It proves the gamma
order \(\beta=2\) is the threshold for this positive ordered-word
compiler.  It does not prove that the exact heat--Wilson bridge
belongs to this class.  Hence exact all-order RGSC, the Wilson
aggregate strong card, all-order strong MTA, WUFC, CPM, cutoff
removal, reflection positivity, Osterwalder--Schrader
reconstruction, a continuum Yang--Mills measure, and a positive
Hamiltonian gap remain open.  The full Yang--Mills mass gap has not
been proved.
\end{firewall}

\begin{theorem}[V38 result ledger]
\label{thm:newV38-results}
V38:
\begin{enumerate}[label=\textup{(V38.\arabic*)},leftmargin=6.5em]
\item identifies \((2q)!\) decay as the missing coefficient class
      not covered by the V34 no-go;
\item constructs and normally sums the infinite block-local root;
\item proves its exact marked heat forest;
\item proves its weighted-polydisc bound and all-arity score card;
\item proves the general gamma-order arity census and its sharp
      \(\beta=2\) threshold for this compiler; and
\item closes a genuinely nonpolynomial calibrated even-entire
      spectral model.
\end{enumerate}
\end{theorem}

\begin{proof}
Use
\Cref{lem:newV38-normal-convergence,
lem:newV38-endpoint,
thm:newV38-entire-forest,
thm:newV38-polydisc,
cor:newV38-entire-card,
prop:newV38-gamma-order,
prop:newV38-even-model}.
\end{proof}

\paragraph{Parent-version record.}
V38 was formed from
\texttt{Renewal\_Yang\_Mills\_Mass\_Gap\_Research\_Notes\_V37.tex}.
The V37 parent had \(17\,743\) logical source lines,
\(560\,879\) bytes, and SHA--256
\[
 \texttt{9d5fcb54b3593d914777ef97e444d4173984b109e5075cccf3702c1af361a0e2}.
\]
The next compulsory companion audit remains V40.

% =============================================================================
% END APPENDED FIFTH-SERIES V38 MATERIAL
% =============================================================================

% =============================================================================
% BEGIN APPENDED FIFTH-SERIES V39 MATERIAL
% =============================================================================

\chapter{V39: Gaussian coefficient test and the inseparable wrapped block}
\label{chap:newV39-coefficients}

\section{The exact radial Gaussian transform}
\label{sec:newV39-transform}

Let \(Y\sim N(0,I_3)\), let
\[
 K={|\xi|^2\over2},
\]
and define the heat-normalized radial Fourier transform
\begin{equation}
 {\cal T}F(K)
 :=
 e^K\mathbb E[F(Y)e^{i\xi\cdot Y}].
 \label{eq:newV39-transform}
\end{equation}

\begin{lemma}[Laguerre transform of every radial monomial]
\label{lem:newV39-Laguerre}
For every integer \(q\ge0\),
\begin{equation}
 {\cal T}(|Y|^{2q})(K)
 =
 2^qq!L_q^{(1/2)}(K),
 \label{eq:newV39-Laguerre}
\end{equation}
where
\begin{equation}
 L_q^{(1/2)}(K)
 =
 \sum_{k=0}^q
 {(-1)^k\over k!}
 {q+1/2\choose q-k}K^k.
 \label{eq:newV39-Laguerre-coefficients}
\end{equation}
\end{lemma}

\begin{proof}
The Gaussian characteristic function is \(e^{-K}\), and
\[
 \mathbb E[|Y|^{2q}e^{i\xi\cdot Y}]
 =
 (-\Delta_\xi)^qe^{-|\xi|^2/2}.
\]
The radial Laguerre identity in three dimensions gives
\[
 (-\Delta_\xi)^qe^{-|\xi|^2/2}
 =
 2^qq!L_q^{(1/2)}(|\xi|^2/2)e^{-|\xi|^2/2}.
\]
Divide by \(e^{-K}\).  The finite coefficient formula is the
standard defining expansion of the generalized Laguerre
polynomial and can also be checked by induction in \(q\).
\end{proof}

\section{The calibrated cosine remainder retains even-entire decay}
\label{sec:newV39-cosine}

Put \(\theta=\sqrt{2s}|Y|\),
\(\alpha_s=(2s)^{-1}\), and retain
\[
 R(\theta)={\theta^2\over2}-(1-\cos\theta).
\]
Then
\begin{equation}
 \alpha_sR(\sqrt{2s}|Y|)
 =
 \sum_{q=2}^\infty
 {(-1)^q2^{q-1}s^{q-1}\over(2q)!}|Y|^{2q}.
 \label{eq:newV39-R-series}
\end{equation}

Let
\begin{equation}
 \Psi_s^R(K)
 :=
 {\cal T}\bigl(\alpha_sR(\sqrt{2s}|Y|)\bigr)(K)
 -{\cal T}\bigl(\alpha_sR(\sqrt{2s}|Y|)\bigr)(0).
 \label{eq:newV39-PsiR}
\end{equation}
The subtraction is the exact Gaussian centering.

\begin{lemma}[Exact Laguerre series for the action root]
\label{lem:newV39-action-Laguerre}
\begin{equation}
 \Psi_s^R(K)
 =
 \sum_{q=2}^\infty
 {(-1)^q2^{2q-1}s^{q-1}q!\over(2q)!}
 \left[L_q^{(1/2)}(K)-L_q^{(1/2)}(0)\right].
 \label{eq:newV39-action-Laguerre}
\end{equation}
The series converges locally uniformly in \(K\), with all
derivatives.
\end{lemma}

\begin{proof}
Insert \eqref{eq:newV39-R-series} into
\Cref{lem:newV39-Laguerre}.  Normal convergence follows either
from the Gaussian exponential moments on compact complex
\(\xi\)-sets or directly from the \((2q)!\) denominator.
\end{proof}

\begin{theorem}[Double-factorial spectral coefficients of the
cosine core]
\label{thm:newV39-cosine-EEMC}
There are numerical \(A,L,s_0>0\) such that, for
\(0<s\le s_0\),
\begin{equation}
 \Psi_s^R(K)
 =
 \sum_{k=1}^\infty c_k(s)s^{k-1}K^k,
 \qquad
 |c_k(s)|\le {A L^{2k}\over(2k)!}.
 \label{eq:newV39-cosine-EEMC}
\end{equation}
Thus the centered Gaussian spectral multiplier of the complete
calibrated cosine remainder belongs to the V38 even-entire class.
\end{theorem}

\begin{proof}
By
\eqref{eq:newV39-Laguerre-coefficients}, the coefficient of \(K^k\)
in \eqref{eq:newV39-action-Laguerre} is
\begin{equation}
 {1\over k!}
 \sum_{q\ge\max\{2,k\}}
 {(-1)^{q+k}2^{2q-1}s^{q-1}q!\over(2q)!}
 {q+1/2\choose q-k}.
 \label{eq:newV39-ck}
\end{equation}
Write \(q=k+n\).  For \(k\ge2\), divide the absolute value of the
\(n\)-th summand by \(s^{k-1}/(2k)!\).  Apart from the exterior
\(2^{2k-1}\), the resulting factor is
\[
 4^ns^n
 {(2k)!\over k!}
 {(k+n)!\over(2k+2n)!}
 {k+n+1/2\choose n}.
\]
Pair the \(2n\) new denominator factors and use
\[
 {k+r\over(2k+2r-1)(2k+2r)}
 \le {1\over2(k+r)}
 \qquad(1\le r\le n),
\]
together with
\[
 {k+n+1/2\choose n}
 \le{(k+n+1/2)^n\over n!}.
\]
The product is bounded by \(C^n/n!\), uniformly in \(k,n\).
Hence the \(n\)-sum is at most \(e^{Cs}\).  Enlarge \(L\) to
absorb \(2^{2k}\).  The case \(k=1\) starts at \(q=2\), has an
additional factor \(s\), and obeys the same bound after enlarging
\(A\).  This proves \eqref{eq:newV39-cosine-EEMC}.
\end{proof}

\begin{corollary}[The infinite cosine-root forest is closed in the
Gaussian tangent model]
\label{cor:newV39-cosine-card}
The exact nonpolynomial score whose Gaussian heat-normalized
multiplier is \(\Psi_s^R\) satisfies the V38 all-arity raw score
card with an exponential base independent of \(s,m\).
\end{corollary}

\begin{proof}
Apply \Cref{thm:newV39-cosine-EEMC} and
\Cref{cor:newV38-entire-card}.
\end{proof}

\section{The isolated Haar logarithm has the wrong coefficient order}
\label{sec:newV39-Haar}

Recall the V32 inner term
\[
 B(\theta)={\theta^2\over6}-\log{\theta\over\sin\theta}.
\]
For \(|\theta|<\pi\),
\begin{equation}
 B(\theta)
 =
 -\sum_{q=2}^\infty
 {\zeta(2q)\over q\pi^{2q}}\theta^{2q}.
 \label{eq:newV39-B-series}
\end{equation}

\begin{lemma}[Top spectral coefficient of the isolated Haar term]
\label{lem:newV39-Haar-top}
Under the Gaussian transform
\eqref{eq:newV39-transform}, the coefficient of \(K^q\) generated
by the \(q\)-th summand of
\(B(\sqrt{2s}|Y|)\) has absolute value
\begin{equation}
 {\zeta(2q)\over q}
 \left({4\over\pi^2}\right)^q s^q.
 \label{eq:newV39-Haar-top}
\end{equation}
Centering changes only the constant coefficient.
\end{lemma}

\begin{proof}
The \(q\)-th physical coefficient is
\[
 -{\zeta(2q)\over q\pi^{2q}}(2s)^q.
\]
By \Cref{lem:newV39-Laguerre}, the normalized Fourier transform
multiplies it by \(2^qq!L_q^{(1/2)}(K)\).  The top Laguerre
coefficient is \((-1)^q/q!\).  Taking the absolute value gives
\eqref{eq:newV39-Haar-top}.
\end{proof}

\begin{theorem}[The separated Haar series fails EEMC]
\label{thm:newV39-Haar-no-EEMC}
The centered Gaussian spectral transform of
\(B(\sqrt{2s}|Y|)\), treated separately from the wrapped images,
does not satisfy the coefficient bound
\eqref{eq:newV38-double-factorial} uniformly in \(q\), for any
fixed \(A,L\) and any fixed \(s>0\).
\end{theorem}

\begin{proof}
Writing the top coefficient
\eqref{eq:newV39-Haar-top} in the V38 normalization
\(a_q(s)s^{q-1}K^q\) gives
\[
 |a_q(s)|
 =
 {s\zeta(2q)\over q}
 \left({4\over\pi^2}\right)^q.
\]
If EEMC held, then
\[
 {s\zeta(2q)\over q}
 \left({4\over\pi^2}\right)^q
 \le {AL^{2q}\over(2q)!}
\]
for all \(q\).  The ratio of the left side to the right grows
faster than any exponential because \((2q)!/q\) does.  This is
impossible.
\end{proof}

\begin{corollary}[Termwise Haar root payment is not all-arity]
\label{cor:newV39-Haar-word-no-go}
Applying the positive ordered-word WPRAC census separately to
\eqref{eq:newV39-B-series} produces the degree majorant
\[
 Cs\sum_{q\ge2}{1\over q}
 \left({4(1+r_0)^2m^2\over\pi^2}\right)^q,
\]
which diverges once \(m\) is large.  Thus the extra exterior
\(s\) does not repair the all-arity coefficient order.
\end{corollary}

\begin{proof}
Use \eqref{eq:newV39-Haar-top} and the \(m^{2q}\) ordered-index
census in the proof of
\Cref{thm:newV38-polydisc}.  Since
\(\zeta(2q)\) is bounded above and below by positive numerical
constants, the displayed logarithmic series has exactly the stated
radius.
\end{proof}

\section{The Haar and image pieces are inseparable at all orders}
\label{sec:newV39-wrapped}

On the inner half, V32 writes
\[
 L_\alpha(\theta)-L_\alpha(0)
 =
 \alpha R(\theta)+B(\theta)+E_\alpha(\theta),
\]
where \(E_\alpha\) is exponentially small there.  The exact
function \(L_\alpha=\alpha\cos\theta-\log h_s(\theta)\) is smooth
at \(\theta=\pi\), whereas
\(B(\theta)\) has a logarithmic singularity at \(\pi\).

\begin{theorem}[Necessary antipodal cancellation]
\label{thm:newV39-antipodal-cancellation}
The paired-image contribution \(E_\alpha\), continued from the
inner representation to an antipodal neighborhood, cancels the
singular jet of \(B\) at \(\theta=\pi\).  Consequently:
\begin{enumerate}[label=\textup{(\roman*)},leftmargin=3.5em]
\item \(B\) and \(E_\alpha\) cannot both have globally uniform
      analytic derivative bounds when treated separately;
\item exponential smallness of \(E_\alpha\) on
      \(0\le\theta\le\pi/2\) does not permit it to be discarded in
      an arbitrary-degree root estimate; and
\item the EEMC must be tested on the complete wrapped block
      \(B+E_\alpha\), not on its two summands.
\end{enumerate}
\end{theorem}

\begin{proof}
The Wilson density and the heat kernel are strictly positive smooth
central functions on the compact group.  Their logarithmic ratio is
therefore smooth at the antipode.  The cosine term and
\(\alpha R\) are smooth there.  By
\eqref{eq:newV39-B-series}, \(B\) diverges logarithmically as
\(\theta\uparrow\pi\).  The only remaining term in the exact image
decomposition must cancel that singularity, including its complete
jet.  This proves \textup{(i)} and \textup{(iii)}.

For every fixed derivative order, the relevant image correction is
exponentially small on the inner half, as V31--V32 prove.  The order
needed by an all-arity root is unbounded; its derivatives detect the
shrinking transition from the shortest image to the antipodal
paired images.  Discarding that transition before the derivative
sum would leave the non-EEMC coefficients of
\Cref{thm:newV39-Haar-no-EEMC}, proving \textup{(ii)}.
\end{proof}

\begin{assessment}[Append-only correction to the V38 programme]
\label{ass:newV39-correction}
V38 correctly identifies the even-entire sufficient class and asks
for the exact coefficient test.  Its proposed separation of the
Haar and paired-image contributions is unavailable at arbitrary
degree.  The action core passes the sharp test in the Gaussian
tangent transform, but the separated Haar block fails it, and its
failure is canceled only by the complete wrapped heat kernel near
the antipode.  From V39 onward, \(B+E_\alpha\) is one indivisible
root object until after its complex or Fourier variation has been
bounded.
\end{assessment}

\section{The revised exact coefficient gate}
\label{sec:newV39-gate}

\begin{definition}[Wrapped even-entire multiplier card]
\label{def:newV39-WEEMC}
The WEEMC is the assertion that the exact centered
heat-normalized Fourier multiplier of
\[
 \alpha R(\theta)+B(\theta)+E_\alpha(\theta)
\]
admits the V38 representation
\[
 \sum_{q\ge1}a_q(s)s^{q-1}K^q,
 \qquad
 |a_q(s)|\le {AL^{2q}\over(2q)!},
\]
or an equivalent complete weighted-polydisc estimate, with
constants uniform in \(s\).  No separate norm is taken on
\(B\) or \(E_\alpha\).
\end{definition}

\begin{proposition}[WEEMC closes the exact heat-endpoint root]
\label{prop:newV39-WEEMC-sufficient}
If WEEMC holds, the exact centered heat-endpoint bridge score
satisfies WPRAC and the all-arity raw score card.
\end{proposition}

\begin{proof}
WEEMC places the complete exact multiplier in the class of
\Cref{cor:newV38-entire-card}.  Centering removes only the trivial
coefficient and is already built into that theorem.
\end{proof}

\section{V40 programme and claim boundary}
\label{sec:newV39-next}

\begin{programme}[V40: audit and full wrapped-multiplier test]
\label{prog:newV39-V40}
\begin{enumerate}[leftmargin=3em]
\item Perform the compulsory read-only audit of the newest active
      SM--Einstein and Navier--Stokes notes.
\item Start from the exact paired-image numerator, not the local
      series for \(B\), and derive a global contour representation
      of \(-\log h_s\) as a central multiplier.
\item Locate the nearest complex zeros of the complete paired heat
      kernel and determine the resulting spectral coefficient order.
\item Prove WEEMC, or exhibit a coefficient subsequence which
      violates \((2q)!\) decay after the antipodal cancellation has
      already been included.
\item If WEEMC fails, absorb the complete wrapped factor into a new
      exactly solvable endpoint; do not return to separated Haar
      jets.
\item Preserve the exact binary and ternary bridge cards of
      V36--V37 independently of the all-order outcome.
\end{enumerate}
\end{programme}

\begin{firewall}[Claim boundary after V39]
\label{fire:newV39-boundary}
V39 proves the exact radial Gaussian--Laguerre transform, proves
that the complete calibrated cosine remainder has
double-factorial spectral coefficients, and closes its infinite
Gaussian tangent forest.  It proves that the separated Haar
logarithm fails EEMC and that the paired-image term is indispensable
for canceling its antipodal jet.  It does not prove WEEMC for the
complete wrapped score or all-order WPRAC for the exact bridge.
Thus exact all-order RGSC, the Wilson aggregate strong card,
all-order strong MTA, WUFC, CPM, cutoff removal, reflection
positivity, Osterwalder--Schrader reconstruction, a continuum
Yang--Mills measure, and a positive Hamiltonian gap remain open.
The full Yang--Mills mass gap has not been proved.
\end{firewall}

\begin{theorem}[V39 result ledger]
\label{thm:newV39-results}
V39:
\begin{enumerate}[label=\textup{(V39.\arabic*)},leftmargin=6.5em]
\item derives the exact Gaussian--Laguerre transform of every radial
      monomial;
\item proves the centered cosine-action multiplier has the V38
      double-factorial coefficient bound;
\item closes the resulting nonpolynomial Gaussian tangent root;
\item computes the top spectral coefficient of every isolated Haar
      logarithm term;
\item proves the separated Haar series fails EEMC and its
      ordered-word all-arity majorant diverges;
\item proves that the paired images cancel the complete antipodal
      Haar jet; and
\item replaces the separated test by the wrapped WEEMC.
\end{enumerate}
\end{theorem}

\begin{proof}
Use
\Cref{lem:newV39-Laguerre,
lem:newV39-action-Laguerre,
thm:newV39-cosine-EEMC,
cor:newV39-cosine-card,
lem:newV39-Haar-top,
thm:newV39-Haar-no-EEMC,
cor:newV39-Haar-word-no-go,
thm:newV39-antipodal-cancellation,
prop:newV39-WEEMC-sufficient}.
\end{proof}

\paragraph{Parent-version record.}
V39 was formed from
\texttt{Renewal\_Yang\_Mills\_Mass\_Gap\_Research\_Notes\_V38.tex}.
The V38 parent had \(18\,169\) logical source lines,
\(574\,336\) bytes, and SHA--256
\[
 \texttt{7ec255a3091aecf7dfc8e89363040efd7c22ca5d85aa5f5212b00c23cf7008ba}.
\]
The next compulsory companion audit is V40.

% =============================================================================
% END APPENDED FIFTH-SERIES V39 MATERIAL
% =============================================================================

% =============================================================================
% BEGIN APPENDED FIFTH-SERIES V40 MATERIAL
% =============================================================================

\chapter{V40: Complex heat zeros and the limit of local image
separation}
\label{chap:newV40-complex-zeros}

\section{Compulsory companion audit}
\label{sec:newV40-audit}

\begin{audit}[Read-only SM--Einstein and Navier--Stokes audit]
\label{audit:newV40-companions}
Before beginning the V40 calculation, the newest active companion
notes were inspected without modification:
\begin{align*}
 &\texttt{Renewal\_SM\_Einstein\_Continuum\_Research\_Notes\_V54.tex},
 \\
 &\qquad 17\,584\ \hbox{logical source lines},\quad
 675\,306\ \hbox{bytes},                                      \\
 &\qquad
 \operatorname{SHA256}
 =
 \texttt{078f1693a4f5f2451caf69ea9e14595699f0209f535089cdb3bdc5e2dea42c7f},
 \\
 &\texttt{Renewal\_NS\_Stability\_Research\_Notes\_V35.tex},
 \\
 &\qquad 25\,472\ \hbox{logical source lines},\quad
 998\,459\ \hbox{bytes},                                      \\
 &\qquad
 \operatorname{SHA256}
 =
 \texttt{aea7449895cd39b2d866869048033b96a47140d037dd9e3738ebb26f3fa55c01}.
\end{align*}
The SM V54 calculation shows that the proper-time shell
polarization pays an order-one amount on each active scale and
therefore retains a logarithmic divergence until the local Maxwell
piece is subtracted or renormalized.  The NS V35 calculation shows
the analogous structural failure of an absolute slot-by-slot or
octave-by-octave estimate: one constant per scale accumulates a
logarithm, so the signed flux must remain assembled.
\end{audit}

\begin{assessment}[Transfer discipline from the audit]
\label{ass:newV40-audit-transfer}
The companion calculations supply no Yang--Mills heat-kernel
inequality.  Their legitimate contribution is methodological:
keep the complete wrapped root object assembled until after its
cancellation has been estimated.  In particular, V40 does not
take separate absolute norms of the Haar term and the image
correction identified in V39.
\end{assessment}

\section{Zeros of the nearest antipodal pair}
\label{sec:newV40-pair-zeros}

Retain the exact V31 notation
\[
 N_\tau(\theta)
 =
 \sum_{k\in\mathbb Z}
 (\theta+2\pi k)e^{-(\theta+2\pi k)^2/(4\tau)},
 \qquad
 h_\tau(\theta)=c_\tau {N_\tau(\theta)\over\sin\theta}.
\]
Put \(\delta=\pi-\theta\).  The two nearest images are
\begin{align}
 P_\tau(\delta)
 &:=
 N_\tau^{\rm pair}(\pi-\delta)                              \notag\\
 &=
 2e^{-(\pi^2+\delta^2)/(4\tau)}F_\tau(\delta),              \label{eq:newV40-pair}\\
 F_\tau(\delta)
 &:=
 \pi\sinh\left({\pi\delta\over2\tau}\right)
 -\delta\cosh\left({\pi\delta\over2\tau}\right).
 \label{eq:newV40-F}
\end{align}
These are entire functions of the complex variable \(\delta\).

\begin{lemma}[First nontrivial imaginary zero of the pair]
\label{lem:newV40-pair-zero}
For all sufficiently small \(\tau>0\), \(F_\tau\) has a simple
zero \(i y_\tau^{(0)}\), where
\begin{equation}
 y_\tau^{(0)}
 =
 2\tau+{4\over\pi^2}\tau^2+O(\tau^3).
 \label{eq:newV40-y0}
\end{equation}
Its conjugate \(-i y_\tau^{(0)}\) is also a simple zero.
\end{lemma}

\begin{proof}
For \(y>0\), set
\[
 x={\pi y\over2\tau},
 \qquad
 \epsilon={2\tau\over\pi^2}.
\]
Since
\[
 F_\tau(iy)
 =
 i\left[
 \pi\sin x-y\cos x
 \right],
\]
a zero with \(\cos x\ne0\) is equivalent to
\begin{equation}
 \tan x=\epsilon x.
 \label{eq:newV40-tan-equation}
\end{equation}
On \((\pi,3\pi/2)\), the function
\(\tan x-\epsilon x\) starts negative at \(\pi\), tends to
\(+\infty\), and has derivative
\(\sec^2x-\epsilon>0\) when \(\epsilon<1\).  Hence it has exactly
one zero \(x_\tau\) in that interval, and that zero is simple.

Write \(x_\tau=\pi+\eta_\tau\).  Equation
\eqref{eq:newV40-tan-equation} and
\(\tan\eta=\eta+O(\eta^3)\) give
\[
 \eta_\tau
 =
 \epsilon(\pi+\eta_\tau)+O(\eta_\tau^3)
 =
 \epsilon\pi+\epsilon^2\pi+O(\epsilon^3).
\]
Therefore
\[
 y_\tau^{(0)}
 =
 {2\tau\over\pi}x_\tau
 =
 2\tau+{4\over\pi^2}\tau^2+O(\tau^3).
\]
Real Taylor coefficients give the conjugate zero.
\end{proof}

\begin{remark}[The canceled zero at the real antipode]
\label{rem:newV40-real-zero}
The identity \(F_\tau(0)=0\) is not a zero of the heat density:
it is canceled by \(\sin\delta\) in
\(h_\tau(\pi-\delta)\).  In contrast,
\(\sin(i y_\tau^{(0)})\ne0\), so the nonzero imaginary roots found
above are genuine candidate zeros of the complexified density.
\end{remark}

\section{Persistence under all remote images}
\label{sec:newV40-full-zero}

Define the remote remainder
\[
 Q_\tau(\delta)
 :=
 N_\tau(\pi-\delta)-P_\tau(\delta).
\]

\begin{lemma}[Complex cap separation]
\label{lem:newV40-complex-separation}
There are constants \(c,C,\tau_0>0\) such that, for
\(0<\tau\le\tau_0\) and \(|\delta|\le4\tau\),
\begin{align}
 |Q_\tau(\delta)|
 +\tau|\partial_\delta Q_\tau(\delta)|
 &\le
 C\tau^{-C}e^{-9\pi^2/(4\tau)+C\tau},                       \label{eq:newV40-Q-bound}\\
 |P_\tau' (i y_\tau^{(0)})|
 &\ge
 c\tau^{-1}e^{-\pi^2/(4\tau)-C\tau}.                        \label{eq:newV40-Pprime}
\end{align}
\end{lemma}

\begin{proof}
The remote terms occur in pairs with real parts of their
distances at least \(3\pi-O(\tau)\).  For a complex number
\(z=a+ib\),
\[
 \left|e^{-z^2/(4\tau)}\right|
 =
 e^{-(a^2-b^2)/(4\tau)}.
\]
On \(|\delta|\le4\tau\), this bounds every first remote pair by a
polynomial in \(\tau^{-1}\) times
\(e^{-9\pi^2/(4\tau)+C\tau}\); the remaining pairs form a
uniformly summable Gaussian tail.  One derivative adds at most a
factor \(C\tau^{-1}\), which proves
\eqref{eq:newV40-Q-bound}.

At the root of \(F_\tau\), differentiation of
\eqref{eq:newV40-F} gives
\[
 F_\tau'(\delta)
 =
 \left({\pi^2\over2\tau}-1\right)
 \cosh\left({\pi\delta\over2\tau}\right)
 -
 {\pi\delta\over2\tau}
 \sinh\left({\pi\delta\over2\tau}\right).
\]
At \(\delta=i y_\tau^{(0)}\), the first factor is
of order \(\tau^{-1}\), while
\(\cos x_\tau=-1+O(\tau^2)\); the remaining term is bounded.
The exponential prefactor in \eqref{eq:newV40-pair} has modulus
\(e^{-\pi^2/(4\tau)+O(\tau)}\).
This proves \eqref{eq:newV40-Pprime}.
\end{proof}

\begin{theorem}[A full wrapped heat-kernel zero at distance
\(2\tau+O(\tau^2)\)]
\label{thm:newV40-full-zero}
There are complex conjugate zeros
\(\delta_\tau,\overline{\delta_\tau}\) of
\(N_\tau(\pi-\delta)\) satisfying
\begin{equation}
 \delta_\tau
 =
 i\left(2\tau+{4\over\pi^2}\tau^2+O(\tau^3)\right)
 +O(e^{-c/\tau}).
 \label{eq:newV40-full-zero}
\end{equation}
They are not zeros of \(\sin\delta\).  Consequently
\(h_\tau(\pi-\delta)\) has genuine complex zeros there.
\end{theorem}

\begin{proof}
By \Cref{lem:newV40-pair-zero}, the pair zero is simple.
Choose a circle centered at \(i y_\tau^{(0)}\) whose radius is
\(r_\tau=e^{-c_1/\tau}\), with \(c_1>0\) smaller than the
remote-image exponent gap.  Taylor's theorem,
\eqref{eq:newV40-Pprime}, and a possibly smaller \(c_1\) give
\[
 |P_\tau(\delta)|
 \ge
 c r_\tau\tau^{-1}e^{-\pi^2/(4\tau)-C\tau}
\]
on that circle.  The quadratic Taylor remainder is smaller than
half this quantity because \(r_\tau/\tau\to0\).
Equation \eqref{eq:newV40-Q-bound} makes
\(|Q_\tau|<|P_\tau|\) on the same circle.  Rouch\'e's theorem
therefore gives exactly one zero of \(P_\tau+Q_\tau\) inside.
Its displacement from \(i y_\tau^{(0)}\) is
\(O(e^{-c/\tau})\).  Conjugation supplies the second zero.

For small \(\tau\),
\(\delta_\tau\ne0\) and is not a real multiple of \(\pi\).
Thus \(\sin\delta_\tau\ne0\), and the quotient defining
\(h_\tau\) retains the zero.
\end{proof}

\section{The collapsing analytic radius}
\label{sec:newV40-radius}

\begin{theorem}[No \(\tau\)-uniform physical analytic strip]
\label{thm:newV40-no-strip}
Let
\[
 \ell_\tau(\delta)
 :=
 \log h_\tau(\pi-\delta)
\]
denote the real branch near \(\delta=0\), analytically continued
on its maximal zero-free disk.  If \(R_\tau\) is the radius of
convergence of its Taylor series at \(0\), then
\begin{equation}
 R_\tau
 \le
 2\tau+{4\over\pi^2}\tau^2+O(\tau^3).
 \label{eq:newV40-radius}
\end{equation}
Equivalently,
\begin{equation}
 \limsup_{n\to\infty}
 \left(
 {|\ell_\tau^{(n)}(0)|\over n!}
 \right)^{1/n}
 \ge
 {1\over 2\tau+O(\tau^2)}.
 \label{eq:newV40-Cauchy-Hadamard}
\end{equation}
\end{theorem}

\begin{proof}
The real heat kernel is positive and analytic at the antipode
after the common real zero of \(N_\tau\) and \(\sin\theta\) is
removed.  Hence the stated branch exists near \(0\).  A
holomorphic logarithm cannot continue through a zero of its
argument.  The zero in \Cref{thm:newV40-full-zero} lies at modulus
\[
 |\delta_\tau|
 =
 2\tau+{4\over\pi^2}\tau^2+O(\tau^3),
\]
so the Taylor disk cannot have larger radius.  The
Cauchy--Hadamard formula gives
\eqref{eq:newV40-Cauchy-Hadamard}.
\end{proof}

\begin{corollary}[Failure of the proposed physical Cauchy route]
\label{cor:newV40-cauchy-no-go}
No proof of WEEMC can be based on a zero-free complex
\(\delta\)-strip of width bounded below independently of \(\tau\)
for the complete wrapped logarithm.
\end{corollary}

\begin{proof}
\Cref{thm:newV40-no-strip} supplies full heat-kernel zeros at
distance \(2\tau+O(\tau^2)\) from the real antipode.
\end{proof}

\begin{warning}[What the zero theorem does not prove]
\label{warn:newV40-not-spectral-no-go}
The collapsing physical analytic radius does not disprove WEEMC.
WEEMC concerns the heat-normalized \emph{spectral} coefficients
after the complete wrapped integration.  Oscillation, heat
weighting, centering, and Laguerre or character cancellation may
improve those coefficients beyond the pointwise Taylor radius.
V40 rules out one proof mechanism, not the desired spectral
estimate.
\end{warning}

\section{How many derivatives may separate the images?}
\label{sec:newV40-derivative-window}

The fixed-order remote-image estimates of V31--V32 remain correct.
The next proposition quantifies why their order may not be sent to
infinity before the images are recombined.

\begin{proposition}[Finite derivative-order separation window]
\label{prop:newV40-window}
Let \(\rho_\tau\) be the relative remote-image correction in the
V31 cap pairing, with the common zero at \(\delta=0\) removed.
There are \(c,C>0\) such that its real Taylor coefficients on the
cap obey
\begin{equation}
 {|\partial_\delta^n\rho_\tau(\delta)|\over n!}
 \le
 C e^{-c/\tau}\left({C\over\tau}\right)^n
 \label{eq:newV40-window-bound}
\end{equation}
on every closed real subcap on which the corresponding
\(c_0\tau\)-complex disks avoid a pair zero.  In particular, the
right side remains at most \(Ce^{-c/(2\tau)}\) only while
\begin{equation}
 n\log(C/\tau)\le {c\over2\tau}.
 \label{eq:newV40-window-range}
\end{equation}
Thus fixed-order image separation is uniformly exponentially
accurate, but its directly justified range is at most
\[
 n=O\left({1\over\tau\log(1/\tau)}\right),
\]
not all derivative orders.
\end{proposition}

\begin{proof}
The image-distance gap used in
\Cref{lem:newV40-complex-separation} persists on a complex disk of
radius \(c_0\tau\) around each permitted real point.  After the
common zero at the cap tip is factored, the nearest pair has no
zero on the chosen disk.  The quotient of the remote images by
that pair is therefore holomorphic and bounded by
\(Ce^{-c/\tau}\).  Cauchy's derivative estimate on a disk of
radius \(c_0\tau\) gives \eqref{eq:newV40-window-bound}.
Substitution of \eqref{eq:newV40-window-range} yields the stated
exponential remainder.  The pair zeros of
\Cref{lem:newV40-pair-zero} show why disks on the \(\tau\) scale,
rather than a fixed scale, are the natural general bound.
\end{proof}

\begin{assessment}[Exact status of the V39 cancellation]
\label{ass:newV40-cancellation-status}
The V39 conclusion is now quantitative.  The separated Haar and
image descriptions may be used for each fixed jet and within the
window \eqref{eq:newV40-window-range}.  They cannot justify the
unbounded-degree summation required by WPRAC.  Beyond that window,
the complete numerator \(N_\tau\), the denominator
\(\sin\theta\), and the heat-normalized transform must remain one
assembled object.
\end{assessment}

\section{The upstream spectral reformulation}
\label{sec:newV40-upstream}

The physical contour route has reached a genuine obstruction, so
the programme now moves upstream to a family whose parameter
derivative produces the logarithm only after the heat convolution
has been assembled.

\begin{definition}[Heat-power spectral family]
\label{def:newV40-power-family}
For \(p\) in a real neighborhood of \(1\) and an irreducible
\(SU(2)\) representation \(\lambda\), define
\begin{equation}
 \mathscr H_{\tau,p}(\lambda)
 :=
 {1\over d_\lambda}
 \int_{SU(2)}
 h_\tau(g)^p\,
 \overline{\chi_\lambda(g)}\,dg.
 \label{eq:newV40-power-family}
\end{equation}
Since \(h_\tau>0\) on the real group,
\begin{equation}
 \left.\partial_p\mathscr H_{\tau,p}(\lambda)\right|_{p=1}
 =
 {1\over d_\lambda}
 \int_{SU(2)}
 h_\tau(g)\log h_\tau(g)\,
 \overline{\chi_\lambda(g)}\,dg.
 \label{eq:newV40-power-derivative}
\end{equation}
The right side is the assembled spectral logarithm; neither the
Haar logarithm nor an individual image occurs separately.
\end{definition}

\begin{lemma}[Legitimacy of the \(p\)-derivative]
\label{lem:newV40-power-derivative}
For each fixed \(\tau>0\), differentiation under the integral in
\eqref{eq:newV40-power-family} is valid for \(p\) in a compact
subinterval of \((0,\infty)\).
\end{lemma}

\begin{proof}
On the compact group, \(h_\tau\) is smooth and strictly positive.
It therefore has finite positive minimum and maximum.  Both
\(h_\tau^p\) and \(h_\tau^p\log h_\tau\) are uniformly bounded
when \(p\) ranges in the stated compact interval.  Dominated
convergence proves the claim.
\end{proof}

\begin{definition}[Power-family spectral card]
\label{def:newV40-PFSC}
The PFSC is a uniform neighborhood \(I\ni1\) and a
heat-normalized representation of
\(\mathscr H_{\tau,p}(\lambda)\) whose \(p\)-derivative at \(1\)
has either:
\begin{enumerate}[label=\textup{(\roman*)},leftmargin=3.5em]
\item the V38 double-factorial coefficient bound; or
\item directly the complete weighted-polydisc bound WPRAC,
\end{enumerate}
with constants uniform in \(\tau\).
\end{definition}

\begin{proposition}[PFSC is an upstream sufficient gate]
\label{prop:newV40-PFSC-sufficient}
PFSC implies WEEMC for the heat-weighted logarithmic row and hence
closes that row of the exact heat-endpoint bridge-score root.
\end{proposition}

\begin{proof}
Identity \eqref{eq:newV40-power-derivative} is exactly the
heat-weighted logarithmic multiplier before centering.  The
trivial representation supplies the scalar centering term.
Alternative \textup{(i)} enters the V38 even-entire theorem;
alternative \textup{(ii)} is already the required conclusion.
In both cases the V35 exact cumulant construction then supplies
the corresponding bridge-score row.
\end{proof}

\begin{programme}[V41: heat-power interpolation]
\label{prog:newV40-V41}
\begin{enumerate}[leftmargin=3em]
\item Use the heat semigroup and Peter--Weyl expansion to normalize
      \(\mathscr H_{\tau,p}(\lambda)\) uniformly near \(p=1\).
\item Derive the exact \(p=1\) value and first \(p\)-variation
      before taking any absolute coefficient norm.
\item Test whether log-convexity, hypercontractivity, or an exact
      radial integral yields PFSC.  Any invoked theorem must be
      proved in the present normalization or stated as a
      conditional input.
\item If PFSC fails, identify an explicit spectral subsequence
      after the full wrapped integration; pointwise physical
      Taylor coefficients are not a counterexample.
\item Preserve the exact binary and ternary cards and the infinite
      tangent cosine card already proved.
\end{enumerate}
\end{programme}

\begin{firewall}[Claim boundary after V40]
\label{fire:newV40-boundary}
V40 proves the existence and asymptotic location of a conjugate
pair of complex zeros of the complete \(SU(2)\) heat kernel near
the antipode.  It proves that the physical logarithm has analytic
radius at most \(2\tau+O(\tau^2)\), quantifies the finite
derivative window in which remote images may be separated, and
replaces the blocked physical-Cauchy route by the assembled
heat-power spectral family.  It does not prove PFSC, WEEMC,
all-order WPRAC for the exact bridge, or all-order RGSC.  The
Wilson aggregate strong card, all-order strong MTA, WUFC, CPM,
cutoff removal, reflection positivity,
Osterwalder--Schrader reconstruction, a continuum Yang--Mills
measure, and a positive Hamiltonian gap remain open.  The full
Yang--Mills mass gap has not been proved.
\end{firewall}

\begin{theorem}[V40 result ledger]
\label{thm:newV40-results}
V40:
\begin{enumerate}[label=\textup{(V40.\arabic*)},leftmargin=6.5em]
\item completes the compulsory V40 companion audit without
      modifying either companion file;
\item locates the first nontrivial imaginary zero of the exact
      antipodal image pair;
\item proves that the zero persists after all remote images are
      restored;
\item proves collapse of the physical analytic radius to at most
      \(2\tau+O(\tau^2)\);
\item rules out a \(\tau\)-uniform physical zero-free-strip proof
      of WEEMC, without claiming spectral failure;
\item quantifies the derivative-order limit of separated image
      estimates; and
\item introduces the exact heat-power spectral family and PFSC as
      the next upstream sufficient gate.
\end{enumerate}
\end{theorem}

\begin{proof}
Use
\Cref{audit:newV40-companions,
lem:newV40-pair-zero,
lem:newV40-complex-separation,
thm:newV40-full-zero,
thm:newV40-no-strip,
cor:newV40-cauchy-no-go,
prop:newV40-window,
lem:newV40-power-derivative,
prop:newV40-PFSC-sufficient}.
\end{proof}

\paragraph{Parent-version record.}
V40 was formed from
\texttt{Renewal\_Yang\_Mills\_Mass\_Gap\_Research\_Notes\_V39.tex}.
The V39 parent had \(18\,618\) logical source lines,
\(588\,725\) bytes, and SHA--256
\[
 \texttt{026728ddc19aa468c78af58da8568899a375575eea2ddd35e845f0e6393747d6}.
\]
The next compulsory companion audit is V45.

% =============================================================================
% END APPENDED FIFTH-SERIES V40 MATERIAL
% =============================================================================

% =============================================================================
% BEGIN APPENDED FIFTH-SERIES V41 MATERIAL
% =============================================================================

\chapter{V41: Spectral no-go for the logarithmic bridge and an
additive-envelope replacement}
\label{chap:newV41-spectral-no-go}

\section{Central Fourier normalization}
\label{sec:newV41-normalization}

No companion audit is due at V41.  The V40 read-only audit remains
the current companion checkpoint.

Index the irreducible \(SU(2)\) characters by \(n\ge0\), so that
\[
 d_n=n+1,
 \qquad
 c_n=n(n+2),
 \qquad
 h_\tau(\theta)
 =
 \sum_{n\ge0}d_ne^{-\tau c_n}\chi_n(\theta).
\]
For a central integrable function \(f\), put
\begin{equation}
 \widehat f(n)
 :=
 {1\over d_n}\int_{SU(2)}
 f(g)\overline{\chi_n(g)}\,dg.
 \label{eq:newV41-hat}
\end{equation}
Thus the character coefficient is \(d_n\widehat f(n)\), and
\begin{equation}
 \widehat h_\tau(n)=e^{-\tau c_n}.
 \label{eq:newV41-heat-hat}
\end{equation}

\begin{lemma}[Gaussian Peter--Weyl decay forces an entire class
function]
\label{lem:newV41-Gaussian-entire}
Suppose \(f\) is a continuous central function and, for some
\(a,C>0\),
\begin{equation}
 |\widehat f(n)|\le Ce^{-a n^2}
 \qquad(n\ge0).
 \label{eq:newV41-Gaussian-decay}
\end{equation}
Then its real class-angle representative is the restriction of an
entire even \(2\pi\)-periodic function of
\(\theta\in\mathbb C\).
\end{lemma}

\begin{proof}
The character identity
\[
 \chi_n(\theta)={\sin((n+1)\theta)\over\sin\theta}
\]
has removable values at the zeros of \(\sin\theta\) and is an
entire polynomial in \(2\cos\theta\).  On every compact
\(\theta\)-set,
\[
 |\chi_n(\theta)|
 \le
 C_K(n+1)e^{(n+1)|\Im\theta|}.
\]
Consequently
\[
 \sum_{n\ge0}d_n\widehat f(n)\chi_n(\theta)
\]
converges normally on every compact subset of \(\mathbb C\) under
\eqref{eq:newV41-Gaussian-decay}.  Its sum is entire, even, and
\(2\pi\)-periodic.  On the real group it has the Peter--Weyl
coefficients of \(f\); \(L^2\) uniqueness and continuity identify
it with \(f\).
\end{proof}

\section{A simple zero makes the heat logarithmic row non-entire}
\label{sec:newV41-log-row}

Let \(w_\alpha=Z_\alpha^{-1}e^{\alpha\cos\theta}\) be the normalized
Wilson one-link density, and let
\[
 r_{\alpha,t}^{\rm geo}
 =
 {h_\tau^{\,1-t}w_\alpha^{\,t}\over
  \int h_\tau^{\,1-t}w_\alpha^{\,t}\,dU}
\]
be the geometric bridge used in V21--V40.  At its heat endpoint,
the centered score has the form
\[
 S_{\alpha,0}^{\rm geo}(\theta)
 =
 \alpha\cos\theta-\log h_\tau(\theta)-\gamma_{\alpha,\tau}
\]
for a real scalar \(\gamma_{\alpha,\tau}\).  Define its signed
score density
\begin{equation}
 g_{\alpha,\tau}(\theta)
 :=
 h_\tau(\theta)S_{\alpha,0}^{\rm geo}(\theta).
 \label{eq:newV41-g}
\end{equation}

\begin{lemma}[Local logarithmic monodromy]
\label{lem:newV41-monodromy}
If an entire function \(H\) has a simple zero at \(z_0\), then the
real-branch function \(H\log H\), analytically continued from any
zero-free real neighborhood, has no single-valued holomorphic
extension through \(z_0\).
\end{lemma}

\begin{proof}
Locally
\[
 H(z)=(z-z_0)a(z),
 \qquad a(z_0)\ne0.
\]
Continuation once around \(z_0\) changes \(\log H\) by \(2\pi i\),
so it changes \(H\log H\) by \(2\pi iH\), which is not identically
zero on the punctured neighborhood.  A holomorphic extension
would be single-valued, a contradiction.
\end{proof}

\begin{theorem}[The complete geometric score density is not
entire]
\label{thm:newV41-score-not-entire}
For every sufficiently small fixed \(\tau>0\),
\(g_{\alpha,\tau}\) has no entire class-angle extension.
\end{theorem}

\begin{proof}
The Gaussian character series for \(h_\tau\) converges normally on
compact complex sets, so \(h_\tau\) is entire.  By
\Cref{thm:newV40-full-zero} it has a simple nonreal zero
\(\theta_\tau=\pi-\delta_\tau\).  The multiplicity is one because
the Rouch\'e disk in that proof contains one zero counted with
multiplicity.

The terms
\[
 \alpha h_\tau\cos\theta
 \quad\hbox{and}\quad
 -\gamma_{\alpha,\tau}h_\tau
\]
in \eqref{eq:newV41-g} are entire.  The remaining term is
\(-h_\tau\log h_\tau\), which has nontrivial monodromy at
\(\theta_\tau\) by \Cref{lem:newV41-monodromy}.  Entire terms
cannot cancel that monodromy.
\end{proof}

\section{The exact all-spin heat normalization fails}
\label{sec:newV41-heat-no-go}

\begin{theorem}[Heat-normalized singleton score no-go]
\label{thm:newV41-singleton-no-go}
For every sufficiently small fixed \(\tau>0\), every \(L<\infty\),
and every \(C<\infty\), there is an \(n\) such that
\begin{equation}
 e^{\tau c_n}
 |\widehat g_{\alpha,\tau}(n)|
 >
 Ce^{Ln}.
 \label{eq:newV41-singleton-growth}
\end{equation}
In particular, the heat-normalized geometric marked singleton is
not uniformly bounded over the Peter--Weyl sectors and does not
have the V38 even-entire growth.
\end{theorem}

\begin{proof}
If \eqref{eq:newV41-singleton-growth} failed for some \(C,L\), then
\[
 |\widehat g_{\alpha,\tau}(n)|
 \le
 C\exp[-\tau n(n+2)+Ln].
\]
For all sufficiently large \(n\), the right side is at most
\(C'e^{-(\tau/2)n^2}\).  The finitely many remaining coefficients
can be absorbed into \(C'\).  By
\Cref{lem:newV41-Gaussian-entire}, \(g_{\alpha,\tau}\) would have
an entire extension, contradicting
\Cref{thm:newV41-score-not-entire}.

The V38 coefficient class grows after heat normalization at most
like
\[
 \sum_{q\ge0}{A(Ln)^{2q}\over(2q)!}
 \le Ae^{Ln}
\]
when \(K=\tau c_n\) and \(\tau\) is fixed.  It is therefore also
excluded.
\end{proof}

\begin{corollary}[PFSC and geometric-bridge WEEMC are false]
\label{cor:newV41-PFSC-false}
The PFSC proposed in V40 is false for the exact logarithmic row.
The V39 WEEMC and the V35 WPRAC are false when interpreted as
all-representation, heat-normalized cards for the exact geometric
bridge.
\end{corollary}

\begin{proof}
At \(p=1\),
\[
 \left.\partial_p\mathscr H_{\tau,p}(n)\right|_{p=1}
 =
 {1\over d_n}
 \int h_\tau\log h_\tau\,\overline{\chi_n}\,dU.
\]
Adding the Wilson cosine and centering terms produces
\(\widehat g_{\alpha,\tau}(n)\), while those added terms are
entire.  PFSC or WEEMC would give the V38
\(e^{L\sqrt{\tau c_n}}\) bound after heat normalization, contrary
to \Cref{thm:newV41-singleton-no-go}.  WPRAC includes the
one-block marked row and is therefore contradicted already at
arity one.
\end{proof}

\begin{assessment}[Relation to the V26 counterexample]
\label{ass:newV41-v26-relation}
This is not a repetition of the V26 fixed-spin, arity-four payer
counterexample.  V26 proves a missing Wilson clock in the first
saturating tangent tensor.  V41 proves an independent
\emph{all-spin, arity-one} obstruction: inverse heat smoothing
cannot be applied to a logarithmic score density whose complex
continuation branches at a heat-kernel zero.
\end{assessment}

\section{The endpoint mismatch is already visible without a score}
\label{sec:newV41-Wilson-tail}

\begin{lemma}[Exact Wilson multiplier]
\label{lem:newV41-Wilson-multiplier}
For the normalized Wilson density
\[
 w_\alpha(\theta)=Z_\alpha^{-1}e^{\alpha\cos\theta},
\]
its scalar multiplier in sector \(n\) is
\begin{equation}
 q_{\alpha,n}
 :=
 \widehat w_\alpha(n)
 =
 {I_{n+1}(\alpha)\over I_1(\alpha)}
 >0,
 \label{eq:newV41-Wilson-multiplier}
\end{equation}
where \(I_k\) is the modified Bessel function.
\end{lemma}

\begin{proof}
Normalized radial Haar measure is
\((2/\pi)\sin^2\theta\,d\theta\).  Hence
\[
 Z_\alpha
 =
 {2\over\pi}\int_0^\pi
 e^{\alpha\cos\theta}\sin^2\theta\,d\theta
 =
 I_0(\alpha)-I_2(\alpha)
 =
 {2I_1(\alpha)\over\alpha}.
\]
Furthermore,
\begin{align*}
 \int w_\alpha\chi_n\,dU
 &=
 {1\over Z_\alpha}{2\over\pi}
 \int_0^\pi
 e^{\alpha\cos\theta}
 \sin\theta\sin((n+1)\theta)\,d\theta\\
 &=
 {I_n(\alpha)-I_{n+2}(\alpha)\over Z_\alpha}
 =
 {(n+1)I_{n+1}(\alpha)\over I_1(\alpha)}.
\end{align*}
Divide by \(d_n=n+1\).
\end{proof}

\begin{proposition}[The Wilson endpoint is not inverse-heat
bounded]
\label{prop:newV41-Wilson-not-heat}
For every fixed \(\alpha,\tau>0\) and \(L<\infty\),
\begin{equation}
 \sup_{n\ge0}
 e^{\tau c_n-Ln}q_{\alpha,n}
 =
 \infty.
 \label{eq:newV41-Wilson-not-heat}
\end{equation}
\end{proposition}

\begin{proof}
The positive power series for \(I_{n+1}\) gives
\[
 I_{n+1}(\alpha)
 \ge
 {(\alpha/2)^{n+1}\over(n+1)!}.
\]
Therefore the logarithm of the left side along \(n\to\infty\) is
bounded below by
\[
 \tau n(n+2)-Ln
 +(n+1)\log(\alpha/2)
 -\log((n+1)!)
 -\log I_1(\alpha).
\]
Stirling's formula makes this
\(\tau n^2-n\log n+O_{\alpha,\tau,L}(n)\), which tends to
\(+\infty\).
\end{proof}

\begin{firewall}[Retired normalization, preserved finite jets]
\label{fire:newV41-retirement}
The exact all-spin programme may no longer normalize the Wilson
endpoint or the geometric score by the inverse heat multiplier
\(e^{\tau c_n}\).  The finite-jet theorems V33--V38 remain correct,
as do the exact cumulant and forest identities conditional on
their stated hypotheses.  What is retired is the assertion that
their heat-normalized infinite-degree hypotheses can hold for the
exact geometric bridge.
\end{firewall}

\section{An additive bridge removes the logarithmic score}
\label{sec:newV41-additive}

Both endpoint densities are normalized.  Define
\begin{equation}
 r_{\alpha,t}^{\rm add}
 :=
 (1-t)h_\tau+t w_\alpha,
 \qquad 0\le t\le1.
 \label{eq:newV41-additive}
\end{equation}
It is a strictly positive normalized central density.

\begin{lemma}[Exact additive score identity]
\label{lem:newV41-additive-score}
The centered additive score is
\begin{equation}
 \Sigma_{\alpha,t}^{\rm add}
 =
 {w_\alpha-h_\tau\over r_{\alpha,t}^{\rm add}},
 \qquad
 \mathbb E_{r_{\alpha,t}^{\rm add}}\Sigma_{\alpha,t}^{\rm add}=0.
 \label{eq:newV41-additive-score}
\end{equation}
For every bounded matrix observable \(F\),
\begin{equation}
 {d\over dt}\mathbb E_{r_{\alpha,t}^{\rm add}}F
 =
 \int F(w_\alpha-h_\tau)\,dU
 =
 \mathbb E_{r_{\alpha,t}^{\rm add}}
 [F\Sigma_{\alpha,t}^{\rm add}].
 \label{eq:newV41-additive-derivative}
\end{equation}
\end{lemma}

\begin{proof}
Differentiate \eqref{eq:newV41-additive}.  Its integral is one for
every \(t\), so the derivative has integral zero.  Division by the
strictly positive real density gives the displayed score and the
two equalities.
\end{proof}

\begin{corollary}[No logarithmic marked density]
\label{cor:newV41-no-log-density}
The additive score-weighted density is exactly
\begin{equation}
 r_{\alpha,t}^{\rm add}\Sigma_{\alpha,t}^{\rm add}
 =
 w_\alpha-h_\tau,
 \label{eq:newV41-signed-endpoint}
\end{equation}
independent of \(t\).  It is an entire class function and contains
no branch at a complex heat-kernel zero.
\end{corollary}

\begin{proof}
This is \eqref{eq:newV41-additive-score}.  Both endpoint character
series extend entire: the heat series by Gaussian convergence and
the Wilson density directly from \(e^{\alpha\cos\theta}\).
\end{proof}

\section{A positive endpoint-envelope normalization}
\label{sec:newV41-envelope}

For each sector put
\begin{equation}
 p_{\alpha,\tau}(n)
 :=
 e^{-\tau c_n}+q_{\alpha,n}>0.
 \label{eq:newV41-envelope}
\end{equation}
If \(B\) is a tensor block, let
\(\mathcal P_B\) be the positive central operator which acts by
\(p_{\alpha,\tau}(n)\) on every total-spin-\(n/2\) summand of
\(\bigotimes_{i\in B}R_i\).

\begin{theorem}[Contractive envelope-normalized block rows]
\label{thm:newV41-envelope-contractive}
Let
\begin{align*}
 M_{t,B}^{\rm add}
 &:=
 \mathbb E_{r_{\alpha,t}^{\rm add}}
 \bigotimes_{i\in B}D^{R_i}(U),\\
 N_B^{\rm add}
 &:=
 \int (w_\alpha-h_\tau)
 \bigotimes_{i\in B}D^{R_i}(U)\,dU.
\end{align*}
Then, at every matrix amplification,
\begin{align}
 \|\mathcal P_B^{-1}M_{t,B}^{\rm add}\|_{\rm cb}
 &\le1,                                                       \label{eq:newV41-envelope-M}\\
 \|\mathcal P_B^{-1}N_B^{\rm add}\|_{\rm cb}
 &\le1.                                                       \label{eq:newV41-envelope-N}
\end{align}
The bounds are uniform in \(t\), all input representations, and
the block size.
\end{theorem}

\begin{proof}
Central Schur decomposition makes all three operators scalar on
each total-spin sector \(n\).  Their respective eigenvalues are
\[
 (1-t)e^{-\tau c_n}+tq_{\alpha,n},
 \qquad
 q_{\alpha,n}-e^{-\tau c_n},
 \qquad
 p_{\alpha,\tau}(n).
\]
Both endpoint multipliers are positive.  Therefore
\[
 0\le
 {(1-t)e^{-\tau c_n}+tq_{\alpha,n}\over
   e^{-\tau c_n}+q_{\alpha,n}}
 \le1
\]
and
\[
 \left|
 {q_{\alpha,n}-e^{-\tau c_n}\over
  e^{-\tau c_n}+q_{\alpha,n}}
 \right|
 \le1.
\]
Direct sums and spectator amplification preserve these scalar
operator bounds.
\end{proof}

\begin{remark}[Why the envelope does not yet prove the forest card]
\label{rem:newV41-envelope-not-enough}
The heat normalization \(e^{\tau C_B^{\rm tot}}\) is a single
functional calculus with the block-local polynomial gates built in
V33--V38.  The envelope inverse
\(\mathcal P_B^{-1}\) is contractive on complete moments but is not
known to factor compatibly when blocks merge.  Taking norms of
moments and then using Bell-number moment--cumulant inversion would
lose the required exponential arity grade.
\end{remark}

\section{The exact envelope merge cocycle}
\label{sec:newV41-cocycle}

For a partition \(\pi\in\Pi(B)\), define
\begin{equation}
 \mathfrak C_{B,\pi}
 :=
 \mathcal P_B^{-1}
 \bigotimes_{A\in\pi}\mathcal P_A.
 \label{eq:newV41-cocycle}
\end{equation}
All factors are central on their stated tensor blocks; the
expression means their canonical embeddings in the full block
algebra.

\begin{lemma}[Exact normalized cumulant formula]
\label{lem:newV41-normalized-cumulant}
Put
\[
 Y_{t,A}:=\mathcal P_A^{-1}M_{t,A}^{\rm add}.
\]
Then the envelope-normalized raw cumulant is exactly
\begin{equation}
 \mathcal P_B^{-1}\kappa_B^{\,r_{\alpha,t}^{\rm add}}
 =
 \sum_{\pi\in\Pi(B)}
 (-1)^{|\pi|-1}(|\pi|-1)!
 \mathfrak C_{B,\pi}
 \bigotimes_{A\in\pi}Y_{t,A}.
 \label{eq:newV41-normalized-cumulant}
\end{equation}
Its \(t\)-derivative is obtained by replacing exactly one
\(Y_{t,A}\) by
\(\mathcal P_A^{-1}N_A^{\rm add}\) and summing over
\(A\in\pi\).
\end{lemma}

\begin{proof}
Start from ordinary moment--cumulant inversion and substitute
\(M_{t,A}^{\rm add}=\mathcal P_A Y_{t,A}\) in every block.
Multiplication on the left by \(\mathcal P_B^{-1}\) gives
\eqref{eq:newV41-normalized-cumulant}.  The envelope operators are
independent of \(t\), while
\(\partial_tM_{t,A}^{\rm add}=N_A^{\rm add}\) by
\Cref{lem:newV41-additive-score}; the product rule proves the
marked formula.
\end{proof}

\begin{definition}[Envelope merge-cocycle card]
\label{def:newV41-EMCC}
EMCC is the existence of a partition-gate interpolation of
\(\mathfrak C_{B,\pi}\) whose connected forest derivatives,
already multiplied by the complete unmarked or singly marked
contractive rows of
\Cref{thm:newV41-envelope-contractive}, obey the weighted
Pr\"ufer bound
\[
 K^{|B|}
 \left(\prod_{i\in B}b_i\right)
 \left(\sum_{i\in B}b_i\right)^{|B|-2}.
\]
No norm is taken on \(\mathfrak C_{B,\pi}\) separately.
\end{definition}

\begin{proposition}[EMCC is the corrected local sufficient gate]
\label{prop:newV41-EMCC-sufficient}
EMCC for the unmarked and singly marked additive rows implies the
raw Wilson aggregate strong card.
\end{proposition}

\begin{proof}
Apply the connected forest identity to the complete right side of
\eqref{eq:newV41-normalized-cumulant}.  EMCC gives the desired
weighted bound uniformly in \(t\).  The marked identity in
\Cref{lem:newV41-normalized-cumulant} and integration from
\(t=0\) to \(1\) compare the raw heat and Wilson cumulants.  The
heat endpoint card is the V29 result.  Finally use the V16
balancing equivalence.  No inverse heat multiplier is applied to
the Wilson endpoint or to a logarithmic density.
\end{proof}

\begin{programme}[V42: exact envelope merge test]
\label{prog:newV41-V42}
\begin{enumerate}[leftmargin=3em]
\item Diagonalize the binary cocycle
      \(\mathcal P_{12}^{-1}(\mathcal P_1\otimes\mathcal P_2)\)
      on every Clebsch--Gordan sector.
\item Use the exact Bessel formula
      \(q_{\alpha,n}=I_{n+1}(\alpha)/I_1(\alpha)\) and the heat
      Gaussian to determine its sharp high-spin growth.
\item Test whether that growth is paid by a canonical star or
      requires the cocycle to remain multiplied by the signed
      endpoint difference.
\item Compute the binary and ternary complete marked cocycles
      before proposing an all-arity bound.
\item If a separate cocycle norm is exponentially too large,
      retain it inside the signed endpoint difference, in the
      sense required by EMCC; do not return to Bell-number bounds.
\end{enumerate}
\end{programme}

\begin{firewall}[Claim boundary after V41]
\label{fire:newV41-boundary}
V41 proves that the exact heat-normalized logarithmic singleton
score grows faster than every \(e^{Ln}\) along a spectral
subsequence.  It therefore disproves PFSC, geometric-bridge WEEMC,
and exact geometric WPRAC under the all-representation heat
normalization.  It independently proves that the Wilson endpoint
itself is not inverse-heat bounded.  It replaces the geometric
bridge by an exact additive bridge, proves uniform contractivity of
all complete envelope-normalized block moments and marked moments,
and isolates the merge cocycle EMCC.  It does not prove EMCC, the
raw Wilson aggregate strong card, all-order strong MTA, WUFC, CPM,
cutoff removal, reflection positivity,
Osterwalder--Schrader reconstruction, a continuum Yang--Mills
measure, or a positive Hamiltonian gap.  The full Yang--Mills mass
gap has not been proved.
\end{firewall}

\begin{theorem}[V41 result ledger]
\label{thm:newV41-results}
V41:
\begin{enumerate}[label=\textup{(V41.\arabic*)},leftmargin=6.5em]
\item proves that Gaussian character decay forces entire
      class-angle continuation;
\item transfers the V40 heat zero into a logarithmic monodromy of
      the complete geometric score density;
\item proves an all-spin, arity-one heat-normalization no-go;
\item computes the exact positive Wilson Bessel multiplier and
      proves the endpoint inverse-heat mismatch;
\item constructs the exact additive heat--Wilson bridge;
\item proves contractive endpoint-envelope block rows at every
      amplification;
\item derives the exact envelope-normalized cumulant and marked
      cocycle formulas; and
\item isolates EMCC as the corrected upstream local gate.
\end{enumerate}
\end{theorem}

\begin{proof}
Use
\Cref{lem:newV41-Gaussian-entire,
lem:newV41-monodromy,
thm:newV41-score-not-entire,
thm:newV41-singleton-no-go,
cor:newV41-PFSC-false,
lem:newV41-Wilson-multiplier,
prop:newV41-Wilson-not-heat,
lem:newV41-additive-score,
thm:newV41-envelope-contractive,
lem:newV41-normalized-cumulant,
prop:newV41-EMCC-sufficient}.
\end{proof}

\paragraph{Parent-version record.}
V41 was formed from
\texttt{Renewal\_Yang\_Mills\_Mass\_Gap\_Research\_Notes\_V40.tex}.
The V40 parent had \(19\,163\) logical source lines,
\(606\,982\) bytes, and SHA--256
\[
 \texttt{e870e3be5aa42eee1ab1acb35b266cb56fcf303858c0019bd50f1f8b6296aa0e}.
\]
The next compulsory companion audit is V45.

% =============================================================================
% END APPENDED FIFTH-SERIES V41 MATERIAL
% =============================================================================

% =============================================================================
% BEGIN APPENDED FIFTH-SERIES V42 MATERIAL
% =============================================================================

\chapter{V42: Thermal-window correction and the first envelope
merge bounds}
\label{chap:newV42-thermal-correction}

\section{Append-only correction of the V41 scope}
\label{sec:newV42-correction}

No companion audit is due at V42.  Before using the V41
high-spin theorem, its quantifiers must be compared with the
actual WPRAC quantifiers.

\begin{correction}[One block is not one admissible high-spin input]
\label{corr:newV42-one-block}
\Cref{thm:newV41-singleton-no-go} is correct as an
all-representation statement at fixed \(\tau\).  The description
of it in \Cref{ass:newV41-v26-relation} as an
\emph{arity-one} obstruction is too broad for the original local
programme.  WPRAC in \Cref{def:newV35-WPRAC} assumes each input
representation is thermal:
\[
 b_i^2=s(1+C_2(R_i))<1.
\]
At fixed \(s\), a one-input representation in this window cannot
run along the unbounded Peter--Weyl subsequence used in V41.
Thus V41 does not disprove the thermal one-input row.
\end{correction}

\begin{proposition}[The thermal one-input geometric score is
bounded]
\label{prop:newV42-thermal-singleton}
Let \(R\) be thermal and let \(N_{\alpha,t,R}\) be its exact
geometric-bridge marked moment.  Then
\begin{equation}
 \|e^{sC_2(R)}N_{\alpha,t,R}\|_{\rm cb}
 \le Cs
 \label{eq:newV42-thermal-singleton}
\end{equation}
uniformly in \(t\) and \(R\).
\end{proposition}

\begin{proof}
For a unitary representation,
\(\|D^R(U)\|=1\).  The V32 adaptive score row at \(p=1\) gives
\[
 \|N_{\alpha,t,R}\|_{\rm cb}
 \le
 \mathbb E_{\nu_{\alpha,t}}|S_{\alpha,t}|
 \le C\alpha^{-1}\le Cs.
\]
Thermality gives \(sC_2(R)<1\), hence
\(\|e^{sC_2(R)}\|\le e\).  Spectator amplification changes
neither bound.
\end{proof}

\begin{correction}[WPRAC remains open in its stated domain]
\label{corr:newV42-WPRAC-open}
The last sentence of
\Cref{cor:newV41-PFSC-false} is valid only for an
all-representation extension of WPRAC.  It does not refute WPRAC
as stated in V35.  A large tensor block of thermal inputs can
indeed contain arbitrarily high total-spin sectors as the arity
grows, but WPRAC estimates the \emph{connected activity}, not an
isolated marked moment.  Moment--cumulant M\"obius terms may cancel
the high-spin tail.  V41 did not estimate that connected
cancellation.
\end{correction}

\begin{assessment}[What V41 actually retires]
\label{ass:newV42-retirement-scope}
V41 rigorously retires:
\begin{enumerate}[label=\textup{(\roman*)},leftmargin=3.5em]
\item PFSC when it demands one global all-spin heat-normalized
      spectral multiplier of V38 even-entire type;
\item WEEMC under that same global all-spin interpretation; and
\item inverse-heat domination of the Wilson endpoint over all
      representations.
\end{enumerate}
It does not retire the fixed thermal jets, binary or ternary
WPRAC, or the possibility that the complete connected
thermal-input activity satisfies all-arity WPRAC.  The additive
envelope bridge is therefore a parallel corrected route, not yet a
proof that the geometric route must be abandoned.
\end{assessment}

\begin{firewall}[No automatic high-spin-to-connected lift]
\label{fire:newV42-no-automatic-lift}
From V42 onward, an unbounded block moment may be used to disprove
a connected card only after its contribution is shown to survive
the complete partition M\"obius sum on an admissible tensor family.
The presence of that moment on a highest-spin sector is not by
itself sufficient.
\end{firewall}

\section{Binary diagonalization of the envelope cocycle}
\label{sec:newV42-binary}

Retain
\[
 p_{\alpha,s}(R)
 =
 q^H_{s,R}+q^W_{\alpha,R}
 =
 e^{-sC_2(R)}+q^W_{\alpha,R}.
\]
For irreducible inputs \(R_1,R_2\), let \(R\) run over their
Clebsch--Gordan summands.

\begin{lemma}[Exact binary cocycle eigenvalue]
\label{lem:newV42-binary-eigenvalue}
On the total-\(R\) summand of \(R_1\otimes R_2\),
\begin{equation}
 \mathfrak C_{\{1,2\},\{\{1\},\{2\}\}}
 =
 {p_{\alpha,s}(R_1)p_{\alpha,s}(R_2)\over
  p_{\alpha,s}(R)}\,I_R.
 \label{eq:newV42-binary-eigenvalue}
\end{equation}
\end{lemma}

\begin{proof}
For a one-point block, \(\mathcal P_{\{i\}}\) is the scalar
\(p_{\alpha,s}(R_i)\).  On a total-\(R\) summand,
\(\mathcal P_{\{1,2\}}\) is the scalar
\(p_{\alpha,s}(R)\).  Substitute these into
\eqref{eq:newV41-cocycle}.
\end{proof}

\begin{theorem}[Uniform binary cocycle bound on thermal inputs]
\label{thm:newV42-binary-cocycle}
If \(b_i^2=s(1+C_2(R_i))<1\), then
\begin{equation}
 \left\|
 \mathfrak C_{\{1,2\},\{\{1\},\{2\}\}}
 \right\|_{\rm cb}
 \le4e^4.
 \label{eq:newV42-binary-cocycle}
\end{equation}
The constant is uniform in \(\alpha,s\) and the two thermal
representations.
\end{theorem}

\begin{proof}
Every central probability multiplier has modulus at most one; the
heat and Wilson multipliers here are positive.  Hence
\[
 0<p_{\alpha,s}(R_i)\le2.
\]
The total-Casimir row triangle inequality gives
\[
 \sqrt{sC_2(R)}
 \le
 \sqrt{sC_2(R_1)}+\sqrt{sC_2(R_2)}
 \le b_1+b_2<2.
\]
Therefore
\[
 p_{\alpha,s}(R)
 \ge e^{-sC_2(R)}
 \ge e^{-4}.
\]
Insert these bounds in
\eqref{eq:newV42-binary-eigenvalue}.  Central direct sums and
amplification preserve the maximum scalar bound.
\end{proof}

\begin{corollary}[Crude complete binary envelope row]
\label{cor:newV42-binary-row}
With the notation of
\Cref{thm:newV41-envelope-contractive},
\begin{align}
 \left\|
 \mathcal P_{12}^{-1}
 \kappa_2^{\,r_{\alpha,t}^{\rm add}}
 \right\|_{\rm cb}
 &\le1+4e^4,                                                  \label{eq:newV42-binary-unmarked}\\
 \left\|
 \partial_t\left[
 \mathcal P_{12}^{-1}
 \kappa_2^{\,r_{\alpha,t}^{\rm add}}
 \right]
 \right\|_{\rm cb}
 &\le1+8e^4.                                                  \label{eq:newV42-binary-marked}
\end{align}
\end{corollary}

\begin{proof}
For \(B=\{1,2\}\),
\Cref{lem:newV41-normalized-cumulant} reads
\[
 \mathcal P_{12}^{-1}\kappa_2
 =
 Y_{t,12}-\mathfrak C_{12,1|2}Y_{t,1}Y_{t,2}.
\]
Every \(Y\) and every normalized marked replacement has norm at
most one by \Cref{thm:newV41-envelope-contractive}.  Apply
\Cref{thm:newV42-binary-cocycle} and the product rule.
\end{proof}

\begin{warning}[The thermal clocks are still missing]
\label{warn:newV42-binary-clocks}
Bounds \eqref{eq:newV42-binary-unmarked}--%
\eqref{eq:newV42-binary-marked} are uniform but do not yet contain
the required factor \(b_1b_2\).  They prove that the envelope
cocycle itself causes no binary blow-up.  They do not prove binary
EMCC.  The missing gain must come from the complete endpoint
difference and connected subtraction, not from contractivity
alone.
\end{warning}

\section{The remote high-spin cocycle has exponential, not
quadratic, grade}
\label{sec:newV42-high-spin}

It is useful to test whether the envelope merely moves the
inverse-heat \(e^{cn^2}\) failure into the cocycle.  At fixed
\(\alpha,s\), it does not in the far Wilson tail.

\begin{lemma}[Fixed-coupling Bessel tail]
\label{lem:newV42-Bessel-tail}
For fixed \(\alpha>0\),
\begin{equation}
 I_k(\alpha)
 =
 {(\alpha/2)^k\over k!}
 \left(1+O_\alpha(k^{-1})\right)
 \qquad(k\to\infty).
 \label{eq:newV42-Bessel-tail}
\end{equation}
\end{lemma}

\begin{proof}
The defining positive series is
\[
 I_k(\alpha)
 =
 {(\alpha/2)^k\over k!}
 \sum_{\ell\ge0}
 {(\alpha^2/4)^\ell k!\over
  \ell!(k+\ell)!}.
\]
The \(\ell=0\) term is one.  For \(\ell\ge1\),
\[
 {k!\over(k+\ell)!}\le(k+1)^{-\ell},
\]
so the rest is bounded by
\[
 \exp\left({\alpha^2\over4(k+1)}\right)-1
 =
 O_\alpha(k^{-1}).
\]
\end{proof}

\begin{proposition}[Balanced far-tail merge asymptotic]
\label{prop:newV42-far-tail}
Fix \(\alpha,s>0\).  Let \(a,b\to\infty\) with
\(\min(a,b)\to\infty\), and put \(N=a+b\).  On the highest
Clebsch--Gordan sector,
\begin{align}
 {p_{\alpha,s}(a)p_{\alpha,s}(b)\over p_{\alpha,s}(N)}
 &=
 {\alpha\over2I_1(\alpha)}
 {(N+1)!\over(a+1)!(b+1)!}
 (1+o(1))                                                     \label{eq:newV42-tail-ratio}\\
 &\le C_\alpha\,2^{N+2}(1+o(1)).
 \label{eq:newV42-tail-exponential}
\end{align}
Thus the far-tail cocycle has exponential arity grade rather than
the quadratic exponential produced by inverse heat smoothing.
\end{proposition}

\begin{proof}
For fixed \(s,\alpha\),
\[
 {e^{-sC_2(n)}\over q^W_{\alpha,n}}\longrightarrow0
\]
because the numerator has a negative quadratic logarithm in \(n\),
whereas \Cref{lem:newV42-Bessel-tail} gives a negative
\(n\log n+O(n)\) logarithm for the denominator.  Hence
\(p_{\alpha,s}(n)=q^W_{\alpha,n}(1+o(1))\).
Using
\(q^W_{\alpha,n}=I_{n+1}(\alpha)/I_1(\alpha)\) and
\Cref{lem:newV42-Bessel-tail} gives
\eqref{eq:newV42-tail-ratio}.  Finally,
\[
 {(N+1)!\over(a+1)!(b+1)!}
 =
 {1\over N+2}{N+2\choose a+1}
 \le2^{N+2}.
\]
\end{proof}

\begin{assessment}[The unresolved crossover]
\label{ass:newV42-crossover}
\Cref{thm:newV42-binary-cocycle} controls two thermal input
blocks, and \Cref{prop:newV42-far-tail} controls the fixed-coupling
remote Wilson tail.  Neither theorem controls an arbitrary merge
of two growing thermal subblocks uniformly through the crossover
\[
 C_2(R)\asymp\alpha^2
\quad\hbox{or}\quad
 n\asymp\alpha,
\]
nor does either produce the required tree clocks.  That crossover
and the connected endpoint difference are the live EMCC
obligations.
\end{assessment}

\section{The connected object must be tested directly}
\label{sec:newV42-connected}

For a bounded scalar observable \(X\), define the additive-bridge
source partition function
\begin{equation}
 Z_{\alpha,s,t}^X(z)
 :=
 \int_{SU(2)}
 e^{zX(U)}r_{\alpha,t}^{\rm add}(U)\,dU.
 \label{eq:newV42-source-Z}
\end{equation}

\begin{lemma}[Connected and marked source identities]
\label{lem:newV42-source-identities}
Near \(z=0\),
\begin{align}
 \partial_z^m\log Z_{\alpha,s,t}^X(0)
 &=
 \kappa_m^{\,r_{\alpha,t}^{\rm add}}(X,\ldots,X),
 \label{eq:newV42-source-cumulant}\\
 \partial_t\partial_z^m\log Z_{\alpha,s,t}^X(0)
 &=
 \kappa_{m+1}^{\,r_{\alpha,t}^{\rm add}}
 (X,\ldots,X,\Sigma_{\alpha,t}^{\rm add}).
 \label{eq:newV42-source-marked}
\end{align}
\end{lemma}

\begin{proof}
The first identity is the standard cumulant-generating identity,
obtained by differentiating the finite exponential series under
the compact-group integral.  For the second, differentiate in
\(t\) and use
\(\partial_t r_{\alpha,t}^{\rm add}=w_\alpha-h_s
=r_{\alpha,t}^{\rm add}\Sigma_{\alpha,t}^{\rm add}\).
Mixed derivatives commute near the origin because the integrand
is entire in \(z\) and uniformly bounded on compact \(z\)-sets.
\end{proof}

\begin{assessment}[Correct analytic target]
\label{ass:newV42-correct-target}
The complex zeros relevant to connected all-arity estimates are
the zeros of \(Z_{\alpha,s,t}^X(z)\) in the \emph{source}
variable, not the zeros of \(h_s(\theta)\) in the physical class
angle.  A zero-free source disk of radius comparable to the
inverse thermal mark gives factorial Cauchy bounds for the complete
cumulant and automatically includes every partition cancellation.
\end{assessment}

\begin{programme}[V43: source-plane zero-free disk]
\label{prog:newV42-V43}
\begin{enumerate}[leftmargin=3em]
\item Begin with \(X=\cos\theta\), for which the Wilson source
      function is the exact shifted Bessel quotient
      \(Z_{\alpha+z}/Z_\alpha\).
\item Locate its complex zeros and prove a source disk on the
      Wilson scale without using a physical heat-kernel strip.
\item Derive the corresponding all-order scalar Wilson cumulant
      bounds by Cauchy's formula.
\item Treat the heat source function and then the additive convex
      combination, keeping the two endpoint source functions
      assembled when testing zeros.
\item Only after the scalar source test succeeds, seek its
      noncommutative block/polydisc analogue and the missing
      \(b_1b_2\) binary endpoint-difference gain.
\end{enumerate}
\end{programme}

\begin{firewall}[Claim boundary after V42]
\label{fire:newV42-boundary}
V42 corrects the scope of V41: the all-spin one-block no-go does
not itself disprove thermal-input WPRAC or its connected
all-arity activity.  It proves the thermal geometric singleton
row, diagonalizes the envelope binary cocycle, proves its uniform
bound for two thermal inputs, and proves exponential far-tail
merge grade at fixed coupling.  It also identifies the
source-partition zeros as the correct connected analytic target.
It does not prove binary EMCC, all-arity WPRAC, RGSC, the raw
Wilson aggregate strong card, all-order strong MTA, WUFC, CPM,
cutoff removal, reflection positivity,
Osterwalder--Schrader reconstruction, a continuum Yang--Mills
measure, or a positive Hamiltonian gap.  The full Yang--Mills mass
gap has not been proved.
\end{firewall}

\begin{theorem}[V42 result ledger]
\label{thm:newV42-results}
V42:
\begin{enumerate}[label=\textup{(V42.\arabic*)},leftmargin=6.5em]
\item repairs the V41 arity/thermal-window overstatement;
\item proves the geometric marked singleton throughout the thermal
      window;
\item preserves PFSC/WEEMC no-go only in their global all-spin
      form and restores thermal connected WPRAC to open status;
\item diagonalizes and uniformly bounds the binary envelope
      cocycle for thermal inputs;
\item proves the fixed-coupling far-tail cocycle has exponential
      merge grade;
\item isolates the uncontrolled growing-block crossover and the
      missing binary clocks; and
\item replaces physical-angle zeros by source-plane zeros as the
      next connected target.
\end{enumerate}
\end{theorem}

\begin{proof}
Use
\Cref{prop:newV42-thermal-singleton,
thm:newV42-binary-cocycle,
cor:newV42-binary-row,
lem:newV42-Bessel-tail,
prop:newV42-far-tail,
lem:newV42-source-identities}.
\end{proof}

\paragraph{Parent-version record.}
V42 was formed from
\texttt{Renewal\_Yang\_Mills\_Mass\_Gap\_Research\_Notes\_V41.tex}.
The V41 parent had \(19\,796\) logical source lines,
\(626\,272\) bytes, and SHA--256
\[
 \texttt{980514090d33a8c6799bf9aef6a011234e1ba62415f5b58de5e0434f4e5af7dc}.
\]
The next compulsory companion audit is V45.

% =============================================================================
% END APPENDED FIFTH-SERIES V42 MATERIAL
% =============================================================================

% =============================================================================
% BEGIN APPENDED FIFTH-SERIES V43 MATERIAL
% =============================================================================

\chapter{V43: Source-plane zero-free disks and a scalar
all-arity card}
\label{chap:newV43-source-disk}

\section{The exact Wilson source function}
\label{sec:newV43-Wilson-source}

No companion audit is due at V43.  Put
\[
 X(U):={1\over2}\operatorname{tr}U=\cos\theta,
 \qquad
 V(U):=1-X(U)=1-\cos\theta.
\]
Thus \(0\le V\le\theta^2/2\).

\begin{lemma}[Shifted Wilson source]
\label{lem:newV43-shift}
For the Wilson endpoint,
\begin{equation}
 Z_{\alpha}^{W,X}(z)
 :=
 \mathbb E_{w_\alpha}e^{zX}
 =
 {Z_{\alpha+z}\over Z_\alpha},
 \qquad
 Z_\beta={2I_1(\beta)\over\beta},
 \label{eq:newV43-shift}
\end{equation}
where the quotient at \(\beta=0\) is removable.
\end{lemma}

\begin{proof}
Multiplication of \(e^{\alpha\cos\theta}\) by
\(e^{z\cos\theta}\) replaces \(\alpha\) by \(\alpha+z\).
The normalization formula for \(Z_\beta\) was proved in
\Cref{lem:newV41-Wilson-multiplier}.
\end{proof}

\begin{lemma}[Canonical product for the Wilson normalization]
\label{lem:newV43-Bessel-product}
Let \(0<j_{1,1}<j_{1,2}<\cdots\) be the positive zeros of
\(J_1\).  Then
\begin{equation}
 {2I_1(z)\over z}
 =
 \prod_{k=1}^\infty
 \left(1+{z^2\over j_{1,k}^2}\right),
 \label{eq:newV43-Bessel-product}
\end{equation}
normally on compact sets.  In particular, the zeros of
\(Z_z\) are exactly
\[
 z=\pm i j_{1,k},
 \qquad k\ge1,
\]
and they are simple.
\end{lemma}

\begin{proof}
The Bessel equation for \(J_1(\lambda r)\) on \(0<r<1\) is the
regular Sturm--Liouville equation
\[
 -(r u')'+{1\over r}u=\lambda^2r u,
 \qquad u(0)=0.
\]
If the regular solution also obeys \(u(1)=0\), multiplication by
\(\overline u\) and integration gives
\[
 \int_0^1\left(r|u'|^2+{1\over r}|u|^2\right)\,dr
 =
 \lambda^2\int_0^1r|u|^2\,dr.
\]
Thus every nonzero zero parameter has \(\lambda^2>0\) and hence is
real.  ODE uniqueness makes the zeros simple.  Sturm comparison
gives \(j_{1,k}\ge ck\), so
\(\sum_kj_{1,k}^{-2}<\infty\).

The even entire function \(2J_1(z)/z\) has order one, value one at
zero, and zeros exactly \(\pm j_{1,k}\).  Hadamard factorization,
paired over the two signs, therefore gives
\[
 {2J_1(z)\over z}
 =
 \prod_{k\ge1}\left(1-{z^2\over j_{1,k}^2}\right).
\]
Finally \(I_1(z)=-iJ_1(iz)\), which proves
\eqref{eq:newV43-Bessel-product}.
\end{proof}

\begin{corollary}[Exact Wilson source zero set]
\label{cor:newV43-Wilson-zeros}
The zeros of \(Z_{\alpha}^{W,X}(z)\) are
\begin{equation}
 z=-\alpha\pm i j_{1,k},
 \qquad k\ge1.
 \label{eq:newV43-Wilson-zeros}
\end{equation}
Consequently its zero-free disk at the source origin has radius
\[
 \sqrt{\alpha^2+j_{1,1}^2}>\alpha.
\]
\end{corollary}

\begin{proof}
Use \eqref{eq:newV43-shift} and
\Cref{lem:newV43-Bessel-product}.  The denominator is a fixed
positive number.
\end{proof}

\section{Exact Wilson cumulants from the zero product}
\label{sec:newV43-Wilson-cumulants}

\begin{theorem}[Wilson scalar cumulant formula]
\label{thm:newV43-Wilson-cumulants}
For every \(m\ge2\),
\begin{equation}
 \kappa_m^{w_\alpha}(X,\ldots,X)
 =
 (-1)^{m-1}(m-1)!
 \sum_{k\ge1}
 \left[
 {1\over(\alpha-i j_{1,k})^m}
 +
 {1\over(\alpha+i j_{1,k})^m}
 \right].
 \label{eq:newV43-Wilson-cumulants}
\end{equation}
For \(\alpha\ge1\),
\begin{equation}
 \left|
 \kappa_m^{w_\alpha}(X,\ldots,X)
 \right|
 \le
 C(m-1)!\,\alpha^{1-m},
 \qquad m\ge2.
 \label{eq:newV43-Wilson-cumulant-bound}
\end{equation}
\end{theorem}

\begin{proof}
The cumulant is the \(m\)-th derivative at zero of
\[
 \log Z_\alpha^{W,X}(z)
 =
 \log Z_{\alpha+z}-\log Z_\alpha.
\]
Differentiate the normally convergent paired logarithm of
\eqref{eq:newV43-Bessel-product}.  For \(m\ge2\) the resulting
series is absolutely convergent and gives
\eqref{eq:newV43-Wilson-cumulants}.

Sturm comparison also gives \(j_{1,k}\ge ck\).  Hence
\begin{align*}
 \sum_{k\ge1}(\alpha^2+j_{1,k}^2)^{-m/2}
 &\le
 C\sum_{k\ge1}(\alpha^2+k^2)^{-m/2}\\
 &\le
 C\left[
 \alpha^{-m}
 +\int_0^\infty(\alpha^2+x^2)^{-m/2}\,dx
 \right]\\
 &\le C\alpha^{1-m}
\end{align*}
uniformly for \(m\ge2\) and \(\alpha\ge1\).  Apply the triangle
inequality in \eqref{eq:newV43-Wilson-cumulants}.
\end{proof}

\begin{assessment}[Sharp scalar clock count]
\label{ass:newV43-Wilson-clocks}
The power \(\alpha^{1-m}\) in
\eqref{eq:newV43-Wilson-cumulant-bound} is precisely the scalar
tree scale: \(X\) is a bounded functional of the fundamental
representation, and
\[
 b_{\rm f}^{\,m}(mb_{\rm f})^{m-2}
 =
 m^{m-2}b_{\rm f}^{\,2m-2}
 \asymp
 m^{m-2}\alpha^{1-m}.
\]
Since \((m-1)!\le m^{m-2}\) for \(m\ge2\), the exact Wilson
formula proves this one-direction scalar projection of the
all-arity local card.
\end{assessment}

\section{A uniform endpoint exponential moment}
\label{sec:newV43-exp-moment}

The exact Bessel product is special to \(X=\cos\theta\) at the
Wilson endpoint.  A concentration argument gives a common disk
for the heat endpoint and every additive mixture.

\begin{lemma}[Scaled \(V\)-exponential moment]
\label{lem:newV43-V-exp}
Let \(\mu\) be either the calibrated heat endpoint or the Wilson
endpoint, with \(s=\tau_\alpha^\sharp\asymp\alpha^{-1}\).
There are numerical \(\eta_0,C>0\) such that
\begin{equation}
 \mathbb E_\mu
 \exp\left(\eta_0{V\over s}\right)
 +
 \mathbb E_\mu
 \left[
 {V\over s}
 \exp\left(\eta_0{V\over s}\right)
 \right]
 \le C.
 \label{eq:newV43-V-exp}
\end{equation}
The same bound holds for every additive mixture
\((1-t)h_s+t w_\alpha\).
\end{lemma}

\begin{proof}
The endpoint cases \(t=0,1\) of the V30 adaptive concentration
theorem give
\[
 \left(\mathbb E_\mu\theta^{2k}\right)^{1/(2k)}
 \le C\sqrt{ks}.
\]
Since \(V\le\theta^2/2\),
\[
 \mathbb E_\mu\left({V\over s}\right)^k
 \le(Ck)^k.
\]
Using \(k!\ge(k/e)^k\), the exponential power series converges
uniformly when \(\eta_0>0\) is sufficiently small.  Differentiating
that scalar series once in \(\eta_0\) gives the second bound after
decreasing \(\eta_0\) if necessary.  An additive-mixture
expectation is the convex combination of the two endpoint
expectations.
\end{proof}

\section{The common additive source disk}
\label{sec:newV43-common-disk}

For \(\mu\in\{H,W\}\), write
\begin{align}
 A_\mu(z)
 &:=
 e^{-z}Z_\mu^X(z)
 =
 \mathbb E_\mu e^{-zV},                                      \label{eq:newV43-A-endpoint}\\
 A_t(z)
 &:=
 (1-t)A_H(z)+tA_W(z).                                        \label{eq:newV43-A-mixture}
\end{align}
Then
\[
 Z_{\alpha,s,t}^X(z)=e^zA_t(z).
\]

\begin{theorem}[Uniform source-plane zero-free disk]
\label{thm:newV43-zero-free}
There is a numerical \(\eta>0\) such that, uniformly in
\(t\in[0,1]\) and sufficiently large \(\alpha\),
\begin{equation}
 |A_t(z)-1|\le {1\over4}
 \qquad
 \left(|z|\le{\eta\over s}\right).
 \label{eq:newV43-close-one}
\end{equation}
Consequently \(Z_{\alpha,s,t}^X\) has no zero in the disk
\begin{equation}
 |z|\le{\eta\over s}.
 \label{eq:newV43-zero-free}
\end{equation}
\end{theorem}

\begin{proof}
For either endpoint and \(|z|\le\eta/s\),
\begin{align*}
 |A_\mu(z)-1|
 &\le
 |z|\mathbb E_\mu
 \left[Ve^{|z|V}\right]\\
 &\le
 \eta\,
 \mathbb E_\mu
 \left[
 {V\over s}e^{\eta V/s}
 \right].
\end{align*}
Choose \(\eta\le\eta_0\) so small that the last expression is at
most \(1/4\), using \Cref{lem:newV43-V-exp}.  The disk is convex
under the endpoint mixture, so the same bound holds for \(A_t\).
It implies \(|A_t|\ge3/4\).  Since \(e^z\ne0\),
\(Z_{\alpha,s,t}^X=e^zA_t\) is zero-free there.
\end{proof}

\section{Scalar all-arity additive cards}
\label{sec:newV43-scalar-cards}

\begin{theorem}[Uniform scalar cumulant and marked-cumulant bounds]
\label{thm:newV43-scalar-cards}
There is a numerical \(K\) such that, for all \(m\ge2\),
\begin{align}
 \left|
 \kappa_m^{\,r_{\alpha,t}^{\rm add}}(X,\ldots,X)
 \right|
 &\le
 K^m m!s^m,                                                   \label{eq:newV43-scalar-unmarked}\\
 \left|
 \kappa_{m+1}^{\,r_{\alpha,t}^{\rm add}}
 (X,\ldots,X,\Sigma_{\alpha,t}^{\rm add})
 \right|
 &\le
 K^m m!s^m.                                                   \label{eq:newV43-scalar-marked}
\end{align}
Both estimates are uniform in \(t\) and sufficiently large
\(\alpha\).
\end{theorem}

\begin{proof}
On the disk in \Cref{thm:newV43-zero-free}, choose the branch
\[
 \log Z_{\alpha,s,t}^X(z)=z+\log A_t(z)
\]
which vanishes at zero.  Bound
\eqref{eq:newV43-close-one} makes
\(|\log A_t|\le C\) uniformly there.  For \(m\ge2\), the derivative
of the linear term \(z\) vanishes.  Cauchy's estimate on the disk
of radius \(\eta/s\) and
\eqref{eq:newV42-source-cumulant} give
\eqref{eq:newV43-scalar-unmarked}.

Moreover,
\[
 \partial_t\log Z_{\alpha,s,t}^X(z)
 =
 {A_W(z)-A_H(z)\over A_t(z)}.
\]
The numerator has modulus at most \(1/2\) and the denominator at
least \(3/4\) on the same disk.  A second application of Cauchy's
estimate, now using
\eqref{eq:newV42-source-marked}, proves
\eqref{eq:newV43-scalar-marked}.
\end{proof}

\begin{corollary}[Scalar fundamental Pr\"ufer grade]
\label{cor:newV43-Prufer}
After increasing \(K\), both left sides in
\Cref{thm:newV43-scalar-cards} are bounded by
\begin{equation}
 K^m b_{\rm f}^{\,m}(mb_{\rm f})^{m-2}
 \label{eq:newV43-Prufer}
\end{equation}
for \(m\ge2\).
\end{corollary}

\begin{proof}
Write \(b_{\rm f}^2=c_{\rm f}s\) with the fixed
\(c_{\rm f}=1+C_2(\mathrm{fund})>0\).  Since \(s\le1\),
\(s^m\le s^{m-1}\).  Also
\[
 m!\le m^m=m^2m^{m-2}\le4^m m^{m-2}
\]
for \(m\ge2\).  Absorb \(4^m\), \(c_{\rm f}^{1-m}\), and the
constants in \Cref{thm:newV43-scalar-cards} into \(K^m\).
\end{proof}

\begin{assessment}[What cancellation has now been kept]
\label{ass:newV43-cancellation}
\Cref{thm:newV43-scalar-cards} estimates
\(\log Z\) before extracting any Taylor coefficient.  Therefore
the complete partition M\"obius cancellation is already present.
The physical heat-kernel zeros of V40 never enter: the relevant
source function remains within distance \(1/4\) of one throughout
a disk of radius \(c/s\).
\end{assessment}

\section{The noncommutative lift}
\label{sec:newV43-next}

\begin{definition}[Matrix source-disk card]
\label{def:newV43-MSDC}
For thermal representation inputs, MSDC is a
spectator-stable, block-local holomorphic source function whose
logarithmic Fr\'echet derivatives produce the complete tensor
cumulants and marked cumulants and which is invertible on the
weighted operator polydisc with
\[
 |z_i|\le {c\over b_i\sum_jb_j}.
\]
The inverse and logarithmic derivative must be bounded before any
partition or word expansion.
\end{definition}

\begin{assessment}[Scalar evidence, not a matrix theorem]
\label{ass:newV43-not-matrix}
The scalar source \(X=\cos\theta\) commutes with itself and uses a
single disk.  A matrix source
\(\exp(\sum_i z_iD^{R_i}(U))\) is not central and its logarithm is
not defined by scalar nonvanishing.  The scalar proof therefore
does not establish MSDC or the full strong Wilson forest card.
It does identify the correct scale and proves that connected
source cancellation can overcome the false all-spin moment
inference corrected in V42.
\end{assessment}

\begin{programme}[V44: operator source invertibility]
\label{prog:newV43-V44}
\begin{enumerate}[leftmargin=3em]
\item Replace scalar nonvanishing by invertibility in a unital
      Banach algebra and prove the elementary logarithm lemma on a
      norm ball around the identity.
\item Center every unitary matrix input at the identity and use the
      heat/Wilson distance moments to bound the complete
      time-ordered exponential before expanding words.
\item Determine whether the natural polydisc radius is
      \(c/(b_i\sum_jb_j)\), which would reproduce the weighted
      Pr\"ufer grade after multivariate Cauchy.
\item Treat the additive marked derivative as a quotient of the
      two complete endpoint source operators, retaining their
      difference until after inversion.
\item At V45 perform the compulsory SM/NS audit before choosing
      between MSDC and the envelope merge crossover.
\end{enumerate}
\end{programme}

\begin{firewall}[Claim boundary after V43]
\label{fire:newV43-boundary}
V43 locates every Wilson \(X=\cos\theta\) source zero, derives the
exact all-order scalar Wilson cumulant formula, proves a common
source-plane zero-free disk of radius \(c/s\) for the heat,
Wilson, and additive bridge laws, and proves scalar unmarked and
marked all-arity Pr\"ufer-grade bounds.  It does not prove the
matrix source-disk card, EMCC, thermal connected WPRAC, RGSC, the
raw Wilson aggregate strong card, all-order strong MTA, WUFC, CPM,
cutoff removal, reflection positivity,
Osterwalder--Schrader reconstruction, a continuum Yang--Mills
measure, or a positive Hamiltonian gap.  The full Yang--Mills mass
gap has not been proved.
\end{firewall}

\begin{theorem}[V43 result ledger]
\label{thm:newV43-results}
V43:
\begin{enumerate}[label=\textup{(V43.\arabic*)},leftmargin=6.5em]
\item computes the shifted Bessel Wilson source exactly;
\item proves its canonical zero product and source zero set;
\item derives the exact Wilson scalar cumulant formula and its
      sharp \(\alpha^{1-m}\) clock;
\item proves a common scaled exponential moment for
      \(1-\cos\theta\);
\item proves a uniform \(c/s\) zero-free source disk for every
      additive mixture;
\item proves scalar unmarked and marked all-arity cumulant cards at
      Pr\"ufer grade; and
\item isolates Banach-algebra source invertibility as the next
      noncommutative gate.
\end{enumerate}
\end{theorem}

\begin{proof}
Use
\Cref{lem:newV43-shift,
lem:newV43-Bessel-product,
cor:newV43-Wilson-zeros,
thm:newV43-Wilson-cumulants,
lem:newV43-V-exp,
thm:newV43-zero-free,
thm:newV43-scalar-cards,
cor:newV43-Prufer}.
\end{proof}

\paragraph{Parent-version record.}
V43 was formed from
\texttt{Renewal\_Yang\_Mills\_Mass\_Gap\_Research\_Notes\_V42.tex}.
The V42 parent had \(20\,237\) logical source lines,
\(640\,590\) bytes, and SHA--256
\[
 \texttt{a343b548358784f74915622c1b100410b0dc41fd677657ddb6a64c4e2e7758cc}.
\]
The next compulsory companion audit is V45.

% =============================================================================
% END APPENDED FIFTH-SERIES V43 MATERIAL
% =============================================================================

% =============================================================================
% BEGIN APPENDED FIFTH-SERIES V44 MATERIAL
% =============================================================================

\chapter{V44: Operator squarefree sources and the binary matrix
card}
\label{chap:newV44-operator-source}

\section{Squarefree operator source}
\label{sec:newV44-source}

No companion audit is due at V44.  Let
\(\mu_{\alpha,s,t}^{\rm add}\) denote the additive bridge law.
For thermal inputs \(R_1,\ldots,R_m\), define the centered-at-the-
identity matrix increments
\begin{equation}
 A_i(U):=D^{R_i}(U)-I_{R_i}.
 \label{eq:newV44-Ai}
\end{equation}
They act on distinct factors of
\(\mathcal H_m=\bigotimes_{i=1}^mR_i\).  Define the multi-affine
operator source
\begin{equation}
 \mathbb Z_{t,m}(\boldsymbol z)
 :=
 \mathbb E_{\mu_{\alpha,s,t}^{\rm add}}
 \bigotimes_{i=1}^m
 \left(I_{R_i}+z_iA_i(U)\right).
 \label{eq:newV44-Z}
\end{equation}
Thus \(\mathbb Z_{t,m}(\boldsymbol0)=I_{\mathcal H_m}\).

\begin{lemma}[Squarefree logarithm equals the tensor cumulant]
\label{lem:newV44-log-cumulant}
Whenever the Banach-algebra logarithm of
\(\mathbb Z_{t,m}\) is defined near the origin, the coefficient of
\(z_1\cdots z_m\) in
\(\log\mathbb Z_{t,m}(\boldsymbol z)\) is
\begin{equation}
 \kappa_m^{\,\mu_{\alpha,s,t}^{\rm add}}
 (D^{R_1},\ldots,D^{R_m})
 \qquad(m\ge2).
 \label{eq:newV44-log-cumulant}
\end{equation}
The coefficient of \(z_1\cdots z_m\) in
\(\partial_t\log\mathbb Z_{t,m}\) is
\begin{equation}
 \kappa_{m+1}^{\,\mu_{\alpha,s,t}^{\rm add}}
 (D^{R_1},\ldots,D^{R_m},\Sigma_{\alpha,t}^{\rm add}).
 \label{eq:newV44-log-marked}
\end{equation}
\end{lemma}

\begin{proof}
Expand
\[
 \mathbb Z_{t,m}
 =
 I+\sum_{\varnothing\ne B\subseteq[m]}
 z_B M_{t,B}^A,
 \qquad
 z_B:=\prod_{i\in B}z_i,
\]
where \(M_{t,B}^A=\mathbb E\bigotimes_{i\in B}A_i\), canonically
embedded in the full tensor algebra.  In
\[
 \log(I+W)=\sum_{k\ge1}{(-1)^{k-1}\over k}W^k,
\]
the squarefree coefficient \(z_{[m]}\) uses ordered, pairwise
disjoint nonempty blocks whose union is \([m]\).  Every unordered
partition with \(k\) blocks occurs in \(k!\) orders.  Its block
operators have disjoint supports and commute, so its total
coefficient is
\[
 {(-1)^{k-1}\over k}k!
 =
 (-1)^{k-1}(k-1)!.
\]
This is moment--cumulant inversion.  For \(m\ge2\), cumulants are
unchanged by replacing any entry \(D^{R_i}\) by \(D^{R_i}-I\),
which proves \eqref{eq:newV44-log-cumulant}.

Differentiate the finite moments in \(t\) and use
\Cref{lem:newV41-additive-score}.  Differentiating the resulting
cumulant identity gives exactly one score insertion, proving
\eqref{eq:newV44-log-marked}.
\end{proof}

\section{A Banach-algebra zero-free polydisc}
\label{sec:newV44-polydisc}

\begin{lemma}[Geodesic matrix increment]
\label{lem:newV44-geodesic-increment}
There is a convention-dependent numerical \(C\) such that
\begin{equation}
 \|A_i(U)\|_{\rm op}
 \le
 C\theta\sqrt{C_2(R_i)}
 \le
 Cb_i{\theta\over\sqrt s}.
 \label{eq:newV44-geodesic-increment}
\end{equation}
\end{lemma}

\begin{proof}
Write \(U=\exp(\theta X)\) with \(X\) a unit Lie-algebra
direction in the fixed metric convention.  Then
\[
 D^{R_i}(U)-I
 =
 \int_0^\theta
 dD^{R_i}(X)e^{u\,dD^{R_i}(X)}\,du.
\]
The exponential is unitary on the real path and
\(\|dD^{R_i}(X)\|\le C\sqrt{C_2(R_i)}\).  This proves the first
bound.  The second follows from
\(sC_2(R_i)\le b_i^2\).
\end{proof}

\begin{lemma}[Uniform radial exponential moment]
\label{lem:newV44-radial-exp}
For
\[
 Y_s:={\theta\over\sqrt s},
\]
there are numerical \(\lambda_0,C>0\) such that
\begin{equation}
 \sup_{\alpha,t}
 \mathbb E_{\mu_{\alpha,s,t}^{\rm add}}
 e^{\lambda_0Y_s}
 \le C.
 \label{eq:newV44-radial-exp}
\end{equation}
\end{lemma}

\begin{proof}
The heat and Wilson endpoint cases of V30 give
\[
 \|Y_s\|_{L^q}\le C\sqrt q.
\]
Expanding \(e^{\lambda Y_s}\), using
\(k!\ge(k/e)^k\), and taking \(\lambda>0\) sufficiently small
gives a uniformly convergent series.  The additive expectation is
the convex combination of the endpoint expectations.
\end{proof}

\begin{theorem}[Operator source invertibility]
\label{thm:newV44-source-invertible}
There is a numerical \(\eta>0\) such that
\begin{equation}
 \sum_{i=1}^m|z_i|b_i\le\eta
 \label{eq:newV44-simplex}
\end{equation}
implies
\begin{equation}
 \|\mathbb Z_{t,m}(\boldsymbol z)-I\|_{\rm cb}
 \le {1\over4}.
 \label{eq:newV44-close-I}
\end{equation}
Consequently \(\mathbb Z_{t,m}\) is invertible and its principal
Banach-algebra logarithm obeys
\begin{equation}
 \|\log\mathbb Z_{t,m}(\boldsymbol z)\|_{\rm cb}\le C
 \label{eq:newV44-log-bound}
\end{equation}
throughout that weighted simplex, uniformly in \(m,t,\alpha\),
the thermal representations, and all spectator amplifications.
\end{theorem}

\begin{proof}
By \Cref{lem:newV44-geodesic-increment},
\begin{align*}
 \|\mathbb Z_{t,m}(\boldsymbol z)-I\|_{\rm cb}
 &\le
 \mathbb E
 \left[
 \prod_{i=1}^m
 (1+|z_i|\|A_i(U)\|)-1
 \right]\\
 &\le
 \mathbb E
 \left[
 \exp\left(
 CY_s\sum_i|z_i|b_i
 \right)-1
 \right].
\end{align*}
\Cref{lem:newV44-radial-exp} and dominated convergence at
\(\eta=0\) make the last expression at most \(1/4\) after
\(\eta\) is chosen sufficiently small.  The Neumann series gives
invertibility.  The normally convergent series for
\(\log(I+W)\) gives \eqref{eq:newV44-log-bound}.  Tensoring by an
identity spectator does not alter any norm.
\end{proof}

\section{The complete additive marked logarithm}
\label{sec:newV44-marked-log}

Let \(\mathbb Z_{H,m}\) and \(\mathbb Z_{W,m}\) denote the two
endpoint source operators.  Then
\[
 \mathbb Z_{t,m}
 =
 (1-t)\mathbb Z_{H,m}+t\mathbb Z_{W,m},
 \qquad
 \partial_t\mathbb Z_{t,m}
 =
 \mathbb Z_{W,m}-\mathbb Z_{H,m}.
\]

\begin{lemma}[Derivative of a Banach-algebra logarithm]
\label{lem:newV44-Dlog}
If \(\|W\|\le1/4\), then for every \(H\),
\begin{equation}
 \|D\log_{I+W}[H]\|
 \le {1\over1-\|W\|}\|H\|
 \le {4\over3}\|H\|.
 \label{eq:newV44-Dlog}
\end{equation}
\end{lemma}

\begin{proof}
Differentiate the absolutely convergent series for
\(\log(I+W)\):
\[
 D\log_{I+W}[H]
 =
 \sum_{k\ge1}{(-1)^{k-1}\over k}
 \sum_{j=0}^{k-1}W^jHW^{k-1-j}.
\]
The norm of the \(k\)-th line is at most
\(\|W\|^{k-1}\|H\|\).  Sum the geometric series.
\end{proof}

\begin{theorem}[Uniform complete marked-source bound]
\label{thm:newV44-marked-source}
On the weighted simplex \eqref{eq:newV44-simplex},
\begin{equation}
 \|\partial_t\log\mathbb Z_{t,m}(\boldsymbol z)\|_{\rm cb}
 \le C.
 \label{eq:newV44-marked-source}
\end{equation}
\end{theorem}

\begin{proof}
\Cref{thm:newV44-source-invertible} applied at the two endpoints
gives
\[
 \|\mathbb Z_{W,m}-\mathbb Z_{H,m}\|_{\rm cb}
 \le {1\over2}.
\]
Apply \Cref{lem:newV44-Dlog} at
\(W=\mathbb Z_{t,m}-I\).
\end{proof}

\section{Cauchy extraction and the exact binary card}
\label{sec:newV44-Cauchy}

\begin{theorem}[Crude all-arity operator source bound]
\label{thm:newV44-crude-all}
For every \(m\ge2\),
\begin{align}
 \left\|
 \kappa_m^{\,\mu_{\alpha,s,t}^{\rm add}}
 (D^{R_1},\ldots,D^{R_m})
 \right\|_{\rm cb}
 &\le
 (Km)^m\prod_{i=1}^m b_i,                                   \label{eq:newV44-crude-unmarked}\\
 \left\|
 \kappa_{m+1}^{\,\mu_{\alpha,s,t}^{\rm add}}
 (D^{R_1},\ldots,D^{R_m},\Sigma_{\alpha,t}^{\rm add})
 \right\|_{\rm cb}
 &\le
 (Km)^m\prod_{i=1}^m b_i.                                   \label{eq:newV44-crude-marked}
\end{align}
\end{theorem}

\begin{proof}
The polydisc
\[
 |z_i|\le r_i:={\eta\over mb_i}
\]
lies in \eqref{eq:newV44-simplex}.  Apply the Banach-valued
multivariate Cauchy estimate to
\(\log\mathbb Z_{t,m}\) and
\(\partial_t\log\mathbb Z_{t,m}\), using
\Cref{lem:newV44-log-cumulant,
thm:newV44-source-invertible,
thm:newV44-marked-source}.  Since every variable occurs once, no
multi-index factorial appears, and
\[
 \prod_i r_i^{-1}
 =
 \left({m\over\eta}\right)^m\prod_i b_i.
\]
\end{proof}

\begin{corollary}[Exact binary additive matrix cards]
\label{cor:newV44-binary}
For two arbitrary thermal representations,
\begin{align}
 \left\|
 \kappa_2^{\,\mu_{\alpha,s,t}^{\rm add}}
 (D^{R_1},D^{R_2})
 \right\|_{\rm cb}
 &\le K^2b_1b_2,                                             \label{eq:newV44-binary-unmarked}\\
 \left\|
 \kappa_3^{\,\mu_{\alpha,s,t}^{\rm add}}
 (D^{R_1},D^{R_2},\Sigma_{\alpha,t}^{\rm add})
 \right\|_{\rm cb}
 &\le K^2b_1b_2.                                             \label{eq:newV44-binary-marked}
\end{align}
\end{corollary}

\begin{proof}
Set \(m=2\) in
\eqref{eq:newV44-crude-unmarked} and
\eqref{eq:newV44-crude-marked}, and absorb the fixed factor
\((2K)^2\) into the displayed base.
\end{proof}

\begin{assessment}[The V42 binary clocks are recovered]
\label{ass:newV44-binary-recovered}
\Cref{cor:newV44-binary} supplies exactly the
\(b_1b_2\) factors missing from the contractive envelope estimate
\Cref{warn:newV42-binary-clocks}.  The gain comes from centering
each matrix input at the identity \emph{before} forming the
complete operator source, not from a small norm of the envelope
cocycle.
\end{assessment}

\section{The precise all-arity deficit}
\label{sec:newV44-deficit}

\begin{proposition}[Direct source invertibility is insufficient
beyond binary arity]
\label{prop:newV44-deficit}
The bound produced by the weighted-simplex argument is
\[
 (Km)^m\prod_i b_i,
\]
whereas the raw strong target is
\[
 K^m\left(\prod_i b_i\right)
 B^{m-2},
 \qquad B:=\sum_i b_i.
\]
For equal marks \(b_i=b\downarrow0\) at any fixed \(m\ge3\), the
ratio of the first scale to the second diverges at least as
\(b^{-(m-2)}\).  Therefore no change of numerical constants in
\Cref{thm:newV44-crude-all} proves the all-arity card.
\end{proposition}

\begin{proof}
For equal marks, the two scales, without their numerical
exponential bases, are
\[
 m^m b^m
 \quad\hbox{and}\quad
 b^m(mb)^{m-2}=m^{m-2}b^{2m-2}.
\]
Their ratio is \(m^2b^{-(m-2)}\), which is unbounded as
\(b\downarrow0\).
\end{proof}

\begin{assessment}[The missing operation is Gaussian
renormalization]
\label{ass:newV44-Gaussian}
The increment \(D^R(U)-I\) has a leading term linear in the scaled
Lie-algebra variable.  The direct norm proof treats all products
of those linear terms absolutely.  A Gaussian source has a
quadratic logarithm and hence zero connected coefficients above
arity two; the missing \(B^{m-2}\) is exactly the cancellation
discarded by the direct exponential-moment bound.
\end{assessment}

\begin{definition}[Gaussian-renormalized matrix source card]
\label{def:newV44-GRMSC}
Let \(\mathbb L_{t,m}^{[\le2]}(\boldsymbol z)\) be the degree-one
and degree-two Taylor polynomial of
\(\log\mathbb Z_{t,m}(\boldsymbol z)\), and put
\begin{equation}
 \mathbb R_{t,m}(\boldsymbol z)
 :=
 \log\mathbb Z_{t,m}(\boldsymbol z)
 -\mathbb L_{t,m}^{[\le2]}(\boldsymbol z).
 \label{eq:newV44-renormalized-source}
\end{equation}
GRMSC is a complete, spectator-stable analytic bound on
\(\mathbb R_{t,m}\) and \(\partial_t\mathbb R_{t,m}\) which,
after squarefree Cauchy extraction, yields
\[
 K^m\left(\prod_i b_i\right)B^{m-2}.
\]
The subtraction is performed in the logarithm; no noncommuting
factorization of exponentials is assumed.
\end{definition}

\begin{proposition}[GRMSC would close the local card]
\label{prop:newV44-GRMSC-sufficient}
GRMSC implies the raw Wilson aggregate strong card.
\end{proposition}

\begin{proof}
For \(m\ge3\), subtracting Taylor degrees one and two does not
change the squarefree coefficient \(z_1\cdots z_m\).  GRMSC
therefore gives the desired unmarked and marked additive-bridge
cumulant bounds by
\Cref{lem:newV44-log-cumulant}.  At \(m=2\), use
\Cref{cor:newV44-binary}.  Integrate the exact marked identity in
\(t\), use the heat endpoint card of V29, and apply the V16
balancing equivalence.
\end{proof}

\begin{programme}[V45: audit and Gaussian-renormalized source]
\label{prog:newV44-V45}
\begin{enumerate}[leftmargin=3em]
\item Perform the compulsory read-only audit of the newest active
      SM--Einstein and Navier--Stokes notes.
\item Compute the degree-two operator source logarithm exactly and
      compare it with the total-Casimir heat covariance.
\item Write a Duhamel formula for the remainder
      \(\mathbb R_{t,m}\) which keeps the quadratic source
      subtraction inside the complete endpoint expectation.
\item Test the formula first at arity three and four, where the
      tangent Gaussian cancellation is explicit.
\item Do not expand the logarithm into ordered words before the
      quadratic cancellation has been applied.
\end{enumerate}
\end{programme}

\begin{firewall}[Claim boundary after V44]
\label{fire:newV44-boundary}
V44 constructs the operator-valued squarefree source, proves that
its Banach-algebra logarithm generates the exact tensor cumulants
and score cumulants, proves a uniform weighted-simplex
invertibility domain, and closes the full binary additive
unmarked and marked matrix cards.  It also proves that this direct
norm argument misses \(m-2\) thermal factors at higher arity and
isolates GRMSC.  It does not prove GRMSC, EMCC, thermal connected
WPRAC, all-arity RGSC, the raw Wilson aggregate strong card,
all-order strong MTA, WUFC, CPM, cutoff removal, reflection
positivity, Osterwalder--Schrader reconstruction, a continuum
Yang--Mills measure, or a positive Hamiltonian gap.  The full
Yang--Mills mass gap has not been proved.
\end{firewall}

\begin{theorem}[V44 result ledger]
\label{thm:newV44-results}
V44:
\begin{enumerate}[label=\textup{(V44.\arabic*)},leftmargin=6.5em]
\item proves the squarefree Banach-logarithm tensor-cumulant
      identity;
\item proves a uniform weighted-simplex invertibility theorem for
      the complete operator source;
\item proves a uniform bound on its exact additive marked
      logarithm;
\item derives the resulting crude all-arity operator bounds;
\item closes the complete binary additive matrix cards with the
      correct \(b_1b_2\) clocks;
\item proves the direct source argument loses exactly the
      higher-arity Gaussian cancellation; and
\item isolates GRMSC as the next sufficient gate.
\end{enumerate}
\end{theorem}

\begin{proof}
Use
\Cref{lem:newV44-log-cumulant,
lem:newV44-geodesic-increment,
lem:newV44-radial-exp,
thm:newV44-source-invertible,
lem:newV44-Dlog,
thm:newV44-marked-source,
thm:newV44-crude-all,
cor:newV44-binary,
prop:newV44-deficit,
prop:newV44-GRMSC-sufficient}.
\end{proof}

\paragraph{Parent-version record.}
V44 was formed from
\texttt{Renewal\_Yang\_Mills\_Mass\_Gap\_Research\_Notes\_V43.tex}.
The V43 parent had \(20\,721\) logical source lines,
\(654\,906\) bytes, and SHA--256
\[
 \texttt{459244dd1ea06d060f19b1348bb2bcb586d5f5fde3ca3ef50a6539fe28ecbc50}.
\]
The next compulsory companion audit is V45.

% =============================================================================
% END APPENDED FIFTH-SERIES V44 MATERIAL
% =============================================================================

% =============================================================================
% BEGIN APPENDED FIFTH-SERIES V45 MATERIAL
% =============================================================================

\chapter{V45: Companion audit, exact quadratic source, and the
ternary matrix card}
\label{chap:newV45-ternary}

\section{Compulsory companion audit}
\label{sec:newV45-audit}

\begin{audit}[Read-only SM--Einstein and Navier--Stokes audit]
\label{audit:newV45-companions}
The newest active companion notes inspected before this append
were:
\begin{align*}
 &\texttt{Renewal\_SM\_Einstein\_Continuum\_Research\_Notes\_V64.tex},
 \\
 &\qquad 20\,726\ \hbox{logical source lines},\quad
 796\,174\ \hbox{bytes},                                      \\
 &\qquad
 \operatorname{SHA256}
 =
 \texttt{d1401dc4b3e5483814485bd367d4fcc07ef052a8a489c9757ae1bc68972e0829},
 \\
 &\texttt{Renewal\_NS\_Stability\_Research\_Notes\_V42.tex},
 \\
 &\qquad 28\,013\ \hbox{logical source lines},\quad
 1\,122\,447\ \hbox{bytes},                                   \\
 &\qquad
 \operatorname{SHA256}
 =
 \texttt{cfe70f62c7d20cf8d0f3b0d0144899f8deb082bf2689f15a46ea09647e4fc1b5}.
\end{align*}
Neither file was modified.

SM V64 proves two distinct multiscale facts.  Local dyadic
neutral details admit local fluxes but lose a factor equal to the
number of scales unless a stable multilevel flux comparison is
proved.  Spectral Newton shells, by contrast, form an exact
positive ledger; an order-one critical contribution on every
shell produces an unavoidable logarithm.  Once the ledger is
positive, signed-shell cancellation is unavailable.

NS V42 reduces the type-I problem to a normalized ancient class
and proves several uniform bounds, but also proves that the class
contains the finite-Dirichlet stationary Liouville problem of
Leray.  Thus a theorem for the entire abstract class would solve a
known hard subproblem; the remaining opportunity must use the
specific ancestry and flat profile of limits actually generated by
the blow-up construction.
\end{audit}

\begin{assessment}[Audit transfer discipline]
\label{ass:newV45-audit-transfer}
No Newton-shell estimate and no Navier--Stokes Liouville estimate
is imported into Yang--Mills.  The legitimate consequences for
the V44 programme are:
\begin{enumerate}[label=\textup{(\roman*)},leftmargin=3.5em]
\item a positive quadratic source ledger must be paid, not
      rhetorically canceled; and
\item GRMSC should first be proved for the exact heat/Wilson
      sources, not for an overbroad class of all sub-Gaussian
      central measures.
\end{enumerate}
\end{assessment}

\section{The exact degree-two operator logarithm}
\label{sec:newV45-quadratic}

For disjoint inputs write
\[
 C_i=C_2(R_i)I,
 \qquad
 \Omega_{ij}=\sum_aJ_i^aJ_j^a.
\]

\begin{lemma}[Exact heat covariance coefficient]
\label{lem:newV45-heat-covariance}
The coefficient of \(z_iz_j\) in the heat-source logarithm is
\begin{equation}
 \kappa_2^H(D^{R_i},D^{R_j})
 =
 e^{-s(C_i+C_j)}
 \left(e^{-2s\Omega_{ij}}-I\right).
 \label{eq:newV45-heat-covariance}
\end{equation}
For thermal inputs,
\begin{equation}
 \|\kappa_2^H(D^{R_i},D^{R_j})\|_{\rm cb}
 \le Cb_ib_j.
 \label{eq:newV45-heat-covariance-bound}
\end{equation}
\end{lemma}

\begin{proof}
V29 gives
\[
 M_{ij}^H
 =
 e^{-s(C_i+C_j)}e^{-2s\Omega_{ij}},
 \qquad
 M_i^HM_j^H=e^{-s(C_i+C_j)}.
\]
Subtract the product of the singleton moments.  For the norm,
\[
 \|e^{-2s\Omega_{ij}}-I\|
 \le
 2s\|\Omega_{ij}\|e^{2s\|\Omega_{ij}\|}.
\]
The row Cauchy--Schwarz inequality gives
\[
 s\|\Omega_{ij}\|
 \le
 \sqrt{sC_2(R_i)}\sqrt{sC_2(R_j)}
 \le b_ib_j<1.
\]
The scalar exterior heat factor is contractive, proving the bound.
\end{proof}

\begin{lemma}[Exact additive covariance ledger]
\label{lem:newV45-additive-covariance}
Let
\[
 \Delta M_i:=M_i^W-M_i^H.
\]
Then
\begin{equation}
 \kappa_{2,t}^{\rm add}
 =
 (1-t)\kappa_2^H+t\kappa_2^W
 +t(1-t)\Delta M_i\Delta M_j.
 \label{eq:newV45-additive-covariance}
\end{equation}
Consequently
\begin{equation}
 \partial_t\kappa_{2,t}^{\rm add}
 =
 \kappa_2^W-\kappa_2^H
 +(1-2t)\Delta M_i\Delta M_j.
 \label{eq:newV45-additive-covariance-marked}
\end{equation}
\end{lemma}

\begin{proof}
Every additive moment is affine:
\[
 M_{B,t}^{\rm add}=(1-t)M_B^H+tM_B^W.
\]
Insert this for \(B=\{i,j\},\{i\},\{j\}\) into
\(\kappa_2=M_{ij}-M_iM_j\) and expand.  Differentiation gives the
second identity.
\end{proof}

\begin{assessment}[The quadratic term is a paid ledger]
\label{ass:newV45-quadratic-ledger}
The degree-two term subtracted in GRMSC is not an error.  At the
heat endpoint it is exactly the pair Hamiltonian
\eqref{eq:newV45-heat-covariance}; along the additive bridge it is
the exact positive-mixture covariance ledger
\eqref{eq:newV45-additive-covariance}.  V44 already pays it at
scale \(b_ib_j\).  GRMSC asks only for the genuinely connected
degrees three and higher.
\end{assessment}

\section{Odd tangent cancellation at arity three}
\label{sec:newV45-parity}

Except on the Haar-null antipode, write
\[
 U=\exp(\theta X_U),
\]
where \(X_U\) is a unit Lie-algebra direction and inversion sends
\((\theta,X_U)\) to \((\theta,-X_U)\).  Put
\begin{align}
 L_i(U)&:=\theta\,dD^{R_i}(X_U),                              \label{eq:newV45-Li}\\
 Q_i(U)&:=D^{R_i}(U)-I-L_i(U).                               \label{eq:newV45-Qi}
\end{align}

\begin{lemma}[Linear and quadratic increment rows]
\label{lem:newV45-LQ-bounds}
With \(Y_s=\theta/\sqrt s\),
\begin{align}
 \|L_i(U)\|
 &\le Cb_iY_s,                                                \label{eq:newV45-L-bound}\\
 \|Q_i(U)\|
 &\le Cb_i^2Y_s^2.                                           \label{eq:newV45-Q-bound}
\end{align}
For every heat, Wilson, or additive-bridge law,
\begin{equation}
 \mathbb E L_i=0,
 \qquad
 \|Q_i-\mathbb EQ_i\|
 \le Cb_i^2(1+Y_s^2).
 \label{eq:newV45-centered-Q}
\end{equation}
\end{lemma}

\begin{proof}
The generator bound in
\Cref{lem:newV44-geodesic-increment} gives
\eqref{eq:newV45-L-bound}.  The integral Taylor remainder
\[
 e^{\theta A}-I-\theta A
 =
 \int_0^\theta(\theta-u)A^2e^{uA}\,du
\]
has norm at most
\(\theta^2\|A\|^2/2\) for skew-Hermitian \(A\).  Since
\(sC_2(R_i)\le b_i^2\), this proves
\eqref{eq:newV45-Q-bound}.

Every endpoint density, and hence every additive mixture, is
invariant under \(U\mapsto U^{-1}\).  The linear term changes sign,
so its expectation vanishes.  Finally,
\[
 \|\mathbb EQ_i\|
 \le Cb_i^2\mathbb EY_s^2\le Cb_i^2
\]
by V30 concentration, proving
\eqref{eq:newV45-centered-Q}.
\end{proof}

\begin{lemma}[The cubic linear tensor vanishes]
\label{lem:newV45-LLL-zero}
For every additive-bridge parameter \(t\),
\begin{equation}
 \mathbb E_t[L_1(U)\otimes L_2(U)\otimes L_3(U)]=0.
 \label{eq:newV45-LLL-zero}
\end{equation}
\end{lemma}

\begin{proof}
The tensor product changes sign under \(U\mapsto U^{-1}\), while
the law and Haar measure are invariant.
\end{proof}

\begin{theorem}[Uniform ternary unmarked matrix card]
\label{thm:newV45-ternary-unmarked}
For arbitrary thermal \(R_1,R_2,R_3\),
\begin{equation}
 \left\|
 \kappa_3^{\,\mu_{\alpha,s,t}^{\rm add}}
 (D^{R_1},D^{R_2},D^{R_3})
 \right\|_{\rm cb}
 \le
 K^3b_1b_2b_3(b_1+b_2+b_3).
 \label{eq:newV45-ternary-unmarked}
\end{equation}
\end{theorem}

\begin{proof}
Because the entries occupy distinct tensor legs, the third
cumulant is the centered product:
\[
 \kappa_3(D^{R_1},D^{R_2},D^{R_3})
 =
 \mathbb E
 \bigotimes_{i=1}^3
 \left(D^{R_i}-\mathbb ED^{R_i}\right).
\]
By \Cref{lem:newV45-LQ-bounds},
\[
 D^{R_i}-\mathbb ED^{R_i}
 =
 L_i+(Q_i-\mathbb EQ_i).
\]
The term with three \(L_i\)'s vanishes by
\Cref{lem:newV45-LLL-zero}.  Every remaining one of the seven
terms contains a nonempty set \(S\) of centered \(Q\)'s.  Its norm
is bounded by a fixed moment of \(Y_s\) times
\[
 \left(\prod_{i=1}^3b_i\right)\prod_{i\in S}b_i.
\]
All moments involved have order at most six and are uniformly
bounded by V30.  Since \(0<b_i<1\),
\[
 \sum_{\varnothing\ne S\subseteq\{1,2,3\}}
 \prod_{i\in S}b_i
 \le C(b_1+b_2+b_3).
\]
This proves the operator bound.  The proof is unchanged after
spectator amplification.
\end{proof}

\section{The marked ternary card by polynomial interpolation}
\label{sec:newV45-marked}

\begin{lemma}[Banach-valued Markov inequality]
\label{lem:newV45-Markov}
If \(P:[0,1]\to\mathcal B\) is a polynomial of degree at most
\(d\) with values in a Banach space, then
\begin{equation}
 \sup_{0\le t\le1}\|P'(t)\|
 \le
 2d^2\sup_{0\le t\le1}\|P(t)\|.
 \label{eq:newV45-Markov}
\end{equation}
\end{lemma}

\begin{proof}
Fix \(t\) and a norm-one functional
\(\ell\in\mathcal B^*\) which norms \(P'(t)\) up to an arbitrary
\(\varepsilon>0\).  The scalar polynomial
\(\ell(P)\), transported from \([0,1]\) to \([-1,1]\), obeys the
classical Markov inequality.  Undoing the affine change of
variables gives the factor \(2d^2\).  Let
\(\varepsilon\downarrow0\), then take the supremum in \(t\).
\end{proof}

\begin{theorem}[Uniform ternary additive score card]
\label{thm:newV45-ternary-marked}
For arbitrary thermal \(R_1,R_2,R_3\),
\begin{align}
 &\left\|
 \kappa_4^{\,\mu_{\alpha,s,t}^{\rm add}}
 (D^{R_1},D^{R_2},D^{R_3},\Sigma_{\alpha,t}^{\rm add})
 \right\|_{\rm cb}
 \notag\\
 &\qquad\le
 K^3b_1b_2b_3(b_1+b_2+b_3).
 \label{eq:newV45-ternary-marked}
\end{align}
\end{theorem}

\begin{proof}
Each moment of the additive law is affine in \(t\).
Moment--cumulant inversion for three inputs shows that
\[
 P(t):=
 \kappa_3^{\,\mu_{\alpha,s,t}^{\rm add}}
 (D^{R_1},D^{R_2},D^{R_3})
\]
is a Banach-valued polynomial of degree at most three.
\Cref{thm:newV45-ternary-unmarked} bounds it uniformly.
Apply \Cref{lem:newV45-Markov} with \(d=3\).
The exact score derivative identity gives
\[
 P'(t)
 =
 \kappa_4^{\,\mu_{\alpha,s,t}^{\rm add}}
 (D^{R_1},D^{R_2},D^{R_3},\Sigma_{\alpha,t}^{\rm add}).
\]
Absorb the fixed Markov factor into \(K^3\).
\end{proof}

\begin{corollary}[Additive local card through arity three]
\label{cor:newV45-through-three}
The full operator-valued unmarked and marked additive-bridge raw
cards hold for \(m=2\) and \(m=3\), uniformly on all thermal
representations and all matrix amplifications.
\end{corollary}

\begin{proof}
Use \Cref{cor:newV44-binary} for \(m=2\) and
\Cref{thm:newV45-ternary-unmarked,
thm:newV45-ternary-marked} for \(m=3\).
\end{proof}

\begin{assessment}[Why the companion audit changes the next step]
\label{ass:newV45-next-direction}
The ternary gain is an exact signed symmetry cancellation before
norms.  At arity four, however, the linear tangent tensor is even
and its fourth connected component is a covariance-type ledger.
The SM audit warns that an actually positive ledger cannot be
discarded.  The correct V46 task is therefore to compute the
fourth connected tangent tensor and show that Gaussian Wick
subtraction, together with the exact heat/Wilson calibration,
pays two additional thermal marks.  A generic central
sub-Gaussian hypothesis is not enough.
\end{assessment}

\begin{programme}[V46: quartic connected source]
\label{prog:newV45-V46}
\begin{enumerate}[leftmargin=3em]
\item Expand each matrix increment through quadratic Lie-algebra
      order with a cubic remainder controlled uniformly on
      thermal representations.
\item Compute the complete fourth cumulant of the four linear
      terms.  Subtract the three Wick pairings before taking an
      operator norm.
\item Use the V37 score-weighted endpoint comparison, or a direct
      calibrated density expansion, to prove that the residual
      fourth tangent cumulant carries an extra factor \(s\).
\item Pay that factor by
      \(s\le(\sum_{i=1}^4b_i)^2/16\), and show that every Taylor
      remainder term carries at least the same two-mark gain after
      inversion parity is used.
\item Differentiate the resulting degree-four additive polynomial
      in \(t\) only after the unmarked bound is proved.
\end{enumerate}
\end{programme}

\begin{firewall}[Claim boundary after V45]
\label{fire:newV45-boundary}
V45 completes the compulsory companion audit, identifies the exact
quadratic heat and additive covariance ledgers, proves the odd
linear tangent cancellation, and closes the complete ternary
unmarked and marked additive matrix cards.  Together with V44 this
settles the additive local card through arity three.  It does not
prove the quartic card, GRMSC, EMCC, thermal connected WPRAC,
all-arity RGSC, the raw Wilson aggregate strong card, all-order
strong MTA, WUFC, CPM, cutoff removal, reflection positivity,
Osterwalder--Schrader reconstruction, a continuum Yang--Mills
measure, or a positive Hamiltonian gap.  The full Yang--Mills mass
gap has not been proved.
\end{firewall}

\begin{theorem}[V45 result ledger]
\label{thm:newV45-results}
V45:
\begin{enumerate}[label=\textup{(V45.\arabic*)},leftmargin=6.5em]
\item completes the required read-only V45 companion audit;
\item computes the exact quadratic heat-source logarithm;
\item derives the exact additive covariance and marked-covariance
      ledgers;
\item proves uniform linear and quadratic matrix increment rows;
\item proves inversion cancellation of the cubic linear tensor;
\item closes the ternary unmarked matrix card;
\item closes the ternary marked card by Banach-valued polynomial
      interpolation; and
\item fixes the quartic Wick residual as the next acquisition.
\end{enumerate}
\end{theorem}

\begin{proof}
Use
\Cref{audit:newV45-companions,
lem:newV45-heat-covariance,
lem:newV45-additive-covariance,
lem:newV45-LQ-bounds,
lem:newV45-LLL-zero,
thm:newV45-ternary-unmarked,
lem:newV45-Markov,
thm:newV45-ternary-marked,
cor:newV45-through-three}.
\end{proof}

\paragraph{Parent-version record.}
V45 was formed from
\texttt{Renewal\_Yang\_Mills\_Mass\_Gap\_Research\_Notes\_V44.tex}.
The V44 parent had \(21\,215\) logical source lines,
\(669\,788\) bytes, and SHA--256
\[
 \texttt{203ae2186b394fa08726969c85d96fde2e8497e6eb72126d7e3724129c9cdf9a}.
\]
The next compulsory companion audit is V50.

% =============================================================================
% END APPENDED FIFTH-SERIES V45 MATERIAL
% =============================================================================

% =============================================================================
% BEGIN APPENDED FIFTH-SERIES V46 MATERIAL
% =============================================================================

\chapter{V46: Quartic Wick residual and the four-input matrix card}
\label{chap:newV46-quartic}

\section{Quantitative Gaussian fourth-moment comparison}
\label{sec:newV46-Gaussian}

No companion audit is due at V46.  Use the scaled radial vector
\begin{equation}
 Y_s:={\theta\boldsymbol n\over\sqrt{2s}}\in\mathbb R^3,
 \label{eq:newV46-Y}
\end{equation}
where \(\boldsymbol n\) is the axis of \(U\).  Let \(\gamma_3\)
denote the standard centered Gaussian law on \(\mathbb R^3\).

\begin{lemma}[Endpoint moment comparison through degree four]
\label{lem:newV46-moment-comparison}
For every polynomial \(P\) on \(\mathbb R^3\) of degree at most
four,
\begin{equation}
 \left|
 \mathbb E_{\mu_{\alpha,s,t}^{\rm add}}P(Y_s)
 -
 \mathbb E_{\gamma_3}P(Y)
 \right|
 \le
 C_Ps
 \label{eq:newV46-moment-comparison}
\end{equation}
uniformly in \(t\in[0,1]\), where
\(s=\tau_\alpha^\sharp\).
\end{lemma}

\begin{proof}
It suffices to prove the two endpoint bounds, because the additive
expectation is their convex combination.  Split the class-angle
integral at
\[
 |Y_s|=s^{-1/8}.
\]
V30 concentration makes the complement
\(O(e^{-cs^{-1/4}})\) in every fixed polynomial moment.

For the heat endpoint, the \(k=0\) image in
\eqref{eq:newV30-image}, multiplied by radial Haar measure, has
the scaled density
\[
 C_s e^{-|Y_s|^2/2}
 {\sin(\sqrt{2s}|Y_s|)\over\sqrt{2s}|Y_s|}
 \left(1+O(e^{-c/s})\right).
\]
On the inner region the non-Gaussian factor is
\[
 1+O\left(s|Y_s|^2\right),
\]
and its normalization has the same expansion.  Gaussian
domination permits termwise integration and gives an \(O(s)\)
error for every polynomial of degree at most four.

For the Wilson endpoint,
\(\alpha=(2s)^{-1}-1/3\), and
\[
 \alpha(\cos(\sqrt{2s}|Y_s|)-1)
 =
 -{|Y_s|^2\over2}
 +O\left(s(|Y_s|^2+|Y_s|^4)\right).
\]
The radial Haar factor is
\[
 \left(
 {\sin(\sqrt{2s}|Y_s|)\over\sqrt{2s}|Y_s|}
 \right)^2
 =
 1+O(s|Y_s|^2).
\]
Taylor expansion under Gaussian domination, followed by
normalization, again gives \(O(s)\).  This proves the endpoint and
mixture claims.
\end{proof}

\begin{corollary}[Fourth cumulant of the scaled tangent vector]
\label{cor:newV46-Y-cumulant}
For all coordinate indices \(a,b,c,d\in\{1,2,3\}\),
\begin{equation}
 \left|
 \kappa_4^{\,\mu_{\alpha,s,t}^{\rm add}}
 (Y_s^a,Y_s^b,Y_s^c,Y_s^d)
 \right|
 \le Cs.
 \label{eq:newV46-Y-cumulant}
\end{equation}
\end{corollary}

\begin{proof}
Inversion symmetry makes every first and third coordinate moment
zero.  The fourth cumulant is the fourth moment minus the three
products of second moments.  By
\Cref{lem:newV46-moment-comparison}, each moment differs by
\(O(s)\) from its Gaussian value.  The Gaussian fourth moment is
exactly the sum of the three Wick pairings, so its fourth cumulant
vanishes.
\end{proof}

\section{Linear, quadratic, and cubic representation pieces}
\label{sec:newV46-LHT}

Retain the principal-geodesic generator
\[
 G_i(U):=dD^{R_i}(X_U)
\]
from V45 and define
\begin{align}
 L_i(U)&:=\theta G_i(U),                                      \label{eq:newV46-L}\\
 H_i(U)&:={\theta^2\over2}G_i(U)^2,                           \label{eq:newV46-H}\\
 T_i(U)&:=D^{R_i}(U)-I-L_i(U)-H_i(U).                         \label{eq:newV46-T}
\end{align}

\begin{lemma}[Three-level Taylor rows]
\label{lem:newV46-LHT-bounds}
With \(Y_s=\theta/\sqrt{2s}\) in norm notation,
\begin{align}
 \|L_i\|&\le Cb_iY_s,                                        \label{eq:newV46-L-bound}\\
 \|H_i\|&\le Cb_i^2Y_s^2,                                    \label{eq:newV46-H-bound}\\
 \|T_i\|&\le Cb_i^3Y_s^3.                                    \label{eq:newV46-T-bound}
\end{align}
Moreover \(L_i\) is odd and \(H_i\) is even under
\(U\mapsto U^{-1}\).
\end{lemma}

\begin{proof}
The first two bounds follow from
\(\|G_i\|\le C\sqrt{C_2(R_i)}\) and
\(sC_2(R_i)\le b_i^2\).  The integral third-order remainder
\[
 e^{\theta G}-I-\theta G-{\theta^2\over2}G^2
 =
 \int_0^\theta{(\theta-u)^2\over2}G^3e^{uG}\,du
\]
has norm at most
\(\theta^3\|G\|^3/6\), proving the third bound.
Inversion changes \(G_i\) to \(-G_i\) and leaves \(\theta\)
fixed, proving the parity statements.
\end{proof}

\begin{lemma}[Centered nonlinear rows]
\label{lem:newV46-centered-HT}
Put
\[
 \widetilde H_i:=H_i-\mathbb EH_i,
 \qquad
 \widetilde T_i:=T_i-\mathbb ET_i.
\]
Then
\begin{align}
 \|\widetilde H_i\|
 &\le Cb_i^2(1+Y_s^2),                                      \label{eq:newV46-H-centered}\\
 \|\widetilde T_i\|
 &\le Cb_i^3(1+Y_s^3).                                      \label{eq:newV46-T-centered}
\end{align}
The variable \(\widetilde H_i\) is even under inversion.
\end{lemma}

\begin{proof}
Use \Cref{lem:newV46-LHT-bounds} and the uniform fixed moments of
\(Y_s\).  The expectation of an even function is constant and
hence remains even.
\end{proof}

\section{The complete four-linear Wick residual}
\label{sec:newV46-linear-four}

\begin{theorem}[Four-linear connected tensor]
\label{thm:newV46-four-linear}
For four arbitrary thermal representations,
\begin{equation}
 \left\|
 \kappa_4^{\,\mu_{\alpha,s,t}^{\rm add}}
 (L_1,L_2,L_3,L_4)
 \right\|_{\rm cb}
 \le
 Cs\prod_{i=1}^4b_i.
 \label{eq:newV46-four-linear}
\end{equation}
\end{theorem}

\begin{proof}
From \eqref{eq:newV46-Y},
\[
 L_i
 =
 \sqrt{2s}\sum_{a=1}^3Y_s^aJ_i^a
\]
up to the fixed metric normalization already absorbed in \(C\).
Multilinearity gives
\begin{align*}
 \kappa_4(L_1,L_2,L_3,L_4)
 &=
 (2s)^2
 \sum_{a,b,c,d}
 \kappa_4(Y_s^a,Y_s^b,Y_s^c,Y_s^d)
 J_1^aJ_2^bJ_3^cJ_4^d.
\end{align*}
Use \Cref{cor:newV46-Y-cumulant} and
\(\|J_i^a\|\le C\sqrt{C_2(R_i)}\).  The exterior generator scale
satisfies
\[
 s^2\prod_i\sqrt{C_2(R_i)}
 =
 \prod_i\sqrt{sC_2(R_i)}
 \le\prod_i b_i.
\]
The coordinate sum has fixed cardinality.  Amplification
preserves the estimate.
\end{proof}

\section{Every non-linear term pays two further marks}
\label{sec:newV46-nonlinear}

\begin{lemma}[Odd \(HLLL\) cumulant vanishes]
\label{lem:newV46-HLLL-zero}
For any placement of one centered quadratic term and three linear
terms,
\begin{equation}
 \kappa_4(\widetilde H_i,L_j,L_k,L_\ell)=0.
 \label{eq:newV46-HLLL-zero}
\end{equation}
\end{lemma}

\begin{proof}
Expand the cumulant over set partitions.  The total number of odd
linear factors is three.  In every partition, at least one block
contains an odd number of \(L\)'s.  Its moment changes sign under
inversion and is zero.  Therefore every partition product is
zero.
\end{proof}

\begin{lemma}[Fixed fourth-cumulant product bound]
\label{lem:newV46-fixed-cumulant}
Suppose \(X_i\) are centered random operators and
\[
 \|X_i\|\le a_i(1+Y_s^{r_i})
\]
with fixed \(r_i\le3\).  Then
\begin{equation}
 \|\kappa_4(X_1,X_2,X_3,X_4)\|_{\rm cb}
 \le C_{\boldsymbol r}\prod_i a_i.
 \label{eq:newV46-fixed-cumulant}
\end{equation}
\end{lemma}

\begin{proof}
Moment--cumulant inversion contains fifteen set partitions.
Apply the operator triangle inequality and H\"older to every block
moment.  V30 supplies all fixed moments of \(Y_s\) needed here.
The same proof holds after amplification.
\end{proof}

\begin{theorem}[Quartic unmarked additive matrix card]
\label{thm:newV46-quartic-unmarked}
For arbitrary thermal \(R_1,\ldots,R_4\), with
\(B=\sum_{i=1}^4b_i\),
\begin{equation}
 \left\|
 \kappa_4^{\,\mu_{\alpha,s,t}^{\rm add}}
 (D^{R_1},D^{R_2},D^{R_3},D^{R_4})
 \right\|_{\rm cb}
 \le
 K^4\left(\prod_{i=1}^4b_i\right)B^2.
 \label{eq:newV46-quartic-unmarked}
\end{equation}
\end{theorem}

\begin{proof}
Cumulants of order four are invariant under adding constants, and
\[
 D^{R_i}-\mathbb ED^{R_i}
 =
 L_i+\widetilde H_i+\widetilde T_i.
\]
Expand the cumulant multilinearly.

The all-linear term is bounded by
\Cref{thm:newV46-four-linear}.  Since
\(b_i^2=s(1+C_2(R_i))\ge s\),
\[
 B\ge4\sqrt s,
 \qquad
 s\le {B^2\over16}.
\]
Thus this term has the required bound.

The terms with exactly one \(\widetilde H\) and three \(L\)'s
vanish by \Cref{lem:newV46-HLLL-zero}.  A term with exactly one
\(\widetilde T_i\) and three \(L\)'s carries the mark factor
\[
 \left(\prod_{j=1}^4b_j\right)b_i^2
 \le
 \left(\prod_jb_j\right)B^2.
\]
Every term with at least two nonlinear entries carries, beyond
\(\prod_jb_j\), either two distinct additional \(b\)'s or at least
two additional powers of one \(b_i\).  Because every \(b_i<1\),
each such extra monomial is at most \(B^2\).  Apply
\Cref{lem:newV46-fixed-cumulant} to all these finitely many terms
and enlarge \(K\).
\end{proof}

\section{The marked quartic card}
\label{sec:newV46-marked}

\begin{theorem}[Quartic additive score card]
\label{thm:newV46-quartic-marked}
Uniformly in \(t\),
\begin{align}
 &\left\|
 \kappa_5^{\,\mu_{\alpha,s,t}^{\rm add}}
 (D^{R_1},D^{R_2},D^{R_3},D^{R_4},
  \Sigma_{\alpha,t}^{\rm add})
 \right\|_{\rm cb}
 \notag\\
 &\qquad\le
 K^4\left(\prod_{i=1}^4b_i\right)
 \left(\sum_{i=1}^4b_i\right)^2.
 \label{eq:newV46-quartic-marked}
\end{align}
\end{theorem}

\begin{proof}
The unmarked fourth cumulant along the additive bridge is a
Banach-valued polynomial in \(t\) of degree at most four.  Apply
the Banach-valued Markov inequality
\Cref{lem:newV45-Markov} with \(d=4\) to
\Cref{thm:newV46-quartic-unmarked}.  Its \(t\)-derivative is the
displayed score cumulant by the exact additive score identity.
Absorb the fixed factor \(2d^2=32\) into \(K^4\).
\end{proof}

\begin{corollary}[Additive matrix local card through arity four]
\label{cor:newV46-through-four}
The complete unmarked and marked additive-bridge matrix cards hold
for every \(2\le m\le4\), uniformly on thermal representations and
all amplifications.
\end{corollary}

\begin{proof}
Combine
\Cref{cor:newV44-binary,
cor:newV45-through-three,
thm:newV46-quartic-unmarked,
thm:newV46-quartic-marked}.
\end{proof}

\section{The all-arity pattern exposed by the quartic proof}
\label{sec:newV46-pattern}

\begin{assessment}[Required tangent cumulant order]
\label{ass:newV46-tangent-order}
For \(m>4\), an \(O(s)\) comparison with a Gaussian is not enough.
The target requires the scaled linear tangent cumulant to obey
\begin{equation}
 \left\|
 \kappa_m(Y_s,\ldots,Y_s)
 \right\|
 \le
 K^m
 s^{\lceil(m-2)/2\rceil}
 \times\hbox{the appropriate connected graph census}.
 \label{eq:newV46-needed-order}
\end{equation}
This is exactly the quartic-root graph order found in V25--V27:
one quartic perturbation can connect at most four linear leaves,
and further leaves require further perturbation vertices.
\end{assessment}

\begin{definition}[Uniform tangent connected-source card]
\label{def:newV46-UTCSC}
UTCSC is an exact analytic representation of the complete
heat/Wilson/additive scaled source logarithm whose degree-\(m\)
coefficient is a sum over connected graphs with:
\begin{enumerate}[label=\textup{(\roman*)},leftmargin=3.5em]
\item ordinary covariance edges paid by \(b_ib_j\);
\item non-Gaussian root vertices carrying at least one factor
      \(s\) and degree at least four; and
\item total weighted graph sum bounded by
      \(K^m(\prod_i b_i)B^{m-2}\).
\end{enumerate}
The graph expansion must be resummed before operator norms at
unbounded \(m\).
\end{definition}

\begin{programme}[V47: tangent connected-source resummation]
\label{prog:newV46-V47}
\begin{enumerate}[leftmargin=3em]
\item Derive the exact scaled radial density ratio of each endpoint
      to the standard Gaussian, retaining the complete wrapped
      image factor.
\item Introduce one scalar source for that density ratio and matrix
      sources for the representation increments; take the
      logarithm before expanding either source.
\item Recover the V25 quartic-root graph as the first derivative of
      the density-ratio source and identify the exact higher root
      degrees.
\item Seek an exponential-grade bound for the complete connected
      root generating function.  Fixed-degree Taylor comparison is
      insufficient at growing arity.
\item Preserve the additive polynomial interpolation so that an
      unmarked all-arity bound automatically yields the marked one.
\end{enumerate}
\end{programme}

\begin{firewall}[Claim boundary after V46]
\label{fire:newV46-boundary}
V46 proves the calibrated endpoint fourth-moment comparison,
computes the four-linear Wick residual, proves the parity
cancellation of the only one-extra-mark quadratic term, and closes
the complete quartic unmarked and marked additive matrix cards.
Together with V44--V45, the additive local card now holds through
arity four.  It does not prove UTCSC, GRMSC, EMCC, thermal
connected WPRAC, all-arity RGSC, the raw Wilson aggregate strong
card, all-order strong MTA, WUFC, CPM, cutoff removal, reflection
positivity, Osterwalder--Schrader reconstruction, a continuum
Yang--Mills measure, or a positive Hamiltonian gap.  The full
Yang--Mills mass gap has not been proved.
\end{firewall}

\begin{theorem}[V46 result ledger]
\label{thm:newV46-results}
V46:
\begin{enumerate}[label=\textup{(V46.\arabic*)},leftmargin=6.5em]
\item proves \(O(s)\) Gaussian moment comparison through degree
      four for both endpoints and every additive mixture;
\item proves an \(O(s)\) fourth cumulant for the scaled tangent
      vector;
\item establishes uniform linear, quadratic, and cubic matrix
      Taylor rows;
\item computes and bounds the complete four-linear Wick residual;
\item proves the odd \(HLLL\) cumulant vanishes;
\item closes the quartic unmarked matrix card;
\item closes the quartic marked card by additive polynomial
      interpolation; and
\item isolates UTCSC as the all-arity continuation.
\end{enumerate}
\end{theorem}

\begin{proof}
Use
\Cref{lem:newV46-moment-comparison,
cor:newV46-Y-cumulant,
lem:newV46-LHT-bounds,
lem:newV46-centered-HT,
thm:newV46-four-linear,
lem:newV46-HLLL-zero,
lem:newV46-fixed-cumulant,
thm:newV46-quartic-unmarked,
thm:newV46-quartic-marked,
cor:newV46-through-four}.
\end{proof}

\paragraph{Parent-version record.}
V46 was formed from
\texttt{Renewal\_Yang\_Mills\_Mass\_Gap\_Research\_Notes\_V45.tex}.
The V45 parent had \(21\,667\) logical source lines,
\(683\,923\) bytes, and SHA--256
\[
 \texttt{3eda55900b68bbe85a1dc43eb0f51433374bf3559047967899aa7967772d706f}.
\]
The next compulsory companion audit is V50.

% =============================================================================
% END APPENDED FIFTH-SERIES V46 MATERIAL
% =============================================================================

% =============================================================================
% BEGIN APPENDED FIFTH-SERIES V47 MATERIAL
% =============================================================================

\chapter{V47: Exact parity decomposition and the five-input
matrix card}
\label{chap:newV47-quintic}

\section{Moment comparison through degree six}
\label{sec:newV47-moments}

No companion audit is due at V47.

\begin{lemma}[Endpoint moment comparison through degree six]
\label{lem:newV47-moment-six}
For every polynomial \(P\) on \(\mathbb R^3\) of degree at most
six,
\begin{equation}
 \left|
 \mathbb E_{\mu_{\alpha,s,t}^{\rm add}}P(Y_s)
 -
 \mathbb E_{\gamma_3}P(Y)
 \right|
 \le C_Ps
 \label{eq:newV47-moment-six}
\end{equation}
uniformly in \(t\), where \(Y_s\) is
\eqref{eq:newV46-Y}.
\end{lemma}

\begin{proof}
Repeat the proof of
\Cref{lem:newV46-moment-comparison}.  On
\(|Y_s|\le s^{-1/8}\), the heat image/Haar factor and the Wilson
action/Haar factor differ from the standard Gaussian density by
\[
 O\left(s(1+|Y_s|^4)\right)
\]
under a fixed slightly weaker Gaussian majorant.  Multiplication
by a polynomial of degree at most six remains integrable.
V30 concentration makes the complementary contribution
exponentially small with every fixed polynomial weight.
Normalization and additive convex combination preserve the
\(O(s)\) bound.
\end{proof}

\section{Exact odd/even representation increments}
\label{sec:newV47-parity}

Define
\begin{align}
 O_i(U)
 &:=
 {D^{R_i}(U)-D^{R_i}(U^{-1})\over2},                         \label{eq:newV47-O}\\
 E_i(U)
 &:=
 {D^{R_i}(U)+D^{R_i}(U^{-1})\over2}-I.                       \label{eq:newV47-E}
\end{align}
Then
\[
 D^{R_i}(U)-I=O_i(U)+E_i(U),
\]
\(O_i\) is odd and \(E_i\) is even under inversion.

\begin{lemma}[Exact parity rows]
\label{lem:newV47-OE-bounds}
Uniformly on thermal representations,
\begin{align}
 \|O_i\|&\le Cb_iY_s,                                        \label{eq:newV47-O-bound}\\
 \|E_i\|&\le Cb_i^2Y_s^2.                                    \label{eq:newV47-E-bound}
\end{align}
Moreover \(\mathbb EO_i=0\) and
\begin{equation}
 \|E_i-\mathbb EE_i\|
 \le Cb_i^2(1+Y_s^2).
 \label{eq:newV47-E-centered}
\end{equation}
\end{lemma}

\begin{proof}
On the principal geodesic,
\[
 O_i=\sinh(\theta G_i),
 \qquad
 E_i=\cosh(\theta G_i)-I,
\]
with the hyperbolic notation evaluated at the skew-Hermitian
operator \(G_i\).  The spectral theorem gives
\[
 \|\sinh(\theta G_i)\|\le\theta\|G_i\|,
 \qquad
 \|\cosh(\theta G_i)-I\|\le{\theta^2\over2}\|G_i\|^2.
\]
Use \(\sqrt{sC_2(R_i)}\le b_i\) and
\(\theta=\sqrt{2s}Y_s\).  Inversion symmetry gives
\(\mathbb EO_i=0\); V30 fixed moments give the centered even bound.
\end{proof}

Retain the leading terms
\[
 L_i=\theta G_i,
 \qquad
 H_i={\theta^2\over2}G_i^2
\]
and define exact parity-preserving remainders
\begin{align}
 O_i^{[3]}&:=O_i-L_i,                                        \label{eq:newV47-O3}\\
 E_i^{[4]}&:=E_i-H_i.                                        \label{eq:newV47-E4}
\end{align}

\begin{lemma}[Odd cubic and even quartic remainder rows]
\label{lem:newV47-remainders}
\begin{align}
 \|O_i^{[3]}\|
 &\le Cb_i^3Y_s^3,                                          \label{eq:newV47-O3-bound}\\
 \|E_i^{[4]}\|
 &\le Cb_i^4Y_s^4.                                          \label{eq:newV47-E4-bound}
\end{align}
The first remainder is odd and the second is even.
\end{lemma}

\begin{proof}
Use the third-order remainder for \(\sinh\) and the fourth-order
remainder for \(\cosh\) on the imaginary spectrum of \(G_i\):
\[
 \|\sinh(\theta G_i)-\theta G_i\|
 \le{\theta^3\over6}\|G_i\|^3,
\]
\[
 \left\|
 \cosh(\theta G_i)-I-{\theta^2\over2}G_i^2
 \right\|
 \le{\theta^4\over24}\|G_i\|^4.
\]
The mark bounds follow as in
\Cref{lem:newV46-LHT-bounds}; parity is exact from the definitions.
\end{proof}

\section{The leading Gaussian quintic identity}
\label{sec:newV47-Gaussian-zero}

\begin{lemma}[One quadratic and four linear Gaussian entries]
\label{lem:newV47-QLLLL-zero}
Let \(Y\sim\gamma_3\), let \(Q(Y)\) be any polynomial of degree at
most two, and let \(\ell_1,\ldots,\ell_4\) be linear forms.  Then
\begin{equation}
 \kappa_5^{\gamma_3}
 (Q(Y),\ell_1(Y),\ell_2(Y),\ell_3(Y),\ell_4(Y))
 =0.
 \label{eq:newV47-QLLLL-zero}
\end{equation}
\end{lemma}

\begin{proof}
For a source \(u\in\mathbb R^3\), Gaussian exponential tilting
gives
\[
 {\mathbb E[Q(Y)e^{\langle u,Y\rangle}]
  \over
  \mathbb E e^{\langle u,Y\rangle}}
 =
 \mathbb E[Q(Y+u)].
\]
The right side is a polynomial of degree at most two in \(u\).
The joint cumulant in \eqref{eq:newV47-QLLLL-zero} is its fourth
directional derivative in the four \(\ell_j\) directions, hence
vanishes.
\end{proof}

\begin{theorem}[Calibrated \(HLLLL\) residual]
\label{thm:newV47-HLLLL}
For every placement of one quadratic entry and four linear
entries,
\begin{align}
 &\left\|
 \kappa_5^{\,\mu_{\alpha,s,t}^{\rm add}}
 (H_i,L_j,L_k,L_\ell,L_r)
 \right\|_{\rm cb}
 \notag\\
 &\qquad\le
 Cs
 \left(\prod_{a\in\{i,j,k,\ell,r\}}b_a\right)b_i.
 \label{eq:newV47-HLLLL}
\end{align}
\end{theorem}

\begin{proof}
After inserting
\[
 L_j=\sqrt{2s}\sum_aY_s^aJ_j^a,
 \qquad
 H_i=s\sum_{a,b}Y_s^aY_s^bJ_i^aJ_i^b,
\]
the cumulant is a fixed finite contraction of moments of \(Y_s\)
of degree at most six.  Its Gaussian value vanishes by
\Cref{lem:newV47-QLLLL-zero}.  The moment--cumulant polynomial is
Lipschitz in this finite list of moments on their uniformly bounded
range.  Hence \Cref{lem:newV47-moment-six} supplies an exterior
\(Cs\).

The four linear generator scales and one quadratic generator scale
are bounded by
\[
 \left(\prod_a b_a\right)b_i.
\]
The finite coordinate sum and spectator amplification are absorbed
in \(C\).
\end{proof}

\section{Complete quintic card}
\label{sec:newV47-card}

\begin{lemma}[Fixed fifth-cumulant product bound]
\label{lem:newV47-fixed-cumulant}
For five centered random operators satisfying fixed polynomial
envelopes
\[
 \|X_i\|\le a_i(1+Y_s^{r_i}),
 \qquad r_i\le4,
\]
one has
\begin{equation}
 \|\kappa_5(X_1,\ldots,X_5)\|_{\rm cb}
 \le C_{\boldsymbol r}\prod_i a_i.
 \label{eq:newV47-fixed-cumulant}
\end{equation}
\end{lemma}

\begin{proof}
The fifth moment--cumulant formula has finitely many partitions.
Apply the operator triangle inequality and H\"older to every block
moment, then use V30 for the needed fixed radial moments.
\end{proof}

\begin{theorem}[Quintic unmarked additive matrix card]
\label{thm:newV47-quintic-unmarked}
For five arbitrary thermal representations, with
\(B=\sum_{i=1}^5b_i\),
\begin{equation}
 \left\|
 \kappa_5^{\,\mu_{\alpha,s,t}^{\rm add}}
 (D^{R_1},\ldots,D^{R_5})
 \right\|_{\rm cb}
 \le
 K^5\left(\prod_{i=1}^5b_i\right)B^3.
 \label{eq:newV47-quintic-unmarked}
\end{equation}
\end{theorem}

\begin{proof}
For cumulants of order at least two, replace each input by
\[
 D^{R_i}-\mathbb ED^{R_i}
 =
 O_i+(E_i-\mathbb EE_i).
\]
Expand multilinearly.  Inversion kills every term with an odd
number of odd \(O\)'s.  With five inputs, the number of even
entries is therefore odd: one, three, or five.

Terms with at least three even entries carry, beyond
\(\prod_i b_i\), at least three further marks.  Their sum is
bounded by
\[
 C\left(\prod_i b_i\right)B^3
\]
using \Cref{lem:newV47-OE-bounds,
lem:newV47-fixed-cumulant}.

Consider a term with exactly one even entry \(E_i-\mathbb EE_i\)
and four odd entries.  Insert
\[
 E_i=H_i+E_i^{[4]},
 \qquad
 O_j=L_j+O_j^{[3]}.
\]
The leading \(H_iLLLL\) term is bounded by
\Cref{thm:newV47-HLLLL}.  Since every
\(b_a\ge\sqrt s\),
\[
 B\ge5\sqrt s,
 \qquad
 b_i s\le {1\over25}B^3.
\]
It therefore has the target scale.

If \(E_i^{[4]}\) is used, it contributes three additional marks
beyond the base \(b_i\).  If any \(O_j^{[3]}\) is used together
with \(H_i\), the quadratic entry supplies one additional mark and
the cubic odd remainder supplies two, again totaling three.
Additional remainders only improve the mark count because
\(b_i<1\).  Apply
\Cref{lem:newV47-remainders,
lem:newV47-fixed-cumulant} and sum the finitely many terms.
\end{proof}

\begin{theorem}[Quintic additive score card]
\label{thm:newV47-quintic-marked}
Uniformly in \(t\),
\begin{align}
 &\left\|
 \kappa_6^{\,\mu_{\alpha,s,t}^{\rm add}}
 (D^{R_1},\ldots,D^{R_5},\Sigma_{\alpha,t}^{\rm add})
 \right\|_{\rm cb}
 \notag\\
 &\qquad\le
 K^5\left(\prod_{i=1}^5b_i\right)
 \left(\sum_{i=1}^5b_i\right)^3.
 \label{eq:newV47-quintic-marked}
\end{align}
\end{theorem}

\begin{proof}
The unmarked fifth cumulant of the additive law is a
Banach-valued polynomial of degree at most five in \(t\).
Apply \Cref{lem:newV45-Markov} with \(d=5\) to
\Cref{thm:newV47-quintic-unmarked}.  Its derivative is the exact
additive score cumulant.  Absorb the fixed Markov factor into
\(K^5\).
\end{proof}

\begin{corollary}[Additive matrix local card through arity five]
\label{cor:newV47-through-five}
The complete unmarked and marked additive-bridge matrix cards hold
for every \(2\le m\le5\), uniformly on thermal representations and
all amplifications.
\end{corollary}

\begin{proof}
Combine
\Cref{cor:newV46-through-four,
thm:newV47-quintic-unmarked,
thm:newV47-quintic-marked}.
\end{proof}

\section{The first genuinely graph-sensitive arity}
\label{sec:newV47-six}

\begin{assessment}[Why arity six is different]
\label{ass:newV47-six}
At arity five, inversion forces at least one even representation
entry, and that entry supplies one of the three required additional
marks.  The remaining two are supplied by the \(O(s)\) Gaussian
comparison.  At arity six, the all-odd linear term is permitted.
Its Gaussian cumulant vanishes, but a bare \(O(s)\) comparison
supplies only two marks, whereas the target needs four.  The
first-order density perturbation must therefore be kept connected
to at least one additional Gaussian pair edge.  This is the first
arity at which the V27 rooted graph structure, rather than a
fixed-moment comparison alone, is indispensable.
\end{assessment}

\begin{definition}[First-order quartic-root operator card]
\label{def:newV47-FQROC}
FQROC is the assertion that the order-\(s\) term in the exact
scaled endpoint density ratio, inserted into the six-or-more input
operator source logarithm, admits a time-ordered rooted forest in
which:
\begin{enumerate}[label=\textup{(\roman*)},leftmargin=3.5em]
\item the density root has degree at most four;
\item every remaining input component is joined by a Gaussian pair
      row of size \(Cb_ib_j\); and
\item the complete signed/time-ordered sum has exponential arity
      grade.
\end{enumerate}
\end{definition}

\begin{programme}[V48: the arity-six first-order graph]
\label{prog:newV47-V48}
\begin{enumerate}[leftmargin=3em]
\item Compute the order-\(s\) polynomial in each scaled endpoint
      density ratio, including normalization.
\item Compute its joint cumulant with six linear representation
      entries before taking norms.
\item Show algebraically that every nonzero Wick contraction has
      one quartic root and at least one additional input pair.
\item Sum the Lie-algebra coordinate contractions into the
      Casimir pair rows of V28 and retain Duhamel time ordering.
\item Prove the complete arity-six bound; only then formulate the
      repeated-root all-arity induction.
\end{enumerate}
\end{programme}

\begin{firewall}[Claim boundary after V47]
\label{fire:newV47-boundary}
V47 proves exact odd/even representation decomposition, extends
the endpoint Gaussian comparison through degree six, proves the
Gaussian \(QLLLL\) cumulant vanishes, and closes the complete
quintic unmarked and marked additive matrix cards.  The additive
local card now holds through arity five.  It does not prove the
arity-six card, FQROC, UTCSC, GRMSC, EMCC, all-arity RGSC, the raw
Wilson aggregate strong card, all-order strong MTA, WUFC, CPM,
cutoff removal, reflection positivity,
Osterwalder--Schrader reconstruction, a continuum Yang--Mills
measure, or a positive Hamiltonian gap.  The full Yang--Mills mass
gap has not been proved.
\end{firewall}

\begin{theorem}[V47 result ledger]
\label{thm:newV47-results}
V47:
\begin{enumerate}[label=\textup{(V47.\arabic*)},leftmargin=6.5em]
\item extends calibrated Gaussian moment comparison through degree
      six;
\item gives an exact odd/even matrix-increment decomposition with
      uniform thermal rows;
\item isolates odd cubic and even quartic remainders;
\item proves the Gaussian \(QLLLL\) connected tensor is zero;
\item proves the calibrated \(HLLLL\) residual carries \(s\);
\item closes the quintic unmarked matrix card;
\item closes the quintic marked card; and
\item isolates arity six as the first mandatory rooted-graph test.
\end{enumerate}
\end{theorem}

\begin{proof}
Use
\Cref{lem:newV47-moment-six,
lem:newV47-OE-bounds,
lem:newV47-remainders,
lem:newV47-QLLLL-zero,
thm:newV47-HLLLL,
lem:newV47-fixed-cumulant,
thm:newV47-quintic-unmarked,
thm:newV47-quintic-marked,
cor:newV47-through-five}.
\end{proof}

\paragraph{Parent-version record.}
V47 was formed from
\texttt{Renewal\_Yang\_Mills\_Mass\_Gap\_Research\_Notes\_V46.tex}.
The V46 parent had \(22\,142\) logical source lines,
\(698\,448\) bytes, and SHA--256
\[
 \texttt{b0e7e041bf91877453fc84691137484bdf21f477ed8276b12c010d14114bd7be}.
\]
The next compulsory companion audit is V50.

% =============================================================================
% END APPENDED FIFTH-SERIES V47 MATERIAL
% =============================================================================

% =============================================================================
% BEGIN APPENDED FIFTH-SERIES V48 MATERIAL
% =============================================================================

\chapter{V48: First density roots and the six-input matrix card}
\label{chap:newV48-six}

\section{The exact first-order scaled density polynomials}
\label{sec:newV48-density}

No companion audit is due at V48.  Retain
\[
 Y_s={\theta\boldsymbol n\over\sqrt{2s}},
 \qquad r:=|Y_s|,
 \qquad \gamma_3(dY)=(2\pi)^{-3/2}e^{-r^2/2}\,dY.
\]

\begin{definition}[First heat and Wilson density roots]
\label{def:newV48-roots}
Define the centered radial polynomials
\begin{align}
 P_H(Y)
 &:=
 1-{r^2\over3},                                               \label{eq:newV48-PH}\\
 P_W(Y)
 &:=
 {r^4\over12}-{r^2\over3}-{1\over4},                         \label{eq:newV48-PW}\\
 P_t(Y)
 &:=
 (1-t)P_H(Y)+tP_W(Y).                                        \label{eq:newV48-Pt}
\end{align}
Then
\[
 \mathbb E_{\gamma_3}P_H
 =
 \mathbb E_{\gamma_3}P_W
 =
 \mathbb E_{\gamma_3}P_t=0.
\]
\end{definition}

\begin{proof}
Use
\(\mathbb E_{\gamma_3}r^2=3\) and
\(\mathbb E_{\gamma_3}r^4=15\).
\end{proof}

\begin{theorem}[Second-order endpoint density expansion in moments]
\label{thm:newV48-density-expansion}
For every fixed polynomial \(F\) on \(\mathbb R^3\),
\begin{align}
 \mathbb E_HF(Y_s)
 &=
 \mathbb E_\gamma F
 +s\mathbb E_\gamma[FP_H]
 +O_F(s^2),                                                   \label{eq:newV48-heat-expansion}\\
 \mathbb E_WF(Y_s)
 &=
 \mathbb E_\gamma F
 +s\mathbb E_\gamma[FP_W]
 +O_F(s^2),                                                   \label{eq:newV48-Wilson-expansion}\\
 \mathbb E_t^{\rm add}F(Y_s)
 &=
 \mathbb E_\gamma F
 +s\mathbb E_\gamma[FP_t]
 +O_F(s^2),                                                   \label{eq:newV48-add-expansion}
\end{align}
uniformly in \(t\in[0,1]\).
\end{theorem}

\begin{proof}
Split at \(r=s^{-1/10}\).  V30 concentration makes the complement
smaller than every power of \(s\) with any fixed polynomial
weight.

For the heat law, the scaled \(k=0\) image and radial Haar measure
give, before normalization,
\[
 e^{-r^2/2}
 {\sin(\sqrt{2s}r)\over\sqrt{2s}r}
 \left(1+O(e^{-c/s})\right).
\]
The non-Gaussian factor is
\[
 1-{s r^2\over3}+O(s^2r^4).
\]
Its Gaussian mean is \(1-s+O(s^2)\).  Division by this mean gives
\[
 1+s\left(1-{r^2\over3}\right)
 +O\left(s^2(1+r^4)\right),
\]
which proves \eqref{eq:newV48-heat-expansion}.

For the Wilson law,
\[
 \alpha={1\over2s}-{1\over3},
 \qquad
 \theta=\sqrt{2s}\,r.
\]
Taylor expansion gives
\begin{align*}
 \alpha(\cos\theta-1)
 &=
 -{r^2\over2}
 +s\left({r^2\over3}+{r^4\over12}\right)
 +O\left(s^2(r^4+r^6)\right),\\
 \left({\sin\theta\over\theta}\right)^2
 &=
 1-{2s r^2\over3}+O(s^2r^4).
\end{align*}
Thus the unnormalized density ratio to \(\gamma_3\) is
\[
 1+s\left({r^4\over12}-{r^2\over3}\right)
 +O\left(s^2(1+r^8)\right).
\]
The Gaussian mean of the first-order polynomial is \(1/4\).
Normalization subtracts that scalar and yields \(P_W\), proving
\eqref{eq:newV48-Wilson-expansion}.  The additive identity is the
convex combination of the endpoint identities.
\end{proof}

\begin{corollary}[The first bridge-density score root]
\label{cor:newV48-root-difference}
\begin{equation}
 P_W-P_H={r^4-15\over12}.
 \label{eq:newV48-root-difference}
\end{equation}
Thus the first additive endpoint difference is precisely the
centered radial quartic root already found in V25, V32, and V37.
\end{corollary}

\begin{proof}
Subtract \eqref{eq:newV48-PH} from
\eqref{eq:newV48-PW}.
\end{proof}

\section{Gaussian root-degree vanishing}
\label{sec:newV48-root-degree}

\begin{lemma}[One polynomial root cannot see more linear leaves
than its degree]
\label{lem:newV48-one-root}
Let \(P\) be a polynomial of degree at most \(d\) and let
\(\ell_1,\ldots,\ell_m\) be linear forms of a standard Gaussian
vector.  If \(m>d\), then
\begin{equation}
 \kappa_{m+1}^{\gamma_3}
 (P(Y),\ell_1(Y),\ldots,\ell_m(Y))=0.
 \label{eq:newV48-one-root}
\end{equation}
\end{lemma}

\begin{proof}
As in \Cref{lem:newV47-QLLLL-zero}, Gaussian tilting gives
\[
 {\mathbb E[P(Y)e^{\langle u,Y\rangle}]
  \over\mathbb Ee^{\langle u,Y\rangle}}
 =
 \mathbb E P(Y+u),
\]
a polynomial of degree at most \(d\) in \(u\).  The cumulant is its
\(m\)-th derivative in the \(\ell_i\) directions and hence
vanishes.
\end{proof}

\begin{lemma}[Two quadratic roots cannot connect four linear
leaves]
\label{lem:newV48-two-quadratic}
If \(Q_1,Q_2\) have degree at most two and
\(\ell_1,\ldots,\ell_4\) are linear Gaussian forms, then
\begin{equation}
 \kappa_6^{\gamma_3}
 (Q_1,Q_2,\ell_1,\ell_2,\ell_3,\ell_4)=0.
 \label{eq:newV48-two-quadratic}
\end{equation}
\end{lemma}

\begin{proof}
After exponential tilting by \(u\), the joint cumulant of the two
quadratics is
\[
 \operatorname{Cov}_{\gamma_3}
 (Q_1(Y+u),Q_2(Y+u)).
\]
The deterministic quadratic terms in \(u\) disappear from the
covariance, odd Gaussian terms vanish, and the remaining expression
has degree at most two in \(u\).  Four directional derivatives
therefore give zero.
\end{proof}

\section{The six-linear cumulant begins at second order}
\label{sec:newV48-six-linear}

\begin{theorem}[Second-order six-linear tangent row]
\label{thm:newV48-six-linear}
For coordinate indices \(a_1,\ldots,a_6\),
\begin{equation}
 \left|
 \kappa_6^{\,\mu_{\alpha,s,t}^{\rm add}}
 (Y_s^{a_1},\ldots,Y_s^{a_6})
 \right|
 \le Cs^2.
 \label{eq:newV48-six-linear}
\end{equation}
\end{theorem}

\begin{proof}
Moment--cumulant inversion is a polynomial in moments of degree at
most six.  Insert \eqref{eq:newV48-add-expansion} into that finite
polynomial.  The zeroth-order coefficient is the sixth Gaussian
cumulant and vanishes.  Differentiation of a cumulant under the
first-order signed tilt \(P_t\) gives
\[
 \kappa_7^{\gamma_3}
 (Y^{a_1},\ldots,Y^{a_6},P_t(Y)).
\]
The root \(P_t\) has degree at most four, so this coefficient
vanishes by \Cref{lem:newV48-one-root}.  The moment remainders in
\Cref{thm:newV48-density-expansion} are \(O(s^2)\), proving the
claim.
\end{proof}

\begin{corollary}[Six-linear representation row]
\label{cor:newV48-six-linear-matrix}
\begin{equation}
 \left\|
 \kappa_6^{\,\mu_{\alpha,s,t}^{\rm add}}
 (L_1,\ldots,L_6)
 \right\|_{\rm cb}
 \le
 Cs^2\prod_{i=1}^6b_i.
 \label{eq:newV48-six-linear-matrix}
\end{equation}
\end{corollary}

\begin{proof}
Insert
\[
 L_i=\sqrt{2s}\sum_aY_s^aJ_i^a
\]
and argue as in \Cref{thm:newV46-four-linear}.  The generator
scales outside the coordinate cumulant are bounded by
\(\prod_i b_i\); apply
\Cref{thm:newV48-six-linear}.
\end{proof}

\section{Leading nonlinear six-input tensors}
\label{sec:newV48-leading-nonlinear}

Define the odd cubic Taylor term
\begin{equation}
 C_i^{[3]}:={\theta^3\over6}G_i^3
 \label{eq:newV48-C3}
\end{equation}
and the odd fifth-order remainder
\[
 O_i^{[5]}:=O_i-L_i-C_i^{[3]}.
\]

\begin{lemma}[Odd fifth-order remainder]
\label{lem:newV48-O5}
\begin{equation}
 \|O_i^{[5]}\|
 \le Cb_i^5Y_s^5.
 \label{eq:newV48-O5}
\end{equation}
\end{lemma}

\begin{proof}
Use the fifth-order Taylor remainder of \(\sinh(\theta G_i)\) on
the imaginary spectrum and
\(\sqrt{sC_2(R_i)}\le b_i\).
\end{proof}

\begin{lemma}[Leading one-cubic row]
\label{lem:newV48-C3-five-L}
For every placement,
\begin{align}
 &\left\|
 \kappa_6^{\,\mu_{\alpha,s,t}^{\rm add}}
 (C_i^{[3]},L_j,L_k,L_\ell,L_r,L_u)
 \right\|_{\rm cb}
 \notag\\
 &\qquad\le
 Cs
 \left(\prod_{a=1}^6b_a\right)b_i^2.
 \label{eq:newV48-C3-five-L}
\end{align}
\end{lemma}

\begin{proof}
The scaled polynomial degree is eight.  Its Gaussian cumulant
vanishes by \Cref{lem:newV48-one-root}, with the cubic entry as the
degree-three root and the other five entries linear.  The
first-order part of
\Cref{thm:newV48-density-expansion}, applied to the fixed
degree-eight moment--cumulant polynomial, gives an \(O(s)\)
difference from the zero Gaussian value.  The cubic entry carries
two additional \(b_i\)'s beyond the base \(\prod_ab_a\).
\end{proof}

\begin{lemma}[Leading two-quadratic row]
\label{lem:newV48-two-H-four-L}
For every placement,
\begin{align}
 &\left\|
 \kappa_6^{\,\mu_{\alpha,s,t}^{\rm add}}
 (H_i,H_j,L_k,L_\ell,L_r,L_u)
 \right\|_{\rm cb}
 \notag\\
 &\qquad\le
 Cs
 \left(\prod_{a=1}^6b_a\right)b_ib_j.
 \label{eq:newV48-two-H-four-L}
\end{align}
\end{lemma}

\begin{proof}
The Gaussian value is zero by
\Cref{lem:newV48-two-quadratic}.
\Cref{thm:newV48-density-expansion}, applied through degree eight,
gives an \(O(s)\) residual.  Each quadratic entry carries one
additional mark beyond its base mark.
\end{proof}

\section{Complete six-input card}
\label{sec:newV48-card}

\begin{lemma}[Fixed sixth-cumulant product bound]
\label{lem:newV48-fixed-cumulant}
For six centered random operators with fixed polynomial radial
envelopes
\[
 \|X_i\|\le a_i(1+Y_s^{r_i}),
 \qquad r_i\le5,
\]
one has
\[
 \|\kappa_6(X_1,\ldots,X_6)\|_{\rm cb}
 \le C_{\boldsymbol r}\prod_i a_i.
\]
\end{lemma}

\begin{proof}
Apply moment--cumulant inversion, the operator triangle inequality,
H\"older, and V30 fixed radial moments.  The partition family is
finite at arity six.
\end{proof}

\begin{theorem}[Six-input unmarked additive matrix card]
\label{thm:newV48-six-unmarked}
For six arbitrary thermal inputs and
\(B=\sum_{i=1}^6b_i\),
\begin{equation}
 \left\|
 \kappa_6^{\,\mu_{\alpha,s,t}^{\rm add}}
 (D^{R_1},\ldots,D^{R_6})
 \right\|_{\rm cb}
 \le
 K^6\left(\prod_{i=1}^6b_i\right)B^4.
 \label{eq:newV48-six-unmarked}
\end{equation}
\end{theorem}

\begin{proof}
Use the exact odd/even decomposition of V47.  Inversion requires an
even number of odd entries, so the number of even entries is
\(0,2,4\), or \(6\).

With no even entries, expand each odd entry as
\[
 O_i=L_i+C_i^{[3]}+O_i^{[5]}.
\]
The all-linear term is bounded by
\Cref{cor:newV48-six-linear-matrix}.  Since
\[
 B\ge6\sqrt s,
 \qquad
 s^2\le {B^4\over6^4},
\]
it has the target scale.  A term with exactly one cubic entry and
five linear entries is bounded by
\Cref{lem:newV48-C3-five-L}; its extra factor satisfies
\[
 sb_i^2\le sB^2\le {B^4\over36}.
\]
Two cubic entries already supply four extra marks.  One
fifth-order remainder also supplies four; further remainders only
improve the count.

With exactly two even entries, expand each as
\[
 E_i=H_i+E_i^{[4]}.
\]
The leading \(HHLLLL\) term is bounded by
\Cref{lem:newV48-two-H-four-L}, and
\[
 sb_ib_j\le sB^2\le {B^4\over36}.
\]
Replacing one \(H\) by \(E^{[4]}\) supplies three extra marks from
that entry and one from the other quadratic entry, totaling four.
Any odd cubic replacement supplies two additional marks on top of
the two quadratic marks.  All other replacements improve the
count.

Finally, four or six even entries directly supply at least four
additional marks.  Apply
\Cref{lem:newV47-OE-bounds,
lem:newV47-remainders,
lem:newV48-O5,
lem:newV48-fixed-cumulant} to the finitely many remaining terms
and enlarge \(K\).
\end{proof}

\begin{theorem}[Six-input additive score card]
\label{thm:newV48-six-marked}
Uniformly in \(t\),
\begin{align}
 &\left\|
 \kappa_7^{\,\mu_{\alpha,s,t}^{\rm add}}
 (D^{R_1},\ldots,D^{R_6},\Sigma_{\alpha,t}^{\rm add})
 \right\|_{\rm cb}
 \notag\\
 &\qquad\le
 K^6\left(\prod_{i=1}^6b_i\right)
 \left(\sum_{i=1}^6b_i\right)^4.
 \label{eq:newV48-six-marked}
\end{align}
\end{theorem}

\begin{proof}
The unmarked sixth cumulant is a Banach-valued polynomial in the
additive parameter of degree at most six.  Apply
\Cref{lem:newV45-Markov} with \(d=6\) to
\Cref{thm:newV48-six-unmarked}.  Its derivative is the displayed
score cumulant.
\end{proof}

\begin{corollary}[Additive matrix local card through arity six]
\label{cor:newV48-through-six}
The complete unmarked and marked additive-bridge matrix cards hold
for every \(2\le m\le6\), uniformly on thermal representations and
all amplifications.
\end{corollary}

\begin{proof}
Combine
\Cref{cor:newV47-through-five,
thm:newV48-six-unmarked,
thm:newV48-six-marked}.
\end{proof}

\section{What the arity-six calculation changes}
\label{sec:newV48-pattern}

\begin{assessment}[Correction of the anticipated first-order
graph]
\label{ass:newV48-correction}
V47 anticipated one quartic root plus one additional Gaussian pair
at arity six.  For six \emph{linear} leaves this connected graph
does not exist: the root has at most four legs, and a pair joining
the two remaining one-leg leaves would be a disconnected
component.  Equivalently,
\Cref{lem:newV48-one-root} kills the full first-order coefficient.
The six-linear cumulant therefore starts at \(s^2\).  Pair
structure first becomes indispensable when nonlinear input
vertices or repeated density roots can share more than one leg.
\end{assessment}

\begin{definition}[Density-root degree card]
\label{def:newV48-DRDC}
DRDC is a uniform expansion of the scaled endpoint density ratio
in centered Gaussian chaos roots such that the coefficient of
total root order \(q\):
\begin{enumerate}[label=\textup{(\roman*)},leftmargin=3.5em]
\item carries \(s^q\);
\item has total chaos degree at most \(4q+2q_{\rm curv}\), with
      every curvature correction explicitly counted; and
\item after connected Gaussian contraction with the exact
      odd/even representation vertices, obeys the weighted
      Pr\"ufer graph census with exponential arity grade.
\end{enumerate}
\end{definition}

\begin{programme}[V49: repeated density-root power counting]
\label{prog:newV48-V49}
\begin{enumerate}[leftmargin=3em]
\item Expand the normalized heat and Wilson density ratios through
      order \(s^2\) and identify their centered chaos degrees.
\item Prove the general Gaussian leg-count inequality: a connected
      Wick graph on polynomial vertices requires at least
      \(v-1\) contractions.
\item Convert that inequality into the minimal power of \(s\) and
      the minimal number of nonlinear input marks at arbitrary
      fixed arity.
\item Determine whether degree growth alone yields the full target
      or whether a stable repeated-root graph sum is still needed
      as \(m\to\infty\).
\item Prepare the compulsory V50 companion audit.
\end{enumerate}
\end{programme}

\begin{firewall}[Claim boundary after V48]
\label{fire:newV48-boundary}
V48 computes the exact first heat and Wilson density roots, proves
their second-order moment expansion, proves the first-order
six-linear coefficient vanishes, controls the leading cubic and
two-quadratic nonlinear rows, and closes the complete six-input
unmarked and marked additive matrix cards.  The additive local
card now holds through arity six.  It does not prove DRDC, UTCSC,
GRMSC, EMCC, all-arity RGSC, the raw Wilson aggregate strong card,
all-order strong MTA, WUFC, CPM, cutoff removal, reflection
positivity, Osterwalder--Schrader reconstruction, a continuum
Yang--Mills measure, or a positive Hamiltonian gap.  The full
Yang--Mills mass gap has not been proved.
\end{firewall}

\begin{theorem}[V48 result ledger]
\label{thm:newV48-results}
V48:
\begin{enumerate}[label=\textup{(V48.\arabic*)},leftmargin=6.5em]
\item computes the centered first-order heat and Wilson scaled
      density roots;
\item proves a uniform second-order density expansion in every
      fixed polynomial moment;
\item proves the Gaussian one-root degree-vanishing lemma;
\item proves the six-linear tangent cumulant begins at \(s^2\);
\item controls the leading cubic--linear and
      quadratic--quadratic--linear residuals;
\item closes the six-input unmarked matrix card;
\item closes the six-input marked card; and
\item isolates the repeated density-root degree census for V49.
\end{enumerate}
\end{theorem}

\begin{proof}
Use
\Cref{thm:newV48-density-expansion,
cor:newV48-root-difference,
lem:newV48-one-root,
lem:newV48-two-quadratic,
thm:newV48-six-linear,
cor:newV48-six-linear-matrix,
lem:newV48-C3-five-L,
lem:newV48-two-H-four-L,
thm:newV48-six-unmarked,
thm:newV48-six-marked,
cor:newV48-through-six}.
\end{proof}

\paragraph{Parent-version record.}
V48 was formed from
\texttt{Renewal\_Yang\_Mills\_Mass\_Gap\_Research\_Notes\_V47.tex}.
The V47 parent had \(22\,584\) logical source lines,
\(711\,913\) bytes, and SHA--256
\[
 \texttt{eecbd67554c1786709424d96bc86d2e1689a7ea411a4211905ae49800bca6d66}.
\]
The next compulsory companion audit is V50.

% =============================================================================
% END APPENDED FIFTH-SERIES V48 MATERIAL
% =============================================================================

% =============================================================================
% BEGIN APPENDED FIFTH-SERIES V49 MATERIAL
% =============================================================================

\chapter{V49: The arity-eight additive-mixture obstruction}
\label{chap:newV49-no-go}

\section{Gaussian connected leg count}
\label{sec:newV49-leg-count}

No companion audit is due at V49.

\begin{lemma}[Connected Wick leg inequality]
\label{lem:newV49-leg-count}
Let \(P_1,\ldots,P_v\) be polynomials of respective degrees
\(d_1,\ldots,d_v\) in a centered Gaussian vector.  If
\begin{equation}
 \sum_{a=1}^v d_a<2(v-1),
 \label{eq:newV49-leg-condition}
\end{equation}
then
\begin{equation}
 \kappa_v^{\gamma}(P_1,\ldots,P_v)=0.
 \label{eq:newV49-leg-zero}
\end{equation}
\end{lemma}

\begin{proof}
Expand every polynomial into monomials and apply Wick's theorem to
each moment in moment--cumulant inversion.  The partition
M\"obius sum cancels every Wick pairing whose contraction graph on
the \(v\) labelled polynomial vertices is disconnected.  A
connected graph on \(v\) vertices has at least \(v-1\) contraction
edges, and each contraction uses two polynomial legs.  Under
\eqref{eq:newV49-leg-condition} there are too few legs for such a
graph.
\end{proof}

\begin{corollary}[Grade versus nonlinear marks]
\label{cor:newV49-grade-count}
Suppose \(m\) representation-input vertices have total polynomial
degree \(m+\rho\), so \(\rho\) is their number of Taylor marks
beyond the linear terms.  Suppose also that logarithmic density
roots have total perturbative grade \(q\), and each grade-\(r\)
root has degree at most \(2r+2\).  A nonzero connected Gaussian
coefficient must satisfy
\begin{equation}
 \rho+2q\ge m-2.
 \label{eq:newV49-grade-count}
\end{equation}
\end{corollary}

\begin{proof}
If there are \(\ell\) density-root vertices of grades summing to
\(q\), their total degree is at most \(2q+2\ell\).
The total number of vertices is \(m+\ell\).  The contrapositive of
\Cref{lem:newV49-leg-count} gives
\[
 m+\rho+2q+2\ell\ge2(m+\ell-1),
\]
which is exactly \eqref{eq:newV49-grade-count}.
\end{proof}

\section{Endpoint logarithmic density grades}
\label{sec:newV49-endpoints}

\begin{lemma}[No endpoint eighth cumulant through order two]
\label{lem:newV49-endpoint-eight}
For \(X=Y_s^1\),
\begin{equation}
 \kappa_8^H(X,\ldots,X)=O(s^3),
 \qquad
 \kappa_8^W(X,\ldots,X)=O(s^3).
 \label{eq:newV49-endpoint-eight}
\end{equation}
\end{lemma}

\begin{proof}
Relative to the standard Gaussian, the logarithm of the normalized
heat density has the finite expansion
\[
 sV_{H,1}(Y)+s^2V_{H,2}(Y)+O(s^3)
\]
in every fixed polynomially weighted moment, with
\[
 \deg V_{H,1}\le2,
 \qquad
 \deg V_{H,2}\le4.
\]
This follows directly from
\[
 \log{\sin(\sqrt{2s}r)\over\sqrt{2s}r}
 =
 -{sr^2\over3}-{s^2r^4\over45}+O(s^3r^6),
\]
with scalar normalization terms omitted.

For the Wilson density, the corresponding logarithm is
\[
 \alpha(\cos(\sqrt{2s}r)-1)+{r^2\over2}
 +2\log{\sin(\sqrt{2s}r)\over\sqrt{2s}r}
 -\hbox{normalization},
\]
so
\[
 \deg V_{W,1}\le4,
 \qquad
 \deg V_{W,2}\le6.
\]

The order-zero eighth cumulant of a Gaussian linear coordinate is
zero.  At order \(s\), there is one root of degree at most four,
which cannot connect eight linear leaves by
\Cref{lem:newV49-leg-count}.  At order \(s^2\), the heat terms
contain either one degree-four root or two degree-two roots; the
Wilson terms contain either one degree-six root or two degree-four
roots.  In the largest case, two degree-four roots and eight
linear leaves provide \(16\) legs on \(10\) vertices, fewer than
the \(18\) legs required for a connected Wick graph.  All
coefficients through order two vanish.  The fixed-order density
expansions and V30 tails bound the remainder by \(O(s^3)\).
\end{proof}

\section{The endpoint cumulant difference}
\label{sec:newV49-deltaK}

Let
\[
 K_H(u;s):=\log\mathbb E_H e^{uY_s^1},
 \qquad
 K_W(u;s):=\log\mathbb E_W e^{uY_s^1},
\]
and put \(\Delta K=K_W-K_H\).

\begin{lemma}[First endpoint cumulant difference]
\label{lem:newV49-deltaK}
Uniformly for \(u\) in every fixed complex disk,
\begin{equation}
 \Delta K(u;s)
 =
 {s\over12}(u^4+10u^2)+O(s^2).
 \label{eq:newV49-deltaK}
\end{equation}
\end{lemma}

\begin{proof}
By \Cref{thm:newV48-density-expansion}, the first difference of
the endpoint densities is
\[
 s(P_W-P_H)
 =
 {s\over12}(|Y|^4-15).
\]
The first variation of the logarithmic moment generating function
is its centered expectation under the Gaussian exponential tilt:
\[
 {s\over12}
 {\mathbb E_\gamma[(|Y|^4-15)e^{uY_1}]
  \over\mathbb E_\gamma e^{uY_1}}.
\]
Gaussian translation sends \(Y\) to \(Y+ue_1\), and V25 computed
\[
 \mathbb E_\gamma(|Y+ue_1|^4-15)=u^4+10u^2.
\]
Analytic domination on a fixed \(u\)-disk controls the remainder
by \(O(s^2)\).
\end{proof}

\section{The mixture logarithm creates a degree-eight root}
\label{sec:newV49-mixture}

For the additive mixture,
\begin{equation}
 K_t^{\rm add}(u;s)
 =
 \log\left[
 (1-t)e^{K_H(u;s)}+te^{K_W(u;s)}
 \right].
 \label{eq:newV49-Kmix}
\end{equation}

\begin{lemma}[Exact mixture-logarithm expansion]
\label{lem:newV49-mixture-log}
\begin{align}
 K_t^{\rm add}
 &=
 K_H+t\Delta K
 +{t(1-t)\over2}(\Delta K)^2
 +O((\Delta K)^3)                                            \label{eq:newV49-mixture-log}\\
 &=
 K_H+t\Delta K
 +{t(1-t)s^2\over288}u^8
 +\{\hbox{terms of \(u\)-degree at most six at order \(s^2\)}\}
 +O(s^3).
 \label{eq:newV49-u8}
\end{align}
\end{lemma}

\begin{proof}
Factor \(e^{K_H}\) in \eqref{eq:newV49-Kmix}:
\[
 K_t^{\rm add}
 =
 K_H+\log(1-t+te^{\Delta K}).
\]
Taylor expansion at \(\Delta K=0\) gives the first line.  Insert
\eqref{eq:newV49-deltaK}.  The only degree-eight term at order
\(s^2\) is the square of \(su^4/12\), with coefficient
\(t(1-t)/2\), which is \(t(1-t)s^2u^8/288\).
\end{proof}

\begin{theorem}[Nonzero eighth scaled-coordinate cumulant]
\label{thm:newV49-eight-Y}
Uniformly for \(t\in[0,1]\),
\begin{equation}
 \kappa_8^{\,\mu_{\alpha,s,t}^{\rm add}}
 (Y_s^1,\ldots,Y_s^1)
 =
 140\,t(1-t)s^2+O(s^3).
 \label{eq:newV49-eight-Y}
\end{equation}
\end{theorem}

\begin{proof}
The eighth cumulant is
\(\partial_u^8K_t^{\rm add}(0;s)\).
By \Cref{lem:newV49-endpoint-eight}, neither endpoint contributes
at order \(s^2\).  The term \(t\Delta K\) likewise has no
degree-eight order-\(s^2\) coefficient because it is the
difference of the endpoint cumulant generators.  Thus
\eqref{eq:newV49-u8} is the unique contribution, and
\[
 {8!\over288}=140.
\]
\end{proof}

\section{Transfer to a bounded fundamental matrix functional}
\label{sec:newV49-transfer}

Use the V25 bounded fundamental functional
\[
 g(U):={1\over2i}\operatorname{tr}(\sigma_1U)
 =n_1\sin\theta.
\]

\begin{lemma}[Eighth-cumulant tangent transfer]
\label{lem:newV49-transfer}
\begin{equation}
 \kappa_8^{\,\mu_{\alpha,s,t}^{\rm add}}
 (g,\ldots,g)
 =
 2240\,t(1-t)s^6+O(s^7).
 \label{eq:newV49-transfer}
\end{equation}
\end{lemma}

\begin{proof}
Since
\[
 g
 =
 \sqrt{2s}\,Y_s^1
 +O\left(s^{3/2}|Y_s|^3\right),
\]
the all-linear term is
\[
 (2s)^4
 \left[
 140t(1-t)s^2+O(s^3)
 \right],
\]
which gives the displayed leading coefficient.

It remains to bound every term containing a cubic remainder.
If \(k\) of the eight entries use that remainder, the exterior
scale is \(s^{4+k}\).  For \(k=1\), the input polynomial degree is
ten; a Gaussian connected graph on eight input vertices is
impossible, and one first density root of degree four is still
insufficient by \Cref{lem:newV49-leg-count}.  Hence this term
starts at \(s^{4+1+2}=s^7\).  For \(k=2\), the input degree is
twelve; its zeroth-order Gaussian cumulant vanishes, so it starts
at \(s^{4+2+1}=s^7\).  For \(k\ge3\), the exterior scale alone is
at least \(s^7\).  Fixed-order Taylor remainders and V30 moments
sum these finitely many terms to \(O(s^7)\).
\end{proof}

\begin{theorem}[The additive all-arity raw card is false]
\label{thm:newV49-additive-false}
The unmarked additive-bridge raw strong card is false at
\(m=8\), already for eight fundamental inputs.  The marked
additive score card is also false at \(m=8\).
\end{theorem}

\begin{proof}
Apply the fixed norm-one functional
\[
 A\longmapsto{1\over2i}\operatorname{tr}(\sigma_1A)
\]
to every fundamental tensor leg.  At \(t=1/2\),
\Cref{lem:newV49-transfer} gives a scalar lower bound
\[
 \left|\kappa_8(g,\ldots,g)\right|
 \ge cs^6
\]
for all sufficiently small \(s\).  The proposed eight-input
Pr\"ufer target is
\[
 K^8b_{\rm f}^8(8b_{\rm f})^6
 =
 C_Ks^7,
\]
because \(b_{\rm f}^2\asymp s\).  No \(K\) independent of \(s\)
can dominate the lower bound.

Differentiate \eqref{eq:newV49-transfer} in \(t\).  The remainder
remains \(O(s^7)\), either directly or by the degree-eight
Banach-valued Markov inequality.  At \(t=1/4\),
\[
 \kappa_9(g,\ldots,g,\Sigma_{\alpha,t}^{\rm add})
 =
 1120s^6+O(s^7),
\]
again larger than the \(Cs^7\) target.
\end{proof}

\begin{correction}[Retirement of the additive all-arity route]
\label{corr:newV49-additive-retired}
The binary through six-input theorems of V44--V48 remain valid.
The conditional implication
\(\mathrm{GRMSC}\Rightarrow\) raw strong card in V44 also remains
logically valid.  However,
\Cref{thm:newV49-additive-false} proves that GRMSC and EMCC cannot
hold with the stated strong target for the additive bridge.  The
additive bridge is therefore retired as an all-arity comparison
path.
\end{correction}

\section{Return upstream to a log-affine bridge}
\label{sec:newV49-upstream}

\begin{assessment}[Why the geometric bridge avoids this new root]
\label{ass:newV49-geometric}
The additive obstruction is the Bernoulli-mixture term
\[
 {t(1-t)\over2}(\Delta K)^2.
\]
For the geometric bridge, the unnormalized physical log density is
affine:
\[
 (1-t)\log h_s+t\log w_\alpha.
\]
Its scaled density roots therefore have the endpoint grade
\(\deg V_r\le2r+2\), with no degree-eight grade-two mixture square.
The physical complex zeros found in V40 obstruct a global
all-spin Fourier normalization, but V42 correctly leaves the
thermal \emph{connected} geometric card open.  The next attack
returns to that connected log-affine object.
\end{assessment}

\begin{definition}[Thermal connected log-affine card]
\label{def:newV49-TCLAC}
TCLAC is an exponential-arity bound for the complete
thermal-input geometric-bridge source logarithm obtained from:
\begin{enumerate}[label=\textup{(\roman*)},leftmargin=3.5em]
\item endpoint log-density roots of grade \(r\) and degree at most
      \(2r+2\);
\item exact odd/even representation vertices;
\item the connected Gaussian leg inequality
      \(\rho+2q\ge m-2\); and
\item the V28--V29 stable time-ordered Casimir forest,
\end{enumerate}
with no inverse-heat norm on an isolated high-spin moment.
\end{definition}

\begin{programme}[V50: audit and geometric grade expansion]
\label{prog:newV49-V50}
\begin{enumerate}[leftmargin=3em]
\item Perform the compulsory read-only audit of the newest active
      SM--Einstein and Navier--Stokes notes.
\item Derive the normalized geometric scaled log-density through
      arbitrary finite grade and prove
      \(\deg V_r\le2r+2\) uniformly in \(t\).
\item Combine that degree theorem with
      \Cref{cor:newV49-grade-count} to prove the correct thermal
      homogeneity at every fixed arity.
\item Identify the coefficient-growth condition needed to replace
      fixed-arity constants by one exponential base \(K^m\).
\item Use V27--V29 only after the complete grade roots have been
      assembled.
\end{enumerate}
\end{programme}

\begin{firewall}[Claim boundary after V49]
\label{fire:newV49-boundary}
V49 proves the connected Gaussian leg inequality, computes the
unique degree-eight order-\(s^2\) root created by the additive
mixture, derives the exact leading eighth scaled and fundamental
cumulants, and disproves both the unmarked and marked additive
all-arity strong cards.  It preserves the proved additive cards
through arity six and returns the programme to the connected
log-affine geometric bridge.  It does not prove TCLAC, thermal
connected WPRAC, all-arity RGSC, the raw Wilson aggregate strong
card, all-order strong MTA, WUFC, CPM, cutoff removal, reflection
positivity, Osterwalder--Schrader reconstruction, a continuum
Yang--Mills measure, or a positive Hamiltonian gap.  The full
Yang--Mills mass gap has not been proved.
\end{firewall}

\begin{theorem}[V49 result ledger]
\label{thm:newV49-results}
V49:
\begin{enumerate}[label=\textup{(V49.\arabic*)},leftmargin=6.5em]
\item proves the general connected Wick leg inequality;
\item derives the grade/extra-mark condition
      \(\rho+2q\ge m-2\);
\item proves both endpoint eighth linear cumulants vanish through
      order \(s^2\);
\item computes
      \(\Delta K=s(u^4+10u^2)/12+O(s^2)\);
\item isolates the additive degree-eight mixture root;
\item proves
      \(\kappa_8(g,\ldots,g)
       =2240t(1-t)s^6+O(s^7)\);
\item disproves the additive unmarked and marked strong cards at
      arity eight; and
\item defines TCLAC as the corrected log-affine route.
\end{enumerate}
\end{theorem}

\begin{proof}
Use
\Cref{lem:newV49-leg-count,
cor:newV49-grade-count,
lem:newV49-endpoint-eight,
lem:newV49-deltaK,
lem:newV49-mixture-log,
thm:newV49-eight-Y,
lem:newV49-transfer,
thm:newV49-additive-false}.
\end{proof}

\paragraph{Parent-version record.}
V49 was formed from
\texttt{Renewal\_Yang\_Mills\_Mass\_Gap\_Research\_Notes\_V48.tex}.
The V48 parent had \(23\,148\) logical source lines,
\(728\,357\) bytes, and SHA--256
\[
 \texttt{180af15cee5cc06df2646b7f6e007f0f89db24210974f99aa2573c4303d5f38c}.
\]
The next compulsory companion audit is V50.

% =============================================================================
% END APPENDED FIFTH-SERIES V49 MATERIAL
% =============================================================================

% =============================================================================
% BEGIN APPENDED FIFTH-SERIES V50 MATERIAL
% =============================================================================

\chapter{V50: Companion audit and fixed-arity geometric
log-density closure}
\label{chap:newV50-fixed-arity}

\section{Compulsory companion audit}
\label{sec:newV50-audit}

\begin{audit}[Read-only SM--Einstein and Navier--Stokes audit]
\label{audit:newV50-companions}
The newest active companion notes inspected before this append
were:
\begin{align*}
 &\texttt{Renewal\_SM\_Einstein\_Continuum\_Research\_Notes\_V68.tex},
 \\
 &\qquad 22\,282\ \hbox{logical source lines},\quad
 856\,963\ \hbox{bytes},                                      \\
 &\qquad
 \operatorname{SHA256}
 =
 \texttt{28ee777da79529dc69eb877f16f20732820b51cad69ca0e4deaf4b128423ae26},
 \\
 &\texttt{Renewal\_NS\_Stability\_Research\_Notes\_V50.tex},
 \\
 &\qquad 30\,285\ \hbox{logical source lines},\quad
 1\,233\,578\ \hbox{bytes},                                   \\
 &\qquad
 \operatorname{SHA256}
 =
 \texttt{a6ff0ce7d0b350ec67a27b2c686b5c8126084712b751282036adc2d0e268091e}.
\end{align*}
Neither file was modified.

SM V68 proves that a positive sharp-time
Osterwalder--Schrader norm diverges by a local boundary contact.
Subtracting the divergent scalar quadratic form produces a finite
distributional part but does not produce a norm-convergent Hilbert
vector and does not preserve positivity automatically.  The
constructive object must remain positively time-smeared and be
related to the Hamiltonian through a semigroup or a closed form.

NS V50 consolidates the type-I ancient-limit reduction and
separates concentration and dispersion branches.  Its Liouville
gate remains open and contains the stationary Leray problem; the
remaining possible gain lies in the ancestry and trace properties
of ancient solutions actually generated by the blow-up, not in an
arbitrary member of the ambient class.
\end{audit}

\begin{assessment}[Audit transfer discipline]
\label{ass:newV50-audit-transfer}
No SM boundary coefficient and no NS ancient-solution estimate is
transferred to Yang--Mills.  The applicable proof discipline is:
\begin{enumerate}[label=\textup{(\roman*)},leftmargin=3.5em]
\item do not replace a failed complete source vector or operator
      by a scalar finite-part subtraction; and
\item exploit the exact geometric heat--Wilson ancestry rather
      than ask for a theorem on every abstract sub-Gaussian central
      law.
\end{enumerate}
\end{assessment}

\section{Finite-grade endpoint logarithms}
\label{sec:newV50-log-grades}

Let \(\Gamma(dY)=\gamma_3(dY)\) and write the two scaled endpoint
laws as
\[
 d\mu_{\mu,s}(Y)
 =
 {\exp(\mathcal V_{\mu,s}(Y))\over
  \mathbb E_\Gamma e^{\mathcal V_{\mu,s}}}\,\Gamma(dY),
 \qquad \mu\in\{H,W\},
\]
on the principal scaled ball, with the wrapped heat images included
in \(\mathcal V_{H,s}\).  The tails outside that ball are understood
in the moment sense used in V46--V48.

\begin{lemma}[Heat logarithmic grade]
\label{lem:newV50-heat-grade}
For every fixed \(Q\), there are even radial polynomials
\(V_{H,r}\), \(1\le r\le Q\), such that
\begin{equation}
 \mathcal V_{H,s}(Y)
 =
 \sum_{r=1}^Qs^rV_{H,r}(Y)
 +\mathcal E_{H,Q,s}(Y),
 \label{eq:newV50-heat-grade}
\end{equation}
where
\begin{equation}
 \deg V_{H,r}\le2r
 \label{eq:newV50-heat-degree}
\end{equation}
and, for every fixed polynomial \(F\),
\begin{equation}
 \mathbb E_\Gamma
 |F\mathcal E_{H,Q,s}|
 \le C_{F,Q}s^{Q+1}.
 \label{eq:newV50-heat-remainder}
\end{equation}
\end{lemma}

\begin{proof}
On the inner scaled region, the shortest heat image gives
\[
 \mathcal V_{H,s}(Y)
 =
 \log{\sin(\sqrt{2s}|Y|)\over\sqrt{2s}|Y|}
 +O(e^{-c/s})
\]
up to a scalar normalization.  The Taylor series of
\(\log(\sin z/z)\) contains only \(z^{2r}\).  Substitution
\(z=\sqrt{2s}|Y|\) proves
\eqref{eq:newV50-heat-degree} and Taylor's theorem gives the
weighted remainder.  On any fixed inner fraction, the remote image
term and all its fixed derivatives are exponentially small by
V31.  The paired cap and V30 concentration make the remaining real
moment contribution \(O(s^N)\) for every fixed \(N\).  Scalar
normalization terms do not change polynomial degree.
\end{proof}

\begin{lemma}[Wilson logarithmic grade]
\label{lem:newV50-Wilson-grade}
For every fixed \(Q\),
\begin{equation}
 \mathcal V_{W,s}(Y)
 =
 \sum_{r=1}^Qs^rV_{W,r}(Y)
 +\mathcal E_{W,Q,s}(Y),
 \label{eq:newV50-Wilson-grade}
\end{equation}
where the \(V_{W,r}\) are even radial polynomials satisfying
\begin{equation}
 \deg V_{W,r}\le2r+2,
 \label{eq:newV50-Wilson-degree}
\end{equation}
and the analogue of
\eqref{eq:newV50-heat-remainder} holds.
\end{lemma}

\begin{proof}
Relative to the standard Gaussian exponent,
\[
 \mathcal V_{W,s}(Y)
 =
 \alpha(\cos(\sqrt{2s}|Y|)-1)+{|Y|^2\over2}
 +2\log{\sin(\sqrt{2s}|Y|)\over\sqrt{2s}|Y|},
\]
up to normalization, with
\(\alpha=(2s)^{-1}-1/3\).
The cosine term of power \(|Y|^{2k}\) has grade \(s^{k-1}\);
hence grade \(r\) has degree at most \(2r+2\).  The Haar logarithm
at grade \(r\) has degree \(2r\).  Taylor's theorem under the V30
Gaussian tail bound proves the fixed weighted remainder.
\end{proof}

\section{The geometric log-affine grade}
\label{sec:newV50-geometric}

The geometric bridge density is
\[
 d\nu_{\alpha,s,t}^{\rm geo}
 =
 {\exp(\mathcal V_{t,s}^{\rm geo})\over
  \mathbb E_\Gamma e^{\mathcal V_{t,s}^{\rm geo}}}\,d\Gamma,
\qquad
 \mathcal V_{t,s}^{\rm geo}
 =
 (1-t)\mathcal V_{H,s}+t\mathcal V_{W,s}.
\]

\begin{theorem}[Uniform finite-grade geometric expansion]
\label{thm:newV50-geometric-grade}
For every fixed \(Q\),
\begin{equation}
 \mathcal V_{t,s}^{\rm geo}(Y)
 =
 \sum_{r=1}^Qs^rV_{t,r}(Y)
 +\mathcal E_{t,Q,s}(Y),
 \qquad
 V_{t,r}=(1-t)V_{H,r}+tV_{W,r},
 \label{eq:newV50-geometric-grade}
\end{equation}
uniformly in \(t\), with
\begin{equation}
 \deg V_{t,r}\le2r+2
 \label{eq:newV50-geometric-degree}
\end{equation}
and the fixed polynomially weighted remainder
\[
 \mathbb E_\Gamma|F\mathcal E_{t,Q,s}|
 \le C_{F,Q}s^{Q+1}.
\]
Differentiation in \(t\) preserves the same grade and degree
bounds.
\end{theorem}

\begin{proof}
Take the affine combination of
\Cref{lem:newV50-heat-grade,lem:newV50-Wilson-grade}.
Uniformity and the derivative statement are immediate.
\end{proof}

\begin{corollary}[Absence of the additive degree-eight defect]
\label{cor:newV50-no-mixture-square}
At geometric grade two, every single log-density root has degree at
most six.  In particular, the degree-eight root
\[
 {t(1-t)\over2}(P_W-P_H)^2
\]
which caused the V49 additive obstruction is absent.
\end{corollary}

\begin{proof}
The geometric unnormalized log density is affine in \(t\), whereas
the displayed square came from the logarithm of a convex mixture.
Apply \eqref{eq:newV50-geometric-degree} with \(r=2\).
\end{proof}

\section{Connected grade expansion}
\label{sec:newV50-connected}

\begin{lemma}[Cumulant coefficients are connected rooted Wick
graphs]
\label{lem:newV50-rooted-Wick}
Fix polynomial input vertices \(A_1(Y),\ldots,A_m(Y)\).  The
coefficient of \(s^q\) in their cumulant under the formal
geometric density
\[
 \exp\left(\sum_{r\ge1}s^rV_{t,r}(Y)\right)d\Gamma
\]
is a finite sum of Gaussian joint cumulants containing
\(\ell\ge0\) root vertices
\[
 V_{t,r_1},\ldots,V_{t,r_\ell},
 \qquad
 r_1+\cdots+r_\ell=q,
\]
with the usual multiplicity factors.  Every surviving Wick
contraction is connected on the \(m+\ell\) labelled vertices.
\end{lemma}

\begin{proof}
Introduce auxiliary sources \(\lambda_r\) for every density root
and \(z_i\) for every input.  The desired cumulant is the
\(z_1\cdots z_m\) coefficient of the logarithm of the normalized
Gaussian source.  Expanding in the \(\lambda_r=s^r\) variables
gives Gaussian joint cumulants of the stated vertices.  By the
moment--cumulant cancellation used in
\Cref{lem:newV49-leg-count}, only connected Wick contraction
graphs survive.
\end{proof}

\begin{theorem}[Geometric grade/mark inequality]
\label{thm:newV50-grade-mark}
Suppose input \(i\) uses a representation Taylor term of degree
\(k_i\ge1\), and put
\[
 \rho:=\sum_{i=1}^m(k_i-1).
\]
If its geometric density coefficient of total grade \(q\) is
nonzero, then
\begin{equation}
 \boxed{\qquad \rho+2q\ge m-2. \qquad}
 \label{eq:newV50-grade-mark}
\end{equation}
The same statement holds after one \(t\)-derivative, hence for one
geometric score insertion.
\end{theorem}

\begin{proof}
Let the coefficient contain \(\ell\) root vertices of grades
\(r_a\) summing to \(q\).  By
\eqref{eq:newV50-geometric-degree}, their total polynomial degree
is at most
\[
 \sum_{a=1}^\ell(2r_a+2)=2q+2\ell.
\]
The input degree is \(m+\rho\).  A connected Wick graph on
\(m+\ell\) vertices requires at least
\(2(m+\ell-1)\) legs by
\Cref{lem:newV49-leg-count}.  Therefore
\[
 m+\rho+2q+2\ell\ge2m+2\ell-2,
\]
which rearranges to the displayed inequality.  Differentiating in
\(t\) replaces one root by \(V_{W,r}-V_{H,r}\) and does not change
its grade or degree.
\end{proof}

\section{Closure at every fixed arity}
\label{sec:newV50-fixed}

\begin{theorem}[Fixed-arity geometric homogeneity card]
\label{thm:newV50-fixed-card}
For every fixed \(m\ge2\), there is a finite constant \(C_m\),
independent of \(\alpha,t\) and the thermal representations, such
that
\begin{align}
 \left\|
 \kappa_m^{\,\nu_{\alpha,s,t}^{\rm geo}}
 (D^{R_1},\ldots,D^{R_m})
 \right\|_{\rm cb}
 &\le
 C_m\left(\prod_{i=1}^mb_i\right)
 B^{m-2},                                                     \label{eq:newV50-fixed-unmarked}\\
 \left\|
 \kappa_{m+1}^{\,\nu_{\alpha,s,t}^{\rm geo}}
 (D^{R_1},\ldots,D^{R_m},S_{\alpha,t}^{\rm geo})
 \right\|_{\rm cb}
 &\le
 C_m\left(\prod_{i=1}^mb_i\right)
 B^{m-2},                                                     \label{eq:newV50-fixed-marked}
\end{align}
where \(B=\sum_i b_i\).
\end{theorem}

\begin{proof}
First suppose \(B\ge1/2\).  Center every input at the identity.
The V44 operator-source argument, used at this fixed \(m\), gives
a bound \(C_m\prod_i b_i\); the fixed lower bound on \(B\)
absorbs \(B^{m-2}\).  For the geometric marked row, use the V32
fixed-\(p\) score bounds with a fixed H\"older exponent depending
only on \(m\).

Now suppose \(B<1/2\).  Expand every exact representation
increment in its geodesic Taylor series and expand the complete
geometric log density to a finite grade larger than \(m\), using
\Cref{thm:newV50-geometric-grade}.  Gaussian tails justify the
finite termwise cumulant expansion; choose the Taylor and density
remainders beyond the target degree, so their fixed-\(m\) moments
obey the same or a smaller bound.

A term with input excess degree \(\rho\) carries
\[
 \left(\prod_i b_i\right)\times
 \{\hbox{a monomial of total \(b\)-degree \(\rho\)}\},
\]
which is at most
\((\prod_i b_i)B^\rho\).  Density grade \(q\) carries \(s^q\).
Since every \(b_i\ge\sqrt s\),
\[
 s^q\le\left({B\over m}\right)^{2q}\le B^{2q}.
\]
By \Cref{thm:newV50-grade-mark}, every nonzero connected term has
\(\rho+2q\ge m-2\).  As \(B<1\),
\[
 B^{\rho+2q}\le B^{m-2}.
\]
There are finitely many retained terms at fixed \(m\); their
Gaussian moments, coordinate sums, and partition factors form
\(C_m\).  The \(t\)-differentiated expansion has the same
grade/degree count, proving the marked row.
\end{proof}

\begin{corollary}[The geometric arity-eight obstruction is absent]
\label{cor:newV50-eight}
The eight-input geometric bridge cumulant and its score derivative
obey the strong thermal homogeneity
\[
 C_8\left(\prod_{i=1}^8b_i\right)B^6.
\]
Thus the V49 counterexample is specific to the additive mixture,
not to the two endpoints or the log-affine geometric path.
\end{corollary}

\begin{proof}
Set \(m=8\) in
\Cref{thm:newV50-fixed-card}.
\end{proof}

\section{The remaining uniform-arity gate}
\label{sec:newV50-uniform}

\begin{warning}[Fixed-arity closure is not all-arity closure]
\label{warn:newV50-not-uniform}
The proof of \Cref{thm:newV50-fixed-card} permits \(C_m\) to absorb
the number of density roots, representation Taylor terms, Gaussian
pairings, and partition coefficients.  It does not prove
\[
 C_m\le K^m.
\]
Nor may the expansion grade be sent to infinity using the
fixed-order image separation of V31: V40 quantifies the collapse
of that separation at derivative order
\(O((s\log(1/s))^{-1})\).
\end{warning}

\begin{definition}[Uniform geometric connected-grade card]
\label{def:newV50-UGCGC}
UGCGC is a resummation of the complete wrapped geometric source
which:
\begin{enumerate}[label=\textup{(\roman*)},leftmargin=3.5em]
\item preserves the root grade/degree inequality
      \(\rho+2q\ge m-2\);
\item sums all Gaussian pair contractions by the V28--V29 stable
      time-ordered forest before norms;
\item keeps the paired heat images assembled beyond the V40
      derivative window; and
\item replaces \(C_m\) in
      \eqref{eq:newV50-fixed-unmarked}--%
      \eqref{eq:newV50-fixed-marked} by one \(K^m\).
\end{enumerate}
\end{definition}

\begin{programme}[V51: coefficient-growth ledger]
\label{prog:newV50-V51}
\begin{enumerate}[leftmargin=3em]
\item Compute the exact coefficient bounds of the endpoint
      logarithmic roots \(V_{\mu,r}\) on the inner scaled chart.
\item Separate root-grade growth from Wick graph counting and
      representation Taylor factorials.
\item Insert the complete grade generating function into the V29
      time-ordered forest rather than summing connected pairings
      directly.
\item Identify whether the only nonuniformity is the remote-image
      transition; if so, formulate one wrapped real-integral
      estimate which bypasses derivative separation.
\item Preserve the fixed-arity theorem independently of the
      all-arity outcome.
\end{enumerate}
\end{programme}

\begin{firewall}[Claim boundary after V50]
\label{fire:newV50-boundary}
V50 completes the compulsory audit, proves the finite-grade heat,
Wilson, and geometric log-density degree bounds, proves the
connected grade/mark inequality, and closes the correct thermal
homogeneity for every fixed arity on the geometric bridge,
including arity eight.  It does not control the constants by one
exponential base and therefore does not prove UGCGC, thermal
connected WPRAC, all-arity RGSC, the raw Wilson aggregate strong
card, all-order strong MTA, WUFC, CPM, cutoff removal, reflection
positivity, Osterwalder--Schrader reconstruction, a continuum
Yang--Mills measure, or a positive Hamiltonian gap.  The full
Yang--Mills mass gap has not been proved.
\end{firewall}

\begin{theorem}[V50 result ledger]
\label{thm:newV50-results}
V50:
\begin{enumerate}[label=\textup{(V50.\arabic*)},leftmargin=6.5em]
\item completes the required V50 companion audit;
\item proves heat log-density grade \(r\) has degree at most
      \(2r\);
\item proves Wilson and geometric grade \(r\) have degree at most
      \(2r+2\);
\item proves the additive degree-eight square is absent from the
      geometric path;
\item proves the connected rooted-Wick coefficient formula;
\item proves the universal grade/mark inequality
      \(\rho+2q\ge m-2\);
\item proves the unmarked and marked thermal homogeneity cards at
      every fixed arity; and
\item isolates UGCGC as the remaining uniform-arity gate.
\end{enumerate}
\end{theorem}

\begin{proof}
Use
\Cref{audit:newV50-companions,
lem:newV50-heat-grade,
lem:newV50-Wilson-grade,
thm:newV50-geometric-grade,
cor:newV50-no-mixture-square,
lem:newV50-rooted-Wick,
thm:newV50-grade-mark,
thm:newV50-fixed-card,
cor:newV50-eight}.
\end{proof}

\paragraph{Parent-version record.}
V50 was formed from
\texttt{Renewal\_Yang\_Mills\_Mass\_Gap\_Research\_Notes\_V49.tex}.
The V49 parent had \(23\,603\) logical source lines,
\(742\,175\) bytes, and SHA--256
\[
 \texttt{07d56844fa335b237136118d4d933fc681cf814933f04fb4ac3487f1a04db30d}.
\]
The next compulsory companion audit is V55.

% =============================================================================
% END APPENDED FIFTH-SERIES V50 MATERIAL
% =============================================================================

\clearpage
\chapter{Fifth-series V51: exact Hartman--Watson heat subordination of
the Wilson endpoint}
\label{ch:newV51-HW}

\section{Reason for changing the V51 attack}
\label{sec:newV51-context}

\begin{correction}[Inherited fifth-series preamble repair]
\label{corr:newV51-preamble}
The fifth-series preamble declared the common theorem counter but
omitted declarations for the already-used environments
\texttt{audit}, \texttt{correction}, and \texttt{warning}.  V51
adds those three declarations.  This changes no mathematical
statement; it makes the accumulated source self-contained.
\end{correction}

V50 proposed a direct coefficient-growth ledger for the geometric
log-affine bridge.  That remains a valid route to UGCGC, but an
upstream exact identity is available before any inner-chart
coefficient is estimated.  The Wilson Fourier multiplier
\[
 q_{\alpha,n}={I_{n+1}(\alpha)\over I_1(\alpha)}
\]
is a Laplace transform of the \(\SU(2)\) heat multiplier
\(\exp\{-S n(n+2)\}\) against a positive probability law.  Thus the
full Wilson density, including all wrapped images and every
representation, is an exact positive mixture of heat kernels.

This chapter imports only the classical scalar Hartman--Watson
transform theorem.  The tilt, the Wilson multiplier identity, the
Peter--Weyl density identity, the tensor-block functional calculus,
and the forest consequences are proved below.  The same calculation
also reveals a genuine obstruction: the random heat time has a
square-root moment-generating singularity at \(1\).  Consequently
one may not take norms in the Duhamel forest and replace every
random-time moment by a thermal-scale power.  The exact
subordination is therefore a complete-object route, not a license
for a termwise moment estimate.

\section{The imported scalar transform and its Wilson tilt}
\label{sec:newV51-transform}

\begin{theorem}[Hartman--Watson transform and infinite
divisibility]
\label{ext:newV51-HW}
For every \(\alpha>0\), there is a probability measure
\(\mu_\alpha\) on \((0,\infty)\) such that
\begin{equation}
 \int_0^\infty e^{-\nu^2T/2}\,\mu_\alpha(dT)
 = {I_{|\nu|}(\alpha)\over I_0(\alpha)}
 \qquad(\nu\in\mathbb R).
 \label{eq:newV51-HW-transform}
\end{equation}
Moreover, \(\mu_\alpha\) is infinitely divisible and
\begin{equation}
 \psi_\alpha(u):=
 -\log {I_{\sqrt{2u}}(\alpha)\over I_0(\alpha)}
 \label{eq:newV51-psi}
\end{equation}
is a complete Bernstein function on \([0,\infty)\).
\end{theorem}

\paragraph{Source boundary.}
The transform \eqref{eq:newV51-HW-transform} is the
Hartman--Watson theorem: P.~Hartman and G.~S.~Watson,
\emph{Normal distribution functions on spheres and the modified
Bessel functions}, Ann.\ Probab.\ \textbf{2} (1974), 593--607,
DOI \(10.1214/\mathrm{aop}/1176996606\).
Infinite divisibility is proved in P.~Hartman,
\emph{Completely monotone families of solutions of \(n\)-th order
linear differential equations and infinitely divisible
distributions}, Ann.\ Scuola Norm.\ Sup.\ Pisa Cl.\ Sci.\
\textbf{3} (1976), 267--287.  The complete-Bernstein formulation
is a standard consequence of the Bondesson-class refinement.  These
classical scalar statements are imported; none is being claimed as
a new proof in this manuscript.

\begin{definition}[Tilted Hartman--Watson heat time]
\label{def:newV51-tilt}
Put
\[
 Z_\alpha:={I_1(\alpha)\over I_0(\alpha)}
 =\int_0^\infty e^{-T/2}\,\mu_\alpha(dT).
\]
Define the probability measure \(\rho_\alpha\) on \((0,\infty)\)
by
\begin{equation}
 \int f(S)\,\rho_\alpha(dS)
 :=
 {1\over Z_\alpha}
 \int f(T/2)e^{-T/2}\,\mu_\alpha(dT)
 \label{eq:newV51-rho}
\end{equation}
for bounded Borel \(f\).
\end{definition}

\begin{lemma}[Exact Laplace transform of the tilted time]
\label{lem:newV51-Laplace}
For every \(\lambda\ge0\),
\begin{equation}
 L_\alpha(\lambda):=
 \int_0^\infty e^{-\lambda S}\,\rho_\alpha(dS)
 =
 {I_{\sqrt{1+\lambda}}(\alpha)\over I_1(\alpha)}.
 \label{eq:newV51-Laplace}
\end{equation}
In particular, \(\rho_\alpha\) is a probability measure and has
moments of every finite order.
\end{lemma}

\begin{proof}
The normalization in \eqref{eq:newV51-rho} follows from
\eqref{eq:newV51-HW-transform} at \(\nu=1\).  For
\(\lambda\ge0\),
\[
 \int e^{-\lambda S}\rho_\alpha(dS)
 ={1\over Z_\alpha}
 \int e^{-(1+\lambda)T/2}\mu_\alpha(dT)
 ={I_{\sqrt{1+\lambda}}(\alpha)\over I_1(\alpha)}.
\]
For every \(0<a<1\), the same calculation with
\(\lambda=-a\) gives a finite exponential moment
\(\int e^{aS}\rho_\alpha(dS)\).  Hence all polynomial moments are
finite.
\end{proof}

\begin{theorem}[Exact Wilson multiplier subordination]
\label{thm:newV51-multiplier}
Let
\[
 c_n=n(n+2),\qquad n=0,1,2,\ldots ,
\]
be the V29 heat-Casimir normalization.  Then
\begin{equation}
 q_{\alpha,n}
 ={I_{n+1}(\alpha)\over I_1(\alpha)}
 =\int_0^\infty e^{-S c_n}\,\rho_\alpha(dS).
 \label{eq:newV51-q-mixture}
\end{equation}
\end{theorem}

\begin{proof}
Since \(1+c_n=(n+1)^2\), substitute
\(\lambda=c_n\) in \eqref{eq:newV51-Laplace}.
\end{proof}

\section{Complete compact-group and block-operator identities}
\label{sec:newV51-group}

\begin{theorem}[Exact positive heat mixture of the Wilson density]
\label{thm:newV51-density}
Let \(w_\alpha(U)dU\) be the normalized central Wilson probability
measure of V41 and let \(h_S(U)dU\) be normalized \(\SU(2)\) heat
kernel measure with character multipliers \(e^{-Sc_n}\).  Then
\begin{equation}
 w_\alpha(U)dU
 =
 \int_0^\infty h_S(U)dU\,\rho_\alpha(dS)
 \label{eq:newV51-density-mixture}
\end{equation}
as finite Borel measures on \(\SU(2)\).  In particular,
\eqref{eq:newV51-density-mixture} keeps the identity chart, the
antipodal chart, and every wrapped heat image assembled.
\end{theorem}

\begin{proof}
Both sides are central probability measures.  On the irreducible
representation of highest weight \(n\), the Fourier transform of
the right side is the scalar
\[
 \int e^{-Sc_n}\rho_\alpha(dS)=q_{\alpha,n}
\]
by \Cref{thm:newV51-multiplier}; this is the Fourier transform of
the left side by V41.  Matrix coefficients of irreducible
representations span a uniformly dense subalgebra of
\(C(\SU(2))\).  Equality of all Fourier transforms therefore gives
equality of the measures.  Since the Wilson measure has a smooth
density, the heat mixture represents that density almost
everywhere as well.
\end{proof}

\begin{corollary}[Exact tensor-block moment]
\label{cor:newV51-block}
For a nonempty block \(B\), let
\[
 \mathcal D_B^a:=\sum_{i\in B}D_i^a,\qquad
 \mathcal C_B:=\sum_{a=1}^3(\mathcal D_B^a)^2\ge0
\]
on \(\bigotimes_{i\in B}V_{R_i}\).  The Wilson block moment is
\begin{equation}
 M_{\alpha,B}^{W}
 =
 \int_0^\infty e^{-S\mathcal C_B}\rho_\alpha(dS)
 =
 L_\alpha(\mathcal C_B).
 \label{eq:newV51-block}
\end{equation}
The equality holds in operator norm, at every matrix
amplification.
\end{corollary}

\begin{proof}
For fixed \(S\), the heat-kernel tensor moment is
\(e^{-S\mathcal C_B}\) by the V29 tensor heat identity.
Integrate \eqref{eq:newV51-density-mixture}.  The integrand is a
positive contraction, so the Bochner integral exists in the
finite-dimensional tensor space and commutes with every
amplification.
\end{proof}

\begin{warning}[Cumulants do not commute with mixing]
\label{warn:newV51-not-one-time}
Equation \eqref{eq:newV51-density-mixture} does not say that a
Wilson cumulant is the \(\rho_\alpha\)-average of a heat cumulant.
The exact moment--cumulant identity is
\begin{equation}
 \kappa_{\alpha,m}^{W}
 =
 \sum_{\pi\in\Pi_m}
 (-1)^{|\pi|-1}(|\pi|-1)!
 \prod_{B\in\pi}
 \int e^{-S_B\mathcal C_B}\rho_\alpha(dS_B).
 \label{eq:newV51-partition-times}
\end{equation}
Thus a partition has one independent time \(S_B\) per block.
Mixture-induced connected terms are exactly the terms lost if the
mixing integral is moved outside the cumulant.
\end{warning}

\begin{proof}
Insert \eqref{eq:newV51-block} in the operator
moment--cumulant formula.  Products belonging to disjoint blocks
act on disjoint tensor factors, so Fubini gives independent
variables \(S_B\).  No common time can be factored through the
M\"obius sum.
\end{proof}

\section{Bernstein functional calculus}
\label{sec:newV51-Bernstein}

\begin{definition}[Wilson subordination exponent]
\label{def:newV51-Phi}
For \(\lambda\ge0\), set
\begin{equation}
 \Phi_\alpha(\lambda)
 :=
 \psi_\alpha\left({1+\lambda\over2}\right)
 -\psi_\alpha\left({1\over2}\right).
 \label{eq:newV51-Phi}
\end{equation}
\end{definition}

\begin{theorem}[Exact Bernstein exponent]
\label{thm:newV51-Bernstein}
\(\Phi_\alpha\) is a Bernstein function,
\(\Phi_\alpha(0)=0\), and
\begin{align}
 L_\alpha(\lambda)&=e^{-\Phi_\alpha(\lambda)},                \label{eq:newV51-L-exp}\\
 M_{\alpha,B}^{W}&=e^{-\Phi_\alpha(\mathcal C_B)},             \label{eq:newV51-block-Phi}\\
 q_{\alpha,n}&=e^{-\Phi_\alpha(c_n)}.                          \label{eq:newV51-q-Phi}
\end{align}
Consequently there are \(d_\alpha\ge0\) and a positive measure
\(\Pi_\alpha\) satisfying
\(\int(1\wedge x)\Pi_\alpha(dx)<\infty\) such that
\begin{equation}
 \Phi_\alpha(\lambda)
 =
 d_\alpha\lambda+
 \int_{(0,\infty)}(1-e^{-\lambda x})\Pi_\alpha(dx).
 \label{eq:newV51-Levy}
\end{equation}
\end{theorem}

\begin{proof}
The first equality follows directly from
\eqref{eq:newV51-psi} and \eqref{eq:newV51-Laplace}.  Since
\(\psi_\alpha'\) is completely monotone,
\[
 \Phi_\alpha'(\lambda)
 ={1\over2}\psi_\alpha'\left({1+\lambda\over2}\right)
\]
is completely monotone.  Also \(\Phi_\alpha(0)=0\).  Hence
\(\Phi_\alpha\) is Bernstein.  Equations
\eqref{eq:newV51-block-Phi} and \eqref{eq:newV51-q-Phi} follow by
spectral calculus.  The final formula is the
L\'evy--Khintchine representation for a Bernstein function with no
killing term.
\end{proof}

\begin{corollary}[Exact subordinated heat semigroup]
\label{cor:newV51-subordinator}
There is a subordinator \((\mathsf S_t)_{t\ge0}\) with
\[
 \mathbb E e^{-\lambda\mathsf S_t}
 =e^{-t\Phi_\alpha(\lambda)}
\]
such that \(\mathsf S_1\) has law \(\rho_\alpha\).  In particular,
the Wilson endpoint is the time-one kernel of the subordinate heat
semigroup \(e^{-t\Phi_\alpha(-\Delta)}\).
\end{corollary}

\begin{proof}
Apply the standard Bernstein-function construction to
\eqref{eq:newV51-Levy}.  At \(t=1\) its Laplace transform is
\(L_\alpha\), which identifies the law with \(\rho_\alpha\).
Then use \Cref{thm:newV51-density}.
\end{proof}

\begin{corollary}[A spectrally calibrated heat time]
\label{cor:newV51-calibration}
Define
\[
 \tau_\alpha:={\Phi_\alpha(3)\over3}>0.
\]
The Wilson and heat multipliers agree exactly in the fundamental
representation:
\[
 q_{\alpha,1}=e^{-3\tau_\alpha}.
\]
Moreover, concavity of \(\Phi_\alpha\) gives
\begin{align}
 \Phi_\alpha(\lambda)&\ge\tau_\alpha\lambda
 &&(0\le\lambda\le3),                                         \label{eq:newV51-concave-low}\\
 \Phi_\alpha(\lambda)&\le\tau_\alpha\lambda
 &&(\lambda\ge3).                                             \label{eq:newV51-concave-high}
\end{align}
\end{corollary}

\begin{proof}
A nonzero Bernstein function with value zero at the origin is
increasing and concave.  Hence
\(\Phi_\alpha(\lambda)/\lambda\) is nonincreasing on
\((0,\infty)\).  Evaluate that ratio at \(3=c_1\).
\end{proof}

\begin{warning}[The calibrated heat kernel is not a high-spin
majorant]
\label{warn:newV51-calibration}
For every nontrivial \(\SU(2)\) Casimir \(c_n\ge3\),
\eqref{eq:newV51-concave-high} implies
\[
 q_{\alpha,n}\ge e^{-\tau_\alpha c_n}.
\]
Thus the calibrated Wilson multiplier decays no faster than the
calibrated heat multiplier.  This recovers, in positive
subordination language, the direction of the V41 high-spin
inverse-heat obstruction.
\end{warning}

\section{Duhamel derivatives before the random-time integral}
\label{sec:newV51-Duhamel}

\begin{lemma}[Mixed affine Duhamel formula]
\label{lem:newV51-Duhamel}
Let
\[
 A(x)=A_0+\sum_{j=1}^r x_jV_j
\]
be self-adjoint on a finite-dimensional Hilbert space, and suppose
\(A(x)\ge\lambda\,\mathbf1\) at the evaluation point.  Then
\begin{align}
 &\partial_{x_1}\cdots\partial_{x_r}
 L_\alpha(A(x))                                               \notag\\
 &\quad=(-1)^r\int\rho_\alpha(dS)
 \sum_{\sigma\in\mathfrak S_r}
 \int_{0<u_1<\cdots<u_r<S}
 e^{-(S-u_r)A}V_{\sigma(r)}
 e^{-(u_r-u_{r-1})A}\cdots
 V_{\sigma(1)}e^{-u_1A}\,du .
 \label{eq:newV51-Duhamel}
\end{align}
Consequently
\begin{equation}
 \left\|
 \partial_{x_1}\cdots\partial_{x_r}L_\alpha(A(x))
 \right\|
 \le
 \mathbb E_{\rho_\alpha}
 \!\left[S^r e^{-\lambda S}\right]
 \prod_{j=1}^r\|V_j\|.
 \label{eq:newV51-Duhamel-norm}
\end{equation}
Both statements remain true at every matrix amplification.
\end{lemma}

\begin{proof}
For a fixed \(S\), repeated differentiation of \(e^{-SA(x)}\)
gives the time-ordered Dyson formula in
\eqref{eq:newV51-Duhamel}.  The simplex has volume \(S^r/r!\)
and there are \(r!\) permutations.  Write
\(A=\lambda\mathbf1+(A-\lambda\mathbf1)\); the scalar factor
\(e^{-\lambda S}\) commutes through the integrand and every
remaining semigroup factor is contractive.  This proves
\eqref{eq:newV51-Duhamel-norm}.  Interchange with
\(\rho_\alpha\) is justified by the finite moments from
\Cref{lem:newV51-Laplace}.  The same proof applies after tensoring
with an identity.
\end{proof}

\begin{corollary}[Exact damped-time derivative identity]
\label{cor:newV51-damped}
For every \(r\ge0\) and \(\lambda\ge0\),
\begin{equation}
 \mathbb E_{\rho_\alpha}[S^r e^{-\lambda S}]
 =(-1)^rL_\alpha^{(r)}(\lambda).
 \label{eq:newV51-damped}
\end{equation}
\end{corollary}

\begin{proof}
Differentiate the Laplace transform under the integral, using any
smaller positive exponential moment as a dominating function.
\end{proof}

\section{The exact long-time obstruction}
\label{sec:newV51-tail}

\begin{lemma}[Order derivative at the Bessel origin]
\label{lem:newV51-order-derivative}
For \(\alpha>0\),
\begin{equation}
 \left.\partial_\nu I_\nu(\alpha)\right|_{\nu=0}
 =-K_0(\alpha)<0.
 \label{eq:newV51-order-derivative}
\end{equation}
\end{lemma}

\begin{proof}
Use the connection formula
\[
 K_\nu(\alpha)
 ={\pi\over2}\,
 {I_{-\nu}(\alpha)-I_\nu(\alpha)\over\sin(\pi\nu)}.
\]
Let \(\nu\to0\).  The numerator has derivative
\(-2\partial_\nu I_\nu(\alpha)|_{\nu=0}\), while
\(\sin(\pi\nu)=\pi\nu+O(\nu^3)\).  This gives the equality.
Positivity of \(K_0\) follows from
\(K_0(\alpha)=\int_0^\infty e^{-\alpha\cosh u}\,du\).
\end{proof}

\begin{theorem}[Exact moment-generating radius]
\label{thm:newV51-radius}
The moment-generating function of \(S\sim\rho_\alpha\) is
\begin{equation}
 \mathcal M_\alpha(z)
 :=\mathbb E e^{zS}
 ={I_{\sqrt{1-z}}(\alpha)\over I_1(\alpha)}
 \qquad(\Re z<1),
 \label{eq:newV51-MGF}
\end{equation}
with the principal square root.  It is analytic for \(|z|<1\)
and has the nonremovable expansion
\begin{equation}
 \mathcal M_\alpha(z)
 ={I_0(\alpha)\over I_1(\alpha)}
 -{K_0(\alpha)\over I_1(\alpha)}\sqrt{1-z}
 +O(1-z)
 \qquad(z\longrightarrow1).
 \label{eq:newV51-branch}
\end{equation}
Therefore its Taylor series at \(0\) has radius exactly \(1\), and
\begin{equation}
 \limsup_{r\to\infty}
 \left({\mathbb E S^r\over r!}\right)^{1/r}=1.
 \label{eq:newV51-moment-limsup}
\end{equation}
\end{theorem}

\begin{proof}
Equation \eqref{eq:newV51-MGF} is
\eqref{eq:newV51-Laplace} with \(\lambda=-z\), first for real
\(z<1\), then throughout the half-plane by analytic continuation
and dominated convergence.  The map \(z\mapsto\sqrt{1-z}\) is
analytic on \(|z|<1\), and \(I_\nu(\alpha)\) is entire in the order
\(\nu\); hence \(\mathcal M_\alpha\) is analytic in that disk.
Taylor expansion in the order at \(\nu=0\), followed by
\Cref{lem:newV51-order-derivative}, proves
\eqref{eq:newV51-branch}.  Its square-root coefficient is nonzero,
so \(z=1\) is not removable.  The radius is therefore exactly
\(1\).  Since the Taylor coefficient of \(z^r\) is
\(\mathbb E S^r/r!\), Cauchy--Hadamard gives
\eqref{eq:newV51-moment-limsup}.
\end{proof}

\begin{corollary}[Damped moments retain a fixed-scale branch]
\label{cor:newV51-damped-radius}
For every fixed \(\lambda\ge0\),
\begin{equation}
 \limsup_{r\to\infty}
 \left(
 {\mathbb E[S^r e^{-\lambda S}]\over r!}
 \right)^{1/r}
 ={1\over1+\lambda}.
 \label{eq:newV51-damped-limsup}
\end{equation}
\end{corollary}

\begin{proof}
The exponential generating function of the numerator is
\[
 \sum_{r\ge0}{z^r\over r!}
 \mathbb E[S^r e^{-\lambda S}]
 =L_\alpha(\lambda-z).
\]
It is analytic for \(|z|<1+\lambda\).  At
\(z=1+\lambda\), the argument of the Bessel order vanishes and the
same nonzero square-root term as in
\eqref{eq:newV51-branch} occurs.  The Taylor radius is exactly
\(1+\lambda\); apply Cauchy--Hadamard.
\end{proof}

\begin{theorem}[No separated thermal random-time moment card]
\label{thm:newV51-no-moment-card}
Let \(s_\alpha\downarrow0\) be any proposed thermal time scale.
There is no constant \(C\), independent of \(\alpha\) and \(r\),
such that, for all sufficiently large \(\alpha\) and every
\(r\ge0\),
\begin{equation}
 \mathbb E_{\rho_\alpha}S^r
 \le r!(Cs_\alpha)^r.
 \label{eq:newV51-false-moment}
\end{equation}
In particular, the stronger estimate
\(\mathbb E S^r\le(Cs_\alpha)^r\) is impossible.
More generally, after a fixed spectral damping \(\lambda\), a
bound with \(Cs_\alpha\) in place of
\((1+\lambda)^{-1}\) cannot hold at all orders whenever
\(Cs_\alpha<(1+\lambda)^{-1}\).
\end{theorem}

\begin{proof}
Choose \(\alpha\) so large that \(Cs_\alpha<1\).
Estimate \eqref{eq:newV51-false-moment} would give the Taylor
series of \(\mathcal M_\alpha\) radius at least
\((Cs_\alpha)^{-1}>1\), contradicting
\Cref{thm:newV51-radius}.  The damped assertion follows in the
same way from \Cref{cor:newV51-damped-radius}.  The estimate without
\(r!\) would make the moment-generating series entire and is still
stronger.
\end{proof}

\begin{warning}[Where the naive subordinated forest loses
uniform arity]
\label{warn:newV51-naive-forest}
If one applies \Cref{lem:newV51-Duhamel} to \(r=m-1\) BKAR tree
edges and takes the norm before summing the connected partition,
the coefficient contains
\(\mathbb E[S^{m-1}e^{-\lambda S}]\).  By
\eqref{eq:newV51-damped-limsup}, its factorial scale is governed
by \((1+\lambda)^{-(m-1)}\), not by
\(s_\alpha^{m-1}\), unless the assembled sector supplies damping
\(\lambda\gtrsim s_\alpha^{-1}\).

Coupled tensor blocks can contain singlets, so such a lower bound
on \(\lambda\) is not automatic.  The M\"obius partition sum,
time ordering, and the rare long-time branch must therefore remain
assembled until their connected cancellations have been used.
\end{warning}

\section{The exact replacement gate}
\label{sec:newV51-gate}

\begin{definition}[Hartman--Watson subordinated forest card,
HWSFC]
\label{def:newV51-HWSFC}
HWSFC is the following complete-object estimate.  Insert
\eqref{eq:newV51-block} into the full M\"obius sum
\eqref{eq:newV51-partition-times}, apply the V28--V29
time-ordered forest to the entire sum, and retain:
\begin{enumerate}[label=\textup{(\roman*)},leftmargin=3.5em]
\item one Hartman--Watson time for each moment block until the
      M\"obius cancellations have occurred;
\item the heat semigroups, including their spectral damping, until
      after all tree insertions;
\item the identity and antipodal heat images as the single exact
      heat kernel \(h_S\); and
\item the L\'evy exponent
      \eqref{eq:newV51-Levy} without expanding its rare-jump sector
      in derivative order.
\end{enumerate}
The target is a constant \(K\), independent of \(m,\alpha\), for
which the resulting unmarked and one geometrically marked
connected operators are bounded by
\begin{equation}
 K^m\left(\prod_{i=1}^m b_i\right)B^{m-2}.
 \label{eq:newV51-HWSFC-target}
\end{equation}
\end{definition}

\begin{proposition}[What HWSFC would bypass]
\label{prop:newV51-bypass}
HWSFC would prove the raw Wilson endpoint aggregate strong card
without:
\begin{enumerate}[label=\textup{(\alph*)},leftmargin=3em]
\item inverting the heat multiplier at high spin;
\item separating wrapped images in a complex physical strip; or
\item estimating the random-time moments before connected
      cancellation.
\end{enumerate}
It would not by itself prove the marked Wilson--heat comparison
unless the \(\alpha\)- or bridge derivative is incorporated in the
same assembled estimate.
\end{proposition}

\begin{proof}
The first conclusion follows from the exact positive identities
\eqref{eq:newV51-density-mixture} and
\eqref{eq:newV51-partition-times} together with the target
\eqref{eq:newV51-HWSFC-target}.  The three bypassed operations are
absent from that construction.  A comparison theorem is a
derivative statement, whereas HWSFC as stated includes only an
unmarked endpoint unless its marked row is also proved.
\end{proof}

\begin{programme}[V52: rare-jump and connected-cancellation
ledger]
\label{prog:newV51-V52}
\begin{enumerate}[leftmargin=3em]
\item Extract from \(\Phi_\alpha\) a small-time analytic sector and
      a rare-jump sector responsible for the branch at
      \(z=1\).
\item Quantify the \(\alpha\)-dependence of the square-root branch
      coefficient \(K_0(\alpha)/I_1(\alpha)\).
\item Insert the blockwise times into the exact M\"obius sum and
      test which long-time contributions cancel before norms.
\item Formulate the surviving long-time term as a wrapped
      real-integral remainder, not a high physical derivative.
\item If the branch does not cancel, combine its exponentially
      small \(\alpha\)-weight with the V50 grade/leg inequality
      before summing arity.
\end{enumerate}
\end{programme}

\begin{firewall}[Claim boundary after V51]
\label{fire:newV51-boundary}
V51 proves, conditional only on the explicitly imported classical
Hartman--Watson scalar theorem, an exact positive heat
subordination of the complete Wilson endpoint, its tensor-block
functional calculus, its Bernstein exponent, and its
time-ordered Duhamel representation.  It also proves the exact
moment-generating radius and rules out the separated
thermal-scale random-time moment card at unbounded arity.

V51 does not prove HWSFC, UGCGC, thermal connected WPRAC,
all-arity RGSC, the raw Wilson aggregate strong card, all-order
strong MTA, WUFC, CPM, cutoff removal, reflection positivity,
Osterwalder--Schrader reconstruction, a continuum Yang--Mills
measure, or a positive Hamiltonian gap.  The full Yang--Mills mass
gap has not been proved.
\end{firewall}

\begin{theorem}[V51 result ledger]
\label{thm:newV51-results}
V51:
\begin{enumerate}[label=\textup{(V51.\arabic*)},leftmargin=6.5em]
\item constructs the tilted Hartman--Watson time law
      \(\rho_\alpha\);
\item proves its exact transform
      \(I_{\sqrt{1+\lambda}}(\alpha)/I_1(\alpha)\);
\item proves exact Wilson multiplier and full-density heat
      subordination;
\item proves the exact tensor-block and Bernstein functional
      calculi;
\item records the independent block-time structure of Wilson
      cumulants;
\item proves the amplified mixed Duhamel formula;
\item proves the square-root singularity and exact moment radius;
\item disproves the naive all-order thermal random-time moment
      card; and
\item isolates HWSFC as the complete-object replacement gate.
\end{enumerate}
\end{theorem}

\begin{proof}
Use
\Cref{lem:newV51-Laplace,
thm:newV51-multiplier,
thm:newV51-density,
cor:newV51-block,
warn:newV51-not-one-time,
thm:newV51-Bernstein,
cor:newV51-subordinator,
lem:newV51-Duhamel,
thm:newV51-radius,
cor:newV51-damped-radius,
thm:newV51-no-moment-card,
def:newV51-HWSFC}.
\end{proof}

\paragraph{Parent-version record.}
V51 was formed from
\texttt{Renewal\_Yang\_Mills\_Mass\_Gap\_Research\_Notes\_V50.tex}.
The V50 parent had \(24\,092\) logical source lines,
\(758\,240\) bytes, and SHA--256
\[
 \texttt{226f4204f409a5aeae615738fda88bb644f9d465c6418a5d8d2fc0e7370a328f}.
\]
The next compulsory companion audit is V55.

% =============================================================================
% END APPENDED FIFTH-SERIES V51 MATERIAL
% =============================================================================

\clearpage
\chapter{Fifth-series V52: the exact rare-time branch and the
wrapped-image crossover}
\label{ch:newV52-branch}

\section{The even--odd order decomposition}
\label{sec:newV52-evenodd}

\begin{definition}[Entire order-parity functions]
\label{def:newV52-EO}
For fixed \(\alpha>0\), define
\begin{align}
 E_\alpha(w)
 &:={I_{\sqrt w}(\alpha)+I_{-\sqrt w}(\alpha)\over2},
                                                               \label{eq:newV52-E}\\
 O_\alpha(w)
 &:={I_{\sqrt w}(\alpha)-I_{-\sqrt w}(\alpha)\over2\sqrt w},
                                                               \label{eq:newV52-O}
\end{align}
with the values at \(w=0\) defined by continuation.
\end{definition}

\begin{lemma}[Entire continuation and exact connection formula]
\label{lem:newV52-EO-entire}
Both \(E_\alpha\) and \(O_\alpha\) are entire functions of \(w\).
Moreover,
\begin{align}
 I_\nu(\alpha)
 &=E_\alpha(\nu^2)+\nu O_\alpha(\nu^2),                        \label{eq:newV52-order-split}\\
 O_\alpha(w)
 &=-{1\over\pi}\,
 {\sin(\pi\sqrt w)\over\sqrt w}\,
 K_{\sqrt w}(\alpha),                                         \label{eq:newV52-O-K}\\
 O_\alpha(0)&=-K_0(\alpha).                                   \label{eq:newV52-O-zero}
\end{align}
\end{lemma}

\begin{proof}
For fixed positive \(\alpha\), the defining series for
\(I_\nu(\alpha)\) is entire in the order \(\nu\).  Split its order
Taylor series into even and odd powers.  The even coefficients
form an entire function of \(\nu^2\), and the odd coefficients
form \(\nu\) times another entire function of \(\nu^2\).  This
proves \eqref{eq:newV52-order-split} and the entireness assertions.

The connection identity
\[
 I_\nu(\alpha)-I_{-\nu}(\alpha)
 =-{2\over\pi}\sin(\pi\nu)K_\nu(\alpha)
\]
gives \eqref{eq:newV52-O-K}.  Letting \(w\to0\) gives
\eqref{eq:newV52-O-zero}; equivalently use
\Cref{lem:newV51-order-derivative}.
\end{proof}

\begin{theorem}[Exact analytic/branch decomposition]
\label{thm:newV52-exact-split}
The moment-generating function of the V51 tilted time admits the
global slit-plane identity
\begin{equation}
 \mathcal M_\alpha(z)
 =
 \mathcal A_\alpha(z)+\mathcal R_\alpha(z),                    \label{eq:newV52-M-split}
\end{equation}
where
\begin{align}
 \mathcal A_\alpha(z)
 &:={E_\alpha(1-z)\over I_1(\alpha)},                          \label{eq:newV52-A}\\
 \mathcal R_\alpha(z)
 &:={\sqrt{1-z}\,O_\alpha(1-z)\over I_1(\alpha)}               \label{eq:newV52-R}\\
 &=-{1\over\pi I_1(\alpha)}
   \sin\!\bigl(\pi\sqrt{1-z}\bigr)
   K_{\sqrt{1-z}}(\alpha).                                    \label{eq:newV52-R-K}
\end{align}
\(\mathcal A_\alpha\) is entire.  Every finite-plane branch of
\(\mathcal M_\alpha\) is carried by
\(\mathcal R_\alpha\), and
\begin{equation}
 \mathcal R_\alpha(z)
 =-\delta_\alpha\sqrt{1-z}
 +O_\alpha\bigl((1-z)^{3/2}\bigr),\qquad
 \delta_\alpha:={K_0(\alpha)\over I_1(\alpha)}>0.              \label{eq:newV52-delta}
\end{equation}
\end{theorem}

\begin{proof}
Put \(w=1-z\) and \(\nu=\sqrt w\) in
\eqref{eq:newV52-order-split}, then divide by \(I_1(\alpha)\).
Equations \eqref{eq:newV52-R-K} and
\eqref{eq:newV52-delta} follow from
\eqref{eq:newV52-O-K} and
\eqref{eq:newV52-O-zero}.  Since both parity functions are entire,
the first summand is entire in \(z\), while the second is a square
root times an entire function of \(1-z\).
\end{proof}

\begin{warning}[This is not a positive-measure splitting]
\label{warn:newV52-not-positive}
The exact decomposition
\eqref{eq:newV52-M-split} separates analytic structures, not
probability measures.  Neither
\(\mathcal A_\alpha\) nor \(\mathcal R_\alpha\) is asserted to be
completely monotone.  Positivity remains attached to their sum
\(L_\alpha(\lambda)=\mathcal M_\alpha(-\lambda)\).
\end{warning}

\section{Nonperturbative size of the branch coefficient}
\label{sec:newV52-delta}

\begin{lemma}[Elementary Bessel bounds]
\label{lem:newV52-Bessel-bounds}
For every \(\alpha\ge1\),
\begin{align}
 {e^{\alpha-1}\over2\pi\sqrt{2\alpha}}
 \le I_1(\alpha)
 &\le {e^\alpha\over\sqrt{2\pi\alpha}},                        \label{eq:newV52-I-bounds}\\
 {e^{-\alpha-1}\over\sqrt\alpha}
 \le K_0(\alpha)
 &\le e^{-\alpha}\sqrt{\pi\over2\alpha}.                       \label{eq:newV52-K-bounds}
\end{align}
\end{lemma}

\begin{proof}
Use
\begin{equation}
 I_1(\alpha)
 ={\alpha\over\pi}\int_{-1}^{1}
 e^{\alpha t}\sqrt{1-t^2}\,dt.                                \label{eq:newV52-I-integral}
\end{equation}
For the upper bound, put \(y=1-t\), use
\(\sqrt{1-t^2}\le\sqrt{2y}\), and extend the integral to
\([0,\infty)\):
\[
 I_1(\alpha)
 \le{\alpha e^\alpha\sqrt2\over\pi}
 \int_0^\infty e^{-\alpha y}y^{1/2}\,dy
 ={e^\alpha\over\sqrt{2\pi\alpha}}.
\]
For the lower bound, restrict the \(y\)-integral to
\([1/(2\alpha),1/\alpha]\).  On that interval,
\(e^{-\alpha y}\ge e^{-1}\),
\(\sqrt{y(2-y)}\ge\sqrt y\ge(2\alpha)^{-1/2}\),
and the interval has length \(1/(2\alpha)\).  This proves the
first lower bound.

Next use
\[
 K_0(\alpha)=\int_0^\infty e^{-\alpha\cosh u}\,du.
\]
Since \(\cosh u\ge1+u^2/2\), Gaussian integration proves the upper
bound.  On \(0\le u\le\alpha^{-1/2}\le1\),
\(\cosh u\le1+u^2\); hence the integrand is at least
\(e^{-\alpha-1}\), proving the lower bound.
\end{proof}

\begin{theorem}[Sharp exponential scale of the branch]
\label{thm:newV52-delta-bounds}
For every \(\alpha\ge1\),
\begin{equation}
 {\sqrt{2\pi}\over e}\,e^{-2\alpha}
 \le\delta_\alpha
 \le 2e\pi^{3/2}e^{-2\alpha}.                                 \label{eq:newV52-delta-bounds}
\end{equation}
With the fifth-series thermal parameter
\begin{equation}
 s={1\over2\alpha+2/3},                                       \label{eq:newV52-s}
\end{equation}
there are absolute \(0<c_\delta<C_\delta<\infty\) such that
\begin{equation}
 c_\delta e^{-1/s}\le\delta_\alpha\le C_\delta e^{-1/s}.
 \label{eq:newV52-delta-s}
\end{equation}
\end{theorem}

\begin{proof}
Divide the lower bound for \(K_0\) by the upper bound for \(I_1\),
and the upper bound for \(K_0\) by the lower bound for \(I_1\), in
\Cref{lem:newV52-Bessel-bounds}.  Finally
\(2\alpha=s^{-1}-2/3\), so the factor \(e^{2/3}\) is absorbed into
the two absolute constants.
\end{proof}

\section{Exact high-order moment asymptotics}
\label{sec:newV52-moments}

\begin{theorem}[Darboux law for tilted-time moments]
\label{thm:newV52-Darboux}
For fixed \(\alpha>0\), as \(r\to\infty\),
\begin{equation}
 {\mathbb E_{\rho_\alpha}S^r\over r!}
 =
 {\delta_\alpha\over2\sqrt\pi}\,r^{-3/2}
 \left(1+O_\alpha(r^{-1})\right).
 \label{eq:newV52-moment-asymptotic}
\end{equation}
More generally, for every fixed \(\lambda\ge0\),
\begin{equation}
 {\mathbb E_{\rho_\alpha}[S^r e^{-\lambda S}]\over r!}
 =
 {\delta_\alpha\over2\sqrt\pi}\,
 (1+\lambda)^{1/2-r}r^{-3/2}
 \left(1+O_{\alpha,\lambda}(r^{-1})\right).
 \label{eq:newV52-damped-asymptotic}
\end{equation}
\end{theorem}

\begin{proof}
By \Cref{thm:newV52-exact-split}, the entire part contributes
Taylor coefficients smaller than \(R^{-r}\) for every fixed
\(R>0\).  The branch part begins with
\(-\delta_\alpha(1-z)^{1/2}\), and its next nonanalytic term is
of order \((1-z)^{3/2}\).  The binomial coefficient estimate
\[
 [z^r](1-z)^{1/2}
 =-{1\over2\sqrt\pi}\,r^{-3/2}
 \left(1+O(r^{-1})\right)
\]
therefore gives \eqref{eq:newV52-moment-asymptotic}; the next
branch term is \(O_\alpha(r^{-5/2})\).

For the damped moment, take coefficients in
\[
 L_\alpha(\lambda-z)
 =\mathcal M_\alpha(z-\lambda).
\]
Its branch point is \(z=1+\lambda\), and its leading singular term
is
\[
 -\delta_\alpha\sqrt{1+\lambda-z}
 =-\delta_\alpha(1+\lambda)^{1/2}
  \left(1-{z\over1+\lambda}\right)^{1/2}.
\]
The same coefficient estimate proves
\eqref{eq:newV52-damped-asymptotic}.
\end{proof}

\begin{corollary}[No finite local counterterm removes the tail]
\label{cor:newV52-no-finite-counterterm}
If \(G\) is entire, then
\(\mathcal M_\alpha-G\) still has the nonzero square-root
coefficient \(-\delta_\alpha\) at \(z=1\).  In particular, no
finite Taylor, heat-time, or local Gaussian counterterm can extend
the random-time moment radius beyond \(1\).
\end{corollary}

\begin{proof}
An entire function has trivial monodromy about \(z=1\) and cannot
cancel a nonzero half-integer power there.
\end{proof}

\section{Identification with the V40 derivative window}
\label{sec:newV52-crossover}

\begin{definition}[Nonperturbative crossover order]
\label{def:newV52-rstar}
For \(0<s<e^{-1}\), put
\begin{equation}
 r_\star(s):={1\over s\log(1/s)}.
 \label{eq:newV52-rstar}
\end{equation}
\end{definition}

\begin{theorem}[Exact scale matching]
\label{thm:newV52-scale-matching}
The thermal power and the Hartman--Watson branch weight meet at
\begin{equation}
 s^{\,r_\star(s)}=e^{-1/s}.                                   \label{eq:newV52-scale-identity}
\end{equation}
Consequently the V51 rare-time branch overtakes a putative
thermal estimate \(s^r\), up to exponential-base and polynomial
factors, at order
\[
 r\asymp {1\over s\log(1/s)}.
\]
This is the same order as the V40 collapse of complex
identity/antipodal image separation.
\end{theorem}

\begin{proof}
Take logarithms:
\[
 \log s^{r_\star}
 ={1\over s\log(1/s)}\log s=-{1\over s}.
\]
Combine this identity with \eqref{eq:newV52-delta-s} and
\eqref{eq:newV52-moment-asymptotic}.  V40 independently located
the first physical complex-zero obstruction at derivative order
\(O((s\log(1/s))^{-1})\).
\end{proof}

\begin{assessment}[The two obstructions are one wrapped sector]
\label{ass:newV52-one-sector}
V40 saw the obstruction in the physical angle plane: an
antipodal image forces zeros toward the identity chart and shrinks
the all-order analytic radius.  V51--V52 see it in the exact
spectral subordination variable: the odd-in-order Bessel sector
has weight \(e^{-1/s}\) but a fixed branch radius.  The identical
crossover scale in \Cref{thm:newV52-scale-matching} is evidence,
now supported by exact formulae on both sides, that these are two
representations of the same wrapped compact-group contribution.
\end{assessment}

\section{Why blockwise absolute subordination still cannot close}
\label{sec:newV52-no-blockwise}

\begin{proposition}[Scalar affine diagnostic]
\label{prop:newV52-scalar-diagnostic}
Fix \(\lambda\ge0\) and define the scalar affine test
\[
 F_{\alpha,s}(x_1,\ldots,x_r)
 :=
 L_\alpha\left(\lambda+s^{-1}\sum_{j=1}^r x_j\right).
\]
Then
\begin{equation}
 \left|
 \partial_{x_1}\cdots\partial_{x_r}
 F_{\alpha,s}(0)
 \right|
 =
 s^{-r}\mathbb E[S^r e^{-\lambda S}],                         \label{eq:newV52-scalar-derivative}
\end{equation}
and
\begin{equation}
 \limsup_{r\to\infty}
 \left(
 {|\partial_{x_1}\cdots\partial_{x_r}F_{\alpha,s}(0)|\over r!}
 \right)^{1/r}
 ={1\over s(1+\lambda)}.                                      \label{eq:newV52-scalar-radius}
\end{equation}
\end{proposition}

\begin{proof}
Differentiate the affine scalar composition and use
\eqref{eq:newV51-damped}.  Then apply
\eqref{eq:newV51-damped-limsup}.
\end{proof}

\begin{corollary}[No positivity-only HWSFC proof]
\label{cor:newV52-no-blackbox}
No proof which replaces every subordinated block derivative by the
absolute Duhamel norm
\eqref{eq:newV51-Duhamel-norm} can establish HWSFC with a uniform
thermal exponential base.  This failure already occurs for the
one-dimensional positive operator in
\Cref{prop:newV52-scalar-diagnostic}.
\end{corollary}

\begin{proof}
Such a proof would imply a uniform analytic radius in the rescaled
variables.  Equation \eqref{eq:newV52-scalar-radius} shows that the
available radius is only \(s(1+\lambda)\), and hence collapses as
\(s\downarrow0\) for every fixed \(\lambda\).
\end{proof}

\begin{warning}[This is a method no-go, not a cumulant no-go]
\label{warn:newV52-method-only}
The scalar diagnostic contains no M\"obius connected projection.
It disproves the blockwise absolute-value route to HWSFC; it does
not disprove the Wilson connected strong card itself.  A valid
closure must either cancel the branch in the assembled connected
source or bound it nonperturbatively without taking arbitrarily
many derivatives.
\end{warning}

\section{The revised wrapped-sector gate}
\label{sec:newV52-gate}

\begin{definition}[Branch-resolved source card, BRSC]
\label{def:newV52-BRSC}
BRSC is a decomposition of the complete Wilson or geometric
connected source into
\[
 \mathcal K_{\rm conn}
 =\mathcal K_{\rm loc}+\mathcal K_{\rm wrap}
\]
with:
\begin{enumerate}[label=\textup{(\roman*)},leftmargin=3.5em]
\item \(\mathcal K_{\rm loc}\) controlled by the V50
      grade/leg forest with one exponential base through orders
      \(r\lesssim r_\star(s)\);
\item \(\mathcal K_{\rm wrap}\) kept as the exact Bessel
      odd-order expression \eqref{eq:newV52-R-K}, or the equivalent
      paired heat-image real integral, without high
      differentiation;
\item its complete connected norm carrying the nonperturbative
      factor \(e^{-1/s}\); and
\item a direct source-domain bound which converts that factor,
      together with the leg weights, into
      \(K^m(\prod_i b_i)B^{m-2}\) for every \(m\), including
      \(m\gg r_\star(s)\).
\end{enumerate}
\end{definition}

\begin{programme}[V53: local/branch coefficient majorant]
\label{prog:newV52-V53}
\begin{enumerate}[leftmargin=3em]
\item Bound the entire parity sector
      \(\mathcal A_\alpha\) on a source disk whose radius is
      proportional to \(1/s\).
\item Bound the remainder
      \(\mathcal R_\alpha+\delta_\alpha\sqrt{1-z}\) uniformly on a
      slit disk and preserve its factor \(\delta_\alpha\).
\item Derive a two-term moment majorant
      \[
      {\mathbb E S^r\over r!}
      \le (Cs)^r+C\delta_\alpha r^{-3/2}.
      \]
\item Insert that majorant before, rather than after, the
      connected grade/leg resummation.
\item Determine whether the V29 tree factorial and the
      \(B^{m-2}\) tree-volume factor absorb the branch term; if
      not, return to the direct geometric source disk required by
      BRSC(iv).
\end{enumerate}
\end{programme}

\begin{firewall}[Claim boundary after V52]
\label{fire:newV52-boundary}
V52 proves the exact entire/branch decomposition of the
Hartman--Watson transform, elementary two-sided
\(e^{-2\alpha}\) bounds on its branch coefficient, the high-order
moment asymptotic, the exact
\((s\log(1/s))^{-1}\) crossover, and a scalar no-go for any
blockwise absolute-value subordination proof.

V52 does not prove BRSC, HWSFC, UGCGC, thermal connected WPRAC,
all-arity RGSC, the raw Wilson aggregate strong card, all-order
strong MTA, WUFC, CPM, cutoff removal, reflection positivity,
Osterwalder--Schrader reconstruction, a continuum Yang--Mills
measure, or a positive Hamiltonian gap.  The full Yang--Mills mass
gap has not been proved.
\end{firewall}

\begin{theorem}[V52 result ledger]
\label{thm:newV52-results}
V52:
\begin{enumerate}[label=\textup{(V52.\arabic*)},leftmargin=6.5em]
\item splits the Bessel order dependence into two entire parity
      functions;
\item identifies the exact odd-order wrapped branch;
\item proves explicit two-sided bounds for
      \(K_0(\alpha)/I_1(\alpha)\);
\item identifies that coefficient with the scale \(e^{-1/s}\);
\item proves the \(r!r^{-3/2}\) high-moment law;
\item proves exact agreement with the V40 derivative-window
      scale;
\item rules out blockwise absolute HWSFC by a scalar diagnostic;
      and
\item formulates BRSC as the required two-regime replacement.
\end{enumerate}
\end{theorem}

\begin{proof}
Use
\Cref{lem:newV52-EO-entire,
thm:newV52-exact-split,
lem:newV52-Bessel-bounds,
thm:newV52-delta-bounds,
thm:newV52-Darboux,
thm:newV52-scale-matching,
prop:newV52-scalar-diagnostic,
cor:newV52-no-blackbox,
def:newV52-BRSC}.
\end{proof}

\paragraph{Parent-version record.}
V52 was formed from
\texttt{Renewal\_Yang\_Mills\_Mass\_Gap\_Research\_Notes\_V51.tex}.
The V51 parent had \(24\,800\) logical source lines,
\(782\,360\) bytes, and SHA--256
\[
 \texttt{daa6256c65efab59f718c1548018e0e8c7893598e3801b529da7f52a60be1214}.
\]
The next compulsory companion audit is V55.

% =============================================================================
% END APPENDED FIFTH-SERIES V52 MATERIAL
% =============================================================================

\clearpage
\chapter{Fifth-series V53: uniform local coefficients and an exact
wrapped remainder}
\label{ch:newV53-Schlafli}

\section{Schl\"afli separation in the order variable}
\label{sec:newV53-Schlafli}

\begin{lemma}[Classical Schl\"afli representation]
\label{lem:newV53-Schlafli}
For \(\alpha>0\) and complex \(\nu\),
\begin{equation}
 I_\nu(\alpha)
 ={1\over\pi}\int_0^\pi e^{\alpha\cos\theta}
       \cos(\nu\theta)\,d\theta
 -{\sin(\pi\nu)\over\pi}
  \int_0^\infty e^{-\alpha\cosh t-\nu t}\,dt.
 \label{eq:newV53-Schlafli}
\end{equation}
Both integrals converge locally uniformly in \(\nu\).
\end{lemma}

\begin{proof}
Equation \eqref{eq:newV53-Schlafli} is the classical
Schl\"afli integral representation for the modified Bessel
function.  The second integral converges locally uniformly because
\(\cosh t\) grows exponentially in \(t\), and the first is over a
compact interval.  The standard derivation first holds in a real
order half-plane by contour deformation of the generating
integral; local uniform convergence then extends the identity to
all \(\nu\) by analytic continuation.  This special-function
identity is used as an exact representation, not as an asymptotic
formula.
\end{proof}

\begin{definition}[Compact and wrapped transform sectors]
\label{def:newV53-CW}
On the principal sheet \(\nu(z)=\sqrt{1-z}\), define
\begin{align}
 \mathcal C_\alpha(z)
 &:={1\over\pi I_1(\alpha)}
   \int_0^\pi e^{\alpha\cos\theta}
       \cos\!\bigl(\theta\nu(z)\bigr)\,d\theta,                 \label{eq:newV53-C}\\
 \mathcal W_\alpha(z)
 &:=-{\sin(\pi\nu(z))\over\pi I_1(\alpha)}
   \int_0^\infty e^{-\alpha\cosh t-\nu(z)t}\,dt.               \label{eq:newV53-W}
\end{align}
\end{definition}

\begin{theorem}[Exact compact/wrapped splitting]
\label{thm:newV53-CW}
For \(|z|<1\),
\begin{equation}
 \mathcal M_\alpha(z)
 =\mathcal C_\alpha(z)+\mathcal W_\alpha(z).                   \label{eq:newV53-CW}
\end{equation}
\(\mathcal C_\alpha\) is entire in \(z\).  The wrapped term carries
both the odd branch \(\mathcal R_\alpha\) of V52 and an
exponentially small entire even-order correction.
\end{theorem}

\begin{proof}
Insert \(\nu=\sqrt{1-z}\) in
\eqref{eq:newV53-Schlafli} and divide by \(I_1(\alpha)\).  Since
\(\cos(\theta\nu)\) is an even entire function of \(\nu\), it is
an entire function of \(\nu^2=1-z\); this proves the entireness of
\(\mathcal C_\alpha\).

For comparison with V52, apply \eqref{eq:newV53-Schlafli} at
\(\nu\) and \(-\nu\).  The odd-order part of the second integral
is
\[
 -{\sin(\pi\nu)\over\pi I_1(\alpha)}
 \int_0^\infty e^{-\alpha\cosh t}\cosh(\nu t)\,dt
 =-{\sin(\pi\nu)\over\pi I_1(\alpha)}K_\nu(\alpha),
\]
which is exactly \(\mathcal R_\alpha\) in
\eqref{eq:newV52-R-K}.  The remaining hyperbolic contribution is
even in \(\nu\), hence entire in \(z\).
\end{proof}

\section{Uniform coefficient bound for the compact sector}
\label{sec:newV53-compact}

\begin{lemma}[Complex compact-sector maximum]
\label{lem:newV53-C-maximum}
There is an absolute constant \(C_0\) such that, for every
\(\alpha\ge1\) and \(R\ge0\),
\begin{equation}
 \sup_{|z|\le R}|\mathcal C_\alpha(z)|
 \le
 C_0\exp\left\{{\pi^2(1+R)\over8\alpha}\right\}.
 \label{eq:newV53-C-maximum}
\end{equation}
\end{lemma}

\begin{proof}
For \(0\le\theta\le\pi\),
\[
 \cos\theta\le1-{2\theta^2\over\pi^2},
\]
because
\(1-\cos\theta=2\sin^2(\theta/2)\) and
\(\sin u\ge2u/\pi\) on \([0,\pi/2]\).
For \(|z|\le R\),
\[
 |\nu(z)|\le\sqrt{1+R},\qquad
 |\cos(\theta\nu(z))|
 \le e^{\theta\sqrt{1+R}}.
\]
Therefore, with
\(a=2\alpha/\pi^2\) and \(b=\sqrt{1+R}\),
\begin{align*}
 {1\over\pi}
 \int_0^\pi e^{\alpha\cos\theta}
 |\cos(\theta\nu)|\,d\theta
 &\le {e^\alpha\over\pi}
 \int_0^\infty e^{-a\theta^2+b\theta}\,d\theta\\
 &\le e^\alpha\sqrt{\pi\over2\alpha}
       \exp\left\{{\pi^2(1+R)\over8\alpha}\right\}.
\end{align*}
Divide by the lower bound for \(I_1(\alpha)\) in
\eqref{eq:newV52-I-bounds}.  All remaining numerical factors are
absorbed into \(C_0\).
\end{proof}

\begin{theorem}[All-order compact coefficient card]
\label{thm:newV53-C-coeff}
Write
\[
 \mathcal C_\alpha(z)=\sum_{r=0}^\infty c_{\alpha,r}z^r.
\]
There is an absolute \(C_C\) such that, for every
\(\alpha\ge1\) and \(r\ge1\),
\begin{equation}
 |c_{\alpha,r}|
 \le\left({C_C\over\alpha}\right)^r.                           \label{eq:newV53-C-coeff}
\end{equation}
\end{theorem}

\begin{proof}
Use Cauchy's estimate on \(|z|=\alpha\).  By
\Cref{lem:newV53-C-maximum},
\[
 |c_{\alpha,r}|
 \le C_0e^{\pi^2/4}\alpha^{-r}.
\]
Increase \(C_C\ge1\) so that
\(C_0e^{\pi^2/4}\le C_C^r\) for every \(r\ge1\).
\end{proof}

\begin{corollary}[Uniform thermal rescaling of the compact sector]
\label{cor:newV53-C-rescaled}
Let \(s=(2\alpha+2/3)^{-1}\).  There is an absolute \(K_C\) such
that
\begin{equation}
 \left|
 {1\over r!}{d^r\over dx^r}
 \mathcal C_\alpha(x/s)\bigg|_{x=0}
 \right|
 \le K_C^r,\qquad r\ge1.                                      \label{eq:newV53-C-rescaled}
\end{equation}
Thus the compact angular sector has precisely the uniform
exponential coefficient growth required by an all-arity forest.
\end{corollary}

\begin{proof}
For \(\alpha\ge1\),
\[
 {1\over\alpha}
 ={2s\over1-2s/3}\le{8\over3}s.
\]
Now apply \eqref{eq:newV53-C-coeff} and multiply the \(r\)-th
Taylor coefficient by \(s^{-r}\).
\end{proof}

\section{Uniform size of the wrapped sector on strict subdisks}
\label{sec:newV53-wrapped}

\begin{lemma}[Wrapped maximum on a strict disk]
\label{lem:newV53-W-maximum}
For every fixed \(0<\eta<1\), there is a finite constant
\(C_\eta\), independent of \(\alpha\ge1\), such that
\begin{equation}
 \sup_{|z|\le\eta}|\mathcal W_\alpha(z)|
 \le C_\eta e^{-2\alpha}.                                    \label{eq:newV53-W-maximum}
\end{equation}
\end{lemma}

\begin{proof}
Put \(b_\eta=\sqrt{1+\eta}\).  On \(|z|\le\eta\),
\[
 |\nu(z)|\le b_\eta,\qquad
 |\sin(\pi\nu(z))|\le e^{\pi b_\eta}.
\]
Using \(\cosh t\ge1+t^2/2\),
\begin{align*}
 \int_0^\infty
 \left|e^{-\alpha\cosh t-\nu t}\right|dt
 &\le e^{-\alpha}\int_0^\infty
 e^{-\alpha t^2/2+b_\eta t}\,dt\\
 &\le e^{-\alpha}\sqrt{2\pi\over\alpha}
       e^{b_\eta^2/(2\alpha)}.
\end{align*}
Divide by the lower bound for \(I_1(\alpha)\) in
\eqref{eq:newV52-I-bounds}.  The powers of \(\alpha\) cancel and
the remaining factor is at most \(C_\eta e^{-2\alpha}\).
\end{proof}

\begin{corollary}[Wrapped coefficient bound]
\label{cor:newV53-W-coeff}
If
\(\mathcal W_\alpha(z)=\sum_{r\ge0}w_{\alpha,r}z^r\) on
\(|z|<1\), then, for every fixed \(0<\eta<1\),
\begin{equation}
 |w_{\alpha,r}|
 \le C_\eta e^{-2\alpha}\eta^{-r},
 \qquad r\ge0.                                                \label{eq:newV53-W-coeff}
\end{equation}
\end{corollary}

\begin{proof}
Cauchy's coefficient estimate on \(|z|=\eta\), followed by
\Cref{lem:newV53-W-maximum}.
\end{proof}

\begin{theorem}[Uniform two-sector moment majorant]
\label{thm:newV53-two-sector}
There are absolute \(C_1,C_2,C_3\) such that, for
\(\alpha\ge1\) and \(r\ge1\),
\begin{equation}
 {\mathbb E_{\rho_\alpha}S^r\over r!}
 \le
 (C_1s)^r+C_2e^{-1/s}C_3^r.                                  \label{eq:newV53-two-sector}
\end{equation}
One may take \(C_3=2\).
\end{theorem}

\begin{proof}
The left side is the \(r\)-th Taylor coefficient of
\(\mathcal M_\alpha\).  Split it by
\eqref{eq:newV53-CW}.  The compact coefficient is bounded by
\eqref{eq:newV53-C-coeff}, then by a thermal power as in
\Cref{cor:newV53-C-rescaled}.  Apply
\eqref{eq:newV53-W-coeff} with \(\eta=1/2\), and use
\(e^{-2\alpha}=e^{2/3}e^{-1/s}\).  Absorb fixed numerical factors
into \(C_1,C_2\).
\end{proof}

\begin{assessment}[The all-order failure is now localized]
\label{ass:newV53-localized}
After rescaling by \(s^{-r}\), the first term in
\eqref{eq:newV53-two-sector} has one uniform exponential base.
The second becomes
\[
 C_2e^{-1/s}\left({C_3\over s}\right)^r.
\]
Thus every loss of uniform thermal analyticity in the exact
Hartman--Watson representation lies in the wrapped hyperbolic
integral \(\mathcal W_\alpha\).  No compact angular coefficient is
responsible.
\end{assessment}

\section{A strict blockwise no-go and the remaining cancellation}
\label{sec:newV53-no-go}

\begin{theorem}[No all-order absorption of the wrapped coefficient]
\label{thm:newV53-no-absorb}
There do not exist constants \(A,K<\infty\) such that
\begin{equation}
 e^{-1/s}\left({A\over s}\right)^r\le K^r
 \label{eq:newV53-false-absorb}
\end{equation}
for every sufficiently small \(s>0\) and every integer \(r\ge0\).
At
\[
 r>{1\over s\log(A/(Ks))}
\]
the inequality reverses whenever \(A>Ks\).
\end{theorem}

\begin{proof}
Take logarithms of \eqref{eq:newV53-false-absorb}.  It would
require
\[
 -{1\over s}+r\log{A\over Ks}\le0.
\]
For fixed small \(s\), the second term grows without bound as
\(r\to\infty\), giving the displayed threshold.
\end{proof}

\begin{corollary}[BRSC(iv) is indispensable]
\label{cor:newV53-BRSC-needed}
The nonperturbative coefficient \(e^{-1/s}\) cannot be converted
to the all-arity thermal card after \(r\) wrapped derivatives have
already produced \(s^{-r}\).  The wrapped source must be bounded
before high differentiation, exactly as required by BRSC(iv).
\end{corollary}

\begin{proof}
Apply \Cref{thm:newV53-no-absorb} to the wrapped term in
\eqref{eq:newV53-two-sector}.
\end{proof}

\begin{warning}[No contradiction with fixed arity]
\label{warn:newV53-fixed}
For each fixed \(r\), \(e^{-1/s}s^{-r}\) is smaller than every
power of \(s\) as \(s\downarrow0\).  Hence the V50 fixed-arity
card remains valid.  The obstruction occurs only when the arity
is unbounded uniformly in \(s\).
\end{warning}

\section{The wrapped source-domain card}
\label{sec:newV53-WSDC}

\begin{definition}[Wrapped source-domain card, WSDC]
\label{def:newV53-WSDC}
WSDC is an estimate on the complete connected source functional,
before extraction of its Taylor coefficients, with the following
properties:
\begin{enumerate}[label=\textup{(\roman*)},leftmargin=3.5em]
\item the compact sector is summed by
      \Cref{thm:newV53-C-coeff};
\item every occurrence of the hyperbolic integral
      \(\mathcal W_\alpha\) remains unexpanded in derivative
      order;
\item the connected M\"obius/forest projection is applied before
      its norm;
\item on a source polydomain scaled by the physical weights
      \(b_i\), the complete wrapped contribution is at most
      \(K^m e^{-1/s}\prod_i b_i\); and
\item a direct large-arity estimate, rather than Cauchy
      differentiation across the branch at distance \(s\), gains
      the remaining \(B^{m-2}\).
\end{enumerate}
\end{definition}

\begin{proposition}[Sufficiency of WSDC plus the compact forest]
\label{prop:newV53-WSDC-suffices}
If WSDC holds with one exponential base and the V29 tree
resummation is applied to the compact sector using
\eqref{eq:newV53-C-rescaled}, then BRSC and the raw Wilson
aggregate strong card follow.
\end{proposition}

\begin{proof}
Split every exact block transform by
\eqref{eq:newV53-CW}.  The term containing only compact factors
has uniform rescaled coefficients by
\Cref{cor:newV53-C-rescaled}, so the V29 stable tree resummation
applies.  Every remaining term contains a wrapped factor and is
bounded in the assembled source by WSDC.  Summing the two bounds
gives the BRSC target, hence the raw Wilson endpoint card.
\end{proof}

\begin{programme}[V54: physical spectral and singlet test of WSDC]
\label{prog:newV53-V54}
\begin{enumerate}[leftmargin=3em]
\item Insert the exact compact/wrapped split into the tensor-block
      moment \(L_\alpha(\mathcal C_B)\).
\item Decompose \(\mathcal C_B\) into total-spin spectral
      projections and locate the wrapped branch relative to the
      physical spectrum \(0,3,8,\ldots\).
\item Test whether tree directions have nonzero repeated
      compression on the singlet projection.
\item If the singlet compression vanishes after connected
      M\"obius projection, prove the resulting divided-difference
      gap card.
\item If it survives, construct an explicit commuting spin
      diagnostic and return to a direct bounded-group source
      estimate without spectral differentiation.
\end{enumerate}
\end{programme}

\begin{firewall}[Claim boundary after V53]
\label{fire:newV53-boundary}
V53 proves an exact compact-angular/wrapped-hyperbolic splitting,
an all-order thermal coefficient card for the compact sector, a
uniform \(e^{-1/s}\) strict-disk bound for the wrapped sector, and
the two-sector moment majorant
\eqref{eq:newV53-two-sector}.  It proves that the wrapped term
cannot be absorbed after arbitrary thermal rescaling and isolates
WSDC.

V53 does not prove WSDC, BRSC, HWSFC, UGCGC, thermal connected
WPRAC, all-arity RGSC, the raw Wilson aggregate strong card,
all-order strong MTA, WUFC, CPM, cutoff removal, reflection
positivity, Osterwalder--Schrader reconstruction, a continuum
Yang--Mills measure, or a positive Hamiltonian gap.  The full
Yang--Mills mass gap has not been proved.
\end{firewall}

\begin{theorem}[V53 result ledger]
\label{thm:newV53-results}
V53:
\begin{enumerate}[label=\textup{(V53.\arabic*)},leftmargin=6.5em]
\item gives the exact Schl\"afli compact/wrapped split;
\item proves a uniform complex maximum for the compact sector;
\item proves its coefficients are bounded by
      \((C/\alpha)^r\);
\item proves a strict-disk \(e^{-2\alpha}\) wrapped maximum;
\item proves the uniform two-sector moment majorant;
\item localizes every all-order thermal-radius loss to the wrapped
      hyperbolic integral;
\item proves that post-differentiation nonperturbative absorption
      is impossible; and
\item isolates WSDC as the direct-source replacement.
\end{enumerate}
\end{theorem}

\begin{proof}
Use
\Cref{lem:newV53-Schlafli,
thm:newV53-CW,
lem:newV53-C-maximum,
thm:newV53-C-coeff,
cor:newV53-C-rescaled,
lem:newV53-W-maximum,
cor:newV53-W-coeff,
thm:newV53-two-sector,
thm:newV53-no-absorb,
def:newV53-WSDC}.
\end{proof}

\paragraph{Parent-version record.}
V53 was formed from
\texttt{Renewal\_Yang\_Mills\_Mass\_Gap\_Research\_Notes\_V52.tex}.
The V52 parent had \(25\,280\) logical source lines,
\(798\,567\) bytes, and SHA--256
\[
 \texttt{674d2758caf95893ee5534a455afba81fe3e1f7ad1cf1a9d5742ce1c1aa21ff4}.
\]
The next compulsory companion audit is V55.

% =============================================================================
% END APPENDED FIFTH-SERIES V53 MATERIAL
% =============================================================================

\clearpage
\chapter{Fifth-series V54: singlet compression and logarithmic
resummation of a rare wrapped sector}
\label{ch:newV54-log}

\section{Endpoint spectrum and the off-spectrum branch}
\label{sec:newV54-spectrum}

\begin{lemma}[Physical total-spin functional calculus]
\label{lem:newV54-physical-spectrum}
For a tensor block \(B\), decompose
\[
 \bigotimes_{i\in B}V_{R_i}
 =\bigoplus_{n\ge0}\mathcal H_{B,n},
 \qquad
 \mathcal C_B=\sum_{n\ge0}c_nP_{B,n},
 \qquad c_n=n(n+2).
\]
Then
\begin{equation}
 M_{\alpha,B}^{W}
 =L_\alpha(\mathcal C_B)
 =\sum_{n\ge0}{I_{n+1}(\alpha)\over I_1(\alpha)}P_{B,n}.
 \label{eq:newV54-spectral-block}
\end{equation}
The V51 branch point \(\lambda=-1\) lies one unit below the
trivial eigenvalue \(c_0=0\) and four units below the first
nontrivial eigenvalue \(c_1=3\).
\end{lemma}

\begin{proof}
The tensor product is a finite orthogonal sum of \(\SU(2)\)
irreducibles.  On the highest-weight-\(n\) summand the total
Casimir is \(c_n\).  Apply the spectral theorem to
\eqref{eq:newV51-block}, then use
\eqref{eq:newV51-q-mixture}.
\end{proof}

\begin{warning}[Endpoint discreteness is not an interpolation
gap]
\label{warn:newV54-not-gap}
The exact endpoint samples in
\eqref{eq:newV54-spectral-block} do not by themselves remove the
branch obstruction.  The V29 forest interpolates pair couplings,
and the resulting positive quadratic operator generally has
continuous eigenvalue motion.  A proof needs selection rules or
connected cancellation, not merely the list
\(0,3,8,\ldots\) at the endpoint.
\end{warning}

\section{The exact singlet compression test}
\label{sec:newV54-singlet}

\begin{lemma}[Mean subordinated time is thermal]
\label{lem:newV54-mean}
There is an absolute constant \(C_S\) such that, for
\(\alpha\ge1\) and \(s=(2\alpha+2/3)^{-1}\),
\begin{equation}
 \mathbb E_{\rho_\alpha}S\le C_Ss.                            \label{eq:newV54-mean}
\end{equation}
Consequently
\begin{equation}
 \Phi_\alpha(\lambda)\le C_Ss\lambda,\qquad\lambda\ge0.
 \label{eq:newV54-Phi-linear}
\end{equation}
\end{lemma}

\begin{proof}
Set \(r=1\) in \eqref{eq:newV53-two-sector}.  The
nonperturbative term \(e^{-1/s}\) is bounded by an absolute
multiple of \(s\) on the compact range
\(0<s\le3/8\).  This proves \eqref{eq:newV54-mean}.
Since \(\Phi_\alpha\) is concave, vanishes at zero, and
\(\Phi_\alpha'(0)=\mathbb E S\),
\(\Phi_\alpha(\lambda)\le\lambda\Phi_\alpha'(0)\).
\end{proof}

\begin{lemma}[Nonzero scalar compression on a conjugate singlet]
\label{lem:newV54-singlet-compression}
Let \(R\) be a nontrivial irreducible representation with Casimir
\(c_R\), and let \(P_0\) project \(R\otimes R^*\) onto its
one-dimensional invariant subspace.  For
\[
 A_R(x):=C_1+C_2+2x\sum_{a=1}^3D_1^aD_2^a
\]
one has
\begin{equation}
 P_0A_R(x)P_0=2c_R(1-x)P_0.                                  \label{eq:newV54-A-singlet}
\end{equation}
Hence, for every \(r\ge0\),
\begin{equation}
 P_0{d^r\over dx^r}L_\alpha(A_R(x))\bigg|_{x=1}P_0
 =(2c_R)^r\mathbb E S^r\,P_0.                                \label{eq:newV54-repeated-singlet}
\end{equation}
\end{lemma}

\begin{proof}
On the invariant line,
\((D_1+D_2)^2=0\).  Since \(C_1=C_2=c_R\), this gives
\(\sum_aD_1^aD_2^a=-c_R\), proving
\eqref{eq:newV54-A-singlet}.  Differentiate
\[
 L_\alpha(2c_R(1-x))
 =\mathbb E e^{-2c_R(1-x)S}
\]
under the integral.
\end{proof}

\begin{assessment}[What the compression disproves]
\label{ass:newV54-compression}
Equation \eqref{eq:newV54-repeated-singlet} realizes the scalar
V52 moment obstruction inside a physical tensor product.  Thus no
operator theorem based only on positivity and the endpoint
spectral gap can remove it.  The repeated derivative uses the same
pair edge, however, whereas a BKAR tree differentiates each edge
once.  It is therefore a method diagnostic, not a counterexample
to the connected tree card.
\end{assessment}

\begin{theorem}[Exact two-leg singlet connected card]
\label{thm:newV54-two-leg}
Let both external legs carry \(R,R^*\) and put
\(b_R=\sqrt{s\,c_R}\).  On the invariant line,
\begin{equation}
 P_0\kappa_{\alpha,2}^{W}(R,R^*)P_0
 =
 \left(1-e^{-2\Phi_\alpha(c_R)}\right)P_0,                    \label{eq:newV54-kappa2-exact}
\end{equation}
and
\begin{equation}
 \left\|
 P_0\kappa_{\alpha,2}^{W}(R,R^*)P_0
 \right\|
 \le2C_S b_R^2.                                               \label{eq:newV54-kappa2-bound}
\end{equation}
\end{theorem}

\begin{proof}
The joint block moment is the identity on the total singlet.
Each singleton moment is
\(L_\alpha(c_R)\mathbf1=e^{-\Phi_\alpha(c_R)}\mathbf1\).
Subtract their product to obtain
\eqref{eq:newV54-kappa2-exact}.  Now use
\(1-e^{-u}\le u\) and
\eqref{eq:newV54-Phi-linear}:
\[
 1-e^{-2\Phi_\alpha(c_R)}
 \le2\Phi_\alpha(c_R)
 \le2C_Ssc_R=2C_Sb_R^2.
\]
\end{proof}

\begin{corollary}[Connected subtraction restores the first
thermal factor]
\label{cor:newV54-first-factor}
Although the undifferentiated singlet joint moment equals \(1\),
its connected two-leg part is \(O(sc_R)\), with a constant uniform
in \(R\) and \(\alpha\).  Thus singlet survival is compatible with
the strong thermal card once the disconnected product is
subtracted.
\end{corollary}

\section{An abstract logarithmic rare-sector lemma}
\label{sec:newV54-rare-log}

\begin{theorem}[Rare mixture inside a logarithm]
\label{thm:newV54-rare-log}
Let \(q,A>0\), \(B\ge0\), and suppose \(0<s\le s_0\) and
\[
 0<\delta_s\le e^{-q/s}.
\]
Let \(F_0,F_1\) be analytic on a neighbourhood of the closed disk
\[
 D_s:=\{z:|z|\le a/s\},
\]
where \(F_0\) is zero-free and
\begin{equation}
 \left|{F_1(z)\over F_0(z)}\right|
 \le e^{A|z|+B}\qquad(z\in D_s).                              \label{eq:newV54-ratio}
\end{equation}
Choose \(0<a<q/(2A)\).  For all sufficiently small \(s\), the
function
\begin{align}
 G_s(z)
 &:=
 \log\!\left((1-\delta_s)F_0(z)+\delta_sF_1(z)\right)
 -\log F_0(z)-\log(1-\delta_s)                               \notag\\
 &=\log\!\left(
 1+{\delta_s\over1-\delta_s}{F_1(z)\over F_0(z)}
 \right)                                                     \label{eq:newV54-G}
\end{align}
is analytic on \(D_s\), with a coherent logarithm, and there are
constants \(C,c>0\) independent of \(m,s\) such that
\begin{equation}
 {|G_s^{(m)}(0)|\over m!}
 \le C e^{-c/s}(Cs)^m,\qquad m\ge0.                           \label{eq:newV54-log-coeff}
\end{equation}
\end{theorem}

\begin{proof}
For small \(s\), \(\delta_s\le1/2\).  On \(D_s\),
\begin{align*}
 \left|
 {\delta_s\over1-\delta_s}{F_1\over F_0}
 \right|
 &\le
 2\exp\left\{-{q\over s}+{Aa\over s}+B\right\}\\
 &\le C_Be^{-q/(2s)}.
\end{align*}
After decreasing \(s_0\), this is at most \(1/2\).  Hence the
principal series for \(\log(1+w)\) converges uniformly and defines
the coherent logarithm in \eqref{eq:newV54-G}; also
\[
 \sup_{D_s}|G_s|\le C e^{-q/(2s)}.
\]
Cauchy's estimate on the circle of radius \(a/s\) gives
\[
 {|G_s^{(m)}(0)|\over m!}
 \le C e^{-q/(2s)}(s/a)^m.
\]
Rename the constants.
\end{proof}

\begin{corollary}[Bernoulli rare-event model]
\label{cor:newV54-Bernoulli}
If \(X_s\) is Bernoulli with
\(\mathbb P(X_s=1)=\delta_s\le e^{-q/s}\), then every cumulant
\(\kappa_m(X_s)\) obeys
\begin{equation}
 {|\kappa_m(X_s)|\over m!}
 \le C e^{-c/s}(Cs)^m.                                       \label{eq:newV54-Bernoulli-bound}
\end{equation}
The zeros of its moment-generating function are exactly
\begin{equation}
 z_k=\log{1-\delta_s\over\delta_s}+(2k+1)i\pi,
 \qquad k\in\mathbb Z,                                       \label{eq:newV54-Bernoulli-zeros}
\end{equation}
so the connected logarithm has analytic radius
\(\asymp\log(1/\delta_s)\gtrsim s^{-1}\), not radius \(O(1)\).
\end{corollary}

\begin{proof}
Apply \Cref{thm:newV54-rare-log} with
\(F_0=1\), \(F_1=e^z\), \(A=1\), and \(B=0\).
The connected generating function is
\(\log(1-\delta_s+\delta_se^z)\), up to its constant value.
Solving \(1-\delta_s+\delta_se^z=0\) gives
\eqref{eq:newV54-Bernoulli-zeros}.
\end{proof}

\begin{assessment}[Why linearizing the rare sector was fatal]
\label{ass:newV54-linearization}
Expanding the Bernoulli logarithm first in \(\delta_s\) produces
the term \(\delta_s(e^z-1)\), whose \(m\)-th derivative is
\(\delta_s\) and appears to lose against \(s^m\) at large \(m\).
The exact logarithm instead moves its first zero to distance
\(\log(1/\delta_s)\asymp1/s\), yielding
\eqref{eq:newV54-Bernoulli-bound}.  Thus the V53 wrapped
post-differentiation no-go is compatible with an all-order
connected card: the nonperturbative sector must be resummed inside
the logarithm before coefficients are extracted.
\end{assessment}

\section{The positive logarithmic realization gate}
\label{sec:newV54-PLRC}

\begin{definition}[Positive logarithmic rare-sector card, PLRC]
\label{def:newV54-PLRC}
PLRC consists of a positive decomposition of the exact Wilson
source measure or the tilted-time law,
\begin{equation}
 \nu_\alpha=(1-\delta_\alpha^{\rm tail})\nu_\alpha^{\rm loc}
 +\delta_\alpha^{\rm tail}\nu_\alpha^{\rm tail},              \label{eq:newV54-positive-split}
\end{equation}
for which:
\begin{enumerate}[label=\textup{(\roman*)},leftmargin=3.5em]
\item \(\delta_\alpha^{\rm tail}\le Ce^{-q/s}\);
\item the local connected source obeys the V50--V53 compact
      all-arity card;
\item the conditional source transforms satisfy a multivariable,
      completely bounded analogue of
      \eqref{eq:newV54-ratio} on a polydisc of radius \(a/s\);
\item the local conditional transform is zero-free there; and
\item the exact logarithm of the mixture is taken before the
      M\"obius/forest derivatives.
\end{enumerate}
\end{definition}

\begin{proposition}[PLRC implies WSDC]
\label{prop:newV54-PLRC-WSDC}
If PLRC holds with constants uniform in the representation labels
after scaling each source direction by its physical weight, then
the wrapped source-domain card WSDC follows.
\end{proposition}

\begin{proof}
Factor the source transform exactly as in
\eqref{eq:newV54-G}, with \(F_0\) the local conditional transform
and \(F_1\) the tail conditional transform.  Apply the
multivariable completely bounded version of
\Cref{thm:newV54-rare-log}.  Cauchy extraction on the
weight-scaled polydisc supplies one factor \(Cs\) per genuinely
wrapped connected derivative, while the local sector is handled
by PLRC(ii).  Taking the logarithm before differentiation is
PLRC(v), so no term of the forbidden form
\(e^{-1/s}s^{-m}\) is produced.  This is precisely WSDC.
\end{proof}

\begin{warning}[The Schl\"afli split cannot serve as PLRC]
\label{warn:newV54-Schlafli-not-positive}
The compact/wrapped decomposition of V53 is signed and therefore
does not define the positive conditional measures in
\eqref{eq:newV54-positive-split}.  PLRC must instead be obtained
from a positive event of the Hartman--Watson time, or directly from
a positive compact-group partition.  Reinterpreting
\(\mathcal W_\alpha\) itself as a probability tail would be
incorrect.
\end{warning}

\begin{programme}[V55: companion audit and positive-tail test]
\label{prog:newV54-V55}
\begin{enumerate}[leftmargin=3em]
\item Perform the compulsory newest-SM/newest-NS audit before
      extending the argument.
\item Split the positive tilted-time law at a fixed time threshold
      and seek a uniform \(e^{-q/s}\) upper bound on the tail
      probability.
\item Express the two conditional heat-source transforms and
      prove the bounded-observable ratio estimate.
\item Test zero-freeness of the local conditional transform on the
      required \(1/s\) source scale.
\item If time truncation does not provide zero-freeness, return
      upstream to a compact-group geodesic/antipodal positive
      partition and preserve the exact connected logarithm.
\end{enumerate}
\end{programme}

\begin{firewall}[Claim boundary after V54]
\label{fire:newV54-boundary}
V54 proves the exact physical endpoint spectral calculus, a
uniform thermal mean-time bound, nonzero singlet compression, the
correct two-leg connected singlet card, and a rigorous logarithmic
rare-mixture lemma showing how \(e^{-q/s}\) produces an all-order
thermal analytic radius when resummed before differentiation.  It
isolates PLRC as the positive realization required for WSDC.

V54 does not prove PLRC, WSDC, BRSC, HWSFC, UGCGC, thermal
connected WPRAC, all-arity RGSC, the raw Wilson aggregate strong
card, all-order strong MTA, WUFC, CPM, cutoff removal, reflection
positivity, Osterwalder--Schrader reconstruction, a continuum
Yang--Mills measure, or a positive Hamiltonian gap.  The full
Yang--Mills mass gap has not been proved.
\end{firewall}

\begin{theorem}[V54 result ledger]
\label{thm:newV54-results}
V54:
\begin{enumerate}[label=\textup{(V54.\arabic*)},leftmargin=6.5em]
\item locates the V51 branch relative to the physical total-spin
      spectrum;
\item proves \(\mathbb ES=O(s)\) and
      \(\Phi_\alpha(\lambda)\le Cs\lambda\);
\item proves that a conjugate singlet has nonzero scalar
      interpolation compression;
\item proves the exact strong two-leg singlet bound;
\item proves the rare-mixture logarithm lemma;
\item computes the Bernoulli source zeros at distance
      \(\asymp1/s\);
\item explains why connected logarithmic resummation can overcome
      the V53 derivative no-go; and
\item formulates PLRC and proves its sufficiency for WSDC.
\end{enumerate}
\end{theorem}

\begin{proof}
Use
\Cref{lem:newV54-physical-spectrum,
lem:newV54-mean,
lem:newV54-singlet-compression,
thm:newV54-two-leg,
thm:newV54-rare-log,
cor:newV54-Bernoulli,
def:newV54-PLRC,
prop:newV54-PLRC-WSDC}.
\end{proof}

\paragraph{Parent-version record.}
V54 was formed from
\texttt{Renewal\_Yang\_Mills\_Mass\_Gap\_Research\_Notes\_V53.tex}.
The V53 parent had \(25\,718\) logical source lines,
\(812\,998\) bytes, and SHA--256
\[
 \texttt{e34c14c00621502eb74c6b1ed2397dea9c4ca42cfe67fcde6cddebd62178eeaf}.
\]
The next compulsory companion audit is V55.

% =============================================================================
% END APPENDED FIFTH-SERIES V54 MATERIAL
% =============================================================================

\clearpage
\chapter{Fifth-series V55: companion audit and a positive
Hartman--Watson thermal-tail split}
\label{ch:newV55-tail}

\section{Compulsory companion audit}
\label{sec:newV55-audit}

\begin{audit}[Newest SM and NS snapshots]
\label{audit:newV55-companions}
The companion files were inspected read-only.  The stable SM snapshot
used here was
\[
 \texttt{Renewal\_SM\_Einstein\_Continuum\_Research\_Notes\_V75.tex},
\]
with \(24\,775\) logical source lines, \(951\,938\) bytes, and
SHA--256
\[
 \texttt{6678a93c8e0910e32898c18e97dfd3e8d58e91a07132ef2d3654062c3431a0d2}.
\]
The NS programme advanced during the audit.  The final atomic snapshot
read was
\[
 \texttt{Renewal\_NS\_Stability\_Research\_Notes\_V79.tex},
\]
with \(35\,559\) logical source lines, \(1\,487\,910\) bytes, and
SHA--256
\[
 \texttt{1a6210f24b58249c12f23c8b9df2dabbaffab6847cf4f69fed97db2ca91510d9}.
\]
Neither companion file was modified.
\end{audit}

\begin{theorem}[V55 audit transfer]
\label{thm:newV55-audit-transfer}
The useful transfers are structural:
\begin{enumerate}[label=\textup{(A\arabic*)},leftmargin=4.5em]
\item SM V75 confirms that a small or exceptional sector must remain
      assembled until its exact contour, constraint, and positivity
      mechanism has been formed; moments of a prematurely separated
      sector do not replace that construction.
\item NS V79 shows that removing a local concentration window can move
      the defect to compact or spread satellite scales.  For YM,
      truncating the Hartman--Watson time must therefore retain the
      conditional tail inside the exact source logarithm; it may not be
      dropped as negligible.
\item Neither companion provides a Yang--Mills heat-kernel constant,
      a Bessel estimate, a source zero-free domain, or a reflection
      positivity theorem.
\end{enumerate}
\end{theorem}

\begin{proof}
SM V75's conformal trace analysis distinguishes convergence of an
isolated rotated Gaussian from positivity of the complete reflected
observable algebra.  NS V79 distinguishes compact satellites, which
force terminal bumps, from spread satellites, which evade the local
implication.  Both distinctions say that the exceptional sector must be
kept with its global assembly.  Their operators and hypotheses are
unrelated to the Wilson endpoint, so no numerical estimate transfers.
\end{proof}

\section{Antipodal heat positivity}
\label{sec:newV55-antipode}

\begin{theorem}[Compact heat-kernel lower bound, standard input]
\label{thm:newV55-heat-lower}
For the normalized \(\SU(2)\) heat kernel used in V29, there are
constants \(a_H,c_H,t_H>0\) such that
\begin{equation}
 h_t(-\mathbf1)\ge
 c_Ht^{-3/2}e^{-a_H/t},
 \qquad 0<t\le t_H.                                           \label{eq:newV55-heat-lower}
\end{equation}
Furthermore
\begin{equation}
 \inf_{t\ge t_H}h_t(-\mathbf1)>0.                             \label{eq:newV55-heat-large}
\end{equation}
\end{theorem}

\begin{proof}
The first row is the standard two-sided Gaussian lower estimate for the
heat kernel on a compact connected three-dimensional Riemannian
manifold, evaluated at the fixed pair
\((\mathbf1,-\mathbf1)\).  It follows from the near-diagonal parametrix
and a finite semigroup chain along a minimizing geodesic; the constants
absorb the metric normalization and the cut locus.  This classical
compact heat-kernel theorem is an explicit external analytic input.

For the second row, strict positivity and continuity give a positive
minimum on every compact time interval
\([t_H,T]\).  The spectral expansion converges uniformly to the
constant equilibrium density \(1\) as \(t\to\infty\); choose \(T\) so
that the kernel is at least \(1/2\) thereafter.
\end{proof}

\begin{lemma}[Uniform antipodal lower bound above a thermal threshold]
\label{lem:newV55-threshold-lower}
Fix \(L>0\).  There is \(c_L>0\) such that, whenever
\(Ls\le t_H\),
\begin{equation}
 \inf_{t\ge Ls}h_t(-\mathbf1)
 \ge c_L e^{-a_H/(Ls)}.                                      \label{eq:newV55-threshold-lower}
\end{equation}
\end{lemma}

\begin{proof}
For \(Ls\le t\le t_H\), use
\eqref{eq:newV55-heat-lower} and
\[
 t^{-3/2}\ge t_H^{-3/2},\qquad
 e^{-a_H/t}\ge e^{-a_H/(Ls)}.
\]
For \(t\ge t_H\), use \eqref{eq:newV55-heat-large}; since the
exponential on the right is at most \(1\), decrease \(c_L\) to cover
both ranges.
\end{proof}

\section{Exact Wilson antipodal cost}
\label{sec:newV55-Wilson-antipode}

\begin{lemma}[Normalized Wilson density at the antipode]
\label{lem:newV55-W-antipode}
In class-angle coordinates \(0\le\theta\le\pi\), normalized Haar
measure has radial density
\((2/\pi)\sin^2\theta\,d\theta\), and
\begin{equation}
 w_\alpha(\theta)
 ={\alpha\over2I_1(\alpha)}e^{\alpha\cos\theta}.               \label{eq:newV55-W-density}
\end{equation}
Consequently, for \(\alpha\ge1\),
\begin{equation}
 w_\alpha(-\mathbf1)
 ={\alpha e^{-\alpha}\over2I_1(\alpha)}
 \le C_W\alpha^{3/2}e^{-2\alpha}.                             \label{eq:newV55-W-antipode}
\end{equation}
\end{lemma}

\begin{proof}
The Bessel recurrence and the standard integral formula give
\[
 \int_{\SU(2)}e^{\alpha\cos\theta}\,dU
 =I_0(\alpha)-I_2(\alpha)
 ={2I_1(\alpha)\over\alpha},
\]
which proves \eqref{eq:newV55-W-density}.  Evaluate at
\(\theta=\pi\), then divide by the lower bound for \(I_1\) in
\eqref{eq:newV52-I-bounds}; the remaining numerical factor is
absorbed into \(C_W\).
\end{proof}

\section{Positive thermal-time tail theorem}
\label{sec:newV55-positive-tail}

\begin{theorem}[Positive Hartman--Watson tail at scale \(Ls\)]
\label{thm:newV55-positive-tail}
Choose \(L>2a_H\), where \(a_H\) is from
\Cref{thm:newV55-heat-lower}.  For
\[
 S\sim\rho_\alpha,\qquad s=(2\alpha+2/3)^{-1},
\]
there are constants \(C_L,q_L>0\), independent of
\(\alpha\ge1\), such that
\begin{equation}
 \delta_{\alpha,L}
 :=\mathbb P_{\rho_\alpha}\{S\ge Ls\}
 \le C_Le^{-q_L/s}.                                          \label{eq:newV55-positive-tail}
\end{equation}
\end{theorem}

\begin{proof}
By the exact positive mixture
\eqref{eq:newV51-density-mixture},
\[
 w_\alpha(-\mathbf1)
 =\int_0^\infty h_t(-\mathbf1)\rho_\alpha(dt)
 \ge
 \delta_{\alpha,L}\inf_{t\ge Ls}h_t(-\mathbf1).
\]
This pointwise use of the measure identity is legitimate: the standard
off-diagonal Gaussian upper bound is uniform on a neighbourhood of the
antipode for small \(t\), while the large-time spectral expansion is
uniform there.  Hence the heat mixture is continuous at
\(-\mathbf1\), and its almost-everywhere equality with the smooth Wilson
density extends to that point.
For small \(s\), combine
\Cref{lem:newV55-threshold-lower} and
\eqref{eq:newV55-W-antipode}:
\begin{align*}
 \delta_{\alpha,L}
 &\le
 C\alpha^{3/2}
 \exp\left\{-2\alpha+{a_H\over Ls}\right\}\\
 &\le
 Cs^{-3/2}
 \exp\left\{-{1-a_H/L\over s}\right\}.
\end{align*}
Since \(L>2a_H\), the exponent coefficient
\(\gamma_L:=1-a_H/L\) is positive.  Absorb
\(s^{-3/2}\) into
\(\exp\{\gamma_L/(2s)\}\) for sufficiently small \(s\), giving
\eqref{eq:newV55-positive-tail} with
\(q_L=\gamma_L/2\).  On the remaining compact interval of
\(s\)-values, enlarge \(C_L\) and use
\(\delta_{\alpha,L}\le1\).
\end{proof}

\begin{corollary}[Exact positive local/tail heat mixture]
\label{cor:newV55-positive-split}
For all sufficiently large \(\alpha\), define conditional
probability laws
\begin{align}
 \rho_{\alpha,L}^{\rm loc}
 &:={\mathbf1_{\{S<Ls\}}\rho_\alpha\over1-\delta_{\alpha,L}},
                                                               \label{eq:newV55-rho-loc}\\
 \rho_{\alpha,L}^{\rm tail}
 &:={\mathbf1_{\{S\ge Ls\}}\rho_\alpha\over\delta_{\alpha,L}}.
                                                               \label{eq:newV55-rho-tail}
\end{align}
Then the Wilson measure has the exact positive decomposition
\begin{equation}
 w_\alpha dU
 =(1-\delta_{\alpha,L})\,\nu_{\alpha,L}^{\rm loc}
 +\delta_{\alpha,L}\,\nu_{\alpha,L}^{\rm tail},               \label{eq:newV55-measure-split}
\end{equation}
where
\begin{align}
 \nu_{\alpha,L}^{\rm loc}(dU)
 &=\int h_t(U)dU\,\rho_{\alpha,L}^{\rm loc}(dt),               \label{eq:newV55-nu-loc}\\
 \nu_{\alpha,L}^{\rm tail}(dU)
 &=\int h_t(U)dU\,\rho_{\alpha,L}^{\rm tail}(dt).              \label{eq:newV55-nu-tail}
\end{align}
This closes PLRC(i) with a genuinely positive, rather than
Schl\"afli-signed, split.
\end{corollary}

\begin{proof}
Partition the positive integral
\eqref{eq:newV51-density-mixture} into the two time events and
normalize each nonzero part.  Apply
\Cref{thm:newV55-positive-tail}.
\end{proof}

\section{The local conditional all-order time card}
\label{sec:newV55-local-card}

\begin{theorem}[Truncated-time Duhamel card]
\label{thm:newV55-local-Duhamel}
Let \(A(x)=A_0+\sum_{j=1}^r x_jV_j\ge0\) at the evaluation point,
as in \Cref{lem:newV51-Duhamel}.  For the local conditional block
transform
\[
 L_{\alpha,L}^{\rm loc}(A)
 :=\int e^{-tA}\rho_{\alpha,L}^{\rm loc}(dt),
\]
one has
\begin{equation}
 \left\|
 \partial_{x_1}\cdots\partial_{x_r}
 L_{\alpha,L}^{\rm loc}(A(x))
 \right\|
 \le (Ls)^r\prod_{j=1}^r\|V_j\|.                             \label{eq:newV55-local-Duhamel}
\end{equation}
The bound is completely amplified and has one exponential base after
thermal normalization.
\end{theorem}

\begin{proof}
Repeat the ordered-simplex proof of
\Cref{lem:newV51-Duhamel}.  Contractivity removes the semigroup
factors, the sum of \(r!\) simplex integrals has total bound \(t^r\),
and \(t<Ls\) on the support of the conditional law.  Integration gives
\(\mathbb E_{\rm loc}t^r\le(Ls)^r\).  Tensoring with identities changes
nothing.
\end{proof}

\begin{corollary}[Two PLRC rows are closed]
\label{cor:newV55-two-rows}
The positive decomposition
\eqref{eq:newV55-measure-split} closes PLRC(i), and
\Cref{thm:newV55-local-Duhamel} closes the time-ordering part of
PLRC(ii).  No post-differentiation term of size
\(e^{-1/s}s^{-r}\) is required in the local sector.
\end{corollary}

\section{A scalar zero-free criterion for the conditional source}
\label{sec:newV55-zero-free}

\begin{lemma}[Support-diameter zero-free strip]
\label{lem:newV55-zero-free-strip}
Let \(X\) be real with essential range in \([a,b]\) under a
probability law, and put
\[
 F(z)=\mathbb E e^{zX}.
\]
Then \(F\) is zero-free in
\begin{equation}
 |\Im z|<{\pi\over b-a}.                                      \label{eq:newV55-zero-free-strip}
\end{equation}
More precisely, if \(m=(a+b)/2\), then
\begin{equation}
 \left|e^{-i(\Im z)m}F(z)\right|
 \ge
 \cos\left({|\Im z|(b-a)\over2}\right)
 \mathbb E e^{(\Re z)X}.                                     \label{eq:newV55-zero-free-lower}
\end{equation}
\end{lemma}

\begin{proof}
Write \(z=x+iy\).  After multiplication by \(e^{-iym}\), every
integrand has argument \(y(X-m)\), whose absolute value is strictly
less than \(\pi/2\) under \eqref{eq:newV55-zero-free-strip}.
Its real part is at least
\[
 e^{xX}\cos\left({|y|(b-a)\over2}\right)>0.
\]
Integrate and use that the modulus is at least the real part.
\end{proof}

\begin{corollary}[When the rare-log lemma applies at thermal scale]
\label{cor:newV55-oscillation}
If the local conditional observable has support diameter
\[
 b-a\le C_Xs,
\]
then its transform is zero-free on a strip of imaginary width
\(\pi/(C_Xs)\).  On every narrower disk of radius \(a_0/s\), its
inverse and the tail/local transform ratio have exponential growth of
the form required in \eqref{eq:newV54-ratio}, provided the tail
observable has a uniform absolute bound.
\end{corollary}

\begin{proof}
Apply \Cref{lem:newV55-zero-free-strip}.  On a narrower strip the cosine
in \eqref{eq:newV55-zero-free-lower} is bounded below.  If
\(|X|\le M\), then both conditional transforms are bounded above by
\(e^{M|z|}\), while the local denominator is bounded below by a fixed
cosine times \(e^{-M|\Re z|}\).  Their ratio is at most
\(Ce^{2M|z|}\).
\end{proof}

\begin{warning}[Time truncation does not imply support-diameter
truncation]
\label{warn:newV55-time-not-space}
Although \(t<Ls\) gives the exact moment bound
\eqref{eq:newV55-local-Duhamel}, a heat kernel at any positive time has
full support on \(\SU(2)\).  Thus the deterministic
\(O(s)\) support-diameter hypothesis of
\Cref{cor:newV55-oscillation} does not follow from time truncation.
Proving the completely bounded multivariable zero-free row remains a
separate source problem.
\end{warning}

\section{The remaining local-zero and tail-log card}
\label{sec:newV55-LZTC}

\begin{definition}[Local-zero/tail-log card, LZTC]
\label{def:newV55-LZTC}
LZTC requires, for the exact positive split
\eqref{eq:newV55-measure-split}:
\begin{enumerate}[label=\textup{(\roman*)},leftmargin=3.5em]
\item a zero-free completely bounded local conditional source transform
      on the weight-scaled polydomain needed for
      \(K^m(\prod b_i)B^{m-2}\);
\item an inverse bound for that local transform;
\item a tail/local conditional transform-ratio bound growing at most
      exponentially in the source norm;
\item application of the exact mixture logarithm before all source and
      forest derivatives; and
\item retention of the tail conditional measure, in accordance with the
      SM/NS audit, rather than deletion of its exponentially small
      probability.
\end{enumerate}
\end{definition}

\begin{proposition}[LZTC completes PLRC]
\label{prop:newV55-LZTC-PLRC}
The positive tail theorem, the local Duhamel card, and LZTC imply PLRC.
Consequently, by \Cref{prop:newV54-PLRC-WSDC}, they imply WSDC.
\end{proposition}

\begin{proof}
PLRC(i) is \Cref{thm:newV55-positive-tail}.
The local time-ordering row of PLRC(ii) is
\Cref{thm:newV55-local-Duhamel}; its remaining connected source
resummation is LZTC(i).  LZTC(ii)--(iii) are the multivariable hypotheses
of the rare-log theorem, and LZTC(iv) supplies its required order of
operations.  LZTC(v) preserves the exact positive mixture.  These are
exactly the PLRC rows.
\end{proof}

\begin{programme}[V56: local conditional zero-free attack]
\label{prog:newV55-V56}
\begin{enumerate}[leftmargin=3em]
\item Write the local conditional tensor source as a time mixture of
      finite-dimensional heat semigroup exponentials.
\item Test zero-freeness first on commuting Cartan sources and locate
      their exact nearest zeros.
\item Compare those zeros with the support-diameter strip
      \eqref{eq:newV55-zero-free-strip}.
\item Test noncommuting sources by numerical-range accretivity before
      using determinants or scalar traces.
\item If a \(1/s\) zero-free polydomain fails even in a scalar physical
      sector, abandon PLRC and return to the V50 real connected-grade
      forest without complex source division.
\end{enumerate}
\end{programme}

\begin{firewall}[Claim boundary after V55]
\label{fire:newV55-boundary}
V55 completes the compulsory companion audit, proves from antipodal heat
positivity that the positive Hartman--Watson event \(S\ge Ls\) has
probability \(O(e^{-q/s})\), constructs the exact positive
local/tail Wilson decomposition, and proves an all-order local
conditional Duhamel card.  It also proves a scalar support-diameter
zero-free criterion and isolates LZTC.

V55 does not prove LZTC, PLRC, WSDC, BRSC, HWSFC, UGCGC, thermal
connected WPRAC, all-arity RGSC, the raw Wilson aggregate strong card,
all-order strong MTA, WUFC, CPM, cutoff removal, reflection positivity,
Osterwalder--Schrader reconstruction, a continuum Yang--Mills measure,
or a positive Hamiltonian gap.  The full Yang--Mills mass gap has not
been proved.
\end{firewall}

\begin{theorem}[V55 result ledger]
\label{thm:newV55-results}
V55:
\begin{enumerate}[label=\textup{(V55.\arabic*)},leftmargin=6.5em]
\item completes the required SM/NS companion audit;
\item proves a uniform antipodal heat-kernel lower bound above \(Ls\);
\item computes the exact normalized Wilson antipodal density;
\item proves the positive \(e^{-q/s}\) tilted-time tail theorem;
\item constructs the exact positive conditional Wilson split;
\item proves the local truncated-time Duhamel card at every arity;
\item proves a scalar support-diameter zero-free strip;
\item closes two PLRC rows; and
\item isolates LZTC as the remaining logarithmic source gate.
\end{enumerate}
\end{theorem}

\begin{proof}
Use
\Cref{audit:newV55-companions,
thm:newV55-audit-transfer,
thm:newV55-heat-lower,
lem:newV55-threshold-lower,
lem:newV55-W-antipode,
thm:newV55-positive-tail,
cor:newV55-positive-split,
thm:newV55-local-Duhamel,
lem:newV55-zero-free-strip,
def:newV55-LZTC}.
\end{proof}

\paragraph{Parent-version record.}
V55 was formed from
\texttt{Renewal\_Yang\_Mills\_Mass\_Gap\_Research\_Notes\_V54.tex}.
The V54 parent had \(26\,118\) logical source lines,
\(827\,061\) bytes, and SHA--256
\[
 \texttt{cc98f8c79e51e5e90c2a31ebaa734f0c2c852fc4fff433b1d3dbac6265229da6}.
\]
The next compulsory companion audit is V60.

% =============================================================================
% END APPENDED FIFTH-SERIES V55 MATERIAL
% =============================================================================

\clearpage
\chapter{Fifth-series V56: conditional source invertibility and
the surviving homogeneity deficit}
\label{ch:newV56-source}

\section{Reuse of the V44 source mechanism}
\label{sec:newV56-reuse}

The zero-free problem must not be restarted from scratch.  V44 already
proved Banach-algebra invertibility of the complete squarefree operator
source on the weighted simplex
\(\sum_i|z_i|b_i\le\eta\).  The first task is to verify that its radial
exponential-moment proof remains uniform for the positive local
conditional law of V55.

\begin{definition}[Joint subordinated heat law]
\label{def:newV56-joint}
Let \(\widehat\nu_\alpha\) be the probability law on
\((0,\infty)\times\SU(2)\) given by
\begin{equation}
 \widehat\nu_\alpha(dS,dU)
 :=\rho_\alpha(dS)h_S(U)dU.                                  \label{eq:newV56-joint}
\end{equation}
Its \(U\)-marginal is the Wilson measure.  Write
\[
 E_{\rm loc}:=\{S<Ls\},\qquad E_{\rm tail}:=\{S\ge Ls\}.
\]
\end{definition}

\begin{lemma}[Local conditional radial exponential moment]
\label{lem:newV56-local-radial}
Let \(\theta(U)\in[0,\pi]\) be the class distance and
\(Y_s=\theta(U)/\sqrt s\).  There are numerical
\(\lambda_L,C_L>0\) such that, for all sufficiently large
\(\alpha\),
\begin{equation}
 \mathbb E_{\widehat\nu_\alpha}
 \left[e^{\lambda_LY_s}\mid E_{\rm loc}\right]\le C_L.
 \label{eq:newV56-local-radial}
\end{equation}
\end{lemma}

\begin{proof}
The \(U\)-marginal is Wilson, so the V44 endpoint case of
\eqref{eq:newV44-radial-exp} gives
\[
 \mathbb E_{\widehat\nu_\alpha}e^{\lambda_0Y_s}\le C.
\]
By \Cref{thm:newV55-positive-tail},
\(\mathbb P(E_{\rm loc})=1-\delta_{\alpha,L}\ge1/2\) for large
\(\alpha\).  Restrict the numerator to \(E_{\rm loc}\), divide by
that probability, and take \(\lambda_L=\lambda_0\).
\end{proof}

\begin{definition}[Conditional squarefree sources]
\label{def:newV56-sources}
With the V44 increments
\(A_i(U)=D^{R_i}(U)-I_{R_i}\), define
\begin{align}
 \mathbb Z_{{\rm loc},m}(\boldsymbol z)
 &:=
 \mathbb E_{\widehat\nu_\alpha}
 \left[
 \bigotimes_{i=1}^m(I+z_iA_i(U))
 \,\middle|\,E_{\rm loc}\right],                              \label{eq:newV56-Zloc}\\
 \mathbb Z_{{\rm tail},m}(\boldsymbol z)
 &:=
 \mathbb E_{\widehat\nu_\alpha}
 \left[
 \bigotimes_{i=1}^m(I+z_iA_i(U))
 \,\middle|\,E_{\rm tail}\right].                             \label{eq:newV56-Ztail}
\end{align}
Then
\begin{equation}
 \mathbb Z_{W,m}
 =(1-\delta_{\alpha,L})\mathbb Z_{{\rm loc},m}
 +\delta_{\alpha,L}\mathbb Z_{{\rm tail},m}.                  \label{eq:newV56-Zsplit}
\end{equation}
\end{definition}

\begin{theorem}[Local conditional source invertibility]
\label{thm:newV56-local-invertible}
There is \(\eta_L>0\), independent of
\(m,\alpha\) and the thermal representations, such that
\begin{equation}
 X(\boldsymbol z):=\sum_{i=1}^m|z_i|b_i\le\eta_L             \label{eq:newV56-simplex}
\end{equation}
implies
\begin{equation}
 \|\mathbb Z_{{\rm loc},m}(\boldsymbol z)-I\|_{\rm cb}
 \le {1\over8}.                                               \label{eq:newV56-Zloc-close}
\end{equation}
Hence its principal logarithm and inverse are uniformly bounded on
that simplex.
\end{theorem}

\begin{proof}
The V44 geodesic increment estimate gives
\[
 \left\|
 \bigotimes_i(I+z_iA_i(U))-I
 \right\|_{\rm cb}
 \le e^{CY_sX(\boldsymbol z)}-1.
\]
Take the local conditional expectation and apply
\Cref{lem:newV56-local-radial}.  Uniform integrability and dominated
convergence as \(X\downarrow0\) permit \(\eta_L\) to be chosen so
that the expectation is at most \(1/8\).  The Neumann and logarithm
series then converge in the completely amplified Banach algebra.  This
is exactly the V44 mechanism with its endpoint expectation replaced by
the verified local conditional expectation.
\end{proof}

\section{Exact singleton and conjugate-singlet zeros}
\label{sec:newV56-low-zeros}

\begin{definition}[Local conditional multiplier deficit]
\label{def:newV56-aR}
For an irreducible \(R\) with Casimir \(c_R\), set
\begin{equation}
 q_{R}^{\rm loc}
 :=\mathbb E_{\widehat\nu_\alpha}
    [e^{-Sc_R}\mid E_{\rm loc}],
 \qquad
 a_R:=1-q_R^{\rm loc}.
 \label{eq:newV56-aR}
\end{equation}
\end{definition}

\begin{lemma}[Thermal size of the local deficit]
\label{lem:newV56-aR}
If \(b_R=\sqrt{sc_R}<1\), then
\begin{equation}
 0\le a_R\le Lsc_R=Lb_R^2.                                   \label{eq:newV56-aR-bound}
\end{equation}
\end{lemma}

\begin{proof}
On \(E_{\rm loc}\), \(0<S<Ls\).  Use
\(1-e^{-u}\le u\):
\[
 a_R
 =\mathbb E[1-e^{-Sc_R}\mid E_{\rm loc}]
 \le c_R\mathbb E[S\mid E_{\rm loc}]
 \le Lsc_R.
\]
\end{proof}

\begin{proposition}[Exact singleton source zero]
\label{prop:newV56-singleton-zero}
The local conditional singleton source is
\begin{equation}
 \mathbb Z_{{\rm loc},1}(z)
 =(1-a_Rz)I_R.                                                \label{eq:newV56-singleton-source}
\end{equation}
If \(a_R>0\), its sole zero is \(z=a_R^{-1}\), and
\begin{equation}
 |z|\ge {1\over Lb_R^2}.                                     \label{eq:newV56-singleton-zero}
\end{equation}
\end{proposition}

\begin{proof}
Centrality gives
\(\mathbb E_{\rm loc}D^R(U)=q_R^{\rm loc}I_R\).
Insert \(A_R=D^R-I_R\), then use
\eqref{eq:newV56-aR-bound}.
\end{proof}

\begin{theorem}[Exact conjugate-singlet source polynomial]
\label{thm:newV56-pair-polynomial}
Let \(P_0\) project \(R\otimes R^*\) onto its invariant line.  Then
\begin{equation}
 P_0\mathbb Z_{{\rm loc},2}(z_1,z_2)P_0
 =
 \left[
 1-a_R(z_1+z_2)+2a_Rz_1z_2
 \right]P_0.                                                  \label{eq:newV56-pair-polynomial}
\end{equation}
On the diagonal \(z_1=z_2=z\), the two zeros have common modulus
\begin{equation}
 |z|={1\over\sqrt{2a_R}}
 \ge {1\over\sqrt{2L}\,b_R}.                                 \label{eq:newV56-pair-zero}
\end{equation}
\end{theorem}

\begin{proof}
The constant and linear coefficients are \(1\) and
\(-a_R\).  For each \(U\), the vector spanning the invariant line is
fixed by \(D^R(U)\otimes D^{R^*}(U)\).  After projection and
expectation,
\[
 P_0\mathbb E[A_1(U)\otimes A_2(U)]P_0
 =(1-q_R^{\rm loc}-q_R^{\rm loc}+1)P_0=2a_RP_0.
\]
This proves the polynomial.  Its diagonal form is
\(1-2a_Rz+2a_Rz^2\).  For \(0<a_R\le1\), its discriminant is
nonpositive and the product of its conjugate roots is
\((2a_R)^{-1}\); hence each root has the displayed modulus.
Use \eqref{eq:newV56-aR-bound}.
\end{proof}

\begin{corollary}[Correct binary weighted bidisc]
\label{cor:newV56-binary-bidisc}
There is \(r_L>0\) such that
\[
 |z_1|,|z_2|\le {r_L\over b_R}
\]
implies that the scalar in
\eqref{eq:newV56-pair-polynomial} lies in the disk
\(|w-1|\le1/2\), uniformly in \(R,\alpha\).
\end{corollary}

\begin{proof}
On that bidisc, using \(a_R\le Lb_R^2\) and \(b_R<1\),
\[
 |a_R(z_1+z_2)-2a_Rz_1z_2|
 \le2Lr_Lb_R+2Lr_L^2
 \le2L(r_L+r_L^2).
\]
Choose \(r_L\) so that the last expression is at most \(1/2\).
\end{proof}

\begin{assessment}[The first physical zeros have the right scale]
\label{ass:newV56-zero-scale}
The singleton zero is at least of order \(b_R^{-2}\), and the first
connected conjugate-singlet zeros are at least of order \(b_R^{-1}\).
The latter is exactly the V44 weighted-source radius needed for the
binary card.  No extra fixed-distance zero is hidden in the local
conditional singlet sector.
\end{assessment}

\section{Tail-weighted source control without commutative
factorization}
\label{sec:newV56-tail-log}

\begin{lemma}[Exponentially small unnormalized tail source]
\label{lem:newV56-tail-source}
After decreasing \(\eta_L\), there are \(C,q>0\) such that on
\eqref{eq:newV56-simplex},
\begin{equation}
 \left\|
 \delta_{\alpha,L}\mathbb Z_{{\rm tail},m}(\boldsymbol z)
 \right\|_{\rm cb}
 \le Ce^{-q/s}.                                               \label{eq:newV56-tail-source}
\end{equation}
\end{lemma}

\begin{proof}
The left side is at most the unnormalized joint expectation
\[
 \mathbb E_{\widehat\nu_\alpha}
 \left[
 \mathbf1_{E_{\rm tail}}
 e^{CY_sX(\boldsymbol z)}
 \right].
\]
Choose \(\eta_L\) so that \(2C\eta_L\le\lambda_0\) in the V44
Wilson radial exponential moment.  Cauchy--Schwarz gives
\[
 \mathbb P(E_{\rm tail})^{1/2}
 \left(\mathbb Ee^{2C\eta_LY_s}\right)^{1/2}
 \le Ce^{-q_L/(2s)}.
\]
Rename \(q=q_L/2\).
\end{proof}

\begin{theorem}[Noncommutative rare-log correction on the baseline
simplex]
\label{thm:newV56-rare-log}
On a possibly smaller simplex
\eqref{eq:newV56-simplex}, define
\begin{equation}
 \Delta\mathbb L_m(\boldsymbol z)
 :=
 \log\mathbb Z_{W,m}(\boldsymbol z)
 -\log\mathbb Z_{{\rm loc},m}(\boldsymbol z)
 -\log(1-\delta_{\alpha,L})I.                                \label{eq:newV56-DeltaL}
\end{equation}
Then the principal logarithms are well defined and
\begin{equation}
 \sup_{X(\boldsymbol z)\le\eta_L}
 \|\Delta\mathbb L_m(\boldsymbol z)\|_{\rm cb}
 \le Ce^{-q/s},                                               \label{eq:newV56-DeltaL-bound}
\end{equation}
uniformly in \(m\), the thermal representations, and spectator
amplifications.
\end{theorem}

\begin{proof}
Put
\[
 A(\boldsymbol z)
 =(1-\delta_{\alpha,L})\mathbb Z_{{\rm loc},m}(\boldsymbol z),
 \qquad
 T(\boldsymbol z)
 =\delta_{\alpha,L}\mathbb Z_{{\rm tail},m}(\boldsymbol z).
\]
By \Cref{thm:newV56-local-invertible},
\(\|A-I\|\le1/4\) for sufficiently large \(\alpha\), after constants
are reduced.  By \Cref{lem:newV56-tail-source},
\(\|T\|\le Ce^{-q/s}\).  Thus
\(A+uT\) remains in \(\|\,\cdot-I\|<1/2\) for \(0\le u\le1\).
The V44 Fr\'echet logarithm estimate gives
\[
 \|\log(A+T)-\log A\|
 \le\int_0^1\|D\log_{A+uT}[T]\|\,du
 \le C\|T\|.
\]
Because the positive scalar \(1-\delta_{\alpha,L}\) commutes with
\(\mathbb Z_{\rm loc}\) and both spectra lie in the principal
logarithm domain,
\[
 \log A
 =\log(1-\delta_{\alpha,L})I+\log\mathbb Z_{\rm loc}.
\]
This proves \eqref{eq:newV56-DeltaL-bound}.  The argument is
Banach-algebraic and unchanged by amplification; no false
\(\log(AB)=\log A+\log B\) step for noncommuting factors is used.
\end{proof}

\begin{corollary}[Weak local-zero/tail-log card]
\label{cor:newV56-weak-LZTC}
The local source is uniformly zero-free and the exact
noncommutative tail correction is exponentially small on the V44
baseline weighted simplex.  Its squarefree coefficient obeys
\begin{equation}
 \left\|[z_1\cdots z_m]\Delta\mathbb L_m\right\|_{\rm cb}
 \le
 Ce^{-q/s}(Km)^m\prod_{i=1}^m b_i.                            \label{eq:newV56-tail-crude}
\end{equation}
\end{corollary}

\begin{proof}
Use the polydisc
\(|z_i|\le\eta_L/(mb_i)\) inside the simplex and apply the
multivariable Cauchy estimate to
\eqref{eq:newV56-DeltaL-bound}.
\end{proof}

\section{Why the baseline zero-free result is still insufficient}
\label{sec:newV56-deficit}

\begin{proposition}[The tail Cauchy bound misses Gaussian
homogeneity]
\label{prop:newV56-tail-deficit}
Estimate \eqref{eq:newV56-tail-crude} does not imply
\[
 K^m\left(\prod_i b_i\right)B^{m-2}
\]
uniformly in \(m,s\).  For equal minimal thermal marks
\(b_i=\sqrt s\), the ratio of the scale on the right of
\eqref{eq:newV56-tail-crude} to the strong target contains
\begin{equation}
 e^{-q/s}C^m m^2s^{\,1-m/2},                                 \label{eq:newV56-deficit-ratio}
\end{equation}
which is unbounded as \(m\to\infty\) for every fixed sufficiently
small \(s\).
\end{proposition}

\begin{proof}
At equal marks,
\[
 \prod_i b_i=s^{m/2},\qquad
 B^{m-2}=(m\sqrt s)^{m-2}.
\]
Divide the crude scale
\(e^{-q/s}(Cm)^ms^{m/2}\) by
\(K^m m^{m-2}s^{m-1}\).  This gives
\eqref{eq:newV56-deficit-ratio} after changing the exponential
base.  For fixed \(s<1\), its factor
\((C/\sqrt s)^m\) eventually dominates the fixed
nonperturbative prefactor.
\end{proof}

\begin{correction}[The strong LZTC row is not mere invertibility]
\label{corr:newV56-LZTC}
V55's LZTC(i) requires a source domain adapted to the full strong
homogeneity, not merely the baseline V44 invertibility simplex.
\Cref{thm:newV56-local-invertible,thm:newV56-rare-log} close the
zero-free and exact-log operations on that baseline domain, but
\Cref{prop:newV56-tail-deficit} proves that this does not yet close
LZTC, PLRC, or WSDC.  The remaining loss is the same
Gaussian-renormalized \(B^{m-2}\) deficit isolated in V44 and V50.
\end{correction}

\begin{definition}[Gaussian-renormalized local/tail card, GRLTC]
\label{def:newV56-GRLTC}
GRLTC is a bound on
\[
 \log\mathbb Z_{W,m}
 -\bigl[\log\mathbb Z_{{\rm loc},m}\bigr]_{\le2}
\]
in which the local time-ordered forest and the exact tail logarithm
\eqref{eq:newV56-DeltaL} are kept assembled, and squarefree
coefficient extraction yields
\[
 K^m\left(\prod_i b_i\right)B^{m-2}.
\]
No term may be estimated only by
\eqref{eq:newV56-tail-crude}.
\end{definition}

\begin{programme}[V57: merge the positive split with the V50
grade/leg forest]
\label{prog:newV56-V57}
\begin{enumerate}[leftmargin=3em]
\item Apply the V50 rooted-Wick leg inequality to the local
      conditional source, using the exact bound \(S<Ls\) instead of
      fixed-order image separation.
\item Keep the rare-log correction
      \(\Delta\mathbb L_m\) inside the squarefree connected
      coefficient rather than applying the baseline Cauchy estimate.
\item Derive its first two Taylor degrees exactly and subtract them
      before the tree sum.
\item Test whether every tail-connected squarefree monomial either has
      V50 excess degree \(\rho+2q\ge m-2\) or receives the larger
      logarithmic source radius of V54.
\item If neither mechanism applies, construct an explicit physical
      high-arity source counterexample before attempting another
      sufficient card.
\end{enumerate}
\end{programme}

\begin{firewall}[Claim boundary after V56]
\label{fire:newV56-boundary}
V56 reuses rather than duplicates V44 to prove local conditional
source invertibility, computes the exact singleton and
conjugate-singlet source zeros, proves an exponentially small
tail-weighted source bound, and proves a noncommutative
rare-log correction estimate on the baseline weighted simplex.  It also
proves that baseline Cauchy extraction still misses the
\(B^{m-2}\) homogeneity and formulates GRLTC.

V56 does not prove GRLTC, LZTC, PLRC, WSDC, BRSC, HWSFC, UGCGC,
thermal connected WPRAC, all-arity RGSC, the raw Wilson aggregate
strong card, all-order strong MTA, WUFC, CPM, cutoff removal,
reflection positivity, Osterwalder--Schrader reconstruction, a
continuum Yang--Mills measure, or a positive Hamiltonian gap.  The full
Yang--Mills mass gap has not been proved.
\end{firewall}

\begin{theorem}[V56 result ledger]
\label{thm:newV56-results}
V56:
\begin{enumerate}[label=\textup{(V56.\arabic*)},leftmargin=6.5em]
\item proves the local conditional radial exponential moment;
\item imports the V44 mechanism to the exact positive local law;
\item computes the local singleton and pair-singlet zeros;
\item verifies their correct \(b^{-2}\) and \(b^{-1}\) scales;
\item proves the tail-weighted source is \(O(e^{-q/s})\);
\item proves the exact noncommutative rare-log correction bound;
\item derives its baseline squarefree Cauchy estimate;
\item proves that estimate still misses strong homogeneity; and
\item identifies GRLTC as the merged replacement gate.
\end{enumerate}
\end{theorem}

\begin{proof}
Use
\Cref{lem:newV56-local-radial,
thm:newV56-local-invertible,
lem:newV56-aR,
prop:newV56-singleton-zero,
thm:newV56-pair-polynomial,
lem:newV56-tail-source,
thm:newV56-rare-log,
cor:newV56-weak-LZTC,
prop:newV56-tail-deficit,
def:newV56-GRLTC}.
\end{proof}

\paragraph{Parent-version record.}
V56 was formed from
\texttt{Renewal\_Yang\_Mills\_Mass\_Gap\_Research\_Notes\_V55.tex}.
The V55 parent had \(26\,587\) logical source lines,
\(844\,056\) bytes, and SHA--256
\[
 \texttt{9f990713e4c4e256d150a839514ffc0b4149cf4189ed0d63c6688dbdd733ecb7}.
\]
The next compulsory companion audit is V60.

% =============================================================================
% END APPENDED FIFTH-SERIES V56 MATERIAL
% =============================================================================

\clearpage
\chapter{Fifth-series V57: the subordinated-clock cumulant hierarchy}
\label{ch:newV57-clock}

\section{Fixed-order clock asymptotics}
\label{sec:newV57-asymptotics}

\begin{theorem}[Uniform fixed-order Debye logarithm, standard input]
\label{thm:newV57-Debye}
Let \(\mu=4\nu^2\).  Uniformly for \(\nu\) in a fixed compact
neighbourhood of \(1\), as \(\alpha\to\infty\),
\begin{align}
 \log I_\nu(\alpha)
 &=
 \alpha-\frac12\log(2\pi\alpha)
 -{\mu-1\over8\alpha}
 -{\mu-1\over16\alpha^2}                                    \notag\\
 &\quad
 +{(\mu-1)(\mu-25)\over384\alpha^3}
 +O(\alpha^{-4}),                                             \label{eq:newV57-Debye}
\end{align}
and the remainder may be differentiated a fixed number of times in
\(\nu\), uniformly on a smaller compact neighbourhood.
\end{theorem}

\begin{proof}
This is the classical fixed-order Debye expansion with its standard
uniform analytic remainder.  To record the logarithmic coefficients,
start from
\[
 I_\nu(\alpha)
 ={e^\alpha\over\sqrt{2\pi\alpha}}
 \left[
 1-{\mu-1\over8\alpha}
 +{(\mu-1)(\mu-9)\over128\alpha^2}
 -{(\mu-1)(\mu-9)(\mu-25)\over3072\alpha^3}
 +O(\alpha^{-4})
 \right].
\]
Expanding \(\log(1+x)\) gives the first two displayed logarithmic
coefficients.  If \(x=\mu-1\), the cubic numerator is
\[
 -x(x-8)(x-24)+3x^2(x-8)-2x^3
 =8x(x-24),
\]
which gives the coefficient
\((\mu-1)(\mu-25)/384\).  Uniform differentiated remainders follow
from the standard contour proof of the Debye expansion.  This
fixed-order special-function estimate is an explicit external analytic
input.
\end{proof}

\begin{corollary}[Laplace logarithm through cubic order]
\label{cor:newV57-logL}
Uniformly for \(\lambda\) in a fixed neighbourhood of \(0\),
\begin{align}
 \log L_\alpha(\lambda)
 &=
 -{\lambda\over2\alpha}
 -{\lambda\over4\alpha^2}
 +{1\over\alpha^3}
   \left(-{3\lambda\over16}+{\lambda^2\over24}\right)
 +O(\alpha^{-4}).                                             \label{eq:newV57-logL}
\end{align}
\end{corollary}

\begin{proof}
In \eqref{eq:newV57-Debye}, put
\(\nu=\sqrt{1+\lambda}\), so
\(\mu=4+4\lambda\), and subtract the value at \(\nu=1\).
The linear differences in the first two correction terms are
\(-\lambda/(2\alpha)\) and
\(-\lambda/(4\alpha^2)\).  Finally,
\[
 (\mu-1)(\mu-25)
 =(3+4\lambda)(-21+4\lambda)
 =-63-72\lambda+16\lambda^2.
\]
Subtract the constant term and divide by \(384\alpha^3\).
\end{proof}

\begin{theorem}[Thermal mean and enhanced clock variance]
\label{thm:newV57-mean-var}
For the full tilted time \(S\sim\rho_\alpha\),
\begin{align}
 \mathbb ES
 &={1\over2\alpha}+{1\over4\alpha^2}
   +{3\over16\alpha^3}+O(\alpha^{-4}),                        \label{eq:newV57-mean}\\
 \operatorname{Var}(S)
 &={1\over12\alpha^3}+O(\alpha^{-4}).                         \label{eq:newV57-var}
\end{align}
In particular,
\begin{equation}
 \mathbb ES=O(s),\qquad
 \operatorname{Var}(S)=O(s^3).                               \label{eq:newV57-thermal-moments}
\end{equation}
\end{theorem}

\begin{proof}
For a Laplace transform,
\[
 -(\log L_\alpha)'(0)=\mathbb ES,\qquad
 (\log L_\alpha)''(0)=\operatorname{Var}(S).
\]
Differentiate \eqref{eq:newV57-logL}.  Since
\(\alpha^{-1}\asymp s\), the final bounds follow.
\end{proof}

\section{Transfer to the positive local clock}
\label{sec:newV57-local}

\begin{lemma}[Fixed moments lost to the tail are nonperturbative]
\label{lem:newV57-tail-moments}
For each fixed integer \(k\ge0\), there are \(C_k,q_k>0\) such
that
\begin{equation}
 \mathbb E\!\left[S^k\mathbf1_{\{S\ge Ls\}}\right]
 \le C_ks^ke^{-q_k/s}.                                      \label{eq:newV57-tail-moments}
\end{equation}
\end{lemma}

\begin{proof}
By Cauchy--Schwarz,
\[
 \mathbb E[S^k\mathbf1_{\{S\ge Ls\}}]
 \le
 \mathbb P(S\ge Ls)^{1/2}
 (\mathbb ES^{2k})^{1/2}.
\]
The probability is \(O(e^{-q_L/s})\) by
\Cref{thm:newV55-positive-tail}.  For fixed \(2k\),
\Cref{thm:newV53-two-sector} gives
\(\mathbb ES^{2k}\le C_ks^{2k}\), after its
nonperturbative term is absorbed into a fixed power of \(s\).
\end{proof}

\begin{corollary}[Local conditional mean and variance]
\label{cor:newV57-local-mean-var}
For \(S_{\rm loc}\sim\rho_{\alpha,L}^{\rm loc}\),
\begin{align}
 \mathbb ES_{\rm loc}
 &={1\over2\alpha}+{1\over4\alpha^2}
   +{3\over16\alpha^3}+O(\alpha^{-4}),                        \label{eq:newV57-local-mean}\\
 \operatorname{Var}(S_{\rm loc})
 &={1\over12\alpha^3}+O(\alpha^{-4}),                         \label{eq:newV57-local-var}
\end{align}
where an additional \(O(e^{-q/s})\) term has been absorbed in each
remainder.  In particular,
\(\operatorname{Var}(S_{\rm loc})\le Cs^3\).
\end{corollary}

\begin{proof}
The first two conditional raw moments differ from the full raw moments
by \eqref{eq:newV57-tail-moments}, division by
\(1-\delta_{\alpha,L}=1+O(e^{-q/s})\), and an exponentially small
normalization correction.  Apply
\Cref{thm:newV57-mean-var} and form the variance.
\end{proof}

\section{The tangent Gaussian mixture}
\label{sec:newV57-Gaussian}

\begin{definition}[Abstract subordinated tangent Gaussian]
\label{def:newV57-Gaussian}
Let \(Q\) be a positive semidefinite bilinear form on a finite real
vector space.  Conditional on \(S_{\rm loc}\), let \(X\) be centered
Gaussian with covariance \(S_{\rm loc}Q\).  Its source logarithm is
\begin{equation}
 \log\mathbb E e^{\ell(X)}
 =
 K_{\rm loc}\!\left({Q(\ell,\ell)\over2}\right),
 \qquad
 K_{\rm loc}(z):=\log\mathbb E e^{zS_{\rm loc}}.              \label{eq:newV57-Gaussian-log}
\end{equation}
\end{definition}

\begin{lemma}[Exact even Gaussian-mixture cumulants]
\label{lem:newV57-Gaussian-cumulants}
Let \(\varkappa_r(S_{\rm loc})=K_{\rm loc}^{(r)}(0)\).
All odd joint cumulants of \(X\) vanish, and
\begin{equation}
 \kappa_{2r}\bigl(\ell_1(X),\ldots,\ell_{2r}(X)\bigr)
 =
 \varkappa_r(S_{\rm loc})
 \sum_{\mathfrak p\in\mathfrak P_2(2r)}
 \prod_{\{i,j\}\in\mathfrak p}Q(\ell_i,\ell_j),               \label{eq:newV57-Gaussian-cumulants}
\end{equation}
where \(\mathfrak P_2(2r)\) is the set of pairings.
\end{lemma}

\begin{proof}
For sources \(t_i\),
\[
 \log\mathbb E
 \exp\left(\sum_i t_i\ell_i(X)\right)
 =
 K_{\rm loc}\left(
 {1\over2}\sum_{i,j}t_it_jQ(\ell_i,\ell_j)
 \right).
\]
Expand \(K_{\rm loc}\) at zero and differentiate.  Only total even
degrees occur.  At degree \(2r\), differentiating the \(r\)-th power
of the quadratic form gives exactly the sum over pairings, with the
factor \(r!\) canceled by the Taylor denominator of
\(K_{\rm loc}\).
\end{proof}

\begin{theorem}[The first mixing correction has strong quartic
homogeneity]
\label{thm:newV57-quartic}
Put
\[
 d_i:=Q(\ell_i,\ell_i)^{1/2},\qquad b_i:=\sqrt s\,d_i,
\]
and assume \(b_i\ge\sqrt s\).  Then
\begin{equation}
 \left|
 \kappa_4(\ell_1(X),\ldots,\ell_4(X))
 \right|
 \le
 C\left(\prod_{i=1}^4b_i\right)
 \left(\sum_{i=1}^4b_i\right)^2.                             \label{eq:newV57-quartic}
\end{equation}
\end{theorem}

\begin{proof}
Set \(r=2\) in
\eqref{eq:newV57-Gaussian-cumulants}.  By
Cauchy--Schwarz for \(Q\) and
\Cref{cor:newV57-local-mean-var},
\[
 |\kappa_4|
 \le3Cs^3\prod_i d_i
 =3Cs\prod_i b_i.
\]
If \(B=\sum_i b_i\), then \(B\ge4\sqrt s\), so
\(s\le B^2/16\).  This gives
\eqref{eq:newV57-quartic}.
\end{proof}

\begin{assessment}[Why the V49 variance-mixture failure is avoided]
\label{ass:newV57-V49}
An arbitrary mixture of heat times separated by \(O(s)\) has variance
of order \(s^2\), producing a quartic tangent cumulant one thermal power
too large.  The exact Hartman--Watson clock instead has variance
of order \(s^3\), and
\Cref{thm:newV57-quartic} recovers that power.  This does not close
higher arity, where cumulants of the clock rather than raw moments are
required.
\end{assessment}

\section{The all-order clock gate}
\label{sec:newV57-CCH}

\begin{definition}[Clock cumulant hierarchy, CCH]
\label{def:newV57-CCH}
CCH is the existence of an absolute \(K_0\) such that, for every
\(r\ge2\),
\begin{equation}
 |\varkappa_r(S_{\rm loc})|
 \le r!K_0^r s^{\,2r-1},                                    \label{eq:newV57-CCH}
\end{equation}
uniformly for sufficiently large \(\alpha\).  The first cumulant
obeys the separate bound \(\mathbb ES_{\rm loc}\le K_0s\).
\end{definition}

\begin{theorem}[CCH closes the tangent Gaussian all-arity card]
\label{thm:newV57-CCH-suffices}
Assume CCH.  Then, for every \(r\ge2\),
\begin{equation}
 \left|
 \kappa_{2r}(\ell_1(X),\ldots,\ell_{2r}(X))
 \right|
 \le
 K^{2r}\left(\prod_{i=1}^{2r}b_i\right)
 B^{2r-2},\qquad B=\sum_i b_i.                               \label{eq:newV57-Gaussian-strong}
\end{equation}
\end{theorem}

\begin{proof}
There are \((2r-1)!!\) pairings.  Hence
\begin{align*}
 |\kappa_{2r}|
 &\le
 r!K_0^rs^{2r-1}(2r-1)!!\prod_i d_i\\
 &=
 r!K_0^r(2r-1)!!s^{r-1}\prod_i b_i.
\end{align*}
Since every \(b_i\ge\sqrt s\),
\[
 B^{2r-2}\ge(2r)^{2r-2}s^{r-1}.
\]
Finally,
\[
 r!(2r-1)!!={(2r)!\over2^r}
 \le K_1^{2r}(2r)^{2r-2}
\]
after increasing a numerical base \(K_1\); the harmless factor
\((2r)^2/2^r\) is exponentially bounded.  Combine the estimates.
\end{proof}

\begin{proposition}[Analytic meaning of CCH]
\label{prop:newV57-CCH-analytic}
CCH is equivalent at the coefficient level to a thermal
zero-free/analytic scale \(O(s^{-2})\) for the centered local-clock
logarithm:
\begin{equation}
 K_{\rm loc}(z)-z\mathbb ES_{\rm loc}
 =
 \sum_{r\ge2}{\varkappa_r(S_{\rm loc})\over r!}z^r
 \label{eq:newV57-clock-series}
\end{equation}
is absolutely bounded by
\[
 {K_0^2s^3|z|^2\over1-K_0s^2|z|}
\]
whenever \(K_0s^2|z|<1\).
\end{proposition}

\begin{proof}
Insert \eqref{eq:newV57-CCH} in the series and sum the geometric
majorant:
\[
 \sum_{r\ge2}K_0^rs^{2r-1}|z|^r
 =
 {K_0^2s^3|z|^2\over1-K_0s^2|z|}.
\]
\end{proof}

\begin{warning}[Bounded support alone gives only the wrong scale]
\label{warn:newV57-support}
The truncation \(0<S_{\rm loc}<Ls\) alone gives clock cumulants at
best on an \(O(s^{-1})\) zero-free scale.  CCH requires the stronger
\(O(s^{-2})\) scale, reflecting concentration width \(s^{3/2}\) plus
the \(e^{-q/s}\) cost of \(O(s)\) deviations.  Thus
\Cref{thm:newV55-local-Duhamel} cannot prove CCH by itself.
\end{warning}

\begin{definition}[Clock moderate-deviation card, CMDC]
\label{def:newV57-CMDC}
CMDC is a centered exponential estimate for \(S_{\rm loc}\) of the
form
\begin{equation}
 \log\mathbb E
 e^{z(S_{\rm loc}-\mathbb ES_{\rm loc})}
 \le C s^3z^2
 \qquad
 (z\in\mathbb R,\ |z|\le c/s^2),                             \label{eq:newV57-CMDC}
\end{equation}
together with a complex zero-free extension to
\(|z|\le c'/s^2\).
\end{definition}

\begin{proposition}[CMDC implies CCH]
\label{prop:newV57-CMDC-CCH}
If CMDC holds with a uniform complex bound on a disk
\(|z|\le c'/s^2\), then CCH holds.
\end{proposition}

\begin{proof}
Apply Cauchy's estimate to the centered logarithm on a circle of radius
\(\gamma/s^2\), with \(0<\gamma<c'\).  Its supremum is \(O(1/s)\)
by \eqref{eq:newV57-CMDC}.  Therefore
\[
 {|\varkappa_r(S_{\rm loc})|\over r!}
 \le {C/s\over(\gamma/s^2)^r}
 =C\gamma^{-r}s^{2r-1},
\]
which is CCH after changing the exponential base.
\end{proof}

\begin{programme}[V58: Riccati and moderate-deviation attack]
\label{prog:newV57-V58}
\begin{enumerate}[leftmargin=3em]
\item Use the modified-Bessel differential equation to derive an exact
      Riccati equation for
      \(\log L_\alpha(\lambda)\).
\item Rescale \(\lambda\) on the \(s^{-2}\) clock-source scale and
      identify the first singular coefficient.
\item Prove CMDC for the positive truncated law, treating the removed
      event with the V55 \(e^{-q/s}\) probability inside the logarithm.
\item If the full \(s^{-2}\) disk meets a Bessel-order obstruction,
      locate its exact zero and compare it with the truncation-induced
      zero displacement.
\item Once CCH is established, insert
      \Cref{thm:newV57-CCH-suffices} into the V50 excess-degree forest
      to attack GRLTC.
\end{enumerate}
\end{programme}

\begin{firewall}[Claim boundary after V57]
\label{fire:newV57-boundary}
V57 proves the fixed-order tilted-clock mean and variance asymptotics,
transfers them to the positive local clock, proves exact tangent
Gaussian-mixture cumulant formulae, and closes the strong quartic
mixing card.  It formulates CCH, proves that CCH closes the tangent
Gaussian card at every arity, and reduces CCH to CMDC.

V57 does not prove CMDC, CCH, GRLTC, LZTC, PLRC, WSDC, BRSC, HWSFC,
UGCGC, thermal connected WPRAC, all-arity RGSC, the raw Wilson
aggregate strong card, all-order strong MTA, WUFC, CPM, cutoff
removal, reflection positivity, Osterwalder--Schrader reconstruction,
a continuum Yang--Mills measure, or a positive Hamiltonian gap.  The
full Yang--Mills mass gap has not been proved.
\end{firewall}

\begin{theorem}[V57 result ledger]
\label{thm:newV57-results}
V57:
\begin{enumerate}[label=\textup{(V57.\arabic*)},leftmargin=6.5em]
\item records the uniform fixed-order Debye logarithm;
\item proves the clock mean through cubic order;
\item proves the enhanced variance
      \(\operatorname{Var}(S)=1/(12\alpha^3)+O(\alpha^{-4})\);
\item transfers these results to the positive local clock;
\item proves the exact tangent Gaussian-mixture cumulant formula;
\item closes its strong quartic card;
\item formulates CCH;
\item proves CCH sufficient for the tangent Gaussian all-arity card;
      and
\item reduces CCH to the clock moderate-deviation card CMDC.
\end{enumerate}
\end{theorem}

\begin{proof}
Use
\Cref{thm:newV57-Debye,
cor:newV57-logL,
thm:newV57-mean-var,
lem:newV57-tail-moments,
cor:newV57-local-mean-var,
lem:newV57-Gaussian-cumulants,
thm:newV57-quartic,
def:newV57-CCH,
thm:newV57-CCH-suffices,
prop:newV57-CMDC-CCH}.
\end{proof}

\paragraph{Parent-version record.}
V57 was formed from
\texttt{Renewal\_Yang\_Mills\_Mass\_Gap\_Research\_Notes\_V56.tex}.
The V56 parent had \(27\,072\) logical source lines,
\(860\,058\) bytes, and SHA--256
\[
 \texttt{4b480ba75edec24b06d3e0ee70354bb003b6c7b4e70cf94546b8e1c03cd86ab2}.
\]
The next compulsory companion audit is V60.

% =============================================================================
% END APPENDED FIFTH-SERIES V57 MATERIAL
% =============================================================================

\clearpage
\chapter{Fifth-series V58: exact Riccati recursion and the wrapped
homogeneous mode}
\label{ch:newV58-Riccati}

\section{The Bessel equation for the clock logarithm}
\label{sec:newV58-equation}

\begin{definition}[Clock Laplace logarithm]
\label{def:newV58-G}
For \(\lambda\) near \(0\), put
\begin{equation}
 G(\lambda,\alpha)
 :=\log L_\alpha(\lambda)
 =\log I_{\sqrt{1+\lambda}}(\alpha)-\log I_1(\alpha),
 \label{eq:newV58-G}
\end{equation}
with the branch satisfying \(\sqrt{1}=1\).  Also put
\begin{equation}
 R_1(\alpha):={I_1'(\alpha)\over I_1(\alpha)}.
 \label{eq:newV58-R1}
\end{equation}
\end{definition}

\begin{theorem}[Exact clock Riccati equation]
\label{thm:newV58-Riccati}
The logarithm \(G\) satisfies
\begin{equation}
 G_{\alpha\alpha}
 +\left(2R_1(\alpha)+{1\over\alpha}\right)G_\alpha
 +(G_\alpha)^2
 ={\lambda\over\alpha^2}.                                    \label{eq:newV58-Riccati}
\end{equation}
\end{theorem}

\begin{proof}
For \(\nu=\sqrt{1+\lambda}\), the modified-Bessel equation is
\[
 I_\nu''+{1\over\alpha}I_\nu'
 -\left(1+{\nu^2\over\alpha^2}\right)I_\nu=0.
\]
If \(R_\nu=I_\nu'/I_\nu\), division by \(I_\nu\) gives
\[
 R_\nu'+R_\nu^2+{1\over\alpha}R_\nu
 =1+{\nu^2\over\alpha^2}.
\]
Subtract the equation at \(\nu=1\).  Since
\[
 G_\alpha=R_\nu-R_1,\qquad
 R_\nu^2-R_1^2=2R_1G_\alpha+(G_\alpha)^2,
\]
and \(\nu^2-1=\lambda\), the result follows.
\end{proof}

\section{Exact coefficient hierarchy}
\label{sec:newV58-hierarchy}

\begin{definition}[Laplace-log coefficients]
\label{def:newV58-gr}
On \(|\lambda|<1\), write
\begin{equation}
 G(\lambda,\alpha)
 =\sum_{r=1}^\infty g_r(\alpha)\lambda^r.
 \label{eq:newV58-gr}
\end{equation}
If \(\varkappa_r(S)\) is the \(r\)-th cumulant of the full clock, then
\begin{equation}
 g_r(\alpha)={(-1)^r\over r!}\varkappa_r(S).
 \label{eq:newV58-gr-cumulant}
\end{equation}
\end{definition}

\begin{theorem}[Triangular Riccati coefficient system]
\label{thm:newV58-coefficient-system}
Let
\[
 A(\alpha):=2R_1(\alpha)+{1\over\alpha}.
\]
Then
\begin{align}
 g_1''+Ag_1'&={1\over\alpha^2},                               \label{eq:newV58-g1}\\
 g_r''+Ag_r'
 &=-\sum_{p=1}^{r-1}g_p'g_{r-p}',
 \qquad r\ge2.                                                \label{eq:newV58-gr-recursion}
\end{align}
The integrating factor is
\begin{equation}
 M(\alpha):=\alpha I_1(\alpha)^2,\qquad {M'\over M}=A,
 \label{eq:newV58-M}
\end{equation}
so equivalently
\begin{align}
 (Mg_1')'&={I_1(\alpha)^2\over\alpha},                         \label{eq:newV58-g1-div}\\
 (Mg_r')'
 &=-M\sum_{p=1}^{r-1}g_p'g_{r-p}',
 \qquad r\ge2.                                                \label{eq:newV58-gr-div}
\end{align}
\end{theorem}

\begin{proof}
Insert \eqref{eq:newV58-gr} in
\eqref{eq:newV58-Riccati} and compare coefficients of
\(\lambda^r\).  The square contributes the convolution in
\eqref{eq:newV58-gr-recursion}.  Finally
\[
 {d\over d\alpha}\log(\alpha I_1^2)
 ={1\over\alpha}+2{I_1'\over I_1}=A.
\]
Multiplication by \(M\) gives the divergence form.
\end{proof}

\begin{corollary}[The first two rows reproduce V57]
\label{cor:newV58-first-two}
The asymptotic solutions of
\eqref{eq:newV58-g1}--\eqref{eq:newV58-gr-recursion} begin as
\begin{align}
 g_1(\alpha)
 &=-{1\over2\alpha}-{1\over4\alpha^2}
   -{3\over16\alpha^3}+O(\alpha^{-4}),                        \label{eq:newV58-g1-asym}\\
 g_2(\alpha)
 &={1\over24\alpha^3}+O(\alpha^{-4}),                         \label{eq:newV58-g2-asym}
\end{align}
and hence give \eqref{eq:newV57-mean}--\eqref{eq:newV57-var}.
\end{corollary}

\begin{proof}
These are the coefficients of \(\lambda\) and \(\lambda^2\) in
\eqref{eq:newV57-logL}.  Direct substitution in
\eqref{eq:newV58-g1}--\eqref{eq:newV58-gr-recursion}, together with
\(R_1=1-(2\alpha)^{-1}+O(\alpha^{-2})\), verifies the displayed
orders.
\end{proof}

\section{The exact homogeneous solution}
\label{sec:newV58-homogeneous}

\begin{lemma}[Wronskian solution]
\label{lem:newV58-H}
The decaying nonconstant solution of
\[
 h''+Ah'=0
\]
is, up to a scalar multiple,
\begin{equation}
 H(\alpha):={K_1(\alpha)\over I_1(\alpha)}
 =\int_\alpha^\infty {du\over uI_1(u)^2}.                     \label{eq:newV58-H}
\end{equation}
In particular,
\begin{equation}
 H'(\alpha)=-{1\over\alpha I_1(\alpha)^2}=-{1\over M(\alpha)}.
 \label{eq:newV58-H-prime}
\end{equation}
\end{lemma}

\begin{proof}
The modified-Bessel Wronskian is
\[
 I_1K_1'-I_1'K_1=-{1\over\alpha}.
\]
Divide by \(I_1^2\) to obtain
\eqref{eq:newV58-H-prime}.  Since \(M'/M=A\),
\((MH')'=0\), which is the homogeneous equation.  Both the ratio
and the integral tend to zero at infinity, so integration proves
\eqref{eq:newV58-H}.
\end{proof}

\begin{theorem}[Homogeneous ambiguity equals the wrapped scale]
\label{thm:newV58-homogeneous-scale}
For \(\alpha\ge1\), there are absolute \(0<c<C<\infty\) such that
\begin{equation}
 ce^{-2\alpha}\le H(\alpha)\le Ce^{-2\alpha}.                 \label{eq:newV58-H-scale}
\end{equation}
Moreover \(H(\alpha)\) is comparable to the V52 branch coefficient
\(\delta_\alpha=K_0(\alpha)/I_1(\alpha)\).
\end{theorem}

\begin{proof}
The standard elementary large-argument bounds for \(I_1\) and \(K_1\)
give
\[
 I_1(\alpha)\asymp{\!e^\alpha\over\sqrt\alpha},
 \qquad
 K_1(\alpha)\asymp{e^{-\alpha}\over\sqrt\alpha}
\]
uniformly for \(\alpha\ge1\), after constants cover the compact
interval near \(1\).  This proves \eqref{eq:newV58-H-scale}.
The ratio \(K_1/K_0\) is positive and continuous on \([1,\infty)\)
and tends to \(1\); hence it is bounded above and below there.
Dividing by \(I_1\) proves comparability with \(\delta_\alpha\).
\end{proof}

\begin{corollary}[Every Riccati row has a wrapped ambiguity]
\label{cor:newV58-row-ambiguity}
If two solutions of the \(r\)-th coefficient equation have the same
lower-order forcing, their difference is
\begin{equation}
 c_{r,0}+c_{r,1}H(\alpha).                                   \label{eq:newV58-row-ambiguity}
\end{equation}
Requiring decay as \(\alpha\to\infty\) removes \(c_{r,0}\) but does
not determine \(c_{r,1}\).
\end{corollary}

\begin{proof}
Subtract the two inhomogeneous equations.  The difference solves the
homogeneous equation, whose two independent solutions are \(1\) and
\(H\) by \Cref{lem:newV58-H}.  Both \(H\) and every inverse power of
\(\alpha\) decay at infinity, so decay alone does not fix the
coefficient of \(H\).
\end{proof}

\begin{assessment}[Riccati form of the beyond-all-orders sector]
\label{ass:newV58-Stokes}
The inverse-power Debye recursion determines the algebraic part of each
clock cumulant.  At every coefficient order it leaves an undetermined
multiple of \(K_1/I_1\asymp e^{-2\alpha}\).  This is the exact
Riccati manifestation of the V40/V52 wrapped sector: it is invisible
to every finite inverse-power expansion but controls sufficiently high
coefficient order.
\end{assessment}

\section{The full-clock high-order obstruction in logarithmic form}
\label{sec:newV58-full-no-go}

\begin{theorem}[Square-root branch of the clock logarithm]
\label{thm:newV58-log-branch}
As \(\lambda\to-1\) on the principal sheet,
\begin{equation}
 G(\lambda,\alpha)
 =
 \log{I_0(\alpha)\over I_1(\alpha)}
 -{K_0(\alpha)\over I_0(\alpha)}\sqrt{1+\lambda}
 +O(1+\lambda).                                               \label{eq:newV58-log-branch}
\end{equation}
Consequently, for fixed \(\alpha\),
\begin{equation}
 |g_r(\alpha)|
 =
 {K_0(\alpha)\over2\sqrt\pi\,I_0(\alpha)}
 r^{-3/2}\left(1+O_\alpha(r^{-1})\right)                     \label{eq:newV58-gr-tail}
\end{equation}
along the absolute-value asymptotic as \(r\to\infty\).
\end{theorem}

\begin{proof}
At \(\nu=0\),
\[
 I_\nu(\alpha)=I_0(\alpha)-K_0(\alpha)\nu+O(\nu^2)
\]
by \Cref{lem:newV51-order-derivative}.  Put
\(\nu=\sqrt{1+\lambda}\), divide by \(I_1\), and expand the
logarithm about the positive constant \(I_0/I_1\).  This proves
\eqref{eq:newV58-log-branch}.  Darboux coefficient extraction at
the nearest branch point \(\lambda=-1\) is legitimate because the
complete-Bernstein representation from V51 continues
\(L_\alpha=e^{-\Phi_\alpha}\) as a zero-free function on the slit
plane \(\mathbb C\setminus(-\infty,-1]\).  Thus no logarithmic zero
occurs inside \(|\lambda|<1\).  Using
\[
 |[\lambda^r](1+\lambda)^{1/2}|
 ={1\over2\sqrt\pi}r^{-3/2}(1+O(r^{-1})),
\]
gives \eqref{eq:newV58-gr-tail}.
\end{proof}

\begin{corollary}[Full-clock CCH is false]
\label{cor:newV58-full-CCH-false}
For every fixed sufficiently large \(\alpha\), no constant \(K\)
independent of \(r\) can satisfy
\[
 |g_r(\alpha)|\le K^r\alpha^{-(2r-1)}
\qquad(r\ge2).
\]
Equivalently, CCH cannot hold for the untruncated full clock.
\end{corollary}

\begin{proof}
The proposed right side has \(r\)-th root \(K\alpha^{-2}\), whereas
\eqref{eq:newV58-gr-tail} has limiting \(r\)-th root \(1\).
\end{proof}

\begin{warning}[Finite Debye order cannot decide CCH]
\label{warn:newV58-finite-Debye}
For every fixed \(r\), the homogeneous term
\(e^{-2\alpha}\) is smaller than the algebraic target
\(\alpha^{-(2r-1)}\) at large \(\alpha\).  For every fixed
\(\alpha\), it eventually dominates as \(r\to\infty\).
Therefore neither any finite Debye expansion nor fixed-arity
verification can establish the uniform clock hierarchy.
\end{warning}

\section{What positive truncation must cancel}
\label{sec:newV58-truncation}

\begin{lemma}[Exact branch cancellation by the positive tail]
\label{lem:newV58-branch-cancellation}
Let
\begin{align}
 M_{\rm loc}(z)
 &:=\mathbb E[e^{zS}\mid S<Ls],                               \label{eq:newV58-Mloc}\\
 T_{\rm un}(z)
 &:=\mathbb E[e^{zS}\mathbf1_{\{S\ge Ls\}}].                  \label{eq:newV58-Tun}
\end{align}
Then
\begin{equation}
 \mathcal M_\alpha(z)
 =(1-\delta_{\alpha,L})M_{\rm loc}(z)+T_{\rm un}(z).          \label{eq:newV58-mgf-split}
\end{equation}
\(M_{\rm loc}\) is entire of exponential type at most \(Ls\).
Hence the complete square-root monodromy of
\(\mathcal M_\alpha\) at \(z=1\) is carried by
\(T_{\rm un}\) and cancels exactly when the tail is subtracted.
\end{lemma}

\begin{proof}
The identity is the partition of the positive expectation into the
two time events.  Since \(0<S<Ls\) under the local conditional law,
\[
 |M_{\rm loc}(z)|\le e^{Ls|z|},
\]
so it is entire.  Rearranging
\eqref{eq:newV58-mgf-split} shows that the difference between the
full transform and the unnormalized tail transform is entire; their
monodromies are therefore identical.
\end{proof}

\begin{definition}[Riccati wrapped-cancellation card, RWCC]
\label{def:newV58-RWCC}
RWCC requires an all-order coefficient analysis of the exact positive
subtraction in \eqref{eq:newV58-mgf-split} which:
\begin{enumerate}[label=\textup{(\roman*)},leftmargin=3.5em]
\item cancels the \(H(\alpha)\) homogeneous component in every
      Riccati coefficient row;
\item leaves algebraic coefficients bounded by
      \(K^r\alpha^{-(2r-1)}\);
\item proves the resulting local moment-generating function is
      zero-free for \(|z|\le c/s^2\); and
\item performs the subtraction before coefficientwise absolute values.
\end{enumerate}
\end{definition}

\begin{proposition}[RWCC implies CCH]
\label{prop:newV58-RWCC-CCH}
RWCC implies the local-clock cumulant hierarchy CCH.
\end{proposition}

\begin{proof}
RWCC(ii)--(iii) give
\[
 {|\varkappa_r(S_{\rm loc})|\over r!}
 \le K^r\alpha^{-(2r-1)}
 \le (K')^rs^{2r-1}.
\]
This is \eqref{eq:newV57-CCH}; row (i) ensures the full-clock
counterexample of \Cref{cor:newV58-full-CCH-false} has actually been
removed rather than hidden in a fixed-order remainder.
\end{proof}

\begin{programme}[V59: positive-tail Riccati cancellation]
\label{prog:newV58-V59}
\begin{enumerate}[leftmargin=3em]
\item Express the unnormalized tail transform
      \(T_{\rm un}\) by integration by parts against its positive
      distribution function.
\item Use the pointwise heat-kernel comparison at variable group
      distance to derive large-deviation bounds for
      \(S/s\), not only the single threshold \(L\).
\item Convert those bounds into a complex estimate for the centered
      local clock on \(|z|\le c/s^2\).
\item Use Rouch\'e's theorem to prove zero-freeness there and hence
      RWCC.
\item If lower-clock deviations prevent that disk, identify their
      exact Bessel large-order rate before changing the truncation.
\end{enumerate}
\end{programme}

\begin{firewall}[Claim boundary after V58]
\label{fire:newV58-boundary}
V58 derives the exact Riccati equation and triangular coefficient
system for the Hartman--Watson clock logarithm.  It proves that the
decaying homogeneous solution is
\(K_1/I_1\asymp e^{-2\alpha}\), identifies it with the wrapped
sector, proves the logarithmic high-order branch law, and proves that
full-clock CCH is false.  It also proves that positive time truncation
cancels the complete branch and isolates RWCC.

V58 does not prove RWCC, CMDC, CCH, GRLTC, LZTC, PLRC, WSDC, BRSC,
HWSFC, UGCGC, thermal connected WPRAC, all-arity RGSC, the raw Wilson
aggregate strong card, all-order strong MTA, WUFC, CPM, cutoff
removal, reflection positivity, Osterwalder--Schrader reconstruction,
a continuum Yang--Mills measure, or a positive Hamiltonian gap.  The
full Yang--Mills mass gap has not been proved.
\end{firewall}

\begin{theorem}[V58 result ledger]
\label{thm:newV58-results}
V58:
\begin{enumerate}[label=\textup{(V58.\arabic*)},leftmargin=6.5em]
\item derives the exact clock Riccati equation;
\item derives its triangular cumulant-coefficient system;
\item identifies the integrating factor \(\alpha I_1^2\);
\item proves the homogeneous solution is \(K_1/I_1\);
\item proves its exact wrapped scale \(e^{-2\alpha}\);
\item derives the square-root tail of the clock logarithm;
\item disproves CCH for the full clock;
\item proves positive truncation cancels the entire branch; and
\item isolates RWCC as the all-order cancellation gate.
\end{enumerate}
\end{theorem}

\begin{proof}
Use
\Cref{thm:newV58-Riccati,
thm:newV58-coefficient-system,
lem:newV58-H,
thm:newV58-homogeneous-scale,
cor:newV58-row-ambiguity,
thm:newV58-log-branch,
cor:newV58-full-CCH-false,
lem:newV58-branch-cancellation,
def:newV58-RWCC}.
\end{proof}

\paragraph{Parent-version record.}
V58 was formed from
\texttt{Renewal\_Yang\_Mills\_Mass\_Gap\_Research\_Notes\_V57.tex}.
The V57 parent had \(27\,524\) logical source lines,
\(874\,270\) bytes, and SHA--256
\[
 \texttt{a4013d8c2f41744f6842361b96bb3486c5d47a7dbf2776fe9bedd1528fcc62a1}.
\]
The next compulsory companion audit is V60.

% =============================================================================
% END APPENDED FIFTH-SERIES V58 MATERIAL
% =============================================================================

\clearpage
\chapter{Fifth-series V59: positive L\'evy spectral splitting and
closure of the small-clock hierarchy}
\label{ch:newV59-Levy}

\section{Stieltjes measure on the Bessel-order cut}
\label{sec:newV59-Stieltjes}

\begin{definition}[Boundary phase]
\label{def:newV59-Theta}
For \(t>1\), define
\begin{equation}
 \Theta_\alpha(t)
 :=
 \lim_{\varepsilon\downarrow0}
 \Im\Phi_\alpha(-t+i\varepsilon),                             \label{eq:newV59-Theta}
\end{equation}
using the complete-Bernstein continuation of V51.
\end{definition}

\begin{lemma}[The shifted exponent remains complete Bernstein]
\label{lem:newV59-Phi-complete}
\(\Phi_\alpha\) is a complete Bernstein function, and its Stieltjes
measure is supported at spectral parameters \(t\ge1\).
\end{lemma}

\begin{proof}
Write the complete-Bernstein representation of the imported V51
function as
\[
 \psi_\alpha(u)=a+bu+\int_{[0,\infty)}
 {u\over u+r}\,\varrho_\alpha(dr).
\]
With \(u_0=1/2\) and \(h=\lambda/2\),
\[
 {u_0+h\over u_0+h+r}-{u_0\over u_0+r}
 ={hr\over(u_0+r)(u_0+r+h)}
 ={\lambda\over\lambda+(1+2r)}\,{2r\over1+2r}.
\]
Thus
\(\psi_\alpha((1+\lambda)/2)-\psi_\alpha(1/2)\) has a
complete-Bernstein representation with drift \(b/2\) and with the
positive measure obtained by pushing
\((2r/(1+2r))\varrho_\alpha(dr)\) to \(t=1+2r\).  Its support is
contained in \([1,\infty)\).
\end{proof}

\begin{theorem}[Exact Stieltjes representation]
\label{thm:newV59-Stieltjes}
There are \(d_\alpha\ge0\) and a positive measure
\(\sigma_\alpha\) supported on \([1,\infty)\) such that
\begin{equation}
 \Phi_\alpha(\lambda)
 =
 d_\alpha\lambda+
 \int_{[1,\infty)}
 {\lambda\over\lambda+t}\,\sigma_\alpha(dt),
 \qquad\lambda\ge0.                                          \label{eq:newV59-Stieltjes}
\end{equation}
The measure is absolutely continuous on \((1,\infty)\), with
\begin{equation}
 \sigma_\alpha(dt)
 ={\Theta_\alpha(t)\over\pi t}\,dt,\qquad t>1,                \label{eq:newV59-sigma-density}
\end{equation}
and \(\Theta_\alpha(t)\ge0\).
\end{theorem}

\begin{proof}
By \Cref{lem:newV59-Phi-complete}, the complete Bernstein function
\(\Phi_\alpha\), which vanishes at the origin, has the
Stieltjes representation
\[
 \Phi(\lambda)=d\lambda+\int_{(0,\infty)}
 {\lambda\over\lambda+t}\,\sigma(dt).
\]
The V51 continuation of
\[
 \Phi_\alpha(\lambda)
 =-\log{I_{\sqrt{1+\lambda}}(\alpha)\over I_1(\alpha)}
\]
is holomorphic on
\(\mathbb C\setminus(-\infty,-1]\), and it is real analytic on
\((-1,\infty)\).  Stieltjes inversion therefore places
\(\sigma_\alpha\) on \([1,\infty)\).

For \(x>0\), the distributional boundary value of
\(\lambda/(\lambda+t)\) at \(\lambda=-x+i0\) has imaginary part
\(\pi x\delta_x(dt)\).  Hence
\[
 \Im\Phi_\alpha(-x+i0)
 =\pi x\,{d\sigma_\alpha\over dx}(x)
\]
at continuity points.  The Bessel boundary value is continuous for
\(x>1\), giving \eqref{eq:newV59-sigma-density}.  Positivity is also
the Pick property of a complete Bernstein function: it maps the upper
half-plane into itself.  Finally, the finite square-root expansion
\eqref{eq:newV58-log-branch} excludes an atom at \(t=1\), since such
an atom would give a pole proportional to
\(\lambda/(\lambda+1)\).
\end{proof}

\begin{corollary}[Exact L\'evy density and cumulant formula]
\label{cor:newV59-Levy-density}
The L\'evy measure of the full tilted clock has completely monotone
density
\begin{equation}
 \pi_\alpha(x)
 ={1\over\pi}\int_1^\infty
 e^{-tx}\Theta_\alpha(t)\,dt,\qquad x>0.                     \label{eq:newV59-Levy-density}
\end{equation}
For every \(r\ge2\),
\begin{equation}
 {\varkappa_r(S)\over r!}
 =
 \int_{[1,\infty)}{1\over t^r}\sigma_\alpha(dt)
 ={1\over\pi}\int_1^\infty
 {\Theta_\alpha(t)\over t^{r+1}}\,dt.                        \label{eq:newV59-cumulant-phase}
\end{equation}
\end{corollary}

\begin{proof}
Use
\[
 {\lambda\over\lambda+t}
 =\int_0^\infty(1-e^{-\lambda x})t e^{-tx}\,dx
\]
in \eqref{eq:newV59-Stieltjes}.  Thus
\[
 \Pi_\alpha(dx)
 =
 \left(\int t e^{-tx}\sigma_\alpha(dt)\right)dx,
\]
which becomes \eqref{eq:newV59-Levy-density} after
\eqref{eq:newV59-sigma-density}.  A subordinator has
\(\varkappa_r(S)=\int x^r\Pi_\alpha(dx)\) for \(r\ge2\).
Tonelli and
\(\int_0^\infty x^rt e^{-tx}dx=r!t^{-r}\)
give \eqref{eq:newV59-cumulant-phase}.
\end{proof}

\section{Exponential phase suppression below the turning scale}
\label{sec:newV59-phase}

\begin{lemma}[Positive compact part at imaginary order]
\label{lem:newV59-compact-lower}
Let \(y\ge0\) and
\begin{equation}
 A_{\alpha,y}
 :={1\over\pi}\int_0^\pi
 e^{\alpha\cos\theta}\cosh(y\theta)\,d\theta.
 \label{eq:newV59-A}
\end{equation}
For \(\alpha\ge1\),
\begin{equation}
 A_{\alpha,y}\ge c\,{e^\alpha\over\sqrt\alpha}                \label{eq:newV59-A-lower}
\end{equation}
with an absolute \(c>0\), uniformly in \(y\ge0\).
\end{lemma}

\begin{proof}
Restrict the integral to
\(0\le\theta\le\alpha^{-1/2}\).  There
\(\cos\theta\ge1-\theta^2/2\ge1-1/(2\alpha)\) and
\(\cosh(y\theta)\ge1\).  The interval length is
\(\alpha^{-1/2}\), proving the claim.
\end{proof}

\begin{lemma}[Remote part at imaginary order]
\label{lem:newV59-remote}
For \(\nu=iy\), the Schl\"afli formula gives
\begin{equation}
 I_{iy}(\alpha)=A_{\alpha,y}+B_{\alpha,y},                    \label{eq:newV59-Iiy}
\end{equation}
where
\begin{equation}
 |B_{\alpha,y}|
 \le{\sinh(\pi y)\over\pi}K_0(\alpha)
 \le C{e^{\pi y-\alpha}\over\sqrt\alpha}.                    \label{eq:newV59-B-bound}
\end{equation}
\end{lemma}

\begin{proof}
In \eqref{eq:newV53-Schlafli},
\(\cos(iy\theta)=\cosh(y\theta)\) and
\(\sin(i\pi y)=i\sinh(\pi y)\).  The absolute value of the
hyperbolic integral is at most
\[
 \int_0^\infty e^{-\alpha\cosh u}\,du=K_0(\alpha).
\]
Apply the upper bound in \eqref{eq:newV52-K-bounds} and
\(\sinh(\pi y)\le e^{\pi y}/2\).
\end{proof}

\begin{theorem}[Boundary phase is nonperturbative below
\(\alpha^2\)]
\label{thm:newV59-phase-bound}
There is \(C>0\) such that, for all sufficiently large
\(\alpha\) and
\begin{equation}
 1<t\le T_\alpha,
 \qquad
 T_\alpha:=1+{\alpha^2\over4\pi^2},                           \label{eq:newV59-T}
\end{equation}
one has
\begin{equation}
 0\le\Theta_\alpha(t)\le Ce^{-3\alpha/2}.                    \label{eq:newV59-phase-bound}
\end{equation}
\end{theorem}

\begin{proof}
On the upper edge of the cut,
\(\sqrt{1-t+i0}=iy\) with \(y=\sqrt{t-1}\).
The range \eqref{eq:newV59-T} gives
\(0\le y\le\alpha/(2\pi)\).  By
\Cref{lem:newV59-compact-lower,lem:newV59-remote},
\[
 {|B_{\alpha,y}|\over A_{\alpha,y}}
 \le Ce^{-2\alpha+\pi y}
 \le Ce^{-3\alpha/2}.
\]
For large \(\alpha\), the right side is at most \(1/2\), so
\(I_{iy}\) lies in a disk about the positive number
\(A_{\alpha,y}\) and
\[
 |\arg I_{iy}(\alpha)|
 \le2{|B_{\alpha,y}|\over A_{\alpha,y}}.
\]
The branch of
\(\Phi_\alpha=-\log(I_{iy}/I_1)\) is continuous from \(t=1\);
the complete-Bernstein Pick property selects its nonnegative
imaginary part.  Hence
\(\Theta_\alpha(t)=|\arg I_{iy}(\alpha)|\) in this range, proving
\eqref{eq:newV59-phase-bound}.
\end{proof}

\section{Positive small-clock/rare-clock decomposition}
\label{sec:newV59-split}

\begin{definition}[Spectral L\'evy split]
\label{def:newV59-split}
Split
\begin{align}
 \Phi_\alpha^{\rm wr}(\lambda)
 &:=
 \int_{[1,T_\alpha]}
 {\lambda\over\lambda+t}\sigma_\alpha(dt),                   \label{eq:newV59-Phi-wr}\\
 \Phi_\alpha^{\rm sm}(\lambda)
 &:=
 d_\alpha\lambda+
 \int_{(T_\alpha,\infty)}
 {\lambda\over\lambda+t}\sigma_\alpha(dt).                   \label{eq:newV59-Phi-sm}
\end{align}
Let \(S_{\rm wr}\) and \(S_{\rm sm}\) be independent subordinators
at time one with these Laplace exponents.
\end{definition}

\begin{theorem}[Rare clock is compound Poisson]
\label{thm:newV59-rare-Poisson}
\(S_{\rm wr}\) is compound Poisson.  Its jump-count intensity is
\begin{equation}
 \Lambda_{\alpha}^{\rm wr}
 =\sigma_\alpha([1,T_\alpha])
 ={1\over\pi}\int_1^{T_\alpha}
 {\Theta_\alpha(t)\over t}\,dt
 \le Ce^{-q/s}                                                \label{eq:newV59-Lambda}
\end{equation}
for absolute \(C,q>0\).  Hence
\begin{equation}
 \mathbb P(S_{\rm wr}\ne0)
 =1-e^{-\Lambda_\alpha^{\rm wr}}
 \le Ce^{-q/s}.                                               \label{eq:newV59-rare-event}
\end{equation}
\end{theorem}

\begin{proof}
The L\'evy density associated with
\eqref{eq:newV59-Phi-wr} has total mass
\[
 \int_0^\infty\int_{[1,T_\alpha]}
 t e^{-tx}\sigma_\alpha(dt)\,dx
 =\sigma_\alpha([1,T_\alpha]).
\]
Thus the subordinator is compound Poisson with the displayed
intensity.  By \Cref{thm:newV59-phase-bound},
\[
 \Lambda_\alpha^{\rm wr}
 \le Ce^{-3\alpha/2}\log T_\alpha
 \le Ce^{-\alpha}
\]
for large \(\alpha\).  Since
\(\alpha=(2s)^{-1}-1/3\), decrease a numerical \(q>0\) and enlarge
\(C\) to obtain \eqref{eq:newV59-Lambda}.  The probability of at
least one Poisson jump is \(1-e^{-\Lambda}\le\Lambda\).
\end{proof}

\begin{theorem}[Exact independent clock decomposition]
\label{thm:newV59-independent}
The full tilted time satisfies
\begin{equation}
 S\ \stackrel{\rm law}{=}\ S_{\rm sm}+S_{\rm wr},             \label{eq:newV59-clock-sum}
\end{equation}
with independent summands.  Conditional on no wrapped jump,
\(S=S_{\rm sm}\); the complementary event has probability
\(O(e^{-q/s})\).
\end{theorem}

\begin{proof}
By construction,
\(\Phi_\alpha=\Phi_\alpha^{\rm sm}+\Phi_\alpha^{\rm wr}\).
The product of the two independent Laplace transforms is therefore
\(e^{-\Phi_\alpha}=L_\alpha\), which identifies the law.
The conditional statement follows because a compound-Poisson sum is
zero exactly when its jump count is zero.
\end{proof}

\section{Closure of the small-clock cumulant hierarchy}
\label{sec:newV59-CCH}

\begin{theorem}[Small-clock CCH]
\label{thm:newV59-small-CCH}
There is an absolute \(K\) such that
\begin{align}
 \mathbb ES_{\rm sm}&\le Ks,                                  \label{eq:newV59-small-mean}\\
 {\varkappa_r(S_{\rm sm})\over r!}
 &\le K^rs^{2r-1},\qquad r\ge2.                              \label{eq:newV59-small-CCH}
\end{align}
\end{theorem}

\begin{proof}
The first cumulant is
\[
 d_\alpha+
 \int_{(T_\alpha,\infty)}{1\over t}\sigma_\alpha(dt)
 \le\Phi_\alpha'(0)=\mathbb ES\le Cs
\]
by \Cref{lem:newV54-mean}.  For \(r\ge2\), the analogue of
\eqref{eq:newV59-cumulant-phase} gives
\begin{align*}
 {\varkappa_r(S_{\rm sm})\over r!}
 &=\int_{(T_\alpha,\infty)}t^{-r}\sigma_\alpha(dt)\\
 &\le T_\alpha^{-(r-1)}
      \int_{(T_\alpha,\infty)}t^{-1}\sigma_\alpha(dt)\\
 &\le Cs\,T_\alpha^{-(r-1)}.
\end{align*}
Since \(T_\alpha\ge c\alpha^2\ge c's^{-2}\) for
\(\alpha\ge1\), the last line is at most
\(K^rs^{2r-1}\).
\end{proof}

\begin{corollary}[Tangent Gaussian all-arity card for the small
clock]
\label{cor:newV59-tangent}
The subordinated tangent Gaussian driven by \(S_{\rm sm}\) obeys
\[
 \left|
 \kappa_{2r}(\ell_1(X),\ldots,\ell_{2r}(X))
 \right|
 \le
 K^{2r}\left(\prod_i b_i\right)B^{2r-2}
\]
at every arity; all odd cumulants vanish.
\end{corollary}

\begin{proof}
The proof of \Cref{thm:newV57-CCH-suffices} uses only the cumulant
hierarchy and applies verbatim with \(S_{\rm loc}\) replaced by
\(S_{\rm sm}\).  Use \Cref{thm:newV59-small-CCH}.
\end{proof}

\begin{correction}[CCH is closed for the L\'evy small clock, not the
time-conditioned clock]
\label{corr:newV59-CCH}
V57 defined CCH for the conditional variable
\(S\mid S<Ls\).  V59 does not prove that conditional hierarchy.
Instead it supplies the stronger structural replacement
\eqref{eq:newV59-clock-sum}: an independent small-jump clock obeying
CCH plus a compound-Poisson wrapped clock whose nonzero event is
nonperturbative.  This is better adapted to the exact logarithmic
source because independent L\'evy exponents add before derivatives.
\end{correction}

\section{Positive compact-group split induced by the clock}
\label{sec:newV59-group-split}

\begin{corollary}[Exact positive small/wrapped Wilson split]
\label{cor:newV59-Wilson-split}
Let \(N_{\rm wr}\) be the wrapped-clock Poisson count and set
\(\varepsilon_\alpha=1-e^{-\Lambda_\alpha^{\rm wr}}\).
Then
\begin{equation}
 w_\alpha dU
 =(1-\varepsilon_\alpha)\nu_\alpha^{\rm sm}
 +\varepsilon_\alpha\nu_\alpha^{\rm wr},                     \label{eq:newV59-Wilson-split}
\end{equation}
where
\begin{align}
 \nu_\alpha^{\rm sm}(dU)
 &:=\mathbb E[h_{S_{\rm sm}}(U)]\,dU,                         \label{eq:newV59-nu-sm}\\
 \nu_\alpha^{\rm wr}(dU)
 &:=\mathbb E[h_{S_{\rm sm}+S_{\rm wr}}(U)
              \mid N_{\rm wr}\ge1]\,dU.                      \label{eq:newV59-nu-wr}
\end{align}
Both are central probability measures and
\(\varepsilon_\alpha\le Ce^{-q/s}\).
\end{corollary}

\begin{proof}
Use the independent representation
\eqref{eq:newV59-clock-sum} in the exact heat mixture
\eqref{eq:newV51-density-mixture}, then condition on
\(N_{\rm wr}=0\) or \(N_{\rm wr}\ge1\).
\end{proof}

\begin{assessment}[What the L\'evy split gains]
\label{ass:newV59-gain}
The V55 time cutoff bounded all local raw moments but did not prove the
enhanced cumulant powers.  The V59 spectral L\'evy cutoff proves those
powers in one line from the mean because every retained Stieltjes
frequency satisfies \(t\gtrsim s^{-2}\).  The discarded sector is not
signed: it is an independent compound-Poisson clock with exponentially
small nonzero probability.  Thus the tangent Gaussian part of the
all-arity obstruction is closed on the local sector, and all remaining
beyond-all-orders work is confined to the exact compact-group source
conditioned on at least one wrapped jump.
\end{assessment}

\section{The remaining nonlinear merge}
\label{sec:newV59-merge}

\begin{definition}[Small-clock nonlinear grade card, SCNGC]
\label{def:newV59-SCNGC}
SCNGC is the insertion of
\eqref{eq:newV59-small-CCH} into the V50 rooted-Wick
grade/leg forest for the exact small-clock group source, proving that:
\begin{enumerate}[label=\textup{(\roman*)},leftmargin=3.5em]
\item the tangent Gaussian sector is controlled by
      \Cref{cor:newV59-tangent};
\item every nonlinear representation Taylor excess pays the V50
      missing legs;
\item all clock-cumulant partitions are summed with one exponential
      base; and
\item the resulting small-clock connected source obeys
      \(K^m(\prod b_i)B^{m-2}\).
\end{enumerate}
\end{definition}

\begin{definition}[Wrapped Poisson source card, WPSC]
\label{def:newV59-WPSC}
WPSC is an exact connected-logarithm estimate for
\(\nu_\alpha^{\rm wr}\) in
\eqref{eq:newV59-Wilson-split}, retaining the prefactor
\(\varepsilon_\alpha\le Ce^{-q/s}\) until after the complete
compact-group source has been resummed, and yielding the same strong
thermal target without differentiating individual long clock jumps.
\end{definition}

\begin{proposition}[SCNGC plus WPSC closes GRLTC]
\label{prop:newV59-merge}
SCNGC and WPSC imply the Gaussian-renormalized local/tail card GRLTC.
\end{proposition}

\begin{proof}
Split the Wilson measure by
\eqref{eq:newV59-Wilson-split}.  SCNGC supplies the complete local
strong source estimate, including its Gaussian clock mixtures.  WPSC
supplies the exact rare connected-logarithm estimate.  The
noncommutative Fr\'echet-log path of V56 combines the two without a
false logarithmic factorization.  Their sum is precisely GRLTC.
\end{proof}

\begin{programme}[V60: audit and small-clock nonlinear merge]
\label{prog:newV59-V60}
\begin{enumerate}[leftmargin=3em]
\item Perform the compulsory newest-SM/newest-NS audit.
\item Expand the exact small-clock heat source by clock cumulants rather
      than raw clock moments.
\item Prove the rooted connected pairing census with the
      \(s^{2r-1}\) clock weights from
      \eqref{eq:newV59-small-CCH}.
\item Insert the V50 Taylor-excess inequality and test SCNGC at all
      arities.
\item Keep WPSC separate; no rare Poisson jump is to be expanded in
      clock length.
\end{enumerate}
\end{programme}

\begin{firewall}[Claim boundary after V59]
\label{fire:newV59-boundary}
V59 derives the exact Stieltjes boundary-phase representation of the
Hartman--Watson Bernstein exponent, proves exponential phase
suppression below a fixed fraction of the \(\alpha^2\) turning scale,
and constructs an exact independent decomposition into a
small-jump clock and an exponentially rare compound-Poisson wrapped
clock.  It proves CCH for the small clock and therefore closes its
tangent Gaussian all-arity card.

V59 does not prove SCNGC, WPSC, GRLTC, LZTC, PLRC, WSDC, BRSC,
HWSFC, UGCGC, thermal connected WPRAC, all-arity RGSC, the raw Wilson
aggregate strong card, all-order strong MTA, WUFC, CPM, cutoff
removal, reflection positivity, Osterwalder--Schrader reconstruction,
a continuum Yang--Mills measure, or a positive Hamiltonian gap.  The
full Yang--Mills mass gap has not been proved.
\end{firewall}

\begin{theorem}[V59 result ledger]
\label{thm:newV59-results}
V59:
\begin{enumerate}[label=\textup{(V59.\arabic*)},leftmargin=6.5em]
\item derives the exact Stieltjes measure of the Wilson clock exponent;
\item derives the exact L\'evy density and clock-cumulant phase formula;
\item proves the imaginary-order compact/remote decomposition;
\item proves exponential boundary-phase suppression through
      \(t\asymp\alpha^2\);
\item constructs the rare compound-Poisson wrapped clock;
\item proves its nonzero probability is \(O(e^{-q/s})\);
\item constructs the independent small clock;
\item proves small-clock CCH and its tangent Gaussian all-arity card;
      and
\item isolates SCNGC and WPSC as the two remaining nonlinear rows.
\end{enumerate}
\end{theorem}

\begin{proof}
Use
\Cref{thm:newV59-Stieltjes,
cor:newV59-Levy-density,
lem:newV59-compact-lower,
lem:newV59-remote,
thm:newV59-phase-bound,
thm:newV59-rare-Poisson,
thm:newV59-independent,
thm:newV59-small-CCH,
cor:newV59-tangent,
cor:newV59-Wilson-split}.
\end{proof}

\paragraph{Parent-version record.}
V59 was formed from
\texttt{Renewal\_Yang\_Mills\_Mass\_Gap\_Research\_Notes\_V58.tex}.
The V58 parent had \(27\,960\) logical source lines,
\(888\,925\) bytes, and SHA--256
\[
 \texttt{9e07073e181be13c113684a84faff7bc97f6fc8ebb6657fe75aab729b854deda}.
\]
The next compulsory companion audit is V60.

% =============================================================================
% END APPENDED FIFTH-SERIES V59 MATERIAL
% =============================================================================

\clearpage
\chapter{Fifth-series V60: compulsory audit, small-clock jump
hyperedges, and the overlap-resummation gate}
\label{ch:newV60-hyperedges}

\section{Purpose and compulsory companion audit}
\label{sec:newV60-purpose}

V59 replaced the unsuccessful deterministic time cutoff by an exact
L\'evy spectral split.  The retained clock \(S_{\rm sm}\) has
\[
 {\varkappa_r(S_{\rm sm})\over r!}
 \le K^rs^{2r-1},\qquad r\ge2,
\]
whereas the discarded low-spectral clock is an exponentially rare
compound-Poisson variable.  The first purpose of V60 is to insert that
hierarchy into the exact compact-group source without returning to a
finite-order tangent expansion.  The second purpose is to identify,
before taking norms, the precise combinatorial operation still needed
to obtain one exponential arity base.

\begin{audit}[Compulsory V60 companion audit]
\label{audit:newV60-companions}
The newest active companion sources were read atomically and were not
modified:
\begin{align*}
 &\texttt{Renewal\_SM\_Einstein\_Continuum\_Research\_Notes\_V81.tex},\\
 &\qquad 27\,271\ \hbox{logical lines},\quad
 1\,048\,966\ \hbox{bytes},\\
 &\qquad
 \texttt{b6c90b8d33d80f937e1520aa75c3b8396ca95752554bab072cdbb94f797daf67},
 \\
 &\texttt{Renewal\_NS\_Stability\_Research\_Notes\_V96.tex},\\
 &\qquad 39\,160\ \hbox{logical lines},\quad
 1\,666\,700\ \hbox{bytes},\\
 &\qquad
 \texttt{3c5458efb72a294600e2dac8dbd078a20b59be82375603f6da2ecb3f69bf4e7}.
\end{align*}

SM V81 proves that the exact nonlinear source algebra of a
high--high to low merger is an infinite two-index ancestry lattice;
the first few output modes locate the response but are not closed
under the nonlinear powers.  Its frequency-uniform compactness
argument sums that infinite ancestry before passing to the
homogenized equation.

NS V96 transfers a Gaussian-weighted pinching and flux ledger from
rescaling limits to the solution at homogeneous points.  It also
separates the exact rate-free identity from the stronger density
conclusions which need a type-I rate.  In particular, an identity
does not supply a uniform estimate unless the scale ledger pays for
it.
\end{audit}

\begin{assessment}[Transfer to the Yang--Mills attack]
\label{ass:newV60-audit-transfer}
No SM constraint estimate and no NS flux estimate is imported into
Yang--Mills.  The useful proof discipline is nevertheless exact:
\begin{enumerate}[label=\textup{(\roman*)},leftmargin=3.5em]
\item retain the full infinite jump-ancestry algebra rather than a
      finite collection of clock cumulants;
\item attach a scale weight to every ancestry merger; and
\item distinguish an exact endpoint-factorizing identity from the
      uniform summability statement required after differentiation.
\end{enumerate}
The construction below follows these three rules.
\end{assessment}

\section{The small-clock L\'evy measure}
\label{sec:newV60-Levy}

Retain the notation of V59.  Thus
\begin{equation}
 \Phi_\alpha^{\rm sm}(\lambda)
 =
 d_\alpha\lambda+
 \int_{(T_\alpha,\infty)}
 {\lambda\over\lambda+t}\,\sigma_\alpha(dt),
 \qquad
 T_\alpha=1+{\alpha^2\over4\pi^2},                            \label{eq:newV60-Phi}
\end{equation}
and put
\begin{equation}
 \Pi_\alpha^{\rm sm}(dx)
 =
 \left\{
 \int_{(T_\alpha,\infty)}t e^{-tx}\sigma_\alpha(dt)
 \right\}dx.                                                  \label{eq:newV60-Pi}
\end{equation}
Then
\begin{equation}
 \Phi_\alpha^{\rm sm}(\lambda)
 =
 d_\alpha\lambda+
 \int_0^\infty(1-e^{-\lambda x})\Pi_\alpha^{\rm sm}(dx).
 \label{eq:newV60-Levy-Khintchine}
\end{equation}

\begin{lemma}[All positive jump moments at the correct scale]
\label{lem:newV60-jump-moments}
For every real \(p\ge1\),
\begin{equation}
 \int_0^\infty x^p\Pi_\alpha^{\rm sm}(dx)
 \le
 \Gamma(p+1)K^{2p}s^{2p-1}.                                  \label{eq:newV60-jump-moments}
\end{equation}
Moreover \(0\le d_\alpha\le Ks\).
\end{lemma}

\begin{proof}
Tonelli's theorem and
\(\int_0^\infty x^pt e^{-tx}dx=\Gamma(p+1)t^{-p}\)
give
\[
 \int x^p\Pi_\alpha^{\rm sm}(dx)
 =
 \Gamma(p+1)
 \int_{(T_\alpha,\infty)}t^{-p}\sigma_\alpha(dt).
\]
Since \(p\ge1\),
\[
 \int_{(T_\alpha,\infty)}t^{-p}\sigma_\alpha(dt)
 \le
 T_\alpha^{1-p}
 \int_{(T_\alpha,\infty)}t^{-1}\sigma_\alpha(dt).
\]
The last integral and \(d_\alpha\) are nonnegative summands of
\(\mathbb ES_{\rm sm}\le Ks\) from
\Cref{thm:newV59-small-CCH}.  Finally
\(T_\alpha\ge c s^{-2}\); absorb \(c^{1-p}\) into \(K^{2p}\).
\end{proof}

\section{One heat jump on an arbitrary label set}
\label{sec:newV60-one-jump}

For a representation \(R\), write
\(\Delta^R(U):=D^R(U)-I\).  For a nonempty
\(A\subseteq[m]\), let
\begin{equation}
 J_A(x)
 :=
 \mathbb E_{h_x}
 \bigotimes_{i\in A}\Delta^{R_i}(U),                          \label{eq:newV60-JA}
\end{equation}
embedded by identities on the complementary tensor factors.

\begin{lemma}[Dimension-free heat-increment moments]
\label{lem:newV60-heat-increment}
There is a numerical \(C\) such that, for every \(p\ge2\),
\(x>0\), and finite-dimensional unitary representation \(R\) of
\(\SU(2)\),
\begin{equation}
 \left(
 \mathbb E_{h_x}\|\Delta^R(U)\|_{\rm op}^p
 \right)^{1/p}
 \le
 C\sqrt{p\,x\{1+C_2(R)\}}.                                   \label{eq:newV60-heat-increment}
\end{equation}
The same estimate holds at every matrix amplification.
\end{lemma}

\begin{proof}
Give \(\SU(2)\) its fixed bi-invariant metric.  Along a minimizing
geodesic from \(I\) to \(U\), differentiation of \(D^R\) and
\(\|\sum_a v_aJ_R^a\|\le |v|\sqrt{C_2(R)}\) give
\[
 \|D^R(U)-I\|
 \le d(I,U)\sqrt{C_2(R)}.
\]
The same inequality is unchanged by amplification.  Brownian motion
on the fixed three-dimensional compact manifold \(\SU(2)\) obeys
\[
 \left(\mathbb E d(I,U_x)^p\right)^{1/p}
 \le C\sqrt{px}
\]
when \(px\le1\).  This follows, for example, by integrating the
standard Gaussian exit bound
\(\mathbb P\{d(I,U_x)>r\}\le C e^{-r^2/(Cx)}\) up to the injectivity
radius and using compactness beyond it.  When \(px>1\), use
\(\|D^R(U)-I\|\le2\), enlarge \(C\), and use
\(1+C_2(R)\ge1\).  Combining the two regimes proves
\eqref{eq:newV60-heat-increment}.
\end{proof}

\begin{corollary}[Raw same-jump product]
\label{cor:newV60-same-jump-raw}
If \(a=|A|\ge2\), then
\begin{equation}
 \|J_A(x)\|_{\rm cb}
 \le
 (C\sqrt a)^a x^{a/2}
 \prod_{i\in A}\sqrt{1+C_2(R_i)}.                             \label{eq:newV60-same-jump-raw}
\end{equation}
\end{corollary}

\begin{proof}
Take the norm inside the expectation in
\eqref{eq:newV60-JA}, apply H\"older with exponent \(a\) to the
\(a\) factors, and use
\Cref{lem:newV60-heat-increment}.  Tensor amplification does not
change any factor norm.
\end{proof}

\section{The all-arity jump-hyperedge card}
\label{sec:newV60-hyperedge}

Define the small-clock jump interaction
\begin{equation}
 Q_A
 :=
 \int_0^\infty J_A(x)\Pi_\alpha^{\rm sm}(dx),
 \qquad |A|\ge2.                                              \label{eq:newV60-QA}
\end{equation}
For a singleton, define instead
\[
 Q_{\{i\}}
 :=
 \int_0^\infty
 \left(e^{-xC_2(R_i)}-I\right)
 \Pi_\alpha^{\rm sm}(dx);
\]
this is finite because
\(\|e^{-xC}-I\|\le xC\) and \(\int x\Pi(dx)<\infty\).
Put
\[
 b_i:=\sqrt{s\{1+C_2(R_i)\}},
 \qquad B_A:=\sum_{i\in A}b_i.
\]

\begin{theorem}[Strong same-jump hyperedge card]
\label{thm:newV60-hyperedge-card}
There is one numerical \(K\) such that, for every
\(A\subseteq[m]\) with \(a=|A|\ge2\),
\begin{equation}
 \boxed{\qquad
 \|Q_A\|_{\rm cb}
 \le
 K^a\left(\prod_{i\in A}b_i\right)B_A^{a-2}.
 \qquad}                                                      \label{eq:newV60-hyperedge-card}
\end{equation}
The constant is independent of \(a,s\), the representations, and
all matrix amplifications.
\end{theorem}

\begin{proof}
Apply
\Cref{cor:newV60-same-jump-raw,lem:newV60-jump-moments}
with \(p=a/2\):
\begin{align*}
 \|Q_A\|_{\rm cb}
 &\le
 (C\sqrt a)^a
 \prod_{i\in A}\sqrt{1+C_2(R_i)}
 \int_0^\infty x^{a/2}\Pi_\alpha^{\rm sm}(dx)\\
 &\le
 (K\sqrt a)^a\Gamma(a/2+1)
 s^{a-1}
 \prod_{i\in A}{b_i\over\sqrt s}\\
 &=
 (K\sqrt a)^a\Gamma(a/2+1)
 s^{a/2-1}\prod_{i\in A}b_i.                                 \tag{*}
\end{align*}
The elementary Stirling bound gives
\(\Gamma(a/2+1)\le(C\sqrt a)^a\).  Also every
\(b_i\ge\sqrt s\), whence
\[
 s^{(a-2)/2}\le\left({B_A\over a}\right)^{a-2}.
\]
Substitute this in \((*)\).  The remaining numerical factor is at
most
\[
 K^a{a^a\over a^{a-2}}=K^aa^2,
\]
and \(a^2\) is absorbed into a larger exponential base.
\end{proof}

\begin{corollary}[Exact Pr\"ufer form]
\label{cor:newV60-Prufer-form}
If \(\mathfrak T(A)\) is the set of labelled trees on \(A\), then
\begin{equation}
 \|Q_A\|_{\rm cb}
 \le
 K^{|A|}
 \sum_{T\in\mathfrak T(A)}
 \prod_{i\in A}b_i^{d_i(T)}.                                 \label{eq:newV60-Prufer-form}
\end{equation}
\end{corollary}

\begin{proof}
The weighted Cayley--Pr\"ufer identity is
\[
 \sum_{T\in\mathfrak T(A)}
 \prod_{i\in A}b_i^{d_i(T)}
 =
 \left(\prod_{i\in A}b_i\right)B_A^{|A|-2}.
\]
Apply \Cref{thm:newV60-hyperedge-card}.
\end{proof}

\begin{assessment}[The clock-root census is closed one root at a time]
\label{ass:newV60-one-root}
The factorial in the clock cumulant is not discarded.  It appears as
\(\Gamma(a/2+1)\) in the positive jump moment, while the heat
increment moment supplies the complementary \(a^{a/2}\).  Their
combined \(a^a\) is exactly of the size already present in the
Pr\"ufer polynomial \(B_A^{a-2}\), because the thermal floor gives
\(B_A\ge a\sqrt s\).  Thus one L\'evy clock root of arbitrary arity
has one exponential base.  No finite clock-cumulant truncation has
been used.
\end{assessment}

\section{Exact jump ancestry and endpoint factorization}
\label{sec:newV60-ancestry}

For \(B\subseteq[m]\), put
\[
 C_B^{\rm tot}
 =
 \sum_{a=1}^3\left(\sum_{i\in B}J_i^a\right)^2.
\]

\begin{lemma}[Subset decomposition of one heat jump]
\label{lem:newV60-subset}
For every nonempty \(B\subseteq[m]\),
\begin{equation}
 e^{-xC_B^{\rm tot}}-I
 =
 \sum_{\varnothing\ne A\subseteq B}J_A(x).                   \label{eq:newV60-subset}
\end{equation}
Consequently, with
\begin{equation}
 L_B^{\rm sm}:=-\Phi_\alpha^{\rm sm}(C_B^{\rm tot}),          \label{eq:newV60-LB}
\end{equation}
one has the exact finite ancestry expansion
\begin{equation}
 L_B^{\rm sm}
 =
 -d_\alpha C_B^{\rm tot}
 +\sum_{\varnothing\ne A\subseteq B}Q_A.                     \label{eq:newV60-LB-subsets}
\end{equation}
\end{lemma}

\begin{proof}
Under the heat law,
\[
 e^{-xC_B^{\rm tot}}
 =
 \mathbb E_{h_x}\bigotimes_{i\in B}D^{R_i}(U).
\]
Expand every factor as \(I+\Delta^{R_i}(U)\) and take the
expectation.  This gives \eqref{eq:newV60-subset}.  Insert it in the
functional-calculus form of
\eqref{eq:newV60-Levy-Khintchine}.  Every sum is finite, and every
Bochner integral is absolutely convergent by the singleton estimate
and \Cref{cor:newV60-same-jump-raw}; hence sum and integral may be
interchanged.
\end{proof}

For \(A\) with \(|A|\ge2\), define its weighted tree polynomial
\[
 W_A:=\sum_{T\in\mathfrak T(A)}
 \prod_{i\in A}b_i^{d_i(T)}
 =
 \left(\prod_{i\in A}b_i\right)B_A^{|A|-2}
\]
and the gate
\begin{equation}
 \chi_A(\boldsymbol x)
 :=
 {1\over W_A}
 \sum_{T\in\mathfrak T(A)}
 \left(\prod_{i\in A}b_i^{d_i(T)}\right)
 \prod_{ij\in E(T)}x_{ij}.                                  \label{eq:newV60-gate}
\end{equation}
Set \(\chi_{\{i\}}=1\).  Also put
\begin{equation}
 {\cal C}_m(\boldsymbol x)
 =
 \sum_iC_2(R_i)I+
 2\sum_{i<j}x_{ij}\Omega_{ij}.                               \label{eq:newV60-Cx}
\end{equation}

\begin{lemma}[Gate endpoints]
\label{lem:newV60-gate-endpoints}
For \(0\le x_{ij}\le1\), one has \(0\le\chi_A\le1\).
If \(\boldsymbol x^\rho\) is the zero--one matrix of a partition
\(\rho\), then
\begin{equation}
 \chi_A(\boldsymbol x^\rho)
 =
 \begin{cases}
 1,&A\ \hbox{is contained in one block of }\rho,\\
 0,&\hbox{otherwise}.
 \end{cases}                                                  \label{eq:newV60-gate-endpoints}
\end{equation}
\end{lemma}

\begin{proof}
The gate is a convex combination of tree monomials in
\([0,1]\).  If \(A\) lies in one block, all its pair variables are
one.  If it meets two blocks, every spanning tree on \(A\) contains
a cross-block edge, whose variable is zero.
\end{proof}

Define the exact interpolation
\begin{align}
 {\cal V}_m(\boldsymbol x)
 &:=
 -d_\alpha{\cal C}_m(\boldsymbol x)
 +\sum_iQ_{\{i\}}
 +\sum_{\substack{A\subseteq[m]\\|A|\ge2}}
 \chi_A(\boldsymbol x)Q_A,                                  \label{eq:newV60-Vx}\\
 {\cal F}_m(\boldsymbol x)
 &:=e^{{\cal V}_m(\boldsymbol x)}.                            \label{eq:newV60-Fx}
\end{align}

\begin{theorem}[Exact endpoint-factorizing small-clock source]
\label{thm:newV60-endpoints}
For every partition \(\rho\) of \([m]\),
\begin{equation}
 {\cal F}_m(\boldsymbol x^\rho)
 =
 \bigotimes_{B\in\rho}
 e^{-\Phi_\alpha^{\rm sm}(C_B^{\rm tot})}.                   \label{eq:newV60-endpoints}
\end{equation}
In particular, the all-one endpoint is the exact small-clock moment
on the full tensor block.
\end{theorem}

\begin{proof}
At a partition endpoint,
\[
 {\cal C}_m(\boldsymbol x^\rho)
 =
 \sum_{B\in\rho}C_B^{\rm tot}.
\]
By \Cref{lem:newV60-gate-endpoints}, the jump terms in
\eqref{eq:newV60-Vx} are exactly
\[
 \sum_{B\in\rho}
 \sum_{\varnothing\ne A\subseteq B}Q_A.
\]
Together with \eqref{eq:newV60-LB-subsets}, this says
\({\cal V}_m(\boldsymbol x^\rho)=\sum_{B\in\rho}L_B^{\rm sm}\).
Terms belonging to disjoint tensor blocks commute, so exponentiation
factorizes and proves \eqref{eq:newV60-endpoints}.
\end{proof}

\begin{corollary}[Exact BKAR forest for the small-clock endpoint]
\label{cor:newV60-BKAR}
For \(m\ge2\), the exact tensor cumulant under
\(\nu_\alpha^{\rm sm}\) is
\begin{equation}
 \kappa_m^{\nu_\alpha^{\rm sm}}
 (D^{R_1},\ldots,D^{R_m})
 =
 \sum_{T\in\mathfrak T_m}
 \int_{[0,1]^{m-1}}
 \partial_{E(T)}
 {\cal F}_m(\boldsymbol x^T(\boldsymbol u))
 \,d\boldsymbol u.                                          \label{eq:newV60-BKAR}
\end{equation}
\end{corollary}

\begin{proof}
Apply the multivariate BKAR identity to
\({\cal F}_m\).  By \Cref{thm:newV60-endpoints}, every partition
endpoint is the tensor product of the exact small-clock moments on
its blocks.  The partition M\"obius sum is therefore the displayed
tensor cumulant, and disconnected forests cancel exactly as in
\Cref{thm:newV29-heat-forest}.
\end{proof}

\section{Stability and exact derivative payment}
\label{sec:newV60-stability}

\begin{lemma}[Small-\(B\) interaction stability]
\label{lem:newV60-interaction-stability}
There are numerical \(b_0,c>0\) such that, if
\(B=\sum_i b_i\le b_0\), then
\begin{equation}
 \sum_{\substack{A\subseteq[m]\\|A|\ge2}}\|Q_A\|_{\rm cb}
 \le cB^2.                                                    \label{eq:newV60-interaction-stability}
\end{equation}
On every BKAR forest matrix
\(\boldsymbol x^T(\boldsymbol u)\),
\begin{equation}
 \|{\cal F}_m(\boldsymbol x^T(\boldsymbol u))\|_{\rm cb}
 \le e^{cB^2}.                                                \label{eq:newV60-F-stable}
\end{equation}
\end{lemma}

\begin{proof}
By \Cref{thm:newV60-hyperedge-card},
\[
 \sum_{|A|=a}\|Q_A\|
 \le
 K^aB^{a-2}\sum_{|A|=a}\prod_{i\in A}b_i
 \le {K^aB^{2a-2}\over a!}.
\]
Sum \(a\ge2\) and choose \(b_0\) so that \(KB^2\le1/2\).
This proves \eqref{eq:newV60-interaction-stability}.

A BKAR forest matrix is positive semidefinite.  Therefore
\({\cal C}_m(\boldsymbol x^T(\boldsymbol u))\ge0\) by the total
Casimir identity of V28.  Each singleton \(Q_{\{i\}}\le0\), and
all gates lie in \([0,1]\).  Thus the first two terms of
\eqref{eq:newV60-Vx} are negative semidefinite, while the norm of
the remaining interaction is at most \(cB^2\).  Since
\({\cal V}_m\) is self-adjoint,
\(\|e^{{\cal V}_m}\|\le e^{cB^2}\).
\end{proof}

\begin{lemma}[A mixed gate derivative retains a spanning-tree
extension]
\label{lem:newV60-gate-derivative}
Let \(E\) be a set of distinct pair variables and \(|A|\ge2\).
Then
\begin{equation}
 \left\|
 \partial_E\{\chi_A(\boldsymbol x)Q_A\}
 \right\|_{\rm cb}
 \le
 K^{|A|}
 \sum_{\substack{T\in\mathfrak T(A)\\E\subseteq E(T)}}
 \prod_{i\in A}b_i^{d_i(T)}                                 \label{eq:newV60-gate-derivative}
\end{equation}
for \(0\le x_{ij}\le1\), with the convention that the sum is zero
unless every edge of \(E\) lies in \(A\).
\end{lemma}

\begin{proof}
Differentiate \eqref{eq:newV60-gate}.  Only trees containing \(E\)
survive, and the undifferentiated pair variables have modulus at
most one.  Multiply the resulting bound by
\(\|Q_A\|\le K^{|A|}W_A\) from
\Cref{thm:newV60-hyperedge-card}; the normalizing \(W_A\) cancels.
\end{proof}

\begin{assessment}[Exact gain over the V50 fixed-grade proof]
\label{ass:newV60-gain}
\Cref{cor:newV60-BKAR,lem:newV60-interaction-stability,
lem:newV60-gate-derivative} provide an exact, nonperturbative
small-clock forest, a stable exponential on the only difficult
small-\(B\) regime, and a real spanning-tree payment for every
mixed clock-root derivative.  This removes the V34 infinite-jet
problem and the V40 derivative-window problem at the level of a
single jump root.  What remains is the overlap of several such
roots inside the same Duhamel derivative.
\end{assessment}

\section{The Bell-overlap obstruction to termwise differentiation}
\label{sec:newV60-Bell}

\begin{proposition}[Naive Fr\'echet differentiation loses the
exponential base]
\label{prop:newV60-Bell-no-go}
Take \eqref{eq:newV60-gate-derivative} as the only estimate on mixed
derivatives of \({\cal V}_m\), insert it term by term into the
multivariate Fr\'echet formula for
\(\partial_{E(T)}e^{{\cal V}_m}\), and then take norms.  The resulting
positive majorant cannot be bounded by
\(K^m\prod_i b_i^{d_i(T)}\) with one \(K\) independent of \(m\).
\end{proposition}

\begin{proof}
Let \(T\) be the star with centre \(1\) and \(k=m-1\) leaf edges.
In the Fr\'echet--Fa\`a di Bruno expansion, the edge derivatives are
partitioned among derivatives of \({\cal V}_m\).  Every block \(F\)
of a set partition of the star edges is itself contained in the
star tree on the centre and the leaves met by \(F\).  Hence
\eqref{eq:newV60-gate-derivative}, after its extra tree edges are
discarded, permits a contribution proportional to
\(\prod_{e\in F}b_{e_-}b_{e_+}\) for every block.

Multiplying over the blocks of the partition gives the same full
star weight for every set partition.  Thus this termwise positive
majorant contains the Bell number \({\rm Bell}(k)\) times that
weight.  Even the subfamily of partitions into pairs has cardinality
\[
 {k!\over2^{k/2}(k/2)!}
\]
for even \(k\), whose \(k\)-th root is unbounded.  It cannot be
absorbed into \(K^m\).  This is a statement about the indicated
positive majorant, not a lower bound on the exact operator
cumulant.
\end{proof}

\begin{warning}[Do not infer failure of SCNGC]
\label{warn:newV60-not-failure}
The Bell family in
\Cref{prop:newV60-Bell-no-go} is created by differentiating the
outer exponential before resumming overlapping jump roots.  In the
tangent Gaussian specialization, the exact logarithm collapses the
same apparent family to one clock cumulant, and V59 proves the
correct all-arity bound.  Therefore the proposition rules out a
proof order; it does not disprove the small-clock nonlinear card.
\end{warning}

\begin{definition}[Overlap-cumulant resummation card, OCRC]
\label{def:newV60-OCRC}
OCRC is a reorganization of
\eqref{eq:newV60-BKAR} which, before norms:
\begin{enumerate}[label=\textup{(\roman*)},leftmargin=3.5em]
\item groups all jump roots meeting the same representation label
      into the exact convolution/logarithmic word at that label;
\item cancels set partitions disconnected in the joint
      label--jump incidence graph;
\item retains the spanning-tree extensions of
      \eqref{eq:newV60-gate-derivative}; and
\item bounds the complete overlap sum by
      \(K^m\prod_i b_i^{d_i(T)}\) for every outer tree \(T\).
\end{enumerate}
\end{definition}

\begin{proposition}[OCRC is the remaining row of SCNGC]
\label{prop:newV60-OCRC-SCNGC}
OCRC, together with the theorems of V59 and V60, implies SCNGC.
\end{proposition}

\begin{proof}
The tangent row is
\Cref{cor:newV59-tangent}.  The arbitrary-arity one-root weights are
\Cref{thm:newV60-hyperedge-card}.  Exact endpoint factorization and
the connected forest are
\Cref{thm:newV60-endpoints,cor:newV60-BKAR}; stability and derivative
payment are
\Cref{lem:newV60-interaction-stability,
lem:newV60-gate-derivative}.  OCRC supplies precisely the only
missing uniform summation over overlapping roots.  Summing the
fixed-tree estimate over trees with the weighted Pr\"ufer identity
then gives
\[
 K^m\left(\prod_i b_i\right)B^{m-2},
\]
which is SCNGC.
\end{proof}

\begin{programme}[V61: go upstream to the overlap logarithm]
\label{prog:newV60-V61}
\begin{enumerate}[leftmargin=3em]
\item Test OCRC first in the exact abelian scalar source, where the
      logarithm can be evaluated before any jump-word expansion.
\item Express the noncommutative source as a time-ordered convolution
      word of independent group jumps and apply moment--cumulant
      cancellation to the joint label--jump incidence graph.
\item Seek a Penrose allocation on that incidence graph which charges
      an overlap at label \(i\) to the unused Taylor factorial of the
      complete matrix exponential, rather than to another absolute
      clock moment.
\item Recover the V59 tangent formula as the two-leg quotient of the
      proposed resummation.
\item Keep the rare wrapped Poisson source outside this calculation;
      no long jump is to be differentiated in clock length.
\end{enumerate}
\end{programme}

\begin{firewall}[Claim boundary after V60]
\label{fire:newV60-boundary}
V60 performs the compulsory companion audit, proves all positive
moments of the small-clock L\'evy measure at scale \(s^{2p-1}\),
proves a dimension-free heat-increment estimate, and proves the
all-arity strong same-jump hyperedge card.  It constructs an exact
endpoint-factorizing interpolation and BKAR forest for the complete
small-clock compact-group source, proves small-\(B\) stability, and
proves exact spanning-tree payment for every mixed gate derivative.
It also proves that termwise Fr\'echet differentiation produces a
Bell-number majorant and isolates OCRC as the required upstream
resummation.

V60 does not prove OCRC, SCNGC, WPSC, GRLTC, LZTC, PLRC, WSDC,
BRSC, HWSFC, UGCGC, thermal connected WPRAC, all-arity RGSC, the
raw Wilson aggregate strong card, all-order strong MTA, WUFC, CPM,
cutoff removal, reflection positivity,
Osterwalder--Schrader reconstruction, a continuum Yang--Mills
measure, or a positive Hamiltonian gap.  The full Yang--Mills mass
gap has not been proved.
\end{firewall}

\begin{theorem}[V60 result ledger]
\label{thm:newV60-results}
V60:
\begin{enumerate}[label=\textup{(V60.\arabic*)},leftmargin=6.5em]
\item completes the required newest-SM/newest-NS audit;
\item proves the fractional-moment hierarchy of the small-clock
      L\'evy measure;
\item proves the dimension-free compact-group heat-increment bound;
\item proves the strong same-jump hyperedge card at every arity;
\item rewrites that card in exact weighted Pr\"ufer form;
\item constructs the complete jump-subset ancestry algebra;
\item constructs an exact endpoint-factorizing small-clock source;
\item derives its exact connected BKAR forest;
\item proves small-source stability and mixed-derivative tree
      payment; and
\item proves the Bell-overlap no-go for termwise differentiation and
      isolates OCRC.
\end{enumerate}
\end{theorem}

\begin{proof}
Use
\Cref{audit:newV60-companions,
lem:newV60-jump-moments,
lem:newV60-heat-increment,
cor:newV60-same-jump-raw,
thm:newV60-hyperedge-card,
cor:newV60-Prufer-form,
lem:newV60-subset,
lem:newV60-gate-endpoints,
thm:newV60-endpoints,
cor:newV60-BKAR,
lem:newV60-interaction-stability,
lem:newV60-gate-derivative,
prop:newV60-Bell-no-go,
prop:newV60-OCRC-SCNGC}.
\end{proof}

\paragraph{Parent-version record.}
V60 was formed from
\texttt{Renewal\_Yang\_Mills\_Mass\_Gap\_Research\_Notes\_V59.tex}.
The V59 parent had \(28\,510\) logical source lines,
\(907\,215\) bytes, and SHA--256
\[
 \texttt{90b5fb5e59ba9440df71f08d599a01ff4cf33fcab9b7ea1f8d79a2c14089699a}.
\]
The next compulsory companion audit is V65.

% =============================================================================
% END APPENDED FIFTH-SERIES V60 MATERIAL
% =============================================================================

\clearpage
\chapter{Fifth-series V61: the clock-factorial reserve and exact
one-centre overlap resummation}
\label{ch:newV61-overlap}

\section{Why V60's Bell family is not yet an obstruction}
\label{sec:newV61-context}

No companion audit is due at V61.  V60 proved that differentiating
the endpoint-factorizing exponential term by term, while retaining
only the relaxed Pr\"ufer bound on each jump root, creates a Bell
number on a fixed star.  That result is a valid no-go for that proof
order.  It deliberately discarded a stronger factor which appeared
one line earlier in the same-jump proof.  V61 retains that factor and
then performs the overlap sum in the logarithmic source.

\begin{correction}[OCRC need not be proved fixed tree by fixed tree]
\label{corr:newV61-OCRC}
The fixed-tree version of OCRC in V60 is sufficient for SCNGC but is
stronger than the final aggregate card.  A Bell family attached to one
fixed star may be absorbed either by the complete aggregate Pr\"ufer
sum or by an exact source quotient, provided the clock-factorial
reserve is kept.  It is therefore incorrect to promote
\Cref{prop:newV60-Bell-no-go} from a no-go for a positive majorant to
an obstruction to the exact source.
\end{correction}

\section{The clock-factorial reserve}
\label{sec:newV61-reserve}

The drift part of the small clock will be included with the two-label
jump interaction.  Define
\begin{equation}
 R_A:=
 \begin{cases}
 Q_A-2d_\alpha\Omega_{ij},&A=\{i,j\},\\
 Q_A,&|A|\ge3.
 \end{cases}                                                  \label{eq:newV61-RA}
\end{equation}

\begin{lemma}[Unrelaxed same-jump estimate]
\label{lem:newV61-reserve}
There is a numerical \(K\) such that, for
\(A\subseteq[m]\), \(a=|A|\ge2\),
\begin{equation}
 \boxed{\qquad
 \|R_A\|_{\rm cb}
 \le
 (Ka)^a s^{(a-2)/2}\prod_{i\in A}b_i .
 \qquad}                                                      \label{eq:newV61-reserve}
\end{equation}
\end{lemma}

\begin{proof}
Before using \(B_A\ge a\sqrt s\), the proof of
\Cref{thm:newV60-hyperedge-card} gave
\[
 \|Q_A\|_{\rm cb}
 \le
 (K\sqrt a)^a\Gamma(a/2+1)
 s^{a/2-1}\prod_{i\in A}b_i.
\]
The Stirling bound
\(\Gamma(a/2+1)\le(C\sqrt a)^a\) proves
\eqref{eq:newV61-reserve} for \(Q_A\).  If \(A=\{i,j\}\), then
\[
 2d_\alpha\|\Omega_{ij}\|
 \le
 2Ks\sqrt{C_2(R_i)C_2(R_j)}
 \le
 2K b_i b_j
\]
by \Cref{lem:newV60-jump-moments}.  This has the right side of
\eqref{eq:newV61-reserve} at \(a=2\).
\end{proof}

\begin{assessment}[Two compatible payments]
\label{ass:newV61-two-payments}
The relaxed estimate
\[
 \|R_A\|\le K^a
 \left(\prod_{i\in A}b_i\right)B_A^{a-2}
\]
pays a spanning tree inside one root.  The unrelaxed estimate
\eqref{eq:newV61-reserve} instead pays
\((a\sqrt s)^{a-2}\).  The former is best when the marks in \(A\)
are comparable; the latter is strictly stronger when one large mark
is repeated as the centre of many roots.  An overlap proof must be
allowed to choose between them before summing partitions.
\end{assessment}

\section{Direct resummation of the star partition family}
\label{sec:newV61-star}

\begin{lemma}[Weighted Bell-star resummation]
\label{lem:newV61-Bell-star}
Let \(n\ge1\), let a centre have weight
\(\beta\in[\sqrt s,1)\), and let the leaves have weights
\(c_1,\ldots,c_n\in[\sqrt s,1)\).  Put
\[
 B:=\beta+\sum_{j=1}^nc_j.
\]
Suppose that, for every nonempty
\(C\subseteq[n]\), a nonnegative activity obeys
\begin{equation}
 a(C)
 \le
 \{K(|C|+1)\}^{|C|+1}
 s^{(|C|-1)/2}\,
 \beta\prod_{j\in C}c_j.                                    \label{eq:newV61-star-activity}
\end{equation}
Then
\begin{equation}
 \sum_{\mathcal P\in{\rm Part}([n])}
 \prod_{C\in\mathcal P}a(C)
 \le
 K_1^{n+1}\,
 \beta\left(\prod_{j=1}^nc_j\right)B^{n-1}.                  \label{eq:newV61-star-sum}
\end{equation}
\end{lemma}

\begin{proof}
Factor \(\prod_jc_j\).  For a block of size \(r\), the remaining
majorant is
\[
 A_r:=\{K(r+1)\}^{r+1}\beta s^{(r-1)/2}.
\]
The labelled exponential formula gives
\begin{equation}
 \sum_{\mathcal P\in{\rm Part}([n])}
 \prod_{C\in\mathcal P}A_{|C|}
 =
 n![z^n]\exp H(z),
 \qquad
 H(z):=\sum_{r\ge1}{A_r\over r!}z^r.                         \label{eq:newV61-star-EGF}
\end{equation}
For \(n\ge1\), the coefficient is unchanged if
\(\exp H\) is replaced by \(\exp H-1\).

Stirling's lower bound on \(r!\) gives
\[
 {A_r\over r!}
 \le
 \beta(CK)^r r^2s^{(r-1)/2}.
\]
Choose a numerical \(\eta>0\), depending only on \(K\), and set
\[
 R:={\eta n\over B}.
\]
Because every leaf is at least \(\sqrt s\),
\[
 n\sqrt s\le B,
 \qquad
 CK\sqrt s\,R\le CK\eta.
\]
Choose \(\eta\) so that the last quantity is at most \(1/4\).
On \(|z|=R\),
\begin{equation}
 |H(z)|
 \le
 C\beta R
 \sum_{r\ge1}r^2(CK\sqrt sR)^{r-1}
 \le C\beta R
 \le C\eta n.                                                \label{eq:newV61-H-bound}
\end{equation}
Moreover
\[
 |e^{H(z)}-1|
 \le |H(z)|e^{|H(z)|}
 \le C\beta R e^{C\eta n}.
\]
Cauchy's estimate in \eqref{eq:newV61-star-EGF} and
\(n!\le n^n\) now give
\begin{align*}
 n!|[z^n](e^{H}-1)|
 &\le
 C\beta n!R^{1-n}e^{C\eta n}\\
 &\le
 C\beta n^n
 \left({\eta n\over B}\right)^{1-n}e^{C\eta n}\\
 &\le
 K_1^{n+1}\beta B^{n-1}.
\end{align*}
Restore \(\prod_jc_j\).
\end{proof}

\begin{corollary}[The V60 Bell family has the strong aggregate
weight]
\label{cor:newV61-Bell-closed}
Consider all V60 star-overlap terms in which every root contains the
centre, the noncentral labels of distinct roots form a partition of
the \(n\) leaves, and every root is bounded by
\eqref{eq:newV61-reserve}.  Their complete positive majorant is at
most
\begin{equation}
 K^{n+1}\,
 \beta\left(\prod_{j=1}^nc_j\right)B^{n-1}.                  \label{eq:newV61-Bell-closed}
\end{equation}
\end{corollary}

\begin{proof}
Use \eqref{eq:newV61-reserve} on the root
\(\{\mathrm{centre}\}\cup C\), then apply
\Cref{lem:newV61-Bell-star}.
\end{proof}

\begin{assessment}[What happened to the Bell number]
\label{ass:newV61-Bell}
The unweighted number of star partitions is superexponential relative
to a fixed base.  A block with \(r\) leaves, however, carries
\(s^{(r-1)/2}\), while the thermal floor gives
\(n\sqrt s\le B\).  The exponential generating function sums these
two facts simultaneously.  The resulting \(n!\) is exactly absorbed
by the \(n\)-th power of \(B\), leaving the required
\(B^{n-1}\) after the nonempty-source factor is extracted.
\end{assessment}

\section{Exact abelian one-centre source quotient}
\label{sec:newV61-abelian}

The preceding combinatorial lemma still uses a positive majorant.
There is an exact resummation which displays the same mechanism
without naming the individual jump roots.

Let \(G\) be a standard real Gaussian, independent of
\(S_{\rm sm}\), and define the normalized tangent variable
\begin{equation}
 Z_{\rm sm}:=\sqrt{2S_{\rm sm}/s}\,G.                         \label{eq:newV61-Z}
\end{equation}
Its cumulant-generating function is
\begin{equation}
 {\cal K}(z)
 :=
 \log\mathbb E e^{zZ_{\rm sm}}
 =
 -\Phi_\alpha^{\rm sm}(-z^2/s).                              \label{eq:newV61-K}
\end{equation}

\begin{lemma}[Uniform analytic derivative disk]
\label{lem:newV61-K-disk}
There are numerical \(c_0,C>0\) such that
\({\cal K}\) is analytic on
\begin{equation}
 |z|\le {c_0\over\sqrt s}                                    \label{eq:newV61-K-disk}
\end{equation}
and
\begin{equation}
 |{\cal K}'(z)|\le C|z|                                      \label{eq:newV61-K-prime}
\end{equation}
there.
\end{lemma}

\begin{proof}
The Stieltjes measure of
\(\Phi_\alpha^{\rm sm}\) is supported in
\((T_\alpha,\infty)\), so
\(\Phi_\alpha^{\rm sm}(-w)\) is analytic for
\(|w|<T_\alpha\).  Since \(sT_\alpha\ge c/s\), this proves
analyticity on \eqref{eq:newV61-K-disk} after \(c_0\) is decreased.

The cumulant series and
\Cref{thm:newV59-small-CCH} give
\[
 {\cal K}(z)
 =
 {\mathbb ES_{\rm sm}\over s}z^2+
 \sum_{r\ge2}
 {\varkappa_r(S_{\rm sm})\over r!}
 {z^{2r}\over s^r}.
\]
Hence
\[
 |{\cal K}'(z)|
 \le
 C|z|+
 {1\over s|z|}
 \sum_{r\ge2}2r(Ks|z|^2)^r.
\]
Choose \(c_0\) so that \(Ks|z|^2\le1/4\).  The series is then at
most \(Cs^2|z|^4\), and the second term is at most
\(Cs|z|^3\le C|z|\).
\end{proof}

\begin{lemma}[One nonlinear character is an exact tilted ratio]
\label{lem:newV61-ratio}
For real \(\tau,t_1,\ldots,t_n\),
\begin{equation}
 \kappa_{n+1}
 \left(e^{i\tau Z_{\rm sm}},
 t_1Z_{\rm sm},\ldots,t_nZ_{\rm sm}\right)
 =
 \left(\prod_{j=1}^nt_j\right)
 \left.{d^n\over du^n}
 \exp\{{\cal K}(u+i\tau)-{\cal K}(u)\}
 \right|_{u=0}.                                               \label{eq:newV61-ratio}
\end{equation}
\end{lemma}

\begin{proof}
Introduce a source \(a\) for the first variable and one scalar source
\(u\) for \(Z_{\rm sm}\).  Differentiation in \(a\) at zero gives
\[
 \left.\partial_a
 \log\mathbb E
 e^{a e^{i\tau Z_{\rm sm}}+uZ_{\rm sm}}
 \right|_{a=0}
 =
 {\mathbb E e^{(u+i\tau)Z_{\rm sm}}
  \over
  \mathbb E e^{uZ_{\rm sm}}}
 =
 e^{{\cal K}(u+i\tau)-{\cal K}(u)}.
\]
Differentiate \(n\) times in independent sources
\(u_jt_j\) and set them to zero.  Since they enter through
\(u=\sum_ju_jt_j\), this proves
\eqref{eq:newV61-ratio}.
\end{proof}

\begin{theorem}[Exact abelian one-centre all-arity card]
\label{thm:newV61-one-centre}
Let
\[
 \sqrt s\le\beta,c_j<1,
 \qquad
 |\tau|\le C_0\beta,
 \qquad
 |t_j|\le C_0c_j,
 \qquad
 B:=\beta+\sum_{j=1}^nc_j.
\]
Then, after decreasing the fixed thermal constant \(C_0\) if
necessary,
\begin{equation}
 \boxed{\quad
 \left|
 \kappa_{n+1}
 \left(e^{i\tau Z_{\rm sm}},
 t_1Z_{\rm sm},\ldots,t_nZ_{\rm sm}\right)
 \right|
 \le
 K^{n+1}\,
 \beta\left(\prod_{j=1}^nc_j\right)B^{n-1}.
 \quad}                                                       \label{eq:newV61-one-centre}
\end{equation}
\end{theorem}

\begin{proof}
Let
\[
 R={\eta n\over B}
\]
with \(\eta>0\) sufficiently small.  The thermal floor gives
\(n\sqrt s\le B\), hence
\(R\le\eta/\sqrt s\).  Also \(B<n+1\), because all weights are
strictly below one, so \(R\ge\eta n/(n+1)\ge\eta/2\).
For \(|u|=R\), the segment from \(u\) to \(u+i\tau\) remains in
\eqref{eq:newV61-K-disk}.  By
\Cref{lem:newV61-K-disk},
\begin{equation}
 |{\cal K}(u+i\tau)-{\cal K}(u)|
 \le C|\tau|(|u|+|\tau|)
 \le C\beta R+C\beta^2
 \le C\eta n+C.                                               \label{eq:newV61-difference}
\end{equation}
Moreover
\[
 \left|
 e^{{\cal K}(u+i\tau)-{\cal K}(u)}-1
 \right|
 \le
 C\beta R e^{C\eta n}.
\]
Here the harmless \(\beta^2\) term is absorbed by \(\beta R\)
because \(\beta<1\) and \(R\ge\eta/2\).

For \(n\ge1\), subtracting one does not change the \(n\)-th
derivative.  Cauchy's estimate therefore yields
\[
 \left|
 {d^n\over du^n}
 e^{{\cal K}(u+i\tau)-{\cal K}(u)}
 \bigg|_{u=0}
 \right|
 \le
 C\beta n!R^{1-n}e^{C\eta n}
 \le
 K^n\beta B^{n-1}.
\]
The last inequality uses \(n!\le n^n\), exactly as in
\Cref{lem:newV61-Bell-star}.  Multiply by
\(\prod_j|t_j|\le C_0^n\prod_jc_j\) in
\eqref{eq:newV61-ratio}.
\end{proof}

\begin{corollary}[One nonlinear tangent centre is closed]
\label{cor:newV61-one-centre-closed}
In the commuting tangent sector, every connected row containing one
complete representation exponential and arbitrarily many linear
representation leaves has the strong thermal homogeneity with one
exponential arity base.
\end{corollary}

\begin{proof}
Simultaneously diagonalize the commuting generators.  On each joint
weight vector, their scalar coefficients satisfy the thermal bounds
in \Cref{thm:newV61-one-centre}.  Take the supremum over weights;
the estimate is dimension-free.
\end{proof}

\begin{assessment}[Exact meaning of the quotient]
\label{ass:newV61-quotient}
The ratio
\[
 {\mathbb E e^{(u+i\tau)Z_{\rm sm}}\over
  \mathbb E e^{uZ_{\rm sm}}}
\]
contains every small-clock jump that meets the nonlinear centre.
Taking this quotient before derivatives is the abelian logarithmic
counterpart of OCRC(i).  Its analytic radius is not a fixed
perturbative disk: it is \(c/\sqrt s\), exactly large enough because
the \(n\) thermal leaves force \(n\sqrt s\le B\).  This is the
upstream cancellation which the V60 fixed-star majorant could not
see.
\end{assessment}

\section{What remains after the one-centre theorem}
\label{sec:newV61-remaining}

\begin{definition}[Multi-centre word card, MCWC]
\label{def:newV61-MCWC}
MCWC is the noncommutative extension of
\Cref{thm:newV61-one-centre} in which:
\begin{enumerate}[label=\textup{(\roman*)},leftmargin=3.5em]
\item two or more labels retain complete matrix exponentials;
\item independent small-clock jumps generate time-ordered words at
      every such label;
\item words with disconnected label--jump incidence cancel before
      norms; and
\item the remaining connected words obey
      \(K^m(\prod_i b_i)B^{m-2}\).
\end{enumerate}
\end{definition}

\begin{proposition}[Reduction of the live small-clock gate]
\label{prop:newV61-reduction}
The tangent row, every one-root row, the complete Bell-star family,
and the exact abelian one-centre row are closed.  Consequently SCNGC
is reduced to MCWC for incidence structures with at least two
nonlinear centres.
\end{proposition}

\begin{proof}
Use
\Cref{cor:newV59-tangent,
thm:newV60-hyperedge-card,
cor:newV61-Bell-closed,
thm:newV61-one-centre,
cor:newV61-one-centre-closed}.
Every excluded structure has at least two labels at which distinct
jump roots overlap complete representation words; those are exactly
the structures retained in MCWC.
\end{proof}

\begin{programme}[V62: noncommutative incidence words]
\label{prog:newV61-V62}
\begin{enumerate}[leftmargin=3em]
\item Write the exact convolution source with one matrix centre as a
      left--right multiplication quotient and recover
      \eqref{eq:newV61-ratio} on commuting weights.
\item Apply ordered-simplex integration to two nonlinear centres and
      identify the first commutator word absent from the scalar
      quotient.
\item Use the Casimir stability identity to pay that commutator by a
      cross-centre pair mark.
\item Prove MCWC first for acyclic label--jump incidence graphs by
      iterated leaf-centre elimination.
\item Test the first incidence cycle separately; do not estimate it
      by the V60 Bell majorant.
\end{enumerate}
\end{programme}

\begin{firewall}[Claim boundary after V61]
\label{fire:newV61-boundary}
V61 retains the unrelaxed clock-factorial reserve, proves an
exponential generating-function bound for the complete Bell-star
partition family, derives the exact small-clock tilted source
quotient, proves its analytic disk and derivative bound, and closes
the exact abelian one-nonlinear-centre row at every arity.  It
corrects the interpretation of the V60 Bell no-go and reduces the
small-clock nonlinear problem to multi-centre noncommutative
incidence words.

V61 does not prove MCWC, OCRC in its full noncommutative form, SCNGC,
WPSC, GRLTC, LZTC, PLRC, WSDC, BRSC, HWSFC, UGCGC, thermal
connected WPRAC, all-arity RGSC, the raw Wilson aggregate strong
card, all-order strong MTA, WUFC, CPM, cutoff removal, reflection
positivity, Osterwalder--Schrader reconstruction, a continuum
Yang--Mills measure, or a positive Hamiltonian gap.  The full
Yang--Mills mass gap has not been proved.
\end{firewall}

\begin{theorem}[V61 result ledger]
\label{thm:newV61-results}
V61:
\begin{enumerate}[label=\textup{(V61.\arabic*)},leftmargin=6.5em]
\item records that fixed-tree OCRC is sufficient but not necessary;
\item proves the unrelaxed same-jump clock-factorial reserve;
\item proves the weighted Bell-star resummation;
\item closes the complete star-overlap positive majorant;
\item derives the exact abelian tilted source quotient;
\item proves its \(c/\sqrt s\) analytic derivative disk;
\item proves the exact abelian one-centre all-arity card;
\item lifts that card to commuting representation weights; and
\item isolates MCWC as the remaining small-clock nonlinear row.
\end{enumerate}
\end{theorem}

\begin{proof}
Use
\Cref{corr:newV61-OCRC,
lem:newV61-reserve,
lem:newV61-Bell-star,
cor:newV61-Bell-closed,
lem:newV61-K-disk,
lem:newV61-ratio,
thm:newV61-one-centre,
cor:newV61-one-centre-closed,
prop:newV61-reduction}.
\end{proof}

\paragraph{Parent-version record.}
V61 was formed from
\texttt{Renewal\_Yang\_Mills\_Mass\_Gap\_Research\_Notes\_V60.tex}.
The V60 parent had \(29\,241\) logical source lines,
\(931\,694\) bytes, and SHA--256
\[
 \texttt{14ec8cc3515d32c3e2a47d8e8236b5fb922b516ca482353b585c83a32d074500}.
\]
The next compulsory companion audit is V65.

% =============================================================================
% END APPENDED FIFTH-SERIES V61 MATERIAL
% =============================================================================

\clearpage
\chapter{Fifth-series V62: noncommutative Gaussian-centre
normal form and the marked-clock quotient}
\label{ch:newV62-centre}

\section{Scope of the first noncommutative test}
\label{sec:newV62-scope}

No companion audit is due at V62.  V61 closed one complete scalar
character against arbitrarily many collinear tangent leaves.  The
first noncommutative issue is already present when the centre is a
full matrix exponential
\[
 \exp\left(i\sqrt{2t}\sum_{a=1}^3G^aJ_0^a\right),
\]
because the three \(J_0^a\) need not commute.  This chapter treats
that centre without a Taylor truncation.  It then uses the exact law
of total cumulance to identify every additional clock connection.

This is a theorem for the subordinated tangent Gaussian model.  It is
not yet the exact compact-group heat endpoint; the latter remains
inside MCWC.

\section{A concentration ledger for the small clock}
\label{sec:newV62-clock}

Let \(\mu_\alpha=\mathbb ES_{\rm sm}\).

\begin{lemma}[Sub-gamma clock transform]
\label{lem:newV62-subgamma}
For \(0\le z<(2Ks^2)^{-1}\),
\begin{equation}
 \log\mathbb E
 e^{z(S_{\rm sm}-\mu_\alpha)}
 \le
 {K^2s^3z^2\over1-Ks^2z}.                                   \label{eq:newV62-subgamma}
\end{equation}
\end{lemma}

\begin{proof}
The centered logarithm has the absolutely convergent series
\[
 \sum_{r\ge2}{\varkappa_r(S_{\rm sm})\over r!}z^r.
\]
Every cumulant is nonnegative because \(S_{\rm sm}\) is a
subordinator.  Apply
\eqref{eq:newV59-small-CCH} and sum the geometric series:
\[
 \sum_{r\ge2}K^rs^{2r-1}z^r
 =
 {K^2s^3z^2\over1-Ks^2z}.
\]
\end{proof}

\begin{theorem}[All raw moments without a false factorial loss]
\label{thm:newV62-clock-moments}
For every integer \(p\ge1\),
\begin{equation}
 \left(\mathbb ES_{\rm sm}^p\right)^{1/p}
 \le
 Ks\left(1+\sqrt{ps}+ps\right)
 \le
 K s(1+p\sqrt s).                                            \label{eq:newV62-clock-moments}
\end{equation}
\end{theorem}

\begin{proof}
The standard Chernoff optimization applied to
\eqref{eq:newV62-subgamma} gives
\[
 \mathbb P\left\{
 S_{\rm sm}-\mu_\alpha
 \ge C\left(s^{3/2}\sqrt u+s^2u\right)
 \right\}
 \le e^{-u},
 \qquad u\ge0.
\]
Indeed, choose
\[
 z=
 {c\sqrt u\over s^{3/2}+s^2\sqrt u},
\]
which remains below \((2Ks^2)^{-1}\) after \(c\) is fixed
sufficiently small, and substitute in the exponential Markov bound.
Integrating the tail by
\[
 \mathbb E X_+^p
 =
 p\int_0^\infty y^{p-1}\mathbb P\{X>y\}\,dy
\]
gives
\[
 \|(S_{\rm sm}-\mu_\alpha)_+\|_{L^p}
 \le C\left(s^{3/2}\sqrt p+s^2p\right).
\]
The negative part is at most \(\mu_\alpha\le Ks\), since
\(S_{\rm sm}\ge0\).  The triangle inequality proves the first
bound.  Finally
\(\sqrt{ps}\le1+ps\le2(1+p\sqrt s)\) for \(s\le1\), after constants
are enlarged.
\end{proof}

\begin{warning}[Why the moment statement is formulated in root form]
\label{warn:newV62-moments}
The estimate is not
\(\mathbb ES_{\rm sm}^p\le p!(Cs)^p\).  The latter discards
concentration about the mean and recreates a factorial at
\(p\lesssim s^{-1/2}\), exactly where the total thermal weight can
still be of order one.  The root estimate
\eqref{eq:newV62-clock-moments} retains the two crossover scales
\(\sqrt{ps}\) and \(ps\).
\end{warning}

\section{One complete matrix centre at deterministic clock time}
\label{sec:newV62-deterministic}

Let \(G=(G^1,G^2,G^3)\) be standard Gaussian.  On mutually distinct
tensor factors define
\begin{align}
 F_{0,t}(G)
 &:=
 \exp\left(i\sqrt{2t}\sum_{a=1}^3G^aJ_0^a\right),             \label{eq:newV62-F}\\
 L_{j,t}(G)
 &:=
 i\sqrt{2t}\sum_{a=1}^3G^aJ_j^a,\qquad 1\le j\le n.          \label{eq:newV62-L}
\end{align}

\begin{lemma}[Gaussian centre identity]
\label{lem:newV62-Gaussian-centre}
If \(F:\mathbb R^3\to{\cal A}\) is smooth with Gaussian-integrable
derivatives and
\(\ell_j(G)=v_j\cdot G\), then
\begin{equation}
 \kappa_{n+1}^{\gamma_3}
 \left(F(G),\ell_1(G),\ldots,\ell_n(G)\right)
 =
 \mathbb E_{\gamma_3}
 D^nF(G)[v_1,\ldots,v_n].                                   \label{eq:newV62-Gaussian-centre}
\end{equation}
The identity remains valid for matrix- and tensor-valued \(F\).
\end{lemma}

\begin{proof}
Introduce a scalar source \(a\) for \(F\) and a vector source \(u\)
for \(G\).  Differentiation in \(a\) at zero gives
\[
 \left.\partial_a
 \log\mathbb E
 e^{aF(G)+u\cdot G}
 \right|_{a=0}
 =
 {\mathbb E[F(G)e^{u\cdot G}]
  \over\mathbb Ee^{u\cdot G}}
 =
 \mathbb EF(G+u)
\]
by the Cameron--Martin identity.  Directional differentiation in
\(v_1,\ldots,v_n\) at \(u=0\) proves
\eqref{eq:newV62-Gaussian-centre}.  The calculation is entrywise in
finite dimension and therefore applies in the tensor algebra.
\end{proof}

\begin{lemma}[Time-ordered matrix exponential derivative]
\label{lem:newV62-Duhamel-centre}
If \(X\) is skew-adjoint and \(H_1,\ldots,H_n\) are arbitrary
matrices, then
\begin{align}
 D^ne^X[H_1,\ldots,H_n]
 ={}&
 \sum_{\sigma\in{\mathfrak S}_n}
 \int_{\Delta_n}
 e^{t_0X}H_{\sigma(1)}e^{t_1X}\cdots
 H_{\sigma(n)}e^{t_nX}\,d\boldsymbol t,                      \label{eq:newV62-Duhamel}\\
 \left\|D^ne^X[H_1,\ldots,H_n]\right\|
 \le{}&\prod_{j=1}^n\|H_j\|,                                \label{eq:newV62-Duhamel-bound}
\end{align}
where
\(\Delta_n=\{t_j\ge0:\sum_{j=0}^nt_j=1\}\).
\end{lemma}

\begin{proof}
Iterate the first Fr\'echet derivative
\[
 De^X[H]=\int_0^1e^{(1-r)X}He^{rX}\,dr.
\]
This gives \eqref{eq:newV62-Duhamel}.  Every exponential factor is
unitary and
\({\rm vol}(\Delta_n)=1/n!\).  The \(n!\) permutations therefore
cancel the simplex volume in the norm bound, proving
\eqref{eq:newV62-Duhamel-bound}.  No commutation of the \(H_j\) is
used.
\end{proof}

\begin{theorem}[Noncommutative one-centre deterministic card]
\label{thm:newV62-deterministic-centre}
For every \(t>0\),
\begin{align}
 &\left\|
 \kappa_{n+1}^{\gamma_3}
 \left(F_{0,t},L_{1,t},\ldots,L_{n,t}\right)
 \right\|_{\rm cb}
 \notag\\
 &\qquad\le
 (Kt)^n C_2(R_0)^{n/2}
 \prod_{j=1}^n C_2(R_j)^{1/2}.                               \label{eq:newV62-deterministic-centre}
\end{align}
In particular, at \(t=s\),
\begin{equation}
 \left\|\kappa_{n+1}^{\gamma_3}
 (F_{0,s},L_{1,s},\ldots,L_{n,s})\right\|_{\rm cb}
 \le
 K^{n+1}b_0\left(\prod_{j=1}^nb_j\right)B^{n-1}.             \label{eq:newV62-deterministic-strong}
\end{equation}
\end{theorem}

\begin{proof}
Expand each \(L_{j,t}\) in its three coordinate components and
apply \Cref{lem:newV62-Gaussian-centre}.  Every derivative of
\(F_{0,t}\) has directions
\(i\sqrt{2t}J_0^{a_j}\).  By
\Cref{lem:newV62-Duhamel-centre},
\[
 \left\|
 \partial_{a_1}\cdots\partial_{a_n}F_{0,t}(G)
 \right\|
 \le
 (2t)^{n/2}\prod_{j=1}^n\|J_0^{a_j}\|.
\]
The \(n\) leaf coefficients contribute another
\((2t)^{n/2}\prod_j\|J_j^{a_j}\|\).  Sum the at most \(3^n\)
coordinate choices and use
\(\|J_i^a\|\le\sqrt{C_2(R_i)}\).  This proves
\eqref{eq:newV62-deterministic-centre}; amplification changes none
of these norms.

At \(t=s\), the right side is at most
\[
 K^n b_0^n\prod_{j=1}^nb_j
 \le
 K^n b_0\left(\prod_jb_j\right)B^{n-1},
\]
because \(b_0\le B\).
\end{proof}

\begin{assessment}[The first commutators are already included]
\label{ass:newV62-commutators}
For \(n=2\), the two terms in
\eqref{eq:newV62-Duhamel} contain both orders \(H_1H_2\) and
\(H_2H_1\); their difference contains \([H_1,H_2]\).
Thus \Cref{thm:newV62-deterministic-centre} is not a Cartan
calculation.  It controls every nested ordering at the centre by
the simplex volume before norms.  This is the exact Taylor-factorial
reserve anticipated in V60--V61.
\end{assessment}

\section{Averaging the one-block centre over the small clock}
\label{sec:newV62-average}

Let
\[
 {\cal C}_{n}(t)
 :=
 \kappa_{n+1}^{\gamma_3}
 \left(F_{0,t},L_{1,t},\ldots,L_{n,t}\right).
\]

\begin{theorem}[The zero-pair one-centre row]
\label{thm:newV62-zero-pair}
With
\[
 B=b_0+\sum_{j=1}^nb_j,
\]
one has
\begin{equation}
 \left\|\mathbb E\,{\cal C}_n(S_{\rm sm})\right\|_{\rm cb}
 \le
 K^{n+1}b_0\left(\prod_{j=1}^nb_j\right)B^{n-1}.             \label{eq:newV62-zero-pair}
\end{equation}
\end{theorem}

\begin{proof}
By \Cref{thm:newV62-deterministic-centre},
\[
 \|\mathbb E{\cal C}_n(S_{\rm sm})\|
 \le
 K^n\mathbb ES_{\rm sm}^n
 C_2(R_0)^{n/2}
 \prod_{j=1}^nC_2(R_j)^{1/2}.
\]
Use \(C_2(R_i)^{1/2}\le b_i/\sqrt s\) and
\Cref{thm:newV62-clock-moments}:
\begin{equation}
 \|\mathbb E{\cal C}_n(S_{\rm sm})\|
 \le
 K^n b_0^n(1+n\sqrt s)^n\prod_{j=1}^nb_j.                   \label{eq:newV62-zero-pair-mid}
\end{equation}
Put \(y=n\sqrt s\).  The thermal floor on the leaves gives
\(B\ge b_0+y\), while \(b_0<1\).  Therefore
\[
 b_0(1+y)\le b_0+y\le B
\]
and hence
\[
 b_0^{n-1}(1+y)^{n-1}\le B^{n-1}.
\]
The remaining factor \(1+y\le1+n\) is absorbed into \(K^n\).
Substitute in \eqref{eq:newV62-zero-pair-mid}.
\end{proof}

\begin{assessment}[Why raw clock moments do not lose the card here]
\label{ass:newV62-zero-pair}
When \(n\sqrt s\lesssim1\), clock concentration gives
\(\|S_{\rm sm}\|_{L^n}\asymp s\) with one exponential base, so no
factorial appears.  When \(n\sqrt s\gtrsim1\), the moment begins to
grow, but the \(n\) leaf floors have already made
\(B\gtrsim n\sqrt s\); that same \(B^{n-1}\) pays the growth.
The proof treats both regimes in the single inequality
\(b_0(1+n\sqrt s)\le B\).
\end{assessment}

\section{Exact total-cumulance normal form}
\label{sec:newV62-total}

\begin{theorem}[One-centre variance-mixture normal form]
\label{thm:newV62-normal-form}
Consider the subordinated tangent Gaussian with one complete centre
\(F_{0,S_{\rm sm}}\) and \(n\) linear leaves
\(L_{j,S_{\rm sm}}\).  Every term in its law-of-total-cumulance
expansion has the following form:
\begin{enumerate}[label=\textup{(\roman*)},leftmargin=3.5em]
\item choose a subset \(A\subseteq[n]\) to join the centre in one
      conditional Gaussian cumulant
      \({\cal C}_{A}(S_{\rm sm})\);
\item pair all leaves in \([n]\setminus A\); and
\item take one outer clock cumulant of
      \({\cal C}_{A}(S_{\rm sm})\) with one copy of
      \(2S_{\rm sm}\Omega_{ij}\) for every selected pair
      \(\{i,j\}\).
\end{enumerate}
No other conditional block survives.
\end{theorem}

\begin{proof}
Apply the multivariate law of total cumulance with conditioning
variable \(S_{\rm sm}\).  Conditional on \(S_{\rm sm}=t\), the
leaves are linear functions of a centered Gaussian.  A conditional
block containing only leaves therefore vanishes unless it has
exactly two members; in that case it equals
\[
 \mathbb E_{\gamma_3}L_{i,t}L_{j,t}
 =
 -2t\Omega_{ij},
\]
up to the harmless sign convention from the factors \(i\).
Exactly one conditional block contains the centre, and by
\Cref{lem:newV62-Gaussian-centre} it is
\({\cal C}_{A}(t)\).  The outer cumulant is taken over the resulting
functions of \(S_{\rm sm}\), proving the list.
\end{proof}

\begin{corollary}[The zero-pair term is closed]
\label{cor:newV62-zero-pair}
The term \(A=[n]\), for which there is no outer covariance pair, is
exactly the left side of
\eqref{eq:newV62-zero-pair} and obeys the strong thermal target.
\end{corollary}

\begin{proof}
For \(A=[n]\), the outer cumulant has one entry and is expectation.
Apply \Cref{thm:newV62-zero-pair}.
\end{proof}

\begin{definition}[Centre-marked clock quotient card, CMCQC]
\label{def:newV62-CMCQC}
For \(k,r\ge0\), CMCQC is a simultaneous bound on
\begin{equation}
 \kappa_{r+1}^{S_{\rm sm}}
 \left(
 {\cal C}_{k}(S_{\rm sm}),
 S_{\rm sm},\ldots,S_{\rm sm}
 \right)                                                     \label{eq:newV62-CMCQC}
\end{equation}
which, after the \(r\) Gaussian pair contractions are restored and
all choices of the \(k\) centre leaves and \(r\) pairs are summed,
is at most
\[
 K^{k+2r+1}
 b_0\left(\prod_{j=1}^{k+2r}b_j\right)
 B^{k+2r-1}.
\]
The estimate must use the normalized marked clock transform
\[
 {\mathbb E[{\cal C}_k(S_{\rm sm})e^{zS_{\rm sm}}]
  \over
  \mathbb Ee^{zS_{\rm sm}}}
\]
before differentiating \(r\) times.
\end{definition}

\begin{proposition}[CMCQC closes the full one-centre tangent row]
\label{prop:newV62-CMCQC}
If CMCQC holds for all \(k,r\), then every one-complete-centre
cumulant of the subordinated tangent Gaussian satisfies the strong
thermal all-arity card with one exponential base.
\end{proposition}

\begin{proof}
Use the exact normal form
\Cref{thm:newV62-normal-form}.  CMCQC includes the choice of the
centre subset, the Gaussian pairings of its complement, the outer
clock cumulant, and their common aggregate weight.  Summing its
bound gives the asserted full row.  The case \(r=0\) has already
been proved in \Cref{thm:newV62-zero-pair}.
\end{proof}

\begin{programme}[V63: marked clock quotient]
\label{prog:newV62-V63}
\begin{enumerate}[leftmargin=3em]
\item Differentiate
      \(\mathbb E[{\cal C}_k(S)e^{zS}]/\mathbb Ee^{zS}\)
      by one outer clock mark and prove the enhanced variance factor
      \(s^3\), uniformly in \(k\).
\item Use the Stieltjes support \(t>T_\alpha\) to obtain an analytic
      disk of radius \(c/s^2\) for the centered clock tilt.
\item Combine the disk with
      \(\|{\cal C}_k(t)\|\le(Kt)^k\) and optimize the Cauchy radius
      against \(k+2r\).
\item Sum the choices of \(A\) and the pairings before norms through
      the same exponential generating function, rather than
      assigning a factorial to each \(r\).
\item If CMCQC closes, return to the exact heat endpoint and use its
      V29 time-ordered forest in place of the Gaussian centre.
\end{enumerate}
\end{programme}

\begin{firewall}[Claim boundary after V62]
\label{fire:newV62-boundary}
V62 proves a sub-gamma transform and sharp root-moment ledger for the
small clock, proves the exact Gaussian-centre identity, controls all
noncommuting Duhamel orders of one complete matrix exponential
without factorial loss, and closes the averaged zero-pair
one-centre row.  It derives the exact law-of-total-cumulance normal
form and reduces the remaining one-centre variance-mixture terms to
the scalar centre-marked clock quotient CMCQC.

V62 does not prove CMCQC, the full one-centre subordinated tangent
card, MCWC, OCRC in its full noncommutative form, SCNGC, WPSC,
GRLTC, LZTC, PLRC, WSDC, BRSC, HWSFC, UGCGC, thermal connected
WPRAC, all-arity RGSC, the raw Wilson aggregate strong card,
all-order strong MTA, WUFC, CPM, cutoff removal, reflection
positivity, Osterwalder--Schrader reconstruction, a continuum
Yang--Mills measure, or a positive Hamiltonian gap.  The full
Yang--Mills mass gap has not been proved.
\end{firewall}

\begin{theorem}[V62 result ledger]
\label{thm:newV62-results}
V62:
\begin{enumerate}[label=\textup{(V62.\arabic*)},leftmargin=6.5em]
\item proves the centered sub-gamma clock transform;
\item proves the all-order clock root-moment ledger;
\item proves the exact Gaussian one-centre identity;
\item proves the factorial-free time-ordered Duhamel derivative;
\item proves the noncommutative deterministic one-centre card;
\item proves the averaged zero-pair one-centre card;
\item derives the exact total-cumulance normal form;
\item proves that only Gaussian pairs can accompany the centre
      block; and
\item isolates CMCQC as the remaining scalar clock quotient.
\end{enumerate}
\end{theorem}

\begin{proof}
Use
\Cref{lem:newV62-subgamma,
thm:newV62-clock-moments,
lem:newV62-Gaussian-centre,
lem:newV62-Duhamel-centre,
thm:newV62-deterministic-centre,
thm:newV62-zero-pair,
thm:newV62-normal-form,
cor:newV62-zero-pair,
prop:newV62-CMCQC}.
\end{proof}

\paragraph{Parent-version record.}
V62 was formed from
\texttt{Renewal\_Yang\_Mills\_Mass\_Gap\_Research\_Notes\_V61.tex}.
The V61 parent had \(29\,786\) logical source lines,
\(948\,302\) bytes, and SHA--256
\[
 \texttt{b7134497ae7876378c7b1c5f2acb35a2667642a228ec925cd7c40ca70fe00bb2}.
\]
The next compulsory companion audit is V65.

% =============================================================================
% END APPENDED FIFTH-SERIES V62 MATERIAL
% =============================================================================

\clearpage
\chapter{Fifth-series V63: the Poisson covariance identity and
closure of the first marked-clock row}
\label{ch:newV63-covariance}

\section{Entry point}
\label{sec:newV63-entry}

No companion audit is due at V63.  V62 reduced the one-complete-centre
variance-mixture problem to CMCQC.  The first nontrivial row is one
outer clock mark, \(r=1\).  Estimating it by a generic complex Cauchy
bound sees only a sub-gamma scale \(s^{3/2}\).  The exact Poisson
covariance identity instead differentiates the centre once and
integrates \(x^2\) against the L\'evy measure, producing the required
\(s^3\) variance scale.

\section{Differentiating the complete Gaussian centre in clock time}
\label{sec:newV63-centre-derivative}

For \(k\ge1\), retain the V62 notation
\[
 {\cal C}_k(t)
 =
 \kappa_{k+1}^{\gamma_3}
 (F_{0,t},L_{1,t},\ldots,L_{k,t})
\]
and put
\[
 A_k:=
 C_2(R_0)^{k/2}
 \prod_{j=1}^kC_2(R_j)^{1/2}.
\]

\begin{lemma}[Clock derivative of the complete centre block]
\label{lem:newV63-centre-derivative}
For \(t>0\),
\begin{align}
 \|{\cal C}_k(t)\|_{\rm cb}
 &\le (Kt)^kA_k,                                             \label{eq:newV63-centre-size}\\
 \|{\cal C}_k'(t)\|_{\rm cb}
 &\le
 K^kA_k
 \left\{
 kt^{k-1}+\sqrt{C_2(R_0)}\,t^{k-1/2}
 \right\}.                                                   \label{eq:newV63-centre-derivative}
\end{align}
\end{lemma}

\begin{proof}
The first line is
\Cref{thm:newV62-deterministic-centre}.  For the derivative, use
\Cref{lem:newV62-Gaussian-centre} and the time-ordered formula
\eqref{eq:newV62-Duhamel}.  The \(k\) derivatives of the centre and
the \(k\) leaf coefficients give an explicit factor \(t^k\).
Differentiating it costs \(k/t\).

The remaining clock dependence is through
\[
 X_t(G)=i\sqrt{2t}\sum_aG^aJ_0^a.
\]
Its derivative obeys
\[
 \|X_t'(G)\|
 \le
 {C|G|\sqrt{C_2(R_0)}\over\sqrt t}.
\]
One further Fr\'echet derivative of \(e^{X_t}\), inserted anywhere
in the ordered simplex, has norm at most this quantity times the
product of the original \(k\) insertion norms.  The sum of all
orders is again paid by the simplex volume, as in
\Cref{lem:newV62-Duhamel-centre}.  Since
\(\mathbb E|G|\le C\), this contribution is the second term of
\eqref{eq:newV63-centre-derivative}.  All estimates are unchanged
by amplification.
\end{proof}

\section{The exact covariance identity}
\label{sec:newV63-Mecke}

\begin{lemma}[Poisson covariance formula for the small clock]
\label{lem:newV63-Mecke}
Let \(f:[0,\infty)\to{\cal A}\) be continuously differentiable and
of polynomial growth.  Then
\begin{equation}
 \operatorname{Cov}
 \left(f(S_{\rm sm}),S_{\rm sm}\right)
 =
 \int_0^\infty
 x\,\mathbb E
 \left[f(S_{\rm sm}+x)-f(S_{\rm sm})\right]
 \Pi_\alpha^{\rm sm}(dx).                                   \label{eq:newV63-Mecke}
\end{equation}
The equality is a Bochner identity.
\end{lemma}

\begin{proof}
First truncate the L\'evy measure to
\(\varepsilon\le x\le\varepsilon^{-1}\).  The jump part is then
\[
 J_\varepsilon=\int x\,N_\varepsilon(dx)
\]
for a finite Poisson random measure, plus an independent remainder
and the deterministic drift.  The Mecke identity gives
\[
 \mathbb E[f(S)J_\varepsilon]
 -
 \mathbb Ef(S)\,\mathbb EJ_\varepsilon
 =
 \int_{\varepsilon}^{\varepsilon^{-1}}
 x\,\mathbb E[f(S+x)-f(S)]\Pi(dx).
\]
The deterministic drift has zero covariance.  Let
\(\varepsilon\downarrow0\).  Polynomial growth, the mean-value
theorem, \Cref{lem:newV60-jump-moments}, and
\Cref{thm:newV62-clock-moments} give uniform integrability on both
sides.  Dominated convergence proves
\eqref{eq:newV63-Mecke}.
\end{proof}

\begin{theorem}[One marked centre-block covariance]
\label{thm:newV63-centre-covariance}
For every \(k\ge1\),
\begin{align}
 \left\|
 \operatorname{Cov}
 \left({\cal C}_k(S_{\rm sm}),S_{\rm sm}\right)
 \right\|_{\rm cb}
 \le{}&
 K^kA_k\left\{
 ks^{k+2}(1+k\sqrt s)^{k-1}\right.
 \notag\\
 &\left.\qquad
 +\sqrt{C_2(R_0)}\,s^{k+5/2}
  (1+k\sqrt s)^{k-1/2}
 \right\}
 \notag\\
 &+
 (Kk)^{k+3}A_k\left\{
 s^{2k+1}+
 \sqrt{C_2(R_0)}\,s^{2k+2}
 \right\}.                                                   \label{eq:newV63-centre-covariance}
\end{align}
\end{theorem}

\begin{proof}
Apply \Cref{lem:newV63-Mecke} and the mean-value theorem to
\({\cal C}_k\).  By
\eqref{eq:newV63-centre-derivative},
\begin{align}
 \|\operatorname{Cov}({\cal C}_k(S),S)\|
 \le{}&
 K^kA_k\int_0^\infty x^2
 \mathbb E\left[
 k(S+x)^{k-1}\right.
 \notag\\
 &\left.\hspace{7em}
 +\sqrt{C_2(R_0)}(S+x)^{k-1/2}
 \right]\Pi(dx).                                             \label{eq:newV63-cov-bound}
\end{align}
For \(q\ge0\),
\[
 (S+x)^q\le2^{q}(S^q+x^q).
\]
The term with \(S^{k-1}\) is bounded, using independence of the
integration variable, by
\[
 K^k k A_k
 \left\{
 \varkappa_2(S)\,\mathbb ES^{k-1}
 +
 \int x^{k+1}\Pi(dx)
 \right\}.
\]
\Cref{thm:newV59-small-CCH,
thm:newV62-clock-moments,
lem:newV60-jump-moments}
give for the typical part
\[
 Ks^3\{Ks(1+k\sqrt s)\}^{k-1},
\]
whereas the pure-jump part is at most
\[
 \Gamma(k+2)K^{2k+2}s^{2k+1},
\]
with \(\Gamma(k+2)\le(Kk)^{k+2}\).

For the second line of \eqref{eq:newV63-cov-bound}, use
\(\sqrt{C_2(R_0)}\le b_0/\sqrt s\le s^{-1/2}\).
The factor with \(S^{k-1/2}\) is at most
\[
 Ks^3\{Ks(1+k\sqrt s)\}^{k-1/2}s^{-1/2},
\]
and the pure-jump term uses
\[
 \int x^{k+3/2}\Pi(dx)
 \le
 \Gamma(k+5/2)K^{2k+3}s^{2k+2}.
\]
These two terms are bounded by
\[
 K^ks^{k+5/2}(1+k\sqrt s)^{k-1/2}
 +(Kk)^{k+3}s^{2k+2},
\]
respectively.
Substitution in \eqref{eq:newV63-cov-bound} proves the theorem.
\end{proof}

\begin{assessment}[The enhanced power is exact]
\label{ass:newV63-enhanced}
The leading term in
\eqref{eq:newV63-centre-covariance} is
\[
 k\,s^{k-1}\varkappa_2(S_{\rm sm})
 =
 O(ks^{k+2}),
\]
not \(O(s^{k+3/2})\).  The factor \(s^3\) is the exact small-clock
variance from V57--V59.  The high-jump terms are smaller because the
retained L\'evy spectrum begins at \(t\asymp s^{-2}\).
\end{assessment}

\section{Closure of CMCQC at one outer pair}
\label{sec:newV63-r1}

\begin{theorem}[First marked-clock quotient row]
\label{thm:newV63-CMCQC-r1}
CMCQC holds for \(r=1\) and every \(k\ge1\).  More explicitly,
after adjoining two further leaves \(i,j\), contracting them into
\(2S_{\rm sm}\Omega_{ij}\), and summing all choices of the two
paired leaves, the resulting contribution is bounded by
\begin{equation}
 K^{k+3}
 b_0\left(\prod_{\ell=1}^{k+2}b_\ell\right)
 B^{k+1}.                                                    \label{eq:newV63-CMCQC-r1}
\end{equation}
\end{theorem}

\begin{proof}
For one fixed choice of the pair,
\(\|\Omega_{ij}\|\le\sqrt{C_2(R_i)C_2(R_j)}\).
Multiply \eqref{eq:newV63-centre-covariance} by this bound and
convert every Casimir factor by
\(\sqrt{C_2(R_\ell)}\le b_\ell/\sqrt s\).  The two typical-clock
terms are at most
\begin{align}
 &K^kks\,b_0^k(1+k\sqrt s)^{k-1}
 \prod_{\ell=1}^{k+2}b_\ell,                                \label{eq:newV63-fixed-pair}\\
 &K^ks\,b_0^{k+1}(1+k\sqrt s)^{k-1/2}
 \prod_{\ell=1}^{k+2}b_\ell.                                \label{eq:newV63-fixed-pair-centre}
\end{align}
The two high-jump terms are bounded by
\begin{align}
 &(Kk)^{k+3}s^k b_0^k
 \prod_{\ell=1}^{k+2}b_\ell,                                \label{eq:newV63-fixed-pair-high1}\\
 &(Kk)^{k+3}s^{k+1/2}b_0^{k+1}
 \prod_{\ell=1}^{k+2}b_\ell.                                \label{eq:newV63-fixed-pair-high2}
\end{align}
Let \(y=k\sqrt s\).  Since \(b_0<1\) and
\(B\ge b_0+y+2\sqrt s\),
\[
 b_0^{k-1}(1+y)^{k-1}\le B^{k-1}.
\]
Moreover \(ks\le B^2/k\le B^2\).  These inequalities put
\eqref{eq:newV63-fixed-pair} below the target.  For
\eqref{eq:newV63-fixed-pair-centre}, use
\[
 b_0^k(1+y)^{k-1/2}\le B^k,
 \qquad s\le B,
\]
which again gives the target.

For the first high-jump term, \(ks\le B^2/k\) and
\(B<k+3\) give
\[
 k^{k+3}s^kb_0^{k-1}
 \le
 B^{2k}k^{3-k}
 \le K^kB^{k+1}.
\]
The last inequality follows from
\((k+3)^{k-1}k^{3-k}\le K^kk^2\).
The second high-jump term is treated identically:
\[
 k^{k+3}s^{k+1/2}b_0^k
 \le
 B^{2k}k^{3-k}\sqrt s
 \le K^kB^{k+1}.
\]
The polynomial number
\(\binom{k+2}{2}\) of choices of the paired leaves is absorbed into
\(K_1^{k+3}\).  This proves \eqref{eq:newV63-CMCQC-r1}.
\end{proof}

\begin{corollary}[One-centre tangent card through one clock
connection]
\label{cor:newV63-one-centre-r1}
In the normal form of
\Cref{thm:newV62-normal-form}, the zero-pair row and the complete
one-pair row obey the strong all-arity thermal estimate with one
exponential base.
\end{corollary}

\begin{proof}
Use
\Cref{thm:newV62-zero-pair,thm:newV63-CMCQC-r1}.
\end{proof}

\section{Iteration identity and the remaining row}
\label{sec:newV63-iteration}

\begin{lemma}[First difference operator behind higher marks]
\label{lem:newV63-difference}
Define
\[
 (\Delta_xf)(t):=f(t+x)-f(t).
\]
Then the first marked row is
\[
 \operatorname{Cov}(f(S),S)
 =
 \int x\,\mathbb E(\Delta_xf)(S)\Pi(dx).
\]
Every further outer clock mark can be generated by applying the
same Poisson add-one operation to the normalized marked transform;
each genuinely new connection inserts one additional difference
\(\Delta_x\) and one additional factor \(x\).
\end{lemma}

\begin{proof}
The first statement is
\Cref{lem:newV63-Mecke}.  Apply the multivariate Mecke identity to
the Poisson representation of the jump part, with independent
marked copies of the add-one operator.  Centering removes terms in
which a new mark forms a disconnected component.  Thus a connected
new mark acts by one finite difference and carries the added jump
size.  This is an identity first for finite L\'evy truncations and
then in the limit by the argument of
\Cref{lem:newV63-Mecke}.
\end{proof}

\begin{definition}[Iterated difference clock card, IDCC]
\label{def:newV63-IDCC}
IDCC is a connected partition formula for
\eqref{eq:newV62-CMCQC} in which \(r\) outer marks are represented
by iterated add-one differences, and which proves
\begin{equation}
 \left\|
 \kappa_{r+1}
 ({\cal C}_k(S),S,\ldots,S)
 \right\|_{\rm cb}
 \le
 K^{k+r}(k+2r)^{2r}
 s^{k+2r}(1+k\sqrt s)^k.                                    \label{eq:newV63-IDCC}
\end{equation}
\end{definition}

\begin{proposition}[IDCC implies CMCQC]
\label{prop:newV63-IDCC}
IDCC, summed over the choices and pairings in
\Cref{thm:newV62-normal-form}, implies CMCQC for every \(k,r\).
\end{proposition}

\begin{proof}
There are at most
\[
 { (k+2r)!\over k!\,2^rr!}
 \le (k+2r)^{2r}
\]
ways to choose and pair the \(2r\) leaves outside the centre block.
After conversion of the \(k+r\) clock powers and all Casimir factors,
\eqref{eq:newV63-IDCC} leaves \(s^r\).
The thermal floor gives
\[
 s^r\le\left({B\over k+2r}\right)^{2r}.
\]
This cancels the pairing and IDCC polynomial powers.  The remaining
\(b_0^{k-1}(1+k\sqrt s)^k\) is paid by \(B^{k-1}\) as in
\Cref{thm:newV63-CMCQC-r1}; numerical and residual polynomial
factors enter one exponential base.
\end{proof}

\begin{programme}[V64: iterated add-one differences]
\label{prog:newV63-V64}
\begin{enumerate}[leftmargin=3em]
\item Prove the full multivariate Mecke--cumulant formula for
      \(r\) copies of the clock, retaining only partitions connected
      to the centre observable.
\item Bound
      \(\Delta_{x_1}\cdots\Delta_{x_q}{\cal C}_k(t)\)
      by the time-ordered derivatives of
      \Cref{lem:newV63-centre-derivative}.
\item Integrate the resulting monomials with
      \Cref{lem:newV60-jump-moments}; every connected outer mark
      must pay two powers of \(s\).
\item Establish \eqref{eq:newV63-IDCC} and close CMCQC.
\item Prepare the compulsory V65 companion audit independently of
      the V64 outcome.
\end{enumerate}
\end{programme}

\begin{firewall}[Claim boundary after V63]
\label{fire:newV63-boundary}
V63 proves the clock derivative bound for a complete
noncommutative Gaussian centre, derives the exact Poisson covariance
identity, proves the enhanced \(s^{k+2}\) centre covariance, and
closes CMCQC for one outer clock mark at every centre arity.  It
identifies iterated add-one differences as the correct route to all
outer marks and formulates IDCC.

V63 does not prove IDCC, CMCQC for \(r\ge2\), the full one-centre
subordinated tangent card, MCWC, OCRC in its full noncommutative
form, SCNGC, WPSC, GRLTC, LZTC, PLRC, WSDC, BRSC, HWSFC, UGCGC,
thermal connected WPRAC, all-arity RGSC, the raw Wilson aggregate
strong card, all-order strong MTA, WUFC, CPM, cutoff removal,
reflection positivity, Osterwalder--Schrader reconstruction, a
continuum Yang--Mills measure, or a positive Hamiltonian gap.  The
full Yang--Mills mass gap has not been proved.
\end{firewall}

\begin{theorem}[V63 result ledger]
\label{thm:newV63-results}
V63:
\begin{enumerate}[label=\textup{(V63.\arabic*)},leftmargin=6.5em]
\item proves the clock derivative of a complete Gaussian centre;
\item proves the exact small-clock Poisson covariance identity;
\item proves the enhanced one-mark centre covariance;
\item retains both typical-clock and high-jump regimes;
\item closes CMCQC for \(r=1\) at every \(k\);
\item closes the zero- and one-pair rows of the one-centre normal
      form;
\item identifies the iterated add-one operation for higher marks;
      and
\item formulates IDCC as the all-mark acquisition target.
\end{enumerate}
\end{theorem}

\begin{proof}
Use
\Cref{lem:newV63-centre-derivative,
lem:newV63-Mecke,
thm:newV63-centre-covariance,
thm:newV63-CMCQC-r1,
cor:newV63-one-centre-r1,
lem:newV63-difference,
prop:newV63-IDCC}.
\end{proof}

\paragraph{Parent-version record.}
V63 was formed from
\texttt{Renewal\_Yang\_Mills\_Mass\_Gap\_Research\_Notes\_V62.tex}.
The V62 parent had \(30\,285\) logical source lines,
\(964\,238\) bytes, and SHA--256
\[
 \texttt{f8a9f651e7021759adc3429f0fcaddf1e00266cb05695bf26b21b61f72ae3a73}.
\]
The next compulsory companion audit is V65.

% =============================================================================
% END APPENDED FIFTH-SERIES V63 MATERIAL
% =============================================================================

\clearpage
\chapter{Fifth-series V64: the multivariate Mecke--cumulant
formula and the all-marks-on-one-jump sector}
\label{ch:newV64-Mecke}

\section{The exact iterated formula}
\label{sec:newV64-formula}

No companion audit is due at V64.  The V63 Poisson covariance
identity is the first member of a complete partition formula.  Write
\[
 \Delta_xf(t):=f(t+x)-f(t)
\]
and let \({\rm Part}([r])\) denote the set partitions of
\([r]\).

\begin{theorem}[Multivariate Mecke--cumulant formula]
\label{thm:newV64-Mecke-cumulant}
Let \(f:[0,\infty)\to{\cal A}\) have polynomial growth and enough
continuous derivatives for the expressions below.  For every
\(r\ge1\),
\begin{align}
 &\kappa_{r+1}
 \left(f(S_{\rm sm}),S_{\rm sm},\ldots,S_{\rm sm}\right)
 \notag\\
 &\quad=
 \sum_{\pi\in{\rm Part}([r])}
 \int_{(0,\infty)^{|\pi|}}
 \mathbb E\left[
 \prod_{B\in\pi}^{\longrightarrow}\Delta_{x_B}f(S_{\rm sm})
 \right]
 \prod_{B\in\pi}x_B^{|B|}
 \Pi_\alpha^{\rm sm}(dx_B),                                 \label{eq:newV64-Mecke-cumulant}
\end{align}
where the arrow means composition of the commuting difference
operators, not multiplication of their values:
\[
 \prod_{B\in\pi}^{\longrightarrow}\Delta_{x_B}f
 :=
 \Delta_{x_{B_1}}\cdots\Delta_{x_{B_{|\pi|}}}f.
\]
The formula holds as a Bochner identity.
\end{theorem}

\begin{proof}
First truncate the L\'evy measure to a finite measure with compact
support.  Put
\[
 M(z)=\mathbb Ee^{zS},
 \qquad
 F_f(z)={\mathbb E[f(S)e^{zS}]\over M(z)}.
\]
By the defining source identity,
\[
 F_f^{(r)}(0)
 =
 \kappa_{r+1}(f(S),S,\ldots,S).                              \label{eq:newV64-source-derivative}
\]
Under exponential tilting by \(z\), the Poisson jump measure has
intensity \(e^{zx}\Pi(dx)\).  Differentiating an expectation under
that tilted law has two effects:
\begin{enumerate}[label=\textup{(\roman*)},leftmargin=3.5em]
\item differentiating an existing factor \(e^{zx_B}\) multiplies it
      by \(x_B\), adjoining the new derivative label to the existing
      block \(B\);
\item differentiating the tilted Poisson expectation creates a new
      jump \(x\), a new difference \(\Delta_x\), and a singleton
      block.
\end{enumerate}
Starting from \(F_f(0)=\mathbb Ef(S)\), these are exactly the two
operations which recursively construct a set partition when the
new label \(r\) is inserted.  Induction proves
\eqref{eq:newV64-Mecke-cumulant} for the truncation.

For the full small-clock measure, truncate at
\(\varepsilon\le x\le\varepsilon^{-1}\).
The finite differences have polynomial growth, while all positive
jump moments are finite by
\Cref{lem:newV60-jump-moments} and all moments of \(S_{\rm sm}\)
are finite by \Cref{thm:newV62-clock-moments}.  Dominated
convergence therefore removes the truncation.
\end{proof}

\begin{corollary}[The first two rows]
\label{cor:newV64-first-rows}
For \(r=1\), \eqref{eq:newV64-Mecke-cumulant} is
\eqref{eq:newV63-Mecke}.  For \(r=2\),
\begin{align}
 \kappa_3(f(S),S,S)
 ={}&
 \int x^2\mathbb E\Delta_xf(S)\Pi(dx)
 \notag\\
 &+
 \iint xy\,
 \mathbb E\Delta_x\Delta_yf(S)\Pi(dx)\Pi(dy).                \label{eq:newV64-r2}
\end{align}
\end{corollary}

\begin{proof}
List the partitions of one and two labels.
\end{proof}

\begin{assessment}[Connected meaning of the partition blocks]
\label{ass:newV64-blocks}
A block \(B\) in \eqref{eq:newV64-Mecke-cumulant} means that all
outer clock marks in \(B\) land on the same L\'evy jump.  Distinct
blocks mean distinct added jumps, and the iterated difference is
the exact operation which keeps all those jumps connected to the
centre observable.  Thus the formula contains no disconnected
clock partition to be removed later.
\end{assessment}

\section{The diagonal outer partition}
\label{sec:newV64-diagonal}

For \(k,r\ge1\), define the all-marks-on-one-jump term
\begin{equation}
 {\cal D}_{k,r}
 :=
 \int_0^\infty
 x^r\,\mathbb E
 \left[
 {\cal C}_k(S_{\rm sm}+x)-{\cal C}_k(S_{\rm sm})
 \right]\Pi_\alpha^{\rm sm}(dx).                             \label{eq:newV64-Dkr}
\end{equation}

\begin{theorem}[Diagonal marked-clock bound]
\label{thm:newV64-diagonal-bound}
There is a numerical \(K\) such that
\begin{align}
 \|{\cal D}_{k,r}\|_{\rm cb}
 \le{}&
 K^{k+r}A_k(r+1)!\left\{
 ks^{k+2r}(1+k\sqrt s)^{k-1}\right.
 \notag\\
 &\left.\qquad
 +\sqrt{C_2(R_0)}\,s^{k+2r+1/2}
 (1+k\sqrt s)^{k-1/2}
 \right\}
 \notag\\
 &+
 \{K(k+r)\}^{k+r+3}A_k
 \left\{
 s^{2k+2r-1}
 +\sqrt{C_2(R_0)}\,s^{2k+2r}
 \right\}.                                                   \label{eq:newV64-diagonal-bound}
\end{align}
\end{theorem}

\begin{proof}
Apply the mean-value theorem and
\Cref{lem:newV63-centre-derivative}:
\begin{align*}
 \|{\cal D}_{k,r}\|
 \le
 K^kA_k\int x^{r+1}\mathbb E\left[
 k(S+x)^{k-1}
 +\sqrt{C_2(R_0)}(S+x)^{k-1/2}
 \right]\Pi(dx).
\end{align*}
Use
\((S+x)^q\le2^q(S^q+x^q)\).  The two products containing \(S\)
are bounded by
\begin{align*}
 &k\left(\int x^{r+1}\Pi(dx)\right)\mathbb ES^{k-1},\\
 &\sqrt{C_2(R_0)}
 \left(\int x^{r+1}\Pi(dx)\right)\mathbb ES^{k-1/2}.
\end{align*}
By
\Cref{lem:newV60-jump-moments,thm:newV62-clock-moments},
these give the first two terms of
\eqref{eq:newV64-diagonal-bound}.

The two pure-jump terms use
\begin{align*}
 \int x^{k+r}\Pi(dx)
 &\le
 \Gamma(k+r+1)K^{2k+2r}s^{2k+2r-1},\\
 \int x^{k+r+1/2}\Pi(dx)
 &\le
 \Gamma(k+r+3/2)K^{2k+2r+1}s^{2k+2r}.
\end{align*}
The elementary bounds
\(\Gamma(k+r+1),\Gamma(k+r+3/2)
\le\{K(k+r)\}^{k+r+2}\)
give the last line of
\eqref{eq:newV64-diagonal-bound}.
\end{proof}

\section{Summation over leaves and Gaussian pairs}
\label{sec:newV64-pair-sum}

\begin{theorem}[The complete diagonal CMCQC sector]
\label{thm:newV64-diagonal-CMCQC}
Fix \(k,r\ge1\) and \(N=k+2r\).  Choose \(k\) of the \(N\) leaves
to enter the centre block, pair the other \(2r\) leaves, multiply
\({\cal D}_{k,r}\) by the corresponding \(r\) operators
\(-2\Omega_{ij}\), and sum all choices and pairings.  The resulting
operator has norm at most
\begin{equation}
 K^{N+1}
 b_0\left(\prod_{\ell=1}^{N}b_\ell\right)B^{N-1}.            \label{eq:newV64-diagonal-CMCQC}
\end{equation}
\end{theorem}

\begin{proof}
For a fixed choice and pairing,
\[
 \prod_{\{i,j\}}\|\Omega_{ij}\|
 \le
 \prod_{\ell\notin A}\sqrt{C_2(R_\ell)}.
\]
Together with \(A_k\), there are \(k+r\) inverse powers of \(s\)
when all Casimir square roots are converted to thermal marks.

The first typical-clock term in
\eqref{eq:newV64-diagonal-bound} therefore becomes
\[
 K^{N}k(r+1)!\,
 s^r b_0^k(1+k\sqrt s)^{k-1}
 \prod_{\ell=1}^{N}b_\ell.
\]
The number of choices and pairings is
\[
 {N!\over k!\,2^rr!}
 \le {N^{2r}\over2^rr!}.
\]
Since \((r+1)!/r!=r+1\),
\[
 N^{2r}s^r\le B^{2r}
\]
by the thermal floor \(N\sqrt s\le B\), while
\[
 b_0^{k-1}(1+k\sqrt s)^{k-1}\le B^{k-1}.
\]
All residual polynomial factors are absorbed into \(K^N\).  The
second typical-clock term has one extra \(b_0\) and the same \(s^r\);
it is smaller by the identical argument.

For the first pure-jump term, after Casimir conversion the remaining
clock power is \(s^{k+r-1}\).  Use
\(\sqrt s\le B/N\) and \(B<N+1\):
\[
 s^{k+r-1}
 \le
 {B^{2k+2r-2}\over N^{2k+2r-2}}.
\]
Multiplying the choice bound and
\((k+r)^{k+r+3}\), then dividing by the target power
\(B^{N-1}=B^{k+2r-1}\), leaves at most
\[
 K^N\,{N^r\over r!},
\]
after harmless polynomial factors and \(B^{k-1}\le(N+1)^{k-1}\)
are included.  Since \(r\le N/2\),
\[
 {N^r\over r!}\le K^N
\]
by \(r!\ge(r/e)^r\).  The second pure-jump term has an additional
\(\sqrt s\,b_0\) and is no larger.  This proves
\eqref{eq:newV64-diagonal-CMCQC}.
\end{proof}

\begin{corollary}[All outer marks on one jump are closed]
\label{cor:newV64-diagonal-closed}
CMCQC holds at every \(k,r\ge1\) for the partition
\(\pi=\{[r]\}\) in
\eqref{eq:newV64-Mecke-cumulant}.
\end{corollary}

\begin{proof}
That partition is exactly \({\cal D}_{k,r}\).  Apply
\Cref{thm:newV64-diagonal-CMCQC}.
\end{proof}

\section{The remaining off-diagonal difference card}
\label{sec:newV64-off-diagonal}

\begin{definition}[Off-diagonal iterated-difference card, OIDC]
\label{def:newV64-OIDC}
OIDC is the estimate, for every partition
\(\pi\in{\rm Part}([r])\) with \(q=|\pi|\ge2\), of
\begin{equation}
 \int
 \left\|
 \mathbb E
 \Delta_{x_1}\cdots\Delta_{x_q}
 {\cal C}_k(S_{\rm sm})
 \right\|_{\rm cb}
 \prod_{\nu=1}^q x_\nu^{|B_\nu|}
 \Pi(dx_\nu),                                                \label{eq:newV64-OIDC}
\end{equation}
such that the sum over \(\pi\), centre leaves, and Gaussian pairings
obeys the strong thermal target with one exponential base.
\end{definition}

\begin{proposition}[OIDC completes IDCC]
\label{prop:newV64-OIDC}
The diagonal theorem
\Cref{thm:newV64-diagonal-CMCQC} together with OIDC implies IDCC
and hence CMCQC.
\end{proposition}

\begin{proof}
\Cref{thm:newV64-Mecke-cumulant} is the exact sum over all
partitions of the \(r\) clock marks.  The one-block partition is
closed by \Cref{thm:newV64-diagonal-CMCQC}; OIDC closes every
partition with at least two blocks.  Their sum is IDCC.  Apply
\Cref{prop:newV63-IDCC}.
\end{proof}

\begin{assessment}[The bottleneck has moved to distinct jumps]
\label{ass:newV64-bottleneck}
No large moment of a single jump remains uncontrolled: the
all-marks-on-one-jump sector is closed at arbitrary \(k,r\), including
the high-jump tail.  The live issue is the simultaneous finite
difference of the complete matrix centre under two or more
distinct added jumps.  A derivative-by-derivative estimate would
again lose cancellations at \(S=0\); OIDC must retain the complete
finite difference or use the analytic \(t^kH_k(t)\) factorization
only on the small-\(tC_2(R_0)\) region.
\end{assessment}

\begin{programme}[V65: compulsory audit and off-diagonal
factorization]
\label{prog:newV64-V65}
\begin{enumerate}[leftmargin=3em]
\item Perform the compulsory newest-SM/newest-NS audit before
      continuing OIDC.
\item Prove the exact factorization
      \({\cal C}_k(t)=t^kH_k(t)\) and a coefficient bound for \(H_k\)
      on \(tC_2(R_0)\le1\).
\item On the complementary region, keep
      \(\Delta_{x_1}\cdots\Delta_{x_q}\) assembled and use unitarity
      rather than the absolute centre power series.
\item Show that each distinct added jump pays one factor
      \(\int x^{|B|+1}\Pi(dx)\), hence two clock powers per outer
      connection.
\item Sum the set partitions by an exponential generating function
      before taking the final operator norm.
\end{enumerate}
\end{programme}

\begin{firewall}[Claim boundary after V64]
\label{fire:newV64-boundary}
V64 proves the exact multivariate Mecke--cumulant partition formula,
identifies its connected jump meaning, proves the arbitrary
all-marks-on-one-jump bound, and closes that complete diagonal
CMCQC sector after summing all centre subsets and Gaussian pairings.
It isolates OIDC for partitions involving two or more distinct
added jumps.

V64 does not prove OIDC, IDCC, CMCQC in full, the full one-centre
subordinated tangent card, MCWC, OCRC in its full noncommutative
form, SCNGC, WPSC, GRLTC, LZTC, PLRC, WSDC, BRSC, HWSFC, UGCGC,
thermal connected WPRAC, all-arity RGSC, the raw Wilson aggregate
strong card, all-order strong MTA, WUFC, CPM, cutoff removal,
reflection positivity, Osterwalder--Schrader reconstruction, a
continuum Yang--Mills measure, or a positive Hamiltonian gap.  The
full Yang--Mills mass gap has not been proved.
\end{firewall}

\begin{theorem}[V64 result ledger]
\label{thm:newV64-results}
V64:
\begin{enumerate}[label=\textup{(V64.\arabic*)},leftmargin=6.5em]
\item proves the multivariate Mecke--cumulant formula;
\item identifies clock-mark blocks with common L\'evy jumps;
\item records the exact two-mark decomposition;
\item bounds the arbitrary diagonal marked-clock term;
\item keeps typical-clock and high-jump powers separate;
\item sums every choice of centre leaves and Gaussian pairs;
\item closes the all-marks-on-one-jump CMCQC sector; and
\item isolates OIDC as the remaining multi-jump row.
\end{enumerate}
\end{theorem}

\begin{proof}
Use
\Cref{thm:newV64-Mecke-cumulant,
cor:newV64-first-rows,
thm:newV64-diagonal-bound,
thm:newV64-diagonal-CMCQC,
cor:newV64-diagonal-closed,
prop:newV64-OIDC}.
\end{proof}

\paragraph{Parent-version record.}
V64 was formed from
\texttt{Renewal\_Yang\_Mills\_Mass\_Gap\_Research\_Notes\_V63.tex}.
The V63 parent had \(30\,746\) logical source lines,
\(978\,216\) bytes, and SHA--256
\[
 \texttt{f4bfc57c40ff1c54778c1d11b714292942a1a264d368b4871096e685bb0bbad8}.
\]
The next compulsory companion audit is V65.

% =============================================================================
% END APPENDED FIFTH-SERIES V64 MATERIAL
% =============================================================================

\clearpage
\chapter{Fifth-series V65: compulsory audit and the inner
multi-jump centre factorization}
\label{ch:newV65-inner}

\section{Compulsory companion audit}
\label{sec:newV65-audit}

\begin{audit}[Newest active companion snapshots]
\label{audit:newV65-companions}
The following sources were read atomically and were not modified:
\begin{align*}
 &\texttt{Renewal\_SM\_Einstein\_Continuum\_Research\_Notes\_V90.tex},\\
 &\qquad 30\,876\ \hbox{logical lines},\quad
 1\,185\,228\ \hbox{bytes},\\
 &\qquad
 \texttt{d822f8d46f0619a01e352d8f98a9c085f9caebf3cb25f4dc440f9141abb3c784},
 \\
 &\texttt{Renewal\_NS\_Stability\_Research\_Notes\_V19.tex},\\
 &\qquad 4\,995\ \hbox{logical lines},\quad
 252\,441\ \hbox{bytes},\\
 &\qquad
 \texttt{33c2da92a60d5b8a2b69fea3c5c37d4d0d698419cc9eeb1c6334bd571a23a25e}.
\end{align*}

SM V90 proves agreement of exact sharp and proper-time trace
subtractions only after the complete trace insertion is retained.
Subtracting merely its leading asymptotic leaves a finite local
scheme factor.  It also proves that cancellation of a metric beat
does not cancel a conformal coarea trace of a different homogeneous
order.

NS V19 derives an exact Gaussian enstrophy balance whose invariant
mean is positive, but explicitly records that the balance is not a
regularity obstruction.  The signed and unsigned pieces must be
kept distinct.
\end{audit}

\begin{assessment}[V65 transfer discipline]
\label{ass:newV65-transfer}
No numerical SM or NS estimate is imported.  For OIDC the applicable
rules are:
\begin{enumerate}[label=\textup{(\roman*)},leftmargin=3.5em]
\item keep the complete iterated finite difference assembled before
      separating its asymptotic regions;
\item do not infer an off-diagonal estimate from the exact Mecke
      balance alone; and
\item do not ask the already-closed one-jump sector to cancel a
      distinct-jump term of another homogeneous order.
\end{enumerate}
\end{assessment}

\section{Analytic factorization of the centre block}
\label{sec:newV65-factorization}

Recall
\[
 {\cal C}_k(t)
 =
 \kappa_{k+1}^{\gamma_3}
 (F_{0,t},L_{1,t},\ldots,L_{k,t}),
 \qquad
 A_k=C_2(R_0)^{k/2}\prod_{j=1}^kC_2(R_j)^{1/2}.
\]

\begin{theorem}[Exact \(t^k\) centre factorization]
\label{thm:newV65-factorization}
For every \(k\ge1\), there is an operator-valued entire function
\(H_k\) such that
\begin{equation}
 {\cal C}_k(t)=t^kH_k(t).                                    \label{eq:newV65-factorization}
\end{equation}
For every \(q\ge0\) and \(0\le tC_2(R_0)\le c_0\),
\begin{equation}
 \|H_k^{(q)}(t)\|_{\rm cb}
 \le
 K^{k+q}(k+2q)^{2q}
 C_2(R_0)^{k/2+q}
 \prod_{j=1}^kC_2(R_j)^{1/2}.                               \label{eq:newV65-H-derivative}
\end{equation}
\end{theorem}

\begin{proof}
Use the Gaussian-centre identity
\Cref{lem:newV62-Gaussian-centre}.  Each of the \(k\) leaf
coefficients contributes \(\sqrt{2t}\), and each of the \(k\)
directional derivatives of \(F_{0,t}\) contributes another
\(\sqrt{2t}\).  Their product is \(t^k\).

Expand the remaining time-ordered centre exponentials in
\eqref{eq:newV62-Duhamel}.  A term with an odd total number of
additional Gaussian coordinates has zero expectation.  Pairing the
additional \(2\ell\) coordinates gives
\begin{equation}
 H_k(t)=\sum_{\ell\ge0}t^\ell H_{k,\ell},\qquad
 \|H_{k,\ell}\|_{\rm cb}
 \le
 {K^{k+\ell}\over\ell!}
 C_2(R_0)^{k/2+\ell}
 \prod_{j=1}^kC_2(R_j)^{1/2}.                               \label{eq:newV65-H-series}
\end{equation}
Indeed, the exponential coefficient and the ordered-simplex volume
cancel the ordering factorial, the Gaussian pairings leave at most
\((K\ell)^\ell\), and
\((2\ell)!\ge\ell!\,\ell^\ell\) absorbs that factor into
\(K^\ell/\ell!\).  Coordinate sums cost only \(3^{k+2\ell}\),
which changes \(K\).

The series is entire.  Differentiating \(q\) times, using
\((\ell)_q/\ell!=1/(\ell-q)!\), and summing the resulting exponential
series on \(tC_2(R_0)\le c_0\) proves
\eqref{eq:newV65-H-derivative}; the displayed polynomial safely
absorbs the choices of which ordered insertions meet the
\(q\) derivatives.  The estimate remains unchanged under
amplification.
\end{proof}

\begin{corollary}[Inner iterated-difference card]
\label{cor:newV65-inner-difference}
Let \(q\ge1\), \(x_\nu\ge0\), and suppose
\begin{equation}
 C_2(R_0)\left(t+\sum_{\nu=1}^qx_\nu\right)\le c_0.           \label{eq:newV65-inner-region}
\end{equation}
Then
\begin{align}
 &\left\|
 \Delta_{x_1}\cdots\Delta_{x_q}{\cal C}_k(t)
 \right\|_{\rm cb}
 \notag\\
 &\quad\le
 K^{k+q}(k+2q)^{2q}A_k
 \left(\prod_{\nu=1}^qx_\nu\right)
 \sum_{h=0}^{\min(k,q)}
 t^{k-h}C_2(R_0)^{q-h}
 \left(\sum_\nu x_\nu\right)^{q-h}.                          \label{eq:newV65-inner-difference}
\end{align}
\end{corollary}

\begin{proof}
The exact rectangle formula is
\[
 \Delta_{x_1}\cdots\Delta_{x_q}f(t)
 =
 \int_0^{x_1}\cdots\int_0^{x_q}
 f^{(q)}(t+u_1+\cdots+u_q)\,du_1\cdots du_q.
\]
Apply it to \(f(t)=t^kH_k(t)\), use Leibniz's rule, and apply
\eqref{eq:newV65-H-derivative} throughout the rectangle, which lies
in \eqref{eq:newV65-inner-region}.  Bounding
\((t+\sum u)^{k-h}\) by the binomial expansion and enlarging the
displayed polynomial gives
\eqref{eq:newV65-inner-difference}.
\end{proof}

\begin{assessment}[Every distinct inner jump now pays a difference]
\label{ass:newV65-inner-payment}
In \eqref{eq:newV64-OIDC}, the block \(B_\nu\) already contributes
\(x_\nu^{|B_\nu|}\).  The inner factorization adds one further
\(x_\nu\) for every distinct jump.  Thus its jump integral is of
order
\[
 \int x_\nu^{|B_\nu|+1}\Pi(dx_\nu)
 \lesssim
 (|B_\nu|+1)!\,s^{2|B_\nu|+1},
\]
which is precisely two clock powers per outer connection.  The
remaining task is summation of the Leibniz index \(h\), the partition
blocks, and the complementary region in which
\(C_2(R_0)(S+\sum x_\nu)>c_0\).
\end{assessment}

\section{The complementary unitary region}
\label{sec:newV65-outer}

\begin{definition}[Outer-centre assembled card, OCAC]
\label{def:newV65-OCAC}
OCAC is a bound on the contribution to OIDC from
\begin{equation}
 C_2(R_0)\left(S_{\rm sm}+\sum_{\nu=1}^qx_\nu\right)>c_0,     \label{eq:newV65-outer-region}
\end{equation}
obtained with the complete alternating finite difference and the
unitarity of \(F_{0,t}\), without expanding \(H_k\) term by term.
After summing clock partitions and Gaussian pairings it must obey
the strong thermal target.
\end{definition}

\begin{proposition}[Inner card plus OCAC closes OIDC]
\label{prop:newV65-inner-outer}
If the partition sum obtained from
\eqref{eq:newV65-inner-difference} has one exponential arity base
and OCAC holds, then OIDC, IDCC, and CMCQC follow.
\end{proposition}

\begin{proof}
Split every integral in \eqref{eq:newV64-OIDC} by
\eqref{eq:newV65-inner-region} and
\eqref{eq:newV65-outer-region}.  The two hypotheses control the two
pieces.  Apply \Cref{prop:newV64-OIDC} and then
\Cref{prop:newV63-IDCC}.
\end{proof}

\begin{programme}[V66: inner partition sum and outer-centre tail]
\label{prog:newV65-V66}
\begin{enumerate}[leftmargin=3em]
\item Insert \eqref{eq:newV65-inner-difference} into the exact
      Mecke partition formula and sum the block factorials by an
      exponential generating function.
\item Use \(C_2(R_0)S>c_0/2\) or
      \(C_2(R_0)\sum x_\nu>c_0/2\) to split the outer region.
\item Bound the first event with the small-clock Stieltjes support
      and the second with a marked union bound that retains the
      complete alternating difference.
\item Test whether the nonperturbative outer probability can be
      converted into the expanding \(1/s\) source radius before
      taking norms.
\item Keep WPSC separate until the small-clock outer-centre row is
      genuinely closed.
\end{enumerate}
\end{programme}

\begin{firewall}[Claim boundary after V65]
\label{fire:newV65-boundary}
V65 completes the compulsory companion audit, proves the exact
\(t^kH_k(t)\) factorization of the complete Gaussian centre, proves
uniform coefficient derivatives on the inner centre region, and
proves that every distinct inner jump supplies the additional
finite-difference factor required by OIDC.  It isolates OCAC for the
complementary unitary region.

V65 does not prove the complete inner partition sum, OCAC, OIDC,
IDCC, CMCQC in full, the full one-centre subordinated tangent card,
MCWC, OCRC in its full noncommutative form, SCNGC, WPSC, GRLTC,
LZTC, PLRC, WSDC, BRSC, HWSFC, UGCGC, thermal connected WPRAC,
all-arity RGSC, the raw Wilson aggregate strong card, all-order
strong MTA, WUFC, CPM, cutoff removal, reflection positivity,
Osterwalder--Schrader reconstruction, a continuum Yang--Mills
measure, or a positive Hamiltonian gap.  The full Yang--Mills mass
gap has not been proved.
\end{firewall}

\begin{theorem}[V65 result ledger]
\label{thm:newV65-results}
V65:
\begin{enumerate}[label=\textup{(V65.\arabic*)},leftmargin=6.5em]
\item completes the required newest-SM/newest-NS audit;
\item imports only the assemble-before-limiting proof discipline;
\item proves the exact \(t^kH_k(t)\) centre factorization;
\item proves the uniform inner coefficient derivative bound;
\item proves the inner iterated-difference estimate;
\item proves one added factor \(x_\nu\) for each distinct jump;
\item separates the complementary unitary region; and
\item isolates OCAC as the outer-centre card.
\end{enumerate}
\end{theorem}

\begin{proof}
Use
\Cref{audit:newV65-companions,
thm:newV65-factorization,
cor:newV65-inner-difference,
prop:newV65-inner-outer}.
\end{proof}

\paragraph{Parent-version record.}
V65 was formed from
\texttt{Renewal\_Yang\_Mills\_Mass\_Gap\_Research\_Notes\_V64.tex}.
The V64 parent had \(31\,147\) logical source lines,
\(990\,912\) bytes, and SHA--256
\[
 \texttt{321b2bf89d8a0e9fa99926d9258f4f4e899b3bb6624cbfbf9e70fbe14d1bfa7c}.
\]
The next compulsory companion audit is V70.

% =============================================================================
% END APPENDED FIFTH-SERIES V65 MATERIAL
% =============================================================================

\clearpage
\chapter{Fifth-series V66: correction and closure of the complete
inner distinct-jump partition sum}
\label{ch:newV66-inner-sum}

\section{Correction of the V65 rectangle estimate}
\label{sec:newV66-correction}

No companion audit is due at V66.

\begin{correction}[Correct Leibniz power in the inner difference]
\label{corr:newV66-inner}
The conclusion of
\Cref{cor:newV65-inner-difference} is replaced by the following
estimate.  Under
\[
 C_2(R_0)\left(t+\sum_{\nu=1}^qx_\nu\right)\le c_0,
\]
one has
\begin{align}
 &\left\|
 \Delta_{x_1}\cdots\Delta_{x_q}{\cal C}_k(t)
 \right\|_{\rm cb}
 \notag\\
 &\quad\le
 K^{k+q}(k+2q)^{2q}A_k
 \left(\prod_{\nu=1}^qx_\nu\right)
 \sum_{h=0}^{\min(k,q)}
 \left(t+\sum_{\nu=1}^qx_\nu\right)^{k-h}
 C_2(R_0)^{q-h}.                                            \label{eq:newV66-inner-correct}
\end{align}
The additional factor
\((\sum x_\nu)^{q-h}\) in
\eqref{eq:newV65-inner-difference} was not justified and is not
used below.
\end{correction}

\begin{proof}
Use the exact rectangle formula in the proof of
\Cref{cor:newV65-inner-difference}.  In the Leibniz expansion of
\[
 {d^q\over dt^q}\{t^kH_k(t)\},
\]
let \(h\) derivatives hit \(t^k\) and \(q-h\) hit \(H_k\).
The former factor is bounded by
\((Kk)^h(t+\sum x)^{k-h}\); the latter is bounded by
\eqref{eq:newV65-H-derivative} with
\(C_2(R_0)^{q-h}\).  The rectangle volume is
\(\prod_\nu x_\nu\).  Sum \(h\) and absorb binomial and falling
factorial coefficients into the displayed polynomial.
\end{proof}

\section{Clock integration on one fixed partition}
\label{sec:newV66-fixed-partition}

Let \(\pi=\{B_1,\ldots,B_q\}\) be a partition of \([r]\), and put
\(p_\nu=|B_\nu|\), so \(\sum_\nu p_\nu=r\).  Denote by
\({\cal I}^{\rm in}_{k,\pi}\) the contribution to
\eqref{eq:newV64-Mecke-cumulant} restricted to
\[
 C_2(R_0)\left(S_{\rm sm}+\sum_{\nu=1}^qx_\nu\right)\le c_0.
\]

\begin{lemma}[Fixed inner partition clock grade]
\label{lem:newV66-fixed-partition}
For \(k\ge1\),
\begin{align}
 \|{\cal I}^{\rm in}_{k,\pi}\|_{\rm cb}
 \le{}&
 K^{k+r+q}(k+2r)^{2r}A_k
 s^{k+2r}(1+k\sqrt s)^k
 \notag\\
 &\times
 \prod_{\nu=1}^q(p_\nu+1)! .                                \label{eq:newV66-fixed-partition}
\end{align}
\end{lemma}

\begin{proof}
Insert \eqref{eq:newV66-inner-correct}.  Besides the exterior
\(x_\nu^{p_\nu}\) from
\eqref{eq:newV64-Mecke-cumulant}, every difference supplies one
additional \(x_\nu\).

Fix \(h\), expand
\[
 \left(S+\sum_\nu x_\nu\right)^{k-h}
\]
multinomially, and let \(j_0+\sum_\nu j_\nu=k-h\).  The corresponding
clock factor is bounded by
\[
 \mathbb ES^{j_0}
 \prod_{\nu=1}^q
 \int x_\nu^{p_\nu+1+j_\nu}\Pi(dx_\nu).
\]
By
\Cref{thm:newV62-clock-moments,
lem:newV60-jump-moments}, its power of \(s\) is at least
\[
 s^{j_0}
 \prod_\nu s^{2(p_\nu+1+j_\nu)-1}
 =
 s^{j_0+2r+q+2\sum_\nu j_\nu}
 \le
 s^{k-h+2r+q}.
\]
The inequality is in the useful direction because
\(j_0+2\sum j_\nu\ge k-h\) and \(0<s\le1\).
Multiplication by
\[
 C_2(R_0)^{q-h}
 \le b_0^{2(q-h)}s^{-(q-h)}
 \le s^{-(q-h)}
\]
leaves \(s^{k+2r}\).

The raw-clock moment contributes
\((1+k\sqrt s)^k\).  Gamma factors from the jump moments are bounded
by
\[
 (p_\nu+1+j_\nu)!
 \le
 (p_\nu+1)!\{K(k+2r)\}^{j_\nu}.
\]
The multinomial sum, the \(h\)-sum, and all remaining polynomial
factors are absorbed by
\(K^{k+r+q}(k+2r)^{2r}\).  This proves
\eqref{eq:newV66-fixed-partition}.
\end{proof}

\section{The weighted partition generating function}
\label{sec:newV66-partitions}

\begin{lemma}[Block-factorial partition sum]
\label{lem:newV66-partition-sum}
There is a numerical \(K\) such that, for every \(r\ge1\),
\begin{equation}
 \sum_{\pi\in{\rm Part}([r])}
 \prod_{B\in\pi}(|B|+1)!
 \le K^r r!.                                                 \label{eq:newV66-partition-sum}
\end{equation}
\end{lemma}

\begin{proof}
The labelled exponential formula gives
\[
 {1\over r!}
 \sum_{\pi\in{\rm Part}([r])}
 \prod_{B\in\pi}(|B|+1)!
 =
 [z^r]\exp\left(
 \sum_{p\ge1}(p+1)z^p
 \right).
\]
For \(|z|<1\),
\[
 \sum_{p\ge1}(p+1)z^p
 ={z(2-z)\over(1-z)^2}.
\]
Cauchy's estimate on the fixed circle \(|z|=1/2\) bounds the
coefficient by \(K^r\).
\end{proof}

\begin{theorem}[Complete inner OIDC]
\label{thm:newV66-inner-OIDC}
After summing every partition \(\pi\) of the \(r\) clock marks, every
choice of \(k\) centre leaves, and every pairing of the remaining
\(2r\) Gaussian leaves, the complete contribution from
\[
 C_2(R_0)\left(S_{\rm sm}+\sum_{\nu=1}^{|\pi|}x_\nu\right)
 \le c_0
\]
obeys
\begin{equation}
 K^{k+2r+1}
 b_0\left(\prod_{\ell=1}^{k+2r}b_\ell\right)
 B^{k+2r-1}.                                                 \label{eq:newV66-inner-OIDC}
\end{equation}
\end{theorem}

\begin{proof}
Sum \eqref{eq:newV66-fixed-partition} and use
\Cref{lem:newV66-partition-sum}.  The number of choices of the
centre subset and pairings of its complement is
\[
 { (k+2r)!\over k!\,2^rr!}
 \le { (k+2r)^{2r}\over2^rr!}.
\]
The \(r!\) from
\eqref{eq:newV66-partition-sum} cancels its denominator.

Convert the \(k+r\) Casimir square-root pairs into thermal marks.
The factor \(s^{k+2r}\) leaves \(s^r\), and
\[
 (k+2r)^{2r}s^r\le B^{2r}
\]
by the thermal floor.  The remaining centre factor satisfies
\[
 b_0^{k-1}(1+k\sqrt s)^k
 \le K^kB^{k-1};
\]
this is the V62 inequality
\(b_0(1+k\sqrt s)\le b_0+k\sqrt s\le B\), with one residual
factor absorbed exponentially.  All displayed polynomial
overcounts enter the same exponential base, proving
\eqref{eq:newV66-inner-OIDC}.
\end{proof}

\begin{corollary}[OIDC is reduced to OCAC]
\label{cor:newV66-OIDC-reduction}
The complete inner part of OIDC is closed.  Hence OCAC alone is the
remaining row required for OIDC, IDCC, and CMCQC.
\end{corollary}

\begin{proof}
Combine
\Cref{thm:newV66-inner-OIDC,prop:newV65-inner-outer}.
\end{proof}

\section{A low-centre outer-tail estimate}
\label{sec:newV66-tail}

\begin{lemma}[Small-clock tail at the centre threshold]
\label{lem:newV66-centre-tail}
There are numerical \(\beta_0,c,C>0\) such that, if
\(b_0\le\beta_0\), then
\begin{equation}
 \mathbb P\left\{
 C_2(R_0)S_{\rm sm}>{c_0\over2}
 \right\}
 \le
 C\exp\left\{-{c\over s b_0^2}\right\}.                      \label{eq:newV66-centre-tail}
\end{equation}
\end{lemma}

\begin{proof}
Since
\[
 C_2(R_0)\le {b_0^2\over s},
\]
the event implies
\[
 S_{\rm sm}>{c_0s\over2b_0^2}.
\]
Choose \(\beta_0\) so that this threshold is at least
\(2\mathbb ES_{\rm sm}\).  The Chernoff bound from
\Cref{lem:newV62-subgamma} has Bernstein form
\[
 \mathbb P\{S-\mathbb ES>u\}
 \le
 \exp\left\{
 -c\min\left({u^2\over s^3},{u\over s^2}\right)
 \right\}.
\]
At \(u\asymp s/b_0^2\), the second quantity is the smaller one and
equals \(c/(sb_0^2)\), proving
\eqref{eq:newV66-centre-tail}.
\end{proof}

\begin{assessment}[Position of the outer card]
\label{ass:newV66-outer}
For a small centre, the part of OCAC caused by the original clock
\(S_{\rm sm}\) is nonperturbative at the stronger scale
\(e^{-c/(sb_0^2)}\).  For a centre of fixed thermal size, that event
need not be rare, but the centre itself supplies a fixed mark.  The
remaining proof must combine this dichotomy with the event that the
added jumps, rather than \(S_{\rm sm}\), cross the centre threshold.
The estimate must remain assembled across the partition because a
large block size can reach the same threshold without a rare
individual jump.
\end{assessment}

\begin{programme}[V67: OCAC]
\label{prog:newV66-V67}
\begin{enumerate}[leftmargin=3em]
\item Derive an exponentially tilted bound for
      \(\sum_\nu x_\nu\) with the block weights
      \(x_\nu^{|B_\nu|}\Pi(dx_\nu)\) still present.
\item Separate block orders below and above
      \(c/(sb_0^2)\); use the tail estimate only below the crossover.
\item In the high-order regime, convert the Gamma factor directly
      into the available \(B^{2r}\) through the thermal floor.
\item Treat \(b_0\le\beta_0\) and \(b_0>\beta_0\) separately, then
      sum the regimes before the final norm.
\item If OCAC closes, conclude OIDC, IDCC, and CMCQC and return to
      the exact compact-group small-clock source.
\end{enumerate}
\end{programme}

\begin{firewall}[Claim boundary after V66]
\label{fire:newV66-boundary}
V66 corrects the V65 rectangle estimate, proves the fixed-partition
inner clock grade, sums all block factorials by an exact exponential
generating function, and closes the complete inner OIDC contribution.
It also proves a nonperturbative low-centre tail bound and reduces
OIDC to OCAC.

V66 does not prove OCAC, OIDC, IDCC, CMCQC in full, the full
one-centre subordinated tangent card, MCWC, OCRC in its full
noncommutative form, SCNGC, WPSC, GRLTC, LZTC, PLRC, WSDC, BRSC,
HWSFC, UGCGC, thermal connected WPRAC, all-arity RGSC, the raw
Wilson aggregate strong card, all-order strong MTA, WUFC, CPM,
cutoff removal, reflection positivity, Osterwalder--Schrader
reconstruction, a continuum Yang--Mills measure, or a positive
Hamiltonian gap.  The full Yang--Mills mass gap has not been proved.
\end{firewall}

\begin{theorem}[V66 result ledger]
\label{thm:newV66-results}
V66:
\begin{enumerate}[label=\textup{(V66.\arabic*)},leftmargin=6.5em]
\item corrects the inherited inner finite-difference estimate;
\item proves the fixed-partition inner clock grade;
\item proves the block-factorial partition generating function;
\item closes the complete inner OIDC partition sum;
\item reduces OIDC, IDCC, and CMCQC to OCAC;
\item proves the low-centre small-clock outer-tail bound; and
\item isolates the block-order crossover for V67.
\end{enumerate}
\end{theorem}

\begin{proof}
Use
\Cref{corr:newV66-inner,
lem:newV66-fixed-partition,
lem:newV66-partition-sum,
thm:newV66-inner-OIDC,
cor:newV66-OIDC-reduction,
lem:newV66-centre-tail}.
\end{proof}

\paragraph{Parent-version record.}
V66 was formed from
\texttt{Renewal\_Yang\_Mills\_Mass\_Gap\_Research\_Notes\_V65.tex}.
The V65 parent had \(31\,424\) logical source lines,
\(1\,000\,935\) bytes, and SHA--256
\[
 \texttt{c2db1839d1f424f1b2ccb8ff02f96c598aa8607e6956f21433ca827ea82b14e1}.
\]
The next compulsory companion audit is V70.

% =============================================================================
% END APPENDED FIFTH-SERIES V66 MATERIAL
% =============================================================================

\clearpage
\chapter{Fifth-series V67: Gamma domination of block-weighted
small-clock jumps}
\label{ch:newV67-gamma}

\section{Block-weighted jump laws}
\label{sec:newV67-laws}

No companion audit is due at V67.  For an integer \(p\ge1\), put
\begin{equation}
 M_p:=\int_0^\infty x^p\Pi_\alpha^{\rm sm}(dx),\qquad
 \eta_p(dx):={x^p\Pi_\alpha^{\rm sm}(dx)\over M_p}.           \label{eq:newV67-eta}
\end{equation}

\begin{theorem}[Uniform Gamma domination]
\label{thm:newV67-gamma}
For \(0\le\lambda<T_\alpha\),
\begin{equation}
 \int e^{\lambda x}\eta_p(dx)
 \le
 \left(1-{\lambda\over T_\alpha}\right)^{-(p+1)}.             \label{eq:newV67-gamma}
\end{equation}
\end{theorem}

\begin{proof}
Use the high-spectral representation
\[
 \Pi_\alpha^{\rm sm}(dx)
 =
 \int_{(T_\alpha,\infty)}t e^{-tx}\sigma_\alpha(dt)\,dx.
\]
Tonelli gives
\begin{align*}
 M_p
 &=
 \Gamma(p+1)\int_{(T_\alpha,\infty)}t^{-p}\sigma_\alpha(dt),\\
 \int e^{\lambda x}x^p\Pi(dx)
 &=
 \Gamma(p+1)\int_{(T_\alpha,\infty)}
 {t\over(t-\lambda)^{p+1}}\sigma_\alpha(dt).
\end{align*}
For \(t>T_\alpha\),
\[
 {t\over(t-\lambda)^{p+1}}
 =
 t^{-p}(1-\lambda/t)^{-(p+1)}
 \le
 t^{-p}(1-\lambda/T_\alpha)^{-(p+1)}.
\]
Divide by \(M_p\).
\end{proof}

\begin{corollary}[A partition has one Gamma shape]
\label{cor:newV67-partition-gamma}
Let \(\pi=\{B_1,\ldots,B_q\}\) partition \([r]\), and put
\(p_\nu=|B_\nu|\).  If \(X_\nu\) are independent with laws
\(\eta_{p_\nu}\), then
\begin{equation}
 \mathbb E
 e^{\lambda\sum_{\nu=1}^qX_\nu}
 \le
 \left(1-{\lambda\over T_\alpha}\right)^{-(r+q)}.             \label{eq:newV67-sum-mgf}
\end{equation}
Consequently, for \(a>0\),
\begin{equation}
 \mathbb P\left\{\sum_{\nu=1}^qX_\nu>a\right\}
 \le
 2^{r+q}e^{-T_\alpha a/2}.                                  \label{eq:newV67-sum-tail}
\end{equation}
\end{corollary}

\begin{proof}
Multiply \eqref{eq:newV67-gamma}; the exponents add to
\(\sum_\nu(p_\nu+1)=r+q\).  Apply exponential Markov at
\(\lambda=T_\alpha/2\).
\end{proof}

\section{The added-jump outer-centre threshold}
\label{sec:newV67-threshold}

\begin{theorem}[Weighted outer-jump tail]
\label{thm:newV67-weighted-tail}
If \(C_2(R_0)>0\), then for every partition \(\pi\) as above,
\begin{align}
 &\int
 \mathbf1_{\{C_2(R_0)\sum_\nu x_\nu>c_0/2\}}
 \prod_{\nu=1}^q x_\nu^{p_\nu}
 \Pi_\alpha^{\rm sm}(dx_\nu)
 \notag\\
 &\qquad\le
 2^{r+q}
 \exp\left\{-{c\over sb_0^2}\right\}
 \prod_{\nu=1}^qM_{p_\nu}.                                  \label{eq:newV67-weighted-tail}
\end{align}
If \(C_2(R_0)=0\), the left side is zero.
\end{theorem}

\begin{proof}
Normalize each factor by \eqref{eq:newV67-eta} and apply
\eqref{eq:newV67-sum-tail} with
\[
 a={c_0\over2C_2(R_0)}.
\]
Since \(T_\alpha\ge c/s^2\) and
\(C_2(R_0)\le b_0^2/s\),
\[
 T_\alpha a\ge {c\over sb_0^2}.
\]
Restore the factors \(M_{p_\nu}\).
\end{proof}

\begin{lemma}[Nonperturbative factor restores any prescribed number
of differences]
\label{lem:newV67-restore}
For every integer \(q\ge1\),
\begin{equation}
 \exp\left\{-{c\over sb_0^2}\right\}
 \le
 (Kqsb_0^2)^q.                                               \label{eq:newV67-restore}
\end{equation}
\end{lemma}

\begin{proof}
The maximum of \(y^qe^{-y}\) on \(y>0\) is \(q^qe^{-q}\).
Set \(y=c/(sb_0^2)\) and rearrange.
\end{proof}

\begin{corollary}[Recovery of the \(q\) missing inner powers]
\label{cor:newV67-recovery}
The left side of
\eqref{eq:newV67-weighted-tail} is at most
\begin{equation}
 K^{r+q}s^{2r}b_0^{2q}q^q
 \prod_{\nu=1}^q p_\nu!.                                    \label{eq:newV67-recovery}
\end{equation}
\end{corollary}

\begin{proof}
By \Cref{lem:newV60-jump-moments},
\[
 M_{p_\nu}\le p_\nu!K^{p_\nu}s^{2p_\nu-1}.
\]
Their product contributes
\[
 K^rs^{2r-q}\prod_\nu p_\nu!.
\]
Apply \Cref{lem:newV67-restore}; its \(s^q\) restores exactly the
power missing because the outer estimate did not differentiate the
complete centre once for each distinct jump.
\end{proof}

\begin{assessment}[Measure-level part of OCAC is closed]
\label{ass:newV67-measure}
For the event that the added jumps cross the centre threshold, the
block-weighted product measure has a Gamma tail with rate
\(T_\alpha\asymp s^{-2}\).  The threshold is
\(C_2(R_0)^{-1}\gtrsim s/b_0^2\), so its cost is
\(e^{-c/(sb_0^2)}\).  This cost restores all \(q\) factors of \(s\)
which were supplied by finite differences in the inner region.
The remaining issue is to retain this recovery after multiplication
by the polynomial clock size of the complete centre and after the
partition sum.
\end{assessment}

\section{Polynomially marked Gamma tail}
\label{sec:newV67-polynomial}

\begin{lemma}[Tail with a polynomial mark]
\label{lem:newV67-polynomial-tail}
For every integer \(d\ge0\),
\begin{align}
 &\int
 \mathbf1_{\{\sum_\nu x_\nu>a\}}
 \left(\sum_\nu x_\nu\right)^d
 \prod_{\nu=1}^q x_\nu^{p_\nu}\Pi(dx_\nu)
 \notag\\
 &\quad\le
 \left\{
 {K(r+q+d)\over T_\alpha}
 \right\}^{d}
 4^{r+q+d}e^{-T_\alpha a/4}
 \prod_{\nu=1}^qM_{p_\nu}.                                  \label{eq:newV67-polynomial-tail}
\end{align}
\end{lemma}

\begin{proof}
Under the normalized product law, use
\[
 y^d\mathbf1_{\{y>a\}}
 \le
 d!\left({4\over T_\alpha}\right)^d
 e^{T_\alpha y/4}\mathbf1_{\{y>a\}}.
\]
Apply exponential Markov with a further parameter
\(T_\alpha/4\).  The total tilt is \(T_\alpha/2\), so
\eqref{eq:newV67-sum-mgf} gives \(2^{r+q}\), while the indicator
gives \(e^{-T_\alpha a/4}\).  Use
\(d!\le d^d\) and enlarge the numerical base.
\end{proof}

\begin{assessment}[The high-order crossover is visible]
\label{ass:newV67-crossover}
The polynomial mark in
\eqref{eq:newV67-polynomial-tail} costs
\((r+q+d)^dT_\alpha^{-d}\asymp
(r+q+d)^ds^{2d}\), not \(s^d\).
Thus a high centre degree improves rather than worsens the clock
power, until its factorial competes with the tail exponent.  The
crossover occurs when \(r+q+d\asymp1/(sb_0^2)\), exactly the scale
at which \eqref{eq:newV67-restore} should be replaced by direct
thermal-floor payment.
\end{assessment}

\begin{definition}[Outer assembly card after Gamma domination,
GOCAC]
\label{def:newV67-GOCAC}
GOCAC is the insertion of
\Cref{thm:newV67-weighted-tail,
lem:newV67-polynomial-tail}
and the original-clock tail
\eqref{eq:newV66-centre-tail} into the complete alternating centre
difference, followed by the joint partition/pairing sum, with the
two order regimes split at
\[
 k+r+q\asymp{1\over sb_0^2}.
\]
\end{definition}

\begin{proposition}[GOCAC implies OCAC]
\label{prop:newV67-GOCAC}
If the low-order tail recovery and high-order thermal-floor
recovery in GOCAC are each bounded with one exponential arity base,
then OCAC holds.
\end{proposition}

\begin{proof}
The outer region is contained in the union of
\[
 \{C_2(R_0)S>c_0/2\},
 \qquad
 \{C_2(R_0)\sum_\nu x_\nu>c_0/2\}.
\]
The first is controlled at measure level by
\Cref{lem:newV66-centre-tail}; the second by
\Cref{thm:newV67-weighted-tail}.  The polynomial clock factors
created by the complete centre are covered by
\Cref{lem:newV67-polynomial-tail}.  The two hypothesized order
summations therefore cover OCAC.
\end{proof}

\begin{programme}[V68: complete the Gamma-tail assembly]
\label{prog:newV67-V68}
\begin{enumerate}[leftmargin=3em]
\item Bound the complete alternating difference by
      \(2^qK^kA_k(S+\sum x)^k\) only after the outer indicator is
      inserted.
\item Use \eqref{eq:newV67-polynomial-tail} below the crossover and
      \(\sqrt s\le B/(k+2r)\) above it.
\item Sum \(\prod p_\nu!\,q^q\) over set partitions by its bivariate
      exponential generating function.
\item Treat the original-clock tail with the same two-regime split.
\item Close GOCAC and hence OCAC if the bivariate coefficient has
      one exponential base after the available \(B^{2r}\) is used.
\end{enumerate}
\end{programme}

\begin{firewall}[Claim boundary after V67]
\label{fire:newV67-boundary}
V67 proves Gamma domination of every block-weighted small-clock
jump law, the exact Gamma tail for a partition sum, the
nonperturbative added-jump centre-threshold cost, recovery of all
missing distinct-jump powers, and a polynomially marked tail bound.
It isolates GOCAC as the remaining assembly operation.

V67 does not prove GOCAC, OCAC, OIDC, IDCC, CMCQC in full, the full
one-centre subordinated tangent card, MCWC, OCRC in its full
noncommutative form, SCNGC, WPSC, GRLTC, LZTC, PLRC, WSDC, BRSC,
HWSFC, UGCGC, thermal connected WPRAC, all-arity RGSC, the raw
Wilson aggregate strong card, all-order strong MTA, WUFC, CPM,
cutoff removal, reflection positivity, Osterwalder--Schrader
reconstruction, a continuum Yang--Mills measure, or a positive
Hamiltonian gap.  The full Yang--Mills mass gap has not been proved.
\end{firewall}

\begin{theorem}[V67 result ledger]
\label{thm:newV67-results}
V67:
\begin{enumerate}[label=\textup{(V67.\arabic*)},leftmargin=6.5em]
\item proves uniform Gamma domination of block-weighted jumps;
\item proves the partition-sum Gamma tail;
\item proves the outer added-jump threshold estimate;
\item converts its nonperturbative cost into \(q\) powers of \(s\);
\item recovers the complete missing distinct-jump clock grade;
\item proves a polynomially marked Gamma-tail estimate; and
\item isolates the exact low/high-order crossover for GOCAC.
\end{enumerate}
\end{theorem}

\begin{proof}
Use
\Cref{thm:newV67-gamma,
cor:newV67-partition-gamma,
thm:newV67-weighted-tail,
lem:newV67-restore,
cor:newV67-recovery,
lem:newV67-polynomial-tail,
prop:newV67-GOCAC}.
\end{proof}

\paragraph{Parent-version record.}
V67 was formed from
\texttt{Renewal\_Yang\_Mills\_Mass_Gap\_Research\_Notes\_V66.tex}.
The V66 parent had \(31\,768\) logical source lines,
\(1\,011\,198\) bytes, and SHA--256
\[
 \texttt{2d1ecedbdce8c41e050452a93972be70a6f3923960ffc5e3b67be6d9e4e4bd2c}.
\]
The next compulsory companion audit is V70.

% =============================================================================
% END APPENDED FIFTH-SERIES V67 MATERIAL
% =============================================================================

\clearpage
\chapter{Fifth-series V68: factorial tail recovery and closure of
the low-centre outer card}
\label{ch:newV68-low-centre}

\section{The factorial form of nonperturbative recovery}
\label{sec:newV68-factorial}

No companion audit is due at V68.

\begin{lemma}[Factorial recovery]
\label{lem:newV68-factorial-recovery}
For every integer \(q\ge1\),
\begin{equation}
 \exp\left\{-{c\over sb_0^2}\right\}
 \le
 q!\left({Ksb_0^2}\right)^q.                                \label{eq:newV68-factorial-recovery}
\end{equation}
\end{lemma}

\begin{proof}
The exponential series gives \(e^y\ge y^q/q!\).
Set \(y=c/(sb_0^2)\) and invert.
\end{proof}

\begin{lemma}[Exact ordered-block identity]
\label{lem:newV68-ordered-blocks}
For every \(r\ge1\),
\begin{equation}
 \sum_{\pi\in{\rm Part}([r])}
 |\pi|!\prod_{B\in\pi}|B|!
 =
 2^{r-1}r!.                                                  \label{eq:newV68-ordered-blocks}
\end{equation}
\end{lemma}

\begin{proof}
For one block of size \(p\), the labelled exponential generating
function after the weight \(p!\) is inserted is
\[
 H(z)=\sum_{p\ge1}{p!\over p!}z^p={z\over1-z}.
\]
A partition into \(q\) unordered blocks has the factor \(1/q!\);
multiplication by \(q!\) changes its generating function to
\(H(z)^q\).  Therefore
\[
 {1\over r!}
 \sum_{\pi}|\pi|!\prod_{B\in\pi}|B|!
 =
 [z^r]\sum_{q\ge1}H(z)^q
 =
 [z^r]{z\over1-2z}
 =
 2^{r-1}.
\]
\end{proof}

\begin{assessment}[Why the \(q!\) must not be replaced by \(q^q\)]
\label{ass:newV68-factorial}
The coarser V67 bound \((Kqsb_0^2)^q\) is correct but loses the
ordered-block cancellation.  The factorial form
\eqref{eq:newV68-factorial-recovery} combines exactly with the
unordered partition factor \(1/q!\); after summing partitions only
the single \(r!\) in \eqref{eq:newV68-ordered-blocks} remains, and
that cancels the denominator in the Gaussian pairing census.
\end{assessment}

\section{Low-centre outer assembly for added jumps}
\label{sec:newV68-added}

Let \(\pi=\{B_1,\ldots,B_q\}\) partition \([r]\), with
\(p_\nu=|B_\nu|\).  Let
\[
 X=\sum_{\nu=1}^qx_\nu.
\]

\begin{theorem}[Fixed-partition added-jump outer bound]
\label{thm:newV68-added-fixed}
There is a numerical \(\beta_0>0\) such that, if
\(b_0\le\beta_0\), the part of
\eqref{eq:newV64-OIDC} on
\[
 C_2(R_0)X>{c_0\over2}
\]
is bounded by
\begin{align}
 &K^{k+r+q}A_k
 s^{k+2r}(1+(k+r)\sqrt s)^k
 \notag\\
 &\qquad\times
 q!\,b_0^{2q}\prod_{\nu=1}^q p_\nu!.                        \label{eq:newV68-added-fixed}
\end{align}
\end{theorem}

\begin{proof}
The complete alternating difference satisfies
\[
 \left\|
 \Delta_{x_1}\cdots\Delta_{x_q}{\cal C}_k(S)
 \right\|
 \le
 2^q(K(S+X))^kA_k
\]
by \eqref{eq:newV63-centre-size}.  Split
\((S+X)^k\le2^k(S^k+X^k)\).

For the \(S^k\) term, independence,
\Cref{thm:newV62-clock-moments}, and
\Cref{thm:newV67-weighted-tail} give
\[
 K^{k+r+q}A_k
 s^k(1+k\sqrt s)^k
 e^{-c/(sb_0^2)}
 s^{2r-q}\prod_\nu p_\nu!.
\]
Apply \Cref{lem:newV68-factorial-recovery}.

For the \(X^k\) term, use
\Cref{lem:newV67-polynomial-tail} with
\(a=c_0/(2C_2(R_0))\).  Since
\(T_\alpha^{-1}\le Ks^2\), its polynomial factor supplies
\[
 \{K(k+r+q)s^2\}^k.
\]
The same factorial recovery supplies \(s^q q!b_0^{2q}\).
This term contains \(s^{2k+2r}\), hence is smaller than the
right side of \eqref{eq:newV68-added-fixed} after the polynomial is
absorbed into \(1+(k+r)\sqrt s\) and the thermal floor.  Enlarge
\(K\).
\end{proof}

\section{Low-centre assembly for the original-clock tail}
\label{sec:newV68-original}

\begin{lemma}[Polynomial original-clock tail]
\label{lem:newV68-original-tail}
If \(b_0\le\beta_0\), then for every \(d\ge0\),
\begin{equation}
 \mathbb E\left[
 S_{\rm sm}^d
 \mathbf1_{\{C_2(R_0)S_{\rm sm}>c_0/2\}}
 \right]
 \le
 \{Ks(1+d\sqrt s)\}^{d}
 \exp\left\{-{c\over sb_0^2}\right\}.                        \label{eq:newV68-original-tail}
\end{equation}
\end{lemma}

\begin{proof}
By H\"older,
\[
 \mathbb E[S^d\mathbf1_E]
 \le
 (\mathbb ES^{2d})^{1/2}\mathbb P(E)^{1/2}.
\]
Apply
\Cref{thm:newV62-clock-moments,
lem:newV66-centre-tail} and decrease \(c\).
\end{proof}

\begin{theorem}[Fixed-partition original-clock outer bound]
\label{thm:newV68-original-fixed}
If \(b_0\le\beta_0\), the part of
\eqref{eq:newV64-OIDC} on
\[
 C_2(R_0)S_{\rm sm}>{c_0\over2}
\]
obeys the same right side as
\eqref{eq:newV68-added-fixed}.
\end{theorem}

\begin{proof}
Use the crude assembled difference bound from the proof of
\Cref{thm:newV68-added-fixed} and
\((S+X)^k\le2^k(S^k+X^k)\).
For \(S^k\), apply
\Cref{lem:newV68-original-tail}; for \(X^k\), use the unrestricted
Gamma moment following from
\eqref{eq:newV67-sum-mgf}, while the original-clock indicator
factors by independence.  The unnormalized jump measures contribute
\[
 K^rs^{2r-q}\prod_\nu p_\nu!.
\]
Apply
\eqref{eq:newV68-factorial-recovery} to restore \(s^q\), and absorb
the polynomial clock moments exactly as in
\Cref{thm:newV68-added-fixed}.
\end{proof}

\section{Summation of the low-centre outer card}
\label{sec:newV68-sum}

\begin{theorem}[Low-centre OCAC]
\label{thm:newV68-low-OCAC}
For \(b_0\le\beta_0\), the complete OCAC contribution, summed over
all clock partitions, choices of centre leaves, and Gaussian
pairings, obeys
\begin{equation}
 K^{k+2r+1}
 b_0\left(\prod_{\ell=1}^{k+2r}b_\ell\right)
 B^{k+2r-1}.                                                 \label{eq:newV68-low-OCAC}
\end{equation}
\end{theorem}

\begin{proof}
The outer event is contained in the union of the two events treated
in
\Cref{thm:newV68-added-fixed,thm:newV68-original-fixed}.
Sum their common bound over \(\pi\) and use the exact identity
\eqref{eq:newV68-ordered-blocks}.  The choice and pairing count is
\[
 { (k+2r)!\over k!\,2^rr!}
 \le { (k+2r)^{2r}\over2^rr!}.
\]
The \(r!\) from the ordered-block identity cancels this denominator.

After conversion of all Casimir square roots, the clock factor
\(s^{k+2r}\) leaves \(s^r\).  Thus
\[
 (k+2r)^{2r}s^r\le B^{2r}.
\]
As before,
\[
 b_0^{k-1}(1+(k+r)\sqrt s)^k
 \le K^kB^{k-1};
\]
the additional \(b_0^{2q}\le1\) only improves the estimate.
This gives \eqref{eq:newV68-low-OCAC}.
\end{proof}

\begin{corollary}[OIDC is closed for a small centre]
\label{cor:newV68-low-OIDC}
If \(b_0\le\beta_0\), then OIDC, IDCC, and CMCQC hold for every
\(k,r\).
\end{corollary}

\begin{proof}
The inner contribution is
\Cref{thm:newV66-inner-OIDC}; the outer contribution is
\Cref{thm:newV68-low-OCAC}.  Apply
\Cref{prop:newV64-OIDC,prop:newV63-IDCC}.
\end{proof}

\section{The remaining high-centre card}
\label{sec:newV68-high}

\begin{definition}[High-centre unitary card, HCUC]
\label{def:newV68-HCUC}
HCUC is OCAC restricted to \(b_0>\beta_0\), proved without a
small probability for
\(C_2(R_0)S_{\rm sm}>c_0/2\).  It must use the complete unitary
centre word, the Duhamel simplex, and the fixed lower mark
\(b_0>\beta_0\) before the clock partition is summed.
\end{definition}

\begin{assessment}[Why the remaining regime is different]
\label{ass:newV68-high}
When \(b_0>\beta_0\), the threshold
\(C_2(R_0)S\asymp1\) lies at ordinary clock time and is not a tail.
No nonperturbative estimate can be used honestly.  On the other
hand, every repeated centre visit carries a fixed mark, and the
time-ordered derivative of its unitary exponential has total
simplex mass one.  HCUC is therefore a noncommutative word
resummation problem, not a large-deviation problem.
\end{assessment}

\begin{proposition}[HCUC completes the one-centre tangent card]
\label{prop:newV68-HCUC}
HCUC implies OCAC for all centre weights and therefore closes OIDC,
IDCC, CMCQC, and the full one-complete-centre subordinated tangent
card.
\end{proposition}

\begin{proof}
\Cref{thm:newV68-low-OCAC} treats
\(b_0\le\beta_0\); HCUC treats its complement.  Apply
\Cref{cor:newV66-OIDC-reduction,
prop:newV64-OIDC,
prop:newV63-IDCC,
prop:newV62-CMCQC}.
\end{proof}

\begin{programme}[V69: high-centre unitary words]
\label{prog:newV68-V69}
\begin{enumerate}[leftmargin=3em]
\item Return to the ordered-simplex formula
      \eqref{eq:newV62-Duhamel} before applying the outer clock
      partition.
\item Root every distinct added jump at its first visit to the
      centre word and use \(b_0>\beta_0\) to pay repeated visits.
\item Sum word interlacings by simplex volume rather than by
      \(q!\) absolute orderings.
\item Prove HCUC first for two added jumps, including the first
      commutator, and then induct on the number of visits.
\item Prepare the compulsory V70 companion audit.
\end{enumerate}
\end{programme}

\begin{firewall}[Claim boundary after V68]
\label{fire:newV68-boundary}
V68 proves factorial nonperturbative recovery, the exact
ordered-block partition identity, the polynomial original-clock
tail, and the complete low-centre outer assembly.  Together with
V66 it closes OIDC, IDCC, and CMCQC for every small centre.  It
isolates HCUC as the remaining fixed-size-centre word problem.

V68 does not prove HCUC, the full one-centre subordinated tangent
card, MCWC, OCRC in its full noncommutative form, SCNGC, WPSC,
GRLTC, LZTC, PLRC, WSDC, BRSC, HWSFC, UGCGC, thermal connected
WPRAC, all-arity RGSC, the raw Wilson aggregate strong card,
all-order strong MTA, WUFC, CPM, cutoff removal, reflection
positivity, Osterwalder--Schrader reconstruction, a continuum
Yang--Mills measure, or a positive Hamiltonian gap.  The full
Yang--Mills mass gap has not been proved.
\end{firewall}

\begin{theorem}[V68 result ledger]
\label{thm:newV68-results}
V68:
\begin{enumerate}[label=\textup{(V68.\arabic*)},leftmargin=6.5em]
\item proves factorial tail recovery;
\item proves the exact ordered-block identity;
\item controls the added-jump outer event with the complete centre
      polynomial;
\item proves a polynomial original-clock tail;
\item controls the original-clock outer event;
\item closes the complete low-centre OCAC sum;
\item closes OIDC, IDCC, and CMCQC for small centres; and
\item isolates HCUC for fixed-size centres.
\end{enumerate}
\end{theorem}

\begin{proof}
Use
\Cref{lem:newV68-factorial-recovery,
lem:newV68-ordered-blocks,
thm:newV68-added-fixed,
lem:newV68-original-tail,
thm:newV68-original-fixed,
thm:newV68-low-OCAC,
cor:newV68-low-OIDC,
prop:newV68-HCUC}.
\end{proof}

\paragraph{Parent-version record.}
V68 was formed from
\texttt{Renewal\_Yang\_Mills\_Mass_Gap\_Research\_Notes\_V67.tex}.
The V67 parent had \(32\,091\) logical source lines,
\(1\,021\,195\) bytes, and SHA--256
\[
 \texttt{1881a8c884773b89469c6a6f7c64da3c9804bbce908fded509d015284d17feef}.
\]
The next compulsory companion audit is V70.

% =============================================================================
% END APPENDED FIFTH-SERIES V68 MATERIAL
% =============================================================================

\clearpage
\chapter{Fifth-series V69: global heat-semigroup differentiation
and the joint partition--pairing resummation}
\label{ch:newV69-global-centre}

\section{Correction of the inherited summation ledger}
\label{sec:newV69-correction}

No companion audit is due at V69.

\begin{correction}[The V66 inner assembly did not close as written]
\label{corr:newV69-V66}
The proof of \Cref{thm:newV66-inner-OIDC} contains an unabsorbed
superexponential factor.  The fixed-partition estimate
\eqref{eq:newV66-fixed-partition} contains
\((k+2r)^{2r}\), while the subsequent choice-and-pairing census
contributes a second factor bounded by \((k+2r)^{2r}/r!\).
After \eqref{eq:newV66-partition-sum} cancels \(r!\), the proof has
\[
 (k+2r)^{4r}s^r.
\]
The thermal floor pays only
\((k+2r)^{2r}s^r\le B^{2r}\); the remaining
\((k+2r)^{2r}\) cannot be absorbed into \(K^{k+2r}\).
Consequently the proof of \Cref{thm:newV66-inner-OIDC}, and hence
the deduction in \Cref{cor:newV68-low-OIDC}, was incomplete.

The fixed-partition estimate, the low-centre tail estimates, and the
exact identities in V67--V68 are not invalidated by this arithmetic
correction.  Their asserted complete closure is replaced by the
global proof below, which never creates either of the two
\((k+2r)^{2r}\) losses.
\end{correction}

\begin{assessment}[Why a termwise repair is insufficient]
\label{ass:newV69-joint}
Writing \(N=k+2r\), a clock derivative which hits the explicit
factor \(t^k\) creates \((k)_h\).  Bounding that factor before the
clock-partition sum would again produce \(k^h\), which is not
exponential in \(N\) when \(h\) is proportional to \(k\).
The derivative-hit index, the number of distinct jumps, and the
pairing denominator \(r!\) must therefore be summed in one
operation.  This is the same assemble-before-majorizing principle
which was identified in V60 and in the V65 companion audit.
\end{assessment}

\section{A global heat-semigroup representation}
\label{sec:newV69-heat}

Let
\[
 U_0(z):=\exp\left(i\sum_{a=1}^3z_aJ_0^a\right),
 \qquad z\in\mathbb R^3.
\]
Tensor placements and harmless powers of \(i\) are suppressed in
the notation below.

\begin{theorem}[Global Gaussian heat representation of the centre]
\label{thm:newV69-heat-representation}
For every \(k\ge0\), there is a smooth operator-valued function
\(\Psi_k:\mathbb R^3\to{\cal A}\) such that
\begin{equation}
 {\cal C}_k(t)=t^kH_k(t),
 \qquad
 H_k(t)=(e^{t\Delta_{\mathbb R^3}}\Psi_k)(0).                \label{eq:newV69-heat-representation}
\end{equation}
With \(A_0:=1\) and
\[
 A_k=C_2(R_0)^{k/2}
     \prod_{j=1}^kC_2(R_j)^{1/2}
 \quad (k\ge1),
\]
one has, globally for \(t\ge0\) and every \(m\ge0\),
\begin{equation}
 \|H_k^{(m)}(t)\|_{\rm cb}
 \le
 K^{k+m}A_k C_2(R_0)^m.                                    \label{eq:newV69-global-H}
\end{equation}
Consequently, for \(q\ge0\),
\begin{align}
 \|{\cal C}_k^{(q)}(t)\|_{\rm cb}
 \le{}&
 K^{k+q}A_k
 \sum_{h=0}^{\min(k,q)}
 {q\choose h}(k)_h
 t^{k-h}C_2(R_0)^{q-h}.                                    \label{eq:newV69-global-C}
\end{align}
\end{theorem}

\begin{proof}
Apply the Gaussian-centre identity
\Cref{lem:newV62-Gaussian-centre} to
\[
 F_{0,t}(G)=U_0(\sqrt{2t}\,G).
\]
Each of the \(k\) derivatives of \(F_{0,t}\), and each of the
\(k\) linear leaf coefficients, supplies one factor \(\sqrt{2t}\).
Thus, after absorbing \(2^k\) and the leaf matrices into
\(\Psi_k\),
\[
 {\cal C}_k(t)
 =
 t^k\mathbb E\Psi_k(\sqrt{2t}\,G)
 =
 t^k(e^{t\Delta}\Psi_k)(0).
\]
This proves \eqref{eq:newV69-heat-representation}.

The heat equation and the Gaussian representation give
\[
 H_k^{(m)}(t)
 =
 (e^{t\Delta}\Delta^m\Psi_k)(0).
\]
Every summand of \(\Delta^m\Psi_k(z)\) is a scalar-coordinate
derivative of \(U_0(z)\) of total order \(k+2m\).  The first \(k\)
directions are accompanied by the \(k\) leaf generators, and the
remaining \(2m\) directions occur in \(m\) equal-coordinate pairs.
For real \(z\), \(U_0(z)\) is unitary.  Therefore the exact
time-ordered formula \eqref{eq:newV62-Duhamel} bounds each summand
by the product of its generator norms, with no derivative
factorial.  Summing the \(3^{k+m}\) coordinate choices gives
\[
 \|\Delta^m\Psi_k(z)\|_{\rm cb}
 \le K^{k+m}A_kC_2(R_0)^m.
\]
Gaussian averaging is a contraction for the completely bounded
norm, proving \eqref{eq:newV69-global-H}.

Finally differentiate \(t^kH_k(t)\) \(q\) times.  If \(h\)
derivatives hit \(t^k\), Leibniz's rule gives
\({q\choose h}(k)_ht^{k-h}H_k^{(q-h)}(t)\).
Insert \eqref{eq:newV69-global-H} and sum \(h\).
\end{proof}

\begin{assessment}[The unitary gain]
\label{ass:newV69-unitary-gain}
The local power-series proof of V65 bounded \(H_k^{(m)}\) only when
\(tC_2(R_0)\) was small.  Formula
\eqref{eq:newV69-heat-representation} differentiates the Gaussian
heat semigroup instead.  Each clock derivative becomes a Laplacian,
and hence two unitary Duhamel insertions.  Their ordered-simplex
volume cancels their orderings before norms are taken.  This is the
global mechanism anticipated, but not completed, by HCUC in V68.
\end{assessment}

\section{A normalized moment bound for all added jumps}
\label{sec:newV69-moments}

Let \(\pi=\{B_1,\ldots,B_q\}\) partition \([r]\), put
\(p_\nu=|B_\nu|\), and normalize
\[
 \eta_{p_\nu+1}(dx)
 =
 {x^{p_\nu+1}\Pi_\alpha^{\rm sm}(dx)\over M_{p_\nu+1}}
\]
whenever \(M_{p_\nu+1}>0\).  Let \(Y_\nu\) have this law,
independently of the other \(Y\)'s and of \(S_{\rm sm}\).

\begin{lemma}[Joint original-and-added clock moment]
\label{lem:newV69-joint-moment}
Put \(N=k+2r\) and \(L=1+N\sqrt s\).  For
\(0\le d\le k\),
\begin{equation}
 \left\|
 S_{\rm sm}+\sum_{\nu=1}^qY_\nu
 \right\|_{L^d}
 \le KsL,                                                    \label{eq:newV69-joint-moment}
\end{equation}
where the assertion at \(d=0\) means that the zeroth moment is one.
\end{lemma}

\begin{proof}
Apply \Cref{thm:newV67-gamma} with \(p=p_\nu+1\).
The moment generating function of \(\sum_\nu Y_\nu\) is dominated
by that of a Gamma variable with rate \(T_\alpha\) and shape
\[
 \sum_{\nu=1}^q(p_\nu+2)=r+2q.
\]
Chernoff integration, or direct integration of the Gamma majorant,
therefore gives
\[
 \left\|\sum_\nu Y_\nu\right\|_{L^d}
 \le {K(r+2q+d)\over T_\alpha}
 \le KN s^2.
\]
On the other hand,
\Cref{thm:newV62-clock-moments} gives
\[
 \|S_{\rm sm}\|_{L^d}
 \le Ks(1+d\sqrt s).
\]
Since \(d\le k\), \(q\le r\), and \(s\le\sqrt s\), Minkowski's
inequality proves \eqref{eq:newV69-joint-moment}.
\end{proof}

\section{The fixed-partition estimate without a polynomial loss}
\label{sec:newV69-fixed}

For \(a=b_0\) define
\begin{equation}
 {\cal F}_{k,q}(L,a)
 :=
 \sum_{h=0}^{\min(k,q)}
 {q\choose h}(k)_h
 L^{k-h}a^{2(q-h)}.                                         \label{eq:newV69-F}
\end{equation}

\begin{theorem}[Global fixed-partition clock grade]
\label{thm:newV69-fixed}
For every \(k\ge0\), \(r\ge1\), and
\(\pi\in{\rm Part}([r])\),
\begin{align}
 \|{\cal I}_{k,\pi}\|_{\rm cb}
 \le{}&
 K^{k+2r}A_k s^{k+2r}
 \left\{\prod_{\nu=1}^q(p_\nu+1)!\right\}
 {\cal F}_{k,q}(L,b_0),                                     \label{eq:newV69-fixed}
\end{align}
where \({\cal I}_{k,\pi}\) is the complete \(\pi\)-term in
\eqref{eq:newV64-Mecke-cumulant}.  No inner/outer restriction is
present.
\end{theorem}

\begin{proof}
The rectangle formula and \eqref{eq:newV69-global-C} give
\begin{align*}
 &\|\Delta_{x_1}\cdots\Delta_{x_q}{\cal C}_k(S)\|_{\rm cb}\\
 &\quad\le
 K^{k+q}A_k\left(\prod_{\nu=1}^qx_\nu\right)
 \sum_{h=0}^{\min(k,q)}
 {q\choose h}(k)_h
 (S+\sum_\nu x_\nu)^{k-h}
 C_2(R_0)^{q-h}.
\end{align*}
For a fixed \(h\), normalize the measures
\(x_\nu^{p_\nu+1}\Pi(dx_\nu)\) by \(M_{p_\nu+1}\) and apply
\Cref{lem:newV69-joint-moment}.  Also,
\Cref{lem:newV60-jump-moments} gives
\[
 \prod_{\nu=1}^qM_{p_\nu+1}
 \le
 K^{2r+2q}s^{2r+q}
 \prod_{\nu=1}^q(p_\nu+1)!.
\]
Finally
\[
 C_2(R_0)^{q-h}
 \le b_0^{2(q-h)}s^{-(q-h)}.
\]
The total power of \(s\) is exactly
\[
 (2r+q)+(k-h)-(q-h)=k+2r.
\]
The remaining factors are precisely
\eqref{eq:newV69-F}; enlarge the exponential base.
\end{proof}

\section{Exact joint resummation}
\label{sec:newV69-resummation}

Put
\begin{equation}
 {\mathfrak A}(z)
 :=
 \sum_{p\ge1}(p+1)z^p
 =
 {z(2-z)\over(1-z)^2}.                                      \label{eq:newV69-Az}
\end{equation}

\begin{lemma}[Partition and derivative-hit identity]
\label{lem:newV69-joint-EGF}
For every \(r\ge1\),
\begin{align}
 &{1\over r!}
 \sum_{\pi\in{\rm Part}([r])}
 \left\{\prod_{B\in\pi}(|B|+1)!\right\}
 {\cal F}_{k,|\pi|}(L,a)
 \notag\\
 &\qquad=
 [z^r]\,
 \exp\{a^2{\mathfrak A}(z)\}
 \{L+{\mathfrak A}(z)\}^{k}.                                \label{eq:newV69-joint-EGF}
\end{align}
\end{lemma}

\begin{proof}
Partitions with \(q\) blocks have labelled exponential generating
function \({\mathfrak A}(z)^q/q!\).  For fixed \(h\),
\[
 \sum_{q\ge h}
 {{q\choose h}\over q!}
 {\mathfrak A}(z)^q a^{2(q-h)}
 =
 {{\mathfrak A}(z)^h\over h!}
 e^{a^2{\mathfrak A}(z)}.
\]
Since \((k)_h/h!={k\choose h}\), summing \(h\) gives
\[
 e^{a^2{\mathfrak A}(z)}
 \sum_{h=0}^k{k\choose h}
 L^{k-h}{\mathfrak A}(z)^h
 =
 e^{a^2{\mathfrak A}(z)}
 \{L+{\mathfrak A}(z)\}^k.
\]
Extract the coefficient of \(z^r\).
\end{proof}

\begin{lemma}[Thermal Cauchy bound]
\label{lem:newV69-Cauchy}
Assume, as in the thermal source polydisc, that
\[
 \sqrt s\le b_i<1.
\]
Let \(N=k+2r\), \(a=b_0\), \(y=N\sqrt s\), and
\(D=a+y\).  Then
\begin{align}
 &a^k y^{2r}
 [z^r]\,
 e^{a^2{\mathfrak A}(z)}
 \{1+y+{\mathfrak A}(z)\}^{k}
 \notag\\
 &\qquad\le
 K^{N}aD^{N-1}.                                             \label{eq:newV69-Cauchy}
\end{align}
\end{lemma}

\begin{proof}
Fix a sufficiently small numerical \(\delta>0\) and take
\[
 \rho=\delta\left({y\over D}\right)^2.
\]
Then \(0<\rho\le\delta\), and on \(|z|=\rho\),
\[
 |{\mathfrak A}(z)|
 \le {\mathfrak A}(\rho)\le C\rho.
\]
Cauchy's estimate bounds the coefficient by
\[
 \rho^{-r}e^{Ca^2\rho}(1+y+C\rho)^k.                         \label{eq:newV69-Cauchy-mid}
\]
Moreover
\[
 y^{2r}\rho^{-r}=\delta^{-r}D^{2r}.                          \label{eq:newV69-radius-pay}
\]

Suppose first that \(k\ge1\).  Since \(0<a<1\),
\[
 a(1+y)\le a+y=D.
\]
Also \(a,y\le D\), so
\[
 a\rho
 =
 \delta{ay^2\over D^2}
 \le\delta D.
\]
It follows that
\[
 a^{k-1}(1+y+C\rho)^k
 \le K^N D^{k-1};
\]
the single residual factor \(1+y+C\rho\le C(1+N)\) is absorbed
into \(K^N\).  Combine this with
\eqref{eq:newV69-Cauchy-mid}--\eqref{eq:newV69-radius-pay}.

If \(k=0\), the coefficient of \(z^r\) is unchanged when the
constant one is subtracted from the exponential.  Thus
\[
 |e^{a^2{\mathfrak A}(z)}-1|
 \le Ca^2\rho
\]
on the same circle.  After
\eqref{eq:newV69-radius-pay}, the left side is bounded by
\[
 K^N a^2y^2D^{2r-2}
 \le K^N aD^{2r-1},
\]
because \(ay^2/D\le y\le N\), and the residual \(N\) is
exponential-base harmless.  This proves both cases.
\end{proof}

\section{Closure of CMCQC and retirement of HCUC}
\label{sec:newV69-closure}

\begin{theorem}[Global centre-marked clock quotient card]
\label{thm:newV69-CMCQC}
For every \(k,r\ge0\) with \(k+2r\ge1\), after restoring the
\(r\) Gaussian pair
contractions and summing all choices of the \(k\) centre leaves,
all pairings of the remaining \(2r\) leaves, and all clock
partitions, one has
\begin{equation}
 K^{k+2r+1}
 b_0\left(\prod_{\ell=1}^{k+2r}b_\ell\right)
 B^{k+2r-1}.                                                 \label{eq:newV69-CMCQC}
\end{equation}
Thus CMCQC holds globally on the thermal source polydisc.
\end{theorem}

\begin{proof}
The case \(r=0\) is \Cref{thm:newV62-zero-pair}.  Let \(r\ge1\)
and \(N=k+2r\).  For a fixed centre subset and pairing, the
Casimir factors in \(A_k\) and the \(r\) operators
\(\Omega_{ij}\) obey
\[
 A_k\prod_{\{i,j\}}\|\Omega_{ij}\|
 \le
 s^{-(k+r)}b_0^k\prod_{\ell=1}^Nb_\ell.
\]
The number of centre subsets and pairings is
\[
 {N!\over k!\,2^rr!}
 \le {N^{2r}\over2^rr!}.
\]
Insert \Cref{thm:newV69-fixed}.  The factor \(s^{k+2r}\)
then leaves \(s^r\), and the choice census supplies
\[
 N^{2r}s^r=(N\sqrt s)^{2r}=y^{2r}.
\]
Crucially, do not estimate the partition sum separately.  Apply
\Cref{lem:newV69-joint-EGF} with \(a=b_0\), followed by
\Cref{lem:newV69-Cauchy}.  Since every leaf mark is at least
\(\sqrt s\),
\[
 B=b_0+\sum_{\ell=1}^Nb_\ell\ge b_0+N\sqrt s=D.
\]
Substitution gives \eqref{eq:newV69-CMCQC}.
\end{proof}

\begin{corollary}[Full one-complete-centre tangent row]
\label{cor:newV69-one-centre}
Every one-complete-centre cumulant of the subordinated tangent
Gaussian driven by \(S_{\rm sm}\) satisfies the strong thermal
all-arity card with one exponential base.
\end{corollary}

\begin{proof}
Apply \Cref{prop:newV62-CMCQC} to
\Cref{thm:newV69-CMCQC}.
\end{proof}

\begin{correction}[HCUC is unnecessary]
\label{corr:newV69-HCUC}
The high-centre unitary card introduced in
\Cref{def:newV68-HCUC} is retired.  The heat-semigroup identity
\eqref{eq:newV69-heat-representation} is global in
\(tC_2(R_0)\), and \Cref{thm:newV69-CMCQC} treats small and large
centres simultaneously.  The V68 low-centre estimates remain
valid auxiliary tail estimates, but no longer carry the closure
claim.
\end{correction}

\begin{assessment}[What actually cancelled]
\label{ass:newV69-cancellation}
There are three dangerous factors:
\((k)_h\) from differentiating \(t^k\), the number of set
partitions of the outer clock marks, and the number of Gaussian
pairings.  Formula \eqref{eq:newV69-joint-EGF} converts
\((k)_h/h!\) into the binomial coefficient \({k\choose h}\);
the \(h!\) is supplied by the same \(q!\) which is visible only
before the partition and pairing ledgers are separated.  The
adaptive Cauchy radius then converts the sole remaining
\((N\sqrt s)^{2r}\) into \(D^{2r}\).  No arity power is hidden
inside the numerical base.
\end{assessment}

\begin{programme}[V70: compulsory audit and return to the exact
small-clock group source]
\label{prog:newV69-V70}
\begin{enumerate}[leftmargin=3em]
\item Read the newest active SM and NS notes and record immutable
      line, byte, and SHA--256 snapshots before using any lesson.
\item Re-enter the endpoint-factorizing small-clock source of V60;
      do not replace the compact-group endpoint by the tangent
      Gaussian.
\item Classify the first two-complete-centre overlap row and derive
      its exact clock-partition generating function.
\item Test whether the same adaptive Cauchy radius closes the
      multi-centre word card MCWC, including noncommuting centre
      words.
\item Keep the wrapped-Poisson card WPSC separate until the exact
      small-clock nonlinear grade card SCNGC is genuinely closed.
\end{enumerate}
\end{programme}

\begin{firewall}[Claim boundary after V69]
\label{fire:newV69-boundary}
V69 corrects the incomplete V66 summation, proves a global
heat-semigroup representation for every clock derivative of the
complete Gaussian centre, proves the normalized joint-clock moment
bound, derives the exact joint partition--derivative generating
function, and closes CMCQC and the full one-complete-centre
subordinated tangent row with one exponential arity base.

V69 does not prove the multi-centre word card, OCRC in its full
noncommutative form, SCNGC, WPSC, GRLTC, LZTC, PLRC, WSDC, BRSC,
HWSFC, UGCGC, thermal connected WPRAC, all-arity RGSC, the raw
Wilson aggregate strong card, all-order strong MTA, WUFC, CPM,
cutoff removal, reflection positivity, Osterwalder--Schrader
reconstruction, a continuum Yang--Mills measure, or a positive
Hamiltonian gap.  The full Yang--Mills mass gap has not been proved.
\end{firewall}

\begin{theorem}[V69 result ledger]
\label{thm:newV69-results}
V69:
\begin{enumerate}[label=\textup{(V69.\arabic*)},leftmargin=6.5em]
\item identifies and corrects the missing arity factor in V66;
\item proves the global Gaussian heat representation of the centre;
\item proves factorial-free global clock derivatives;
\item proves a normalized original-and-added clock moment bound;
\item proves the sharp fixed-partition clock grade;
\item derives the exact joint partition--derivative EGF;
\item proves the adaptive thermal Cauchy estimate;
\item closes every nontrivial CMCQC row; and
\item closes the full one-complete-centre subordinated tangent row.
\end{enumerate}
\end{theorem}

\begin{proof}
Use
\Cref{corr:newV69-V66,
thm:newV69-heat-representation,
lem:newV69-joint-moment,
thm:newV69-fixed,
lem:newV69-joint-EGF,
lem:newV69-Cauchy,
thm:newV69-CMCQC,
cor:newV69-one-centre,
corr:newV69-HCUC}.
\end{proof}

\paragraph{Parent-version record.}
V69 was formed from
\texttt{Renewal\_Yang\_Mills\_Mass\_Gap\_Research\_Notes\_V68.tex}.
The V68 parent had \(32\,449\) logical source lines,
\(1\,031\,935\) bytes, and SHA--256
\[
 \texttt{1581b543565f838a8438cbc50d080cb38704598ed3c6b2f6c63bf631aa0c706b}.
\]
The next compulsory companion audit is V70.

% =============================================================================
% END APPENDED FIFTH-SERIES V69 MATERIAL
% =============================================================================

\clearpage
\chapter{Fifth-series V70: compulsory audit and the acyclic
two-centre incidence row}
\label{ch:newV70-two-centre}

\section{Compulsory companion audit}
\label{sec:newV70-audit}

\begin{audit}[V70 immutable companion snapshots]
\label{audit:newV70-companions}
The newest active companion sources were read without modification.
During the audit the parallel NS programme replaced V29 by V30, so
the inventory was refreshed and the final active source was read.
The snapshots used below are
\begin{align*}
 &\texttt{Renewal\_SM\_Einstein\_Continuum\_Research\_Notes\_V94.tex},\\
 &\qquad 32\,941\ \hbox{logical lines},\quad
 1\,264\,229\ \hbox{bytes},\\
 &\qquad
 \texttt{87ce0ae7bdd5645bc5cb1579aab1511c4a2e52c59d984c1c9bfaa8d6ac771003},
 \\
 &\texttt{Renewal\_NS\_Stability\_Research\_Notes\_V30.tex},\\
 &\qquad 7\,557\ \hbox{logical lines},\quad
 384\,630\ \hbox{bytes},\\
 &\qquad
 \texttt{4035533fc8d9a4264bf27f72e4d243f8f971bb39ed22f65c34d78b3662065e13}.
\end{align*}

SM V94 expands the Wilson kernel square into an exact free positive
part and finitely many covariant-commutator remainders.  Only after
that algebraic decomposition are plaquette norms inserted.  This
proves a gap on the declared admissible class while keeping
measure concentration, the Yukawa floor, and the determinant phase
as separate gates.

NS V30 consolidates the rate-free blow-up ledger.  Its normalized
profile equation makes the Rayleigh quotient exact, but it also
records that division by a vanishing Gaussian norm amplifies the
defect and does not yield compactness or a limiting eigenfunction.
Exact normalization is therefore informative without being a
uniform estimate.
\end{audit}

\begin{assessment}[Audit transfer to the YM programme]
\label{ass:newV70-audit-transfer}
No SM lattice constant and no NS analytic estimate is imported into
Yang--Mills.  Two proof disciplines are used:
\begin{enumerate}[label=\textup{(\roman*)},leftmargin=3.5em]
\item keep the complete ordered operator word until its positive or
      commutator structure has been isolated; and
\item do not normalize each overlap activity separately when the
      normalizing quantity can be small.
\end{enumerate}
Accordingly, the two-centre calculation below first contracts the
complete Duhamel word and then sums the unnormalized star activities.
\end{assessment}

\section{Stable ordered-simplex words}
\label{sec:newV70-Duhamel}

\begin{lemma}[Dissipative Duhamel word contraction]
\label{lem:newV70-Duhamel}
Let \(V=V^*\) and suppose \(V\preceq\varepsilon I\).  For arbitrary
operators \(A_1,\ldots,A_p\),
\begin{equation}
 \|D^pe^V[A_1,\ldots,A_p]\|_{\rm cb}
 \le
 e^\varepsilon\prod_{\nu=1}^p\|A_\nu\|_{\rm cb}.             \label{eq:newV70-Duhamel}
\end{equation}
The constant contains no \(p!\).
\end{lemma}

\begin{proof}
The exact Fr\'echet formula is
\[
 D^pe^V[A_1,\ldots,A_p]
 =
 \sum_{\sigma\in\mathfrak S_p}
 \int_{\Delta_p}
 e^{t_0V}A_{\sigma(1)}e^{t_1V}\cdots
 A_{\sigma(p)}e^{t_pV}\,d\boldsymbol t.
\]
Since \(\sum_jt_j=1\),
\[
 \prod_{j=0}^p\|e^{t_jV}\|
 \le e^{\varepsilon\sum_jt_j}=e^\varepsilon.
\]
There are \(p!\) orders and
\({\rm vol}(\Delta_p)=1/p!\).  Take norms after this cancellation.
The argument is unchanged by matrix amplification.
\end{proof}

\begin{corollary}[Application on the small-clock BKAR matrices]
\label{cor:newV70-Duhamel-BKAR}
There is a numerical \(b_*>0\) such that, whenever
\(B=\sum_i b_i\le b_*\), every ordered word obtained by
differentiating
\({\cal F}_m(\boldsymbol x)=e^{{\cal V}_m(\boldsymbol x)}\)
on a BKAR forest matrix obeys
\begin{equation}
 \|D^pe^{{\cal V}_m}[A_1,\ldots,A_p]\|_{\rm cb}
 \le C\prod_{\nu=1}^p\|A_\nu\|_{\rm cb}.                    \label{eq:newV70-Duhamel-BKAR}
\end{equation}
\end{corollary}

\begin{proof}
The proof of
\Cref{lem:newV60-interaction-stability} gives
\({\cal V}_m\preceq cB^2I\) on every BKAR forest matrix.
Apply \Cref{lem:newV70-Duhamel} with
\(\varepsilon=cB^2\) and decrease \(b_*\) if necessary.
\end{proof}

\section{The two-centre incidence class}
\label{sec:newV70-incidence}

\begin{definition}[Active label--root incidence graph]
\label{def:newV70-incidence}
For a term in a mixed derivative of
\({\cal F}_m\), its active incidence graph is the bipartite graph
whose vertices are the representation labels and the occurrences
of nonsingleton root activities \(R_A\), with an edge \(iA\) exactly
when \(i\in A\).  A label is a nonlinear centre when it meets at
least two active roots.
\end{definition}

\begin{definition}[Two-centre dumbbell incidence class, TDIC]
\label{def:newV70-TDIC}
TDIC consists of connected acyclic active incidence graphs having
exactly two nonlinear centres.  The complete TDIC contribution
includes every Duhamel ordering and every partition of the
noncentral labels among the active roots.
\end{definition}

\begin{lemma}[Exact structure of a TDIC graph]
\label{lem:newV70-TDIC-structure}
Let the two nonlinear centres be \(c_1,c_2\).  Then:
\begin{enumerate}[label=\textup{(\roman*)},leftmargin=3.5em]
\item there is a unique active root \(A_*\) containing both
      \(c_1\) and \(c_2\);
\item every other active root contains exactly one of
      \(c_1,c_2\); and
\item if
      \(L_*=A_*\setminus\{c_1,c_2\}\), the noncentral label sets of
      roots attached to \(c_i\) form a partition of a set \(L_i\),
      and \(L_*,L_1,L_2\) are pairwise disjoint.
\end{enumerate}
\end{lemma}

\begin{proof}
The unique path between \(c_1\) and \(c_2\) cannot contain an
intermediate label: such a label has degree at least two and would
be a third nonlinear centre.  Hence one root \(A_*\) is adjacent to
both centres.  A second such root would create a four-cycle, proving
uniqueness.

Any other root must be connected to the graph.  If it contained
neither centre, the path joining it to \(A_*\) would contain an
intermediate label of degree at least two.  If it contained both
centres, uniqueness would fail.  Thus it contains exactly one
centre.  Finally a noncentral label has degree one; consequently it
occurs in exactly one active root, proving all disjointness and
partition assertions.
\end{proof}

\section{Acyclic two-centre closure}
\label{sec:newV70-closure}

\begin{theorem}[TDIC strong aggregate card]
\label{thm:newV70-TDIC}
On the difficult small-total-mark region \(B\le b_*\), the complete
positive majorant of all TDIC terms on \(m\) labels is bounded by
\begin{equation}
 K^m\left(\prod_{i=1}^mb_i\right)B^{m-2}.                   \label{eq:newV70-TDIC}
\end{equation}
In particular, Duhamel orderings and the two independent
centre-star partition families create no Bell or word factorial.
\end{theorem}

\begin{proof}
First fix \(c_1,c_2\) and a decomposition
\[
 [m]\setminus\{c_1,c_2\}=L_*\mathbin{\dot\cup}L_1
 \mathbin{\dot\cup}L_2
\]
as in \Cref{lem:newV70-TDIC-structure}.  Put
\(n_i=|L_i|\) and \(\ell=|L_*|\).  Since \(c_i\) is nonlinear,
the attached family at \(c_i\) is nonempty; hence \(n_i\ge1\).

If \(\ell\ge1\), the shared root equals \(Q_{A_*}\) and is bounded
by the relaxed same-jump card
\Cref{thm:newV60-hyperedge-card}.  If \(\ell=0\), it is
\(R_{\{c_1,c_2\}}=Q_{\{c_1,c_2\}}-2d_\alpha\Omega_{c_1c_2}\);
the \(a=2\) case in the proof of
\Cref{lem:newV61-reserve} gives the same pair bound.  Thus in both
cases
\begin{equation}
 \|R_{A_*}\|_{\rm cb}
 \le
 K^{\ell+2}
 b_{c_1}b_{c_2}
 \left(\prod_{j\in L_*}b_j\right)B^\ell.                   \label{eq:newV70-shared}
\end{equation}
For the roots attached only at \(c_i\), their noncentral label sets
partition \(L_i\).  Apply
\Cref{cor:newV61-Bell-closed} with centre weight \(b_{c_i}\):
\begin{equation}
 \sum_{\mathcal P_i\in{\rm Part}(L_i)}
 \prod_{C\in\mathcal P_i}
 \|R_{\{c_i\}\cup C}\|_{\rm cb}
 \le
 K^{n_i+1}b_{c_i}
 \left(\prod_{j\in L_i}b_j\right)B^{n_i-1}.                 \label{eq:newV70-star-i}
\end{equation}
The replacement of each local sum of marks by the global \(B\)
only enlarges the right side.

By \Cref{cor:newV70-Duhamel-BKAR}, all orderings of the shared root
and the two star families cost only one numerical factor times the
product of
\eqref{eq:newV70-shared} and the two instances of
\eqref{eq:newV70-star-i}.  That product is at most
\begin{align*}
 K^m
 b_{c_1}^2b_{c_2}^2
 \left(\prod_{j\ne c_1,c_2}b_j\right)
 B^{\ell+n_1+n_2-2}.
\end{align*}
Use
\[
 b_{c_1}b_{c_2}\le B^2
\]
on the repeated centre factors.  Since
\(m=2+\ell+n_1+n_2\), the result is
\[
 K^m\left(\prod_{i=1}^mb_i\right)B^{m-2}.
\]

There are at most \(m^2\) ordered choices of the centres and at most
\(3^{m-2}\) decompositions into \(L_*,L_1,L_2\).  Both are absorbed
into a larger \(K^m\).  This proves
\eqref{eq:newV70-TDIC}.
\end{proof}

\begin{corollary}[The first genuinely multi-centre row is closed]
\label{cor:newV70-TDIC-closed}
The connected acyclic MCWC sector with exactly two nonlinear centres
obeys the strong aggregate card in the small-\(B\) regime.
\end{corollary}

\begin{proof}
\Cref{lem:newV70-TDIC-structure} exhausts that sector, and
\Cref{thm:newV70-TDIC} bounds its complete positive majorant.
\end{proof}

\begin{assessment}[Where noncommutativity first remains]
\label{ass:newV70-cycle}
The ordered matrix words themselves are not the obstruction in TDIC:
\Cref{lem:newV70-Duhamel} sums their permutation orders by simplex
volume.  Acyclic incidence also is not the obstruction: the unique
shared root reduces the remaining overlap to two Bell-star sums.
The first unresolved geometry is therefore an incidence cycle.  On a
cycle, two or more shared roots meet the same two centres, so there is
no unique root at which to cut the graph, and multiplying two relaxed
root cards repeats both centre marks without a spare \(B^2\) for
every repetition.
\end{assessment}

\begin{programme}[V71: the first incidence cycle]
\label{prog:newV70-V71}
\begin{enumerate}[leftmargin=3em]
\item Treat the four-cycle with centres \(c_1,c_2\) and two shared
      roots as one ordered convolution word, not as a product of two
      relaxed root norms.
\item Expand only the commutator difference between the two possible
      shared-root orders and retain the symmetric word intact.
\item Test whether the unrelaxed clock factor
      \(s^{(|A|-2)/2}\) pays the repeated centre marks after the
      cycle is cut once.
\item Formulate a cycle-opening identity on the bipartite incidence
      graph which reduces every unicyclic word to one commutator
      insertion plus an acyclic TDIC word.
\item Do not promote the two-centre result to SCNGC until all
      incidence cycles and arbitrarily many centres are controlled.
\end{enumerate}
\end{programme}

\begin{firewall}[Claim boundary after V70]
\label{fire:newV70-boundary}
V70 performs the compulsory SM/NS audit, proves a factorial-free
Duhamel contraction for stable self-adjoint source words, classifies
the connected acyclic incidence graphs with exactly two nonlinear
centres, and proves their complete strong aggregate card in the
small-total-mark regime.

V70 does not prove the unicyclic two-centre row, MCWC in full, OCRC
in its full noncommutative form, SCNGC, WPSC, GRLTC, LZTC, PLRC,
WSDC, BRSC, HWSFC, UGCGC, thermal connected WPRAC, all-arity RGSC,
the raw Wilson aggregate strong card, all-order strong MTA, WUFC,
CPM, cutoff removal, reflection positivity,
Osterwalder--Schrader reconstruction, a continuum Yang--Mills
measure, or a positive Hamiltonian gap.  The full Yang--Mills mass
gap has not been proved.
\end{firewall}

\begin{theorem}[V70 result ledger]
\label{thm:newV70-results}
V70:
\begin{enumerate}[label=\textup{(V70.\arabic*)},leftmargin=6.5em]
\item performs the compulsory active-SM/active-NS audit;
\item imports only two proof-order disciplines from that audit;
\item proves stable dissipative Duhamel word contraction;
\item defines the active label--root incidence graph;
\item classifies every acyclic graph with two nonlinear centres;
\item reduces it to one shared root and two disjoint star families;
\item closes both star families by the V61 generating function;
\item proves the TDIC strong aggregate card; and
\item isolates the first incidence cycle as the next obstruction.
\end{enumerate}
\end{theorem}

\begin{proof}
Use
\Cref{audit:newV70-companions,
lem:newV70-Duhamel,
cor:newV70-Duhamel-BKAR,
lem:newV70-TDIC-structure,
thm:newV70-TDIC,
cor:newV70-TDIC-closed}.
\end{proof}

\paragraph{Parent-version record.}
V70 was formed from
\texttt{Renewal\_Yang\_Mills\_Mass\_Gap\_Research\_Notes\_V69.tex}.
The V69 parent had \(32\,997\) logical source lines,
\(1\,048\,936\) bytes, and SHA--256
\[
 \texttt{656d1ad027b31bc4d5eb78dc11a51c2b13d2f4ea397355d4671aaa9773f0d49e}.
\]
The next compulsory companion audit is V75.

% =============================================================================
% END APPENDED FIFTH-SERIES V70 MATERIAL
% =============================================================================

\clearpage
\chapter{Fifth-series V71: double-centre Bell resummation and
closure of the complete two-centre incidence sector}
\label{ch:newV71-double-centre}

\section{Why the first incidence cycle is still summable}
\label{sec:newV71-entry}

No companion audit is due at V71.

V70 closed the acyclic two-centre class by cutting at its unique
root shared by both nonlinear centres.  If several roots are shared,
the incidence graph contains cycles and there is no unique cut.
Nevertheless, the shared roots have disjoint noncentral label sets:
otherwise a repeated noncentral label would itself be a third
nonlinear centre.  They therefore form a labelled set partition,
and the unrelaxed root estimate of V61 supplies one more half-power
of \(s\) per noncentral label than the one-centre Bell family.

\section{The double-centre partition lemma}
\label{sec:newV71-partition}

\begin{lemma}[Double-centre Bell resummation]
\label{lem:newV71-double-Bell}
Let two centres have weights
\(\beta_1,\beta_2\in[\sqrt s,1)\), let
\(c_1,\ldots,c_n\in[\sqrt s,1)\), \(n\ge1\), and put
\[
 B=\beta_1+\beta_2+\sum_{j=1}^nc_j,
 \qquad \beta=\beta_1\beta_2.
\]
Suppose that every nonempty \(C\subseteq[n]\) has an activity
satisfying
\begin{equation}
 a(C)
 \le
 \{K(|C|+2)\}^{|C|+2}
 s^{|C|/2}\,
 \beta\prod_{j\in C}c_j.                                   \label{eq:newV71-double-activity}
\end{equation}
Then
\begin{equation}
 \sum_{\mathcal P\in{\rm Part}([n])}
 \prod_{C\in\mathcal P}a(C)
 \le
 K_1^{n+2}\,
 \beta\left(\prod_{j=1}^nc_j\right)B^n.                    \label{eq:newV71-double-Bell}
\end{equation}
\end{lemma}

\begin{proof}
Factor \(\prod_jc_j\).  For a block of size \(r\), put
\[
 A_r=\{K(r+2)\}^{r+2}\beta s^{r/2}.
\]
The labelled exponential formula gives
\[
 \sum_{\mathcal P}\prod_{C\in\mathcal P}A_{|C|}
 =
 n![z^n]\{\exp H(z)-1\},
 \qquad
 H(z)=\sum_{r\ge1}{A_r\over r!}z^r.                         \label{eq:newV71-double-EGF}
\]
Stirling's lower bound gives
\[
 {A_r\over r!}
 \le
 \beta K_0^r(r+2)^2s^{r/2}.
\]
Choose a sufficiently small numerical \(\eta>0\) and use the
Cauchy radius
\[
 R={\eta n\over B}.
\]
The leaf floor gives \(n\sqrt s\le B\), and therefore
\(K_0\sqrt sR\le K_0\eta\).  On \(|z|=R\),
\[
 |H(z)|
 \le
 C\beta\sqrt sR
 \sum_{r\ge1}(r+2)^2(K_0\sqrt sR)^{r-1}
 \le C\beta\sqrt sR.
\]
After decreasing \(\eta\), this is bounded by a numerical constant,
so
\[
 |\exp H(z)-1|\le C\beta\sqrt sR.
\]
Cauchy's estimate and \(n!\le n^n\) now give
\begin{align*}
 n!|[z^n](e^H-1)|
 &\le
 C\beta\sqrt s\,n!R^{1-n}\\
 &\le
 K_1^n\beta\,n\sqrt s\,B^{n-1}\\
 &\le
 K_1^n\beta B^n.
\end{align*}
Restore \(\prod_jc_j\).
\end{proof}

\begin{corollary}[All nonempty roots shared by two centres]
\label{cor:newV71-shared-family}
Fix centres \(c_1,c_2\) and a noncentral label set \(L\).
Sum every family of active roots
\[
 \{c_1,c_2\}\cup C,
 \qquad C\in\mathcal P,
\]
where \(\mathcal P\) ranges over all partitions of \(L\).
The complete positive majorant is at most
\begin{equation}
 K^{|L|+2}
 b_{c_1}b_{c_2}
 \left(\prod_{j\in L}b_j\right)B^{|L|}.                    \label{eq:newV71-shared-family}
\end{equation}
The same estimate holds if the optional bare pair root
\(R_{\{c_1,c_2\}}\) is also present, and for \(L=\varnothing\)
the family consists of that bare root alone.
\end{corollary}

\begin{proof}
For \(C\ne\varnothing\), apply
\Cref{lem:newV61-reserve} to the root of size \(|C|+2\):
\[
 \|R_{\{c_1,c_2\}\cup C}\|_{\rm cb}
 \le
 \{K(|C|+2)\}^{|C|+2}
 s^{|C|/2}b_{c_1}b_{c_2}\prod_{j\in C}b_j.
\]
This is \eqref{eq:newV71-double-activity}; apply
\Cref{lem:newV71-double-Bell}.  The \(a=2\) case of
\Cref{lem:newV61-reserve} bounds the optional bare root by
\(K^2b_{c_1}b_{c_2}\).  Since
\(b_{c_1}b_{c_2}<1\), adjoining it only changes the numerical base.
The same two-label estimate proves the case \(L=\varnothing\).
\end{proof}

\begin{assessment}[The additional clock half-power]
\label{ass:newV71-half-power}
For a root with one centre and \(r\) leaves, V61's activity contains
\(s^{(r-1)/2}\); extracting the nonempty source factor leaves the
single centre mark.  For a root with two centres and \(r\) leaves,
the power is \(s^{r/2}\).  The factor \(n\sqrt s\le B\) in the last
line of the proof is exactly what upgrades the shared family from
\(B^{n-1}\) to \(B^n\), the correct homogeneity for two retained
centres.
\end{assessment}

\section{Classification with two nonlinear centres}
\label{sec:newV71-classification}

\begin{lemma}[General two-centre incidence structure]
\label{lem:newV71-two-structure}
Let a connected active incidence graph have exactly two nonlinear
centre labels \(c_1,c_2\).  Then:
\begin{enumerate}[label=\textup{(\roman*)},leftmargin=3.5em]
\item every active root contains at least one centre;
\item the roots containing both centres form a nonempty family;
\item apart from a possible bare pair root, their noncentral label
      sets form a partition of a set \(L_*\);
\item roots containing \(c_i\) but not the other centre have
      noncentral label sets forming a partition of a set \(L_i\);
      and
\item \(L_*,L_1,L_2\) are pairwise disjoint and exhaust all
      noncentral labels.
\end{enumerate}
Conversely, every such collection has connected incidence graph and
exactly the two declared nonlinear centres, provided each declared
centre meets at least two roots.
\end{lemma}

\begin{proof}
A noncentral label has incidence degree one.  Hence a root containing
neither \(c_1\) nor \(c_2\) could not be connected to the rest of the
graph: any path from it to a centre would first pass through a
noncentral label of degree at least two.  This proves (i).

If no root contained both centres, a path between them would contain
an intermediate label of degree at least two, contrary to the
assumption that there are only two nonlinear centres.  Thus the
shared family is nonempty.  Two different roots cannot share a
noncentral label, since that label would have degree at least two.
This proves all partition and disjointness statements.  The converse
is immediate from the construction.
\end{proof}

\section{Closure of the complete two-centre sector}
\label{sec:newV71-closure}

\begin{theorem}[Complete two-centre MCWC sector]
\label{thm:newV71-two-centre}
On \(B\le b_*\), the complete positive majorant of every connected
MCWC incidence structure with exactly two nonlinear centres is
bounded by
\begin{equation}
 K^m\left(\prod_{i=1}^mb_i\right)B^{m-2}.                   \label{eq:newV71-two-centre}
\end{equation}
Arbitrarily many roots shared by the two centres and every resulting
incidence cycle are included.
\end{theorem}

\begin{proof}
Fix the two centres and the decomposition
\[
 [m]\setminus\{c_1,c_2\}
 =
 L_*\mathbin{\dot\cup}L_1\mathbin{\dot\cup}L_2
\]
from \Cref{lem:newV71-two-structure}.  Write
\(n_*=|L_*|\), \(n_i=|L_i|\), and
\(\delta_i=\mathbf1_{\{n_i>0\}}\).

The full family of shared roots is bounded by
\Cref{cor:newV71-shared-family}:
\[
 K^{n_*+2}b_{c_1}b_{c_2}
 \left(\prod_{j\in L_*}b_j\right)B^{n_*}.
\]
If \(n_i>0\), the roots containing only the centre \(c_i\) are
bounded, after their complete partition sum, by
\Cref{cor:newV61-Bell-closed}:
\[
 K^{n_i+1}b_{c_i}
 \left(\prod_{j\in L_i}b_j\right)B^{n_i-1}.
\]
For \(n_i=0\), insert the factor one.  All Duhamel orderings of
these root insertions are bounded by
\Cref{cor:newV70-Duhamel-BKAR} without an ordering factorial.

After the common factor \(\prod_{i=1}^mb_i\) is removed, the only
repeated centre marks are
\[
 b_{c_1}^{\delta_1}b_{c_2}^{\delta_2}
 \le B^{\delta_1+\delta_2}.
\]
The total power of \(B\) is therefore
\[
 n_*+(n_1-\delta_1)+(n_2-\delta_2)
 +\delta_1+\delta_2
 =
 n_*+n_1+n_2=m-2.
\]
Finally sum the at most \(m^2\) centre choices and
\(3^{m-2}\) decompositions of the noncentral labels.  These change
only the exponential base and prove
\eqref{eq:newV71-two-centre}.
\end{proof}

\begin{corollary}[Two-centre incidence cycles are retired]
\label{cor:newV71-two-retired}
Within the small-\(B\) source, the first unresolved MCWC sector must
have at least three nonlinear centres.  A cycle involving only two
centres is already contained in
\Cref{thm:newV71-two-centre}.
\end{corollary}

\begin{proof}
V70 treats the acyclic subfamily, and
\Cref{thm:newV71-two-centre} treats every number of shared roots.
\end{proof}

\section{The next obstruction: a core of three or more centres}
\label{sec:newV71-core}

\begin{definition}[Nonlinear-centre core]
\label{def:newV71-core}
Contract every degree-one noncentral label in an active incidence
graph into the root which contains it.  The remaining bipartite
multihypergraph on the nonlinear centre labels and active roots is
the nonlinear-centre core.  V61 closes a core with one centre and
V71 closes every connected core with two centres.
\end{definition}

\begin{assessment}[Why three centres are different]
\label{ass:newV71-three}
With three centres, a root may meet one, two, or all three centres.
Pairwise shared-root families need not have disjoint leaf sets at the
level of centre labels, and the core can contain a triangle of three
pair-shared families.  Applying the two-centre card independently on
its three edges repeats every centre twice.  Those repeated marks
suggest the weighted spanning-tree polynomial on the centre core,
but taking all three pair families independently also retains the
cycle.  The next proof must perform a Penrose tree allocation on the
centre core before multiplying the two-centre Bell bounds.
\end{assessment}

\begin{programme}[V72: centre-core tree allocation]
\label{prog:newV71-V72}
\begin{enumerate}[leftmargin=3em]
\item Prove a weighted Penrose allocation for a connected
      centre core whose hyperedge activities have already been
      resummed over their degree-one leaves.
\item Treat the three-centre triangle explicitly and identify the
      factor which deletes one of its three pair-shared families
      from the spanning-tree payment.
\item Retain roots meeting all three centres as hyperedges; do not
      split them into three pair activities.
\item Determine whether the unrelaxed reserve supplies the
      tree-excess factor for every core cycle.
\item If the centre-core tree allocation closes, combine it with
      V61 and V71 to attack arbitrary-centre MCWC.
\end{enumerate}
\end{programme}

\begin{firewall}[Claim boundary after V71]
\label{fire:newV71-boundary}
V71 proves the double-centre Bell generating-function estimate,
classifies every connected active incidence structure with exactly
two nonlinear centres, and closes its complete positive majorant,
including all two-centre incidence cycles, in the small-total-mark
regime.

V71 does not prove the three-or-more-centre core card, MCWC in full,
OCRC in its full noncommutative form, SCNGC, WPSC, GRLTC, LZTC,
PLRC, WSDC, BRSC, HWSFC, UGCGC, thermal connected WPRAC,
all-arity RGSC, the raw Wilson aggregate strong card, all-order
strong MTA, WUFC, CPM, cutoff removal, reflection positivity,
Osterwalder--Schrader reconstruction, a continuum Yang--Mills
measure, or a positive Hamiltonian gap.  The full Yang--Mills mass
gap has not been proved.
\end{firewall}

\begin{theorem}[V71 result ledger]
\label{thm:newV71-results}
V71:
\begin{enumerate}[label=\textup{(V71.\arabic*)},leftmargin=6.5em]
\item proves the double-centre Bell resummation;
\item obtains the correct extra \(B\)-power from the clock floor;
\item closes every family of roots shared by two centres;
\item classifies all connected two-centre incidence graphs;
\item combines shared and one-centre root families without
      repeated-centre loss;
\item closes every two-centre incidence cycle;
\item closes the complete exactly-two-centre MCWC sector; and
\item isolates the centre core with at least three vertices.
\end{enumerate}
\end{theorem}

\begin{proof}
Use
\Cref{lem:newV71-double-Bell,
cor:newV71-shared-family,
lem:newV71-two-structure,
thm:newV71-two-centre,
cor:newV71-two-retired}.
\end{proof}

\paragraph{Parent-version record.}
V71 was formed from
\texttt{Renewal\_Yang\_Mills\_Mass\_Gap\_Research\_Notes\_V70.tex}.
The V70 parent had \(33\,342\) logical source lines,
\(1\,061\,784\) bytes, and SHA--256
\[
 \texttt{2ce7f53627fb326fb66fcb394020698886be4002e9c668baf60ed5a7dcab61d9}.
\]
The next compulsory companion audit is V75.

% =============================================================================
% END APPENDED FIFTH-SERIES V71 MATERIAL
% =============================================================================

\clearpage
\chapter{Fifth-series V72: effective pair edges and the
connected centre-core tree bound}
\label{ch:newV72-pair-core}

\section{The pair-core sector}
\label{sec:newV72-entry}

No companion audit is due at V72.

\begin{definition}[Pair-core incidence class, PCIC]
\label{def:newV72-PCIC}
Fix a set \({\cal C}\) of \(p\ge2\) nonlinear centre labels.
PCIC consists of the connected active incidence structures for
which every root meeting at least two centres meets exactly two.
Roots meeting one centre are allowed as decorations, and every
noncentral label has incidence degree one.  All root multiplicities,
all partitions of the noncentral labels, and all Duhamel orderings
are included.
\end{definition}

The last condition is automatic once \({\cal C}\) is the complete
set of nonlinear centres.  It is displayed because it permits one
labelled exponential generating function to assign every remaining
label exactly once.

\section{A connected graph bound}
\label{sec:newV72-graph}

\begin{lemma}[Connected graphs are paid by one spanning tree]
\label{lem:newV72-connected-graph}
Let \(w_{ij}\ge0\) for \(1\le i<j\le p\).  Then
\begin{equation}
 \sum_{\substack{G\ {\rm connected}\\V(G)=[p]}}
 \prod_{ij\in E(G)}w_{ij}
 \le
 \exp\left\{\sum_{i<j}w_{ij}\right\}
 \sum_{T\in\mathfrak T_p}\prod_{ij\in E(T)}w_{ij}.           \label{eq:newV72-connected-graph}
\end{equation}
\end{lemma}

\begin{proof}
Assign to every connected graph one of its spanning trees.  For a
fixed tree \(T\), enlarge the assigned family to all graphs
containing \(T\).  Its total weight is
\[
 \left(\prod_{e\in T}w_e\right)
 \prod_{e\notin T}(1+w_e)
 \le
 \left(\prod_{e\in T}w_e\right)
 \exp\left\{\sum_ew_e\right\}.
\]
Sum \(T\).  The enlargement may count a graph more than once, which
is in the useful direction.
\end{proof}

\begin{corollary}[Weighted pair-core tree]
\label{cor:newV72-weighted-tree}
If \(w_{ij}\le Kb_ib_j\) and
\(\sum_{i\in{\cal C}}b_i=B_{\cal C}\), then, provided
\(\sum_{i<j}w_{ij}\le C\),
\begin{equation}
 \sum_{G\ {\rm connected}}\prod_{ij\in E(G)}w_{ij}
 \le
 K_1^p
 \left(\prod_{i\in{\cal C}}b_i\right)
 B_{\cal C}^{p-2}.                                         \label{eq:newV72-weighted-tree}
\end{equation}
\end{corollary}

\begin{proof}
Use \Cref{lem:newV72-connected-graph} and the weighted
Cayley--Pr\"ufer identity
\[
 \sum_{T\in\mathfrak T_p}
 \prod_{ij\in E(T)}b_ib_j
 =
 \left(\prod_{i\in{\cal C}}b_i\right)
 B_{\cal C}^{p-2}.
\]
\end{proof}

\section{All roots supported on one centre pair}
\label{sec:newV72-effective-edge}

Let \(e=\{i,j\}\subset{\cal C}\).  After the marks of its
noncentral labels have been factored out, define the majorant
\begin{align}
 H_e(z)
 &:=
 \sum_{r\ge1}
 { \{K(r+2)\}^{r+2}s^{r/2}b_ib_j\over r!}\,z^r,             \label{eq:newV72-He}\\
 w_e^{(0)}
 &:=K^2b_ib_j,                                               \label{eq:newV72-w0}\\
 W_e(z)
 &:=
 (1+w_e^{(0)})e^{H_e(z)}-1.                                 \label{eq:newV72-We}
\end{align}
The constant term represents the optional bare root
\(R_{\{i,j\}}\); the exponential partitions the labels assigned to
the pair among arbitrarily many nonempty shared roots.

\begin{lemma}[Effective pair edge on the adaptive circle]
\label{lem:newV72-effective-edge}
Let \(n\ge1\), let \(B\) be the sum of all centre and noncentral
marks, and put
\[
 R={\eta n\over B}
\]
with a sufficiently small numerical \(\eta>0\).  On \(|z|=R\),
\begin{equation}
 |H_e(z)|\le Cb_ib_j\sqrt sR,
 \qquad
 |W_e(z)|\le K_0b_ib_j.                                    \label{eq:newV72-effective-edge}
\end{equation}
The same second bound holds at \(z=0\) when \(n=0\).
\end{lemma}

\begin{proof}
The proof of \Cref{lem:newV71-double-Bell} gives
\[
 { \{K(r+2)\}^{r+2}s^{r/2}b_ib_j\over r!}
 \le
 b_ib_jK_1^r(r+2)^2s^{r/2}.
\]
Since the \(n\) noncentral marks are each at least \(\sqrt s\),
\(n\sqrt s\le B\), and hence \(K_1\sqrt sR\le K_1\eta\).
Summing the geometric polynomial series proves the first bound.

Now
\[
 |W_e|
 \le
 w_e^{(0)}e^{|H_e|}+e^{|H_e|}-1.
\]
The first term is at most \(Kb_ib_j\).  The second is at most
\(C|H_e|\), and \(\sqrt sR\le\eta\); enlarge \(K_0\).
At \(z=0\), \(H_e(0)=0\) and
\(W_e(0)=w_e^{(0)}\).
\end{proof}

\begin{assessment}[Why the pair allocation is simultaneous]
\label{ass:newV72-simultaneous}
Assigning each noncentral label to one of \(O(p^2)\) centre pairs
before estimating would cost \(p^{2n}\).  In
\eqref{eq:newV72-We}, multiplication of the edge generating
functions and one final coefficient extraction performs exactly
the same labelled allocation with its multinomial denominators
still present.
\end{assessment}

\section{One-centre decorations}
\label{sec:newV72-decorations}

For \(i\in{\cal C}\), define
\begin{equation}
 H_i^{(1)}(z)
 :=
 \sum_{r\ge1}
 { \{K(r+1)\}^{r+1}s^{(r-1)/2}b_i\over r!}\,z^r.            \label{eq:newV72-H1}
\end{equation}
This is the V61 Bell-star majorant for all roots which meet \(i\)
and no other centre.

\begin{lemma}[All vertex decorations on the adaptive circle]
\label{lem:newV72-decorations}
With the radius of
\Cref{lem:newV72-effective-edge},
\begin{equation}
 \left|
 \exp\left\{\sum_{i\in{\cal C}}H_i^{(1)}(z)\right\}
 \right|
 \le e^{C\eta n}.                                           \label{eq:newV72-decorations}
\end{equation}
\end{lemma}

\begin{proof}
The coefficient estimate in the proof of
\Cref{lem:newV61-Bell-star} gives, on \(|z|=R\),
\[
 |H_i^{(1)}(z)|\le Cb_iR.
\]
Therefore
\[
 \sum_i|H_i^{(1)}(z)|
 \le CB_{\cal C}R
 \le C\eta n,
\]
because \(B_{\cal C}\le B\).  Exponentiate.
\end{proof}

\section{The complete pair-core coefficient}
\label{sec:newV72-coefficient}

Define
\begin{equation}
 {\cal G}_{\cal C}(z)
 :=
 \exp\left\{\sum_{i\in{\cal C}}H_i^{(1)}(z)\right\}
 \sum_{\substack{G\ {\rm connected}\\V(G)={\cal C}}}
 \prod_{e\in E(G)}W_e(z).                                   \label{eq:newV72-G}
\end{equation}

\begin{lemma}[Pair-core coefficient bound]
\label{lem:newV72-coefficient}
Let \(n\) be the number of noncentral labels.  If \(B\le b_*\), with
\(b_*>0\) sufficiently small, then
\begin{equation}
 n!\left|[z^n]{\cal G}_{\cal C}(z)\right|
 \le
 K^{p+n}
 \left(\prod_{i\in{\cal C}}b_i\right)
 B^{p+n-2}.                                                 \label{eq:newV72-coefficient}
\end{equation}
For \(n=0\), the left side means \(|{\cal G}_{\cal C}(0)|\).
\end{lemma}

\begin{proof}
First take \(n\ge1\) and the adaptive radius \(R=\eta n/B\).
By \Cref{lem:newV72-effective-edge},
\[
 \sum_{e\subset{\cal C}}|W_e(z)|
 \le
 K_0\sum_{i<j}b_ib_j
 \le {K_0\over2}B_{\cal C}^2.
\]
Choose \(b_*\) so that this is bounded by a fixed numerical
constant.  Apply
\Cref{cor:newV72-weighted-tree} pointwise on \(|z|=R\), then
\Cref{lem:newV72-decorations}:
\[
 |{\cal G}_{\cal C}(z)|
 \le
 e^{C\eta n}K^p
 \left(\prod_{i\in{\cal C}}b_i\right)
 B_{\cal C}^{p-2}.
\]
Cauchy's estimate and \(n!\le n^n\) give
\[
 n!R^{-n}e^{C\eta n}
 \le
 K^nB^n.
\]
Use \(B_{\cal C}\le B\).  This proves
\eqref{eq:newV72-coefficient}.

For \(n=0\), all \(H\)'s vanish and
\(W_e(0)=K^2b_ib_j\).  Apply
\Cref{cor:newV72-weighted-tree} directly.
\end{proof}

\begin{theorem}[PCIC strong aggregate card]
\label{thm:newV72-PCIC}
On \(B\le b_*\), the complete positive majorant of PCIC, summed over
all choices of the centre set, all noncentral-label allocations,
all root partitions, all centre-core cycles, and all Duhamel
orders, obeys
\begin{equation}
 K^m\left(\prod_{\ell=1}^mb_\ell\right)B^{m-2}.              \label{eq:newV72-PCIC}
\end{equation}
\end{theorem}

\begin{proof}
Fix a centre set \({\cal C}\) of size \(p\) and factor the mark of
each of the \(n=m-p\) noncentral labels.  The labelled exponential
formula shows that their complete allocation among one-centre
decorations and centre-pair root families is exactly majorized by
\[
 n![z^n]{\cal G}_{\cal C}(z).
\]
Indeed, \(e^{H_i^{(1)}}\) partitions labels assigned to vertex \(i\),
while \(W_e\) partitions labels assigned to the occupied edge \(e\);
the connected graph sum enforces connectivity of the centre core.
Apply \Cref{lem:newV72-coefficient} and restore the noncentral
marks.  Every operator ordering is already paid by
\Cref{cor:newV70-Duhamel-BKAR}.

Finally sum the at most \(2^m\) choices of \({\cal C}\).  This only
changes the exponential base and proves
\eqref{eq:newV72-PCIC}.
\end{proof}

\begin{corollary}[All pair-supported centre cores are closed]
\label{cor:newV72-PCIC-closed}
The first unresolved MCWC incidence structure contains at least one
active root meeting three or more nonlinear centres.
\end{corollary}

\begin{proof}
One-centre cores are closed by V61, two-centre cores by V71, and
every connected core of arbitrary size whose roots meet at most two
centres is PCIC and is closed by
\Cref{thm:newV72-PCIC}.
\end{proof}

\begin{assessment}[The role of centre-core cycles]
\label{ass:newV72-cycles}
Cycles themselves do not cause the remaining loss.  After all roots
on a fixed centre pair are resummed into \(W_{ij}\), the total edge
activity is \(O(B^2)\); non-tree edges exponentiate at stable cost,
and one weighted Cayley tree pays connectivity.  The unresolved
object is a genuine centre hyperedge, not an ordinary graph cycle.
\end{assessment}

\begin{programme}[V73: genuine centre hyperedges]
\label{prog:newV72-V73}
\begin{enumerate}[leftmargin=3em]
\item For a fixed centre subset \(C\), \(|C|\ge3\), resum all roots
      whose centre intersection is exactly \(C\) and whose remaining
      labels have degree one.
\item Derive the effective hyperedge activity with its retained
      factor \(s^{(|C|-2)/2}\prod_{i\in C}b_i\).
\item Prove a connected hypergraph-to-incidence-tree allocation
      compatible with ordinary pair edges from V72.
\item Test the single all-centre hyperedge and the three pair edges
      on three centres in the same coefficient formula.
\item Keep the small-\(B\) stability and Duhamel word contraction
      outside the hypergraph combinatorics.
\end{enumerate}
\end{programme}

\begin{firewall}[Claim boundary after V72]
\label{fire:newV72-boundary}
V72 resums all roots supported on a fixed centre pair into an
effective edge, sums every noncentral-label allocation by one global
exponential generating function, reduces all connected centre graphs
to weighted Cayley trees, and closes PCIC for arbitrarily many
nonlinear centres in the small-total-mark regime.

V72 does not prove the genuine centre-hyperedge row, MCWC in full,
OCRC in its full noncommutative form, SCNGC, WPSC, GRLTC, LZTC,
PLRC, WSDC, BRSC, HWSFC, UGCGC, thermal connected WPRAC,
all-arity RGSC, the raw Wilson aggregate strong card, all-order
strong MTA, WUFC, CPM, cutoff removal, reflection positivity,
Osterwalder--Schrader reconstruction, a continuum Yang--Mills
measure, or a positive Hamiltonian gap.  The full Yang--Mills mass
gap has not been proved.
\end{firewall}

\begin{theorem}[V72 result ledger]
\label{thm:newV72-results}
V72:
\begin{enumerate}[label=\textup{(V72.\arabic*)},leftmargin=6.5em]
\item proves the connected-graph-to-tree majorant;
\item constructs the effective pair-edge generating function;
\item proves its \(O(b_ib_j)\) adaptive-circle bound;
\item resums all one-centre decorations simultaneously;
\item avoids the false \(p^n\) label-allocation factor;
\item proves the pair-core coefficient bound;
\item closes PCIC at arbitrary centre arity; and
\item isolates genuine centre hyperedges as the next row.
\end{enumerate}
\end{theorem}

\begin{proof}
Use
\Cref{lem:newV72-connected-graph,
cor:newV72-weighted-tree,
lem:newV72-effective-edge,
lem:newV72-decorations,
lem:newV72-coefficient,
thm:newV72-PCIC,
cor:newV72-PCIC-closed}.
\end{proof}

\paragraph{Parent-version record.}
V72 was formed from
\texttt{Renewal\_Yang\_Mills\_Mass\_Gap\_Research\_Notes\_V71.tex}.
The V71 parent had \(33\,704\) logical source lines,
\(1\,074\,277\) bytes, and SHA--256
\[
 \texttt{b35194e3415a3b026d9c2ab0f225552e59e0c88a47119b72cd75e350f9bba737}.
\]
The next compulsory companion audit is V75.

% =============================================================================
% END APPENDED FIFTH-SERIES V72 MATERIAL
% =============================================================================

\clearpage
\chapter{Fifth-series V73: effective centre hyperedges and the
one-hyperedge rooted-forest card}
\label{ch:newV73-hyperedge}

\section{A fixed retained centre set}
\label{sec:newV73-entry}

No companion audit is due at V73.

Let \(C\) be a set of \(d\ge2\) nonlinear centre labels, with
\[
 \beta_C:=\prod_{i\in C}b_i,
 \qquad
 B_C:=\sum_{i\in C}b_i.
\]
We first sum every root whose intersection with the nonlinear-centre
set is exactly \(C\).

\section{The multi-centre Bell lemma}
\label{sec:newV73-Bell}

\begin{lemma}[Fixed-\(d\)-centre Bell resummation]
\label{lem:newV73-d-Bell}
Let \(n\ge1\), let the \(d\) centre weights and \(n\) leaf weights
belong to \([\sqrt s,1)\), and let \(B\) be their total sum.
Suppose that every nonempty \(A\subseteq[n]\) obeys
\begin{equation}
 a_C(A)
 \le
 \{K(|A|+d)\}^{|A|+d}
 s^{(|A|+d-2)/2}\,
 \beta_C\prod_{j\in A}c_j.                                 \label{eq:newV73-d-activity}
\end{equation}
If \(B\le b_*\), with \(b_*>0\) sufficiently small, then
\begin{equation}
 \sum_{\mathcal P\in{\rm Part}([n])}
 \prod_{A\in\mathcal P}a_C(A)
 \le
 K_1^{n+d}\,
 \beta_C\left(\prod_{j=1}^nc_j\right)
 B_C^{d-2}B^n.                                              \label{eq:newV73-d-Bell}
\end{equation}
\end{lemma}

\begin{proof}
Factor the leaf marks and put
\[
 A_r=
 \{K(r+d)\}^{r+d}
 s^{(r+d-2)/2}\beta_C,
 \qquad
 H_C(z)=\sum_{r\ge1}{A_r\over r!}z^r.
\]
By \(r!\ge(r/e)^r\),
\begin{align*}
 {A_r\over r!}
 &\le
 K^{r+d}(r+d)^d
 s^{(r+d-2)/2}\beta_C.
\end{align*}
Indeed,
\((1+d/r)^r\le e^d\), and its exponential factor is included
in \(K^{r+d}\).

Choose \(R=\eta n/B\), where \(\eta>0\) is numerical and small.
Then \(\sqrt sR\le\eta\).  Also
\[
 (r+d)^d\le d^de^r.
\]
After decreasing \(\eta\) once more, summation on \(|z|=R\) gives
\begin{align}
 |H_C(z)|
 &\le
 K^d d^d s^{(d-2)/2}\beta_C\sqrt sR
 \notag\\
 &\le
 K_0^d\beta_CB_C^{d-2}\sqrt sR.                             \label{eq:newV73-H-bound}
\end{align}
The last step uses
\[
 d^d s^{(d-2)/2}
 =
 d^2(d\sqrt s)^{d-2}
 \le d^2B_C^{d-2},
\]
and absorbs \(d^2\) into \(K_0^d\).

By the arithmetic--geometric mean inequality,
\[
 \beta_C\le(B_C/d)^d.
\]
Hence \(K_0^d\beta_CB_C^{d-2}\) is uniformly small when \(B\le b_*\),
after \(b_*\) is fixed in terms of \(K_0\).  Therefore
\[
 |e^{H_C(z)}-1|
 \le
 K_2^d\beta_CB_C^{d-2}\sqrt sR.
\]
The labelled exponential formula and Cauchy's estimate yield
\begin{align*}
 \sum_{\mathcal P}\prod_{A\in\mathcal P}A_{|A|}
 &=
 n![z^n]\{e^{H_C(z)}-1\}\\
 &\le
 K^{n+d}\beta_CB_C^{d-2}
 n\sqrt s\,B^{n-1}\\
 &\le
 K^{n+d}\beta_CB_C^{d-2}B^n.
\end{align*}
Restore the leaf marks.
\end{proof}

\begin{assessment}[Interpolation between one and two centres]
\label{ass:newV73-interpolation}
At \(d=1\), the core factor \(B_C^{-1}\) is not meaningful and the
separate V61 nonempty-source argument gives the one-centre star.
At \(d=2\), \Cref{lem:newV73-d-Bell} is the V71 double-centre
mechanism.  For \(d\ge3\), the retained factor
\(s^{(d-2)/2}\) is first converted into
\(B_C^{d-2}\); it is exactly the internal spanning-tree degree of a
\(d\)-centre hyperedge.
\end{assessment}

\section{The effective centre hyperedge}
\label{sec:newV73-effective}

Define the bare-root majorant
\begin{equation}
 w_C^{(0)}
 :=
 (Kd)^d s^{(d-2)/2}\beta_C                                  \label{eq:newV73-w0}
\end{equation}
and let \(H_C(z)\) be the series in the preceding proof.  Put
\begin{equation}
 W_C(z):=(1+w_C^{(0)})e^{H_C(z)}-1.                          \label{eq:newV73-WC}
\end{equation}

\begin{lemma}[Strong effective hyperedge]
\label{lem:newV73-effective}
On the adaptive circle \(|z|=\eta n/B\), \(n\ge1\), and on
\(B\le b_*\),
\begin{equation}
 |W_C(z)|
 \le
 K^d\beta_CB_C^{d-2}.                                      \label{eq:newV73-effective}
\end{equation}
The same estimate holds at \(z=0\) when \(n=0\).
\end{lemma}

\begin{proof}
The floor computation in
\eqref{eq:newV73-H-bound} gives
\[
 w_C^{(0)}
 \le K^d\beta_CB_C^{d-2}.
\]
Moreover
\[
 |W_C|
 \le
 w_C^{(0)}e^{|H_C|}+e^{|H_C|}-1.
\]
Use \eqref{eq:newV73-H-bound},
\(\sqrt sR\le\eta\), and the smallness established in the proof of
\Cref{lem:newV73-d-Bell}.  At \(z=0\), \(H_C(0)=0\).
\end{proof}

\begin{corollary}[Internal tree form]
\label{cor:newV73-internal-tree}
The effective hyperedge satisfies
\begin{equation}
 |W_C(z)|
 \le
 K^d
 \sum_{T\in\mathfrak T(C)}
 \prod_{i\in C}b_i^{d_i(T)}.                               \label{eq:newV73-internal-tree}
\end{equation}
\end{corollary}

\begin{proof}
Use the weighted Cayley--Pr\"ufer identity
\[
 \sum_{T\in\mathfrak T(C)}
 \prod_{i\in C}b_i^{d_i(T)}
 =
 \beta_CB_C^{d-2}
\]
in \eqref{eq:newV73-effective}.
\end{proof}

\section{Pair graphs attached to a preconnected root set}
\label{sec:newV73-rooted}

\begin{lemma}[Rooted pair-graph majorant]
\label{lem:newV73-rooted-graph}
Let \({\cal C}\) be a set of \(p\) centres and let
\(\varnothing\ne C\subseteq{\cal C}\) be declared internally
connected.  Suppose \(w_{ij}\le Kb_ib_j\) and
\(\sum_{i<j}w_{ij}\le C_0\).  Then
\begin{align}
 &\sum_{\substack{G\ {\rm on}\ {\cal C}\\
                  G\ {\rm connects\ every\ vertex\ to}\ C}}
 \prod_{ij\in E(G)}w_{ij}
 \notag\\
 &\qquad\le
 K_1^{p-|C|}
 \left(\prod_{i\in{\cal C}\setminus C}b_i\right)
 B_{\cal C}^{p-|C|}.                                       \label{eq:newV73-rooted-graph}
\end{align}
Here all vertices of \(C\) are contracted to one root before
connectivity is tested.
\end{lemma}

\begin{proof}
Assign to every admissible graph a spanning forest in which each
component contains exactly one vertex of \(C\).  Enlarge the assigned
family to all graphs containing that forest; the unused edges cost at
most
\(\exp\{\sum w_{ij}\}\le e^{C_0}\), as in
\Cref{lem:newV72-connected-graph}.

Orient the forest toward its roots.  Drop acyclicity and all
restrictions on parent choices.  Every
\(i\in{\cal C}\setminus C\) then contributes at most
\[
 Kb_i\sum_{j\in{\cal C}}b_j=Kb_iB_{\cal C}.
\]
The product of these bounds proves
\eqref{eq:newV73-rooted-graph}.
\end{proof}

\section{One genuine hyperedge and arbitrary pair roots}
\label{sec:newV73-one-hyperedge}

\begin{definition}[One genuine hyperedge class, OGHIC]
\label{def:newV73-OGHIC}
OGHIC consists of PCIC enlarged by one occupied effective centre
subset \(C_*\) with \(|C_*|\ge3\).  Arbitrarily many original roots
with centre intersection \(C_*\) are included inside \(W_{C_*}\);
all other core roots meet one or two centres.
\end{definition}

\begin{theorem}[OGHIC strong aggregate card]
\label{thm:newV73-OGHIC}
On \(B\le b_*\), the complete positive majorant of OGHIC, including
all centre choices, noncentral-label allocations, root partitions,
pair-core cycles, and Duhamel orders, obeys
\begin{equation}
 K^m\left(\prod_{\ell=1}^mb_\ell\right)B^{m-2}.              \label{eq:newV73-OGHIC}
\end{equation}
\end{theorem}

\begin{proof}
Fix a centre set \({\cal C}\), \(|{\cal C}|=p\), and
\(C_*\subseteq{\cal C}\), \(|C_*|=d\ge3\).  Use the same global
label variable \(z\) for:
\begin{enumerate}[label=\textup{(\roman*)},leftmargin=3.5em]
\item one-centre decorations \(e^{\sum_iH_i^{(1)}(z)}\);
\item all effective pair edges \(W_{ij}(z)\); and
\item the effective hyperedge \(W_{C_*}(z)\).
\end{enumerate}
The labelled product formula then assigns every noncentral label
to exactly one root family; no preliminary allocation is made.

Let \(n=m-p\).  If \(n\ge1\), take
\(R=\eta n/B\).  On \(|z|=R\),
\Cref{lem:newV73-effective} bounds the hyperedge by
\[
 K^d\beta_{C_*}B_{C_*}^{d-2}.
\]
By \Cref{lem:newV72-effective-edge}, every pair edge is bounded by
\(Kb_ib_j\), and their total activity is \(O(B_{\cal C}^2)\).
Apply \Cref{lem:newV73-rooted-graph} after contracting \(C_*\).
Multiplication by the hyperedge gives
\begin{align*}
 &K^p
 \beta_{C_*}B_{C_*}^{d-2}
 \left(\prod_{i\in{\cal C}\setminus C_*}b_i\right)
 B_{\cal C}^{p-d}\\
 &\qquad\le
 K^p\left(\prod_{i\in{\cal C}}b_i\right)
 B^{p-2}.
\end{align*}
The vertex-decoration factor is at most \(e^{C\eta n}\) by
\Cref{lem:newV72-decorations}.  Cauchy's estimate supplies
\[
 n!R^{-n}e^{C\eta n}\le K^nB^n.
\]
Restore the \(n\) noncentral marks.

For \(n=0\), evaluate all generating functions at zero and use the
same rooted-graph estimate.  Finally sum the at most \(2^m\) centre
sets and at most \(2^p\) choices of \(C_*\).  Duhamel orderings cost
only the constant in
\Cref{cor:newV70-Duhamel-BKAR}.  All enumeration factors enter
one exponential base, proving \eqref{eq:newV73-OGHIC}.
\end{proof}

\begin{corollary}[The remaining core has at least two genuine
hyperedges]
\label{cor:newV73-remaining}
After V72--V73, an unresolved MCWC core must contain at least two
distinct occupied centre subsets of cardinality at least three.
\end{corollary}

\begin{proof}
If none is present, apply \Cref{thm:newV72-PCIC}; if exactly one is
present, apply \Cref{thm:newV73-OGHIC}.
\end{proof}

\begin{assessment}[What contraction achieved]
\label{ass:newV73-contraction}
An effective \(d\)-centre hyperedge already carries its internal
tree degree \(B_C^{d-2}\).  Contracting it turns the remaining pair
problem into a rooted forest, which contributes exactly one
\(b_iB\) for every centre outside \(C\).  Thus
\[
 (d-2)+(p-d)=p-2.
\]
The equality is the centre-core analogue of the thermal degree
ledger used throughout the programme.
\end{assessment}

\begin{programme}[V74: several genuine centre hyperedges]
\label{prog:newV73-V74}
\begin{enumerate}[leftmargin=3em]
\item Form the hyperedge-intersection graph of the occupied effective
      centre subsets \(C\), \(|C|\ge3\).
\item Prove the card first when those subsets have an acyclic
      incidence graph, by successive contraction and the rooted
      pair-forest lemma.
\item For an intersection cycle, retain one internal tree expansion
      \eqref{eq:newV73-internal-tree} per hyperedge and test whether
      surplus internal edges exponentiate under
      \(\sum_C|W_C|=O(B^2)\).
\item Sum the choices of centre subsets before replacing local
      \(B_C\)'s by the global \(B\).
\item If all effective-hyperedge cores close, combine V72--V74 into
      MCWC and return to SCNGC.
\end{enumerate}
\end{programme}

\begin{firewall}[Claim boundary after V73]
\label{fire:newV73-boundary}
V73 proves the fixed-\(d\)-centre Bell resummation, constructs a
strong effective centre hyperedge, proves the rooted pair-graph
majorant, and closes every core having one genuine centre hyperedge
and otherwise arbitrary pair-supported roots and vertex decorations.

V73 does not prove cores with several genuine centre hyperedges,
MCWC in full, OCRC in its full noncommutative form, SCNGC, WPSC,
GRLTC, LZTC, PLRC, WSDC, BRSC, HWSFC, UGCGC, thermal connected
WPRAC, all-arity RGSC, the raw Wilson aggregate strong card,
all-order strong MTA, WUFC, CPM, cutoff removal, reflection
positivity, Osterwalder--Schrader reconstruction, a continuum
Yang--Mills measure, or a positive Hamiltonian gap.  The full
Yang--Mills mass gap has not been proved.
\end{firewall}

\begin{theorem}[V73 result ledger]
\label{thm:newV73-results}
V73:
\begin{enumerate}[label=\textup{(V73.\arabic*)},leftmargin=6.5em]
\item proves the fixed-\(d\)-centre Bell lemma;
\item constructs an effective centre hyperedge;
\item proves its internal weighted-tree form;
\item proves the rooted pair-graph majorant;
\item performs all noncentral-label allocations globally;
\item closes OGHIC;
\item permits arbitrary pair cycles and vertex decorations; and
\item isolates cores with at least two genuine hyperedges.
\end{enumerate}
\end{theorem}

\begin{proof}
Use
\Cref{lem:newV73-d-Bell,
lem:newV73-effective,
cor:newV73-internal-tree,
lem:newV73-rooted-graph,
thm:newV73-OGHIC,
cor:newV73-remaining}.
\end{proof}

\paragraph{Parent-version record.}
V73 was formed from
\texttt{Renewal\_Yang\_Mills\_Mass\_Gap\_Research\_Notes\_V72.tex}.
The V72 parent had \(34\,083\) logical source lines,
\(1\,086\,501\) bytes, and SHA--256
\[
 \texttt{926374cef6dfd909b294935122ee6789b05bdfe58f286454bc02261129c1c69a}.
\]
The next compulsory companion audit is V75.

% =============================================================================
% END APPENDED FIFTH-SERIES V73 MATERIAL
% =============================================================================

\clearpage
\chapter{Fifth-series V74: the centre-core excess ledger and
reserve-preserving hyperedge stars}
\label{ch:newV74-core-ledger}

\section{Fixed connected centre cores}
\label{sec:newV74-fixed}

No companion audit is due at V74.

Let \({\cal H}\) be a connected simple hypergraph on a set
\({\cal C}\) of \(p\ge2\) nonlinear centres.  Write
\[
 E=|{\cal H}|,
 \qquad d_A=|A|\quad(A\in{\cal H}),
\]
and define its incidence excess
\begin{equation}
 g({\cal H})
 :=
 \sum_{A\in{\cal H}}(d_A-1)-p+1
 =
 \sum_Ad_A-p-E+1.                                           \label{eq:newV74-excess}
\end{equation}
Connectedness gives \(g({\cal H})\ge0\); it is the cyclomatic
number of the bipartite incidence graph.

\begin{lemma}[Exact thermal degree of a fixed core]
\label{lem:newV74-fixed-degree}
Suppose every occupied effective hyperedge obeys
\[
 w_A\le K^{d_A}
 \left(\prod_{i\in A}b_i\right)B_A^{d_A-2}.
\]
Then
\begin{align}
 \prod_{A\in{\cal H}}w_A
 \le{}&
 K^{\sum_Ad_A}
 \left(\prod_{i\in{\cal C}}b_i\right)
 B_{\cal C}^{p-2+2g({\cal H})}.                             \label{eq:newV74-fixed-degree}
\end{align}
Moreover
\begin{equation}
 E\le p+g-1,
 \qquad
 \sum_Ad_A\le2p+2g-2.                                      \label{eq:newV74-size-excess}
\end{equation}
\end{lemma}

\begin{proof}
Let \(\deg_{\cal H}(i)\) be the number of hyperedges containing
\(i\).  After one \(b_i\) per centre is factored, the repeated centre
marks obey
\[
 \prod_i b_i^{\deg_{\cal H}(i)-1}
 \le
 B_{\cal C}^{\sum_Ad_A-p}.
\]
Also
\[
 \prod_AB_A^{d_A-2}
 \le
 B_{\cal C}^{\sum_Ad_A-2E}.
\]
The sum of these two exponents is
\[
 2\sum_Ad_A-p-2E
 =
 p-2+2g({\cal H}),
\]
which proves \eqref{eq:newV74-fixed-degree}.

Since every \(d_A\ge2\),
\[
 g=\sum_A(d_A-1)-p+1\ge E-p+1,
\]
giving the first inequality in
\eqref{eq:newV74-size-excess}.  Substitute it into
\(\sum_Ad_A=p+E+g-1\) to obtain the second.
\end{proof}

\begin{corollary}[Every fixed incidence cycle has a smallness
reserve]
\label{cor:newV74-cycle-reserve}
After \(b_*>0\) is chosen so that \(K^2b_*^2\le1/2\), every fixed
connected centre core on \(B_{\cal C}\le b_*\) obeys
\begin{equation}
 \prod_{A\in{\cal H}}w_A
 \le
 K_1^p\,2^{-g({\cal H})}
 \left(\prod_{i\in{\cal C}}b_i\right)
 B_{\cal C}^{p-2}.                                         \label{eq:newV74-cycle-reserve}
\end{equation}
\end{corollary}

\begin{proof}
By \eqref{eq:newV74-size-excess},
\[
 K^{\sum d_A}B_{\cal C}^{2g}
 \le
 K^{2p}(K^2B_{\cal C}^2)^g.
\]
Apply \Cref{lem:newV74-fixed-degree}.
\end{proof}

\section{Why the fixed-core estimate cannot yet be summed}
\label{sec:newV74-warning}

\begin{warning}[Acyclic incidence still contains a Bell family]
\label{warn:newV74-Bell}
The factor \(2^{-g}\) controls cycles but does not sum all cores.
Let one centre \(c\) be joined to \(p-1\) other centre labels by an
ordinary star.  Partition its \(p-1\) star edges arbitrarily, and
replace each block by the hyperedge consisting of \(c\) and the
leaves in that block.  Every resulting incidence graph is a tree,
so \(g=0\), and there are \({\rm Bell}(p-1)\) such cores.
Therefore summing
\eqref{eq:newV74-cycle-reserve} term by term loses the exponential
base.  The unrelaxed \(s\)-reserve must be retained until this
partition has been summed.
\end{warning}

\begin{proof}
Two hyperedges in the construction meet only at \(c\); hence the
bipartite incidence graph is connected and has no cycle.  Distinct
set partitions give distinct hyperedge families.  The Bell-number
growth was quantified in
\Cref{prop:newV60-Bell-no-go}.
\end{proof}

\section{The effective hyperedge retains the unrelaxed reserve}
\label{sec:newV74-reserve}

\begin{lemma}[Reserve-preserving effective hyperedge]
\label{lem:newV74-effective-reserve}
Retain the notation \(W_C(z)\) of
\eqref{eq:newV73-WC}.  On the adaptive circle
\(|z|=\eta n/B\), and on \(B\le b_*\),
\begin{equation}
 |W_C(z)|
 \le
 (K|C|)^{|C|}
 s^{(|C|-2)/2}\prod_{i\in C}b_i.                            \label{eq:newV74-effective-reserve}
\end{equation}
The same estimate holds at \(z=0\).
\end{lemma}

\begin{proof}
Put \(d=|C|\).  The bare term
\eqref{eq:newV73-w0} is already the right side.
The proof of \eqref{eq:newV73-H-bound}, before converting
\(d\sqrt s\) into \(B_C\), gives
\[
 |H_C(z)|
 \le
 (K_0d)^d s^{(d-2)/2}
 \left(\prod_{i\in C}b_i\right)\sqrt s\,|z|.
\]
On the adaptive circle, \(\sqrt s|z|\le\eta\).
The arithmetic--geometric mean bound used in V73 and the smallness
of \(B\) make both the bare activity and \(H_C\) uniformly small.
Insert these estimates into
\[
 |W_C|\le w_C^{(0)}e^{|H_C|}+e^{|H_C|}-1
\]
and enlarge \(K\).  At \(z=0\), \(H_C(0)=0\).
\end{proof}

\section{The complete centre-core star}
\label{sec:newV74-star}

\begin{definition}[Centre-core star class, CCSC]
\label{def:newV74-CCSC}
Fix a centre \(c\).  CCSC consists of effective hyperedges whose
centre sets are
\[
 \{c\}\cup A,\qquad A\in\mathcal P,
\]
where \(\mathcal P\) is an arbitrary partition of all other
nonlinear centres.  Degree-one noncentral labels may be assigned
among these effective hyperedges or among one-centre decorations.
All root multiplicities and Duhamel orderings are included.
\end{definition}

\begin{theorem}[CCSC strong aggregate card]
\label{thm:newV74-CCSC}
On \(B\le b_*\), the complete CCSC positive majorant on \(m\)
labels obeys
\begin{equation}
 K^m\left(\prod_{\ell=1}^mb_\ell\right)B^{m-2}.              \label{eq:newV74-CCSC}
\end{equation}
\end{theorem}

\begin{proof}
Fix a set \({\cal C}\) of \(p\) nonlinear centres, the star centre
\(c\in{\cal C}\), and let \(n=m-p\) be the number of noncentral
labels.  Use one labelled variable \(z\) for every noncentral label.
If \(n\ge1\), take the adaptive circle \(|z|=\eta n/B\).

For a nonempty block
\(A\subseteq{\cal C}\setminus\{c\}\),
\Cref{lem:newV74-effective-reserve} gives
\[
 |W_{\{c\}\cup A}(z)|
 \le
 \{K(|A|+1)\}^{|A|+1}
 s^{(|A|-1)/2}
 b_c\prod_{i\in A}b_i.
\]
This is exactly the activity hypothesis
\eqref{eq:newV61-star-activity}.  Therefore
\Cref{lem:newV61-Bell-star}, applied to the \(p-1\) centre leaves,
gives
\begin{equation}
 \sum_{\mathcal P\in{\rm Part}({\cal C}\setminus\{c\})}
 \prod_{A\in\mathcal P}|W_{\{c\}\cup A}(z)|
 \le
 K^p\left(\prod_{i\in{\cal C}}b_i\right)
 B_{\cal C}^{p-2}.                                         \label{eq:newV74-core-star}
\end{equation}
The bound is uniform on the adaptive circle.  All one-centre
decorations contribute the factor in
\Cref{lem:newV72-decorations}, at most \(e^{C\eta n}\).
Thus Cauchy's estimate supplies
\[
 n!|z|^{-n}e^{C\eta n}\le K^nB^n.
\]
Restore the noncentral marks and use
\(B_{\cal C}\le B\).

For \(n=0\), evaluate at \(z=0\) and use
\eqref{eq:newV74-core-star} directly.  Sum the at most \(2^m\)
centre sets and at most \(m\) choices of \(c\).  Finally apply
\Cref{cor:newV70-Duhamel-BKAR} to all operator orderings.
\end{proof}

\begin{corollary}[The Bell hypertree star is closed]
\label{cor:newV74-star-closed}
The superexponential family exhibited in
\Cref{warn:newV74-Bell} obeys the strong aggregate card once the
unrelaxed reserve is retained through the effective-hyperedge
coefficient sum.
\end{corollary}

\begin{proof}
It is the nondecorated subfamily of
\Cref{thm:newV74-CCSC}.
\end{proof}

\section{The remaining multi-hyperedge operation}
\label{sec:newV74-remaining}

\begin{definition}[Centre-incidence resummation card, CIRC]
\label{def:newV74-CIRC}
CIRC is the joint summation of all connected effective-hyperedge
cores which:
\begin{enumerate}[label=\textup{(\roman*)},leftmargin=3.5em]
\item retains \eqref{eq:newV74-effective-reserve} until all
      incidence-tree partitions have been summed;
\item uses the \(B^{2g}\) surplus of
      \eqref{eq:newV74-fixed-degree} only for incidence cycles;
\item assigns each noncentral label once through the global
      coefficient variable; and
\item yields
      \(K^p(\prod_{i\in{\cal C}}b_i)B_{\cal C}^{p-2}\)
      before noncentral decorations are restored.
\end{enumerate}
\end{definition}

\begin{proposition}[CIRC completes the small-\(B\) MCWC core]
\label{prop:newV74-CIRC}
CIRC, together with V72--V74, closes every centre-core incidence
structure in the small-\(B\) regime and therefore closes the
small-\(B\) part of MCWC.
\end{proposition}

\begin{proof}
PCIC is \Cref{thm:newV72-PCIC}; a unique genuine hyperedge is
\Cref{thm:newV73-OGHIC}; a centre-core star is
\Cref{thm:newV74-CCSC}.  Every remaining core has several genuine
hyperedges and non-star incidence.  CIRC is exactly their aggregate
bound, with the same vertex decorations and Duhamel contraction
already proved in V70--V73.
\end{proof}

\begin{assessment}[Position before the V75 audit]
\label{ass:newV74-position}
There is no fixed-core power-counting obstruction: every incidence
cycle improves the thermal degree by two.  There is also no star
enumeration obstruction: the retained clock reserve closes the
complete Bell family.  The live question is whether a non-star
incidence hypertree can be pruned recursively without using the same
centre mark in two independent Bell sums.  This is a combinatorial
allocation question, not a new analytic estimate on the Wilson
clock.
\end{assessment}

\begin{programme}[V75: compulsory audit and hypertree pruning]
\label{prog:newV74-V75}
\begin{enumerate}[leftmargin=3em]
\item Perform the compulsory newest-SM/newest-NS audit and record
      immutable source snapshots.
\item Root the bipartite incidence hypertree at one centre and prune
      a deepest hyperedge whose child centres carry disjoint
      descendant cores.
\item Apply the V61 star EGF once at each pruned parent, retaining
      the parent mark for its unique edge toward the root.
\item Prove that the pruning order is independent of choices, or
      define a canonical order which prevents double payment.
\item Treat incidence cycles only after the hypertree recursion has
      been closed; use the \(B^{2g}\) reserve for surplus hyperedges.
\end{enumerate}
\end{programme}

\begin{firewall}[Claim boundary after V74]
\label{fire:newV74-boundary}
V74 proves the exact thermal degree and two-power cycle reserve of
every fixed centre core, proves that relaxed fixed-core bounds cannot
be summed because of an acyclic Bell family, proves that effective
hyperedges retain the unrelaxed clock reserve, and closes the
complete centre-core star family.

V74 does not prove CIRC, non-star multi-hyperedge cores, MCWC in
full, OCRC in its full noncommutative form, SCNGC, WPSC, GRLTC,
LZTC, PLRC, WSDC, BRSC, HWSFC, UGCGC, thermal connected WPRAC,
all-arity RGSC, the raw Wilson aggregate strong card, all-order
strong MTA, WUFC, CPM, cutoff removal, reflection positivity,
Osterwalder--Schrader reconstruction, a continuum Yang--Mills
measure, or a positive Hamiltonian gap.  The full Yang--Mills mass
gap has not been proved.
\end{firewall}

\begin{theorem}[V74 result ledger]
\label{thm:newV74-results}
V74:
\begin{enumerate}[label=\textup{(V74.\arabic*)},leftmargin=6.5em]
\item defines the centre-core incidence excess;
\item proves the exact fixed-core thermal degree;
\item proves two extra thermal powers per incidence cycle;
\item identifies the acyclic Bell no-go for relaxed enumeration;
\item proves the reserve-preserving effective-hyperedge bound;
\item closes the complete centre-core star;
\item formulates CIRC; and
\item isolates canonical non-star hypertree pruning for V75.
\end{enumerate}
\end{theorem}

\begin{proof}
Use
\Cref{lem:newV74-fixed-degree,
cor:newV74-cycle-reserve,
warn:newV74-Bell,
lem:newV74-effective-reserve,
thm:newV74-CCSC,
cor:newV74-star-closed,
prop:newV74-CIRC}.
\end{proof}

\paragraph{Parent-version record.}
V74 was formed from
\texttt{Renewal\_Yang\_Mills\_Mass\_Gap\_Research\_Notes\_V73.tex}.
The V73 parent had \(34\,482\) logical source lines,
\(1\,098\,777\) bytes, and SHA--256
\[
 \texttt{4c07e87513baa701a7d6aae292cb09174c8e8c2c0a03689b3c9f4c8e4c14754f}.
\]
The next compulsory companion audit is V75.

% =============================================================================
% END APPENDED FIFTH-SERIES V74 MATERIAL
% =============================================================================

\clearpage
\chapter{Fifth-series V75: compulsory audit and the canonical
rooted-hypertree quotient}
\label{ch:newV75-hypertree}

\section{Compulsory companion audit}
\label{sec:newV75-audit}

\begin{audit}[V75 immutable companion snapshots]
\label{audit:newV75-companions}
The newest active companion sources were read without modification:
\begin{align*}
 &\texttt{Renewal\_SM\_Einstein\_Continuum\_Research\_Notes\_V99.tex},\\
 &\qquad 35\,122\ \hbox{logical lines},\quad
 1\,349\,056\ \hbox{bytes},\\
 &\qquad
 \texttt{53cbaa0772dba73ae19618614c0d47d98197436759668ef6743d26570c74b28e},
 \\
 &\texttt{Renewal\_NS\_Stability\_Research\_Notes\_V41.tex},\\
 &\qquad 10\,071\ \hbox{logical lines},\quad
 512\,596\ \hbox{bytes},\\
 &\qquad
 \texttt{ef30fd70b900b16fa0d50e36048bfa0249ff47263d4803a62422df63763bae40}.
\end{align*}

SM V99 proves two exact Ginsparg--Wilson intertwiners, assembles the
four chiral kinetic and improved-mass corners into the standard
massive overlap operator, and obtains its singular-value floor.
The same note proves that treating an unimproved mass insertion as a
perturbation reaches the endpoint of the available margin and
certifies no positive floor.

NS V41 decomposes the invariant blow-up identities mode by mode,
proves an explicit flux bound and a \(1/k\) upper bound for the mean
energy in one Hermite mode, and computes the first dipole row.
It also proves that the modewise identities form an unclosed
hierarchy: the nonlinearity raises the moment order.
\end{audit}

\begin{assessment}[Audit transfer]
\label{ass:newV75-audit-transfer}
No SM or NS estimate is imported.  The relevant proof discipline is:
\begin{enumerate}[label=\textup{(\roman*)},leftmargin=3.5em]
\item assemble the exact object which has the desired margin before
      comparing norms; and
\item do not infer an aggregate hierarchy bound from valid
      componentwise bounds.
\end{enumerate}
Thus V75 does not promote the V74 fixed-core estimate to CIRC.  It
constructs the exact rooted combinatorial quotient which must still
be summed.
\end{assessment}

\section{Canonical rooting of an incidence hypertree}
\label{sec:newV75-rooting}

\begin{definition}[Incidence hypertree]
\label{def:newV75-hypertree}
An incidence hypertree on a centre set \({\cal C}\) is a connected
simple hypergraph whose bipartite centre--hyperedge incidence graph
is a tree.  Equivalently, its excess
\eqref{eq:newV74-excess} is zero.
\end{definition}

\begin{theorem}[Rooted-tree and child-partition bijection]
\label{thm:newV75-bijection}
Fix a root centre \(o\in{\cal C}\).  Incidence hypertrees on
\({\cal C}\) are in bijection with the following data:
\begin{enumerate}[label=\textup{(\roman*)},leftmargin=3.5em]
\item an ordinary rooted tree \(T\) on \({\cal C}\), oriented toward
      \(o\); and
\item for every centre \(i\), a set partition
      \(\mathcal P_i\) of its children
      \({\rm Ch}_T(i)\).
\end{enumerate}
A block \(A\in\mathcal P_i\) represents the hyperedge
\(\{i\}\cup A\).
\end{theorem}

\begin{proof}
Root the bipartite incidence tree at \(o\).  Every hyperedge vertex
has a unique adjacent centre on the path to \(o\); call it the parent
centre.  All its other adjacent centres are its children.  Replace
that hyperedge vertex and its incidence edges by ordinary edges from
the parent centre to each child centre.  Since the incidence graph
was a tree, the result is a connected ordinary graph with
\[
 \sum_{A\in{\cal H}}(|A|-1)=p-1
\]
edges, hence an ordinary tree.  Hyperedges with the same parent
partition that parent's child set.

Conversely, replace every block \(A\in\mathcal P_i\) by one
hyperedge vertex adjacent to \(i\) and all vertices of \(A\).
The resulting incidence graph is connected.  Its number of
incidence edges is
\[
 \sum_i\sum_{A\in\mathcal P_i}(|A|+1)
 =
 (p-1)+\sum_i|\mathcal P_i|,
\]
which is one less than its number
\(p+\sum_i|\mathcal P_i|\) of vertices.  It is therefore a tree.
The constructions are inverse.
\end{proof}

\section{Exact local child-partition sum}
\label{sec:newV75-local}

For a rooted ordinary tree \(T\), put
\[
 q_i:=|{\rm Ch}_T(i)|,
 \qquad
 I(T):=\{i:q_i>0\},
 \qquad
 B_i^T:=b_i+\sum_{j\in{\rm Ch}_T(i)}b_j.
\]

\begin{lemma}[All hyperedge blockings over one rooted tree]
\label{lem:newV75-fixed-tree}
Sum, for each \(i\), every partition of
\({\rm Ch}_T(i)\), and bound the corresponding effective
hyperedges by the reserve
\eqref{eq:newV74-effective-reserve}.  Then the complete blocking
sum over the fixed rooted tree \(T\) is at most
\begin{align}
 K^{2p}
 \left(\prod_{j\ne o}b_j\right)
 \left(\prod_{i\in I(T)}b_i\right)
 \prod_{i\in I(T)}(B_i^T)^{q_i-1}.                          \label{eq:newV75-fixed-tree}
\end{align}
In particular,
\begin{equation}
 \textup{\eqref{eq:newV75-fixed-tree}}
 \le
 K^{2p}\left(\prod_{i\in{\cal C}}b_i\right)
 B_{\cal C}^{p-2}.                                         \label{eq:newV75-fixed-tree-strong}
\end{equation}
\end{lemma}

\begin{proof}
Fix \(i\in I(T)\).  For a nonempty child block
\(A\subseteq{\rm Ch}_T(i)\),
\eqref{eq:newV74-effective-reserve} is
\[
 \{K(|A|+1)\}^{|A|+1}
 s^{(|A|-1)/2}b_i\prod_{j\in A}b_j.
\]
Apply \Cref{lem:newV61-Bell-star} with centre \(i\) and its
\(q_i\) children.  The complete local partition sum is bounded by
\[
 K^{q_i+1}b_i
 \left(\prod_{j\in{\rm Ch}_T(i)}b_j\right)
 (B_i^T)^{q_i-1}.
\]
Multiply this over \(i\in I(T)\).  Every nonroot vertex occurs
exactly once as a child, so the product of all child marks is
\(\prod_{j\ne o}b_j\).  Also
\[
 \sum_{i\in I(T)}(q_i+1)
 =
 (p-1)+|I(T)|\le2p-2,
\]
which proves \eqref{eq:newV75-fixed-tree}.

For the second assertion, the root \(o\) belongs to \(I(T)\).
After one mark for every centre is removed, the repeated parent
marks are
\[
 \prod_{i\in I(T)\setminus\{o\}}b_i
 \le B_{\cal C}^{|I(T)|-1}.
\]
Moreover
\[
 \sum_{i\in I(T)}(q_i-1)
 =
 p-1-|I(T)|.
\]
Replacing every \(B_i^T\) by \(B_{\cal C}\) gives the total exponent
\[
 (|I(T)|-1)+(p-1-|I(T)|)=p-2.
\]
\end{proof}

\begin{assessment}[What is exact and what is not]
\label{ass:newV75-fixed-not-summed}
\Cref{thm:newV75-bijection} is an exact decomposition, and
\Cref{lem:newV75-fixed-tree} sums every hyperedge blocking above one
ordinary rooted tree.  The final relaxation
\eqref{eq:newV75-fixed-tree-strong} is uniform in \(T\), but summing
it over the \(p^{p-2}\) ordinary trees would be invalid.  The local
sums \(B_i^T\) and repeated parent marks must remain present during
the tree sum.
\end{assessment}

\section{The failed parent-map relaxation}
\label{sec:newV75-parent-map}

\begin{proposition}[Dropping acyclicity recreates a
superexponential factor]
\label{prop:newV75-parent-no-go}
If one replaces the rooted trees in
\eqref{eq:newV75-fixed-tree} by arbitrary parent maps and only then
uses \(B_i^T\le B_{\cal C}\), the resulting majorant contains
\begin{equation}
 \max_{1\le h\le p-1}
 {h^{p-1}\over(h-1)!},                                     \label{eq:newV75-parent-loss}
\end{equation}
whose \(p\)-th root is unbounded.  Hence that relaxation cannot
prove CIRC.
\end{proposition}

\begin{proof}
Let \(I\) be the set of active parents, require \(o\in I\), and put
\(|I|=h\).  There are at most \(h^{p-1}\) maps of the \(p-1\)
nonroot vertices into \(I\).  After
\eqref{eq:newV75-fixed-tree} is relaxed to the global \(B_{\cal C}\),
summing the repeated parent marks over all \(I\) gives at best
\[
 \sum_{\substack{I\ni o\\|I|=h}}
 \prod_{i\in I\setminus\{o\}}b_i
 \le {B_{\cal C}^{h-1}\over(h-1)!}.
\]
This cancels the corresponding power of \(B_{\cal C}\) but leaves
\eqref{eq:newV75-parent-loss}.

Take \(h=\lfloor p/2\rfloor\).  Stirling's formula gives
\[
 {h^{p-1}\over(h-1)!}
 \ge
 c^p h^{p-h},
\]
whose \(p\)-th root grows like a positive power of \(p\).
\end{proof}

\section{The exact remaining hypertree quotient}
\label{sec:newV75-quotient}

\begin{definition}[Rooted hypertree quotient card, RHQC]
\label{def:newV75-RHQC}
RHQC is the estimate
\begin{align}
 &\sum_{\substack{T\in\mathfrak T({\cal C})\\T\ {\rm rooted\ at}\ o}}
 \left(\prod_{j\ne o}b_j\right)
 \left(\prod_{i\in I(T)}b_i\right)
 \prod_{i\in I(T)}(B_i^T)^{q_i-1}
 \notag\\
 &\qquad\le
 K^p
 \left(\prod_{i\in{\cal C}}b_i\right)
 B_{\cal C}^{p-2}.                                         \label{eq:newV75-RHQC}
\end{align}
The estimate must use the acyclic rooted-tree structure; the
parent-map relaxation of
\Cref{prop:newV75-parent-no-go} is forbidden.
\end{definition}

\begin{proposition}[RHQC closes incidence hypertrees]
\label{prop:newV75-RHQC-closes}
RHQC implies the CIRC bound for every incidence-hypertree core.
\end{proposition}

\begin{proof}
Use the bijection
\Cref{thm:newV75-bijection}, apply
\Cref{lem:newV75-fixed-tree} to sum all child partitions over each
ordinary rooted tree, and then apply
\eqref{eq:newV75-RHQC}.  The numerical factor \(K^{2p}\) changes
only the exponential base.  The global noncentral-label coefficient
and Duhamel orderings are then restored exactly as in V72--V74.
\end{proof}

\begin{assessment}[Position after the audit]
\label{ass:newV75-position}
The fixed-hypertree power count and every local Bell partition are
closed.  The live row is now the single explicit sum
\eqref{eq:newV75-RHQC}.  This is narrower than CIRC: it contains no
clock analysis, no operator word, no noncentral-label allocation,
and no incidence cycle.  It is a weighted rooted-tree inequality with
a child-set-dependent local sum.
\end{assessment}

\begin{programme}[V76: matrix-tree or Lagrange route to RHQC]
\label{prog:newV75-V76}
\begin{enumerate}[leftmargin=3em]
\item Rewrite the left side of \eqref{eq:newV75-RHQC} by directed
      Prüfer codes, retaining \(B_i^T\) rather than replacing it by
      \(B_{\cal C}\).
\item In parallel, seek a determinant representation by the directed
      matrix-tree theorem after introducing one auxiliary source per
      possible child.
\item Test RHQC first for paths, stars, and depth-two trees with
      heterogeneous marks.
\item If the exact determinant has mixed signs, use the rooted-tree
      recursion rather than an absolute cofactor expansion.
\item Only after RHQC is proved, spend the \(B^{2g}\) reserve on
      surplus hyperedges of cyclic cores.
\end{enumerate}
\end{programme}

\begin{firewall}[Claim boundary after V75]
\label{fire:newV75-boundary}
V75 performs the compulsory SM/NS audit, proves the exact bijection
between incidence hypertrees and rooted ordinary trees with child
partitions, sums all hyperedge blockings over each fixed rooted tree,
and reduces the aggregate hypertree problem to the explicit RHQC
inequality.  It also proves that dropping acyclicity produces a
superexponential parent-map majorant.

V75 does not prove RHQC, CIRC, non-star hypertree aggregation, MCWC
in full, OCRC in its full noncommutative form, SCNGC, WPSC, GRLTC,
LZTC, PLRC, WSDC, BRSC, HWSFC, UGCGC, thermal connected WPRAC,
all-arity RGSC, the raw Wilson aggregate strong card, all-order
strong MTA, WUFC, CPM, cutoff removal, reflection positivity,
Osterwalder--Schrader reconstruction, a continuum Yang--Mills
measure, or a positive Hamiltonian gap.  The full Yang--Mills mass
gap has not been proved.
\end{firewall}

\begin{theorem}[V75 result ledger]
\label{thm:newV75-results}
V75:
\begin{enumerate}[label=\textup{(V75.\arabic*)},leftmargin=6.5em]
\item performs the compulsory active-SM/active-NS audit;
\item imports only exact-assembly and hierarchy-boundary discipline;
\item proves the rooted-tree/child-partition bijection;
\item sums every hyperedge blocking over one rooted tree;
\item proves the strong card for each fixed rooted tree;
\item proves the parent-map relaxation is superexponential;
\item formulates RHQC; and
\item proves RHQC is sufficient for every incidence hypertree.
\end{enumerate}
\end{theorem}

\begin{proof}
Use
\Cref{audit:newV75-companions,
thm:newV75-bijection,
lem:newV75-fixed-tree,
prop:newV75-parent-no-go,
prop:newV75-RHQC-closes}.
\end{proof}

\paragraph{Parent-version record.}
V75 was formed from
\texttt{Renewal\_Yang\_Mills\_Mass\_Gap\_Research\_Notes\_V74.tex}.
The V74 parent had \(34\,855\) logical source lines,
\(1\,110\,936\) bytes, and SHA--256
\[
 \texttt{5424cc6c7c96795590ac0ce8ae546f22f9670ba417d57cef422d2b2523fc8f15}.
\]
The next compulsory companion audit is V80.

% =============================================================================
% END APPENDED FIFTH-SERIES V75 MATERIAL
% =============================================================================

\clearpage
\chapter{Fifth-series V76: balanced centre marks and the
coefficientwise hypertree target}
\label{ch:newV76-balanced}

\section{The normalized rooted-hypertree polynomial}
\label{sec:newV76-polynomial}

Keep the root \(o\in{\cal C}\) fixed, write \(p=|{\cal C}|\ge2\),
and put \(B=\sum_{i\in{\cal C}}b_i\).  The quotient left after the
common product of all centre marks has been removed from
\eqref{eq:newV75-RHQC} is
\begin{equation}
 {\cal Q}_{p,o}(\mathbf b)
 :=
 \sum_{T\in\mathfrak T_o({\cal C})}
 \left(\prod_{i\in I(T)\setminus\{o\}}b_i\right)
 \prod_{i\in I(T)}
 \left(b_i+\sum_{j\in{\rm Ch}_T(i)}b_j\right)^{q_i-1}.
 \label{eq:newV76-Q}
\end{equation}
Thus RHQC is precisely
\begin{equation}
 {\cal Q}_{p,o}(\mathbf b)\le K^p B^{p-2}.
 \label{eq:newV76-Q-target}
\end{equation}

\begin{lemma}[Positivity, degree, and coefficient mass]
\label{lem:newV76-Q-basic}
The polynomial \({\cal Q}_{p,o}\) has nonnegative coefficients and
is homogeneous of degree \(p-2\).  Moreover, if
\[
 a(T):=\prod_{i\in I(T)}(q_i+1)^{q_i-1},
\]
then
\begin{equation}
 {\cal Q}_{p,o}(\mathbf 1)
 =
 \sum_{T\in\mathfrak T_o({\cal C})}a(T).
 \label{eq:newV76-coefficient-mass}
\end{equation}
\end{lemma}

\begin{proof}
Every factor in \eqref{eq:newV76-Q} is a polynomial with
nonnegative coefficients.  If \(h=|I(T)|\), its summand has degree
\[
 (h-1)+\sum_{i\in I(T)}(q_i-1)
 =
 (h-1)+(p-1-h)=p-2.
\]
At \(\mathbf b=\mathbf1\), the local sum at \(i\) is \(q_i+1\);
this gives \eqref{eq:newV76-coefficient-mass}.
\end{proof}

\section{A Prüfer estimate with the correct \(p^{p-2}\) scale}
\label{sec:newV76-prufer}

\begin{lemma}[Weighted child-count enumeration]
\label{lem:newV76-child-count}
For a fixed root \(o\) and \(p\ge2\),
\begin{equation}
 \sum_{T\in\mathfrak T_o({\cal C})}a(T)
 \le
 (8e^2)^p p^{p-2}.
 \label{eq:newV76-child-count-bound}
\end{equation}
\end{lemma}

\begin{proof}
For a rooted tree let \(q_i\) be the number of children of \(i\).
Then
\[
 q_o\ge1,\qquad q_i\ge0\ (i\ne o),\qquad
 \sum_iq_i=p-1.
\]
The undirected degrees are \(d_o=q_o\) and
\(d_i=q_i+1\) for \(i\ne o\).  The Prüfer degree formula therefore
shows that the number of labelled trees with the prescribed child
counts is
\begin{equation}
 { (p-2)!\over
   (q_o-1)!\prod_{i\ne o}q_i!}.
 \label{eq:newV76-degree-count}
\end{equation}

For \(r\ge0\), set
\[
 u_0=1,\qquad
 u_r={(r+1)^{r-1}\over r!}\quad(r\ge1),
\]
and, for \(r\ge1\), set
\[
 v_r={(r+1)^{r-1}\over(r-1)!}=r u_r.
\]
The elementary Stirling lower bound \(r!\ge(r/e)^r\) and
\((1+1/r)^r\le e\) give
\[
 u_r\le e^{2r}\quad(r\ge0),
 \qquad
 v_r\le(2e^2)^r\quad(r\ge1),
\]
where \(r\le2^r\) was used in the second estimate.  Consequently,
for every admissible child-count vector,
\[
 v_{q_o}\prod_{i\ne o}u_{q_i}
 \le(2e^2)^{p-1}.
\]

After writing \(q_o'=q_o-1\), the number of admissible vectors is
\[
 {2p-3\choose p-1}\le4^{p-1}.
\]
Multiplying \eqref{eq:newV76-degree-count} by \(a(T)\), summing the
child-count vectors, and using \((p-2)!\le p^{p-2}\) yields
\[
 \sum_Ta(T)
 \le
 (p-2)!(8e^2)^{p-1}
 \le(8e^2)^p p^{p-2}.
\]
\end{proof}

\section{Closure of the balanced-mark sector}
\label{sec:newV76-balanced}

\begin{theorem}[Balanced-mark RHQC]
\label{thm:newV76-balanced-RHQC}
Assume that, for some \(\Lambda\ge1\),
\begin{equation}
 \max_{i\in{\cal C}}b_i\le{\Lambda B\over p}.
 \label{eq:newV76-balance}
\end{equation}
Then
\begin{equation}
 {\cal Q}_{p,o}(\mathbf b)
 \le(8e^2\Lambda)^p B^{p-2}.
 \label{eq:newV76-balanced-Q}
\end{equation}
Hence RHQC, and therefore the complete incidence-hypertree core
card, holds uniformly in every fixed-\(\Lambda\) balanced-mark
sector.
\end{theorem}

\begin{proof}
By \Cref{lem:newV76-Q-basic}, each tree summand is a homogeneous
polynomial of degree \(p-2\) with nonnegative coefficients.
Under \eqref{eq:newV76-balance}, every monomial in that summand is
at most
\[
 \left({\Lambda B\over p}\right)^{p-2}
\]
times its coefficient.  Therefore
\[
 {\cal Q}_{p,o}(\mathbf b)
 \le
 \left({\Lambda B\over p}\right)^{p-2}
 {\cal Q}_{p,o}(\mathbf1).
\]
Apply \Cref{lem:newV76-child-count}; since \(\Lambda\ge1\),
\[
 {\cal Q}_{p,o}(\mathbf b)
 \le
 (8e^2)^p\Lambda^{p-2}B^{p-2}
 \le(8e^2\Lambda)^pB^{p-2}.
\]
The last assertion follows from
\Cref{prop:newV75-RHQC-closes}, with the pre-existing numerical
factor \(K^{2p}\) absorbed into the exponential base.
\end{proof}

\begin{corollary}[Equal centre marks]
\label{cor:newV76-equal}
If \(b_i=b\) for every \(i\in{\cal C}\), then
\[
 {\cal Q}_{p,o}(b,\ldots,b)
 \le(8e^2)^p(pb)^{p-2}.
\]
Thus arbitrary incidence hypertrees are closed in the equal-mark
model.
\end{corollary}

\begin{proof}
Here \(B=pb\) and \eqref{eq:newV76-balance} holds with
\(\Lambda=1\).
\end{proof}

\section{Exact concentration endpoint and the next quotient}
\label{sec:newV76-coefficients}

For a multi-index
\(\alpha=(\alpha_i)_{i\in{\cal C}}\in{\mathbb N}^{\cal C}\), write
\(\mathbf b^\alpha=\prod_i b_i^{\alpha_i}\) and
\[
 {p-2\choose\alpha}:={(p-2)!\over\prod_i\alpha_i!}
 \quad\hbox{when}\quad |\alpha|=p-2.
\]

\begin{proposition}[Pure-power coefficients]
\label{prop:newV76-pure-power}
For \(p\ge3\) and \(h\in{\cal C}\),
\begin{equation}
 [b_h^{p-2}]\,{\cal Q}_{p,o}
 =
 \begin{cases}
  1,&h=o,\\
  2^{p-2},&h\ne o.
 \end{cases}
 \label{eq:newV76-pure-power}
\end{equation}
\end{proposition}

\begin{proof}
Suppose first that \(h=o\).  A contribution to \(b_o^{p-2}\)
cannot contain a mandatory factor \(b_i\) from a nonroot internal
vertex.  Hence \(I(T)=\{o\}\), so \(T\) is the unique star rooted at
\(o\).  Selecting \(b_o\) from every factor of
\(B_o^{p-2}\) gives coefficient one.

Now let \(h\ne o\).  Every nonroot internal vertex supplies its
mandatory mark in \eqref{eq:newV76-Q}; hence the only possible
nonroot internal vertex is \(h\).  The root star contributes one.
Otherwise \(I(T)=\{o,h\}\).  Then \(h\) is a child of \(o\), and
each of the other \(p-2\) vertices is a leaf attached either to
\(o\) or to \(h\), with at least one attached to \(h\).  There are
\(2^{p-2}-1\) such trees.  If \(R\) and \(S\) are respectively the
leaves attached to \(o\) and \(h\), the corresponding quotient
summand is
\[
 b_h(B_o^T)^{|R|}(B_h^T)^{|S|-1}.
\]
The mark \(b_h\) occurs in both displayed local sums, and choosing
it in every factor contributes coefficient one.  Adding the root
star gives \(2^{p-2}\).
\end{proof}

\begin{definition}[Coefficientwise quotient domination card, CQDC]
\label{def:newV76-CQDC}
CQDC is the coefficient estimate
\begin{equation}
 [\mathbf b^\alpha]\,{\cal Q}_{p,o}
 \le
 K^p {p-2\choose\alpha}
 \qquad(|\alpha|=p-2),
 \label{eq:newV76-CQDC}
\end{equation}
uniformly in \(p,o,\alpha\).
\end{definition}

\begin{proposition}[CQDC implies full RHQC]
\label{prop:newV76-CQDC-implies}
CQDC implies \eqref{eq:newV76-Q-target} for arbitrary nonnegative
centre marks.
\end{proposition}

\begin{proof}
Positivity and homogeneity give
\[
 {\cal Q}_{p,o}(\mathbf b)
 =
 \sum_{|\alpha|=p-2}
 [\mathbf b^\alpha]{\cal Q}_{p,o}\,\mathbf b^\alpha.
\]
Apply \eqref{eq:newV76-CQDC} and the multinomial theorem:
\[
 {\cal Q}_{p,o}(\mathbf b)
 \le
 K^p\sum_{|\alpha|=p-2}
 {p-2\choose\alpha}\mathbf b^\alpha
 =
 K^pB^{p-2}.
\]
\end{proof}

\begin{assessment}[What the endpoint calculation says]
\label{ass:newV76-endpoint}
The balanced theorem uses only total coefficient mass and is
therefore sharp enough away from concentration but not at a
heterogeneous endpoint.  The exact coefficient
\eqref{eq:newV76-pure-power} shows two facts.  First, concentration
does not itself create a superexponential obstruction.  Second, a
coefficientwise proof must allow an exponential multiplicity of at
least \(2^{p-2}\); an injectivity claim with constant one would be
false.  CQDC is now the exact purely combinatorial live row.
\end{assessment}

\begin{programme}[V77: decorated Prüfer fibres]
\label{prog:newV76-V77}
\begin{enumerate}[leftmargin=3em]
\item Expand each local power in \eqref{eq:newV76-Q}; regard a
      monomial contribution as a rooted tree with \(p-2\) marked
      local choices.
\item Map that decorated tree to a Prüfer word with the same mark
      content \(\alpha\).
\item Prove that every word has at most \(K^p\) preimages.  The
      \(2^{p-2}\) family in
      \Cref{prop:newV76-pure-power} is the calibration case.
\item If direct deletion is not finite-state, stratify by the
      support of \(\alpha\) and keep the forest induced by that
      support rather than replacing it by an arbitrary parent map.
\item Only after CQDC is obtained, return to the \(B^{2g}\) reserve
      for cyclic centre cores.
\end{enumerate}
\end{programme}

\begin{firewall}[Claim boundary after V76]
\label{fire:newV76-boundary}
V76 rewrites RHQC as a positive homogeneous quotient polynomial,
proves the correctly scaled Prüfer coefficient-mass estimate, and
closes RHQC for every fixed-balance centre-mark sector, including
equal marks.  It also computes the exact pure-power coefficients,
formulates CQDC, and proves that CQDC would imply full RHQC.

V76 does not prove CQDC, full heterogeneous RHQC, CIRC, cyclic
centre-core aggregation, MCWC in full, OCRC in its full
noncommutative form, SCNGC, WPSC, GRLTC, LZTC, PLRC, WSDC, BRSC,
HWSFC, UGCGC, thermal connected WPRAC, all-arity RGSC, the raw
Wilson aggregate strong card, all-order strong MTA, WUFC, CPM,
cutoff removal, reflection positivity, Osterwalder--Schrader
reconstruction, a continuum Yang--Mills measure, or a positive
Hamiltonian gap.  The full Yang--Mills mass gap has not been proved.
\end{firewall}

\begin{theorem}[V76 result ledger]
\label{thm:newV76-results}
V76:
\begin{enumerate}[label=\textup{(V76.\arabic*)},leftmargin=6.5em]
\item defines the normalized rooted-hypertree quotient polynomial;
\item proves its positivity and exact degree \(p-2\);
\item proves an exponential-times-\(p^{p-2}\) Prüfer mass bound;
\item closes RHQC for uniformly balanced centre marks;
\item closes the equal-mark incidence-hypertree model;
\item computes every pure-power coefficient exactly;
\item formulates CQDC; and
\item proves CQDC is sufficient for full RHQC.
\end{enumerate}
\end{theorem}

\begin{proof}
Use
\Cref{lem:newV76-Q-basic,
lem:newV76-child-count,
thm:newV76-balanced-RHQC,
cor:newV76-equal,
prop:newV76-pure-power,
prop:newV76-CQDC-implies}.
\end{proof}

\paragraph{Parent-version record.}
V76 was formed from
\texttt{Renewal\_Yang\_Mills\_Mass\_Gap\_Research\_Notes\_V75.tex}.
The V75 parent had \(35\,212\) logical source lines,
\(1\,123\,415\) bytes, and SHA--256
\[
 \texttt{0c44941e3b7448031f6c81f9ba36844ac17b4d6e856c747bb87941e0fa3c1100}.
\]
The next compulsory companion audit is V80.

% =============================================================================
% END APPENDED FIFTH-SERIES V76 MATERIAL
% =============================================================================

\clearpage
\chapter{Fifth-series V77: two-level Cayley gluing and full RHQC}
\label{ch:newV77-gluing}

\section{Local Cayley trees}
\label{sec:newV77-local}

For a rooted ordinary tree \(T\), retain the notation
\[
 S_i^T:=\{i\}\cup{\rm Ch}_T(i),
 \qquad |S_i^T|=q_i+1
 \quad(i\in I(T)).
\]
The weighted Cayley--Pr\"ufer identity, already used in
\Cref{cor:newV60-Prufer-form}, gives the polynomial identity
\begin{equation}
 \left(\sum_{v\in S}b_v\right)^{|S|-2}
 =
 \sum_{L\in\mathfrak T(S)}
 \prod_{v\in S}b_v^{d_L(v)-1}
 \qquad(|S|\ge2).
 \label{eq:newV77-local-Cayley}
\end{equation}
Therefore every local factor in \eqref{eq:newV76-Q} is exactly a
sum over an ordinary tree \(L_i\) on \(S_i^T\), rather than merely
an upper bound.

\begin{lemma}[Gluing the local Cayley trees]
\label{lem:newV77-gluing}
Fix \(T\in\mathfrak T_o({\cal C})\) and choose
\(L_i\in\mathfrak T(S_i^T)\) for every \(i\in I(T)\).
Then
\[
 U:=\bigcup_{i\in I(T)}L_i
\]
is an ordinary tree on \({\cal C}\).  Moreover,
\begin{equation}
 \left(\prod_{i\in I(T)\setminus\{o\}}b_i\right)
 \prod_{i\in I(T)}
 \prod_{v\in S_i^T}b_v^{d_{L_i}(v)-1}
 =
 \prod_{v\in{\cal C}}b_v^{d_U(v)-1}.
 \label{eq:newV77-weight-gluing}
\end{equation}
\end{lemma}

\begin{proof}
Two distinct closed child stars \(S_i^T,S_j^T\) intersect in at
most one vertex.  Hence the local trees have disjoint edge sets.
Their union is connected: the root star connects \(o\) to all
first-generation vertices, and each later local tree attaches at
its own root to the already connected preceding generations.
Furthermore,
\[
 |E(U)|=\sum_{i\in I(T)}q_i=p-1.
\]
Thus \(U\) is a tree.

Let \(m_v\) be the number of local child stars containing \(v\).
The root has \(m_o=1\).  A nonroot leaf of \(T\) occurs only in its
parent's star and also has \(m_v=1\).  A nonroot internal vertex
occurs in its parent's star and its own star, so \(m_v=2\).
Disjointness of the local edge sets gives
\[
 \sum_{i:\,v\in S_i^T}\{d_{L_i}(v)-1\}
 =d_U(v)-m_v.
\]
The mandatory factor in the first product on the left of
\eqref{eq:newV77-weight-gluing} adds one precisely in the last
case.  The resulting exponent is \(d_U(v)-1\) in all three cases.
\end{proof}

\section{The fibre over one global tree}
\label{sec:newV77-fibre}

\begin{definition}[Cayley-gluing fibre]
\label{def:newV77-fibre}
For \(U\in\mathfrak T({\cal C})\), let \(M_o(U)\) be the number of
pairs
\[
 \left(T,(L_i)_{i\in I(T)}\right)
\]
from \Cref{lem:newV77-gluing} whose local edge union is \(U\).
\end{definition}

\begin{lemma}[Sharp exponential fibre bound]
\label{lem:newV77-fibre-bound}
For every \(U\in\mathfrak T({\cal C})\),
\begin{equation}
 M_o(U)\le2^{p-2}.
 \label{eq:newV77-fibre-bound}
\end{equation}
\end{lemma}

\begin{proof}
A point of the fibre partitions \(E(U)\) into the nonempty,
connected edge sets \(E(L_i)\).  At a vertex \(v\), at most two
parts occur: the part belonging to the parent star of \(v\), and,
when \(v\) is internal in \(T\), the part belonging to its own
star.  The root occurs in only its own part.

This global edge partition determines the original pair.  Indeed,
its connected parts form a tree under intersection.  Root that
part-intersection tree at the unique part containing \(o\).
Every nonroot part meets its parent part in one vertex, and that
vertex is the root label \(i\) of the corresponding \(L_i\).
Every other vertex of that part is a child of \(i\) in \(T\).
The part itself then recovers \(L_i\).  Thus the map from fibre
points to such global edge partitions is injective.

It remains to count a larger class containing all these partitions.
At every vertex \(v\), partition its \(d_U(v)\) incident edges
according to their global edge part.  There are at most two local
classes.  A global partition is determined by all these local
partitions: two edges belong to the same connected global part
exactly when they can be joined by a chain of locally equivalent
incident edges.  The number of set partitions of a \(d\)-element
set into at most two nonempty classes is
\[
 {d\brace1}+{d\brace2}
 =
 1+{2^d-2\over2}=2^{d-1}.
\]
Consequently,
\[
 M_o(U)
 \le
 \prod_{v\in{\cal C}}2^{d_U(v)-1}
 =
 2^{\,2(p-1)-p}
 =
 2^{p-2}.
\]
\end{proof}

\begin{remark}[Calibration]
\label{rem:newV77-calibration}
When \(U\) is the star centred at a nonroot vertex \(h\), the fibre
has exactly \(2^{p-2}\) elements: each other nonroot vertex chooses
whether it is attached at the first or second rooted level.  Thus
\eqref{eq:newV77-fibre-bound} is sharp and recovers
\Cref{prop:newV76-pure-power}.
\end{remark}

\section{Coefficientwise domination and heterogeneous closure}
\label{sec:newV77-closure}

\begin{theorem}[Exact global-tree representation]
\label{thm:newV77-global-tree}
For arbitrary nonnegative marks,
\begin{equation}
 {\cal Q}_{p,o}(\mathbf b)
 =
 \sum_{U\in\mathfrak T({\cal C})}
 M_o(U)\prod_{v\in{\cal C}}b_v^{d_U(v)-1}.
 \label{eq:newV77-global-tree}
\end{equation}
\end{theorem}

\begin{proof}
Insert \eqref{eq:newV77-local-Cayley} in every local factor of
\eqref{eq:newV76-Q}.  The resulting terms are exactly the pairs
in \Cref{def:newV77-fibre}; their weights are transformed by
\Cref{lem:newV77-gluing}.  Group equal output trees \(U\).
\end{proof}

\begin{theorem}[CQDC and full heterogeneous RHQC]
\label{thm:newV77-CQDC-RHQC}
For every \(p\ge2\), root \(o\), and nonnegative centre marks,
\begin{equation}
 {\cal Q}_{p,o}(\mathbf b)
 \le2^{p-2}B^{p-2}.
 \label{eq:newV77-full-RHQC}
\end{equation}
In particular, CQDC holds with exponential base \(2\), and RHQC
holds without a balance hypothesis.
\end{theorem}

\begin{proof}
Combine \Cref{thm:newV77-global-tree,lem:newV77-fibre-bound}:
\[
 {\cal Q}_{p,o}(\mathbf b)
 \le
 2^{p-2}
 \sum_{U\in\mathfrak T({\cal C})}
 \prod_vb_v^{d_U(v)-1}.
\]
The weighted Cayley--Pr\"ufer identity evaluates the last sum as
\(B^{p-2}\), proving \eqref{eq:newV77-full-RHQC}.

The same argument is coefficientwise.  Since
\[
 \sum_U\prod_vb_v^{d_U(v)-1}=B^{p-2},
\]
the coefficient of \(\mathbf b^\alpha\) on the tree sum is
\({p-2\choose\alpha}\).  Hence
\[
 [\mathbf b^\alpha]{\cal Q}_{p,o}
 \le
 2^{p-2}{p-2\choose\alpha}
 \le
 2^p{p-2\choose\alpha},
\]
which is \eqref{eq:newV76-CQDC} with \(K=2\).
\end{proof}

\begin{corollary}[All incidence-hypertree cores close]
\label{cor:newV77-hypertrees-close}
The complete sum of every connected centre core whose incidence
graph is a tree obeys the strong centre-card target, uniformly in
the number of centres, the arities of all hyperedges, and the
heterogeneity of the centre marks.
\end{corollary}

\begin{proof}
\Cref{thm:newV77-CQDC-RHQC} proves RHQC.  Apply
\Cref{prop:newV75-RHQC-closes}; all local hyperedge blockings,
noncentral-label allocations, and Duhamel orderings were already
included there.
\end{proof}

\begin{assessment}[The acyclic obstruction is removed]
\label{ass:newV77-position}
The Bell obstruction in \Cref{warn:newV74-Bell} arose only after
the local hyperedge reserve was replaced by a fixed-core global
power.  The two-level Cayley construction keeps that reserve until
all hyperedge blockings have become local trees.  Their union has
exactly \(p-1\) edges, hence one global Cayley tree; only an
edge-partition fibre remains, and
\Cref{lem:newV77-fibre-bound} pays it with the sharp exponential
factor \(2^{p-2}\).  No balance condition and no parent-map
relaxation is used.
\end{assessment}

\section{The next core: incidence cycles}
\label{sec:newV77-cycles}

\begin{programme}[V78: mark and cut surplus incidences]
\label{prog:newV77-V78}
\begin{enumerate}[leftmargin=3em]
\item For a connected core of excess \(g\ge1\), choose a canonical
      set of \(g\) surplus incidence edges whose deletion leaves an
      incidence hypertree.
\item Retain the \(B^{2g}\) reserve of
      \Cref{cor:newV74-cycle-reserve} until the surplus attachments
      are marked.
\item Apply the V77 two-level Cayley gluing to the residual
      hypertree, now with \(2g\) distinguished attachment slots.
\item Prove a coefficientwise attachment bound of the form
      \(K^{p+g}B^{2g}\), rather than counting all possible cuts
      independently.
\item Sum \(g\) only after the Wilson-side small-\(B\) factor has
      been restored.
\end{enumerate}
\end{programme}

\begin{firewall}[Claim boundary after V77]
\label{fire:newV77-boundary}
V77 proves the exact local-to-global Cayley gluing identity, the
sharp \(2^{p-2}\) fibre bound, CQDC, full heterogeneous RHQC, and
the strong aggregate card for every incidence-hypertree centre
core.

V77 does not prove cyclic CIRC, the surplus-incidence attachment
sum, MCWC in full, OCRC in its full noncommutative form, SCNGC,
WPSC, GRLTC, LZTC, PLRC, WSDC, BRSC, HWSFC, UGCGC, thermal
connected WPRAC, all-arity RGSC, the raw Wilson aggregate strong
card, all-order strong MTA, WUFC, CPM, cutoff removal, reflection
positivity, Osterwalder--Schrader reconstruction, a continuum
Yang--Mills measure, or a positive Hamiltonian gap.  The full
Yang--Mills mass gap has not been proved.
\end{firewall}

\begin{theorem}[V77 result ledger]
\label{thm:newV77-results}
V77:
\begin{enumerate}[label=\textup{(V77.\arabic*)},leftmargin=6.5em]
\item expands every local child-star power into weighted Cayley
      trees;
\item proves that all local trees glue to one global tree;
\item proves exact preservation of the Prüfer monomial;
\item identifies the gluing fibre with constrained connected
      edge partitions;
\item proves the sharp fibre bound \(2^{p-2}\);
\item proves the exact global-tree representation;
\item proves CQDC and full heterogeneous RHQC; and
\item closes every incidence-hypertree centre core.
\end{enumerate}
\end{theorem}

\begin{proof}
Use
\Cref{lem:newV77-gluing,
lem:newV77-fibre-bound,
thm:newV77-global-tree,
thm:newV77-CQDC-RHQC,
cor:newV77-hypertrees-close}.
\end{proof}

\paragraph{Parent-version record.}
V77 was formed from
\texttt{Renewal\_Yang\_Mills\_Mass\_Gap\_Research\_Notes\_V76.tex}.
The V76 parent had \(35\,574\) logical source lines,
\(1\,134\,257\) bytes, and SHA--256
\[
 \texttt{06c211fb95198be15e8e6776a18374af164182b38114101d4677db294329244f}.
\]
The next compulsory companion audit is V80.

% =============================================================================
% END APPENDED FIFTH-SERIES V77 MATERIAL
% =============================================================================

\clearpage
\chapter{Fifth-series V78: canonical incidence skeletons and
closure of CIRC}
\label{ch:newV78-CIRC}

\section{Two pointwise effective-hyperedge proxies}
\label{sec:newV78-proxies}

Fix a centre set \({\cal C}\), put \(p=|{\cal C}|\ge2\), and write
\[
 B_{\cal C}:=\sum_{i\in{\cal C}}b_i.
\]
Every centre mark obeys \(b_i\ge\sqrt s\), hence
\begin{equation}
 p\sqrt s\le B_{\cal C}.                                   \label{eq:newV78-floor}
\end{equation}
For \(C\subseteq{\cal C}\), \(d=|C|\ge2\), introduce
\begin{align}
 {\mathfrak r}(C)
 &:=(Kd)^d s^{(d-2)/2}\prod_{i\in C}b_i,                    \label{eq:newV78-reserve-proxy}\\
 {\mathfrak q}(C)
 &:=K^d\left(\prod_{i\in C}b_i\right)B_C^{d-2}.             \label{eq:newV78-strong-proxy}
\end{align}
After one enlargement of \(K\),
\Cref{lem:newV74-effective-reserve,lem:newV73-effective} give,
uniformly on the adaptive coefficient circle and also at \(z=0\),
\begin{equation}
 |W_C(z)|\le\min\{{\mathfrak r}(C),{\mathfrak q}(C)\}.
 \label{eq:newV78-two-proxies}
\end{equation}

\begin{lemma}[Summed whole-hyperedge activity]
\label{lem:newV78-whole-gas}
If \(B_{\cal C}\le b_*\), then
\begin{equation}
 \sum_{\substack{C\subseteq{\cal C}\\|C|\ge2}}
 {\mathfrak q}(C)
 \le C B_{\cal C}^2.                                      \label{eq:newV78-whole-gas}
\end{equation}
Consequently the positive sum over an arbitrary set of whole
surplus hyperedges is at most \(\exp(CB_{\cal C}^2)\).
\end{lemma}

\begin{proof}
For \(d\ge2\),
\[
 \sum_{|C|=d}\prod_{i\in C}b_i
 \le {B_{\cal C}^d\over d!}.
\]
Since \(B_C\le B_{\cal C}\),
\[
 \sum_{|C|=d}{\mathfrak q}(C)
 \le {K^dB_{\cal C}^{2d-2}\over d!}.
\]
Sum \(d\ge2\), decreasing \(b_*\) so that \(KB_{\cal C}^2\le1\).
For a set \({\cal F}\) of surplus hyperedges, positivity gives
\[
 \sum_{\cal F}\prod_{C\in{\cal F}}{\mathfrak q}(C)
 =
 \prod_C(1+{\mathfrak q}(C))
 \le
 \exp\left\{\sum_C{\mathfrak q}(C)\right\}.
\]
\end{proof}

\section{Extending one skeleton hyperedge}
\label{sec:newV78-extension}

\begin{lemma}[All surplus incidences on one retained hyperedge]
\label{lem:newV78-extension}
Let \(C\subseteq{\cal C}\), \(d=|C|\ge2\).  Then
\begin{equation}
 \sum_{J\subseteq{\cal C}\setminus C}
 {{\mathfrak r}(C\cup J)\over{\mathfrak r}(C)}
 \le \exp(CB_{\cal C}^2).
 \label{eq:newV78-extension}
\end{equation}
\end{lemma}

\begin{proof}
Put \(k=|J|\) and \(n=d+k\).  Directly from
\eqref{eq:newV78-reserve-proxy},
\begin{align*}
 {{\mathfrak r}(C\cup J)\over{\mathfrak r}(C)}
 &=
 K^k{n^n\over d^d}s^{k/2}\prod_{j\in J}b_j\\
 &\le
 (eKp\sqrt s)^k\prod_{j\in J}b_j.
\end{align*}
Indeed,
\[
 {n^n\over d^d}
 =
 n^k\left(1+{k\over d}\right)^d
 \le p^ke^k.
\]
Therefore
\begin{align*}
 \sum_{J\subseteq{\cal C}\setminus C}
 {{\mathfrak r}(C\cup J)\over{\mathfrak r}(C)}
 &\le
 \prod_{j\notin C}(1+eKp\sqrt s\,b_j)\\
 &\le
 \exp(eKp\sqrt s\,B_{\cal C})
 \le
 \exp(eK B_{\cal C}^2),
\end{align*}
where \eqref{eq:newV78-floor} was used last.
\end{proof}

\begin{assessment}[Why incidence additions do not cost
\(p^{|J|}\)]
\label{ass:newV78-extension}
The size ratio of the unrelaxed reserve contributes at most
\((Kp)^{|J|}\), but the same added incidence contributes
\(\sqrt s\,b_j\).  Summing the distinct added centres produces the
product in \eqref{eq:newV78-extension}, and
\(p\sqrt s\le B_{\cal C}\) converts the apparent arity loss into
the genuine cycle degree \(B_{\cal C}^2\).  No incidence is chosen
independently after its centre mark has been discarded.
\end{assessment}

\section{Canonical incidence-skeleton decomposition}
\label{sec:newV78-skeleton}

\begin{lemma}[Hypertree skeleton of every connected core]
\label{lem:newV78-skeleton}
Let \({\cal H}\) be a connected simple hypergraph on
\({\cal C}\), all of whose hyperedges have cardinality at least
two.  There is an injective assignment
\begin{equation}
 {\cal H}\longmapsto
 \left(
 {\cal H}_0,\ (J_C)_{C\in{\cal H}_0},\ {\cal F}
 \right)                                                   \label{eq:newV78-skeleton-data}
\end{equation}
with the following properties:
\begin{enumerate}[label=\textup{(\roman*)},leftmargin=3.5em]
\item \({\cal H}_0\) is an incidence hypertree on all \(p\)
      centres and \(|{\cal H}_0|\le p-1\);
\item \(J_C\subseteq{\cal C}\setminus C\);
\item \({\cal F}\) is a set of whole surplus hyperedges; and
\item
      \[
      {\cal H}
      =
      \{C\cup J_C:C\in{\cal H}_0\}\cup{\cal F}.
      \]
\end{enumerate}
\end{lemma}

\begin{proof}
Fix total orders on the centre and hyperedge labels, and let
\(\tau({\cal H})\) be the lexicographically first spanning tree of
the bipartite incidence graph.  Let \({\cal F}\) consist of those
hyperedge vertices having degree one in \(\tau({\cal H})\).
Delete these leaf vertices from \(\tau({\cal H})\).

The remaining graph is a tree and still contains every centre
vertex.  To see the latter assertion, a centre lost from the
nontrivial remaining component would have been adjacent in the
tree only to hyperedge leaves; then it and those leaves would form
a component separate from every other centre, contrary to
connectedness and \(p\ge2\).  Every retained hyperedge vertex has
at least two remaining incident centre vertices.  For such a
vertex \(A\), put
\[
 C_A:=N_{\tau({\cal H})}(A),
 \qquad
 J_{C_A}:=A\setminus C_A.
\]
The family \({\cal H}_0=\{C_A:A\notin{\cal F}\}\) has the remaining
tree as its incidence graph.  Its members are distinct: two
retained hyperedge vertices with the same centre set of size at
least two would create a four-cycle in that tree.  Thus
\({\cal H}_0\) is an incidence hypertree.  Since
\(\sum_{C\in{\cal H}_0}(|C|-1)=p-1\), it has at most \(p-1\)
hyperedges.

The original family is reconstructed from the displayed data.
Because the spanning tree was selected deterministically, the
assignment is injective.
\end{proof}

\section{Full centre-incidence resummation}
\label{sec:newV78-CIRC}

\begin{theorem}[CIRC]
\label{thm:newV78-CIRC}
On \(B_{\cal C}\le b_*\), uniformly on the adaptive coefficient
circle and at \(z=0\),
\begin{equation}
 \sum_{\substack{{\cal H}\ {\rm connected\ simple}\\
                  {\cal H}\ {\rm on}\ {\cal C}}}
 \prod_{A\in{\cal H}}|W_A(z)|
 \le
 K^p\left(\prod_{i\in{\cal C}}b_i\right)
 B_{\cal C}^{p-2}.
 \label{eq:newV78-CIRC}
\end{equation}
Thus the centre-incidence resummation card
\Cref{def:newV74-CIRC} holds.
\end{theorem}

\begin{proof}
Apply \Cref{lem:newV78-skeleton} to each connected core and use the
two different bounds in \eqref{eq:newV78-two-proxies}: use the
reserve proxy on every retained-and-extended skeleton hyperedge
and the strong proxy on every whole surplus hyperedge.  Positivity
and injectivity bound the left side of
\eqref{eq:newV78-CIRC} by
\begin{align*}
 &\sum_{{\cal H}_0\ {\rm incidence\ hypertree}}
   \prod_{C\in{\cal H}_0}{\mathfrak r}(C)
   \prod_{C\in{\cal H}_0}
   \left\{
    \sum_{J\subseteq{\cal C}\setminus C}
    {{\mathfrak r}(C\cup J)\over{\mathfrak r}(C)}
   \right\}\\
 &\qquad\times
   \sum_{\cal F}\prod_{A\in{\cal F}}{\mathfrak q}(A).
\end{align*}
By \Cref{lem:newV78-extension,lem:newV78-whole-gas} and
\(|{\cal H}_0|\le p-1\), the last two factors are at most
\[
 \exp(CpB_{\cal C}^2)\exp(CB_{\cal C}^2)\le K_1^p.
\]
The V77 incidence-hypertree theorem, specifically
\Cref{cor:newV77-hypertrees-close}, bounds the first sum by
\[
 K_2^p\left(\prod_{i\in{\cal C}}b_i\right)
 B_{\cal C}^{p-2}.
\]
Absorb all numerical exponentials into one base.
\end{proof}

\begin{corollary}[Small-mark MCWC core]
\label{cor:newV78-MCWC}
All connected centre-core incidence structures in the difficult
small-\(B\) regime of MCWC obey the strong thermal aggregate card.
This includes arbitrary numbers and arities of genuine hyperedges,
arbitrary pair-supported cycles, every noncentral-label allocation,
and every Duhamel ordering already incorporated in V70--V74.
\end{corollary}

\begin{proof}
Apply \Cref{prop:newV74-CIRC} to
\Cref{thm:newV78-CIRC}.  The global noncentral-label coefficient
variable and the stable Duhamel contraction are respectively those
of V72--V74 and \Cref{cor:newV70-Duhamel-BKAR}.
\end{proof}

\begin{assessment}[The fixed-excess ledger is no longer needed]
\label{ass:newV78-position}
\Cref{cor:newV74-cycle-reserve} correctly showed two thermal powers
per fixed incidence cycle, but summing cores by excess would still
require a separate census of all cut choices.  The canonical
skeleton proof sums more directly.  Surplus incidences are summed
inside each retained hyperedge by
\Cref{lem:newV78-extension}; wholly surplus hyperedges exponentiate
by \Cref{lem:newV78-whole-gas}.  The incidence excess is therefore
summed to all orders without a \(p^{2g}\) or \(g!\) loss.
\end{assessment}

\begin{programme}[V79: promote the core closure through the
nonlinear dependency chain]
\label{prog:newV78-V79}
\begin{enumerate}[leftmargin=3em]
\item Recheck the exact reduction from the difficult small-\(B\)
      MCWC core to full MCWC; prove the complementary estimate
      explicitly if it is only implicit upstream.
\item Verify that every OCRC word class is represented by the
      effective-hyperedge core used in CIRC, with no omitted
      singleton or disconnected logarithmic sector.
\item If that audit passes, invoke
      \Cref{prop:newV60-OCRC-SCNGC,prop:newV61-reduction} to close
      SCNGC.
\item Keep the wrapped Poisson card WPSC separate; CIRC concerns
      only the small-clock source.
\item Then attack WPSC at the exact connected logarithm, preserving
      the factor \(e^{-q/s}\).
\end{enumerate}
\end{programme}

\begin{firewall}[Claim boundary after V78]
\label{fire:newV78-boundary}
V78 proves a canonical incidence-hypertree skeleton decomposition,
sums all surplus incidences on retained hyperedges, exponentiates
all whole surplus hyperedges, proves CIRC, and closes every
small-\(B\) MCWC centre core.

V78 does not yet promote this result to full MCWC, OCRC, or SCNGC;
it does not prove WPSC, GRLTC, LZTC, PLRC, WSDC, BRSC, HWSFC,
UGCGC, thermal connected WPRAC, all-arity RGSC, the raw Wilson
aggregate strong card, all-order strong MTA, WUFC, CPM, cutoff
removal, reflection positivity, Osterwalder--Schrader
reconstruction, a continuum Yang--Mills measure, or a positive
Hamiltonian gap.  The full Yang--Mills mass gap has not been proved.
\end{firewall}

\begin{theorem}[V78 result ledger]
\label{thm:newV78-results}
V78:
\begin{enumerate}[label=\textup{(V78.\arabic*)},leftmargin=6.5em]
\item defines reserve and strong proxies for each effective
      hyperedge;
\item proves the whole-surplus-hyperedge gas has activity
      \(O(B_{\cal C}^2)\);
\item sums all incidence extensions of one skeleton hyperedge;
\item proves the canonical incidence-skeleton decomposition;
\item proves that every skeleton has at most \(p-1\) retained
      hyperedges;
\item proves full CIRC;
\item sums all incidence excesses without a cut census; and
\item closes the difficult small-\(B\) MCWC centre core.
\end{enumerate}
\end{theorem}

\begin{proof}
Use
\Cref{lem:newV78-whole-gas,
lem:newV78-extension,
lem:newV78-skeleton,
thm:newV78-CIRC,
cor:newV78-MCWC}.
\end{proof}

\paragraph{Parent-version record.}
V78 was formed from
\texttt{Renewal\_Yang\_Mills\_Mass\_Gap\_Research\_Notes\_V77.tex}.
The V77 parent had \(35\,887\) logical source lines,
\(1\,144\,602\) bytes, and SHA--256
\[
 \texttt{d9416f63f155cc90a99bbc6c684301368998ab361683bd9ca71b29782a83d224}.
\]
The next compulsory companion audit is V80.

% =============================================================================
% END APPENDED FIFTH-SERIES V78 MATERIAL
% =============================================================================

\clearpage
\chapter{Fifth-series V79: promotion audit and the intermediate
thermal wedge}
\label{ch:newV79-promotion}

\section{What CIRC proves and what it does not}
\label{sec:newV79-audit}

No companion audit is due at V79.  The logical chain inherited from
V60--V61 and V74 is
\[
 \text{CIRC on a declared mark region}
 \Longrightarrow
 \text{MCWC on that region}
 \Longrightarrow
 \text{the corresponding SCNGC sector}.
\]
\Cref{thm:newV78-CIRC} establishes the first implication only for
the difficult small-total-mark region \(B\le b_*\).  The
fixed-arity theorem \Cref{thm:newV50-fixed-card} cannot fill the
complement at all arities because its constant is \(C_m\), with no
proof that \(C_m\le K^m\).  We therefore establish below the part of
the complement which follows uniformly from the complete operator
source.

\section{The small-clock source has the V44 analytic disk}
\label{sec:newV79-source}

For the exact small-clock probability measure
\(\nu_\alpha^{\rm sm}\) of
\eqref{eq:newV59-nu-sm}, set
\[
 Y_s(U):={\theta(U)\over\sqrt s}.
\]

\begin{lemma}[Small-clock radial exponential moment]
\label{lem:newV79-radial}
After restricting to the continuum regime \(0<s\le s_0\), there
are numerical \(\lambda_0,C>0\) such that
\begin{equation}
 \mathbb E_{\nu_\alpha^{\rm sm}}e^{\lambda_0Y_s}\le C.
 \label{eq:newV79-radial}
\end{equation}
\end{lemma}

\begin{proof}
The positive Wilson split \eqref{eq:newV59-Wilson-split} gives, for
every nonnegative \(f\),
\[
 (1-\varepsilon_\alpha)
 \mathbb E_{\nu_\alpha^{\rm sm}}f
 \le
 \mathbb E_{w_\alpha dU}f.
\]
By \eqref{eq:newV59-rare-event}, decrease \(s_0\) so that
\(\varepsilon_\alpha\le1/2\).  The Wilson endpoint case of
\Cref{lem:newV44-radial-exp} bounds the right side with
\(f=e^{\lambda_0Y_s}\).  Divide by
\(1-\varepsilon_\alpha\).
\end{proof}

Define the squarefree small-clock operator source
\begin{equation}
 \mathbb Z_m^{\rm sm}(\mathbf z)
 :=
 \mathbb E_{\nu_\alpha^{\rm sm}}
 \bigotimes_{i=1}^m
 \{I_{R_i}+z_i(D^{R_i}(U)-I_{R_i})\}.
 \label{eq:newV79-source}
\end{equation}

\begin{theorem}[Crude all-arity small-clock source card]
\label{thm:newV79-crude}
There is a numerical \(K\) such that, for every \(m\ge2\),
\begin{equation}
 \left\|
 \kappa_m^{\nu_\alpha^{\rm sm}}
 (D^{R_1},\ldots,D^{R_m})
 \right\|_{\rm cb}
 \le
 (Km)^m\prod_{i=1}^mb_i.
 \label{eq:newV79-crude}
\end{equation}
The estimate is spectator-stable.
\end{theorem}

\begin{proof}
Repeat the source argument of V44 with
\(\mathbb Z_m^{\rm sm}\).  By
\Cref{lem:newV44-geodesic-increment,lem:newV79-radial},
\[
 \|\mathbb Z_m^{\rm sm}(\mathbf z)-I\|_{\rm cb}
 \le
 \mathbb E_{\nu_\alpha^{\rm sm}}
 \left[
  \exp\left(CY_s\sum_i|z_i|b_i\right)-1
 \right].
\]
Hence there is a numerical \(\eta>0\) for which
\(\sum_i|z_i|b_i\le\eta\) implies that the right side is at most
\(1/4\).  The principal Banach logarithm is then uniformly
bounded.  The squarefree logarithm argument in
\Cref{lem:newV44-log-cumulant} applies to every probability
measure and identifies the coefficient of \(z_1\cdots z_m\) with
the displayed tensor cumulant.  Cauchy extraction on
\[
 |z_i|={\eta\over mb_i}
\]
gives \eqref{eq:newV79-crude}.  Tensoring by a spectator identity
does not change the norm.
\end{proof}

\section{The high-average-mark sector}
\label{sec:newV79-high}

\begin{theorem}[Uniform closure when \(B\) is proportional to
arity]
\label{thm:newV79-high}
Fix \(\delta>0\).  On
\begin{equation}
 B:=\sum_{i=1}^mb_i\ge\delta m,                              \label{eq:newV79-high-region}
\end{equation}
the exact small-clock connected source obeys
\begin{equation}
 \left\|
 \kappa_m^{\nu_\alpha^{\rm sm}}
 (D^{R_1},\ldots,D^{R_m})
 \right\|_{\rm cb}
 \le
 K_\delta^m
 \left(\prod_{i=1}^mb_i\right)B^{m-2}.
 \label{eq:newV79-high}
\end{equation}
\end{theorem}

\begin{proof}
By \Cref{thm:newV79-crude}, it is enough to compare
\(m^m\) with \(B^{m-2}\).  From
\eqref{eq:newV79-high-region},
\[
 m^m=m^2m^{m-2}
 \le
 m^2\delta^{-(m-2)}B^{m-2}.
\]
Since \(m^2\le4^m\) for \(m\ge2\), all remaining factors enter
one base \(K_\delta\).
\end{proof}

\begin{corollary}[Two disjoint uniform regimes]
\label{cor:newV79-two-regimes}
The strong all-arity small-clock source card is proved in both
\[
 B\le b_*
 \qquad\hbox{and}\qquad
 B\ge\delta m.
\]
The first region is supplied by
\Cref{cor:newV78-MCWC}; the second by
\Cref{thm:newV79-high}.
\end{corollary}

\begin{proof}
Only the stated citations are used.  No fixed-arity constant is
promoted to an exponential base.
\end{proof}

\section{The genuine remaining wedge}
\label{sec:newV79-wedge}

\begin{definition}[Intermediate thermal wedge]
\label{def:newV79-wedge}
For fixed numerical \(b_*,\delta>0\), the intermediate wedge is
\begin{equation}
 b_*<B<\delta m.                                             \label{eq:newV79-wedge}
\end{equation}
\end{definition}

\begin{warning}[CIRC alone does not yet prove SCNGC]
\label{warn:newV79-no-promotion}
The current theorem chain contains no uniform all-arity estimate on
\eqref{eq:newV79-wedge}.  Therefore
\Cref{thm:newV78-CIRC} may not yet be promoted to full MCWC,
OCRC, or SCNGC.
\end{warning}

\begin{proof}
\Cref{thm:newV78-CIRC} assumes \(B\le b_*\).
\Cref{thm:newV79-high} assumes \(B\ge\delta m\).
For arbitrarily large \(m\), the interval
\((b_*,\delta m)\) is nonempty.  The direct source estimate there
still contains \(m^m\), and
\Cref{warn:newV50-not-uniform} forbids replacing the missing
all-arity argument by the fixed-\(m\) constant \(C_m\).
\end{proof}

\begin{definition}[Thermal-block renormalization card, TBRC]
\label{def:newV79-TBRC}
TBRC is an exact block-level source or BKAR construction on the
intermediate wedge with the following properties:
\begin{enumerate}[label=\textup{(\roman*)},leftmargin=3.5em]
\item the representation labels are partitioned into thermal blocks
      having total mark at most \(b_*\);
\item V78 CIRC is applied inside each occupied block before norms;
\item connections between occupied blocks are paid by one weighted
      tree on the block marks; and
\item summing the block partitions and the block tree yields
      \(K^m(\prod_i b_i)B^{m-2}\), with no Bell or \(m^m\) factor.
\end{enumerate}
\end{definition}

\begin{proposition}[TBRC is sufficient for the missing MCWC sector]
\label{prop:newV79-TBRC}
TBRC, together with
\Cref{cor:newV79-two-regimes}, proves full MCWC.  The inherited
reductions then prove SCNGC.
\end{proposition}

\begin{proof}
The three regions \(B\le b_*\),
\(b_*<B<\delta m\), and \(B\ge\delta m\) exhaust all mark vectors.
The first and third are
\Cref{cor:newV79-two-regimes}; TBRC is exactly the second.
Then use
\Cref{prop:newV74-CIRC,prop:newV61-reduction,
prop:newV60-OCRC-SCNGC}.
\end{proof}

\begin{assessment}[Upstream correction]
\label{ass:newV79-position}
The V60 phrase \emph{the only difficult small-\(B\) regime} was
not itself a uniform complementary theorem.  V79 replaces that
phrase by two proved regions and one explicit wedge.  The remaining
problem is no longer centre-hypergraph enumeration: V77--V78 closed
that operation.  It is the construction of an exact block source
which permits local CIRC resummations to be joined without losing
the squarefree logarithmic cancellation.
\end{assessment}

\begin{programme}[V80: compulsory audit and TBRC]
\label{prog:newV79-V80}
\begin{enumerate}[leftmargin=3em]
\item Perform the compulsory read-only audit of the newest active
      SM and NS notes and record immutable snapshots.
\item Search those notes only for an exact block-assembly or
      concentration mechanism transferable at the level of proof
      structure.
\item Construct a deterministic thermal partition with block sums
      in \([b_*/2,b_*]\), apart from one remainder block.
\item Test a block-level BKAR interpolation whose endpoints are the
      exact small-clock sources on unions of blocks.
\item Prove the block-edge derivative has one product of block
      marks before applying the weighted Cayley identity.
\end{enumerate}
\end{programme}

\begin{firewall}[Claim boundary after V79]
\label{fire:newV79-boundary}
V79 transfers the uniform radial exponential moment to the exact
small-clock source, proves its crude all-arity operator-source
bound, closes the high-average-mark region \(B\ge\delta m\), and
identifies the exact intermediate wedge left between that theorem
and V78 CIRC.  It formulates TBRC and proves its sufficiency.

V79 does not prove TBRC, full MCWC, OCRC, SCNGC, WPSC, GRLTC,
LZTC, PLRC, WSDC, BRSC, HWSFC, UGCGC, thermal connected WPRAC,
all-arity RGSC, the raw Wilson aggregate strong card, all-order
strong MTA, WUFC, CPM, cutoff removal, reflection positivity,
Osterwalder--Schrader reconstruction, a continuum Yang--Mills
measure, or a positive Hamiltonian gap.  The full Yang--Mills mass
gap has not been proved.
\end{firewall}

\begin{theorem}[V79 result ledger]
\label{thm:newV79-results}
V79:
\begin{enumerate}[label=\textup{(V79.\arabic*)},leftmargin=6.5em]
\item audits the exact CIRC-to-SCNGC dependency chain;
\item transfers the Wilson radial exponential moment to
      \(\nu_\alpha^{\rm sm}\);
\item proves the crude all-arity small-clock source card;
\item closes the region \(B\ge\delta m\);
\item combines it with the V78 region \(B\le b_*\);
\item proves that the intermediate wedge remains nonempty;
\item formulates TBRC; and
\item proves TBRC would complete MCWC and SCNGC.
\end{enumerate}
\end{theorem}

\begin{proof}
Use
\Cref{lem:newV79-radial,
thm:newV79-crude,
thm:newV79-high,
cor:newV79-two-regimes,
warn:newV79-no-promotion,
prop:newV79-TBRC}.
\end{proof}

\paragraph{Parent-version record.}
V79 was formed from
\texttt{Renewal\_Yang\_Mills\_Mass\_Gap\_Research\_Notes\_V78.tex}.
The V78 parent had \(36\,228\) logical source lines,
\(1\,156\,245\) bytes, and SHA--256
\[
 \texttt{b1f4fe1cb8938214593309a610889aa8862bde4a8b8dc6774450c2a257c18415}.
\]
The next compulsory companion audit is V80.

% =============================================================================
% END APPENDED FIFTH-SERIES V79 MATERIAL
% =============================================================================

\clearpage
\chapter{Fifth-series V80: compulsory audit and the assembled
block-edge card}
\label{ch:newV80-blocks}

\section{Compulsory companion audit}
\label{sec:newV80-audit}

\begin{audit}[V80 immutable companion snapshots]
\label{audit:newV80-companions}
The newest active companion sources were read without modification:
\begin{align*}
 &\texttt{Renewal\_SM\_Einstein\_Continuum\_Research\_Notes\_V5.tex},\\
 &\qquad 2\,266\ \hbox{logical lines},\quad
 86\,825\ \hbox{bytes},\\
 &\qquad
 \texttt{57f0163fe00429d6e7d37dedf60ef1188ceff4d3baf73f3ce36a728eee2f91da},
 \\
 &\texttt{Renewal\_NS\_Stability\_Research\_Notes\_V56.tex},\\
 &\qquad 13\,168\ \hbox{logical lines},\quad
 672\,039\ \hbox{bytes},\\
 &\qquad
 \texttt{56b16436e25ecf817ea38c9ddecc943d1f1b81c33ecc9ab34c425aef8792fb08}.
\end{align*}

SM V5 proves an exact finite-cutoff relative Ward identity after the
full and massive-base chiral blocks are placed in the same
determinant line with a common covariance law.  It transfers the
assembled identity coefficientwise to a unique heat asymptotic
scale and separately controls the entire large-proper-time sign
tail.  It does not infer either statement from termwise trace
bounds.

NS V56 proves a closed local enstrophy inequality using a bounded
overlap partition of unity, couples it to a global inequality
controlled by maximal local enstrophy, and obtains a post-window
delay of order \(M^{-1}\).  It also records that repeated delays
with geometrically growing constants can be summable and therefore
do not by themselves exclude a later spike.
\end{audit}

\begin{assessment}[Audit transfer to TBRC]
\label{ass:newV80-audit-transfer}
No SM or NS numerical estimate is imported.  Two proof disciplines
are retained:
\begin{enumerate}[label=\textup{(\roman*)},leftmargin=3.5em]
\item the block Ward or cumulant cancellation must be obtained from
      one exact assembled source before coefficient norms are taken;
      and
\item local block estimates must have bounded ownership of every
      label and edge, while the growth of constants is recorded in
      one global ledger rather than hidden in an iteration.
\end{enumerate}
Accordingly, the construction below uses one deterministic thermal
partition and reduces TBRC to one assembled block-edge identity.
\end{assessment}

\section{A deterministic thermal partition}
\label{sec:newV80-partition}

Fix
\[
 a:=\min\{b_*/4,1/4\}.
\]
Call a label heavy if \(b_i>a\), and light otherwise.

\begin{lemma}[Greedy thermal blocks]
\label{lem:newV80-blocks}
There is a deterministic partition
\(\mathcal P=\{P_1,\ldots,P_r\}\) of \([m]\) such that, with
\(w_P:=\sum_{i\in P}b_i\),
\begin{enumerate}[label=\textup{(\roman*)},leftmargin=3.5em]
\item every heavy label is a singleton block;
\item every nonsingleton block and every light singleton obeys
      \(w_P\le2a\le b_*/2\);
\item at most one light block has \(w_P<a\); and
\item
      \[
      r\le {B\over a}+1.
      \]
\end{enumerate}
On the intermediate wedge \eqref{eq:newV79-wedge}, \(r\ge2\).
\end{lemma}

\begin{proof}
Put every heavy label in its own block.  Read the light labels in
their fixed label order, accumulating them until their sum first
reaches \(a\), then close that block and restart.  Since each added
mark is at most \(a\), a completed light block has weight in
\([a,2a]\).  The final incomplete block, if present, has weight
less than \(a\).

Every heavy singleton and every completed light block has weight at
least \(a\).  Their disjoint weights sum to at most \(B\), proving
the bound on \(r\), with one added for the possible remainder.
If \(r=1\) and \(m\ge2\), that block cannot be a heavy singleton;
it is therefore a light block of weight at most \(2a\le b_*/2\),
contrary to \(B>b_*\) on the intermediate wedge.
\end{proof}

\begin{assessment}[Local ownership]
\label{ass:newV80-local}
Every nonsingleton block lies in the V78 CIRC region.  A heavy
singleton has no internal connected resummation to perform.  Thus
all within-block nonlinear words are controlled.  What is not yet
proved is that the exact \(m\)-label cumulant can be assembled from
these local sources with one cross-block tree and no extra
Leonov--Shiryaev partition sum.
\end{assessment}

\section{The exact block-edge target}
\label{sec:newV80-ABEC}

For a tree \(T\) on the block set \(\mathcal P\), write
\(d_T(P)\) for its block degree and \(n_P=|P|\).

\begin{definition}[Assembled block-edge card, ABEC]
\label{def:newV80-ABEC}
ABEC is an exact representation
\begin{equation}
 \kappa_m^{\nu_\alpha^{\rm sm}}
 (D^{R_1},\ldots,D^{R_m})
 =
 \sum_{T\in\mathfrak T(\mathcal P)}{\cal A}_{\mathcal P,T}
 \label{eq:newV80-ABEC-representation}
\end{equation}
formed from one block-level endpoint-factorizing source, together
with the bound
\begin{align}
 \|{\cal A}_{\mathcal P,T}\|_{\rm cb}
 \le{}&
 K^m
 \prod_{P\in\mathcal P}
 \left[
  \left(\prod_{i\in P}b_i\right)w_P^{n_P-2}
 \right]
 \prod_{P\in\mathcal P}w_P^{d_T(P)}.
 \label{eq:newV80-ABEC-bound}
\end{align}
For a singleton, the bracket is interpreted as
\(b_iw_P^{-1}=1\).  The representation must be derived before
taking norms; a cumulant-of-products formula followed by absolute
partition summation is not ABEC.
\end{definition}

\begin{theorem}[ABEC closes the intermediate wedge]
\label{thm:newV80-ABEC}
ABEC implies TBRC and the strong all-arity source card on
\eqref{eq:newV79-wedge}.
\end{theorem}

\begin{proof}
Sum \eqref{eq:newV80-ABEC-bound} over block trees and use the
weighted Cayley--Pr\"ufer identity:
\begin{equation}
 \sum_{T\in\mathfrak T(\mathcal P)}
 \prod_{P\in\mathcal P}w_P^{d_T(P)}
 =
 \left(\prod_{P\in\mathcal P}w_P\right)B^{r-2}.
 \label{eq:newV80-block-Cayley}
\end{equation}
Therefore
\begin{align*}
 \|\kappa_m^{\nu_\alpha^{\rm sm}}\|_{\rm cb}
 &\le
 K^m
 \left(\prod_{i=1}^mb_i\right)
 \left(\prod_{P\in\mathcal P}w_P^{n_P-1}\right)
 B^{r-2}\\
 &\le
 K^m
 \left(\prod_{i=1}^mb_i\right)
 B^{\,\sum_P(n_P-1)+r-2}\\
 &=
 K^m
 \left(\prod_{i=1}^mb_i\right)B^{m-2}.
\end{align*}
The deterministic partition creates no partition census.  Its
small nonsingleton blocks use V78 CIRC internally, while heavy
singletons contribute only their cross-block degree.  These are
exactly the rows required by TBRC.
\end{proof}

\begin{corollary}[ABEC would close SCNGC]
\label{cor:newV80-SCNGC}
ABEC, together with V78--V79, proves full MCWC and hence SCNGC.
\end{corollary}

\begin{proof}
Apply \Cref{thm:newV80-ABEC} and then
\Cref{prop:newV79-TBRC}.
\end{proof}

\section{Why the naive block cumulant is not ABEC}
\label{sec:newV80-no-go}

\begin{warning}[Cumulants of block products contain extra
indecomposable partitions]
\label{warn:newV80-products}
Let
\[
 X_P(U):=\bigotimes_{i\in P}D^{R_i}(U).
\]
The cumulant
\(\kappa_r(X_{P_1},\ldots,X_{P_r})\) is not equal to the
\(m\)-label cumulant in
\eqref{eq:newV80-ABEC-representation}.  Expanding it into cumulants
of the individual labels produces every partition of \([m]\) that
is connected relative to \(\mathcal P\), not only the one-block
partition.  Taking norms at that stage recreates a Bell family.
\end{warning}

\begin{proof}
This is the moment--cumulant formula applied first to the \(r\)
products and then to the individual factors.  A partition of
\([m]\) survives exactly when its join with the block partition
\(\mathcal P\) is the one-block partition.  There are generally
many such partitions.  For two blocks of \(k\) labels, all perfect
matchings across the blocks survive, already giving
\(k!\) terms.  Hence a termwise absolute estimate cannot satisfy
ABEC with one exponential base.
\end{proof}

\begin{assessment}[Exact remaining analytic operation]
\label{ass:newV80-position}
The thermal partition and the block-tree power ledger are complete.
The unresolved operation is not a combinatorial sum: it is an exact
source interpolation whose block forest formula cancels the
indecomposable product partitions before norms.  The SM audit shows
the same order at a relative Ward identity; the NS audit shows why
repeated local estimates without a global ownership formula are
insufficient.
\end{assessment}

\begin{programme}[V81: a shared-clock block interpolation]
\label{prog:newV80-V81}
\begin{enumerate}[leftmargin=3em]
\item Use the L\'evy--Khintchine representation of
      \(\Phi_\alpha^{\rm sm}\) to couple one clock to all blocks at
      the all-one endpoint and independent clocks to partition
      components at the other endpoints.
\item Construct this coupling as a positive Poisson marking of
      clock jumps, rather than as a signed M\"obius sum of block
      products.
\item Differentiate one block edge and prove that it forces one
      marked jump meeting both endpoint blocks, supplying
      \(w_Pw_Q\).
\item Keep all within-block jump roots inside the exact V78
      effective hyperedges before applying the Duhamel norm.
\item If positivity fails for a general forest matrix, go upstream
      to a hierarchical coalescent interpolation indexed by the
      BKAR edge order.
\end{enumerate}
\end{programme}

\begin{firewall}[Claim boundary after V80]
\label{fire:newV80-boundary}
V80 performs the compulsory SM/NS audit, constructs a deterministic
thermal partition with controlled block count, formulates ABEC,
proves that ABEC closes the intermediate wedge and hence SCNGC, and
proves that a naive cumulant-of-block-products expansion recreates a
factorial family.

V80 does not prove ABEC, TBRC, full MCWC, OCRC, SCNGC, WPSC, GRLTC,
LZTC, PLRC, WSDC, BRSC, HWSFC, UGCGC, thermal connected WPRAC,
all-arity RGSC, the raw Wilson aggregate strong card, all-order
strong MTA, WUFC, CPM, cutoff removal, reflection positivity,
Osterwalder--Schrader reconstruction, a continuum Yang--Mills
measure, or a positive Hamiltonian gap.  The full Yang--Mills mass
gap has not been proved.
\end{firewall}

\begin{theorem}[V80 result ledger]
\label{thm:newV80-results}
V80:
\begin{enumerate}[label=\textup{(V80.\arabic*)},leftmargin=6.5em]
\item performs the compulsory active-SM/active-NS audit;
\item imports only exact-assembly and bounded-ownership discipline;
\item constructs deterministic heavy and light thermal blocks;
\item proves the block count \(r\le B/a+1\);
\item formulates ABEC;
\item proves the exact block-tree power ledger;
\item proves ABEC would close TBRC, MCWC, and SCNGC; and
\item proves the naive product-cumulant route has a factorial
      obstruction.
\end{enumerate}
\end{theorem}

\begin{proof}
Use
\Cref{audit:newV80-companions,
lem:newV80-blocks,
thm:newV80-ABEC,
cor:newV80-SCNGC,
warn:newV80-products}.
\end{proof}

\paragraph{Parent-version record.}
V80 was formed from
\texttt{Renewal\_Yang\_Mills\_Mass\_Gap\_Research\_Notes\_V79.tex}.
The V79 parent had \(36\,540\) logical source lines,
\(1\,166\,363\) bytes, and SHA--256
\[
 \texttt{4bddd0e3571af30e56649e11fd5209c2f809f6070953904d836f82e47abd6823}.
\]
The next compulsory companion audit is V85.

% =============================================================================
% END APPENDED FIFTH-SERIES V80 MATERIAL
% =============================================================================

\clearpage
\chapter{Fifth-series V81: a positive percolative block source}
\label{ch:newV81-percolation}

\section{Percolation cluster intensities}
\label{sec:newV81-clusters}

Let \(\mathcal P\) be the deterministic thermal partition of
\Cref{lem:newV80-blocks}.  Put \(r=|\mathcal P|\), and give every
edge \(e=\{P,Q\}\) of the complete graph on \(\mathcal P\) an
independent Bernoulli state with opening probability
\(x_e\in[0,1]\).  Write \(\pi(\omega)\) for the partition into open
components.  For nonempty
\(\mathscr C\subseteq\mathcal P\), define
\begin{align}
 \lambda_{\mathscr C}(\mathbf x)
 &:=
 \mathbb P_{\mathbf x}
 \{\mathscr C\text{ is a component of }\pi(\omega)\},
 \label{eq:newV81-lambda}\\
 q_{PQ}(\mathbf x)
 &:=
 \mathbb P_{\mathbf x}\{P\leftrightarrow Q\},
 \qquad q_{PP}:=1.                                         \label{eq:newV81-q}
\end{align}

\begin{lemma}[Positive partition ledger]
\label{lem:newV81-partition-ledger}
The functions \(\lambda_{\mathscr C}\) and \(q_{PQ}\) are
multi-affine polynomials in the edge variables, and
\begin{align}
 &\lambda_{\mathscr C}\ge0,\qquad
 \sum_{\mathscr C\ni P}\lambda_{\mathscr C}=1,
 \label{eq:newV81-marginal-ledger}\\
 &q_{PQ}=\sum_{\mathscr C\ni P,Q}\lambda_{\mathscr C}.
 \label{eq:newV81-connectivity-ledger}
\end{align}
The matrix \(q(\mathbf x)\) is positive semidefinite.  At the
zero--one point \(\mathbf x^\rho\) of a partition \(\rho\) of
\(\mathcal P\),
\begin{equation}
 \lambda_{\mathscr C}(\mathbf x^\rho)
 =
 \mathbf1_{\{\mathscr C\in\rho\}},
 \qquad
 q(\mathbf x^\rho)=\mathbf x^\rho.
 \label{eq:newV81-partition-endpoint}
\end{equation}
\end{lemma}

\begin{proof}
Every event in the definitions is a finite union of Bernoulli edge
configurations, proving polynomiality and nonnegativity.  For a
fixed block \(P\), exactly one random component contains it; this
proves \eqref{eq:newV81-marginal-ledger}.  Summing the same
component indicator over components containing both endpoints
gives \eqref{eq:newV81-connectivity-ledger}.

For a deterministic partition \(\pi\), its connectivity matrix is
the Gram matrix of the component-indicator vectors and is positive
semidefinite.  The matrix \(q(\mathbf x)\) is the expectation of
that matrix, hence is positive semidefinite.  At
\(\mathbf x^\rho\), all within-part edges are open and all
cross-part edges are closed almost surely, proving the endpoint
identities.
\end{proof}

\section{The shared-cluster small-clock law}
\label{sec:newV81-law}

For
\(\mathscr C\subseteq\mathcal P\), let
\[
 \langle\mathscr C\rangle:=\bigcup_{P\in\mathscr C}P
\]
be its label union.  Construct a probability law
\(\mu_{\mathbf x}^{\rm perc}\) on
\(\SU(2)^{\mathcal P}\) as follows:
\begin{enumerate}[label=\textup{(\roman*)},leftmargin=3.5em]
\item its Gaussian drift has block covariance
      \(d_\alpha q(\mathbf x)\);
\item independently for every nonempty
      \(\mathscr C\subseteq\mathcal P\), use a Poisson random
      measure of heat jumps with intensity
      \(\lambda_{\mathscr C}(\mathbf x)
        \Pi_\alpha^{\rm sm}(dy)\);
\item a jump of size \(y\) and type \(\mathscr C\) applies the same
      heat increment of law \(h_y\) to every coordinate in
      \(\mathscr C\), and the identity to the other coordinates.
\end{enumerate}
This is a standard finite-type L\'evy construction; infinite jump
activity is interpreted by the L\'evy--Khintchine limit.

\begin{theorem}[Positive percolative interpolation]
\label{thm:newV81-positive}
For every \(\mathbf x\in[0,1]^{\binom r2}\),
\(\mu_{\mathbf x}^{\rm perc}\) is a probability law with central
coordinate marginals.  In
the tensor representation carried by all labels, its moment
operator is \(e^{{\cal V}^{\rm perc}(\mathbf x)}\), where
\begin{align}
 {\cal V}^{\rm perc}(\mathbf x)
 ={}&
 -d_\alpha{\cal C}_m(q(\mathbf x))
 \notag\\
 &+
 \sum_{\varnothing\ne\mathscr C\subseteq\mathcal P}
 \lambda_{\mathscr C}(\mathbf x)
 \int_0^\infty
 \left(e^{-yC_{\langle\mathscr C\rangle}^{\rm tot}}-I\right)
 \Pi_\alpha^{\rm sm}(dy).                                  \label{eq:newV81-V}
\end{align}
Moreover
\begin{equation}
 {\cal V}^{\rm perc}(\mathbf x)\preceq0,
 \qquad
 \|e^{{\cal V}^{\rm perc}(\mathbf x)}\|_{\rm cb}\le1.
 \label{eq:newV81-contraction}
\end{equation}
Every coordinate marginal is exactly
\(\nu_\alpha^{\rm sm}\).
\end{theorem}

\begin{proof}
\Cref{lem:newV81-partition-ledger} makes the Gaussian covariance
positive semidefinite and every Poisson intensity nonnegative, so
the construction is a probability law.  The
L\'evy--Khintchine formula in the tensor representation is exactly
\eqref{eq:newV81-V}.  Since \(q(\mathbf x)\succeq0\),
\({\cal C}_m(q(\mathbf x))\succeq0\).  Also
\(C_{\langle\mathscr C\rangle}^{\rm tot}\succeq0\), whence every
jump summand \(e^{-yC}-I\preceq0\).  This proves
\eqref{eq:newV81-contraction}; matrix amplification changes none
of the order statements.

For a fixed coordinate \(P\), the Gaussian variance is
\(q_{PP}=1\), and
\(\sum_{\mathscr C\ni P}\lambda_{\mathscr C}=1\).  Its projected
L\'evy triplet is therefore exactly
\((d_\alpha,\Pi_\alpha^{\rm sm})\), proving the marginal assertion.
\end{proof}

\begin{corollary}[Exact partition endpoints]
\label{cor:newV81-endpoints}
At a partition endpoint \(\mathbf x^\rho\),
\begin{equation}
 \mu_{\mathbf x^\rho}^{\rm perc}
 =
 \bigotimes_{\mathscr C\in\rho}
 \nu_{\alpha,\langle\mathscr C\rangle}^{\rm sm},
 \label{eq:newV81-law-endpoint}
\end{equation}
where one common small-clock group variable is used by all labels
in \(\langle\mathscr C\rangle\), and different
\(\mathscr C\)'s are independent.  At the all-one endpoint this
is the original common small-clock source on all \(m\) labels.
\end{corollary}

\begin{proof}
Use \eqref{eq:newV81-partition-endpoint} in
\eqref{eq:newV81-V}.  The exponent becomes
\[
 \sum_{\mathscr C\in\rho}
 \left\{
  -d_\alpha C_{\langle\mathscr C\rangle}^{\rm tot}
  +\int_0^\infty
   (e^{-yC_{\langle\mathscr C\rangle}^{\rm tot}}-I)
   \Pi_\alpha^{\rm sm}(dy)
 \right\}.
\]
The summands act on disjoint tensor factors and are precisely the
small-clock logarithms \eqref{eq:newV60-LB}.
\end{proof}

\section{The full squarefree logarithm}
\label{sec:newV81-log}

If \(i\in P\), write \(P(i)=P\).  Define
\begin{equation}
 \mathbb Z_{\mathbf x}(\mathbf z)
 :=
 \mathbb E_{\mu_{\mathbf x}^{\rm perc}}
 \bigotimes_{i=1}^m
 \{I_{R_i}+z_i(D^{R_i}(U_{P(i)})-I_{R_i})\}.
 \label{eq:newV81-source}
\end{equation}

\begin{lemma}[Uniform logarithmic source disk]
\label{lem:newV81-log-disk}
There is a numerical \(\eta>0\), independent of
\(m,r,\mathbf x\), such that
\begin{equation}
 \sum_i|z_i|b_i\le\eta
 \quad\Longrightarrow\quad
 \|\mathbb Z_{\mathbf x}(\mathbf z)-I\|_{\rm cb}\le{1\over4}.
 \label{eq:newV81-log-disk}
\end{equation}
Hence the principal Banach logarithm is uniformly defined and
bounded there.
\end{lemma}

\begin{proof}
Put
\[
 a_P:=\sum_{i\in P}|z_i|b_i,
 \qquad A:=\sum_Pa_P.
\]
By \Cref{lem:newV44-geodesic-increment},
\[
 \|\mathbb Z_{\mathbf x}(\mathbf z)-I\|_{\rm cb}
 \le
 \mathbb E
 \left[
  \exp\left(C\sum_Pa_PY_s(U_P)\right)-1
 \right].
\]
Every coordinate marginal is \(\nu_\alpha^{\rm sm}\), so
\Cref{lem:newV79-radial} and generalized H\"older give, whenever
\(CA\le\lambda_0\),
\[
 \mathbb E\exp\left(C\sum_Pa_PY_s(U_P)\right)
 \le
 \prod_{P:a_P>0}
 \left\{\mathbb Ee^{CAY_s(U_P)}\right\}^{a_P/A}
 \le C_0.
\]
Using \(e^u-1\le ue^u\), the same H\"older argument with a slightly
larger exponential parameter bounds the preceding difference by
\(C_1A\).  Choose \(\eta\) so that \(C_1\eta\le1/4\).
The logarithm series then converges normally.
\end{proof}

Define the completely assembled squarefree coefficient
\begin{equation}
 {\cal G}(\mathbf x)
 :=
 [z_1\cdots z_m]\log\mathbb Z_{\mathbf x}(\mathbf z).
 \label{eq:newV81-G}
\end{equation}

\begin{lemma}[Disconnected faces vanish after the logarithm]
\label{lem:newV81-face-zero}
If the support graph of \(\mathbf x\) is disconnected on the block
set \(\mathcal P\), then
\begin{equation}
 {\cal G}(\mathbf x)=0.                                    \label{eq:newV81-face-zero}
\end{equation}
At the all-one point,
\begin{equation}
 {\cal G}(\mathbf1)
 =
 \kappa_m^{\nu_\alpha^{\rm sm}}
 (D^{R_1},\ldots,D^{R_m}).                                 \label{eq:newV81-all-one}
\end{equation}
\end{lemma}

\begin{proof}
If the support graph has components \(\rho\), no percolation
cluster crosses \(\rho\), and the covariance \(q(\mathbf x)\) is
block diagonal across \(\rho\).  Thus the law and source factor:
\[
 \mathbb Z_{\mathbf x}
 =
 \bigotimes_{\mathscr C\in\rho}
 \mathbb Z_{\mathbf x,\mathscr C}.
\]
The factors act on disjoint tensor supports, commute, and lie in
the logarithmic disk of
\Cref{lem:newV81-log-disk}.  Therefore
\[
 \log\mathbb Z_{\mathbf x}
 =
 \sum_{\mathscr C\in\rho}
 \log\mathbb Z_{\mathbf x,\mathscr C}.
\]
No summand contains every one of the \(m\) squarefree variables,
proving \eqref{eq:newV81-face-zero}.  At \(\mathbf1\), use
\Cref{cor:newV81-endpoints} and the general squarefree logarithm
identity \Cref{lem:newV44-log-cumulant}.
\end{proof}

\section{Exact block-tree representation}
\label{sec:newV81-tree}

For a forest \(F\) on \(\mathcal P\) and
\(\mathbf u\in[0,1]^{E(F)}\), let
\(\mathbf x^F(\mathbf u)\) be the BKAR matrix: its \(PQ\) entry is
the minimum edge parameter on the \(F\)-path from \(P\) to \(Q\),
or zero when no path exists.

\begin{theorem}[ABEC representation row]
\label{thm:newV81-ABEC-representation}
The exact cumulant admits
\begin{equation}
 \kappa_m^{\nu_\alpha^{\rm sm}}
 (D^{R_1},\ldots,D^{R_m})
 =
 \sum_{T\in\mathfrak T(\mathcal P)}
 \int_{[0,1]^{r-1}}
 \partial_{E(T)}{\cal G}(\mathbf x^T(\mathbf u))
 \,d\mathbf u.
 \label{eq:newV81-ABEC-representation}
\end{equation}
In particular, the representation row
\eqref{eq:newV80-ABEC-representation} of ABEC holds with
\[
 {\cal A}_{\mathcal P,T}
 :=
 \int_{[0,1]^{r-1}}
 \partial_{E(T)}{\cal G}(\mathbf x^T(\mathbf u))
 \,d\mathbf u.
\]
\end{theorem}

\begin{proof}
The cluster weights and connectivity probabilities are polynomials
in \(\mathbf x\); every finite-dimensional moment and the logarithm
on \eqref{eq:newV81-log-disk} are therefore smooth.  Apply the
BKAR forest formula to \({\cal G}(\mathbf1)\).

If \(F\) is disconnected, then throughout the face on which all
edges between distinct \(F\)-components vanish,
\({\cal G}=0\) by \Cref{lem:newV81-face-zero}.  The derivative
\(\partial_{E(F)}\) uses only edges within those components, so its
forest term also vanishes.  For \(r\ge2\), the only connected
forests are trees.  Use \eqref{eq:newV81-all-one}.
\end{proof}

\section{One pivotal block edge}
\label{sec:newV81-pivotal}

For disjoint nonempty block families
\(\mathscr A,\mathscr B\), define the exact merger
\begin{equation}
 {\cal M}(\mathscr A,\mathscr B)
 :=
 L_{\langle\mathscr A\cup\mathscr B\rangle}^{\rm sm}
 -L_{\langle\mathscr A\rangle}^{\rm sm}
 -L_{\langle\mathscr B\rangle}^{\rm sm}.
 \label{eq:newV81-merger}
\end{equation}

\begin{lemma}[Russo merger formula]
\label{lem:newV81-Russo}
Let \(e=\{P,Q\}\).  Sample all percolation edges other than \(e\)
and keep \(e\) closed.  If \(P,Q\) lie in distinct components
\(\mathscr A,\mathscr B\), then opening \(e\) changes the moment
exponent by \({\cal M}(\mathscr A,\mathscr B)\); if they are
already connected, it changes nothing.  Consequently
\begin{equation}
 \partial_e{\cal V}^{\rm perc}(\mathbf x)
 =
 \mathbb E_{\mathbf x,\ne e}
 \left[
  \mathbf1_{\{P\nleftrightarrow Q\}}
  {\cal M}(\mathscr A,\mathscr B)
 \right].
 \label{eq:newV81-Russo}
\end{equation}
Moreover
\begin{align}
 {\cal M}(\mathscr A,\mathscr B)
 ={}&
 -2d_\alpha
 \sum_{\substack{i\in\langle\mathscr A\rangle\\
                 j\in\langle\mathscr B\rangle}}
 \Omega_{ij}
 \notag\\
 &+
 \sum_{\substack{\varnothing\ne S\subseteq
                  \langle\mathscr A\cup\mathscr B\rangle\\
                  S\cap\langle\mathscr A\rangle\ne\varnothing\\
                  S\cap\langle\mathscr B\rangle\ne\varnothing}}
 Q_S.                                                       \label{eq:newV81-merger-roots}
\end{align}
\end{lemma}

\begin{proof}
Conditioning independent Bernoulli percolation on \(e=1\) or
\(e=0\) is the elementary Russo derivative for a multi-affine
expectation.  When the endpoint components are distinct, opening
\(e\) replaces the two independent cluster exponents by their
union exponent.  This gives
\eqref{eq:newV81-Russo}.

Expand each \(L^{\rm sm}\) by
\eqref{eq:newV60-LB-subsets}.  Singleton and same-side jump roots
cancel.  The total Casimir identity leaves exactly the displayed
cross term.
\end{proof}

\begin{assessment}[What V81 removes]
\label{ass:newV81-position}
The interpolation is a genuine probability law for every edge
matrix, and its complete moment exponent is dissipative.  Thus the
intermediate wedge no longer suffers the unstable
\(e^{cB^2}\) estimate of the gated V60 interpolation.  Taking the
squarefree logarithm before BKAR also removes the factorial family
of \Cref{warn:newV80-products}.  Every first block derivative is an
exact merger meeting both sides of its pivotal edge.

The remaining row is quantitative: prove
\eqref{eq:newV80-ABEC-bound} for the mixed derivatives in
\eqref{eq:newV81-ABEC-representation}.  Higher percolation
derivatives and Fr\'echet derivatives of the assembled logarithm
must be resummed before norms.
\end{assessment}

\begin{programme}[V82: pivotal forests and the ABEC bound]
\label{prog:newV81-V82}
\begin{enumerate}[leftmargin=3em]
\item Expand mixed percolation derivatives by conditioning all
      forest edges simultaneously; record the partition formed
      when those edges are closed.
\item Prove that survival forces every forest edge to be pivotal
      in the induced quotient, so the merger roots form a connected
      block-incidence structure.
\item Apply V78 CIRC inside light thermal blocks before estimating
      the merger operators.
\item Use the dissipative logarithmic resolvent on
      \eqref{eq:newV81-log-disk} to contract all ordered words
      without a factorial.
\item Match the two endpoint-block marks per pivotal edge with
      \(\prod_Pw_P^{d_T(P)}\).
\end{enumerate}
\end{programme}

\begin{firewall}[Claim boundary after V81]
\label{fire:newV81-boundary}
V81 constructs a positive percolative shared-clock law with exact
small-clock marginals, proves contraction for every interpolation
matrix, proves exact factorization at every partition endpoint,
constructs a uniform logarithmic source disk, proves disconnected
squarefree logarithms vanish, derives the exact block-tree ABEC
representation, and proves the one-edge Russo merger formula.

V81 does not prove the ABEC derivative bound, ABEC, TBRC, full
MCWC, OCRC, SCNGC, WPSC, GRLTC, LZTC, PLRC, WSDC, BRSC, HWSFC,
UGCGC, thermal connected WPRAC, all-arity RGSC, the raw Wilson
aggregate strong card, all-order strong MTA, WUFC, CPM, cutoff
removal, reflection positivity, Osterwalder--Schrader
reconstruction, a continuum Yang--Mills measure, or a positive
Hamiltonian gap.  The full Yang--Mills mass gap has not been proved.
\end{firewall}

\begin{theorem}[V81 result ledger]
\label{thm:newV81-results}
V81:
\begin{enumerate}[label=\textup{(V81.\arabic*)},leftmargin=6.5em]
\item defines the nonnegative percolation cluster intensities;
\item proves the induced connectivity matrix is positive
      semidefinite;
\item constructs the percolative shared-clock L\'evy law;
\item proves exact marginals, endpoint factorization, and
      contraction;
\item proves a uniform block-source logarithmic disk;
\item proves disconnected full squarefree logarithms vanish;
\item proves the exact ABEC block-tree representation; and
\item proves the pivotal-edge Russo merger formula.
\end{enumerate}
\end{theorem}

\begin{proof}
Use
\Cref{lem:newV81-partition-ledger,
thm:newV81-positive,
cor:newV81-endpoints,
lem:newV81-log-disk,
lem:newV81-face-zero,
thm:newV81-ABEC-representation,
lem:newV81-Russo}.
\end{proof}

\paragraph{Parent-version record.}
V81 was formed from
\texttt{Renewal\_Yang\_Mills\_Mass\_Gap\_Research\_Notes\_V80.tex}.
The V80 parent had \(36\,851\) logical source lines,
\(1\,177\,460\) bytes, and SHA--256
\[
 \texttt{7b027912315e23e817f463a08735b5ba9b5d2d5e8f71efc90f9facae42232d30}.
\]
The next compulsory companion audit is V85.

% =============================================================================
% END APPENDED FIFTH-SERIES V81 MATERIAL
% =============================================================================

\clearpage
\chapter{Fifth-series V82: simultaneous Russo quotients and the
shortcut-cycle correction}
\label{ch:newV82-Russo}

\section{The exact percolation gates seen by label roots}
\label{sec:newV82-gates}

For a nonempty label set \(S\subseteq[m]\), define its block
support
\[
 \beta(S):=\{P\in\mathcal P:S\cap P\ne\varnothing\}.
\]
For a nonempty block family
\(\mathscr D\subseteq\mathcal P\), put
\begin{equation}
 g_{\mathscr D}(\mathbf x)
 :=
 \mathbb P_{\mathbf x}
 \{\mathscr D\text{ is contained in one open component}\}
 =
 \sum_{\mathscr C\supseteq\mathscr D}
 \lambda_{\mathscr C}(\mathbf x).
 \label{eq:newV82-gate}
\end{equation}
Thus \(g_{\{P\}}=1\) and
\(g_{\{P,Q\}}=q_{PQ}\).

\begin{lemma}[Exact gated root expansion]
\label{lem:newV82-gated-expansion}
The percolative moment exponent of V81 has the exact expansion
\begin{align}
 {\cal V}^{\rm perc}(\mathbf x)
 ={}&
 -d_\alpha
 \left\{
  \sum_iC_2(R_i)I+
  2\sum_{i<j}
  g_{\{P(i),P(j)\}}(\mathbf x)\Omega_{ij}
 \right\}
 \notag\\
 &+
 \sum_{\varnothing\ne S\subseteq[m]}
 g_{\beta(S)}(\mathbf x)Q_S.
 \label{eq:newV82-gated-expansion}
\end{align}
All gates are nonnegative multi-affine polynomials and are at most
one on the real cube.
\end{lemma}

\begin{proof}
For the Gaussian part, two labels share covariance exactly when
their block coordinates lie in the same percolation component;
the expectation of that indicator is \(q_{P(i)P(j)}\).

For a jump cluster \(\mathscr C\), apply the exact subset
decomposition \eqref{eq:newV60-subset} to the labels in
\(\langle\mathscr C\rangle\).  A root \(Q_S\) occurs precisely for
clusters containing every block in \(\beta(S)\).  Its summed
intensity is
\(\sum_{\mathscr C\supseteq\beta(S)}\lambda_{\mathscr C}
=g_{\beta(S)}\).  The final assertion follows directly from the
Bernoulli probability definition.
\end{proof}

\begin{lemma}[Finite-difference derivative of a gate]
\label{lem:newV82-gate-difference}
Let \(E\) be a nonempty set of \(k\) block edges.  Then
\begin{equation}
 \partial_Eg_{\mathscr D}(\mathbf x)
 =
 \sum_{F\subseteq E}
 (-1)^{k-|F|}
 g_{\mathscr D}(\mathbf x_{E=\mathbf1_F}),
 \label{eq:newV82-gate-difference}
\end{equation}
where the variables outside \(E\) are unchanged.  In particular,
\begin{equation}
 |\partial_Eg_{\mathscr D}(\mathbf x)|
 \le2^{k-1}.                                                \label{eq:newV82-gate-derivative}
\end{equation}
\end{lemma}

\begin{proof}
A Bernoulli expectation is affine in each opening probability.
Iterate the one-variable identity
\(\partial_ex=f|_{x_e=1}-f|_{x_e=0}\).
There are \(2^{k-1}\) positive and \(2^{k-1}\) negative terms in
\eqref{eq:newV82-gate-difference}, all in \([0,1]\), which proves
\eqref{eq:newV82-gate-derivative}.
\end{proof}

\section{Condition all differentiated edges at once}
\label{sec:newV82-simultaneous}

For a percolation partition \(\rho\) of \(\mathcal P\), define
\begin{equation}
 {\cal H}(\rho)
 :=
 \sum_{\mathscr C\in\rho}
 L_{\langle\mathscr C\rangle}^{\rm sm}.
 \label{eq:newV82-component-H}
\end{equation}
By the proof of \Cref{thm:newV81-positive},
\begin{equation}
 {\cal V}^{\rm perc}(\mathbf x)
 =
 \mathbb E_{\mathbf x}{\cal H}(\pi(\omega)).
 \label{eq:newV82-V-expectation}
\end{equation}

Fix an edge set \(E\).  Condition all edges outside \(E\), keep the
edges of \(E\) closed, and let \(\rho_0\) be the resulting component
partition.  Contract every part of \(\rho_0\).  The images of the
edges in \(E\) form a multigraph
\(\Gamma_E(\rho_0)\); an edge whose endpoints lie in one part
becomes a loop.

\begin{theorem}[Simultaneous Russo quotient]
\label{thm:newV82-simultaneous-Russo}
For every nonempty \(E\),
\begin{equation}
 \partial_E{\cal V}^{\rm perc}(\mathbf x)
 =
 \mathbb E_{\mathbf x,\notin E}
 \left[
  \Delta_E{\cal H}(\rho_0)
 \right],
 \label{eq:newV82-simultaneous-Russo}
\end{equation}
where
\begin{equation}
 \Delta_E{\cal H}(\rho_0)
 :=
 \sum_{F\subseteq E}(-1)^{|E|-|F|}
 {\cal H}(\rho_0\vee F).                                   \label{eq:newV82-component-difference}
\end{equation}
Here \(\rho_0\vee F\) is obtained by joining the contracted
components along the images of the edges of \(F\).

The difference in \eqref{eq:newV82-component-difference} vanishes
unless \(\Gamma_E(\rho_0)\) is loop-free and its edge-bearing part
is connected.  Consequently every surviving quotient has at most
\(|E|+1\) vertices, and percolation differentiation creates no
factor larger than \(2^{|E|-1}\) before the operator roots are
estimated.
\end{theorem}

\begin{proof}
Apply the mixed Bernoulli finite-difference identity to
\eqref{eq:newV82-V-expectation}, conditioning the undifferentiated
edges.  This gives
\eqref{eq:newV82-simultaneous-Russo} and
\eqref{eq:newV82-component-difference}.

If an edge of \(E\) becomes a loop, opening or closing it leaves
\(\rho_0\vee F\) unchanged; its one-variable difference is zero.
Suppose instead that the edge-bearing multigraph is the disjoint
union of two nonempty vertex-disjoint edge families
\(E_1,E_2\).  The component functional
\({\cal H}\) changes additively on those two disjoint vertex sets.
After one difference in every edge of \(E_1\), the result is
independent of the variables in \(E_2\), so the remaining mixed
difference is zero.  This proves the connectedness condition.
A connected loop-free multigraph with \(|E|\) edges has at most
\(|E|+1\) vertices.  The last assertion is the same positive-versus-
negative count as in \eqref{eq:newV82-gate-derivative}.
\end{proof}

\section{The acyclic quotient is an exact spanning gate}
\label{sec:newV82-tree-quotient}

For a block support \(\mathscr D\), let
\(\overline{\mathscr D}\) be the set of contracted
\(\rho_0\)-vertices which meet \(\mathscr D\).

\begin{lemma}[Tree-quotient gate criterion]
\label{lem:newV82-tree-criterion}
Suppose \(\Gamma_E(\rho_0)\) is an ordinary tree.  Let
\(\tau_{\Gamma}(\overline{\mathscr D})\) be the minimal subtree
spanning \(\overline{\mathscr D}\).  Then, conditionally on the
outside edges,
\begin{equation}
 \Delta_E
 \mathbf1_{\{\mathscr D\text{ lies in one component}\}}
 =
 \begin{cases}
 1,&\tau_{\Gamma}(\overline{\mathscr D})=\Gamma_E(\rho_0),\\
 0,&\text{otherwise}.
 \end{cases}                                                \label{eq:newV82-tree-criterion}
\end{equation}
\end{lemma}

\begin{proof}
On a tree, the vertices in
\(\overline{\mathscr D}\) are connected after the edge set \(F\)
is opened exactly when every edge of their minimal spanning subtree
belongs to \(F\).  Thus the indicator is the squarefree monomial
\[
 \prod_{e\in\tau_{\Gamma}(\overline{\mathscr D})}
 \mathbf1_{\{e\in F\}}.
\]
Its full mixed difference in all variables of \(E\) is zero unless
that monomial contains every edge, and is one when it does.
\end{proof}

\begin{corollary}[Exact root content of an acyclic quotient]
\label{cor:newV82-tree-roots}
Under the hypothesis of
\Cref{lem:newV82-tree-criterion}, the conditioned value of
\(\Delta_E{\cal H}(\rho_0)\) is
\begin{align}
 &-2d_\alpha
 \sum_{\substack{i<j\\
 \tau_\Gamma(\overline{\{P(i),P(j)\}})=\Gamma}}
 \Omega_{ij}
 \notag\\
 &\quad+
 \sum_{\substack{\varnothing\ne S\subseteq[m]\\
 \tau_\Gamma(\overline{\beta(S)})=\Gamma}}
 Q_S.                                                       \label{eq:newV82-tree-roots}
\end{align}
Thus every surviving drift pair or jump root spans the entire
acyclic quotient.  No absolute inclusion--exclusion factor remains
in this sector.
\end{corollary}

\begin{proof}
Apply \Cref{lem:newV82-tree-criterion} termwise to the exact gated
expansion \eqref{eq:newV82-gated-expansion}.
\end{proof}

\section{Why edgewise pivotality is false}
\label{sec:newV82-shortcut}

\begin{warning}[A non-tree shortcut produces a cyclic quotient]
\label{warn:newV82-shortcut}
Take three blocks and differentiate in the path edges
\[
 E=\{\{1,2\},\{2,3\}\}.
\]
Condition the undifferentiated edge \(\{1,3\}\) to be open.
With both differentiated edges closed, the contracted vertices are
\(\{1,3\}\) and \(\{2\}\), joined by two parallel edge images.
Then
\begin{equation}
 \Delta_E{\cal H}
 =
 -{\cal M}(\{1,3\},\{2\})\ne0
 \quad\hbox{in general}.                                   \label{eq:newV82-shortcut}
\end{equation}
Hence it is false that every differentiated forest edge must be
individually pivotal.
\end{warning}

\begin{proof}
With neither path edge open,
\({\cal H}=L_{\{1,3\}}^{\rm sm}+L_{\{2\}}^{\rm sm}\).
Opening either one path edge, or both, produces the one-component
value \(L_{\{1,2,3\}}^{\rm sm}\).  Therefore
\[
 \Delta_E{\cal H}
 =
 L_{\{1,2,3\}}^{\rm sm}
 -L_{\{1,2,3\}}^{\rm sm}
 -L_{\{1,2,3\}}^{\rm sm}
 +L_{\{1,3\}}^{\rm sm}+L_{\{2\}}^{\rm sm},
\]
which is \eqref{eq:newV82-shortcut}.
\end{proof}

\begin{assessment}[Steiner-block ownership]
\label{ass:newV82-Steiner}
In the shortcut example, a differentiated connectivity gate for a
root supported on blocks \(1\) and \(2\) can have value \(-1\);
block \(3\) is only a Steiner component of the open shortcut.
Its thermal mark must therefore come from the squarefree background
coefficient inside \(\log\mathbb Z_{\mathbf x}\), not from the
single differentiated root.  Any proof which assigns two endpoint
marks independently to every Russo derivative is invalid.
\end{assessment}

\begin{definition}[Shortcut-cycle ownership card, SCOC]
\label{def:newV82-SCOC}
SCOC is the following bound inside
\(\partial_{E(T)}{\cal G}(\mathbf x^T(\mathbf u))\):
after conditioning undifferentiated percolation edges and forming
the connected quotient \(\Gamma_{E(T)}(\rho_0)\),
\begin{enumerate}[label=\textup{(\roman*)},leftmargin=3.5em]
\item the acyclic quotient terms use
      \eqref{eq:newV82-tree-roots};
\item every quotient cycle is combined with the squarefree
      background roots which contain its Steiner blocks before
      norms;
\item the resulting connected block-incidence core pays
      \(\prod_Pw_P^{d_T(P)}\); and
\item all quotient excesses and Duhamel/logarithmic orders cost at
      most \(K^m\).
\end{enumerate}
\end{definition}

\begin{proposition}[SCOC is the remaining percolative row of ABEC]
\label{prop:newV82-SCOC}
If SCOC holds after the within-block CIRC resummations of V78, then
the bound \eqref{eq:newV80-ABEC-bound} holds.
\end{proposition}

\begin{proof}
\Cref{thm:newV82-simultaneous-Russo} exhausts every conditioned
mixed derivative of the moment exponent and removes disconnected
quotients.  \Cref{cor:newV82-tree-roots} treats quotient excess
zero.  SCOC supplies the cyclic quotient terms with the same
prescribed tree monomial.  The uniform logarithmic disk
\Cref{lem:newV81-log-disk} and the dissipative moment contraction
\Cref{thm:newV81-positive} control the assembled outer resolvents;
the remaining local label roots are exactly the V78 CIRC
structures inside each light block.  Multiplying the stated vertex
and tree factors gives \eqref{eq:newV80-ABEC-bound}.
\end{proof}

\begin{assessment}[Position after the correction]
\label{ass:newV82-position}
V82 replaces the false edgewise-pivotal programme by an exact
connected-quotient theorem.  Acyclic quotients have coefficient
one and an explicit spanning-root criterion.  The sole new
operation is ownership of Steiner blocks on quotient cycles.
This is narrower than the original mixed-derivative problem: all
percolation finite differences, disconnected quotients, and
acyclic quotient gates have been evaluated exactly.
\end{assessment}

\begin{programme}[V83: cycle cutting with background ownership]
\label{prog:newV82-V83}
\begin{enumerate}[leftmargin=3em]
\item For a connected quotient multigraph \(\Gamma\), choose its
      canonical spanning tree and mark the quotient excess.
\item Expand only those background gates needed to attach every
      Steiner block created by a surplus quotient edge.
\item Use the V78 canonical skeleton decomposition on the joint
      graph of differentiated merger roots and background roots.
\item Prove that each surplus quotient edge either contributes two
      block marks or removes one independent background component;
      do not count the two choices separately.
\item Sum the quotient excess before applying the logarithmic
      resolvent norm.
\end{enumerate}
\end{programme}

\begin{firewall}[Claim boundary after V82]
\label{fire:newV82-boundary}
V82 derives the exact percolation gates for every label root,
proves the simultaneous Russo quotient formula, removes all
disconnected and looped quotients, proves the exact spanning-root
criterion for acyclic quotients, and gives a concrete shortcut
counterexample to edgewise pivotality.  It isolates SCOC as the
remaining percolative row of ABEC.

V82 does not prove SCOC, the ABEC derivative bound, ABEC, TBRC,
full MCWC, OCRC, SCNGC, WPSC, GRLTC, LZTC, PLRC, WSDC, BRSC,
HWSFC, UGCGC, thermal connected WPRAC, all-arity RGSC, the raw
Wilson aggregate strong card, all-order strong MTA, WUFC, CPM,
cutoff removal, reflection positivity, Osterwalder--Schrader
reconstruction, a continuum Yang--Mills measure, or a positive
Hamiltonian gap.  The full Yang--Mills mass gap has not been proved.
\end{firewall}

\begin{theorem}[V82 result ledger]
\label{thm:newV82-results}
V82:
\begin{enumerate}[label=\textup{(V82.\arabic*)},leftmargin=6.5em]
\item expands the percolative exponent with exact connectivity
      gates;
\item proves exponential, factorial-free mixed gate derivatives;
\item proves the simultaneous Russo quotient identity;
\item proves only connected loop-free quotient multigraphs survive;
\item proves the exact tree-quotient spanning criterion;
\item computes the acyclic quotient root content;
\item disproves edgewise pivotality by an explicit shortcut; and
\item isolates SCOC and proves its sufficiency for the ABEC bound.
\end{enumerate}
\end{theorem}

\begin{proof}
Use
\Cref{lem:newV82-gated-expansion,
lem:newV82-gate-difference,
thm:newV82-simultaneous-Russo,
lem:newV82-tree-criterion,
cor:newV82-tree-roots,
warn:newV82-shortcut,
prop:newV82-SCOC}.
\end{proof}

\paragraph{Parent-version record.}
V82 was formed from
\texttt{Renewal\_Yang\_Mills\_Mass\_Gap\_Research\_Notes\_V81.tex}.
The V81 parent had \(37\,338\) logical source lines,
\(1\,193\,983\) bytes, and SHA--256
\[
 \texttt{3bf7c95adc8dcd03bfef48aaff1cd5f277ed9724ab9e90e4442eb7047366ca2c}.
\]
The next compulsory companion audit is V85.

% =============================================================================
% END APPENDED FIFTH-SERIES V82 MATERIAL
% =============================================================================

\clearpage
\chapter{Fifth-series V83: aggregate ABEC and the Kruskal
shortcut identity}
\label{ch:newV83-Kruskal}

\section{The termwise block-tree target is stronger than necessary}
\label{sec:newV83-aggregate}

For the deterministic thermal partition \(\mathcal P\), set
\begin{equation}
 v_P:=
 \left(\prod_{i\in P}b_i\right)w_P^{n_P-2}.
 \label{eq:newV83-vertex-weight}
\end{equation}
For a heavy singleton this equals one, as in
\Cref{def:newV80-ABEC}.  Retain the exact tree terms
\({\cal A}_{\mathcal P,T}\) of
\Cref{thm:newV81-ABEC-representation}.

\begin{definition}[Aggregate assembled block-edge card, AABEC]
\label{def:newV83-AABEC}
AABEC is the estimate
\begin{equation}
 \sum_{T\in\mathfrak T(\mathcal P)}
 \|{\cal A}_{\mathcal P,T}\|_{\rm cb}
 \le
 K^m
 \left(\prod_{P\in\mathcal P}v_P\right)
 \left(\prod_{P\in\mathcal P}w_P\right)B^{r-2}.
 \label{eq:newV83-AABEC}
\end{equation}
Unlike \eqref{eq:newV80-ABEC-bound}, AABEC does not require the
input BKAR tree \(T\) to own the same tree monomial after
percolation shortcuts have been exposed.
\end{definition}

\begin{theorem}[AABEC is sufficient]
\label{thm:newV83-AABEC}
AABEC proves the strong all-arity small-clock source card on the
intermediate wedge and therefore implies TBRC.
\end{theorem}

\begin{proof}
Use the exact representation
\eqref{eq:newV81-ABEC-representation} and
\eqref{eq:newV83-AABEC}.  Since \(w_P\le B\),
\begin{align*}
 \left(\prod_Pv_P\right)
 \left(\prod_Pw_P\right)B^{r-2}
 &=
 \left(\prod_{i=1}^mb_i\right)
 \left(\prod_Pw_P^{n_P-1}\right)B^{r-2}\\
 &\le
 \left(\prod_{i=1}^mb_i\right)
 B^{\,m-r+r-2}.
\end{align*}
This is the strong target.  The deterministic local blocks and the
same exact source assembly are the rows of TBRC.
\end{proof}

\begin{proposition}[Relation to the original ABEC]
\label{prop:newV83-old-ABEC}
The termwise ABEC bound \eqref{eq:newV80-ABEC-bound} implies AABEC,
but AABEC alone is sufficient for
\Cref{cor:newV80-SCNGC}.
\end{proposition}

\begin{proof}
Sum \eqref{eq:newV80-ABEC-bound} over \(T\) and apply the block
Cayley identity \eqref{eq:newV80-block-Cayley}.  Sufficiency is
\Cref{thm:newV83-AABEC,prop:newV79-TBRC}.
\end{proof}

\begin{correction}[The live target is aggregate]
\label{corr:newV83-target}
The termwise tree monomial in V80 is retained as a sufficient
strong form, but it is retired as the required target.  In
particular, the Steiner-block observation of
\Cref{ass:newV82-Steiner} is not a defect if the corresponding
operator ancestry is reassigned to a different output spanning
tree before the sum over input BKAR trees is estimated.
\end{correction}

\section{Kruskal partition of unity}
\label{sec:newV83-Kruskal}

Let \(G\) be a connected simple graph on \(\mathcal P\).  If
\(T\subseteq G\) is a spanning tree and
\(\mathbf u\in[0,1]^{E(T)}\), set
\[
 x_e^T(\mathbf u)
 :=
 \min_{f\in{\rm path}_T(e)}u_f
 \qquad(e\in E(G)\setminus E(T)).
\]

\begin{theorem}[Maximum-spanning-tree identity]
\label{thm:newV83-Kruskal}
For every finite connected simple graph \(G\),
\begin{equation}
 \sum_{\substack{T\subseteq G\\T\ {\rm spanning\ tree}}}
 \int_{[0,1]^{r-1}}
 \prod_{e\in E(G)\setminus E(T)}
 x_e^T(\mathbf u)\,d\mathbf u
 =
 1.
 \label{eq:newV83-Kruskal}
\end{equation}
\end{theorem}

\begin{proof}
Give every edge of \(G\) an independent uniform random weight
\(U_e\in[0,1]\).  With probability one all weights are distinct,
so Kruskal's decreasing-weight algorithm produces a unique maximum
spanning tree.

Fix \(T\subseteq G\) and condition on \(U_f=u_f\) for
\(f\in E(T)\).  The cycle criterion says that \(T\) is the maximum
spanning tree exactly when
\[
 U_e\le\min_{f\in{\rm path}_T(e)}u_f
 \quad\text{for every }e\notin T.
\]
The conditional probability of these independent inequalities is
the product in \eqref{eq:newV83-Kruskal}.  Integrating the tree
weights gives the probability that the maximum spanning tree is
\(T\).  Sum the mutually exclusive exhaustive events.
\end{proof}

\begin{corollary}[A fixed shortcut graph costs no tree census]
\label{cor:newV83-shortcut}
Suppose a positive-majorant ancestry has a fixed connected
percolation shortcut graph \(G\), and its term in the input BKAR
tree \(T\subseteq G\) contains the interpolation factor
\[
 \prod_{e\in E(G)\setminus E(T)}x_e^T(\mathbf u).
\]
Then summing all compatible input trees and their BKAR parameters
costs exactly one.  If a Russo finite difference has first been
split into signs, its at most \(2^{r-1}\) sign choices enlarge only
the exponential base.
\end{corollary}

\begin{proof}
The first assertion is
\Cref{thm:newV83-Kruskal}.  Since \(r\le m\), the second factor is
at most \(2^m\).
\end{proof}

\begin{assessment}[The triangle shortcut is reallocated]
\label{ass:newV83-triangle}
In \Cref{warn:newV82-shortcut}, the open edge \(\{1,3\}\) and the
two differentiated path edges form the triangle \(G=K_3\).
The three input spanning trees each receive its Kruskal region;
their integrals sum to one.  Thus the operator roots may pay the
output tree \(\{\{1,2\},\{1,3\}\}\), even when the input derivative
tree was \(\{\{1,2\},\{2,3\}\}\).  No comparison of \(w_1\) and
\(w_2\) is needed.
\end{assessment}

\section{Normalized percolation configurations must be retained}
\label{sec:newV83-normalization}

\begin{warning}[Do not sum shortcut graphs after dropping closed
edges]
\label{warn:newV83-normalization}
The identity \eqref{eq:newV83-Kruskal} treats one fixed connected
open graph.  Summing all open graphs after replacing every closed
edge factor \(1-x_e\) by one would create as many as
\(2^{\binom r2}\) terms and is invalid.  The Bernoulli
configuration probability must be retained until the open-graph
Kruskal allocation has been made; only then may the normalized
sum over configurations be used.
\end{warning}

\begin{proof}
There are \(2^{\binom r2}\) simple graphs on \(r\) labelled
vertices, whose \(r\)-th root is unbounded.  In contrast,
\[
 \sum_{\omega}
 \prod_{e:\omega_e=1}x_e
 \prod_{e:\omega_e=0}(1-x_e)=1.
\]
Dropping the closed factors before the sum destroys this
normalization.
\end{proof}

\begin{definition}[Conditioned percolation ancestry card, CPAC]
\label{def:newV83-CPAC}
CPAC is the following operation in the mixed derivative of
\({\cal G}\):
\begin{enumerate}[label=\textup{(\roman*)},leftmargin=3.5em]
\item retain the complete Bernoulli probability of every
      undifferentiated-edge configuration;
\item combine its Russo quotient roots with the squarefree
      background logarithm before norms;
\item apply the V78 CIRC resummation to every light thermal block
      and the trivial vertex card to every heavy singleton;
\item expand the resulting connected block ancestry into an output
      graph \(G\) with the vertex weight \(v_P\) and one incident
      \(w_P\) per output-tree edge; and
\item use \Cref{thm:newV83-Kruskal} before summing the normalized
      percolation configurations.
\end{enumerate}
The required conclusion is the aggregate estimate
\eqref{eq:newV83-AABEC}.
\end{definition}

\begin{proposition}[CPAC is the remaining row of AABEC]
\label{prop:newV83-CPAC}
CPAC implies AABEC and hence full MCWC and SCNGC.
\end{proposition}

\begin{proof}
The exact tree representation is
\Cref{thm:newV81-ABEC-representation}.  Connected quotient support
is \Cref{thm:newV82-simultaneous-Russo}.  For each fixed connected
output graph, \Cref{thm:newV83-Kruskal} sums every compatible input
tree at unit cost.  The percolation probabilities then sum to one,
and the CPAC vertex and output-tree weights give
\eqref{eq:newV83-AABEC}.  Apply
\Cref{thm:newV83-AABEC,prop:newV79-TBRC}.
\end{proof}

\begin{assessment}[Position after aggregate reallocation]
\label{ass:newV83-position}
SCOC is absorbed into the more precise CPAC.  Quotient cycles and
Steiner shortcuts themselves are no longer a counting bottleneck:
the BKAR tree sum pays them by the exact Kruskal partition of unity.
The remaining analytic operation is the dressed squarefree
block-vertex estimate in CPAC(iii)--(iv), with the full normalized
percolation probability retained.
\end{assessment}

\begin{programme}[V84: dressed block vertices]
\label{prog:newV83-V84}
\begin{enumerate}[leftmargin=3em]
\item Write the squarefree coefficient inside one light block with
      \(d\) prescribed external merger incidences.
\item Add \(d\) formal port labels and extract their multilinear
      coefficient from the V78 CIRC polynomial before setting port
      marks to the adjacent block weights.
\item Prove the local factor
      \(K^{n_P}v_Pw_P^{d}\) uniformly in \(d\), including the
      singleton convention.
\item Insert these dressed vertices into the conditioned
      percolation ancestry without expanding the Bernoulli
      normalization.
\item If the port coefficient cannot be extracted from the V78
      aggregate inequality, return to the V77 coefficientwise
      Cayley representation rather than differentiating the V78
      scalar bound.
\end{enumerate}
\end{programme}

\begin{firewall}[Claim boundary after V83]
\label{fire:newV83-boundary}
V83 replaces the unnecessarily termwise ABEC target by AABEC,
proves AABEC is sufficient, proves the exact Kruskal
maximum-spanning-tree identity, and shows that every fixed
percolation shortcut graph costs no tree census.  It also proves
that Bernoulli closed-edge normalization cannot be dropped and
isolates CPAC's dressed block-vertex estimate as the remaining row.

V83 does not prove CPAC, AABEC, ABEC, TBRC, full MCWC, OCRC,
SCNGC, WPSC, GRLTC, LZTC, PLRC, WSDC, BRSC, HWSFC, UGCGC,
thermal connected WPRAC, all-arity RGSC, the raw Wilson aggregate
strong card, all-order strong MTA, WUFC, CPM, cutoff removal,
reflection positivity, Osterwalder--Schrader reconstruction, a
continuum Yang--Mills measure, or a positive Hamiltonian gap.  The
full Yang--Mills mass gap has not been proved.
\end{firewall}

\begin{theorem}[V83 result ledger]
\label{thm:newV83-results}
V83:
\begin{enumerate}[label=\textup{(V83.\arabic*)},leftmargin=6.5em]
\item defines the aggregate block-edge card AABEC;
\item proves AABEC suffices for the intermediate wedge;
\item retires the termwise input-tree monomial as a requirement;
\item proves the exact Kruskal tree partition of unity;
\item pays every fixed shortcut graph at unit tree-sum cost;
\item reallocates the triangle shortcut to its output tree;
\item proves the closed-edge normalization no-go; and
\item isolates CPAC as the remaining dressed-vertex operation.
\end{enumerate}
\end{theorem}

\begin{proof}
Use
\Cref{thm:newV83-AABEC,
prop:newV83-old-ABEC,
thm:newV83-Kruskal,
cor:newV83-shortcut,
warn:newV83-normalization,
prop:newV83-CPAC}.
\end{proof}

\paragraph{Parent-version record.}
V83 was formed from
\texttt{Renewal\_Yang\_Mills\_Mass\_Gap\_Research\_Notes\_V82.tex}.
The V82 parent had \(37\,745\) logical source lines,
\(1\,208\,484\) bytes, and SHA--256
\[
 \texttt{cee404577eec967dac992ea9b5f606a61f4029c1112f8bf892cf66be0eb5ca7f}.
\]
The next compulsory companion audit is V85.

% =============================================================================
% END APPENDED FIFTH-SERIES V83 MATERIAL
% =============================================================================

\clearpage
\chapter{Fifth-series V84: ported Cayley vertices and the
remaining ancestry allocation}
\label{ch:newV84-ports}

\section{Labeled external ports}
\label{sec:newV84-ports}

Let \(P\in\mathcal P\) be a thermal block, put
\[
 n=|P|,\qquad w=w_P=\sum_{i\in P}b_i,
\]
and let \(D\) be a finite set of labeled external ports,
\(|D|=d\).  An anchor map is any function \(a:D\to P\).
For \(L\in\mathfrak T(P)\), use
\[
 \omega_P(L):=\prod_{i\in P}b_i^{d_L(i)}.
\]
When \(n=1\), \(\mathfrak T(P)\) contains the single empty tree and
\(\omega_P(L)=1\).

\begin{theorem}[Ported Cayley identity]
\label{thm:newV84-ported-Cayley}
For every \(n\ge1\) and \(d\ge0\),
\begin{align}
 &\sum_{L\in\mathfrak T(P)}
  \sum_{a:D\to P}
  \omega_P(L)\prod_{\ell\in D}b_{a(\ell)}
 \notag\\
 &\qquad=
 \left(\prod_{i\in P}b_i\right)w^{n-2+d}
 =
 v_Pw^d.                                                    \label{eq:newV84-ported-Cayley}
\end{align}
\end{theorem}

\begin{proof}
For \(n\ge2\), the weighted Cayley--Pr\"ufer identity gives
\[
 \sum_{L\in\mathfrak T(P)}\omega_P(L)
 =
 \left(\prod_{i\in P}b_i\right)w^{n-2}.
\]
Independently,
\[
 \sum_{a:D\to P}\prod_{\ell\in D}b_{a(\ell)}
 =
 \prod_{\ell\in D}\left(\sum_{i\in P}b_i\right)
 =
 w^d.
\]
Multiply.  If \(n=1\), the left side is \(b_i^d\), while the
right side is \(b_i b_i^{d-1}=b_i^d\), including \(d=0\) under the
singleton convention \(v_P=1\).
\end{proof}

\begin{corollary}[No port-choice arity loss]
\label{cor:newV84-no-port-loss}
The \(n^d\) possible anchor maps must not be counted without their
weights.  Their complete weighted sum is exactly \(w^d\), with no
constant depending on \(n\) or \(d\).
\end{corollary}

\begin{proof}
It is the second factor in the proof of
\Cref{thm:newV84-ported-Cayley}.
\end{proof}

\section{Why the ports are labeled}
\label{sec:newV84-labeled}

\begin{lemma}[Port-degree ledger]
\label{lem:newV84-port-degree}
Let \(U\) be an output tree on the block set \(\mathcal P\).
For each block \(P\), label its port set by the incident edges of
\(U\):
\[
 D_P(U):=\{e\in E(U):P\in e\}.
\]
Then
\begin{equation}
 |D_P(U)|=d_U(P),
 \qquad
 \sum_{P\in\mathcal P}|D_P(U)|=2r-2.                       \label{eq:newV84-port-degree}
\end{equation}
Thus any fibre bounded by \(K^{n_P+|D_P(U)|}\) at block \(P\)
costs at most \(K^{3m}\) globally.
\end{lemma}

\begin{proof}
The first two identities are the handshaking lemma for a tree.
Also \(\sum_Pn_P=m\) and \(r\le m\), so
\[
 \sum_P\{n_P+|D_P(U)|\}=m+2r-2\le3m-2.
\]
\end{proof}

\begin{assessment}[Ports are half-edges, not indistinguishable
insertions]
\label{ass:newV84-labeled}
Every port is owned by a specified neighboring block in the output
tree.  Therefore no division by \(d!\), and no multiplication by
the number of unordered port partitions, belongs in the local
ledger.  Repeated anchors at one representation label are allowed;
they supply repeated powers of that label's mark exactly as in the
Pr\"ufer degree monomial.
\end{assessment}

\section{A tree-valued form of the V77 core}
\label{sec:newV84-tree-valued}

\begin{proposition}[Incidence-hypertree output remains tree-valued]
\label{prop:newV84-tree-valued}
Let \(P\) consist only of nonlinear-centre labels and suppose its
effective-hyperedge incidence graph is a tree.  Before the final
weighted Cayley sum is evaluated, the complete V77 positive
majorant has the form
\begin{equation}
 K^n
 \sum_{L\in\mathfrak T(P)}
 M(L)\prod_{i\in P}b_i^{d_L(i)},
 \qquad
 0\le M(L)\le2^{n-2}.                                      \label{eq:newV84-tree-valued}
\end{equation}
\end{proposition}

\begin{proof}
In \Cref{thm:newV77-global-tree}, the quotient after one common
factor \(b_i\) per centre is
\[
 \sum_LM(L)\prod_i b_i^{d_L(i)-1}.
\]
Restore \(\prod_i b_i\).  The omitted local effective-hyperedge
constants have total arity at most \(2n-2\) by the incidence-tree
degree identity and therefore enter \(K^n\).  The fibre estimate is
\Cref{lem:newV77-fibre-bound}.
\end{proof}

\begin{corollary}[Ported incidence-hypertree centre]
\label{cor:newV84-ported-hypertree}
If the centre-core ancestry of
\Cref{prop:newV84-tree-valued} also carries \(d\) already-owned,
labeled external half-edges, each with an explicit anchor factor
\(b_{a(\ell)}\), then summing the local output trees and all anchor
maps costs at most
\begin{equation}
 K^{n+d}v_Pw_P^d.                                          \label{eq:newV84-ported-hypertree}
\end{equation}
\end{corollary}

\begin{proof}
Use \(M(L)\le2^{n-2}\) in
\eqref{eq:newV84-tree-valued}, then apply
\Cref{thm:newV84-ported-Cayley}.  The factor \(2^{n-2}\) is
absorbed in the exponential base; allowing \(K^d\) accommodates
the declared ownership constants of the ports.
\end{proof}

\begin{assessment}[Scope of the proved ported sector]
\label{ass:newV84-scope}
\Cref{cor:newV84-ported-hypertree} proves the desired local vertex
factor when the external merger expansion has already exposed one
anchor mark per port and the remaining local centre incidence is
acyclic.  V78 closes cyclic and whole-surplus centre cores after
the scalar Cayley sum, but its published bound does not retain the
output tree and anchor map coefficientwise.  That refinement may
not be inferred by differentiating the scalar V78 inequality.
\end{assessment}

\section{The exact remaining allocation card}
\label{sec:newV84-allocation}

\begin{definition}[Ported ancestry allocation card, PAAC]
\label{def:newV84-PAAC}
PAAC is an allocation of every conditioned CPAC operator ancestry,
before norms, to:
\begin{enumerate}[label=\textup{(\roman*)},leftmargin=3.5em]
\item one connected output block graph, followed by its Kruskal
      output tree \(U\);
\item for every block \(P\), one local Cayley tree
      \(L_P\in\mathfrak T(P)\);
\item an anchor map
      \(a_P:D_P(U)\to P\); and
\item a fibre label taking at most
      \(K^{n_P+d_U(P)}\) values locally,
\end{enumerate}
such that the ancestry norm before the sums is bounded by
\begin{equation}
 K^m
 \prod_{P\in\mathcal P}
 \left\{
  \omega_P(L_P)
  \prod_{\ell\in D_P(U)}b_{a_P(\ell)}
 \right\}.
 \label{eq:newV84-PAAC}
\end{equation}
The normalized Bernoulli probability is retained as required by
\Cref{warn:newV83-normalization}.
\end{definition}

\begin{theorem}[PAAC closes CPAC]
\label{thm:newV84-PAAC}
PAAC implies CPAC, AABEC, TBRC, full MCWC, and SCNGC.
\end{theorem}

\begin{proof}
For a fixed output tree \(U\), sum
\eqref{eq:newV84-PAAC} over every local tree and anchor map.
\Cref{thm:newV84-ported-Cayley} gives
\[
 K^m\prod_{P\in\mathcal P}v_Pw_P^{d_U(P)}.
\]
\Cref{lem:newV84-port-degree} absorbs the allocation fibres into a
larger \(K^m\).  Sum the output trees with the weighted Cayley
identity to obtain the right side of
\eqref{eq:newV83-AABEC}.  The Kruskal input-tree sum and the
Bernoulli configuration sum cost one by
\Cref{thm:newV83-Kruskal,warn:newV83-normalization}.  Hence CPAC
and AABEC hold.  Apply
\Cref{thm:newV83-AABEC,prop:newV79-TBRC}.
\end{proof}

\begin{proposition}[Heavy-singleton PAAC]
\label{prop:newV84-heavy}
PAAC holds at every heavy singleton block with fibre one.
\end{proposition}

\begin{proof}
If \(P=\{i\}\), the local tree is empty and every incident port has
the unique anchor \(i\).  Its weight in
\eqref{eq:newV84-PAAC} is \(b_i^{d_U(P)}=w_P^{d_U(P)}\), exactly
the singleton case of
\eqref{eq:newV84-ported-Cayley}.
\end{proof}

\begin{proposition}[Acyclic exposed-anchor PAAC]
\label{prop:newV84-acyclic-PAAC}
PAAC holds at a light block whenever:
\begin{enumerate}[label=\textup{(\roman*)},leftmargin=3.5em]
\item its internal nonlinear-centre incidence graph is a tree; and
\item every external output-tree half-edge has already exposed one
      anchor label in the block and one factor of that anchor's
      mark.
\end{enumerate}
\end{proposition}

\begin{proof}
Use \Cref{cor:newV84-ported-hypertree}.  Its fibre
\(2^{n_P-2}\) satisfies the PAAC exponential-fibre requirement.
\end{proof}

\begin{assessment}[The remaining local cases]
\label{ass:newV84-position}
The port power and all anchor choices are now summed exactly.
Heavy vertices and acyclic light vertices with exposed anchors
satisfy PAAC.  The remaining cases are precise:
\begin{enumerate}[label=\textup{(\roman*)},leftmargin=3.5em]
\item a cyclic or whole-surplus V78 core must retain its output
      Cayley tree rather than only its scalar sum; and
\item a mixed derivative of the full squarefree logarithm must
      expose one anchor mark for every incident output-tree
      half-edge.
\end{enumerate}
These are coefficientwise ancestry statements; neither follows
from an already-summed norm inequality.
\end{assessment}

\begin{programme}[V85: compulsory audit and tree-valued CIRC]
\label{prog:newV84-V85}
\begin{enumerate}[leftmargin=3em]
\item Perform the compulsory newest-SM/newest-NS audit.
\item Revisit the canonical skeleton proof of V78 without summing
      the V77 output Cayley tree.
\item Attach every skeleton extension and whole-surplus hyperedge
      to a canonical edge or vertex of that output tree.
\item Prove that the attachment fibre is exponential while the
      tree monomial is unchanged or improved.
\item Then return to the full-logarithm derivative and prove the
      exposed-anchor row of PAAC.
\end{enumerate}
\end{programme}

\begin{firewall}[Claim boundary after V84]
\label{fire:newV84-boundary}
V84 proves the exact ported Cayley identity, the global port-degree
ledger, the tree-valued V77 incidence-hypertree form, the complete
ported acyclic-centre sector, and the heavy-singleton vertex.  It
formulates PAAC and proves that PAAC closes CPAC, AABEC, TBRC,
MCWC, and SCNGC.

V84 does not prove cyclic tree-valued CIRC, the general
exposed-anchor row, PAAC, CPAC, AABEC, TBRC, full MCWC, OCRC,
SCNGC, WPSC, GRLTC, LZTC, PLRC, WSDC, BRSC, HWSFC, UGCGC,
thermal connected WPRAC, all-arity RGSC, the raw Wilson aggregate
strong card, all-order strong MTA, WUFC, CPM, cutoff removal,
reflection positivity, Osterwalder--Schrader reconstruction, a
continuum Yang--Mills measure, or a positive Hamiltonian gap.  The
full Yang--Mills mass gap has not been proved.
\end{firewall}

\begin{theorem}[V84 result ledger]
\label{thm:newV84-results}
V84:
\begin{enumerate}[label=\textup{(V84.\arabic*)},leftmargin=6.5em]
\item proves the exact ported Cayley identity;
\item proves that anchor choices cost \(w_P^d\), not \(n_P^d\);
\item proves the global port-degree ledger;
\item retains the V77 output Cayley tree coefficientwise;
\item closes the ported acyclic centre-core sector;
\item formulates PAAC and proves its sufficiency;
\item closes every heavy-singleton block vertex; and
\item closes acyclic light vertices with exposed anchors.
\end{enumerate}
\end{theorem}

\begin{proof}
Use
\Cref{thm:newV84-ported-Cayley,
lem:newV84-port-degree,
prop:newV84-tree-valued,
cor:newV84-ported-hypertree,
thm:newV84-PAAC,
prop:newV84-heavy,
prop:newV84-acyclic-PAAC}.
\end{proof}

\paragraph{Parent-version record.}
V84 was formed from
\texttt{Renewal\_Yang\_Mills\_Mass\_Gap\_Research\_Notes\_V83.tex}.
The V83 parent had \(38\,057\) logical source lines,
\(1\,219\,669\) bytes, and SHA--256
\[
 \texttt{05561b4f94b06d6e7e1bbbe6e5b8e116bd68664a0e4b732de34fb94f6eacb4e1}.
\]
The next compulsory companion audit is V85.

% =============================================================================
% END APPENDED FIFTH-SERIES V84 MATERIAL
% =============================================================================

\clearpage
\chapter{Fifth-series V85: compulsory audit and tree-valued CIRC}
\label{ch:newV85-tree-CIRC}

\section{Compulsory companion audit}
\label{sec:newV85-audit}

\begin{audit}[V85 immutable companion snapshots]
\label{audit:newV85-companions}
The newest active companion sources were read without modification:
\begin{align*}
 &\texttt{Renewal\_SM\_Einstein\_Continuum\_Research\_Notes\_V9.tex},\\
 &\qquad 3\,991\ \hbox{logical lines},\quad
 155\,101\ \hbox{bytes},\\
 &\qquad
 \texttt{7fbe1d2664eb4bc6995c96554b1f852625c0b54d4fa30f3e1de7b6a0c66f4f41},
 \\
 &\texttt{Renewal\_NS\_Stability\_Research\_Notes\_V60.tex},\\
 &\qquad 14\,557\ \hbox{logical lines},\quad
 749\,203\ \hbox{bytes},\\
 &\qquad
 \texttt{2c13df20c489b7df7cc3c88654156e3ac5d0edc39660ee61ef180722e71c8305}.
\end{align*}

SM V9 derives the exact overlap heat Hessian as the sum of its two
separated Duhamel insertions and the contact term, proves the
one-point functional vanishes, and obtains exact lattice
transversality before any norm estimate.  It proves explicitly that
estimating the separated and contact terms independently destroys
the Ward cancellation and cannot exclude a mass tensor.

NS V60 consolidates a bounded-overlap local trap and a quantitative
Morrey floor, then pinches the energy of a concentrating ball between
two quantities of the same parabolic order.  Its local energy
identity has an error of that same order and therefore does not
retain enough information to connect successive windows; the
coefficientwise inter-window row remains open.
\end{audit}

\begin{assessment}[Audit transfer]
\label{ass:newV85-audit-transfer}
No companion estimate is imported.  The proof-order lessons are:
\begin{enumerate}[label=\textup{(\roman*)},leftmargin=3.5em]
\item Gaussian drift pairs and jump mergers must be assembled before
      anchor ownership is asserted, just as the SM contact and
      separated terms must be assembled before transversality; and
\item a scalar bound of the correct total degree need not retain the
      coefficientwise tree needed at the next scale, just as the NS
      ball-energy pinch does not propagate a cluster between windows.
\end{enumerate}
V85 therefore revisits the V78 proof itself and retains its V77
output tree.
\end{assessment}

\section{Centre-only output trees survive the V78 skeleton}
\label{sec:newV85-centres}

Let \({\cal C}\) be a centre set of cardinality \(p\ge2\), and put
\(B_{\cal C}=\sum_{i\in{\cal C}}b_i\le b_*\).  This section is at
zero noncentral-label variable, so every effective hyperedge is
\(W_A(0)\).

\begin{theorem}[Tree-valued centre-only CIRC]
\label{thm:newV85-tree-CIRC}
After one enlargement of the numerical base,
\begin{equation}
 \sum_{\substack{{\cal H}\ {\rm connected\ simple}\\
                  {\cal H}\ {\rm on}\ {\cal C}}}
 \prod_{A\in{\cal H}}|W_A(0)|
 \le
 K^p
 \sum_{L\in\mathfrak T({\cal C})}
 \prod_{i\in{\cal C}}b_i^{d_L(i)}.
 \label{eq:newV85-tree-CIRC}
\end{equation}
Thus cyclic cores, incidence extensions, and wholly surplus
hyperedges do not require the final Cayley sum to be evaluated.
\end{theorem}

\begin{proof}
Use the canonical skeleton injection
\Cref{lem:newV78-skeleton}.  For its base incidence hypertree,
\Cref{prop:newV84-tree-valued} gives
\[
 K_0^p
 \sum_{L\in\mathfrak T({\cal C})}
 M(L)\prod_i b_i^{d_L(i)},
 \qquad M(L)\le2^{p-2}.
\]
Absorb \(M(L)\) into \(K_1^p\), but do not perform the sum over
\(L\).

For each retained skeleton hyperedge \(C\), sum every incidence
extension \(J\) by \Cref{lem:newV78-extension}.  This multiplies
the preceding tree-valued expression by at most
\(\exp(CB_{\cal C}^2)\).  There are at most \(p-1\) skeleton
hyperedges, so all extensions cost
\[
 \exp(CpB_{\cal C}^2)\le K_2^p.
\]
This factor is independent of \(L\).

Sum the set of wholly surplus hyperedges by
\Cref{lem:newV78-whole-gas}; its factor
\(\exp(CB_{\cal C}^2)\) is also independent of \(L\) and is
absorbed into \(K_3^p\).  The resulting expression is exactly
\eqref{eq:newV85-tree-CIRC}.
\end{proof}

\begin{corollary}[The scalar V78 theorem is recovered]
\label{cor:newV85-scalar}
Evaluating the final tree sum in
\eqref{eq:newV85-tree-CIRC} gives
\[
 K^p\left(\prod_{i\in{\cal C}}b_i\right)
 B_{\cal C}^{p-2},
\]
the centre-only case of \Cref{thm:newV78-CIRC}.
\end{corollary}

\begin{proof}
Apply the weighted Cayley--Pr\"ufer identity.
\end{proof}

\begin{assessment}[Surplus roots do not alter the output tree]
\label{ass:newV85-surplus}
The V78 incidence extensions and wholly surplus hyperedges were
summed by positive gases after the canonical skeleton had already
produced its V77 output tree.  Their bounds are scalar factors
independent of that tree.  Therefore the same output tree can own
the complete cyclic core; no new tree choice and no surplus-root
attachment fibre is required.
\end{assessment}

\section{Ported cyclic centre cores}
\label{sec:newV85-ported}

\begin{corollary}[Full centre-core PAAC with exposed anchors]
\label{cor:newV85-centre-PAAC}
Let \(D\) be a set of \(d\) labeled external ports.  Suppose each
port has already exposed an anchor \(a(\ell)\in{\cal C}\) and the
factor \(b_{a(\ell)}\).  Then the complete connected centre-core
majorant, including arbitrary incidence excess, summed over all
output trees and anchor maps, is at most
\begin{equation}
 K^{p+d}
 \left(\prod_{i\in{\cal C}}b_i\right)
 B_{\cal C}^{p-2+d}.
 \label{eq:newV85-centre-PAAC}
\end{equation}
\end{corollary}

\begin{proof}
Multiply the tree-valued right side of
\eqref{eq:newV85-tree-CIRC} by the anchor weights and sum
\(a:D\to{\cal C}\).  Apply the ported Cayley identity
\Cref{thm:newV84-ported-Cayley}.  The optional \(K^d\) records the
uniform port-ownership constants.
\end{proof}

\begin{corollary}[Cyclic centre row of PAAC]
\label{cor:newV85-cyclic-PAAC}
PAAC holds at a light block for every centre-only local ancestry,
acyclic or cyclic, once every incident output-tree half-edge has
exposed one anchor mark in that block.
\end{corollary}

\begin{proof}
Use \Cref{cor:newV85-centre-PAAC}; the local allocation fibre is
already included in its exponential base.
\end{proof}

\section{The noncentral-label restoration cannot be inferred}
\label{sec:newV85-noncentral}

\begin{warning}[The adaptive Cauchy factor is not tree-valued]
\label{warn:newV85-Cauchy}
In V72--V78, \(q\) noncentral labels were restored from one global
coefficient variable by the bound
\[
 q!R^{-q}e^{C\eta q}\le K^qB^q.
\]
This is sufficient after the centre Cayley sum has been evaluated,
but the scalar power \(B^q\) does not specify how the noncentral
labels attach to an output tree.  It therefore cannot be used as a
proof of PAAC for a full thermal block.
\end{warning}

\begin{proof}
The monomial \(B^q\) includes arbitrary parent choices for every
one of the \(q\) labeled factors, including choices which form
directed cycles among the noncentral labels.  A tree-valued
restoration must cancel or reallocate those choices before norms.
The scalar inequality contains no such allocation data.
\end{proof}

\begin{definition}[Tree-valued label-restoration card, TLRC]
\label{def:newV85-TLRC}
TLRC is a coefficientwise refinement of the global noncentral-label
generating function which, for a fixed centre-core output tree:
\begin{enumerate}[label=\textup{(\roman*)},leftmargin=3.5em]
\item restores every noncentral label exactly once;
\item attaches all restored labels to produce an ordinary tree on
      the full thermal block;
\item retains every already exposed external port anchor; and
\item has fibre at most \(K^{n+d}\), where \(n\) is the full block
      size and \(d\) its port degree.
\end{enumerate}
\end{definition}

\begin{proposition}[TLRC completes the local PAAC row]
\label{prop:newV85-TLRC}
TLRC, together with
\Cref{cor:newV85-cyclic-PAAC,prop:newV84-heavy}, proves PAAC.
\end{proposition}

\begin{proof}
At a heavy singleton use \Cref{prop:newV84-heavy}.  At a light
block, \Cref{cor:newV85-cyclic-PAAC} supplies a tree-valued core on
all nonlinear centres and preserves the exposed port anchors.
TLRC attaches every remaining label and yields the local tree and
anchor data of \Cref{def:newV84-PAAC}.  The fibre exponents sum to
at most a constant times \(m\) by
\Cref{lem:newV84-port-degree}.
\end{proof}

\begin{assessment}[Position after tree-valued CIRC]
\label{ass:newV85-position}
Cyclic centre cores no longer obstruct PAAC: their V77 output tree
survives every V78 surplus sum.  The remaining local operation is
now only TLRC, followed by the separate exposed-anchor theorem for
the complete assembled logarithmic derivative.  The audit forbids
claiming either from an already-summed scalar estimate.
\end{assessment}

\begin{programme}[V86: coefficientwise noncentral leaves]
\label{prog:newV85-V86}
\begin{enumerate}[leftmargin=3em]
\item Return to the exact global EGF before its adaptive Cauchy
      estimate and give every noncentral label its own variable.
\item Expand each effective hyperedge into its V77 local tree before
      extracting the squarefree coefficient.
\item Use a Prüfer word to attach noncentral labels to the
      centre-core output tree, and bound the fibre coefficientwise.
\item Prove the pure-star and path attachment sectors first; test
      whether a directed cycle of noncentral parent choices cancels
      in the full squarefree logarithm.
\item Only after TLRC is proved, address anchor exposure in the
      assembled drift-plus-jump merger.
\end{enumerate}
\end{programme}

\begin{firewall}[Claim boundary after V85]
\label{fire:newV85-boundary}
V85 performs the compulsory SM/NS audit, proves tree-valued
centre-only CIRC, retains the V77 output tree through all V78
surplus sums, and closes cyclic centre-core PAAC whenever external
anchors are exposed.  It proves that the old adaptive Cauchy factor
does not retain enough information and isolates TLRC.

V85 does not prove TLRC, general PAAC, CPAC, AABEC, TBRC, full
MCWC, OCRC, SCNGC, WPSC, GRLTC, LZTC, PLRC, WSDC, BRSC, HWSFC,
UGCGC, thermal connected WPRAC, all-arity RGSC, the raw Wilson
aggregate strong card, all-order strong MTA, WUFC, CPM, cutoff
removal, reflection positivity, Osterwalder--Schrader
reconstruction, a continuum Yang--Mills measure, or a positive
Hamiltonian gap.  The full Yang--Mills mass gap has not been proved.
\end{firewall}

\begin{theorem}[V85 result ledger]
\label{thm:newV85-results}
V85:
\begin{enumerate}[label=\textup{(V85.\arabic*)},leftmargin=6.5em]
\item performs the compulsory active-SM/active-NS audit;
\item imports only assembled-cancellation and coefficient-retention
      discipline;
\item proves tree-valued centre-only CIRC;
\item retains the same output tree through every surplus-root sum;
\item recovers the scalar V78 theorem as a corollary;
\item closes ported cyclic centre cores with exposed anchors;
\item proves the global Cauchy factor is not tree-valued; and
\item formulates TLRC and proves its sufficiency for local PAAC.
\end{enumerate}
\end{theorem}

\begin{proof}
Use
\Cref{audit:newV85-companions,
thm:newV85-tree-CIRC,
cor:newV85-scalar,
cor:newV85-centre-PAAC,
cor:newV85-cyclic-PAAC,
warn:newV85-Cauchy,
prop:newV85-TLRC}.
\end{proof}

\paragraph{Parent-version record.}
V85 was formed from
\texttt{Renewal\_Yang\_Mills\_Mass\_Gap\_Research\_Notes\_V84.tex}.
The V84 parent had \(38\,397\) logical source lines,
\(1\,231\,133\) bytes, and SHA--256
\[
 \texttt{13401335946752fb68ec4a64466b21cff9cd42b39c5ea618e4e1b8a56199b88e}.
\]
The next compulsory companion audit is V90.

% =============================================================================
% END APPENDED FIFTH-SERIES V85 MATERIAL
% =============================================================================

\clearpage
\chapter{Fifth-series V86: normalized root occurrences on all
thermal labels}
\label{ch:newV86-occurrences}

\section{Return before the centre contraction}
\label{sec:newV86-return}

No companion audit is due at V86.

V85 formulated TLRC after the degree-one labels had been contracted
into effective centre hyperedges.  That formulation is sufficient,
but it is not the only possible proof order.  The present note
returns to the active root occurrences before that contraction and
treats every thermal label as a vertex of the incidence
hypergraph.  The V77--V78 combinatorics itself does not require a
vertex to have incidence degree at least two.

There is, however, one normalization issue.  Different Duhamel
blocks may select the same root support.  A connected
\emph{simple} support hypergraph is already covered by V85, whereas
a repeated support must first be summed with the factorial supplied
by an assembled analytic type source.  The next theorem records
exactly what is obtained when that factorial is present.

Let \(X\) be a finite label set, \(n=|X|\ge2\), and put
\[
 B_X:=\sum_{i\in X}b_i,\qquad b_i\ge\sqrt s .
\]
For \(A\subseteq X\), \(|A|=a\ge2\), use the two root proxies
\begin{align}
 {\mathfrak r}_X(A)
 &:=(Ka)^a s^{(a-2)/2}\prod_{i\in A}b_i,                    \label{eq:newV86-r}\\
 {\mathfrak q}_X(A)
 &:=K^a\left(\prod_{i\in A}b_i\right)B_A^{a-2}.             \label{eq:newV86-q}
\end{align}
Both bounds hold for the assembled drift-plus-jump root \(R_A\) by
\Cref{lem:newV61-reserve,thm:newV60-hyperedge-card}; at \(a=2\)
the drift term is included in
\Cref{lem:newV61-reserve}.

\section{Exact multiplicity coarsening}
\label{sec:newV86-coarsening}

\begin{definition}[Exponentially normalized support multihypergraph]
\label{def:newV86-normalized-multi}
Let \(u_A\ge0\) for every \(A\subseteq X\), \(|A|\ge2\).
A multiplicity vector is
\[
 \boldsymbol\nu=(\nu_A)_{A\subseteq X,\ |A|\ge2},
 \qquad \nu_A\in{\mathbb N}_0.
\]
Its occupied simple support is
\[
 {\rm supp}\,\boldsymbol\nu:=\{A:\nu_A\ge1\}.
\]
It is connected when this occupied support is a connected
hypergraph on all vertices of \(X\).  Its normalized weight is
\begin{equation}
 {\sf w}(\boldsymbol\nu)
 :=
 \prod_A{u_A^{\nu_A}\over \nu_A!}.                          \label{eq:newV86-normalized-weight}
\end{equation}
\end{definition}

\begin{lemma}[Exact support coarsening]
\label{lem:newV86-exact-coarsening}
For arbitrary nonnegative activities,
\begin{equation}
 \sum_{\substack{\boldsymbol\nu\\
                  {\rm supp}\,\boldsymbol\nu\ {\rm connected}}}
 {\sf w}(\boldsymbol\nu)
 =
 \sum_{\substack{{\cal H}\ {\rm connected\ simple}\\
                  {\cal H}\ {\rm on}\ X}}
 \prod_{A\in{\cal H}}\bigl(e^{u_A}-1\bigr).                 \label{eq:newV86-exact-coarsening}
\end{equation}
\end{lemma}

\begin{proof}
Fix the occupied support \({\cal H}\).  Multiplicities at distinct
supports are independent, and
\[
 \sum_{\nu_A\ge1}{u_A^{\nu_A}\over\nu_A!}=e^{u_A}-1.
\]
Multiplication over \(A\in{\cal H}\), followed by summation over
the connected simple supports, proves
\eqref{eq:newV86-exact-coarsening}.  All sums are finite products
of nonnegative entire series, so no rearrangement issue occurs.
\end{proof}

\begin{lemma}[The two proxies survive exponentiation]
\label{lem:newV86-proxy-inheritance}
There is a numerical \(b_*>0\) such that, if \(B_X\le b_*\) and
\[
 0\le u_A\le
 \min\{{\mathfrak r}_X(A),{\mathfrak q}_X(A)\},
                                                               \label{eq:newV86-two-proxy}
\]
then, after one enlargement of \(K\),
\begin{equation}
 e^{u_A}-1
 \le
 \min\{{\mathfrak r}_X(A),{\mathfrak q}_X(A)\}               \label{eq:newV86-exp-proxy}
\end{equation}
with the enlarged proxies on the right.
\end{lemma}

\begin{proof}
The strong proxy and the elementary-symmetric bound give
\[
 u_A\le
 K^a\left(\prod_{i\in A}b_i\right)B_A^{a-2}
 \le
 K^a{B_A^{2a-2}\over a^a}.
\]
Decrease \(b_*\) so the last expression is at most \(1/2\) for
every \(a\ge2\).  Hence \(e^{u_A}-1\le2u_A\).  Since \(a\ge2\),
the factor two is absorbed into a fixed enlargement of the
exponential base in both
\eqref{eq:newV86-r} and \eqref{eq:newV86-q}.
\end{proof}

\section{A tree-valued theorem on the full label set}
\label{sec:newV86-full-tree}

\begin{theorem}[Normalized full-label occurrence theorem]
\label{thm:newV86-full-label}
Suppose \(B_X\le b_*\) and the activities obey
\eqref{eq:newV86-two-proxy}.  Then
\begin{equation}
 \sum_{\substack{\boldsymbol\nu\\
                  {\rm supp}\,\boldsymbol\nu\ {\rm connected}}}
 {\sf w}(\boldsymbol\nu)
 \le
 K^n
 \sum_{L\in\mathfrak T(X)}
 \prod_{i\in X}b_i^{d_L(i)}.                               \label{eq:newV86-full-label}
\end{equation}
In particular, every degree-one label already belongs to the
ordinary output tree; no separate scalar Cauchy restoration is
used.
\end{theorem}

\begin{proof}
By \Cref{lem:newV86-exact-coarsening}, the left side is the
connected simple-hypergraph sum with effective activity
\(\widetilde u_A=e^{u_A}-1\).  By
\Cref{lem:newV86-proxy-inheritance}, every
\(\widetilde u_A\) obeys the reserve and strong proxies used in
V78.

Apply the canonical incidence-skeleton injection
\Cref{lem:newV78-skeleton} with \(X\), rather than only the
nonlinear-centre set, as its vertex set.  On the retained
incidence hypertree, the proof of
\Cref{prop:newV84-tree-valued} is purely combinatorial and gives
\[
 K_0^n
 \sum_{L\in\mathfrak T(X)}
 M(L)\prod_{i\in X}b_i^{d_L(i)},
 \qquad M(L)\le2^{n-2}.
\]
Every incidence extension costs
\(\exp(CB_X^2)\) by
\Cref{lem:newV78-extension}; there are at most \(n-1\) retained
hyperedges.  Every wholly surplus support costs
\(\exp(CB_X^2)\) by
\Cref{lem:newV78-whole-gas}.  These factors are independent of
the output tree.  Absorb them and \(2^{n-2}\) into \(K^n\).
This proves \eqref{eq:newV86-full-label}.
\end{proof}

\begin{corollary}[Scalar strong form]
\label{cor:newV86-full-label-scalar}
Under the hypotheses of \Cref{thm:newV86-full-label},
\begin{equation}
 \sum_{\substack{\boldsymbol\nu\\
                  {\rm supp}\,\boldsymbol\nu\ {\rm connected}}}
 {\sf w}(\boldsymbol\nu)
 \le
 K^n\left(\prod_{i\in X}b_i\right)B_X^{n-2}.                \label{eq:newV86-full-label-scalar}
\end{equation}
\end{corollary}

\begin{proof}
Evaluate the final tree sum in
\eqref{eq:newV86-full-label} by the weighted
Cayley--Pr\"ufer identity.
\end{proof}

\begin{corollary}[Ported normalized occurrence vertex]
\label{cor:newV86-ported}
Let \(D\) be a set of \(d\) labeled external ports.  If every port
has already exposed an anchor \(a(\ell)\in X\) and its factor
\(b_{a(\ell)}\), then the normalized connected occurrence sum,
summed over its output tree and all anchor maps, is at most
\begin{equation}
 K^{n+d}
 \left(\prod_{i\in X}b_i\right)B_X^{n-2+d}.                 \label{eq:newV86-ported}
\end{equation}
\end{corollary}

\begin{proof}
Multiply the tree-valued right side of
\eqref{eq:newV86-full-label} by the \(d\) anchor marks, sum the
anchor maps, and apply
\Cref{thm:newV84-ported-Cayley}.  The optional \(K^d\) pays the
uniform port-ownership constants.
\end{proof}

\begin{assessment}[What has been gained]
\label{ass:newV86-gain}
The centre/noncentral distinction is unnecessary after root
multiplicities have been normalized by support.  Incidence leaves
are ordinary vertices in the V77 gluing, and incidence cycles are
ordinary V78 surplus data.  Thus the normalized occurrence problem
produces the full local Cayley tree in one step and avoids the
directed-parent cycle identified in
\Cref{warn:newV85-Cauchy}.
\end{assessment}

\section{Where the factorial comes from}
\label{sec:newV86-Duhamel}

\begin{lemma}[Noncommutative assembled-type coefficient]
\label{lem:newV86-type-factorial}
Let \(V=V^*\preceq\varepsilon I\), let
\(\{H_\alpha\}_{\alpha\in I}\) be a finite operator family, and
let \(\boldsymbol\nu\in{\mathbb N}_0^I\), with
\(q=\sum_\alpha\nu_\alpha\).  Then
\begin{equation}
 \left\|
 [\boldsymbol z^{\boldsymbol\nu}]
 \exp\left(V+\sum_{\alpha\in I}z_\alpha H_\alpha\right)
 \right\|_{\rm cb}
 \le
 e^\varepsilon
 \prod_{\alpha\in I}
 {\|H_\alpha\|_{\rm cb}^{\nu_\alpha}\over\nu_\alpha!}.
                                                               \label{eq:newV86-type-factorial}
\end{equation}
No commutativity among the \(H_\alpha\) is assumed.
\end{lemma}

\begin{proof}
The coefficient is the sum of the time-ordered simplex integrals
over all distinct words containing \(\nu_\alpha\) copies of each
letter \(H_\alpha\).  There are
\[
 {q!\over\prod_\alpha\nu_\alpha!}
\]
such words, while the simplex has volume \(1/q!\).
As in \Cref{lem:newV70-Duhamel}, the product of all intervening
semigroup norms is at most \(e^\varepsilon\).  Taking norms after
the word count and simplex volume have been combined gives
\eqref{eq:newV86-type-factorial}.
\end{proof}

\begin{corollary}[Already assembled supports satisfy the theorem]
\label{cor:newV86-assembled}
If the exact local source is first written in the analytic form
\[
 \exp\left\{V+\sum_{\substack{A\subseteq X\\|A|\ge2}}
 z_AH_A\right\},
\]
with one type variable per root support and
\[
 \|H_A\|_{\rm cb}
 \le\min\{{\mathfrak r}_X(A),{\mathfrak q}_X(A)\},
\]
then its connected support coefficients obey
\Cref{thm:newV86-full-label,cor:newV86-ported}.
\end{corollary}

\begin{proof}
Apply \Cref{lem:newV86-type-factorial}; the resulting positive
coefficient majorant is exactly
\eqref{eq:newV86-normalized-weight}.  Then apply
\Cref{thm:newV86-full-label,cor:newV86-ported}.
\end{proof}

\section{The exact remaining coarsening operation}
\label{sec:newV86-card}

In the BKAR derivative, an insertion is not initially indexed only
by \(A\).  It is indexed by a pair \((A,E)\), where \(E\) is the
nonempty set of outer-tree edge derivatives assigned to that
Duhamel block.  The estimate
\Cref{lem:newV60-gate-derivative} retains a spanning tree on \(A\)
which contains \(E\).  Distinct blocks carry disjoint \(E\)'s, but
several of them may have the same support \(A\).  Therefore
\Cref{lem:newV86-type-factorial} cannot simply be applied with
type \(A\): first the edge-decorated insertions must be assembled
without forgetting which external half-edge each block owns.

\begin{definition}[Anchored support-coarsening card, ASCC]
\label{def:newV86-ASCC}
ASCC is a reorganization of the exact partitioned BKAR--Duhamel
derivative, before norms, with the following properties:
\begin{enumerate}[label=\textup{(\roman*)},leftmargin=3.5em]
\item all insertions \((A,E)\) having the same label support \(A\)
      become coefficients of one analytic support source \(z_A\);
\item a multiplicity \(\nu_A\) consequently carries
      \(1/\nu_A!\) after the ordered simplex is summed;
\item every outer-tree half-edge retains one declared local anchor
      mark \(b_i\), including when several edge sets have been
      coarsened into the same support; and
\item the assembled support activity still obeys both proxies
      \eqref{eq:newV86-r}--\eqref{eq:newV86-q}, with only one
      exponential base.
\end{enumerate}
\end{definition}

\begin{proposition}[ASCC closes PAAC]
\label{prop:newV86-ASCC}
ASCC implies PAAC, hence also CPAC, AABEC, TBRC, full MCWC, and
SCNGC.
\end{proposition}

\begin{proof}
At a light block, ASCC(i)--(ii) put the connected occurrence sum
in the hypotheses of
\Cref{thm:newV86-full-label}.  ASCC(iii) supplies the labeled
ports required by \Cref{cor:newV86-ported}, and ASCC(iv) supplies
its two-proxy hypothesis.  Thus
\[
 K^{n_P+d_U(P)}v_Pw_P^{d_U(P)}
\]
pays the complete local vertex, which is
\eqref{eq:newV84-PAAC}.  Heavy singletons are
\Cref{prop:newV84-heavy}.  Hence PAAC holds.  Invoke
\Cref{thm:newV84-PAAC}.
\end{proof}

\begin{corollary}[Two sectors have no multiplicity obstruction]
\label{cor:newV86-closed-sectors}
The multiplicity-normalization part of ASCC is automatic for:
\begin{enumerate}[label=\textup{(\roman*)},leftmargin=3.5em]
\item every ancestry in which no two edge-decorated Duhamel blocks
      have the same support; and
\item every repeated-support ancestry which is already presented
      as one analytic type source before the outer-edge
      coefficient is extracted.
\end{enumerate}
In either sector, PAAC follows if the separate exposed-anchor and
two-proxy requirements ASCC(iii)--(iv) are verified.
\end{corollary}

\begin{proof}
In (i), every support multiplicity is zero or one, so no
multiplicity normalization is needed and the simple-support proof
inside \Cref{thm:newV86-full-label} applies once its two proxies are
available.  In (ii), the required factorial is
\Cref{lem:newV86-type-factorial}.  If ASCC(iii)--(iv) also hold,
\Cref{cor:newV86-ported} gives the PAAC vertex.
\end{proof}

\begin{warning}[The missing factorial may not be inserted after
norms]
\label{warn:newV86-factorial}
The set partition of the outer-tree derivative edges is unlabeled,
but its blocks are edge-decorated insertions \((A,E)\), not copies
of a pre-existing support type \(A\).  Replacing their positive
sum by \(\prod_A1/\nu_A!\) without an exact support-source identity
would be circular.  If norms are taken first, the star family of
\Cref{prop:newV60-Bell-no-go} reappears.
\end{warning}

\begin{assessment}[Position after the upstream return]
\label{ass:newV86-position}
TLRC has been replaced by a more precise route.  The entire
full-label hypergraph, including degree-one labels and incidence
cycles, has a tree-valued strong bound once support multiplicities
are exponentially normalized.  Noncommutativity does not destroy
that factorial after the support types are assembled.  The only
unproved local operation is now ASCC: coarsening the
edge-decorated gate derivatives into those support types while
retaining their anchors and both root proxies.
\end{assessment}

\begin{programme}[V87: polarized support sources]
\label{prog:newV86-V87}
\begin{enumerate}[leftmargin=3em]
\item Introduce one polarization variable \(y_e\) for each local
      outer-tree edge and one type variable \(z_A\) for each root
      support.
\item Write the complete squarefree
      \(\prod_ey_e\)-coefficient before applying the Fr\'echet
      norm, so the partition into edge sets is generated rather
      than counted.
\item Prove ASCC first when a support is repeated twice; retain the
      two distinguished differentiated tree edges in the
      normalized gate.
\item Treat the repeated pair-support row separately, because its
      normalized gate is affine and two nonempty derivatives of
      one copy vanish.
\item If the polarized identity produces contact terms, assemble
      them with the drift derivative before assigning anchors, in
      accordance with the V85 audit.
\end{enumerate}
\end{programme}

\begin{firewall}[Claim boundary after V86]
\label{fire:newV86-boundary}
V86 proves exact multiplicity coarsening for exponentially
normalized supports, proves that both root proxies survive this
coarsening, proves the tree-valued full-label occurrence theorem,
adds labeled ports, and derives the noncommutative assembled-type
factorial.  It closes the distinct-support and already-assembled
support sectors and reduces the remaining local problem to ASCC.

V86 does not prove the exposed-anchor or two-proxy rows for either
of the preceding sectors.  It does not prove ASCC, TLRC in its
original fixed-core form,
general PAAC, CPAC, AABEC, TBRC, full MCWC, OCRC, SCNGC, WPSC,
GRLTC, LZTC, PLRC, WSDC, BRSC, HWSFC, UGCGC, thermal connected
WPRAC, all-arity RGSC, the raw Wilson aggregate strong card,
all-order strong MTA, WUFC, CPM, cutoff removal, reflection
positivity, Osterwalder--Schrader reconstruction, a continuum
Yang--Mills measure, or a positive Hamiltonian gap.  The full
Yang--Mills mass gap has not been proved.
\end{firewall}

\begin{theorem}[V86 result ledger]
\label{thm:newV86-results}
V86:
\begin{enumerate}[label=\textup{(V86.\arabic*)},leftmargin=6.5em]
\item returns to root occurrences on the full thermal label set;
\item proves the exact exponential multiplicity coarsening;
\item proves inheritance of the reserve and strong proxies;
\item proves the normalized full-label tree-valued theorem;
\item proves its ported version;
\item proves the noncommutative assembled-type factorial;
\item closes the multiplicity subproblem for distinct and
      already-assembled support sectors; and
\item formulates ASCC and proves its sufficiency for PAAC.
\end{enumerate}
\end{theorem}

\begin{proof}
Use
\Cref{lem:newV86-exact-coarsening,
lem:newV86-proxy-inheritance,
thm:newV86-full-label,
cor:newV86-ported,
lem:newV86-type-factorial,
cor:newV86-assembled,
prop:newV86-ASCC,
cor:newV86-closed-sectors}.
\end{proof}

\paragraph{Parent-version record.}
V86 was formed from
\texttt{Renewal\_Yang\_Mills\_Mass\_Gap\_Research\_Notes\_V85.tex}.
The V85 parent had \(38\,707\) logical source lines,
\(1\,242\,836\) bytes, and SHA--256
\[
 \texttt{4d46962277ff998d08b99f8c6e3686cfb03e80e32c22d082524c9b6f955deab7}.
\]
The next compulsory companion audit is V90.

% =============================================================================
% END APPENDED FIFTH-SERIES V86 MATERIAL
% =============================================================================

\clearpage
\chapter{Fifth-series V87: polarized support sources and
weighted-tree gate coefficients}
\label{ch:newV87-polarization}

\section{Two independent polarizations}
\label{sec:newV87-two-polarizations}

No companion audit is due at V87.

Fix a finite outer-edge set \(F_0\).  Give every \(e\in F_0\) a
squarefree polarization variable \(y_e\), and give every root
support \(A\) a type variable \(z_A\).  These variables serve
different purposes:
\[
 y_e\quad\hbox{records ownership of the edge derivative},\qquad
 z_A\quad\hbox{records multiplicity of the support}.
\]
They must not be identified before the relevant coefficients have
been extracted.

\begin{lemma}[Doubly polarized Duhamel majorant]
\label{lem:newV87-double-polarization}
Let \(V=V^*\preceq\varepsilon I\), and for each type
\(\alpha\) let
\[
 H_\alpha(\boldsymbol y)
 =
 \sum_{E\subseteq F_0}H_{\alpha,E}\boldsymbol y^E
\]
be an operator polynomial.  Suppose nonnegative scalar
coefficients \(h_{\alpha,E}\) satisfy
\(\|H_{\alpha,E}\|_{\rm cb}\le h_{\alpha,E}\), and put
\[
 h_\alpha(\boldsymbol y)
 :=
 \sum_{E\subseteq F_0}h_{\alpha,E}\boldsymbol y^E.
\]
For every \(F\subseteq F_0\) and multiplicity vector
\(\boldsymbol\nu\), one has the coefficientwise bound
\begin{align}
 &\left\|
 [\boldsymbol y^F\boldsymbol z^{\boldsymbol\nu}]
 \exp\left\{
 V+\sum_\alpha z_\alpha H_\alpha(\boldsymbol y)
 \right\}
 \right\|_{\rm cb}
 \notag\\
 &\qquad\le
 e^\varepsilon
 [\boldsymbol y^F]
 \prod_\alpha
 {h_\alpha(\boldsymbol y)^{\nu_\alpha}\over\nu_\alpha!}.
                                                               \label{eq:newV87-double-polarization}
\end{align}
Only squarefree \(\boldsymbol y\)-coefficients are asserted.
\end{lemma}

\begin{proof}
Fix \(q=\sum_\alpha\nu_\alpha\).  Expand the coefficient in the
type variables by the ordered-simplex formula.  As in
\Cref{lem:newV86-type-factorial}, the words with the prescribed
type multiplicities and the simplex volume give
\(\prod_\alpha1/\nu_\alpha!\) after norms.

For each word, extract \(\boldsymbol y^F\).  The product rule
assigns disjoint subsets of \(F\) to the \(q\) ordered factors.
Replacing each operator coefficient by \(h_{\alpha,E}\) gives
exactly the coefficient of the scalar product on the right.
Positivity permits the sums to be interchanged.
\end{proof}

\begin{corollary}[The support factorial precedes edge extraction]
\label{cor:newV87-order}
For the exact BKAR source, put the undifferentiated source at
\(\boldsymbol y=0\) into \(V\), and introduce \(z_A\) on the
increment
\begin{equation}
 H_A^\circ(\boldsymbol y)
 :=
 \{\chi_A(\boldsymbol x+\boldsymbol y)
       -\chi_A(\boldsymbol x)\}R_A.                         \label{eq:newV87-gate-increment}
\end{equation}
Do this before extracting \(\prod_{e\in F_0}y_e\).  Then every
repeated support \(A\) retains the factor \(1/\nu_A!\), every
explicit occurrence owns a nonempty edge set, and the
\(y\)-coefficient assigns each outer edge to exactly one Duhamel
block.
\end{corollary}

\begin{proof}
The constant coefficient of
\eqref{eq:newV87-gate-increment} is zero.  Apply
\Cref{lem:newV87-double-polarization} with type \(\alpha=A\).
Consequently no explicit factor receives the empty edge set.  A
squarefree monomial cannot assign one edge to two factors, and the
type factorial is already outside the subsequent
\(\boldsymbol y\)-coefficient.
\end{proof}

\section{Exact coefficients of the normalized tree gate}
\label{sec:newV87-gate}

Fix \(A\), \(a=|A|\ge2\), and recall
\[
 W_A=
 \sum_{T\in\mathfrak T(A)}
 \prod_{i\in A}b_i^{d_T(i)}
 =
 \left(\prod_{i\in A}b_i\right)B_A^{a-2}.
\]
For a forest \(E\) on \(A\), let \({\cal K}(E)\) be its set of
vertex components, including isolated vertices, and put
\[
 k:=|E|,\qquad c:=|{\cal K}(E)|=a-k.
\]

\begin{lemma}[Weighted Cayley trees containing a fixed forest]
\label{lem:newV87-forest-extension}
If \(E\) is a forest on \(A\), then
\begin{align}
 &\sum_{\substack{T\in\mathfrak T(A)\\E\subseteq E(T)}}
 \prod_{i\in A}b_i^{d_T(i)}
 \notag\\
 &\qquad=
 \left(\prod_{\{i,j\}\in E}b_ib_j\right)
 \left(\prod_{C\in{\cal K}(E)}B_C\right)
 B_A^{c-2}.                                                  \label{eq:newV87-forest-extension}
\end{align}
If \(E\) contains a cycle, the left side is zero.
\end{lemma}

\begin{proof}
The cyclic assertion is immediate.  Suppose \(E\) is a forest and
contract each component \(C\in{\cal K}(E)\) to one vertex.  An
edge between two contracted components \(C,D\) may choose arbitrary
endpoints \(i\in C\), \(j\in D\); its complete weight is
\[
 \sum_{i\in C}\sum_{j\in D}b_ib_j=B_CB_D.
\]
Thus the sum over the edges added between components is the
weighted Cayley polynomial on the \(c\) contracted vertices:
\[
 \left(\prod_C B_C\right)B_A^{c-2}.
\]
The already fixed edges contribute
\(\prod_{\{i,j\}\in E}b_ib_j\), proving
\eqref{eq:newV87-forest-extension}.  The same formula covers
\(c=1\): the contracted tree is empty and
\((\prod_CB_C)B_A^{-1}=1\).
\end{proof}

\begin{definition}[Forest gate coefficient]
\label{def:newV87-gamma}
For a forest \(E\) on \(A\), define
\begin{equation}
 \gamma_A(E)
 :=
 {1\over W_A}
 \sum_{\substack{T\in\mathfrak T(A)\\E\subseteq E(T)}}
 \prod_{i\in A}b_i^{d_T(i)}.                               \label{eq:newV87-gamma}
\end{equation}
Set \(\gamma_A(E)=0\) when \(E\) contains a cycle.
\end{definition}

\begin{corollary}[Exact normalized forest coefficient]
\label{cor:newV87-gamma-formula}
For a forest \(E\),
\begin{equation}
 \gamma_A(E)
 =
 {1\over B_A^k}
 {\displaystyle
  \left(\prod_{\{i,j\}\in E}b_ib_j\right)
  \left(\prod_{C\in{\cal K}(E)}B_C\right)
  \over
  \displaystyle\prod_{i\in A}b_i}.                         \label{eq:newV87-gamma-formula}
\end{equation}
Moreover, on every \(0\le x_e\le1\),
\begin{equation}
 \left|
 [\boldsymbol y^E]\,
 \chi_A(\boldsymbol x+\boldsymbol y)
 \right|
 \le\gamma_A(E).                                            \label{eq:newV87-gate-coeff}
\end{equation}
\end{corollary}

\begin{proof}
Divide \eqref{eq:newV87-forest-extension} by \(W_A\) and use
\(c=a-k\).  For \eqref{eq:newV87-gate-coeff}, expand every tree
monomial in \(\chi_A(\boldsymbol x+\boldsymbol y)\).  A surviving
tree contains \(E\); its remaining \(x\)-factors have modulus at
most one.
\end{proof}

\section{A product bound for forest coefficients}
\label{sec:newV87-product}

\begin{lemma}[Component-root expansion]
\label{lem:newV87-component-root}
For every tree component \(C\) of \(E\),
\begin{equation}
 B_C\prod_{v\in C}b_v^{d_{E|C}(v)-1}
 \le
 \prod_{\{i,j\}\in E|C}(b_i+b_j).                           \label{eq:newV87-component-root}
\end{equation}
The left side is interpreted as one when \(C\) is an isolated
vertex.
\end{lemma}

\begin{proof}
For a nontrivial component, expand \(B_C=\sum_{r\in C}b_r\).
For each \(r\), orient every edge of the component toward \(r\).
The resulting term
\[
 b_r\prod_{v\in C}b_v^{d(v)-1}
\]
is the monomial in
\(\prod_{\{i,j\}\in E|C}(b_i+b_j)\) obtained by choosing, on each
edge, the endpoint closer to \(r\).  Different roots give distinct
endpoint-choice monomials.  Their sum is therefore bounded by the
complete positive expansion.  For an isolated vertex both sides
equal one.
\end{proof}

\begin{corollary}[One-edge-probability domination]
\label{cor:newV87-product-gamma}
For every forest \(E\) on \(A\),
\begin{equation}
 \gamma_A(E)
 \le
 \prod_{\{i,j\}\in E}{b_i+b_j\over B_A}.                    \label{eq:newV87-product-gamma}
\end{equation}
In particular, each factor on the right lies in \([0,1]\).
\end{corollary}

\begin{proof}
Regroup the numerator of
\eqref{eq:newV87-gamma-formula} component by component.  Its
component factor after division by the vertex marks is exactly the
left side of \eqref{eq:newV87-component-root}.  Apply that lemma
and retain \(B_A^{-k}\).
\end{proof}

\begin{remark}[Exact one-edge marginal]
\label{rem:newV87-one-edge}
For \(E=\{\{i,j\}\}\), equality holds:
\[
 \gamma_A(E)={b_i+b_j\over B_A}.
\]
Thus the right side of
\eqref{eq:newV87-product-gamma} is the product of the one-edge
inclusion probabilities of the weighted Cayley tree.  No
probabilistic negative-dependence theorem is needed; the
deterministic component-root proof gives precisely the direction
used below.
\end{remark}

\section{Repeated supports collapse coefficientwise}
\label{sec:newV87-repeated}

\begin{lemma}[High-arity root activity is smaller than
\(a^{-a}\)]
\label{lem:newV87-small-activity}
After decreasing \(b_*>0\), every \(a\ge3\), \(B_A\le b_*\), and
activity satisfying the strong proxy obeys
\begin{equation}
 0\le u_A
 \le K^a\left(\prod_{i\in A}b_i\right)B_A^{a-2}
 \le a^{-a}.                                                \label{eq:newV87-small-activity}
\end{equation}
\end{lemma}

\begin{proof}
The arithmetic--geometric mean inequality gives
\[
 u_A\le {K^aB_A^{2a-2}\over a^a}.
\]
Choose \(b_*\) so that
\(K^ab_*^{2a-2}\le1\) for every \(a\ge3\).  It is enough to
enforce this at the finitely many first values and
\(Kb_*^2<1\), after which the left side decreases
geometrically.
\end{proof}

\begin{lemma}[Small-activity Touchard bound]
\label{lem:newV87-Touchard}
Let \(k\ge1\) and \(0\le u\le k^{-k}\).  Then
\begin{equation}
 \sum_{\pi\in{\rm Part}([k])}u^{|\pi|}
 \le2u.                                                      \label{eq:newV87-Touchard}
\end{equation}
\end{lemma}

\begin{proof}
There is one one-block partition.  Every other partition has at
least two blocks, and the total number of set partitions is at
most \(k^k\).  Hence
\[
 \sum_\pi u^{|\pi|}
 \le u+k^ku^2\le2u.
\]
\end{proof}

\begin{theorem}[Coefficientwise repeated-support collapse]
\label{thm:newV87-repeated-collapse}
Let \(A\) have \(a\ge3\) labels, let \(F\) be a nonempty forest on
\(A\), and put \(k=|F|\le a-1\).  Suppose one occurrence of support
\(A\) has activity \(u_A\) satisfying both root proxies.  In the
doubly polarized Duhamel expansion, sum every partition of the
edges of \(F\) among an arbitrary positive number of repeated
copies of that same support.  Its complete positive coefficient is
at most
\begin{equation}
 2u_A
 \prod_{\{i,j\}\in F}{b_i+b_j\over B_A}.                    \label{eq:newV87-repeated-collapse}
\end{equation}
The factor two is independent of \(a,k\), and the multiplicity.
\end{theorem}

\begin{proof}
By \Cref{lem:newV87-double-polarization}, the support multiplicity
is exponentially normalized.  Expanding its squarefree
\(\boldsymbol y^F\)-coefficient gives a set partition
\(\pi\) of \(F\): each block \(E\in\pi\) is the nonempty edge set
assigned to one occurrence.  Every block is a forest.  By
\Cref{cor:newV87-gamma-formula,cor:newV87-product-gamma}, the term for
\(\pi\) is at most
\[
 u_A^{|\pi|}
 \prod_{E\in\pi}
 \prod_{\{i,j\}\in E}{b_i+b_j\over B_A}
 =
 u_A^{|\pi|}
 \prod_{\{i,j\}\in F}{b_i+b_j\over B_A}.
\]
Since \(k\le a-1\),
\(\eqref{eq:newV87-small-activity}\) implies
\[
 u_A\le a^{-a}\le k^{-k}.
\]
Sum \(\pi\) and apply \Cref{lem:newV87-Touchard}.
\end{proof}

\begin{corollary}[A pair support cannot repeat as a differentiated
block]
\label{cor:newV87-pair}
If \(A=\{i,j\}\), at most one Duhamel block with support \(A\) can
own a nonempty outer-edge derivative set.  Its normalized gate is
\(\chi_A=x_{ij}\), and the unique possible derivative is
\(\partial_{ij}\chi_A=1\).
\end{corollary}

\begin{proof}
There is only one edge variable on \(A\), and the edge sets assigned
to distinct Duhamel blocks are nonempty and disjoint.  Hence two
such blocks cannot both occur.  The displayed gate formula follows
directly from \eqref{eq:newV60-gate}.
\end{proof}

\begin{assessment}[ASCC multiplicities are closed]
\label{ass:newV87-ASCC}
\Cref{cor:newV87-order,thm:newV87-repeated-collapse,
cor:newV87-pair} prove ASCC(i)--(ii) without inserting a factorial
after norms.  They also retain, for every owned forest edge, the
explicit endpoint-choice factor
\[
 {b_i+b_j\over B_A}
 =
 {b_i\over B_A}+{b_j\over B_A}.
\]
The support activity itself changes only from \(u_A\) to \(2u_A\),
so both proxies survive after the numerical base is enlarged.
What remains is to convert these normalized endpoint choices into
the unnormalized local port marks required by PAAC without losing
the reserve when \(k\) is close to \(a-1\).
\end{assessment}

\section{The remaining normalized-anchor conversion}
\label{sec:newV87-anchor}

\begin{definition}[Normalized-anchor conversion card, NACC]
\label{def:newV87-NACC}
NACC is an allocation of every factor on the right side of
\eqref{eq:newV87-repeated-collapse} which:
\begin{enumerate}[label=\textup{(\roman*)},leftmargin=3.5em]
\item expands each \(b_i+b_j\) into one chosen endpoint anchor;
\item cancels the resulting powers of \(B_A^{-1}\) only against
      clock reserve or an already retained Cayley-tree factor;
\item preserves one raw mark \(b_{a(\ell)}\) for every external
      local port; and
\item has fibre at most \(K^{a+k}\), uniformly up to
      \(k=a-1\).
\end{enumerate}
\end{definition}

\begin{proposition}[NACC completes ASCC]
\label{prop:newV87-NACC}
NACC, together with the results of V87, proves ASCC and therefore
PAAC.
\end{proposition}

\begin{proof}
The type-source organization and its factorial are
\Cref{cor:newV87-order}.  Repeated supports are reduced to one
two-proxy activity by
\Cref{thm:newV87-repeated-collapse}; pair supports are
\Cref{cor:newV87-pair}.  NACC supplies the missing raw anchors with
an exponential fibre.  These are precisely ASCC(i)--(iv).
Apply \Cref{prop:newV86-ASCC}.
\end{proof}

\begin{assessment}[A natural split for NACC]
\label{ass:newV87-split}
If \(k\le a-2\), then
\[
 {s^{(a-2)/2}\over B_A^k}
 \le
 {s^{(a-2-k)/2}\over a^k},
\]
so every normalized endpoint choice consumes at most one available
clock half-power.  The saturated case \(k=a-1\) cannot use this
inequality: it would create \(s^{-1/2}\).  In that case \(F\) is
already a spanning tree on \(A\), and the strong gate estimate
retains its complete Cayley monomial.  V88 will keep these two rows
separate.
\end{assessment}

\begin{programme}[V88: raw anchors from normalized forest edges]
\label{prog:newV87-V88}
\begin{enumerate}[leftmargin=3em]
\item For \(k\le a-2\), orient every owned forest component and use
      one reserve half-power per selected endpoint denominator.
\item Sum the endpoint orientations before replacing local
      \(B_A\)'s by the block weight.
\item For \(k=a-1\), abandon the reserve proxy and use the exact
      spanning-tree coefficient in
      \Cref{lem:newV87-forest-extension}.
\item Prove that the two rows meet at \(k=a-2\) with one uniform
      exponential fibre.
\item Insert the resulting raw anchors into the ported Cayley
      identity and test the full NACC conclusion.
\end{enumerate}
\end{programme}

\begin{firewall}[Claim boundary after V87]
\label{fire:newV87-boundary}
V87 proves the doubly polarized Duhamel coefficient bound, proves
the exact weighted-Cayley extension formula for a prescribed
forest, bounds every forest gate coefficient by its one-edge
endpoint probabilities, and collapses every repeated high-arity
support coefficientwise to one occurrence with a factor two.  It
also proves that differentiated pair supports never repeat and
closes the multiplicity rows ASCC(i)--(ii).  It formulates NACC as
the remaining joint anchor/reserve operation.

V87 does not prove NACC, ASCC(iii)--(iv), TLRC in its original
fixed-core form, general PAAC, CPAC, AABEC, TBRC, full MCWC, OCRC,
SCNGC, WPSC, GRLTC, LZTC, PLRC, WSDC, BRSC, HWSFC, UGCGC,
thermal connected WPRAC, all-arity RGSC, the raw Wilson aggregate
strong card, all-order strong MTA, WUFC, CPM, cutoff removal,
reflection positivity, Osterwalder--Schrader reconstruction, a
continuum Yang--Mills measure, or a positive Hamiltonian gap.  The
full Yang--Mills mass gap has not been proved.
\end{firewall}

\begin{theorem}[V87 result ledger]
\label{thm:newV87-results}
V87:
\begin{enumerate}[label=\textup{(V87.\arabic*)},leftmargin=6.5em]
\item proves the doubly polarized Duhamel majorant;
\item retains the support factorial before edge extraction;
\item proves the exact fixed-forest Cayley formula;
\item proves product domination by one-edge gate probabilities;
\item proves the small-activity Touchard estimate;
\item collapses every repeated high-arity support;
\item proves differentiated pair supports do not repeat; and
\item formulates NACC and proves its sufficiency for ASCC.
\end{enumerate}
\end{theorem}

\begin{proof}
Use
\Cref{lem:newV87-double-polarization,
cor:newV87-order,
lem:newV87-forest-extension,
cor:newV87-product-gamma,
lem:newV87-Touchard,
thm:newV87-repeated-collapse,
cor:newV87-pair,
prop:newV87-NACC}.
\end{proof}

\paragraph{Parent-version record.}
V87 was formed from
\texttt{Renewal\_Yang\_Mills\_Mass\_Gap\_Research\_Notes\_V86.tex}.
The V86 parent had \(39\,173\) logical source lines,
\(1\,260\,061\) bytes, and SHA--256
\[
 \texttt{7f0fee3fdb55ecee856db5e31133b637c2b0bad14d73c515aaad65cae693c128}.
\]
The next compulsory companion audit is V90.

% =============================================================================
% END APPENDED FIFTH-SERIES V87 MATERIAL
% =============================================================================

\clearpage
\chapter{Fifth-series V88: block-compatible Cayley trees and the
component-coalescence defect}
\label{ch:newV88-compatible}

\section{Homogeneity correction to NACC}
\label{sec:newV88-correction}

No companion audit is due at V88.

The normalized endpoint factor in V87 is a probability and has
thermal degree zero.  It records which endpoint is selected, but
does not itself supply the additional raw mark required by a PAAC
port.  This distinction is forced by scaling.

\begin{lemma}[Scale degree of the V87 coefficient]
\label{lem:newV88-scaling}
Fix positive reference data
\((\widehat s,\widehat b_i)\), and for \(0<\lambda\le1\) put
\[
 s_\lambda=\lambda^2\widehat s,\qquad
 b_{i,\lambda}=\lambda\widehat b_i.
\]
Then the hypotheses \(b_i\ge\sqrt s\) are preserved.  Both root
proxies
\eqref{eq:newV86-r}--\eqref{eq:newV86-q} have degree
\(2a-2\):
\[
 {\mathfrak r}_{\lambda}(A)
 =\lambda^{2a-2}{\mathfrak r}_{1}(A),
 \qquad
 {\mathfrak q}_{\lambda}(A)
 =\lambda^{2a-2}{\mathfrak q}_{1}(A).
                                                               \label{eq:newV88-proxy-scale}
\]
Every gate coefficient \(\gamma_A(E)\), and every factor
\((b_i+b_j)/B_A\), has degree zero.
\end{lemma}

\begin{proof}
The product of the \(a\) marks contributes \(\lambda^a\);
\(s^{(a-2)/2}\) and \(B_A^{a-2}\) each contribute
\(\lambda^{a-2}\).  This proves
\eqref{eq:newV88-proxy-scale}.  Formula
\eqref{eq:newV87-gamma-formula} has equal total mark degree in its
numerator and denominator.  The one-edge quotient is manifestly
scale invariant.
\end{proof}

\begin{proposition}[A normalized endpoint is not a raw port]
\label{prop:newV88-no-NACC}
No bound derived only from
\eqref{eq:newV87-repeated-collapse} can, uniformly in the thermal
scale, replace \(k\ge1\) normalized endpoint factors by \(k\) raw
anchor marks while leaving the full root-tree weight unchanged.
Consequently NACC, as a standalone conversion of the V87
single-support majorant, is retired.
\end{proposition}

\begin{proof}
The nonzero left side of
\eqref{eq:newV87-repeated-collapse} has degree \(2a-2\) under
\Cref{lem:newV88-scaling}.  A full root-tree factor has the same
degree.  Multiplying it by \(k\) additional raw anchors gives
degree \(2a-2+k\).  An asserted uniform inequality from the former
to the latter would have the form
\[
 c\lambda^{2a-2}
 \le C^{a+k}\lambda^{2a-2+k}
\]
for fixed positive reference data.  Division by
\(\lambda^{2a-2}\) and passage to \(\lambda\downarrow0\) is a
contradiction.
\end{proof}

\begin{correction}[The port must be cut from a global tree]
\label{corr:newV88-NACC}
\Cref{prop:newV87-NACC} remains a valid conditional implication,
but the V87 coefficient does not establish its hypothesis.  A
cross-block root edge already carries two endpoint marks; one must
be assigned to each incident block.  It is incorrect to keep the
entire root tree inside one block and then ask a normalized gate
probability to create an additional port mark.  The root tree must
first be cut along the thermal partition.
\end{correction}

\begin{remark}[Pair calibration]
\label{rem:newV88-pair}
For \(A=\{i,j\}\),
\(\partial_{ij}\chi_A=1\), and the pair root has weight
\(Kb_ib_j\).  If \(i\) and \(j\) lie in different blocks, this is
exactly one raw anchor \(b_i\) at the first endpoint and one raw
anchor \(b_j\) at the second.  There is no additional local
pair-tree factor to be multiplied by these anchors.  This
two-label example already forces
\Cref{corr:newV88-NACC}.
\end{remark}

\section{Cutting one global Cayley tree by blocks}
\label{sec:newV88-cut}

Let \(X\) be partitioned into nonempty thermal blocks
\(\mathcal P\).  Write
\[
 r:=|\mathcal P|,
 \qquad
 n_P:=|P|,
 \qquad
 w_P:=\sum_{i\in P}b_i.
\]
For a tree \(S\in\mathfrak T(X)\), let \(S[P]\) be its induced
forest on \(P\), let \(c_P(S)\) be the number of its components,
and let \(E_\times(S)\) be the cross-block edges.

\begin{lemma}[Component quotient of a global tree]
\label{lem:newV88-component-quotient}
Contract every component of every induced forest \(S[P]\) to one
vertex.  Retain every cross-block edge.  The resulting graph
\(\widehat U_S\) is a tree.  In particular,
\begin{equation}
 |E_\times(S)|
 =
 \sum_{P\in\mathcal P}c_P(S)-1.                             \label{eq:newV88-cross-count}
\end{equation}
\end{lemma}

\begin{proof}
Contracting connected subtrees of a tree neither disconnects it nor
creates a cycle.  The contracted graph is therefore a tree.  Its
vertex set consists of the
\(\sum_Pc_P(S)\) within-block components, giving
\eqref{eq:newV88-cross-count}.
\end{proof}

\begin{lemma}[Exact tree-weight cut]
\label{lem:newV88-weight-cut}
For every \(S\in\mathfrak T(X)\),
\begin{equation}
 \prod_{i\in X}b_i^{d_S(i)}
 =
 \left\{
 \prod_{P\in\mathcal P}
 \prod_{i\in P}b_i^{d_{S[P]}(i)}
 \right\}
 \left\{
 \prod_{\{i,j\}\in E_\times(S)}b_ib_j
 \right\}.                                                  \label{eq:newV88-weight-cut}
\end{equation}
Thus every cross-block edge exposes one raw endpoint anchor in
each of its two blocks.
\end{lemma}

\begin{proof}
At each vertex \(i\), split its degree in \(S\) into its
within-block degree and its number of incident cross-block edges.
The two factors on the right reproduce the corresponding powers of
\(b_i\).
\end{proof}

\begin{definition}[Block-compatible tree]
\label{def:newV88-compatible}
A global tree \(S\) is block-compatible when every \(S[P]\) is
connected, equivalently \(c_P(S)=1\) for every block.
\end{definition}

\begin{corollary}[Compatible quotient is a block tree]
\label{cor:newV88-compatible-quotient}
If \(S\) is block-compatible, contraction of each whole block gives
an ordinary tree \(U_S\) on \(\mathcal P\).  There is at most one
edge of \(S\) between any fixed pair of blocks.
\end{corollary}

\begin{proof}
Now \(\widehat U_S\) in
\Cref{lem:newV88-component-quotient} has one vertex per block, so
it is \(U_S\).  Parallel edges would be a two-cycle in the
contracted multigraph, impossible because \(\widehat U_S\) is a
tree.
\end{proof}

\section{Exact hierarchical Cayley identity}
\label{sec:newV88-hierarchical}

Let
\[
 \mathfrak T_{\mathcal P}^{\rm comp}(X)
 :=
 \{S\in\mathfrak T(X):S\hbox{ is block-compatible}\}.
\]

\begin{theorem}[Bijection for compatible global trees]
\label{thm:newV88-compatible-bijection}
Block-compatible trees are in bijection with data
\begin{equation}
 \left(
 U,\ (L_P)_{P\in\mathcal P},\
 (a_{P,e})_{\substack{P\in\mathcal P\\e\in D_P(U)}}
 \right),                                                   \label{eq:newV88-compatible-data}
\end{equation}
where:
\begin{enumerate}[label=\textup{(\roman*)},leftmargin=3.5em]
\item \(U\) is a tree on the block set;
\item \(L_P\) is a tree on \(P\); and
\item \(a_{P,e}\in P\) is the endpoint in \(P\) of the global
      edge lying over \(e\in E(U)\).
\end{enumerate}
\end{theorem}

\begin{proof}
Given a compatible \(S\), use
\Cref{cor:newV88-compatible-quotient}, set \(L_P=S[P]\), and
record the two endpoints of every cross-block edge.  This gives
\eqref{eq:newV88-compatible-data}.

Conversely, take the disjoint union of the local trees \(L_P\) and,
for each \(e=\{P,Q\}\in E(U)\), add the edge
\(\{a_{P,e},a_{Q,e}\}\).  The result is connected.  Its edge count
is
\[
 \sum_P(n_P-1)+(r-1)=|X|-1,
\]
so it is a tree.  Its induced graph on every block is \(L_P\).
The two constructions are inverse.
\end{proof}

\begin{theorem}[Exact hierarchical weighted Cayley sum]
\label{thm:newV88-hierarchical-Cayley}
With \(v_P\) as in \eqref{eq:newV83-vertex-weight},
\begin{align}
 &\sum_{S\in\mathfrak T_{\mathcal P}^{\rm comp}(X)}
 \prod_{i\in X}b_i^{d_S(i)}
 \notag\\
 &\qquad=
 \sum_{U\in\mathfrak T(\mathcal P)}
 \prod_{P\in\mathcal P}
 \left\{
  \sum_{L_P\in\mathfrak T(P)}
  \sum_{a_P:D_P(U)\to P}
  \prod_{i\in P}b_i^{d_{L_P}(i)}
  \prod_{e\in D_P(U)}b_{a_P(e)}
 \right\}
 \notag\\
 &\qquad=
 \left(\prod_{P\in\mathcal P}v_P\right)
 \left(\prod_{P\in\mathcal P}w_P\right)
 B_X^{r-2}.                                                  \label{eq:newV88-hierarchical-Cayley}
\end{align}
\end{theorem}

\begin{proof}
The first equality is the bijection
\Cref{thm:newV88-compatible-bijection} together with the exact
weight cut \eqref{eq:newV88-weight-cut}.  For fixed \(U\), apply
the ported Cayley identity
\Cref{thm:newV84-ported-Cayley} in every block.  This gives
\[
 \prod_Pv_Pw_P^{d_U(P)}.
\]
The weighted Cayley identity on the block set then gives
\[
 \sum_U\prod_Pw_P^{d_U(P)}
 =
 \left(\prod_Pw_P\right)B_X^{r-2}.
\]
\end{proof}

\begin{corollary}[Compatible-tree PAAC sector]
\label{cor:newV88-compatible-PAAC}
Every operator ancestry which is allocated, with fibre at most
\(K^{|X|}\), to block-compatible global Cayley trees with their
tree monomials satisfies the PAAC weight and hence AABEC.
\end{corollary}

\begin{proof}
For a fixed compatible tree, the data in
\Cref{thm:newV88-compatible-bijection} are exactly the local trees,
raw anchors, and output block tree required by
\Cref{def:newV84-PAAC}.  Sum them by
\Cref{thm:newV88-hierarchical-Cayley} and invoke
\Cref{thm:newV84-PAAC}.
\end{proof}

\section{Why the complete full-label tree sum is insufficient}
\label{sec:newV88-insufficient}

\begin{proposition}[All trees cannot be replaced by compatible
trees uniformly]
\label{prop:newV88-ratio}
The ratio of the complete weighted Cayley sum to the compatible
sum is
\begin{equation}
 {B_X^{|X|-2}\prod_{i\in X}b_i
  \over
  \left(\prod_Pv_P\right)\left(\prod_Pw_P\right)B_X^{r-2}}
 =
 {B_X^{|X|-r}\over
  \prod_{P\in\mathcal P}w_P^{n_P-1}}.                       \label{eq:newV88-ratio}
\end{equation}
It is not bounded by \(K^{|X|}\) uniformly over the thermal data,
even for three labels.
\end{proposition}

\begin{proof}
The numerator is the weighted Cayley--Pr\"ufer identity and the
denominator is
\Cref{thm:newV88-hierarchical-Cayley}; cancellation gives
\eqref{eq:newV88-ratio}.

For the explicit three-label example, take
\[
 P=\{1,2\},\qquad Q=\{3\},\qquad
 b_1=b_2=\epsilon,\qquad b_3=a_0
\]
with \(a_0>0\) fixed and \(\epsilon\downarrow0\), scaling \(s\) so
that the thermal floor holds.  Then \(r=2\), and the ratio is
\[
 {2\epsilon+a_0\over2\epsilon},
\]
which is unbounded.  This is allowed for the possible final
incomplete light block in
\Cref{lem:newV80-blocks}.
\end{proof}

\begin{correction}[Scope of the V86 full-label theorem]
\label{corr:newV88-full-label}
\Cref{thm:newV86-full-label} is a valid tree-valued occurrence
theorem, but it does not by itself imply PAAC after a nontrivial
thermal partition.  It sums all global Cayley trees, whereas PAAC
sums precisely the block-compatible family in
\Cref{thm:newV88-hierarchical-Cayley}.  The ratio
\eqref{eq:newV88-ratio} forbids discarding this distinction.
\end{correction}

\section{The compatibility defect is an exact cycle excess}
\label{sec:newV88-defect}

Identify all component vertices of
\(\widehat U_S\) which belong to the same block, retaining parallel
cross edges.  Denote the resulting connected block multigraph by
\(G_S\).

\begin{theorem}[Defect--excess identity]
\label{thm:newV88-defect-excess}
The cyclomatic excess of \(G_S\) is
\begin{equation}
 g(G_S)
 :=
 |E(G_S)|-|\mathcal P|+1
 =
 \sum_{P\in\mathcal P}\{c_P(S)-1\}.                         \label{eq:newV88-defect-excess}
\end{equation}
Thus every failure of within-block connectedness creates exactly
one quotient-cycle unit, and there are no other units.
\end{theorem}

\begin{proof}
By \eqref{eq:newV88-cross-count},
\[
 |E(G_S)|=|E_\times(S)|
 =\sum_Pc_P(S)-1.
\]
Since \(G_S\) is connected on \(r\) block vertices, subtract
\(r-1\) to obtain
\eqref{eq:newV88-defect-excess}.
\end{proof}

\begin{assessment}[Combinatorial exchange behind the defect]
\label{ass:newV88-exchange}
Choose a spanning tree \(U\) of \(G_S\).  Exactly
\(g(G_S)\) cross edges are then surplus.  Inside each block \(P\),
exactly \(c_P(S)-1\) new edges are needed to connect the components
of \(S[P]\).  The totals agree by
\eqref{eq:newV88-defect-excess}.  After adding those internal edges
and deleting the surplus cross edges, one obtains local trees and a
block tree.  This is a combinatorial basis exchange.  It is not yet
a weighted estimate: comparing a deleted cross-edge monomial with
an added internal-edge monomial termwise is invalid for
heterogeneous marks.
\end{assessment}

\begin{definition}[Block-compatibility allocation card, BCAC]
\label{def:newV88-BCAC}
BCAC is an aggregate allocation of the exact conditioned
percolation ancestry to block-compatible global Cayley trees such
that:
\begin{enumerate}[label=\textup{(\roman*)},leftmargin=3.5em]
\item the component quotient \(\widehat U_S\) and the normalized
      Bernoulli configuration are retained before the allocation;
\item a Kruskal tree of \(G_S\) selects the surplus cross edges;
\item the squarefree background source coalesces the
      \(c_P(S)\) components inside each block before norms; and
\item all choices and operator orderings have total fibre at most
      \(K^{|X|}\).
\end{enumerate}
\end{definition}

\begin{proposition}[BCAC is the corrected local-to-block bridge]
\label{prop:newV88-BCAC}
BCAC implies PAAC, CPAC, AABEC, TBRC, full MCWC, and SCNGC.
\end{proposition}

\begin{proof}
BCAC allocates every ancestry to the compatible sector of
\Cref{cor:newV88-compatible-PAAC}.  The latter proves PAAC.
Apply \Cref{thm:newV84-PAAC}.
\end{proof}

\begin{assessment}[Position after the correction]
\label{ass:newV88-position}
The repeated-support factorial and gate coefficient results of V87
remain valid.  They solve multiplicity before the thermal partition.
The missing thermal powers cannot be supplied by normalized
endpoint probabilities.  They must arise when each global root
tree is cut across blocks and combined with the background roots
which coalesce its within-block components.  The exact number of
required coalescences is the quotient excess
\eqref{eq:newV88-defect-excess}.
\end{assessment}

\begin{programme}[V89: weighted component coalescence]
\label{prog:newV88-V89}
\begin{enumerate}[leftmargin=3em]
\item Condition on one global tree \(S\), its component tree
      \(\widehat U_S\), and the full Bernoulli shortcut
      configuration.
\item Apply the V83 Kruskal identity to \(G_S\) before comparing
      any edge weights.
\item For each surplus cross edge, expose the exact squarefree
      background root which joins the corresponding two components
      inside a repeated block.
\item Sum all internal coalescing edges with the rooted-forest
      formula \Cref{lem:newV87-forest-extension}, rather than
      choosing endpoints without their weights.
\item Prove BCAC first for \(g(G_S)=1\); use that row to determine
      whether the general excess exponentiates.
\end{enumerate}
\end{programme}

\begin{firewall}[Claim boundary after V88]
\label{fire:newV88-boundary}
V88 proves the homogeneity obstruction to standalone NACC, corrects
the location from which raw port marks must arise, proves the exact
component quotient and tree-weight cut, proves the compatible-tree
bijection and hierarchical weighted Cayley identity, closes the
block-compatible PAAC sector, and identifies incompatibility
exactly with quotient cycle excess.  It formulates BCAC as the
corrected local-to-block bridge.

V88 does not prove BCAC, weighted component coalescence, PAAC in
general, CPAC, AABEC, TBRC, full MCWC, OCRC, SCNGC, WPSC, GRLTC,
LZTC, PLRC, WSDC, BRSC, HWSFC, UGCGC, thermal connected WPRAC,
all-arity RGSC, the raw Wilson aggregate strong card, all-order
strong MTA, WUFC, CPM, cutoff removal, reflection positivity,
Osterwalder--Schrader reconstruction, a continuum Yang--Mills
measure, or a positive Hamiltonian gap.  The full Yang--Mills mass
gap has not been proved.
\end{firewall}

\begin{theorem}[V88 result ledger]
\label{thm:newV88-results}
V88:
\begin{enumerate}[label=\textup{(V88.\arabic*)},leftmargin=6.5em]
\item proves the scale degree of root and gate factors;
\item proves normalized endpoints cannot create raw ports;
\item proves the exact component quotient of a global tree;
\item proves the exact tree-weight cut;
\item proves the compatible-tree data bijection;
\item proves the hierarchical weighted Cayley identity;
\item closes the compatible-tree PAAC sector; and
\item proves the compatibility-defect/quotient-excess identity.
\end{enumerate}
\end{theorem}

\begin{proof}
Use
\Cref{lem:newV88-scaling,
prop:newV88-no-NACC,
lem:newV88-component-quotient,
lem:newV88-weight-cut,
thm:newV88-compatible-bijection,
thm:newV88-hierarchical-Cayley,
cor:newV88-compatible-PAAC,
thm:newV88-defect-excess}.
\end{proof}

\paragraph{Parent-version record.}
V88 was formed from
\texttt{Renewal\_Yang\_Mills\_Mass\_Gap\_Research\_Notes\_V87.tex}.
The V87 parent had \(39\,681\) logical source lines,
\(1\,277\,244\) bytes, and SHA--256
\[
 \texttt{ac1148fa6b3fab623fcfb66641b2c67ccb1d9a78fdf9d2ccdcb90d5cb71a8de5}.
\]
The next compulsory companion audit is V90.

% =============================================================================
% END APPENDED FIFTH-SERIES V88 MATERIAL
% =============================================================================

\clearpage
\chapter{Fifth-series V89: one-defect weighted component
coalescence}
\label{ch:newV89-one-defect}

\section{A crossing edge from a background Cayley tree}
\label{sec:newV89-crossing}

No companion audit is due at V89.

Let \(P\) be a light nonsingleton block.  Thus
\[
 w_P=\sum_{i\in P}b_i\le b_*/2<1.
\]
Fix a bipartition \(P=C_1\mathbin{\dot\cup}C_2\) into two nonempty
sets, and fix a total order on the edges of the complete graph on
\(P\).

\begin{lemma}[Tree allocation to the first crossing edge]
\label{lem:newV89-first-crossing}
For each \(L\in\mathfrak T(P)\), let \(e(L)\) be the first edge of
\(L\) crossing \(C_1|C_2\).  Then, for every
\(e=\{i,j\}\) with \(i\in C_1\), \(j\in C_2\),
\begin{equation}
 \sum_{\substack{L\in\mathfrak T(P)\\e(L)=e}}
 \prod_{v\in P}b_v^{d_L(v)}
 \le b_ib_j.                                                \label{eq:newV89-first-crossing}
\end{equation}
Consequently,
\begin{equation}
 \sum_{L\in\mathfrak T(P)}
 \prod_{v\in P}b_v^{d_L(v)}
 \le
 \sum_{\substack{i\in C_1\\j\in C_2}}b_ib_j
 =
 B_{C_1}B_{C_2}.                                            \label{eq:newV89-cut-sum}
\end{equation}
\end{lemma}

\begin{proof}
Every spanning tree crosses the nontrivial cut, so \(e(L)\) is
defined.  Enlarge the family assigned to \(e\) to all trees
containing \(e\).  By
\Cref{lem:newV87-forest-extension} with one prescribed edge,
\begin{align*}
 \sum_{\substack{L\in\mathfrak T(P)\\e\in E(L)}}
 \prod_{v\in P}b_v^{d_L(v)}
 &=
 b_ib_j(b_i+b_j)
 \left(\prod_{v\in P\setminus\{i,j\}}b_v\right)
 w_P^{|P|-3}.
\end{align*}
If \(|P|=2\), the factor after \(b_ib_j\) equals one under the
contracted-tree convention.  If \(|P|\ge3\), then
\[
 b_i+b_j\le w_P<1,\qquad
 \prod_{v\ne i,j}b_v\le1,\qquad
 w_P^{|P|-3}\le1.
\]
This proves \eqref{eq:newV89-first-crossing}.  Sum it over the
crossing edges to obtain \eqref{eq:newV89-cut-sum}.
\end{proof}

\begin{assessment}[Why no endpoint census appears]
\label{ass:newV89-crossing}
The chosen edge retains its exact weight \(b_ib_j\).  All remaining
tree edges are summed before being discarded, through
\Cref{lem:newV87-forest-extension}.  Counting possible trees after
dropping their weights would recreate a Cayley factor; the
allocation in \Cref{lem:newV89-first-crossing} avoids it.
\end{assessment}

\section{The one-defect exchange}
\label{sec:newV89-exchange}

Let \(S\in\mathfrak T(X)\) have compatibility defect one:
\[
 \sum_{Q\in\mathcal P}\{c_Q(S)-1\}=1.                       \label{eq:newV89-one-defect}
\]
There is then a unique block \(P\) such that \(S[P]\) has two
components, say \(C_1,C_2\); every other induced block graph is a
tree.

\begin{lemma}[Unique quotient cycle]
\label{lem:newV89-cycle}
The block multigraph \(G_S\) is unicyclic.  Its unique cycle is the
image of the path in the component tree \(\widehat U_S\) joining
the component vertices \(C_1\) and \(C_2\).
\end{lemma}

\begin{proof}
By \Cref{thm:newV88-defect-excess}, \(G_S\) is connected with
cyclomatic excess one, hence unicyclic.  The component quotient
\(\widehat U_S\) is a tree by
\Cref{lem:newV88-component-quotient}.  Identifying its two vertices
\(C_1,C_2\) creates exactly the cycle given by their unique path.
\end{proof}

Fix a deterministic rule selecting one cross edge \(f(S)\) on this
cycle, for example the first in a fixed edge order.

\begin{lemma}[Component exchange produces a compatible tree]
\label{lem:newV89-exchange}
Let \(e=\{i,j\}\) with \(i\in C_1\), \(j\in C_2\), and put
\begin{equation}
 S':=S-f(S)+e.                                               \label{eq:newV89-exchanged-tree}
\end{equation}
Then \(S'\) is a block-compatible global tree.  Moreover
\begin{equation}
 \left(\prod_{v\in X}b_v^{d_S(v)}\right)b_ib_j
 =
 \left(\prod_{v\in X}b_v^{d_{S'}(v)}\right)
 b_{f_-}b_{f_+}
 \le
 \prod_{v\in X}b_v^{d_{S'}(v)},                             \label{eq:newV89-weight-exchange}
\end{equation}
where \(f(S)=\{f_-,f_+\}\).
\end{lemma}

\begin{proof}
Adding \(e\) to the tree \(S\) creates the cycle consisting of
\(e\) and the path from \(C_1\) to \(C_2\) in
\(\widehat U_S\).  The edge \(f(S)\) lies on that path.  Deleting
it therefore leaves a tree.

The new edge joins the two components of \(S[P]\), while the
deleted edge is cross-block.  Thus \(S'[P]\) is connected, and no
other induced block tree was changed.  Hence \(S'\) is compatible.

The first equality in \eqref{eq:newV89-weight-exchange} follows by
cancelling all unchanged edge factors.  Every thermal mark is
strictly less than one, so \(b_{f_-}b_{f_+}\le1\), proving the
inequality.
\end{proof}

\begin{lemma}[The one-defect exchange fibre is exponential]
\label{lem:newV89-fibre}
The map from a valid pair \((S,e)\) to the compatible tree \(S'\)
in \eqref{eq:newV89-exchanged-tree} has fibre at most
\(|X|^4\), hence at most \(K^{|X|}\) for one numerical \(K\).
\end{lemma}

\begin{proof}
Given \(S'\), a preimage is determined by:
\begin{enumerate}[label=\textup{(\roman*)},leftmargin=3.5em]
\item the added within-block edge \(e\in E(S')\); and
\item the deleted cross-block edge \(f\) on \(X\).
\end{enumerate}
There are fewer than \(|X|^2\) choices for each.  They determine
\(S=S'-e+f\), after which validity and the deterministic
cycle-selection rule can only reduce the count.  Finally
\(|X|^4\le16^{|X|}\) for \(|X|\ge2\).
\end{proof}

\section{A fresh background tree closes one defect}
\label{sec:newV89-fresh}

\begin{definition}[Fresh tree-valued background factor]
\label{def:newV89-fresh}
For a fixed differentiated ancestry tree \(S\) and its bad block
\(P\), a fresh tree-valued background factor is a disjoint
squarefree-background contribution whose complete positive
majorant, before norms and before choosing an edge, is
\begin{equation}
 K^{|P|}
 \sum_{L\in\mathfrak T(P)}
 \prod_{i\in P}b_i^{d_L(i)}.                               \label{eq:newV89-fresh}
\end{equation}
Fresh means that none of its coefficient labels or root
occurrences has already been used to construct \(S\).
\end{definition}

\begin{theorem}[Fresh-background one-defect BCAC]
\label{thm:newV89-one-defect-BCAC}
Every defect-one ancestry which carries a fresh tree-valued
background factor in its unique bad block satisfies BCAC with one
exponential fibre.
\end{theorem}

\begin{proof}
Fix \(S\), its bad block \(P=C_1\dot\cup C_2\), and the remaining
conditioned data.  Apply
\Cref{lem:newV89-first-crossing} to the fresh background tree sum.
Allocate it to a crossing edge
\(e=\{i,j\}\), retaining the factor \(b_ib_j\).

Use \Cref{lem:newV89-exchange}.  The product of the ancestry
tree weight and the retained crossing-edge weight is bounded by
the compatible-tree weight of \(S'\).  The fibre is exponential by
\Cref{lem:newV89-fibre}; the background allocation itself was
performed by a positive weighted sum and introduces no tree
census.  All constants enter \(K^{|X|}\).  This is BCAC.
\end{proof}

\begin{corollary}[One-defect PAAC sector]
\label{cor:newV89-one-defect-PAAC}
The fresh-background defect-one sector satisfies PAAC, CPAC,
AABEC, TBRC, full MCWC, and SCNGC.
\end{corollary}

\begin{proof}
Use
\Cref{thm:newV89-one-defect-BCAC,prop:newV88-BCAC}.
\end{proof}

\begin{warning}[Freshness is not automatic]
\label{warn:newV89-freshness}
\Cref{thm:newV85-tree-CIRC} produces a tree-valued background
majorant after a connected centre source has been isolated.  It
does not prove that this source is disjoint from the differentiated
root occurrences already used to construct \(S\).  Reusing the
same occurrence in both factors would square its thermal payment
and is invalid.  Thus
\Cref{thm:newV89-one-defect-BCAC} is a proved sector theorem, not a
proof of the complete defect-one ancestry.
\end{warning}

\section{The exact fresh-coalescence card}
\label{sec:newV89-card}

\begin{definition}[Fresh component-coalescence card, FCCC]
\label{def:newV89-FCCC}
FCCC is a coefficientwise split of the exact conditioned
squarefree logarithm such that, for every incompatible global tree
\(S\):
\begin{enumerate}[label=\textup{(\roman*)},leftmargin=3.5em]
\item the root occurrences allocated to \(S\) and the occurrences
      allocated to component coalescence are disjoint;
\item in each block \(P\), the latter occurrences supply a
      tree-valued factor connecting the \(c_P(S)\) components of
      \(S[P]\);
\item every coalescing edge is retained with its two endpoint
      marks before norms; and
\item all splits, edge choices, and Duhamel orders have aggregate
      fibre at most \(K^{|X|}\).
\end{enumerate}
\end{definition}

\begin{proposition}[FCCC completes BCAC]
\label{prop:newV89-FCCC}
FCCC implies BCAC and therefore PAAC, CPAC, AABEC, TBRC, full
MCWC, and SCNGC.
\end{proposition}

\begin{proof}
For each block \(P\), use the supplied tree-valued factor to select
\(c_P(S)-1\) edges which connect the components of \(S[P]\).
The total number of added edges is
\[
 \sum_P\{c_P(S)-1\}=g(G_S)
\]
by \Cref{thm:newV88-defect-excess}.  Choose a spanning tree of
\(G_S\) by the normalized V83 Kruskal allocation and delete the
same number of surplus cross edges.  The resulting global tree is
block-compatible.

For every exchange, the product of the added edge weight and the
old global-tree weight equals the new global-tree weight times the
deleted cross-edge weight, as in
\eqref{eq:newV89-weight-exchange}.  The latter is at most one.
FCCC(iv), the Kruskal identity, and the normalized Bernoulli
configuration prevent a superexponential inverse fibre.  Hence
BCAC holds.  Apply \Cref{prop:newV88-BCAC}.
\end{proof}

\begin{assessment}[What remains after the one-cycle theorem]
\label{ass:newV89-position}
The weighted exchange itself is no longer an obstruction.  A fresh
within-block crossing edge dominates the deleted cross edge without
any comparison of heterogeneous endpoint weights, because the
deleted factor is at most one.  The live issue is squarefree
freshness: the exact logarithm must allocate disjoint root
occurrences simultaneously to the global differentiated tree and
to the component-coalescing background forest.
\end{assessment}

\begin{programme}[V90: compulsory audit and fresh component
forests]
\label{prog:newV89-V90}
\begin{enumerate}[leftmargin=3em]
\item Perform the compulsory newest-SM/newest-NS audit and record
      immutable snapshots.
\item Return to the exact squarefree logarithmic coefficient before
      the V78 positive majorant and color every root occurrence
      either differentiated or background.
\item Prove the two-color coefficient identity without duplicating
      an occurrence.
\item For one bad block with two components, test whether connected
      cancellation forces a background root path across their cut.
\item If it does not, enlarge the exchange object so a mixed
      differentiated/background path, rather than a fresh
      background tree, supplies the coalescing edge.
\end{enumerate}
\end{programme}

\begin{firewall}[Claim boundary after V89]
\label{fire:newV89-boundary}
V89 proves a weighted allocation of background Cayley trees to cut
edges, proves the exact one-defect component exchange and its
exponential fibre, and closes the entire defect-one sector whenever
one fresh tree-valued background factor is present.  It formulates
FCCC and proves its sufficiency for BCAC.

V89 does not prove freshness for the exact squarefree logarithm,
FCCC, general BCAC, general PAAC, CPAC, AABEC, TBRC, full MCWC,
OCRC, SCNGC, WPSC, GRLTC, LZTC, PLRC, WSDC, BRSC, HWSFC, UGCGC,
thermal connected WPRAC, all-arity RGSC, the raw Wilson aggregate
strong card, all-order strong MTA, WUFC, CPM, cutoff removal,
reflection positivity, Osterwalder--Schrader reconstruction, a
continuum Yang--Mills measure, or a positive Hamiltonian gap.  The
full Yang--Mills mass gap has not been proved.
\end{firewall}

\begin{theorem}[V89 result ledger]
\label{thm:newV89-results}
V89:
\begin{enumerate}[label=\textup{(V89.\arabic*)},leftmargin=6.5em]
\item allocates a background Cayley sum to one cut edge;
\item proves the unique-cycle description at defect one;
\item proves the weighted component exchange;
\item proves its inverse fibre is exponential;
\item closes fresh-background defect-one BCAC and PAAC;
\item proves why freshness may not be inferred;
\item formulates FCCC; and
\item proves FCCC suffices for the complete local chain.
\end{enumerate}
\end{theorem}

\begin{proof}
Use
\Cref{lem:newV89-first-crossing,
lem:newV89-cycle,
lem:newV89-exchange,
lem:newV89-fibre,
thm:newV89-one-defect-BCAC,
cor:newV89-one-defect-PAAC,
warn:newV89-freshness,
prop:newV89-FCCC}.
\end{proof}

\paragraph{Parent-version record.}
V89 was formed from
\texttt{Renewal\_Yang\_Mills\_Mass\_Gap\_Research\_Notes\_V88.tex}.
The V88 parent had \(40\,180\) logical source lines,
\(1\,294\,372\) bytes, and SHA--256
\[
 \texttt{a6c27830e1be1ec89a6ae0b280ab7466d3f9d0beb071131e81626444740512af}.
\]
The next compulsory companion audit is V90.

% =============================================================================
% END APPENDED FIFTH-SERIES V89 MATERIAL
% =============================================================================

\clearpage
\chapter{Fifth-series V90: compulsory audit and the
local--transversal split}
\label{ch:newV90-transversal}

\section{Compulsory companion audit}
\label{sec:newV90-audit}

\begin{audit}[V90 immutable companion snapshots]
\label{audit:newV90-companions}
The newest active companion sources were read without modification:
\begin{align*}
 &\texttt{Renewal\_SM\_Einstein\_Continuum\_Research\_Notes\_V16.tex},\\
 &\qquad 6\,989\ \hbox{logical lines},\quad
 274\,001\ \hbox{bytes},\\
 &\qquad
 \texttt{20a43ab5feecb2df6a27575fc95b7b6ab0c6bc765ad6ff3c77ac2ba8b3bbb4b7},
 \\
 &\texttt{Renewal\_NS\_Stability\_Research\_Notes\_V71.tex},\\
 &\qquad 17\,240\ \hbox{logical lines},\quad
 895\,976\ \hbox{bytes},\\
 &\qquad
 \texttt{91b7b07cacd6910c0d89ea400ceef62505cbe724dffcc1209a730c2152ea3d38}.
\end{align*}

SM V16 proves a radial-gauge, factorial-free local parametrix for
the smooth-background overlap heat density.  It identifies the
continuum \(F^2\) coefficient and integrable mesoscopic errors, but
states explicitly that fixed admissibility does not control the
rough-field complement and therefore does not imply interacting
measure concentration.

NS V71 rescales at the sub-parabolic length
\(\lambda/\mathfrak Y\), obtains viscosity-only local constants,
and improves the Morrey floor from an exponential in
\(\mathfrak Y^3\) to one in \(\mathfrak Y^2\).  It also states that
this adapted-scale configuration has not been connected to the
preceding window configuration; that inter-scale transport is the
next open row.
\end{audit}

\begin{assessment}[Audit transfer]
\label{ass:newV90-audit-transfer}
No companion estimate is imported.  Two proof-boundary lessons
apply:
\begin{enumerate}[label=\textup{(\roman*)},leftmargin=3.5em]
\item a theorem on locally regular objects does not control a
      complementary rough or transversal sector; and
\item choosing the scale which normalizes a local estimate does not
      itself connect objects living at different scales or in
      different blocks.
\end{enumerate}
Accordingly, V90 does not infer fresh component edges from local
CIRC.  It first tests whether the block-compatible PAAC target can
even follow from the available root proxies in the transversal
sector.
\end{assessment}

\section{Freshness is not forced by connectedness}
\label{sec:newV90-freshness}

\begin{lemma}[Exact two-color occurrence identity]
\label{lem:newV90-two-color}
Let \(\mathscr R\) be a finite family of distinguishable root
occurrences, and give every occurrence \(\rho\) two color variables
\(z_\rho^{\rm d},z_\rho^{\rm b}\).  In a squarefree coefficient,
the substitution
\[
 z_\rho=z_\rho^{\rm d}+z_\rho^{\rm b}
\]
expands exactly into disjoint two-colorings of the selected
occurrences: every selected occurrence is differentiated or
background, and no occurrence has both colors.
\end{lemma}

\begin{proof}
Squarefreeness permits the exponent of \(z_\rho\) to be only zero
or one.  If it is one, the displayed substitution contributes
exactly its two linear summands.  Multiplication over the selected
occurrences gives all colorings once and gives no term containing
\(z_\rho^{\rm d}z_\rho^{\rm b}\).
\end{proof}

\begin{proposition}[Connectedness does not imply a fresh crossing]
\label{prop:newV90-no-freshness}
The connected squarefree logarithm and
\Cref{lem:newV90-two-color} do not imply FCCC.  A connected colored
incidence structure can have no background occurrence, and its
differentiated color can have positive block-compatibility defect.
\end{proposition}

\begin{proof}
Take one differentiated root whose support meets every label, and
take no background root.  Its label--root incidence graph is
connected.  Expand its strong one-root payment into a global Cayley
tree \(S\).  Choose a thermal partition for which \(S[P]\) is
disconnected in some block \(P\); such examples already occur on
three labels, with \(P=\{1,2\}\) and
\(S=\{\{1,3\},\{2,3\}\}\).  The defect is positive, but the
background color is empty and supplies no fresh crossing.
\end{proof}

\begin{correction}[Scope of the V89 freshness route]
\label{corr:newV90-freshness}
\Cref{thm:newV89-one-defect-BCAC} remains valid for the sector in
which a fresh background tree is actually present.  Neither
logarithmic connectedness nor the two-color squarefree identity
proves that hypothesis.  FCCC therefore cannot be promoted from a
conditional card to a general theorem by connectivity alone.
\end{correction}

\section{A proxy-level no-go for universal PAAC}
\label{sec:newV90-no-go}

Retain the deterministic threshold
\(a=\min\{b_*/4,1/4\}\) of
\Cref{lem:newV80-blocks}.  Fix an integer \(k\ge2\), to be chosen
below, and for each \(r\ge2\) put
\[
 n=kr,\qquad
 b_i={a\over k},\qquad
 s=\left({a\over k}\right)^2
 \quad(1\le i\le n).
                                                               \label{eq:newV90-balanced-data}
\]
In the fixed label order, the greedy partition has \(r\) completed
light blocks, each of size \(k\) and weight \(a\).

\begin{lemma}[The balanced family lies in the intermediate wedge]
\label{lem:newV90-wedge}
Choose \(k>a/\delta\).  For all sufficiently large \(r\), the data
\eqref{eq:newV90-balanced-data} obey
\[
 b_*<B=ra<\delta n,
\]
and hence lie in the intermediate wedge
\eqref{eq:newV79-wedge}.
\end{lemma}

\begin{proof}
The right inequality is \(ra<\delta kr\), which is the choice of
\(k\).  The left inequality holds once \(r>b_*/a\).
\end{proof}

\begin{lemma}[Compatible-tree weight in the balanced family]
\label{lem:newV90-compatible-weight}
For the greedy partition of
\eqref{eq:newV90-balanced-data},
\begin{align}
 {\cal C}_{n,r}
 &:=
 \sum_{S\in\mathfrak T_{\mathcal P}^{\rm comp}([n])}
 \prod_{i=1}^n b_i^{d_S(i)}
 \notag\\
 &=
 \left({a\over k}\right)^n
 a^{\,n-r}(ra)^{r-2}.                                     \label{eq:newV90-compatible-weight}
\end{align}
\end{lemma}

\begin{proof}
Apply \Cref{thm:newV88-hierarchical-Cayley}.  Every block has
\[
 v_P=\left({a\over k}\right)^k a^{k-2},
 \qquad w_P=a.
\]
Multiplying the \(r\) vertex factors, the \(r\) block marks, and
\(B^{r-2}=(ra)^{r-2}\) gives
\eqref{eq:newV90-compatible-weight}.
\end{proof}

\begin{theorem}[One transversal root defeats a proxy-only PAAC
deduction]
\label{thm:newV90-proxy-no-go}
Let a formal one-root activity on \([n]\) be assigned its strong
proxy
\begin{equation}
 u_{[n]}
 :=
 K_0^n\left(\prod_{i=1}^nb_i\right)B^{n-2}.                \label{eq:newV90-one-root-proxy}
\end{equation}
For the balanced data, this activity also obeys the reserve proxy
after one fixed enlargement of its numerical base.  Nevertheless,
\begin{equation}
 {u_{[n]}\over{\cal C}_{n,r}}
 =
 K_0^n r^{\,n-r}.                                          \label{eq:newV90-proxy-gap}
\end{equation}
For fixed \(k>1\), the factor \(r^{n-r}\) cannot be absorbed into
\(K^n\) uniformly in \(r\).  Therefore the two known one-root
proxies alone cannot prove a universal allocation to
block-compatible trees.
\end{theorem}

\begin{proof}
For equal marks \(b_i=\sqrt s\), the reserve proxy is
\[
 (K_1n)^ns^{(n-2)/2}\prod_i b_i
 =
 (K_1n)^n\left({a\over k}\right)^{2n-2}.
\]
The strong proxy is
\[
 K_0^n n^{n-2}\left({a\over k}\right)^{2n-2},
\]
so the latter is bounded by the former after a fixed choice of
bases.  Thus
\eqref{eq:newV90-one-root-proxy} is compatible with both available
proxy hypotheses.

Divide \eqref{eq:newV90-one-root-proxy} by
\eqref{eq:newV90-compatible-weight}.  All mark powers cancel and
\[
 {B^{n-2}\over a^{n-r}(ra)^{r-2}}
 =
 {(ra)^{n-2}\over a^{n-r}(ra)^{r-2}}
 =
 r^{n-r}.
\]
This proves \eqref{eq:newV90-proxy-gap}.  Since \(n=kr\),
\[
 \{r^{n-r}\}^{1/n}=r^{1-1/k}\longrightarrow\infty.
\]
\end{proof}

\begin{warning}[This is a proof-method no-go]
\label{warn:newV90-not-lower}
\Cref{thm:newV90-proxy-no-go} is not a lower bound on the exact
operator \(R_{[n]}\).  It proves that the established upper proxies
permit a transversal one-root majorant which is much larger than
the complete PAAC tree sum.  Hence a PAAC proof needs additional
exact cancellation, or it must keep such a transversal root as a
global object rather than force it through local block vertices.
\end{warning}

\begin{correction}[BCAC and PAAC are not the universal proxy target]
\label{corr:newV90-PAAC}
The implications
\[
 {\rm BCAC}\Longrightarrow{\rm PAAC}
 \Longrightarrow{\rm AABEC}
\]
remain correct.  What is retired is the claim that the root-proxy
programme V60--V89 can establish BCAC for every transversal
ancestry.  The compatible-tree sum is genuinely too small for that
deduction on \eqref{eq:newV90-balanced-data}.
\end{correction}

\section{The corrected hybrid target}
\label{sec:newV90-hybrid}

\begin{definition}[Local and transversal root occurrences]
\label{def:newV90-local-transversal}
Relative to the deterministic thermal partition \(\mathcal P\), a
root occurrence of support \(A\) is:
\[
 \begin{cases}
 \textup{local},&A\subseteq P\text{ for some }P\in\mathcal P,\\
 \textup{transversal},&A\text{ meets at least two blocks}.
 \end{cases}
\]
The two classes are disjoint and exhaust all occurrences.
\end{definition}

\begin{definition}[Hybrid local--transversal card, HLTC]
\label{def:newV90-HLTC}
HLTC is an exact reorganization of the intermediate-wedge
squarefree logarithm which:
\begin{enumerate}[label=\textup{(\roman*)},leftmargin=3.5em]
\item applies the small-\(B\), tree-valued V78 CIRC only to
      connected components made entirely of local occurrences;
\item retains every transversal occurrence with its full
      all-label support and its reserve/strong minimum;
\item constructs one connected quotient hypergraph whose vertices
      are the local components and whose hyperedges are the
      transversal occurrences;
\item sums that quotient without forcing its hyperedges into the
      block-compatible PAAC family; and
\item yields directly
      \[
      K^m\left(\prod_{i=1}^mb_i\right)B^{m-2}
      \]
      after the normalized Bernoulli and Duhamel sums.
\end{enumerate}
\end{definition}

\begin{proposition}[HLTC is sufficient]
\label{prop:newV90-HLTC}
HLTC, together with the already proved low- and high-total-mark
regions, proves full MCWC and SCNGC.
\end{proposition}

\begin{proof}
On the intermediate wedge, HLTC(v) is the strong target in
\Cref{def:newV61-MCWC}.  The regions
\(B\le b_*\) and \(B\ge\delta m\) are
\Cref{cor:newV79-two-regimes}.  Apply
\Cref{prop:newV79-TBRC} with its intermediate-wedge conclusion
replaced by the identical conclusion supplied directly by HLTC,
then use
\Cref{prop:newV61-reduction,prop:newV60-OCRC-SCNGC}.
\end{proof}

\begin{corollary}[The single-root transversal sector is closed]
\label{cor:newV90-single-root}
An ancestry consisting of one transversal root and no other
nonsingleton occurrence obeys the full strong all-label target,
without PAAC.
\end{corollary}

\begin{proof}
Apply the strong same-jump card
\Cref{thm:newV60-hyperedge-card} to its full support.  Singleton
decorations are contracted by the stable source exactly as in
V60--V70.  No thermal-block factorization is made.
\end{proof}

\begin{assessment}[The new bottleneck]
\label{ass:newV90-position}
The local CIRC and compatible-tree results remain useful, but only
for the occurrences which are actually local.  A transversal root
can carry the global Cayley degree which is absent from the product
of local block weights.  The next problem is therefore not fresh
component coalescence.  It is a hybrid connected-hypergraph sum on
local components, retaining each transversal support until the
global strong degree has been paid.
\end{assessment}

\begin{definition}[One-transversal-root attachment card, OTRAC]
\label{def:newV90-OTRAC}
OTRAC is the HLTC estimate for ancestries having exactly one
transversal root, with arbitrary local connected occurrence
components attached to its support.  The required allocation must
retain the transversal root's full strong polynomial and attach
each local component by a rooted forest, with total fibre
\(K^m\).
\end{definition}

\begin{proposition}[OTRAC is the first genuine HLTC row]
\label{prop:newV90-OTRAC}
If OTRAC holds, every intermediate-wedge ancestry with exactly one
transversal root obeys the full strong target.  No PAAC or
block-compatible-tree estimate is needed in this sector.
\end{proposition}

\begin{proof}
This is the defining OTRAC conclusion, together with the local
tree-valued CIRC factors in HLTC(i).  The quotient hypergraph has
one transversal hyperedge, so there is no transversal incidence
cycle.
\end{proof}

\begin{programme}[V91: one transversal root with local
attachments]
\label{prog:newV90-V91}
\begin{enumerate}[leftmargin=3em]
\item Fix the support \(A_*\) of the transversal root and contract
      it to a rooted set.
\item In every block, decompose the local CIRC output into the
      components which meet \(A_*\) and those which must attach to
      it.
\item Apply the rooted weighted-forest identity before replacing
      any local component weight by \(w_P\).
\item Prove OTRAC first when \(A_*\) meets every local component;
      then attach the remaining components by the V73 rooted-graph
      mechanism.
\item Keep the full factor \(B_{A_*}^{|A_*|-2}\); do not replace it
      by a product of block weights.
\end{enumerate}
\end{programme}

\begin{firewall}[Claim boundary after V90]
\label{fire:newV90-boundary}
V90 performs the compulsory SM/NS audit, proves the exact two-color
occurrence identity, proves that connectedness does not imply a
fresh background crossing, and gives an explicit intermediate-wedge
balanced family for which the available one-root proxies exceed the
entire compatible-tree sum by \(r^{n-r}\).  It corrects the claimed
scope of the BCAC/PAAC route, defines HLTC, proves its sufficiency,
and closes the isolated single-transversal-root sector.

V90 does not prove OTRAC with local attachments, HLTC, full MCWC,
OCRC, SCNGC, WPSC, GRLTC, LZTC, PLRC, WSDC, BRSC, HWSFC, UGCGC,
thermal connected WPRAC, all-arity RGSC, the raw Wilson aggregate
strong card, all-order strong MTA, WUFC, CPM, cutoff removal,
reflection positivity, Osterwalder--Schrader reconstruction, a
continuum Yang--Mills measure, or a positive Hamiltonian gap.  The
full Yang--Mills mass gap has not been proved.
\end{firewall}

\begin{theorem}[V90 result ledger]
\label{thm:newV90-results}
V90:
\begin{enumerate}[label=\textup{(V90.\arabic*)},leftmargin=6.5em]
\item performs the compulsory active-SM/active-NS audit;
\item proves the exact disjoint two-color occurrence identity;
\item disproves freshness as a consequence of connectedness;
\item constructs the balanced greedy-block test family;
\item proves the \(r^{n-r}\) proxy/compatible-tree gap;
\item corrects the universal scope of BCAC and PAAC;
\item formulates HLTC and proves its sufficiency; and
\item isolates OTRAC as its first nontrivial row.
\end{enumerate}
\end{theorem}

\begin{proof}
Use
\Cref{audit:newV90-companions,
lem:newV90-two-color,
prop:newV90-no-freshness,
lem:newV90-wedge,
lem:newV90-compatible-weight,
thm:newV90-proxy-no-go,
prop:newV90-HLTC,
prop:newV90-OTRAC}.
\end{proof}

\paragraph{Parent-version record.}
V90 was formed from
\texttt{Renewal\_Yang\_Mills\_Mass\_Gap\_Research\_Notes\_V89.tex}.
The V89 parent had \(40\,534\) logical source lines,
\(1\,307\,468\) bytes, and SHA--256
\[
 \texttt{25b71fe773bda6c99a19eab59fad95d7a07bf115781fef444d45c986cc04e3e1}.
\]
The next compulsory companion audit is V95.

% =============================================================================
% END APPENDED FIFTH-SERIES V90 MATERIAL
% =============================================================================

\clearpage
\chapter{Fifth-series V91: one transversal root and
single-anchor local components}
\label{ch:newV91-one-transversal}

\section{The component structure with one transversal root}
\label{sec:newV91-structure}

No companion audit is due at V91.

Fix a connected root-occurrence ancestry on a label set \(X\) and
the deterministic thermal partition \(\mathcal P\).  Suppose it has
exactly one transversal occurrence, with support \(A_*\), and all
other nonsingleton occurrences are local.  Form the incidence
subgraph containing only the local occurrences and their label
vertices; its nontrivial connected components will be called local
components.

\begin{lemma}[Every local component meets the transversal support]
\label{lem:newV91-meets}
Every local component \(C\) satisfies
\begin{equation}
 C\cap A_*\ne\varnothing.                                  \label{eq:newV91-meets}
\end{equation}
Moreover, two distinct local components have disjoint intersections
with \(A_*\).
\end{lemma}

\begin{proof}
A local component lies in one thermal block.  If it did not meet
\(A_*\), no local occurrence could join it to another block, and
the unique transversal occurrence would not meet it.  It would be
a separate component of the full incidence graph, contrary to
connectedness.

If two local components shared a label \(i\in A_*\), their incidence
graphs would be joined through the common label vertex \(i\), so
they would be one local component.  Thus the intersections are
disjoint.
\end{proof}

\begin{definition}[Single-anchor local sector]
\label{def:newV91-single-anchor}
The one-transversal ancestry is in the single-anchor local sector
when every local component \(C\) obeys
\[
 |C\cap A_*|=1.
\]
Write \(r(C)\) for this unique anchor label.
\end{definition}

\section{The exact rooted-forest polynomial}
\label{sec:newV91-rooted-forest}

Let \(Y\) be a finite set with nonempty declared root set
\(R\subseteq Y\), and put \(N=Y\setminus R\),
\(q=|N|\).  Let \(\mathfrak F_R(Y)\) be the forests on \(Y\) in
which every component contains exactly one vertex of \(R\).
Use the edge weight \(b_ib_j\).

\begin{theorem}[Weighted rooted-forest identity]
\label{thm:newV91-rooted-forest}
If \(q\ge1\), then
\begin{equation}
 \sum_{F\in\mathfrak F_R(Y)}
 \prod_{\{i,j\}\in E(F)}b_ib_j
 =
 B_R\left(\prod_{i\in N}b_i\right)B_Y^{q-1}.                \label{eq:newV91-rooted-forest}
\end{equation}
If \(q=0\), the forest is empty and the sum is one.
\end{theorem}

\begin{proof}
Orient every forest edge toward the unique root in its component.
Each nonroot \(i\in N\) then has one outgoing parent edge.  Factor
the child mark \(b_i\) from that edge.  The remaining weight of the
oriented edge \(i\to j\) is \(b_j\).

The directed matrix-forest theorem says that the sum of these
parent weights is the determinant of the \(N\times N\) matrix
\[
 M_{ij}
 =
 \begin{cases}
 B_Y-b_i,&i=j,\\
 -b_j,&i\ne j.
 \end{cases}
\]
Equivalently,
\[
 M=B_YI-\boldsymbol 1\,\boldsymbol b_N^{\,T}.
\]
The matrix determinant lemma gives
\[
 \det M
 =
 B_Y^q\left(1-{B_N\over B_Y}\right)
 =
 B_RB_Y^{q-1}.
\]
Restoring \(\prod_{i\in N}b_i\) proves
\eqref{eq:newV91-rooted-forest}.  The \(q=0\) statement is the
definition.
\end{proof}

\begin{corollary}[A connected tree polynomial is dominated by the
rooted forest]
\label{cor:newV91-tree-to-forest}
Suppose \(B_Y\le1\).  Then
\begin{equation}
 \sum_{T\in\mathfrak T(Y)}
 \prod_{i\in Y}b_i^{d_T(i)}
 \le
 \sum_{F\in\mathfrak F_R(Y)}
 \prod_{\{i,j\}\in E(F)}b_ib_j.                             \label{eq:newV91-tree-to-forest}
\end{equation}
\end{corollary}

\begin{proof}
Put \(h=|R|\).  If \(q\ge1\), divide the weighted Cayley identity
by \eqref{eq:newV91-rooted-forest}.  The ratio is
\[
 {\left(\prod_{i\in R}b_i\right)B_Y^{h-1}\over B_R}.
\]
By arithmetic--geometric mean,
\[
 {\prod_{i\in R}b_i\over B_R}
 \le {B_R^{h-1}\over h^h}.
\]
Since \(B_R,B_Y\le1\), the ratio is at most one.  If \(q=0\), the
tree sum is
\((\prod_{i\in R}b_i)B_R^{h-2}\le1\), again by the same
arithmetic--geometric mean estimate, with the cases \(h=1,2\)
read directly.
\end{proof}

\begin{remark}[One root gives equality]
\label{rem:newV91-one-root}
If \(R=\{r\}\), every tree on \(Y\) is a rooted forest with unique
root \(r\), and
\[
 \sum_{T\in\mathfrak T(Y)}
 \prod_i b_i^{d_T(i)}
 =
 b_r\left(\prod_{i\ne r}b_i\right)B_Y^{|Y|-2}.
\]
Thus no relaxation is made in the single-anchor sector.
\end{remark}

\section{Gluing rooted attachments to the transversal tree}
\label{sec:newV91-gluing}

Let \(\mathcal C\) be the set of local components.  Labels of
\(X\setminus A_*\) belong to exactly one \(C\in\mathcal C\).
An anchor in \(A_*\) which has no local occurrence is allowed to
belong to no component.

\begin{lemma}[Tree gluing at distinct anchors]
\label{lem:newV91-tree-gluing}
Let \(T_*\in\mathfrak T(A_*)\).  For every
\(C\in\mathcal C\), suppose
\(C\cap A_*=\{r(C)\}\), and let \(L_C\in\mathfrak T(C)\).
Then
\begin{equation}
 U:=T_*\cup\bigcup_{C\in\mathcal C}L_C                   \label{eq:newV91-glued-tree}
\end{equation}
is a tree on \(X\), and
\begin{equation}
 \left(\prod_{i\in A_*}b_i^{d_{T_*}(i)}\right)
 \prod_{C\in\mathcal C}
 \left(\prod_{i\in C}b_i^{d_{L_C}(i)}\right)
 =
 \prod_{i\in X}b_i^{d_U(i)}.                               \label{eq:newV91-glued-weight}
\end{equation}
\end{lemma}

\begin{proof}
Each \(L_C\) meets the previously constructed graph at exactly its
anchor \(r(C)\).  Attaching a tree at one vertex preserves
connectedness and creates no cycle.  Iteration proves that \(U\) is
a tree.  At an anchor, the degrees from \(T_*\) and its unique
local component add; all other vertices occur in one factor.
This proves \eqref{eq:newV91-glued-weight}.
\end{proof}

\begin{lemma}[The complete gluing data are recovered from the
output tree]
\label{lem:newV91-gluing-fibre}
For fixed \(A_*\) and thermal partition, the map from the complete
data
\[
 \bigl(T_*,\mathcal C,(L_C)_{C\in\mathcal C}\bigr)
 \longmapsto U
\]
in \eqref{eq:newV91-glued-tree} is injective, even when the
local-component label sets vary.
\end{lemma}

\begin{proof}
No local tree contains two vertices of \(A_*\).  Hence the induced
tree \(U[A_*]\) is exactly \(T_*\).  Delete its edges.  Every
remaining nontrivial component contains one anchor; its vertex set
recovers the corresponding local-component label set \(C\), and
its edges recover \(L_C\).  An isolated anchor corresponds to no
nontrivial local component.  Thus all input data are recovered.
\end{proof}

\section{Closure of the single-anchor OTRAC row}
\label{sec:newV91-closure}

\begin{theorem}[Single-anchor one-transversal-root card]
\label{thm:newV91-single-anchor}
Assume the local background occurrences have first been
exponentially normalized by support as in V86--V87.  On every light
block, the complete positive majorant of the one-transversal,
single-anchor local sector obeys
\begin{equation}
 K^{|X|}
 \left(\prod_{i\in X}b_i\right)B_X^{|X|-2}.                 \label{eq:newV91-single-anchor}
\end{equation}
All arities of the transversal root and all local incidence
cycles are included.
\end{theorem}

\begin{proof}
Expand the unique transversal root by the strong Pr\"ufer form
\Cref{cor:newV60-Prufer-form}; for a pair support include the drift
term by \Cref{lem:newV61-reserve}:
\[
 K^{|A_*|}
 \sum_{T_*\in\mathfrak T(A_*)}
 \prod_{i\in A_*}b_i^{d_{T_*}(i)}.
\]
For each local component \(C\), its total mark is at most that of
its light block and hence at most \(b_*/2\).  The normalized
full-label tree theorem
\Cref{thm:newV86-full-label} gives
\[
 K^{|C|}
 \sum_{L_C\in\mathfrak T(C)}
 \prod_{i\in C}b_i^{d_{L_C}(i)}.
\]
The V78 surplus-root factors are already included and are
independent of \(L_C\).

Multiply these expansions.  By
\Cref{lem:newV91-tree-gluing,lem:newV91-gluing-fibre}, every term
is the weight of one global tree \(U\), with no duplicate for fixed
component data.  The sum of constant exponents is
\[
 |A_*|+\sum_C|C|
 =
 |X|+|\mathcal C|
 \le2|X|,
\]
because distinct local components have distinct anchors by
\Cref{lem:newV91-meets}.  Absorb this into one base and enlarge the
resulting subset of global trees to all
\(\mathfrak T(X)\).  The weighted Cayley identity gives
\eqref{eq:newV91-single-anchor}.
\end{proof}

\begin{corollary}[Single-anchor OTRAC]
\label{cor:newV91-OTRAC}
OTRAC holds for every ancestry in which each local connected
component meets the unique transversal root support at exactly one
label.
\end{corollary}

\begin{proof}
The structural reduction is \Cref{lem:newV91-meets}; the strong
bound is \Cref{thm:newV91-single-anchor}.
\end{proof}

\begin{assessment}[Why the output tree is global]
\label{ass:newV91-global}
The proof does not factor the transversal root through the thermal
blocks.  Its tree \(T_*\) remains a global core.  A local component
is attached only at its unique common label, so its local Cayley
tree can be glued without a cycle.  The final weighted Cayley sum
is therefore the full strong target, not the smaller
block-compatible PAAC sum ruled out in V90.
\end{assessment}

\section{Several anchors in one local component}
\label{sec:newV91-multi-anchor}

If \(R_C:=C\cap A_*\) has at least two labels, a local output tree
on \(C\) can connect two vertices already joined by \(T_*\), and
the union need not be a tree.  The correct acyclic replacement is a
forest in \(\mathfrak F_{R_C}(C)\).  Corollary
\ref{cor:newV91-tree-to-forest} proves that the total local tree
polynomial is small enough for this replacement.  What it does not
control is the number of ways in which the anchor set \(A_*\) is
partitioned among several multi-anchor local incidence components.

\begin{definition}[Multi-anchor surplus attachment card, MSAC]
\label{def:newV91-MSAC}
MSAC is the aggregate allocation, for one fixed transversal root,
of every multi-anchor local component \(C\) to a forest in
\(\mathfrak F_{R_C}(C)\) such that:
\begin{enumerate}[label=\textup{(\roman*)},leftmargin=3.5em]
\item every nontransversal label is restored exactly once;
\item the local-component partition of \(A_*\) is summed before its
      tree weights are relaxed;
\item after gluing to one transversal tree \(T_*\), every term is
      one global Cayley tree on \(X\); and
\item the total inverse fibre is at most \(K^{|X|}\).
\end{enumerate}
\end{definition}

\begin{proposition}[MSAC completes OTRAC]
\label{prop:newV91-MSAC}
MSAC, together with \Cref{thm:newV91-single-anchor}, proves OTRAC
for arbitrary local attachments.
\end{proposition}

\begin{proof}
Use \Cref{thm:newV91-single-anchor} on all single-anchor
components.  On every remaining component use MSAC.  A rooted
forest has one component at each anchor, so its union with the
transversal tree is acyclic and connected.  MSAC(iv) absorbs all
component groupings; the final global Cayley sum gives the strong
target.
\end{proof}

\begin{assessment}[The next combinatorial object]
\label{ass:newV91-position}
The one-transversal programme is now reduced to multi-anchor local
components.  Each such component is surplus relative to the
already connected transversal support.  The scalar comparison
\eqref{eq:newV91-tree-to-forest} pays it, but a proof of MSAC must
retain the partition of its anchor labels until the rooted forests
have been glued.  This is the rooted-forest analogue of the V77
Bell-to-Cayley quotient.
\end{assessment}

\begin{programme}[V92: multi-anchor surplus forests]
\label{prog:newV91-V92}
\begin{enumerate}[leftmargin=3em]
\item Root the transversal tree \(T_*\) and orient every local
      rooted forest toward its anchor set.
\item Encode a multi-anchor local component by the minimal subtree
      of \(T_*\) spanning its anchors and by its nontransversal
      rooted forest.
\item Sum anchor partitions with the V77 edge-partition fibre,
      retaining the small local block weight through the sum.
\item Treat first the case in which every multi-anchor component
      has two anchors; test a matching family on a transversal
      star.
\item Use the V78 whole-surplus gas only after every
      nontransversal label has been assigned to an output forest.
\end{enumerate}
\end{programme}

\begin{firewall}[Claim boundary after V91]
\label{fire:newV91-boundary}
V91 proves the exact weighted rooted-forest identity, proves a
connected Cayley polynomial is dominated by the corresponding
rooted-forest polynomial on a light block, classifies local
components around one transversal root, and closes the complete
single-anchor local sector by an injective global-tree gluing.  It
formulates MSAC as the remaining one-transversal operation.

V91 does not prove MSAC, general OTRAC, HLTC, full MCWC, OCRC,
SCNGC, WPSC, GRLTC, LZTC, PLRC, WSDC, BRSC, HWSFC, UGCGC,
thermal connected WPRAC, all-arity RGSC, the raw Wilson aggregate
strong card, all-order strong MTA, WUFC, CPM, cutoff removal,
reflection positivity, Osterwalder--Schrader reconstruction, a
continuum Yang--Mills measure, or a positive Hamiltonian gap.  The
full Yang--Mills mass gap has not been proved.
\end{firewall}

\begin{theorem}[V91 result ledger]
\label{thm:newV91-results}
V91:
\begin{enumerate}[label=\textup{(V91.\arabic*)},leftmargin=6.5em]
\item proves every local component meets the transversal support;
\item proves distinct components have disjoint anchor sets;
\item proves the exact weighted rooted-forest identity;
\item compares connected tree and rooted-forest polynomials;
\item proves the distinct-anchor global-tree gluing;
\item proves its injectivity;
\item closes the single-anchor OTRAC sector; and
\item formulates MSAC and proves its sufficiency.
\end{enumerate}
\end{theorem}

\begin{proof}
Use
\Cref{lem:newV91-meets,
thm:newV91-rooted-forest,
cor:newV91-tree-to-forest,
lem:newV91-tree-gluing,
lem:newV91-gluing-fibre,
thm:newV91-single-anchor,
cor:newV91-OTRAC,
prop:newV91-MSAC}.
\end{proof}

\paragraph{Parent-version record.}
V91 was formed from
\texttt{Renewal\_Yang\_Mills\_Mass\_Gap\_Research\_Notes\_V90.tex}.
The V90 parent had \(40\,965\) logical source lines,
\(1\,323\,103\) bytes, and SHA--256
\[
 \texttt{30a4f916583435a5f76f134ba5b43b2be0c9422d7bdab37188b4ba958fc32102}.
\]
The next compulsory companion audit is V95.

% =============================================================================
% END APPENDED FIFTH-SERIES V91 MATERIAL
% =============================================================================

\clearpage
\chapter{Fifth-series V92: the multi-anchor surplus gas and
closure of one-transversal-root attachments}
\label{ch:newV92-multi-anchor}

\section{Retaining the tree-to-forest ratio}
\label{sec:newV92-ratio}

No companion audit is due at V92.

Let \(C\) be a local connected component in a light block \(P\),
let
\[
 R:=C\cap A_*,\qquad h:=|R|\ge2,\qquad
 N:=C\setminus R,\qquad q:=|N|,
\]
and put \(X_R:=B_C\).  The weighted connected-tree polynomial on
\(C\) and the rooted-forest polynomial of V91 are denoted by
\[
 {\cal T}(C):=
 \sum_{L\in\mathfrak T(C)}
 \prod_{i\in C}b_i^{d_L(i)},
 \qquad
 \Phi_R(C):=
 \sum_{F\in\mathfrak F_R(C)}
 \prod_{\{i,j\}\in E(F)}b_ib_j.
\]

\begin{lemma}[Exact multi-anchor ratio]
\label{lem:newV92-exact-ratio}
For every \(q\ge0\),
\begin{equation}
 {\cal T}(C)=\theta_R(C)\Phi_R(C),                           \label{eq:newV92-exact-ratio}
\end{equation}
where
\begin{equation}
 \theta_R(C)
 :=
 {\displaystyle\prod_{r\in R}b_r\over B_R}
 B_C^{h-1}.                                                 \label{eq:newV92-theta}
\end{equation}
For \(q=0\), \(\Phi_R(C)=1\) and the same formula is valid.
\end{lemma}

\begin{proof}
If \(q\ge1\), the weighted Cayley and rooted-forest identities give
\begin{align*}
 {\cal T}(C)
 &=
 \left(\prod_{r\in R}b_r\right)
 \left(\prod_{i\in N}b_i\right)B_C^{q+h-2},\\
 \Phi_R(C)
 &=
 B_R\left(\prod_{i\in N}b_i\right)B_C^{q-1}.
\end{align*}
Their ratio is \eqref{eq:newV92-theta}.  If \(q=0\), then
\[
 {\cal T}(R)=
 \left(\prod_{r\in R}b_r\right)B_R^{h-2}
 =
 {\prod_{r\in R}b_r\over B_R}B_R^{h-1},
\]
while the rooted forest has no edges and weight one.
\end{proof}

\begin{assessment}[Do not replace \(\theta\) by one]
\label{ass:newV92-theta}
Corollary \ref{cor:newV91-tree-to-forest} used only
\(\theta_R(C)\le1\).  That relaxation is sufficient for one fixed
component but loses the factor which pays the set partition of
several anchor groups.  V92 retains
\eqref{eq:newV92-theta} until every such grouping has been summed.
\end{assessment}

\section{Output branches and their masses}
\label{sec:newV92-branches}

Fix one transversal core tree \(T_*\) on \(A_*\).  After every local
component has been replaced by a rooted forest, its union with
\(T_*\) is a global tree.  Delete the core edges \(E(T_*)\).
For every anchor \(r\in A_*\), let \(D_r\) be the resulting tree
component containing \(r\), allowing \(D_r=\{r\}\), and set
\begin{equation}
 x_r:=B_{D_r}=\sum_{i\in D_r}b_i.                           \label{eq:newV92-branch-mass}
\end{equation}
If \(I_P:=A_*\cap P\), then
\begin{equation}
 b_r\le x_r,\qquad
 \sum_{r\in I_P}x_r\le w_P.                                \label{eq:newV92-branch-budget}
\end{equation}
For a group \(R\subseteq I_P\), the original local component is the
union of the branch sets \(D_r\), \(r\in R\), and therefore
\[
 B_C=X_R:=\sum_{r\in R}x_r.
\]

\begin{lemma}[Branch data determine the component once the anchor
group is known]
\label{lem:newV92-component-recovery}
Fix the output global tree, its core \(T_*\), and an anchor group
\(R\subseteq I_P\).  Then the label set and rooted forest of the
corresponding local component are uniquely determined:
\[
 C_R=\bigcup_{r\in R}D_r,
 \qquad
 F_R=\bigcup_{r\in R}D_r.
\]
Here the second union means the union of the branch edge sets.
\end{lemma}

\begin{proof}
Deleting \(E(T_*)\) recovers every \(D_r\) and all of its edges.
A rooted forest for the group \(R\) is precisely their disjoint
union, because it has one component at each root in \(R\).  Thus
only the grouping of the anchors was forgotten.
\end{proof}

\section{The anchor-group surplus estimate}
\label{sec:newV92-gas}

\begin{lemma}[Summed multi-anchor activity]
\label{lem:newV92-summed-theta}
Let \(I\) be finite, let \(0<b_i\le x_i\), and suppose
\[
 \sum_{i\in I}x_i\le w\le1.
\]
For \(R\subseteq I\), \(|R|=h\ge2\), put
\begin{equation}
 \theta(R)
 :=
 {\prod_{i\in R}b_i\over\sum_{i\in R}b_i}
 \left(\sum_{i\in R}x_i\right)^{h-1}.                      \label{eq:newV92-abstract-theta}
\end{equation}
Then
\begin{equation}
 \sum_{\substack{R\subseteq I\\|R|\ge2}}\theta(R)
 \le
 \sum_{h\ge2}{w^{2h-2}\over(h-1)!}
 =
 w^2e^{w^2}
 \le e\,w^2.                                                \label{eq:newV92-summed-theta}
\end{equation}
\end{lemma}

\begin{proof}
Fix \(R\), put \(X_R=\sum_{i\in R}x_i\) and
\(B_R=\sum_{i\in R}b_i\).  Since \(B_R\ge b_j\) for every
\(j\in R\),
\[
 {\prod_{i\in R}b_i\over B_R}
 \le
 \prod_{i\in R\setminus\{j\}}b_i.
\]
Average these inequalities with the probabilities \(x_j/X_R\).
Multiplication by \(X_R^{h-1}\) gives
\begin{equation}
 \theta(R)
 \le
 X_R^{h-2}
 \sum_{j\in R}x_j
 \prod_{i\in R\setminus\{j\}}b_i
 \le
 w^{h-2}
 \sum_{j\in R}x_j
 \prod_{i\in R\setminus\{j\}}b_i.                          \label{eq:newV92-theta-split}
\end{equation}
Sum \eqref{eq:newV92-theta-split} over all \(h\)-element subsets.
For fixed \(j\), the elementary-symmetric bound gives
\[
 \sum_{\substack{S\subseteq I\setminus\{j\}\\|S|=h-1}}
 \prod_{i\in S}b_i
 \le {(\sum_{i\in I}b_i)^{h-1}\over(h-1)!}
 \le {w^{h-1}\over(h-1)!}.
\]
Therefore
\[
 \sum_{|R|=h}\theta(R)
 \le
 {w^{h-2}\over(h-1)!}
 \left(\sum_jx_j\right)w^{h-1}
 \le {w^{2h-2}\over(h-1)!}.
\]
Sum \(h\ge2\).
\end{proof}

\begin{corollary}[All anchor partitions cost a small gas]
\label{cor:newV92-partition-gas}
Under the hypotheses of
\Cref{lem:newV92-summed-theta},
\begin{equation}
 \sum_{\pi\in{\rm Part}(I)}
 \prod_{\substack{R\in\pi\\|R|\ge2}}\theta(R)
 \le \exp(e\,w^2).                                         \label{eq:newV92-partition-gas}
\end{equation}
Singleton groups have weight one.
\end{corollary}

\begin{proof}
Enlarge the disjoint family of nonsingleton blocks of a partition
to an arbitrary set of nonsingleton subsets.  Positivity gives
\[
 \sum_\pi\prod_{\substack{R\in\pi\\|R|\ge2}}\theta(R)
 \le
 \prod_{\substack{R\subseteq I\\|R|\ge2}}\{1+\theta(R)\}
 \le
 \exp\left\{\sum_{|R|\ge2}\theta(R)\right\}.
\]
Apply \Cref{lem:newV92-summed-theta}.
\end{proof}

\begin{corollary}[The product of all light-block gases is
exponential in the label number]
\label{cor:newV92-global-gas}
For the deterministic thermal partition,
\begin{equation}
 \prod_{\substack{P\in\mathcal P\\P\ {\rm light}}}
 \exp(e\,w_P^2)
 \le K^m.                                                    \label{eq:newV92-global-gas}
\end{equation}
\end{corollary}

\begin{proof}
Every light block satisfies \(w_P\le b_*/2\), hence
\[
 \sum_{P\ {\rm light}}w_P^2
 \le {b_*\over2}\sum_Pw_P
 ={b_*\over2}B
 \le {b_*\over2}m,
\]
because every \(b_i<1\).  Exponentiate and enlarge the numerical
base.
\end{proof}

\section{MSAC and OTRAC are closed}
\label{sec:newV92-closure}

\begin{theorem}[Multi-anchor surplus attachment theorem]
\label{thm:newV92-MSAC}
MSAC holds.  More precisely, after the local normalized occurrence
sum of V86--V87, all multi-anchor local components around one
transversal root can be replaced by rooted forests and their
complete forgotten anchor grouping has aggregate fibre at most
\(K^m\).
\end{theorem}

\begin{proof}
For each local connected component \(C\), V86 supplies the
tree-valued majorant \(K^{|C|}{\cal T}(C)\).  If
\(|R_C|=1\), use the equality in
\Cref{rem:newV91-one-root}.  If \(|R_C|\ge2\), use the exact
factorization
\[
 {\cal T}(C)=\theta_{R_C}(C)\Phi_{R_C}(C)
\]
from \Cref{lem:newV92-exact-ratio}, retaining the scalar
\(\theta\).

Fix the resulting global output tree and its transversal core
\(T_*\).  By
\Cref{lem:newV92-component-recovery}, the only forgotten data in
one block are the groups in a partition of \(I_P=A_*\cap P\).
For a group \(R\), its exact scalar is
\eqref{eq:newV92-abstract-theta}, with the branch masses
\eqref{eq:newV92-branch-mass}.  The budget
\eqref{eq:newV92-branch-budget} permits
\Cref{cor:newV92-partition-gas}.  Multiply the block bounds and use
\Cref{cor:newV92-global-gas}.

All local V86 fibres cost
\[
 K^{\sum_C|C|}
 \le K^{m+|A_*|}
 \le K^{2m},
\]
because nonanchor labels are disjoint and distinct local components
have disjoint anchor sets.  Thus the complete grouping and local
tree fibre is \(K^m\) after the base is enlarged.  Each nonanchor
label occurs in exactly one rooted branch, and gluing these branches
to \(T_*\) gives one global tree.  These are MSAC(i)--(iv).
\end{proof}

\begin{theorem}[Full one-transversal-root attachment card]
\label{thm:newV92-OTRAC}
OTRAC holds for arbitrary local connected components.  Hence every
intermediate-wedge ancestry with exactly one transversal root obeys
\begin{equation}
 K^m\left(\prod_{i=1}^mb_i\right)B^{m-2}.                   \label{eq:newV92-OTRAC}
\end{equation}
\end{theorem}

\begin{proof}
Single-anchor components are
\Cref{thm:newV91-single-anchor}; multi-anchor components are
\Cref{thm:newV92-MSAC}.  Apply
\Cref{prop:newV91-MSAC}.  After all rooted forests are glued to the
transversal tree, enlarge their output family to all global Cayley
trees and evaluate the weighted sum.
\end{proof}

\begin{corollary}[The first hybrid row is complete]
\label{cor:newV92-first-hybrid}
The exactly-one-transversal-root sector of HLTC is closed, with no
block-compatible PAAC hypothesis and no fresh-background
assumption.
\end{corollary}

\begin{proof}
Use \Cref{thm:newV92-OTRAC,prop:newV90-OTRAC}.
\end{proof}

\begin{assessment}[What paid the Bell grouping]
\label{ass:newV92-payment}
A multi-anchor local component is surplus because its anchors are
already connected by the transversal root.  Cutting its local tree
into a rooted forest loses the grouping of those anchors.  The lost
grouping is paid by the exact ratio \(\theta_R(C)\), not by counting
partitions.  Inequality \eqref{eq:newV92-theta-split} distributes
one anchor denominator against the branch masses, and the remaining
elementary-symmetric sum makes the complete gas \(O(w_P^2)\).
\end{assessment}

\section{Several transversal roots}
\label{sec:newV92-several}

With two or more transversal roots, the transversal supports
themselves form a connected hypergraph only after local components
are included.  A local component may meet several transversal
supports and can be essential rather than surplus.  The one-core
argument cannot declare all of its multi-anchor connections
discardable.

\begin{definition}[Transversal-incidence card, TIC]
\label{def:newV92-TIC}
TIC is the remaining HLTC operation for two or more transversal
root occurrences.  It must:
\begin{enumerate}[label=\textup{(\roman*)},leftmargin=3.5em]
\item form the exact bipartite incidence graph whose left vertices
      are transversal roots and whose right data are local connected
      components;
\item select a canonical incidence hypertree before treating any
      local component as surplus;
\item use rooted-forest output for every component meeting an
      already connected set of transversal anchors;
\item retain the strong polynomial of every essential transversal
      root; and
\item sum incidence excess and repeated supports with one
      exponential fibre, yielding the global strong target.
\end{enumerate}
\end{definition}

\begin{proposition}[TIC completes HLTC]
\label{prop:newV92-TIC}
TIC, together with \Cref{thm:newV92-OTRAC}, proves HLTC and
therefore full MCWC and SCNGC.
\end{proposition}

\begin{proof}
Zero transversal roots are local CIRC.  Exactly one is
\Cref{thm:newV92-OTRAC}.  TIC treats every remaining connected
ancestry.  Thus HLTC holds; apply
\Cref{prop:newV90-HLTC}.
\end{proof}

\begin{programme}[V93: the transversal incidence hypertree]
\label{prog:newV92-V93}
\begin{enumerate}[leftmargin=3em]
\item Contract every local connected component to its set of
      incident transversal roots and root the resulting bipartite
      graph at one transversal occurrence.
\item Prove the incidence-hypertree row first, using one essential
      rooted-forest attachment at each pruning step.
\item Retain each transversal root's full label support and strong
      Cayley polynomial during the pruning.
\item Apply the V77 two-level tree gluing to the quotient only after
      local branch masses have been assigned.
\item Leave quotient incidence cycles for the V78 surplus gas after
      the hypertree row is coefficientwise closed.
\end{enumerate}
\end{programme}

\begin{firewall}[Claim boundary after V92]
\label{fire:newV92-boundary}
V92 proves the exact tree-to-rooted-forest ratio, reconstructs local
components from output branches and anchor groups, proves the
summed multi-anchor activity is \(O(w_P^2)\), and exponentiates all
anchor partitions without a Bell loss.  It proves MSAC and closes
OTRAC for arbitrary local attachments, completing the
exactly-one-transversal-root row of HLTC.  It formulates TIC for
several transversal roots.

V92 does not prove TIC, HLTC, full MCWC, OCRC, SCNGC, WPSC, GRLTC,
LZTC, PLRC, WSDC, BRSC, HWSFC, UGCGC, thermal connected WPRAC,
all-arity RGSC, the raw Wilson aggregate strong card, all-order
strong MTA, WUFC, CPM, cutoff removal, reflection positivity,
Osterwalder--Schrader reconstruction, a continuum Yang--Mills
measure, or a positive Hamiltonian gap.  The full Yang--Mills mass
gap has not been proved.
\end{firewall}

\begin{theorem}[V92 result ledger]
\label{thm:newV92-results}
V92:
\begin{enumerate}[label=\textup{(V92.\arabic*)},leftmargin=6.5em]
\item proves the exact multi-anchor tree/forest ratio;
\item defines the output-branch mass budget;
\item proves component recovery from an anchor group;
\item proves the summed \(\theta\)-activity bound;
\item proves the anchor-partition gas estimate;
\item proves its global exponential ledger;
\item proves MSAC and full OTRAC; and
\item formulates TIC and proves its sufficiency for HLTC.
\end{enumerate}
\end{theorem}

\begin{proof}
Use
\Cref{lem:newV92-exact-ratio,
lem:newV92-component-recovery,
lem:newV92-summed-theta,
cor:newV92-partition-gas,
cor:newV92-global-gas,
thm:newV92-MSAC,
thm:newV92-OTRAC,
prop:newV92-TIC}.
\end{proof}

\paragraph{Parent-version record.}
V92 was formed from
\texttt{Renewal\_Yang\_Mills\_Mass\_Gap\_Research\_Notes\_V91.tex}.
The V91 parent had \(41\,379\) logical source lines,
\(1\,337\,500\) bytes, and SHA--256
\[
 \texttt{65a24353e41176b134b9d084b908abfe94af2b92a5765e2d17efd5d89395f38f}.
\]
The next compulsory companion audit is V95.

% =============================================================================
% END APPENDED FIFTH-SERIES V92 MATERIAL
% =============================================================================

\clearpage
\chapter{Fifth-series V93: primitive-root incidence hypertrees
across thermal blocks}
\label{ch:newV93-transversal-hypertree}

\section{Stable Duhamel contraction in the intermediate wedge}
\label{sec:newV93-contraction}

No companion audit is due at V93.

The small-\(B\) contraction in V70 used
\({\cal V}\preceq cB^2I\).  The percolative source constructed in
V81 is stronger for the block problem: at every interpolation
matrix,
\[
 {\cal V}^{\rm perc}(\boldsymbol x)\preceq0.
\]
Thus the Duhamel word itself creates no total-mark loss in the
intermediate wedge.

\begin{lemma}[Global percolative Duhamel contraction]
\label{lem:newV93-Duhamel}
Let \(A_1,\ldots,A_q\) be arbitrary operator insertions obtained by
differentiating the percolative moment exponent
\eqref{eq:newV81-V}.  Then, on every interpolation matrix,
\begin{equation}
 \left\|
 D^qe^{{\cal V}^{\rm perc}}
 [A_1,\ldots,A_q]
 \right\|_{\rm cb}
 \le
 \prod_{\nu=1}^q\|A_\nu\|_{\rm cb}.                         \label{eq:newV93-Duhamel}
\end{equation}
The constant is independent of \(q,m,B\).
\end{lemma}

\begin{proof}
By \Cref{thm:newV81-positive},
\({\cal V}^{\rm perc}\preceq0\) at every matrix amplification.
Apply the ordered-simplex proof of
\Cref{lem:newV70-Duhamel} with \(\varepsilon=0\).  The \(q!\)
orders cancel the simplex volume \(1/q!\).
\end{proof}

\begin{assessment}[Analytic and combinatorial gates separate]
\label{ass:newV93-separation}
The global semigroup is contractive, so the remaining intermediate
wedge problem is the positive summation of the differentiated root
insertions and their percolation gates.  No factor
\(\exp(CB^2)\) is needed for the Duhamel word.
\end{assessment}

\section{Primitive root insertions retain the reserve}
\label{sec:newV93-root-reserve}

\begin{definition}[Primitive active-root hypergraph]
\label{def:newV93-primitive}
After expanding each nonempty Duhamel derivative block, form the
bipartite graph whose left vertices are all representation labels
used by nonsingleton insertions and whose right vertices are the
distinguishable primitive occurrences of \(R_A\).  Join \(i\) to
the occurrence when \(i\in A\).  This is the primitive active-root
hypergraph.  Singleton dissipative decorations remain in the
contractive background.
\end{definition}

\begin{lemma}[Differentiated gates do not enlarge one root beyond
an exponential base]
\label{lem:newV93-gate-reserve}
After the exact support polarization of V87, an occurrence of
support \(A\), \(|A|=a\ge2\), with a nonempty forest \(E\) of
owned label-edge derivatives has norm at most
\begin{equation}
 K^{a+|E|}
 (Ka)^a s^{(a-2)/2}\prod_{i\in A}b_i.                      \label{eq:newV93-gate-reserve}
\end{equation}
Across a squarefree outer tree, the product of all extra
\(K^{|E|}\) factors is at most \(K^m\).
\end{lemma}

\begin{proof}
The coefficient of the normalized gate is bounded by
\(\gamma_A(E)\le1\) by
\Cref{cor:newV87-gamma-formula}.  Multiply by the unrelaxed root
estimate \eqref{eq:newV61-reserve}.  Percolation finite-difference
sign splittings and port ownership enlarge the numerical constant
once per owned edge.  The owned edge sets of distinct Duhamel
blocks are disjoint subsets of an outer tree, so
\[
 \sum |E|\le m-1.
\]
\end{proof}

\begin{remark}[Repeated supports are cyclic]
\label{rem:newV93-repeat}
Two primitive occurrences with the same support \(A\), \(|A|\ge2\),
create a four-cycle in the incidence graph.  Therefore a primitive
incidence hypertree is automatically simple by support.  The
repeated-support collapse of V87 is needed only in the cyclic
sector.
\end{remark}

\section{The all-label incidence-hypertree theorem}
\label{sec:newV93-hypertree}

\begin{definition}[Primitive-root ancestry hypertree, PRAH]
\label{def:newV93-PRAH}
PRAH is the sector in which the primitive active-root incidence
graph is connected and is a tree after singleton decorations have
been removed.  Local and transversal roots are both allowed, with
arbitrary arities.
\end{definition}

\begin{lemma}[V77 is independent of total smallness at the
reserve level]
\label{lem:newV93-V77-uniform}
Let \(X\) have \(m\ge2\) labels with \(b_i\ge\sqrt s\), with no
upper assumption on \(B_X\).  Assign every hyperedge \(A\) the
reserve activity
\[
 (K|A|)^{|A|}s^{(|A|-2)/2}\prod_{i\in A}b_i.
\]
Then the complete sum over incidence hypertrees on \(X\) is at most
\begin{equation}
 K^m\left(\prod_{i\in X}b_i\right)B_X^{m-2}.                \label{eq:newV93-V77-uniform}
\end{equation}
\end{lemma}

\begin{proof}
Apply the rooted-hypertree/child-partition bijection
\Cref{thm:newV75-bijection}.  The local Bell sums use only the
reserve activity and the floor \(b_i\ge\sqrt s\), as in
\Cref{lem:newV75-fixed-tree}.  The resulting rooted quotient is
RHQC.  Its proof
\Cref{thm:newV77-CQDC-RHQC} is a positive polynomial identity for
arbitrary nonnegative marks and yields
\[
 2^{m-2}\sum_{U\in\mathfrak T(X)}
 \prod_i b_i^{d_U(i)-1}.
\]
Restoring one mark per label and applying the weighted Cayley
identity gives \eqref{eq:newV93-V77-uniform}.  No step in this
combinatorial argument assumes \(B_X\le b_*\); that assumption in
V74--V78 was used to derive and sum effective cyclic
centre-hyperedges, not in RHQC itself.
\end{proof}

\begin{theorem}[PRAH strong aggregate card]
\label{thm:newV93-PRAH}
On the intermediate wedge, the complete positive majorant of all
PRAH terms obeys
\begin{equation}
 K^m\left(\prod_{i=1}^mb_i\right)B^{m-2}.                   \label{eq:newV93-PRAH}
\end{equation}
This includes arbitrary mixtures and arities of local and
transversal primitive roots.
\end{theorem}

\begin{proof}
Use \Cref{lem:newV93-Duhamel} before taking the product of insertion
norms.  Apply \Cref{lem:newV93-gate-reserve} to every primitive
root.  The total gate and sign cost is \(K^m\).  The remaining sum
is exactly the reserve-weighted incidence-hypertree sum of
\Cref{lem:newV93-V77-uniform}.  Singleton dissipative terms are
part of the contraction and create no connected hyperedge choice.
\end{proof}

\begin{corollary}[Acyclic TIC is closed]
\label{cor:newV93-acyclic-TIC}
TIC holds for every ancestry whose primitive local--transversal
incidence graph has excess zero.
\end{corollary}

\begin{proof}
Such an ancestry is PRAH; apply \Cref{thm:newV93-PRAH}.
\end{proof}

\begin{assessment}[Why block compatibility is absent]
\label{ass:newV93-no-block}
PRAH is summed on the original label set.  The V77 output tree may
cross and re-enter a thermal block, but this is harmless because
the conclusion is the global strong polynomial required by HLTC,
not the smaller block-compatible PAAC target.  The thermal
partition is used for the exact percolative source and its
contraction, not imposed on the output Cayley tree.
\end{assessment}

\section{The precise remaining cycle sector}
\label{sec:newV93-cycles}

For a connected primitive occurrence hypergraph \({\cal H}\) on
\(X\), define its incidence excess by
\begin{equation}
 g({\cal H})
 :=
 \sum_{A\in{\cal H}}(|A|-1)-|X|+1.                         \label{eq:newV93-excess}
\end{equation}
This is the cyclomatic number of its bipartite incidence graph,
including multiplicity of distinguishable occurrences.

\begin{lemma}[Already closed cycles]
\label{lem:newV93-closed-cycles}
The following cyclic sectors are already closed:
\begin{enumerate}[label=\textup{(\roman*)},leftmargin=3.5em]
\item every cycle contained entirely in one light block;
\item every ancestry with exactly one transversal root, including
      arbitrary local cycles; and
\item every primitive incidence structure with \(g=0\).
\end{enumerate}
\end{lemma}

\begin{proof}
Item (i) is local tree-valued CIRC on \(w_P\le b_*/2\), using
\Cref{thm:newV85-tree-CIRC,thm:newV86-full-label}.  Item (ii) is
\Cref{thm:newV92-OTRAC}.  Item (iii) is
\Cref{thm:newV93-PRAH}.
\end{proof}

\begin{corollary}[First unresolved transversal cycle]
\label{cor:newV93-first-cycle}
An unresolved TIC ancestry has at least two transversal roots and
positive primitive incidence excess after all purely local cyclic
components have been resummed.
\end{corollary}

\begin{proof}
This is the complement of the three cases in
\Cref{lem:newV93-closed-cycles}.
\end{proof}

\begin{warning}[The V78 global surplus estimate is unavailable]
\label{warn:newV93-no-global-V78}
At fixed excess, the relaxed strong ledger contributes
\[
 B^{m-2+2g}.
\]
In V78 the extra \(B^{2g}\) was summable because \(B\le b_*\).
On the intermediate wedge \(B\) may grow with \(m\), so neither
\(\exp(CB^2)\) nor \(\exp(CmB^2)\) can be absorbed into \(K^m\).
Thus transversal cycles cannot be closed by applying the V78
surplus gas globally.
\end{warning}

\begin{definition}[Transversal-cycle normalization card, TCNC]
\label{def:newV93-TCNC}
TCNC is the cyclic row of TIC.  It must:
\begin{enumerate}[label=\textup{(\roman*)},leftmargin=3.5em]
\item choose a canonical primitive incidence hypertree before
      relaxing any cycle marks;
\item retain the full normalized Bernoulli configuration of every
      undifferentiated block edge;
\item pair each surplus transversal incidence with either a
      Kruskal shortcut factor or a repeated-support type
      normalization;
\item sum the quotient-cycle choices before using
      \(B_A\le B\); and
\item leave at most \(K^m\) times the PRAH tree sum.
\end{enumerate}
\end{definition}

\begin{proposition}[TCNC completes TIC]
\label{prop:newV93-TCNC}
TCNC, together with \Cref{cor:newV93-acyclic-TIC}, proves TIC,
HLTC, full MCWC, and SCNGC.
\end{proposition}

\begin{proof}
The primitive incidence excess is either zero or positive.  The
first case is \Cref{cor:newV93-acyclic-TIC}; TCNC treats the second.
Thus TIC holds.  Apply
\Cref{prop:newV92-TIC,prop:newV90-HLTC}.
\end{proof}

\begin{assessment}[Position after the hypertree closure]
\label{ass:newV93-position}
All analytic Duhamel orders and all acyclic primitive
local--transversal incidences are now closed at arbitrary total
mark.  The remaining intermediate-wedge operation is genuinely
cyclic and genuinely transversal.  Its smallness cannot come from
the global total mark; it must come from the normalized
percolation/Kruskal probability or from exact repeated-support
symmetrization.
\end{assessment}

\begin{programme}[V94: the first transversal incidence cycle]
\label{prog:newV93-V94}
\begin{enumerate}[leftmargin=3em]
\item Classify the minimal four-cycle formed by two transversal
      root occurrences sharing two label or local-component
      vertices.
\item Condition all other percolation edges and write the exact
      two-root simultaneous Russo quotient.
\item Apply the V83 maximum-spanning-tree identity before splitting
      signs or dropping closed-edge factors.
\item Test whether the one surplus incidence is paid by one
      shortcut probability with no \(B^2\) relaxation.
\item If the four-cycle closes, formulate its canonical contraction
      as the induction step for arbitrary transversal excess.
\end{enumerate}
\end{programme}

\begin{firewall}[Claim boundary after V93]
\label{fire:newV93-boundary}
V93 proves global percolative Duhamel contraction on the
intermediate wedge, proves differentiated primitive roots retain
the unrelaxed reserve, proves the V77 hypertree bound requires no
smallness of the total mark at the reserve level, and closes every
primitive local--transversal incidence hypertree.  It identifies
the only remaining TIC sector as positive transversal incidence
excess and formulates TCNC.

V93 does not prove TCNC, cyclic TIC, HLTC, full MCWC, OCRC, SCNGC,
WPSC, GRLTC, LZTC, PLRC, WSDC, BRSC, HWSFC, UGCGC, thermal
connected WPRAC, all-arity RGSC, the raw Wilson aggregate strong
card, all-order strong MTA, WUFC, CPM, cutoff removal, reflection
positivity, Osterwalder--Schrader reconstruction, a continuum
Yang--Mills measure, or a positive Hamiltonian gap.  The full
Yang--Mills mass gap has not been proved.
\end{firewall}

\begin{theorem}[V93 result ledger]
\label{thm:newV93-results}
V93:
\begin{enumerate}[label=\textup{(V93.\arabic*)},leftmargin=6.5em]
\item proves intermediate-wedge Duhamel contraction;
\item proves the differentiated primitive-root reserve;
\item proves repeated supports are necessarily cyclic;
\item proves the V77 reserve hypertree theorem at arbitrary \(B\);
\item closes PRAH;
\item closes acyclic TIC;
\item classifies the already closed cyclic sectors; and
\item formulates TCNC and proves its sufficiency.
\end{enumerate}
\end{theorem}

\begin{proof}
Use
\Cref{lem:newV93-Duhamel,
lem:newV93-gate-reserve,
lem:newV93-V77-uniform,
thm:newV93-PRAH,
cor:newV93-acyclic-TIC,
lem:newV93-closed-cycles,
cor:newV93-first-cycle,
prop:newV93-TCNC}.
\end{proof}

\paragraph{Parent-version record.}
V93 was formed from
\texttt{Renewal\_Yang\_Mills\_Mass\_Gap\_Research\_Notes\_V92.tex}.
The V92 parent had \(41\,819\) logical source lines,
\(1\,351\,944\) bytes, and SHA--256
\[
 \texttt{7dd3394106fba96942c280cb1983ce0e5c79bc1523f4334e4a1e6aff875975ac}.
\]
The next compulsory companion audit is V95.

% =============================================================================
% END APPENDED FIFTH-SERIES V93 MATERIAL
% =============================================================================

\clearpage
\chapter{Fifth-series V94: fusion of a locally overlapped
transversal four-cycle}
\label{ch:newV94-fusion}

\section{A two-root union--intersection inequality}
\label{sec:newV94-inequality}

No companion audit is due at V94.

For finite supports \(A,B\), put
\[
 a:=|A|,\qquad b:=|B|,\qquad
 I:=A\cap B,\quad h:=|I|\ge2,\qquad
 C:=A\cup B,\quad c:=|C|=a+b-h.
\]
Write \(\mathfrak r\) and \(\mathfrak q\) for the reserve and strong
proxies, with their numerical bases enlarged uniformly when
needed.

\begin{lemma}[Entropy inequality for union and intersection]
\label{lem:newV94-entropy}
For integers \(a,b,h,c\) as above,
\begin{equation}
 a^ab^b\le h^hc^c.                                         \label{eq:newV94-entropy}
\end{equation}
\end{lemma}

\begin{proof}
Write \(a=h+x\), \(b=h+y\), so \(c=h+x+y\), with
\(x,y\ge0\).  Let \(f(t)=t\log t\).  Since \(f'\) is increasing,
\[
 f(h+x+y)-f(h+x)
 =
 \int_0^y f'(h+x+t)\,dt
 \ge
 \int_0^y f'(h+t)\,dt
 =
 f(h+y)-f(h).
\]
Rearrange and exponentiate.
\end{proof}

\begin{lemma}[Two-root fusion]
\label{lem:newV94-fusion}
\begin{equation}
 \mathfrak r(A)\mathfrak r(B)
 \le \mathfrak r(C)\mathfrak r(I).                         \label{eq:newV94-fusion}
\end{equation}
Thus both the union and intersection retain the clock reserve, with
the same numerical base.
\end{lemma}

\begin{proof}
The numerical bases cancel because \(a+b=c+h\).  The left side is
\[
 a^ab^b
 s^{(a+b-4)/2}
 \left(\prod_{i\in C}b_i\right)
 \left(\prod_{i\in I}b_i\right).
\]
The product \(\mathfrak r(C)\mathfrak r(I)\) is
\[
 c^ch^h
 s^{(c+h-4)/2}
 \left(\prod_{i\in C}b_i\right)
 \left(\prod_{i\in I}b_i\right).
\]
The clock exponents agree because \(a+b=c+h\), and the remaining
arity inequality is exactly \Cref{lem:newV94-entropy}.
\end{proof}

\begin{corollary}[The residual overlap has a strong proxy]
\label{cor:newV94-overlap-strong}
For \(h\ge2\),
\begin{equation}
 \mathfrak r(I)
 \le K^h\mathfrak q(I).                                    \label{eq:newV94-overlap-strong}
\end{equation}
\end{corollary}

\begin{proof}
Use \(B_I\ge h\sqrt s\):
\[
 h^hs^{(h-2)/2}
 \le h^2B_I^{h-2}\le4^hB_I^{h-2}.
\]
Restore the common product of marks.
\end{proof}

\begin{corollary}[Optional strong form on the union]
\label{cor:newV94-union-strong}
After enlarging the numerical base,
\begin{equation}
 \mathfrak r(A)\mathfrak r(B)
 \le K^{a+b}\mathfrak q(C)\mathfrak r(I).                  \label{eq:newV94-union-strong}
\end{equation}
\end{corollary}

\begin{proof}
The floor \(B_C\ge c\sqrt s\) gives
\(\mathfrak r(C)\le K^c\mathfrak q(C)\), exactly as in
\Cref{cor:newV94-overlap-strong}.  Apply
\Cref{lem:newV94-fusion}.
\end{proof}

\begin{assessment}[Thermal degree is conserved]
\label{ass:newV94-degree}
The two input supports have total size \(a+b\), while the union and
intersection have total size \(c+h=a+b\).  Thus the fusion neither
creates nor discards an arity exponent or a clock half-power.  The
union remains eligible for the V77 reserve hypertree sum, while the
intersection is a residual overlap activity.
\end{assessment}

\section{Local overlaps form a small residual gas}
\label{sec:newV94-local-gas}

\begin{lemma}[All residual overlaps in one light block]
\label{lem:newV94-local-overlaps}
Let \(P\) be light.  Then
\begin{equation}
 \sum_{\substack{I\subseteq P\\|I|\ge2}}\mathfrak r(I)
 \le Cw_P^2.                                                \label{eq:newV94-local-overlaps}
\end{equation}
An arbitrary ordered list of such residual overlaps has total
weight at most a numerical constant, after \(b_*\) is decreased.
\end{lemma}

\begin{proof}
By \Cref{cor:newV94-overlap-strong}, enlarge the base and dominate
each reserve activity by its strong proxy.  The whole-hyperedge
gas estimate \Cref{lem:newV78-whole-gas}, applied on \(P\), gives
\eqref{eq:newV94-local-overlaps}.

For an ordered list, positivity gives the geometric bound
\[
 \sum_{\ell\ge0}
 \left(
  \sum_{\substack{I\subseteq P\\|I|\ge2}}\mathfrak r(I)
 \right)^\ell
 \le {1\over1-Cw_P^2}\le2
\]
after \(b_*\) is chosen so that \(C(b_*/2)^2\le1/2\).
For exponentially normalized repeated supports, the smaller
exponential gas follows from V86--V87.
\end{proof}

\begin{corollary}[All blockwise residual gases cost \(K^m\)]
\label{cor:newV94-global-local-gas}
The product of the ordered residual-overlap gas bounds over the
light blocks is at most \(K^m\).
\end{corollary}

\begin{proof}
There are at most \(m\) blocks, and each factor is at most two.
\end{proof}

\section{The locally overlapped fundamental cycle}
\label{sec:newV94-class}

\begin{definition}[Locally overlapped fundamental-cycle class,
LOFC]
\label{def:newV94-LOFC}
LOFC consists of connected primitive ancestries with a canonical
distinguished pair of transversal supports \(A,B\) such that:
\begin{enumerate}[label=\textup{(\roman*)},leftmargin=3.5em]
\item \(I=A\cap B\) has at least two labels and is contained in one
      light block;
\item replacing the two occurrence vertices \(A,B\) by one formal
      union occurrence \(C=A\cup B\) produces a primitive incidence
      hypertree \({\cal H}_0\); and
\item \(C\) is not already the support of another occurrence in
      \({\cal H}_0\).
\end{enumerate}
The distinguished pair is chosen by a fixed lexicographic rule.
\end{definition}

\begin{lemma}[Fusion preserves connected support and removes the
fundamental excess]
\label{lem:newV94-support-fusion}
For an LOFC ancestry, \({\cal H}_0\) is connected on the same label
set.  The replacement removes \(h-1\) units of primitive incidence
excess:
\begin{equation}
 g({\cal H})-g({\cal H}_0)=h-1.                             \label{eq:newV94-excess-drop}
\end{equation}
\end{lemma}

\begin{proof}
Every incidence formerly attached to \(A\) or \(B\) is attached to
their union \(C\), so no path is lost.  The contribution of the two
old hyperedges to
\(\sum_D(|D|-1)\) is \(a+b-2\); the contribution of \(C\) is
\[
 c-1=a+b-h-1.
\]
Their difference is \(h-1\).
\end{proof}

\begin{lemma}[Inverse split fibre]
\label{lem:newV94-split-fibre}
For fixed union \(C\) and intersection \(I\), the number of ordered
pairs \((A,B)\) satisfying
\[
 A\cup B=C,\qquad A\cap B=I
\]
is at most \(2^{|C|}\).
\end{lemma}

\begin{proof}
Every label of \(C\setminus I\) belongs to exactly one of
\(A\setminus I\) and \(B\setminus I\).  Its two choices determine
the ordered pair.  Restrictions that both roots be transversal or
distinct only reduce the number.
\end{proof}

\begin{theorem}[LOFC strong aggregate card]
\label{thm:newV94-LOFC}
On the intermediate wedge, the complete positive majorant of LOFC
obeys
\begin{equation}
 K^m\left(\prod_{i=1}^mb_i\right)B^{m-2}.                   \label{eq:newV94-LOFC}
\end{equation}
\end{theorem}

\begin{proof}
Apply the contractive Duhamel bound
\Cref{lem:newV93-Duhamel} and the primitive reserve
\Cref{lem:newV93-gate-reserve}.  On the distinguished pair, use
\Cref{lem:newV94-fusion}.  This replaces its positive weight by a
reserve union-root factor on \(C\) and a residual reserve factor on
\(I\).

By definition, the fused ancestry \({\cal H}_0\) is an incidence
hypertree, and every one of its primitive activities, including
the union support \(C\), retains the reserve.  Apply
\Cref{thm:newV93-PRAH} to obtain the global strong tree sum.

Sum \(I\) inside its light block by
\Cref{lem:newV94-local-overlaps}; this factor is numerical.
The inverse split costs at most \(2^{|C|}\) by
\Cref{lem:newV94-split-fibre}.  Choices of the canonical pair,
signs, and owned derivative edges are polynomial or exponential in
\(m\) and enter \(K^m\).  The final global Cayley sum is
\((\prod_i b_i)B^{m-2}\).
\end{proof}

\begin{corollary}[The first locally overlapped transversal cycle is
closed]
\label{cor:newV94-first-cycle}
Every primitive four-cycle whose two root vertices have a local
overlap and whose fusion leaves an incidence hypertree satisfies
TCNC.
\end{corollary}

\begin{proof}
It is LOFC; apply \Cref{thm:newV94-LOFC}.
\end{proof}

\begin{assessment}[What fusion did and did not prove]
\label{ass:newV94-fusion-scope}
Fusion closes a genuine transversal incidence cycle without using
the global \(B^{2g}\) factor.  Its residual cycle degree is localized
to \(I\) and summed at the light-block scale.  The proof does not
cover an overlap \(I\) meeting two or more thermal blocks, nor a
fusion whose union collides with another occupied support or leaves
a further transversal cycle.
\end{assessment}

\section{The remaining transversal-overlap card}
\label{sec:newV94-remaining}

\begin{definition}[Transversal-overlap contraction card, TOCC]
\label{def:newV94-TOCC}
TOCC is the remaining TCNC operation in which the canonical
cycle pair \(A,B\) has \(I=A\cap B\) meeting at least two thermal
blocks, or its union \(C=A\cup B\) remains in a positive-excess
transversal core.  It must:
\begin{enumerate}[label=\textup{(\roman*)},leftmargin=3.5em]
\item retain the union--intersection fusion
      \eqref{eq:newV94-fusion};
\item treat the residual \(I\) as a smaller transversal root,
      not as a light-block gas;
\item contract a monotone complexity, such as the ordered pair
      \((g,\sum |A|)\), without an arity-dependent inverse fibre;
\item preserve the normalized percolation/Kruskal factors; and
\item terminate in PRAH, OTRAC, or LOFC with total cost \(K^m\).
\end{enumerate}
\end{definition}

\begin{proposition}[TOCC completes TCNC]
\label{prop:newV94-TOCC}
TOCC, together with
\Cref{cor:newV93-acyclic-TIC,cor:newV94-first-cycle}, proves TCNC,
TIC, HLTC, full MCWC, and SCNGC.
\end{proposition}

\begin{proof}
The canonical primitive cycle has either local or transversal
overlap.  A local overlap whose fusion ends at a hypertree is LOFC.
All other positive-excess cases are TOCC.  The zero-excess case is
acyclic TIC.  Thus TCNC holds; apply
\Cref{prop:newV93-TCNC}.
\end{proof}

\begin{programme}[V95: compulsory audit and transversal overlap]
\label{prog:newV94-V95}
\begin{enumerate}[leftmargin=3em]
\item Perform the compulsory newest-SM/newest-NS audit.
\item Test fusion on the minimal transversal overlap consisting of
      two pair-block paths sharing two thermal blocks.
\item Order the pair \((C,I)\) by support span and prove whether
      replacing \(A,B\) by \(C,I\) decreases a lexicographic
      cycle complexity after the residual is normalized.
\item Use the V83 Kruskal identity to sum choices when \(I\) follows
      an already open block shortcut.
\item If fusion can cycle between union and intersection, return to
      the exact percolation quotient before applying
      \eqref{eq:newV94-fusion}.
\end{enumerate}
\end{programme}

\begin{firewall}[Claim boundary after V94]
\label{fire:newV94-boundary}
V94 proves the union--intersection entropy inequality, the
two-root fusion estimate, strong domination of the residual
overlap, and the local residual-overlap gas.  It closes every
locally overlapped fundamental transversal cycle whose fusion
leaves an incidence hypertree, and formulates TOCC for genuinely
transversal overlaps and persistent excess.

V94 does not prove TOCC, full TCNC, cyclic TIC, HLTC, full MCWC,
OCRC, SCNGC, WPSC, GRLTC, LZTC, PLRC, WSDC, BRSC, HWSFC, UGCGC,
thermal connected WPRAC, all-arity RGSC, the raw Wilson aggregate
strong card, all-order strong MTA, WUFC, CPM, cutoff removal,
reflection positivity, Osterwalder--Schrader reconstruction, a
continuum Yang--Mills measure, or a positive Hamiltonian gap.  The
full Yang--Mills mass gap has not been proved.
\end{firewall}

\begin{theorem}[V94 result ledger]
\label{thm:newV94-results}
V94:
\begin{enumerate}[label=\textup{(V94.\arabic*)},leftmargin=6.5em]
\item proves the union--intersection entropy inequality;
\item proves the two-root fusion bound;
\item proves residual reserve-to-strong domination;
\item proves the local overlap gas;
\item proves the excess drop under fusion;
\item proves the inverse split fibre;
\item closes LOFC; and
\item formulates TOCC and proves its sufficiency.
\end{enumerate}
\end{theorem}

\begin{proof}
Use
\Cref{lem:newV94-entropy,
lem:newV94-fusion,
cor:newV94-overlap-strong,
lem:newV94-local-overlaps,
lem:newV94-support-fusion,
lem:newV94-split-fibre,
thm:newV94-LOFC,
prop:newV94-TOCC}.
\end{proof}

\paragraph{Parent-version record.}
V94 was formed from
\texttt{Renewal\_Yang\_Mills\_Mass\_Gap\_Research\_Notes\_V93.tex}.
The V93 parent had \(42\,180\) logical source lines,
\(1\,365\,160\) bytes, and SHA--256
\[
 \texttt{ecb8a8145a4022dc25662956fac775ccf320cb8edd1c0a04dbcececec6cca47f}.
\]
The next compulsory companion audit is V95.

% =============================================================================
% END APPENDED FIFTH-SERIES V94 MATERIAL
% =============================================================================

\clearpage
\chapter{Fifth-series V95: compulsory audit, submacroscopic
transversal overlaps, and two-laminar uncrossing}
\label{ch:newV95-uncrossing}

\section{Compulsory companion audit}
\label{sec:newV95-audit}

The newest settled companion snapshots read for this audit were:
\begin{center}
\small
\begin{tabular}{p{0.24\textwidth}p{0.69\textwidth}}
Standard Model--Einstein &
\texttt{Renewal\_SM\_Einstein\_Continuum\_Research\_Notes\_V20.tex},
\(339\,111\) bytes, \(8\,657\) logical source lines, SHA--256
\texttt{4737F516E6586F92C8651784BAC453085027C134F04A70D5F4D53A02A489CF71}.
\\ [0.4em]
Navier--Stokes &
\texttt{Renewal\_NS\_Stability\_Research\_Notes\_V88.tex},
\(1\,079\,335\) bytes, \(20\,577\) logical source lines, SHA--256
\texttt{C3466AB4BB2B58F2EFA1991F8655C2C75551A6831441BC5B4A63B38253F61026}.
\end{tabular}
\end{center}

\begin{audit}[Standard Model--Einstein V20]
\label{audit:newV95-SM}
SM V20 keeps the exact factor \(1-a^2m^2/4\) inside the
massive-overlap square until the complete lattice quadratic
response has been subtracted.  Only after that exact subtraction
does it identify the quartic continuum coefficient.  The
transferable rule is structural: a normalization must remain
attached to the full object whose lower-order contribution it
cancels.  Here this confirms that the union--intersection reserve
of \Cref{lem:newV94-fusion} and the closed-edge Bernoulli factors
of \Cref{warn:newV83-normalization} must be retained through the
complete overlap sum.  No SM heat-kernel coefficient is imported.
\end{audit}

\begin{audit}[Navier--Stokes V88]
\label{audit:newV95-NS}
NS V88 records an audited theorem manifest and a separate import
ledger after re-deriving its preceding blocks.  In particular, it
does not promote exact integral identities or termwise profile
facts to a global regularity estimate.  The corresponding warning
here is that a termwise uncrossing inequality is not an aggregate
cycle bound until its inverse fibres have been summed.  No
Navier--Stokes estimate is imported.
\end{audit}

\begin{audit}[Direction selected by the audit]
\label{audit:newV95-direction}
V95 therefore keeps three ledgers separate:
\begin{enumerate}[label=\textup{(\roman*)},leftmargin=3.5em]
\item the termwise reserve inequality;
\item the sum over overlap supports and inverse splits; and
\item the normalized percolation/Kruskal configuration sum.
\end{enumerate}
The first two can be closed when the overlap arity is
submacroscopic.  The remaining high-arity family is then
uncrossed without claiming that its aggregate fibre has already
been controlled.
\end{audit}

\section{A global overlap series at fixed arity}
\label{sec:newV95-overlap-series}

Let \(C\) be a support of size at most \(m\), put
\(B_C=\sum_{i\in C}b_i\), and recall that \(0<b_i<1\).
For \(h\ge2\), let
\[
 {\cal R}_h(C)
 :=
 \sum_{\substack{I\subseteq C\\|I|=h}}\mathfrak r(I).
\]

\begin{lemma}[Exact elementary-symmetric reduction]
\label{lem:newV95-elementary}
After fixing the reserve base \(K_0\),
\begin{equation}
 {\cal R}_h(C)
 =
 (K_0h)^h s^{(h-2)/2}e_h((b_i)_{i\in C})
 \le
 (eK_0)^h s^{(h-2)/2}B_C^h.                \label{eq:newV95-elementary}
\end{equation}
\end{lemma}

\begin{proof}
The equality is the definition of the reserve activity summed
over all \(h\)-subsets.  Expanding \(B_C^h\), every squarefree
monomial \(\prod_{i\in I}b_i\) occurs \(h!\) times, so
\[
 e_h((b_i)_{i\in C})\le {B_C^h\over h!}.
\]
The elementary bound \(h!\ge(h/e)^h\) gives
\((K_0h)^h/h!\le(eK_0)^h\).
\end{proof}

\begin{definition}[Submacroscopic overlap threshold]
\label{def:newV95-Hm}
For \(m\ge2\), set
\begin{equation}
 H_m:=\max\left\{2,\left\lfloor{m\over\log(m+e)}\right\rfloor\right\}.
                                                        \label{eq:newV95-Hm}
\end{equation}
\end{definition}

\begin{lemma}[Submacroscopic overlaps cost only an exponential
base]
\label{lem:newV95-submacro-gas}
There is a numerical \(K\), independent of \(m,s,C\), such that
\begin{equation}
 \sum_{h=2}^{\min\{H_m,|C|\}}{\cal R}_h(C)\le K^m.
                                                        \label{eq:newV95-submacro-gas}
\end{equation}
\end{lemma}

\begin{proof}
The clock regime has \(0<s\le1\), so the clock factor in
\eqref{eq:newV95-elementary} is at most one.  Also \(B_C\le|C|\le
m\).  Hence the left side is at most
\[
 \sum_{h=2}^{H_m}(eK_0m)^h.
\]
The terms with \(m\) in a fixed bounded range are absorbed by
enlarging \(K\).  For all remaining \(m\),
\[
 H_m\log(eK_0m)
 \le
 {m\log(eK_0m)\over\log(m+e)}
 \le c_0m
\]
with a numerical \(c_0\).  There are at most \(m\) summands, and
\(m\le2^m\).  This proves \eqref{eq:newV95-submacro-gas}.
\end{proof}

\begin{remark}[The two-label transversal overlap]
\label{rem:newV95-two-label}
For \(h=2\), the estimate is particularly transparent:
\[
 {\cal R}_2(C)
 \le K^2\sum_{\{i,j\}\subseteq C}b_ib_j
 \le {K^2\over2}B_C^2
 \le K^m.
\]
Thus an overlap consisting of one label in each of two thermal
blocks is not itself a global \(B^2\) obstruction when it occurs
only once.  The obstruction in \Cref{warn:newV93-no-global-V78}
concerns an arbitrary gas of surplus overlaps.
\end{remark}

\section{Fundamental cycles with submacroscopic overlap}
\label{sec:newV95-SOFC}

\begin{definition}[Submacroscopic-overlap fundamental-cycle
class, SOFC]
\label{def:newV95-SOFC}
SOFC is the class obtained from LOFC in
\Cref{def:newV94-LOFC} by replacing condition (i) with
\[
 2\le |A\cap B|\le H_m,
\]
without requiring \(A\cap B\) to lie in one thermal block.
Conditions (ii)--(iii), including the canonical distinguished
pair and the requirement that its union replacement be a
primitive incidence hypertree, are unchanged.
\end{definition}

\begin{theorem}[SOFC strong aggregate card]
\label{thm:newV95-SOFC}
On the intermediate wedge, the complete positive majorant of
SOFC obeys
\begin{equation}
 K^m\left(\prod_{i=1}^mb_i\right)B^{m-2}.                 \label{eq:newV95-SOFC}
\end{equation}
\end{theorem}

\begin{proof}
Fix the fused support \(C=A\cup B\).  Use the contractive
Duhamel estimate and primitive reserve exactly as in
\Cref{thm:newV94-LOFC}.  By
\Cref{lem:newV94-fusion}, the distinguished pair is bounded by
\(\mathfrak r(C)\mathfrak r(I)\), where \(I=A\cap B\).
The union-replaced ancestry is a primitive incidence hypertree,
so its complete sum, with \(C\) carrying its reserve, is bounded
by \Cref{thm:newV93-PRAH}.

For fixed \(C,I\), the inverse ordered split costs at most
\(2^{|C|}\le2^m\) by \Cref{lem:newV94-split-fibre}.  Now sum all
permitted residual intersections \(I\subseteq C\) by
\Cref{lem:newV95-submacro-gas}.  Both costs enter the exponential
base.  The distinguished pair is canonical, while the number of
root occurrences in the fused hypertree is at most \(m-1\);
all remaining finite sign, port, and occurrence choices are also
bounded by \(K^m\).  The PRAH output is the right side of
\eqref{eq:newV95-SOFC}.
\end{proof}

\begin{corollary}[The minimal genuinely transversal four-cycle is
closed]
\label{cor:newV95-minimal-transversal}
Suppose two transversal root occurrences share exactly two
labels in distinct thermal blocks, and replacing them by their
union leaves a primitive incidence hypertree with no union-support
collision.  Then the ancestry satisfies TCNC.
\end{corollary}

\begin{proof}
Its overlap has arity two and hence is at most \(H_m\).  It is an
SOFC ancestry, so apply \Cref{thm:newV95-SOFC}.
\end{proof}

\begin{assessment}[New location of the fundamental-cycle
bottleneck]
\label{ass:newV95-SOFC-scope}
For a single fundamental fusion, neither locality of the overlap
nor a Kruskal shortcut factor is needed below arity \(H_m\).
Only overlaps with
\[
 |A\cap B|>H_m
\]
remain in this row.  They contain a macroscopic amount of
incidence data relative to the exponential ledger, even though
their cardinality may still be \(o(m)\).
\end{assessment}

\section{Reserve-preserving two-laminar uncrossing}
\label{sec:newV95-laminar}

Two supports \(A,B\) are called \emph{properly two-crossing} when
\[
 |A\cap B|\ge2,\qquad A\setminus B\ne\varnothing,\qquad
 B\setminus A\ne\varnothing.
\]
A multiset of supports is \emph{two-laminar} when it has no
properly two-crossing pair.  Thus any two of its supports are
nested or meet in at most one label.

\begin{lemma}[Uncrossing invariants and strict potential gain]
\label{lem:newV95-potential}
Let a finite multiset \({\cal F}\) contain a properly two-crossing
pair \(A,B\), and form
\[
 {\cal F}'
 =
 ({\cal F}\setminus\{A,B\})\cup\{A\cap B,A\cup B\},
\]
with multiplicities.  Then:
\begin{enumerate}[label=\textup{(\roman*)},leftmargin=3.5em]
\item for every label \(i\), its incidence degree
      \(d_i({\cal F})=\#\{D\in{\cal F}:i\in D\}\) is unchanged;
\item the number of support occurrences, the total incidence
      count, and the incidence excess are unchanged;
\item
\begin{equation}
 \prod_{D\in{\cal F}}\mathfrak r(D)
 \le
 \prod_{D\in{\cal F}'}\mathfrak r(D);                    \label{eq:newV95-product-uncross}
\end{equation}
and
\item for \(\Psi({\cal F})=\sum_{D\in{\cal F}}|D|^2\),
\begin{equation}
 \Psi({\cal F}')-\Psi({\cal F})
 =
 2|A\setminus B|\,|B\setminus A|>0.                       \label{eq:newV95-potential}
\end{equation}
\end{enumerate}
\end{lemma}

\begin{proof}
The pointwise indicator identity
\[
 {\bf1}_A+{\bf1}_B
 =
 {\bf1}_{A\cap B}+{\bf1}_{A\cup B}
\]
proves (i), and hence preserves the total incidence count.
There are two supports before and after, so the incidence excess
\(\sum_D(|D|-1)-|X|+1\) is also unchanged.  This proves (ii).
Item (iii) is \Cref{lem:newV94-fusion}, with all other factors
cancelled.

For (iv), write \(h=|A\cap B|\),
\(x=|A\setminus B|>0\), and \(y=|B\setminus A|>0\).
The changed part of the potential is
\[
 h^2+(h+x+y)^2-(h+x)^2-(h+y)^2=2xy.
\]
\end{proof}

\begin{theorem}[Finite two-laminar normal form]
\label{thm:newV95-laminar}
Every finite multiset \({\cal F}\) of nonsingleton supports admits
a finite sequence of proper two-uncrossings ending in a
two-laminar multiset \({\cal L}\).  Along the sequence the label
degrees and incidence excess are fixed and
\begin{equation}
 \prod_{D\in{\cal F}}\mathfrak r(D)
 \le
 \prod_{D\in{\cal L}}\mathfrak r(D).                      \label{eq:newV95-laminar}
\end{equation}
If \({\cal F}\) has \(q\) occurrences on \(m\) labels, every such
sequence has fewer than \(qm^2/2+1\) steps.
\end{theorem}

\begin{proof}
Repeatedly uncross any properly two-crossing pair.  All new
supports remain nonsingleton because their intersection has at
least two labels.  By \Cref{lem:newV95-potential}, the
integer-valued potential strictly increases by at least two at
each step, while
\[
 \Psi\le qm^2.
\]
Thus the procedure terminates in fewer than \(qm^2/2+1\) steps.
At termination no properly two-crossing pair remains, which is
exactly two-laminarity.  The invariant statements and
\eqref{eq:newV95-laminar} follow by iterating
\Cref{lem:newV95-potential}.
\end{proof}

\begin{corollary}[Geometry of the unresolved high-overlap sector]
\label{cor:newV95-high-laminar}
Term by term, every unresolved family of high-arity transversal
overlaps is reserve-dominated by a two-laminar family with the
same incidence degrees and the same positive excess.
\end{corollary}

\begin{proof}
Apply \Cref{thm:newV95-laminar}.  Supports intersecting in at most
one label create no bipartite incidence four-cycle; every
remaining multi-incidence relation is nested.
\end{proof}

\begin{warning}[Termwise uncrossing is not the aggregate card]
\label{warn:newV95-fibre}
Theorem~\ref{thm:newV95-laminar} does not prove TOCC.  A sequence
may have \(O(qm^2)\) steps, and multiplying the one-step inverse
split bound \(2^m\) over those steps is superexponential.  Nor
does the V83 Kruskal identity repair this by itself: that identity
sums input spanning trees for one fixed open shortcut graph,
whereas uncrossing changes the primitive support family.  The
normalized Bernoulli configuration and the support fibre must be
summed jointly before either is discarded.
\end{warning}

\section{The remaining high-overlap nested card}
\label{sec:newV95-HONC}

\begin{definition}[High-overlap nested-chain card, HONC]
\label{def:newV95-HONC}
HONC is the following remaining operation.
\begin{enumerate}[label=\textup{(\roman*)},leftmargin=3.5em]
\item Start with the canonical positive-excess transversal core
      after local CIRC and SOFC have been removed.
\item Retain the complete normalized Bernoulli configuration and
      its V83 Kruskal allocation.
\item Sum all primitive support families having the same
      incidence-degree vector before applying the two-laminar
      domination.
\item Bound the resulting nested-support fibre by \(K^m\), using
      the degree vector rather than a history of inverse splits.
\item End with the PRAH strong polynomial.
\end{enumerate}
\end{definition}

\begin{proposition}[HONC completes the fundamental overlap row]
\label{prop:newV95-HONC}
HONC, together with \Cref{thm:newV95-SOFC}, closes every
fundamental two-root transversal overlap whose union replacement
is a primitive incidence hypertree and does not collide with an
occupied support.
\end{proposition}

\begin{proof}
The overlap arity is at most \(H_m\) or larger.  The first case is
SOFC.  In the second, \Cref{cor:newV95-high-laminar} gives the
termwise nested geometry, and HONC supplies precisely the missing
aggregate fibre estimate.
\end{proof}

\begin{programme}[V96: degree-vector fibre before nested chains]
\label{prog:newV95-V96}
\begin{enumerate}[leftmargin=3em]
\item For \(q\) distinguishable root occurrences, encode a support
      family by its \(m\times q\) zero--one incidence matrix.
\item Factor out the incidence-hypertree owner selected by PRAH
      before counting any surplus row choices.
\item Derive the exact number of row patterns with prescribed
      degrees and test it against the normalized percolation
      probabilities.
\item Prove the complete lattice-sorting identity for supports
      only after the fibre is expressed by degree vectors.
\item If the unrestricted degree fibre is superexponential,
      identify the missing Duhamel or exponential-generating
      denominator rather than weakening the target.
\end{enumerate}
\end{programme}

\begin{firewall}[Claim boundary after V95]
\label{fire:newV95-boundary}
V95 audits SM V20 and NS V88, proves the exact fixed-arity overlap
series, proves that all overlap arities through
\(H_m=O(m/\log m)\) cost only \(K^m\), and closes the corresponding
submacroscopic fundamental transversal cycles without a locality
assumption.  It also proves finite reserve-preserving
two-laminar uncrossing with exact incidence-degree and excess
invariants.

V95 does not prove the aggregate fibre in
\Cref{warn:newV95-fibre}, HONC, all fundamental high-overlap
cycles, support collisions, persistent-excess TOCC, full TCNC,
cyclic TIC, HLTC, full MCWC, OCRC, SCNGC, WPSC, GRLTC, LZTC,
PLRC, WSDC, BRSC, HWSFC, UGCGC, thermal connected WPRAC,
all-arity RGSC, the raw Wilson aggregate strong card, all-order
strong MTA, WUFC, CPM, cutoff removal, reflection positivity,
Osterwalder--Schrader reconstruction, a continuum Yang--Mills
measure, or a positive Hamiltonian gap.  The full Yang--Mills
mass gap has not been proved.
\end{firewall}

\begin{theorem}[V95 result ledger]
\label{thm:newV95-results}
V95:
\begin{enumerate}[label=\textup{(V95.\arabic*)},leftmargin=6.5em]
\item performs the compulsory SM/NS audit;
\item proves the fixed-arity overlap series estimate;
\item proves the submacroscopic overlap gas;
\item closes SOFC and the minimal transversal four-cycle;
\item proves the uncrossing invariants and strict potential gain;
\item proves the finite two-laminar normal form;
\item isolates the aggregate fibre obstruction; and
\item formulates HONC.
\end{enumerate}
\end{theorem}

\begin{proof}
Use
\Cref{audit:newV95-SM,
audit:newV95-NS,
lem:newV95-elementary,
lem:newV95-submacro-gas,
thm:newV95-SOFC,
cor:newV95-minimal-transversal,
lem:newV95-potential,
thm:newV95-laminar,
warn:newV95-fibre,
prop:newV95-HONC}.
\end{proof}

\paragraph{Parent-version record.}
V95 was formed from
\texttt{Renewal\_Yang\_Mills\_Mass\_Gap\_Research\_Notes\_V94.tex}.
The V94 parent had \(42\,562\) logical source lines,
\(1\,377\,918\) bytes, and SHA--256
\[
 \texttt{111e389c2fe4ca66d3688b642dc6f3c5ec5b18c70f007fd7bad580258654e345}.
\]
The next compulsory companion audit is V100.

% =============================================================================
% END APPENDED FIFTH-SERIES V95 MATERIAL
% =============================================================================

\clearpage
\chapter{Fifth-series V96: a submacroscopic-excess theorem and
the dense degree-fibre obstruction}
\label{ch:newV96-excess}

\section{The canonical skeleton is exactly excess-graded}
\label{sec:newV96-grading}

No companion audit is due at V96.

Let \(X\) have \(m\ge2\) labels, total mark \(B\), and let
\({\cal H}\) be a connected simple hypergraph of primitive root
supports.  Give every support \(A\) an activity \(u_A\) obeying
\begin{equation}
 0\le u_A\le\min\{\mathfrak r(A),\mathfrak q(A)\},          \label{eq:newV96-two-proxy}
\end{equation}
with the reserve and strong proxies of V86.  Apply the canonical
incidence-skeleton injection
\Cref{lem:newV78-skeleton}:
\[
 {\cal H}\longmapsto
 \left({\cal H}_0,(J_C)_{C\in{\cal H}_0},{\cal F}\right).
\]

\begin{lemma}[Exact surplus-grade identity]
\label{lem:newV96-grade}
Define
\begin{equation}
 \kappa
 :=
 \sum_{C\in{\cal H}_0}|J_C|
 +\sum_{A\in{\cal F}}(|A|-1).                              \label{eq:newV96-grade}
\end{equation}
Then
\begin{equation}
 \kappa=g({\cal H}).                                       \label{eq:newV96-grade-equals}
\end{equation}
\end{lemma}

\begin{proof}
The skeleton is an incidence hypertree on \(X\), so
\[
 \sum_{C\in{\cal H}_0}(|C|-1)=m-1.
\]
Replacing \(C\) by \(C\cup J_C\) adds \(|J_C|\) incidences
without adding a support vertex.  Adding a whole surplus support
\(A\) adds \(|A|\) incidences and one support vertex, hence
\(|A|-1\) units of excess.  Substitution into
\[
 g({\cal H})=\sum_{A\in{\cal H}}(|A|-1)-m+1
\]
gives \eqref{eq:newV96-grade-equals}.
\end{proof}

\section{Coefficientwise extension and whole-support series}
\label{sec:newV96-series}

\begin{lemma}[One skeleton support, graded]
\label{lem:newV96-one-extension}
For every retained support \(C\),
\begin{equation}
 \sum_{J\subseteq X\setminus C}
 {\mathfrak r(C\cup J)\over\mathfrak r(C)}t^{|J|}
 \preceq
 \exp(KB^2t),                                              \label{eq:newV96-one-extension}
\end{equation}
where \(\preceq\) denotes coefficientwise domination of formal
power series with nonnegative coefficients.
\end{lemma}

\begin{proof}
The proof of \Cref{lem:newV78-extension}, without setting \(t=1\),
gives for \(|J|=k\)
\[
 {\mathfrak r(C\cup J)\over\mathfrak r(C)}
 \le(eKm\sqrt s)^k\prod_{j\in J}b_j.
\]
Summing the \(k\)-subsets and using
\(e_k((b_j)_{j\notin C})\le B^k/k!\) gives
\[
 [t^k]\text{(left side)}
 \le{(eKm\sqrt s\,B)^k\over k!}
 \le{(eKB^2)^k\over k!},
\]
because \(m\sqrt s\le B\).  Enlarge \(K\).
\end{proof}

\begin{corollary}[All retained-support extensions]
\label{cor:newV96-all-extensions}
For a fixed skeleton \({\cal H}_0\),
\begin{equation}
 \prod_{C\in{\cal H}_0}
 \left\{
  \sum_{J\subseteq X\setminus C}
  {\mathfrak r(C\cup J)\over\mathfrak r(C)}t^{|J|}
 \right\}
 \preceq
 \exp(KmB^2t).                                             \label{eq:newV96-all-extensions}
\end{equation}
Consequently its coefficient of grade \(k\) is at most
\((KmB^2)^k/k!\).
\end{corollary}

\begin{proof}
There are at most \(m-1\) supports in an incidence hypertree.
Multiply \Cref{lem:newV96-one-extension}.
\end{proof}

\begin{lemma}[Graded whole-surplus-support gas]
\label{lem:newV96-whole-graded}
The formal series over sets of whole surplus supports satisfies
\begin{equation}
 [t^\ell]\!
 \prod_{\substack{A\subseteq X\\|A|\ge2}}
 \left(1+\mathfrak q(A)t^{|A|-1}\right)
 \le (KB^2)^\ell K^\ell
 \qquad(\ell\ge0).                                        \label{eq:newV96-whole-graded}
\end{equation}
\end{lemma}

\begin{proof}
For \(d\ge2\), the elementary-symmetric estimate gives
\begin{equation}
 \sum_{|A|=d}\mathfrak q(A)
 \le {K^dB^{2d-2}\over d!}.                               \label{eq:newV96-qsum}
\end{equation}
Put \(z=K_1B^2t\).  Using \(1+x\le e^x\), the product in
\eqref{eq:newV96-whole-graded} is coefficientwise dominated,
after enlarging \(K_1\), by
\[
 \exp\left\{
  K_1\sum_{j\ge1}{z^j\over(j+1)!}
 \right\}.
\]
In the coefficient of \(z^\ell\), write \(n_j\) for the number
of selected terms of degree \(j\).  Then
\(\sum_jjn_j=\ell\), so
\(\sum_j(j+1)n_j\le2\ell\).  Discarding the factorial
denominators bounds each contribution by \(K_1^{2\ell}\).
The number of integer partitions of \(\ell\) is at most the
number \(2^{\ell-1}\) of compositions.  Hence the coefficient is
at most \(K_2^\ell\), proving the claim after rescaling \(z\).
\end{proof}

\section{All simple cores of submacroscopic excess}
\label{sec:newV96-low-excess}

\begin{theorem}[Fixed-excess simple-core estimate]
\label{thm:newV96-fixed-excess}
For every \(g\ge0\),
\begin{equation}
 \sum_{\substack{{\cal H}\ {\rm connected,\ simple}\\
                  {\cal H}\ {\rm on}\ X,\ g({\cal H})=g}}
 \prod_{A\in{\cal H}}u_A
 \le
 K^m(KmB^2)^g
 \left(\prod_{i\in X}b_i\right)B^{m-2}.                   \label{eq:newV96-fixed-excess}
\end{equation}
\end{theorem}

\begin{proof}
Use the injective skeleton data of
\Cref{lem:newV78-skeleton}.  On each retained skeleton support
use the reserve proxy; on every whole surplus support use the
strong proxy.  By \Cref{lem:newV96-grade}, the coefficient of
total formal grade \(g\) selects exactly the cores of incidence
excess \(g\).

For a fixed skeleton, convolve
\Cref{cor:newV96-all-extensions} with
\Cref{lem:newV96-whole-graded}.  The coefficient of grade \(g\)
is at most
\[
 \sum_{k+\ell=g}{(KmB^2)^k\over k!}(KB^2)^\ell K^\ell
 \le (K_3mB^2)^g.
\]
The remaining skeleton sum is precisely the
reserve-weighted PRAH sum and is bounded by
\Cref{thm:newV93-PRAH}.  Absorb the convolution and gate bases
into \(K\).
\end{proof}

\begin{definition}[Submacroscopic excess threshold]
\label{def:newV96-Gm}
Set
\begin{equation}
 G_m:=\left\lfloor{m\over\log(m+e)}\right\rfloor.          \label{eq:newV96-Gm}
\end{equation}
\end{definition}

\begin{theorem}[Submacroscopic-excess TCNC for simple supports]
\label{thm:newV96-low-excess}
On the intermediate wedge, the complete positive majorant of all
connected simple primitive support hypergraphs satisfying
\(0\le g({\cal H})\le G_m\) obeys
\begin{equation}
 K^m\left(\prod_{i=1}^mb_i\right)B^{m-2}.                 \label{eq:newV96-low-excess}
\end{equation}
\end{theorem}

\begin{proof}
On the intermediate wedge, \(B<m\).  Therefore, for
\(g\le G_m\),
\[
 (KmB^2)^g
 \le(Km^3)^g
 =
 \exp\{g\log(Km^3)\}
 \le K_1^m,
\]
because
\[
 {m\log(Km^3)\over\log(m+e)}\le c_0m.
\]
There are at most \(m\) grades in the sum.  Apply
\Cref{thm:newV96-fixed-excess} and absorb this last factor into
the exponential base.
\end{proof}

\begin{corollary}[Location of the remaining simple cyclic sector]
\label{cor:newV96-high-excess}
Every unresolved simple-support TIC ancestry has
\begin{equation}
 g({\cal H})>G_m,                                         \label{eq:newV96-high-excess}
\end{equation}
after the previously closed local, one-transversal-root, and
acyclic rows have been removed.
\end{corollary}

\begin{proof}
The rows of excess at most \(G_m\) are
\Cref{thm:newV96-low-excess}; the earlier exclusions are
\Cref{lem:newV93-closed-cycles}.
\end{proof}

\begin{assessment}[Why this is stronger than the V95 fundamental
row]
\label{ass:newV96-strength}
Theorem~\ref{thm:newV96-low-excess} permits arbitrary persistent
cycles and support collisions, provided the primitive supports
are distinct and the total incidence excess is at most \(G_m\).
It does not choose a distinguished pair and does not require one
fusion to leave a hypertree.  The high-excess condition
\eqref{eq:newV96-high-excess} is now compulsory for every
remaining simple core.
\end{assessment}

\section{Exact degree fibres}
\label{sec:newV96-degree}

For \(q\) distinguishable root occurrences, encode their supports
by an \(m\times q\) zero--one matrix \(M\), with
\(M_{ij}=1\) when label \(i\) belongs to root \(j\).  Its row sums
are the label-incidence degrees \(d_i\).

\begin{lemma}[Degree-fibre census]
\label{lem:newV96-degree-census}
Without column restrictions, the number of incidence matrices
with prescribed row sums \(0\le d_i\le q\) is exactly
\begin{equation}
 N(\boldsymbol d)=\prod_{i=1}^m{q\choose d_i}.             \label{eq:newV96-degree-census}
\end{equation}
\end{lemma}

\begin{proof}
In row \(i\), choose independently the \(d_i\) root columns that
contain label \(i\).
\end{proof}

\begin{lemma}[Boolean sorting of a common-core family]
\label{lem:newV96-Boolean-sort}
Suppose all \(q\) supports contain a fixed set \(R\) with
\(|R|\ge2\).  Apply any comparison sorting network to the ordered
supports, replacing a compared pair by its intersection and
union.  The output is the nested chain
\begin{equation}
 L_j=\{i:d_i\ge q-j+1\},\qquad 1\le j\le q,                \label{eq:newV96-level-chain}
\end{equation}
and
\begin{equation}
 \prod_{\nu=1}^q\mathfrak r(A_\nu)
 \le
 \prod_{j=1}^q\mathfrak r(L_j).                           \label{eq:newV96-sorted-reserve}
\end{equation}
\end{lemma}

\begin{proof}
For each label \(i\), its membership bits are acted on by the
ordinary Boolean comparator
\((x,y)\mapsto(\min\{x,y\},\max\{x,y\})\).  A sorting network
therefore leaves \(q-d_i\) zeros followed by \(d_i\) ones, which
is \eqref{eq:newV96-level-chain}.  Every intermediate support
contains \(R\), so every comparator intersection has size at
least two.  Apply \Cref{lem:newV94-fusion} at every comparator
and multiply.
\end{proof}

\begin{theorem}[Dense common-core degree fibre]
\label{thm:newV96-dense-fibre}
For every sufficiently large even \(q\), put \(m=q+2\).  There is
a fixed degree vector and a family of at least
\begin{equation}
 {1\over2}{q\choose q/2}^{q}                              \label{eq:newV96-dense-count}
\end{equation}
ordered \(q\)-tuples of pairwise distinct, connected,
nonsingleton supports on \(m\) labels, all containing the same
two-label core.  Every tuple has the same sorted output:
\[
 \underbrace{R,\ldots,R}_{q/2\ {\rm times}},
 \quad
 \underbrace{X,\ldots,X}_{q/2\ {\rm times}}.
\]
Even after division by the two output type factorials
\((q/2)!^2\), the fibre has size \(\exp(cq^2)\).
\end{theorem}

\begin{proof}
Let \(R=\{1,2\}\), put both core labels in all \(q\) supports, and
for each of the remaining \(q\) labels independently choose
exactly \(q/2\) support columns.  This gives
\({q\choose q/2}^q\) matrices with degrees
\[
 d_1=d_2=q,\qquad d_i=q/2\quad(i\ge3).
\]
All columns contain \(R\), hence every support is nonsingleton
and the incidence graph is connected.

Fix two columns.  In one noncore row their bits agree in
\[
 {q-2\choose q/2-2}+{q-2\choose q/2}
\]
of the \({q\choose q/2}\) choices.  The ratio is
\[
 {q-2\over2(q-1)}< {1\over2}.
\]
Thus the proportion of matrices in which this column pair agrees
in every noncore row is less than \(2^{-q}\).  A union bound over
\({q\choose2}\) pairs is less than \(1/2\) for all sufficiently
large \(q\).  At least half the matrices therefore have pairwise
distinct columns, proving \eqref{eq:newV96-dense-count}.

Equation~\eqref{eq:newV96-level-chain} gives the displayed
two-level chain.  Finally,
\({q\choose q/2}\ge2^q/(q+1)\), whereas
\((q/2)!\le(q/2)^{q/2}\).  Hence the logarithm of the quotient of
\eqref{eq:newV96-dense-count} by \((q/2)!^2\) is
\[
 q^2\log2-O(q\log q)\ge cq^2.
\]
\end{proof}

\begin{corollary}[Degree-only HONC is impossible]
\label{cor:newV96-HONC-no-go}
The HONC prescription of summing a positive post-norm support
fibre using only its degree vector and then applying laminar
domination cannot have a uniform \(K^m\) bound.  The
repeated-support type factorial of V86 is also insufficient after
the sorting map.
\end{corollary}

\begin{proof}
All inputs in \Cref{thm:newV96-dense-fibre} have the same degree
vector, are already connected, and have distinct support types,
so their input type factorials equal one.  They all map to one
reserve-dominating chain.  The fibre remains
\(\exp(cm^2)\) even after the two possible output factorials are
granted.  This exceeds \(K^m\) for every fixed \(K\).
\end{proof}

\begin{correction}[HONC is retired]
\label{corr:newV96-HONC}
Definition~\ref{def:newV95-HONC} is retired as a proposed
sufficient operation.  Its termwise two-laminar reduction remains
correct, but item (iv), interpreted as a positive degree-only
fibre bound, is false by \Cref{cor:newV96-HONC-no-go}.
\end{correction}

\section{The required upstream source card}
\label{sec:newV96-SISC}

\begin{definition}[Simultaneous incidence-source card, SISC]
\label{def:newV96-SISC}
SISC is an operator-level identity, before norms and before
support sorting, which:
\begin{enumerate}[label=\textup{(\roman*)},leftmargin=3.5em]
\item retains the analytic variables that generate membership of
      each label in each primitive root;
\item combines all distinct support families with the normalized
      Bernoulli/Kruskal configuration and the exact Duhamel type
      coefficient;
\item produces a joint row normalization strong enough that the
      dense fibre of \Cref{thm:newV96-dense-fibre} has total mass
      at most \(K^m\);
\item preserves the reserve and exposed-anchor ledgers; and
\item only then applies the two-laminar or skeleton reduction.
\end{enumerate}
\end{definition}

\begin{proposition}[SISC completes the remaining simple row]
\label{prop:newV96-SISC}
If SISC yields the PRAH strong polynomial after summing all
simple cores with \(g>G_m\), then
\Cref{thm:newV96-low-excess} and SISC prove TCNC for all
simple-support primitive ancestries.
\end{proposition}

\begin{proof}
Split by incidence excess at \(G_m\).  The lower row is
\Cref{thm:newV96-low-excess}; SISC treats the complement.
\end{proof}

\begin{programme}[V97: return to the exact simultaneous source]
\label{prog:newV96-V97}
\begin{enumerate}[leftmargin=3em]
\item Write the coefficient of a prescribed distinct support
      family in the exact analytic Duhamel source before norms.
\item Verify explicitly that the V86 type factorial is one on the
      dense distinct-support family.
\item Locate every remaining factorial or Bernoulli probability
      in the label-membership variables.
\item Test whether the percolative Russo quotient supplies a
      normalized row law; if it does not, state the corresponding
      no-go theorem.
\item Replace SISC by the smallest exact upstream identity that
      survives this test.
\end{enumerate}
\end{programme}

\begin{firewall}[Claim boundary after V96]
\label{fire:newV96-boundary}
V96 proves that the V78 canonical skeleton is exactly graded by
incidence excess, derives coefficientwise extension and
whole-support series, proves a fixed-excess global estimate, and
closes every connected simple primitive core with
\(g\le m/\log(m+e)\) on the intermediate wedge.  It also proves
the exact degree-fibre census, the common-core Boolean sorting
identity, and an \(\exp(cm^2)\) dense fibre that rules out the
degree-only HONC proposal.

V96 does not prove SISC, the simple high-excess row, arbitrary
repeated transversal supports, full TOCC, full TCNC, cyclic TIC,
HLTC, full MCWC, OCRC, SCNGC, WPSC, GRLTC, LZTC, PLRC, WSDC,
BRSC, HWSFC, UGCGC, thermal connected WPRAC, all-arity RGSC, the
raw Wilson aggregate strong card, all-order strong MTA, WUFC,
CPM, cutoff removal, reflection positivity,
Osterwalder--Schrader reconstruction, a continuum Yang--Mills
measure, or a positive Hamiltonian gap.  The full Yang--Mills
mass gap has not been proved.
\end{firewall}

\begin{theorem}[V96 result ledger]
\label{thm:newV96-results}
V96:
\begin{enumerate}[label=\textup{(V96.\arabic*)},leftmargin=6.5em]
\item proves the skeleton surplus-grade identity;
\item proves graded extension and whole-support bounds;
\item proves the fixed-excess simple-core estimate;
\item closes all simple cores of submacroscopic excess;
\item proves the exact degree-fibre census;
\item proves common-core Boolean sorting;
\item proves the dense degree-fibre obstruction; and
\item retires HONC and formulates SISC.
\end{enumerate}
\end{theorem}

\begin{proof}
Use
\Cref{lem:newV96-grade,
lem:newV96-one-extension,
cor:newV96-all-extensions,
lem:newV96-whole-graded,
thm:newV96-fixed-excess,
thm:newV96-low-excess,
lem:newV96-degree-census,
lem:newV96-Boolean-sort,
thm:newV96-dense-fibre,
cor:newV96-HONC-no-go,
prop:newV96-SISC}.
\end{proof}

\paragraph{Parent-version record.}
V96 was formed from
\texttt{Renewal\_Yang\_Mills\_Mass\_Gap\_Research\_Notes\_V95.tex}.
The V95 parent had \(43\,026\) logical source lines,
\(1\,394\,545\) bytes, and SHA--256
\[
 \texttt{c43dc30eb18f14c95a31c5e052870087a575e1ffc0360e8e7a27d426cb4bf1a1}.
\]
The next compulsory companion audit is V100.

% =============================================================================
% END APPENDED FIFTH-SERIES V96 MATERIAL
% =============================================================================

\clearpage
\chapter{Fifth-series V97: exact optional-support collapse at a
saturated forest gate}
\label{ch:newV97-support-collapse}

\section{Optional labels are contractions before norms}
\label{sec:newV97-optional}

No companion audit is due at V97.

For a heat-group element \(U\), write
\[
 \Delta_i(U):=D^{R_i}(U)-I.
\]
As in V60, absent tensor factors are identities and
\[
 Q_A
 =
 \int_0^\infty
 \mathbb E_{h_x}
 \bigotimes_{i\in A}\Delta_i(U)\,
 \Pi_\alpha^{\rm sm}(dx),
 \qquad |A|\ge2.
\]

\begin{lemma}[Exact optional-support identity]
\label{lem:newV97-optional}
Let \(F\subseteq X\), \(|F|\ge2\), and choose
\(0\le\lambda_j\le1\) for \(j\notin F\).  Then
\begin{align}
 &\sum_{F\subseteq A\subseteq X}
 \left(\prod_{j\in A\setminus F}\lambda_j\right)Q_A
 \notag\\
 &\quad=
 \int_0^\infty\mathbb E_{h_x}
 \left[
  \bigotimes_{i\in F}\Delta_i(U)
  \otimes
  \bigotimes_{j\notin F}
  \{I+\lambda_j\Delta_j(U)\}
 \right]\Pi_\alpha^{\rm sm}(dx).                           \label{eq:newV97-optional}
\end{align}
Moreover every optional factor is a contraction:
\begin{equation}
 \|I+\lambda_j\Delta_j(U)\|
 =
 \|(1-\lambda_j)I+\lambda_jD^{R_j}(U)\|
 \le1.                                                     \label{eq:newV97-contraction}
\end{equation}
\end{lemma}

\begin{proof}
Expand the tensor product over \(j\notin F\).  Choosing the
\(\lambda_j\Delta_j\) term at precisely the labels in
\(A\setminus F\) gives the summand indexed by \(A\).
The finite sum may be moved through the heat expectation and
L\'evy integral; absolute integrability follows from the two
forced increments and the V60 same-jump bound.  This proves
\eqref{eq:newV97-optional}.  Equation
\eqref{eq:newV97-contraction} follows because it is a convex
combination of two contractions, one of which is unitary.
\end{proof}

\begin{corollary}[The unrestricted membership cube costs no
\(2^{m-|F|}\)]
\label{cor:newV97-cube}
The completely summed optional-support expression in
\eqref{eq:newV97-optional} is bounded by the same raw
same-jump integral as the forced set \(F\), independently of the
number of optional labels.
\end{corollary}

\begin{proof}
Take the norm inside the integrals and use
\eqref{eq:newV97-contraction}.  Only the factors
\(\Delta_i\), \(i\in F\), remain.
\end{proof}

\begin{assessment}[Resolution of the dense incidence cube]
\label{ass:newV97-cube}
The \(\exp(cm^2)\) family of V96 was created by taking absolute
values before summing the optional membership bit of each label
in each root.  At one jump occurrence, the exact two choices are
\(I\) and \(\Delta_j\), and their sum is \(D^{R_j}(U)\), not two
positive terms.  Thus the exact same-jump algebra has the
normalization required by SISC.  It remains to transport this
identity through the nonconstant interpolation gate.
\end{assessment}

\section{Laplace representation of a saturated forest
coefficient}
\label{sec:newV97-Laplace}

Let \(E\) be a nonempty forest, let
\[
 F:=V(E),\qquad f:=|F|,\qquad k:=|E|\ge1,
\]
and let \({\cal K}_F(E)\) denote its nontrivial components on
\(F\).  Define
\begin{equation}
 c_E:=
 { \displaystyle
   \left(\prod_{\{i,j\}\in E}b_ib_j\right)
   \left(\prod_{C\in{\cal K}_F(E)}B_C\right)
  \over
  \displaystyle\prod_{i\in F}b_i}.                        \label{eq:newV97-cE}
\end{equation}

\begin{lemma}[Optional dependence of the exact forest
probability]
\label{lem:newV97-gamma-optional}
For every \(A\supseteq F\),
\begin{equation}
 \gamma_A(E)={c_E\over B_A^k}.                             \label{eq:newV97-gamma-optional}
\end{equation}
\end{lemma}

\begin{proof}
Apply the exact formula
\eqref{eq:newV87-gamma-formula}.  Every vertex of
\(A\setminus F\) is an isolated component of \(E\) on \(A\).
Its component factor \(b_j\) cancels the same factor in
\(\prod_{i\in A}b_i\).  What remains is exactly
\eqref{eq:newV97-gamma-optional}.
\end{proof}

\begin{theorem}[Saturated forest support-collapse identity]
\label{thm:newV97-saturated-collapse}
For every nonempty forest \(E\),
\begin{equation}
 \left\|
  \sum_{F\subseteq A\subseteq X}\gamma_A(E)Q_A
 \right\|_{\rm cb}
 \le
 \gamma_F(E)\,\mathfrak r(F)
 \le\mathfrak r(F),                                       \label{eq:newV97-saturated-collapse}
\end{equation}
after one uniform enlargement of the reserve base.
\end{theorem}

\begin{proof}
For \(k\ge1\),
\[
 {1\over B_A^k}
 =
 {1\over(k-1)!}
 \int_0^\infty t^{k-1}e^{-tB_A}\,dt.
\]
Use \Cref{lem:newV97-gamma-optional}, interchange the finite
support sum with the \(t\)-integral, and factor
\(e^{-tB_A}=e^{-tB_F}\prod_{j\in A\setminus F}e^{-tb_j}\).
Lemma~\ref{lem:newV97-optional}, with
\(\lambda_j=e^{-tb_j}\), and
\eqref{eq:newV97-contraction} show that the norm is at most
\[
 {c_E\over(k-1)!}
 \int_0^\infty t^{k-1}e^{-tB_F}\,dt
 \int_0^\infty
 \mathbb E_{h_x}
 \prod_{i\in F}\|\Delta_i(U)\|
 \Pi_\alpha^{\rm sm}(dx).
\]
The first integral is \((k-1)!/B_F^k\).  The proof of the
unrelaxed same-jump estimate
\Cref{lem:newV61-reserve}, applied only to the forced increments,
bounds the second by \(\mathfrak r(F)\).  Finally
\[
 {c_E\over B_F^k}=\gamma_F(E)\le1
\]
by \Cref{cor:newV87-gamma-formula}.
\end{proof}

\begin{corollary}[The drift-corrected root also collapses]
\label{cor:newV97-R-collapse}
With \(R_A\) defined by \eqref{eq:newV61-RA},
\begin{equation}
 \left\|
  \sum_{F\subseteq A\subseteq X}\gamma_A(E)R_A
 \right\|_{\rm cb}
 \le K^f f^f s^{(f-2)/2}\prod_{i\in F}b_i.                 \label{eq:newV97-R-collapse}
\end{equation}
\end{corollary}

\begin{proof}
If \(f\ge3\), every \(A\supseteq F\) has arity at least three,
so \(R_A=Q_A\) and
\Cref{thm:newV97-saturated-collapse} applies.  If \(f=2\), then
\(E\) is the unique edge on \(F=\{i,j\}\),
\(\gamma_F(E)=1\), and only the term \(A=F\) receives the drift
correction \(-2d_\alpha\Omega_{ij}\).  Its norm is at most
\(Kb_ib_j\) by the proof of \Cref{lem:newV61-reserve}.
Combine the two bounds.
\end{proof}

\begin{assessment}[What the saturated theorem pays]
\label{ass:newV97-saturated}
For a fixed owned forest \(E\), all labels not incident to \(E\)
have been summed exactly, with no support census and no total-mark
power.  The result is one reserve activity on the forced endpoint
set \(F=V(E)\), with the exact residual probability
\(\gamma_F(E)\) still available.  This is the operator-level
collapse that the degree-only HONC route could not supply.
\end{assessment}

\section{Why the known factorials and shortcut identity were not
the missing normalization}
\label{sec:newV97-not-factorial}

\begin{lemma}[Distinct Duhamel types have no type factorial]
\label{lem:newV97-distinct}
For \(q\) pairwise distinct support types
\(A_1,\ldots,A_q\), the V86 noncommutative assembled-type bound
contains
\[
 \prod_{\nu=1}^q{1\over1!}=1.
\]
Thus neither the ordered simplex nor support-type coarsening
suppresses a dense distinct-support fibre.
\end{lemma}

\begin{proof}
Set \(\nu_{A_\ell}=1\) in
\eqref{eq:newV86-type-factorial}.  The \(q!\) distinct time
orders exactly cancel the simplex volume \(1/q!\).
\end{proof}

\begin{lemma}[The dense fibre survives prescribed owned star
edges]
\label{lem:newV97-owned-dense}
In the construction of
\Cref{thm:newV96-dense-fibre}, require noncore label \(j+2\) to
belong to support column \(j\), and assign the distinct outer-tree
edge \(\{1,j+2\}\) to that root.  There remain
\(\exp(cq^2)\) pairwise-distinct support families.  With equal
marks, the product of their one-edge gate probabilities loses at
most \(\exp(Cq\log q)\), so a termwise positive gate majorant is
still superexponential.
\end{lemma}

\begin{proof}
Each noncore row now chooses \(q/2-1\) further columns among
\(q-1\), giving
\[
 {q-1\choose q/2-1}^{q}
 =
 2^{-q}{q\choose q/2}^{q}
 =
 \exp(q^2\log2-O(q\log q)).
\]
The same union-bound argument as in
\Cref{thm:newV96-dense-fibre} retains a fixed positive fraction
with distinct columns: for a fixed column pair, every row whose
prescribed column is outside the pair makes the two bits equal
with probability at most \(2/3\), while the two exceptional rows
only decrease the bound.  Thus the collision probability is at
most \(q^2(2/3)^{q-2}\).

For equal marks and any support \(A\),
\[
 \gamma_A(\{\{1,j+2\}\})={2\over|A|}\ge {2\over q+2}.
\]
The product over \(q\) roots is therefore at least
\((2/(q+2))^q=\exp(-O(q\log q))\), which cannot offset the
\(\exp(cq^2)\) census.
\end{proof}

\begin{assessment}[Kruskal has a different job]
\label{ass:newV97-Kruskal}
The owned edges in \Cref{lem:newV97-owned-dense} lie in one fixed
outer star.  There is no shortcut-tree census to pay: the
corresponding Kruskal allocation has total mass one.  V83 remains
essential for changing output block trees, but it cannot replace
the optional-support contraction
\Cref{lem:newV97-optional}.
\end{assessment}

\section{The unsaturated interpolation defect}
\label{sec:newV97-defect}

For \(0\le\boldsymbol x\le1\), write
\begin{equation}
 c_A(E;\boldsymbol x)
 :=
 [\boldsymbol y^E]\,
 \chi_A(\boldsymbol x+\boldsymbol y).                      \label{eq:newV97-cAx}
\end{equation}
The gate coefficients are nonnegative and
\[
 0\le c_A(E;\boldsymbol x)\le\gamma_A(E),
 \qquad
 c_A(E;\boldsymbol 1)=\gamma_A(E).
\]

\begin{warning}[Scalar coefficient domination does not preserve
the operator cancellation]
\label{warn:newV97-monotonicity}
The pointwise inequality
\(c_A(E;\boldsymbol x)\le\gamma_A(E)\) does not permit
\[
 \left\|\sum_Ac_A(E;\boldsymbol x)Q_A\right\|
 \le
 \left\|\sum_A\gamma_A(E)Q_A\right\|.
\]
Consequently
\Cref{thm:newV97-saturated-collapse} cannot be inserted at a
general BKAR interpolation matrix by monotonicity.
\end{warning}

\begin{proof}
Operator norm is not monotone under independently increasing
nonnegative scalar coefficients of sign-indefinite operators.
For example, take \(Q_1=I\), \(Q_2=-I\),
\((c_1,c_2)=(1,0)\), and
\((\gamma_1,\gamma_2)=(1,1)\).  The two norms are respectively
one and zero.
\end{proof}

\begin{definition}[Interpolation-gate mixture card, IGMC]
\label{def:newV97-IGMC}
IGMC is an exact representation, at every BKAR interpolation
matrix, of the family
\(\{c_A(E;\boldsymbol x):A\supseteq V(E)\}\) as a positive
mixture of optional-label laws for which the factors at label
\(j\) have the contraction form
\[
 (1-\lambda_j)I+\lambda_jD^{R_j}(U),
 \qquad 0\le\lambda_j\le1.
\]
The mixture must retain the owned forest, normalized Bernoulli
configuration, and Kruskal variables, and have total variation
at most \(K^{|V(E)|+|E|}\).
\end{definition}

\begin{proposition}[IGMC supplies the support-fibre part of SISC]
\label{prop:newV97-IGMC}
IGMC and \Cref{cor:newV97-R-collapse} sum every optional
root-support membership before norms with an exponential cost
depending only on forced forest data.  In particular, they remove
the dense degree fibre of \Cref{thm:newV96-dense-fibre} from
SISC; the remaining data are the disjoint owned-edge partition
and its component-coalescence defect.
\end{proposition}

\begin{proof}
Apply the optional-support contraction inside each positive
mixture component, then integrate the mixture.  Every root
occurrence owns a nonempty edge set, and the owned sets are
disjoint subsets of an outer tree.  Thus their total edge count
is at most \(m-1\), so the mixture costs multiply to at most
\(K^m\).  After optional labels have disappeared, a root is
supported on \(V(E)\).  The excess of these forced supports is
exactly the component-coalescence defect of V88: if \(E_\nu\)
has \(c_\nu\) nontrivial components, then
\[
 \sum_\nu\{|V(E_\nu)|-1\}-m+1
 =
 \sum_\nu(c_\nu-1),
\]
because \(\sum_\nu|E_\nu|=m-1\) for a partition of the outer
tree.  No other support fibre remains.
\end{proof}

\begin{programme}[V98: Bernoulli lift of the unsaturated gate]
\label{prog:newV97-V98}
\begin{enumerate}[leftmargin=3em]
\item Represent
      \(c_A(E;\boldsymbol x)\) as the probability that a weighted
      Cayley tree contains \(E\) and all its remaining edges are
      open in an independent Bernoulli graph.
\item Condition first on the open graph, not on the support.
\item Test whether the matrix-tree contraction of the optional
      vertices admits the Laplace mixture required by IGMC.
\item Prove the partition-endpoint row separately if the
      continuous interpolation row fails.
\item Keep the V88 component-coalescence defect explicit after
      optional supports have collapsed.
\end{enumerate}
\end{programme}

\begin{firewall}[Claim boundary after V97]
\label{fire:newV97-boundary}
V97 proves the exact optional-support identity and contraction,
derives the Laplace representation of normalized forest
probabilities, and proves that every saturated forest gate sums
all optional support labels to one reserve activity on its forced
endpoint set.  It proves the same for drift-corrected roots and
shows that neither distinct-type Duhamel coefficients nor
one-edge gate probabilities pay the dense fibre termwise.  It
isolates the unsaturated interpolation coefficient as the exact
remaining support-collapse defect and formulates IGMC.

V97 does not prove IGMC, the unsaturated support collapse, the
forced-forest component-coalescence row, SISC, the simple
high-excess row, arbitrary repeated transversal supports, full
TOCC, full TCNC, cyclic TIC, HLTC, full MCWC, OCRC, SCNGC, WPSC,
GRLTC, LZTC, PLRC, WSDC, BRSC, HWSFC, UGCGC, thermal connected
WPRAC, all-arity RGSC, the raw Wilson aggregate strong card,
all-order strong MTA, WUFC, CPM, cutoff removal, reflection
positivity, Osterwalder--Schrader reconstruction, a continuum
Yang--Mills measure, or a positive Hamiltonian gap.  The full
Yang--Mills mass gap has not been proved.
\end{firewall}

\begin{theorem}[V97 result ledger]
\label{thm:newV97-results}
V97:
\begin{enumerate}[label=\textup{(V97.\arabic*)},leftmargin=6.5em]
\item proves the exact optional-support contraction;
\item derives the optional form of the forest probability;
\item proves saturated forest support collapse;
\item includes the drift-corrected pair root;
\item proves the distinct-type factorial is absent;
\item proves the owned-star dense fibre persists termwise;
\item isolates the unsaturated monotonicity obstruction; and
\item formulates IGMC and its sufficiency for the support fibre.
\end{enumerate}
\end{theorem}

\begin{proof}
Use
\Cref{lem:newV97-optional,
cor:newV97-cube,
lem:newV97-gamma-optional,
thm:newV97-saturated-collapse,
cor:newV97-R-collapse,
lem:newV97-distinct,
lem:newV97-owned-dense,
warn:newV97-monotonicity,
prop:newV97-IGMC}.
\end{proof}

\paragraph{Parent-version record.}
V97 was formed from
\texttt{Renewal\_Yang\_Mills\_Mass\_Gap\_Research\_Notes\_V96.tex}.
The V96 parent had \(43\,520\) logical source lines,
\(1\,411\,283\) bytes, and SHA--256
\[
 \texttt{c821c1aa09f8b8316146bf027849f00304eec8181b76abfef50e1da88de682fc}.
\]
The next compulsory companion audit is V100.

% =============================================================================
% END APPENDED FIFTH-SERIES V97 MATERIAL
% =============================================================================

\clearpage
\chapter{Fifth-series V98: Bernoulli lifting and the uniform
interpolation support collapse}
\label{ch:newV98-Bernoulli}

\section{The exact probabilistic meaning of a gate coefficient}
\label{sec:newV98-lift}

No companion audit is due at V98.

For a support \(A\), let \(\mu_A\) be the weighted Cayley-tree
law
\begin{equation}
 \mu_A(T)
 :=
 {1\over W_A}\prod_{i\in A}b_i^{d_T(i)},
 \qquad T\in\mathfrak T(A).                                \label{eq:newV98-mu}
\end{equation}
The weighted Cayley identity says that \(\mu_A\) is a probability
law.

\begin{theorem}[Bernoulli lift of the differentiated gate]
\label{thm:newV98-Bernoulli}
Fix a forest \(E\) on \(A\).  Independently sample
\(\mathsf T_A\sim\mu_A\), and sample a Bernoulli graph
\(\omega_{\boldsymbol x}\) in which every pair edge \(e\) is open
independently with probability \(x_e\).  Then
\begin{equation}
 c_A(E;\boldsymbol x)
 =
 \mathbb P\left\{
 E\subseteq\mathsf T_A,\quad
 \mathsf T_A\setminus E\subseteq\omega_{\boldsymbol x}
 \right\}.                                                 \label{eq:newV98-Bernoulli}
\end{equation}
Equivalently, conditioning on the open graph,
\begin{equation}
 c_A(E;\boldsymbol x)
 =
 \mathbb E_{\omega_{\boldsymbol x}}
 \mu_A\{T:E\subseteq T\subseteq E\cup\omega_{\boldsymbol x}\}.
                                                               \label{eq:newV98-conditioned}
\end{equation}
\end{theorem}

\begin{proof}
Differentiate the tree polynomial defining \(\chi_A\):
\[
 c_A(E;\boldsymbol x)
 =
 \sum_{\substack{T\in\mathfrak T(A)\\E\subseteq T}}
 \mu_A(T)\prod_{e\in T\setminus E}x_e.
\]
For fixed \(T\), the product is exactly the probability that all
edges of \(T\setminus E\) are open.  Sum over \(T\), and then
condition on \(\omega_{\boldsymbol x}\).
\end{proof}

\begin{corollary}[Normalization and saturation]
\label{cor:newV98-normalization}
For every interpolation matrix,
\[
 0\le c_A(E;\boldsymbol x)\le\gamma_A(E),
 \qquad
 c_A(E;\boldsymbol 1)=\gamma_A(E).
\]
The loss between the two coefficients is an exact Bernoulli
connectivity event, not an arity constant.
\end{corollary}

\begin{proof}
Drop, or impose, the second event in
\eqref{eq:newV98-Bernoulli}.
\end{proof}

\begin{assessment}[What must remain normalized]
\label{ass:newV98-normalized}
Equation~\eqref{eq:newV98-conditioned} makes precise why neither
the closed-edge Bernoulli factors nor the weighted Cayley-tree law
may be discarded before optional supports are summed.  IGMC is a
projective-consistency statement for these two probability laws
as the support \(A\) varies.
\end{assessment}

\section{Complete collapse on the uniform interpolation line}
\label{sec:newV98-uniform}

Let \(E\) be a nonempty forest on \(F=V(E)\), let
\(f=|F|\), \(k=|E|\), and let \(c(E)=f-k\) be its number of
nontrivial components.

\begin{lemma}[Uniform gate factorization]
\label{lem:newV98-uniform-gate}
If \(x_{ij}=\lambda\in[0,1]\) for all distinct \(i,j\), then for
every \(A\supseteq F\),
\begin{equation}
 c_A(E;\lambda\boldsymbol 1)
 =
 \lambda^{|A|-1-k}\gamma_A(E)
 =
 \lambda^{c(E)-1}
 \lambda^{|A\setminus F|}\gamma_A(E),                      \label{eq:newV98-uniform-gate}
\end{equation}
with the convention \(0^0=1\).
\end{lemma}

\begin{proof}
Every tree on \(A\) containing \(E\) has exactly
\(|A|-1-k\) undifferentiated edges.  Hence every monomial in the
definition of \(c_A\) has the same power of \(\lambda\).  The
remaining weighted tree sum is \(\gamma_A(E)\).
Finally
\[
 |A|-1-k=|A\setminus F|+f-1-k
 =|A\setminus F|+c(E)-1.
\]
\end{proof}

\begin{theorem}[Uniform interpolation support collapse]
\label{thm:newV98-uniform-collapse}
For \(0\le\lambda\le1\),
\begin{equation}
 \left\|
  \sum_{F\subseteq A\subseteq X}
  c_A(E;\lambda\boldsymbol 1)R_A
 \right\|_{\rm cb}
 \le
 K^f\lambda^{c(E)-1}
 f^fs^{(f-2)/2}\prod_{i\in F}b_i.                          \label{eq:newV98-uniform-collapse}
\end{equation}
The constant is independent of \(m\) and of the number of
optional labels.
\end{theorem}

\begin{proof}
Insert \Cref{lem:newV98-uniform-gate} and the Laplace formula
used in \Cref{thm:newV97-saturated-collapse}.  After
\(B_A^{-k}\) is represented by the \(t\)-integral, every optional
label now has weight
\[
 \lambda e^{-tb_j}\in[0,1].
\]
Apply \Cref{lem:newV97-optional}; its optional factor is
\[
 I+\lambda e^{-tb_j}\Delta_j(U)
 =
 (1-\lambda e^{-tb_j})I
 +\lambda e^{-tb_j}D^{R_j}(U),
\]
a contraction.  The \(t\)-integral leaves
\(\gamma_F(E)\le1\), and the forced heat-increment integral is
the reserve bound on \(F\).  The prefactor
\(\lambda^{c(E)-1}\) remains.  The pair drift correction is
handled exactly as in \Cref{cor:newV97-R-collapse}.
\end{proof}

\begin{corollary}[IGMC on the one-level hierarchy]
\label{cor:newV98-one-level-IGMC}
IGMC holds whenever the BKAR interpolation matrix has one
off-diagonal level.  In that row the support-fibre cost is one,
apart from the reserve base on the forced endpoints.
\end{corollary}

\begin{proof}
The proof of \Cref{thm:newV98-uniform-collapse} is the required
positive product-law representation, with the Laplace parameter
\(t\) as mixing variable.
\end{proof}

\section{Partition endpoints and support localization}
\label{sec:newV98-endpoints}

Let \(\rho\) be a partition of \(X\), and let
\(\boldsymbol x^\rho\) be its zero--one connectivity matrix.

\begin{lemma}[No optional vertex may enter a fresh partition
block]
\label{lem:newV98-endpoint-localization}
If \(A\supseteq F\) contains a label \(j\) whose \(\rho\)-block
does not meet \(F\), then
\begin{equation}
 c_A(E;\boldsymbol x^\rho)=0.                              \label{eq:newV98-endpoint-zero}
\end{equation}
More generally, this coefficient is nonzero only if the graph
formed by \(E\) together with the complete graphs on the sets
\(A\cap P\), \(P\in\rho\), is connected.
\end{lemma}

\begin{proof}
By \Cref{thm:newV98-Bernoulli}, every undifferentiated edge of the
sampled tree must lie inside a \(\rho\)-block.  All cross-block
edges of that tree must therefore belong to \(E\).  A block
containing \(j\) but no vertex of \(F=V(E)\) has no edge of \(E\)
incident to it and cannot connect to the rest of the tree.  The
same argument gives the stated connectivity criterion.
\end{proof}

\begin{assessment}[Endpoint gain]
\label{ass:newV98-endpoint}
At a partition endpoint, optional support labels are confined to
blocks already touched by the owned forest.  For thermal
partitions this converts the unrestricted global membership cube
into a product of local light-block problems plus forced heavy
singletons.  The continuous BKAR interpolation cannot be replaced
by its endpoints, so this is a boundary row rather than IGMC in
full.
\end{assessment}

\section{Why arbitrary cube matrices are the wrong test}
\label{sec:newV98-cube-no-go}

\begin{lemma}[A non-ultrametric gate need not be a positive
optional product law]
\label{lem:newV98-nonultrametric}
Take four equal marks, \(E=\{\{1,2\}\}\), and set
\[
 x_{13}=x_{23}=x_{14}=x_{24}=1,\qquad x_{34}=0.
\]
For optional labels \(3,4\), the normalized coefficients are
\[
 c_\varnothing=1,\qquad
 c_{\{3\}}=c_{\{4\}}={2\over3},\qquad
 c_{\{3,4\}}={1\over4}.
\]
They cannot have a representation
\[
 c_J=\int\prod_{j\in J}\lambda_j\,d\mu(\boldsymbol\lambda)
\]
with a positive measure of mass one on \([0,1]^2\).
\end{lemma}

\begin{proof}
On three vertices, the two trees containing \(12\) give
\(2/3\).  On four vertices, a tree containing \(12\) and using
only the four displayed open edges chooses one edge incident to
3 and one incident to 4, giving four of the sixteen Cayley trees;
hence \(1/4\).

For \(0\le\lambda_3,\lambda_4\le1\),
\[
 (1-\lambda_3)(1-\lambda_4)\ge0.
\]
Any positive product-moment representation would therefore
require
\[
 c_{\{3,4\}}\ge c_{\{3\}}+c_{\{4\}}-c_\varnothing={1\over3},
\]
contrary to \(c_{\{3,4\}}=1/4\).
\end{proof}

\begin{remark}[The counterexample is not a BKAR matrix]
\label{rem:newV98-not-BKAR}
The data in \Cref{lem:newV98-nonultrametric} violate the
ultrametric inequality
\[
 x_{34}\ge\min\{x_{31},x_{14}\}.
\]
Thus they rule out extending IGMC to the whole entrywise cube but
do not disprove IGMC on the BKAR locus.
\end{remark}

\section{The exact hierarchy of a BKAR interpolation matrix}
\label{sec:newV98-ultrametric}

\begin{lemma}[Nested partition representation]
\label{lem:newV98-nested}
Let \(T_0\) be the input BKAR tree and
\[
 x_{ij}^{T_0}(\boldsymbol u)
 =
 \min_{e\in{\rm path}_{T_0}(i,j)}u_e.
\]
For \(t\in[0,1]\), let \(\rho_t\) be the component partition of
the forest \(\{e\in T_0:u_e\ge t\}\).  Then
\begin{equation}
 x_{ij}^{T_0}(\boldsymbol u)
 =
 \int_0^1{\bf1}_{\{i,j\ {\rm in\ one\ block\ of}\ \rho_t\}}\,dt.
                                                               \label{eq:newV98-layer-cake}
\end{equation}
The partitions \(\rho_t\) are nested and change at no more than
\(m-1\) levels.
\end{lemma}

\begin{proof}
The vertices \(i,j\) lie in one component at level \(t\) exactly
when every edge on their unique \(T_0\)-path has \(u_e\ge t\).
This is equivalent to
\(t\le\min_{e\in{\rm path}}u_e=x_{ij}\).  Integrating the
indicator proves \eqref{eq:newV98-layer-cake}.  Monotonicity in
\(t\) gives nesting, and only the \(m-1\) edge values can change
the partition.
\end{proof}

\begin{lemma}[Top-level affine decomposition]
\label{lem:newV98-top-layer}
Let \(\lambda=\min_{e\in T_0}u_e<1\), remove all edges with value
\(\lambda\), and let \(\rho\) be the resulting component
partition.  Then
\begin{equation}
 \boldsymbol x^{T_0}
 =
 \lambda\boldsymbol 1
 +(1-\lambda)\boldsymbol z^\rho,                            \label{eq:newV98-top-layer}
\end{equation}
where \(z_{ij}=0\) across distinct blocks of \(\rho\), while
inside each block \(z_{ij}=(x_{ij}-\lambda)/(1-\lambda)\) is
again a BKAR ultrametric matrix generated by the corresponding
subtree.
\end{lemma}

\begin{proof}
Every path between distinct \(\rho\)-blocks crosses an edge of
minimal value \(\lambda\), so its minimum is \(\lambda\).
Inside one block every path edge is larger than \(\lambda\);
the affine rescaling preserves path minima and gives the stated
subtree matrix.
\end{proof}

\begin{definition}[Ultrametric layer-collapse card, ULCC]
\label{def:newV98-ULCC}
ULCC is an induction on the nested decomposition
\eqref{eq:newV98-top-layer} which expands the multi-affine gate
only after:
\begin{enumerate}[label=\textup{(\roman*)},leftmargin=3.5em]
\item the uniform \(\lambda\)-part has been summed by
      \Cref{thm:newV98-uniform-collapse};
\item every remainder edge has been assigned to one smaller
      \(\rho\)-block;
\item optional labels in untouched blocks have been removed by
      \Cref{lem:newV98-endpoint-localization}; and
\item the component-coalescence choices across the hierarchy are
      summed with the V83 Kruskal normalization.
\end{enumerate}
Its required total layer fibre is \(K^m\), not \(K\) per
root--layer pair.
\end{definition}

\begin{proposition}[ULCC implies IGMC on the BKAR locus]
\label{prop:newV98-ULCC}
ULCC proves IGMC and therefore supplies the support-fibre part of
SISC.
\end{proposition}

\begin{proof}
Iterate \Cref{lem:newV98-top-layer}.  Every terminal block is a
one-level or partition-endpoint row, covered respectively by
\Cref{thm:newV98-uniform-collapse} and
\Cref{lem:newV98-endpoint-localization}.  ULCC supplies the
uniform aggregate fibre for assigning remainder edges and
coalescing components.  This is exactly the mixture required in
\Cref{def:newV97-IGMC}; apply
\Cref{prop:newV97-IGMC}.
\end{proof}

\begin{programme}[V99: the two-level ultrametric recursion]
\label{prog:newV98-V99}
\begin{enumerate}[leftmargin=3em]
\item Expand one top-level decomposition
      \eqref{eq:newV98-top-layer} for a prescribed owned forest.
\item Group all remainder edges by child block before summing
      optional supports.
\item Apply the exact weighted-forest contraction in each child
      block and retain every cross-block edge weight.
\item Prove ULCC first when the hierarchy has two nontrivial
      levels.
\item If the layer assignment has a Bell fibre, return to the
      Bernoulli graph in \eqref{eq:newV98-conditioned} before
      expanding the affine decomposition.
\end{enumerate}
\end{programme}

\begin{firewall}[Claim boundary after V98]
\label{fire:newV98-boundary}
V98 proves the exact Bernoulli lift of every differentiated gate,
proves complete optional-support collapse for every uniform
interpolation matrix, localizes optional supports at partition
endpoints, and proves the nested layer and affine decomposition
of every BKAR interpolation matrix.  It also gives a
non-ultrametric counterexample showing why positivity cannot be
asserted on the whole entrywise cube, and formulates ULCC for the
actual BKAR locus.

V98 does not prove ULCC, full IGMC, the general unsaturated
support collapse, the forced-forest component-coalescence row,
SISC, the simple high-excess row, arbitrary repeated transversal
supports, full TOCC, full TCNC, cyclic TIC, HLTC, full MCWC,
OCRC, SCNGC, WPSC, GRLTC, LZTC, PLRC, WSDC, BRSC, HWSFC,
UGCGC, thermal connected WPRAC, all-arity RGSC, the raw Wilson
aggregate strong card, all-order strong MTA, WUFC, CPM, cutoff
removal, reflection positivity, Osterwalder--Schrader
reconstruction, a continuum Yang--Mills measure, or a positive
Hamiltonian gap.  The full Yang--Mills mass gap has not been
proved.
\end{firewall}

\begin{theorem}[V98 result ledger]
\label{thm:newV98-results}
V98:
\begin{enumerate}[label=\textup{(V98.\arabic*)},leftmargin=6.5em]
\item proves the Bernoulli lift of gate coefficients;
\item proves uniform-line optional-support collapse;
\item proves IGMC on the one-level hierarchy;
\item proves partition-endpoint support localization;
\item proves the arbitrary-cube product-law no-go;
\item proves every BKAR matrix is nested ultrametric;
\item proves its top-level affine recursion; and
\item formulates ULCC and proves its sufficiency.
\end{enumerate}
\end{theorem}

\begin{proof}
Use
\Cref{thm:newV98-Bernoulli,
cor:newV98-normalization,
lem:newV98-uniform-gate,
thm:newV98-uniform-collapse,
cor:newV98-one-level-IGMC,
lem:newV98-endpoint-localization,
lem:newV98-nonultrametric,
lem:newV98-nested,
lem:newV98-top-layer,
prop:newV98-ULCC}.
\end{proof}

\paragraph{Parent-version record.}
V98 was formed from
\texttt{Renewal\_Yang\_Mills\_Mass\_Gap\_Research\_Notes\_V97.tex}.
The V97 parent had \(43\,965\) logical source lines,
\(1\,426\,288\) bytes, and SHA--256
\[
 \texttt{ae22a8125aa0ca1f151ce840a8048a423673ab55686f4253165a802073d0cc8f}.
\]
The next compulsory companion audit is V100.

% =============================================================================
% END APPENDED FIFTH-SERIES V98 MATERIAL
% =============================================================================

\clearpage
\chapter{Fifth-series V99: exact expansion of one ultrametric
layer and removal of all optional support labels}
\label{ch:newV99-layer}

\section{The affine layer expansion has no Bell fibre}
\label{sec:newV99-expansion}

No companion audit is due at V99.

Fix a support \(A\), a nonempty owned forest \(E\) on \(A\), and
write
\[
 a:=|A|,\qquad k:=|E|.
\]
Suppose an interpolation matrix has the top-layer decomposition
of \Cref{lem:newV98-top-layer},
\begin{equation}
 \boldsymbol x
 =
 \lambda\boldsymbol 1+(1-\lambda)\boldsymbol z^\rho,
 \qquad 0\le\lambda\le1,                                  \label{eq:newV99-layer}
\end{equation}
where \(z_{ij}=0\) across distinct child blocks of \(\rho\).

\begin{theorem}[Exact one-layer forest expansion]
\label{thm:newV99-layer-expansion}
One has
\begin{align}
 c_A(E;\boldsymbol x)
 &=
 \sum_{\substack{
       H\subseteq\binom A2\setminus E\\
       E\cup H\ {\rm a\ forest}}}
 \lambda^{a-1-k-|H|}
 (1-\lambda)^{|H|}
 \boldsymbol z^H
 \gamma_A(E\cup H).                                      \label{eq:newV99-layer-expansion}
\end{align}
Terms with a negative exponent do not occur.
\end{theorem}

\begin{proof}
From the differentiated gate polynomial,
\[
 c_A(E;\boldsymbol x)
 =
 \sum_{\substack{T\in\mathfrak T(A)\\E\subseteq T}}
 \mu_A(T)\prod_{e\in T\setminus E}
 \{\lambda+(1-\lambda)z_e\}.
\]
Expand the last product.  Let \(H\subseteq T\setminus E\) be the
set of edges from which the \(z\)-term is selected.  Since \(T\)
is a tree, \(E\cup H\) is a forest, and the other
\(a-1-k-|H|\) edges contribute \(\lambda\).  Conversely, after
\(H\) is fixed, summing all \(T\supseteq E\cup H\) gives
\(\gamma_A(E\cup H)\) by definition.  This proves the formula.
\end{proof}

\begin{lemma}[Layer assignments form a subprobability]
\label{lem:newV99-layer-probability}
For a fixed \(T\supseteq E\),
\begin{align}
 &\sum_{H\subseteq T\setminus E}
 \lambda^{|T\setminus E|-|H|}
 (1-\lambda)^{|H|}\boldsymbol z^H
 \notag\\
 &\qquad=
 \prod_{e\in T\setminus E}
 \{\lambda+(1-\lambda)z_e\}
 \le1.                                                     \label{eq:newV99-layer-probability}
\end{align}
Thus expansion of one layer introduces neither a Bell number nor
an independent forest census.
\end{lemma}

\begin{proof}
This is the binomial product identity edge by edge.  Every
\(z_e\) lies in \([0,1]\), so each factor is at most one.
\end{proof}

\begin{assessment}[The correct order of the layer sum]
\label{ass:newV99-order}
The sets \(H\) must be summed while they are still choices inside
one sampled Cayley tree \(T\).  Counting all forests \(H\) after
forgetting \(T\) would destroy
\eqref{eq:newV99-layer-probability}, exactly as dropping closed
Bernoulli edges destroyed the normalization in V83.
\end{assessment}

\section{Optional support labels disappear after one layer}
\label{sec:newV99-collapse}

For a term \(H\) in \eqref{eq:newV99-layer-expansion}, put
\[
 G:=E\cup H,\qquad F_G:=V(G),\qquad
 c(G):=|F_G|-|G|.
\]
Since \(G\) has no isolated vertices on \(F_G\), \(c(G)\) is its
number of nontrivial components.

\begin{lemma}[Separation of the optional-label power]
\label{lem:newV99-power}
For every \(A\supseteq F_G\),
\begin{equation}
 |A|-1-k-|H|
 =
 |A\setminus F_G|+c(G)-1.                                 \label{eq:newV99-power}
\end{equation}
\end{lemma}

\begin{proof}
Use \(|G|=k+|H|\) and
\(c(G)=|F_G|-|G|\).
\end{proof}

\begin{theorem}[Post-collapse one-layer forest bound]
\label{thm:newV99-post-collapse}
Let the support sum range over every \(A\) containing
\(V(E)\).  Then
\begin{align}
 &\left\|
 \sum_{A\supseteq V(E)}
 c_A(E;\boldsymbol x)R_A
 \right\|_{\rm cb}
 \notag\\
 &\quad\le
 \sum_{\substack{
       H\subseteq\binom X2\setminus E\\
       E\cup H\ {\rm a\ forest}}}
 \lambda^{c(E\cup H)-1}
 (1-\lambda)^{|H|}
 \boldsymbol z^H
 \gamma_{F_{E\cup H}}(E\cup H)\,
 \mathfrak r(F_{E\cup H}),                                \label{eq:newV99-post-collapse}
\end{align}
after one uniform enlargement of the reserve base.  In
particular, the right side contains no optional support label.
\end{theorem}

\begin{proof}
Insert \Cref{thm:newV99-layer-expansion} and interchange the two
finite sums before taking norms.  Fix \(H\), set \(G=E\cup H\),
and use \Cref{lem:newV99-power}.  Its contribution is
\[
 \lambda^{c(G)-1}(1-\lambda)^{|H|}\boldsymbol z^H
 \sum_{A\supseteq F_G}
 \lambda^{|A\setminus F_G|}\gamma_A(G)R_A.
\]
Repeat the Laplace proof of
\Cref{thm:newV98-uniform-collapse}, now with forced forest \(G\).
Every optional factor has parameter
\(\lambda e^{-tb_j}\in[0,1]\), so
\Cref{lem:newV97-optional} makes it a contraction.  The resulting
forced term is bounded by
\(\gamma_{F_G}(G)\mathfrak r(F_G)\), including the pair drift
correction.  Take the triangle inequality only after each fixed
\(H\) support sum has collapsed.  This gives
\eqref{eq:newV99-post-collapse}.
\end{proof}

\begin{corollary}[The dense degree fibre is absent after one
ultrametric layer]
\label{cor:newV99-no-dense}
For each fixed owned forest \(E\), the right side of
\eqref{eq:newV99-post-collapse} is indexed by completion forests
inside child blocks, not by arbitrary incidence rows
\(A\supseteq V(E)\).  Hence the
\(\exp(cm^2)\) optional-membership fibre of V96 has been removed
before norms.
\end{corollary}

\begin{proof}
Every label outside \(F_{E\cup H}\) has already been summed in
the contraction factor
\((1-\lambda e^{-tb_j})I+\lambda e^{-tb_j}D^{R_j}(U)\).
\end{proof}

\begin{assessment}[Two complementary small factors]
\label{ass:newV99-beta}
Each surviving completion forest carries the exact beta-type
ledger
\[
 \lambda^{c(G)-1}(1-\lambda)^{|H|}.
\]
Many components are suppressed when \(\lambda\) is small; many
chosen child-layer edges are suppressed when \(\lambda\) is
close to one.  Neither factor may be bounded by one before the
forest sum is performed.
\end{assessment}

\section{The remaining forest projection}
\label{sec:newV99-projection}

\begin{lemma}[Connected completion row]
\label{lem:newV99-connected-row}
If \(G=E\cup H\) is a tree on \(F_G\), then
\begin{equation}
 \gamma_{F_G}(G)
 =
 {1\over W_{F_G}}
 \prod_{i\in F_G}b_i^{d_G(i)},                             \label{eq:newV99-tree-probability}
\end{equation}
and the component factor in
\eqref{eq:newV99-post-collapse} is
\(\lambda^0=1\).
\end{lemma}

\begin{proof}
A spanning tree on \(F_G\) containing the tree \(G\) must equal
\(G\).  Apply the definition of \(\gamma\).  Connectedness gives
\(c(G)=1\).
\end{proof}

\begin{lemma}[Empty child-layer row]
\label{lem:newV99-empty-row}
The term \(H=\varnothing\) in
\eqref{eq:newV99-post-collapse} is exactly the uniform-line
collapse of V98:
\begin{equation}
 \lambda^{c(E)-1}\gamma_{V(E)}(E)\mathfrak r(V(E)).         \label{eq:newV99-empty-row}
\end{equation}
\end{lemma}

\begin{proof}
Set \(H=\varnothing\) in the theorem.
\end{proof}

\begin{definition}[Forest-projection card, FPC]
\label{def:newV99-FPC}
FPC is the aggregate estimate of the right side of
\eqref{eq:newV99-post-collapse}, simultaneously over:
\begin{enumerate}[label=\textup{(\roman*)},leftmargin=3.5em]
\item the disjoint owned forests \(E_\nu\) partitioning the outer
      BKAR tree edges;
\item all child-block completion forests \(H_\nu\);
\item the nested layer parameters;
\item the retained probabilities
      \(\gamma_{F_{G_\nu}}(G_\nu)\); and
\item the V88 component-coalescence exchanges.
\end{enumerate}
The beta ledgers and the Bernoulli/Kruskal probabilities must be
summed before they are dropped.  The target is the PRAH strong
polynomial times \(K^m\).
\end{definition}

\begin{proposition}[FPC proves ULCC]
\label{prop:newV99-FPC}
FPC proves ULCC, hence IGMC and the support-fibre part of SISC.
\end{proposition}

\begin{proof}
The exact top-layer expansion is
\Cref{thm:newV99-layer-expansion}; its edge assignments have unit
mass by \Cref{lem:newV99-layer-probability}.  All optional support
labels are removed by
\Cref{thm:newV99-post-collapse}.  FPC supplies the remaining
uniform forest projection and coalescence bound.  Iterate on the
strictly smaller child blocks of
\Cref{lem:newV98-top-layer}, proving ULCC.  Then apply
\Cref{prop:newV98-ULCC,prop:newV97-IGMC}.
\end{proof}

\begin{warning}[Do not sum completion forests by their raw
number]
\label{warn:newV99-raw-H}
For \(\lambda=1/2\), the scalar beta factor alone assigns
\((1/2)^{c(G)-1+|H|}\).  A raw count of labelled forests can still
be superexponential if their Cayley probability
\(\gamma_{F_G}(G)\) is dropped.  Thus
\eqref{eq:newV99-layer-probability} and
\eqref{eq:newV99-tree-probability} must be used before any forest
census.
\end{warning}

\section{Exact relation to the V88 defect}
\label{sec:newV99-defect}

\begin{lemma}[Forced-support excess after optional collapse]
\label{lem:newV99-forced-excess}
Let the nonempty edge sets \(E_1,\ldots,E_q\) partition an outer
tree on \(m\) labels.  If \(c_\nu\) is the number of nontrivial
components of \(E_\nu\), then the incidence excess of the forced
support family \(V(E_\nu)\) is
\begin{equation}
 g_{\rm forced}
 =
 \sum_{\nu=1}^q(c_\nu-1).                                  \label{eq:newV99-forced-excess}
\end{equation}
\end{lemma}

\begin{proof}
If \(k_\nu=|E_\nu|\), then
\(|V(E_\nu)|=k_\nu+c_\nu\).  Since
\(\sum_\nu k_\nu=m-1\),
\begin{align*}
 g_{\rm forced}
 &=
 \sum_\nu\{|V(E_\nu)|-1\}-m+1\\
 &=
 (m-1)+\sum_\nu(c_\nu-1)-m+1.
\end{align*}
\end{proof}

\begin{assessment}[No new defect was created]
\label{ass:newV99-no-new-defect}
The optional-support collapse converts the arbitrary primitive
incidence excess into the already identified component defect:
disconnected pieces inside one owned edge block.  This is the
same combinatorial quantity as the V88
compatibility-defect/quotient-excess identity after each global
tree is cut by the relevant hierarchy.  FPC must therefore be
proved jointly with component coalescence; treating them as two
independent gases would count the same defect twice.
\end{assessment}

\begin{programme}[V100: compulsory audit and beta--forest
projection]
\label{prog:newV99-V100}
\begin{enumerate}[leftmargin=3em]
\item Perform the compulsory newest-SM/newest-NS audit.
\item Sum the connected completion row by the exact weighted
      Cayley probability, retaining the beta ledger.
\item Test the two-component row with one child-layer edge and
      the V89 weighted exchange.
\item State the precise FPC remainder after this first
      beta--forest calculation.
\item Freeze the completed V100 as the next \texttt{\_f5}
      archive and start a new active V1 with context, a rigorous
      result manifest, and the live bottleneck.
\end{enumerate}
\end{programme}

\begin{firewall}[Claim boundary after V99]
\label{fire:newV99-boundary}
V99 proves the exact affine expansion of one BKAR ultrametric
layer, proves its edge assignments form a subprobability, and
uses the V97 Laplace contraction to remove every optional support
label at that layer.  It leaves a positive forest projection with
the exact beta ledger
\(\lambda^{c-1}(1-\lambda)^{|H|}\), normalized forest
probability, and reserve activity.  It proves the forced-support
excess is precisely the previously identified component defect
and formulates FPC.

V99 does not prove FPC, ULCC, full IGMC, the general
component-coalescence row, SISC, the simple high-excess row,
arbitrary repeated transversal supports, full TOCC, full TCNC,
cyclic TIC, HLTC, full MCWC, OCRC, SCNGC, WPSC, GRLTC, LZTC,
PLRC, WSDC, BRSC, HWSFC, UGCGC, thermal connected WPRAC,
all-arity RGSC, the raw Wilson aggregate strong card, all-order
strong MTA, WUFC, CPM, cutoff removal, reflection positivity,
Osterwalder--Schrader reconstruction, a continuum Yang--Mills
measure, or a positive Hamiltonian gap.  The full Yang--Mills
mass gap has not been proved.
\end{firewall}

\begin{theorem}[V99 result ledger]
\label{thm:newV99-results}
V99:
\begin{enumerate}[label=\textup{(V99.\arabic*)},leftmargin=6.5em]
\item proves the exact one-layer forest expansion;
\item proves the layer-assignment subprobability;
\item separates the optional-label power;
\item proves post-collapse removal of all optional supports;
\item identifies the beta-type forest ledger;
\item isolates the connected and empty layer rows;
\item formulates FPC and proves its sufficiency; and
\item identifies forced excess with component defect.
\end{enumerate}
\end{theorem}

\begin{proof}
Use
\Cref{thm:newV99-layer-expansion,
lem:newV99-layer-probability,
lem:newV99-power,
thm:newV99-post-collapse,
lem:newV99-connected-row,
lem:newV99-empty-row,
prop:newV99-FPC,
lem:newV99-forced-excess}.
\end{proof}

\paragraph{Parent-version record.}
V99 was formed from
\texttt{Renewal\_Yang\_Mills\_Mass\_Gap\_Research\_Notes\_V98.tex}.
The V98 parent had \(44\,404\) logical source lines,
\(1\,441\,088\) bytes, and SHA--256
\[
 \texttt{d6e4a8a89a38e4bf7cc687ff7dadc58e982762b53d68b832922fd8e1b31bc945}.
\]
The next compulsory companion audit is V100.

% =============================================================================
% END APPENDED FIFTH-SERIES V99 MATERIAL
% =============================================================================

\clearpage
\chapter{Fifth-series V100: companion audit, minimum-edge
disintegration, and fixed-vertex forest projection}
\label{ch:newV100-projection}

\section{Compulsory companion audit}
\label{sec:newV100-audit}

The newest settled companion snapshots read for this audit were:
\begin{center}
\small
\begin{tabular}{p{0.24\textwidth}p{0.69\textwidth}}
Standard Model--Einstein &
\texttt{Renewal\_SM\_Einstein\_Continuum\_Research\_Notes\_V24.tex},
\(397\,734\) bytes, \(10\,081\) logical source lines, SHA--256
\texttt{70117850268C2BCD4C17E720384D2E0252FADA71CC348FB2E5B833A84A688432}.
\\ [0.4em]
Navier--Stokes &
\texttt{Renewal\_NS\_Stability\_Research\_Notes\_V15.tex},
\(147\,056\) bytes, \(2\,878\) logical source lines, SHA--256
\texttt{D0433B54C1409B3EE28D870A43A33A43DFF76B7426C182B788290B53D82BEC5F}.
\end{tabular}
\end{center}

\begin{audit}[Standard Model--Einstein V24]
\label{audit:newV100-SM}
SM V24 proves that complete flat overlap data do not determine the
mixed \(R|H|^2\) coefficient: changing an unobserved curvature
endomorphism leaves the flat operator fixed and changes the mixed
owner.  It therefore records the missing metric lattice operator
as an explicit card rather than interpolating it from special
rows.  The corresponding lesson here is that the uniform-line and
partition-endpoint support collapses do not determine the
interior BKAR gate; the actual ultrametric layer recursion must be
proved.  No SM coefficient is imported.
\end{audit}

\begin{audit}[Navier--Stokes V15]
\label{audit:newV100-NS}
NS V15 records that localization of the type-I cross term to a
finite configuration neighbourhood does not create a sign; the
localized gate remains a genuine mean inequality.  Analogously,
removing optional support labels in V99 does not by itself bound
the remaining forced forest.  Its normalized Cayley probability
and beta weights must be summed.  No Navier--Stokes estimate is
imported.
\end{audit}

\begin{audit}[Direction selected by the audit]
\label{audit:newV100-direction}
V100 keeps the actual BKAR integration measure.  It first
disintegrates that measure at its minimum tree-edge parameter,
then sums all layer-completion edges which use no new forced
vertex.  This proves an interior interpolation row, rather than
extrapolating from the two boundary rows.
\end{audit}

\section{Exact minimum-edge disintegration of the BKAR cube}
\label{sec:newV100-minimum}

Let \(T_0\) be a tree on \(m\ge2\) labels.  For
\(\boldsymbol u\in[0,1]^{E(T_0)}\), let
\(\boldsymbol x^{T_0}(\boldsymbol u)\) be its path-minimum
matrix.

\begin{lemma}[Minimum-edge coordinates]
\label{lem:newV100-minimum}
Outside a null set there is a unique edge \(e_*\) minimizing the
\(u\)-coordinates.  Put
\[
 \lambda:=u_{e_*},\qquad
 v_f:={u_f-\lambda\over1-\lambda}
 \quad(f\ne e_*).
\]
If \(\rho(e_*)\) is the two-component partition of
\(T_0-e_*\), then
\begin{equation}
 \boldsymbol x^{T_0}(\boldsymbol u)
 =
 \lambda\boldsymbol 1
 +(1-\lambda)
 \boldsymbol z^{T_0-e_*}(\boldsymbol v),                  \label{eq:newV100-minimum-matrix}
\end{equation}
where \(\boldsymbol z\) is zero across the two components and is
the path-minimum matrix of the corresponding subtree inside each
component.  The coordinate Jacobian is
\begin{equation}
 d\boldsymbol u=(1-\lambda)^{m-2}\,
 d\lambda\,d\boldsymbol v.                                \label{eq:newV100-Jacobian}
\end{equation}
\end{lemma}

\begin{proof}
On the region where \(e_*\) is the unique minimum, every other
coordinate ranges from \(\lambda\) to one, so the displayed
change of variables maps that region bijectively to
\([0,1]\times[0,1]^{m-2}\).  A path crossing \(e_*\) has minimum
\(\lambda\).  A path inside one component has edge coordinates
\(\lambda+(1-\lambda)v_f\), whose minimum is
\(\lambda+(1-\lambda)\min v_f\).  This proves
\eqref{eq:newV100-minimum-matrix}.  The \(m-2\) affine
coordinates give \eqref{eq:newV100-Jacobian}.
\end{proof}

\begin{theorem}[Normalized minimum-edge disintegration]
\label{thm:newV100-disintegration}
For every integrable function \(\Phi\) of the path-minimum
matrix,
\begin{align}
 &\int_{[0,1]^{m-1}}
 \Phi(\boldsymbol x^{T_0}(\boldsymbol u))\,d\boldsymbol u
 \notag\\
 &\quad=
 \sum_{e_*\in E(T_0)}
 \int_0^1(1-\lambda)^{m-2}\,d\lambda
 \int_{[0,1]^{m-2}}
 \Phi\!\left(
  \lambda\boldsymbol1+
  (1-\lambda)\boldsymbol z^{T_0-e_*}(\boldsymbol v)
 \right)d\boldsymbol v.                                  \label{eq:newV100-disintegration}
\end{align}
In particular,
\begin{equation}
 \sum_{e_*\in E(T_0)}
 \int_0^1(1-\lambda)^{m-2}\,d\lambda=1.                   \label{eq:newV100-normalized-minimum}
\end{equation}
\end{theorem}

\begin{proof}
Partition the cube, up to the tie null set, by its unique
minimizing edge and apply
\Cref{lem:newV100-minimum}.  For \(\Phi=1\), the right side is
\((m-1)/(m-1)=1\), proving the normalization.
\end{proof}

\begin{assessment}[The top-layer choice is already paid]
\label{ass:newV100-minimum}
Choosing the edge which creates the top two child blocks costs
\(m-1\), but its minimum-coordinate integral is exactly
\(1/(m-1)\).  Keeping the Jacobian therefore prevents a new
layer census.  Bounding \((1-\lambda)^{m-2}\) by one before
summing \(e_*\) would lose this normalization.
\end{assessment}

\section{Fixed-vertex forest projection}
\label{sec:newV100-fixed}

Let \(E\) be a forest on a fixed vertex set \(F\), with no
isolated vertices, and put
\[
 c:=c(E)=|F|-|E|.
\]
For a child-layer matrix \(0\le\boldsymbol z\le1\), define
\begin{align}
 {\cal P}_{F,E}(\lambda,\boldsymbol z)
 &:=
 \sum_{\substack{
       H\subseteq\binom F2\setminus E\\
       E\cup H\ {\rm a\ forest}}}
 \lambda^{c(E\cup H)-1}
 (1-\lambda)^{|H|}
 \boldsymbol z^H\,
 \gamma_F(E\cup H).                                      \label{eq:newV100-P}
\end{align}

\begin{theorem}[Fixed-vertex projection has unit Cayley mass]
\label{thm:newV100-fixed-projection}
For every \(0\le\lambda\le1\),
\begin{equation}
 {\cal P}_{F,E}(\lambda,\boldsymbol z)
 \le\gamma_F(E)\le1.                                     \label{eq:newV100-fixed-projection}
\end{equation}
If \(z_e=1\) for every edge which can occur in a tree containing
\(E\), then the first inequality is equality.
\end{theorem}

\begin{proof}
Expand
\[
 \gamma_F(E\cup H)
 =
 \sum_{\substack{T\in\mathfrak T(F)\\E\cup H\subseteq T}}
 \mu_F(T)
\]
and interchange the two finite positive sums.  For a fixed
\(T\supseteq E\), the set \(T\setminus E\) has \(c-1\) edges.
Every \(H\subseteq T\setminus E\) reduces the number of
components by \(|H|\), so
\[
 c(E\cup H)-1=c-1-|H|.
\]
Its complete contribution is
\begin{align*}
 \mu_F(T)
 \sum_{H\subseteq T\setminus E}
 \lambda^{c-1-|H|}(1-\lambda)^{|H|}\boldsymbol z^H
 &=
 \mu_F(T)
 \prod_{e\in T\setminus E}
 \{\lambda+(1-\lambda)z_e\}\\
 &\le\mu_F(T).
\end{align*}
Sum \(T\supseteq E\) to obtain \(\gamma_F(E)\).  If all relevant
\(z_e\)'s equal one, every product equals one.
\end{proof}

\begin{corollary}[Exact two-component beta interpolation]
\label{cor:newV100-two-component}
If \(E\) has two components, then
\begin{align}
 &\lambda\gamma_F(E)
 +(1-\lambda)
 \sum_{\substack{e\ {\rm joining}\\
                  \text{the two components}}}
 z_e\,\mu_F(E\cup\{e\})
 \le\gamma_F(E).                                         \label{eq:newV100-two-component}
\end{align}
Equality holds when all displayed \(z_e=1\).
\end{corollary}

\begin{proof}
Here every tree containing \(E\) is \(E\cup\{e\}\) for a unique
crossing edge \(e\).  Specialize
\Cref{thm:newV100-fixed-projection}.
\end{proof}

\begin{corollary}[Fixed-vertex row of the post-collapse bound]
\label{cor:newV100-fixed-row}
In \eqref{eq:newV99-post-collapse}, sum every completion
\(H\) satisfying \(V(E\cup H)=V(E)=F\) before dropping its beta
or Cayley factors.  Their total is at most
\begin{equation}
 \gamma_F(E)\mathfrak r(F)\le\mathfrak r(F).              \label{eq:newV100-fixed-row}
\end{equation}
\end{corollary}

\begin{proof}
The reserve activity \(\mathfrak r(F)\) is common to all these
terms.  Apply \Cref{thm:newV100-fixed-projection}.
\end{proof}

\begin{theorem}[Integrated fixed-vertex FPC row]
\label{thm:newV100-integrated-fixed}
After summing all minimum-edge choices and their BKAR parameters,
the fixed-vertex completion row at the top layer costs at most
one copy of its pre-layer reserve majorant.  Iterating the same
statement down the child trees creates no depth-dependent
constant.
\end{theorem}

\begin{proof}
At a fixed top decomposition, use
\Cref{cor:newV100-fixed-row}.  Integrate in the child parameters.
Then sum the top minimizing edge and \(\lambda\) with
\eqref{eq:newV100-normalized-minimum}.  The bound is unchanged.
Each recursive child has strictly fewer tree edges; repeat.  At
every recursive node the corresponding minimum-edge
disintegration has total mass one, so no factor depending on
depth is accumulated.
\end{proof}

\begin{assessment}[The first interior component-coalescence row
is closed]
\label{ass:newV100-fixed-closed}
Theorem~\ref{thm:newV100-fixed-projection} includes every number
of components of \(E\) and every forest of child-layer edges
which joins those components without introducing a new vertex.
Thus the V89 two-component exchange is not merely a one-defect
phenomenon at this layer: all fixed-vertex coalescences are paid
by one conditional Cayley probability.
\end{assessment}

\section{The precise remaining vertex-accretion row}
\label{sec:newV100-accretion}

\begin{definition}[Vertex-accreting forest-projection card,
VAFPC]
\label{def:newV100-VAFPC}
VAFPC is the part of FPC in which at least one child-layer edge
\(H\) is incident to a label outside the pre-layer forced set
\(V(E)\).  It must:
\begin{enumerate}[label=\textup{(\roman*)},leftmargin=3.5em]
\item reveal the new forced vertices component by component;
\item retain the exact factor
      \(\lambda^{c(E\cup H)-1}(1-\lambda)^{|H|}\);
\item retain \(\gamma_{V(E\cup H)}(E\cup H)\) until all rooted
      child forests have been summed;
\item use the minimum-edge Jacobians before counting hierarchy
      nodes; and
\item return one reserve activity on the pre-layer forced data,
      or an equivalent block-compatible Cayley tree, with total
      cost \(K^m\).
\end{enumerate}
\end{definition}

\begin{proposition}[VAFPC completes FPC]
\label{prop:newV100-VAFPC}
VAFPC, together with
\Cref{thm:newV100-integrated-fixed}, proves FPC, ULCC, IGMC, and
the support-fibre part of SISC.
\end{proposition}

\begin{proof}
Every completion forest \(H\) either leaves the forced vertex set
fixed or introduces a new forced vertex.  The first row is
\Cref{thm:newV100-integrated-fixed}; the second is VAFPC.
Therefore FPC holds.  Apply
\Cref{prop:newV99-FPC,prop:newV98-ULCC,prop:newV97-IGMC}.
\end{proof}

\begin{programme}[Sixth-series V1 after the mandatory freeze]
\label{prog:newV100-restart}
Freeze this completed fifth-series V100 as
\texttt{Renewal\_Yang\_Mills\_Mass\_Gap\_Research\_Notes\_V100\_f5.tex}.
Start a new active V1 containing:
\begin{enumerate}[leftmargin=3em]
\item the context and claim boundary of the mass-gap programme;
\item an audited manifest of the principal proved results;
\item the current nonlinear dependency chain;
\item the VAFPC bottleneck in its exact beta--Cayley form; and
\item the first vertex-accretion calculation, beginning with one
      rooted child tree.
\end{enumerate}
\end{programme}

\begin{firewall}[Claim boundary after V100]
\label{fire:newV100-boundary}
V100 audits SM V24 and NS V15, proves the exact normalized
minimum-edge disintegration of the BKAR cube, proves that every
fixed-vertex forest projection has unit conditional Cayley mass,
closes all fixed-vertex component-coalescence rows through every
ultrametric depth, and isolates VAFPC as the remaining
vertex-accretion operation.

V100 does not prove VAFPC, full FPC, ULCC, full IGMC, SISC, the
simple high-excess row, arbitrary repeated transversal supports,
full TOCC, full TCNC, cyclic TIC, HLTC, full MCWC, OCRC, SCNGC,
WPSC, GRLTC, LZTC, PLRC, WSDC, BRSC, HWSFC, UGCGC, thermal
connected WPRAC, all-arity RGSC, the raw Wilson aggregate strong
card, all-order strong MTA, WUFC, CPM, cutoff removal, reflection
positivity, Osterwalder--Schrader reconstruction, a continuum
Yang--Mills measure, or a positive Hamiltonian gap.  The full
Yang--Mills mass gap has not been proved.
\end{firewall}

\begin{theorem}[V100 result ledger]
\label{thm:newV100-results}
V100:
\begin{enumerate}[label=\textup{(V100.\arabic*)},leftmargin=7em]
\item performs the compulsory SM/NS audit;
\item proves minimum-edge coordinates and their Jacobian;
\item proves normalized BKAR minimum-edge disintegration;
\item proves the fixed-vertex projection theorem;
\item proves the exact two-component beta interpolation;
\item closes the integrated fixed-vertex FPC row;
\item isolates VAFPC; and
\item records the mandatory fifth-series freeze and restart.
\end{enumerate}
\end{theorem}

\begin{proof}
Use
\Cref{audit:newV100-SM,
audit:newV100-NS,
lem:newV100-minimum,
thm:newV100-disintegration,
thm:newV100-fixed-projection,
cor:newV100-two-component,
thm:newV100-integrated-fixed,
prop:newV100-VAFPC}.
\end{proof}

\paragraph{Parent-version record.}
V100 was formed from
\texttt{Renewal\_Yang\_Mills\_Mass\_Gap\_Research\_Notes\_V99.tex}.
The V99 parent had \(44\,802\) logical source lines,
\(1\,454\,235\) bytes, and SHA--256
\[
 \texttt{6cfcbe207cc33f3f51c9fd87af7e4abeee71ebe3d69c69db143daac828392271}.
\]
This version is the terminal active version of the fifth series
and is frozen as archive \texttt{\_f5}.

% =============================================================================
% END APPENDED FIFTH-SERIES V100 MATERIAL
% =============================================================================

\end{document}
